\newpage
\begin{problem}{}{}
% Subject: varc
\textbf{Instructions}

Instructions   

**The passage below is accompanied by a set of six questions. Choose the best answer to each question.**
Understanding where you are in the world is a basic survival skill, which is why we, like most species come hard-wired with specialised brain areas to create cognitive maps of our surroundings. Where humans are unique, though, with the possible exception of honeybees, is that we try to communicate this understanding of the world with others. We have a long history of doing this by drawing maps — the earliest versions yet discovered were scrawled on cave walls 14,000 years ago. Human cultures have been drawing them on stone tablets, papyrus, paper and now computer screens ever since.
Given such a long history of human map-making, it is perhaps surprising that it is only within the last few hundred years that north has been consistently considered to be at the top. In fact, for much of human history, north almost never appeared at the top, according to Jerry Brotton, a map historian... "North was rarely put at the top for the simple fact that north is where darkness comes from," he says. "West is also very unlikely to be put at the top because west is where the sun disappears."
Confusingly, early Chinese maps seem to buck this trend. But, Brotton, says, even though they did have compasses at the time, that isn't the reason that they placed north at the top. Early Chinese compasses were actually oriented to point south, which was considered to be more desirable than deepest darkest north. But in Chinese maps, the Emperor, who lived in the north of the country was always put at the top of the map, with everyone else, his loyal subjects, looking up towards him. "In Chinese culture the Emperor looks south because it's where the winds come from, it's a good direction. North is not very good but you are in a position of subjection to the emperor, so you look up to him," says Brotton.
Given that each culture has a very different idea of who, or what, they should look up to it's perhaps not surprising that there is very little consistency in which way early maps pointed. In ancient Egyptian times the top of the world was east, the position of sunrise. Early Islamic maps favoured south at the top because most of the early Muslim cultures were north of Mecca, so they imagined looking up (south) towards it. Christian maps from the same era (called Mappa Mundi) put east at the top, towards the Garden of Eden and with Jerusalem in the centre.
So when did everyone get together and decide that north was the top? It's tempting to put it down to European explorers like Christopher Columbus and Ferdinand Megellan, who were navigating by the North Star. But Brotton argues that these early explorers didn't think of the world like that at all. "When Columbus describes the world it is in accordance with east being at the top, he says. "Columbus says he is going towards paradise, so his mentality is from a medieval mappa mundi." We've got to remember, adds Brotton, that at the time, "no one knows what they are doing and where they are going."

\vspace{1em}
\textbf{Question 1}

Which one of the following best describes what the passage is trying to do?

\begin{enumerate}[label=\Alph*.]
    \item It questions an explanation about how maps are designed.
    \item It corrects the misconception about the way maps are designed.
    \item It critiques a methodology used to create maps.
    \item It explores some myths about maps.
[](https://cracku.in/1-which-one-of-the-following-best-describes-what-the-x-cat-2017-shift-1)
\end{enumerate}
\vspace{1em}
\textbf{Solution:}

Option A: He gives explanation/reasoning behind the various methods adopted in history but does not question them. 
Option C: The author doesnt discuss about the merits/ demirits or evaluate a method of map-making. It is involved more with the history of the method.
Option D speaks about myth. While most of the data are quoted from history, they cannot be misrepresented as a myth.
The author starts the passage by talking about the history of map making. The author then mentions how north was never put at the top in ancient times.   
"Given such a long history of human map-making, it is perhaps surprising that it is only within the last few hundred years that north has been consistently considered to be at the top. In fact, for much of human history, north almost never appeared at the top".
He implicitly means that people have a preemptive notion of maps facing the north. He goes on to mention that it was considered a bad direction. He says that north being put at the top is a fairly recent phenomenon. He then goes on to discuss why different people started putting north at the top. He mentions that the reasons for different people putting north at the top were different from what people think now. 
He cites various examples to show that north' s presence in top is more recent and due to varied factors.  
Hence, he is trying to clear certain misconceptions about why north is put at the top in the maps. Thus, option B is the most suitable answer.

\end{problem}

\newpage
\begin{problem}{}{}
% Subject: varc
\textbf{Instructions}

Instructions   

**The passage below is accompanied by a set of six questions. Choose the best answer to each question.**
Understanding where you are in the world is a basic survival skill, which is why we, like most species come hard-wired with specialised brain areas to create cognitive maps of our surroundings. Where humans are unique, though, with the possible exception of honeybees, is that we try to communicate this understanding of the world with others. We have a long history of doing this by drawing maps — the earliest versions yet discovered were scrawled on cave walls 14,000 years ago. Human cultures have been drawing them on stone tablets, papyrus, paper and now computer screens ever since.
Given such a long history of human map-making, it is perhaps surprising that it is only within the last few hundred years that north has been consistently considered to be at the top. In fact, for much of human history, north almost never appeared at the top, according to Jerry Brotton, a map historian... "North was rarely put at the top for the simple fact that north is where darkness comes from," he says. "West is also very unlikely to be put at the top because west is where the sun disappears."
Confusingly, early Chinese maps seem to buck this trend. But, Brotton, says, even though they did have compasses at the time, that isn't the reason that they placed north at the top. Early Chinese compasses were actually oriented to point south, which was considered to be more desirable than deepest darkest north. But in Chinese maps, the Emperor, who lived in the north of the country was always put at the top of the map, with everyone else, his loyal subjects, looking up towards him. "In Chinese culture the Emperor looks south because it's where the winds come from, it's a good direction. North is not very good but you are in a position of subjection to the emperor, so you look up to him," says Brotton.
Given that each culture has a very different idea of who, or what, they should look up to it's perhaps not surprising that there is very little consistency in which way early maps pointed. In ancient Egyptian times the top of the world was east, the position of sunrise. Early Islamic maps favoured south at the top because most of the early Muslim cultures were north of Mecca, so they imagined looking up (south) towards it. Christian maps from the same era (called Mappa Mundi) put east at the top, towards the Garden of Eden and with Jerusalem in the centre.
So when did everyone get together and decide that north was the top? It's tempting to put it down to European explorers like Christopher Columbus and Ferdinand Megellan, who were navigating by the North Star. But Brotton argues that these early explorers didn't think of the world like that at all. "When Columbus describes the world it is in accordance with east being at the top, he says. "Columbus says he is going towards paradise, so his mentality is from a medieval mappa mundi." We've got to remember, adds Brotton, that at the time, "no one knows what they are doing and where they are going."

\vspace{1em}
\textbf{Question 2}

Early maps did NOT put north at the top for all the following reasons EXCEPT

\begin{enumerate}[label=\Alph*.]
    \item North was the source of darkness.
    \item South was favoured by some emperors.
    \item East and south were more important for religious reasons for some civilisations.
    \item East was considered by some civilisations to be a more positive direction.
[](https://cracku.in/2-early-maps-did-not-put-north-at-the-top-for-all-th-x-cat-2017-shift-1)
\end{enumerate}
\vspace{1em}
\textbf{Solution:}

The passage mentions that the Chinese put North at the top of the map because the emperor would live in the North and he preferred to look towards South. Hence, the fact that South was preferred by some emperors is not a reason why North was put at the top. Hence, option B is false. All other options are mentioned in the passage.

\end{problem}

\newpage
\begin{problem}{}{}
% Subject: varc
\textbf{Instructions}

Instructions   

**The passage below is accompanied by a set of six questions. Choose the best answer to each question.**
Understanding where you are in the world is a basic survival skill, which is why we, like most species come hard-wired with specialised brain areas to create cognitive maps of our surroundings. Where humans are unique, though, with the possible exception of honeybees, is that we try to communicate this understanding of the world with others. We have a long history of doing this by drawing maps — the earliest versions yet discovered were scrawled on cave walls 14,000 years ago. Human cultures have been drawing them on stone tablets, papyrus, paper and now computer screens ever since.
Given such a long history of human map-making, it is perhaps surprising that it is only within the last few hundred years that north has been consistently considered to be at the top. In fact, for much of human history, north almost never appeared at the top, according to Jerry Brotton, a map historian... "North was rarely put at the top for the simple fact that north is where darkness comes from," he says. "West is also very unlikely to be put at the top because west is where the sun disappears."
Confusingly, early Chinese maps seem to buck this trend. But, Brotton, says, even though they did have compasses at the time, that isn't the reason that they placed north at the top. Early Chinese compasses were actually oriented to point south, which was considered to be more desirable than deepest darkest north. But in Chinese maps, the Emperor, who lived in the north of the country was always put at the top of the map, with everyone else, his loyal subjects, looking up towards him. "In Chinese culture the Emperor looks south because it's where the winds come from, it's a good direction. North is not very good but you are in a position of subjection to the emperor, so you look up to him," says Brotton.
Given that each culture has a very different idea of who, or what, they should look up to it's perhaps not surprising that there is very little consistency in which way early maps pointed. In ancient Egyptian times the top of the world was east, the position of sunrise. Early Islamic maps favoured south at the top because most of the early Muslim cultures were north of Mecca, so they imagined looking up (south) towards it. Christian maps from the same era (called Mappa Mundi) put east at the top, towards the Garden of Eden and with Jerusalem in the centre.
So when did everyone get together and decide that north was the top? It's tempting to put it down to European explorers like Christopher Columbus and Ferdinand Megellan, who were navigating by the North Star. But Brotton argues that these early explorers didn't think of the world like that at all. "When Columbus describes the world it is in accordance with east being at the top, he says. "Columbus says he is going towards paradise, so his mentality is from a medieval mappa mundi." We've got to remember, adds Brotton, that at the time, "no one knows what they are doing and where they are going."

\vspace{1em}
\textbf{Question 3}

According to the passage, early Chinese maps placed north at the top because

\begin{enumerate}[label=\Alph*.]
    \item the Chinese invented the compass and were aware of magnetic north.
    \item they wanted to show respect to the emperor.
    \item the Chinese emperor appreciated the winds from the south.
    \item north was considered the most desirable direction.
[](https://cracku.in/3-according-to-the-passage-early-chinese-maps-placed-x-cat-2017-shift-1)
\end{enumerate}
\vspace{1em}
\textbf{Solution:}

We can straightaway eliminate options A and D. The passage states that the Chinese compasses pointed to magnetic south and that south was considered a more desirable direction. While option C is true, that is not the reason why North was placed at the top. 
The passage states that the emperor lived in the north and hence maps depicted him as above his subjects. Thus, north was placed at the top of the map to show respect for the emperor. Thus, option B is correct.

\end{problem}

\newpage
\begin{problem}{}{}
% Subject: varc
\textbf{Instructions}

Instructions   

**The passage below is accompanied by a set of six questions. Choose the best answer to each question.**
Understanding where you are in the world is a basic survival skill, which is why we, like most species come hard-wired with specialised brain areas to create cognitive maps of our surroundings. Where humans are unique, though, with the possible exception of honeybees, is that we try to communicate this understanding of the world with others. We have a long history of doing this by drawing maps — the earliest versions yet discovered were scrawled on cave walls 14,000 years ago. Human cultures have been drawing them on stone tablets, papyrus, paper and now computer screens ever since.
Given such a long history of human map-making, it is perhaps surprising that it is only within the last few hundred years that north has been consistently considered to be at the top. In fact, for much of human history, north almost never appeared at the top, according to Jerry Brotton, a map historian... "North was rarely put at the top for the simple fact that north is where darkness comes from," he says. "West is also very unlikely to be put at the top because west is where the sun disappears."
Confusingly, early Chinese maps seem to buck this trend. But, Brotton, says, even though they did have compasses at the time, that isn't the reason that they placed north at the top. Early Chinese compasses were actually oriented to point south, which was considered to be more desirable than deepest darkest north. But in Chinese maps, the Emperor, who lived in the north of the country was always put at the top of the map, with everyone else, his loyal subjects, looking up towards him. "In Chinese culture the Emperor looks south because it's where the winds come from, it's a good direction. North is not very good but you are in a position of subjection to the emperor, so you look up to him," says Brotton.
Given that each culture has a very different idea of who, or what, they should look up to it's perhaps not surprising that there is very little consistency in which way early maps pointed. In ancient Egyptian times the top of the world was east, the position of sunrise. Early Islamic maps favoured south at the top because most of the early Muslim cultures were north of Mecca, so they imagined looking up (south) towards it. Christian maps from the same era (called Mappa Mundi) put east at the top, towards the Garden of Eden and with Jerusalem in the centre.
So when did everyone get together and decide that north was the top? It's tempting to put it down to European explorers like Christopher Columbus and Ferdinand Megellan, who were navigating by the North Star. But Brotton argues that these early explorers didn't think of the world like that at all. "When Columbus describes the world it is in accordance with east being at the top, he says. "Columbus says he is going towards paradise, so his mentality is from a medieval mappa mundi." We've got to remember, adds Brotton, that at the time, "no one knows what they are doing and where they are going."

\vspace{1em}
\textbf{Question 4}

It can be inferred from the passage that European explorers like Columbus and Megellan

\begin{enumerate}[label=\Alph*.]
    \item set the precedent for north-up maps
    \item navigated by the compass
    \item used an eastward orientation for religious reasons
    \item navigated with the help of early maps
[](https://cracku.in/4-it-can-be-inferred-from-the-passage-that-european--x-cat-2017-shift-1)
\end{enumerate}
\vspace{1em}
\textbf{Solution:}

We can straightaway eliminate option A. The author says that though one might think that the trend of north-up maps was set by these explorers, this is in fact not true. Options B and D are also incorrect. The passage says that the explorers navigated with the help of the North Star. 
In the passage, it is given that "When Columbus describes the world it is in accordance with east being at the top, he says. "Columbus says he is going towards paradise, so his mentality is from a medieval mappa mundi." Hence, we can infer that statement C is true.

\end{problem}

\newpage
\begin{problem}{}{}
% Subject: varc
\textbf{Instructions}

Instructions   

**The passage below is accompanied by a set of six questions. Choose the best answer to each question.**
Understanding where you are in the world is a basic survival skill, which is why we, like most species come hard-wired with specialised brain areas to create cognitive maps of our surroundings. Where humans are unique, though, with the possible exception of honeybees, is that we try to communicate this understanding of the world with others. We have a long history of doing this by drawing maps — the earliest versions yet discovered were scrawled on cave walls 14,000 years ago. Human cultures have been drawing them on stone tablets, papyrus, paper and now computer screens ever since.
Given such a long history of human map-making, it is perhaps surprising that it is only within the last few hundred years that north has been consistently considered to be at the top. In fact, for much of human history, north almost never appeared at the top, according to Jerry Brotton, a map historian... "North was rarely put at the top for the simple fact that north is where darkness comes from," he says. "West is also very unlikely to be put at the top because west is where the sun disappears."
Confusingly, early Chinese maps seem to buck this trend. But, Brotton, says, even though they did have compasses at the time, that isn't the reason that they placed north at the top. Early Chinese compasses were actually oriented to point south, which was considered to be more desirable than deepest darkest north. But in Chinese maps, the Emperor, who lived in the north of the country was always put at the top of the map, with everyone else, his loyal subjects, looking up towards him. "In Chinese culture the Emperor looks south because it's where the winds come from, it's a good direction. North is not very good but you are in a position of subjection to the emperor, so you look up to him," says Brotton.
Given that each culture has a very different idea of who, or what, they should look up to it's perhaps not surprising that there is very little consistency in which way early maps pointed. In ancient Egyptian times the top of the world was east, the position of sunrise. Early Islamic maps favoured south at the top because most of the early Muslim cultures were north of Mecca, so they imagined looking up (south) towards it. Christian maps from the same era (called Mappa Mundi) put east at the top, towards the Garden of Eden and with Jerusalem in the centre.
So when did everyone get together and decide that north was the top? It's tempting to put it down to European explorers like Christopher Columbus and Ferdinand Megellan, who were navigating by the North Star. But Brotton argues that these early explorers didn't think of the world like that at all. "When Columbus describes the world it is in accordance with east being at the top, he says. "Columbus says he is going towards paradise, so his mentality is from a medieval mappa mundi." We've got to remember, adds Brotton, that at the time, "no one knows what they are doing and where they are going."

\vspace{1em}
\textbf{Question 5}

Which one of the following about the northern orientation of modern maps is asserted in the passage?

\begin{enumerate}[label=\Alph*.]
    \item The biggest contributory factor was the understanding of magnetic north.
    \item The biggest contributory factor was the role of European explorers.
    \item The biggest contributory factor was the influence of Christian maps.
    \item The biggest contributory factor is not stated in the passage.
[](https://cracku.in/5-which-one-of-the-following-about-the-northern-orie-x-cat-2017-shift-1)
\end{enumerate}
\vspace{1em}
\textbf{Solution:}

In the passage, the author discusses how North was traditionally not put on the top of early maps. The author explicitly refutes the role of the compass and of European explorers in placing North at the top of maps. Hence, we can eliminate options A and B. The author says that East was placed at the top of Christian maps. Hence, option C is also incorrect. Thought the author counters all known explanations as to why North was placed on the top, he does not offer any explanation of his own. Hence, option D is correct.

\end{problem}

\newpage
\begin{problem}{}{}
% Subject: varc
\textbf{Instructions}

Instructions   

**The passage below is accompanied by a set of six questions. Choose the best answer to each question.**
Understanding where you are in the world is a basic survival skill, which is why we, like most species come hard-wired with specialised brain areas to create cognitive maps of our surroundings. Where humans are unique, though, with the possible exception of honeybees, is that we try to communicate this understanding of the world with others. We have a long history of doing this by drawing maps — the earliest versions yet discovered were scrawled on cave walls 14,000 years ago. Human cultures have been drawing them on stone tablets, papyrus, paper and now computer screens ever since.
Given such a long history of human map-making, it is perhaps surprising that it is only within the last few hundred years that north has been consistently considered to be at the top. In fact, for much of human history, north almost never appeared at the top, according to Jerry Brotton, a map historian... "North was rarely put at the top for the simple fact that north is where darkness comes from," he says. "West is also very unlikely to be put at the top because west is where the sun disappears."
Confusingly, early Chinese maps seem to buck this trend. But, Brotton, says, even though they did have compasses at the time, that isn't the reason that they placed north at the top. Early Chinese compasses were actually oriented to point south, which was considered to be more desirable than deepest darkest north. But in Chinese maps, the Emperor, who lived in the north of the country was always put at the top of the map, with everyone else, his loyal subjects, looking up towards him. "In Chinese culture the Emperor looks south because it's where the winds come from, it's a good direction. North is not very good but you are in a position of subjection to the emperor, so you look up to him," says Brotton.
Given that each culture has a very different idea of who, or what, they should look up to it's perhaps not surprising that there is very little consistency in which way early maps pointed. In ancient Egyptian times the top of the world was east, the position of sunrise. Early Islamic maps favoured south at the top because most of the early Muslim cultures were north of Mecca, so they imagined looking up (south) towards it. Christian maps from the same era (called Mappa Mundi) put east at the top, towards the Garden of Eden and with Jerusalem in the centre.
So when did everyone get together and decide that north was the top? It's tempting to put it down to European explorers like Christopher Columbus and Ferdinand Megellan, who were navigating by the North Star. But Brotton argues that these early explorers didn't think of the world like that at all. "When Columbus describes the world it is in accordance with east being at the top, he says. "Columbus says he is going towards paradise, so his mentality is from a medieval mappa mundi." We've got to remember, adds Brotton, that at the time, "no one knows what they are doing and where they are going."

\vspace{1em}
\textbf{Question 6}

The role of natural phenomena in influencing map-making conventions is seen most clearly in

\begin{enumerate}[label=\Alph*.]
    \item early Egyptian maps
    \item early Islamic maps
    \item early Chinese maps
    \item early Christian maps
[](https://cracku.in/6-the-role-of-natural-phenomena-in-influencing-map-m-x-cat-2017-shift-1)
\end{enumerate}
\vspace{1em}
\textbf{Solution:}

According to the passage, early Egyptian maps placed the East at the top because that was the position of sunrise. Hence, we can say that natural phenomena dictated the map-making convention in this case. Thus, option A is correct.
Options B and D are incorrect as the conventions were decided by religious factors and not natural phenomena. Option C also can be eliminated as the orientation was a result of their desire to honour their emperor. Hence, the answer is option A.

Instructions   

**The passage below is accompanied by a set of six questions. Choose the best answer to each question.**
I used a smartphone GPS to find my way through the cobblestoned maze of Geneva's Old Town, in search of a handmade machine that changed the world more than any other invention. Near a 13th-century cathedral in this Swiss city on the shores of a lovely lake, I found what I was looking for: a Gutenberg printing press. "This was the Internet of its day — at least as influential as the iPhone," said Gabriel de Montmollin, the director of the Museum of the Reformation, toying with the replica of Johann Gutenberg's great invention.
Before the invention of the printing press, it used to take four monks up to a year to produce a single book. With the advance in movable type in 15th-century Europe, one press could crank out 3,000 pages a day. Before long, average people could travel to places that used to be unknown to them — with maps! Medical information passed more freely and quickly, diminishing the sway of quacks. The printing press offered the prospect that tyrants would never be able to kill a book or suppress an idea. Gutenberg's brainchild broke the monopoly that clerics had on scripture. And later, stirred by pamphlets from a version of that same press, the American colonies rose up against a king and gave birth to a nation.
So, a question in the summer of this 10th anniversary of the iPhone: has the device that is perhaps the most revolutionary of all time given us a single magnificent idea? Nearly every advancement of the written word through new technology has also advanced humankind. Sure, you can say the iPhone changed everything. By putting the world's recorded knowledge in the palm of a hand, it revolutionized work, dining, travel and socializing. It made us more narcissistic — here's more of me doing cool stuff! — and it unleashed an army of awful trolls. We no longer have the patience to sit through a baseball game without that reach to the pocket. And one more casualty of Apple selling more than a billion phones in a decade's time: daydreaming has become a lost art.
For all of that, I'm still waiting to see if the iPhone can do what the printing press did for religion and democracy...the Geneva museum makes a strong case that the printing press opened more minds than anything else...it's hard to imagine the French or American revolutions without those enlightened voices in print...
Not long after Steve Jobs introduced his iPhone, he said the bound book was probably headed for history's attic. Not so fast. After a period of rapid growth in e-books, something closer to the medium for Chaucer's volumes has made a great comeback
The hope of the iPhone, and the Internet in general, was that it would free people in closed societies. But the failure of the Arab Spring, and the continued suppression of ideas in North Korea, China and Iran, has not borne that out. The iPhone is still young. It has certainly been "one of the most important, world-changing and successful products in. history," as Apple C.E.O. Tim Cook said. But I'm not sure if the world changed for the better with the iPhone — as it did with the printing press — or merely changed.

\end{problem}

\newpage
\begin{problem}{}{}
% Subject: varc
\textbf{Instructions}

Instructions   

**The passage below is accompanied by a set of six questions. Choose the best answer to each question.**
Understanding where you are in the world is a basic survival skill, which is why we, like most species come hard-wired with specialised brain areas to create cognitive maps of our surroundings. Where humans are unique, though, with the possible exception of honeybees, is that we try to communicate this understanding of the world with others. We have a long history of doing this by drawing maps — the earliest versions yet discovered were scrawled on cave walls 14,000 years ago. Human cultures have been drawing them on stone tablets, papyrus, paper and now computer screens ever since.
Given such a long history of human map-making, it is perhaps surprising that it is only within the last few hundred years that north has been consistently considered to be at the top. In fact, for much of human history, north almost never appeared at the top, according to Jerry Brotton, a map historian... "North was rarely put at the top for the simple fact that north is where darkness comes from," he says. "West is also very unlikely to be put at the top because west is where the sun disappears."
Confusingly, early Chinese maps seem to buck this trend. But, Brotton, says, even though they did have compasses at the time, that isn't the reason that they placed north at the top. Early Chinese compasses were actually oriented to point south, which was considered to be more desirable than deepest darkest north. But in Chinese maps, the Emperor, who lived in the north of the country was always put at the top of the map, with everyone else, his loyal subjects, looking up towards him. "In Chinese culture the Emperor looks south because it's where the winds come from, it's a good direction. North is not very good but you are in a position of subjection to the emperor, so you look up to him," says Brotton.
Given that each culture has a very different idea of who, or what, they should look up to it's perhaps not surprising that there is very little consistency in which way early maps pointed. In ancient Egyptian times the top of the world was east, the position of sunrise. Early Islamic maps favoured south at the top because most of the early Muslim cultures were north of Mecca, so they imagined looking up (south) towards it. Christian maps from the same era (called Mappa Mundi) put east at the top, towards the Garden of Eden and with Jerusalem in the centre.
So when did everyone get together and decide that north was the top? It's tempting to put it down to European explorers like Christopher Columbus and Ferdinand Megellan, who were navigating by the North Star. But Brotton argues that these early explorers didn't think of the world like that at all. "When Columbus describes the world it is in accordance with east being at the top, he says. "Columbus says he is going towards paradise, so his mentality is from a medieval mappa mundi." We've got to remember, adds Brotton, that at the time, "no one knows what they are doing and where they are going."

\vspace{1em}
\textbf{Question 7}

The printing press has been likened to the Internet for which one of the following reasons?

\begin{enumerate}[label=\Alph*.]
    \item It enabled rapid access to new information and the sharing of new ideas.
    \item It represented new and revolutionary technology compared to the past.
    \item It encouraged reading among people by giving them access to thousands of books.
    \item It gave people access to pamphlets and literature in several languages.
[](https://cracku.in/7-the-printing-press-has-been-likened-to-the-interne-x-cat-2017-shift-1)
\end{enumerate}
\vspace{1em}
\textbf{Solution:}

In the first passage, the author mentions printing press as the internet of its day. Immediately, in the next paragraph he elucidates how printing press helped in spreading ideas and information. Thus, the author likened the printing press to the internet because it enabled access to new information and sharing of ideas.
Hence, option A is the correct answer.

\end{problem}

\newpage
\begin{problem}{}{}
% Subject: varc
\textbf{Instructions}

Instructions   

**The passage below is accompanied by a set of six questions. Choose the best answer to each question.**
Understanding where you are in the world is a basic survival skill, which is why we, like most species come hard-wired with specialised brain areas to create cognitive maps of our surroundings. Where humans are unique, though, with the possible exception of honeybees, is that we try to communicate this understanding of the world with others. We have a long history of doing this by drawing maps — the earliest versions yet discovered were scrawled on cave walls 14,000 years ago. Human cultures have been drawing them on stone tablets, papyrus, paper and now computer screens ever since.
Given such a long history of human map-making, it is perhaps surprising that it is only within the last few hundred years that north has been consistently considered to be at the top. In fact, for much of human history, north almost never appeared at the top, according to Jerry Brotton, a map historian... "North was rarely put at the top for the simple fact that north is where darkness comes from," he says. "West is also very unlikely to be put at the top because west is where the sun disappears."
Confusingly, early Chinese maps seem to buck this trend. But, Brotton, says, even though they did have compasses at the time, that isn't the reason that they placed north at the top. Early Chinese compasses were actually oriented to point south, which was considered to be more desirable than deepest darkest north. But in Chinese maps, the Emperor, who lived in the north of the country was always put at the top of the map, with everyone else, his loyal subjects, looking up towards him. "In Chinese culture the Emperor looks south because it's where the winds come from, it's a good direction. North is not very good but you are in a position of subjection to the emperor, so you look up to him," says Brotton.
Given that each culture has a very different idea of who, or what, they should look up to it's perhaps not surprising that there is very little consistency in which way early maps pointed. In ancient Egyptian times the top of the world was east, the position of sunrise. Early Islamic maps favoured south at the top because most of the early Muslim cultures were north of Mecca, so they imagined looking up (south) towards it. Christian maps from the same era (called Mappa Mundi) put east at the top, towards the Garden of Eden and with Jerusalem in the centre.
So when did everyone get together and decide that north was the top? It's tempting to put it down to European explorers like Christopher Columbus and Ferdinand Megellan, who were navigating by the North Star. But Brotton argues that these early explorers didn't think of the world like that at all. "When Columbus describes the world it is in accordance with east being at the top, he says. "Columbus says he is going towards paradise, so his mentality is from a medieval mappa mundi." We've got to remember, adds Brotton, that at the time, "no one knows what they are doing and where they are going."

\vspace{1em}
\textbf{Question 8}

According to the passage, the invention of the printing press did all of the following EXCEPT

\begin{enumerate}[label=\Alph*.]
    \item promoted the spread of enlightened political views across countries.
    \item gave people direct access to authentic medical information and religious texts.
    \item shortened the time taken to produce books and pamphlets.
    \item enabled people to perform various tasks simultaneously.
[](https://cracku.in/8-according-to-the-passage-the-invention-of-the-prin-x-cat-2017-shift-1)
\end{enumerate}
\vspace{1em}
\textbf{Solution:}

From the lines "Medical information passed more freely and quickly, diminishing the sway of quacks" and "Gutenberg's brainchild broke the monopoly that clerics had on scripture" , we can infer that option B is true. 
From the lines "And later, stirred by pamphlets from a version of that same press, the American colonies rose up against a king and gave birth to a nation" and "it's hard to imagine the French or American revolutions without those enlightened voices in print", we can infer option A is true.
From the lines "Before the invention of the printing press, it used to take four monks up to a year to produce a single book. With the advance in movable type in 15th-century Europe, one press could crank out 3,000 pages a day", we can infer that option C is true.
Option D has not been stated nor implied anywhere in the passage.

\end{problem}

\newpage
\begin{problem}{}{}
% Subject: varc
\textbf{Instructions}

Instructions   

**The passage below is accompanied by a set of six questions. Choose the best answer to each question.**
Understanding where you are in the world is a basic survival skill, which is why we, like most species come hard-wired with specialised brain areas to create cognitive maps of our surroundings. Where humans are unique, though, with the possible exception of honeybees, is that we try to communicate this understanding of the world with others. We have a long history of doing this by drawing maps — the earliest versions yet discovered were scrawled on cave walls 14,000 years ago. Human cultures have been drawing them on stone tablets, papyrus, paper and now computer screens ever since.
Given such a long history of human map-making, it is perhaps surprising that it is only within the last few hundred years that north has been consistently considered to be at the top. In fact, for much of human history, north almost never appeared at the top, according to Jerry Brotton, a map historian... "North was rarely put at the top for the simple fact that north is where darkness comes from," he says. "West is also very unlikely to be put at the top because west is where the sun disappears."
Confusingly, early Chinese maps seem to buck this trend. But, Brotton, says, even though they did have compasses at the time, that isn't the reason that they placed north at the top. Early Chinese compasses were actually oriented to point south, which was considered to be more desirable than deepest darkest north. But in Chinese maps, the Emperor, who lived in the north of the country was always put at the top of the map, with everyone else, his loyal subjects, looking up towards him. "In Chinese culture the Emperor looks south because it's where the winds come from, it's a good direction. North is not very good but you are in a position of subjection to the emperor, so you look up to him," says Brotton.
Given that each culture has a very different idea of who, or what, they should look up to it's perhaps not surprising that there is very little consistency in which way early maps pointed. In ancient Egyptian times the top of the world was east, the position of sunrise. Early Islamic maps favoured south at the top because most of the early Muslim cultures were north of Mecca, so they imagined looking up (south) towards it. Christian maps from the same era (called Mappa Mundi) put east at the top, towards the Garden of Eden and with Jerusalem in the centre.
So when did everyone get together and decide that north was the top? It's tempting to put it down to European explorers like Christopher Columbus and Ferdinand Megellan, who were navigating by the North Star. But Brotton argues that these early explorers didn't think of the world like that at all. "When Columbus describes the world it is in accordance with east being at the top, he says. "Columbus says he is going towards paradise, so his mentality is from a medieval mappa mundi." We've got to remember, adds Brotton, that at the time, "no one knows what they are doing and where they are going."

\vspace{1em}
\textbf{Question 9}

Steve Jobs predicted which one of the following with the introduction of the iPhone?

\begin{enumerate}[label=\Alph*.]
    \item People would switch from reading on the Internet to reading on their iPhones.
    \item People would lose interest in historical and traditional classics.
    \item Reading printed books would become a thing of the past.
    \item The production of e-books would eventually fall.
[](https://cracku.in/9-steve-jobs-predicted-which-one-of-the-following-wi-x-cat-2017-shift-1)
\end{enumerate}
\vspace{1em}
\textbf{Solution:}

Refer to the following lines - "Not long after Steve Jobs introduced his iPhone, he said the bound book was probably headed for history's attic". Thus, we can infer that Steve Jobs predicted that reading printed books would become a thing of the past. Hence, option C is the right answer.

\end{problem}

\newpage
\begin{problem}{}{}
% Subject: varc
\textbf{Instructions}

Instructions   

**The passage below is accompanied by a set of six questions. Choose the best answer to each question.**
Understanding where you are in the world is a basic survival skill, which is why we, like most species come hard-wired with specialised brain areas to create cognitive maps of our surroundings. Where humans are unique, though, with the possible exception of honeybees, is that we try to communicate this understanding of the world with others. We have a long history of doing this by drawing maps — the earliest versions yet discovered were scrawled on cave walls 14,000 years ago. Human cultures have been drawing them on stone tablets, papyrus, paper and now computer screens ever since.
Given such a long history of human map-making, it is perhaps surprising that it is only within the last few hundred years that north has been consistently considered to be at the top. In fact, for much of human history, north almost never appeared at the top, according to Jerry Brotton, a map historian... "North was rarely put at the top for the simple fact that north is where darkness comes from," he says. "West is also very unlikely to be put at the top because west is where the sun disappears."
Confusingly, early Chinese maps seem to buck this trend. But, Brotton, says, even though they did have compasses at the time, that isn't the reason that they placed north at the top. Early Chinese compasses were actually oriented to point south, which was considered to be more desirable than deepest darkest north. But in Chinese maps, the Emperor, who lived in the north of the country was always put at the top of the map, with everyone else, his loyal subjects, looking up towards him. "In Chinese culture the Emperor looks south because it's where the winds come from, it's a good direction. North is not very good but you are in a position of subjection to the emperor, so you look up to him," says Brotton.
Given that each culture has a very different idea of who, or what, they should look up to it's perhaps not surprising that there is very little consistency in which way early maps pointed. In ancient Egyptian times the top of the world was east, the position of sunrise. Early Islamic maps favoured south at the top because most of the early Muslim cultures were north of Mecca, so they imagined looking up (south) towards it. Christian maps from the same era (called Mappa Mundi) put east at the top, towards the Garden of Eden and with Jerusalem in the centre.
So when did everyone get together and decide that north was the top? It's tempting to put it down to European explorers like Christopher Columbus and Ferdinand Megellan, who were navigating by the North Star. But Brotton argues that these early explorers didn't think of the world like that at all. "When Columbus describes the world it is in accordance with east being at the top, he says. "Columbus says he is going towards paradise, so his mentality is from a medieval mappa mundi." We've got to remember, adds Brotton, that at the time, "no one knows what they are doing and where they are going."

\vspace{1em}
\textbf{Question 10}

"I'm still waiting to see if the iPhone can do what the printing press did for religion and democracy." The author uses which one of the following to indicate his uncertainty?

\begin{enumerate}[label=\Alph*.]
    \item The rise of religious groups in many parts of the world.
    \item The expansion in trolling and narcissism among users of the Internet.
    \item The continued suppression of free speech in closed societies.
    \item The decline in reading habits among those who use the device.
[](https://cracku.in/10-im-still-waiting-to-see-if-the-iphone-can-do-what--x-cat-2017-shift-1)
\end{enumerate}
\vspace{1em}
\textbf{Solution:}

The author says that the iPhone has not fulfilled its potential as a piece of revolutionary technology. He goes on to say that the hope was that the iPhone could help in liberating people in closed societies. However, the failure of the Arab spring and continued suppression in places like North Korea shows that this has not happened. Hence, the author uses the continued suppression of free speech in closed societies to indicate why he is still uncertain about the potential of the iPhone. Hence, option C is correct.

\end{problem}

\newpage
\begin{problem}{}{}
% Subject: varc
\textbf{Instructions}

Instructions   

**The passage below is accompanied by a set of six questions. Choose the best answer to each question.**
Understanding where you are in the world is a basic survival skill, which is why we, like most species come hard-wired with specialised brain areas to create cognitive maps of our surroundings. Where humans are unique, though, with the possible exception of honeybees, is that we try to communicate this understanding of the world with others. We have a long history of doing this by drawing maps — the earliest versions yet discovered were scrawled on cave walls 14,000 years ago. Human cultures have been drawing them on stone tablets, papyrus, paper and now computer screens ever since.
Given such a long history of human map-making, it is perhaps surprising that it is only within the last few hundred years that north has been consistently considered to be at the top. In fact, for much of human history, north almost never appeared at the top, according to Jerry Brotton, a map historian... "North was rarely put at the top for the simple fact that north is where darkness comes from," he says. "West is also very unlikely to be put at the top because west is where the sun disappears."
Confusingly, early Chinese maps seem to buck this trend. But, Brotton, says, even though they did have compasses at the time, that isn't the reason that they placed north at the top. Early Chinese compasses were actually oriented to point south, which was considered to be more desirable than deepest darkest north. But in Chinese maps, the Emperor, who lived in the north of the country was always put at the top of the map, with everyone else, his loyal subjects, looking up towards him. "In Chinese culture the Emperor looks south because it's where the winds come from, it's a good direction. North is not very good but you are in a position of subjection to the emperor, so you look up to him," says Brotton.
Given that each culture has a very different idea of who, or what, they should look up to it's perhaps not surprising that there is very little consistency in which way early maps pointed. In ancient Egyptian times the top of the world was east, the position of sunrise. Early Islamic maps favoured south at the top because most of the early Muslim cultures were north of Mecca, so they imagined looking up (south) towards it. Christian maps from the same era (called Mappa Mundi) put east at the top, towards the Garden of Eden and with Jerusalem in the centre.
So when did everyone get together and decide that north was the top? It's tempting to put it down to European explorers like Christopher Columbus and Ferdinand Megellan, who were navigating by the North Star. But Brotton argues that these early explorers didn't think of the world like that at all. "When Columbus describes the world it is in accordance with east being at the top, he says. "Columbus says he is going towards paradise, so his mentality is from a medieval mappa mundi." We've got to remember, adds Brotton, that at the time, "no one knows what they are doing and where they are going."

\vspace{1em}
\textbf{Question 11}

The author attributes the French and American revolutions to the invention of the printing press because

\begin{enumerate}[label=\Alph*.]
    \item maps enabled large numbers of Europeans to travel and settle in the American continent.
    \item the rapid spread of information exposed people to new ideas on freedom and democracy.
    \item it encouraged religious freedom among the people by destroying the monopoly of religious leaders on the scriptures.
    \item it made available revolutionary strategies and opinions to the people.
[](https://cracku.in/11-the-author-attributes-the-french-and-american-revo-x-cat-2017-shift-1)
\end{enumerate}
\vspace{1em}
\textbf{Solution:}

Refer to the line "it's hard to imagine the French or American revolutions without those enlightened voices in print". Hence, from this we can straightaway eliminate options A and C. Between options B and D, B correctly captures the point made by the author. The printing press allowed the spread of enlightened voices and as a result people were exposed to new ideas on freedom and democracy.   

Option D slightly distorts what is given in the passage. The passage does not mention any revolutionary "strategies". Hence, we can eliminate this option.
Thus, the answer is option B.

\end{problem}

\newpage
\begin{problem}{}{}
% Subject: varc
\textbf{Instructions}

Instructions   

**The passage below is accompanied by a set of six questions. Choose the best answer to each question.**
Understanding where you are in the world is a basic survival skill, which is why we, like most species come hard-wired with specialised brain areas to create cognitive maps of our surroundings. Where humans are unique, though, with the possible exception of honeybees, is that we try to communicate this understanding of the world with others. We have a long history of doing this by drawing maps — the earliest versions yet discovered were scrawled on cave walls 14,000 years ago. Human cultures have been drawing them on stone tablets, papyrus, paper and now computer screens ever since.
Given such a long history of human map-making, it is perhaps surprising that it is only within the last few hundred years that north has been consistently considered to be at the top. In fact, for much of human history, north almost never appeared at the top, according to Jerry Brotton, a map historian... "North was rarely put at the top for the simple fact that north is where darkness comes from," he says. "West is also very unlikely to be put at the top because west is where the sun disappears."
Confusingly, early Chinese maps seem to buck this trend. But, Brotton, says, even though they did have compasses at the time, that isn't the reason that they placed north at the top. Early Chinese compasses were actually oriented to point south, which was considered to be more desirable than deepest darkest north. But in Chinese maps, the Emperor, who lived in the north of the country was always put at the top of the map, with everyone else, his loyal subjects, looking up towards him. "In Chinese culture the Emperor looks south because it's where the winds come from, it's a good direction. North is not very good but you are in a position of subjection to the emperor, so you look up to him," says Brotton.
Given that each culture has a very different idea of who, or what, they should look up to it's perhaps not surprising that there is very little consistency in which way early maps pointed. In ancient Egyptian times the top of the world was east, the position of sunrise. Early Islamic maps favoured south at the top because most of the early Muslim cultures were north of Mecca, so they imagined looking up (south) towards it. Christian maps from the same era (called Mappa Mundi) put east at the top, towards the Garden of Eden and with Jerusalem in the centre.
So when did everyone get together and decide that north was the top? It's tempting to put it down to European explorers like Christopher Columbus and Ferdinand Megellan, who were navigating by the North Star. But Brotton argues that these early explorers didn't think of the world like that at all. "When Columbus describes the world it is in accordance with east being at the top, he says. "Columbus says he is going towards paradise, so his mentality is from a medieval mappa mundi." We've got to remember, adds Brotton, that at the time, "no one knows what they are doing and where they are going."

\vspace{1em}
\textbf{Question 12}

The main conclusion of the passage is that the new technology has

\begin{enumerate}[label=\Alph*.]
    \item some advantages, but these are outweighed by its disadvantages.
    \item so far not proved as successful as the printing press in opening people's minds.
    \item been disappointing because it has changed society too rapidly.
    \item been more wasteful than the printing press because people spend more time daydreaming or surfing.
[](https://cracku.in/12-the-main-conclusion-of-the-passage-is-that-the-new-x-cat-2017-shift-1)
\end{enumerate}
\vspace{1em}
\textbf{Solution:}

The main point of the passage is that unlike the Gutenberg printing press, the iPhone has in comparison done nothing to make the society more liberated or enlightened. This point has been accurately captured by option B.
The author is not weighing the advantages or disadvantages of new technology. Hence, we can eliminate option A. 
The author does not say that the society has rapidly changed as a result of new technology. In fact, he says that nothing really has changed as a result of it.
The author says that people are no longer daydreaming as a result of new technology. Hence, option D, which contradicts what is given in the passage, can be eliminated.

Instructions   

**The passage below is accompanied by a set of six questions. Choose the best answer to each question.**
This year alone, more than 8,600 stores could close, according to industry estimates, many of them the brand-name anchor outlets that real estate developers once stumbled over themselves to court. Already there have been 5,300 retail closings this year... Sears Holdings — which owns Kmart — said in March that there's "substantial doubt" it can stay in business altogether, and will close 300 stores this year. So far this year, nine national retail chains have filed for bankruptcy.
Local jobs are a major casualty of what analysts are calling, with only a hint of hyperbole, the retail apocalypse. Since 2002, department stores have lost 448,000 jobs, a 25% decline, while the number of store closures this year is on pace to surpass the worst depths of the Great Recession. The growth of online retailers, meanwhile, has failed to offset those losses, with the e-commerce sector adding just 178,000 jobs over the past 15 years. Some of those jobs can be found in the massive distribution centers Amazon has opened across the country, often not too far from malls the company helped shutter.
But those are workplaces, not gathering places. The mall is both. And in the 61 years since the first enclosed one opened in suburban Minneapolis, the shopping mall has been where a huge swath of middle-class America went for far more than shopping. It was the home of first jobs and blind dates, the place for family photos and ear piercings, where goths and grandmothers could somehow walk through the same doors and find something they all liked. Sure, the food was lousy for you and the oceans of parking lots encouraged car- heavy development, something now scorned by contemporary planners. But for better or worse, the mall has been America's public square for the last 60 years.
So what happens when it disappears?
Think of your mall. Or think of the one you went to as a kid. Think of the perfume clouds in the department stores. The fountains splashing below the skylights. The cinnamon wafting from the food court. As far back as ancient Greece, societies have congregated around a central marketplace. In medieval Europe, they were outside cathedrals. For half of the 20th century and almost 20 years into the new one, much of America has found their agora on the terrazzo between Orange Julius and Sbarro, Waldenbooks and the Gap, Sunglass Hut and Hot Topic.
That mall was an ecosystem unto itself, a combination. of community and commercialism peddling everything you needed and everything you didn' t: Magic Eye posters, wind catchers, Air Jordans....
A growing number of Americans, however, don't see the need to go to any Macy's at all. Our digital lives are frictionless and ruthlessly efficient, with retail and romance available at a click. Malls were designed for leisure, abundance, ambling. You parked and planned to spend some time. Today, much of that time has been given over to busier lives and second jobs and apps that let you swipe right instead of haunt the food court. Malls, says Harvard business professor Leonard Schlesinger, "were built for patterns of social interaction that increasingly don't exist."

\end{problem}

\newpage
\begin{problem}{}{}
% Subject: varc
\textbf{Instructions}

Instructions   

**The passage below is accompanied by a set of six questions. Choose the best answer to each question.**
Understanding where you are in the world is a basic survival skill, which is why we, like most species come hard-wired with specialised brain areas to create cognitive maps of our surroundings. Where humans are unique, though, with the possible exception of honeybees, is that we try to communicate this understanding of the world with others. We have a long history of doing this by drawing maps — the earliest versions yet discovered were scrawled on cave walls 14,000 years ago. Human cultures have been drawing them on stone tablets, papyrus, paper and now computer screens ever since.
Given such a long history of human map-making, it is perhaps surprising that it is only within the last few hundred years that north has been consistently considered to be at the top. In fact, for much of human history, north almost never appeared at the top, according to Jerry Brotton, a map historian... "North was rarely put at the top for the simple fact that north is where darkness comes from," he says. "West is also very unlikely to be put at the top because west is where the sun disappears."
Confusingly, early Chinese maps seem to buck this trend. But, Brotton, says, even though they did have compasses at the time, that isn't the reason that they placed north at the top. Early Chinese compasses were actually oriented to point south, which was considered to be more desirable than deepest darkest north. But in Chinese maps, the Emperor, who lived in the north of the country was always put at the top of the map, with everyone else, his loyal subjects, looking up towards him. "In Chinese culture the Emperor looks south because it's where the winds come from, it's a good direction. North is not very good but you are in a position of subjection to the emperor, so you look up to him," says Brotton.
Given that each culture has a very different idea of who, or what, they should look up to it's perhaps not surprising that there is very little consistency in which way early maps pointed. In ancient Egyptian times the top of the world was east, the position of sunrise. Early Islamic maps favoured south at the top because most of the early Muslim cultures were north of Mecca, so they imagined looking up (south) towards it. Christian maps from the same era (called Mappa Mundi) put east at the top, towards the Garden of Eden and with Jerusalem in the centre.
So when did everyone get together and decide that north was the top? It's tempting to put it down to European explorers like Christopher Columbus and Ferdinand Megellan, who were navigating by the North Star. But Brotton argues that these early explorers didn't think of the world like that at all. "When Columbus describes the world it is in accordance with east being at the top, he says. "Columbus says he is going towards paradise, so his mentality is from a medieval mappa mundi." We've got to remember, adds Brotton, that at the time, "no one knows what they are doing and where they are going."

\vspace{1em}
\textbf{Question 13}

The central idea of this passage is that:

\begin{enumerate}[label=\Alph*.]
    \item the closure of malls has affected the economic and social life of middle-class America.
    \item Is the advantages of malls outweigh their disadvantages.
    \item malls used to perform a social function that has been lost.
    \item malls are closing down because people have found alternate ways to shop.
[](https://cracku.in/13-the-central-idea-of-this-passage-is-that-x-cat-2017-shift-1)
\end{enumerate}
\vspace{1em}
\textbf{Solution:}

The author argues that malls were more than just shopping places. Towards the end of the passage, the author tries to invoke a sense of nostalgia by stating how malls used to play an important role in everyone's life and how that function is getting lost.   
Let us evaluate the options.  
  
Option B states that the advantages of the malls outweigh the disadvantages. The author has not mentioned anything about the disadvantages of the malls. Therefore, we can easily eliminate option B.   
  
Option D states that malls are closing down since people have found an alternate way to shop. Though the option is true, the main point of the author is not that malls are closing down. The author is concerned about the fact that malls were places of congregation and the closure of malls takes a social function away with them. Therefore, we can eliminate option D as well.   
  
Options A and C are close. Option A states that the closure of malls has affected the social and economic life of America. However, the author states that the closure of malls is reflective of the changing social structure of America. Option A gets the relationship backwards and hence, option A can be eliminated.   
  
Option C states that malls used to perform a social function that has been lost. This seems to be the main point that the author is trying to emphasize through the paragraph. The author states how malls used to be places of congregation and how the new generation finds no need to go to malls. Therefore, option C is the right answer.

\end{problem}

\newpage
\begin{problem}{}{}
% Subject: varc
\textbf{Instructions}

Instructions   

**The passage below is accompanied by a set of six questions. Choose the best answer to each question.**
Understanding where you are in the world is a basic survival skill, which is why we, like most species come hard-wired with specialised brain areas to create cognitive maps of our surroundings. Where humans are unique, though, with the possible exception of honeybees, is that we try to communicate this understanding of the world with others. We have a long history of doing this by drawing maps — the earliest versions yet discovered were scrawled on cave walls 14,000 years ago. Human cultures have been drawing them on stone tablets, papyrus, paper and now computer screens ever since.
Given such a long history of human map-making, it is perhaps surprising that it is only within the last few hundred years that north has been consistently considered to be at the top. In fact, for much of human history, north almost never appeared at the top, according to Jerry Brotton, a map historian... "North was rarely put at the top for the simple fact that north is where darkness comes from," he says. "West is also very unlikely to be put at the top because west is where the sun disappears."
Confusingly, early Chinese maps seem to buck this trend. But, Brotton, says, even though they did have compasses at the time, that isn't the reason that they placed north at the top. Early Chinese compasses were actually oriented to point south, which was considered to be more desirable than deepest darkest north. But in Chinese maps, the Emperor, who lived in the north of the country was always put at the top of the map, with everyone else, his loyal subjects, looking up towards him. "In Chinese culture the Emperor looks south because it's where the winds come from, it's a good direction. North is not very good but you are in a position of subjection to the emperor, so you look up to him," says Brotton.
Given that each culture has a very different idea of who, or what, they should look up to it's perhaps not surprising that there is very little consistency in which way early maps pointed. In ancient Egyptian times the top of the world was east, the position of sunrise. Early Islamic maps favoured south at the top because most of the early Muslim cultures were north of Mecca, so they imagined looking up (south) towards it. Christian maps from the same era (called Mappa Mundi) put east at the top, towards the Garden of Eden and with Jerusalem in the centre.
So when did everyone get together and decide that north was the top? It's tempting to put it down to European explorers like Christopher Columbus and Ferdinand Megellan, who were navigating by the North Star. But Brotton argues that these early explorers didn't think of the world like that at all. "When Columbus describes the world it is in accordance with east being at the top, he says. "Columbus says he is going towards paradise, so his mentality is from a medieval mappa mundi." We've got to remember, adds Brotton, that at the time, "no one knows what they are doing and where they are going."

\vspace{1em}
\textbf{Question 14}

Why does the author say in paragraph 2, 'the massive distribution centers Amazon has opened across the country, often not too far from malls the company helped shutter'?

\begin{enumerate}[label=\Alph*.]
    \item To highlight the irony of the situation.
    \item To indicate that malls and distribution centres are located in the same area.
    \item To show that Amazon is helping certain brands go online.
    \item To indicate that the shopping habits of the American middle class have changed.
[](https://cracku.in/14-why-does-the-author-say-in-paragraph-2-the-massive-x-cat-2017-shift-1)
\end{enumerate}
\vspace{1em}
\textbf{Solution:}

No where has it been mentioned that Amazon is helping brands go online. Therefore, we can eliminate option C.   
  
Option D states that the purpose of the line is to indicate that the shopping habits of the middle class America has changed. However, the author talks about shopping habits towards the end of the passage. The author has not introduced the topic of 'shopping habits' when the given line has been mentioned. Therefore, we can eliminate option D as well.   
  
Option B states that the author uses the line to indicate that the malls and distribution centres are located in the same area. However, the author states that the distribution centres are replacing the malls. He does not intend to convey that they co-exist. Therefore, we can eliminate option B as well.   
  
Option A states that the author uses the line to indicate the irony of the situation. Option A captures the fact that the author finds it ironic that distribution centres are springing up near the places where the malls once existed. Therefore, option A is the right answer.

\end{problem}

\newpage
\begin{problem}{}{}
% Subject: varc
\textbf{Instructions}

Instructions   

**The passage below is accompanied by a set of six questions. Choose the best answer to each question.**
Understanding where you are in the world is a basic survival skill, which is why we, like most species come hard-wired with specialised brain areas to create cognitive maps of our surroundings. Where humans are unique, though, with the possible exception of honeybees, is that we try to communicate this understanding of the world with others. We have a long history of doing this by drawing maps — the earliest versions yet discovered were scrawled on cave walls 14,000 years ago. Human cultures have been drawing them on stone tablets, papyrus, paper and now computer screens ever since.
Given such a long history of human map-making, it is perhaps surprising that it is only within the last few hundred years that north has been consistently considered to be at the top. In fact, for much of human history, north almost never appeared at the top, according to Jerry Brotton, a map historian... "North was rarely put at the top for the simple fact that north is where darkness comes from," he says. "West is also very unlikely to be put at the top because west is where the sun disappears."
Confusingly, early Chinese maps seem to buck this trend. But, Brotton, says, even though they did have compasses at the time, that isn't the reason that they placed north at the top. Early Chinese compasses were actually oriented to point south, which was considered to be more desirable than deepest darkest north. But in Chinese maps, the Emperor, who lived in the north of the country was always put at the top of the map, with everyone else, his loyal subjects, looking up towards him. "In Chinese culture the Emperor looks south because it's where the winds come from, it's a good direction. North is not very good but you are in a position of subjection to the emperor, so you look up to him," says Brotton.
Given that each culture has a very different idea of who, or what, they should look up to it's perhaps not surprising that there is very little consistency in which way early maps pointed. In ancient Egyptian times the top of the world was east, the position of sunrise. Early Islamic maps favoured south at the top because most of the early Muslim cultures were north of Mecca, so they imagined looking up (south) towards it. Christian maps from the same era (called Mappa Mundi) put east at the top, towards the Garden of Eden and with Jerusalem in the centre.
So when did everyone get together and decide that north was the top? It's tempting to put it down to European explorers like Christopher Columbus and Ferdinand Megellan, who were navigating by the North Star. But Brotton argues that these early explorers didn't think of the world like that at all. "When Columbus describes the world it is in accordance with east being at the top, he says. "Columbus says he is going towards paradise, so his mentality is from a medieval mappa mundi." We've got to remember, adds Brotton, that at the time, "no one knows what they are doing and where they are going."

\vspace{1em}
\textbf{Question 15}

In the first paragraph, the phrase "real estate developers once stumbled over themselves to court..." suggests that they

\begin{enumerate}[label=\Alph*.]
    \item took brand-name anchor outlets to court.
    \item no longer pursue brand-name anchor outlets.
    \item collaborated with one another to get brand-name anchor outlets.
    \item were eager to get brand-name anchor outlets to set up shop in their mall.
[](https://cracku.in/15-in-the-first-paragraph-the-phrase-real-estate-deve-x-cat-2017-shift-1)
\end{enumerate}
\vspace{1em}
\textbf{Solution:}

In the first paragraph, the author states that "many of them the brand-name anchor outlets that real estate developers once stumbled over themselves to court", indicating that the brand outlets used to be sought after once but are no longer in demand. The author states that the real estate developers used to fight each other over these outlets earlier. Therefore, we can infer that the real estate developers no longer show the kind of enthusiasm they did (since the outlets are closing down) and hence, option B is the right answer.

\end{problem}

\newpage
\begin{problem}{}{}
% Subject: varc
\textbf{Instructions}

Instructions   

**The passage below is accompanied by a set of six questions. Choose the best answer to each question.**
Understanding where you are in the world is a basic survival skill, which is why we, like most species come hard-wired with specialised brain areas to create cognitive maps of our surroundings. Where humans are unique, though, with the possible exception of honeybees, is that we try to communicate this understanding of the world with others. We have a long history of doing this by drawing maps — the earliest versions yet discovered were scrawled on cave walls 14,000 years ago. Human cultures have been drawing them on stone tablets, papyrus, paper and now computer screens ever since.
Given such a long history of human map-making, it is perhaps surprising that it is only within the last few hundred years that north has been consistently considered to be at the top. In fact, for much of human history, north almost never appeared at the top, according to Jerry Brotton, a map historian... "North was rarely put at the top for the simple fact that north is where darkness comes from," he says. "West is also very unlikely to be put at the top because west is where the sun disappears."
Confusingly, early Chinese maps seem to buck this trend. But, Brotton, says, even though they did have compasses at the time, that isn't the reason that they placed north at the top. Early Chinese compasses were actually oriented to point south, which was considered to be more desirable than deepest darkest north. But in Chinese maps, the Emperor, who lived in the north of the country was always put at the top of the map, with everyone else, his loyal subjects, looking up towards him. "In Chinese culture the Emperor looks south because it's where the winds come from, it's a good direction. North is not very good but you are in a position of subjection to the emperor, so you look up to him," says Brotton.
Given that each culture has a very different idea of who, or what, they should look up to it's perhaps not surprising that there is very little consistency in which way early maps pointed. In ancient Egyptian times the top of the world was east, the position of sunrise. Early Islamic maps favoured south at the top because most of the early Muslim cultures were north of Mecca, so they imagined looking up (south) towards it. Christian maps from the same era (called Mappa Mundi) put east at the top, towards the Garden of Eden and with Jerusalem in the centre.
So when did everyone get together and decide that north was the top? It's tempting to put it down to European explorers like Christopher Columbus and Ferdinand Megellan, who were navigating by the North Star. But Brotton argues that these early explorers didn't think of the world like that at all. "When Columbus describes the world it is in accordance with east being at the top, he says. "Columbus says he is going towards paradise, so his mentality is from a medieval mappa mundi." We've got to remember, adds Brotton, that at the time, "no one knows what they are doing and where they are going."

\vspace{1em}
\textbf{Question 16}

The author calls the mall an ecosystem unto itself because

\begin{enumerate}[label=\Alph*.]
    \item people of all ages and from all walks of life went there.
    \item people could shop as well as eat in one place.
    \item it was a commercial space as well as a gathering place.
    \item it sold things that were needed as well as those that were not.
[](https://cracku.in/16-the-author-calls-the-mall-an-ecosystem-unto-itself-x-cat-2017-shift-1)
\end{enumerate}
\vspace{1em}
\textbf{Solution:}

The author calls malls an ecosystem since they served as places of congregation. The author compares the market places from ancient Greece to emphasize how malls, just like the market places through out the history, served as the places of social interaction as well. Only option C captures the fact that malls serves both as commercial and social centres and hence, option C is the right answer.

\end{problem}

\newpage
\begin{problem}{}{}
% Subject: varc
\textbf{Instructions}

Instructions   

**The passage below is accompanied by a set of six questions. Choose the best answer to each question.**
Understanding where you are in the world is a basic survival skill, which is why we, like most species come hard-wired with specialised brain areas to create cognitive maps of our surroundings. Where humans are unique, though, with the possible exception of honeybees, is that we try to communicate this understanding of the world with others. We have a long history of doing this by drawing maps — the earliest versions yet discovered were scrawled on cave walls 14,000 years ago. Human cultures have been drawing them on stone tablets, papyrus, paper and now computer screens ever since.
Given such a long history of human map-making, it is perhaps surprising that it is only within the last few hundred years that north has been consistently considered to be at the top. In fact, for much of human history, north almost never appeared at the top, according to Jerry Brotton, a map historian... "North was rarely put at the top for the simple fact that north is where darkness comes from," he says. "West is also very unlikely to be put at the top because west is where the sun disappears."
Confusingly, early Chinese maps seem to buck this trend. But, Brotton, says, even though they did have compasses at the time, that isn't the reason that they placed north at the top. Early Chinese compasses were actually oriented to point south, which was considered to be more desirable than deepest darkest north. But in Chinese maps, the Emperor, who lived in the north of the country was always put at the top of the map, with everyone else, his loyal subjects, looking up towards him. "In Chinese culture the Emperor looks south because it's where the winds come from, it's a good direction. North is not very good but you are in a position of subjection to the emperor, so you look up to him," says Brotton.
Given that each culture has a very different idea of who, or what, they should look up to it's perhaps not surprising that there is very little consistency in which way early maps pointed. In ancient Egyptian times the top of the world was east, the position of sunrise. Early Islamic maps favoured south at the top because most of the early Muslim cultures were north of Mecca, so they imagined looking up (south) towards it. Christian maps from the same era (called Mappa Mundi) put east at the top, towards the Garden of Eden and with Jerusalem in the centre.
So when did everyone get together and decide that north was the top? It's tempting to put it down to European explorers like Christopher Columbus and Ferdinand Megellan, who were navigating by the North Star. But Brotton argues that these early explorers didn't think of the world like that at all. "When Columbus describes the world it is in accordance with east being at the top, he says. "Columbus says he is going towards paradise, so his mentality is from a medieval mappa mundi." We've got to remember, adds Brotton, that at the time, "no one knows what they are doing and where they are going."

\vspace{1em}
\textbf{Question 17}

Why does the author say that the mall has been America's public square?

\begin{enumerate}[label=\Alph*.]
    \item Malls did not bar anybody from entering the space.
    \item Malls were a great place to shop for a huge section of the middle class.
    \item Malls were a hangout place where families grew close to each other.
    \item Malls were a great place for everyone to gather and interact.
[](https://cracku.in/17-why-does-the-author-say-that-the-mall-has-been-ame-x-cat-2017-shift-1)
\end{enumerate}
\vspace{1em}
\textbf{Solution:}

The main point that the author places is that malls served as places of social interaction and cannot be reduced to merely places of commercial activity. Therefore, the author compares the malls as America's public square to show the bustling social life that the malls hosted. Therefore, option D is the right answer.

\end{problem}

\newpage
\begin{problem}{}{}
% Subject: varc
\textbf{Instructions}

Instructions   

**The passage below is accompanied by a set of six questions. Choose the best answer to each question.**
Understanding where you are in the world is a basic survival skill, which is why we, like most species come hard-wired with specialised brain areas to create cognitive maps of our surroundings. Where humans are unique, though, with the possible exception of honeybees, is that we try to communicate this understanding of the world with others. We have a long history of doing this by drawing maps — the earliest versions yet discovered were scrawled on cave walls 14,000 years ago. Human cultures have been drawing them on stone tablets, papyrus, paper and now computer screens ever since.
Given such a long history of human map-making, it is perhaps surprising that it is only within the last few hundred years that north has been consistently considered to be at the top. In fact, for much of human history, north almost never appeared at the top, according to Jerry Brotton, a map historian... "North was rarely put at the top for the simple fact that north is where darkness comes from," he says. "West is also very unlikely to be put at the top because west is where the sun disappears."
Confusingly, early Chinese maps seem to buck this trend. But, Brotton, says, even though they did have compasses at the time, that isn't the reason that they placed north at the top. Early Chinese compasses were actually oriented to point south, which was considered to be more desirable than deepest darkest north. But in Chinese maps, the Emperor, who lived in the north of the country was always put at the top of the map, with everyone else, his loyal subjects, looking up towards him. "In Chinese culture the Emperor looks south because it's where the winds come from, it's a good direction. North is not very good but you are in a position of subjection to the emperor, so you look up to him," says Brotton.
Given that each culture has a very different idea of who, or what, they should look up to it's perhaps not surprising that there is very little consistency in which way early maps pointed. In ancient Egyptian times the top of the world was east, the position of sunrise. Early Islamic maps favoured south at the top because most of the early Muslim cultures were north of Mecca, so they imagined looking up (south) towards it. Christian maps from the same era (called Mappa Mundi) put east at the top, towards the Garden of Eden and with Jerusalem in the centre.
So when did everyone get together and decide that north was the top? It's tempting to put it down to European explorers like Christopher Columbus and Ferdinand Megellan, who were navigating by the North Star. But Brotton argues that these early explorers didn't think of the world like that at all. "When Columbus describes the world it is in accordance with east being at the top, he says. "Columbus says he is going towards paradise, so his mentality is from a medieval mappa mundi." We've got to remember, adds Brotton, that at the time, "no one knows what they are doing and where they are going."

\vspace{1em}
\textbf{Question 18}

The author describes 'Perfume clouds in the department stores' in order to

\begin{enumerate}[label=\Alph*.]
    \item evoke memories by painting a picture of malls.
    \item describe the smells and sights of malls.
    \item emphasise that all brands were available under one roof.
    \item show that malls smelt good because of the various stores and food court.
[](https://cracku.in/18-the-author-describes-perfume-clouds-in-the-departm-x-cat-2017-shift-1)
\end{enumerate}
\vspace{1em}
\textbf{Solution:}

In the last 2 paragraphs, the author tries to invoke a sense of nostalgia by describing the malls. He tries to make the readers feel the ambiance and vibrance of the malls by taking us down the memory lane. Option B adopts an objective tone. The author does not merely describe a mall. He tries to invoke our memories of the malls. Options C and D can be eliminated since the purpose of the author is not to emphasize that all brands were available at one place or to establish the reason for the sweet smell that the malls carried. The author tries to make an emotional appeal and only option A captures this point. Therefore, option A is the right answer.

Instructions   

**The passage below is accompanied by a set of three questions. Choose the best answer to each question.**
Scientists have long recognised the incredible diversity within a species. But they thought it reflected evolutionary changes that unfolded imperceptibly, over millions of years. That divergence between populations within a species was enforced, according to Ernst Mayr, the great evolutionary biologist of the 1940s, when a population was separated from the rest of the species by a mountain range or a desert, preventing breeding across the divide over geologic scales of time. Without the separation, gene flow was relentless. But as the separation persisted, the isolated population grew apart and speciation occurred.
In the mid-1960s, the biologist Paul Ehrlich — author of The Population Bomb (1968) — and his Stanford University colleague Peter Raven challenged Mayr's ideas about speciation. They had studied checkerspot butterflies living in the Jasper Ridge Biological Preserve in California, and it soon became clear that they were not examining a single population. Through years of capturing, marking and then recapturing the butterflies, they were able to prove that within the population, spread over just 50 acres of suitable checkerspot habitat, there were three groups that rarely interacted despite their very close proximity.
Among other ideas, Ehrlich and Raven argued in a now classic paper from 1969 that gene flow was not as predictable and ubiquitous as Mayr and his cohort maintained, and thus evolutionary divergence between neighbouring groups in a population was probably common. They also asserted that isolation and gene flow were less important to evolutionary divergence than natural selection (when factors such as mate choice, weather, disease or predation cause better-adapted individuals to survive and pass on their successful genetic traits). For example, Ehrlich and Raven suggested that, without the force of natural selection, an isolated population would remain unchanged and that, in other scenarios, natural selection could be strong enough to overpower gene flow...

\end{problem}

\newpage
\begin{problem}{}{}
% Subject: varc
\textbf{Instructions}

Instructions   

**The passage below is accompanied by a set of six questions. Choose the best answer to each question.**
Understanding where you are in the world is a basic survival skill, which is why we, like most species come hard-wired with specialised brain areas to create cognitive maps of our surroundings. Where humans are unique, though, with the possible exception of honeybees, is that we try to communicate this understanding of the world with others. We have a long history of doing this by drawing maps — the earliest versions yet discovered were scrawled on cave walls 14,000 years ago. Human cultures have been drawing them on stone tablets, papyrus, paper and now computer screens ever since.
Given such a long history of human map-making, it is perhaps surprising that it is only within the last few hundred years that north has been consistently considered to be at the top. In fact, for much of human history, north almost never appeared at the top, according to Jerry Brotton, a map historian... "North was rarely put at the top for the simple fact that north is where darkness comes from," he says. "West is also very unlikely to be put at the top because west is where the sun disappears."
Confusingly, early Chinese maps seem to buck this trend. But, Brotton, says, even though they did have compasses at the time, that isn't the reason that they placed north at the top. Early Chinese compasses were actually oriented to point south, which was considered to be more desirable than deepest darkest north. But in Chinese maps, the Emperor, who lived in the north of the country was always put at the top of the map, with everyone else, his loyal subjects, looking up towards him. "In Chinese culture the Emperor looks south because it's where the winds come from, it's a good direction. North is not very good but you are in a position of subjection to the emperor, so you look up to him," says Brotton.
Given that each culture has a very different idea of who, or what, they should look up to it's perhaps not surprising that there is very little consistency in which way early maps pointed. In ancient Egyptian times the top of the world was east, the position of sunrise. Early Islamic maps favoured south at the top because most of the early Muslim cultures were north of Mecca, so they imagined looking up (south) towards it. Christian maps from the same era (called Mappa Mundi) put east at the top, towards the Garden of Eden and with Jerusalem in the centre.
So when did everyone get together and decide that north was the top? It's tempting to put it down to European explorers like Christopher Columbus and Ferdinand Megellan, who were navigating by the North Star. But Brotton argues that these early explorers didn't think of the world like that at all. "When Columbus describes the world it is in accordance with east being at the top, he says. "Columbus says he is going towards paradise, so his mentality is from a medieval mappa mundi." We've got to remember, adds Brotton, that at the time, "no one knows what they are doing and where they are going."

\vspace{1em}
\textbf{Question 19}

Which of the following best sums up Ehrlich and Raven's argument in their classic 1969 paper?

\begin{enumerate}[label=\Alph*.]
    \item Ernst Mayr was wrong in identifying physical separation as the cause of species diversity.
    \item Checkerspot butterflies in the 50-acre Jasper Ridge Preserve formed three groups that rarely interacted with each other.
    \item While a factor, isolation was not as important to speciation as natural selection.
    \item Gene flow is less common and more erratic than Mayr and his colleagues claimed.
[](https://cracku.in/19-which-of-the-following-best-sums-up-ehrlich-and-ra-x-cat-2017-shift-1)
\end{enumerate}
\vspace{1em}
\textbf{Solution:}

In the last paragraph, the author uses the line "They also asserted that isolation and gene flow were less important to evolutionary divergence than natural selection". Also, through the example of checkerspot butterflies, the author brings to light the fact that isolation, though a factor, is not very dominant in influencing speciation. Therefore, option C is the right answer.

\end{problem}

\newpage
\begin{problem}{}{}
% Subject: varc
\textbf{Instructions}

Instructions   

**The passage below is accompanied by a set of six questions. Choose the best answer to each question.**
Understanding where you are in the world is a basic survival skill, which is why we, like most species come hard-wired with specialised brain areas to create cognitive maps of our surroundings. Where humans are unique, though, with the possible exception of honeybees, is that we try to communicate this understanding of the world with others. We have a long history of doing this by drawing maps — the earliest versions yet discovered were scrawled on cave walls 14,000 years ago. Human cultures have been drawing them on stone tablets, papyrus, paper and now computer screens ever since.
Given such a long history of human map-making, it is perhaps surprising that it is only within the last few hundred years that north has been consistently considered to be at the top. In fact, for much of human history, north almost never appeared at the top, according to Jerry Brotton, a map historian... "North was rarely put at the top for the simple fact that north is where darkness comes from," he says. "West is also very unlikely to be put at the top because west is where the sun disappears."
Confusingly, early Chinese maps seem to buck this trend. But, Brotton, says, even though they did have compasses at the time, that isn't the reason that they placed north at the top. Early Chinese compasses were actually oriented to point south, which was considered to be more desirable than deepest darkest north. But in Chinese maps, the Emperor, who lived in the north of the country was always put at the top of the map, with everyone else, his loyal subjects, looking up towards him. "In Chinese culture the Emperor looks south because it's where the winds come from, it's a good direction. North is not very good but you are in a position of subjection to the emperor, so you look up to him," says Brotton.
Given that each culture has a very different idea of who, or what, they should look up to it's perhaps not surprising that there is very little consistency in which way early maps pointed. In ancient Egyptian times the top of the world was east, the position of sunrise. Early Islamic maps favoured south at the top because most of the early Muslim cultures were north of Mecca, so they imagined looking up (south) towards it. Christian maps from the same era (called Mappa Mundi) put east at the top, towards the Garden of Eden and with Jerusalem in the centre.
So when did everyone get together and decide that north was the top? It's tempting to put it down to European explorers like Christopher Columbus and Ferdinand Megellan, who were navigating by the North Star. But Brotton argues that these early explorers didn't think of the world like that at all. "When Columbus describes the world it is in accordance with east being at the top, he says. "Columbus says he is going towards paradise, so his mentality is from a medieval mappa mundi." We've got to remember, adds Brotton, that at the time, "no one knows what they are doing and where they are going."

\vspace{1em}
\textbf{Question 20}

All of the following statements are true according to the passage EXCEPT

\begin{enumerate}[label=\Alph*.]
    \item Gene flow contributes to evolutionary divergence.
    \item The Population Bomb questioned dominant ideas about species diversity.
    \item Evolutionary changes unfold imperceptibly over time.
    \item Checkerspot butterflies are known to exhibit speciation while living in close proximity.
[](https://cracku.in/20-all-of-the-following-statements-are-true-according-x-cat-2017-shift-1)
\end{enumerate}
\vspace{1em}
\textbf{Solution:}

In the second line of the first paragraph, the author mentions that evolutionary changes evolve imperceptibly over time.   
In the third paragraph, the author mentions that "Isolation and gene flow were less important to evolutionary divergence than natural selection". Therefore, we can infer that gene flow contributes to evolutionary divergence.  
The author explains that 3 species of checkerspot butterflies living within 50 acres hardly interact with each other. Therefore, we can infer option D. 
The author mentioned the book "the population bomb" as an introduction to the biologist Paul Ehrlich.  
Option B cannot be inferred from the passage and hence, it is the right answer.

\end{problem}

\newpage
\begin{problem}{}{}
% Subject: varc
\textbf{Instructions}

Instructions   

**The passage below is accompanied by a set of six questions. Choose the best answer to each question.**
Understanding where you are in the world is a basic survival skill, which is why we, like most species come hard-wired with specialised brain areas to create cognitive maps of our surroundings. Where humans are unique, though, with the possible exception of honeybees, is that we try to communicate this understanding of the world with others. We have a long history of doing this by drawing maps — the earliest versions yet discovered were scrawled on cave walls 14,000 years ago. Human cultures have been drawing them on stone tablets, papyrus, paper and now computer screens ever since.
Given such a long history of human map-making, it is perhaps surprising that it is only within the last few hundred years that north has been consistently considered to be at the top. In fact, for much of human history, north almost never appeared at the top, according to Jerry Brotton, a map historian... "North was rarely put at the top for the simple fact that north is where darkness comes from," he says. "West is also very unlikely to be put at the top because west is where the sun disappears."
Confusingly, early Chinese maps seem to buck this trend. But, Brotton, says, even though they did have compasses at the time, that isn't the reason that they placed north at the top. Early Chinese compasses were actually oriented to point south, which was considered to be more desirable than deepest darkest north. But in Chinese maps, the Emperor, who lived in the north of the country was always put at the top of the map, with everyone else, his loyal subjects, looking up towards him. "In Chinese culture the Emperor looks south because it's where the winds come from, it's a good direction. North is not very good but you are in a position of subjection to the emperor, so you look up to him," says Brotton.
Given that each culture has a very different idea of who, or what, they should look up to it's perhaps not surprising that there is very little consistency in which way early maps pointed. In ancient Egyptian times the top of the world was east, the position of sunrise. Early Islamic maps favoured south at the top because most of the early Muslim cultures were north of Mecca, so they imagined looking up (south) towards it. Christian maps from the same era (called Mappa Mundi) put east at the top, towards the Garden of Eden and with Jerusalem in the centre.
So when did everyone get together and decide that north was the top? It's tempting to put it down to European explorers like Christopher Columbus and Ferdinand Megellan, who were navigating by the North Star. But Brotton argues that these early explorers didn't think of the world like that at all. "When Columbus describes the world it is in accordance with east being at the top, he says. "Columbus says he is going towards paradise, so his mentality is from a medieval mappa mundi." We've got to remember, adds Brotton, that at the time, "no one knows what they are doing and where they are going."

\vspace{1em}
\textbf{Question 21}

The author discusses Mayr, Ehrlich and Raven to demonstrate that

\begin{enumerate}[label=\Alph*.]
    \item evolution is a sensitive and controversial topic.
    \item Ehrlich and Raven's ideas about evolutionary divergence are widely accepted by scientists.
    \item the causes of speciation are debated by scientists.
    \item checkerspot butterflies offer the best example of Ehrlich and Raven's ideas about speciation.
[](https://cracku.in/21-the-author-discusses-mayr-ehrlich-and-raven-to-dem-x-cat-2017-shift-1)
\end{enumerate}
\vspace{1em}
\textbf{Solution:}

The author provides 'Checkerspot butterflies' as an example to drive home his point. The primary intention of the author is not to discuss checkerspot butterflies. Therefore, we can eliminate option D. Option B can be eliminated as well since the author has not mentioned that  
the theories are widely accepted.   
  
The author explains the contrasting views of the scientists to show that speciation is a debated topic. Option A states that evolution is a controversial topic. The passage deals with speciation. We cannot generalize speciation to evolution. Also, the intention of the author is to show the differing views among the scientists rather than to establish that the topic is controversial. Therefore, option C is the right answer.

Instructions   

**The passage below is accompanied by a set of three questions. Choose the best answer to each question.**
Do sports mega events like the summer Olympic Games benefit the host city economically? It depends, but the prospects are less than rosy. The trick is converting...several billion dollars in operating costs during the 17-day fiesta of the Games into a basis for long-term economic returns. These days, the summer Olympic Games themselves generate total revenue of $4 billion to $5 billion, but the lion's share of this goes to the International Olympics Committee, the National Olympics Committees and the International Sports Federations. Any economic benefit would have to flow from the value of the Games as an advertisement for the city, the new transportation and communications infrastructure that was created for the Games, or the ongoing use of the new facilities.
Evidence suggests that the advertising effect is far from certain. The infrastructure benefit depends on the initial condition of the city and the effectiveness of the planning. The facilities benefit is dubious at best for buildings such as velodromes or natatoriums and problematic for 100,000-seat Olympic stadiums. The latter require a conversion plan for future use, the former are usually doomed to near vacancy. Hosting the summer Games generally requires 30-plus sports venues and dozens of training centers. Today, the Bird's Nest in Beijing sits virtually empty, while the Olympic Stadium in Sydney costs some $30 million a year to operate.
Part of the problem is that Olympics planning takes place in a frenzied and time-pressured atmosphere of intense competition with the other prospective host cities — not optimal conditions for contemplating the future shape of an urban landscape. Another part of the problem is that urban land is generally scarce and growing scarcer. The new facilities often stand for decades or longer. Even if they have future use, are they the best use of precious urban real estate?
Further, cities must consider the human cost. Residential areas often are razed and citizens relocated (without adequate preparation or compensation). Life is made more hectic and congested. There are, after all, other productive uses that can be made of vanishing fiscal resources.

\end{problem}

\newpage
\begin{problem}{}{}
% Subject: varc
\textbf{Instructions}

Instructions   

**The passage below is accompanied by a set of six questions. Choose the best answer to each question.**
Understanding where you are in the world is a basic survival skill, which is why we, like most species come hard-wired with specialised brain areas to create cognitive maps of our surroundings. Where humans are unique, though, with the possible exception of honeybees, is that we try to communicate this understanding of the world with others. We have a long history of doing this by drawing maps — the earliest versions yet discovered were scrawled on cave walls 14,000 years ago. Human cultures have been drawing them on stone tablets, papyrus, paper and now computer screens ever since.
Given such a long history of human map-making, it is perhaps surprising that it is only within the last few hundred years that north has been consistently considered to be at the top. In fact, for much of human history, north almost never appeared at the top, according to Jerry Brotton, a map historian... "North was rarely put at the top for the simple fact that north is where darkness comes from," he says. "West is also very unlikely to be put at the top because west is where the sun disappears."
Confusingly, early Chinese maps seem to buck this trend. But, Brotton, says, even though they did have compasses at the time, that isn't the reason that they placed north at the top. Early Chinese compasses were actually oriented to point south, which was considered to be more desirable than deepest darkest north. But in Chinese maps, the Emperor, who lived in the north of the country was always put at the top of the map, with everyone else, his loyal subjects, looking up towards him. "In Chinese culture the Emperor looks south because it's where the winds come from, it's a good direction. North is not very good but you are in a position of subjection to the emperor, so you look up to him," says Brotton.
Given that each culture has a very different idea of who, or what, they should look up to it's perhaps not surprising that there is very little consistency in which way early maps pointed. In ancient Egyptian times the top of the world was east, the position of sunrise. Early Islamic maps favoured south at the top because most of the early Muslim cultures were north of Mecca, so they imagined looking up (south) towards it. Christian maps from the same era (called Mappa Mundi) put east at the top, towards the Garden of Eden and with Jerusalem in the centre.
So when did everyone get together and decide that north was the top? It's tempting to put it down to European explorers like Christopher Columbus and Ferdinand Megellan, who were navigating by the North Star. But Brotton argues that these early explorers didn't think of the world like that at all. "When Columbus describes the world it is in accordance with east being at the top, he says. "Columbus says he is going towards paradise, so his mentality is from a medieval mappa mundi." We've got to remember, adds Brotton, that at the time, "no one knows what they are doing and where they are going."

\vspace{1em}
\textbf{Question 22}

The central point in the first paragraph is that the economic benefits of the Olympic Games

\begin{enumerate}[label=\Alph*.]
    \item are shared equally among the three organising committees.
    \item accrue mostly through revenue from advertisements and ticket sales.
    \item accrue to host cities, if at all, only in the long term.
    \item are usually eroded by expenditure incurred by the host city.
[](https://cracku.in/22-the-central-point-in-the-first-paragraph-is-that-t-x-cat-2017-shift-1)
\end{enumerate}
\vspace{1em}
\textbf{Solution:}

The author expresses his views that hosting Olympics cost a lot and the financial prospects are not that good. Towards the end of the first paragraph, he uses the line "Any economic benefit would have to flow from the value of the Games as an advertisement for the city, the new transportation and communications infrastructure that was created for the Games, or the ongoing use of the new facilities". From this line, we can infer that the author states that there are any benefits, it should be from advertisements and the use of the facilities (implying that the other streams are unlikely to yield any revenue).  
  
It has not been mentioned that the 3 committees share the profit equally. Even if they do, it is not the main point that the author intends to convey through the first paragraph. Therefore, we can eliminate option A.   
  
Option B states that the revenue accrues through advertisement and ticket sales. No information has been provided about 'ticket sales' and hence, option B can be eliminated.  
  
Option D states that the revenues are usually eroded by the expenditure incurred by host city. The second paragraph builds on revenue from 'advertising and infrastructure' implying that the primary purpose of the first paragraph was to introduce these topics rather than to lament that the expenditure of the host city erodes the revenue. Therefore, option C is the right answer.

\end{problem}

\newpage
\begin{problem}{}{}
% Subject: varc
\textbf{Instructions}

Instructions   

**The passage below is accompanied by a set of six questions. Choose the best answer to each question.**
Understanding where you are in the world is a basic survival skill, which is why we, like most species come hard-wired with specialised brain areas to create cognitive maps of our surroundings. Where humans are unique, though, with the possible exception of honeybees, is that we try to communicate this understanding of the world with others. We have a long history of doing this by drawing maps — the earliest versions yet discovered were scrawled on cave walls 14,000 years ago. Human cultures have been drawing them on stone tablets, papyrus, paper and now computer screens ever since.
Given such a long history of human map-making, it is perhaps surprising that it is only within the last few hundred years that north has been consistently considered to be at the top. In fact, for much of human history, north almost never appeared at the top, according to Jerry Brotton, a map historian... "North was rarely put at the top for the simple fact that north is where darkness comes from," he says. "West is also very unlikely to be put at the top because west is where the sun disappears."
Confusingly, early Chinese maps seem to buck this trend. But, Brotton, says, even though they did have compasses at the time, that isn't the reason that they placed north at the top. Early Chinese compasses were actually oriented to point south, which was considered to be more desirable than deepest darkest north. But in Chinese maps, the Emperor, who lived in the north of the country was always put at the top of the map, with everyone else, his loyal subjects, looking up towards him. "In Chinese culture the Emperor looks south because it's where the winds come from, it's a good direction. North is not very good but you are in a position of subjection to the emperor, so you look up to him," says Brotton.
Given that each culture has a very different idea of who, or what, they should look up to it's perhaps not surprising that there is very little consistency in which way early maps pointed. In ancient Egyptian times the top of the world was east, the position of sunrise. Early Islamic maps favoured south at the top because most of the early Muslim cultures were north of Mecca, so they imagined looking up (south) towards it. Christian maps from the same era (called Mappa Mundi) put east at the top, towards the Garden of Eden and with Jerusalem in the centre.
So when did everyone get together and decide that north was the top? It's tempting to put it down to European explorers like Christopher Columbus and Ferdinand Megellan, who were navigating by the North Star. But Brotton argues that these early explorers didn't think of the world like that at all. "When Columbus describes the world it is in accordance with east being at the top, he says. "Columbus says he is going towards paradise, so his mentality is from a medieval mappa mundi." We've got to remember, adds Brotton, that at the time, "no one knows what they are doing and where they are going."

\vspace{1em}
\textbf{Question 23}

Sports facilities built for the Olympics are not fully utilised after the Games are over because

\begin{enumerate}[label=\Alph*.]
    \item their scale and the costs of operating them are large.
    \item their location away from the city centre usually limits easy access.
    \item the authorities do not adapt them to local conditions.
    \item they become outdated having being built with little planning and under time pressure.
[](https://cracku.in/23-sports-facilities-built-for-the-olympics-are-not-f-x-cat-2017-shift-1)
\end{enumerate}
\vspace{1em}
\textbf{Solution:}

In the second paragraph, the author mentions that the facilities are not fully utilized due to their sheer capacity. 'The facilities benefit is dubious at best for buildings such as velodromes or natatoriums and problematic for 100,000-seat Olympic stadiums. The latter require a conversion plan for future use, the former are usually doomed to near vacancy'. Nowhere has it been mentioned that the facilities become outdated or that they are located far from the city. The author mentions in the second paragraph that the structures are expensive to maintain.  
  
Therefore, option A is the right answer.

\end{problem}

\newpage
\begin{problem}{}{}
% Subject: varc
\textbf{Instructions}

Instructions   

**The passage below is accompanied by a set of six questions. Choose the best answer to each question.**
Understanding where you are in the world is a basic survival skill, which is why we, like most species come hard-wired with specialised brain areas to create cognitive maps of our surroundings. Where humans are unique, though, with the possible exception of honeybees, is that we try to communicate this understanding of the world with others. We have a long history of doing this by drawing maps — the earliest versions yet discovered were scrawled on cave walls 14,000 years ago. Human cultures have been drawing them on stone tablets, papyrus, paper and now computer screens ever since.
Given such a long history of human map-making, it is perhaps surprising that it is only within the last few hundred years that north has been consistently considered to be at the top. In fact, for much of human history, north almost never appeared at the top, according to Jerry Brotton, a map historian... "North was rarely put at the top for the simple fact that north is where darkness comes from," he says. "West is also very unlikely to be put at the top because west is where the sun disappears."
Confusingly, early Chinese maps seem to buck this trend. But, Brotton, says, even though they did have compasses at the time, that isn't the reason that they placed north at the top. Early Chinese compasses were actually oriented to point south, which was considered to be more desirable than deepest darkest north. But in Chinese maps, the Emperor, who lived in the north of the country was always put at the top of the map, with everyone else, his loyal subjects, looking up towards him. "In Chinese culture the Emperor looks south because it's where the winds come from, it's a good direction. North is not very good but you are in a position of subjection to the emperor, so you look up to him," says Brotton.
Given that each culture has a very different idea of who, or what, they should look up to it's perhaps not surprising that there is very little consistency in which way early maps pointed. In ancient Egyptian times the top of the world was east, the position of sunrise. Early Islamic maps favoured south at the top because most of the early Muslim cultures were north of Mecca, so they imagined looking up (south) towards it. Christian maps from the same era (called Mappa Mundi) put east at the top, towards the Garden of Eden and with Jerusalem in the centre.
So when did everyone get together and decide that north was the top? It's tempting to put it down to European explorers like Christopher Columbus and Ferdinand Megellan, who were navigating by the North Star. But Brotton argues that these early explorers didn't think of the world like that at all. "When Columbus describes the world it is in accordance with east being at the top, he says. "Columbus says he is going towards paradise, so his mentality is from a medieval mappa mundi." We've got to remember, adds Brotton, that at the time, "no one knows what they are doing and where they are going."

\vspace{1em}
\textbf{Question 24}

The author feels that the Games place a burden on the host city for all of the following reasons EXCEPT that

\begin{enumerate}[label=\Alph*.]
    \item they divert scarce urban land from more productive uses.
    \item they involve the demolition of residential structures to accommodate sports facilities and infrastructure.
    \item the finances used to fund the Games could be better used for other purposes.
    \item the influx of visitors during the Games places a huge strain on the urban infrastructure.  

[](https://cracku.in/24-the-author-feels-that-the-games-place-a-burden-on--x-cat-2017-shift-1)
\end{enumerate}
\vspace{1em}
\textbf{Solution:}

We can infer option A from the line 'Even if they have future use, are they the best use of precious urban real estate?'. The author feels that the urban real estate can be put to better use than hosting Olympics.  
  
The author mentions that the fiscal resources can be used for more productive pursuits than hosting Olympics in the last line of the passage.   
We can infer option B from the line "Residential areas often are razed and citizens relocated ".  
  
Option D cannot be inferred from the passage and hence, it is the right answer.

Instructions   

For the following questions answer them individually

\end{problem}

\newpage
\begin{problem}{}{}
% Subject: varc
\textbf{Instructions}

Instructions   

**The passage below is accompanied by a set of six questions. Choose the best answer to each question.**
Understanding where you are in the world is a basic survival skill, which is why we, like most species come hard-wired with specialised brain areas to create cognitive maps of our surroundings. Where humans are unique, though, with the possible exception of honeybees, is that we try to communicate this understanding of the world with others. We have a long history of doing this by drawing maps — the earliest versions yet discovered were scrawled on cave walls 14,000 years ago. Human cultures have been drawing them on stone tablets, papyrus, paper and now computer screens ever since.
Given such a long history of human map-making, it is perhaps surprising that it is only within the last few hundred years that north has been consistently considered to be at the top. In fact, for much of human history, north almost never appeared at the top, according to Jerry Brotton, a map historian... "North was rarely put at the top for the simple fact that north is where darkness comes from," he says. "West is also very unlikely to be put at the top because west is where the sun disappears."
Confusingly, early Chinese maps seem to buck this trend. But, Brotton, says, even though they did have compasses at the time, that isn't the reason that they placed north at the top. Early Chinese compasses were actually oriented to point south, which was considered to be more desirable than deepest darkest north. But in Chinese maps, the Emperor, who lived in the north of the country was always put at the top of the map, with everyone else, his loyal subjects, looking up towards him. "In Chinese culture the Emperor looks south because it's where the winds come from, it's a good direction. North is not very good but you are in a position of subjection to the emperor, so you look up to him," says Brotton.
Given that each culture has a very different idea of who, or what, they should look up to it's perhaps not surprising that there is very little consistency in which way early maps pointed. In ancient Egyptian times the top of the world was east, the position of sunrise. Early Islamic maps favoured south at the top because most of the early Muslim cultures were north of Mecca, so they imagined looking up (south) towards it. Christian maps from the same era (called Mappa Mundi) put east at the top, towards the Garden of Eden and with Jerusalem in the centre.
So when did everyone get together and decide that north was the top? It's tempting to put it down to European explorers like Christopher Columbus and Ferdinand Megellan, who were navigating by the North Star. But Brotton argues that these early explorers didn't think of the world like that at all. "When Columbus describes the world it is in accordance with east being at the top, he says. "Columbus says he is going towards paradise, so his mentality is from a medieval mappa mundi." We've got to remember, adds Brotton, that at the time, "no one knows what they are doing and where they are going."

\vspace{1em}
\textbf{Question 25}

**The passage given below is followed by four summaries. Choose the option that best captures the author' s position.**  
To me, a "classic" means precisely the opposite of what my predecessors understood: a work is classical by reason of its resistance to contemporaneity and supposed universality, by reason of its capacity to indicate human particularity and difference in that past epoch. The classic is not what tells me about shared humanity — or, more truthfully put, what lets me recognize myself as already present in the past, what nourishes in me the illusion that everything has been like me and has existed only to prepare the way for me. Instead, the classic is what gives access to radically different forms of human consciousness for any given generation of readers, and thereby expands for them the range of possibilities of what it means to be a human being.

\begin{enumerate}[label=\Alph*.]
    \item A classic is able to focus on the contemporary human condition and a unified experience of human consciousness.
    \item A classical work seeks to resist particularity and temporal difference even as it focuses on a common humanity.
    \item A classic is a work exploring the new, going beyond the universal, the contemporary, and the notion of a unified human consciousness.
    \item A classic is a work that provides access to a universal experience of the human race as opposed to radically different forms of human consciousness.
[](https://cracku.in/25-the-passage-given-below-is-followed-by-four-summar-x-cat-2017-shift-1)
\end{enumerate}
\vspace{1em}
\textbf{Solution:}

The author states that a classic is not which puts him at the centre of the universe but one which gives access to radically different forms of human consciousness.   
  
Let us evaluate the options.  
  
Option A states that a classic should focus on unified human experience. The author mentions the exact opposite in the paragraph. Therefore, we can eliminate option A. We can eliminate option D too since it mentions the polar opposite of what that is mentioned in the paragraph. The author is of the view that a classic should go beyond providing a unified human experience and expose one to radically different forms of human consciousness.   
  
We can eliminate option B since it states that a classic focuses on common humanity. Only option C captures the essence of the given paragraph and hence, option C is the right answer.

\end{problem}

\newpage
\begin{problem}{}{}
% Subject: varc
\textbf{Instructions}

Instructions   

**The passage below is accompanied by a set of six questions. Choose the best answer to each question.**
Understanding where you are in the world is a basic survival skill, which is why we, like most species come hard-wired with specialised brain areas to create cognitive maps of our surroundings. Where humans are unique, though, with the possible exception of honeybees, is that we try to communicate this understanding of the world with others. We have a long history of doing this by drawing maps — the earliest versions yet discovered were scrawled on cave walls 14,000 years ago. Human cultures have been drawing them on stone tablets, papyrus, paper and now computer screens ever since.
Given such a long history of human map-making, it is perhaps surprising that it is only within the last few hundred years that north has been consistently considered to be at the top. In fact, for much of human history, north almost never appeared at the top, according to Jerry Brotton, a map historian... "North was rarely put at the top for the simple fact that north is where darkness comes from," he says. "West is also very unlikely to be put at the top because west is where the sun disappears."
Confusingly, early Chinese maps seem to buck this trend. But, Brotton, says, even though they did have compasses at the time, that isn't the reason that they placed north at the top. Early Chinese compasses were actually oriented to point south, which was considered to be more desirable than deepest darkest north. But in Chinese maps, the Emperor, who lived in the north of the country was always put at the top of the map, with everyone else, his loyal subjects, looking up towards him. "In Chinese culture the Emperor looks south because it's where the winds come from, it's a good direction. North is not very good but you are in a position of subjection to the emperor, so you look up to him," says Brotton.
Given that each culture has a very different idea of who, or what, they should look up to it's perhaps not surprising that there is very little consistency in which way early maps pointed. In ancient Egyptian times the top of the world was east, the position of sunrise. Early Islamic maps favoured south at the top because most of the early Muslim cultures were north of Mecca, so they imagined looking up (south) towards it. Christian maps from the same era (called Mappa Mundi) put east at the top, towards the Garden of Eden and with Jerusalem in the centre.
So when did everyone get together and decide that north was the top? It's tempting to put it down to European explorers like Christopher Columbus and Ferdinand Megellan, who were navigating by the North Star. But Brotton argues that these early explorers didn't think of the world like that at all. "When Columbus describes the world it is in accordance with east being at the top, he says. "Columbus says he is going towards paradise, so his mentality is from a medieval mappa mundi." We've got to remember, adds Brotton, that at the time, "no one knows what they are doing and where they are going."

\vspace{1em}
\textbf{Question 26}

**The passage given below is followed by four summaries. Choose the option that best captures the author' s position.  
**  
A translator of literary works needs a secure hold upon the two languages involved, supported by a good measure of familiarity with the two cultures. For an Indian translating works in an Indian language into English, finding satisfactory equivalents in a generalized western culture of practices and symbols in the original would be less difficult than gaining fluent control of contemporary English. When a westerner works on texts in Indian languages the interpretation of cultural elements will be the major challenge, rather than control over the grammar and essential vocabulary of the language concerned. It is much easier to remedy lapses in language in a text translated into English, than flaws of content. Since it is easier for an Indian to learn the English language than it is for a Briton or American to comprehend Indian culture, translations of Indian texts is better left to Indians.

\begin{enumerate}[label=\Alph*.]
    \item While translating, the Indian and the westerner face the same challenges but they have different skill profiles and the former has the advantage.
    \item As preserving cultural meanings is the essence of literary translation Indians' knowledge of the local culture outweighs the initial disadvantage of lower fluency in English.
    \item Indian translators should translate Indian texts into English as their work is less likely to pose cultural problems which are harder to address than the quality of language.
    \item Westerners might be good at gaining reasonable fluency in new languages, but as understanding the culture reflected in literature is crucial, Indians remain better placed.
[](https://cracku.in/26-the-passage-given-below-is-followed-by-four-summar-x-cat-2017-shift-1)
\end{enumerate}
\vspace{1em}
\textbf{Solution:}

Let us note down the important points put down by the author.   
  
Indians have better knowledge of their culture. A westerner might be fluent in the language but will find it hard to relate to the culture. Indians, on the other hand, might be less fluent in the language but will be able to preserve the culture when a text is translated. Therefore, Indians should translate Indian texts.   
  
Let us evaluate the options now.   
  
Option A states that Indians and Westerners face the same challenges but they have different skill sets. Indians and Westerners face different challenges while translating the text. Indians face difficulty in the language and westerners face difficulty in relating to the culture. Therefore, we can eliminate option A.  
  
Option D fails to capture the fact that the primary intention of the paragraph is not to pit Indians against westerners but to suggest that Indians should translate Indian texts. Also, it does not capture the fact that Indians will retain the advantage only when translating the Indian texts. Therefore, we can eliminate option D.   
  
Option B, though true, fails to capture the India-centric angle that the paragraph adopts. The paragraph places huge emphasis on the term 'Indian texts' and only option C manages to capture this fact. Also, only option C captures the fact that it is easier to remedy errors in the language than to fix errors in the interpretation of culture. Therefore, option C is the right answer.

\end{problem}

\newpage
\begin{problem}{}{}
% Subject: varc
\textbf{Instructions}

Instructions   

**The passage below is accompanied by a set of six questions. Choose the best answer to each question.**
Understanding where you are in the world is a basic survival skill, which is why we, like most species come hard-wired with specialised brain areas to create cognitive maps of our surroundings. Where humans are unique, though, with the possible exception of honeybees, is that we try to communicate this understanding of the world with others. We have a long history of doing this by drawing maps — the earliest versions yet discovered were scrawled on cave walls 14,000 years ago. Human cultures have been drawing them on stone tablets, papyrus, paper and now computer screens ever since.
Given such a long history of human map-making, it is perhaps surprising that it is only within the last few hundred years that north has been consistently considered to be at the top. In fact, for much of human history, north almost never appeared at the top, according to Jerry Brotton, a map historian... "North was rarely put at the top for the simple fact that north is where darkness comes from," he says. "West is also very unlikely to be put at the top because west is where the sun disappears."
Confusingly, early Chinese maps seem to buck this trend. But, Brotton, says, even though they did have compasses at the time, that isn't the reason that they placed north at the top. Early Chinese compasses were actually oriented to point south, which was considered to be more desirable than deepest darkest north. But in Chinese maps, the Emperor, who lived in the north of the country was always put at the top of the map, with everyone else, his loyal subjects, looking up towards him. "In Chinese culture the Emperor looks south because it's where the winds come from, it's a good direction. North is not very good but you are in a position of subjection to the emperor, so you look up to him," says Brotton.
Given that each culture has a very different idea of who, or what, they should look up to it's perhaps not surprising that there is very little consistency in which way early maps pointed. In ancient Egyptian times the top of the world was east, the position of sunrise. Early Islamic maps favoured south at the top because most of the early Muslim cultures were north of Mecca, so they imagined looking up (south) towards it. Christian maps from the same era (called Mappa Mundi) put east at the top, towards the Garden of Eden and with Jerusalem in the centre.
So when did everyone get together and decide that north was the top? It's tempting to put it down to European explorers like Christopher Columbus and Ferdinand Megellan, who were navigating by the North Star. But Brotton argues that these early explorers didn't think of the world like that at all. "When Columbus describes the world it is in accordance with east being at the top, he says. "Columbus says he is going towards paradise, so his mentality is from a medieval mappa mundi." We've got to remember, adds Brotton, that at the time, "no one knows what they are doing and where they are going."

\vspace{1em}
\textbf{Question 27}

**The passage given below is followed by four summaries. Choose the option that best captures the author' s position.  
**  
For each of the past three years, temperatures have hit peaks not seen since the birth of meteorology, and probably not for more than 110,000 years. The amount of carbon dioxide in the air is at its highest level in 4 million years. This does not cause storms like Harvey — there have always been storms and hurricanes along the Gulf of Mexico — but it makes them wetter and more powerful. As the seas warm, they evaporate more easily and provide energy to storm fronts. As the air above them warms, it holds more water vapour. For every half a degree Celsius in warming, there is about a 3% increase in atmospheric moisture content. Scientists call this the Clausius-Clapeyron equation. This means the skies fill more quickly and have more to dump. The storm surge was greater because sea levels have risen 20 cm as a result of more than 100 years of human- related global warming which has melted glaciers and thermally expanded the volume of seawater.

\begin{enumerate}[label=\Alph*.]
    \item The storm Harvey is one of the regular, annual ones from the Gulf of Mexico; global warming and Harvey are unrelated phenomena.
    \item Global warming does not breed storms but makes them more destructive; the Clausius-Clapeyron equation, though it predicts potential increase in atmospheric moisture content, cannot predict the scale of damage storms might wreck.
    \item Global warming melts glaciers, resulting in seawater volume expansion; this enables more water vapour to fill the air above faster. Thus, modern storms contain more destructive energy.
    \item It is naive to think that rising sea levels and the force of tropical storms are unrelated; Harvey was destructive as global warming has armed it with more moisture content, but this may not be true of all storms.
[](https://cracku.in/27-the-passage-given-below-is-followed-by-four-summar-x-cat-2017-shift-1)
\end{enumerate}
\vspace{1em}
\textbf{Solution:}

Let us note down the important points in the given paragraph.  
  
Global warming does not cause storms but make them more powerful. Due to the increase in the temperature, the air can absorb more moisture. This relationship (the change in the ability to absorb water with the increase in the temperature) is given by the Clausius-Clapeyron equation.  
  
Let us evaluate the options.  
  
The author provides storm Harvey as an example to illustrate how increased temperatures can arm the storms with more power. Harvey is not the central theme of the given paragraph. We can eliminate options A and D since option D places much emphasis on storm Harvey and option A states that there is no relationship between the increase in temperature and the power of storms.   
  
Option B states that the Clausius-Clapeyron equation cannot predict the quantum of destruction that a storm might cause. This point is totally out of context with respect to what that is being discussed in the paragraph. Therefore, we can eliminate option B as well.   
  
Option C precisely explains the mechanism through which global warming makes the modern storms more destructive. Therefore, option C is the right answer.

\end{problem}

\newpage
\begin{problem}{}{}
% Subject: varc
\textbf{Instructions}

Instructions   

**The passage below is accompanied by a set of six questions. Choose the best answer to each question.**
Understanding where you are in the world is a basic survival skill, which is why we, like most species come hard-wired with specialised brain areas to create cognitive maps of our surroundings. Where humans are unique, though, with the possible exception of honeybees, is that we try to communicate this understanding of the world with others. We have a long history of doing this by drawing maps — the earliest versions yet discovered were scrawled on cave walls 14,000 years ago. Human cultures have been drawing them on stone tablets, papyrus, paper and now computer screens ever since.
Given such a long history of human map-making, it is perhaps surprising that it is only within the last few hundred years that north has been consistently considered to be at the top. In fact, for much of human history, north almost never appeared at the top, according to Jerry Brotton, a map historian... "North was rarely put at the top for the simple fact that north is where darkness comes from," he says. "West is also very unlikely to be put at the top because west is where the sun disappears."
Confusingly, early Chinese maps seem to buck this trend. But, Brotton, says, even though they did have compasses at the time, that isn't the reason that they placed north at the top. Early Chinese compasses were actually oriented to point south, which was considered to be more desirable than deepest darkest north. But in Chinese maps, the Emperor, who lived in the north of the country was always put at the top of the map, with everyone else, his loyal subjects, looking up towards him. "In Chinese culture the Emperor looks south because it's where the winds come from, it's a good direction. North is not very good but you are in a position of subjection to the emperor, so you look up to him," says Brotton.
Given that each culture has a very different idea of who, or what, they should look up to it's perhaps not surprising that there is very little consistency in which way early maps pointed. In ancient Egyptian times the top of the world was east, the position of sunrise. Early Islamic maps favoured south at the top because most of the early Muslim cultures were north of Mecca, so they imagined looking up (south) towards it. Christian maps from the same era (called Mappa Mundi) put east at the top, towards the Garden of Eden and with Jerusalem in the centre.
So when did everyone get together and decide that north was the top? It's tempting to put it down to European explorers like Christopher Columbus and Ferdinand Megellan, who were navigating by the North Star. But Brotton argues that these early explorers didn't think of the world like that at all. "When Columbus describes the world it is in accordance with east being at the top, he says. "Columbus says he is going towards paradise, so his mentality is from a medieval mappa mundi." We've got to remember, adds Brotton, that at the time, "no one knows what they are doing and where they are going."

\vspace{1em}
\textbf{Question 28}

The five sentences labelled 1, 2, 3, 4, 5) given in this question, when properly sequenced, form a coherent paragraph. Each sentence is labelled with a number. Decide on. the proper order for the sentences and key in this sequence of five numbers as your answer.  
1. The process of handing down implies not a passive transfer, but some contestation in defining what exactly is to be handed down.  
2. Wherever Western scholars have worked on the Indian past, the selection is even more apparent and the inventing of a tradition much more recognizable.  
3. Every generation selects what it requires from the past and makes its innovations, some more than others.   
4. It is now a truism to say that traditions are not handed down unchanged, but are invented.  
5. Just as life has death as its opposite, so is tradition by default the opposite of innovation.
Backspace  
789  
456  
123  
0.-  
Clear All  

Submit
[](https://cracku.in/28-the-five-sentences-labelled-1-2-3-4-5-given-in-thi-x-cat-2017-shift-1)

\begin{enumerate}[label=\Alph*.]
\end{enumerate}
\vspace{1em}
\textbf{Solution:}

A quick glance at all the five questions suggest that the passage is discussing tradition and innovation. Hence, 5 is the best opening sentence as it sets the context for the discussion. 4 and 1 form a pair as both of them talk about handing down. Similarly, 3, 2 also form a pair as both talk about selections. Hence, the correct logical order of the given sentences will be 54132.

\end{problem}

\newpage
\begin{problem}{}{}
% Subject: varc
\textbf{Instructions}

Instructions   

**The passage below is accompanied by a set of six questions. Choose the best answer to each question.**
Understanding where you are in the world is a basic survival skill, which is why we, like most species come hard-wired with specialised brain areas to create cognitive maps of our surroundings. Where humans are unique, though, with the possible exception of honeybees, is that we try to communicate this understanding of the world with others. We have a long history of doing this by drawing maps — the earliest versions yet discovered were scrawled on cave walls 14,000 years ago. Human cultures have been drawing them on stone tablets, papyrus, paper and now computer screens ever since.
Given such a long history of human map-making, it is perhaps surprising that it is only within the last few hundred years that north has been consistently considered to be at the top. In fact, for much of human history, north almost never appeared at the top, according to Jerry Brotton, a map historian... "North was rarely put at the top for the simple fact that north is where darkness comes from," he says. "West is also very unlikely to be put at the top because west is where the sun disappears."
Confusingly, early Chinese maps seem to buck this trend. But, Brotton, says, even though they did have compasses at the time, that isn't the reason that they placed north at the top. Early Chinese compasses were actually oriented to point south, which was considered to be more desirable than deepest darkest north. But in Chinese maps, the Emperor, who lived in the north of the country was always put at the top of the map, with everyone else, his loyal subjects, looking up towards him. "In Chinese culture the Emperor looks south because it's where the winds come from, it's a good direction. North is not very good but you are in a position of subjection to the emperor, so you look up to him," says Brotton.
Given that each culture has a very different idea of who, or what, they should look up to it's perhaps not surprising that there is very little consistency in which way early maps pointed. In ancient Egyptian times the top of the world was east, the position of sunrise. Early Islamic maps favoured south at the top because most of the early Muslim cultures were north of Mecca, so they imagined looking up (south) towards it. Christian maps from the same era (called Mappa Mundi) put east at the top, towards the Garden of Eden and with Jerusalem in the centre.
So when did everyone get together and decide that north was the top? It's tempting to put it down to European explorers like Christopher Columbus and Ferdinand Megellan, who were navigating by the North Star. But Brotton argues that these early explorers didn't think of the world like that at all. "When Columbus describes the world it is in accordance with east being at the top, he says. "Columbus says he is going towards paradise, so his mentality is from a medieval mappa mundi." We've got to remember, adds Brotton, that at the time, "no one knows what they are doing and where they are going."

\vspace{1em}
\textbf{Question 29}

**The five sentences labelled (1, 2, 3, 4, 5) given in this question, when properly sequenced, form a coherent paragraph. Each sentence is labelled with a number. Decide on the proper order for the sentences and key in this sequence of five numbers as your answer.**
1. Scientists have for the first time managed to edit genes in a human embryo to repair a genetic mutation, fuelling hopes that such procedures may one day be available outside laboratory conditions.  
2. The cardiac disease causes sudden death in otherwise healthy young athletes and affects about one in 500 people overall.  
3. Correcting the mutation in the gene would not only ensure that the child is healthy but also prevents transmission of the mutation to future generations.  
4. It is caused by a mutation in a particular gene and a child will suffer from the condition even if it inherits only one copy of the mutated gene.  
5. In results announced in Nature this week, scientists fixed a mutation that thickens the heart muscle, a condition called hypertrophic cardiomyopathy.
Backspace  
789  
456  
123  
0.-  
Clear All  

Submit
[](https://cracku.in/29-the-five-sentences-labelled-1-2-3-4-5-given-in-thi-x-cat-2017-shift-1)

\begin{enumerate}[label=\Alph*.]
\end{enumerate}
\vspace{1em}
\textbf{Solution:}

After reading all the given sentences, we know that the paragraph is about the gene editing and a specific disease where the process has been used effectively. Statement 1 is the opening sentences introducing the topic that how for the first time scientists have managed to edit gene. Statement 5 provides the specifics of the case where the gene editing has taken place. 'The cardiac disease' mentioned in statement 2 refers to hypertrophic cardiomyopathy discussed in statement 5. So, statement 2 must follow statement 5. Statement 4 explains the causes of the of the disease mentioned in statement 2. Statement 3 is a conclusion about the result of the process explained earlier. Thus, the correct order is 1 - 5 - 2 - 4 - 3.  
Hence, 15243 is the correct answer.

\end{problem}

\newpage
\begin{problem}{}{}
% Subject: varc
\textbf{Instructions}

Instructions   

**The passage below is accompanied by a set of six questions. Choose the best answer to each question.**
Understanding where you are in the world is a basic survival skill, which is why we, like most species come hard-wired with specialised brain areas to create cognitive maps of our surroundings. Where humans are unique, though, with the possible exception of honeybees, is that we try to communicate this understanding of the world with others. We have a long history of doing this by drawing maps — the earliest versions yet discovered were scrawled on cave walls 14,000 years ago. Human cultures have been drawing them on stone tablets, papyrus, paper and now computer screens ever since.
Given such a long history of human map-making, it is perhaps surprising that it is only within the last few hundred years that north has been consistently considered to be at the top. In fact, for much of human history, north almost never appeared at the top, according to Jerry Brotton, a map historian... "North was rarely put at the top for the simple fact that north is where darkness comes from," he says. "West is also very unlikely to be put at the top because west is where the sun disappears."
Confusingly, early Chinese maps seem to buck this trend. But, Brotton, says, even though they did have compasses at the time, that isn't the reason that they placed north at the top. Early Chinese compasses were actually oriented to point south, which was considered to be more desirable than deepest darkest north. But in Chinese maps, the Emperor, who lived in the north of the country was always put at the top of the map, with everyone else, his loyal subjects, looking up towards him. "In Chinese culture the Emperor looks south because it's where the winds come from, it's a good direction. North is not very good but you are in a position of subjection to the emperor, so you look up to him," says Brotton.
Given that each culture has a very different idea of who, or what, they should look up to it's perhaps not surprising that there is very little consistency in which way early maps pointed. In ancient Egyptian times the top of the world was east, the position of sunrise. Early Islamic maps favoured south at the top because most of the early Muslim cultures were north of Mecca, so they imagined looking up (south) towards it. Christian maps from the same era (called Mappa Mundi) put east at the top, towards the Garden of Eden and with Jerusalem in the centre.
So when did everyone get together and decide that north was the top? It's tempting to put it down to European explorers like Christopher Columbus and Ferdinand Megellan, who were navigating by the North Star. But Brotton argues that these early explorers didn't think of the world like that at all. "When Columbus describes the world it is in accordance with east being at the top, he says. "Columbus says he is going towards paradise, so his mentality is from a medieval mappa mundi." We've got to remember, adds Brotton, that at the time, "no one knows what they are doing and where they are going."

\vspace{1em}
\textbf{Question 30}

The five sentences labelled (1, 2, 3, 4, 5) given in this question, when properly sequenced, form a coherent paragraph. Each sentence is labelled with a number. Decide on. the proper order for the sentences and key in this sequence of five numbers as your answer.  
1. The study suggests that the disease did not spread with such intensity, but that it may have driven human migrations across Europe and Asia.  
2. The oldest sample came from an individual who lived in southeast Russia about 5,000 years ago.  
3. The ages of the skeletons correspond to a time of mass exodus from today's Russia and Ukraine into western Europe and central Asia, suggesting that a pandemic could have driven these migrations.  
4. In the analysis of fragments of DNA from 101 Bronze Age skeletons for sequences from Yersinia pestis, the bacterium that causes the disease, seven tested positive.  
5. DNA from Bronze Age human skeletons indicate that the black plague could have emerged as early as 3,000 BCE, long before the epidemic that swept through Europe in the rnid-1300s.
Backspace  
789  
456  
123  
0.-  
Clear All  

Submit
[](https://cracku.in/30-the-five-sentences-labelled-1-2-3-4-5-given-in-thi-x-cat-2017-shift-1)

\begin{enumerate}[label=\Alph*.]
\end{enumerate}
\vspace{1em}
\textbf{Solution:}

On carefully reading the sentences, we see that the paragraph is about the disease - black plaque and its impact on migration. Sentence 5 introduces the subject. Hence, it should be the opening sentence of the paragraph. Sentence 4 should follow 5 as 4 talks about the bronze age skeletons, which are mentioned in 5. Sentence 1 should follow 4 as 'the study' refers to the analysis of skeletons mentioned in 4. Sentences 2 and 3 form a pair as 2 mentions that the oldest sample was found in Russia and 3 mentions that this could have been the reason for mass exodus. 54123 forms a coherent paragraph.

\end{problem}

\newpage
\begin{problem}{}{}
% Subject: varc
\textbf{Instructions}

Instructions   

**The passage below is accompanied by a set of six questions. Choose the best answer to each question.**
Understanding where you are in the world is a basic survival skill, which is why we, like most species come hard-wired with specialised brain areas to create cognitive maps of our surroundings. Where humans are unique, though, with the possible exception of honeybees, is that we try to communicate this understanding of the world with others. We have a long history of doing this by drawing maps — the earliest versions yet discovered were scrawled on cave walls 14,000 years ago. Human cultures have been drawing them on stone tablets, papyrus, paper and now computer screens ever since.
Given such a long history of human map-making, it is perhaps surprising that it is only within the last few hundred years that north has been consistently considered to be at the top. In fact, for much of human history, north almost never appeared at the top, according to Jerry Brotton, a map historian... "North was rarely put at the top for the simple fact that north is where darkness comes from," he says. "West is also very unlikely to be put at the top because west is where the sun disappears."
Confusingly, early Chinese maps seem to buck this trend. But, Brotton, says, even though they did have compasses at the time, that isn't the reason that they placed north at the top. Early Chinese compasses were actually oriented to point south, which was considered to be more desirable than deepest darkest north. But in Chinese maps, the Emperor, who lived in the north of the country was always put at the top of the map, with everyone else, his loyal subjects, looking up towards him. "In Chinese culture the Emperor looks south because it's where the winds come from, it's a good direction. North is not very good but you are in a position of subjection to the emperor, so you look up to him," says Brotton.
Given that each culture has a very different idea of who, or what, they should look up to it's perhaps not surprising that there is very little consistency in which way early maps pointed. In ancient Egyptian times the top of the world was east, the position of sunrise. Early Islamic maps favoured south at the top because most of the early Muslim cultures were north of Mecca, so they imagined looking up (south) towards it. Christian maps from the same era (called Mappa Mundi) put east at the top, towards the Garden of Eden and with Jerusalem in the centre.
So when did everyone get together and decide that north was the top? It's tempting to put it down to European explorers like Christopher Columbus and Ferdinand Megellan, who were navigating by the North Star. But Brotton argues that these early explorers didn't think of the world like that at all. "When Columbus describes the world it is in accordance with east being at the top, he says. "Columbus says he is going towards paradise, so his mentality is from a medieval mappa mundi." We've got to remember, adds Brotton, that at the time, "no one knows what they are doing and where they are going."

\vspace{1em}
\textbf{Question 31}

The five sentences labelled (1, 2, 3, 4, 5) given in this question, when properly sequenced, form a coherent paragraph. Each sentence is labelled with a number. Decide on the proper order for the sentences and key in this sequence of five numbers as your answer.  
1. This visual turn in social media has merely accentuated this announcing instinct of ours, enabling us with easy-to-create, easy-to-share, easy-to-store and easy-to-consume platforms, gadgets and apps.  
2. There is absolutely nothing new about us framing the vision of who we are or what we want, visually or otherwise, in our Facebook page, for example.  
3. Turning the pages of most family albums, which belong to a period well before the digital dissemination of self-created and self-curated moments and images, would reconfirm the basic instinct of documenting our presence in a particular space, on a significant occasion, with others who matter.  
4. We are empowered to book our faces and act as celebrities within the confinement of our respective friend lists, and communicate our activities, companionship and locations with minimal clicks and touches.  
5. What is unprecedented is not the desire to put out newsfeeds related to the self, but the ease with which this broadcast operation can now be executed, often provoking (un)anticipated responses from beyond one' s immediate location.
Backspace  
789  
456  
123  
0.-  
Clear All  

Submit
[](https://cracku.in/31-the-five-sentences-labelled-1-2-3-4-5-given-in-thi-x-cat-2017-shift-1)

\begin{enumerate}[label=\Alph*.]
\end{enumerate}
\vspace{1em}
\textbf{Solution:}

3 is the starting sentence of the given paragraph as it introduces us the topic.  
2 follows 3 as it further talks about what's given in 3 as can be understood by "There is absolutely nothing new about us framing the vision" and this vision talked about in 3.  
1 follows 2 as "This visual turn" has it's precedent is in 2. The visual turn in 1 is explained in 2.  
5 concludes the summary and it also follows 4 as "What is unprecedented is not the desire to put out newsfeeds related to the self" given in 5 is explained in 4.  
Hence, the correct order is 32145.

\end{problem}

\newpage
\begin{problem}{}{}
% Subject: varc
\textbf{Instructions}

Instructions   

**The passage below is accompanied by a set of six questions. Choose the best answer to each question.**
Understanding where you are in the world is a basic survival skill, which is why we, like most species come hard-wired with specialised brain areas to create cognitive maps of our surroundings. Where humans are unique, though, with the possible exception of honeybees, is that we try to communicate this understanding of the world with others. We have a long history of doing this by drawing maps — the earliest versions yet discovered were scrawled on cave walls 14,000 years ago. Human cultures have been drawing them on stone tablets, papyrus, paper and now computer screens ever since.
Given such a long history of human map-making, it is perhaps surprising that it is only within the last few hundred years that north has been consistently considered to be at the top. In fact, for much of human history, north almost never appeared at the top, according to Jerry Brotton, a map historian... "North was rarely put at the top for the simple fact that north is where darkness comes from," he says. "West is also very unlikely to be put at the top because west is where the sun disappears."
Confusingly, early Chinese maps seem to buck this trend. But, Brotton, says, even though they did have compasses at the time, that isn't the reason that they placed north at the top. Early Chinese compasses were actually oriented to point south, which was considered to be more desirable than deepest darkest north. But in Chinese maps, the Emperor, who lived in the north of the country was always put at the top of the map, with everyone else, his loyal subjects, looking up towards him. "In Chinese culture the Emperor looks south because it's where the winds come from, it's a good direction. North is not very good but you are in a position of subjection to the emperor, so you look up to him," says Brotton.
Given that each culture has a very different idea of who, or what, they should look up to it's perhaps not surprising that there is very little consistency in which way early maps pointed. In ancient Egyptian times the top of the world was east, the position of sunrise. Early Islamic maps favoured south at the top because most of the early Muslim cultures were north of Mecca, so they imagined looking up (south) towards it. Christian maps from the same era (called Mappa Mundi) put east at the top, towards the Garden of Eden and with Jerusalem in the centre.
So when did everyone get together and decide that north was the top? It's tempting to put it down to European explorers like Christopher Columbus and Ferdinand Megellan, who were navigating by the North Star. But Brotton argues that these early explorers didn't think of the world like that at all. "When Columbus describes the world it is in accordance with east being at the top, he says. "Columbus says he is going towards paradise, so his mentality is from a medieval mappa mundi." We've got to remember, adds Brotton, that at the time, "no one knows what they are doing and where they are going."

\vspace{1em}
\textbf{Question 32}

**Five sentences related to a topic are given below. Four of them can be put together to form a meaningful and coherent short paragraph. Identify the odd one out.**
l. People who study children's language spend a lot of time watching how babies react to the speech they hear around them.  
2. They make films of adults and babies interacting, and examine them very carefully to see whether the babies show any signs of understanding what the adults say.  
3. They believe that babies begin to react to language from the very moment they are born.  
4. Sometimes the signs are very subtle — slight movements of the baby's eyes or the head or the hands.  
5. You'd never notice them if you were just sitting with the child, but by watching a recording over and over, you can spot them.
Backspace  
789  
456  
123  
0.-  
Clear All  

Submit
[](https://cracku.in/32-five-sentences-related-to-a-topic-are-given-below--x-cat-2017-shift-1)

\begin{enumerate}[label=\Alph*.]
\end{enumerate}
\vspace{1em}
\textbf{Solution:}

After reading all the sentences, we know that the paragraph is about the children's language and the signs that they show. Statement 1 is the opening sentence as it introduces us to the method adopted to study children's language. Statement 2 further explains the method how people study the signs given by children. Statements 4 and 5 are about the signs mentioned in statement 2. Thus, all the four statements are related to the methodology adopted by people to study children's language. Therefore, these 4 sentences form a paragraph.  
Statement 3 is about the reaction of children to a certain language. So, statement 3 is about a different topic and does not fit in the paragraph.  
Hence, 3 is the correct answer.

\end{problem}

\newpage
\begin{problem}{}{}
% Subject: varc
\textbf{Instructions}

Instructions   

**The passage below is accompanied by a set of six questions. Choose the best answer to each question.**
Understanding where you are in the world is a basic survival skill, which is why we, like most species come hard-wired with specialised brain areas to create cognitive maps of our surroundings. Where humans are unique, though, with the possible exception of honeybees, is that we try to communicate this understanding of the world with others. We have a long history of doing this by drawing maps — the earliest versions yet discovered were scrawled on cave walls 14,000 years ago. Human cultures have been drawing them on stone tablets, papyrus, paper and now computer screens ever since.
Given such a long history of human map-making, it is perhaps surprising that it is only within the last few hundred years that north has been consistently considered to be at the top. In fact, for much of human history, north almost never appeared at the top, according to Jerry Brotton, a map historian... "North was rarely put at the top for the simple fact that north is where darkness comes from," he says. "West is also very unlikely to be put at the top because west is where the sun disappears."
Confusingly, early Chinese maps seem to buck this trend. But, Brotton, says, even though they did have compasses at the time, that isn't the reason that they placed north at the top. Early Chinese compasses were actually oriented to point south, which was considered to be more desirable than deepest darkest north. But in Chinese maps, the Emperor, who lived in the north of the country was always put at the top of the map, with everyone else, his loyal subjects, looking up towards him. "In Chinese culture the Emperor looks south because it's where the winds come from, it's a good direction. North is not very good but you are in a position of subjection to the emperor, so you look up to him," says Brotton.
Given that each culture has a very different idea of who, or what, they should look up to it's perhaps not surprising that there is very little consistency in which way early maps pointed. In ancient Egyptian times the top of the world was east, the position of sunrise. Early Islamic maps favoured south at the top because most of the early Muslim cultures were north of Mecca, so they imagined looking up (south) towards it. Christian maps from the same era (called Mappa Mundi) put east at the top, towards the Garden of Eden and with Jerusalem in the centre.
So when did everyone get together and decide that north was the top? It's tempting to put it down to European explorers like Christopher Columbus and Ferdinand Megellan, who were navigating by the North Star. But Brotton argues that these early explorers didn't think of the world like that at all. "When Columbus describes the world it is in accordance with east being at the top, he says. "Columbus says he is going towards paradise, so his mentality is from a medieval mappa mundi." We've got to remember, adds Brotton, that at the time, "no one knows what they are doing and where they are going."

\vspace{1em}
\textbf{Question 33}

**Five sentences related to a topic are given below. Four of them can be put together to form a meaningful and coherent short paragraph. Identify the odd one out.  
  
** 1. Neuroscientists have just begun studying exercise's impact within brain cells — on the genes themselves.   
2. Even there, in the roots of our biology, they' ve found signs of the body's influence on the mind.  
3. It turns out that moving our muscles produces proteins that travel through the bloodstream and into the brain, where they play pivotal roles in the mechanisms of our highest thought processes.  
4. In today's technology-driven, plasma-screened-in world, it's easy to forget that we are born movers-animals, in fact — because we' ve engineered movement right out of our lives.  
5. It's only in the past few years that neuroscientists have begun to describe these factors and how they work, and each new discovery adds awe-inspiring depth to the picture.
Backspace  
789  
456  
123  
0.-  
Clear All  

Submit
[](https://cracku.in/33-five-sentences-related-to-a-topic-are-given-below--x-cat-2017-shift-1)

\begin{enumerate}[label=\Alph*.]
\end{enumerate}
\vspace{1em}
\textbf{Solution:}

After reading all the sentences, we know that the passage is about the study of the effect of exercise on the mind by the neuroscientists. Statement 1 is the opening sentence as it introduces the main idea and suggests that neuroscientists have started studying the effect of exercise on the mind. Statement 2 discusses the findings of the study mentioned in statement 1. Statement 3 further elaborates the finding and explains the reasons behind the effect of exercise on brain cells. Statement 5 is a conclusion sentence based on the study mentioned in the three sentences mentioned earlier. Thus, 1235 forms a meaningful paragraph.  
Statement 4 focuses on our ignorance of movements due to the widespread use of technology. Other four sentences are about the relationship between exercise and brain cells. However, statement 4 is about a different topic. Therefore, statement 4 is an odd sentence which does not fit in the paragraph.  
Hence, 4 is the correct answer.

\end{problem}

\newpage
\begin{problem}{}{}
% Subject: varc
\textbf{Instructions}

Instructions   

**The passage below is accompanied by a set of six questions. Choose the best answer to each question.**
Understanding where you are in the world is a basic survival skill, which is why we, like most species come hard-wired with specialised brain areas to create cognitive maps of our surroundings. Where humans are unique, though, with the possible exception of honeybees, is that we try to communicate this understanding of the world with others. We have a long history of doing this by drawing maps — the earliest versions yet discovered were scrawled on cave walls 14,000 years ago. Human cultures have been drawing them on stone tablets, papyrus, paper and now computer screens ever since.
Given such a long history of human map-making, it is perhaps surprising that it is only within the last few hundred years that north has been consistently considered to be at the top. In fact, for much of human history, north almost never appeared at the top, according to Jerry Brotton, a map historian... "North was rarely put at the top for the simple fact that north is where darkness comes from," he says. "West is also very unlikely to be put at the top because west is where the sun disappears."
Confusingly, early Chinese maps seem to buck this trend. But, Brotton, says, even though they did have compasses at the time, that isn't the reason that they placed north at the top. Early Chinese compasses were actually oriented to point south, which was considered to be more desirable than deepest darkest north. But in Chinese maps, the Emperor, who lived in the north of the country was always put at the top of the map, with everyone else, his loyal subjects, looking up towards him. "In Chinese culture the Emperor looks south because it's where the winds come from, it's a good direction. North is not very good but you are in a position of subjection to the emperor, so you look up to him," says Brotton.
Given that each culture has a very different idea of who, or what, they should look up to it's perhaps not surprising that there is very little consistency in which way early maps pointed. In ancient Egyptian times the top of the world was east, the position of sunrise. Early Islamic maps favoured south at the top because most of the early Muslim cultures were north of Mecca, so they imagined looking up (south) towards it. Christian maps from the same era (called Mappa Mundi) put east at the top, towards the Garden of Eden and with Jerusalem in the centre.
So when did everyone get together and decide that north was the top? It's tempting to put it down to European explorers like Christopher Columbus and Ferdinand Megellan, who were navigating by the North Star. But Brotton argues that these early explorers didn't think of the world like that at all. "When Columbus describes the world it is in accordance with east being at the top, he says. "Columbus says he is going towards paradise, so his mentality is from a medieval mappa mundi." We've got to remember, adds Brotton, that at the time, "no one knows what they are doing and where they are going."

\vspace{1em}
\textbf{Question 34}

**Five sentences related to a topic are given below. Four of them can be put together to form a meaningful and coherent short paragraph. Identify the odd one out.**  
1. The water that made up ancient lakes and perhaps an ocean was lost.  
2. Particles from the Sun collided with molecules in the atmosphere, knocking them into space or giving them an electric charge that caused them to be swept away by the solar wind.  
3. Most of the planet's remaining water is now frozen or buried, but clues over the past decade suggested that some liquid water, a presumed necessity for life, might survive in underground aquifers.  
4. Data from NASA's MAVEN orbiter show that solar storms stripped away most of Mars's once-thick atmosphere.  
5. A recent study reveals how Mars lost much of its early water, while another indicates that some liquid water remains.  

Backspace  
789  
456  
123  
0.-  
Clear All  

Submit
[](https://cracku.in/34-five-sentences-related-to-a-topic-are-given-below--x-cat-2017-shift-1)

\begin{enumerate}[label=\Alph*.]
\end{enumerate}
\vspace{1em}
\textbf{Solution:}

On reading the sentences, we can infer that the paragraph is about how Mars lost most of its water.   
Sentence 5 states that Mars has lost much of its water according to one study but some liquid water remains according to another. Therefore, the rest of the paragraph should explain about these 2 findings.   
Sentences 4 and 2 together substantiate the first study. They try to explain how a solar storm swept away Mar's water content. Both the sentences talk about the solar storm. Sentence 3 talks about the water that is remaining on the planet.   
Sentences 5423 can be put together into a coherent paragraph.  
  
Therefore, sentence 1 is the odd one out.
[](https://cracku.in/downloads/1883)
  *  
  *

\end{problem}

\newpage
\begin{problem}{}{}
% Subject: varc
\textbf{Instructions}

Instructions   

**The passage below is accompanied by a set of six questions. Choose the best answer to each question.**
Cities are the true fonts of creativity... With their diverse populations, dense social networks, and public spaces where people can meet spontaneously and serendipitously, they spark and catalyze new ideas. With their infrastructure for finance, organization and trade, they allow those ideas to be swiftly actualized.
As for what staunches creativity, that's easy, if ironic. It's the very institutions that we build to manage, exploit and perpetuate the fruits of creativity — our big bureaucracies, and sad to say, too many of our schools. Creativity is disruptive; schools and organizations are regimented, standardized and stultifying.
The education expert Sir Ken Robinson points to a 1968 study reporting on a group of 1,600 children who were tested over time for their ability to think in out-of-the-box ways. When the children were between 3 and 5 years old, 98 percent achieved positive scores. When they were 8 to 10, only 32 percent passed the same test, and only 10 percent at 13 to 15. When 280,000 25-year-olds took the test, just 2 percent passed. By the time we are adults, our creativity has been wrung out of us.
I once asked the great urbanist Jane Jacobs what makes some places more creative than others. She said, essentially, that the question was an easy one. All cities, she said, were filled with creative people; that's our default state as people. But some cities had more than their shares of leaders, people and institutions that blocked out that creativity. She called them "squelchers."
Creativity (or the lack of it) follows the same general contours of the great socio-economic divide — our rising inequality — that plagues us. According to my own estimates, roughly a third of us across the United States, and perhaps as much as half of us in our most creative cities — are able to do work which engages our creative faculties to some extent, whether as artists, musicians, writers, techies, innovators, entrepreneurs, doctors, lawyers, journalists or educators — those of us who work with our minds. That leaves a group that I term "the other 66 percent," who toil in low-wage rote and rotten jobs — if they have jobs at all — in which their creativity is subjugated, ignored or wasted.
Creativity itself is not in danger. It's flourishing is all around us — in science and technology, arts and culture, in our rapidly revitalizing cities. But we still have a long way to go if we want to build a truly creative society that supports and rewards the creativity of each and every one of us.

\vspace{1em}
\textbf{Question 1}

In the author's view, cities promote human creativity for all the following reasons EXCEPT that they

\begin{enumerate}[label=\Alph*.]
    \item contain spaces that enable people to meet and share new ideas.
    \item expose people to different and novel ideas, because they are home to varied groups of people.
    \item provide the financial and institutional networks that enable ideas to become reality.
    \item provide access to cultural activities that promote new and creative ways of thinking.
[](https://cracku.in/1-in-the-authors-view-cities-promote-human-creativit-x-cat-2017-shift-2)
\end{enumerate}
\vspace{1em}
\textbf{Solution:}

In the paragraph starting with 'cities are true fronts of creativity', author mentions that cities have diverse population.   
The author also mentions that cities provide the space where people can meet and share ideas. Then, the author discusses the financial and organizational infrastructure that cities provide for ideas to flourish.  
No where has it been mentioned that cities provide access to cultural activities. We cannot infer option D from the passage. 
Therefore, option D is the right answer.

\end{problem}

\newpage
\begin{problem}{}{}
% Subject: varc
\textbf{Instructions}

Instructions   

**The passage below is accompanied by a set of six questions. Choose the best answer to each question.**
Cities are the true fonts of creativity... With their diverse populations, dense social networks, and public spaces where people can meet spontaneously and serendipitously, they spark and catalyze new ideas. With their infrastructure for finance, organization and trade, they allow those ideas to be swiftly actualized.
As for what staunches creativity, that's easy, if ironic. It's the very institutions that we build to manage, exploit and perpetuate the fruits of creativity — our big bureaucracies, and sad to say, too many of our schools. Creativity is disruptive; schools and organizations are regimented, standardized and stultifying.
The education expert Sir Ken Robinson points to a 1968 study reporting on a group of 1,600 children who were tested over time for their ability to think in out-of-the-box ways. When the children were between 3 and 5 years old, 98 percent achieved positive scores. When they were 8 to 10, only 32 percent passed the same test, and only 10 percent at 13 to 15. When 280,000 25-year-olds took the test, just 2 percent passed. By the time we are adults, our creativity has been wrung out of us.
I once asked the great urbanist Jane Jacobs what makes some places more creative than others. She said, essentially, that the question was an easy one. All cities, she said, were filled with creative people; that's our default state as people. But some cities had more than their shares of leaders, people and institutions that blocked out that creativity. She called them "squelchers."
Creativity (or the lack of it) follows the same general contours of the great socio-economic divide — our rising inequality — that plagues us. According to my own estimates, roughly a third of us across the United States, and perhaps as much as half of us in our most creative cities — are able to do work which engages our creative faculties to some extent, whether as artists, musicians, writers, techies, innovators, entrepreneurs, doctors, lawyers, journalists or educators — those of us who work with our minds. That leaves a group that I term "the other 66 percent," who toil in low-wage rote and rotten jobs — if they have jobs at all — in which their creativity is subjugated, ignored or wasted.
Creativity itself is not in danger. It's flourishing is all around us — in science and technology, arts and culture, in our rapidly revitalizing cities. But we still have a long way to go if we want to build a truly creative society that supports and rewards the creativity of each and every one of us.

\vspace{1em}
\textbf{Question 2}

The author uses 'ironic' in the third paragraph to point out that

\begin{enumerate}[label=\Alph*.]
    \item people need social contact rather than isolation to nurture their creativity.
    \item institutions created to promote creativity eventually stifle it.
    \item the larger the creative population in a city, the more likely it is to be stifled.
    \item large bureaucracies and institutions are the inevitable outcome of successful cities.
[](https://cracku.in/2-the-author-uses-ironic-in-the-third-paragraph-to-p-x-cat-2017-shift-2)
\end{enumerate}
\vspace{1em}
\textbf{Solution:}

'Irony' is a term used to define an activity defeating its very purpose. Therefore, the answer must be along similar lines - a method or activity that stifles its purpose.  
  
In the passage (1968 survey), the author describes how schools and colleges, the institutions that were supposed to foster creativity, stifle it. Also, in the paragraph preceding the paragraph about survey, the author mentions explicitly that the institutes created to promote creativity stifle it. Therefore, option B is the right answer.

\end{problem}

\newpage
\begin{problem}{}{}
% Subject: varc
\textbf{Instructions}

Instructions   

**The passage below is accompanied by a set of six questions. Choose the best answer to each question.**
Cities are the true fonts of creativity... With their diverse populations, dense social networks, and public spaces where people can meet spontaneously and serendipitously, they spark and catalyze new ideas. With their infrastructure for finance, organization and trade, they allow those ideas to be swiftly actualized.
As for what staunches creativity, that's easy, if ironic. It's the very institutions that we build to manage, exploit and perpetuate the fruits of creativity — our big bureaucracies, and sad to say, too many of our schools. Creativity is disruptive; schools and organizations are regimented, standardized and stultifying.
The education expert Sir Ken Robinson points to a 1968 study reporting on a group of 1,600 children who were tested over time for their ability to think in out-of-the-box ways. When the children were between 3 and 5 years old, 98 percent achieved positive scores. When they were 8 to 10, only 32 percent passed the same test, and only 10 percent at 13 to 15. When 280,000 25-year-olds took the test, just 2 percent passed. By the time we are adults, our creativity has been wrung out of us.
I once asked the great urbanist Jane Jacobs what makes some places more creative than others. She said, essentially, that the question was an easy one. All cities, she said, were filled with creative people; that's our default state as people. But some cities had more than their shares of leaders, people and institutions that blocked out that creativity. She called them "squelchers."
Creativity (or the lack of it) follows the same general contours of the great socio-economic divide — our rising inequality — that plagues us. According to my own estimates, roughly a third of us across the United States, and perhaps as much as half of us in our most creative cities — are able to do work which engages our creative faculties to some extent, whether as artists, musicians, writers, techies, innovators, entrepreneurs, doctors, lawyers, journalists or educators — those of us who work with our minds. That leaves a group that I term "the other 66 percent," who toil in low-wage rote and rotten jobs — if they have jobs at all — in which their creativity is subjugated, ignored or wasted.
Creativity itself is not in danger. It's flourishing is all around us — in science and technology, arts and culture, in our rapidly revitalizing cities. But we still have a long way to go if we want to build a truly creative society that supports and rewards the creativity of each and every one of us.

\vspace{1em}
\textbf{Question 3}

The central idea of this passage is that

\begin{enumerate}[label=\Alph*.]
    \item social interaction is necessary to nurture creativity.
    \item creativity and ideas are gradually declining in all societies.
    \item the creativity divide is widening in societies in line with socio-economic trends.
    \item more people should work in jobs that engage their creative faculties.
[](https://cracku.in/3-the-central-idea-of-this-passage-is-that-x-cat-2017-shift-2)
\end{enumerate}
\vspace{1em}
\textbf{Solution:}

The entire passage revolves around how cities provide grounds for creativity to flourish and how our education system stifles it.   

Option B states that creativity and ideas are gradually declining. But, in the last paragraph, the author mentions that 'Creativity itself is not in danger'. Therefore, we can rule out option B.  
  
Option D states that more people must engage in creative jobs. But it cannot be said to be the central idea of the passage. As we have discussed, the passage revolves around social interaction and creativity divide. Therefore, we can eliminate option D too.  
  
Options A and C are close. But, the author describes creativity divide more as an effect than the problem itself. Barring the last 2 paragraphs, the author describes about the importance of social interaction and how the lack of it kills creativity. Since the question is about the central idea, option A can be deemed a better fit than option C. 
Therefore, option A is the right answer.

\end{problem}

\newpage
\begin{problem}{}{}
% Subject: varc
\textbf{Instructions}

Instructions   

**The passage below is accompanied by a set of six questions. Choose the best answer to each question.**
Cities are the true fonts of creativity... With their diverse populations, dense social networks, and public spaces where people can meet spontaneously and serendipitously, they spark and catalyze new ideas. With their infrastructure for finance, organization and trade, they allow those ideas to be swiftly actualized.
As for what staunches creativity, that's easy, if ironic. It's the very institutions that we build to manage, exploit and perpetuate the fruits of creativity — our big bureaucracies, and sad to say, too many of our schools. Creativity is disruptive; schools and organizations are regimented, standardized and stultifying.
The education expert Sir Ken Robinson points to a 1968 study reporting on a group of 1,600 children who were tested over time for their ability to think in out-of-the-box ways. When the children were between 3 and 5 years old, 98 percent achieved positive scores. When they were 8 to 10, only 32 percent passed the same test, and only 10 percent at 13 to 15. When 280,000 25-year-olds took the test, just 2 percent passed. By the time we are adults, our creativity has been wrung out of us.
I once asked the great urbanist Jane Jacobs what makes some places more creative than others. She said, essentially, that the question was an easy one. All cities, she said, were filled with creative people; that's our default state as people. But some cities had more than their shares of leaders, people and institutions that blocked out that creativity. She called them "squelchers."
Creativity (or the lack of it) follows the same general contours of the great socio-economic divide — our rising inequality — that plagues us. According to my own estimates, roughly a third of us across the United States, and perhaps as much as half of us in our most creative cities — are able to do work which engages our creative faculties to some extent, whether as artists, musicians, writers, techies, innovators, entrepreneurs, doctors, lawyers, journalists or educators — those of us who work with our minds. That leaves a group that I term "the other 66 percent," who toil in low-wage rote and rotten jobs — if they have jobs at all — in which their creativity is subjugated, ignored or wasted.
Creativity itself is not in danger. It's flourishing is all around us — in science and technology, arts and culture, in our rapidly revitalizing cities. But we still have a long way to go if we want to build a truly creative society that supports and rewards the creativity of each and every one of us.

\vspace{1em}
\textbf{Question 4}

Jane Jacobs believed that cities that are more creative

\begin{enumerate}[label=\Alph*.]
    \item have to struggle to retain their creativity.
    \item have to 'squelch' unproductive people and promote creative ones.
    \item have leaders and institutions that do not block creativity.
    \item typically do not start off as creative hubs.
[](https://cracku.in/4-jane-jacobs-believed-that-cities-that-are-more-cre-x-cat-2017-shift-2)
\end{enumerate}
\vspace{1em}
\textbf{Solution:}

In the passage, the author clearly describes that Jane Jacobs attributes creativity to the type of leaders. From the paragraph about 'squelchers', we can infer that Jane Jacobs holds leaders responsible for the creativity of the people. Therefore, option C is the right answer.

\end{problem}

\newpage
\begin{problem}{}{}
% Subject: varc
\textbf{Instructions}

Instructions   

**The passage below is accompanied by a set of six questions. Choose the best answer to each question.**
Cities are the true fonts of creativity... With their diverse populations, dense social networks, and public spaces where people can meet spontaneously and serendipitously, they spark and catalyze new ideas. With their infrastructure for finance, organization and trade, they allow those ideas to be swiftly actualized.
As for what staunches creativity, that's easy, if ironic. It's the very institutions that we build to manage, exploit and perpetuate the fruits of creativity — our big bureaucracies, and sad to say, too many of our schools. Creativity is disruptive; schools and organizations are regimented, standardized and stultifying.
The education expert Sir Ken Robinson points to a 1968 study reporting on a group of 1,600 children who were tested over time for their ability to think in out-of-the-box ways. When the children were between 3 and 5 years old, 98 percent achieved positive scores. When they were 8 to 10, only 32 percent passed the same test, and only 10 percent at 13 to 15. When 280,000 25-year-olds took the test, just 2 percent passed. By the time we are adults, our creativity has been wrung out of us.
I once asked the great urbanist Jane Jacobs what makes some places more creative than others. She said, essentially, that the question was an easy one. All cities, she said, were filled with creative people; that's our default state as people. But some cities had more than their shares of leaders, people and institutions that blocked out that creativity. She called them "squelchers."
Creativity (or the lack of it) follows the same general contours of the great socio-economic divide — our rising inequality — that plagues us. According to my own estimates, roughly a third of us across the United States, and perhaps as much as half of us in our most creative cities — are able to do work which engages our creative faculties to some extent, whether as artists, musicians, writers, techies, innovators, entrepreneurs, doctors, lawyers, journalists or educators — those of us who work with our minds. That leaves a group that I term "the other 66 percent," who toil in low-wage rote and rotten jobs — if they have jobs at all — in which their creativity is subjugated, ignored or wasted.
Creativity itself is not in danger. It's flourishing is all around us — in science and technology, arts and culture, in our rapidly revitalizing cities. But we still have a long way to go if we want to build a truly creative society that supports and rewards the creativity of each and every one of us.

\vspace{1em}
\textbf{Question 5}

The 1968 study is used here to show that

\begin{enumerate}[label=\Alph*.]
    \item as they get older, children usually learn to be more creative.
    \item schooling today does not encourage creative thinking in children.
    \item the more children learn, the less creative they become.
    \item technology today prevents children from being creative.
[](https://cracku.in/5-the-1968-study-is-used-here-to-show-that-x-cat-2017-shift-2)
\end{enumerate}
\vspace{1em}
\textbf{Solution:}

There has been no talk about technology in the entire passage. Therefore, we can eliminate option B straight away.  
Also, option A states that children become more creative as they get older. However, the exact opposite has been discussed in the passage. Therefore, we can eliminate option A too.  
Among options B and C, option C attributes reduction in creativity to learning more. But, in the paragraph about 'what staunches creativity', the author mentions that institutions that were created to promote creativity stifle it. He then produces the 1968 study as a validation of the argument. Therefore, the author implies that schools and colleges stifle creativity.
Hence, option B is the right answer.

\end{problem}

\newpage
\begin{problem}{}{}
% Subject: varc
\textbf{Instructions}

Instructions   

**The passage below is accompanied by a set of six questions. Choose the best answer to each question.**
Cities are the true fonts of creativity... With their diverse populations, dense social networks, and public spaces where people can meet spontaneously and serendipitously, they spark and catalyze new ideas. With their infrastructure for finance, organization and trade, they allow those ideas to be swiftly actualized.
As for what staunches creativity, that's easy, if ironic. It's the very institutions that we build to manage, exploit and perpetuate the fruits of creativity — our big bureaucracies, and sad to say, too many of our schools. Creativity is disruptive; schools and organizations are regimented, standardized and stultifying.
The education expert Sir Ken Robinson points to a 1968 study reporting on a group of 1,600 children who were tested over time for their ability to think in out-of-the-box ways. When the children were between 3 and 5 years old, 98 percent achieved positive scores. When they were 8 to 10, only 32 percent passed the same test, and only 10 percent at 13 to 15. When 280,000 25-year-olds took the test, just 2 percent passed. By the time we are adults, our creativity has been wrung out of us.
I once asked the great urbanist Jane Jacobs what makes some places more creative than others. She said, essentially, that the question was an easy one. All cities, she said, were filled with creative people; that's our default state as people. But some cities had more than their shares of leaders, people and institutions that blocked out that creativity. She called them "squelchers."
Creativity (or the lack of it) follows the same general contours of the great socio-economic divide — our rising inequality — that plagues us. According to my own estimates, roughly a third of us across the United States, and perhaps as much as half of us in our most creative cities — are able to do work which engages our creative faculties to some extent, whether as artists, musicians, writers, techies, innovators, entrepreneurs, doctors, lawyers, journalists or educators — those of us who work with our minds. That leaves a group that I term "the other 66 percent," who toil in low-wage rote and rotten jobs — if they have jobs at all — in which their creativity is subjugated, ignored or wasted.
Creativity itself is not in danger. It's flourishing is all around us — in science and technology, arts and culture, in our rapidly revitalizing cities. But we still have a long way to go if we want to build a truly creative society that supports and rewards the creativity of each and every one of us.

\vspace{1em}
\textbf{Question 6}

The author's conclusions about the most 'creative cities' in the US (paragraph 6) are based on his assumption that

\begin{enumerate}[label=\Alph*.]
    \item people who work with their hands are not doing creative work.
    \item more than half the population works in non-creative jobs.
    \item only artists, musicians, writers, and so on should be valued in a society.
    \item most cities ignore or waste the creativity of low-wage workers.
[](https://cracku.in/6-the-authors-conclusions-about-the-most-creative-ci-x-cat-2017-shift-2)
\end{enumerate}
\vspace{1em}
\textbf{Solution:}

In the paragraph regarding creative cities, the author makes a remark that the creativity of only those people who work with their mind are utilized. Therefore, we can infer that the author thinks that the creativity of people who do not work with their minds (who work with their hands) is not utilized. Therefore, option A is the right answer.

Instructions   

**The passage below is accompanied by a set of six questions. Choose the best answer to each question.**
  
During the frigid season... it's often necessary to nestle under a blanket to try to stay warm. The temperature difference between the blanket and the air outside is so palpable that we often have trouble leaving our warm refuge. Many plants and animals similarly hunker down, relying on snow cover for safety from winter's harsh conditions. The small area between the snowpack and the ground, called the subnivium... might be the most important ecosystem that you have never heard of.
The subnivium is so well-insulated and stable that its temperature holds steady at around 32 degree Fahrenheit (0 degree Celsius). Although that might still sound cold, a constant temperature of 32 degree Fahrenheit can often be 30 to 40 degrees warmer than the air temperature during the peak of winter. Because of this large temperature difference, a wide variety of species...depend on the subnivium for winter protection.
For many organisms living in temperate and Arctic regions, the difference between being under the snow or outside it is a matter of life and death. Consequently, disruptions to the subnivium brought about by climate change will affect everything from population dynamics to nutrient cycling through the ecosystem.
The formation and stability of the subnivium requires more than a few flurries. Winter ecologists have suggested that eight inches of snow is necessary to develop a stable layer of insulation. Depth is not the only factor, however. More accurately, the stability of the subnivium depends on the interaction between snow depth and snow density. Imagine being under a stack of blankets that are all flattened and pressed together. When compressed, the blankets essentially form one compacted layer. In contrast, when they are lightly placed on top of one another, their insulative capacity increases because the air pockets between them trap heat. Greater depths of low-density snow are therefore better at insulating the ground.
Both depth and density of snow are sensitive to temperature. Scientists are now beginning to explore how climate change will affect the subnivium, as well as the species that depend on it. At first glance, warmer winters seem beneficial for species that have difficulty surviving subzero temperatures; however, as with most ecological phenomena, the consequences are not so straightforward. Research has shown that the snow season (the period when snow is more likely than rain) has become shorter since l970. When rain falls on snow, it increases the density of the snow and reduces its insulative capacity. Therefore, even though winters are expected to become warmer overall from future climate change, the subnivium will tend to become colder and more variable with less protection from the above-ground temperatures.
The effects of a colder subnivium are complex... For example, shrubs such as crowberry and alpine azalea that grow along the forest floor tend to block the wind and so retain higher depths of snow around them. This captured snow helps to keep soils insulated and in turn increases plant decomposition and nutrient release. In field experiments, researchers removed a portion. of the snow cover to investigate the importance of the subnivium's insulation. They found that soil frost in the snow-free area resulted in damage to plant roots and sometimes even the death of the plant.

\end{problem}

\newpage
\begin{problem}{}{}
% Subject: varc
\textbf{Instructions}

Instructions   

**The passage below is accompanied by a set of six questions. Choose the best answer to each question.**
Cities are the true fonts of creativity... With their diverse populations, dense social networks, and public spaces where people can meet spontaneously and serendipitously, they spark and catalyze new ideas. With their infrastructure for finance, organization and trade, they allow those ideas to be swiftly actualized.
As for what staunches creativity, that's easy, if ironic. It's the very institutions that we build to manage, exploit and perpetuate the fruits of creativity — our big bureaucracies, and sad to say, too many of our schools. Creativity is disruptive; schools and organizations are regimented, standardized and stultifying.
The education expert Sir Ken Robinson points to a 1968 study reporting on a group of 1,600 children who were tested over time for their ability to think in out-of-the-box ways. When the children were between 3 and 5 years old, 98 percent achieved positive scores. When they were 8 to 10, only 32 percent passed the same test, and only 10 percent at 13 to 15. When 280,000 25-year-olds took the test, just 2 percent passed. By the time we are adults, our creativity has been wrung out of us.
I once asked the great urbanist Jane Jacobs what makes some places more creative than others. She said, essentially, that the question was an easy one. All cities, she said, were filled with creative people; that's our default state as people. But some cities had more than their shares of leaders, people and institutions that blocked out that creativity. She called them "squelchers."
Creativity (or the lack of it) follows the same general contours of the great socio-economic divide — our rising inequality — that plagues us. According to my own estimates, roughly a third of us across the United States, and perhaps as much as half of us in our most creative cities — are able to do work which engages our creative faculties to some extent, whether as artists, musicians, writers, techies, innovators, entrepreneurs, doctors, lawyers, journalists or educators — those of us who work with our minds. That leaves a group that I term "the other 66 percent," who toil in low-wage rote and rotten jobs — if they have jobs at all — in which their creativity is subjugated, ignored or wasted.
Creativity itself is not in danger. It's flourishing is all around us — in science and technology, arts and culture, in our rapidly revitalizing cities. But we still have a long way to go if we want to build a truly creative society that supports and rewards the creativity of each and every one of us.

\vspace{1em}
\textbf{Question 7}

The purpose of this passage is to

\begin{enumerate}[label=\Alph*.]
    \item introduce readers to a relatively unknown ecosystem: the subnivium.
    \item explain how the subnivium works to provide shelter and food to several species.
    \item outline the effects of climate change on the subnivium.
    \item draw an analogy between the effect of blankets on humans and of snow cover on species living in the subnivium.
[](https://cracku.in/7-the-purpose-of-this-passage-is-to-x-cat-2017-shift-2)
\end{enumerate}
\vspace{1em}
\textbf{Solution:}

The entire passage revolves around the effects of climate change on subnivium.  
We can eliminate option D directly as it talks about a small illustration. It cannot be said to be the purpose of the passage. Options A and B emphasize subnivium as the subject. However, the passage is about the effects of climate change on subnivium rather than subnivium itself. Throughout the passage, the author discusses the effects of various climatic changes and how it affects the subnivium.  
  
Therefore, option C is the right answer.

\end{problem}

\newpage
\begin{problem}{}{}
% Subject: varc
\textbf{Instructions}

Instructions   

**The passage below is accompanied by a set of six questions. Choose the best answer to each question.**
Cities are the true fonts of creativity... With their diverse populations, dense social networks, and public spaces where people can meet spontaneously and serendipitously, they spark and catalyze new ideas. With their infrastructure for finance, organization and trade, they allow those ideas to be swiftly actualized.
As for what staunches creativity, that's easy, if ironic. It's the very institutions that we build to manage, exploit and perpetuate the fruits of creativity — our big bureaucracies, and sad to say, too many of our schools. Creativity is disruptive; schools and organizations are regimented, standardized and stultifying.
The education expert Sir Ken Robinson points to a 1968 study reporting on a group of 1,600 children who were tested over time for their ability to think in out-of-the-box ways. When the children were between 3 and 5 years old, 98 percent achieved positive scores. When they were 8 to 10, only 32 percent passed the same test, and only 10 percent at 13 to 15. When 280,000 25-year-olds took the test, just 2 percent passed. By the time we are adults, our creativity has been wrung out of us.
I once asked the great urbanist Jane Jacobs what makes some places more creative than others. She said, essentially, that the question was an easy one. All cities, she said, were filled with creative people; that's our default state as people. But some cities had more than their shares of leaders, people and institutions that blocked out that creativity. She called them "squelchers."
Creativity (or the lack of it) follows the same general contours of the great socio-economic divide — our rising inequality — that plagues us. According to my own estimates, roughly a third of us across the United States, and perhaps as much as half of us in our most creative cities — are able to do work which engages our creative faculties to some extent, whether as artists, musicians, writers, techies, innovators, entrepreneurs, doctors, lawyers, journalists or educators — those of us who work with our minds. That leaves a group that I term "the other 66 percent," who toil in low-wage rote and rotten jobs — if they have jobs at all — in which their creativity is subjugated, ignored or wasted.
Creativity itself is not in danger. It's flourishing is all around us — in science and technology, arts and culture, in our rapidly revitalizing cities. But we still have a long way to go if we want to build a truly creative society that supports and rewards the creativity of each and every one of us.

\vspace{1em}
\textbf{Question 8}

All of the following statements are true EXCEPT

\begin{enumerate}[label=\Alph*.]
    \item Snow depth and Snow density both influence the stability of the subnivium.
    \item Climate change has some positive effects on the subnivium.
    \item The subnivium maintains a steady temperature that can be 30 to 40 degrees warmer than the winter air temperature.
    \item Researchers have established the adverse effects of dwindling snow cover on the subnivium.
[](https://cracku.in/8-all-of-the-following-statements-are-true-except-x-cat-2017-shift-2)
\end{enumerate}
\vspace{1em}
\textbf{Solution:}

The author mentions that 'Both depth and density of snow are sensitive to temperature.' Therefore, we can easily eliminate option A.  
Option C talks about the insulating properties of subnivium which has been explicitly mentioned in the passage - 'Although that might still sound cold, a constant temperature of 32°F can often be 30 to 40 degrees warmer than the air temperature during the peak of winter.' Therefore, we can eliminate option C too.  
  
Option D states that researchers have established the adverse effects of the dwindling snow cover in subnivium. From the line starting with 'research has shown that...', we can infer that the effects of the dwindling snow cover on subnivium has been established.  
  
The entire passage does not discuss any positive effect of climate change on the subnivium. Therefore, we can say that option B is the right answer.

\end{problem}

\newpage
\begin{problem}{}{}
% Subject: varc
\textbf{Instructions}

Instructions   

**The passage below is accompanied by a set of six questions. Choose the best answer to each question.**
Cities are the true fonts of creativity... With their diverse populations, dense social networks, and public spaces where people can meet spontaneously and serendipitously, they spark and catalyze new ideas. With their infrastructure for finance, organization and trade, they allow those ideas to be swiftly actualized.
As for what staunches creativity, that's easy, if ironic. It's the very institutions that we build to manage, exploit and perpetuate the fruits of creativity — our big bureaucracies, and sad to say, too many of our schools. Creativity is disruptive; schools and organizations are regimented, standardized and stultifying.
The education expert Sir Ken Robinson points to a 1968 study reporting on a group of 1,600 children who were tested over time for their ability to think in out-of-the-box ways. When the children were between 3 and 5 years old, 98 percent achieved positive scores. When they were 8 to 10, only 32 percent passed the same test, and only 10 percent at 13 to 15. When 280,000 25-year-olds took the test, just 2 percent passed. By the time we are adults, our creativity has been wrung out of us.
I once asked the great urbanist Jane Jacobs what makes some places more creative than others. She said, essentially, that the question was an easy one. All cities, she said, were filled with creative people; that's our default state as people. But some cities had more than their shares of leaders, people and institutions that blocked out that creativity. She called them "squelchers."
Creativity (or the lack of it) follows the same general contours of the great socio-economic divide — our rising inequality — that plagues us. According to my own estimates, roughly a third of us across the United States, and perhaps as much as half of us in our most creative cities — are able to do work which engages our creative faculties to some extent, whether as artists, musicians, writers, techies, innovators, entrepreneurs, doctors, lawyers, journalists or educators — those of us who work with our minds. That leaves a group that I term "the other 66 percent," who toil in low-wage rote and rotten jobs — if they have jobs at all — in which their creativity is subjugated, ignored or wasted.
Creativity itself is not in danger. It's flourishing is all around us — in science and technology, arts and culture, in our rapidly revitalizing cities. But we still have a long way to go if we want to build a truly creative society that supports and rewards the creativity of each and every one of us.

\vspace{1em}
\textbf{Question 9}

Based on this extract, the author would support which one of the following actions?

\begin{enumerate}[label=\Alph*.]
    \item The use of snow machines in winter to ensure snow cover of at least eight inches.
    \item Government action to curb climate change.
    \item Adding nutrients to the soil in winter.
    \item Planting more shrubs in areas of short snow season.
[](https://cracku.in/9-based-on-this-extract-the-author-would-support-whi-x-cat-2017-shift-2)
\end{enumerate}
\vspace{1em}
\textbf{Solution:}

The author mentions in the passage that the quality of snow also plays a vital role. Therefore, maintaining 8 inches of snow with a machine will not fix the problem. Moreover, the option feels too shallow and unsustainable. Therefore, we can eliminate option A.  

Option C also feels shallow and unrealistic. Moreover, it has not been mentioned that adding nutrients will fix the issue. 
Option D suggests planting shrubs. But, in the last paragraph the author mentions that the effects are multilayered and complex. Options A, C, and D try to address the symptom than attacking the cause. Option B offers a more viable solution and addresses the cause of the issue rather than its manifestation. Therefore, the author is most likely to agree with option B and hence, option B is the right answer.

\end{problem}

\newpage
\begin{problem}{}{}
% Subject: varc
\textbf{Instructions}

Instructions   

**The passage below is accompanied by a set of six questions. Choose the best answer to each question.**
Cities are the true fonts of creativity... With their diverse populations, dense social networks, and public spaces where people can meet spontaneously and serendipitously, they spark and catalyze new ideas. With their infrastructure for finance, organization and trade, they allow those ideas to be swiftly actualized.
As for what staunches creativity, that's easy, if ironic. It's the very institutions that we build to manage, exploit and perpetuate the fruits of creativity — our big bureaucracies, and sad to say, too many of our schools. Creativity is disruptive; schools and organizations are regimented, standardized and stultifying.
The education expert Sir Ken Robinson points to a 1968 study reporting on a group of 1,600 children who were tested over time for their ability to think in out-of-the-box ways. When the children were between 3 and 5 years old, 98 percent achieved positive scores. When they were 8 to 10, only 32 percent passed the same test, and only 10 percent at 13 to 15. When 280,000 25-year-olds took the test, just 2 percent passed. By the time we are adults, our creativity has been wrung out of us.
I once asked the great urbanist Jane Jacobs what makes some places more creative than others. She said, essentially, that the question was an easy one. All cities, she said, were filled with creative people; that's our default state as people. But some cities had more than their shares of leaders, people and institutions that blocked out that creativity. She called them "squelchers."
Creativity (or the lack of it) follows the same general contours of the great socio-economic divide — our rising inequality — that plagues us. According to my own estimates, roughly a third of us across the United States, and perhaps as much as half of us in our most creative cities — are able to do work which engages our creative faculties to some extent, whether as artists, musicians, writers, techies, innovators, entrepreneurs, doctors, lawyers, journalists or educators — those of us who work with our minds. That leaves a group that I term "the other 66 percent," who toil in low-wage rote and rotten jobs — if they have jobs at all — in which their creativity is subjugated, ignored or wasted.
Creativity itself is not in danger. It's flourishing is all around us — in science and technology, arts and culture, in our rapidly revitalizing cities. But we still have a long way to go if we want to build a truly creative society that supports and rewards the creativity of each and every one of us.

\vspace{1em}
\textbf{Question 10}

In paragraph 6, the author provides the examples of crowberry and alpine azalea to demonstrate that

\begin{enumerate}[label=\Alph*.]
    \item Despite frigid temperatures, several species survive in temperate and Arctic regions.
    \item Due to frigid temperatures in the temperate and Arctic regions, plant species that survive tend to be shrubs rather than trees.
    \item The crowberry and alpine azalea are abundant in temperate and Arctic regions.
    \item The stability of the subnivium depends on several interrelated factors, including shrubs on the forest floor.
[](https://cracku.in/10-in-paragraph-6-the-author-provides-the-examples-of-x-cat-2017-shift-2)
\end{enumerate}
\vspace{1em}
\textbf{Solution:}

The reason for the inclusion of the shrubs must be in line with the central idea of the passage. Options A and C are too general and hence, can be ruled out easily. Option B states that plants that tend to survive turn out to be shrubs. But, it has not been mentioned anywhere in the passage.  
The last paragraph clearly mentions that the effects of colder subnivium are multilayered and interrelated. The shrubs tend to prove the point. The paragraph discusses the effect on the shrubs in detail, adding substance to the statement.  
  
Therefore, option D is the right answer.

\end{problem}

\newpage
\begin{problem}{}{}
% Subject: varc
\textbf{Instructions}

Instructions   

**The passage below is accompanied by a set of six questions. Choose the best answer to each question.**
Cities are the true fonts of creativity... With their diverse populations, dense social networks, and public spaces where people can meet spontaneously and serendipitously, they spark and catalyze new ideas. With their infrastructure for finance, organization and trade, they allow those ideas to be swiftly actualized.
As for what staunches creativity, that's easy, if ironic. It's the very institutions that we build to manage, exploit and perpetuate the fruits of creativity — our big bureaucracies, and sad to say, too many of our schools. Creativity is disruptive; schools and organizations are regimented, standardized and stultifying.
The education expert Sir Ken Robinson points to a 1968 study reporting on a group of 1,600 children who were tested over time for their ability to think in out-of-the-box ways. When the children were between 3 and 5 years old, 98 percent achieved positive scores. When they were 8 to 10, only 32 percent passed the same test, and only 10 percent at 13 to 15. When 280,000 25-year-olds took the test, just 2 percent passed. By the time we are adults, our creativity has been wrung out of us.
I once asked the great urbanist Jane Jacobs what makes some places more creative than others. She said, essentially, that the question was an easy one. All cities, she said, were filled with creative people; that's our default state as people. But some cities had more than their shares of leaders, people and institutions that blocked out that creativity. She called them "squelchers."
Creativity (or the lack of it) follows the same general contours of the great socio-economic divide — our rising inequality — that plagues us. According to my own estimates, roughly a third of us across the United States, and perhaps as much as half of us in our most creative cities — are able to do work which engages our creative faculties to some extent, whether as artists, musicians, writers, techies, innovators, entrepreneurs, doctors, lawyers, journalists or educators — those of us who work with our minds. That leaves a group that I term "the other 66 percent," who toil in low-wage rote and rotten jobs — if they have jobs at all — in which their creativity is subjugated, ignored or wasted.
Creativity itself is not in danger. It's flourishing is all around us — in science and technology, arts and culture, in our rapidly revitalizing cities. But we still have a long way to go if we want to build a truly creative society that supports and rewards the creativity of each and every one of us.

\vspace{1em}
\textbf{Question 11}

Which one of the following statements can be inferred from the passage?

\begin{enumerate}[label=\Alph*.]
    \item In an ecosystem, altering any one element has a ripple effect on all others.
    \item Climate change affects temperate and Artie regions more than equatorial or arid ones.
    \item A compact layer of wool is warmer than a similarly compact layer of goose down.
    \item The loss of the subnivium, while tragic, will affect only temperate and Artic regions.
[](https://cracku.in/11-which-one-of-the-following-statements-can-be-infer-x-cat-2017-shift-2)
\end{enumerate}
\vspace{1em}
\textbf{Solution:}

Options B and D mention that it will be the arctic and the temperate regions that will be affected. Though we do not know the effect of climate change on the tropical regions, we cannot claim that there will be no effects. The passage does not give us sufficient information to make that claim. Therefore, we can rule out options B and D. 
Option C states that a compact layer of wool is warmer than a similarly compact layer of goose down. Again, the passage does not provide us with sufficient information to substantiate this claim. We do not have sufficient details to compare 2 different materials. Therefore, option C can be ruled out as well.  
  
Option A talks about ripple effect. The entire passage is about how the effects of climate change are interrelated. Ripple effect also discusses the same. Therefore, option A is the right answer.

\end{problem}

\newpage
\begin{problem}{}{}
% Subject: varc
\textbf{Instructions}

Instructions   

**The passage below is accompanied by a set of six questions. Choose the best answer to each question.**
Cities are the true fonts of creativity... With their diverse populations, dense social networks, and public spaces where people can meet spontaneously and serendipitously, they spark and catalyze new ideas. With their infrastructure for finance, organization and trade, they allow those ideas to be swiftly actualized.
As for what staunches creativity, that's easy, if ironic. It's the very institutions that we build to manage, exploit and perpetuate the fruits of creativity — our big bureaucracies, and sad to say, too many of our schools. Creativity is disruptive; schools and organizations are regimented, standardized and stultifying.
The education expert Sir Ken Robinson points to a 1968 study reporting on a group of 1,600 children who were tested over time for their ability to think in out-of-the-box ways. When the children were between 3 and 5 years old, 98 percent achieved positive scores. When they were 8 to 10, only 32 percent passed the same test, and only 10 percent at 13 to 15. When 280,000 25-year-olds took the test, just 2 percent passed. By the time we are adults, our creativity has been wrung out of us.
I once asked the great urbanist Jane Jacobs what makes some places more creative than others. She said, essentially, that the question was an easy one. All cities, she said, were filled with creative people; that's our default state as people. But some cities had more than their shares of leaders, people and institutions that blocked out that creativity. She called them "squelchers."
Creativity (or the lack of it) follows the same general contours of the great socio-economic divide — our rising inequality — that plagues us. According to my own estimates, roughly a third of us across the United States, and perhaps as much as half of us in our most creative cities — are able to do work which engages our creative faculties to some extent, whether as artists, musicians, writers, techies, innovators, entrepreneurs, doctors, lawyers, journalists or educators — those of us who work with our minds. That leaves a group that I term "the other 66 percent," who toil in low-wage rote and rotten jobs — if they have jobs at all — in which their creativity is subjugated, ignored or wasted.
Creativity itself is not in danger. It's flourishing is all around us — in science and technology, arts and culture, in our rapidly revitalizing cities. But we still have a long way to go if we want to build a truly creative society that supports and rewards the creativity of each and every one of us.

\vspace{1em}
\textbf{Question 12}

In paragraph 1, the author uses blankets as a device to

\begin{enumerate}[label=\Alph*.]
    \item evoke the bitter cold of winter in the minds of readers.
    \item explain how blankets work to keep us warm.
    \item draw an analogy between blankets and the snow pack.
    \item alert readers to the fatal effects of excessive exposure to the cold.
[](https://cracku.in/12-in-paragraph-1-the-author-uses-blankets-as-a-devic-x-cat-2017-shift-2)
\end{enumerate}
\vspace{1em}
\textbf{Solution:}

In the passage, author uses the example to explain how having some spaces between layers increases the insulating property. He then uses the same logic to explain the effects of increase in density of snow on subnivium. Therefore, the author uses the example to draw an analogy. Therefore, option C is the right answer.

Instructions   

**The passage below is accompanied by a set of six questions. Choose the best answer to each question.**
Some of it reflects the remarkable abilities of Elon Musk, the company's founder, as a salesman, engineer, and a man able to get the most out his factory workers and the governments he deals with...Mr Musk is selling a dream that the world wants to believe in. 
This last may be the most important factor in the story. The private car is...a device of immense practical help and economic significance, but at the same time a theatre for myths of unattainable self-fulfilment. The one thing you will never see in a car advertisement is traffic, even though that is the element in which drivers spend their lives. Every single driver in a traffic jam is trying to escape from it, yet it is the inevitable consequence of mass car ownership.
The sleek and swift electric car is at one level merely the most contemporary fantasy of autonomy and power. But it might also disrupt our exterior landscapes nearly as much as the fossil fuel-engined car did in the last century. Electrical cars would of course pollute far less than fossil fuel-driven ones; instead of oil reserves, the rarest materials for batteries would make undeserving despots and their dynasties fantastically rich. Petrol stations would disappear. The air in cities would once more be breathable and their streets as quiet as those of Venice. This isn't an unmixed good. Cars that were as silent as bicycles would still be as dangerous as they are now to anyone they hit without audible warning.
The dream goes further than that. The electric cars of the future will be so thoroughly equipped with sensors and reaction mechanisms that they will never hit anyone. Just as brakes don't let you skid today, the steering wheel of tomorrow will swerve you away from danger before you have even noticed it...
This is where the fantasy of autonomy comes full circle. The logical outcome of cars which need no driver is that they will become cars which need no owner either. Instead, they will work as taxis do, summoned at will but only for the journeys we actually need. This the future towards which Uber...is working. The ultimate development of the private car will be to reinvent public transport. Traffic jams will be abolished only when the private car becomes a public utility. What then will happen to our fantasies of independence? We' ll all have to take to electrically powered bicycles.

\end{problem}

\newpage
\begin{problem}{}{}
% Subject: varc
\textbf{Instructions}

Instructions   

**The passage below is accompanied by a set of six questions. Choose the best answer to each question.**
Cities are the true fonts of creativity... With their diverse populations, dense social networks, and public spaces where people can meet spontaneously and serendipitously, they spark and catalyze new ideas. With their infrastructure for finance, organization and trade, they allow those ideas to be swiftly actualized.
As for what staunches creativity, that's easy, if ironic. It's the very institutions that we build to manage, exploit and perpetuate the fruits of creativity — our big bureaucracies, and sad to say, too many of our schools. Creativity is disruptive; schools and organizations are regimented, standardized and stultifying.
The education expert Sir Ken Robinson points to a 1968 study reporting on a group of 1,600 children who were tested over time for their ability to think in out-of-the-box ways. When the children were between 3 and 5 years old, 98 percent achieved positive scores. When they were 8 to 10, only 32 percent passed the same test, and only 10 percent at 13 to 15. When 280,000 25-year-olds took the test, just 2 percent passed. By the time we are adults, our creativity has been wrung out of us.
I once asked the great urbanist Jane Jacobs what makes some places more creative than others. She said, essentially, that the question was an easy one. All cities, she said, were filled with creative people; that's our default state as people. But some cities had more than their shares of leaders, people and institutions that blocked out that creativity. She called them "squelchers."
Creativity (or the lack of it) follows the same general contours of the great socio-economic divide — our rising inequality — that plagues us. According to my own estimates, roughly a third of us across the United States, and perhaps as much as half of us in our most creative cities — are able to do work which engages our creative faculties to some extent, whether as artists, musicians, writers, techies, innovators, entrepreneurs, doctors, lawyers, journalists or educators — those of us who work with our minds. That leaves a group that I term "the other 66 percent," who toil in low-wage rote and rotten jobs — if they have jobs at all — in which their creativity is subjugated, ignored or wasted.
Creativity itself is not in danger. It's flourishing is all around us — in science and technology, arts and culture, in our rapidly revitalizing cities. But we still have a long way to go if we want to build a truly creative society that supports and rewards the creativity of each and every one of us.

\vspace{1em}
\textbf{Question 13}

Which of the following statements best reflects the author's argument?

\begin{enumerate}[label=\Alph*.]
    \item Hybrid and electric vehicles signal the end of the age of internal combustion engines.
    \item Elon Musk is a remarkably gifted salesman.
    \item The private car represents an unattainable myth of independence.
    \item The future Uber car will be environmentally friendlier than even the Tesla.
[](https://cracku.in/13-which-of-the-following-statements-best-reflects-th-x-cat-2017-shift-2)
\end{enumerate}
\vspace{1em}
\textbf{Solution:}

The main argument around which the passage revolves can be referred from following statemenst made by the author "...The private car is...a device of immense practical ... but at the same time a theatre for myths of unattainable self-fulfilment... ". In the later part, the author substanitiates this argument. Hence, option C, which states the same, must be the right answer.

\end{problem}

\newpage
\begin{problem}{}{}
% Subject: varc
\textbf{Instructions}

Instructions   

**The passage below is accompanied by a set of six questions. Choose the best answer to each question.**
Cities are the true fonts of creativity... With their diverse populations, dense social networks, and public spaces where people can meet spontaneously and serendipitously, they spark and catalyze new ideas. With their infrastructure for finance, organization and trade, they allow those ideas to be swiftly actualized.
As for what staunches creativity, that's easy, if ironic. It's the very institutions that we build to manage, exploit and perpetuate the fruits of creativity — our big bureaucracies, and sad to say, too many of our schools. Creativity is disruptive; schools and organizations are regimented, standardized and stultifying.
The education expert Sir Ken Robinson points to a 1968 study reporting on a group of 1,600 children who were tested over time for their ability to think in out-of-the-box ways. When the children were between 3 and 5 years old, 98 percent achieved positive scores. When they were 8 to 10, only 32 percent passed the same test, and only 10 percent at 13 to 15. When 280,000 25-year-olds took the test, just 2 percent passed. By the time we are adults, our creativity has been wrung out of us.
I once asked the great urbanist Jane Jacobs what makes some places more creative than others. She said, essentially, that the question was an easy one. All cities, she said, were filled with creative people; that's our default state as people. But some cities had more than their shares of leaders, people and institutions that blocked out that creativity. She called them "squelchers."
Creativity (or the lack of it) follows the same general contours of the great socio-economic divide — our rising inequality — that plagues us. According to my own estimates, roughly a third of us across the United States, and perhaps as much as half of us in our most creative cities — are able to do work which engages our creative faculties to some extent, whether as artists, musicians, writers, techies, innovators, entrepreneurs, doctors, lawyers, journalists or educators — those of us who work with our minds. That leaves a group that I term "the other 66 percent," who toil in low-wage rote and rotten jobs — if they have jobs at all — in which their creativity is subjugated, ignored or wasted.
Creativity itself is not in danger. It's flourishing is all around us — in science and technology, arts and culture, in our rapidly revitalizing cities. But we still have a long way to go if we want to build a truly creative society that supports and rewards the creativity of each and every one of us.

\vspace{1em}
\textbf{Question 14}

The author points out all of the following about electric cars EXCEPT

\begin{enumerate}[label=\Alph*.]
    \item Their reliance on rare materials for batteries will support despotic rule.
    \item They will reduce air and noise pollution.
    \item They will not decrease the number of traffic jams.
    \item They will ultimately undermine rather than further driver autonomy.
[](https://cracku.in/14-the-author-points-out-all-of-the-following-about-e-x-cat-2017-shift-2)
\end{enumerate}
\vspace{1em}
\textbf{Solution:}

Refer to the following lines, " ...the rarest materials for batteries would make undeserving despots and their dynasties fantastically rich..." Thus, the author states that relying on rare materials will support despotic rule. Hence, option A is correct and can be eliminated.
Refer to the following lines, "...The air in cities would once more be breathable and their streets as quiet as .. Cars that were as silent as bicycles ...". Hence, option B is correct and can be eliminated.
Refer to the following lines, "...traffic jam is trying to escape from it, yet it is the inevitable consequence of mass car ownership..." Hence, the author points out that the problem of traffic jam which persists with electric car as it is a consequence of mass car ownership. Hence, option C is correct and can be eliminated.
Option D, which states that electric cars will undermine the driver autonomy, cannot be inferred from the passage. Hence, option D is the right choice.

\end{problem}

\newpage
\begin{problem}{}{}
% Subject: varc
\textbf{Instructions}

Instructions   

**The passage below is accompanied by a set of six questions. Choose the best answer to each question.**
Cities are the true fonts of creativity... With their diverse populations, dense social networks, and public spaces where people can meet spontaneously and serendipitously, they spark and catalyze new ideas. With their infrastructure for finance, organization and trade, they allow those ideas to be swiftly actualized.
As for what staunches creativity, that's easy, if ironic. It's the very institutions that we build to manage, exploit and perpetuate the fruits of creativity — our big bureaucracies, and sad to say, too many of our schools. Creativity is disruptive; schools and organizations are regimented, standardized and stultifying.
The education expert Sir Ken Robinson points to a 1968 study reporting on a group of 1,600 children who were tested over time for their ability to think in out-of-the-box ways. When the children were between 3 and 5 years old, 98 percent achieved positive scores. When they were 8 to 10, only 32 percent passed the same test, and only 10 percent at 13 to 15. When 280,000 25-year-olds took the test, just 2 percent passed. By the time we are adults, our creativity has been wrung out of us.
I once asked the great urbanist Jane Jacobs what makes some places more creative than others. She said, essentially, that the question was an easy one. All cities, she said, were filled with creative people; that's our default state as people. But some cities had more than their shares of leaders, people and institutions that blocked out that creativity. She called them "squelchers."
Creativity (or the lack of it) follows the same general contours of the great socio-economic divide — our rising inequality — that plagues us. According to my own estimates, roughly a third of us across the United States, and perhaps as much as half of us in our most creative cities — are able to do work which engages our creative faculties to some extent, whether as artists, musicians, writers, techies, innovators, entrepreneurs, doctors, lawyers, journalists or educators — those of us who work with our minds. That leaves a group that I term "the other 66 percent," who toil in low-wage rote and rotten jobs — if they have jobs at all — in which their creativity is subjugated, ignored or wasted.
Creativity itself is not in danger. It's flourishing is all around us — in science and technology, arts and culture, in our rapidly revitalizing cities. But we still have a long way to go if we want to build a truly creative society that supports and rewards the creativity of each and every one of us.

\vspace{1em}
\textbf{Question 15}

According to the author, the main reason for Tesla's remarkable sales is that

\begin{enumerate}[label=\Alph*.]
    \item in the long run, the Tesla is more cost effective than fossil fuel-driven cars.
    \item the US government has announced a tax subsidy for Tesla buyers.
    \item the company is rapidly upscaling the number of specialised charging stations for customer convenience.
    \item people believe in the autonomy represented by private cars.
[](https://cracku.in/15-according-to-the-author-the-main-reason-for-teslas-x-cat-2017-shift-2)
\end{enumerate}
\vspace{1em}
\textbf{Solution:}

Refer to following lines, "...Mr Musk is selling a dream that the world ...This last may be the most important factor in the story. The private car..". Since people believe in the autonomy represented by private cars, Tesla had remarkable sales. 
Option A is nowhere stated as the reason for high sales for Tesla. Hence, it can be eliminated.
Nowhere in the passage has the author mentioned about tax subsidy to Tesla buyers. Hence, option B can be eliminated.
The author states that there are limited number of specialised charging stations. But the author doesn't state that there will be significant increase in charging stations rapidly which would increase the sales. Hence, option C can be eliminated.
Option D correctly highlights the main reason for Tesla's remarkable sales. Hence, option D is the right choice.

\end{problem}

\newpage
\begin{problem}{}{}
% Subject: varc
\textbf{Instructions}

Instructions   

**The passage below is accompanied by a set of six questions. Choose the best answer to each question.**
Cities are the true fonts of creativity... With their diverse populations, dense social networks, and public spaces where people can meet spontaneously and serendipitously, they spark and catalyze new ideas. With their infrastructure for finance, organization and trade, they allow those ideas to be swiftly actualized.
As for what staunches creativity, that's easy, if ironic. It's the very institutions that we build to manage, exploit and perpetuate the fruits of creativity — our big bureaucracies, and sad to say, too many of our schools. Creativity is disruptive; schools and organizations are regimented, standardized and stultifying.
The education expert Sir Ken Robinson points to a 1968 study reporting on a group of 1,600 children who were tested over time for their ability to think in out-of-the-box ways. When the children were between 3 and 5 years old, 98 percent achieved positive scores. When they were 8 to 10, only 32 percent passed the same test, and only 10 percent at 13 to 15. When 280,000 25-year-olds took the test, just 2 percent passed. By the time we are adults, our creativity has been wrung out of us.
I once asked the great urbanist Jane Jacobs what makes some places more creative than others. She said, essentially, that the question was an easy one. All cities, she said, were filled with creative people; that's our default state as people. But some cities had more than their shares of leaders, people and institutions that blocked out that creativity. She called them "squelchers."
Creativity (or the lack of it) follows the same general contours of the great socio-economic divide — our rising inequality — that plagues us. According to my own estimates, roughly a third of us across the United States, and perhaps as much as half of us in our most creative cities — are able to do work which engages our creative faculties to some extent, whether as artists, musicians, writers, techies, innovators, entrepreneurs, doctors, lawyers, journalists or educators — those of us who work with our minds. That leaves a group that I term "the other 66 percent," who toil in low-wage rote and rotten jobs — if they have jobs at all — in which their creativity is subjugated, ignored or wasted.
Creativity itself is not in danger. It's flourishing is all around us — in science and technology, arts and culture, in our rapidly revitalizing cities. But we still have a long way to go if we want to build a truly creative society that supports and rewards the creativity of each and every one of us.

\vspace{1em}
\textbf{Question 16}

The author comes to the conclusion that

\begin{enumerate}[label=\Alph*.]
    \item car drivers will no longer own cars but will have to use public transport.
    \item cars will be controlled by technology that is more efficient than car drivers.
    \item car drivers dream of autonomy but the future may be public transport.
    \item electrically powered bicycles are the only way to achieve autonomy in transportation.
[](https://cracku.in/16-the-author-comes-to-the-conclusion-that-x-cat-2017-shift-2)
\end{enumerate}
\vspace{1em}
\textbf{Solution:}

Refer to the following lines, "...Traffic jams will be abolished only when the private car becomes a public utility...". According to the author the fantasy of autonomy comes full circle. Thus the author states that car drivers want autonomy but public transport will be the future as only then the traffic problem will be solved. 
Option A, which states that car drivers will no longer own cars, is too extreme and can be ruled out.
Option B specifies something that is out of the scope of the paragraph. Hence, option B is incorrect. 
Option C completely reflects what the author says in the end. Hence, it is the right choice.
Option D is again too extreme and cannot be found in the passage. Also, the author mentions electric bicycles just to provide an illustration. Hence,it can be eliminated. 
Hence, option C is the right answer.

\end{problem}

\newpage
\begin{problem}{}{}
% Subject: varc
\textbf{Instructions}

Instructions   

**The passage below is accompanied by a set of six questions. Choose the best answer to each question.**
Cities are the true fonts of creativity... With their diverse populations, dense social networks, and public spaces where people can meet spontaneously and serendipitously, they spark and catalyze new ideas. With their infrastructure for finance, organization and trade, they allow those ideas to be swiftly actualized.
As for what staunches creativity, that's easy, if ironic. It's the very institutions that we build to manage, exploit and perpetuate the fruits of creativity — our big bureaucracies, and sad to say, too many of our schools. Creativity is disruptive; schools and organizations are regimented, standardized and stultifying.
The education expert Sir Ken Robinson points to a 1968 study reporting on a group of 1,600 children who were tested over time for their ability to think in out-of-the-box ways. When the children were between 3 and 5 years old, 98 percent achieved positive scores. When they were 8 to 10, only 32 percent passed the same test, and only 10 percent at 13 to 15. When 280,000 25-year-olds took the test, just 2 percent passed. By the time we are adults, our creativity has been wrung out of us.
I once asked the great urbanist Jane Jacobs what makes some places more creative than others. She said, essentially, that the question was an easy one. All cities, she said, were filled with creative people; that's our default state as people. But some cities had more than their shares of leaders, people and institutions that blocked out that creativity. She called them "squelchers."
Creativity (or the lack of it) follows the same general contours of the great socio-economic divide — our rising inequality — that plagues us. According to my own estimates, roughly a third of us across the United States, and perhaps as much as half of us in our most creative cities — are able to do work which engages our creative faculties to some extent, whether as artists, musicians, writers, techies, innovators, entrepreneurs, doctors, lawyers, journalists or educators — those of us who work with our minds. That leaves a group that I term "the other 66 percent," who toil in low-wage rote and rotten jobs — if they have jobs at all — in which their creativity is subjugated, ignored or wasted.
Creativity itself is not in danger. It's flourishing is all around us — in science and technology, arts and culture, in our rapidly revitalizing cities. But we still have a long way to go if we want to build a truly creative society that supports and rewards the creativity of each and every one of us.

\vspace{1em}
\textbf{Question 17}

In paragraphs 5 and 6, the author provides the example of Uber to argue that

\begin{enumerate}[label=\Alph*.]
    \item in the future, electric cars will be equipped with mechanisms that prevent collisions.
    \item in the future, traffic jams will not exist.
    \item in the future, the private car will be transformed into a form of public transport.
    \item in the future, Uber rides will outstrip Tesla sales.
[](https://cracku.in/17-in-paragraphs-5-and-6-the-author-provides-the-exam-x-cat-2017-shift-2)
\end{enumerate}
\vspace{1em}
\textbf{Solution:}

In paragraph 5 and 6 the author states that instead of cars having owners they'll work as taxis do, call the taxis at will and use only for journey which we actually need. According to the author this is the future towards which Uber is working. 
Option A, which states that electric cars will have mechanisms to prevent collisions, is out of context. Hence, it is incorrect and can be eliminated.
Option B, which states that future will definitely have no trafiic jams, is incorrect. Hence, option B can be eliminated. 
Option C captures the points which we discussed. Hence, option C is the right answer. 
The passage doesn't give any comparison about Uber rides and Tesla sales. Hence, option D is incorrect and can be eliminated.

\end{problem}

\newpage
\begin{problem}{}{}
% Subject: varc
\textbf{Instructions}

Instructions   

**The passage below is accompanied by a set of six questions. Choose the best answer to each question.**
Cities are the true fonts of creativity... With their diverse populations, dense social networks, and public spaces where people can meet spontaneously and serendipitously, they spark and catalyze new ideas. With their infrastructure for finance, organization and trade, they allow those ideas to be swiftly actualized.
As for what staunches creativity, that's easy, if ironic. It's the very institutions that we build to manage, exploit and perpetuate the fruits of creativity — our big bureaucracies, and sad to say, too many of our schools. Creativity is disruptive; schools and organizations are regimented, standardized and stultifying.
The education expert Sir Ken Robinson points to a 1968 study reporting on a group of 1,600 children who were tested over time for their ability to think in out-of-the-box ways. When the children were between 3 and 5 years old, 98 percent achieved positive scores. When they were 8 to 10, only 32 percent passed the same test, and only 10 percent at 13 to 15. When 280,000 25-year-olds took the test, just 2 percent passed. By the time we are adults, our creativity has been wrung out of us.
I once asked the great urbanist Jane Jacobs what makes some places more creative than others. She said, essentially, that the question was an easy one. All cities, she said, were filled with creative people; that's our default state as people. But some cities had more than their shares of leaders, people and institutions that blocked out that creativity. She called them "squelchers."
Creativity (or the lack of it) follows the same general contours of the great socio-economic divide — our rising inequality — that plagues us. According to my own estimates, roughly a third of us across the United States, and perhaps as much as half of us in our most creative cities — are able to do work which engages our creative faculties to some extent, whether as artists, musicians, writers, techies, innovators, entrepreneurs, doctors, lawyers, journalists or educators — those of us who work with our minds. That leaves a group that I term "the other 66 percent," who toil in low-wage rote and rotten jobs — if they have jobs at all — in which their creativity is subjugated, ignored or wasted.
Creativity itself is not in danger. It's flourishing is all around us — in science and technology, arts and culture, in our rapidly revitalizing cities. But we still have a long way to go if we want to build a truly creative society that supports and rewards the creativity of each and every one of us.

\vspace{1em}
\textbf{Question 18}

In paragraph 6, the author mentions electrically powered bicycles to argue that

\begin{enumerate}[label=\Alph*.]
    \item if Elon Musk were a true visionary, he would invest funds in developing electric bicycles.
    \item our fantasies of autonomy might unexpectedly require us to consider electric bicycles.
    \item in terms of environmental friendliness and safety, electric bicycles rather than electric cars are the future.
    \item electric buses are the best form of public transport.
[](https://cracku.in/18-in-paragraph-6-the-author-mentions-electrically-po-x-cat-2017-shift-2)
\end{enumerate}
\vspace{1em}
\textbf{Solution:}

According to the author we'll have no traffic jams only when cars become public utility. And, for us(people) to have independence, we'll have to start using electrically powered bicycles. Hence, option B, which states this, is the right answer. 
Option A is nowhere mentioned in the passage. Hence, can be eliminated.
Author doesn't give any comparison, as mentioned in option C, between electric powered bicycle and electric cars. Hence, option C can be eliminated. 
The author nowhere mentions that electric buses are the best form of public transport. Hence, option D can be eliminated. 
Hence, option B is the right choice.

Instructions   

**The passage below is accompanied by a set of three questions. Choose the best answer to each question.**
Typewriters are the epitome of a technology that has been comprehensively rendered obsolete by the digital age. The ink comes off the ribbon, they weigh a ton, and second thoughts are a disaster. But they are also personal, portable and, above all, private. Type a document and lock it away and more or less the only way anyone else can get it is if you give it to them. That is why the Russians have decided to go back to typewriters in some government offices, and why in the US, some departments have never abandoned them. Yet it is not just their resistance to algorithms and secret surveillance that keeps typewriter production lines — well one, at least — in business (the last British one closed a year ago). Nor is it only the nostalgic appeal of the metal body and the stout well-defined keys that make them popular on eBay. A typewriter demands something particular: attentiveness. By the time the paper is loaded, the ribbon tightened, the carriage returned, the spacing and the margins set, there's a big premium on hitting the right key. That means sorting out ideas, pulling together a kind of order and organising details before actually striking off. There can be no thinking on screen with a typewriter. Nor are there any easy distractions. No online shopping. No urgent emails. No Twitter. No need even for electricity — perfect for writing in a remote hideaway. The thinking process is accompanied by the encouraging clack of keys, and the ratchet of the carriage return. Ping!

\end{problem}

\newpage
\begin{problem}{}{}
% Subject: varc
\textbf{Instructions}

Instructions   

**The passage below is accompanied by a set of six questions. Choose the best answer to each question.**
Cities are the true fonts of creativity... With their diverse populations, dense social networks, and public spaces where people can meet spontaneously and serendipitously, they spark and catalyze new ideas. With their infrastructure for finance, organization and trade, they allow those ideas to be swiftly actualized.
As for what staunches creativity, that's easy, if ironic. It's the very institutions that we build to manage, exploit and perpetuate the fruits of creativity — our big bureaucracies, and sad to say, too many of our schools. Creativity is disruptive; schools and organizations are regimented, standardized and stultifying.
The education expert Sir Ken Robinson points to a 1968 study reporting on a group of 1,600 children who were tested over time for their ability to think in out-of-the-box ways. When the children were between 3 and 5 years old, 98 percent achieved positive scores. When they were 8 to 10, only 32 percent passed the same test, and only 10 percent at 13 to 15. When 280,000 25-year-olds took the test, just 2 percent passed. By the time we are adults, our creativity has been wrung out of us.
I once asked the great urbanist Jane Jacobs what makes some places more creative than others. She said, essentially, that the question was an easy one. All cities, she said, were filled with creative people; that's our default state as people. But some cities had more than their shares of leaders, people and institutions that blocked out that creativity. She called them "squelchers."
Creativity (or the lack of it) follows the same general contours of the great socio-economic divide — our rising inequality — that plagues us. According to my own estimates, roughly a third of us across the United States, and perhaps as much as half of us in our most creative cities — are able to do work which engages our creative faculties to some extent, whether as artists, musicians, writers, techies, innovators, entrepreneurs, doctors, lawyers, journalists or educators — those of us who work with our minds. That leaves a group that I term "the other 66 percent," who toil in low-wage rote and rotten jobs — if they have jobs at all — in which their creativity is subjugated, ignored or wasted.
Creativity itself is not in danger. It's flourishing is all around us — in science and technology, arts and culture, in our rapidly revitalizing cities. But we still have a long way to go if we want to build a truly creative society that supports and rewards the creativity of each and every one of us.

\vspace{1em}
\textbf{Question 19}

Which one of the following best describes what the passage is trying to do?

\begin{enumerate}[label=\Alph*.]
    \item It describes why people continue to use typewriters even in the digital age.
    \item It argues that typewriters will continue to be used even though they are an obsolete technology.
    \item It highlights the personal benefits of using typewriters.
    \item It shows that computers offer fewer options than typewriters.
[](https://cracku.in/19-which-one-of-the-following-best-describes-what-the-x-cat-2017-shift-2)
\end{enumerate}
\vspace{1em}
\textbf{Solution:}

The passage starts by introducing typewriters. The author later states that some goverment offices in Russia are going back to typewriters and that in the US some offices still use them. The author then goes on to give reasons for the same by highlighting positive aspects of the typewriter. 
Option A correctly describes what the passage is trying to do. Hence, it is the right choice.
The author doesn't state that use of typewriters will be perennial. Hence, it can be eliminated.
The main aim of the passage is not to highlight personal benefit. Hence, it can be eliminated.
Option D is out of scope as the author nowhere states or highlights that computers offer less options than typewriters. Hence, it can be eliminated.
Hence, option A is the right answer.

\end{problem}

\newpage
\begin{problem}{}{}
% Subject: varc
\textbf{Instructions}

Instructions   

**The passage below is accompanied by a set of six questions. Choose the best answer to each question.**
Cities are the true fonts of creativity... With their diverse populations, dense social networks, and public spaces where people can meet spontaneously and serendipitously, they spark and catalyze new ideas. With their infrastructure for finance, organization and trade, they allow those ideas to be swiftly actualized.
As for what staunches creativity, that's easy, if ironic. It's the very institutions that we build to manage, exploit and perpetuate the fruits of creativity — our big bureaucracies, and sad to say, too many of our schools. Creativity is disruptive; schools and organizations are regimented, standardized and stultifying.
The education expert Sir Ken Robinson points to a 1968 study reporting on a group of 1,600 children who were tested over time for their ability to think in out-of-the-box ways. When the children were between 3 and 5 years old, 98 percent achieved positive scores. When they were 8 to 10, only 32 percent passed the same test, and only 10 percent at 13 to 15. When 280,000 25-year-olds took the test, just 2 percent passed. By the time we are adults, our creativity has been wrung out of us.
I once asked the great urbanist Jane Jacobs what makes some places more creative than others. She said, essentially, that the question was an easy one. All cities, she said, were filled with creative people; that's our default state as people. But some cities had more than their shares of leaders, people and institutions that blocked out that creativity. She called them "squelchers."
Creativity (or the lack of it) follows the same general contours of the great socio-economic divide — our rising inequality — that plagues us. According to my own estimates, roughly a third of us across the United States, and perhaps as much as half of us in our most creative cities — are able to do work which engages our creative faculties to some extent, whether as artists, musicians, writers, techies, innovators, entrepreneurs, doctors, lawyers, journalists or educators — those of us who work with our minds. That leaves a group that I term "the other 66 percent," who toil in low-wage rote and rotten jobs — if they have jobs at all — in which their creativity is subjugated, ignored or wasted.
Creativity itself is not in danger. It's flourishing is all around us — in science and technology, arts and culture, in our rapidly revitalizing cities. But we still have a long way to go if we want to build a truly creative society that supports and rewards the creativity of each and every one of us.

\vspace{1em}
\textbf{Question 20}

According to the passage, some governments still use typewriters because:

\begin{enumerate}[label=\Alph*.]
    \item they do not want to abandon old technologies that may be useful in the future.
    \item they want to ensure that typewriter production lines remain in business.
    \item they like the nostalgic appeal of typewriter.
    \item they can control who reads the document.
[](https://cracku.in/20-according-to-the-passage-some-governments-still-us-x-cat-2017-shift-2)
\end{enumerate}
\vspace{1em}
\textbf{Solution:}

Refer to the sentence, " ... the only way anyone else can get it is if you give it to them. That is why the Russians have decided to go back to typewriters in some government offices..." Hence, the government uses typewriters to control who views the document as the only way someone can read the document is by physically accessing it to them. Option D, which highlights this, is the right answer.

\end{problem}

\newpage
\begin{problem}{}{}
% Subject: varc
\textbf{Instructions}

Instructions   

**The passage below is accompanied by a set of six questions. Choose the best answer to each question.**
Cities are the true fonts of creativity... With their diverse populations, dense social networks, and public spaces where people can meet spontaneously and serendipitously, they spark and catalyze new ideas. With their infrastructure for finance, organization and trade, they allow those ideas to be swiftly actualized.
As for what staunches creativity, that's easy, if ironic. It's the very institutions that we build to manage, exploit and perpetuate the fruits of creativity — our big bureaucracies, and sad to say, too many of our schools. Creativity is disruptive; schools and organizations are regimented, standardized and stultifying.
The education expert Sir Ken Robinson points to a 1968 study reporting on a group of 1,600 children who were tested over time for their ability to think in out-of-the-box ways. When the children were between 3 and 5 years old, 98 percent achieved positive scores. When they were 8 to 10, only 32 percent passed the same test, and only 10 percent at 13 to 15. When 280,000 25-year-olds took the test, just 2 percent passed. By the time we are adults, our creativity has been wrung out of us.
I once asked the great urbanist Jane Jacobs what makes some places more creative than others. She said, essentially, that the question was an easy one. All cities, she said, were filled with creative people; that's our default state as people. But some cities had more than their shares of leaders, people and institutions that blocked out that creativity. She called them "squelchers."
Creativity (or the lack of it) follows the same general contours of the great socio-economic divide — our rising inequality — that plagues us. According to my own estimates, roughly a third of us across the United States, and perhaps as much as half of us in our most creative cities — are able to do work which engages our creative faculties to some extent, whether as artists, musicians, writers, techies, innovators, entrepreneurs, doctors, lawyers, journalists or educators — those of us who work with our minds. That leaves a group that I term "the other 66 percent," who toil in low-wage rote and rotten jobs — if they have jobs at all — in which their creativity is subjugated, ignored or wasted.
Creativity itself is not in danger. It's flourishing is all around us — in science and technology, arts and culture, in our rapidly revitalizing cities. But we still have a long way to go if we want to build a truly creative society that supports and rewards the creativity of each and every one of us.

\vspace{1em}
\textbf{Question 21}

The writer praises typewriters for all the following reasons EXCEPT

\begin{enumerate}[label=\Alph*.]
    \item Unlike computers, they can only be used for typing.
    \item You cannot revise what you have typed on a typewriter.
    \item Typewriters are noisier than computers.
    \item Typewriters are messier to use than computers.
[](https://cracku.in/21-the-writer-praises-typewriters-for-all-the-followi-x-cat-2017-shift-2)
\end{enumerate}
\vspace{1em}
\textbf{Solution:}

The author states that, "..Nor are there any easy distractions...". Hence, the only thing one can do using typewriter is write, unlike computers. Hence, option A is correct.
Refer to following lines "... there's a big premium on hitting the right key...". Thus, there's premium attached on hitting right keys because as you cannot revise what you've typed on the typewriter. Hence, option B is correct.
Refer to following lines "...thinking process is accompanied by the encouraging clack of keys...". Hence, typewriters are noisier than computers. So, option C is correct. 
Nowhere in the passage does the author state or highlight that typewriters are messiar than the computers. Hence, option D is not the reason why the author praises typewriters. 
Thus, option D is the right choice.

Instructions   

**The passage below is accompanied by a set of three questions. Choose the best answer to each question.**
Despite their fierce reputation, Vikings may not have always been the plunderers and pillagers popular culture imagines them to be. In fact, they got their start trading in northern European markets, researchers suggest.
Combs carved from animal antlers, as well as comb manufacturing waste and raw antler material has turned up at three archaeological sites in Denmark, including a medieval marketplace in the city of Ribe. A team of researchers from Denmark and the U.K. hoped to identify the species of animal to which the antlers once belonged by analyzing collagen proteins in the samples and comparing them across the animal kingdom, Laura Geggel reports for LiveScience. Somewhat surprisingly, molecular analysis of the artifacts revealed that some combs and other material had been carved from reindeer antlers.... Given that reindeer (Rangifer tarandus) don't live in Denmark, the researchers posit that it arrived on Viking ships from Norway. Antler craftsmanship, in the form of decorative combs, was part of Viking culture. Such combs served as symbols of good health, Geggel writes. The fact that the animals shed their antlers also made them easy to collect from the large herds that inhabited Norway.
Since the artifacts were found in marketplace areas at each site it's more likely that the Norsemen came to trade rather than pillage. Most of the artifacts also date to the 780s, but some are as old as 725. That predates the beginning of Viking raids on Great Britain by about 70 years. (Traditionally, the so-called "Viking Age" began with these raids in 793 and ended with the Norman conquest of Great Britain in l066.) Archaeologists had suspected that the Vikings had experience with long maritime voyages [that] might have preceded their raiding days. Beyond Norway, these combs would have been a popular industry in Scandinavia as well: It's possible that the antler combs represent a larger trade network, where the Norsemen supplied raw material to craftsmen in Denmark and elsewhere.

\end{problem}

\newpage
\begin{problem}{}{}
% Subject: varc
\textbf{Instructions}

Instructions   

**The passage below is accompanied by a set of six questions. Choose the best answer to each question.**
Cities are the true fonts of creativity... With their diverse populations, dense social networks, and public spaces where people can meet spontaneously and serendipitously, they spark and catalyze new ideas. With their infrastructure for finance, organization and trade, they allow those ideas to be swiftly actualized.
As for what staunches creativity, that's easy, if ironic. It's the very institutions that we build to manage, exploit and perpetuate the fruits of creativity — our big bureaucracies, and sad to say, too many of our schools. Creativity is disruptive; schools and organizations are regimented, standardized and stultifying.
The education expert Sir Ken Robinson points to a 1968 study reporting on a group of 1,600 children who were tested over time for their ability to think in out-of-the-box ways. When the children were between 3 and 5 years old, 98 percent achieved positive scores. When they were 8 to 10, only 32 percent passed the same test, and only 10 percent at 13 to 15. When 280,000 25-year-olds took the test, just 2 percent passed. By the time we are adults, our creativity has been wrung out of us.
I once asked the great urbanist Jane Jacobs what makes some places more creative than others. She said, essentially, that the question was an easy one. All cities, she said, were filled with creative people; that's our default state as people. But some cities had more than their shares of leaders, people and institutions that blocked out that creativity. She called them "squelchers."
Creativity (or the lack of it) follows the same general contours of the great socio-economic divide — our rising inequality — that plagues us. According to my own estimates, roughly a third of us across the United States, and perhaps as much as half of us in our most creative cities — are able to do work which engages our creative faculties to some extent, whether as artists, musicians, writers, techies, innovators, entrepreneurs, doctors, lawyers, journalists or educators — those of us who work with our minds. That leaves a group that I term "the other 66 percent," who toil in low-wage rote and rotten jobs — if they have jobs at all — in which their creativity is subjugated, ignored or wasted.
Creativity itself is not in danger. It's flourishing is all around us — in science and technology, arts and culture, in our rapidly revitalizing cities. But we still have a long way to go if we want to build a truly creative society that supports and rewards the creativity of each and every one of us.

\vspace{1em}
\textbf{Question 22}

The primary purpose of the passage is:

\begin{enumerate}[label=\Alph*.]
    \item to explain the presence of reindeer antler combs in Denmark.
    \item to contradict the widely-accepted beginning date for the Viking Age in Britain, and propose an alternate one.
    \item to challenge the popular perception of Vikings as raiders by using evidence that suggests their early trade relations with Europe.
    \item to argue that besides being violent pillagers, Vikings were also skilled craftsmen and efficient traders.
[](https://cracku.in/22-the-primary-purpose-of-the-passage-is-x-cat-2017-shift-2)
\end{enumerate}
\vspace{1em}
\textbf{Solution:}

The passage revolves around how vikings did not start out as pillagers but as traders. The intention of the author seems to dispel the notion that the Vikings were pillagers. The combs have been used just as an illustration to prove the author's hypothesis. Therefore, option A can be ruled out.  
  
Option B states that the purpose was to change the period of Viking age. However, the passage does not hint any such intention. The author cites that the combs had made their way to Britain before the Viking age to substantiate the fact that Vikings were traders before they became pillagers.  
Therefore, we can rule out option B too.  
  
Option D states that despite being pillagers, Vikings were efficient traders and craftsmen. However, the passage talks about a period prior to which Vikings turned pillagers. Therefore, we can eliminate option D too.  
  
Option C states that the purpose of the passage is to dispel the notion that Vikings were pillagers. This seems the most appropriate option as the passage tries to establish the fact that Vikings started out as traders. Therefore, option C is the right answer.

\end{problem}

\newpage
\begin{problem}{}{}
% Subject: varc
\textbf{Instructions}

Instructions   

**The passage below is accompanied by a set of six questions. Choose the best answer to each question.**
Cities are the true fonts of creativity... With their diverse populations, dense social networks, and public spaces where people can meet spontaneously and serendipitously, they spark and catalyze new ideas. With their infrastructure for finance, organization and trade, they allow those ideas to be swiftly actualized.
As for what staunches creativity, that's easy, if ironic. It's the very institutions that we build to manage, exploit and perpetuate the fruits of creativity — our big bureaucracies, and sad to say, too many of our schools. Creativity is disruptive; schools and organizations are regimented, standardized and stultifying.
The education expert Sir Ken Robinson points to a 1968 study reporting on a group of 1,600 children who were tested over time for their ability to think in out-of-the-box ways. When the children were between 3 and 5 years old, 98 percent achieved positive scores. When they were 8 to 10, only 32 percent passed the same test, and only 10 percent at 13 to 15. When 280,000 25-year-olds took the test, just 2 percent passed. By the time we are adults, our creativity has been wrung out of us.
I once asked the great urbanist Jane Jacobs what makes some places more creative than others. She said, essentially, that the question was an easy one. All cities, she said, were filled with creative people; that's our default state as people. But some cities had more than their shares of leaders, people and institutions that blocked out that creativity. She called them "squelchers."
Creativity (or the lack of it) follows the same general contours of the great socio-economic divide — our rising inequality — that plagues us. According to my own estimates, roughly a third of us across the United States, and perhaps as much as half of us in our most creative cities — are able to do work which engages our creative faculties to some extent, whether as artists, musicians, writers, techies, innovators, entrepreneurs, doctors, lawyers, journalists or educators — those of us who work with our minds. That leaves a group that I term "the other 66 percent," who toil in low-wage rote and rotten jobs — if they have jobs at all — in which their creativity is subjugated, ignored or wasted.
Creativity itself is not in danger. It's flourishing is all around us — in science and technology, arts and culture, in our rapidly revitalizing cities. But we still have a long way to go if we want to build a truly creative society that supports and rewards the creativity of each and every one of us.

\vspace{1em}
\textbf{Question 23}

The evidence - "Most of the artifacts also date to the 780s, but some are as old as 725" — has been used in the passage to argue that:

\begin{enumerate}[label=\Alph*.]
    \item the beginning date of the Viking Age should be changed from 793 to 725.
    \item the Viking raids started as early as 725.
    \item some of the antler artifacts found in Denmark and Great Britain could have come from Scandinavia.
    \item the Vikings' trade relations with Europe pre-dates the Viking raids.
[](https://cracku.in/23-the-evidence-most-of-the-artifacts-also-date-to-th-x-cat-2017-shift-2)
\end{enumerate}
\vspace{1em}
\textbf{Solution:}

The author mentions the statement to imply that the Vikings had trade relations with the British before the Viking age. The Viking age started in 793, whereas the artifacts predate this period. Therefore, the intention of the line "Most of the artifacts also date to the 780s, but some are as old as 725" is to emphasize that Vikings had trade relations. Therefore, option D is the right answer.

\end{problem}

\newpage
\begin{problem}{}{}
% Subject: varc
\textbf{Instructions}

Instructions   

**The passage below is accompanied by a set of six questions. Choose the best answer to each question.**
Cities are the true fonts of creativity... With their diverse populations, dense social networks, and public spaces where people can meet spontaneously and serendipitously, they spark and catalyze new ideas. With their infrastructure for finance, organization and trade, they allow those ideas to be swiftly actualized.
As for what staunches creativity, that's easy, if ironic. It's the very institutions that we build to manage, exploit and perpetuate the fruits of creativity — our big bureaucracies, and sad to say, too many of our schools. Creativity is disruptive; schools and organizations are regimented, standardized and stultifying.
The education expert Sir Ken Robinson points to a 1968 study reporting on a group of 1,600 children who were tested over time for their ability to think in out-of-the-box ways. When the children were between 3 and 5 years old, 98 percent achieved positive scores. When they were 8 to 10, only 32 percent passed the same test, and only 10 percent at 13 to 15. When 280,000 25-year-olds took the test, just 2 percent passed. By the time we are adults, our creativity has been wrung out of us.
I once asked the great urbanist Jane Jacobs what makes some places more creative than others. She said, essentially, that the question was an easy one. All cities, she said, were filled with creative people; that's our default state as people. But some cities had more than their shares of leaders, people and institutions that blocked out that creativity. She called them "squelchers."
Creativity (or the lack of it) follows the same general contours of the great socio-economic divide — our rising inequality — that plagues us. According to my own estimates, roughly a third of us across the United States, and perhaps as much as half of us in our most creative cities — are able to do work which engages our creative faculties to some extent, whether as artists, musicians, writers, techies, innovators, entrepreneurs, doctors, lawyers, journalists or educators — those of us who work with our minds. That leaves a group that I term "the other 66 percent," who toil in low-wage rote and rotten jobs — if they have jobs at all — in which their creativity is subjugated, ignored or wasted.
Creativity itself is not in danger. It's flourishing is all around us — in science and technology, arts and culture, in our rapidly revitalizing cities. But we still have a long way to go if we want to build a truly creative society that supports and rewards the creativity of each and every one of us.

\vspace{1em}
\textbf{Question 24}

All of the following hold true for Vikings EXCEPT

\begin{enumerate}[label=\Alph*.]
    \item Vikings brought reindeer from Norway to Denmark for trade purposes.
    \item Before becoming the raiders of northern Europe, Vikings had trade relations with European nations.
    \item Antler combs, regarded by the Vikings as a symbol of good health, were part of the Viking culture.
    \item Vikings, once upon a time, had trade relations with Denmark and Scandinavia.
[](https://cracku.in/24-all-of-the-following-hold-true-for-vikings-except-x-cat-2017-shift-2)
\end{enumerate}
\vspace{1em}
\textbf{Solution:}

In the passage, it has been mentioned that "Such combs served as symbols of good health, Geggel writes" . Therefore, we can infer option C and hence, it can be eliminated.  
Option D states that "Vikings, once upon a time, had trade relations with Denmark and Scandinavia". The last paragraph mentions that "Beyond Norway, these combs would have been a popular industry in Scandinavia as well. It’s possible that the antler combs represent a larger trade network, where the Norsemen supplied raw material to craftsmen in Denmark and elsewhere". Therefore, we can rule out option D too.  
Option B states that the Vikings had trade relations with Northern Europe. This is the very theme of the passage. Hence, we can eliminate option B too.  
  
Option A states that "Vikings brought reindeer from Norway to Denmark for trade purposes". However, the passage only mentions that they brought the antlers - not the Reindeers itself. Therefore, option A is the right answer.

Instructions   

For the following questions answer them individually

\end{problem}

\newpage
\begin{problem}{}{}
% Subject: varc
\textbf{Instructions}

Instructions   

**The passage below is accompanied by a set of six questions. Choose the best answer to each question.**
Cities are the true fonts of creativity... With their diverse populations, dense social networks, and public spaces where people can meet spontaneously and serendipitously, they spark and catalyze new ideas. With their infrastructure for finance, organization and trade, they allow those ideas to be swiftly actualized.
As for what staunches creativity, that's easy, if ironic. It's the very institutions that we build to manage, exploit and perpetuate the fruits of creativity — our big bureaucracies, and sad to say, too many of our schools. Creativity is disruptive; schools and organizations are regimented, standardized and stultifying.
The education expert Sir Ken Robinson points to a 1968 study reporting on a group of 1,600 children who were tested over time for their ability to think in out-of-the-box ways. When the children were between 3 and 5 years old, 98 percent achieved positive scores. When they were 8 to 10, only 32 percent passed the same test, and only 10 percent at 13 to 15. When 280,000 25-year-olds took the test, just 2 percent passed. By the time we are adults, our creativity has been wrung out of us.
I once asked the great urbanist Jane Jacobs what makes some places more creative than others. She said, essentially, that the question was an easy one. All cities, she said, were filled with creative people; that's our default state as people. But some cities had more than their shares of leaders, people and institutions that blocked out that creativity. She called them "squelchers."
Creativity (or the lack of it) follows the same general contours of the great socio-economic divide — our rising inequality — that plagues us. According to my own estimates, roughly a third of us across the United States, and perhaps as much as half of us in our most creative cities — are able to do work which engages our creative faculties to some extent, whether as artists, musicians, writers, techies, innovators, entrepreneurs, doctors, lawyers, journalists or educators — those of us who work with our minds. That leaves a group that I term "the other 66 percent," who toil in low-wage rote and rotten jobs — if they have jobs at all — in which their creativity is subjugated, ignored or wasted.
Creativity itself is not in danger. It's flourishing is all around us — in science and technology, arts and culture, in our rapidly revitalizing cities. But we still have a long way to go if we want to build a truly creative society that supports and rewards the creativity of each and every one of us.

\vspace{1em}
\textbf{Question 25}

The passage given below is followed by four summaries. Choose the option that best captures the author' s position.  
North American walnut sphinx moth caterpillars (Amorpha juglandis) look like easy meals for birds, but they have a trick up their sleeves — they produce whistles that sound like bird alarm calls, scaring potential predators away. At first, scientists suspected birds were simply startled by the loud noise. But a new study suggests a more sophisticated mechanism: the caterpillar's whistle appears to mimic a bird alarm call, sending avian predators scrambling for cover. When pecked by a bird, the caterpillars whistle by compressing their bodies like an accordion and forcing air out through specialized holes in their sides. The whistles are impressively loud — they have been measured at over 80 dB from 5 cm away from the caterpillar — considering they are made by a two-inch long insect.

\begin{enumerate}[label=\Alph*.]
    \item North American walnut sphinx moth caterpillars will whistle periodically to ward off predator birds - they have a specialized vocal tract that helps them whistle.
    \item North American walnut sphinx moth caterpillars can whistle very loudly; the loudness of their whistles is shocking as they are very small insects.
    \item The North American walnut sphinx moth caterpillars, in a case of acoustic deception, produce whistles that mimic bird alarm calls to defend themselves.
    \item North American. walnut sphinx moth caterpillars, in. a case of deception and camouflage, produce whistles that mimic bird alarm calls to defend themselves.
[](https://cracku.in/25-the-passage-given-below-is-followed-by-four-summar-x-cat-2017-shift-2)
\end{enumerate}
\vspace{1em}
\textbf{Solution:}

According to the paragraph, the North American walnut sphinx moth caterpillars produce whistles which are extremely loud considering their size. These whistles appear to mimic bird(predator) alarm calls which scares them to look for cover. Thus, these sounds act as acoustic deception and help the insect to defend themselves against predators. 
Option A mentions about vocal tracts which is out of scope. Hence, it can be eliminated.
Option B though correct, fails to mention the use of sound to defend against the predators. Hence, it can be eliminated.
Option C captures all the main points and hence is right choice. 
Option D mentions 'camouflage' which is also out of context. Hence, it can be eliminated. 
Hence, option C is the right answer.

\end{problem}

\newpage
\begin{problem}{}{}
% Subject: varc
\textbf{Instructions}

Instructions   

**The passage below is accompanied by a set of six questions. Choose the best answer to each question.**
Cities are the true fonts of creativity... With their diverse populations, dense social networks, and public spaces where people can meet spontaneously and serendipitously, they spark and catalyze new ideas. With their infrastructure for finance, organization and trade, they allow those ideas to be swiftly actualized.
As for what staunches creativity, that's easy, if ironic. It's the very institutions that we build to manage, exploit and perpetuate the fruits of creativity — our big bureaucracies, and sad to say, too many of our schools. Creativity is disruptive; schools and organizations are regimented, standardized and stultifying.
The education expert Sir Ken Robinson points to a 1968 study reporting on a group of 1,600 children who were tested over time for their ability to think in out-of-the-box ways. When the children were between 3 and 5 years old, 98 percent achieved positive scores. When they were 8 to 10, only 32 percent passed the same test, and only 10 percent at 13 to 15. When 280,000 25-year-olds took the test, just 2 percent passed. By the time we are adults, our creativity has been wrung out of us.
I once asked the great urbanist Jane Jacobs what makes some places more creative than others. She said, essentially, that the question was an easy one. All cities, she said, were filled with creative people; that's our default state as people. But some cities had more than their shares of leaders, people and institutions that blocked out that creativity. She called them "squelchers."
Creativity (or the lack of it) follows the same general contours of the great socio-economic divide — our rising inequality — that plagues us. According to my own estimates, roughly a third of us across the United States, and perhaps as much as half of us in our most creative cities — are able to do work which engages our creative faculties to some extent, whether as artists, musicians, writers, techies, innovators, entrepreneurs, doctors, lawyers, journalists or educators — those of us who work with our minds. That leaves a group that I term "the other 66 percent," who toil in low-wage rote and rotten jobs — if they have jobs at all — in which their creativity is subjugated, ignored or wasted.
Creativity itself is not in danger. It's flourishing is all around us — in science and technology, arts and culture, in our rapidly revitalizing cities. But we still have a long way to go if we want to build a truly creative society that supports and rewards the creativity of each and every one of us.

\vspace{1em}
\textbf{Question 26}

The passage given below is followed by four summaries. Choose the option that best captures the author's position.  
  
Both Socrates and Bacon were very good at asking useful questions. In fact, Socrates is largely credited with coming up with a way of asking questions, 'the Socratic method,' which itself is at the core of the 'scientific method,' popularised by Bacon. The Socratic method disproves arguments by finding exceptions to them, and can therefore lead your opponent to a point where they admit something that contradicts their original position. In common with Socrates, Bacon stressed it was as important to disprove a theory as it was to prove one — and real-world observation and experimentation were key to achieving both aims. Bacon also saw science as a collaborative affair, with scientists working together, challenging each other.

\begin{enumerate}[label=\Alph*.]
    \item Both Socrates and Bacon advocated clever questioning of the opponents to disprove their arguments and theories.
    \item Both Socrates and Bacon advocated challenging arguments and theories by observation and experimentation.
    \item Both Socrates and Bacon advocated confirming arguments and theories by finding exceptions.
    \item Both Socrates and Bacon advocated examining arguments and theories from both sides to prove them.
[](https://cracku.in/26-the-passage-given-below-is-followed-by-four-summar-x-cat-2017-shift-2)
\end{enumerate}
\vspace{1em}
\textbf{Solution:}

According to the paragraph, Socrates and Bacon were good at asking questions. The Socratic method works in a way by finding exceptions to the arguments of the opponent, which makes the opponent to agree on something that contradicts their original position. In a similar way, Bacon stressed that it was important to disprove theory as it is to prove it. Thus both Socrates and Bacon stressed on examining arguments from both ends - to prove as well as disprove. 
Option A, which speaks only about disproving of arguments, can be eliminated.
Option B talks only about examining and observation. Hence, it can be eliminated.
Option C talks only about confirming of arguments and not the other way. Hence, it can be eliminated.
Option D captures the main points which we discussed earlier. 
Hence, option D is the right answer.

\end{problem}

\newpage
\begin{problem}{}{}
% Subject: varc
\textbf{Instructions}

Instructions   

**The passage below is accompanied by a set of six questions. Choose the best answer to each question.**
Cities are the true fonts of creativity... With their diverse populations, dense social networks, and public spaces where people can meet spontaneously and serendipitously, they spark and catalyze new ideas. With their infrastructure for finance, organization and trade, they allow those ideas to be swiftly actualized.
As for what staunches creativity, that's easy, if ironic. It's the very institutions that we build to manage, exploit and perpetuate the fruits of creativity — our big bureaucracies, and sad to say, too many of our schools. Creativity is disruptive; schools and organizations are regimented, standardized and stultifying.
The education expert Sir Ken Robinson points to a 1968 study reporting on a group of 1,600 children who were tested over time for their ability to think in out-of-the-box ways. When the children were between 3 and 5 years old, 98 percent achieved positive scores. When they were 8 to 10, only 32 percent passed the same test, and only 10 percent at 13 to 15. When 280,000 25-year-olds took the test, just 2 percent passed. By the time we are adults, our creativity has been wrung out of us.
I once asked the great urbanist Jane Jacobs what makes some places more creative than others. She said, essentially, that the question was an easy one. All cities, she said, were filled with creative people; that's our default state as people. But some cities had more than their shares of leaders, people and institutions that blocked out that creativity. She called them "squelchers."
Creativity (or the lack of it) follows the same general contours of the great socio-economic divide — our rising inequality — that plagues us. According to my own estimates, roughly a third of us across the United States, and perhaps as much as half of us in our most creative cities — are able to do work which engages our creative faculties to some extent, whether as artists, musicians, writers, techies, innovators, entrepreneurs, doctors, lawyers, journalists or educators — those of us who work with our minds. That leaves a group that I term "the other 66 percent," who toil in low-wage rote and rotten jobs — if they have jobs at all — in which their creativity is subjugated, ignored or wasted.
Creativity itself is not in danger. It's flourishing is all around us — in science and technology, arts and culture, in our rapidly revitalizing cities. But we still have a long way to go if we want to build a truly creative society that supports and rewards the creativity of each and every one of us.

\vspace{1em}
\textbf{Question 27}

**The passage given below is followed by four summaries. Choose the option that best captures the author' s position.**  
A fundamental property of language is that it is slippery and messy and more liquid than solid, a gelatinous mass that changes shape to fit. As Wittgenstein would remind us, "usage has no sharp boundary." Oftentimes, the only way to determine the meaning of a word is to examine how it is used. This insight is often described as the "meaning is use" doctrine. There are differences between the "meaning is use" doctrine and a dictionary-first theory of meaning. "The dictionary's careful fixing of words to definitions, like butterflies pinned under glass, can suggest that this is how language works. The definitions can seem to ensure and fix the meaning of words, just as the gold standard can back a country's currency." What Wittgenstein found in the circulation of ordinary language, however, was a free-floating currency of meaning. The value of each word arises out of the exchange. The lexicographer abstracts a meaning from that exchange, which is then set within the conventions of the dictionary definition.

\begin{enumerate}[label=\Alph*.]
    \item Dictionary definitions are like 'gold standards' — artificial, theoretical and dogmatic. Actual meaning of words is their free-exchange value.
    \item Language is already slippery; given this, accounting for 'meaning in use' will only exasperate the problem. That is why lexicographers 'fix' meanings.
    \item Meaning is dynamic; definitions are static. The 'meaning in use' theory helps us understand that definitions of words are culled from their meaning in exchange and use and not vice versa.
    \item The meaning of words in dictionaries is clear, fixed and less dangerous and ambiguous than the meaning that arises when words are exchanged between people.
[](https://cracku.in/27-the-passage-given-below-is-followed-by-four-summar-x-cat-2017-shift-2)
\end{enumerate}
\vspace{1em}
\textbf{Solution:}

According to the paragraph, language is like a gelatinous mass that changes shape to fit. Also, many times the only way to find meaning of word is to examine how it is used. It is stated that definitions are fixed for the word by dictionary.Wittgenstein found that circulation of ordinary language was a free-floating currency of meaning. So the meanings are dynamic. Thus, the value of word arises from the exchange and then the lexicographer abstracts meaning from that exchange. Thus, definitions are picked up from the meaning in use. 
Option A, which states that definitions are like dogmatic, cannot be found in the paragraph. Hence, it can be eliminated.
The paragraph doesn't talk about why lexicographers fix meanings. Hence, option B can be eliminated.
Option C covers all the main points. Hence, it is the right choice.
The purpose of the passage is not to compare meaning of words in dictionaries with meaning which arises from exchange. Hence, option D can be eliminated.
Hence, option C is the right choice.

\end{problem}

\newpage
\begin{problem}{}{}
% Subject: varc
\textbf{Instructions}

Instructions   

**The passage below is accompanied by a set of six questions. Choose the best answer to each question.**
Cities are the true fonts of creativity... With their diverse populations, dense social networks, and public spaces where people can meet spontaneously and serendipitously, they spark and catalyze new ideas. With their infrastructure for finance, organization and trade, they allow those ideas to be swiftly actualized.
As for what staunches creativity, that's easy, if ironic. It's the very institutions that we build to manage, exploit and perpetuate the fruits of creativity — our big bureaucracies, and sad to say, too many of our schools. Creativity is disruptive; schools and organizations are regimented, standardized and stultifying.
The education expert Sir Ken Robinson points to a 1968 study reporting on a group of 1,600 children who were tested over time for their ability to think in out-of-the-box ways. When the children were between 3 and 5 years old, 98 percent achieved positive scores. When they were 8 to 10, only 32 percent passed the same test, and only 10 percent at 13 to 15. When 280,000 25-year-olds took the test, just 2 percent passed. By the time we are adults, our creativity has been wrung out of us.
I once asked the great urbanist Jane Jacobs what makes some places more creative than others. She said, essentially, that the question was an easy one. All cities, she said, were filled with creative people; that's our default state as people. But some cities had more than their shares of leaders, people and institutions that blocked out that creativity. She called them "squelchers."
Creativity (or the lack of it) follows the same general contours of the great socio-economic divide — our rising inequality — that plagues us. According to my own estimates, roughly a third of us across the United States, and perhaps as much as half of us in our most creative cities — are able to do work which engages our creative faculties to some extent, whether as artists, musicians, writers, techies, innovators, entrepreneurs, doctors, lawyers, journalists or educators — those of us who work with our minds. That leaves a group that I term "the other 66 percent," who toil in low-wage rote and rotten jobs — if they have jobs at all — in which their creativity is subjugated, ignored or wasted.
Creativity itself is not in danger. It's flourishing is all around us — in science and technology, arts and culture, in our rapidly revitalizing cities. But we still have a long way to go if we want to build a truly creative society that supports and rewards the creativity of each and every one of us.

\vspace{1em}
\textbf{Question 28}

The five sentences labelled (1, 2, 3, 4, 5) given in this question, when properly sequenced, form a coherent paragraph. Each sentence is labelled with a number. Decide on the proper order for the sentences and key in this sequence of **FIVE NUMBERS** as your answer.
1: The implications of retelling of Indian stories, hence, takes on new meaning in a modern India.  
2. The stories we tell reflect the world around us.  
3. We cannot help but retell the stories that we value — after all, they are never quite right for us — in our time. 
4. And even if we manage to get them quite right, they are only right for us — other people living around us will have different reasons for telling similar stories.
5. As soon as we capture a story, the world we were trying to capture has changed.
Backspace  
789  
456  
123  
0.-  
Clear All  

Submit
[](https://cracku.in/28-the-five-sentences-labelled-1-2-3-4-5-given-in-thi-x-cat-2017-shift-2)

\begin{enumerate}[label=\Alph*.]
\end{enumerate}
\vspace{1em}
\textbf{Solution:}

Sentence 2, which introduces the topic of what stories tell, must be the starting sentence. Sentence 5 elaborates on sentence 2. Hence sentence 5 logically follows sentence 2. According to the sentence 3, the stories we retell are never quite right for us. Sentence 4 elaborates on what if we get the stories quite right. Hence, 3-4 forms a pair which must come after sentence 5. Sentence 1 which concludes the topics of discussion must be the ending sentence. Hence, 25341 is the right answer.

\end{problem}

\newpage
\begin{problem}{}{}
% Subject: varc
\textbf{Instructions}

Instructions   

**The passage below is accompanied by a set of six questions. Choose the best answer to each question.**
Cities are the true fonts of creativity... With their diverse populations, dense social networks, and public spaces where people can meet spontaneously and serendipitously, they spark and catalyze new ideas. With their infrastructure for finance, organization and trade, they allow those ideas to be swiftly actualized.
As for what staunches creativity, that's easy, if ironic. It's the very institutions that we build to manage, exploit and perpetuate the fruits of creativity — our big bureaucracies, and sad to say, too many of our schools. Creativity is disruptive; schools and organizations are regimented, standardized and stultifying.
The education expert Sir Ken Robinson points to a 1968 study reporting on a group of 1,600 children who were tested over time for their ability to think in out-of-the-box ways. When the children were between 3 and 5 years old, 98 percent achieved positive scores. When they were 8 to 10, only 32 percent passed the same test, and only 10 percent at 13 to 15. When 280,000 25-year-olds took the test, just 2 percent passed. By the time we are adults, our creativity has been wrung out of us.
I once asked the great urbanist Jane Jacobs what makes some places more creative than others. She said, essentially, that the question was an easy one. All cities, she said, were filled with creative people; that's our default state as people. But some cities had more than their shares of leaders, people and institutions that blocked out that creativity. She called them "squelchers."
Creativity (or the lack of it) follows the same general contours of the great socio-economic divide — our rising inequality — that plagues us. According to my own estimates, roughly a third of us across the United States, and perhaps as much as half of us in our most creative cities — are able to do work which engages our creative faculties to some extent, whether as artists, musicians, writers, techies, innovators, entrepreneurs, doctors, lawyers, journalists or educators — those of us who work with our minds. That leaves a group that I term "the other 66 percent," who toil in low-wage rote and rotten jobs — if they have jobs at all — in which their creativity is subjugated, ignored or wasted.
Creativity itself is not in danger. It's flourishing is all around us — in science and technology, arts and culture, in our rapidly revitalizing cities. But we still have a long way to go if we want to build a truly creative society that supports and rewards the creativity of each and every one of us.

\vspace{1em}
\textbf{Question 29}

The five sentences labelled (1, 2, 3, 4, 5) given in this question, when properly sequenced, form a coherent paragraph. Each sentence is labelled with a number. Decide on the proper order for the sentences and **key in this sequence of FIVE NUMBERS** as your answer.  
  
1. Before plants can take life from atmosphere, nitrogen must undergo transformations similar to ones that food undergoes in our digestive machinery.  
2. In its aerial form nitrogen is insoluble, unusable and is in need of transformation.  
3. Lightning starts the series of chemical reactions that need to happen to nitrogen, ultimately helping it nourish our earth.  
4. Nitrogen — an essential food for plants — is an abundant resource, with about 22 million tons of it floating over each square mile of earth.  
5. One of the most dramatic examples in nature of ill wind that blows goodness is lightning.
Backspace  
789  
456  
123  
0.-  
Clear All  

Submit
[](https://cracku.in/29-the-five-sentences-labelled-1-2-3-4-5-given-in-thi-x-cat-2017-shift-2)

\begin{enumerate}[label=\Alph*.]
\end{enumerate}
\vspace{1em}
\textbf{Solution:}

On closely reading the sentences, we see that the topic of discussion is how chemical reactions started by lightning affects nitrogen in the air which ultimately help nourish the earth. Hence, sentence 5 which introduces lightning to us must be the starting sentence. Sentence 3 further elaborates on chemical reactions, started by lightning, which affects nitrogen. Hence, sentence 3 logically follows sentence 5. Sentence 4 states how nitrogen helps in nourishing the earth. Hence, sentence 4 must follow sentence 3. Sentence 2 states why in its natural form nitrogen is unusable and sentence 1 further elaborates on why nitrogen must undergo transformation for plants to use it. Hence sentences 2 and 1 form a pair which must follow sentence 4. Hence, the correct sequence 53421.

\end{problem}

\newpage
\begin{problem}{}{}
% Subject: varc
\textbf{Instructions}

Instructions   

**The passage below is accompanied by a set of six questions. Choose the best answer to each question.**
Cities are the true fonts of creativity... With their diverse populations, dense social networks, and public spaces where people can meet spontaneously and serendipitously, they spark and catalyze new ideas. With their infrastructure for finance, organization and trade, they allow those ideas to be swiftly actualized.
As for what staunches creativity, that's easy, if ironic. It's the very institutions that we build to manage, exploit and perpetuate the fruits of creativity — our big bureaucracies, and sad to say, too many of our schools. Creativity is disruptive; schools and organizations are regimented, standardized and stultifying.
The education expert Sir Ken Robinson points to a 1968 study reporting on a group of 1,600 children who were tested over time for their ability to think in out-of-the-box ways. When the children were between 3 and 5 years old, 98 percent achieved positive scores. When they were 8 to 10, only 32 percent passed the same test, and only 10 percent at 13 to 15. When 280,000 25-year-olds took the test, just 2 percent passed. By the time we are adults, our creativity has been wrung out of us.
I once asked the great urbanist Jane Jacobs what makes some places more creative than others. She said, essentially, that the question was an easy one. All cities, she said, were filled with creative people; that's our default state as people. But some cities had more than their shares of leaders, people and institutions that blocked out that creativity. She called them "squelchers."
Creativity (or the lack of it) follows the same general contours of the great socio-economic divide — our rising inequality — that plagues us. According to my own estimates, roughly a third of us across the United States, and perhaps as much as half of us in our most creative cities — are able to do work which engages our creative faculties to some extent, whether as artists, musicians, writers, techies, innovators, entrepreneurs, doctors, lawyers, journalists or educators — those of us who work with our minds. That leaves a group that I term "the other 66 percent," who toil in low-wage rote and rotten jobs — if they have jobs at all — in which their creativity is subjugated, ignored or wasted.
Creativity itself is not in danger. It's flourishing is all around us — in science and technology, arts and culture, in our rapidly revitalizing cities. But we still have a long way to go if we want to build a truly creative society that supports and rewards the creativity of each and every one of us.

\vspace{1em}
\textbf{Question 30}

The five sentences (labelled 1, 2, 3, 4, 5) given in this question, when properly sequenced, form a coherent paragraph. Each sentence is labelled with a number. Decide on the proper order for the sentences and key in this sequence of five numbers as your answer.  
1. This has huge implications for the health care system as it operates today, where depleted resources and time lead to patients rotating in and out of doctor's offices, oftentimes receiving minimal care or concern (what is commonly referred to as "bed side manner") from doctors.  
2. The placebo effect is when an individual's medical condition or pain shows signs of improvement based on a fake intervention that has been presented to them as a real one and used to be regularly dismissed by researchers as a psychological effect.  
3. The placebo effect is not solely based on believing in treatment, however, as the clinical setting in which treatments are administered is also paramount.  
4. That the mind has the power to trigger biochemical changes because the individual believes that a given drug or intervention will be effective could empower chronic patients through the notion of our bodies' capacity for self-healing.  
5. Placebo effects are now studied not just as foils for "real" interventions but as a potential portal into the self-healing powers of the body.
Backspace  
789  
456  
123  
0.-  
Clear All  

Submit
[](https://cracku.in/30-the-five-sentences-labelled-1-2-3-4-5-given-in-thi-x-cat-2017-shift-2)

\begin{enumerate}[label=\Alph*.]
\end{enumerate}
\vspace{1em}
\textbf{Solution:}

On closely reading the sentences, we can see that the passage is about placebo effect. Sentence 2, which introduces the placebo effect, must be the starting sentence. Sentence 5 states that placebo effect are now not just studied for real interventions, as stated in sentence 2, but as potential portal into self healing power. Hence, sentence 5 logically follows sentence 2. Sentence 4 which elaborates on self healing must follow sentence 5. Sentence 3 makes a point that apart from the belief in the treatment, the clinical setting also has a role to play. Sentence 1 elaborates on this. Hence, 31 is a pair which must follow sentence 4. Thus, the correct order is 25431.

\end{problem}

\newpage
\begin{problem}{}{}
% Subject: varc
\textbf{Instructions}

Instructions   

**The passage below is accompanied by a set of six questions. Choose the best answer to each question.**
Cities are the true fonts of creativity... With their diverse populations, dense social networks, and public spaces where people can meet spontaneously and serendipitously, they spark and catalyze new ideas. With their infrastructure for finance, organization and trade, they allow those ideas to be swiftly actualized.
As for what staunches creativity, that's easy, if ironic. It's the very institutions that we build to manage, exploit and perpetuate the fruits of creativity — our big bureaucracies, and sad to say, too many of our schools. Creativity is disruptive; schools and organizations are regimented, standardized and stultifying.
The education expert Sir Ken Robinson points to a 1968 study reporting on a group of 1,600 children who were tested over time for their ability to think in out-of-the-box ways. When the children were between 3 and 5 years old, 98 percent achieved positive scores. When they were 8 to 10, only 32 percent passed the same test, and only 10 percent at 13 to 15. When 280,000 25-year-olds took the test, just 2 percent passed. By the time we are adults, our creativity has been wrung out of us.
I once asked the great urbanist Jane Jacobs what makes some places more creative than others. She said, essentially, that the question was an easy one. All cities, she said, were filled with creative people; that's our default state as people. But some cities had more than their shares of leaders, people and institutions that blocked out that creativity. She called them "squelchers."
Creativity (or the lack of it) follows the same general contours of the great socio-economic divide — our rising inequality — that plagues us. According to my own estimates, roughly a third of us across the United States, and perhaps as much as half of us in our most creative cities — are able to do work which engages our creative faculties to some extent, whether as artists, musicians, writers, techies, innovators, entrepreneurs, doctors, lawyers, journalists or educators — those of us who work with our minds. That leaves a group that I term "the other 66 percent," who toil in low-wage rote and rotten jobs — if they have jobs at all — in which their creativity is subjugated, ignored or wasted.
Creativity itself is not in danger. It's flourishing is all around us — in science and technology, arts and culture, in our rapidly revitalizing cities. But we still have a long way to go if we want to build a truly creative society that supports and rewards the creativity of each and every one of us.

\vspace{1em}
\textbf{Question 31}

The **five sentences** (labelled 1, 2, 3, 4, 5) given in this question, when properly sequenced, form a coherent paragraph. Each sentence is labelled with a number. Decide on the proper order for the sentences and key in this sequence of **five numbers** as your answer.  
1. Johnson treated English very practically, as a living language, with many different shades of meaning and adopted his definitions on the principle of English common law — according to precedent.  
2. Masking a profound inner torment, Johnson found solace in compiling the words of a language that was, in its coarse complexity and comprehensive genius, the precise analogue of his character.  
3. Samuel Johnson was a pioneer who raised common sense to heights of genius, and a man of robust popular instincts whose watchwords were clarity, precision and simplicity.  
4. The 18th century English reader, in the new world of global trade and global warfare, needed a dictionary with authoritative acts of definition of words of a language that was becoming seeded throughout the first British empire by a vigorous and practical champion.  
5. The Johnson who challenged Bishop Berkeley's solipsist theory of the nonexistence of matter by kicking a large stone ("I refute it thus") is the same Johnson for whom language must have a daily practical use.
Backspace  
789  
456  
123  
0.-  
Clear All  

Submit
[](https://cracku.in/31-the-five-sentences-labelled-1-2-3-4-5-given-in-thi-x-cat-2017-shift-2)

\begin{enumerate}[label=\Alph*.]
\end{enumerate}
\vspace{1em}
\textbf{Solution:}

Sentence 4 should be the opening sentence since it talks about the need for a dictionary in the 18th century. The other 4 statements talk about Samuel Johnson.   
3 must follow 4 since it introduces the subject, Samuel Johnson. Only sentence 3 contains the full name of Samuel Johnson.   
3 should be followed by 5 since it describes Johnson's character. 5 plays the role of a general introduction and hence, it should be placed before any specific detail regarding Johnson's contribution to the dictionary is introduced.  
Out of sentences 1 and 2, 1 should precede 2 since it establishes that Johnson worked on English and sentence 2 explains the innate connection between Johnson and the language (English).  
43512 is the correct order.

\end{problem}

\newpage
\begin{problem}{}{}
% Subject: varc
\textbf{Instructions}

Instructions   

**The passage below is accompanied by a set of six questions. Choose the best answer to each question.**
Cities are the true fonts of creativity... With their diverse populations, dense social networks, and public spaces where people can meet spontaneously and serendipitously, they spark and catalyze new ideas. With their infrastructure for finance, organization and trade, they allow those ideas to be swiftly actualized.
As for what staunches creativity, that's easy, if ironic. It's the very institutions that we build to manage, exploit and perpetuate the fruits of creativity — our big bureaucracies, and sad to say, too many of our schools. Creativity is disruptive; schools and organizations are regimented, standardized and stultifying.
The education expert Sir Ken Robinson points to a 1968 study reporting on a group of 1,600 children who were tested over time for their ability to think in out-of-the-box ways. When the children were between 3 and 5 years old, 98 percent achieved positive scores. When they were 8 to 10, only 32 percent passed the same test, and only 10 percent at 13 to 15. When 280,000 25-year-olds took the test, just 2 percent passed. By the time we are adults, our creativity has been wrung out of us.
I once asked the great urbanist Jane Jacobs what makes some places more creative than others. She said, essentially, that the question was an easy one. All cities, she said, were filled with creative people; that's our default state as people. But some cities had more than their shares of leaders, people and institutions that blocked out that creativity. She called them "squelchers."
Creativity (or the lack of it) follows the same general contours of the great socio-economic divide — our rising inequality — that plagues us. According to my own estimates, roughly a third of us across the United States, and perhaps as much as half of us in our most creative cities — are able to do work which engages our creative faculties to some extent, whether as artists, musicians, writers, techies, innovators, entrepreneurs, doctors, lawyers, journalists or educators — those of us who work with our minds. That leaves a group that I term "the other 66 percent," who toil in low-wage rote and rotten jobs — if they have jobs at all — in which their creativity is subjugated, ignored or wasted.
Creativity itself is not in danger. It's flourishing is all around us — in science and technology, arts and culture, in our rapidly revitalizing cities. But we still have a long way to go if we want to build a truly creative society that supports and rewards the creativity of each and every one of us.

\vspace{1em}
\textbf{Question 32}

Five sentences related to a topic are given below. Four of them can be put together to form a meaningful and coherent short paragraph. Identify the odd one out.  
1. Although we are born with the gift of language, research shows that we are surprisingly unskilled when it comes to communicating with others.  
2. We must carefully orchestrate our speech if we want to achieve our goals and bring our dreams to fruition.   
3. We often choose our words without thought, oblivious of the emotional effects they can have on others.  
4. We talk more than we need to, ignoring the effect we are having on those listening to us.  
5. We listen poorly, without realizing it, and we often fail to pay attention to the subtle meanings conveyed by facial expressions, body gestures, and the tone and cadence of our voice.
Backspace  
789  
456  
123  
0.-  
Clear All  

Submit
[](https://cracku.in/32-five-sentences-related-to-a-topic-are-given-below--x-cat-2017-shift-2)

\begin{enumerate}[label=\Alph*.]
\end{enumerate}
\vspace{1em}
\textbf{Solution:}

The paragraph is about how inefficient we are when it comes to usage of language. Sentence 2 stands out as an imperative or instructive statement whereas the other statements simply elaborate on the point that we are unskilled at language usage. Therefore, option 2 is the odd sentence.

\end{problem}

\newpage
\begin{problem}{}{}
% Subject: varc
\textbf{Instructions}

Instructions   

**The passage below is accompanied by a set of six questions. Choose the best answer to each question.**
Cities are the true fonts of creativity... With their diverse populations, dense social networks, and public spaces where people can meet spontaneously and serendipitously, they spark and catalyze new ideas. With their infrastructure for finance, organization and trade, they allow those ideas to be swiftly actualized.
As for what staunches creativity, that's easy, if ironic. It's the very institutions that we build to manage, exploit and perpetuate the fruits of creativity — our big bureaucracies, and sad to say, too many of our schools. Creativity is disruptive; schools and organizations are regimented, standardized and stultifying.
The education expert Sir Ken Robinson points to a 1968 study reporting on a group of 1,600 children who were tested over time for their ability to think in out-of-the-box ways. When the children were between 3 and 5 years old, 98 percent achieved positive scores. When they were 8 to 10, only 32 percent passed the same test, and only 10 percent at 13 to 15. When 280,000 25-year-olds took the test, just 2 percent passed. By the time we are adults, our creativity has been wrung out of us.
I once asked the great urbanist Jane Jacobs what makes some places more creative than others. She said, essentially, that the question was an easy one. All cities, she said, were filled with creative people; that's our default state as people. But some cities had more than their shares of leaders, people and institutions that blocked out that creativity. She called them "squelchers."
Creativity (or the lack of it) follows the same general contours of the great socio-economic divide — our rising inequality — that plagues us. According to my own estimates, roughly a third of us across the United States, and perhaps as much as half of us in our most creative cities — are able to do work which engages our creative faculties to some extent, whether as artists, musicians, writers, techies, innovators, entrepreneurs, doctors, lawyers, journalists or educators — those of us who work with our minds. That leaves a group that I term "the other 66 percent," who toil in low-wage rote and rotten jobs — if they have jobs at all — in which their creativity is subjugated, ignored or wasted.
Creativity itself is not in danger. It's flourishing is all around us — in science and technology, arts and culture, in our rapidly revitalizing cities. But we still have a long way to go if we want to build a truly creative society that supports and rewards the creativity of each and every one of us.

\vspace{1em}
\textbf{Question 33}

Five sentences related to a topic are given below. Four of them can be put together to form a meaningful and coherent short paragraph. Identify the odd one out.  
1: Over the past fortnight, one of its finest champions managed to pull off a similar impression.  
2. Wimbledon's greatest illusion is the sense of timelessness it evokes.  
3. At 35 years and 342 days, Roger Federer became the oldest man to win the singles title in the Open Era — a full 14 years after he first claimed the title as a scruffy, pony-tailed upstart.  
4. Once he had survived the opening week, the second week witnessed the range of a rested Federer's genius.   
5. Given that his method isn't reliant on explosive athleticism or muscular ball-striking, both vulnerable to decay, there is cause to believe that Federer will continue to enchant for a while longer.
Backspace  
789  
456  
123  
0.-  
Clear All  

Submit
[](https://cracku.in/33-five-sentences-related-to-a-topic-are-given-below--x-cat-2017-shift-2)

\begin{enumerate}[label=\Alph*.]
\end{enumerate}
\vspace{1em}
\textbf{Solution:}

The paragraph is about the timelessness of the Wimbledon and Roger Federer. Sentences 1, 2 and 3 are clearly the part of the paragraph.   
Wimbledon - one of its finest champion - Roger Federer  
Now, sentences 4 and 5 are close. Sentence 4 talks about a specific event. Whereas, sentence 5 is inline with the idea of timelessness and hence, can be used as the concluding sentence of the paragraph. Moreover, sentence 4 appears hanging (Missing some previous statement). Therefore, sentence 4 is the odd one out.

\end{problem}

\newpage
\begin{problem}{}{}
% Subject: varc
\textbf{Instructions}

Instructions   

**The passage below is accompanied by a set of six questions. Choose the best answer to each question.**
Cities are the true fonts of creativity... With their diverse populations, dense social networks, and public spaces where people can meet spontaneously and serendipitously, they spark and catalyze new ideas. With their infrastructure for finance, organization and trade, they allow those ideas to be swiftly actualized.
As for what staunches creativity, that's easy, if ironic. It's the very institutions that we build to manage, exploit and perpetuate the fruits of creativity — our big bureaucracies, and sad to say, too many of our schools. Creativity is disruptive; schools and organizations are regimented, standardized and stultifying.
The education expert Sir Ken Robinson points to a 1968 study reporting on a group of 1,600 children who were tested over time for their ability to think in out-of-the-box ways. When the children were between 3 and 5 years old, 98 percent achieved positive scores. When they were 8 to 10, only 32 percent passed the same test, and only 10 percent at 13 to 15. When 280,000 25-year-olds took the test, just 2 percent passed. By the time we are adults, our creativity has been wrung out of us.
I once asked the great urbanist Jane Jacobs what makes some places more creative than others. She said, essentially, that the question was an easy one. All cities, she said, were filled with creative people; that's our default state as people. But some cities had more than their shares of leaders, people and institutions that blocked out that creativity. She called them "squelchers."
Creativity (or the lack of it) follows the same general contours of the great socio-economic divide — our rising inequality — that plagues us. According to my own estimates, roughly a third of us across the United States, and perhaps as much as half of us in our most creative cities — are able to do work which engages our creative faculties to some extent, whether as artists, musicians, writers, techies, innovators, entrepreneurs, doctors, lawyers, journalists or educators — those of us who work with our minds. That leaves a group that I term "the other 66 percent," who toil in low-wage rote and rotten jobs — if they have jobs at all — in which their creativity is subjugated, ignored or wasted.
Creativity itself is not in danger. It's flourishing is all around us — in science and technology, arts and culture, in our rapidly revitalizing cities. But we still have a long way to go if we want to build a truly creative society that supports and rewards the creativity of each and every one of us.

\vspace{1em}
\textbf{Question 34}

**Five sentences related to a topic are given below. Four of them can be put together to form a meaningful and coherent short paragraph. Identify the odd one out.**
1. Those geometric symbols and aerodynamic swooshes are more than just skin deep.  
2. The Commonwealth Bank logo — a yellow diamond, with a black chunk sliced out in one corner — is so recognisable that the bank doesn't even use its full name in its advertising.  
3. It's not just logos with hidden shapes; sometimes brands will have meanings or stories within them that are deliberately vague or lost in time, urging you to delve deeper to solve the riddle.  
4. Graphic designers embed cryptic references because it adds a story to the brand; they want people to spend more time with a brand and have that idea that they are an insider if they can understand the hidden message.   
5. But the CommBank logo has more to it than meets the eye, as squirrelled away in that diamond is the Southern Cross constellation.
Backspace  
789  
456  
123  
0.-  
Clear All  

Submit
[](https://cracku.in/34-five-sentences-related-to-a-topic-are-given-below--x-cat-2017-shift-2)

\begin{enumerate}[label=\Alph*.]
\end{enumerate}
\vspace{1em}
\textbf{Solution:}

After reading all the sentences, we know that the paragraph is talking about the logos and brands and the purpose behind them. Statement 2 is the opening sentence which tells us about the Commonwealth Bank's logo. Statement 5 provides more information about the logo, mentioned in statement 2, that there is more to what is visible to eyes about the Bank's logo. Statement 3 gives another insight that it's not only logos, but brand names are also created with hidden meanings. Statement 4 provides the reasoning behind passing subtle messages through brands and logos. Thus, 2534 forms a meaningful paragraph.  
Statement 1 describes something as geometric symbols and aerodynamic swooshes which does not connect well with the other four sentences. Thus, statement 1 does not fit in the paragraph.  
Hence, option 1 is the correct answer.
[](https://cracku.in/downloads/1882)
  *  
  *

\end{problem}

\newpage
\begin{problem}{}{}
% Subject: varc
\textbf{Instructions}

Instructions   

**The passage below is accompanied by a set of questions. Choose the best answer to each question.**
The word ‘anarchy’ comes from the Greek '_anarkhia_ ', meaning contrary to authority or without a ruler, and was used in a derogatory sense until 1840, when it was adopted by Pierre-Joseph Proudhon to describe his political and social ideology. Proudhon argued that organization without government was both possible and desirable. In the evolution of political ideas, anarchism can be seen as an ultimate projection of both liberalism and socialism, and the differing strands of anarchist thought can be related to their emphasis on one or the other of these.
Historically, anarchism arose not only as an explanation of the gulf between the rich and the poor in any community, and of the reason why the poor have been obliged to fight for their share of a common inheritance, but as a radical answer to the question ‘What went wrong?’ that followed the ultimate outcome of the French Revolution. It had ended not only with a reign of terror and the emergence of a newly rich ruling caste, but with a new adored emperor, Napoleon Bonaparte, strutting through his conquered territories.
The anarchists and their precursors were unique on the political Left in affirming that workers and peasants, grasping the chance that arose to bring an end to centuries of exploitation and tyranny, were inevitably betrayed by the new class of politicians, whose first priority was to re-establish a centralized state power. After every revolutionary uprising, usually won at a heavy cost for ordinary populations, the new rulers had no hesitation in applying violence and terror, a secret police, and a professional army to maintain their control.
For anarchists the state itself is the enemy, and they have applied the same interpretation to the outcome of every revolution of the 19th and 20th centuries. This is not merely because every state keeps a watchful and sometimes punitive eye on its dissidents, but because every state protects the privileges of the powerful.
The mainstream of anarchist propaganda for more than a century has been anarchist- communism, which argues that property in land, natural resources, and the means of production should be held in mutual control by local communities, federating for innumerable joint purposes with other communes. It differs from state socialism in opposing the concept of any central authority. Some anarchists prefer to distinguish between anarchist-communism and collectivist anarchism in order to stress the obviously desirable freedom of an individual or family to possess the resources needed for living, while not implying the right to own the resources needed by others. . . .
There are, unsurprisingly, several traditions of individualist anarchism, one of them deriving from the ‘conscious egoism’ of the German writer Max Stirner (1806-56), and another from a remarkable series of 19th-century American figures who argued that in protecting our own autonomy and associating with others for common advantages, we are promoting the good of all. These thinkers differed from free-market liberals in their absolute mistrust of American capitalism, and in their emphasis on mutualism.

\vspace{1em}
\textbf{Question 1}

The author makes all of the following arguments in the passage, EXCEPT:

\begin{enumerate}[label=\Alph*.]
    \item The failure of the French Revolution was because of its betrayal by the new class of politicians who emerged from it.
    \item For anarchists, the state is the enemy because all states apply violence and terror to maintain their control.
    \item Individualist anarchism is actually constituted of many streams, all of which focus on the autonomy of the individual.
    \item The popular perception of anarchism as espousing lawlessness and violence comes from a mainstream mistrust of collectivism.
[](https://cracku.in/1-the-author-makes-all-of-the-following-arguments-in-x-cat-2020-slot-1)
\end{enumerate}
\vspace{1em}
\textbf{Solution:}

_Option A_ : This element is discussed in the third paragraph: {".._.The anarchists and their precursors were unique on the political Left in affirming that workers and peasants, grasping the chance that arose to bring an end to centuries of exploitation and tyranny, were inevitably betrayed by the new class of politicians, whose first priority was to re-establish a centralized state power._.."} Prior to making this claim, the failure of the French revolution is highlighted, and this segment is subsequently tied to it. Hence, Option A is definitely a point that is made in the passage.
_Option B:_ This has been outlined in the fourth paragraph: {".._.For anarchists, the state itself is the enemy, and they have applied the same interpretation to the outcome of every revolution of the 19th and 20th centuries. This is not merely because every state keeps a watchful and sometimes punitive eye on its dissidents, but because every state protects the privileges of the powerful.._. "}
_Option C_ : The last two paragraphs present the perception associated with individualist-anarchism: {".._.desirable freedom of an individual or family to possess the resources needed for living, while not implying the right to own the resources needed by others. . . .__There are, unsurprisingly, several traditions of individualist anarchism, one of them deriving from the ‘conscious egoism’ of the German writer Max Stirner (1806-56), and another from a remarkable series of 19th-century American figures who argued that in protecting our own autonomy and associating with others for common advantages, we are promoting the good of all.._."} Thus, we can evidently understand the statement in C to be true.
_Option D_ , however, is not inferable from the discussion undertaken in the passage. The author does not pinpoint "a mainstream mistrust of collectivism" as the source of the misconception of anarchism "espousing lawlessness and violence". No causal element is discussed, and thus, the statement in D is not an argument that is made in the passage. 
Hence, Option D is the correct answer.

\end{problem}

\newpage
\begin{problem}{}{}
% Subject: varc
\textbf{Instructions}

Instructions   

**The passage below is accompanied by a set of questions. Choose the best answer to each question.**
The word ‘anarchy’ comes from the Greek '_anarkhia_ ', meaning contrary to authority or without a ruler, and was used in a derogatory sense until 1840, when it was adopted by Pierre-Joseph Proudhon to describe his political and social ideology. Proudhon argued that organization without government was both possible and desirable. In the evolution of political ideas, anarchism can be seen as an ultimate projection of both liberalism and socialism, and the differing strands of anarchist thought can be related to their emphasis on one or the other of these.
Historically, anarchism arose not only as an explanation of the gulf between the rich and the poor in any community, and of the reason why the poor have been obliged to fight for their share of a common inheritance, but as a radical answer to the question ‘What went wrong?’ that followed the ultimate outcome of the French Revolution. It had ended not only with a reign of terror and the emergence of a newly rich ruling caste, but with a new adored emperor, Napoleon Bonaparte, strutting through his conquered territories.
The anarchists and their precursors were unique on the political Left in affirming that workers and peasants, grasping the chance that arose to bring an end to centuries of exploitation and tyranny, were inevitably betrayed by the new class of politicians, whose first priority was to re-establish a centralized state power. After every revolutionary uprising, usually won at a heavy cost for ordinary populations, the new rulers had no hesitation in applying violence and terror, a secret police, and a professional army to maintain their control.
For anarchists the state itself is the enemy, and they have applied the same interpretation to the outcome of every revolution of the 19th and 20th centuries. This is not merely because every state keeps a watchful and sometimes punitive eye on its dissidents, but because every state protects the privileges of the powerful.
The mainstream of anarchist propaganda for more than a century has been anarchist- communism, which argues that property in land, natural resources, and the means of production should be held in mutual control by local communities, federating for innumerable joint purposes with other communes. It differs from state socialism in opposing the concept of any central authority. Some anarchists prefer to distinguish between anarchist-communism and collectivist anarchism in order to stress the obviously desirable freedom of an individual or family to possess the resources needed for living, while not implying the right to own the resources needed by others. . . .
There are, unsurprisingly, several traditions of individualist anarchism, one of them deriving from the ‘conscious egoism’ of the German writer Max Stirner (1806-56), and another from a remarkable series of 19th-century American figures who argued that in protecting our own autonomy and associating with others for common advantages, we are promoting the good of all. These thinkers differed from free-market liberals in their absolute mistrust of American capitalism, and in their emphasis on mutualism.

\vspace{1em}
\textbf{Question 2}

The author believes that the new ruling class of politicians betrayed the principles of the French Revolution, but does not specify in what way. In the context of the passage, which statement below is the likeliest explanation of that betrayal?

\begin{enumerate}[label=\Alph*.]
    \item The new ruling class was constituted mainly of anarchists who were against the destructive impact of the Revolution on the market.
    \item The anarchists did not want a new ruling class, but were not politically strong enough to stop them.
    \item The new ruling class struck a deal with the old ruling class to share power between them.
    \item The new ruling class rode to power on the strength of the workers’ revolutionary anger, but then turned to oppress that very class.
[](https://cracku.in/2-the-author-believes-that-the-new-ruling-class-of-p-x-cat-2020-slot-1)
\end{enumerate}
\vspace{1em}
\textbf{Solution:}

The third paragraph renders us with sufficient information concerning the manner of betrayal. The failure of the Frech Revolution is being attributed to the deceit of the new class of politicians who not only prioritised the re-establishment and centralisation of authority but also unleashed a reign of terror on the ordinary masses. The author hints at this re-installation of power and the atrocities committed on the common people, who are the very drivers of the revolution, being the desertion of the principles of the French Revolution by the new ruling class of politicians. Any statement that aligns with this understanding would be a suitable answer. Option A brings in the concept of an impact on the market, which is not mentioned in the passage. Additionally, it does not signify betrayal of any kind. Same can be said about the statement in Option B. Option C is a strong candidate that symbolises betrayal. However, can we appropriately associate it with the underlying principles of the revolution? The author's primary focus is on the cost borne by the ordinary populations, and despite the revolution being realised by their sacrifices, they are still the ones being subjected to inequity and torment. Although C highlights betrayal, it does not truly coincide with the idea being conveyed. The afore-mentioned element is appropriately captured by Option D: it highlights how the new class of politicians utilised the indignation of the masses/workers for the revolution - to topple authority and rise to power - and then stab them in the back by turning to oppress them.

\end{problem}

\newpage
\begin{problem}{}{}
% Subject: varc
\textbf{Instructions}

Instructions   

**The passage below is accompanied by a set of questions. Choose the best answer to each question.**
The word ‘anarchy’ comes from the Greek '_anarkhia_ ', meaning contrary to authority or without a ruler, and was used in a derogatory sense until 1840, when it was adopted by Pierre-Joseph Proudhon to describe his political and social ideology. Proudhon argued that organization without government was both possible and desirable. In the evolution of political ideas, anarchism can be seen as an ultimate projection of both liberalism and socialism, and the differing strands of anarchist thought can be related to their emphasis on one or the other of these.
Historically, anarchism arose not only as an explanation of the gulf between the rich and the poor in any community, and of the reason why the poor have been obliged to fight for their share of a common inheritance, but as a radical answer to the question ‘What went wrong?’ that followed the ultimate outcome of the French Revolution. It had ended not only with a reign of terror and the emergence of a newly rich ruling caste, but with a new adored emperor, Napoleon Bonaparte, strutting through his conquered territories.
The anarchists and their precursors were unique on the political Left in affirming that workers and peasants, grasping the chance that arose to bring an end to centuries of exploitation and tyranny, were inevitably betrayed by the new class of politicians, whose first priority was to re-establish a centralized state power. After every revolutionary uprising, usually won at a heavy cost for ordinary populations, the new rulers had no hesitation in applying violence and terror, a secret police, and a professional army to maintain their control.
For anarchists the state itself is the enemy, and they have applied the same interpretation to the outcome of every revolution of the 19th and 20th centuries. This is not merely because every state keeps a watchful and sometimes punitive eye on its dissidents, but because every state protects the privileges of the powerful.
The mainstream of anarchist propaganda for more than a century has been anarchist- communism, which argues that property in land, natural resources, and the means of production should be held in mutual control by local communities, federating for innumerable joint purposes with other communes. It differs from state socialism in opposing the concept of any central authority. Some anarchists prefer to distinguish between anarchist-communism and collectivist anarchism in order to stress the obviously desirable freedom of an individual or family to possess the resources needed for living, while not implying the right to own the resources needed by others. . . .
There are, unsurprisingly, several traditions of individualist anarchism, one of them deriving from the ‘conscious egoism’ of the German writer Max Stirner (1806-56), and another from a remarkable series of 19th-century American figures who argued that in protecting our own autonomy and associating with others for common advantages, we are promoting the good of all. These thinkers differed from free-market liberals in their absolute mistrust of American capitalism, and in their emphasis on mutualism.

\vspace{1em}
\textbf{Question 3}

Which one of the following best expresses the similarity between American individualist anarchists and free-market liberals as well as the difference between the former and the latter?

\begin{enumerate}[label=\Alph*.]
    \item Both prioritise individual autonomy; but the former also emphasise mutual dependence, while the latter do not do so.
    \item Both reject the regulatory power of the state; but the former favour a people’s state, while the latter favour state intervention in markets.
    \item Both have sophisticated arguments for capitalism; but the former argue for a morally upright capitalism, while the latter argue that the market is the only morality
    \item Both are founded on the moral principles of altruism; but the latter conceive of the market as a force too mystical for the former to comprehend.
[](https://cracku.in/3-which-one-of-the-following-best-expresses-the-simi-x-cat-2020-slot-1)
\end{enumerate}
\vspace{1em}
\textbf{Solution:}

We can derive our understanding of these two groups from the following: {"._..There are, unsurprisingly, several traditions of individualist anarchism, one of them deriving from the ‘conscious egoism’ of the German writer Max Stirner (1806-56), and another from a remarkable series of 19th-century American figures who argued that in**protecting our own autonomy** and **associating with others for common advantages** , we are promoting the good of all. These thinkers **differed** from free-market liberals in their absolute **mistrust of American capitalism** , and in their **emphasis on mutualism**._.."} 
It is evident that both the groups prioritize individual autonomy; however, the latter does not put emphasis on mutualism. The mistrust of American capitalism is an additional facet marking the difference {but not captured in any of the options}. _Option A_ aptly highlights the difference between the two factions, and is, hence, the correct answer.
_Option B_ : Preferences concerning state configuration or involvement are not discussed in the passage, and hence, we can eliminate this option.
_Option C_ : Perspectives regarding the nature of capitalism are not stated in the passage, and thus, we can effortlessly discard this option
_Option D_ : None of the elements stated in this option can be understood from the passage, and therefore, can be rejected.

\end{problem}

\newpage
\begin{problem}{}{}
% Subject: varc
\textbf{Instructions}

Instructions   

**The passage below is accompanied by a set of questions. Choose the best answer to each question.**
The word ‘anarchy’ comes from the Greek '_anarkhia_ ', meaning contrary to authority or without a ruler, and was used in a derogatory sense until 1840, when it was adopted by Pierre-Joseph Proudhon to describe his political and social ideology. Proudhon argued that organization without government was both possible and desirable. In the evolution of political ideas, anarchism can be seen as an ultimate projection of both liberalism and socialism, and the differing strands of anarchist thought can be related to their emphasis on one or the other of these.
Historically, anarchism arose not only as an explanation of the gulf between the rich and the poor in any community, and of the reason why the poor have been obliged to fight for their share of a common inheritance, but as a radical answer to the question ‘What went wrong?’ that followed the ultimate outcome of the French Revolution. It had ended not only with a reign of terror and the emergence of a newly rich ruling caste, but with a new adored emperor, Napoleon Bonaparte, strutting through his conquered territories.
The anarchists and their precursors were unique on the political Left in affirming that workers and peasants, grasping the chance that arose to bring an end to centuries of exploitation and tyranny, were inevitably betrayed by the new class of politicians, whose first priority was to re-establish a centralized state power. After every revolutionary uprising, usually won at a heavy cost for ordinary populations, the new rulers had no hesitation in applying violence and terror, a secret police, and a professional army to maintain their control.
For anarchists the state itself is the enemy, and they have applied the same interpretation to the outcome of every revolution of the 19th and 20th centuries. This is not merely because every state keeps a watchful and sometimes punitive eye on its dissidents, but because every state protects the privileges of the powerful.
The mainstream of anarchist propaganda for more than a century has been anarchist- communism, which argues that property in land, natural resources, and the means of production should be held in mutual control by local communities, federating for innumerable joint purposes with other communes. It differs from state socialism in opposing the concept of any central authority. Some anarchists prefer to distinguish between anarchist-communism and collectivist anarchism in order to stress the obviously desirable freedom of an individual or family to possess the resources needed for living, while not implying the right to own the resources needed by others. . . .
There are, unsurprisingly, several traditions of individualist anarchism, one of them deriving from the ‘conscious egoism’ of the German writer Max Stirner (1806-56), and another from a remarkable series of 19th-century American figures who argued that in protecting our own autonomy and associating with others for common advantages, we are promoting the good of all. These thinkers differed from free-market liberals in their absolute mistrust of American capitalism, and in their emphasis on mutualism.

\vspace{1em}
\textbf{Question 4}

Of the following sets of concepts, identify the set that is conceptually closest to the concerns of the passage.

\begin{enumerate}[label=\Alph*.]
    \item Anarchism, State, Individual, Freedom.
    \item Revolution, State, Strike, Egoism.
    \item Revolution, State, Protection, Liberals.
    \item Anarchism, Betrayal, Power, State.
[](https://cracku.in/4-of-the-following-sets-of-concepts-identify-the-set-x-cat-2020-slot-1)
\end{enumerate}
\vspace{1em}
\textbf{Solution:}

One has to be mindful of the question: the set that is **conceptually closest** to the **concerns** of the passage.
The passage is clearly about 'Anarchism', and the discussion pertaining to 'Revolution' is a sub-element. This elementary understanding helps us pare down the viable choices to Options A and D. 'Betrayal' is again a minor sub-element used to supplement the core argument concerning the authority of the 'State' and the justified antagonism of the anarchists towards it. Additionally, towards the end of the discussion, individual autonomy and freedom serve as the focal points. Thus, of the two likely choices, Option A is more appropriate and conceptually closest to the concerns of the passage.

\end{problem}

\newpage
\begin{problem}{}{}
% Subject: varc
\textbf{Instructions}

Instructions   

**The passage below is accompanied by a set of questions. Choose the best answer to each question.**
The word ‘anarchy’ comes from the Greek '_anarkhia_ ', meaning contrary to authority or without a ruler, and was used in a derogatory sense until 1840, when it was adopted by Pierre-Joseph Proudhon to describe his political and social ideology. Proudhon argued that organization without government was both possible and desirable. In the evolution of political ideas, anarchism can be seen as an ultimate projection of both liberalism and socialism, and the differing strands of anarchist thought can be related to their emphasis on one or the other of these.
Historically, anarchism arose not only as an explanation of the gulf between the rich and the poor in any community, and of the reason why the poor have been obliged to fight for their share of a common inheritance, but as a radical answer to the question ‘What went wrong?’ that followed the ultimate outcome of the French Revolution. It had ended not only with a reign of terror and the emergence of a newly rich ruling caste, but with a new adored emperor, Napoleon Bonaparte, strutting through his conquered territories.
The anarchists and their precursors were unique on the political Left in affirming that workers and peasants, grasping the chance that arose to bring an end to centuries of exploitation and tyranny, were inevitably betrayed by the new class of politicians, whose first priority was to re-establish a centralized state power. After every revolutionary uprising, usually won at a heavy cost for ordinary populations, the new rulers had no hesitation in applying violence and terror, a secret police, and a professional army to maintain their control.
For anarchists the state itself is the enemy, and they have applied the same interpretation to the outcome of every revolution of the 19th and 20th centuries. This is not merely because every state keeps a watchful and sometimes punitive eye on its dissidents, but because every state protects the privileges of the powerful.
The mainstream of anarchist propaganda for more than a century has been anarchist- communism, which argues that property in land, natural resources, and the means of production should be held in mutual control by local communities, federating for innumerable joint purposes with other communes. It differs from state socialism in opposing the concept of any central authority. Some anarchists prefer to distinguish between anarchist-communism and collectivist anarchism in order to stress the obviously desirable freedom of an individual or family to possess the resources needed for living, while not implying the right to own the resources needed by others. . . .
There are, unsurprisingly, several traditions of individualist anarchism, one of them deriving from the ‘conscious egoism’ of the German writer Max Stirner (1806-56), and another from a remarkable series of 19th-century American figures who argued that in protecting our own autonomy and associating with others for common advantages, we are promoting the good of all. These thinkers differed from free-market liberals in their absolute mistrust of American capitalism, and in their emphasis on mutualism.

\vspace{1em}
\textbf{Question 5}

According to the passage, what is the one idea that is common to all forms of anarchism?

\begin{enumerate}[label=\Alph*.]
    \item They all derive from the work of Pierre-Joseph Proudhon.
    \item They are all opposed to the centralisation of power in the state.
    \item They all focus on the primacy of the power of the individual.
    \item There is no idea common to all forms of anarchism; that is why it is anarchic.
[](https://cracku.in/5-according-to-the-passage-what-is-the-one-idea-that-x-cat-2020-slot-1)
\end{enumerate}
\vspace{1em}
\textbf{Solution:}

_Option A_ : The various derivations of anarchism cannot be associated with Pierre-Joseph Proudhon. We don't have the requisite information to support this claim.
_Option B_ : This has been presented as the motto of every anarchist faction: wherein the state is the enemy {opposition to the centralization of power}. While there are additional beliefs associated with different schools of anarchist thought, this antagonism towards centralised power forms their core. 
_Option C_ : This belief has been shown to be limited to individualist anarchism {based on the information from the passage}; cannot be stated as a common attribute.
_Option D_ : This assertion would be a gross understatement; the latter half of the option is especially a bit odd {no such idea has been implied}.
Hence, of the given choices, Option B is the correct answer.

Instructions   

**The passage below is accompanied by a set of questions. Choose the best answer to each question.**
Few realise that the government of China, governing an empire of some 60 million people during the Tang dynasty (618-907), implemented a complex financial system that recognised grain, coins and textiles as money. . . . Coins did have certain advantages: they were durable, recognisable and provided a convenient medium of exchange, especially for smaller transactions. However, there were also disadvantages. A continuing shortage of copper meant that government mints could not produce enough coins for the entire empire, to the extent that for most of the dynasty’s history, coins constituted only a tenth of the money supply. One of the main objections to calls for taxes to be paid in coin was that peasant producers who could weave cloth or grow grain - the other two major currencies of the Tang - would not be able to produce coins, and therefore would not be able to pay their taxes. . . .
As coins had advantages and disadvantages, so too did textiles. If in circulation for a long period of time, they could show signs of wear and tear. Stained, faded and torn bolts of textiles had less value than a brand new bolt. Furthermore, a full bolt had a particular value. If consumers cut textiles into smaller pieces to buy or sell something worth less than a full bolt, that, too, greatly lessened the value of the textiles. Unlike coins, textiles could not be used for small transactions; as [an official] noted, textiles could not “be exchanged by the foot and the inch” . . .
But textiles had some advantages over coins. For a start, textile production was widespread and there were fewer problems with the supply of textiles. For large transactions, textiles weighed less than their equivalent in coins since a string of coins . . . could weigh as much as 4 kg. Furthermore, the dimensions of a bolt of silk held remarkably steady from the third to the tenth century: 56 cm wide and 12 m long . . . The values of different textiles were also more stable than the fluctuating values of coins. . . .
The government also required the use of textiles for large transactions. Coins, on the other hand, were better suited for smaller transactions, and possibly, given the costs of transporting coins, for a more local usage. Grain, because it rotted easily, was not used nearly as much as coins and textiles, but taxpayers were required to pay grain to the government as a share of their annual tax obligations, and official salaries were expressed in weights of grain. . . .
In actuality, our own currency system today has some similarities even as it is changing in front of our eyes. . . . We have cash - coins for small transactions like paying for parking at a meter, and banknotes for other items; cheques and debit/credit cards for other, often larger, types of payments. At the same time, we are shifting to electronic banking and making payments online. Some young people never use cash [and] do not know how to write a cheque . . .

\end{problem}

\newpage
\begin{problem}{}{}
% Subject: varc
\textbf{Instructions}

Instructions   

**The passage below is accompanied by a set of questions. Choose the best answer to each question.**
The word ‘anarchy’ comes from the Greek '_anarkhia_ ', meaning contrary to authority or without a ruler, and was used in a derogatory sense until 1840, when it was adopted by Pierre-Joseph Proudhon to describe his political and social ideology. Proudhon argued that organization without government was both possible and desirable. In the evolution of political ideas, anarchism can be seen as an ultimate projection of both liberalism and socialism, and the differing strands of anarchist thought can be related to their emphasis on one or the other of these.
Historically, anarchism arose not only as an explanation of the gulf between the rich and the poor in any community, and of the reason why the poor have been obliged to fight for their share of a common inheritance, but as a radical answer to the question ‘What went wrong?’ that followed the ultimate outcome of the French Revolution. It had ended not only with a reign of terror and the emergence of a newly rich ruling caste, but with a new adored emperor, Napoleon Bonaparte, strutting through his conquered territories.
The anarchists and their precursors were unique on the political Left in affirming that workers and peasants, grasping the chance that arose to bring an end to centuries of exploitation and tyranny, were inevitably betrayed by the new class of politicians, whose first priority was to re-establish a centralized state power. After every revolutionary uprising, usually won at a heavy cost for ordinary populations, the new rulers had no hesitation in applying violence and terror, a secret police, and a professional army to maintain their control.
For anarchists the state itself is the enemy, and they have applied the same interpretation to the outcome of every revolution of the 19th and 20th centuries. This is not merely because every state keeps a watchful and sometimes punitive eye on its dissidents, but because every state protects the privileges of the powerful.
The mainstream of anarchist propaganda for more than a century has been anarchist- communism, which argues that property in land, natural resources, and the means of production should be held in mutual control by local communities, federating for innumerable joint purposes with other communes. It differs from state socialism in opposing the concept of any central authority. Some anarchists prefer to distinguish between anarchist-communism and collectivist anarchism in order to stress the obviously desirable freedom of an individual or family to possess the resources needed for living, while not implying the right to own the resources needed by others. . . .
There are, unsurprisingly, several traditions of individualist anarchism, one of them deriving from the ‘conscious egoism’ of the German writer Max Stirner (1806-56), and another from a remarkable series of 19th-century American figures who argued that in protecting our own autonomy and associating with others for common advantages, we are promoting the good of all. These thinkers differed from free-market liberals in their absolute mistrust of American capitalism, and in their emphasis on mutualism.

\vspace{1em}
\textbf{Question 6}

During the Tang period, which one of the following would not be an economically sound decision for a small purchase in the local market that is worth one-eighth of a bolt of cloth?

\begin{enumerate}[label=\Alph*.]
    \item Making the payment with the appropriate weight of grain.
    \item Paying with a faded bolt of cloth that has approximately the same value.
    \item Using coins issued by the government to make the payment.
    \item Cutting one-eighth of the fabric from a new bolt to pay the amount.
[](https://cracku.in/6-during-the-tang-period-which-one-of-the-following--x-cat-2020-slot-1)
\end{enumerate}
\vspace{1em}
\textbf{Solution:}

Options A and C are economically sound decisions. Option B might be a viable choice, as well. However, Option D would not be an economically sound decision for a small purchase in the local market that is worth one-eighth of a bolt of cloth. Concerning small transactions, the author highlights the limitations that stemmed from a textile form of currency: "..._Furthermore, a full bolt had a particular value. If consumers cut textiles into smaller pieces to buy or sell something worth less than a full bolt, that, too, greatly lessened the value of the textiles_..." Thus, cutting the cloth would greatly diminish its value {and consequently cannot be considered as a sensible decision}. Hence, Option D is the correct answer.

\end{problem}

\newpage
\begin{problem}{}{}
% Subject: varc
\textbf{Instructions}

Instructions   

**The passage below is accompanied by a set of questions. Choose the best answer to each question.**
The word ‘anarchy’ comes from the Greek '_anarkhia_ ', meaning contrary to authority or without a ruler, and was used in a derogatory sense until 1840, when it was adopted by Pierre-Joseph Proudhon to describe his political and social ideology. Proudhon argued that organization without government was both possible and desirable. In the evolution of political ideas, anarchism can be seen as an ultimate projection of both liberalism and socialism, and the differing strands of anarchist thought can be related to their emphasis on one or the other of these.
Historically, anarchism arose not only as an explanation of the gulf between the rich and the poor in any community, and of the reason why the poor have been obliged to fight for their share of a common inheritance, but as a radical answer to the question ‘What went wrong?’ that followed the ultimate outcome of the French Revolution. It had ended not only with a reign of terror and the emergence of a newly rich ruling caste, but with a new adored emperor, Napoleon Bonaparte, strutting through his conquered territories.
The anarchists and their precursors were unique on the political Left in affirming that workers and peasants, grasping the chance that arose to bring an end to centuries of exploitation and tyranny, were inevitably betrayed by the new class of politicians, whose first priority was to re-establish a centralized state power. After every revolutionary uprising, usually won at a heavy cost for ordinary populations, the new rulers had no hesitation in applying violence and terror, a secret police, and a professional army to maintain their control.
For anarchists the state itself is the enemy, and they have applied the same interpretation to the outcome of every revolution of the 19th and 20th centuries. This is not merely because every state keeps a watchful and sometimes punitive eye on its dissidents, but because every state protects the privileges of the powerful.
The mainstream of anarchist propaganda for more than a century has been anarchist- communism, which argues that property in land, natural resources, and the means of production should be held in mutual control by local communities, federating for innumerable joint purposes with other communes. It differs from state socialism in opposing the concept of any central authority. Some anarchists prefer to distinguish between anarchist-communism and collectivist anarchism in order to stress the obviously desirable freedom of an individual or family to possess the resources needed for living, while not implying the right to own the resources needed by others. . . .
There are, unsurprisingly, several traditions of individualist anarchism, one of them deriving from the ‘conscious egoism’ of the German writer Max Stirner (1806-56), and another from a remarkable series of 19th-century American figures who argued that in protecting our own autonomy and associating with others for common advantages, we are promoting the good of all. These thinkers differed from free-market liberals in their absolute mistrust of American capitalism, and in their emphasis on mutualism.

\vspace{1em}
\textbf{Question 7}

When discussing textiles as currency in the Tang period, the author uses the words “steady” and “stable” to indicate all of the following EXCEPT:

\begin{enumerate}[label=\Alph*.]
    \item reliable transportation.
    \item reliable quality.
    \item reliable measurements.
    \item reliable supply.
[](https://cracku.in/7-when-discussing-textiles-as-currency-in-the-tang-p-x-cat-2020-slot-1)
\end{enumerate}
\vspace{1em}
\textbf{Solution:}

The author evidently highlights the stability of utilizing textiles given its abundant supply{".._.textile production was widespread and there were fewer problems with the supply of textiles.._."}, the steady standard of measurement {".._.the dimensions of a bolt of silk held remarkably steady from the third to the tenth century: 56 cm wide and 12 m long._.."} and reliable quality {"._..Furthermore, a full bolt had a particular value. If consumers ..._ "}. Transportation is not a dominant variable that is considered in the discussion pertaining to textiles. Hence, Option A is the correct choice.

\end{problem}

\newpage
\begin{problem}{}{}
% Subject: varc
\textbf{Instructions}

Instructions   

**The passage below is accompanied by a set of questions. Choose the best answer to each question.**
The word ‘anarchy’ comes from the Greek '_anarkhia_ ', meaning contrary to authority or without a ruler, and was used in a derogatory sense until 1840, when it was adopted by Pierre-Joseph Proudhon to describe his political and social ideology. Proudhon argued that organization without government was both possible and desirable. In the evolution of political ideas, anarchism can be seen as an ultimate projection of both liberalism and socialism, and the differing strands of anarchist thought can be related to their emphasis on one or the other of these.
Historically, anarchism arose not only as an explanation of the gulf between the rich and the poor in any community, and of the reason why the poor have been obliged to fight for their share of a common inheritance, but as a radical answer to the question ‘What went wrong?’ that followed the ultimate outcome of the French Revolution. It had ended not only with a reign of terror and the emergence of a newly rich ruling caste, but with a new adored emperor, Napoleon Bonaparte, strutting through his conquered territories.
The anarchists and their precursors were unique on the political Left in affirming that workers and peasants, grasping the chance that arose to bring an end to centuries of exploitation and tyranny, were inevitably betrayed by the new class of politicians, whose first priority was to re-establish a centralized state power. After every revolutionary uprising, usually won at a heavy cost for ordinary populations, the new rulers had no hesitation in applying violence and terror, a secret police, and a professional army to maintain their control.
For anarchists the state itself is the enemy, and they have applied the same interpretation to the outcome of every revolution of the 19th and 20th centuries. This is not merely because every state keeps a watchful and sometimes punitive eye on its dissidents, but because every state protects the privileges of the powerful.
The mainstream of anarchist propaganda for more than a century has been anarchist- communism, which argues that property in land, natural resources, and the means of production should be held in mutual control by local communities, federating for innumerable joint purposes with other communes. It differs from state socialism in opposing the concept of any central authority. Some anarchists prefer to distinguish between anarchist-communism and collectivist anarchism in order to stress the obviously desirable freedom of an individual or family to possess the resources needed for living, while not implying the right to own the resources needed by others. . . .
There are, unsurprisingly, several traditions of individualist anarchism, one of them deriving from the ‘conscious egoism’ of the German writer Max Stirner (1806-56), and another from a remarkable series of 19th-century American figures who argued that in protecting our own autonomy and associating with others for common advantages, we are promoting the good of all. These thinkers differed from free-market liberals in their absolute mistrust of American capitalism, and in their emphasis on mutualism.

\vspace{1em}
\textbf{Question 8}

In the context of the passage, which one of the following can be inferred with regard to the use of currency during the Tang era?

\begin{enumerate}[label=\Alph*.]
    \item Currency that deteriorated easily was not used for official work.
    \item Copper coins were more valuable and durable than textiles.
    \item Grains were the most used currency because of government requirements.
    \item Currency usage was similar to that of modern times.
[](https://cracku.in/8-in-the-context-of-the-passage-which-one-of-the-fol-x-cat-2020-slot-1)
\end{enumerate}
\vspace{1em}
\textbf{Solution:}

Let us inspect the individual statements:
_Option A_ : The author does not make any such claim. The following is stated in this regard: "._..Grain, because it rotted easily, was not used nearly as much as coins and textiles._..". Thus, the marginal usage of grains is presented relative to that of coins and textiles. Therefore, it cannot be understood that perishable currencies such as grains were not utilised for official work. 
_Option B_ : The author simply states that ".._.Coins did have certain advantages: they were durable, recognisable and provided a convenient medium of exchange, especially for smaller transactions.._." Presenting them as being more valuable than textile currency would be incorrect. 
_Option C_ : This would be imprecise as the author portrays the following: ".._.Grain, because it rotted easily, was not used nearly as much as coins and textiles, but taxpayers were required to pay grain to the government as a share of their annual tax obligations, and official salaries were expressed in weights of grain._.." The statement here deviates from this depiction.
_Option D_ : This statement can be understood from the final paragraph, wherein the author states: "._..In actuality, our own currency system today has some similarities even as it is changing in front of our eyes. . . . We have cash - coins for small transactions like paying for parking at a meter, and banknotes for other items; cheques._..". Hence, Option D is a valid inference

\end{problem}

\newpage
\begin{problem}{}{}
% Subject: varc
\textbf{Instructions}

Instructions   

**The passage below is accompanied by a set of questions. Choose the best answer to each question.**
The word ‘anarchy’ comes from the Greek '_anarkhia_ ', meaning contrary to authority or without a ruler, and was used in a derogatory sense until 1840, when it was adopted by Pierre-Joseph Proudhon to describe his political and social ideology. Proudhon argued that organization without government was both possible and desirable. In the evolution of political ideas, anarchism can be seen as an ultimate projection of both liberalism and socialism, and the differing strands of anarchist thought can be related to their emphasis on one or the other of these.
Historically, anarchism arose not only as an explanation of the gulf between the rich and the poor in any community, and of the reason why the poor have been obliged to fight for their share of a common inheritance, but as a radical answer to the question ‘What went wrong?’ that followed the ultimate outcome of the French Revolution. It had ended not only with a reign of terror and the emergence of a newly rich ruling caste, but with a new adored emperor, Napoleon Bonaparte, strutting through his conquered territories.
The anarchists and their precursors were unique on the political Left in affirming that workers and peasants, grasping the chance that arose to bring an end to centuries of exploitation and tyranny, were inevitably betrayed by the new class of politicians, whose first priority was to re-establish a centralized state power. After every revolutionary uprising, usually won at a heavy cost for ordinary populations, the new rulers had no hesitation in applying violence and terror, a secret police, and a professional army to maintain their control.
For anarchists the state itself is the enemy, and they have applied the same interpretation to the outcome of every revolution of the 19th and 20th centuries. This is not merely because every state keeps a watchful and sometimes punitive eye on its dissidents, but because every state protects the privileges of the powerful.
The mainstream of anarchist propaganda for more than a century has been anarchist- communism, which argues that property in land, natural resources, and the means of production should be held in mutual control by local communities, federating for innumerable joint purposes with other communes. It differs from state socialism in opposing the concept of any central authority. Some anarchists prefer to distinguish between anarchist-communism and collectivist anarchism in order to stress the obviously desirable freedom of an individual or family to possess the resources needed for living, while not implying the right to own the resources needed by others. . . .
There are, unsurprisingly, several traditions of individualist anarchism, one of them deriving from the ‘conscious egoism’ of the German writer Max Stirner (1806-56), and another from a remarkable series of 19th-century American figures who argued that in protecting our own autonomy and associating with others for common advantages, we are promoting the good of all. These thinkers differed from free-market liberals in their absolute mistrust of American capitalism, and in their emphasis on mutualism.

\vspace{1em}
\textbf{Question 9}

According to the passage, the modern currency system shares all the following features with that of the Tang, EXCEPT that:

\begin{enumerate}[label=\Alph*.]
    \item it uses different materials as currency.
    \item it uses different currencies for different situations.
    \item its currencies fluctuate in value over time.
    \item it is undergoing transformation.
[](https://cracku.in/9-according-to-the-passage-the-modern-currency-syste-x-cat-2020-slot-1)
\end{enumerate}
\vspace{1em}
\textbf{Solution:}

Let us pay heed to the following excerpt: "_...In actuality, our own currency system today has some similarities even as it is changing in front of our eyes. . . . We have cash - coins for small transactions like paying for parking at a meter, and banknotes for other items; cheques and debit/credit cards for other, often larger, types of payments. At the same time, we are shifting to electronic banking and making payments online. Some young people never use cash [and] do not know how to write a cheque_..." 
Based on the above, we can infer Options A, B and C as the similarities that are implied. Concerning Option D, while the author does mention the present currency system as a changing/dynamic one, the same cannot be said about the Tang dynasty. The author does not highlight this element as a similarity. Hence, Option D is the correct answer.

Instructions   

**The passage below is accompanied by a set of questions. Choose the best answer to each question.**
Vocabulary used in speech or writing organizes itself in seven parts of speech (eight, if you count interjections such as Oh! and Gosh! and Fuhgeddaboudit!). Communication composed of these parts of speech must be organized by rules of grammar upon which we agree. When these rules break down, confusion and misunderstanding result. Bad grammar produces bad sentences. My favorite example from Strunk and White is this one: “As a mother of five, with another one on the way, my ironing board is always up.”
Nouns and verbs are the two indispensable parts of writing. Without one of each, no group of words can be a sentence, since a sentence is, by definition, a group of words containing a subject (noun) and a predicate (verb); these strings of words begin with a capital letter, end with a period, and combine to make a complete thought which starts in the writer’s head and then leaps to the reader’s.
Must you write complete sentences each time, every time? Perish the thought. If your work consists only of fragments and floating clauses, the Grammar Police aren’t going to come and take you away. Even William Strunk, that Mussolini of rhetoric, recognized the delicious pliability of language. “It is an old observation,” he writes, “that the best writers sometimes disregard the rules of rhetoric.” Yet he goes on to add this thought, which I urge you to consider: “Unless he is certain of doing well, [the writer] will probably do best to follow the rules.”
The telling clause here is Unless he is certain of doing well. If you don’t have a rudimentary grasp of how the parts of speech translate into coherent sentences, how can you be certain that you are doing well? How will you know if you’re doing ill, for that matter? The answer, of course, is that you can’t, you won’t. One who does grasp the rudiments of grammar finds a comforting simplicity at its heart, where there need be only nouns, the words that name, and verbs, the words that act.
Take any noun, put it with any verb, and you have a sentence. It never fails. Rocks explode. Jane transmits. Mountains float. These are all perfect sentences. Many such thoughts make little rational sense, but even the stranger ones (Plums deify!) have a kind of poetic weight that’s nice. The simplicity of noun-verb construction is useful—at the very least it can provide a safety net for your writing. Strunk and White caution against too many simple sentences in a row, but simple sentences provide a path you can follow when you fear getting lost in the tangles of rhetoric—all those restrictive and nonrestrictive clauses, those modifying phrases, those appositives and compound-complex sentences. If you start to freak out at the sight of such unmapped territory (unmapped by you, at least), just remind yourself that rocks explode, Jane transmits, mountains float, and plums deify. Grammar is . . . the pole you grab to get your thoughts up on their feet and walking.

\end{problem}

\newpage
\begin{problem}{}{}
% Subject: varc
\textbf{Instructions}

Instructions   

**The passage below is accompanied by a set of questions. Choose the best answer to each question.**
The word ‘anarchy’ comes from the Greek '_anarkhia_ ', meaning contrary to authority or without a ruler, and was used in a derogatory sense until 1840, when it was adopted by Pierre-Joseph Proudhon to describe his political and social ideology. Proudhon argued that organization without government was both possible and desirable. In the evolution of political ideas, anarchism can be seen as an ultimate projection of both liberalism and socialism, and the differing strands of anarchist thought can be related to their emphasis on one or the other of these.
Historically, anarchism arose not only as an explanation of the gulf between the rich and the poor in any community, and of the reason why the poor have been obliged to fight for their share of a common inheritance, but as a radical answer to the question ‘What went wrong?’ that followed the ultimate outcome of the French Revolution. It had ended not only with a reign of terror and the emergence of a newly rich ruling caste, but with a new adored emperor, Napoleon Bonaparte, strutting through his conquered territories.
The anarchists and their precursors were unique on the political Left in affirming that workers and peasants, grasping the chance that arose to bring an end to centuries of exploitation and tyranny, were inevitably betrayed by the new class of politicians, whose first priority was to re-establish a centralized state power. After every revolutionary uprising, usually won at a heavy cost for ordinary populations, the new rulers had no hesitation in applying violence and terror, a secret police, and a professional army to maintain their control.
For anarchists the state itself is the enemy, and they have applied the same interpretation to the outcome of every revolution of the 19th and 20th centuries. This is not merely because every state keeps a watchful and sometimes punitive eye on its dissidents, but because every state protects the privileges of the powerful.
The mainstream of anarchist propaganda for more than a century has been anarchist- communism, which argues that property in land, natural resources, and the means of production should be held in mutual control by local communities, federating for innumerable joint purposes with other communes. It differs from state socialism in opposing the concept of any central authority. Some anarchists prefer to distinguish between anarchist-communism and collectivist anarchism in order to stress the obviously desirable freedom of an individual or family to possess the resources needed for living, while not implying the right to own the resources needed by others. . . .
There are, unsurprisingly, several traditions of individualist anarchism, one of them deriving from the ‘conscious egoism’ of the German writer Max Stirner (1806-56), and another from a remarkable series of 19th-century American figures who argued that in protecting our own autonomy and associating with others for common advantages, we are promoting the good of all. These thinkers differed from free-market liberals in their absolute mistrust of American capitalism, and in their emphasis on mutualism.

\vspace{1em}
\textbf{Question 10}

Which one of the following statements, if false, could be seen as supporting the arguments in the passage?

\begin{enumerate}[label=\Alph*.]
    \item An understanding of grammar helps a writer decide if she/he is writing well or not.
    \item Regarding grammar, women writers tend to be more attentive to method and  
accuracy.
    \item It has been observed that writers sometimes disregard the rules of rhetoric.
    \item Perish the thought that complete sentences necessarily need nouns and verbs!
[](https://cracku.in/10-which-one-of-the-following-statements-if-false-cou-x-cat-2020-slot-1)
\end{enumerate}
\vspace{1em}
\textbf{Solution:}

We need to find a statement, which if **false** , aligns with the discussion/serves as a supporting argument. Let us inspect the individual options:
_Option A_ : In its current form, this statement is in tune with the author's assertion. However, if false, it is antagonistic to the claim being made in the passage. Hence, we can eliminate this option.
_Options B and C_ : Regardless of whether these statements are true or false, they do very little to further the idea presented in the passage.
_Option D_ : This statement implies that complete sentences do not need nouns and verbs. However, in the passage, the author says otherwise; thus, if this sentence is false, it perfectly aligns with the argument made in the passage. Therefore, Option D, if false, could be seen as a supplementary argument.

\end{problem}

\newpage
\begin{problem}{}{}
% Subject: varc
\textbf{Instructions}

Instructions   

**The passage below is accompanied by a set of questions. Choose the best answer to each question.**
The word ‘anarchy’ comes from the Greek '_anarkhia_ ', meaning contrary to authority or without a ruler, and was used in a derogatory sense until 1840, when it was adopted by Pierre-Joseph Proudhon to describe his political and social ideology. Proudhon argued that organization without government was both possible and desirable. In the evolution of political ideas, anarchism can be seen as an ultimate projection of both liberalism and socialism, and the differing strands of anarchist thought can be related to their emphasis on one or the other of these.
Historically, anarchism arose not only as an explanation of the gulf between the rich and the poor in any community, and of the reason why the poor have been obliged to fight for their share of a common inheritance, but as a radical answer to the question ‘What went wrong?’ that followed the ultimate outcome of the French Revolution. It had ended not only with a reign of terror and the emergence of a newly rich ruling caste, but with a new adored emperor, Napoleon Bonaparte, strutting through his conquered territories.
The anarchists and their precursors were unique on the political Left in affirming that workers and peasants, grasping the chance that arose to bring an end to centuries of exploitation and tyranny, were inevitably betrayed by the new class of politicians, whose first priority was to re-establish a centralized state power. After every revolutionary uprising, usually won at a heavy cost for ordinary populations, the new rulers had no hesitation in applying violence and terror, a secret police, and a professional army to maintain their control.
For anarchists the state itself is the enemy, and they have applied the same interpretation to the outcome of every revolution of the 19th and 20th centuries. This is not merely because every state keeps a watchful and sometimes punitive eye on its dissidents, but because every state protects the privileges of the powerful.
The mainstream of anarchist propaganda for more than a century has been anarchist- communism, which argues that property in land, natural resources, and the means of production should be held in mutual control by local communities, federating for innumerable joint purposes with other communes. It differs from state socialism in opposing the concept of any central authority. Some anarchists prefer to distinguish between anarchist-communism and collectivist anarchism in order to stress the obviously desirable freedom of an individual or family to possess the resources needed for living, while not implying the right to own the resources needed by others. . . .
There are, unsurprisingly, several traditions of individualist anarchism, one of them deriving from the ‘conscious egoism’ of the German writer Max Stirner (1806-56), and another from a remarkable series of 19th-century American figures who argued that in protecting our own autonomy and associating with others for common advantages, we are promoting the good of all. These thinkers differed from free-market liberals in their absolute mistrust of American capitalism, and in their emphasis on mutualism.

\vspace{1em}
\textbf{Question 11}

“Take any noun, put it with any verb, and you have a sentence. It never fails. Rocks explode. Jane transmits. Mountains float.” None of the following statements can be seen as similar EXCEPT:

\begin{enumerate}[label=\Alph*.]
    \item Take any vegetable, put some spices in it, and you have a dish
    \item Take an apple tree, plant it in a field, and you have an orchard.
    \item A group of nouns arranged in a row becomes a sentence.
    \item A collection of people with the same sports equipment is a sports team.
[](https://cracku.in/11-take-any-noun-put-it-with-any-verb-and-you-have-a--x-cat-2020-slot-1)
\end{enumerate}
\vspace{1em}
\textbf{Solution:}

The author intends to highlight that the rudimentary combination of a noun and a verb serves as the simplest form of expression; two basic yet immensely significant entities coupled together that is representative of a broader and perhaps, complex group of entities {a sentence}. Option A is closest to such a relationship: vegetables and spices (two elements) combined to represent a larger group - 'dishes'. The same cannot be said about the rest of the options, since they evidently deviate from the core message being conveyed.

\end{problem}

\newpage
\begin{problem}{}{}
% Subject: varc
\textbf{Instructions}

Instructions   

**The passage below is accompanied by a set of questions. Choose the best answer to each question.**
The word ‘anarchy’ comes from the Greek '_anarkhia_ ', meaning contrary to authority or without a ruler, and was used in a derogatory sense until 1840, when it was adopted by Pierre-Joseph Proudhon to describe his political and social ideology. Proudhon argued that organization without government was both possible and desirable. In the evolution of political ideas, anarchism can be seen as an ultimate projection of both liberalism and socialism, and the differing strands of anarchist thought can be related to their emphasis on one or the other of these.
Historically, anarchism arose not only as an explanation of the gulf between the rich and the poor in any community, and of the reason why the poor have been obliged to fight for their share of a common inheritance, but as a radical answer to the question ‘What went wrong?’ that followed the ultimate outcome of the French Revolution. It had ended not only with a reign of terror and the emergence of a newly rich ruling caste, but with a new adored emperor, Napoleon Bonaparte, strutting through his conquered territories.
The anarchists and their precursors were unique on the political Left in affirming that workers and peasants, grasping the chance that arose to bring an end to centuries of exploitation and tyranny, were inevitably betrayed by the new class of politicians, whose first priority was to re-establish a centralized state power. After every revolutionary uprising, usually won at a heavy cost for ordinary populations, the new rulers had no hesitation in applying violence and terror, a secret police, and a professional army to maintain their control.
For anarchists the state itself is the enemy, and they have applied the same interpretation to the outcome of every revolution of the 19th and 20th centuries. This is not merely because every state keeps a watchful and sometimes punitive eye on its dissidents, but because every state protects the privileges of the powerful.
The mainstream of anarchist propaganda for more than a century has been anarchist- communism, which argues that property in land, natural resources, and the means of production should be held in mutual control by local communities, federating for innumerable joint purposes with other communes. It differs from state socialism in opposing the concept of any central authority. Some anarchists prefer to distinguish between anarchist-communism and collectivist anarchism in order to stress the obviously desirable freedom of an individual or family to possess the resources needed for living, while not implying the right to own the resources needed by others. . . .
There are, unsurprisingly, several traditions of individualist anarchism, one of them deriving from the ‘conscious egoism’ of the German writer Max Stirner (1806-56), and another from a remarkable series of 19th-century American figures who argued that in protecting our own autonomy and associating with others for common advantages, we are promoting the good of all. These thinkers differed from free-market liberals in their absolute mistrust of American capitalism, and in their emphasis on mutualism.

\vspace{1em}
\textbf{Question 12}

Which one of the following quotes best captures the main concern of the passage?

\begin{enumerate}[label=\Alph*.]
    \item “The telling clause here is Unless he is certain of doing well.”
    \item “Bad grammar produces bad sentences.”
    \item “Strunk and White caution against too many simple sentences in a row, but simple sentences provide a path you can follow when you fear getting lost in the tangles of rhetoric . . .”
    \item “Nouns and verbs are the two indispensable parts of writing. Without one of each, no group of words can be a sentence . . .”
[](https://cracku.in/12-which-one-of-the-following-quotes-best-captures-th-x-cat-2020-slot-1)
\end{enumerate}
\vspace{1em}
\textbf{Solution:}

The author begins by highlighting the necessity for a set of codes (enabled by grammar) to organize communication and avoid confusion. He then proceeds to present supplementary arguments in this regard (elements associated with rhetoric and its specialists) and emphasizes how even proper, intentional simplification can be attained only through a firm grasp of the rudiments of grammar. Therefore, it can be observed that grammar is the focal point here and the correct choice should definitely align with this. Option B aptly captures the main concern raised in the passage. Options A and C fail to include the idea revolving around grammar and instead focuses on the additives. Option D is close; however, the mention of nouns and verbs is with the intention to supplement the idea highlighted in B. These simply serve as an illustration to emphasize the significance of grammar. Thus, between the two options B and D, Option B is the suitable choice.

\end{problem}

\newpage
\begin{problem}{}{}
% Subject: varc
\textbf{Instructions}

Instructions   

**The passage below is accompanied by a set of questions. Choose the best answer to each question.**
The word ‘anarchy’ comes from the Greek '_anarkhia_ ', meaning contrary to authority or without a ruler, and was used in a derogatory sense until 1840, when it was adopted by Pierre-Joseph Proudhon to describe his political and social ideology. Proudhon argued that organization without government was both possible and desirable. In the evolution of political ideas, anarchism can be seen as an ultimate projection of both liberalism and socialism, and the differing strands of anarchist thought can be related to their emphasis on one or the other of these.
Historically, anarchism arose not only as an explanation of the gulf between the rich and the poor in any community, and of the reason why the poor have been obliged to fight for their share of a common inheritance, but as a radical answer to the question ‘What went wrong?’ that followed the ultimate outcome of the French Revolution. It had ended not only with a reign of terror and the emergence of a newly rich ruling caste, but with a new adored emperor, Napoleon Bonaparte, strutting through his conquered territories.
The anarchists and their precursors were unique on the political Left in affirming that workers and peasants, grasping the chance that arose to bring an end to centuries of exploitation and tyranny, were inevitably betrayed by the new class of politicians, whose first priority was to re-establish a centralized state power. After every revolutionary uprising, usually won at a heavy cost for ordinary populations, the new rulers had no hesitation in applying violence and terror, a secret police, and a professional army to maintain their control.
For anarchists the state itself is the enemy, and they have applied the same interpretation to the outcome of every revolution of the 19th and 20th centuries. This is not merely because every state keeps a watchful and sometimes punitive eye on its dissidents, but because every state protects the privileges of the powerful.
The mainstream of anarchist propaganda for more than a century has been anarchist- communism, which argues that property in land, natural resources, and the means of production should be held in mutual control by local communities, federating for innumerable joint purposes with other communes. It differs from state socialism in opposing the concept of any central authority. Some anarchists prefer to distinguish between anarchist-communism and collectivist anarchism in order to stress the obviously desirable freedom of an individual or family to possess the resources needed for living, while not implying the right to own the resources needed by others. . . .
There are, unsurprisingly, several traditions of individualist anarchism, one of them deriving from the ‘conscious egoism’ of the German writer Max Stirner (1806-56), and another from a remarkable series of 19th-century American figures who argued that in protecting our own autonomy and associating with others for common advantages, we are promoting the good of all. These thinkers differed from free-market liberals in their absolute mistrust of American capitalism, and in their emphasis on mutualism.

\vspace{1em}
\textbf{Question 13}

All of the following statements can be inferred from the passage EXCEPT that:

\begin{enumerate}[label=\Alph*.]
    \item sentences do not always have to be complete.
    \item the subject-predicate relation is the same as the noun-verb relation.
    \item the primary purpose of grammar is to ensure that sentences remain simple.
    \item “Grammar Police” is a metaphor for critics who focus on linguistic rules.
[](https://cracku.in/13-all-of-the-following-statements-can-be-inferred-fr-x-cat-2020-slot-1)
\end{enumerate}
\vspace{1em}
\textbf{Solution:}

_Option A_ : This can be inferred from "..._Must you write complete sentences each time, every time? Perish the thought. If your work consists only of fragments and floating clauses, the Grammar Police aren’t going to come and take you away._.." The author presents an illustration to show how a simple combination of a noun and verb forms a complete expression.
_Option B_ : Based on the limited information available from the passage, we can make this inference from "._..no group of words can be a sentence, since a sentence is, by definition, a group of words containing a subject (noun) and a predicate (verb)._.."
_Option C_ : The author does not make any such claim. Grammar serves as a mechanism to organize communication and avoid confusion. However, the "primary purpose" of grammar is not to ensure that sentences remain "simple". Since we cannot infer this statement from the passage, it is the correct answer.
_Option D_ : Although not explicitly mentioned, we can understand the sentiment behind using the term "Grammar Police". The author is using this as a metaphor for the strong adherents of the grammatic rules (who are perhaps swift to judge and criticize).
Therefore, we can infer all of the statements except for the one in Option C.

\end{problem}

\newpage
\begin{problem}{}{}
% Subject: varc
\textbf{Instructions}

Instructions   

**The passage below is accompanied by a set of questions. Choose the best answer to each question.**
The word ‘anarchy’ comes from the Greek '_anarkhia_ ', meaning contrary to authority or without a ruler, and was used in a derogatory sense until 1840, when it was adopted by Pierre-Joseph Proudhon to describe his political and social ideology. Proudhon argued that organization without government was both possible and desirable. In the evolution of political ideas, anarchism can be seen as an ultimate projection of both liberalism and socialism, and the differing strands of anarchist thought can be related to their emphasis on one or the other of these.
Historically, anarchism arose not only as an explanation of the gulf between the rich and the poor in any community, and of the reason why the poor have been obliged to fight for their share of a common inheritance, but as a radical answer to the question ‘What went wrong?’ that followed the ultimate outcome of the French Revolution. It had ended not only with a reign of terror and the emergence of a newly rich ruling caste, but with a new adored emperor, Napoleon Bonaparte, strutting through his conquered territories.
The anarchists and their precursors were unique on the political Left in affirming that workers and peasants, grasping the chance that arose to bring an end to centuries of exploitation and tyranny, were inevitably betrayed by the new class of politicians, whose first priority was to re-establish a centralized state power. After every revolutionary uprising, usually won at a heavy cost for ordinary populations, the new rulers had no hesitation in applying violence and terror, a secret police, and a professional army to maintain their control.
For anarchists the state itself is the enemy, and they have applied the same interpretation to the outcome of every revolution of the 19th and 20th centuries. This is not merely because every state keeps a watchful and sometimes punitive eye on its dissidents, but because every state protects the privileges of the powerful.
The mainstream of anarchist propaganda for more than a century has been anarchist- communism, which argues that property in land, natural resources, and the means of production should be held in mutual control by local communities, federating for innumerable joint purposes with other communes. It differs from state socialism in opposing the concept of any central authority. Some anarchists prefer to distinguish between anarchist-communism and collectivist anarchism in order to stress the obviously desirable freedom of an individual or family to possess the resources needed for living, while not implying the right to own the resources needed by others. . . .
There are, unsurprisingly, several traditions of individualist anarchism, one of them deriving from the ‘conscious egoism’ of the German writer Max Stirner (1806-56), and another from a remarkable series of 19th-century American figures who argued that in protecting our own autonomy and associating with others for common advantages, we are promoting the good of all. These thinkers differed from free-market liberals in their absolute mistrust of American capitalism, and in their emphasis on mutualism.

\vspace{1em}
\textbf{Question 14}

Inferring from the passage, the author could be most supportive of which one of the following practices?

\begin{enumerate}[label=\Alph*.]
    \item The availability of language software that will standardise the rules of grammar as an aid to writers.
    \item A campaign demanding that a writer’s creative license should allow the breaking of grammatical rules.
    \item A Creative Writing course that focuses on how to avoid the use of rhetoric.
    \item The critique of standardised rules of punctuation and capitalisation.
[](https://cracku.in/14-inferring-from-the-passage-the-author-could-be-mos-x-cat-2020-slot-1)
\end{enumerate}
\vspace{1em}
\textbf{Solution:}

In this question, we need to identify the statement that coincides with the aspects discussed by the author. Putting ourselves in the author's shoes, we know that grammar (and the necessity to appropriately learn it) is the primary idea that needs to be conveyed. Writers with a substantial understanding of the elementary rules in grammar can appreciate the "_comforting simplicity at its heart_ " (is what the author claims). Thus, the author will surely support any stance that concurs with the above.
_Option A_ : This will definitely supplement the assertion made by the author. It will enable writers with the requisite understanding of the standard governing rules - a significant necessity highlighted by the author. He is bound to favor such an action. 
_Option B_ : The author would endorse such drastic measures (does not match the tone). The pliability of grammatical rules is a noticeable comment made by the author (but only for those well-versed with it). Hence, we can eliminate this option.
_Option C_ : The author does not portray any view that promotes the eschewal of rhetoric. This again deviates from the discussion in the passage and can be scrapped as the correct choice.
_Option D_ : The focus is broad: on grammar (not just on punctuation and capitalization). 
It is evident that Option A is the only sensible statement that the author would support here.

Instructions   

**The passage below is accompanied by a set of questions. Choose the best answer to each question.**
In the late 1960s, while studying the northern-elephant-seal population along the coasts of Mexico and California, Burney Le Boeuf and his colleagues couldn’t help but notice that the threat calls of males at some sites sounded different from those of males at other sites. . . . That was the first time dialects were documented in a nonhuman mammal. . . .
All the northern elephant seals that exist today are descendants of the small herd that survived on Isla Guadalupe [after the near extinction of the species in the nineteenth century]. As that tiny population grew, northern elephant seals started to recolonize former breeding locations. It was precisely on the more recently colonized islands where Le Boeuf found that the tempos of the male vocal displays showed stronger differences to the ones from Isla Guadalupe, the founder colony.
In order to test the reliability of these dialects over time, Le Boeuf and other researchers visited Año Nuevo Island in California—the island where males showed the slowest pulse rates in their calls—every winter from 1968 to 1972. “What we found is that the pulse rate increased, but it still remained relatively slow compared to the other colonies we had measured in the past” Le Boeuf told me.
At the individual level, the pulse of the calls stayed the same: A male would maintain his vocal signature throughout his lifetime. But the average pulse rate was changing. Immigration could have been responsible for this increase, as in the early 1970s, 43 percent of the males on Año Nuevo had come from southern rookeries that had a faster pulse rate. This led Le Boeuf and his collaborator, Lewis Petrinovich, to deduce that the dialects were, perhaps, a result of isolation over time, after the breeding sites had been recolonized. For instance, the first settlers of Año Nuevo could have had, by chance, calls with low pulse rates. At other sites, where the scientists found faster pulse rates, the opposite would have happened—seals with faster rates would have happened to arrive first.
As the population continued to expand and the islands kept on receiving immigrants from the original population, the calls in all locations would have eventually regressed to the average pulse rate of the founder colony. In the decades that followed, scientists noticed that the geographical variations reported in 1969 were not obvious anymore. . . . In the early 2010s, while studying northern elephant seals on Año Nuevo Island, [researcher Caroline] Casey noticed, too, that what Le Boeuf had heard decades ago was not what she heard now. . . . By performing more sophisticated statistical analyses on both sets of data, [Casey and Le Boeuf] confirmed that dialects existed back then but had vanished. Yet there are other differences between the males from the late 1960s and their great-great-grandsons: Modern males exhibit more individual diversity, and their calls are more complex. While 50 years ago the drumming pattern was quite simple and the dialects denoted just a change in tempo, Casey explained, the calls recorded today have more complex structures, sometimes featuring doublets or triplets. . . .

\end{problem}

\newpage
\begin{problem}{}{}
% Subject: varc
\textbf{Instructions}

Instructions   

**The passage below is accompanied by a set of questions. Choose the best answer to each question.**
The word ‘anarchy’ comes from the Greek '_anarkhia_ ', meaning contrary to authority or without a ruler, and was used in a derogatory sense until 1840, when it was adopted by Pierre-Joseph Proudhon to describe his political and social ideology. Proudhon argued that organization without government was both possible and desirable. In the evolution of political ideas, anarchism can be seen as an ultimate projection of both liberalism and socialism, and the differing strands of anarchist thought can be related to their emphasis on one or the other of these.
Historically, anarchism arose not only as an explanation of the gulf between the rich and the poor in any community, and of the reason why the poor have been obliged to fight for their share of a common inheritance, but as a radical answer to the question ‘What went wrong?’ that followed the ultimate outcome of the French Revolution. It had ended not only with a reign of terror and the emergence of a newly rich ruling caste, but with a new adored emperor, Napoleon Bonaparte, strutting through his conquered territories.
The anarchists and their precursors were unique on the political Left in affirming that workers and peasants, grasping the chance that arose to bring an end to centuries of exploitation and tyranny, were inevitably betrayed by the new class of politicians, whose first priority was to re-establish a centralized state power. After every revolutionary uprising, usually won at a heavy cost for ordinary populations, the new rulers had no hesitation in applying violence and terror, a secret police, and a professional army to maintain their control.
For anarchists the state itself is the enemy, and they have applied the same interpretation to the outcome of every revolution of the 19th and 20th centuries. This is not merely because every state keeps a watchful and sometimes punitive eye on its dissidents, but because every state protects the privileges of the powerful.
The mainstream of anarchist propaganda for more than a century has been anarchist- communism, which argues that property in land, natural resources, and the means of production should be held in mutual control by local communities, federating for innumerable joint purposes with other communes. It differs from state socialism in opposing the concept of any central authority. Some anarchists prefer to distinguish between anarchist-communism and collectivist anarchism in order to stress the obviously desirable freedom of an individual or family to possess the resources needed for living, while not implying the right to own the resources needed by others. . . .
There are, unsurprisingly, several traditions of individualist anarchism, one of them deriving from the ‘conscious egoism’ of the German writer Max Stirner (1806-56), and another from a remarkable series of 19th-century American figures who argued that in protecting our own autonomy and associating with others for common advantages, we are promoting the good of all. These thinkers differed from free-market liberals in their absolute mistrust of American capitalism, and in their emphasis on mutualism.

\vspace{1em}
\textbf{Question 15}

All of the following can be inferred from Le Boeuf’s study as described in the passage EXCEPT that:

\begin{enumerate}[label=\Alph*.]
    \item male northern elephant seals might not have exhibited dialects had they not become nearly extinct in the nineteenth century.
    \item changes in population and migration had no effect on the call pulse rate of individual male northern elephant seals.
    \item the influx of new northern elephant seals into Año Nuevo Island would have soon made the call pulse rate of its male seals exceed that of those at Isla Guadalupe.
    \item the average call pulse rate of male northern elephant seals at Año Nuevo Island increased from the early 1970s till the disappearance of dialects.
[](https://cracku.in/15-all-of-the-following-can-be-inferred-from-le-boeuf-x-cat-2020-slot-1)
\end{enumerate}
\vspace{1em}
\textbf{Solution:}

Let us inspect the individual statements:
_Option A_ : This can be inferred from the second paragraph [".._.It was precisely on the more recently colonized islands where Le Boeuf found that the tempos of the male vocal displays showed stronger differences to the ones from Isla Guadalupe, the founder colony..._ "]. The inception of the eventual dialects has been indirectly attributed to the dynamic changes that occurred as a result of the near extinction of the elephant seals. Hence, we can infer Option A from the passage.
_Option B_ : This can be inferred from the fourth paragraph: ["._..At the individual level, the pulse of the calls stayed the same: A male would maintain his vocal signature throughout his lifetime._.."] The changing variables have little to no effect of the individual vocal signature of the elephant seals. Thus, Option B can be inferred from the passage.
_Option C_ : No such claim is made in the passage. 
_Option D_ : This can be inferred from [".._.could have been responsible for this increase, as in the early 1970s, 43 per cent of the males on Año Nuevo had come from southern rookeries that had a faster pulse rate._.."] and the discussion at the beginning of the final paragraph: [".._.As the population continued to expand and the islands kept on receiving immigrants from the original population, the calls in all locations would have eventually regressed to the average pulse rate of the founder colony.._.]. 
Therefore, of the given statements, we cannot infer Option C.

\end{problem}

\newpage
\begin{problem}{}{}
% Subject: varc
\textbf{Instructions}

Instructions   

**The passage below is accompanied by a set of questions. Choose the best answer to each question.**
The word ‘anarchy’ comes from the Greek '_anarkhia_ ', meaning contrary to authority or without a ruler, and was used in a derogatory sense until 1840, when it was adopted by Pierre-Joseph Proudhon to describe his political and social ideology. Proudhon argued that organization without government was both possible and desirable. In the evolution of political ideas, anarchism can be seen as an ultimate projection of both liberalism and socialism, and the differing strands of anarchist thought can be related to their emphasis on one or the other of these.
Historically, anarchism arose not only as an explanation of the gulf between the rich and the poor in any community, and of the reason why the poor have been obliged to fight for their share of a common inheritance, but as a radical answer to the question ‘What went wrong?’ that followed the ultimate outcome of the French Revolution. It had ended not only with a reign of terror and the emergence of a newly rich ruling caste, but with a new adored emperor, Napoleon Bonaparte, strutting through his conquered territories.
The anarchists and their precursors were unique on the political Left in affirming that workers and peasants, grasping the chance that arose to bring an end to centuries of exploitation and tyranny, were inevitably betrayed by the new class of politicians, whose first priority was to re-establish a centralized state power. After every revolutionary uprising, usually won at a heavy cost for ordinary populations, the new rulers had no hesitation in applying violence and terror, a secret police, and a professional army to maintain their control.
For anarchists the state itself is the enemy, and they have applied the same interpretation to the outcome of every revolution of the 19th and 20th centuries. This is not merely because every state keeps a watchful and sometimes punitive eye on its dissidents, but because every state protects the privileges of the powerful.
The mainstream of anarchist propaganda for more than a century has been anarchist- communism, which argues that property in land, natural resources, and the means of production should be held in mutual control by local communities, federating for innumerable joint purposes with other communes. It differs from state socialism in opposing the concept of any central authority. Some anarchists prefer to distinguish between anarchist-communism and collectivist anarchism in order to stress the obviously desirable freedom of an individual or family to possess the resources needed for living, while not implying the right to own the resources needed by others. . . .
There are, unsurprisingly, several traditions of individualist anarchism, one of them deriving from the ‘conscious egoism’ of the German writer Max Stirner (1806-56), and another from a remarkable series of 19th-century American figures who argued that in protecting our own autonomy and associating with others for common advantages, we are promoting the good of all. These thinkers differed from free-market liberals in their absolute mistrust of American capitalism, and in their emphasis on mutualism.

\vspace{1em}
\textbf{Question 16}

Which one of the following conditions, if true, could have ensured that male northern elephant seal dialects did not disappear?

\begin{enumerate}[label=\Alph*.]
    \item Besides Isla Guadalupe, there was one more surviving colony with the same average male call tempo from which no migration took place.
    \item The call tempo of individual male seals in host colonies changed to match the average call tempo of immigrant male seals.
    \item Besides Isla Guadalupe, there was one more founder colony with the same average male call tempo from which male seals migrated to various other colonies.
    \item The call tempo of individual immigrant male seals changed to match the average tempo of resident male seals in the host colony.
[](https://cracku.in/16-which-one-of-the-following-conditions-if-true-coul-x-cat-2020-slot-1)
\end{enumerate}
\vspace{1em}
\textbf{Solution:}

A noticeable clue to ensure that the male northern elephant seal dialects did not disappear is presented in the penultimate paragraph: [".._.At the individual level, the pulse of the calls stayed the same: A male would maintain his vocal signature throughout his lifetime. But the average pulse rate was changing. Immigration could have been responsible for this increase, as in the early 1970s, 43 per cent of the males on Año Nuevo had come from southern rookeries that had a faster pulse rate._.. "]. The loss in the dialect is due to the influx of seals with a faster pulse rate and "as the population continued to expand and the islands kept on receiving immigrants from the original population, the calls in all locations would have eventually regressed to the average pulse rate of the founder colony." Thus, if the individual pulse rate of the immigrants varies or adapts to the existing population, this could preserve the dialect in a particular region. The statement in Option D reflects this specific idea and helps sustain the existing dialect in a given population. Options A and C do little to contribute to the cause of preventing the disappearance of the dialects. Option B aligns with the discussion in the passage and is responsible for the regression of the dialects. Hence, Option D is the correct answer.

\end{problem}

\newpage
\begin{problem}{}{}
% Subject: varc
\textbf{Instructions}

Instructions   

**The passage below is accompanied by a set of questions. Choose the best answer to each question.**
The word ‘anarchy’ comes from the Greek '_anarkhia_ ', meaning contrary to authority or without a ruler, and was used in a derogatory sense until 1840, when it was adopted by Pierre-Joseph Proudhon to describe his political and social ideology. Proudhon argued that organization without government was both possible and desirable. In the evolution of political ideas, anarchism can be seen as an ultimate projection of both liberalism and socialism, and the differing strands of anarchist thought can be related to their emphasis on one or the other of these.
Historically, anarchism arose not only as an explanation of the gulf between the rich and the poor in any community, and of the reason why the poor have been obliged to fight for their share of a common inheritance, but as a radical answer to the question ‘What went wrong?’ that followed the ultimate outcome of the French Revolution. It had ended not only with a reign of terror and the emergence of a newly rich ruling caste, but with a new adored emperor, Napoleon Bonaparte, strutting through his conquered territories.
The anarchists and their precursors were unique on the political Left in affirming that workers and peasants, grasping the chance that arose to bring an end to centuries of exploitation and tyranny, were inevitably betrayed by the new class of politicians, whose first priority was to re-establish a centralized state power. After every revolutionary uprising, usually won at a heavy cost for ordinary populations, the new rulers had no hesitation in applying violence and terror, a secret police, and a professional army to maintain their control.
For anarchists the state itself is the enemy, and they have applied the same interpretation to the outcome of every revolution of the 19th and 20th centuries. This is not merely because every state keeps a watchful and sometimes punitive eye on its dissidents, but because every state protects the privileges of the powerful.
The mainstream of anarchist propaganda for more than a century has been anarchist- communism, which argues that property in land, natural resources, and the means of production should be held in mutual control by local communities, federating for innumerable joint purposes with other communes. It differs from state socialism in opposing the concept of any central authority. Some anarchists prefer to distinguish between anarchist-communism and collectivist anarchism in order to stress the obviously desirable freedom of an individual or family to possess the resources needed for living, while not implying the right to own the resources needed by others. . . .
There are, unsurprisingly, several traditions of individualist anarchism, one of them deriving from the ‘conscious egoism’ of the German writer Max Stirner (1806-56), and another from a remarkable series of 19th-century American figures who argued that in protecting our own autonomy and associating with others for common advantages, we are promoting the good of all. These thinkers differed from free-market liberals in their absolute mistrust of American capitalism, and in their emphasis on mutualism.

\vspace{1em}
\textbf{Question 17}

Which one of the following best sums up the overall history of transformation of male northern elephant seal calls?

\begin{enumerate}[label=\Alph*.]
    \item Owing to migrations in the aftermath of near species extinction, the average call pulse rates in the recolonised breeding locations exhibited a gradual increase until they matched the tempo at the founding colony.
    \item The calls have transformed from exhibiting simple composition, less individual variety, and great regional variety to complex composition, great individual variety, and less regional variety.
    \item Owing to migrations in the aftermath of near species extinction, the calls have transformed from exhibiting complex composition, less individual variety, and great regional variety to simple composition, less individual variety, and great regional variety.
    \item The calls have transformed from exhibiting simple composition, great individual variety, and less regional variety to complex composition, less individual variety, and great regional variety.
[](https://cracku.in/17-which-one-of-the-following-best-sums-up-the-overal-x-cat-2020-slot-1)
\end{enumerate}
\vspace{1em}
\textbf{Solution:}

The following excerpt serves as an essential source for comparing the difference in the attributes of the elephant seals: ["._..Yet there are other differences between the males from the late 1960s and their great-great-grandsons: Modern males exhibit more individual diversity, and their calls are more complex. While 50 years ago the drumming pattern was quite simple and the dialects denoted just a change in tempo, Casey explained, the calls recorded today have more complex structures, sometimes featuring doublets or triplets. . . ._ "]
In the late 1960s, the elephant seal calls were marked by having a simple drumming pattern which later transformed into calls with marked individual diversity and sophistication. Additionally, the dialects that were present in the 1960s were not evident during the study undertaken in the early 2010s, thereby indicating a decrease in the regional variations in the calls. These elements are aptly captured in Option B.

\end{problem}

\newpage
\begin{problem}{}{}
% Subject: varc
\textbf{Instructions}

Instructions   

**The passage below is accompanied by a set of questions. Choose the best answer to each question.**
The word ‘anarchy’ comes from the Greek '_anarkhia_ ', meaning contrary to authority or without a ruler, and was used in a derogatory sense until 1840, when it was adopted by Pierre-Joseph Proudhon to describe his political and social ideology. Proudhon argued that organization without government was both possible and desirable. In the evolution of political ideas, anarchism can be seen as an ultimate projection of both liberalism and socialism, and the differing strands of anarchist thought can be related to their emphasis on one or the other of these.
Historically, anarchism arose not only as an explanation of the gulf between the rich and the poor in any community, and of the reason why the poor have been obliged to fight for their share of a common inheritance, but as a radical answer to the question ‘What went wrong?’ that followed the ultimate outcome of the French Revolution. It had ended not only with a reign of terror and the emergence of a newly rich ruling caste, but with a new adored emperor, Napoleon Bonaparte, strutting through his conquered territories.
The anarchists and their precursors were unique on the political Left in affirming that workers and peasants, grasping the chance that arose to bring an end to centuries of exploitation and tyranny, were inevitably betrayed by the new class of politicians, whose first priority was to re-establish a centralized state power. After every revolutionary uprising, usually won at a heavy cost for ordinary populations, the new rulers had no hesitation in applying violence and terror, a secret police, and a professional army to maintain their control.
For anarchists the state itself is the enemy, and they have applied the same interpretation to the outcome of every revolution of the 19th and 20th centuries. This is not merely because every state keeps a watchful and sometimes punitive eye on its dissidents, but because every state protects the privileges of the powerful.
The mainstream of anarchist propaganda for more than a century has been anarchist- communism, which argues that property in land, natural resources, and the means of production should be held in mutual control by local communities, federating for innumerable joint purposes with other communes. It differs from state socialism in opposing the concept of any central authority. Some anarchists prefer to distinguish between anarchist-communism and collectivist anarchism in order to stress the obviously desirable freedom of an individual or family to possess the resources needed for living, while not implying the right to own the resources needed by others. . . .
There are, unsurprisingly, several traditions of individualist anarchism, one of them deriving from the ‘conscious egoism’ of the German writer Max Stirner (1806-56), and another from a remarkable series of 19th-century American figures who argued that in protecting our own autonomy and associating with others for common advantages, we are promoting the good of all. These thinkers differed from free-market liberals in their absolute mistrust of American capitalism, and in their emphasis on mutualism.

\vspace{1em}
\textbf{Question 18}

From the passage it can be inferred that the call pulse rate of male northern elephant seals in the southern rookeries was faster because:

\begin{enumerate}[label=\Alph*.]
    \item the male northern elephant seals of Isla Guadalupe with faster call pulse rates might have been the original settlers of the southern rookeries.
    \item a large number of male northern elephant seals migrated from the southern rookeries to Año Nuevo Island in the early 1970s.
    \item a large number of male northern elephant seals from Año Nuevo Island might have migrated to the southern rookeries to recolonise them.
    \item the calls of male northern elephant seals in the southern rookeries have more sophisticated structures, containing doublets and triplets.
[](https://cracku.in/18-from-the-passage-it-can-be-inferred-that-the-call--x-cat-2020-slot-1)
\end{enumerate}
\vspace{1em}
\textbf{Solution:}

We can make a direct inference based on the following excerpt from the fourth paragraph: [".._.This led Le Boeuf and his collaborator, Lewis Petrinovich, to deduce that the dialects were, perhaps, a result of isolation over time, after the breeding sites had been recolonised. For instance, the first settlers of Año Nuevo could have had, by chance, calls with low pulse rates. At other sites, where the scientists found faster pulse rates, the opposite would have happened—seals with faster rates would have happened to arrive first.._."]
Based on the above information, the only reason behind the call pulse rate of male northern elephant seals in the southern rookeries being faster would be because the male northern elephant seals of Isla Guadalupe with faster call pulse rates might have been the original settlers of this region. Option A correctly highlights this reason and is, hence, the correct answer.

Instructions   

For the following questions answer them individually

\end{problem}

\newpage
\begin{problem}{}{}
% Subject: varc
\textbf{Instructions}

Instructions   

**The passage below is accompanied by a set of questions. Choose the best answer to each question.**
The word ‘anarchy’ comes from the Greek '_anarkhia_ ', meaning contrary to authority or without a ruler, and was used in a derogatory sense until 1840, when it was adopted by Pierre-Joseph Proudhon to describe his political and social ideology. Proudhon argued that organization without government was both possible and desirable. In the evolution of political ideas, anarchism can be seen as an ultimate projection of both liberalism and socialism, and the differing strands of anarchist thought can be related to their emphasis on one or the other of these.
Historically, anarchism arose not only as an explanation of the gulf between the rich and the poor in any community, and of the reason why the poor have been obliged to fight for their share of a common inheritance, but as a radical answer to the question ‘What went wrong?’ that followed the ultimate outcome of the French Revolution. It had ended not only with a reign of terror and the emergence of a newly rich ruling caste, but with a new adored emperor, Napoleon Bonaparte, strutting through his conquered territories.
The anarchists and their precursors were unique on the political Left in affirming that workers and peasants, grasping the chance that arose to bring an end to centuries of exploitation and tyranny, were inevitably betrayed by the new class of politicians, whose first priority was to re-establish a centralized state power. After every revolutionary uprising, usually won at a heavy cost for ordinary populations, the new rulers had no hesitation in applying violence and terror, a secret police, and a professional army to maintain their control.
For anarchists the state itself is the enemy, and they have applied the same interpretation to the outcome of every revolution of the 19th and 20th centuries. This is not merely because every state keeps a watchful and sometimes punitive eye on its dissidents, but because every state protects the privileges of the powerful.
The mainstream of anarchist propaganda for more than a century has been anarchist- communism, which argues that property in land, natural resources, and the means of production should be held in mutual control by local communities, federating for innumerable joint purposes with other communes. It differs from state socialism in opposing the concept of any central authority. Some anarchists prefer to distinguish between anarchist-communism and collectivist anarchism in order to stress the obviously desirable freedom of an individual or family to possess the resources needed for living, while not implying the right to own the resources needed by others. . . .
There are, unsurprisingly, several traditions of individualist anarchism, one of them deriving from the ‘conscious egoism’ of the German writer Max Stirner (1806-56), and another from a remarkable series of 19th-century American figures who argued that in protecting our own autonomy and associating with others for common advantages, we are promoting the good of all. These thinkers differed from free-market liberals in their absolute mistrust of American capitalism, and in their emphasis on mutualism.

\vspace{1em}
\textbf{Question 19}

**Five jumbled up sentences, related to a topic, are given below. Four of them can be put together to form a coherent paragraph. Identify the odd one out and key in the number of the sentence as your answer:**  
  
1. For feminists, the question of how we read is inextricably linked with the question of what we read.  
2. Elaine Showalter’s critique of the literary curriculum is exemplary of this work.  
3. Androcentric literature structures the reading experience differently depending on the gender of the reader.  
4. The documentation of this realization was one of the earliest tasks undertaken by feminist critics.  
5. More specifically, the feminist inquiry into the activity of reading begins with the realization that the literary canon is androcentric, and that this has a profoundly damaging effect on women readers.
Backspace  
789  
456  
123  
0.-  
Clear All  

Submit
[](https://cracku.in/19-five-jumbled-up-sentences-related-to-a-topic-are-g-x-cat-2020-slot-1)

\begin{enumerate}[label=\Alph*.]
\end{enumerate}
\vspace{1em}
\textbf{Solution:}

The passage is focused on delineating the mechanism of reading as perceived by feminists (and the criticism associated with it). The arrangement (1)-(5)-(4)-(2) forms a coherent paragraph and Statement (3) stands out like a sore thumb. While the passage describes the elements associated with a feminist perspective, (3) brings in a description of androcentric literature that does not align with the context.

\end{problem}

\newpage
\begin{problem}{}{}
% Subject: varc
\textbf{Instructions}

Instructions   

**The passage below is accompanied by a set of questions. Choose the best answer to each question.**
The word ‘anarchy’ comes from the Greek '_anarkhia_ ', meaning contrary to authority or without a ruler, and was used in a derogatory sense until 1840, when it was adopted by Pierre-Joseph Proudhon to describe his political and social ideology. Proudhon argued that organization without government was both possible and desirable. In the evolution of political ideas, anarchism can be seen as an ultimate projection of both liberalism and socialism, and the differing strands of anarchist thought can be related to their emphasis on one or the other of these.
Historically, anarchism arose not only as an explanation of the gulf between the rich and the poor in any community, and of the reason why the poor have been obliged to fight for their share of a common inheritance, but as a radical answer to the question ‘What went wrong?’ that followed the ultimate outcome of the French Revolution. It had ended not only with a reign of terror and the emergence of a newly rich ruling caste, but with a new adored emperor, Napoleon Bonaparte, strutting through his conquered territories.
The anarchists and their precursors were unique on the political Left in affirming that workers and peasants, grasping the chance that arose to bring an end to centuries of exploitation and tyranny, were inevitably betrayed by the new class of politicians, whose first priority was to re-establish a centralized state power. After every revolutionary uprising, usually won at a heavy cost for ordinary populations, the new rulers had no hesitation in applying violence and terror, a secret police, and a professional army to maintain their control.
For anarchists the state itself is the enemy, and they have applied the same interpretation to the outcome of every revolution of the 19th and 20th centuries. This is not merely because every state keeps a watchful and sometimes punitive eye on its dissidents, but because every state protects the privileges of the powerful.
The mainstream of anarchist propaganda for more than a century has been anarchist- communism, which argues that property in land, natural resources, and the means of production should be held in mutual control by local communities, federating for innumerable joint purposes with other communes. It differs from state socialism in opposing the concept of any central authority. Some anarchists prefer to distinguish between anarchist-communism and collectivist anarchism in order to stress the obviously desirable freedom of an individual or family to possess the resources needed for living, while not implying the right to own the resources needed by others. . . .
There are, unsurprisingly, several traditions of individualist anarchism, one of them deriving from the ‘conscious egoism’ of the German writer Max Stirner (1806-56), and another from a remarkable series of 19th-century American figures who argued that in protecting our own autonomy and associating with others for common advantages, we are promoting the good of all. These thinkers differed from free-market liberals in their absolute mistrust of American capitalism, and in their emphasis on mutualism.

\vspace{1em}
\textbf{Question 20}

The passage given below is followed by four alternate summaries. Choose the option that best captures the essence of the passage.  
  
For nearly a century most psychologists have embraced one view of intelligence. Individuals are born with more or less intelligence potential (I.Q.); this potential is heavily influenced by heredity and difficult to alter; experts in measurement can determine a person’s intelligence early in life, currently from paper-and-pencil measures, perhaps eventually from examining the brain in action or even scrutinizing his/her genome. Recently, criticism of this conventional wisdom has mounted. Biologists ask if speaking of a single entity called “intelligence” is coherent and question the validity of measures used to estimate heritability of a trait in humans, who, unlike plants or animals, are not conceived and bred under controlled conditions.

\begin{enumerate}[label=\Alph*.]
    \item Biologists have questioned the long-standing view that ‘intelligence’ is a single entity and the attempts to estimate it's heritability.
    \item Biologists have criticised that conventional wisdom that individuals are born with more or less intelligence potential.
    \item Biologists have started questioning psychologists' view of 'intelligence' as a measurable immutable characteristic of an individual.
    \item Biologists have questioned the view that ‘intelligence’ is a single entity and the ways in which what is inherited
[](https://cracku.in/20-the-passage-given-below-is-followed-by-four-altern-x-cat-2020-slot-1)
\end{enumerate}
\vspace{1em}
\textbf{Solution:}

The paragraph highlights the following:
1. The validity of the ubiquitous perspective held by psychologists {of intelligence being a measurable, unalterable entity that is greatly influenced by heredity} is now being questioned by biologists. 
2. The dubiety concerning the aspect of intelligence being hereditary {given the fact that "_humans, who, unlike plants or animals, are not conceived and bred under controlled conditions._ "}
Thus, a statement capturing these elements is bound to be the answer. _Option A_ aptly encompasses these two key points.
_Option B_ : Calling the widely -held perspective as conventional wisdom would be inappropriate. Additionally, the statement here fails to capture point (2).
_Option C_ : Although close, it misses out on the second half of the discussion.
_Option D_ : This option might appear confusing, given that it touches upon both the key elements. However, it is unspecific and comes across as a bit odd {"_ways in which what is inherited_ " doesn't make sense}. Between Options A and D, A is definitely the better choice.

\end{problem}

\newpage
\begin{problem}{}{}
% Subject: varc
\textbf{Instructions}

Instructions   

**The passage below is accompanied by a set of questions. Choose the best answer to each question.**
The word ‘anarchy’ comes from the Greek '_anarkhia_ ', meaning contrary to authority or without a ruler, and was used in a derogatory sense until 1840, when it was adopted by Pierre-Joseph Proudhon to describe his political and social ideology. Proudhon argued that organization without government was both possible and desirable. In the evolution of political ideas, anarchism can be seen as an ultimate projection of both liberalism and socialism, and the differing strands of anarchist thought can be related to their emphasis on one or the other of these.
Historically, anarchism arose not only as an explanation of the gulf between the rich and the poor in any community, and of the reason why the poor have been obliged to fight for their share of a common inheritance, but as a radical answer to the question ‘What went wrong?’ that followed the ultimate outcome of the French Revolution. It had ended not only with a reign of terror and the emergence of a newly rich ruling caste, but with a new adored emperor, Napoleon Bonaparte, strutting through his conquered territories.
The anarchists and their precursors were unique on the political Left in affirming that workers and peasants, grasping the chance that arose to bring an end to centuries of exploitation and tyranny, were inevitably betrayed by the new class of politicians, whose first priority was to re-establish a centralized state power. After every revolutionary uprising, usually won at a heavy cost for ordinary populations, the new rulers had no hesitation in applying violence and terror, a secret police, and a professional army to maintain their control.
For anarchists the state itself is the enemy, and they have applied the same interpretation to the outcome of every revolution of the 19th and 20th centuries. This is not merely because every state keeps a watchful and sometimes punitive eye on its dissidents, but because every state protects the privileges of the powerful.
The mainstream of anarchist propaganda for more than a century has been anarchist- communism, which argues that property in land, natural resources, and the means of production should be held in mutual control by local communities, federating for innumerable joint purposes with other communes. It differs from state socialism in opposing the concept of any central authority. Some anarchists prefer to distinguish between anarchist-communism and collectivist anarchism in order to stress the obviously desirable freedom of an individual or family to possess the resources needed for living, while not implying the right to own the resources needed by others. . . .
There are, unsurprisingly, several traditions of individualist anarchism, one of them deriving from the ‘conscious egoism’ of the German writer Max Stirner (1806-56), and another from a remarkable series of 19th-century American figures who argued that in protecting our own autonomy and associating with others for common advantages, we are promoting the good of all. These thinkers differed from free-market liberals in their absolute mistrust of American capitalism, and in their emphasis on mutualism.

\vspace{1em}
\textbf{Question 21}

The four sentences (labelled 1, 2, 3, 4) below, when properly sequenced would yield a coherent paragraph. Decide on the proper sequencing of the order of the sentences and key in the sequence of the four numbers as your answer:  
  
1. Relying on narrative structure alone, indigenous significances of nineteenth-century San folktales are hard to determine.  
2. Using their supernatural potency, benign shamans transcend the levels of the San cosmos in order to deal with social conflict and to protect material resources and enjoy a measure of respect that sets them apart from ordinary people.  
3. Selected tales reveal that they deal with a form of spiritual conflict that has social implications and concern conflict between people and living or dead malevolent shamans.  
4. Meaning can be elicited, and the tales contextualized, by probing beneath the narrative of verbatim, original-language records and exploring the connotations of highly significant words and phrases.
Backspace  
789  
456  
123  
0.-  
Clear All  

Submit
[](https://cracku.in/21-the-four-sentences-labelled-1-2-3-4-below-when-pro-x-cat-2020-slot-1)

\begin{enumerate}[label=\Alph*.]
\end{enumerate}
\vspace{1em}
\textbf{Solution:}

Statement (1) introduces the predicament we face: the difficulty in understanding the indigenous significance of the nineteenth century San folktales by simply probing the narrative structure. The author indicates that there are additional variables that need to be considered. The aspect of evaluating elements beyond the "_the narrative of verbatim_ " as highlighted in Statement (4). The outcome of such a methodology of exploring these narratives is then presented in the Statements (3) followed by (2). Statement (3) highlights dead-malevolent shamans while (2) talks about benign shamans. Hence, we obtain (1)-(4)-(3)-(2) as a coherent paragraph.

\end{problem}

\newpage
\begin{problem}{}{}
% Subject: varc
\textbf{Instructions}

Instructions   

**The passage below is accompanied by a set of questions. Choose the best answer to each question.**
The word ‘anarchy’ comes from the Greek '_anarkhia_ ', meaning contrary to authority or without a ruler, and was used in a derogatory sense until 1840, when it was adopted by Pierre-Joseph Proudhon to describe his political and social ideology. Proudhon argued that organization without government was both possible and desirable. In the evolution of political ideas, anarchism can be seen as an ultimate projection of both liberalism and socialism, and the differing strands of anarchist thought can be related to their emphasis on one or the other of these.
Historically, anarchism arose not only as an explanation of the gulf between the rich and the poor in any community, and of the reason why the poor have been obliged to fight for their share of a common inheritance, but as a radical answer to the question ‘What went wrong?’ that followed the ultimate outcome of the French Revolution. It had ended not only with a reign of terror and the emergence of a newly rich ruling caste, but with a new adored emperor, Napoleon Bonaparte, strutting through his conquered territories.
The anarchists and their precursors were unique on the political Left in affirming that workers and peasants, grasping the chance that arose to bring an end to centuries of exploitation and tyranny, were inevitably betrayed by the new class of politicians, whose first priority was to re-establish a centralized state power. After every revolutionary uprising, usually won at a heavy cost for ordinary populations, the new rulers had no hesitation in applying violence and terror, a secret police, and a professional army to maintain their control.
For anarchists the state itself is the enemy, and they have applied the same interpretation to the outcome of every revolution of the 19th and 20th centuries. This is not merely because every state keeps a watchful and sometimes punitive eye on its dissidents, but because every state protects the privileges of the powerful.
The mainstream of anarchist propaganda for more than a century has been anarchist- communism, which argues that property in land, natural resources, and the means of production should be held in mutual control by local communities, federating for innumerable joint purposes with other communes. It differs from state socialism in opposing the concept of any central authority. Some anarchists prefer to distinguish between anarchist-communism and collectivist anarchism in order to stress the obviously desirable freedom of an individual or family to possess the resources needed for living, while not implying the right to own the resources needed by others. . . .
There are, unsurprisingly, several traditions of individualist anarchism, one of them deriving from the ‘conscious egoism’ of the German writer Max Stirner (1806-56), and another from a remarkable series of 19th-century American figures who argued that in protecting our own autonomy and associating with others for common advantages, we are promoting the good of all. These thinkers differed from free-market liberals in their absolute mistrust of American capitalism, and in their emphasis on mutualism.

\vspace{1em}
\textbf{Question 22}

The passage given below is followed by four alternate summaries. Choose the option that best captures the essence of the passage.  
  
As Soviet power declined, the world became to some extent multipolar, and Europe strove to define an independent identity. What a journey Europe has undertaken to reach this point. It had in every century changed its internal structure and invented new ways of thinking about the nature of international order. Now at the culmination of an era, Europe, in order to participate in it, felt obliged to set aside the political mechanisms through which it had conducted its affairs for three and a half centuries. Impelled also by the desire to cushion the emergent unification of Germany, the new European Union established a common currency in 2002 and a formal political structure in 2004. It proclaimed a Europe united, whole, and free, adjusting its differences by peaceful mechanisms.

\begin{enumerate}[label=\Alph*.]
    \item Europe has consistently changed in keeping with the changing world order and that has culminated in a united Europe.
    \item The establishment of a formal political structure in Europe was hastened by the unification of Germany and the emergence of a multipolar world.
    \item Europe has consistently changed its internal structure to successfully adapt to the changing world order.
    \item Europe has chosen to lower political and economic heterogeneity, in order to adapt itself to an emerging multi-polar world.
[](https://cracku.in/22-the-passage-given-below-is-followed-by-four-altern-x-cat-2020-slot-1)
\end{enumerate}
\vspace{1em}
\textbf{Solution:}

The passage begins by highlighting Europe's continual attempt to adapt itself in a multipolar world by striving to be a dynamic entity- nationally and internationally{"_changed its internal structure and invented new ways of thinking about the nature of international order_ "}. Post this, the author portrays how certain stimuli in the modern world has lead Europe to review its political components {"_set aside the political mechanisms through which it had conducted its affairs for three and a half centuries_ "} and to make changes in its economic structure {"_established a common currency_ "}. Thus, the passage presents two key elements: (1) the fact that Europe has consistently tried to adapt to a changing world and (2) the manner in which Europe has attempted to achieve that in the existing multi-polar setup. Option D correctly highlights these points.
Option A: This misses out on the point (2). Furthermore, directly ascribing the unification of Europe to its attempt to rapidly change would be incorrect. 
Option B: This does not fully capture the essence of the passage and misses out on point 1.
Option C: Although the statement in this option captures point 1, it misses out on point 2.
Hence, Option D is the correct answer.

\end{problem}

\newpage
\begin{problem}{}{}
% Subject: varc
\textbf{Instructions}

Instructions   

**The passage below is accompanied by a set of questions. Choose the best answer to each question.**
The word ‘anarchy’ comes from the Greek '_anarkhia_ ', meaning contrary to authority or without a ruler, and was used in a derogatory sense until 1840, when it was adopted by Pierre-Joseph Proudhon to describe his political and social ideology. Proudhon argued that organization without government was both possible and desirable. In the evolution of political ideas, anarchism can be seen as an ultimate projection of both liberalism and socialism, and the differing strands of anarchist thought can be related to their emphasis on one or the other of these.
Historically, anarchism arose not only as an explanation of the gulf between the rich and the poor in any community, and of the reason why the poor have been obliged to fight for their share of a common inheritance, but as a radical answer to the question ‘What went wrong?’ that followed the ultimate outcome of the French Revolution. It had ended not only with a reign of terror and the emergence of a newly rich ruling caste, but with a new adored emperor, Napoleon Bonaparte, strutting through his conquered territories.
The anarchists and their precursors were unique on the political Left in affirming that workers and peasants, grasping the chance that arose to bring an end to centuries of exploitation and tyranny, were inevitably betrayed by the new class of politicians, whose first priority was to re-establish a centralized state power. After every revolutionary uprising, usually won at a heavy cost for ordinary populations, the new rulers had no hesitation in applying violence and terror, a secret police, and a professional army to maintain their control.
For anarchists the state itself is the enemy, and they have applied the same interpretation to the outcome of every revolution of the 19th and 20th centuries. This is not merely because every state keeps a watchful and sometimes punitive eye on its dissidents, but because every state protects the privileges of the powerful.
The mainstream of anarchist propaganda for more than a century has been anarchist- communism, which argues that property in land, natural resources, and the means of production should be held in mutual control by local communities, federating for innumerable joint purposes with other communes. It differs from state socialism in opposing the concept of any central authority. Some anarchists prefer to distinguish between anarchist-communism and collectivist anarchism in order to stress the obviously desirable freedom of an individual or family to possess the resources needed for living, while not implying the right to own the resources needed by others. . . .
There are, unsurprisingly, several traditions of individualist anarchism, one of them deriving from the ‘conscious egoism’ of the German writer Max Stirner (1806-56), and another from a remarkable series of 19th-century American figures who argued that in protecting our own autonomy and associating with others for common advantages, we are promoting the good of all. These thinkers differed from free-market liberals in their absolute mistrust of American capitalism, and in their emphasis on mutualism.

\vspace{1em}
\textbf{Question 23}

The passage given below is followed by four alternate summaries. Choose the option that best captures the essence of the passage.  
  
For years, movies and television series like Crime Scene Investigation (CSI) paint an unrealistic picture of the “science of voices.” In the 1994 movie Clear and Present Danger an expert listens to a brief recorded utterance and declares that the speaker is “Cuban, aged 35 to 45, educated in the [...] eastern United States.” The recording is then fed to a supercomputer that matches the voice to that of a suspect, concluding that the probability of correct identification is 90%. This sequence sums up a good number of misimpressions about forensic phonetics, which have led to errors in real- life justice. Indeed, that movie scene exemplifies the so-called “CSI effect”—the phenomenon in which judges hold unrealistic expectations of the capabilities of forensic science.

\begin{enumerate}[label=\Alph*.]
    \item Although voice recognition is often presented as evidence in legal cases, its scientific basis can be shaky.
    \item Voice recognition as used in many movies to identify criminals has been used to identify criminals in real life also.
    \item Movies and televisions have led to the belief that the use of forensic phonetics in legal investigations is robust and fool proof.
    \item Voice recognition has started to feature prominently in crime-scene intelligence investigations because of movies and television series.
[](https://cracku.in/23-the-passage-given-below-is-followed-by-four-altern-x-cat-2020-slot-1)
\end{enumerate}
\vspace{1em}
\textbf{Solution:}

The passage is primarily about the misconception regarding the capacity of forensic phonetics stemming from its portrayal in movies and television. This aspect is correctly captured by the point stated in C. The author does not question the "scientific basis" of the evidence (based on voice recognition) that is presented in legal cases. Hence, Option A can be eliminated. The claim made in Option B cannot be understood from the passage. Same can be said about Option D. Hence, Option C is the correct answer.

\end{problem}

\newpage
\begin{problem}{}{}
% Subject: varc
\textbf{Instructions}

Instructions   

**The passage below is accompanied by a set of questions. Choose the best answer to each question.**
The word ‘anarchy’ comes from the Greek '_anarkhia_ ', meaning contrary to authority or without a ruler, and was used in a derogatory sense until 1840, when it was adopted by Pierre-Joseph Proudhon to describe his political and social ideology. Proudhon argued that organization without government was both possible and desirable. In the evolution of political ideas, anarchism can be seen as an ultimate projection of both liberalism and socialism, and the differing strands of anarchist thought can be related to their emphasis on one or the other of these.
Historically, anarchism arose not only as an explanation of the gulf between the rich and the poor in any community, and of the reason why the poor have been obliged to fight for their share of a common inheritance, but as a radical answer to the question ‘What went wrong?’ that followed the ultimate outcome of the French Revolution. It had ended not only with a reign of terror and the emergence of a newly rich ruling caste, but with a new adored emperor, Napoleon Bonaparte, strutting through his conquered territories.
The anarchists and their precursors were unique on the political Left in affirming that workers and peasants, grasping the chance that arose to bring an end to centuries of exploitation and tyranny, were inevitably betrayed by the new class of politicians, whose first priority was to re-establish a centralized state power. After every revolutionary uprising, usually won at a heavy cost for ordinary populations, the new rulers had no hesitation in applying violence and terror, a secret police, and a professional army to maintain their control.
For anarchists the state itself is the enemy, and they have applied the same interpretation to the outcome of every revolution of the 19th and 20th centuries. This is not merely because every state keeps a watchful and sometimes punitive eye on its dissidents, but because every state protects the privileges of the powerful.
The mainstream of anarchist propaganda for more than a century has been anarchist- communism, which argues that property in land, natural resources, and the means of production should be held in mutual control by local communities, federating for innumerable joint purposes with other communes. It differs from state socialism in opposing the concept of any central authority. Some anarchists prefer to distinguish between anarchist-communism and collectivist anarchism in order to stress the obviously desirable freedom of an individual or family to possess the resources needed for living, while not implying the right to own the resources needed by others. . . .
There are, unsurprisingly, several traditions of individualist anarchism, one of them deriving from the ‘conscious egoism’ of the German writer Max Stirner (1806-56), and another from a remarkable series of 19th-century American figures who argued that in protecting our own autonomy and associating with others for common advantages, we are promoting the good of all. These thinkers differed from free-market liberals in their absolute mistrust of American capitalism, and in their emphasis on mutualism.

\vspace{1em}
\textbf{Question 24}

The four sentences (labelled 1, 2, 3, 4) below, when properly sequenced would yield a coherent paragraph. Decide on the proper sequencing of the order of the sentences and key in the sequence of the four numbers as your answer:  
  
1. Tensions and sometimes conflict remain an issue in and between the 11 states in South East Asia (Brunei Darussalam, Cambodia, Indonesia, Laos, Malaysia, Myanmar, the Philippines, Singapore, Thailand, Timor-Leste and Vietnam).  
2. China’s rise as a regional military power and its claims in the South China Sea have become an increasingly pressing security concern for many South East Asian states.  
3. Since the 1990s, the security environment of South East Asia has seen both continuity and profound changes.  
4. These concerns cause states from outside the region to take an active interest in South East Asian security.
Backspace  
789  
456  
123  
0.-  
Clear All  

Submit
[](https://cracku.in/24-the-four-sentences-labelled-1-2-3-4-below-when-pro-x-cat-2020-slot-1)

\begin{enumerate}[label=\Alph*.]
\end{enumerate}
\vspace{1em}
\textbf{Solution:}

Statement (3) introduces the core message: the continuing concerns in the region of South-East Asia. It highlights that since the 1990s there have been dynamic changes in the security environment of this region. Statment (1) adds details by listing the countries present in this region and the existing environment rife with tension/conflict {a transition from the 1990s into the present world}. Statement (2) continues on this line and cites an event in this regard: China becoming a pressing security concern for the regional countries. Statement (4) puts a cap on this discussion by indicating that the persisting regional instability attracts entities from outside the region. Hence, (3)-(1)-(2)-(4) forms a coherent arrangement.

\end{problem}

\newpage
\begin{problem}{}{}
% Subject: varc
\textbf{Instructions}

Instructions   

**The passage below is accompanied by a set of questions. Choose the best answer to each question.**
The word ‘anarchy’ comes from the Greek '_anarkhia_ ', meaning contrary to authority or without a ruler, and was used in a derogatory sense until 1840, when it was adopted by Pierre-Joseph Proudhon to describe his political and social ideology. Proudhon argued that organization without government was both possible and desirable. In the evolution of political ideas, anarchism can be seen as an ultimate projection of both liberalism and socialism, and the differing strands of anarchist thought can be related to their emphasis on one or the other of these.
Historically, anarchism arose not only as an explanation of the gulf between the rich and the poor in any community, and of the reason why the poor have been obliged to fight for their share of a common inheritance, but as a radical answer to the question ‘What went wrong?’ that followed the ultimate outcome of the French Revolution. It had ended not only with a reign of terror and the emergence of a newly rich ruling caste, but with a new adored emperor, Napoleon Bonaparte, strutting through his conquered territories.
The anarchists and their precursors were unique on the political Left in affirming that workers and peasants, grasping the chance that arose to bring an end to centuries of exploitation and tyranny, were inevitably betrayed by the new class of politicians, whose first priority was to re-establish a centralized state power. After every revolutionary uprising, usually won at a heavy cost for ordinary populations, the new rulers had no hesitation in applying violence and terror, a secret police, and a professional army to maintain their control.
For anarchists the state itself is the enemy, and they have applied the same interpretation to the outcome of every revolution of the 19th and 20th centuries. This is not merely because every state keeps a watchful and sometimes punitive eye on its dissidents, but because every state protects the privileges of the powerful.
The mainstream of anarchist propaganda for more than a century has been anarchist- communism, which argues that property in land, natural resources, and the means of production should be held in mutual control by local communities, federating for innumerable joint purposes with other communes. It differs from state socialism in opposing the concept of any central authority. Some anarchists prefer to distinguish between anarchist-communism and collectivist anarchism in order to stress the obviously desirable freedom of an individual or family to possess the resources needed for living, while not implying the right to own the resources needed by others. . . .
There are, unsurprisingly, several traditions of individualist anarchism, one of them deriving from the ‘conscious egoism’ of the German writer Max Stirner (1806-56), and another from a remarkable series of 19th-century American figures who argued that in protecting our own autonomy and associating with others for common advantages, we are promoting the good of all. These thinkers differed from free-market liberals in their absolute mistrust of American capitalism, and in their emphasis on mutualism.

\vspace{1em}
\textbf{Question 25}

Five jumbled up sentences, related to a topic, are given below. Four of them can be put together to form a coherent paragraph. Identify the odd one out and key in the number of the sentence as your answer:  
  
1. Talk was the most common way for enslaved men and women to subvert the rules of their bondage, to gain more agency than they were supposed to have.  
2. Even in conditions of extreme violence and unfreedom, their words remained ubiquitous, ephemeral, irrepressible, and potentially transgressive.  
3. Slaves came from societies in which oaths, orations, and invocations carried great potency, both between people and as a connection to the all-powerful spirit world.  
4. Freedom of speech and the power to silence may have been preeminent markers of white liberty in Colonies, but at the same time, slavery depended on dialogue: slaves could never be completely muted.  
5. Slave-owners obsessed over slave talk, though they could never control it, yet feared its power to bind and inspire—for, as everyone knew, oaths, whispers, and secret conversations bred conspiracy and revolt.
Backspace  
789  
456  
123  
0.-  
Clear All  

Submit
[](https://cracku.in/25-five-jumbled-up-sentences-related-to-a-topic-are-g-x-cat-2020-slot-1)

\begin{enumerate}[label=\Alph*.]
\end{enumerate}
\vspace{1em}
\textbf{Solution:}

Statements (1), (2), (4) and (5) discuss the aspect of how talking/freedom of speech was a significant facet associated with the slaves. They emphasise how this element was central to slavery { "slaves could never be completely muted"}. Contrarily, Statment (3) goes on a tangential route that associates the potency of "oaths, orations, and invocations" based on the origin of the slaves {discusses "..._a connection to the all-powerful spirit world._.." which is clearly out of place}. Thus, Statement (3) is the odd-one-out here.

\end{problem}

\newpage
\begin{problem}{}{}
% Subject: varc
\textbf{Instructions}

Instructions   

**The passage below is accompanied by a set of questions. Choose the best answer to each question.**
The word ‘anarchy’ comes from the Greek '_anarkhia_ ', meaning contrary to authority or without a ruler, and was used in a derogatory sense until 1840, when it was adopted by Pierre-Joseph Proudhon to describe his political and social ideology. Proudhon argued that organization without government was both possible and desirable. In the evolution of political ideas, anarchism can be seen as an ultimate projection of both liberalism and socialism, and the differing strands of anarchist thought can be related to their emphasis on one or the other of these.
Historically, anarchism arose not only as an explanation of the gulf between the rich and the poor in any community, and of the reason why the poor have been obliged to fight for their share of a common inheritance, but as a radical answer to the question ‘What went wrong?’ that followed the ultimate outcome of the French Revolution. It had ended not only with a reign of terror and the emergence of a newly rich ruling caste, but with a new adored emperor, Napoleon Bonaparte, strutting through his conquered territories.
The anarchists and their precursors were unique on the political Left in affirming that workers and peasants, grasping the chance that arose to bring an end to centuries of exploitation and tyranny, were inevitably betrayed by the new class of politicians, whose first priority was to re-establish a centralized state power. After every revolutionary uprising, usually won at a heavy cost for ordinary populations, the new rulers had no hesitation in applying violence and terror, a secret police, and a professional army to maintain their control.
For anarchists the state itself is the enemy, and they have applied the same interpretation to the outcome of every revolution of the 19th and 20th centuries. This is not merely because every state keeps a watchful and sometimes punitive eye on its dissidents, but because every state protects the privileges of the powerful.
The mainstream of anarchist propaganda for more than a century has been anarchist- communism, which argues that property in land, natural resources, and the means of production should be held in mutual control by local communities, federating for innumerable joint purposes with other communes. It differs from state socialism in opposing the concept of any central authority. Some anarchists prefer to distinguish between anarchist-communism and collectivist anarchism in order to stress the obviously desirable freedom of an individual or family to possess the resources needed for living, while not implying the right to own the resources needed by others. . . .
There are, unsurprisingly, several traditions of individualist anarchism, one of them deriving from the ‘conscious egoism’ of the German writer Max Stirner (1806-56), and another from a remarkable series of 19th-century American figures who argued that in protecting our own autonomy and associating with others for common advantages, we are promoting the good of all. These thinkers differed from free-market liberals in their absolute mistrust of American capitalism, and in their emphasis on mutualism.

\vspace{1em}
\textbf{Question 26}

The four sentences (labelled 1, 2, 3, 4) below, when properly sequenced would yield a coherent paragraph. Decide on the proper sequencing of the order of the sentences and key in the sequence of the four numbers as your answer:  
  
1. Man has used poisons for assassination purposes ever since the dawn of civilization, against individual enemies but also occasionally against armies.  
2. These dangers were soon recognized, and resulted in two international declarations—in 1874 in Brussels and in 1899 in The Hague—that prohibited the use of poisoned weapons.  
3. The foundation of microbiology by Louis Pasteur and Robert Koch offered new prospects for those interested in biological weapons because it allowed agents to be chosen and designed on a rational basis.  
4. Though treaties were all made in good faith, they contained no means of control, and so fail
Backspace  
789  
456  
123  
0.-  
Clear All  

Submit
[](https://cracku.in/26-the-four-sentences-labelled-1-2-3-4-below-when-pro-x-cat-2020-slot-1)

\begin{enumerate}[label=\Alph*.]
\end{enumerate}
\vspace{1em}
\textbf{Solution:}

Statement (1) serves as a general introduction to the topic: an inclination towards using poison {or associated agents}. The author mentions that the utility of this element has existed since the dawn of civilization. Statment (3) supplements this idea by mentioning that certain advancements served as a new avenue for the zealots of this subject. Statement (2) delineates on the safeguard that was put into place once the threat posed by this domain was realized. Statement (4) pinpoints the flaw in this safeguard. Hence, we notice that (1)-(3)-(2)-(4) forms a logical arrangement.
[](https://cracku.in/downloads/11132)
  *  
  *

\end{problem}

\newpage
\begin{problem}{}{}
% Subject: varc
\textbf{Instructions}

Instructions   

The passage below is accompanied by a set of questions. Choose the best answer to each question.
The claims advanced here may be condensed into two assertions: [first, that visual] culture is what images, acts of seeing, and attendant intellectual, emotional, and perceptual sensibilities do to build, maintain, or transform the worlds in which people live. [And second, that the] study of visual culture is the analysis and interpretation of images and the ways of seeing (or gazes) that configure the agents, practices, conceptualities, and institutions that put images to work. . . .
Accordingly, the study of visual culture should be characterized by several concerns. First, scholars of visual culture need to examine any and all imagery - high and low, art and non-art. . . . They must not restrict themselves to objects of a particular beauty or aesthetic value. Indeed, any kind of imagery may be found to offer up evidence of the visual construction of reality. . . .
Second, the study of visual culture must scrutinize visual practice as much as images themselves, asking what images do when they are put to use. If scholars engaged in this enterprise inquire what makes an image beautiful or why this image or that constitutes a masterpiece or a work of genius, they should do so with the purpose of investigating an artist’s or a work’s contribution to the experience of beauty, taste, value, or genius. No amount of social analysis can account fully for the existence of Michelangelo or Leonardo. They were unique creators of images that changed the way their contemporaries thought and felt and have continued to shape the history of art, artists, museums, feeling, and aesthetic value. But study of the critical, artistic, and popular reception of works by such artists as Michelangelo and Leonardo can shed important light on the meaning of these artists and their works for many different people. And the history of meaning-making has a great deal to do with how scholars as well as lay audiences today understand these artists and their achievements.
Third, scholars studying visual culture might properly focus their interpretative work on life worlds by examining images, practices, visual technologies, taste, and artistic style as constitutive of social relations. The task is to understand how artifacts contribute to the construction of a world. . . . Important methodological implications follow: ethnography and reception studies become productive forms of gathering information, since these move beyond the image as a closed and fixed meaning-event. . . .
Fourth, scholars may learn a great deal when they scrutinize the constituents of vision, that is, the structures of perception as a physiological process as well as the epistemological frameworks informing a system of visual representation. Vision is a socially and a biologically constructed operation, depending on the design of the human body and how it engages the interpretive devices developed by a culture in order to see intelligibly. . . . Seeing . . . operates on the foundation of covenants with images that establish the conditions for meaningful visual experience.
Finally, the scholar of visual culture seeks to regard images as evidence for explanation, not as epiphenomena.

\vspace{1em}
\textbf{Question 1}

“Seeing . . . operates on the foundation of covenants with images that establish the conditions for meaningful visual experience.” In light of the passage, which one of the following statements best conveys the meaning of this sentence?

\begin{enumerate}[label=\Alph*.]
    \item Sight becomes a meaningful visual experience because of covenants of meaningfulness that we establish with the images we see.
    \item Images are meaningful visual experiences when they have a foundation of covenants seeing them.
    \item Sight as a meaningful visual experience is possible when there is a foundational condition established in images of covenants.
    \item The way we experience sight is through images operated on by meaningful covenants.
[](https://cracku.in/1-seeing-operates-on-the-foundation-of-covenants-wit-x-cat-2020-slot-2)
\end{enumerate}
\vspace{1em}
\textbf{Solution:}

Let's look at the options one by one.
Option B is "Images are meaningful visual experiences when they have a foundation of covenants seeing them." This is a distorted option. The actual statement states that sight becomes a meaning visual experience when images are associated with covenants . There is nowhere any discussion about the meaningfulness of images
Option C is "Sight as a meaningful visual experience is possible when there is a foundational condition established in images of covenants." Again a twisted option. There is nothing said about the possibility of sight as a visual experience. There could be other cases too in which meaningful visual experience could be established
Option D is "The way we experience sight is through images operated on by meaningful covenants." Entirely out of context option. The statement is about meaningful visual experience not about meaningful covenants.
Option A correctly encapsulates all the points. Hence it is the correct answer.

\end{problem}

\newpage
\begin{problem}{}{}
% Subject: varc
\textbf{Instructions}

Instructions   

The passage below is accompanied by a set of questions. Choose the best answer to each question.
The claims advanced here may be condensed into two assertions: [first, that visual] culture is what images, acts of seeing, and attendant intellectual, emotional, and perceptual sensibilities do to build, maintain, or transform the worlds in which people live. [And second, that the] study of visual culture is the analysis and interpretation of images and the ways of seeing (or gazes) that configure the agents, practices, conceptualities, and institutions that put images to work. . . .
Accordingly, the study of visual culture should be characterized by several concerns. First, scholars of visual culture need to examine any and all imagery - high and low, art and non-art. . . . They must not restrict themselves to objects of a particular beauty or aesthetic value. Indeed, any kind of imagery may be found to offer up evidence of the visual construction of reality. . . .
Second, the study of visual culture must scrutinize visual practice as much as images themselves, asking what images do when they are put to use. If scholars engaged in this enterprise inquire what makes an image beautiful or why this image or that constitutes a masterpiece or a work of genius, they should do so with the purpose of investigating an artist’s or a work’s contribution to the experience of beauty, taste, value, or genius. No amount of social analysis can account fully for the existence of Michelangelo or Leonardo. They were unique creators of images that changed the way their contemporaries thought and felt and have continued to shape the history of art, artists, museums, feeling, and aesthetic value. But study of the critical, artistic, and popular reception of works by such artists as Michelangelo and Leonardo can shed important light on the meaning of these artists and their works for many different people. And the history of meaning-making has a great deal to do with how scholars as well as lay audiences today understand these artists and their achievements.
Third, scholars studying visual culture might properly focus their interpretative work on life worlds by examining images, practices, visual technologies, taste, and artistic style as constitutive of social relations. The task is to understand how artifacts contribute to the construction of a world. . . . Important methodological implications follow: ethnography and reception studies become productive forms of gathering information, since these move beyond the image as a closed and fixed meaning-event. . . .
Fourth, scholars may learn a great deal when they scrutinize the constituents of vision, that is, the structures of perception as a physiological process as well as the epistemological frameworks informing a system of visual representation. Vision is a socially and a biologically constructed operation, depending on the design of the human body and how it engages the interpretive devices developed by a culture in order to see intelligibly. . . . Seeing . . . operates on the foundation of covenants with images that establish the conditions for meaningful visual experience.
Finally, the scholar of visual culture seeks to regard images as evidence for explanation, not as epiphenomena.

\vspace{1em}
\textbf{Question 2}

All of the following statements may be considered valid inferences from the passage, EXCEPT:

\begin{enumerate}[label=\Alph*.]
    \item visual culture is not just about how we see, but also about how our visual practices can impact and change the world.
    \item artifacts are meaningful precisely because they help to construct the meanings of the world for us.
    \item understanding the structures of perception is an important part of understanding how visual cultures work.
    \item studying visual culture requires institutional structures without which the structures of perception cannot be analysed.
[](https://cracku.in/2-all-of-the-following-statements-may-be-considered--x-cat-2020-slot-2)
\end{enumerate}
\vspace{1em}
\textbf{Solution:}

Option D states "studying visual culture requires institutional structures without which the structures of perception cannot be analysed".
Look at the penultimate paragraph of the question," Vision is a socially and a biologically constructed operation, depending on the design of the human body and how it engages the interpretive devices developed by a culture in order to see intelligibly"
Studying visual culture thus depends on the design of human body and interpretative devices developed by the culture. Nowhere it is mentioned that ,without institutional structures of culture vision can't be analysed as it also depends on the design of human body. Hence this is a wrong inference. Remaining all three options are correct.

\end{problem}

\newpage
\begin{problem}{}{}
% Subject: varc
\textbf{Instructions}

Instructions   

The passage below is accompanied by a set of questions. Choose the best answer to each question.
The claims advanced here may be condensed into two assertions: [first, that visual] culture is what images, acts of seeing, and attendant intellectual, emotional, and perceptual sensibilities do to build, maintain, or transform the worlds in which people live. [And second, that the] study of visual culture is the analysis and interpretation of images and the ways of seeing (or gazes) that configure the agents, practices, conceptualities, and institutions that put images to work. . . .
Accordingly, the study of visual culture should be characterized by several concerns. First, scholars of visual culture need to examine any and all imagery - high and low, art and non-art. . . . They must not restrict themselves to objects of a particular beauty or aesthetic value. Indeed, any kind of imagery may be found to offer up evidence of the visual construction of reality. . . .
Second, the study of visual culture must scrutinize visual practice as much as images themselves, asking what images do when they are put to use. If scholars engaged in this enterprise inquire what makes an image beautiful or why this image or that constitutes a masterpiece or a work of genius, they should do so with the purpose of investigating an artist’s or a work’s contribution to the experience of beauty, taste, value, or genius. No amount of social analysis can account fully for the existence of Michelangelo or Leonardo. They were unique creators of images that changed the way their contemporaries thought and felt and have continued to shape the history of art, artists, museums, feeling, and aesthetic value. But study of the critical, artistic, and popular reception of works by such artists as Michelangelo and Leonardo can shed important light on the meaning of these artists and their works for many different people. And the history of meaning-making has a great deal to do with how scholars as well as lay audiences today understand these artists and their achievements.
Third, scholars studying visual culture might properly focus their interpretative work on life worlds by examining images, practices, visual technologies, taste, and artistic style as constitutive of social relations. The task is to understand how artifacts contribute to the construction of a world. . . . Important methodological implications follow: ethnography and reception studies become productive forms of gathering information, since these move beyond the image as a closed and fixed meaning-event. . . .
Fourth, scholars may learn a great deal when they scrutinize the constituents of vision, that is, the structures of perception as a physiological process as well as the epistemological frameworks informing a system of visual representation. Vision is a socially and a biologically constructed operation, depending on the design of the human body and how it engages the interpretive devices developed by a culture in order to see intelligibly. . . . Seeing . . . operates on the foundation of covenants with images that establish the conditions for meaningful visual experience.
Finally, the scholar of visual culture seeks to regard images as evidence for explanation, not as epiphenomena.

\vspace{1em}
\textbf{Question 3}

Which one of the following best describes the word “epiphenomena” in the last sentence of the passage?

\begin{enumerate}[label=\Alph*.]
    \item Overarching collections of images.
    \item Visual phenomena of epic proportions.
    \item Phenomena amenable to analysis.
    \item Phenomena supplemental to the evidence.
[](https://cracku.in/3-which-one-of-the-following-best-describes-the-word-x-cat-2020-slot-2)
\end{enumerate}
\vspace{1em}
\textbf{Solution:}

This is a vocab question. Epiphenomena means "a secondary effect or byproduct". The option "Phenomena supplemental to the evidence", is the closest one. Hence it is the correct answer.

\end{problem}

\newpage
\begin{problem}{}{}
% Subject: varc
\textbf{Instructions}

Instructions   

The passage below is accompanied by a set of questions. Choose the best answer to each question.
The claims advanced here may be condensed into two assertions: [first, that visual] culture is what images, acts of seeing, and attendant intellectual, emotional, and perceptual sensibilities do to build, maintain, or transform the worlds in which people live. [And second, that the] study of visual culture is the analysis and interpretation of images and the ways of seeing (or gazes) that configure the agents, practices, conceptualities, and institutions that put images to work. . . .
Accordingly, the study of visual culture should be characterized by several concerns. First, scholars of visual culture need to examine any and all imagery - high and low, art and non-art. . . . They must not restrict themselves to objects of a particular beauty or aesthetic value. Indeed, any kind of imagery may be found to offer up evidence of the visual construction of reality. . . .
Second, the study of visual culture must scrutinize visual practice as much as images themselves, asking what images do when they are put to use. If scholars engaged in this enterprise inquire what makes an image beautiful or why this image or that constitutes a masterpiece or a work of genius, they should do so with the purpose of investigating an artist’s or a work’s contribution to the experience of beauty, taste, value, or genius. No amount of social analysis can account fully for the existence of Michelangelo or Leonardo. They were unique creators of images that changed the way their contemporaries thought and felt and have continued to shape the history of art, artists, museums, feeling, and aesthetic value. But study of the critical, artistic, and popular reception of works by such artists as Michelangelo and Leonardo can shed important light on the meaning of these artists and their works for many different people. And the history of meaning-making has a great deal to do with how scholars as well as lay audiences today understand these artists and their achievements.
Third, scholars studying visual culture might properly focus their interpretative work on life worlds by examining images, practices, visual technologies, taste, and artistic style as constitutive of social relations. The task is to understand how artifacts contribute to the construction of a world. . . . Important methodological implications follow: ethnography and reception studies become productive forms of gathering information, since these move beyond the image as a closed and fixed meaning-event. . . .
Fourth, scholars may learn a great deal when they scrutinize the constituents of vision, that is, the structures of perception as a physiological process as well as the epistemological frameworks informing a system of visual representation. Vision is a socially and a biologically constructed operation, depending on the design of the human body and how it engages the interpretive devices developed by a culture in order to see intelligibly. . . . Seeing . . . operates on the foundation of covenants with images that establish the conditions for meaningful visual experience.
Finally, the scholar of visual culture seeks to regard images as evidence for explanation, not as epiphenomena.

\vspace{1em}
\textbf{Question 4}

“No amount of social analysis can account fully for the existence of Michelangelo or Leonardo.” In light of the passage, which one of the following interpretations of this sentence is the most accurate?

\begin{enumerate}[label=\Alph*.]
    \item Socially existing beings cannot be analysed, unlike the art of Michelangelo or Leonardo which can.
    \item Social analytical accounts of people like Michelangelo or Leonardo cannot explain their genius.
    \item Michelangelo or Leonardo cannot be subjected to social analysis because of their genius.
    \item No analyses exist of Michelangelo’s or Leonardo’s social accounts.
[](https://cracku.in/4-no-amount-of-social-analysis-can-account-fully-for-x-cat-2020-slot-2)
\end{enumerate}
\vspace{1em}
\textbf{Solution:}

Let's look at the options one by one.
Option A states that ,"Socially existing beings cannot be analysed, unlike the art of Michelangelo or Leonardo which can." Twisted option. The excerpt from the passage states that "no amount of social analysis is enough for Michelangelo and Leonardo, because they were such vast artists" However other beings could be socially analysed because not everyone is like Michelangelo or Leonardo.
Option C states that ,"Michelangelo or Leonardo cannot be subjected to social analysis because of their genius." This is an entirely wrong option. These artists can be subjected to social analysis, but nothing will do justice to them.
Option D states that ,"No analyses exist of Michelangelo’s or Leonardo’s social accounts.". This is beyond the scope of the pasaage, as nothing has been mentioned about this.
Option B is the correct answer.

\end{problem}

\newpage
\begin{problem}{}{}
% Subject: varc
\textbf{Instructions}

Instructions   

The passage below is accompanied by a set of questions. Choose the best answer to each question.
The claims advanced here may be condensed into two assertions: [first, that visual] culture is what images, acts of seeing, and attendant intellectual, emotional, and perceptual sensibilities do to build, maintain, or transform the worlds in which people live. [And second, that the] study of visual culture is the analysis and interpretation of images and the ways of seeing (or gazes) that configure the agents, practices, conceptualities, and institutions that put images to work. . . .
Accordingly, the study of visual culture should be characterized by several concerns. First, scholars of visual culture need to examine any and all imagery - high and low, art and non-art. . . . They must not restrict themselves to objects of a particular beauty or aesthetic value. Indeed, any kind of imagery may be found to offer up evidence of the visual construction of reality. . . .
Second, the study of visual culture must scrutinize visual practice as much as images themselves, asking what images do when they are put to use. If scholars engaged in this enterprise inquire what makes an image beautiful or why this image or that constitutes a masterpiece or a work of genius, they should do so with the purpose of investigating an artist’s or a work’s contribution to the experience of beauty, taste, value, or genius. No amount of social analysis can account fully for the existence of Michelangelo or Leonardo. They were unique creators of images that changed the way their contemporaries thought and felt and have continued to shape the history of art, artists, museums, feeling, and aesthetic value. But study of the critical, artistic, and popular reception of works by such artists as Michelangelo and Leonardo can shed important light on the meaning of these artists and their works for many different people. And the history of meaning-making has a great deal to do with how scholars as well as lay audiences today understand these artists and their achievements.
Third, scholars studying visual culture might properly focus their interpretative work on life worlds by examining images, practices, visual technologies, taste, and artistic style as constitutive of social relations. The task is to understand how artifacts contribute to the construction of a world. . . . Important methodological implications follow: ethnography and reception studies become productive forms of gathering information, since these move beyond the image as a closed and fixed meaning-event. . . .
Fourth, scholars may learn a great deal when they scrutinize the constituents of vision, that is, the structures of perception as a physiological process as well as the epistemological frameworks informing a system of visual representation. Vision is a socially and a biologically constructed operation, depending on the design of the human body and how it engages the interpretive devices developed by a culture in order to see intelligibly. . . . Seeing . . . operates on the foundation of covenants with images that establish the conditions for meaningful visual experience.
Finally, the scholar of visual culture seeks to regard images as evidence for explanation, not as epiphenomena.

\vspace{1em}
\textbf{Question 5}

Which set of keywords below most closely captures the arguments of the passage?

\begin{enumerate}[label=\Alph*.]
    \item Visual Culture, Aesthetic Value, Lay Audience, Visual Experience.
    \item Visual Construction of Reality, Work of Genius, Ethnography, Epiphenomena.
    \item Imagery, Visual Practices, Lifeworlds, Structures of Perception.
    \item Scholars, Social Analysis, Michelangelo and Leonardo, Interpretive Devices.
[](https://cracku.in/5-which-set-of-keywords-below-most-closely-captures--x-cat-2020-slot-2)
\end{enumerate}
\vspace{1em}
\textbf{Solution:}

This is a fact based easy question. Read the passage carefully. Imagery can be inferred from the second paragraph. And from the subsequent paragraphs you can also infer Visual Practices, Lifeworlds and Structures of Perception(penultimate paragraph).
Hence Option C is the correct answer.

Instructions   

The passage below is accompanied by a set of questions. Choose the best answer to each question.
174 incidents of piracy were reported to the International Maritime Bureau last year, with Somali pirates responsible for only three. The rest ranged from the discreet theft of coils of rope in the Yellow Sea to the notoriously ferocious Nigerian gunmen attacking and hijacking oil tankers in the Gulf of Guinea, as well as armed robbery off Singapore and the Venezuelan coast and kidnapping in the Sundarbans in the Bay of Bengal. For [Dr. Peter] Lehr, an expert on modern-day piracy, the phenomenon’s history should be a source of instruction rather than entertainment, piracy past offering lessons for piracy present. . . .
But . . . where does piracy begin or end? According to St Augustine, a corsair captain once told Alexander the Great that in the forceful acquisition of power and wealth at sea, the difference between an emperor and a pirate was simply one of scale. By this logic, European empire-builders were the most successful pirates of all time. A more eclectic history might have included the conquistadors, Vasco da Gama and the East India Company. But Lehr sticks to the disorganized small fry, making comparisons with the renegades of today possible.
The main motive for piracy has always been a combination of need and greed. Why toil away as a starving peasant in the 16th century when a successful pirate made up to £4,000 on each raid? Anyone could turn to freebooting if the rewards were worth the risk . . . .
Increased globalisation has done more to encourage piracy than suppress it. European colonialism weakened delicate balances of power, leading to an influx of opportunists on the high seas. A rise in global shipping has meant rich pickings for freebooters. Lehr writes: “It quickly becomes clear that in those parts of the world that have not profited from globalisation and modernisation, and where abject poverty and the daily struggle for survival are still a reality, the root causes of piracy are still the same as they were a couple of hundred years ago.” . . .
Modern pirate prevention has failed. After the French yacht Le Gonant was ransomed for $2 million in 2008, opportunists from all over Somalia flocked to the coast for a piece of the action. . . . A consistent rule, even today, is there are never enough warships to patrol pirate-infested waters. Such ships are costly and only solve the problem temporarily; Somali piracy is bound to return as soon as the warships are withdrawn. Robot shipping, eliminating hostages, has been proposed as a possible solution; but as Lehr points out, this will only make pirates switch their targets to smaller carriers unable to afford the technology.
His advice isn’t new. Proposals to end illegal fishing are often advanced but they are difficult to enforce. Investment in local welfare put a halt to Malaysian piracy in the 1970s, but was dependent on money somehow filtering through a corrupt bureaucracy to the poor on the periphery. Diplomatic initiatives against piracy are plagued by mutual distrust: the Russians execute pirates, while the EU and US are reluctant to capture them for fear they’ll claim asylum.

\end{problem}

\newpage
\begin{problem}{}{}
% Subject: varc
\textbf{Instructions}

Instructions   

The passage below is accompanied by a set of questions. Choose the best answer to each question.
The claims advanced here may be condensed into two assertions: [first, that visual] culture is what images, acts of seeing, and attendant intellectual, emotional, and perceptual sensibilities do to build, maintain, or transform the worlds in which people live. [And second, that the] study of visual culture is the analysis and interpretation of images and the ways of seeing (or gazes) that configure the agents, practices, conceptualities, and institutions that put images to work. . . .
Accordingly, the study of visual culture should be characterized by several concerns. First, scholars of visual culture need to examine any and all imagery - high and low, art and non-art. . . . They must not restrict themselves to objects of a particular beauty or aesthetic value. Indeed, any kind of imagery may be found to offer up evidence of the visual construction of reality. . . .
Second, the study of visual culture must scrutinize visual practice as much as images themselves, asking what images do when they are put to use. If scholars engaged in this enterprise inquire what makes an image beautiful or why this image or that constitutes a masterpiece or a work of genius, they should do so with the purpose of investigating an artist’s or a work’s contribution to the experience of beauty, taste, value, or genius. No amount of social analysis can account fully for the existence of Michelangelo or Leonardo. They were unique creators of images that changed the way their contemporaries thought and felt and have continued to shape the history of art, artists, museums, feeling, and aesthetic value. But study of the critical, artistic, and popular reception of works by such artists as Michelangelo and Leonardo can shed important light on the meaning of these artists and their works for many different people. And the history of meaning-making has a great deal to do with how scholars as well as lay audiences today understand these artists and their achievements.
Third, scholars studying visual culture might properly focus their interpretative work on life worlds by examining images, practices, visual technologies, taste, and artistic style as constitutive of social relations. The task is to understand how artifacts contribute to the construction of a world. . . . Important methodological implications follow: ethnography and reception studies become productive forms of gathering information, since these move beyond the image as a closed and fixed meaning-event. . . .
Fourth, scholars may learn a great deal when they scrutinize the constituents of vision, that is, the structures of perception as a physiological process as well as the epistemological frameworks informing a system of visual representation. Vision is a socially and a biologically constructed operation, depending on the design of the human body and how it engages the interpretive devices developed by a culture in order to see intelligibly. . . . Seeing . . . operates on the foundation of covenants with images that establish the conditions for meaningful visual experience.
Finally, the scholar of visual culture seeks to regard images as evidence for explanation, not as epiphenomena.

\vspace{1em}
\textbf{Question 6}

“Why toil away as a starving peasant in the 16th century when a successful pirate made up to £4,000 on each raid?” In this sentence, the author’s tone can best be described as being:

\begin{enumerate}[label=\Alph*.]
    \item facetious, about the hardships of peasant life in medieval England.
    \item analytical, to explain the contrasts between peasant and pirate life in medieval England.
    \item ironic, about the reasons why so many took to piracy in medieval times.
    \item indignant, at the scale of wealth successful pirates could amass in medieval times.
[](https://cracku.in/6-why-toil-away-as-a-starving-peasant-in-the-16th-ce-x-cat-2020-slot-2)
\end{enumerate}
\vspace{1em}
\textbf{Solution:}

Option A states that the author is "facetious". Facetious means, treating serious issues with inappropriate humor. Author is nowhere mocking the hard life of peasants. On the contrary, he is actually admitting it and suggesting it as the reason why so many people became pirates in the first place. Hence this option is incorrect.
Option B states analytical. The author is nowhere contrasting the lives of peasant and pirates. He has not portrayed the lives of pirates in a pleasant light. Ultimately they also have to face surveillance and other dangers, thus the contrast in the option is wrong.
Indignant is completely wrong. The author is not angry at the pirates , for amassing huge wealth. Entirely out of context.
Option C talks about irony, which is actually correct. The author is ironical that an honest peasant has to toil day and night and still has to sleep empty stomach but on the other hand a pirate could easily amass twice the fortunes of peasant, without breaking much sweat. This entirely captures the essence of the mentioned line. Hence this answer is correct.

\end{problem}

\newpage
\begin{problem}{}{}
% Subject: varc
\textbf{Instructions}

Instructions   

The passage below is accompanied by a set of questions. Choose the best answer to each question.
The claims advanced here may be condensed into two assertions: [first, that visual] culture is what images, acts of seeing, and attendant intellectual, emotional, and perceptual sensibilities do to build, maintain, or transform the worlds in which people live. [And second, that the] study of visual culture is the analysis and interpretation of images and the ways of seeing (or gazes) that configure the agents, practices, conceptualities, and institutions that put images to work. . . .
Accordingly, the study of visual culture should be characterized by several concerns. First, scholars of visual culture need to examine any and all imagery - high and low, art and non-art. . . . They must not restrict themselves to objects of a particular beauty or aesthetic value. Indeed, any kind of imagery may be found to offer up evidence of the visual construction of reality. . . .
Second, the study of visual culture must scrutinize visual practice as much as images themselves, asking what images do when they are put to use. If scholars engaged in this enterprise inquire what makes an image beautiful or why this image or that constitutes a masterpiece or a work of genius, they should do so with the purpose of investigating an artist’s or a work’s contribution to the experience of beauty, taste, value, or genius. No amount of social analysis can account fully for the existence of Michelangelo or Leonardo. They were unique creators of images that changed the way their contemporaries thought and felt and have continued to shape the history of art, artists, museums, feeling, and aesthetic value. But study of the critical, artistic, and popular reception of works by such artists as Michelangelo and Leonardo can shed important light on the meaning of these artists and their works for many different people. And the history of meaning-making has a great deal to do with how scholars as well as lay audiences today understand these artists and their achievements.
Third, scholars studying visual culture might properly focus their interpretative work on life worlds by examining images, practices, visual technologies, taste, and artistic style as constitutive of social relations. The task is to understand how artifacts contribute to the construction of a world. . . . Important methodological implications follow: ethnography and reception studies become productive forms of gathering information, since these move beyond the image as a closed and fixed meaning-event. . . .
Fourth, scholars may learn a great deal when they scrutinize the constituents of vision, that is, the structures of perception as a physiological process as well as the epistemological frameworks informing a system of visual representation. Vision is a socially and a biologically constructed operation, depending on the design of the human body and how it engages the interpretive devices developed by a culture in order to see intelligibly. . . . Seeing . . . operates on the foundation of covenants with images that establish the conditions for meaningful visual experience.
Finally, the scholar of visual culture seeks to regard images as evidence for explanation, not as epiphenomena.

\vspace{1em}
\textbf{Question 7}

“A more eclectic history might have included the conquistadors, Vasco da Gama and the East India Company. But Lehr sticks to the disorganised small fry . . .” From this statement we can infer that the author believes that:

\begin{enumerate}[label=\Alph*.]
    \item Lehr does not assign adequate blame to empire builders for their past deeds.
    \item the disorganised piracy of today is no match for the organised piracy of the past.
    \item Vasco da Gama and the East India Company laid the ground for modern piracy.
    \item colonialism should be considered an organised form of piracy.
[](https://cracku.in/7-a-more-eclectic-history-might-have-included-the-co-x-cat-2020-slot-2)
\end{enumerate}
\vspace{1em}
\textbf{Solution:}

Option B and Option C are factually out of scope. Nothing has been said about who laid the groundwork for modern piracy. Neither has any comparisons made between the piracy of today and yesteryears,
In this statement , the author is obviously not assigning any blame to Vasco Da Gama and East India company. He just wants to highlight their roles in early pirates history.
Option D can be inferred from the passage.

\end{problem}

\newpage
\begin{problem}{}{}
% Subject: varc
\textbf{Instructions}

Instructions   

The passage below is accompanied by a set of questions. Choose the best answer to each question.
The claims advanced here may be condensed into two assertions: [first, that visual] culture is what images, acts of seeing, and attendant intellectual, emotional, and perceptual sensibilities do to build, maintain, or transform the worlds in which people live. [And second, that the] study of visual culture is the analysis and interpretation of images and the ways of seeing (or gazes) that configure the agents, practices, conceptualities, and institutions that put images to work. . . .
Accordingly, the study of visual culture should be characterized by several concerns. First, scholars of visual culture need to examine any and all imagery - high and low, art and non-art. . . . They must not restrict themselves to objects of a particular beauty or aesthetic value. Indeed, any kind of imagery may be found to offer up evidence of the visual construction of reality. . . .
Second, the study of visual culture must scrutinize visual practice as much as images themselves, asking what images do when they are put to use. If scholars engaged in this enterprise inquire what makes an image beautiful or why this image or that constitutes a masterpiece or a work of genius, they should do so with the purpose of investigating an artist’s or a work’s contribution to the experience of beauty, taste, value, or genius. No amount of social analysis can account fully for the existence of Michelangelo or Leonardo. They were unique creators of images that changed the way their contemporaries thought and felt and have continued to shape the history of art, artists, museums, feeling, and aesthetic value. But study of the critical, artistic, and popular reception of works by such artists as Michelangelo and Leonardo can shed important light on the meaning of these artists and their works for many different people. And the history of meaning-making has a great deal to do with how scholars as well as lay audiences today understand these artists and their achievements.
Third, scholars studying visual culture might properly focus their interpretative work on life worlds by examining images, practices, visual technologies, taste, and artistic style as constitutive of social relations. The task is to understand how artifacts contribute to the construction of a world. . . . Important methodological implications follow: ethnography and reception studies become productive forms of gathering information, since these move beyond the image as a closed and fixed meaning-event. . . .
Fourth, scholars may learn a great deal when they scrutinize the constituents of vision, that is, the structures of perception as a physiological process as well as the epistemological frameworks informing a system of visual representation. Vision is a socially and a biologically constructed operation, depending on the design of the human body and how it engages the interpretive devices developed by a culture in order to see intelligibly. . . . Seeing . . . operates on the foundation of covenants with images that establish the conditions for meaningful visual experience.
Finally, the scholar of visual culture seeks to regard images as evidence for explanation, not as epiphenomena.

\vspace{1em}
\textbf{Question 8}

We can deduce that the author believes that piracy can best be controlled in the long run:

\begin{enumerate}[label=\Alph*.]
    \item if we eliminate poverty and income disparities in affected regions.
    \item through international cooperation in enforcing stringent deterrents.
    \item through lucrative welfare schemes to improve the lives of people in affected regions.
    \item through the extensive deployment of technology to track ships and cargo.
[](https://cracku.in/8-we-can-deduce-that-the-author-believes-that-piracy-x-cat-2020-slot-2)
\end{enumerate}
\vspace{1em}
\textbf{Solution:}

Look at the 4th paragraph of the passage. " It quickly becomes clear that in those parts of the world that have not profited from globalisation and modernisation, and where abject poverty and the daily struggle for survival are still a reality, the root causes of piracy are still the same as they were a couple of hundred years ago." Hence to best control the problem of piracy one has to solve this root problem that is eliminate income disparities and poverty in affected regions. Hence 1 is the correct answer. Remaining options are solution but not the one which can eliminate piracy in the long run.

\end{problem}

\newpage
\begin{problem}{}{}
% Subject: varc
\textbf{Instructions}

Instructions   

The passage below is accompanied by a set of questions. Choose the best answer to each question.
The claims advanced here may be condensed into two assertions: [first, that visual] culture is what images, acts of seeing, and attendant intellectual, emotional, and perceptual sensibilities do to build, maintain, or transform the worlds in which people live. [And second, that the] study of visual culture is the analysis and interpretation of images and the ways of seeing (or gazes) that configure the agents, practices, conceptualities, and institutions that put images to work. . . .
Accordingly, the study of visual culture should be characterized by several concerns. First, scholars of visual culture need to examine any and all imagery - high and low, art and non-art. . . . They must not restrict themselves to objects of a particular beauty or aesthetic value. Indeed, any kind of imagery may be found to offer up evidence of the visual construction of reality. . . .
Second, the study of visual culture must scrutinize visual practice as much as images themselves, asking what images do when they are put to use. If scholars engaged in this enterprise inquire what makes an image beautiful or why this image or that constitutes a masterpiece or a work of genius, they should do so with the purpose of investigating an artist’s or a work’s contribution to the experience of beauty, taste, value, or genius. No amount of social analysis can account fully for the existence of Michelangelo or Leonardo. They were unique creators of images that changed the way their contemporaries thought and felt and have continued to shape the history of art, artists, museums, feeling, and aesthetic value. But study of the critical, artistic, and popular reception of works by such artists as Michelangelo and Leonardo can shed important light on the meaning of these artists and their works for many different people. And the history of meaning-making has a great deal to do with how scholars as well as lay audiences today understand these artists and their achievements.
Third, scholars studying visual culture might properly focus their interpretative work on life worlds by examining images, practices, visual technologies, taste, and artistic style as constitutive of social relations. The task is to understand how artifacts contribute to the construction of a world. . . . Important methodological implications follow: ethnography and reception studies become productive forms of gathering information, since these move beyond the image as a closed and fixed meaning-event. . . .
Fourth, scholars may learn a great deal when they scrutinize the constituents of vision, that is, the structures of perception as a physiological process as well as the epistemological frameworks informing a system of visual representation. Vision is a socially and a biologically constructed operation, depending on the design of the human body and how it engages the interpretive devices developed by a culture in order to see intelligibly. . . . Seeing . . . operates on the foundation of covenants with images that establish the conditions for meaningful visual experience.
Finally, the scholar of visual culture seeks to regard images as evidence for explanation, not as epiphenomena.

\vspace{1em}
\textbf{Question 9}

The author ascribes the rise in piracy today to all of the following factors EXCEPT:

\begin{enumerate}[label=\Alph*.]
    \item decreased surveillance of the high seas.
    \item the high rewards via ransoms for successful piracy attempts.
    \item the growth in international shipping with globalisation.
    \item colonialism’s disruption of historic ties among countries.
[](https://cracku.in/9-the-author-ascribes-the-rise-in-piracy-today-to-al-x-cat-2020-slot-2)
\end{enumerate}
\vspace{1em}
\textbf{Solution:}

The author in the passage that. "surveillance is there, but they are costly and will never be enough to cover the entire sea" He nowhere mentions that surveillance has decreased, actually it is not possible to map the entire sea. Hence this option is incorrect. Remaining all the options are correct.

Instructions   

The passage below is accompanied by a set of questions. Choose the best answer to each question.
Aggression is any behavior that is directed toward injuring, harming, or inflicting pain on another living being or group of beings. Generally, the victim(s) of aggression must wish to avoid such behavior in order for it to be considered true aggression. Aggression is also categorized according to its ultimate intent. Hostile aggression is an aggressive act that results from anger, and is intended to inflict pain or injury because of that anger. Instrumental aggression is an aggressive act that is regarded as a means to an end other than pain or injury. For example, an enemy combatant may be subjected to torture in order to extract useful intelligence, though those inflicting the torture may have no real feelings of anger or animosity toward their subject. The concept of aggression is very broad, and includes many categories of behavior (e.g., verbal aggression, street crime, child abuse, spouse abuse, group conflict, war, etc.). A number of theories and models of aggression have arisen to explain these diverse forms of behavior, and these theories/models tend to be categorized according to their specific focus. The most common system of categorization groups the various approaches to aggression into three separate areas, based upon the three key variables that are present whenever any aggressive act or set of acts is committed. The first variable is the aggressor him/herself. The second is the social situation or circumstance in which the aggressive act(s) occur. The third variable is the target or victim of aggression.
Regarding theories and research on the aggressor, the fundamental focus is on the factors that lead an individual (or group) to commit aggressive acts. At the most basic level, some argue that aggressive urges and actions are the result of inborn, biological factors. Sigmund Freud (1930) proposed that all individuals are born with a death instinct that predisposes us to a variety of aggressive behaviors, including suicide (self directed aggression) and mental illness (possibly due to an unhealthy or unnatural suppression of aggressive urges). Other influential perspectives supporting a biological basis for aggression conclude that humans evolved with an abnormally low neural inhibition of aggressive impulses (in comparison to other species), and that humans possess a powerful instinct for property accumulation and territorialism. It is proposed that this instinct accounts for hostile behaviors ranging from minor street crime to world wars. Hormonal factors also appear to play a significant role in fostering aggressive tendencies. For example, the hormone testosterone has been shown to increase aggressive behaviors when injected into animals. Men and women convicted of violent crimes also possess significantly higher levels of testosterone than men and women convicted of non violent crimes. Numerous studies comparing different age groups, racial/ethnic groups, and cultures also indicate that men, overall, are more likely to engage in a variety of aggressive behaviors (e.g., sexual assault, aggravated assault, etc.) than women. One explanation for higher levels of aggression in men is based on the assumption that, on average, men have higher levels of testosterone than women.

\end{problem}

\newpage
\begin{problem}{}{}
% Subject: varc
\textbf{Instructions}

Instructions   

The passage below is accompanied by a set of questions. Choose the best answer to each question.
The claims advanced here may be condensed into two assertions: [first, that visual] culture is what images, acts of seeing, and attendant intellectual, emotional, and perceptual sensibilities do to build, maintain, or transform the worlds in which people live. [And second, that the] study of visual culture is the analysis and interpretation of images and the ways of seeing (or gazes) that configure the agents, practices, conceptualities, and institutions that put images to work. . . .
Accordingly, the study of visual culture should be characterized by several concerns. First, scholars of visual culture need to examine any and all imagery - high and low, art and non-art. . . . They must not restrict themselves to objects of a particular beauty or aesthetic value. Indeed, any kind of imagery may be found to offer up evidence of the visual construction of reality. . . .
Second, the study of visual culture must scrutinize visual practice as much as images themselves, asking what images do when they are put to use. If scholars engaged in this enterprise inquire what makes an image beautiful or why this image or that constitutes a masterpiece or a work of genius, they should do so with the purpose of investigating an artist’s or a work’s contribution to the experience of beauty, taste, value, or genius. No amount of social analysis can account fully for the existence of Michelangelo or Leonardo. They were unique creators of images that changed the way their contemporaries thought and felt and have continued to shape the history of art, artists, museums, feeling, and aesthetic value. But study of the critical, artistic, and popular reception of works by such artists as Michelangelo and Leonardo can shed important light on the meaning of these artists and their works for many different people. And the history of meaning-making has a great deal to do with how scholars as well as lay audiences today understand these artists and their achievements.
Third, scholars studying visual culture might properly focus their interpretative work on life worlds by examining images, practices, visual technologies, taste, and artistic style as constitutive of social relations. The task is to understand how artifacts contribute to the construction of a world. . . . Important methodological implications follow: ethnography and reception studies become productive forms of gathering information, since these move beyond the image as a closed and fixed meaning-event. . . .
Fourth, scholars may learn a great deal when they scrutinize the constituents of vision, that is, the structures of perception as a physiological process as well as the epistemological frameworks informing a system of visual representation. Vision is a socially and a biologically constructed operation, depending on the design of the human body and how it engages the interpretive devices developed by a culture in order to see intelligibly. . . . Seeing . . . operates on the foundation of covenants with images that establish the conditions for meaningful visual experience.
Finally, the scholar of visual culture seeks to regard images as evidence for explanation, not as epiphenomena.

\vspace{1em}
\textbf{Question 10}

All of the following statements can be seen as logically implied by the arguments of the passage EXCEPT:

\begin{enumerate}[label=\Alph*.]
    \item a common theory of aggression is that it is the result of an abnormally low neural regulation of testosterone.
    \item if the alleged aggressive act is not sought to be avoided, it cannot really be considered aggression.
    \item Freud’s theory of aggression proposes that aggression results from the suppression of aggressive urges.
    \item the Freudian theory of suicide as self-inflicted aggression implies that an aggressive act need not be sought to be avoided in order for it to be considered aggression.
[](https://cracku.in/10-all-of-the-following-statements-can-be-seen-as-log-x-cat-2020-slot-2)
\end{enumerate}
\vspace{1em}
\textbf{Solution:}

Look at the last paragraph of the passage," Other influential perspectives supporting a biological basis for aggression conclude that humans evolved with an abnormally low neural inhibition of aggressive impulses". While the first option talks about neural regulation of testosterone, which is factually incorrect. Hence this is the answer.

\end{problem}

\newpage
\begin{problem}{}{}
% Subject: varc
\textbf{Instructions}

Instructions   

The passage below is accompanied by a set of questions. Choose the best answer to each question.
The claims advanced here may be condensed into two assertions: [first, that visual] culture is what images, acts of seeing, and attendant intellectual, emotional, and perceptual sensibilities do to build, maintain, or transform the worlds in which people live. [And second, that the] study of visual culture is the analysis and interpretation of images and the ways of seeing (or gazes) that configure the agents, practices, conceptualities, and institutions that put images to work. . . .
Accordingly, the study of visual culture should be characterized by several concerns. First, scholars of visual culture need to examine any and all imagery - high and low, art and non-art. . . . They must not restrict themselves to objects of a particular beauty or aesthetic value. Indeed, any kind of imagery may be found to offer up evidence of the visual construction of reality. . . .
Second, the study of visual culture must scrutinize visual practice as much as images themselves, asking what images do when they are put to use. If scholars engaged in this enterprise inquire what makes an image beautiful or why this image or that constitutes a masterpiece or a work of genius, they should do so with the purpose of investigating an artist’s or a work’s contribution to the experience of beauty, taste, value, or genius. No amount of social analysis can account fully for the existence of Michelangelo or Leonardo. They were unique creators of images that changed the way their contemporaries thought and felt and have continued to shape the history of art, artists, museums, feeling, and aesthetic value. But study of the critical, artistic, and popular reception of works by such artists as Michelangelo and Leonardo can shed important light on the meaning of these artists and their works for many different people. And the history of meaning-making has a great deal to do with how scholars as well as lay audiences today understand these artists and their achievements.
Third, scholars studying visual culture might properly focus their interpretative work on life worlds by examining images, practices, visual technologies, taste, and artistic style as constitutive of social relations. The task is to understand how artifacts contribute to the construction of a world. . . . Important methodological implications follow: ethnography and reception studies become productive forms of gathering information, since these move beyond the image as a closed and fixed meaning-event. . . .
Fourth, scholars may learn a great deal when they scrutinize the constituents of vision, that is, the structures of perception as a physiological process as well as the epistemological frameworks informing a system of visual representation. Vision is a socially and a biologically constructed operation, depending on the design of the human body and how it engages the interpretive devices developed by a culture in order to see intelligibly. . . . Seeing . . . operates on the foundation of covenants with images that establish the conditions for meaningful visual experience.
Finally, the scholar of visual culture seeks to regard images as evidence for explanation, not as epiphenomena.

\vspace{1em}
\textbf{Question 11}

The author identifies three essential factors according to which theories of aggression are most commonly categorised. Which of the following options is closest to the factors identified by the author?

\begin{enumerate}[label=\Alph*.]
    \item Aggressor - Circumstances of aggression - Victim.
    \item Psychologically - Sociologically - Medically.
    \item Hostile - Instrumental - Hormonal.
    \item Extreme - Moderate - Mild.
[](https://cracku.in/11-the-author-identifies-three-essential-factors-acco-x-cat-2020-slot-2)
\end{enumerate}
\vspace{1em}
\textbf{Solution:}

The author explains that theories of aggression are often categorized based on three main factors: the aggressor (the person or group committing the aggression), the circumstances surrounding the aggressive act, and the victim (the person or group being targeted). These factors help to understand and explain different types of aggression.

\end{problem}

\newpage
\begin{problem}{}{}
% Subject: varc
\textbf{Instructions}

Instructions   

The passage below is accompanied by a set of questions. Choose the best answer to each question.
The claims advanced here may be condensed into two assertions: [first, that visual] culture is what images, acts of seeing, and attendant intellectual, emotional, and perceptual sensibilities do to build, maintain, or transform the worlds in which people live. [And second, that the] study of visual culture is the analysis and interpretation of images and the ways of seeing (or gazes) that configure the agents, practices, conceptualities, and institutions that put images to work. . . .
Accordingly, the study of visual culture should be characterized by several concerns. First, scholars of visual culture need to examine any and all imagery - high and low, art and non-art. . . . They must not restrict themselves to objects of a particular beauty or aesthetic value. Indeed, any kind of imagery may be found to offer up evidence of the visual construction of reality. . . .
Second, the study of visual culture must scrutinize visual practice as much as images themselves, asking what images do when they are put to use. If scholars engaged in this enterprise inquire what makes an image beautiful or why this image or that constitutes a masterpiece or a work of genius, they should do so with the purpose of investigating an artist’s or a work’s contribution to the experience of beauty, taste, value, or genius. No amount of social analysis can account fully for the existence of Michelangelo or Leonardo. They were unique creators of images that changed the way their contemporaries thought and felt and have continued to shape the history of art, artists, museums, feeling, and aesthetic value. But study of the critical, artistic, and popular reception of works by such artists as Michelangelo and Leonardo can shed important light on the meaning of these artists and their works for many different people. And the history of meaning-making has a great deal to do with how scholars as well as lay audiences today understand these artists and their achievements.
Third, scholars studying visual culture might properly focus their interpretative work on life worlds by examining images, practices, visual technologies, taste, and artistic style as constitutive of social relations. The task is to understand how artifacts contribute to the construction of a world. . . . Important methodological implications follow: ethnography and reception studies become productive forms of gathering information, since these move beyond the image as a closed and fixed meaning-event. . . .
Fourth, scholars may learn a great deal when they scrutinize the constituents of vision, that is, the structures of perception as a physiological process as well as the epistemological frameworks informing a system of visual representation. Vision is a socially and a biologically constructed operation, depending on the design of the human body and how it engages the interpretive devices developed by a culture in order to see intelligibly. . . . Seeing . . . operates on the foundation of covenants with images that establish the conditions for meaningful visual experience.
Finally, the scholar of visual culture seeks to regard images as evidence for explanation, not as epiphenomena.

\vspace{1em}
\textbf{Question 12}

“[A]n enemy combatant may be subjected to torture in order to extract useful intelligence, though those inflicting the torture may have no real feelings of anger or animosity toward their subject.” Which one of the following best explicates the larger point being made by the author here?

\begin{enumerate}[label=\Alph*.]
    \item In certain kinds of aggression, inflicting pain is not the objective, and is no more than a utilitarian means to achieve another end.
    \item Information revealed by subjecting an enemy combatant to torture is not always reliable because of the animosity involved.
    \item When an enemy combatant refuses to reveal information, the use of torture can sometimes involve real feelings of hostility.
    \item The use of torture to extract information is most effective when the torturer is not emotionally involved in the torture.
[](https://cracku.in/12-an-enemy-combatant-may-be-subjected-to-torture-in--x-cat-2020-slot-2)
\end{enumerate}
\vspace{1em}
\textbf{Solution:}

Let's look at the options one by one.
Option B is a distorted option. The sentence is not about the reliability of the information. Its about the inherent aggression involved.
Option C is also incorrect. Nothing has been mentioned about when the enemy refuses to reveal information. Hence this is beyond the scope of the passage.
On similar lines option D is also incorrect. Nothing has been said about the most effective method to extract information.
Option A correctly captures the essence. Sometimes the aggressor has no intention to inflict pain, its just a means to an utilitarian end. Hence this is the correct option.

\end{problem}

\newpage
\begin{problem}{}{}
% Subject: varc
\textbf{Instructions}

Instructions   

The passage below is accompanied by a set of questions. Choose the best answer to each question.
The claims advanced here may be condensed into two assertions: [first, that visual] culture is what images, acts of seeing, and attendant intellectual, emotional, and perceptual sensibilities do to build, maintain, or transform the worlds in which people live. [And second, that the] study of visual culture is the analysis and interpretation of images and the ways of seeing (or gazes) that configure the agents, practices, conceptualities, and institutions that put images to work. . . .
Accordingly, the study of visual culture should be characterized by several concerns. First, scholars of visual culture need to examine any and all imagery - high and low, art and non-art. . . . They must not restrict themselves to objects of a particular beauty or aesthetic value. Indeed, any kind of imagery may be found to offer up evidence of the visual construction of reality. . . .
Second, the study of visual culture must scrutinize visual practice as much as images themselves, asking what images do when they are put to use. If scholars engaged in this enterprise inquire what makes an image beautiful or why this image or that constitutes a masterpiece or a work of genius, they should do so with the purpose of investigating an artist’s or a work’s contribution to the experience of beauty, taste, value, or genius. No amount of social analysis can account fully for the existence of Michelangelo or Leonardo. They were unique creators of images that changed the way their contemporaries thought and felt and have continued to shape the history of art, artists, museums, feeling, and aesthetic value. But study of the critical, artistic, and popular reception of works by such artists as Michelangelo and Leonardo can shed important light on the meaning of these artists and their works for many different people. And the history of meaning-making has a great deal to do with how scholars as well as lay audiences today understand these artists and their achievements.
Third, scholars studying visual culture might properly focus their interpretative work on life worlds by examining images, practices, visual technologies, taste, and artistic style as constitutive of social relations. The task is to understand how artifacts contribute to the construction of a world. . . . Important methodological implications follow: ethnography and reception studies become productive forms of gathering information, since these move beyond the image as a closed and fixed meaning-event. . . .
Fourth, scholars may learn a great deal when they scrutinize the constituents of vision, that is, the structures of perception as a physiological process as well as the epistemological frameworks informing a system of visual representation. Vision is a socially and a biologically constructed operation, depending on the design of the human body and how it engages the interpretive devices developed by a culture in order to see intelligibly. . . . Seeing . . . operates on the foundation of covenants with images that establish the conditions for meaningful visual experience.
Finally, the scholar of visual culture seeks to regard images as evidence for explanation, not as epiphenomena.

\vspace{1em}
\textbf{Question 13}

The author discusses all of the following arguments in the passage EXCEPT that:

\begin{enumerate}[label=\Alph*.]
    \item men in general are believed to be more hormonally driven to exhibit violence than women.
    \item several studies indicate that aggression may have roots in the biological condition of humanity.
    \item aggression in most societies is kept under control through moderating the death instinct identified by Freud.
    \item the nature of aggression can vary depending on several factors, including intent.
[](https://cracku.in/13-the-author-discusses-all-of-the-following-argument-x-cat-2020-slot-2)
\end{enumerate}
\vspace{1em}
\textbf{Solution:}

Option C is a twisted option. It distorts what is given in the passage. The passage talks about the identification of death instinct by Freud. It doesn't talk about its moderation and keeping the resulting aggression in check. Hence this option is completely out of context and is incorrect.

Instructions   

**The passage below is accompanied by a set of questions. Choose the best answer to each question.**
In a low-carbon world, renewable energy technologies are hot business. For investors looking to redirect funds, wind turbines and solar panels, among other technologies, seem a straightforward choice. But renewables need to be further scrutinized before being championed as forging a path toward a low-carbon future. Both the direct and indirect impacts of renewable energy must be examined to ensure that a climate-smart future does not intensify social and environmental harm. As renewable energy production requires land, water, and labor, among other inputs, it imposes costs on people and the environment. Hydropower projects, for instance, have led to community dispossession and exclusion . . . Renewable energy supply chains are also intertwined with mining, and their technologies contribute to growing levels of electronic waste . Furthermore, although renewable energy can be produced and distributed through small-scale, local systems, such an approach might not generate the high returns on investment needed to attract capital.
Although an emerging sector, renewables are enmeshed in long-standing resource extraction through their dependence on minerals and metals . . . Scholars document the negative consequences of mining . . . even for mining operations that commit to socially responsible practices[:] “many of the world’s largest reservoirs of minerals like cobalt, copper, lithium, [and] rare earth minerals”—the ones needed for renewable technologies—“are found in fragile states and under communities of marginalized peoples in Africa, Asia, and Latin America.” Since the demand for metals and minerals will increase substantially in a renewable-powered future . . . this intensification could exacerbate the existing consequences of extractive activities.
Among the connections between climate change and waste, O’Neill . . . highlights that “devices developed to reduce our carbon footprint, such as lithium batteries for hybrid and electric cars or solar panels[,] become potentially dangerous electronic waste at the end of their productive life.” The disposal of toxic waste has long perpetuated social injustice through the flows of waste to the Global South and to marginalized communities in the Global North . . .
While renewable energy is a more recent addition to financial portfolios, investments in the sector must be considered in light of our understanding of capital accumulation. As agricultural finance reveals, the concentration of control of corporate activity facilitates profit generation. For some climate activists, the promise of renewables rests on their ability not only to reduce emissions but also to provide distributed, democratized access to energy . . . But Burke and Stephens . . . caution that “renewable energy systems offer a possibility but not a certainty for more democratic energy futures.” Small-scale, distributed forms of energy are only highly profitable to institutional investors if control is consolidated somewhere in the financial chain. Renewable energy can be produced at the household or neighborhood level. However, such small-scale, localized production is unlikely to generate high returns for investors. For financial growth to be sustained and expanded by the renewable sector, production and trade in renewable energy technologies will need to be highly concentrated, and large asset management firms will likely drive those developments.

\end{problem}

\newpage
\begin{problem}{}{}
% Subject: varc
\textbf{Instructions}

Instructions   

The passage below is accompanied by a set of questions. Choose the best answer to each question.
The claims advanced here may be condensed into two assertions: [first, that visual] culture is what images, acts of seeing, and attendant intellectual, emotional, and perceptual sensibilities do to build, maintain, or transform the worlds in which people live. [And second, that the] study of visual culture is the analysis and interpretation of images and the ways of seeing (or gazes) that configure the agents, practices, conceptualities, and institutions that put images to work. . . .
Accordingly, the study of visual culture should be characterized by several concerns. First, scholars of visual culture need to examine any and all imagery - high and low, art and non-art. . . . They must not restrict themselves to objects of a particular beauty or aesthetic value. Indeed, any kind of imagery may be found to offer up evidence of the visual construction of reality. . . .
Second, the study of visual culture must scrutinize visual practice as much as images themselves, asking what images do when they are put to use. If scholars engaged in this enterprise inquire what makes an image beautiful or why this image or that constitutes a masterpiece or a work of genius, they should do so with the purpose of investigating an artist’s or a work’s contribution to the experience of beauty, taste, value, or genius. No amount of social analysis can account fully for the existence of Michelangelo or Leonardo. They were unique creators of images that changed the way their contemporaries thought and felt and have continued to shape the history of art, artists, museums, feeling, and aesthetic value. But study of the critical, artistic, and popular reception of works by such artists as Michelangelo and Leonardo can shed important light on the meaning of these artists and their works for many different people. And the history of meaning-making has a great deal to do with how scholars as well as lay audiences today understand these artists and their achievements.
Third, scholars studying visual culture might properly focus their interpretative work on life worlds by examining images, practices, visual technologies, taste, and artistic style as constitutive of social relations. The task is to understand how artifacts contribute to the construction of a world. . . . Important methodological implications follow: ethnography and reception studies become productive forms of gathering information, since these move beyond the image as a closed and fixed meaning-event. . . .
Fourth, scholars may learn a great deal when they scrutinize the constituents of vision, that is, the structures of perception as a physiological process as well as the epistemological frameworks informing a system of visual representation. Vision is a socially and a biologically constructed operation, depending on the design of the human body and how it engages the interpretive devices developed by a culture in order to see intelligibly. . . . Seeing . . . operates on the foundation of covenants with images that establish the conditions for meaningful visual experience.
Finally, the scholar of visual culture seeks to regard images as evidence for explanation, not as epiphenomena.

\vspace{1em}
\textbf{Question 14}

Which one of the following statements best captures the main argument of the last paragraph of the passage?

\begin{enumerate}[label=\Alph*.]
    \item The development of the renewable energy sector is a double-edged sword.
    \item Renewable energy systems are not democratic unless they are corporate-controlled.
    \item Renewable energy produced at the household or neighbourhood level is more efficient than mass-produced forms of energy.
    \item Most forms of renewable energy are not profitable investments for institutional investors.
[](https://cracku.in/14-which-one-of-the-following-statements-best-capture-x-cat-2020-slot-2)
\end{enumerate}
\vspace{1em}
\textbf{Solution:}

The last paragraph talks about how renewable energy can be produced at local or household level but it also talks about the roadblocks in terms of financing. Hence it portrays two sides of the coin, talking about both the positives and negatives. Thus we can fairly conclude that renewable energy is a double edged sword. Option1 is the correct answer.

\end{problem}

\newpage
\begin{problem}{}{}
% Subject: varc
\textbf{Instructions}

Instructions   

The passage below is accompanied by a set of questions. Choose the best answer to each question.
The claims advanced here may be condensed into two assertions: [first, that visual] culture is what images, acts of seeing, and attendant intellectual, emotional, and perceptual sensibilities do to build, maintain, or transform the worlds in which people live. [And second, that the] study of visual culture is the analysis and interpretation of images and the ways of seeing (or gazes) that configure the agents, practices, conceptualities, and institutions that put images to work. . . .
Accordingly, the study of visual culture should be characterized by several concerns. First, scholars of visual culture need to examine any and all imagery - high and low, art and non-art. . . . They must not restrict themselves to objects of a particular beauty or aesthetic value. Indeed, any kind of imagery may be found to offer up evidence of the visual construction of reality. . . .
Second, the study of visual culture must scrutinize visual practice as much as images themselves, asking what images do when they are put to use. If scholars engaged in this enterprise inquire what makes an image beautiful or why this image or that constitutes a masterpiece or a work of genius, they should do so with the purpose of investigating an artist’s or a work’s contribution to the experience of beauty, taste, value, or genius. No amount of social analysis can account fully for the existence of Michelangelo or Leonardo. They were unique creators of images that changed the way their contemporaries thought and felt and have continued to shape the history of art, artists, museums, feeling, and aesthetic value. But study of the critical, artistic, and popular reception of works by such artists as Michelangelo and Leonardo can shed important light on the meaning of these artists and their works for many different people. And the history of meaning-making has a great deal to do with how scholars as well as lay audiences today understand these artists and their achievements.
Third, scholars studying visual culture might properly focus their interpretative work on life worlds by examining images, practices, visual technologies, taste, and artistic style as constitutive of social relations. The task is to understand how artifacts contribute to the construction of a world. . . . Important methodological implications follow: ethnography and reception studies become productive forms of gathering information, since these move beyond the image as a closed and fixed meaning-event. . . .
Fourth, scholars may learn a great deal when they scrutinize the constituents of vision, that is, the structures of perception as a physiological process as well as the epistemological frameworks informing a system of visual representation. Vision is a socially and a biologically constructed operation, depending on the design of the human body and how it engages the interpretive devices developed by a culture in order to see intelligibly. . . . Seeing . . . operates on the foundation of covenants with images that establish the conditions for meaningful visual experience.
Finally, the scholar of visual culture seeks to regard images as evidence for explanation, not as epiphenomena.

\vspace{1em}
\textbf{Question 15}

Which one of the following statements, if true, could be an accurate inference from the first paragraph of the passage?

\begin{enumerate}[label=\Alph*.]
    \item The author has reservations about the consequences of renewable energy systems
    \item The author has reservations about the consequences of non-renewable energy systems.
    \item The author does not think renewable energy systems can be as efficient as non-renewable energy systems.
    \item The author’s only reservation is about the profitability of renewable energy systems.
[](https://cracku.in/15-which-one-of-the-following-statements-if-true-coul-x-cat-2020-slot-2)
\end{enumerate}
\vspace{1em}
\textbf{Solution:}

Look at these lines from the first paragraph of the passage, "But renewables need to be further scrutinized before being championed as forging a path toward a low-carbon future. Both the direct and indirect impacts of renewable energy must be examined to ensure that a climate-smart future does not intensify social and environmental harm. As renewable energy production requires land, water, and labor, among other inputs, it imposes costs on people and the environment."
These lines clearly indicate that the author has reservations about the renewable energy systems.
Option (2) is factually incorrect. Option 3 is beyond the scope of the passage. No comparison has been made as such. Option (4) is partially correct. The author has reservations about profitability and various other factors too which is accurately captured by the first option.

\end{problem}

\newpage
\begin{problem}{}{}
% Subject: varc
\textbf{Instructions}

Instructions   

The passage below is accompanied by a set of questions. Choose the best answer to each question.
The claims advanced here may be condensed into two assertions: [first, that visual] culture is what images, acts of seeing, and attendant intellectual, emotional, and perceptual sensibilities do to build, maintain, or transform the worlds in which people live. [And second, that the] study of visual culture is the analysis and interpretation of images and the ways of seeing (or gazes) that configure the agents, practices, conceptualities, and institutions that put images to work. . . .
Accordingly, the study of visual culture should be characterized by several concerns. First, scholars of visual culture need to examine any and all imagery - high and low, art and non-art. . . . They must not restrict themselves to objects of a particular beauty or aesthetic value. Indeed, any kind of imagery may be found to offer up evidence of the visual construction of reality. . . .
Second, the study of visual culture must scrutinize visual practice as much as images themselves, asking what images do when they are put to use. If scholars engaged in this enterprise inquire what makes an image beautiful or why this image or that constitutes a masterpiece or a work of genius, they should do so with the purpose of investigating an artist’s or a work’s contribution to the experience of beauty, taste, value, or genius. No amount of social analysis can account fully for the existence of Michelangelo or Leonardo. They were unique creators of images that changed the way their contemporaries thought and felt and have continued to shape the history of art, artists, museums, feeling, and aesthetic value. But study of the critical, artistic, and popular reception of works by such artists as Michelangelo and Leonardo can shed important light on the meaning of these artists and their works for many different people. And the history of meaning-making has a great deal to do with how scholars as well as lay audiences today understand these artists and their achievements.
Third, scholars studying visual culture might properly focus their interpretative work on life worlds by examining images, practices, visual technologies, taste, and artistic style as constitutive of social relations. The task is to understand how artifacts contribute to the construction of a world. . . . Important methodological implications follow: ethnography and reception studies become productive forms of gathering information, since these move beyond the image as a closed and fixed meaning-event. . . .
Fourth, scholars may learn a great deal when they scrutinize the constituents of vision, that is, the structures of perception as a physiological process as well as the epistemological frameworks informing a system of visual representation. Vision is a socially and a biologically constructed operation, depending on the design of the human body and how it engages the interpretive devices developed by a culture in order to see intelligibly. . . . Seeing . . . operates on the foundation of covenants with images that establish the conditions for meaningful visual experience.
Finally, the scholar of visual culture seeks to regard images as evidence for explanation, not as epiphenomena.

\vspace{1em}
\textbf{Question 16}

Which one of the following statements, if false, could be seen as best supporting the arguments in the passage?

\begin{enumerate}[label=\Alph*.]
    \item Renewable energy systems are as expensive as non-renewable energy systems.
    \item Renewable energy systems have little or no environmental impact.
    \item Renewable energy systems are not as profitable as non-renewable energy systems.
    \item The production and distribution of renewable energy through small-scale, local systems is not economically sustainable.
[](https://cracku.in/16-which-one-of-the-following-statements-if-false-cou-x-cat-2020-slot-2)
\end{enumerate}
\vspace{1em}
\textbf{Solution:}

Here we have to look at those options which if false will support the author's arguments. Look at the option 2. It states that renewable energy have little or no environmental impact. Its negation will be renewable energy have considerable environmental impact, this is what the author states throughout the passage, about the harmful impact of renewable energy. Hence option 2 is correct.

\end{problem}

\newpage
\begin{problem}{}{}
% Subject: varc
\textbf{Instructions}

Instructions   

The passage below is accompanied by a set of questions. Choose the best answer to each question.
The claims advanced here may be condensed into two assertions: [first, that visual] culture is what images, acts of seeing, and attendant intellectual, emotional, and perceptual sensibilities do to build, maintain, or transform the worlds in which people live. [And second, that the] study of visual culture is the analysis and interpretation of images and the ways of seeing (or gazes) that configure the agents, practices, conceptualities, and institutions that put images to work. . . .
Accordingly, the study of visual culture should be characterized by several concerns. First, scholars of visual culture need to examine any and all imagery - high and low, art and non-art. . . . They must not restrict themselves to objects of a particular beauty or aesthetic value. Indeed, any kind of imagery may be found to offer up evidence of the visual construction of reality. . . .
Second, the study of visual culture must scrutinize visual practice as much as images themselves, asking what images do when they are put to use. If scholars engaged in this enterprise inquire what makes an image beautiful or why this image or that constitutes a masterpiece or a work of genius, they should do so with the purpose of investigating an artist’s or a work’s contribution to the experience of beauty, taste, value, or genius. No amount of social analysis can account fully for the existence of Michelangelo or Leonardo. They were unique creators of images that changed the way their contemporaries thought and felt and have continued to shape the history of art, artists, museums, feeling, and aesthetic value. But study of the critical, artistic, and popular reception of works by such artists as Michelangelo and Leonardo can shed important light on the meaning of these artists and their works for many different people. And the history of meaning-making has a great deal to do with how scholars as well as lay audiences today understand these artists and their achievements.
Third, scholars studying visual culture might properly focus their interpretative work on life worlds by examining images, practices, visual technologies, taste, and artistic style as constitutive of social relations. The task is to understand how artifacts contribute to the construction of a world. . . . Important methodological implications follow: ethnography and reception studies become productive forms of gathering information, since these move beyond the image as a closed and fixed meaning-event. . . .
Fourth, scholars may learn a great deal when they scrutinize the constituents of vision, that is, the structures of perception as a physiological process as well as the epistemological frameworks informing a system of visual representation. Vision is a socially and a biologically constructed operation, depending on the design of the human body and how it engages the interpretive devices developed by a culture in order to see intelligibly. . . . Seeing . . . operates on the foundation of covenants with images that establish the conditions for meaningful visual experience.
Finally, the scholar of visual culture seeks to regard images as evidence for explanation, not as epiphenomena.

\vspace{1em}
\textbf{Question 17}

All of the following statements, if true, could be seen as supporting the arguments in the passage, EXCEPT:

\begin{enumerate}[label=\Alph*.]
    \item The example of agricultural finance helps us to see how to concentrate corporate activity in the renewable energy sector.
    \item One reason for the perpetuation of social injustice lies in the problem of the disposal of toxic waste.
    \item Marginalised people in Africa, Asia and Latin America have often been the main sufferers of corporate mineral extraction projects
    \item The possible negative impacts of renewable energy need to be studied before it can be offered as a financial investment opportunity.
[](https://cracku.in/17-all-of-the-following-statements-if-true-could-be-s-x-cat-2020-slot-2)
\end{enumerate}
\vspace{1em}
\textbf{Solution:}

Look at the last line of the passage," For financial growth to be sustained and expanded by the renewable sector, production and trade in renewable energy technologies will need to be highly concentrated, and large asset management firms will likely drive those developments " Nowhere it is mentioned that the negative impacts of renewable energy needs to be studied , for it to be financially viable. Hence option 4 is incorrect.

\end{problem}

\newpage
\begin{problem}{}{}
% Subject: varc
\textbf{Instructions}

Instructions   

The passage below is accompanied by a set of questions. Choose the best answer to each question.
The claims advanced here may be condensed into two assertions: [first, that visual] culture is what images, acts of seeing, and attendant intellectual, emotional, and perceptual sensibilities do to build, maintain, or transform the worlds in which people live. [And second, that the] study of visual culture is the analysis and interpretation of images and the ways of seeing (or gazes) that configure the agents, practices, conceptualities, and institutions that put images to work. . . .
Accordingly, the study of visual culture should be characterized by several concerns. First, scholars of visual culture need to examine any and all imagery - high and low, art and non-art. . . . They must not restrict themselves to objects of a particular beauty or aesthetic value. Indeed, any kind of imagery may be found to offer up evidence of the visual construction of reality. . . .
Second, the study of visual culture must scrutinize visual practice as much as images themselves, asking what images do when they are put to use. If scholars engaged in this enterprise inquire what makes an image beautiful or why this image or that constitutes a masterpiece or a work of genius, they should do so with the purpose of investigating an artist’s or a work’s contribution to the experience of beauty, taste, value, or genius. No amount of social analysis can account fully for the existence of Michelangelo or Leonardo. They were unique creators of images that changed the way their contemporaries thought and felt and have continued to shape the history of art, artists, museums, feeling, and aesthetic value. But study of the critical, artistic, and popular reception of works by such artists as Michelangelo and Leonardo can shed important light on the meaning of these artists and their works for many different people. And the history of meaning-making has a great deal to do with how scholars as well as lay audiences today understand these artists and their achievements.
Third, scholars studying visual culture might properly focus their interpretative work on life worlds by examining images, practices, visual technologies, taste, and artistic style as constitutive of social relations. The task is to understand how artifacts contribute to the construction of a world. . . . Important methodological implications follow: ethnography and reception studies become productive forms of gathering information, since these move beyond the image as a closed and fixed meaning-event. . . .
Fourth, scholars may learn a great deal when they scrutinize the constituents of vision, that is, the structures of perception as a physiological process as well as the epistemological frameworks informing a system of visual representation. Vision is a socially and a biologically constructed operation, depending on the design of the human body and how it engages the interpretive devices developed by a culture in order to see intelligibly. . . . Seeing . . . operates on the foundation of covenants with images that establish the conditions for meaningful visual experience.
Finally, the scholar of visual culture seeks to regard images as evidence for explanation, not as epiphenomena.

\vspace{1em}
\textbf{Question 18}

Based on the passage, we can infer that the author would be most supportive of which one of the following practices?

\begin{enumerate}[label=\Alph*.]
    \item More stringent global policies and regulations to ensure a more just system of toxic waste disposal.
    \item The study of the coexistence of marginalised people with their environments.
    \item Encouragement for the development of more environment-friendly carbon-based fuels.
    \item The localised, small-scale development of renewable energy systems.
[](https://cracku.in/18-based-on-the-passage-we-can-infer-that-the-author--x-cat-2020-slot-2)
\end{enumerate}
\vspace{1em}
\textbf{Solution:}

Let's look at the options one by one. Option 2 states "The study of the coexistence of marginalised people with their environments." This passage is about the impact of renewable energy on marginalized people, not about their coexistence. Hence ignore this option.
Option 3 states "Encouragement for the development of more environment-friendly carbon-based fuels". In the third paragraph, the author talks about reducing carbon footprint. Hence he will never support carbon based fuels. This option is incorrect.
Option 4 is completely negated by the last paragraph. The author talks about the financial viability of the localised small scale production. Hence this option too is incorrect.
Option 1 is completely correct."many of the world’s largest reservoirs of minerals like cobalt, copper, lithium, [and] rare earth minerals”—the ones needed for renewable technologies—“are found in fragile states and under communities of marginalized peoples in Africa, Asia, and Latin America.” Since the demand for metals and minerals will increase substantially in a renewable-powered future . . . this intensification could exacerbate the existing consequences of extractive activities." 
This intensification could exacerbate the existing consequences, i.e. if propar disposal of toxic materials is not done it can make the conditions worse. Option 1 is the correct answer.

Instructions   

For the following questions answer them individually

\end{problem}

\newpage
\begin{problem}{}{}
% Subject: varc
\textbf{Instructions}

Instructions   

The passage below is accompanied by a set of questions. Choose the best answer to each question.
The claims advanced here may be condensed into two assertions: [first, that visual] culture is what images, acts of seeing, and attendant intellectual, emotional, and perceptual sensibilities do to build, maintain, or transform the worlds in which people live. [And second, that the] study of visual culture is the analysis and interpretation of images and the ways of seeing (or gazes) that configure the agents, practices, conceptualities, and institutions that put images to work. . . .
Accordingly, the study of visual culture should be characterized by several concerns. First, scholars of visual culture need to examine any and all imagery - high and low, art and non-art. . . . They must not restrict themselves to objects of a particular beauty or aesthetic value. Indeed, any kind of imagery may be found to offer up evidence of the visual construction of reality. . . .
Second, the study of visual culture must scrutinize visual practice as much as images themselves, asking what images do when they are put to use. If scholars engaged in this enterprise inquire what makes an image beautiful or why this image or that constitutes a masterpiece or a work of genius, they should do so with the purpose of investigating an artist’s or a work’s contribution to the experience of beauty, taste, value, or genius. No amount of social analysis can account fully for the existence of Michelangelo or Leonardo. They were unique creators of images that changed the way their contemporaries thought and felt and have continued to shape the history of art, artists, museums, feeling, and aesthetic value. But study of the critical, artistic, and popular reception of works by such artists as Michelangelo and Leonardo can shed important light on the meaning of these artists and their works for many different people. And the history of meaning-making has a great deal to do with how scholars as well as lay audiences today understand these artists and their achievements.
Third, scholars studying visual culture might properly focus their interpretative work on life worlds by examining images, practices, visual technologies, taste, and artistic style as constitutive of social relations. The task is to understand how artifacts contribute to the construction of a world. . . . Important methodological implications follow: ethnography and reception studies become productive forms of gathering information, since these move beyond the image as a closed and fixed meaning-event. . . .
Fourth, scholars may learn a great deal when they scrutinize the constituents of vision, that is, the structures of perception as a physiological process as well as the epistemological frameworks informing a system of visual representation. Vision is a socially and a biologically constructed operation, depending on the design of the human body and how it engages the interpretive devices developed by a culture in order to see intelligibly. . . . Seeing . . . operates on the foundation of covenants with images that establish the conditions for meaningful visual experience.
Finally, the scholar of visual culture seeks to regard images as evidence for explanation, not as epiphenomena.

\vspace{1em}
\textbf{Question 19}

The passage given below is followed by four alternate summaries. Choose the option that best captures the essence of the passage.
All humans make decisions based on one or a combination of two factors. This is either intuition or information. Decisions made through intuition are usually fast, people don’t even think about the problem. It is quite philosophical, meaning that someone who made a decision based on intuition will have difficulty explaining the reasoning behind it. The decision-maker would often utilize her senses in drawing conclusions, which again is based on some experience in the field of study. On the other side of the spectrum, we have decisions made based on information. These decisions are rational — it is based on facts and figures, which unfortunately also means that it can be quite slow. The decision-maker would frequently use reports, analyses, and indicators to form her conclusion. This methodology results in accurate, quantifiable decisions, meaning that a person can clearly explain the rationale behind it.

\begin{enumerate}[label=\Alph*.]
    \item It is better to make decisions based on information because it is more accurate, and the rationale behind it can be explained.
    \item Decisions based on intuition and information result in differential speed and ability to provide a rationale.
    \item While decisions based on intuition can be made fast, the reasons that led to these cannot be spelt out.
    \item We make decisions based on intuition or information on the basis of the time available.
[](https://cracku.in/19-the-passage-given-below-is-followed-by-four-altern-x-cat-2020-slot-2)
\end{enumerate}
\vspace{1em}
\textbf{Solution:}

Let's look at the options one by one.
The author has not made any comparisons between the two methodologies, about which method is better. Hence optionA is incorrect.
Option 3 is incorrect too because it doesnot summarise the para accurately. Nothing is mentioned about decisions based on information.
Although time is a factor, nowhere it is mentioned that time is the factor in determining the method of our decision making. Hence Option4 is incorrect too.
Option 2 correctly summarises the passage by differentiating between the two methods. Hence it is correct.

\end{problem}

\newpage
\begin{problem}{}{}
% Subject: varc
\textbf{Instructions}

Instructions   

The passage below is accompanied by a set of questions. Choose the best answer to each question.
The claims advanced here may be condensed into two assertions: [first, that visual] culture is what images, acts of seeing, and attendant intellectual, emotional, and perceptual sensibilities do to build, maintain, or transform the worlds in which people live. [And second, that the] study of visual culture is the analysis and interpretation of images and the ways of seeing (or gazes) that configure the agents, practices, conceptualities, and institutions that put images to work. . . .
Accordingly, the study of visual culture should be characterized by several concerns. First, scholars of visual culture need to examine any and all imagery - high and low, art and non-art. . . . They must not restrict themselves to objects of a particular beauty or aesthetic value. Indeed, any kind of imagery may be found to offer up evidence of the visual construction of reality. . . .
Second, the study of visual culture must scrutinize visual practice as much as images themselves, asking what images do when they are put to use. If scholars engaged in this enterprise inquire what makes an image beautiful or why this image or that constitutes a masterpiece or a work of genius, they should do so with the purpose of investigating an artist’s or a work’s contribution to the experience of beauty, taste, value, or genius. No amount of social analysis can account fully for the existence of Michelangelo or Leonardo. They were unique creators of images that changed the way their contemporaries thought and felt and have continued to shape the history of art, artists, museums, feeling, and aesthetic value. But study of the critical, artistic, and popular reception of works by such artists as Michelangelo and Leonardo can shed important light on the meaning of these artists and their works for many different people. And the history of meaning-making has a great deal to do with how scholars as well as lay audiences today understand these artists and their achievements.
Third, scholars studying visual culture might properly focus their interpretative work on life worlds by examining images, practices, visual technologies, taste, and artistic style as constitutive of social relations. The task is to understand how artifacts contribute to the construction of a world. . . . Important methodological implications follow: ethnography and reception studies become productive forms of gathering information, since these move beyond the image as a closed and fixed meaning-event. . . .
Fourth, scholars may learn a great deal when they scrutinize the constituents of vision, that is, the structures of perception as a physiological process as well as the epistemological frameworks informing a system of visual representation. Vision is a socially and a biologically constructed operation, depending on the design of the human body and how it engages the interpretive devices developed by a culture in order to see intelligibly. . . . Seeing . . . operates on the foundation of covenants with images that establish the conditions for meaningful visual experience.
Finally, the scholar of visual culture seeks to regard images as evidence for explanation, not as epiphenomena.

\vspace{1em}
\textbf{Question 20}

The passage given below is followed by four alternate summaries. Choose the option that best captures the essence of the passage.
The rural-urban continuum and the heterogeneity of urban settings pose an obvious challenge to identifying urban areas and measuring urbanization rates in a consistent way within and across countries. An objective methodology for distinguishing between urban and rural areas that is based on one or two metrics with fixed thresholds may not adequately capture the wide diversity of places. A richer combination of criteria would better describe the multifaceted nature of a city’s function and its environment, but the joint interpretation of these criteria may require an element of human judgment.

\begin{enumerate}[label=\Alph*.]
    \item With the diversity of urban landscapes, measurable criteria for defining urban areas may need to be supplemented with human judgement.
    \item Current methodologies used to define urban and rural areas are no longer relevant to our being able to study trends in urbanisation.
    \item The difficulty of accurately identifying urban areas means that we need to create a rich combination of criteria that can be applied to all urban areas.
    \item Distinguishing between urban and rural areas might call for some judgement on the objective methodology being used to define a city’s functions.
[](https://cracku.in/20-the-passage-given-below-is-followed-by-four-altern-x-cat-2020-slot-2)
\end{enumerate}
\vspace{1em}
\textbf{Solution:}

Option B is incorrect, as it mentions "is no longer relevant" whereas the author says may no longer be relevant.
Option C is incorrect too, as the passage talks about distinguishing between urban/rural, not about accurately identifying rural areas.
Option D is distorted. Judgement would be required on the richer criteria , not on the objective methodology, as mentioned in this option.
Option A is correct as it correctly captures the essence of the passage.

\end{problem}

\newpage
\begin{problem}{}{}
% Subject: varc
\textbf{Instructions}

Instructions   

The passage below is accompanied by a set of questions. Choose the best answer to each question.
The claims advanced here may be condensed into two assertions: [first, that visual] culture is what images, acts of seeing, and attendant intellectual, emotional, and perceptual sensibilities do to build, maintain, or transform the worlds in which people live. [And second, that the] study of visual culture is the analysis and interpretation of images and the ways of seeing (or gazes) that configure the agents, practices, conceptualities, and institutions that put images to work. . . .
Accordingly, the study of visual culture should be characterized by several concerns. First, scholars of visual culture need to examine any and all imagery - high and low, art and non-art. . . . They must not restrict themselves to objects of a particular beauty or aesthetic value. Indeed, any kind of imagery may be found to offer up evidence of the visual construction of reality. . . .
Second, the study of visual culture must scrutinize visual practice as much as images themselves, asking what images do when they are put to use. If scholars engaged in this enterprise inquire what makes an image beautiful or why this image or that constitutes a masterpiece or a work of genius, they should do so with the purpose of investigating an artist’s or a work’s contribution to the experience of beauty, taste, value, or genius. No amount of social analysis can account fully for the existence of Michelangelo or Leonardo. They were unique creators of images that changed the way their contemporaries thought and felt and have continued to shape the history of art, artists, museums, feeling, and aesthetic value. But study of the critical, artistic, and popular reception of works by such artists as Michelangelo and Leonardo can shed important light on the meaning of these artists and their works for many different people. And the history of meaning-making has a great deal to do with how scholars as well as lay audiences today understand these artists and their achievements.
Third, scholars studying visual culture might properly focus their interpretative work on life worlds by examining images, practices, visual technologies, taste, and artistic style as constitutive of social relations. The task is to understand how artifacts contribute to the construction of a world. . . . Important methodological implications follow: ethnography and reception studies become productive forms of gathering information, since these move beyond the image as a closed and fixed meaning-event. . . .
Fourth, scholars may learn a great deal when they scrutinize the constituents of vision, that is, the structures of perception as a physiological process as well as the epistemological frameworks informing a system of visual representation. Vision is a socially and a biologically constructed operation, depending on the design of the human body and how it engages the interpretive devices developed by a culture in order to see intelligibly. . . . Seeing . . . operates on the foundation of covenants with images that establish the conditions for meaningful visual experience.
Finally, the scholar of visual culture seeks to regard images as evidence for explanation, not as epiphenomena.

\vspace{1em}
\textbf{Question 21}

Five jumbled up sentences, related to a topic, are given below. Four of them can be put together to form a coherent paragraph. Identify the odd one out and key in the number of the sentence as your answer:  
  
1. You can observe the truth of this in every e-business model ever constructed: monopolise and protect data.  
2. Economists and technologists believe that a new kind of capitalism is being created - different from industrial capitalism as was merchant capitalism.  
3. In 1962, Kenneth Arrow, the guru of mainstream economics, said that in a free market economy the purpose of inventing things is to create intellectual property rights.  
4. There is, alongside the world of monopolised information and surveillance, a different dynamic growing up: information as a social good, incapable of being owned or exploited or priced.  
5. Yet information is abundant. Information goods are freely replicable. Once a thing is made, it can be copied and pasted infinitely.
Backspace  
789  
456  
123  
0.-  
Clear All  

Submit
[](https://cracku.in/21-five-jumbled-up-sentences-related-to-a-topic-are-g-x-cat-2020-slot-2)

\begin{enumerate}[label=\Alph*.]
\end{enumerate}
\vspace{1em}
\textbf{Solution:}

We notice that Statement (3) serves as an introduction by touching upon the topic of intellectual property rights. Statement (2) continues along this line and mentions how 'monopolizing and protecting data' {an idea perhaps associated with the aforementioned intellectual property rights} can be observed in certain business models. Thus, the discussion revolves around data/information and its place in the market system. Statement (5) mentions that despite the monopolization, "information is abundant". And this brings us to the description in Statement (4) which portrays information differently - as a social good. Hence, although a bit discontiguous, 3-1-5-4 seems to talk about the same subject while Statement (2) {on "Merchant Capitalism"} diverges from it. Therefore, (2) is the odd-one-out.

\end{problem}

\newpage
\begin{problem}{}{}
% Subject: varc
\textbf{Instructions}

Instructions   

The passage below is accompanied by a set of questions. Choose the best answer to each question.
The claims advanced here may be condensed into two assertions: [first, that visual] culture is what images, acts of seeing, and attendant intellectual, emotional, and perceptual sensibilities do to build, maintain, or transform the worlds in which people live. [And second, that the] study of visual culture is the analysis and interpretation of images and the ways of seeing (or gazes) that configure the agents, practices, conceptualities, and institutions that put images to work. . . .
Accordingly, the study of visual culture should be characterized by several concerns. First, scholars of visual culture need to examine any and all imagery - high and low, art and non-art. . . . They must not restrict themselves to objects of a particular beauty or aesthetic value. Indeed, any kind of imagery may be found to offer up evidence of the visual construction of reality. . . .
Second, the study of visual culture must scrutinize visual practice as much as images themselves, asking what images do when they are put to use. If scholars engaged in this enterprise inquire what makes an image beautiful or why this image or that constitutes a masterpiece or a work of genius, they should do so with the purpose of investigating an artist’s or a work’s contribution to the experience of beauty, taste, value, or genius. No amount of social analysis can account fully for the existence of Michelangelo or Leonardo. They were unique creators of images that changed the way their contemporaries thought and felt and have continued to shape the history of art, artists, museums, feeling, and aesthetic value. But study of the critical, artistic, and popular reception of works by such artists as Michelangelo and Leonardo can shed important light on the meaning of these artists and their works for many different people. And the history of meaning-making has a great deal to do with how scholars as well as lay audiences today understand these artists and their achievements.
Third, scholars studying visual culture might properly focus their interpretative work on life worlds by examining images, practices, visual technologies, taste, and artistic style as constitutive of social relations. The task is to understand how artifacts contribute to the construction of a world. . . . Important methodological implications follow: ethnography and reception studies become productive forms of gathering information, since these move beyond the image as a closed and fixed meaning-event. . . .
Fourth, scholars may learn a great deal when they scrutinize the constituents of vision, that is, the structures of perception as a physiological process as well as the epistemological frameworks informing a system of visual representation. Vision is a socially and a biologically constructed operation, depending on the design of the human body and how it engages the interpretive devices developed by a culture in order to see intelligibly. . . . Seeing . . . operates on the foundation of covenants with images that establish the conditions for meaningful visual experience.
Finally, the scholar of visual culture seeks to regard images as evidence for explanation, not as epiphenomena.

\vspace{1em}
\textbf{Question 22}

The four sentences (labelled 1, 2, 3, 4) below, when properly sequenced would yield a coherent paragraph. Decide on the proper sequencing of the order of the sentences and key in the sequence of the four numbers as your answer:  
1. But the attention of the layman, not surprisingly, has been captured by the atom bomb, although there is at least a chance that it may never be used again.  
2. Of all the changes introduced by man into the household of nature, [controlled] large-scale nuclear fission is undoubtedly the most dangerous and most profound.  
3. The danger to humanity created by the so-called peaceful uses of atomic energy may, however, be much greater.  
4. The resultant ionizing radiation has become the most serious agent of pollution of the environment and the greatest threat to man’s survival on earth.
Backspace  
789  
456  
123  
0.-  
Clear All  

Submit
[](https://cracku.in/22-the-four-sentences-labelled-1-2-3-4-below-when-pro-x-cat-2020-slot-2)

\begin{enumerate}[label=\Alph*.]
\end{enumerate}
\vspace{1em}
\textbf{Solution:}

Of all the sentences, sentence 2 will be the starting point because it introduces the concept of nuclear fission. Sentence 4 continues this concept by talking about the "resulting ionizing radiation". Out of sentence 1 and 3, 1 will come first because it introduces the concept of atom bomb and sentence 3 further explores it by talking about how the impact of atomic energy could be much greater. Hence the correct order is 2413.

\end{problem}

\newpage
\begin{problem}{}{}
% Subject: varc
\textbf{Instructions}

Instructions   

The passage below is accompanied by a set of questions. Choose the best answer to each question.
The claims advanced here may be condensed into two assertions: [first, that visual] culture is what images, acts of seeing, and attendant intellectual, emotional, and perceptual sensibilities do to build, maintain, or transform the worlds in which people live. [And second, that the] study of visual culture is the analysis and interpretation of images and the ways of seeing (or gazes) that configure the agents, practices, conceptualities, and institutions that put images to work. . . .
Accordingly, the study of visual culture should be characterized by several concerns. First, scholars of visual culture need to examine any and all imagery - high and low, art and non-art. . . . They must not restrict themselves to objects of a particular beauty or aesthetic value. Indeed, any kind of imagery may be found to offer up evidence of the visual construction of reality. . . .
Second, the study of visual culture must scrutinize visual practice as much as images themselves, asking what images do when they are put to use. If scholars engaged in this enterprise inquire what makes an image beautiful or why this image or that constitutes a masterpiece or a work of genius, they should do so with the purpose of investigating an artist’s or a work’s contribution to the experience of beauty, taste, value, or genius. No amount of social analysis can account fully for the existence of Michelangelo or Leonardo. They were unique creators of images that changed the way their contemporaries thought and felt and have continued to shape the history of art, artists, museums, feeling, and aesthetic value. But study of the critical, artistic, and popular reception of works by such artists as Michelangelo and Leonardo can shed important light on the meaning of these artists and their works for many different people. And the history of meaning-making has a great deal to do with how scholars as well as lay audiences today understand these artists and their achievements.
Third, scholars studying visual culture might properly focus their interpretative work on life worlds by examining images, practices, visual technologies, taste, and artistic style as constitutive of social relations. The task is to understand how artifacts contribute to the construction of a world. . . . Important methodological implications follow: ethnography and reception studies become productive forms of gathering information, since these move beyond the image as a closed and fixed meaning-event. . . .
Fourth, scholars may learn a great deal when they scrutinize the constituents of vision, that is, the structures of perception as a physiological process as well as the epistemological frameworks informing a system of visual representation. Vision is a socially and a biologically constructed operation, depending on the design of the human body and how it engages the interpretive devices developed by a culture in order to see intelligibly. . . . Seeing . . . operates on the foundation of covenants with images that establish the conditions for meaningful visual experience.
Finally, the scholar of visual culture seeks to regard images as evidence for explanation, not as epiphenomena.

\vspace{1em}
\textbf{Question 23}

The passage given below is followed by four alternate summaries. Choose the option that best captures the essence of the passage.  
  
With the Treaty of Westphalia, the papacy had been confined to ecclesiastical functions, and the doctrine of sovereign equality reigned. What political theory could then explain the origin and justify the functions of secular political order? In his Leviathan, published in 1651, three years after the Peace of Westphalia, Thomas Hobbes provided such a theory. He imagined a “state of nature” in the past when the absence of authority produced a “war of all against all.” To escape such intolerable insecurity, he theorized, people delivered their rights to a sovereign power in return for the sovereign’s provision of security for all within the state’s border. The sovereign state’s monopoly on power was established as the only way to overcome the perpetual fear of violent death and war.

\begin{enumerate}[label=\Alph*.]
    \item Thomas Hobbes theorized the voluntary surrender of rights by people as essential for emergence of sovereign states.
    \item Thomas Hobbes theorized the emergence of sovereign states as a form of transactional governance to limit the power of the papacy.
    \item Thomas Hobbes theorized that sovereign states emerged out of people’s voluntary desire to overcome the sense of insecurity and establish the doctrine of sovereign equality.
    \item Thomas Hobbes theorized the emergence of sovereign states based on a transactional relationship between people and sovereign state that was necessitated by a sense of insecurity of the people.
[](https://cracku.in/23-the-passage-given-below-is-followed-by-four-altern-x-cat-2020-slot-2)
\end{enumerate}
\vspace{1em}
\textbf{Solution:}

Let's look at the options one by one. Option 1 talks about voluntary surrender of rights. But the passage talks about "transfer" of rights, not "surrender" of rights.
The main point is not about "powers of papacy", option 2 is inconsequential.
Option 3 does not cover one of the main points of the passage, the transfer of rights, between people and the sovereign power. It's an incomplete option.
Option 4 is the correct summary, it talks about the transactional relationship i.e give and take or transfer of rights between people and sovereign government.

\end{problem}

\newpage
\begin{problem}{}{}
% Subject: varc
\textbf{Instructions}

Instructions   

The passage below is accompanied by a set of questions. Choose the best answer to each question.
The claims advanced here may be condensed into two assertions: [first, that visual] culture is what images, acts of seeing, and attendant intellectual, emotional, and perceptual sensibilities do to build, maintain, or transform the worlds in which people live. [And second, that the] study of visual culture is the analysis and interpretation of images and the ways of seeing (or gazes) that configure the agents, practices, conceptualities, and institutions that put images to work. . . .
Accordingly, the study of visual culture should be characterized by several concerns. First, scholars of visual culture need to examine any and all imagery - high and low, art and non-art. . . . They must not restrict themselves to objects of a particular beauty or aesthetic value. Indeed, any kind of imagery may be found to offer up evidence of the visual construction of reality. . . .
Second, the study of visual culture must scrutinize visual practice as much as images themselves, asking what images do when they are put to use. If scholars engaged in this enterprise inquire what makes an image beautiful or why this image or that constitutes a masterpiece or a work of genius, they should do so with the purpose of investigating an artist’s or a work’s contribution to the experience of beauty, taste, value, or genius. No amount of social analysis can account fully for the existence of Michelangelo or Leonardo. They were unique creators of images that changed the way their contemporaries thought and felt and have continued to shape the history of art, artists, museums, feeling, and aesthetic value. But study of the critical, artistic, and popular reception of works by such artists as Michelangelo and Leonardo can shed important light on the meaning of these artists and their works for many different people. And the history of meaning-making has a great deal to do with how scholars as well as lay audiences today understand these artists and their achievements.
Third, scholars studying visual culture might properly focus their interpretative work on life worlds by examining images, practices, visual technologies, taste, and artistic style as constitutive of social relations. The task is to understand how artifacts contribute to the construction of a world. . . . Important methodological implications follow: ethnography and reception studies become productive forms of gathering information, since these move beyond the image as a closed and fixed meaning-event. . . .
Fourth, scholars may learn a great deal when they scrutinize the constituents of vision, that is, the structures of perception as a physiological process as well as the epistemological frameworks informing a system of visual representation. Vision is a socially and a biologically constructed operation, depending on the design of the human body and how it engages the interpretive devices developed by a culture in order to see intelligibly. . . . Seeing . . . operates on the foundation of covenants with images that establish the conditions for meaningful visual experience.
Finally, the scholar of visual culture seeks to regard images as evidence for explanation, not as epiphenomena.

\vspace{1em}
\textbf{Question 24}

Five jumbled up sentences, related to a topic, are given below. Four of them can be put together to form a coherent paragraph. Identify the odd one out and key in the number of the sentence as your answer:  
  
1. The victim’s trauma after assault rarely gets the attention that we lavish on the moment of damage that divided the survivor from a less encumbered past.  
2. One thing we often do with narratives of sexual assault is sort their respective parties into different temporalities: it seems we are interested in perpetrators’ futures and victims’ pasts.  
3. One result is that we don’t have much of a vocabulary for what happens in a victim’s life after the painful past has been excavated, even when our shared language gestures toward the future, as the term “survivor” does.  
4. Even the most charitable questions asked about the victims seem to focus on the past, in pursuit of understanding or of corroboration of painful details.  
5. As more and more stories of sexual assault have been made public in the last two years, the genre of their telling has exploded --- crimes have a tendency to become not just stories but genres.
Backspace  
789  
456  
123  
0.-  
Clear All  

Submit
[](https://cracku.in/24-five-jumbled-up-sentences-related-to-a-topic-are-g-x-cat-2020-slot-2)

\begin{enumerate}[label=\Alph*.]
\end{enumerate}
\vspace{1em}
\textbf{Solution:}

The subject here is about the narratives of sexual assault and the subsequent treatment of the victims. The paragraph begins with Statement (5) which talks about sexual assault narratives turning into a genre. Statement (2) comments on a stark element within these narratives: the general interest in perpetrators’ futures and victims’ pasts. Statement (3) delineates on the outcome of the latter: the focus on a victim's past' and Statement (1) further adds to this point. Although Statement (4) appears on a similar topic, we cannot place it in the above arrangement since it diverges into the subject of asking questions to the survivors {an element that is yet to be discussed}. Thus, (4) is the odd-one-out.

\end{problem}

\newpage
\begin{problem}{}{}
% Subject: varc
\textbf{Instructions}

Instructions   

The passage below is accompanied by a set of questions. Choose the best answer to each question.
The claims advanced here may be condensed into two assertions: [first, that visual] culture is what images, acts of seeing, and attendant intellectual, emotional, and perceptual sensibilities do to build, maintain, or transform the worlds in which people live. [And second, that the] study of visual culture is the analysis and interpretation of images and the ways of seeing (or gazes) that configure the agents, practices, conceptualities, and institutions that put images to work. . . .
Accordingly, the study of visual culture should be characterized by several concerns. First, scholars of visual culture need to examine any and all imagery - high and low, art and non-art. . . . They must not restrict themselves to objects of a particular beauty or aesthetic value. Indeed, any kind of imagery may be found to offer up evidence of the visual construction of reality. . . .
Second, the study of visual culture must scrutinize visual practice as much as images themselves, asking what images do when they are put to use. If scholars engaged in this enterprise inquire what makes an image beautiful or why this image or that constitutes a masterpiece or a work of genius, they should do so with the purpose of investigating an artist’s or a work’s contribution to the experience of beauty, taste, value, or genius. No amount of social analysis can account fully for the existence of Michelangelo or Leonardo. They were unique creators of images that changed the way their contemporaries thought and felt and have continued to shape the history of art, artists, museums, feeling, and aesthetic value. But study of the critical, artistic, and popular reception of works by such artists as Michelangelo and Leonardo can shed important light on the meaning of these artists and their works for many different people. And the history of meaning-making has a great deal to do with how scholars as well as lay audiences today understand these artists and their achievements.
Third, scholars studying visual culture might properly focus their interpretative work on life worlds by examining images, practices, visual technologies, taste, and artistic style as constitutive of social relations. The task is to understand how artifacts contribute to the construction of a world. . . . Important methodological implications follow: ethnography and reception studies become productive forms of gathering information, since these move beyond the image as a closed and fixed meaning-event. . . .
Fourth, scholars may learn a great deal when they scrutinize the constituents of vision, that is, the structures of perception as a physiological process as well as the epistemological frameworks informing a system of visual representation. Vision is a socially and a biologically constructed operation, depending on the design of the human body and how it engages the interpretive devices developed by a culture in order to see intelligibly. . . . Seeing . . . operates on the foundation of covenants with images that establish the conditions for meaningful visual experience.
Finally, the scholar of visual culture seeks to regard images as evidence for explanation, not as epiphenomena.

\vspace{1em}
\textbf{Question 25}

The four sentences (labelled 1, 2, 3, 4) below, when properly sequenced would yield a coherent paragraph. Decide on the proper sequencing of the order of the sentences and key in the sequence of the four numbers as your answer:  
  
1. It also has four movable auxiliary telescopes 1.8 m in diameter.  
2. Completed in 2006, the Very Large Telescope (VLT) has four reflecting telescopes, 8.2 m in diameter that can observe objects 4 billion times weaker than can normally be seen with the naked eye.  
3. This configuration enables one to distinguish an astronaut on the Moon.  
4. When these are combined with the large telescopes, they produce what is called interferometry: a simulation of the power of a mirror 16 m in diameter and the resolution of a telescope of 200 m.
Backspace  
789  
456  
123  
0.-  
Clear All  

Submit
[](https://cracku.in/25-the-four-sentences-labelled-1-2-3-4-below-when-pro-x-cat-2020-slot-2)

\begin{enumerate}[label=\Alph*.]
\end{enumerate}
\vspace{1em}
\textbf{Solution:}

This is a very easy parajumble question. Sentence 2 introduces the topic, VLT. Option 4 further explores it by discussing about the telescopes. 43 forms a natural block as one provides the use and other the application. Hence 2143 is the correct order

\end{problem}

\newpage
\begin{problem}{}{}
% Subject: varc
\textbf{Instructions}

Instructions   

The passage below is accompanied by a set of questions. Choose the best answer to each question.
The claims advanced here may be condensed into two assertions: [first, that visual] culture is what images, acts of seeing, and attendant intellectual, emotional, and perceptual sensibilities do to build, maintain, or transform the worlds in which people live. [And second, that the] study of visual culture is the analysis and interpretation of images and the ways of seeing (or gazes) that configure the agents, practices, conceptualities, and institutions that put images to work. . . .
Accordingly, the study of visual culture should be characterized by several concerns. First, scholars of visual culture need to examine any and all imagery - high and low, art and non-art. . . . They must not restrict themselves to objects of a particular beauty or aesthetic value. Indeed, any kind of imagery may be found to offer up evidence of the visual construction of reality. . . .
Second, the study of visual culture must scrutinize visual practice as much as images themselves, asking what images do when they are put to use. If scholars engaged in this enterprise inquire what makes an image beautiful or why this image or that constitutes a masterpiece or a work of genius, they should do so with the purpose of investigating an artist’s or a work’s contribution to the experience of beauty, taste, value, or genius. No amount of social analysis can account fully for the existence of Michelangelo or Leonardo. They were unique creators of images that changed the way their contemporaries thought and felt and have continued to shape the history of art, artists, museums, feeling, and aesthetic value. But study of the critical, artistic, and popular reception of works by such artists as Michelangelo and Leonardo can shed important light on the meaning of these artists and their works for many different people. And the history of meaning-making has a great deal to do with how scholars as well as lay audiences today understand these artists and their achievements.
Third, scholars studying visual culture might properly focus their interpretative work on life worlds by examining images, practices, visual technologies, taste, and artistic style as constitutive of social relations. The task is to understand how artifacts contribute to the construction of a world. . . . Important methodological implications follow: ethnography and reception studies become productive forms of gathering information, since these move beyond the image as a closed and fixed meaning-event. . . .
Fourth, scholars may learn a great deal when they scrutinize the constituents of vision, that is, the structures of perception as a physiological process as well as the epistemological frameworks informing a system of visual representation. Vision is a socially and a biologically constructed operation, depending on the design of the human body and how it engages the interpretive devices developed by a culture in order to see intelligibly. . . . Seeing . . . operates on the foundation of covenants with images that establish the conditions for meaningful visual experience.
Finally, the scholar of visual culture seeks to regard images as evidence for explanation, not as epiphenomena.

\vspace{1em}
\textbf{Question 26}

The four sentences (labelled 1, 2, 3, 4) below, when properly sequenced would yield a coherent paragraph. Decide on the proper sequencing of the order of the sentences and key in the sequence of the four numbers as your answer:  
1. While you might think that you see or are aware of all the changes that happen in your immediate environment, there is simply too much information for your brain to fully process everything.  
2. Psychologists use the term ‘change blindness’ to describe this tendency of people to be blind to changes though they are in the immediate environment.  
3. It cannot be aware of every single thing that happens in the world around you.  
4. Sometimes big shifts happen in front of your eyes and you are not at all aware of these changes.
Backspace  
789  
456  
123  
0.-  
Clear All  

Submit
[](https://cracku.in/26-the-four-sentences-labelled-1-2-3-4-below-when-pro-x-cat-2020-slot-2)

\begin{enumerate}[label=\Alph*.]
\end{enumerate}
\vspace{1em}
\textbf{Solution:}

1 and 3 forms a natural block , as the "It" in sentence 3 refers to the "brain" discussed in statement 1. Again 4 and 2 forms a natural block, as the "this tendency" in statement 2 , refers to the tendency of the people ignoring things happening in front of their eyes as discussed in statement 4. Hence the correct order is 1342
[](https://cracku.in/downloads/11131)
  *  
  *

\end{problem}

\newpage
\begin{problem}{}{}
% Subject: varc
\textbf{Instructions}

Instructions   

The passage below is accompanied by a set of questions. Choose the best answer to each question.
Mode of transportation affects the travel experience and thus can produce new types of travel writing and perhaps even new “identities.” Modes of transportation determine the types and duration of social encounters; affect the organization and passage of space and time; . . . and also affect perception and knowledge—how and what the traveler comes to know and write about. The completion of the first U.S. transcontinental highway during the 1920s . . . for example, inaugurated a new genre of travel literature about the United States—the automotive or road narrative. Such narratives highlight the experiences of mostly male protagonists “discovering themselves” on their journeys, emphasizing the independence of road travel and the value of rural folk traditions.
Travel writing’s relationship to empire building— as a type of “colonialist discourse”—has drawn the most attention from academicians. Close connections have been observed between European (and American) political, economic, and administrative goals for the colonies and their manifestations in the cultural practice of writing travel books. Travel writers’ descriptions of foreign places have been analyzed as attempts to validate, promote, or challenge the ideologies and practices of colonial or imperial domination and expansion. Mary Louise Pratt’s study of the genres and conventions of 18th- and 19th-century exploration narratives about South America and Africa (e.g., the “monarch of all I survey” trope) offered ways of thinking about travel writing as embedded within relations of power between metropole and periphery, as did Edward Said’s theories of representation and cultural imperialism. Particularly Said’s book, Orientalism, helped scholars understand ways in which representations of people in travel texts were intimately bound up with notions of self, in this case, that the Occident defined itself through essentialist, ethnocentric, and racist representations of the Orient. Said’s work became a model for demonstrating cultural forms of imperialism in travel texts, showing how the political, economic, or administrative fact of dominance relies on legitimating discourses such as those articulated through travel writing. . . .
Feminist geographers’ studies of travel writing challenge the masculinist history of geography by questioning who and what are relevant subjects of geographic study and, indeed, what counts as geographic knowledge itself. Such questions are worked through ideological constructs that posit men as explorers and women as travelers—or, conversely, men as travelers and women as tied to the home. Studies of Victorian women who were professional travel writers, tourists, wives of colonial administrators, and other (mostly) elite women who wrote narratives about their experiences abroad during the 19th century have been particularly revealing. From a “liberal” feminist perspective, travel presented one means toward female liberation for middle- and upper-class Victorian women. Many studies from the 1970s onward demonstrated the ways in which women’s gendered identities were negotiated differently “at home” than they were “away,” thereby showing women’s self-development through travel. The more recent post structural turn in studies of Victorian travel writing has focused attention on women’s diverse and fragmented identities as they narrated their travel experiences, emphasizing women’s sense of themselves as women in new locations, but only as they worked through their ties to nation, class, whiteness, and colonial and imperial power structures

\vspace{1em}
\textbf{Question 1}

From the passage, we can infer that feminist scholars’ understanding of the experiences of Victorian women travellers is influenced by all of the following EXCEPT scholars':

\begin{enumerate}[label=\Alph*.]
    \item awareness of gender issues in Victorian society
    \item knowledge of class tensions in Victorian society
    \item perspective that they bring to their research
    \item awareness of the ways in which identity is formed
[](https://cracku.in/1-from-the-passage-we-can-infer-that-feminist-schola-x-cat-2020-slot-3)
\end{enumerate}
\vspace{1em}
\textbf{Solution:}

The aspects that guided/influenced the understanding of feminist scholars who examined Victorian women's travelling experiences need to be found out. We understand that the attempt made by feminist scholars was to "challenge the masculinist history of geography by (1) questioning who and what are relevant subjects of geographic study, and (2) what counts as geographic knowledge itself." And considering the role of women in travel writing during the Victorian era enabled these scholars to inspect new perspectives and achieve the aforementioned objectives {"..._such questions are worked through ideological constructs that posit men as explorers and women as travelers—or, conversely, men as travelers and women as tied to the home._.."}[_Option C_]. The varied viewpoints offered by women during this period originated from the difference in the gendered identities, implicitly indicating that the presence of an inequality/ineuity in the gender roles _[Option A_]. Additionally, studies concerning the manner in which travel altered a woman's gendered identities were also available; this further shaped the feminist scholars' understanding of the travelling experiences of Victorian women [_Option D_]. It has not been presented or implied that the knowledge of "class" tensions, as stated in _Option B_ , was an imperative element that influenced scholars' understanding. Hence, Option B is the correct answer.

\end{problem}

\newpage
\begin{problem}{}{}
% Subject: varc
\textbf{Instructions}

Instructions   

The passage below is accompanied by a set of questions. Choose the best answer to each question.
Mode of transportation affects the travel experience and thus can produce new types of travel writing and perhaps even new “identities.” Modes of transportation determine the types and duration of social encounters; affect the organization and passage of space and time; . . . and also affect perception and knowledge—how and what the traveler comes to know and write about. The completion of the first U.S. transcontinental highway during the 1920s . . . for example, inaugurated a new genre of travel literature about the United States—the automotive or road narrative. Such narratives highlight the experiences of mostly male protagonists “discovering themselves” on their journeys, emphasizing the independence of road travel and the value of rural folk traditions.
Travel writing’s relationship to empire building— as a type of “colonialist discourse”—has drawn the most attention from academicians. Close connections have been observed between European (and American) political, economic, and administrative goals for the colonies and their manifestations in the cultural practice of writing travel books. Travel writers’ descriptions of foreign places have been analyzed as attempts to validate, promote, or challenge the ideologies and practices of colonial or imperial domination and expansion. Mary Louise Pratt’s study of the genres and conventions of 18th- and 19th-century exploration narratives about South America and Africa (e.g., the “monarch of all I survey” trope) offered ways of thinking about travel writing as embedded within relations of power between metropole and periphery, as did Edward Said’s theories of representation and cultural imperialism. Particularly Said’s book, Orientalism, helped scholars understand ways in which representations of people in travel texts were intimately bound up with notions of self, in this case, that the Occident defined itself through essentialist, ethnocentric, and racist representations of the Orient. Said’s work became a model for demonstrating cultural forms of imperialism in travel texts, showing how the political, economic, or administrative fact of dominance relies on legitimating discourses such as those articulated through travel writing. . . .
Feminist geographers’ studies of travel writing challenge the masculinist history of geography by questioning who and what are relevant subjects of geographic study and, indeed, what counts as geographic knowledge itself. Such questions are worked through ideological constructs that posit men as explorers and women as travelers—or, conversely, men as travelers and women as tied to the home. Studies of Victorian women who were professional travel writers, tourists, wives of colonial administrators, and other (mostly) elite women who wrote narratives about their experiences abroad during the 19th century have been particularly revealing. From a “liberal” feminist perspective, travel presented one means toward female liberation for middle- and upper-class Victorian women. Many studies from the 1970s onward demonstrated the ways in which women’s gendered identities were negotiated differently “at home” than they were “away,” thereby showing women’s self-development through travel. The more recent post structural turn in studies of Victorian travel writing has focused attention on women’s diverse and fragmented identities as they narrated their travel experiences, emphasizing women’s sense of themselves as women in new locations, but only as they worked through their ties to nation, class, whiteness, and colonial and imperial power structures

\vspace{1em}
\textbf{Question 2}

American travel literature of the 1920s:

\begin{enumerate}[label=\Alph*.]
    \item showed participation in local traditions.
    \item developed the male protagonists’ desire for independence
    \item presented travellers’ discovery of their identity as different from others.
    \item celebrated the freedom that travel gives.
[](https://cracku.in/2-american-travel-literature-of-the-1920s-x-cat-2020-slot-3)
\end{enumerate}
\vspace{1em}
\textbf{Solution:}

We can zero-in on the answer based on the following excerpt: {.._.The completion of the first U.S. transcontinental highway during the 1920s . . . for example, inaugurated a new genre of travel literature about the United States—the automotive or road narrative. Such narratives highlight the experiences of mostly male protagonists “discovering themselves” on their journeys, emphasizing the independence of road travel and the value of rural folk traditions.._.}
_Option A_ : talks about participation in local traditions which is not mentioned or implied.
_Option B_ : is a distorted comment that does not align with the idea discussed. The author states that road journeys enabled the male protagonist's experience of "discovering themselves"; the phrase "desire for independence" would be incorrect in this regard.
_Option C_ : is not even remotely discussed/implied.
_Option D_ : It is mentioned that the inauguration of a transcontinental highway during the 1920s paved the way for a new genre that emphasised the freedom attached to such road travelling enterprises. Hence, it implicitly depicted travel as an experience celebrating an individual's independence {..._emphasizing the independence_...}. 
Hence, of the given options, Option D aptly captures the characteristics of American travel literature of the 1920s.

\end{problem}

\newpage
\begin{problem}{}{}
% Subject: varc
\textbf{Instructions}

Instructions   

The passage below is accompanied by a set of questions. Choose the best answer to each question.
Mode of transportation affects the travel experience and thus can produce new types of travel writing and perhaps even new “identities.” Modes of transportation determine the types and duration of social encounters; affect the organization and passage of space and time; . . . and also affect perception and knowledge—how and what the traveler comes to know and write about. The completion of the first U.S. transcontinental highway during the 1920s . . . for example, inaugurated a new genre of travel literature about the United States—the automotive or road narrative. Such narratives highlight the experiences of mostly male protagonists “discovering themselves” on their journeys, emphasizing the independence of road travel and the value of rural folk traditions.
Travel writing’s relationship to empire building— as a type of “colonialist discourse”—has drawn the most attention from academicians. Close connections have been observed between European (and American) political, economic, and administrative goals for the colonies and their manifestations in the cultural practice of writing travel books. Travel writers’ descriptions of foreign places have been analyzed as attempts to validate, promote, or challenge the ideologies and practices of colonial or imperial domination and expansion. Mary Louise Pratt’s study of the genres and conventions of 18th- and 19th-century exploration narratives about South America and Africa (e.g., the “monarch of all I survey” trope) offered ways of thinking about travel writing as embedded within relations of power between metropole and periphery, as did Edward Said’s theories of representation and cultural imperialism. Particularly Said’s book, Orientalism, helped scholars understand ways in which representations of people in travel texts were intimately bound up with notions of self, in this case, that the Occident defined itself through essentialist, ethnocentric, and racist representations of the Orient. Said’s work became a model for demonstrating cultural forms of imperialism in travel texts, showing how the political, economic, or administrative fact of dominance relies on legitimating discourses such as those articulated through travel writing. . . .
Feminist geographers’ studies of travel writing challenge the masculinist history of geography by questioning who and what are relevant subjects of geographic study and, indeed, what counts as geographic knowledge itself. Such questions are worked through ideological constructs that posit men as explorers and women as travelers—or, conversely, men as travelers and women as tied to the home. Studies of Victorian women who were professional travel writers, tourists, wives of colonial administrators, and other (mostly) elite women who wrote narratives about their experiences abroad during the 19th century have been particularly revealing. From a “liberal” feminist perspective, travel presented one means toward female liberation for middle- and upper-class Victorian women. Many studies from the 1970s onward demonstrated the ways in which women’s gendered identities were negotiated differently “at home” than they were “away,” thereby showing women’s self-development through travel. The more recent post structural turn in studies of Victorian travel writing has focused attention on women’s diverse and fragmented identities as they narrated their travel experiences, emphasizing women’s sense of themselves as women in new locations, but only as they worked through their ties to nation, class, whiteness, and colonial and imperial power structures

\vspace{1em}
\textbf{Question 3}

From the passage, it can be inferred that scholars argue that Victorian women  
experienced self-development through their travels because:

\begin{enumerate}[label=\Alph*.]
    \item they developed a feminist perspective of the world.
    \item they were from the progressive middle- and upper-classes of society.
    \item they were on a quest to discover their diverse identities.
    \item their identity was redefined when they were away from home.
[](https://cracku.in/3-from-the-passage-it-can-be-inferred-that-scholars--x-cat-2020-slot-3)
\end{enumerate}
\vspace{1em}
\textbf{Solution:}

The question requires us to probe the reason behind why the scholars argue that Victorian women experienced self-development through their travels. Let us pay heed to the following segment from the passage: {.._.Many studies from the 1970s onward demonstrated the ways in which women’s gendered identities were negotiated differently “at home” than they were “away,” thereby showing women’s self-development through travel._..}. It is highlighted that travelling {being "away"} enabled women's identities to be "_negotiated differently_ "{highlighting transformation/reconfiguration}, which in turn was the cause of their self-development. Option D is closest to this understanding. Options A, B and C fail to present the reason and are invalid statements/inferences.

\end{problem}

\newpage
\begin{problem}{}{}
% Subject: varc
\textbf{Instructions}

Instructions   

The passage below is accompanied by a set of questions. Choose the best answer to each question.
Mode of transportation affects the travel experience and thus can produce new types of travel writing and perhaps even new “identities.” Modes of transportation determine the types and duration of social encounters; affect the organization and passage of space and time; . . . and also affect perception and knowledge—how and what the traveler comes to know and write about. The completion of the first U.S. transcontinental highway during the 1920s . . . for example, inaugurated a new genre of travel literature about the United States—the automotive or road narrative. Such narratives highlight the experiences of mostly male protagonists “discovering themselves” on their journeys, emphasizing the independence of road travel and the value of rural folk traditions.
Travel writing’s relationship to empire building— as a type of “colonialist discourse”—has drawn the most attention from academicians. Close connections have been observed between European (and American) political, economic, and administrative goals for the colonies and their manifestations in the cultural practice of writing travel books. Travel writers’ descriptions of foreign places have been analyzed as attempts to validate, promote, or challenge the ideologies and practices of colonial or imperial domination and expansion. Mary Louise Pratt’s study of the genres and conventions of 18th- and 19th-century exploration narratives about South America and Africa (e.g., the “monarch of all I survey” trope) offered ways of thinking about travel writing as embedded within relations of power between metropole and periphery, as did Edward Said’s theories of representation and cultural imperialism. Particularly Said’s book, Orientalism, helped scholars understand ways in which representations of people in travel texts were intimately bound up with notions of self, in this case, that the Occident defined itself through essentialist, ethnocentric, and racist representations of the Orient. Said’s work became a model for demonstrating cultural forms of imperialism in travel texts, showing how the political, economic, or administrative fact of dominance relies on legitimating discourses such as those articulated through travel writing. . . .
Feminist geographers’ studies of travel writing challenge the masculinist history of geography by questioning who and what are relevant subjects of geographic study and, indeed, what counts as geographic knowledge itself. Such questions are worked through ideological constructs that posit men as explorers and women as travelers—or, conversely, men as travelers and women as tied to the home. Studies of Victorian women who were professional travel writers, tourists, wives of colonial administrators, and other (mostly) elite women who wrote narratives about their experiences abroad during the 19th century have been particularly revealing. From a “liberal” feminist perspective, travel presented one means toward female liberation for middle- and upper-class Victorian women. Many studies from the 1970s onward demonstrated the ways in which women’s gendered identities were negotiated differently “at home” than they were “away,” thereby showing women’s self-development through travel. The more recent post structural turn in studies of Victorian travel writing has focused attention on women’s diverse and fragmented identities as they narrated their travel experiences, emphasizing women’s sense of themselves as women in new locations, but only as they worked through their ties to nation, class, whiteness, and colonial and imperial power structures

\vspace{1em}
\textbf{Question 4}

According to the passage, Said’s book, “Orientalism”:

\begin{enumerate}[label=\Alph*.]
    \item argued that cultural imperialism was more significant than colonial domination.
    \item explained the difference between the representation of people and the actual fact
    \item illustrated how narrow minded and racist westerners were
    \item demonstrated how cultural imperialism was used to justify colonial domination.
[](https://cracku.in/4-according-to-the-passage-saids-book-orientalism-x-cat-2020-slot-3)
\end{enumerate}
\vspace{1em}
\textbf{Solution:}

A direct inference from the excerpt: {.._.Said’s work became a model for demonstrating cultural forms of imperialism in travel texts, showing how the political, economic, or administrative fact of dominance relies on legitimating discourses such as those articulated through travel writing. ._ .}. It is stated that Said's work rendered scholars with an understanding of "cultural imperialism" and the manner in which it was used to justify colonial domination {or similar pursuits thereof}. Option B aptly captures this aspect. Options A, B and C either diverge from the discussion or are distorted comments.

\end{problem}

\newpage
\begin{problem}{}{}
% Subject: varc
\textbf{Instructions}

Instructions   

The passage below is accompanied by a set of questions. Choose the best answer to each question.
Mode of transportation affects the travel experience and thus can produce new types of travel writing and perhaps even new “identities.” Modes of transportation determine the types and duration of social encounters; affect the organization and passage of space and time; . . . and also affect perception and knowledge—how and what the traveler comes to know and write about. The completion of the first U.S. transcontinental highway during the 1920s . . . for example, inaugurated a new genre of travel literature about the United States—the automotive or road narrative. Such narratives highlight the experiences of mostly male protagonists “discovering themselves” on their journeys, emphasizing the independence of road travel and the value of rural folk traditions.
Travel writing’s relationship to empire building— as a type of “colonialist discourse”—has drawn the most attention from academicians. Close connections have been observed between European (and American) political, economic, and administrative goals for the colonies and their manifestations in the cultural practice of writing travel books. Travel writers’ descriptions of foreign places have been analyzed as attempts to validate, promote, or challenge the ideologies and practices of colonial or imperial domination and expansion. Mary Louise Pratt’s study of the genres and conventions of 18th- and 19th-century exploration narratives about South America and Africa (e.g., the “monarch of all I survey” trope) offered ways of thinking about travel writing as embedded within relations of power between metropole and periphery, as did Edward Said’s theories of representation and cultural imperialism. Particularly Said’s book, Orientalism, helped scholars understand ways in which representations of people in travel texts were intimately bound up with notions of self, in this case, that the Occident defined itself through essentialist, ethnocentric, and racist representations of the Orient. Said’s work became a model for demonstrating cultural forms of imperialism in travel texts, showing how the political, economic, or administrative fact of dominance relies on legitimating discourses such as those articulated through travel writing. . . .
Feminist geographers’ studies of travel writing challenge the masculinist history of geography by questioning who and what are relevant subjects of geographic study and, indeed, what counts as geographic knowledge itself. Such questions are worked through ideological constructs that posit men as explorers and women as travelers—or, conversely, men as travelers and women as tied to the home. Studies of Victorian women who were professional travel writers, tourists, wives of colonial administrators, and other (mostly) elite women who wrote narratives about their experiences abroad during the 19th century have been particularly revealing. From a “liberal” feminist perspective, travel presented one means toward female liberation for middle- and upper-class Victorian women. Many studies from the 1970s onward demonstrated the ways in which women’s gendered identities were negotiated differently “at home” than they were “away,” thereby showing women’s self-development through travel. The more recent post structural turn in studies of Victorian travel writing has focused attention on women’s diverse and fragmented identities as they narrated their travel experiences, emphasizing women’s sense of themselves as women in new locations, but only as they worked through their ties to nation, class, whiteness, and colonial and imperial power structures

\vspace{1em}
\textbf{Question 5}

From the passage, we can infer that travel writing is most similar to:

\begin{enumerate}[label=\Alph*.]
    \item political journalism
    \item feminist writing
    \item autobiographical writing.
    \item historical fiction.
[](https://cracku.in/5-from-the-passage-we-can-infer-that-travel-writing--x-cat-2020-slot-3)
\end{enumerate}
\vspace{1em}
\textbf{Solution:}

Since travel writing involves the presentation and recounting of personal travel experiences and/or perspective of the world, the closest category to this would be autobiographical writing [Option C]. Political journalism and feminist writing can be hurled out the window. Associating travel literature to historical fiction would be inappropriate as well. Hence, of the given choices, Option C is the correct answer.

Instructions   

**The passage below is accompanied by a set of questions. Choose the best answer to each question.**
Although one of the most contested concepts in political philosophy, human nature is something on which most people seem to agree. By and large, according to Rutger Bregman in his new book Humankind, we have a rather pessimistic view - not of ourselves exactly, but of everyone else. We see other people as selfish, untrustworthy and dangerous and therefore we behave towards them with defensiveness and suspicion. This was how the 17th-century philosopher Thomas Hobbes conceived our natural state to be, believing that all that stood between us and violent anarchy was a strong state and firm leadership. But in following Hobbes, argues Bregman, we ensure that the negative view we have of human nature is reflected back at us. He instead puts his faith in Jean-Jacques Rousseau, the 18th-century French thinker, who famously declared that man was born free and it was civilisation - with its coercive powers, social classes and restrictive laws - that put him in chains.
Hobbes and Rousseau are seen as the two poles of the human nature argument and it’s no surprise that Bregman strongly sides with the Frenchman. He takes Rousseau’s intuition and paints a picture of a prelapsarian idyll in which, for the better part of 300,000 years, Homo sapiens lived a fulfilling life in harmony with nature . . . Then we discovered agriculture and for the next 10,000 years it was all property, war, greed and injustice. . . . 
It was abandoning our nomadic lifestyle and then domesticating animals, says Bregman, that brought about infectious diseases such as measles, smallpox, tuberculosis, syphilis, malaria, cholera and plague. This may be true, but what Bregman never really seems to get to grips with is that pathogens were not the only things that grew with agriculture - so did the number of humans. It’s one thing to maintain friendly relations and a property-less mode of living when you’re 30 or 40 hunter-gatherers following the food. But life becomes a great deal more complex and knowledge far more extensive when there are settlements of many thousands. “Civilisation has become synonymous with peace and progress and wilderness with war and decline,” writes Bregman. “In reality, for most of human existence, it was the other way around.” Whereas traditional history depicts the collapse of civilisations as “dark ages” in which everything gets worse, modern scholars, he claims, see them more as a reprieve, in which the enslaved gain their freedom and culture flourishes. Like much else in this book, the truth is probably somewhere between the two stated positions.
In any case, the fear of civilisational collapse, Bregman believes, is unfounded. It’s the result of what the Dutch biologist Frans de Waal calls “veneer theory” - the idea that just below the surface, our bestial nature is waiting to break out. . . . There’s a great deal of reassuring human decency to be taken from this bold and thought-provoking book and a wealth of evidence in support of the contention that the sense of who we are as a species has been deleteriously distorted. But it seems equally misleading to offer the false choice of Rousseau and Hobbes when, clearly, humanity encompasses both.

\end{problem}

\newpage
\begin{problem}{}{}
% Subject: varc
\textbf{Instructions}

Instructions   

The passage below is accompanied by a set of questions. Choose the best answer to each question.
Mode of transportation affects the travel experience and thus can produce new types of travel writing and perhaps even new “identities.” Modes of transportation determine the types and duration of social encounters; affect the organization and passage of space and time; . . . and also affect perception and knowledge—how and what the traveler comes to know and write about. The completion of the first U.S. transcontinental highway during the 1920s . . . for example, inaugurated a new genre of travel literature about the United States—the automotive or road narrative. Such narratives highlight the experiences of mostly male protagonists “discovering themselves” on their journeys, emphasizing the independence of road travel and the value of rural folk traditions.
Travel writing’s relationship to empire building— as a type of “colonialist discourse”—has drawn the most attention from academicians. Close connections have been observed between European (and American) political, economic, and administrative goals for the colonies and their manifestations in the cultural practice of writing travel books. Travel writers’ descriptions of foreign places have been analyzed as attempts to validate, promote, or challenge the ideologies and practices of colonial or imperial domination and expansion. Mary Louise Pratt’s study of the genres and conventions of 18th- and 19th-century exploration narratives about South America and Africa (e.g., the “monarch of all I survey” trope) offered ways of thinking about travel writing as embedded within relations of power between metropole and periphery, as did Edward Said’s theories of representation and cultural imperialism. Particularly Said’s book, Orientalism, helped scholars understand ways in which representations of people in travel texts were intimately bound up with notions of self, in this case, that the Occident defined itself through essentialist, ethnocentric, and racist representations of the Orient. Said’s work became a model for demonstrating cultural forms of imperialism in travel texts, showing how the political, economic, or administrative fact of dominance relies on legitimating discourses such as those articulated through travel writing. . . .
Feminist geographers’ studies of travel writing challenge the masculinist history of geography by questioning who and what are relevant subjects of geographic study and, indeed, what counts as geographic knowledge itself. Such questions are worked through ideological constructs that posit men as explorers and women as travelers—or, conversely, men as travelers and women as tied to the home. Studies of Victorian women who were professional travel writers, tourists, wives of colonial administrators, and other (mostly) elite women who wrote narratives about their experiences abroad during the 19th century have been particularly revealing. From a “liberal” feminist perspective, travel presented one means toward female liberation for middle- and upper-class Victorian women. Many studies from the 1970s onward demonstrated the ways in which women’s gendered identities were negotiated differently “at home” than they were “away,” thereby showing women’s self-development through travel. The more recent post structural turn in studies of Victorian travel writing has focused attention on women’s diverse and fragmented identities as they narrated their travel experiences, emphasizing women’s sense of themselves as women in new locations, but only as they worked through their ties to nation, class, whiteness, and colonial and imperial power structures

\vspace{1em}
\textbf{Question 6}

None of the following views is expressed in the passage EXCEPT that:

\begin{enumerate}[label=\Alph*.]
    \item most people agree with Hobbes’ pessimistic view of human nature as being intrinsically untrustworthy and selfish.
    \item Hobbes and Rousseau disagreed on the fundamental nature of humans, but both believed in the need for a strong state.
    \item Bregman agrees with Hobbes that firm leadership is needed to ensure property rights and regulate strife.
    \item the author of the review believes in the veneer theory of human nature.
[](https://cracku.in/6-none-of-the-following-views-is-expressed-in-the-pa-x-cat-2020-slot-3)
\end{enumerate}
\vspace{1em}
\textbf{Solution:}

We need to find a viewpoint that is presented in the passage. Let us inspect the individual options: 
Option A: The introductory lines of the passage helps us infer this: {..._Although one of the most contested concepts in political philosophy, human nature is something on which most people seem to agree. By and large, according to Rutger Bregman in his new book Humankind, we have a rather pessimistic view - not of ourselves exactly, but of everyone else. We see other people as selfish, untrustworthy and dangerous and therefore we behave towards them with defensiveness and suspicion_...}
Option B: The author calls the viewpoints of Hobbes and Rousseau as polar opposites {"_Hobbes and Rousseau are seen as the two poles of the human nature argument_ "} and does not present a similarity, especially any comment of the form: "_both believed in the need for a strong state._ " Thus, we can eliminate this option.
Option C: No such view has been presented. 
Option D: The author's opinion of Frans de Waal's “veneer theory” is not evident/not highlighted. Hence, we can eliminate this option.
Thus, of the given statements, Option A is the correct answer.

\end{problem}

\newpage
\begin{problem}{}{}
% Subject: varc
\textbf{Instructions}

Instructions   

The passage below is accompanied by a set of questions. Choose the best answer to each question.
Mode of transportation affects the travel experience and thus can produce new types of travel writing and perhaps even new “identities.” Modes of transportation determine the types and duration of social encounters; affect the organization and passage of space and time; . . . and also affect perception and knowledge—how and what the traveler comes to know and write about. The completion of the first U.S. transcontinental highway during the 1920s . . . for example, inaugurated a new genre of travel literature about the United States—the automotive or road narrative. Such narratives highlight the experiences of mostly male protagonists “discovering themselves” on their journeys, emphasizing the independence of road travel and the value of rural folk traditions.
Travel writing’s relationship to empire building— as a type of “colonialist discourse”—has drawn the most attention from academicians. Close connections have been observed between European (and American) political, economic, and administrative goals for the colonies and their manifestations in the cultural practice of writing travel books. Travel writers’ descriptions of foreign places have been analyzed as attempts to validate, promote, or challenge the ideologies and practices of colonial or imperial domination and expansion. Mary Louise Pratt’s study of the genres and conventions of 18th- and 19th-century exploration narratives about South America and Africa (e.g., the “monarch of all I survey” trope) offered ways of thinking about travel writing as embedded within relations of power between metropole and periphery, as did Edward Said’s theories of representation and cultural imperialism. Particularly Said’s book, Orientalism, helped scholars understand ways in which representations of people in travel texts were intimately bound up with notions of self, in this case, that the Occident defined itself through essentialist, ethnocentric, and racist representations of the Orient. Said’s work became a model for demonstrating cultural forms of imperialism in travel texts, showing how the political, economic, or administrative fact of dominance relies on legitimating discourses such as those articulated through travel writing. . . .
Feminist geographers’ studies of travel writing challenge the masculinist history of geography by questioning who and what are relevant subjects of geographic study and, indeed, what counts as geographic knowledge itself. Such questions are worked through ideological constructs that posit men as explorers and women as travelers—or, conversely, men as travelers and women as tied to the home. Studies of Victorian women who were professional travel writers, tourists, wives of colonial administrators, and other (mostly) elite women who wrote narratives about their experiences abroad during the 19th century have been particularly revealing. From a “liberal” feminist perspective, travel presented one means toward female liberation for middle- and upper-class Victorian women. Many studies from the 1970s onward demonstrated the ways in which women’s gendered identities were negotiated differently “at home” than they were “away,” thereby showing women’s self-development through travel. The more recent post structural turn in studies of Victorian travel writing has focused attention on women’s diverse and fragmented identities as they narrated their travel experiences, emphasizing women’s sense of themselves as women in new locations, but only as they worked through their ties to nation, class, whiteness, and colonial and imperial power structures

\vspace{1em}
\textbf{Question 7}

According to the passage, the “collapse of civilisations” is viewed by Bregman as:

\begin{enumerate}[label=\Alph*.]
    \item a time that enables changes in societies and cultures.
    \item a sign of regression in society’s trajectory.
    \item resulting from a breakdown in the veneer of human nature.
    \item a temporary phase which can be rectified by social action.
[](https://cracku.in/7-according-to-the-passage-the-collapse-of-civilisat-x-cat-2020-slot-3)
\end{enumerate}
\vspace{1em}
\textbf{Solution:}

Bregman considers the aftermath of civilizational collapse as a period that allows for certain changes or alterations in the society {.._.“Civilisation has become synonymous with peace and progress and wilderness with war and decline,” writes Bregman. “In reality, for most of human existence, it was the other way around.” Whereas traditional history depicts the collapse of civilisations as “dark ages” in which everything gets worse, modern scholars, he claims, see them more as a reprieve, in which the enslaved gain their freedom and culture flourishes._.. }. Option A correctly captures this point. Options B, C and D are either not stated or distorted interpretations.

\end{problem}

\newpage
\begin{problem}{}{}
% Subject: varc
\textbf{Instructions}

Instructions   

The passage below is accompanied by a set of questions. Choose the best answer to each question.
Mode of transportation affects the travel experience and thus can produce new types of travel writing and perhaps even new “identities.” Modes of transportation determine the types and duration of social encounters; affect the organization and passage of space and time; . . . and also affect perception and knowledge—how and what the traveler comes to know and write about. The completion of the first U.S. transcontinental highway during the 1920s . . . for example, inaugurated a new genre of travel literature about the United States—the automotive or road narrative. Such narratives highlight the experiences of mostly male protagonists “discovering themselves” on their journeys, emphasizing the independence of road travel and the value of rural folk traditions.
Travel writing’s relationship to empire building— as a type of “colonialist discourse”—has drawn the most attention from academicians. Close connections have been observed between European (and American) political, economic, and administrative goals for the colonies and their manifestations in the cultural practice of writing travel books. Travel writers’ descriptions of foreign places have been analyzed as attempts to validate, promote, or challenge the ideologies and practices of colonial or imperial domination and expansion. Mary Louise Pratt’s study of the genres and conventions of 18th- and 19th-century exploration narratives about South America and Africa (e.g., the “monarch of all I survey” trope) offered ways of thinking about travel writing as embedded within relations of power between metropole and periphery, as did Edward Said’s theories of representation and cultural imperialism. Particularly Said’s book, Orientalism, helped scholars understand ways in which representations of people in travel texts were intimately bound up with notions of self, in this case, that the Occident defined itself through essentialist, ethnocentric, and racist representations of the Orient. Said’s work became a model for demonstrating cultural forms of imperialism in travel texts, showing how the political, economic, or administrative fact of dominance relies on legitimating discourses such as those articulated through travel writing. . . .
Feminist geographers’ studies of travel writing challenge the masculinist history of geography by questioning who and what are relevant subjects of geographic study and, indeed, what counts as geographic knowledge itself. Such questions are worked through ideological constructs that posit men as explorers and women as travelers—or, conversely, men as travelers and women as tied to the home. Studies of Victorian women who were professional travel writers, tourists, wives of colonial administrators, and other (mostly) elite women who wrote narratives about their experiences abroad during the 19th century have been particularly revealing. From a “liberal” feminist perspective, travel presented one means toward female liberation for middle- and upper-class Victorian women. Many studies from the 1970s onward demonstrated the ways in which women’s gendered identities were negotiated differently “at home” than they were “away,” thereby showing women’s self-development through travel. The more recent post structural turn in studies of Victorian travel writing has focused attention on women’s diverse and fragmented identities as they narrated their travel experiences, emphasizing women’s sense of themselves as women in new locations, but only as they worked through their ties to nation, class, whiteness, and colonial and imperial power structures

\vspace{1em}
\textbf{Question 8}

According to the author, the main reason why Bregman contrasts life in pre-agricultural societies with agricultural societies is to:

\begin{enumerate}[label=\Alph*.]
    \item advocate the promotion of less complex societies as a basis for greater security and prosperity
    \item highlight the enormous impact that settled farming had on population growth
    \item make the argument that an environmentally conscious lifestyle is a more harmonious way of living.
    \item bolster his argument that people are basically decent, but progress as we know it can make them selfish.
[](https://cracku.in/8-according-to-the-author-the-main-reason-why-bregma-x-cat-2020-slot-3)
\end{enumerate}
\vspace{1em}
\textbf{Solution:}

Bregman disagrees with Hobbes' standpoint of humans being inherently selfish or bestial and instead takes Rousseau's side. He asserts that civilizational progress caused by the post-agricultural setup is responsible for the negative/undesired circumstances. In this regard, he presents the contrasting picture of pre and post-agricultural societies {attaches the image of "a prelapsarian idyll" to the nomadic lifestyle, while considers the discovery of agriculture as a misevent}. Thus, this depiction supplements "his argument that people are basically decent, but progress as we know it can make them selfish." _Option D_ is the appropriate answer. 
_Option A_ : The aspect of complexity is not the primary focal point. Thus, this option can be eliminated.
_Option B_ : This diverges from the discussion onto a new line of discussion: "impact that settled farming had on population growth". Hence, we can discard this choice as well.
_Option C_ : Again, the focus is not on the environment; hence, we can scrap off this option.

\end{problem}

\newpage
\begin{problem}{}{}
% Subject: varc
\textbf{Instructions}

Instructions   

The passage below is accompanied by a set of questions. Choose the best answer to each question.
Mode of transportation affects the travel experience and thus can produce new types of travel writing and perhaps even new “identities.” Modes of transportation determine the types and duration of social encounters; affect the organization and passage of space and time; . . . and also affect perception and knowledge—how and what the traveler comes to know and write about. The completion of the first U.S. transcontinental highway during the 1920s . . . for example, inaugurated a new genre of travel literature about the United States—the automotive or road narrative. Such narratives highlight the experiences of mostly male protagonists “discovering themselves” on their journeys, emphasizing the independence of road travel and the value of rural folk traditions.
Travel writing’s relationship to empire building— as a type of “colonialist discourse”—has drawn the most attention from academicians. Close connections have been observed between European (and American) political, economic, and administrative goals for the colonies and their manifestations in the cultural practice of writing travel books. Travel writers’ descriptions of foreign places have been analyzed as attempts to validate, promote, or challenge the ideologies and practices of colonial or imperial domination and expansion. Mary Louise Pratt’s study of the genres and conventions of 18th- and 19th-century exploration narratives about South America and Africa (e.g., the “monarch of all I survey” trope) offered ways of thinking about travel writing as embedded within relations of power between metropole and periphery, as did Edward Said’s theories of representation and cultural imperialism. Particularly Said’s book, Orientalism, helped scholars understand ways in which representations of people in travel texts were intimately bound up with notions of self, in this case, that the Occident defined itself through essentialist, ethnocentric, and racist representations of the Orient. Said’s work became a model for demonstrating cultural forms of imperialism in travel texts, showing how the political, economic, or administrative fact of dominance relies on legitimating discourses such as those articulated through travel writing. . . .
Feminist geographers’ studies of travel writing challenge the masculinist history of geography by questioning who and what are relevant subjects of geographic study and, indeed, what counts as geographic knowledge itself. Such questions are worked through ideological constructs that posit men as explorers and women as travelers—or, conversely, men as travelers and women as tied to the home. Studies of Victorian women who were professional travel writers, tourists, wives of colonial administrators, and other (mostly) elite women who wrote narratives about their experiences abroad during the 19th century have been particularly revealing. From a “liberal” feminist perspective, travel presented one means toward female liberation for middle- and upper-class Victorian women. Many studies from the 1970s onward demonstrated the ways in which women’s gendered identities were negotiated differently “at home” than they were “away,” thereby showing women’s self-development through travel. The more recent post structural turn in studies of Victorian travel writing has focused attention on women’s diverse and fragmented identities as they narrated their travel experiences, emphasizing women’s sense of themselves as women in new locations, but only as they worked through their ties to nation, class, whiteness, and colonial and imperial power structures

\vspace{1em}
\textbf{Question 9}

The author has differing views from Bregman regarding:

\begin{enumerate}[label=\Alph*.]
    \item the role of agriculture in the advancement of knowledge.
    \item the role of pathogens in the spread of infectious diseases.
    \item a property-less mode of living being socially harmonious.
    \item a civilised society being coercive and unjust.
[](https://cracku.in/9-the-author-has-differing-views-from-bregman-regard-x-cat-2020-slot-3)
\end{enumerate}
\vspace{1em}
\textbf{Solution:}

At the end of the passage, the author states the following: {..._There’s a great deal of reassuring human decency to be taken from this bold and thought-provoking book and a wealth of evidence in support of the contention that the sense of who we are as a species has been deleteriously distorted. But it seems equally misleading to offer the false choice of Rousseau and Hobbes when, clearly, humanity encompasses both.._.} Thus, he does not truly agree with Bregman's portrayal of the civilized society. Option D correctly captures this disagreement.

Instructions   

The passage below is accompanied by a set of questions. Choose the best answer to each question.
I’ve been following the economic crisis for more than two years now. I began working on the subject as part of the background to a novel, and soon realized that I had stumbled across the most interesting story I’ve ever found. While I was beginning to work on it, the British bank Northern Rock blew up, and it became clear that, as I wrote at the time, “If our laws are not extended to control the new kinds of super-powerful, super-complex, and potentially super risky investment vehicles, they will one day cause a financial disaster of global-systemic proportions.” . . . I was both right and too late, because all the groundwork for the crisis had already been done—though the sluggishness of the world’s governments, in not preparing for the great unraveling of autumn 2008, was then and still is stupefying. But this is the first reason why I wrote this book: because what’s happened is extraordinarily interesting. It is an absolutely amazing story, full of human interest and drama, one whose byways of mathematics, economics, and psychology are both central to the story of the last decades and mysteriously unknown to the general public. We have heard a lot about “the two cultures” of science and the arts—we heard a particularly large amount about it in 2009, because it was the fiftieth anniversary of the speech during which C. P. Snow first used the phrase. But I’m not sure the idea of a huge gap between science and the arts is as true as it was half a century ago—it’s certainly true, for instance, that a general reader who wants to pick up an education in the fundamentals of science will find it easier than ever before. It seems to me that there is a much bigger gap between the world of finance and that of the general public and that there is a need to narrow that gap, if the financial industry is not to be a kind of priesthood, administering to its own mysteries and feared and resented by the rest of us. Many bright, literate people have no idea about all sorts of economic basics, of a type that financial insiders take as elementary facts of how the world works. I am an outsider to finance and economics, and my hope is that I can talk across that gulf.
My need to understand is the same as yours, whoever you are. That’s one of the strangest ironies of this story: after decades in which the ideology of the Western world was personally and economically individualistic, we’ve suddenly been hit by a crisis which shows in the starkest terms that whether we like it or not—and there are large parts of it that you would have to be crazy to like—we’re all in this together. The aftermath of the crisis is going to dominate the economics and politics of our societies for at least a decade to come and perhaps longer.

\end{problem}

\newpage
\begin{problem}{}{}
% Subject: varc
\textbf{Instructions}

Instructions   

The passage below is accompanied by a set of questions. Choose the best answer to each question.
Mode of transportation affects the travel experience and thus can produce new types of travel writing and perhaps even new “identities.” Modes of transportation determine the types and duration of social encounters; affect the organization and passage of space and time; . . . and also affect perception and knowledge—how and what the traveler comes to know and write about. The completion of the first U.S. transcontinental highway during the 1920s . . . for example, inaugurated a new genre of travel literature about the United States—the automotive or road narrative. Such narratives highlight the experiences of mostly male protagonists “discovering themselves” on their journeys, emphasizing the independence of road travel and the value of rural folk traditions.
Travel writing’s relationship to empire building— as a type of “colonialist discourse”—has drawn the most attention from academicians. Close connections have been observed between European (and American) political, economic, and administrative goals for the colonies and their manifestations in the cultural practice of writing travel books. Travel writers’ descriptions of foreign places have been analyzed as attempts to validate, promote, or challenge the ideologies and practices of colonial or imperial domination and expansion. Mary Louise Pratt’s study of the genres and conventions of 18th- and 19th-century exploration narratives about South America and Africa (e.g., the “monarch of all I survey” trope) offered ways of thinking about travel writing as embedded within relations of power between metropole and periphery, as did Edward Said’s theories of representation and cultural imperialism. Particularly Said’s book, Orientalism, helped scholars understand ways in which representations of people in travel texts were intimately bound up with notions of self, in this case, that the Occident defined itself through essentialist, ethnocentric, and racist representations of the Orient. Said’s work became a model for demonstrating cultural forms of imperialism in travel texts, showing how the political, economic, or administrative fact of dominance relies on legitimating discourses such as those articulated through travel writing. . . .
Feminist geographers’ studies of travel writing challenge the masculinist history of geography by questioning who and what are relevant subjects of geographic study and, indeed, what counts as geographic knowledge itself. Such questions are worked through ideological constructs that posit men as explorers and women as travelers—or, conversely, men as travelers and women as tied to the home. Studies of Victorian women who were professional travel writers, tourists, wives of colonial administrators, and other (mostly) elite women who wrote narratives about their experiences abroad during the 19th century have been particularly revealing. From a “liberal” feminist perspective, travel presented one means toward female liberation for middle- and upper-class Victorian women. Many studies from the 1970s onward demonstrated the ways in which women’s gendered identities were negotiated differently “at home” than they were “away,” thereby showing women’s self-development through travel. The more recent post structural turn in studies of Victorian travel writing has focused attention on women’s diverse and fragmented identities as they narrated their travel experiences, emphasizing women’s sense of themselves as women in new locations, but only as they worked through their ties to nation, class, whiteness, and colonial and imperial power structures

\vspace{1em}
\textbf{Question 10}

Which one of the following, if true, would be an accurate inference from the first  
sentence of the passage?

\begin{enumerate}[label=\Alph*.]
    \item The author’s preoccupation with the economic crisis is not less than two years old.
    \item The author is preoccupied with the economic crisis because he is being followed.
    \item The economic crisis outlasted the author’s preoccupation with it.
    \item The author has witnessed many economic crises by travelling a lot for two years
[](https://cracku.in/10-which-one-of-the-following-if-true-would-be-an-acc-x-cat-2020-slot-3)
\end{enumerate}
\vspace{1em}
\textbf{Solution:}

This is a direct inference question that requires minimal effort and can be accurately answered using the option-elimination mechanism. The first sentence of the passage is as follows: " .._.I’ve been following the economic crisis for more than two years now._..". It is evident that the author has been following {events/information associated with} the economic crisis for at least two years if not more. Thus, Option A is a sensible inference to draw from this statement. Options B, C and D are inane interpretations of the same and can be effortlessly eliminated.

\end{problem}

\newpage
\begin{problem}{}{}
% Subject: varc
\textbf{Instructions}

Instructions   

The passage below is accompanied by a set of questions. Choose the best answer to each question.
Mode of transportation affects the travel experience and thus can produce new types of travel writing and perhaps even new “identities.” Modes of transportation determine the types and duration of social encounters; affect the organization and passage of space and time; . . . and also affect perception and knowledge—how and what the traveler comes to know and write about. The completion of the first U.S. transcontinental highway during the 1920s . . . for example, inaugurated a new genre of travel literature about the United States—the automotive or road narrative. Such narratives highlight the experiences of mostly male protagonists “discovering themselves” on their journeys, emphasizing the independence of road travel and the value of rural folk traditions.
Travel writing’s relationship to empire building— as a type of “colonialist discourse”—has drawn the most attention from academicians. Close connections have been observed between European (and American) political, economic, and administrative goals for the colonies and their manifestations in the cultural practice of writing travel books. Travel writers’ descriptions of foreign places have been analyzed as attempts to validate, promote, or challenge the ideologies and practices of colonial or imperial domination and expansion. Mary Louise Pratt’s study of the genres and conventions of 18th- and 19th-century exploration narratives about South America and Africa (e.g., the “monarch of all I survey” trope) offered ways of thinking about travel writing as embedded within relations of power between metropole and periphery, as did Edward Said’s theories of representation and cultural imperialism. Particularly Said’s book, Orientalism, helped scholars understand ways in which representations of people in travel texts were intimately bound up with notions of self, in this case, that the Occident defined itself through essentialist, ethnocentric, and racist representations of the Orient. Said’s work became a model for demonstrating cultural forms of imperialism in travel texts, showing how the political, economic, or administrative fact of dominance relies on legitimating discourses such as those articulated through travel writing. . . .
Feminist geographers’ studies of travel writing challenge the masculinist history of geography by questioning who and what are relevant subjects of geographic study and, indeed, what counts as geographic knowledge itself. Such questions are worked through ideological constructs that posit men as explorers and women as travelers—or, conversely, men as travelers and women as tied to the home. Studies of Victorian women who were professional travel writers, tourists, wives of colonial administrators, and other (mostly) elite women who wrote narratives about their experiences abroad during the 19th century have been particularly revealing. From a “liberal” feminist perspective, travel presented one means toward female liberation for middle- and upper-class Victorian women. Many studies from the 1970s onward demonstrated the ways in which women’s gendered identities were negotiated differently “at home” than they were “away,” thereby showing women’s self-development through travel. The more recent post structural turn in studies of Victorian travel writing has focused attention on women’s diverse and fragmented identities as they narrated their travel experiences, emphasizing women’s sense of themselves as women in new locations, but only as they worked through their ties to nation, class, whiteness, and colonial and imperial power structures

\vspace{1em}
\textbf{Question 11}

Which one of the following best captures the main argument of the last paragraph of  
the passage?

\begin{enumerate}[label=\Alph*.]
    \item The aftermath of the crisis will strengthen the central ideology of individualism in the  
Western world
    \item Whoever you are, you would be crazy to think that there is no crisis.
    \item In the decades to come, other ideologies will emerge in the aftermath of the crisis.
    \item The ideology of individualism must be set aside in order to deal with the crisis.
[](https://cracku.in/11-which-one-of-the-following-best-captures-the-main--x-cat-2020-slot-3)
\end{enumerate}
\vspace{1em}
\textbf{Solution:}

In the final paragraph, the author highlights the crisis as an ironical situation: a group of individualistic entities facing an issue with collective impact and thereby, needs to be dealt with "together" or by shifting from the existing self-centred setup {".._.That’s one of the strangest ironies of this story: after decades in which the ideology of the Western world was personally and economically individualistic, we’ve suddenly been hit by a crisis which shows in the starkest terms that whether we like it or not—and there are large parts of it that you would have to be crazy to like—we’re all in this together...._ "}. Option D is the closest choice that captures this element. Option A is contrary to the point presented in the concluding para. Options B and C are either divergent to the point made or merely distorted comments. Hence, Option D is the correct answer.

\end{problem}

\newpage
\begin{problem}{}{}
% Subject: varc
\textbf{Instructions}

Instructions   

The passage below is accompanied by a set of questions. Choose the best answer to each question.
Mode of transportation affects the travel experience and thus can produce new types of travel writing and perhaps even new “identities.” Modes of transportation determine the types and duration of social encounters; affect the organization and passage of space and time; . . . and also affect perception and knowledge—how and what the traveler comes to know and write about. The completion of the first U.S. transcontinental highway during the 1920s . . . for example, inaugurated a new genre of travel literature about the United States—the automotive or road narrative. Such narratives highlight the experiences of mostly male protagonists “discovering themselves” on their journeys, emphasizing the independence of road travel and the value of rural folk traditions.
Travel writing’s relationship to empire building— as a type of “colonialist discourse”—has drawn the most attention from academicians. Close connections have been observed between European (and American) political, economic, and administrative goals for the colonies and their manifestations in the cultural practice of writing travel books. Travel writers’ descriptions of foreign places have been analyzed as attempts to validate, promote, or challenge the ideologies and practices of colonial or imperial domination and expansion. Mary Louise Pratt’s study of the genres and conventions of 18th- and 19th-century exploration narratives about South America and Africa (e.g., the “monarch of all I survey” trope) offered ways of thinking about travel writing as embedded within relations of power between metropole and periphery, as did Edward Said’s theories of representation and cultural imperialism. Particularly Said’s book, Orientalism, helped scholars understand ways in which representations of people in travel texts were intimately bound up with notions of self, in this case, that the Occident defined itself through essentialist, ethnocentric, and racist representations of the Orient. Said’s work became a model for demonstrating cultural forms of imperialism in travel texts, showing how the political, economic, or administrative fact of dominance relies on legitimating discourses such as those articulated through travel writing. . . .
Feminist geographers’ studies of travel writing challenge the masculinist history of geography by questioning who and what are relevant subjects of geographic study and, indeed, what counts as geographic knowledge itself. Such questions are worked through ideological constructs that posit men as explorers and women as travelers—or, conversely, men as travelers and women as tied to the home. Studies of Victorian women who were professional travel writers, tourists, wives of colonial administrators, and other (mostly) elite women who wrote narratives about their experiences abroad during the 19th century have been particularly revealing. From a “liberal” feminist perspective, travel presented one means toward female liberation for middle- and upper-class Victorian women. Many studies from the 1970s onward demonstrated the ways in which women’s gendered identities were negotiated differently “at home” than they were “away,” thereby showing women’s self-development through travel. The more recent post structural turn in studies of Victorian travel writing has focused attention on women’s diverse and fragmented identities as they narrated their travel experiences, emphasizing women’s sense of themselves as women in new locations, but only as they worked through their ties to nation, class, whiteness, and colonial and imperial power structures

\vspace{1em}
\textbf{Question 12}

According to the passage, the author is likely to be supportive of which one of the  
following programmes?

\begin{enumerate}[label=\Alph*.]
    \item Economic policies that are more sensitively calibrated to the fluctuations of the  
market.
    \item An educational curriculum that promotes economic research
    \item An educational curriculum that promotes developing financial literacy in the masses.
    \item The complete nationalisation of all financial institutions
[](https://cracku.in/12-according-to-the-passage-the-author-is-likely-to-b-x-cat-2020-slot-3)
\end{enumerate}
\vspace{1em}
\textbf{Solution:}

A predominant idea discussed by the author is regarding the lack of financial literacy that could be truly beneficial to our understanding of the world. This is emphasised via the following excerpt: 
{._..It seems to me that there is a much bigger gap between the world of finance and that of the general public and that there is a need to narrow that gap, if the financial industry is not to be a kind of priesthood, administering to its own mysteries and feared and resented by the rest of us. Many bright, literate people have no idea about all sorts of economic basics, of a type that financial insiders take as elementary facts of how the world works. I am an outsider to finance and economics, and my hope is that I can talk across that gulf...._} 
Option C aligns with this concern and is, consequently, an idea that the author is bound to support. Option B appears as another likely candidate; however, economic research is not one of the focal points mentioned. Options A and D contain elements that are not discussed or are opposite to the author's ideas. Hence, Option C is a programme that the author is most likely to be supportive of.

\end{problem}

\newpage
\begin{problem}{}{}
% Subject: varc
\textbf{Instructions}

Instructions   

The passage below is accompanied by a set of questions. Choose the best answer to each question.
Mode of transportation affects the travel experience and thus can produce new types of travel writing and perhaps even new “identities.” Modes of transportation determine the types and duration of social encounters; affect the organization and passage of space and time; . . . and also affect perception and knowledge—how and what the traveler comes to know and write about. The completion of the first U.S. transcontinental highway during the 1920s . . . for example, inaugurated a new genre of travel literature about the United States—the automotive or road narrative. Such narratives highlight the experiences of mostly male protagonists “discovering themselves” on their journeys, emphasizing the independence of road travel and the value of rural folk traditions.
Travel writing’s relationship to empire building— as a type of “colonialist discourse”—has drawn the most attention from academicians. Close connections have been observed between European (and American) political, economic, and administrative goals for the colonies and their manifestations in the cultural practice of writing travel books. Travel writers’ descriptions of foreign places have been analyzed as attempts to validate, promote, or challenge the ideologies and practices of colonial or imperial domination and expansion. Mary Louise Pratt’s study of the genres and conventions of 18th- and 19th-century exploration narratives about South America and Africa (e.g., the “monarch of all I survey” trope) offered ways of thinking about travel writing as embedded within relations of power between metropole and periphery, as did Edward Said’s theories of representation and cultural imperialism. Particularly Said’s book, Orientalism, helped scholars understand ways in which representations of people in travel texts were intimately bound up with notions of self, in this case, that the Occident defined itself through essentialist, ethnocentric, and racist representations of the Orient. Said’s work became a model for demonstrating cultural forms of imperialism in travel texts, showing how the political, economic, or administrative fact of dominance relies on legitimating discourses such as those articulated through travel writing. . . .
Feminist geographers’ studies of travel writing challenge the masculinist history of geography by questioning who and what are relevant subjects of geographic study and, indeed, what counts as geographic knowledge itself. Such questions are worked through ideological constructs that posit men as explorers and women as travelers—or, conversely, men as travelers and women as tied to the home. Studies of Victorian women who were professional travel writers, tourists, wives of colonial administrators, and other (mostly) elite women who wrote narratives about their experiences abroad during the 19th century have been particularly revealing. From a “liberal” feminist perspective, travel presented one means toward female liberation for middle- and upper-class Victorian women. Many studies from the 1970s onward demonstrated the ways in which women’s gendered identities were negotiated differently “at home” than they were “away,” thereby showing women’s self-development through travel. The more recent post structural turn in studies of Victorian travel writing has focused attention on women’s diverse and fragmented identities as they narrated their travel experiences, emphasizing women’s sense of themselves as women in new locations, but only as they worked through their ties to nation, class, whiteness, and colonial and imperial power structures

\vspace{1em}
\textbf{Question 13}

All of the following, if true, could be seen as supporting the arguments in the passage, EXCEPT:

\begin{enumerate}[label=\Alph*.]
    \item The failure of economic systems does not necessarily mean the failure of their ideologies.
    \item The difficulty with understanding financial matters is that they have become so arcane.
    \item The story of the economic crisis is also one about international relations, global financial security, and mass psychology.
    \item Economic crises could be averted by changing prevailing ideologies and beliefs.
[](https://cracku.in/13-all-of-the-following-if-true-could-be-seen-as-supp-x-cat-2020-slot-3)
\end{enumerate}
\vspace{1em}
\textbf{Solution:}

Let us inspect the individual options:
Option A: If true, this statement could be antithetical to the point put forth in the concluding paragraph: the author believes that an individualistic ideology isn't the right way forward; instead, we need to deal with such crises collectively. Thus, this is a conflicting viewpoint and thereby, the correct answer. 
Option B: If true, aligns with the author's point in the first paragraph: {.._.It seems to me that there is a much bigger gap between the world of finance and that of the general public and that there is a need to narrow that gap, if the financial industry is not to be a kind of priesthood, administering to its own mysteries and feared and resented by the rest of us._..} 
Options C and D: If true, this is in tune with the author's claim {made in the final segment} about dealing with such crises together. 
Hence, of the given statements, Options A deviates from the author's argument and thus, is the correct answer.

\end{problem}

\newpage
\begin{problem}{}{}
% Subject: varc
\textbf{Instructions}

Instructions   

The passage below is accompanied by a set of questions. Choose the best answer to each question.
Mode of transportation affects the travel experience and thus can produce new types of travel writing and perhaps even new “identities.” Modes of transportation determine the types and duration of social encounters; affect the organization and passage of space and time; . . . and also affect perception and knowledge—how and what the traveler comes to know and write about. The completion of the first U.S. transcontinental highway during the 1920s . . . for example, inaugurated a new genre of travel literature about the United States—the automotive or road narrative. Such narratives highlight the experiences of mostly male protagonists “discovering themselves” on their journeys, emphasizing the independence of road travel and the value of rural folk traditions.
Travel writing’s relationship to empire building— as a type of “colonialist discourse”—has drawn the most attention from academicians. Close connections have been observed between European (and American) political, economic, and administrative goals for the colonies and their manifestations in the cultural practice of writing travel books. Travel writers’ descriptions of foreign places have been analyzed as attempts to validate, promote, or challenge the ideologies and practices of colonial or imperial domination and expansion. Mary Louise Pratt’s study of the genres and conventions of 18th- and 19th-century exploration narratives about South America and Africa (e.g., the “monarch of all I survey” trope) offered ways of thinking about travel writing as embedded within relations of power between metropole and periphery, as did Edward Said’s theories of representation and cultural imperialism. Particularly Said’s book, Orientalism, helped scholars understand ways in which representations of people in travel texts were intimately bound up with notions of self, in this case, that the Occident defined itself through essentialist, ethnocentric, and racist representations of the Orient. Said’s work became a model for demonstrating cultural forms of imperialism in travel texts, showing how the political, economic, or administrative fact of dominance relies on legitimating discourses such as those articulated through travel writing. . . .
Feminist geographers’ studies of travel writing challenge the masculinist history of geography by questioning who and what are relevant subjects of geographic study and, indeed, what counts as geographic knowledge itself. Such questions are worked through ideological constructs that posit men as explorers and women as travelers—or, conversely, men as travelers and women as tied to the home. Studies of Victorian women who were professional travel writers, tourists, wives of colonial administrators, and other (mostly) elite women who wrote narratives about their experiences abroad during the 19th century have been particularly revealing. From a “liberal” feminist perspective, travel presented one means toward female liberation for middle- and upper-class Victorian women. Many studies from the 1970s onward demonstrated the ways in which women’s gendered identities were negotiated differently “at home” than they were “away,” thereby showing women’s self-development through travel. The more recent post structural turn in studies of Victorian travel writing has focused attention on women’s diverse and fragmented identities as they narrated their travel experiences, emphasizing women’s sense of themselves as women in new locations, but only as they worked through their ties to nation, class, whiteness, and colonial and imperial power structures

\vspace{1em}
\textbf{Question 14}

Which one of the following, if false, could be seen as supporting the author’s claims?

\begin{enumerate}[label=\Alph*.]
    \item The huge gap between science and the arts has steadily narrowed over time.
    \item The economic crisis was not a failure of collective action to rectify economic problems.
    \item Most people are yet to gain any real understanding of the workings of the financial world.
    \item The global economic crisis lasted for more than two years.
[](https://cracku.in/14-which-one-of-the-following-if-false-could-be-seen--x-cat-2020-slot-3)
\end{enumerate}
\vspace{1em}
\textbf{Solution:}

On skimming through the statements, it is evident that Options A and D do very little, if at all anything, to support the author's claim. Option C, in its current form, aligns with the author's assertion; however, if false, it is opposite to the author's argument about financial literacy {on negating the statement, it indicates that most people do have an understanding of the workings of the financial world}. Option B, if false, is in line with the assertions made in the concluding paragraph.

Instructions   

**The passage below is accompanied by a set of questions. Choose the best answer to each question.**
[There is] a curious new reality: Human contact is becoming a luxury good. As more screens appear in the lives of the poor, screens are disappearing from the lives of the rich. The richer you are, the more you spend to be off-screen. . . . 
The joy — at least at first — of the internet revolution was its democratic nature. Facebook is the same Facebook whether you are rich or poor. Gmail is the same Gmail. And it’s all free. There is something mass market and unappealing about that. And as studies show that time on these advertisement-support platforms is unhealthy, it all starts to seem déclassé, like drinking soda or smoking cigarettes, which wealthy people do less than poor people. The wealthy can afford to opt out of having their data and their attention sold as a product. The poor and middle class don’t have the same kind of resources to make that happen.
Screen exposure starts young. And children who spent more than two hours a day looking at a screen got lower scores on thinking and language tests, according to early results of a landmark study on brain development of more than 11,000 children that the National Institutes of Health is supporting. Most disturbingly, the study is finding that the brains of children who spend a lot of time on screens are different. For some kids, there is premature thinning of their cerebral cortex. In adults, one study found an association between screen time and depression. . . .
Tech companies worked hard to get public schools to buy into programs that required schools to have one laptop per student, arguing that it would better prepare children for their screen-based future. But this idea isn’t how the people who actually build the screen-based future raise their own children. In Silicon Valley, time on screens is increasingly seen as unhealthy. Here, the popular elementary school is the local Waldorf School, which promises a back-to-nature, nearly screen-free education. So as wealthy kids are growing up with less screen time, poor kids are growing up with more. How comfortable someone is with human engagement could become a new class marker.
Human contact is, of course, not exactly like organic food . . . . But with screen time, there has been a concerted effort on the part of Silicon Valley behemoths to confuse the public. The poor and the middle class are told that screens are good and important for them and their children. There are fleets of psychologists and neuroscientists on staff at big tech companies working to hook eyes and minds to the screen as fast as possible and for as long as possible. And so human contact is rare. . . . 
There is a small movement to pass a “right to disconnect” bill, which would allow workers to turn their phones off, but for now, a worker can be punished for going offline and not being available. There is also the reality that in our culture of increasing isolation, in which so many of the traditional gathering places and social structures have disappeared, screens are filling a crucial void.

\end{problem}

\newpage
\begin{problem}{}{}
% Subject: varc
\textbf{Instructions}

Instructions   

The passage below is accompanied by a set of questions. Choose the best answer to each question.
Mode of transportation affects the travel experience and thus can produce new types of travel writing and perhaps even new “identities.” Modes of transportation determine the types and duration of social encounters; affect the organization and passage of space and time; . . . and also affect perception and knowledge—how and what the traveler comes to know and write about. The completion of the first U.S. transcontinental highway during the 1920s . . . for example, inaugurated a new genre of travel literature about the United States—the automotive or road narrative. Such narratives highlight the experiences of mostly male protagonists “discovering themselves” on their journeys, emphasizing the independence of road travel and the value of rural folk traditions.
Travel writing’s relationship to empire building— as a type of “colonialist discourse”—has drawn the most attention from academicians. Close connections have been observed between European (and American) political, economic, and administrative goals for the colonies and their manifestations in the cultural practice of writing travel books. Travel writers’ descriptions of foreign places have been analyzed as attempts to validate, promote, or challenge the ideologies and practices of colonial or imperial domination and expansion. Mary Louise Pratt’s study of the genres and conventions of 18th- and 19th-century exploration narratives about South America and Africa (e.g., the “monarch of all I survey” trope) offered ways of thinking about travel writing as embedded within relations of power between metropole and periphery, as did Edward Said’s theories of representation and cultural imperialism. Particularly Said’s book, Orientalism, helped scholars understand ways in which representations of people in travel texts were intimately bound up with notions of self, in this case, that the Occident defined itself through essentialist, ethnocentric, and racist representations of the Orient. Said’s work became a model for demonstrating cultural forms of imperialism in travel texts, showing how the political, economic, or administrative fact of dominance relies on legitimating discourses such as those articulated through travel writing. . . .
Feminist geographers’ studies of travel writing challenge the masculinist history of geography by questioning who and what are relevant subjects of geographic study and, indeed, what counts as geographic knowledge itself. Such questions are worked through ideological constructs that posit men as explorers and women as travelers—or, conversely, men as travelers and women as tied to the home. Studies of Victorian women who were professional travel writers, tourists, wives of colonial administrators, and other (mostly) elite women who wrote narratives about their experiences abroad during the 19th century have been particularly revealing. From a “liberal” feminist perspective, travel presented one means toward female liberation for middle- and upper-class Victorian women. Many studies from the 1970s onward demonstrated the ways in which women’s gendered identities were negotiated differently “at home” than they were “away,” thereby showing women’s self-development through travel. The more recent post structural turn in studies of Victorian travel writing has focused attention on women’s diverse and fragmented identities as they narrated their travel experiences, emphasizing women’s sense of themselves as women in new locations, but only as they worked through their ties to nation, class, whiteness, and colonial and imperial power structures

\vspace{1em}
\textbf{Question 15}

Which of the following statements about the negative effects of screen time is the author least likely to endorse?

\begin{enumerate}[label=\Alph*.]
    \item It is shown to have adverse effects on young children’s learning
    \item It increases human contact as it fills an isolation void.
    \item It can cause depression in viewers.
    \item It is designed to be addictive.
[](https://cracku.in/15-which-of-the-following-statements-about-the-negati-x-cat-2020-slot-3)
\end{enumerate}
\vspace{1em}
\textbf{Solution:}

_Option A_ : This author will agree with this assertion. It has been mentioned in the third para: {.._. For some kids, there is premature thinning of their cerebral cortex..._}.
_Option B_ : Reduction in human contact or social engagement is one of the negative impacts of screen time that the author highlights in the passage. Option B is antithetical to this presentation and is, hence, not a claim that the author is likely to endorse. 
_Option C_ : The author states this element in the third para: {.._. In adults, one study found an association between screen time and depression..._}  
  
_Option D_ : A point along similar lines has been presented in the fifth para: {.._.There are fleets of psychologists and neuroscientists on staff at big tech companies working to hook eyes and minds to the screen as fast as possible and for as long as possible. And so human contact is rare._..}
Hence, Option B is the correct answer.

\end{problem}

\newpage
\begin{problem}{}{}
% Subject: varc
\textbf{Instructions}

Instructions   

The passage below is accompanied by a set of questions. Choose the best answer to each question.
Mode of transportation affects the travel experience and thus can produce new types of travel writing and perhaps even new “identities.” Modes of transportation determine the types and duration of social encounters; affect the organization and passage of space and time; . . . and also affect perception and knowledge—how and what the traveler comes to know and write about. The completion of the first U.S. transcontinental highway during the 1920s . . . for example, inaugurated a new genre of travel literature about the United States—the automotive or road narrative. Such narratives highlight the experiences of mostly male protagonists “discovering themselves” on their journeys, emphasizing the independence of road travel and the value of rural folk traditions.
Travel writing’s relationship to empire building— as a type of “colonialist discourse”—has drawn the most attention from academicians. Close connections have been observed between European (and American) political, economic, and administrative goals for the colonies and their manifestations in the cultural practice of writing travel books. Travel writers’ descriptions of foreign places have been analyzed as attempts to validate, promote, or challenge the ideologies and practices of colonial or imperial domination and expansion. Mary Louise Pratt’s study of the genres and conventions of 18th- and 19th-century exploration narratives about South America and Africa (e.g., the “monarch of all I survey” trope) offered ways of thinking about travel writing as embedded within relations of power between metropole and periphery, as did Edward Said’s theories of representation and cultural imperialism. Particularly Said’s book, Orientalism, helped scholars understand ways in which representations of people in travel texts were intimately bound up with notions of self, in this case, that the Occident defined itself through essentialist, ethnocentric, and racist representations of the Orient. Said’s work became a model for demonstrating cultural forms of imperialism in travel texts, showing how the political, economic, or administrative fact of dominance relies on legitimating discourses such as those articulated through travel writing. . . .
Feminist geographers’ studies of travel writing challenge the masculinist history of geography by questioning who and what are relevant subjects of geographic study and, indeed, what counts as geographic knowledge itself. Such questions are worked through ideological constructs that posit men as explorers and women as travelers—or, conversely, men as travelers and women as tied to the home. Studies of Victorian women who were professional travel writers, tourists, wives of colonial administrators, and other (mostly) elite women who wrote narratives about their experiences abroad during the 19th century have been particularly revealing. From a “liberal” feminist perspective, travel presented one means toward female liberation for middle- and upper-class Victorian women. Many studies from the 1970s onward demonstrated the ways in which women’s gendered identities were negotiated differently “at home” than they were “away,” thereby showing women’s self-development through travel. The more recent post structural turn in studies of Victorian travel writing has focused attention on women’s diverse and fragmented identities as they narrated their travel experiences, emphasizing women’s sense of themselves as women in new locations, but only as they worked through their ties to nation, class, whiteness, and colonial and imperial power structures

\vspace{1em}
\textbf{Question 16}

The statement “The richer you are, the more you spend to be off-screen” is supported by which other line from the passage?

\begin{enumerate}[label=\Alph*.]
    \item “Gmail is the same Gmail. And it’s all free.”
    \item “How comfortable someone is with human engagement could become a new class marker.”
    \item “. . . screens are filling a crucial void.”
    \item “. . . studies show that time on these advertisement-support platforms is unhealthy .
[](https://cracku.in/16-the-statement-the-richer-you-are-the-more-you-spen-x-cat-2020-slot-3)
\end{enumerate}
\vspace{1em}
\textbf{Solution:}

At the very beginning, the author highlights the disparity in the screen time with regard to the wealthy and the common masses: {.._.As more screens appear in the lives of the poor, screens are disappearing from the lives of the rich. The richer you are, the more you spend to be off-screen_...}. The statement in Option B supplements this assertion by further highlighting this observed difference in activity between the rich and the common: {..._As more screens appear in the lives of the poor, screens are disappearing from the lives of the rich. The richer you are, the more you spend to be off-screen._..}. None of the other options can be attached to the given statements. Hence, Option B is the correct answer.

\end{problem}

\newpage
\begin{problem}{}{}
% Subject: varc
\textbf{Instructions}

Instructions   

The passage below is accompanied by a set of questions. Choose the best answer to each question.
Mode of transportation affects the travel experience and thus can produce new types of travel writing and perhaps even new “identities.” Modes of transportation determine the types and duration of social encounters; affect the organization and passage of space and time; . . . and also affect perception and knowledge—how and what the traveler comes to know and write about. The completion of the first U.S. transcontinental highway during the 1920s . . . for example, inaugurated a new genre of travel literature about the United States—the automotive or road narrative. Such narratives highlight the experiences of mostly male protagonists “discovering themselves” on their journeys, emphasizing the independence of road travel and the value of rural folk traditions.
Travel writing’s relationship to empire building— as a type of “colonialist discourse”—has drawn the most attention from academicians. Close connections have been observed between European (and American) political, economic, and administrative goals for the colonies and their manifestations in the cultural practice of writing travel books. Travel writers’ descriptions of foreign places have been analyzed as attempts to validate, promote, or challenge the ideologies and practices of colonial or imperial domination and expansion. Mary Louise Pratt’s study of the genres and conventions of 18th- and 19th-century exploration narratives about South America and Africa (e.g., the “monarch of all I survey” trope) offered ways of thinking about travel writing as embedded within relations of power between metropole and periphery, as did Edward Said’s theories of representation and cultural imperialism. Particularly Said’s book, Orientalism, helped scholars understand ways in which representations of people in travel texts were intimately bound up with notions of self, in this case, that the Occident defined itself through essentialist, ethnocentric, and racist representations of the Orient. Said’s work became a model for demonstrating cultural forms of imperialism in travel texts, showing how the political, economic, or administrative fact of dominance relies on legitimating discourses such as those articulated through travel writing. . . .
Feminist geographers’ studies of travel writing challenge the masculinist history of geography by questioning who and what are relevant subjects of geographic study and, indeed, what counts as geographic knowledge itself. Such questions are worked through ideological constructs that posit men as explorers and women as travelers—or, conversely, men as travelers and women as tied to the home. Studies of Victorian women who were professional travel writers, tourists, wives of colonial administrators, and other (mostly) elite women who wrote narratives about their experiences abroad during the 19th century have been particularly revealing. From a “liberal” feminist perspective, travel presented one means toward female liberation for middle- and upper-class Victorian women. Many studies from the 1970s onward demonstrated the ways in which women’s gendered identities were negotiated differently “at home” than they were “away,” thereby showing women’s self-development through travel. The more recent post structural turn in studies of Victorian travel writing has focused attention on women’s diverse and fragmented identities as they narrated their travel experiences, emphasizing women’s sense of themselves as women in new locations, but only as they worked through their ties to nation, class, whiteness, and colonial and imperial power structures

\vspace{1em}
\textbf{Question 17}

The author claims that Silicon Valley tech companies have tried to “confuse the public” by:

\begin{enumerate}[label=\Alph*.]
    \item pushing for greater privacy while working with advertisement-support platforms to mine data.
    \item concealing the findings of psychologists and neuroscientists on screen-time use from the public.
    \item developing new work-efficiency programmes while lobbying for the “right to disconnect” bill.
    \item promoting screen time in public schools while opting for a screen-free education for their own children.
[](https://cracku.in/17-the-author-claims-that-silicon-valley-tech-compani-x-cat-2020-slot-3)
\end{enumerate}
\vspace{1em}
\textbf{Solution:}

The following excerpt from the fourth paragraph throws light into this matter: {.._.Tech companies worked hard to get public schools to buy into programs that required schools to have one laptop per student, arguing that it would better prepare children for their screen-based future. But this idea isn’t how the people who actually build the screen-based future raise their own children..._}. It is understood that though tech companies manipulate public schools into engaging in a process involving more screen time, they avoid a similar course of activity when it comes to their own children {whom they subject to a screen-free education and upbringing}. Option D aptly captures this two-facedness. Option B is a distorted interpretation, while Options A and C cannot be inferred from the passage. 
Hence, Option D is the correct answer.

\end{problem}

\newpage
\begin{problem}{}{}
% Subject: varc
\textbf{Instructions}

Instructions   

The passage below is accompanied by a set of questions. Choose the best answer to each question.
Mode of transportation affects the travel experience and thus can produce new types of travel writing and perhaps even new “identities.” Modes of transportation determine the types and duration of social encounters; affect the organization and passage of space and time; . . . and also affect perception and knowledge—how and what the traveler comes to know and write about. The completion of the first U.S. transcontinental highway during the 1920s . . . for example, inaugurated a new genre of travel literature about the United States—the automotive or road narrative. Such narratives highlight the experiences of mostly male protagonists “discovering themselves” on their journeys, emphasizing the independence of road travel and the value of rural folk traditions.
Travel writing’s relationship to empire building— as a type of “colonialist discourse”—has drawn the most attention from academicians. Close connections have been observed between European (and American) political, economic, and administrative goals for the colonies and their manifestations in the cultural practice of writing travel books. Travel writers’ descriptions of foreign places have been analyzed as attempts to validate, promote, or challenge the ideologies and practices of colonial or imperial domination and expansion. Mary Louise Pratt’s study of the genres and conventions of 18th- and 19th-century exploration narratives about South America and Africa (e.g., the “monarch of all I survey” trope) offered ways of thinking about travel writing as embedded within relations of power between metropole and periphery, as did Edward Said’s theories of representation and cultural imperialism. Particularly Said’s book, Orientalism, helped scholars understand ways in which representations of people in travel texts were intimately bound up with notions of self, in this case, that the Occident defined itself through essentialist, ethnocentric, and racist representations of the Orient. Said’s work became a model for demonstrating cultural forms of imperialism in travel texts, showing how the political, economic, or administrative fact of dominance relies on legitimating discourses such as those articulated through travel writing. . . .
Feminist geographers’ studies of travel writing challenge the masculinist history of geography by questioning who and what are relevant subjects of geographic study and, indeed, what counts as geographic knowledge itself. Such questions are worked through ideological constructs that posit men as explorers and women as travelers—or, conversely, men as travelers and women as tied to the home. Studies of Victorian women who were professional travel writers, tourists, wives of colonial administrators, and other (mostly) elite women who wrote narratives about their experiences abroad during the 19th century have been particularly revealing. From a “liberal” feminist perspective, travel presented one means toward female liberation for middle- and upper-class Victorian women. Many studies from the 1970s onward demonstrated the ways in which women’s gendered identities were negotiated differently “at home” than they were “away,” thereby showing women’s self-development through travel. The more recent post structural turn in studies of Victorian travel writing has focused attention on women’s diverse and fragmented identities as they narrated their travel experiences, emphasizing women’s sense of themselves as women in new locations, but only as they worked through their ties to nation, class, whiteness, and colonial and imperial power structures

\vspace{1em}
\textbf{Question 18}

The author is least likely to agree with the view that the increase in screen-time is fuelled by the fact that:

\begin{enumerate}[label=\Alph*.]
    \item screens provide social contact in an increasingly isolating world
    \item there is a growth in computer-based teaching in public schools
    \item some workers face punitive action if they are not online
    \item with falling costs, people are streaming more content on their devices
[](https://cracku.in/18-the-author-is-least-likely-to-agree-with-the-view--x-cat-2020-slot-3)
\end{enumerate}
\vspace{1em}
\textbf{Solution:}

_Option A_ : The author has already discussed a point along similar lines towards the end of the passage: {._..There is also the reality that in our culture of increasing isolation, in which so many of the traditional gathering places and social structures have disappeared, screens are filling a crucial void..._}. Hence, the author definitely considers this as one of the causes behind the increase in screen time.
_Option B_ : The discussion in passage 4 and 5 highlights how the tech companies have {perhaps successfully} convinced schools to integrate a screen-based educational culture to prepare students for a "screen-based future". Thus, this is another factor that the author attributes to the rise in screen time.
_Option C_ : This aspect has been discussed towards the end of the passage: {..._There is a small movement to pass a “right to disconnect” bill, which would allow workers to turn their phones off, but for now, a worker can be punished for going offline and not being available..._}. Hence, the author considers this as another factor contributing to the increase in screen time. 
_Option D_ : There is no discussion mentioning or detailing the elements presented in D. Thus, we cannot conclusively comment on whether the author is likely to agree with this point. 
Hence, Option D is the correct answer.

Instructions   

For the following questions answer them individually

\end{problem}

\newpage
\begin{problem}{}{}
% Subject: varc
\textbf{Instructions}

Instructions   

The passage below is accompanied by a set of questions. Choose the best answer to each question.
Mode of transportation affects the travel experience and thus can produce new types of travel writing and perhaps even new “identities.” Modes of transportation determine the types and duration of social encounters; affect the organization and passage of space and time; . . . and also affect perception and knowledge—how and what the traveler comes to know and write about. The completion of the first U.S. transcontinental highway during the 1920s . . . for example, inaugurated a new genre of travel literature about the United States—the automotive or road narrative. Such narratives highlight the experiences of mostly male protagonists “discovering themselves” on their journeys, emphasizing the independence of road travel and the value of rural folk traditions.
Travel writing’s relationship to empire building— as a type of “colonialist discourse”—has drawn the most attention from academicians. Close connections have been observed between European (and American) political, economic, and administrative goals for the colonies and their manifestations in the cultural practice of writing travel books. Travel writers’ descriptions of foreign places have been analyzed as attempts to validate, promote, or challenge the ideologies and practices of colonial or imperial domination and expansion. Mary Louise Pratt’s study of the genres and conventions of 18th- and 19th-century exploration narratives about South America and Africa (e.g., the “monarch of all I survey” trope) offered ways of thinking about travel writing as embedded within relations of power between metropole and periphery, as did Edward Said’s theories of representation and cultural imperialism. Particularly Said’s book, Orientalism, helped scholars understand ways in which representations of people in travel texts were intimately bound up with notions of self, in this case, that the Occident defined itself through essentialist, ethnocentric, and racist representations of the Orient. Said’s work became a model for demonstrating cultural forms of imperialism in travel texts, showing how the political, economic, or administrative fact of dominance relies on legitimating discourses such as those articulated through travel writing. . . .
Feminist geographers’ studies of travel writing challenge the masculinist history of geography by questioning who and what are relevant subjects of geographic study and, indeed, what counts as geographic knowledge itself. Such questions are worked through ideological constructs that posit men as explorers and women as travelers—or, conversely, men as travelers and women as tied to the home. Studies of Victorian women who were professional travel writers, tourists, wives of colonial administrators, and other (mostly) elite women who wrote narratives about their experiences abroad during the 19th century have been particularly revealing. From a “liberal” feminist perspective, travel presented one means toward female liberation for middle- and upper-class Victorian women. Many studies from the 1970s onward demonstrated the ways in which women’s gendered identities were negotiated differently “at home” than they were “away,” thereby showing women’s self-development through travel. The more recent post structural turn in studies of Victorian travel writing has focused attention on women’s diverse and fragmented identities as they narrated their travel experiences, emphasizing women’s sense of themselves as women in new locations, but only as they worked through their ties to nation, class, whiteness, and colonial and imperial power structures

\vspace{1em}
\textbf{Question 19}

**The four sentences (labelled 1, 2, 3, 4) below, when properly sequenced would yield a coherent paragraph. Decide on the proper sequencing of the order of the sentences and key in the sequence of the four numbers as your answer:**  
1. Complex computational elements of the CNS are organized according to a “nested” hierarchic criterion; the organization is not permanent and can change dynamically from moment to moment as they carry out a computational task.  
2. Echolocation in bats exemplifies adaptation produced by natural selection; a function not produced by natural selection for its current use is exaptation -- feathers might have originally arisen in the context of selection for insulation.  
3. From a structural standpoint, consistent with exaptation, the living organism is organized as a complex of “Russian Matryoshka Dolls” -- smaller structures are contained within larger ones in multiple layers.  
4. The exaptation concept, and the Russian-doll organization concept of living beings deduced from studies on evolution of the various apparatuses in mammals, can be applied for the most complex human organ: the central nervous system (CNS).
Backspace  
789  
456  
123  
0.-  
Clear All  

Submit
[](https://cracku.in/19-the-four-sentences-labelled-1-2-3-4-below-when-pro-x-cat-2020-slot-3)

\begin{enumerate}[label=\Alph*.]
\end{enumerate}
\vspace{1em}
\textbf{Solution:}

Statement (2) introduces us to certain evolutionary influences in the bodily mechanisms of animals, especially mammals (bats). Statement (4) continues on the observed influence on the various apparatuses in mammals, specifically in the case of the most complex human organ: the central nervous system (CNS). The author mentions that the example of exaptation and Russian dolls serve to assist in our understanding of this complex organ (CNS). The significance of the Russian dolls is elaborated in Statement (3), specifically with regard to its structural implications: "  _smaller structures are contained within larger ones in multiple layers_ "—this portryas the presence of some structural hierarchy that can be observed in such complex apparatuses. Statement (1) highlights how the complex elements in the CNS are organised in a “nested hierarchic criterion", thereby, serving as a continuation to (3). Therefore, (2)-(4)-(3)-(1) forms a coherent paragraph.

\end{problem}

\newpage
\begin{problem}{}{}
% Subject: varc
\textbf{Instructions}

Instructions   

The passage below is accompanied by a set of questions. Choose the best answer to each question.
Mode of transportation affects the travel experience and thus can produce new types of travel writing and perhaps even new “identities.” Modes of transportation determine the types and duration of social encounters; affect the organization and passage of space and time; . . . and also affect perception and knowledge—how and what the traveler comes to know and write about. The completion of the first U.S. transcontinental highway during the 1920s . . . for example, inaugurated a new genre of travel literature about the United States—the automotive or road narrative. Such narratives highlight the experiences of mostly male protagonists “discovering themselves” on their journeys, emphasizing the independence of road travel and the value of rural folk traditions.
Travel writing’s relationship to empire building— as a type of “colonialist discourse”—has drawn the most attention from academicians. Close connections have been observed between European (and American) political, economic, and administrative goals for the colonies and their manifestations in the cultural practice of writing travel books. Travel writers’ descriptions of foreign places have been analyzed as attempts to validate, promote, or challenge the ideologies and practices of colonial or imperial domination and expansion. Mary Louise Pratt’s study of the genres and conventions of 18th- and 19th-century exploration narratives about South America and Africa (e.g., the “monarch of all I survey” trope) offered ways of thinking about travel writing as embedded within relations of power between metropole and periphery, as did Edward Said’s theories of representation and cultural imperialism. Particularly Said’s book, Orientalism, helped scholars understand ways in which representations of people in travel texts were intimately bound up with notions of self, in this case, that the Occident defined itself through essentialist, ethnocentric, and racist representations of the Orient. Said’s work became a model for demonstrating cultural forms of imperialism in travel texts, showing how the political, economic, or administrative fact of dominance relies on legitimating discourses such as those articulated through travel writing. . . .
Feminist geographers’ studies of travel writing challenge the masculinist history of geography by questioning who and what are relevant subjects of geographic study and, indeed, what counts as geographic knowledge itself. Such questions are worked through ideological constructs that posit men as explorers and women as travelers—or, conversely, men as travelers and women as tied to the home. Studies of Victorian women who were professional travel writers, tourists, wives of colonial administrators, and other (mostly) elite women who wrote narratives about their experiences abroad during the 19th century have been particularly revealing. From a “liberal” feminist perspective, travel presented one means toward female liberation for middle- and upper-class Victorian women. Many studies from the 1970s onward demonstrated the ways in which women’s gendered identities were negotiated differently “at home” than they were “away,” thereby showing women’s self-development through travel. The more recent post structural turn in studies of Victorian travel writing has focused attention on women’s diverse and fragmented identities as they narrated their travel experiences, emphasizing women’s sense of themselves as women in new locations, but only as they worked through their ties to nation, class, whiteness, and colonial and imperial power structures

\vspace{1em}
\textbf{Question 20}

**The four sentences (labelled 1, 2, 3, 4) below, when properly sequenced would yield a coherent paragraph. Decide on the proper sequencing of the order of the sentences and key in the sequence of the four numbers as your answer:**  
1. It advocated a conservative approach to antitrust enforcement that espouses faith in efficient markets and voiced suspicion regarding the merits of judicial intervention to correct anticompetitive practices.  
2. Many industries have consistently gained market share, the lion’s share - without any official concern; the most successful technology companies have grown into veritable titans, on the premise that they advance ‘public interest’.  
3. That the new anticompetitive risks posed by tech giants like Google, Facebook, and Amazon, necessitate new legal solutions could be attributed to the dearth of enforcement actions against monopolies and the few cases challenging mergers in the USA.  
4. The criterion of ‘consumer welfare standard’ and the principle that antitrust law should serve consumer interests and that it should protect competition rather than individual competitors was an antitrust law introduced by, and named after, the 'Chicago school'.
Backspace  
789  
456  
123  
0.-  
Clear All  

Submit
[](https://cracku.in/20-the-four-sentences-labelled-1-2-3-4-below-when-pro-x-cat-2020-slot-3)

\begin{enumerate}[label=\Alph*.]
\end{enumerate}
\vspace{1em}
\textbf{Solution:}

Statement (4) opens the discussion by mentioning an "antitrust law" and a few essential features attached to it. Statement (1) further describes this provision {"._..It advocated.._." }: the move to support an efficient market and curb anti-competitive practises. Statement (2) then elaborates on the latter point of anticompetitive practices and subsequently, highlights the relevance of this provision. Statement (3) further emphasises the necessity of "new legal solutions" to deal with the elements discussed earlier {in (2) and (3)}. Hence, (4)-(1)-(2)-(3) forms a logical arrangement.

\end{problem}

\newpage
\begin{problem}{}{}
% Subject: varc
\textbf{Instructions}

Instructions   

The passage below is accompanied by a set of questions. Choose the best answer to each question.
Mode of transportation affects the travel experience and thus can produce new types of travel writing and perhaps even new “identities.” Modes of transportation determine the types and duration of social encounters; affect the organization and passage of space and time; . . . and also affect perception and knowledge—how and what the traveler comes to know and write about. The completion of the first U.S. transcontinental highway during the 1920s . . . for example, inaugurated a new genre of travel literature about the United States—the automotive or road narrative. Such narratives highlight the experiences of mostly male protagonists “discovering themselves” on their journeys, emphasizing the independence of road travel and the value of rural folk traditions.
Travel writing’s relationship to empire building— as a type of “colonialist discourse”—has drawn the most attention from academicians. Close connections have been observed between European (and American) political, economic, and administrative goals for the colonies and their manifestations in the cultural practice of writing travel books. Travel writers’ descriptions of foreign places have been analyzed as attempts to validate, promote, or challenge the ideologies and practices of colonial or imperial domination and expansion. Mary Louise Pratt’s study of the genres and conventions of 18th- and 19th-century exploration narratives about South America and Africa (e.g., the “monarch of all I survey” trope) offered ways of thinking about travel writing as embedded within relations of power between metropole and periphery, as did Edward Said’s theories of representation and cultural imperialism. Particularly Said’s book, Orientalism, helped scholars understand ways in which representations of people in travel texts were intimately bound up with notions of self, in this case, that the Occident defined itself through essentialist, ethnocentric, and racist representations of the Orient. Said’s work became a model for demonstrating cultural forms of imperialism in travel texts, showing how the political, economic, or administrative fact of dominance relies on legitimating discourses such as those articulated through travel writing. . . .
Feminist geographers’ studies of travel writing challenge the masculinist history of geography by questioning who and what are relevant subjects of geographic study and, indeed, what counts as geographic knowledge itself. Such questions are worked through ideological constructs that posit men as explorers and women as travelers—or, conversely, men as travelers and women as tied to the home. Studies of Victorian women who were professional travel writers, tourists, wives of colonial administrators, and other (mostly) elite women who wrote narratives about their experiences abroad during the 19th century have been particularly revealing. From a “liberal” feminist perspective, travel presented one means toward female liberation for middle- and upper-class Victorian women. Many studies from the 1970s onward demonstrated the ways in which women’s gendered identities were negotiated differently “at home” than they were “away,” thereby showing women’s self-development through travel. The more recent post structural turn in studies of Victorian travel writing has focused attention on women’s diverse and fragmented identities as they narrated their travel experiences, emphasizing women’s sense of themselves as women in new locations, but only as they worked through their ties to nation, class, whiteness, and colonial and imperial power structures

\vspace{1em}
\textbf{Question 21}

**Five jumbled up sentences, related to a topic, are given below. Four of them can be put together to form a coherent paragraph. Identify the odd one out and key in the number of the sentence as your answer:**  
1. Machine learning models are prone to learning human-like biases from the training data that feeds these algorithms.  
2. Hate speech detection is part of the on-going effort against oppressive and abusive language on social media.  
3. The current automatic detection models miss out on something vital: context.  
4. It uses complex algorithms to flag racist or violent speech faster and better than human beings alone.  
5. For instance, algorithms struggle to determine if group identifiers like "gay" or "black" are used in offensive or prejudiced ways because they're trained on imbalanced datasets with unusually high rates of hate speech.
Backspace  
789  
456  
123  
0.-  
Clear All  

Submit
[](https://cracku.in/21-five-jumbled-up-sentences-related-to-a-topic-are-g-x-cat-2020-slot-3)

\begin{enumerate}[label=\Alph*.]
\end{enumerate}
\vspace{1em}
\textbf{Solution:}

On reading the statements, the arrangement (2)-(4)-(1)-(5) can be linked to form a paragraph, while Statement (3) stands out. Statements (2) and (4) talk about hate speech detection and the algorithms involved, while Statements (1) and (5) indicate the issue associated with the aforementioned algorithms. Hence, (3) is the odd one out.

\end{problem}

\newpage
\begin{problem}{}{}
% Subject: varc
\textbf{Instructions}

Instructions   

The passage below is accompanied by a set of questions. Choose the best answer to each question.
Mode of transportation affects the travel experience and thus can produce new types of travel writing and perhaps even new “identities.” Modes of transportation determine the types and duration of social encounters; affect the organization and passage of space and time; . . . and also affect perception and knowledge—how and what the traveler comes to know and write about. The completion of the first U.S. transcontinental highway during the 1920s . . . for example, inaugurated a new genre of travel literature about the United States—the automotive or road narrative. Such narratives highlight the experiences of mostly male protagonists “discovering themselves” on their journeys, emphasizing the independence of road travel and the value of rural folk traditions.
Travel writing’s relationship to empire building— as a type of “colonialist discourse”—has drawn the most attention from academicians. Close connections have been observed between European (and American) political, economic, and administrative goals for the colonies and their manifestations in the cultural practice of writing travel books. Travel writers’ descriptions of foreign places have been analyzed as attempts to validate, promote, or challenge the ideologies and practices of colonial or imperial domination and expansion. Mary Louise Pratt’s study of the genres and conventions of 18th- and 19th-century exploration narratives about South America and Africa (e.g., the “monarch of all I survey” trope) offered ways of thinking about travel writing as embedded within relations of power between metropole and periphery, as did Edward Said’s theories of representation and cultural imperialism. Particularly Said’s book, Orientalism, helped scholars understand ways in which representations of people in travel texts were intimately bound up with notions of self, in this case, that the Occident defined itself through essentialist, ethnocentric, and racist representations of the Orient. Said’s work became a model for demonstrating cultural forms of imperialism in travel texts, showing how the political, economic, or administrative fact of dominance relies on legitimating discourses such as those articulated through travel writing. . . .
Feminist geographers’ studies of travel writing challenge the masculinist history of geography by questioning who and what are relevant subjects of geographic study and, indeed, what counts as geographic knowledge itself. Such questions are worked through ideological constructs that posit men as explorers and women as travelers—or, conversely, men as travelers and women as tied to the home. Studies of Victorian women who were professional travel writers, tourists, wives of colonial administrators, and other (mostly) elite women who wrote narratives about their experiences abroad during the 19th century have been particularly revealing. From a “liberal” feminist perspective, travel presented one means toward female liberation for middle- and upper-class Victorian women. Many studies from the 1970s onward demonstrated the ways in which women’s gendered identities were negotiated differently “at home” than they were “away,” thereby showing women’s self-development through travel. The more recent post structural turn in studies of Victorian travel writing has focused attention on women’s diverse and fragmented identities as they narrated their travel experiences, emphasizing women’s sense of themselves as women in new locations, but only as they worked through their ties to nation, class, whiteness, and colonial and imperial power structures

\vspace{1em}
\textbf{Question 22}

**Five jumbled up sentences, related to a topic, are given below. Four of them can be put together to form a coherent paragraph. Identify the odd one out and key in the number of the sentence as your answer:**  
1. The logic of displaying one’s inner qualities through outward appearance was based on a distinction between being a woman and being feminine.  
2. 'Appearance' became a signifier of conduct - to look was to be and conformity to the feminine ideal was measured by how well women could use the tools of the fashion and beauty industries.  
3. The makeover-centric media sets out subtly and not-so-subtly, ‘good’ and ‘bad’ ways to be a woman, layering these over inequalities of race and class.  
4. The denigration of working-class women and women of colour often centres on their perceived failure to embody feminine beauty.  
5. ‘Woman’ was considered a biological category, but femininity was a ‘process’ by which women became specific kinds of women.
Backspace  
789  
456  
123  
0.-  
Clear All  

Submit
[](https://cracku.in/22-five-jumbled-up-sentences-related-to-a-topic-are-g-x-cat-2020-slot-3)

\begin{enumerate}[label=\Alph*.]
\end{enumerate}
\vspace{1em}
\textbf{Solution:}

Statement (1) talks about how the "logic" of determining a woman's inner quality boiled down to the distinction between the perception of "being a woman and being feminine". Statement (5) highlights the difference in this understanding: the former being a 'biological category' and latter being a 'process'. Statement (2) continues on the manner in which the measure of "feminine ideal" was dependent on a woman's appearance. Statement (4) continues on this line by presenting how the incapacity to meet up to this ideal led to the denigration of working-class women and women of colour. We notice that Statement (3) is the odd one out here.

\end{problem}

\newpage
\begin{problem}{}{}
% Subject: varc
\textbf{Instructions}

Instructions   

The passage below is accompanied by a set of questions. Choose the best answer to each question.
Mode of transportation affects the travel experience and thus can produce new types of travel writing and perhaps even new “identities.” Modes of transportation determine the types and duration of social encounters; affect the organization and passage of space and time; . . . and also affect perception and knowledge—how and what the traveler comes to know and write about. The completion of the first U.S. transcontinental highway during the 1920s . . . for example, inaugurated a new genre of travel literature about the United States—the automotive or road narrative. Such narratives highlight the experiences of mostly male protagonists “discovering themselves” on their journeys, emphasizing the independence of road travel and the value of rural folk traditions.
Travel writing’s relationship to empire building— as a type of “colonialist discourse”—has drawn the most attention from academicians. Close connections have been observed between European (and American) political, economic, and administrative goals for the colonies and their manifestations in the cultural practice of writing travel books. Travel writers’ descriptions of foreign places have been analyzed as attempts to validate, promote, or challenge the ideologies and practices of colonial or imperial domination and expansion. Mary Louise Pratt’s study of the genres and conventions of 18th- and 19th-century exploration narratives about South America and Africa (e.g., the “monarch of all I survey” trope) offered ways of thinking about travel writing as embedded within relations of power between metropole and periphery, as did Edward Said’s theories of representation and cultural imperialism. Particularly Said’s book, Orientalism, helped scholars understand ways in which representations of people in travel texts were intimately bound up with notions of self, in this case, that the Occident defined itself through essentialist, ethnocentric, and racist representations of the Orient. Said’s work became a model for demonstrating cultural forms of imperialism in travel texts, showing how the political, economic, or administrative fact of dominance relies on legitimating discourses such as those articulated through travel writing. . . .
Feminist geographers’ studies of travel writing challenge the masculinist history of geography by questioning who and what are relevant subjects of geographic study and, indeed, what counts as geographic knowledge itself. Such questions are worked through ideological constructs that posit men as explorers and women as travelers—or, conversely, men as travelers and women as tied to the home. Studies of Victorian women who were professional travel writers, tourists, wives of colonial administrators, and other (mostly) elite women who wrote narratives about their experiences abroad during the 19th century have been particularly revealing. From a “liberal” feminist perspective, travel presented one means toward female liberation for middle- and upper-class Victorian women. Many studies from the 1970s onward demonstrated the ways in which women’s gendered identities were negotiated differently “at home” than they were “away,” thereby showing women’s self-development through travel. The more recent post structural turn in studies of Victorian travel writing has focused attention on women’s diverse and fragmented identities as they narrated their travel experiences, emphasizing women’s sense of themselves as women in new locations, but only as they worked through their ties to nation, class, whiteness, and colonial and imperial power structures

\vspace{1em}
\textbf{Question 23}

**The passage given below is followed by four alternate summaries. Choose the option that best captures the essence of the passage.**
The dominant hypotheses in modern science believe that language evolved to allow humans to exchange factual information about the physical world. But an alternative view is that language evolved, in modern humans at least, to facilitate social bonding. It increased our ancestors’ chances of survival by enabling them to hunt more successfully or to cooperate more extensively. Language meant that things could be explained and that plans and past experiences could be shared efficiently.

\begin{enumerate}[label=\Alph*.]
    \item Since its origin, language has been continuously evolving to higher forms, from being used to identify objects to ensuring human survival by enabling our ancestors to bond and cooperate.
    \item From the belief that humans invented language to process factual information, scholars now think that language was the outcome of the need to ensure social cohesion and thus human survival.
    \item Most believe that language originated from a need to articulate facts, but others think it emerged from the need to promote social cohesion and cooperation, thus enabling human survival.
    \item Experts are challenging the narrow view of the origin of language, as being merely used to describe facts and label objects, to being necessary to promote more complex interactions among humans
[](https://cracku.in/23-the-passage-given-below-is-followed-by-four-altern-x-cat-2020-slot-3)
\end{enumerate}
\vspace{1em}
\textbf{Solution:}

One predominant viewpoint: language originated to exchange factual information
An alternative viewpoint: language originated to facilitate social bonding and consequently, to ensure human survival.
The summary needs to highlight these two core viewpoints. _Option C_ does this without deviating from the discussion.
_Option A_ : The evolution of language is not the focal point here; the views held in this regard are. {"language has been continuously evolving to higher forms"} Thus, we can eliminate this option since it comes across as a misrepresentation.
_Option B_ : This is a trap wherein the statement captures both the core viewpoints but there is a distortion involved: "...From the belief ..." to "...scholars now..." indicates a shift in the viewpoint. However, this is not the case - the author simply states two prevalent perspectives on the subject. 
_Option D_ : is again a distortion since experts are not "challenging any views; the author simply highlights the presence of two viewpoints {no conflict presented}
Hence, Option C is the correct answer.

\end{problem}

\newpage
\begin{problem}{}{}
% Subject: varc
\textbf{Instructions}

Instructions   

The passage below is accompanied by a set of questions. Choose the best answer to each question.
Mode of transportation affects the travel experience and thus can produce new types of travel writing and perhaps even new “identities.” Modes of transportation determine the types and duration of social encounters; affect the organization and passage of space and time; . . . and also affect perception and knowledge—how and what the traveler comes to know and write about. The completion of the first U.S. transcontinental highway during the 1920s . . . for example, inaugurated a new genre of travel literature about the United States—the automotive or road narrative. Such narratives highlight the experiences of mostly male protagonists “discovering themselves” on their journeys, emphasizing the independence of road travel and the value of rural folk traditions.
Travel writing’s relationship to empire building— as a type of “colonialist discourse”—has drawn the most attention from academicians. Close connections have been observed between European (and American) political, economic, and administrative goals for the colonies and their manifestations in the cultural practice of writing travel books. Travel writers’ descriptions of foreign places have been analyzed as attempts to validate, promote, or challenge the ideologies and practices of colonial or imperial domination and expansion. Mary Louise Pratt’s study of the genres and conventions of 18th- and 19th-century exploration narratives about South America and Africa (e.g., the “monarch of all I survey” trope) offered ways of thinking about travel writing as embedded within relations of power between metropole and periphery, as did Edward Said’s theories of representation and cultural imperialism. Particularly Said’s book, Orientalism, helped scholars understand ways in which representations of people in travel texts were intimately bound up with notions of self, in this case, that the Occident defined itself through essentialist, ethnocentric, and racist representations of the Orient. Said’s work became a model for demonstrating cultural forms of imperialism in travel texts, showing how the political, economic, or administrative fact of dominance relies on legitimating discourses such as those articulated through travel writing. . . .
Feminist geographers’ studies of travel writing challenge the masculinist history of geography by questioning who and what are relevant subjects of geographic study and, indeed, what counts as geographic knowledge itself. Such questions are worked through ideological constructs that posit men as explorers and women as travelers—or, conversely, men as travelers and women as tied to the home. Studies of Victorian women who were professional travel writers, tourists, wives of colonial administrators, and other (mostly) elite women who wrote narratives about their experiences abroad during the 19th century have been particularly revealing. From a “liberal” feminist perspective, travel presented one means toward female liberation for middle- and upper-class Victorian women. Many studies from the 1970s onward demonstrated the ways in which women’s gendered identities were negotiated differently “at home” than they were “away,” thereby showing women’s self-development through travel. The more recent post structural turn in studies of Victorian travel writing has focused attention on women’s diverse and fragmented identities as they narrated their travel experiences, emphasizing women’s sense of themselves as women in new locations, but only as they worked through their ties to nation, class, whiteness, and colonial and imperial power structures

\vspace{1em}
\textbf{Question 24}

**The passage given below is followed by four alternate summaries. Choose the option that best captures the essence of the passage.**
Aesthetic political representation urges us to realize that ‘the representative has autonomy with regard to the people represented’ but autonomy then is not an excuse to abandon one’s responsibility. Aesthetic autonomy requires cultivation of ‘disinterestedness’ on the part of actors which is not indifference. To have disinterestedness, that is, to have comportment towards the beautiful that is devoid of all ulterior references to use - requires a kind of aesthetic commitment; it is the liberation of ourselves for the release of what has proper worth only in itself.

\begin{enumerate}[label=\Alph*.]
    \item Aesthetic political representation advocates autonomy for the representatives manifested through disinterestedness which itself is different from indifference.
    \item Aesthetic political representation advocates autonomy for the representatives drawing from disinterestedness, which itself is different from indifference.
    \item Disinterestedness is different from indifference as the former means a non-subjective evaluation of things which is what constitutes aesthetic political representation
    \item Disinterestedness, as distinct from indifference, is the basis of political representation.
[](https://cracku.in/24-the-passage-given-below-is-followed-by-four-altern-x-cat-2020-slot-3)
\end{enumerate}
\vspace{1em}
\textbf{Solution:}

The paragraph discusses two essential elements: it begins by presenting the facet of autonomy enjoyed by the representative in Aesthetic political representation and then highlights the cultivation of "disinterestedness" in this regard. Additionally, the author distinctly identifies the aforementioned concept as being not the same as that of "indifference". Post this, towards the end. The author presents the reason behind this assertion. Option B correctly captures these two aspects without distorting the overall meaning.
Option A: The author does not claim that the autonomy "manifested" through disinterestedness. 
Option C: The statement here contains added elements which cannot be inferred from the passage. 
Option D: This alternative fails to capture the essence of the discussion and describes a single component. {'political representation' might again be incorrect}
Hence, of the given summaries, Option B aptly captures the substance of the passage.

\end{problem}

\newpage
\begin{problem}{}{}
% Subject: varc
\textbf{Instructions}

Instructions   

The passage below is accompanied by a set of questions. Choose the best answer to each question.
Mode of transportation affects the travel experience and thus can produce new types of travel writing and perhaps even new “identities.” Modes of transportation determine the types and duration of social encounters; affect the organization and passage of space and time; . . . and also affect perception and knowledge—how and what the traveler comes to know and write about. The completion of the first U.S. transcontinental highway during the 1920s . . . for example, inaugurated a new genre of travel literature about the United States—the automotive or road narrative. Such narratives highlight the experiences of mostly male protagonists “discovering themselves” on their journeys, emphasizing the independence of road travel and the value of rural folk traditions.
Travel writing’s relationship to empire building— as a type of “colonialist discourse”—has drawn the most attention from academicians. Close connections have been observed between European (and American) political, economic, and administrative goals for the colonies and their manifestations in the cultural practice of writing travel books. Travel writers’ descriptions of foreign places have been analyzed as attempts to validate, promote, or challenge the ideologies and practices of colonial or imperial domination and expansion. Mary Louise Pratt’s study of the genres and conventions of 18th- and 19th-century exploration narratives about South America and Africa (e.g., the “monarch of all I survey” trope) offered ways of thinking about travel writing as embedded within relations of power between metropole and periphery, as did Edward Said’s theories of representation and cultural imperialism. Particularly Said’s book, Orientalism, helped scholars understand ways in which representations of people in travel texts were intimately bound up with notions of self, in this case, that the Occident defined itself through essentialist, ethnocentric, and racist representations of the Orient. Said’s work became a model for demonstrating cultural forms of imperialism in travel texts, showing how the political, economic, or administrative fact of dominance relies on legitimating discourses such as those articulated through travel writing. . . .
Feminist geographers’ studies of travel writing challenge the masculinist history of geography by questioning who and what are relevant subjects of geographic study and, indeed, what counts as geographic knowledge itself. Such questions are worked through ideological constructs that posit men as explorers and women as travelers—or, conversely, men as travelers and women as tied to the home. Studies of Victorian women who were professional travel writers, tourists, wives of colonial administrators, and other (mostly) elite women who wrote narratives about their experiences abroad during the 19th century have been particularly revealing. From a “liberal” feminist perspective, travel presented one means toward female liberation for middle- and upper-class Victorian women. Many studies from the 1970s onward demonstrated the ways in which women’s gendered identities were negotiated differently “at home” than they were “away,” thereby showing women’s self-development through travel. The more recent post structural turn in studies of Victorian travel writing has focused attention on women’s diverse and fragmented identities as they narrated their travel experiences, emphasizing women’s sense of themselves as women in new locations, but only as they worked through their ties to nation, class, whiteness, and colonial and imperial power structures

\vspace{1em}
\textbf{Question 25}

**The passage given below is followed by four alternate summaries. Choose the option that best captures the essence of the passage.**
Brown et al. (2001) suggest that ‘metabolic theory may provide a conceptual foundation for much of ecology just as genetic theory provides a foundation for much of evolutionary biology’. One of the successes of genetic theory is the diversity of theoretical approaches and models that have been developed and applied. A Web of Science (v. 5.9. Thomson Reuters) search on genetic* + theor* + evol* identifies more than 12000 publications between 2005 and 2012. Considering only the 10 most-cited papers within this 12000 publication set, genetic theory can be seen to focus on genome dynamics, phylogenetic inference, game theory and the regulation of gene expression. There is no one fundamental genetic equation, but rather a wide array of genetic models, ranging from simple to complex, with differing inputs and outputs, and divergent areas of application, loosely connected to each other through the shared conceptual foundation of heritable variation.

\begin{enumerate}[label=\Alph*.]
    \item Genetic theory has evolved to spawn a wide range of theoretical models and applications but Metabolic theory need not evolve in a similar manner in the field of ecology
    \item Genetic theory has a wide range of theoretical approaches and application and is foundational to evolutionary biology and Metabolic theory has the potential to do the same for ecology
    \item Genetic theory has a wide range of theoretical approaches and applications and Metabolic theory must have the same in the field of ecology
    \item Genetic theory provides an example of how a range of theoretical approaches and applications can make a theory successful.
[](https://cracku.in/25-the-passage-given-below-is-followed-by-four-altern-x-cat-2020-slot-3)
\end{enumerate}
\vspace{1em}
\textbf{Solution:}

There are two key points discussed in the passage:
1. The prospect of "metabolic theory" being foundational to the field of ecology; the same as is the case in (2)
2. Genetic theory being the conceptual basis of evolutionary biology {given the diverse and extensive theoretical approaches and models available}. 
Thus, the summary needs to capture both these points. _Option B_ fulfils this requirement.
_Option A_ : is a distorted claim since it is not implied in the passage; the author does not assert that "metabolic theory need not evolve in a similar manner". 
_Option C_ : is again a misinterpretation because the author does not claim that metabolic theory "must" contribute in a similar fashion. Instead, the focus is on the "potential" of this theory.
_Option D_ : is divergent since the author does not discuss the "success" of a theory.
Hence, Option B is the correct answer.

\end{problem}

\newpage
\begin{problem}{}{}
% Subject: varc
\textbf{Instructions}

Instructions   

The passage below is accompanied by a set of questions. Choose the best answer to each question.
Mode of transportation affects the travel experience and thus can produce new types of travel writing and perhaps even new “identities.” Modes of transportation determine the types and duration of social encounters; affect the organization and passage of space and time; . . . and also affect perception and knowledge—how and what the traveler comes to know and write about. The completion of the first U.S. transcontinental highway during the 1920s . . . for example, inaugurated a new genre of travel literature about the United States—the automotive or road narrative. Such narratives highlight the experiences of mostly male protagonists “discovering themselves” on their journeys, emphasizing the independence of road travel and the value of rural folk traditions.
Travel writing’s relationship to empire building— as a type of “colonialist discourse”—has drawn the most attention from academicians. Close connections have been observed between European (and American) political, economic, and administrative goals for the colonies and their manifestations in the cultural practice of writing travel books. Travel writers’ descriptions of foreign places have been analyzed as attempts to validate, promote, or challenge the ideologies and practices of colonial or imperial domination and expansion. Mary Louise Pratt’s study of the genres and conventions of 18th- and 19th-century exploration narratives about South America and Africa (e.g., the “monarch of all I survey” trope) offered ways of thinking about travel writing as embedded within relations of power between metropole and periphery, as did Edward Said’s theories of representation and cultural imperialism. Particularly Said’s book, Orientalism, helped scholars understand ways in which representations of people in travel texts were intimately bound up with notions of self, in this case, that the Occident defined itself through essentialist, ethnocentric, and racist representations of the Orient. Said’s work became a model for demonstrating cultural forms of imperialism in travel texts, showing how the political, economic, or administrative fact of dominance relies on legitimating discourses such as those articulated through travel writing. . . .
Feminist geographers’ studies of travel writing challenge the masculinist history of geography by questioning who and what are relevant subjects of geographic study and, indeed, what counts as geographic knowledge itself. Such questions are worked through ideological constructs that posit men as explorers and women as travelers—or, conversely, men as travelers and women as tied to the home. Studies of Victorian women who were professional travel writers, tourists, wives of colonial administrators, and other (mostly) elite women who wrote narratives about their experiences abroad during the 19th century have been particularly revealing. From a “liberal” feminist perspective, travel presented one means toward female liberation for middle- and upper-class Victorian women. Many studies from the 1970s onward demonstrated the ways in which women’s gendered identities were negotiated differently “at home” than they were “away,” thereby showing women’s self-development through travel. The more recent post structural turn in studies of Victorian travel writing has focused attention on women’s diverse and fragmented identities as they narrated their travel experiences, emphasizing women’s sense of themselves as women in new locations, but only as they worked through their ties to nation, class, whiteness, and colonial and imperial power structures

\vspace{1em}
\textbf{Question 26}

**The four sentences (labelled 1, 2, 3, 4) below, when properly sequenced would yield a coherent paragraph. Decide on the proper sequencing of the order of the sentences and key in the sequence of the four numbers as your answer:**  
1. Each one personified a different aspect of good fortune.  
2. The others were versions of popular Buddhist gods, Hindu gods and Daoist gods.  
3. Seven popular Japanese deities, the Shichi Fukujin, were considered to bring good luck and happiness.  
4. Although they were included in the Shinto pantheon, only two of them, Daikoku and Ebisu, were indigenous Japanese gods.
Backspace  
789  
456  
123  
0.-  
Clear All  

Submit
[](https://cracku.in/26-the-four-sentences-labelled-1-2-3-4-below-when-pro-x-cat-2020-slot-3)

\begin{enumerate}[label=\Alph*.]
\end{enumerate}
\vspace{1em}
\textbf{Solution:}

Statement (3) opens the paragraph by introducing the subject: seven popular Japanese deities who bring good luck. Statement (1) then comments on the aspect of good fortune followed by statements (4) and (2). Statement (4) clarifies how only two of these seven entities qualify as indigenous Japanese gods while Statement (2) comments on the origin/background of the rest. Hence, (3)-(1)-(4)-(2) forms a coherent arrangement. 
[](https://cracku.in/downloads/11130)
  *  
  *

\end{problem}

\newpage
\begin{problem}{}{}
% Subject: varc
\textbf{Instructions}

Instructions   

**The passage below is accompanied by a set of questions. Choose the best answer to each question.**
We cannot travel outside our neighbourhood without passports. We must wear the same plain clothes. We must exchange our houses every ten years. We cannot avoid labour. We all go to bed at the same time . . . We have religious freedom, but we cannot deny that the soul dies with the body, since ‘but for the fear of punishment, they would have nothing but contempt for the laws and customs of society'. . . . In More’s time, for much of the population, given the plenty and security on offer, such restraints would not have seemed overly unreasonable. For modern readers, however, Utopia appears to rely upon relentless transparency, the repression of variety, and the curtailment of privacy. Utopia provides security: but at what price? In both its external and internal relations, indeed, it seems perilously dystopian.
Such a conclusion might be fortified by examining selectively the tradition which follows More on these points. This often portrays societies where . . . 'it would be almost impossible for man to be depraved, or wicked'. . . . This is achieved both through institutions and mores, which underpin the common life. . . . The passions are regulated and inequalities of wealth and distinction are minimized. Needs, vanity, and emulation are restrained, often by prizing equality and holding riches in contempt. The desire for public power is curbed. Marriage and sexual intercourse are often controlled: in Tommaso Campanella’s The City of the Sun (1623), the first great literary utopia after More’s, relations are forbidden to men before the age of twenty-one and women before nineteen. Communal child-rearing is normal; for Campanella, this commences at age two. Greater simplicity of life, ‘living according to nature’, is often a result: the desire for simplicity and purity are closely related. People become more alike in appearance, opinion, and outlook than they often have been. Unity, order, and homogeneity thus prevail at the cost of individuality and diversity. This model, as J. C. Davis demonstrates, dominated early modern utopianism. . . . And utopian homogeneity remains a familiar theme well into the twentieth century.
Given these considerations, it is not unreasonable to take as our starting point here the hypothesis that utopia and dystopia evidently share more in common than is often supposed. Indeed, they might be twins, the progeny of the same parents. Insofar as this proves to be the case, my linkage of both here will be uncomfortably close for some readers. Yet we should not mistake this argument for the assertion that all utopias are, or tend to produce, dystopias. Those who defend this proposition will find that their association here is not nearly close enough. For we have only to acknowledge the existence of thousands of successful intentional communities in which a cooperative ethos predominates and where harmony without coercion is the rule to set aside such an assertion. Here the individual’s submersion in the group is consensual (though this concept is not unproblematic). It results not in enslavement but voluntary submission to group norms. Harmony is achieved without . . . harming others.

\vspace{1em}
\textbf{Question 1}

Following from the passage, which one of the following may be seen as a characteristic of a utopian society?

\begin{enumerate}[label=\Alph*.]
    \item A society without any laws to restrain one’s individuality.
    \item A society where public power is earned through merit rather than through privilege.
    \item Institutional surveillance of every individual to ensure his/her security and welfare.
    \item The regulation of homogeneity through promoting competitive heterogeneity.
[](https://cracku.in/1-following-from-the-passage-which-one-of-the-follow-x-cat-2021-slot-1)
\end{enumerate}
\vspace{1em}
\textbf{Solution:}

_Option A_ : The author does not discuss a utopian narrative that involves a society without laws or social structure. Instead, he talks about regulations that curb individuality and promote homogeneity. Hence, we can eliminate Option A as a potential choice.
_Option B_ : The second paragraph relays the following idea: "_The passions are regulated, and inequalities of wealth and distinction are minimized. Needs, vanity, and emulation are restrained, often by prizing equality and holding riches in contempt. The desire for public power is curbed_ " Given that public power is not looked at favourably, the entire debate on the mechanism to attain this facet (public power) becomes irrelevant. Thus, we can eliminate Option B.
_Option C_ : "_Such a conclusion might be fortified by examining selectively the tradition which follows More on these points. This often portrays societies where . . . 'it would be almost impossible for man to be depraved, or wicked'. . . . This is achieved both through**institutions** and mores, which underpin the common life_" From the above excerpt, it is clarified that the utopian doctrines are enforced by institutions present in the society; additionally, the author mentions in the preceding paragraph that security in a utopian setting is attained through the curtailment of privacy. Hence, Option C is likely to be the correct choice. 
_Option D_ : There is no mention of regulating homogeneity through promoting competitive heterogeneity. We can, therefore, eliminate this choice.
Hence, Option C is the correct choice.

\end{problem}

\newpage
\begin{problem}{}{}
% Subject: varc
\textbf{Instructions}

Instructions   

**The passage below is accompanied by a set of questions. Choose the best answer to each question.**
We cannot travel outside our neighbourhood without passports. We must wear the same plain clothes. We must exchange our houses every ten years. We cannot avoid labour. We all go to bed at the same time . . . We have religious freedom, but we cannot deny that the soul dies with the body, since ‘but for the fear of punishment, they would have nothing but contempt for the laws and customs of society'. . . . In More’s time, for much of the population, given the plenty and security on offer, such restraints would not have seemed overly unreasonable. For modern readers, however, Utopia appears to rely upon relentless transparency, the repression of variety, and the curtailment of privacy. Utopia provides security: but at what price? In both its external and internal relations, indeed, it seems perilously dystopian.
Such a conclusion might be fortified by examining selectively the tradition which follows More on these points. This often portrays societies where . . . 'it would be almost impossible for man to be depraved, or wicked'. . . . This is achieved both through institutions and mores, which underpin the common life. . . . The passions are regulated and inequalities of wealth and distinction are minimized. Needs, vanity, and emulation are restrained, often by prizing equality and holding riches in contempt. The desire for public power is curbed. Marriage and sexual intercourse are often controlled: in Tommaso Campanella’s The City of the Sun (1623), the first great literary utopia after More’s, relations are forbidden to men before the age of twenty-one and women before nineteen. Communal child-rearing is normal; for Campanella, this commences at age two. Greater simplicity of life, ‘living according to nature’, is often a result: the desire for simplicity and purity are closely related. People become more alike in appearance, opinion, and outlook than they often have been. Unity, order, and homogeneity thus prevail at the cost of individuality and diversity. This model, as J. C. Davis demonstrates, dominated early modern utopianism. . . . And utopian homogeneity remains a familiar theme well into the twentieth century.
Given these considerations, it is not unreasonable to take as our starting point here the hypothesis that utopia and dystopia evidently share more in common than is often supposed. Indeed, they might be twins, the progeny of the same parents. Insofar as this proves to be the case, my linkage of both here will be uncomfortably close for some readers. Yet we should not mistake this argument for the assertion that all utopias are, or tend to produce, dystopias. Those who defend this proposition will find that their association here is not nearly close enough. For we have only to acknowledge the existence of thousands of successful intentional communities in which a cooperative ethos predominates and where harmony without coercion is the rule to set aside such an assertion. Here the individual’s submersion in the group is consensual (though this concept is not unproblematic). It results not in enslavement but voluntary submission to group norms. Harmony is achieved without . . . harming others.

\vspace{1em}
\textbf{Question 2}

All of the following arguments are made in the passage EXCEPT that:

\begin{enumerate}[label=\Alph*.]
    \item in More’s time, there was plenty and security, so people did not need restraints that could appear unreasonable.
    \item there have been thousands of communities where homogeneity and stability have been achieved through choice, rather than by force.
    \item the tradition of utopian literature has often shown societies in which it would be nearly impossible for anyone to be sinful or criminal.
    \item in early modern utopianism, the stability of utopian societies was seen to be achieved only with individuals surrendering their sense of self.
[](https://cracku.in/2-all-of-the-following-arguments-are-made-in-the-pas-x-cat-2021-slot-1)
\end{enumerate}
\vspace{1em}
\textbf{Solution:}

_Option A_ : The statement here appears to be a distortion. The author says: "_In More’s time, for much of the population, given the plenty and security on offer, such restraints would not have seemed overly unreasonable._ " It is being conveyed that the form of restrictions discussed at the beginning of the passage would not seem unreasonable to the citizens/members of More's utopia. However, the author feels that the opinions of modern readers would be drastically different. 
_Option B_ : "_For we have only to acknowledge the existence of thousands of successful intentional communities in which a cooperative ethos predominates and where harmony without coercion is the rule to set aside such an assertion. Here the individual’s submersion in the group is consensual (though this concept is not unproblematic). It results not in enslavement but voluntary submission to group norms. Harmony is achieved without . . . harming others._ " Towards the end of the discussion, the author indicates that homogeneity and stability (that often constitute a utopian universe) need not be achieved via coercion. Members of many communities voluntarily concede to the group's norms at the cost of their individuality.  _  
_
_Option C_ : The statement here correlates to the assertion made by the author in the second paragraph: "_Such a conclusion might be fortified by examining selectively the tradition which follows More on these points. This often portrays societies where . . . 'it would be almost impossible for man to be depraved, or wicked'. . . . This is achieved both through institutions and mores, which underpin the common life._ " 
_Option D_ : The introductory segment of the discussion highlights the restraints placed on individual freedom in a utopian society. Furthermore, the author mentions that many most utopian narratives in literary history are built on a premise devoid of individuality or diversity, as stated in the following excerpt: "_People become more alike in appearance, opinion, and outlook than they often have been. Unity, order, and homogeneity thus prevail at the cost of individuality and diversity. This model, as J. C. Davis demonstrates, dominated early modern utopianism. . . . And utopian homogeneity remains a familiar theme well into the twentieth century._ "
Hence, Option A is the correct choice.

\end{problem}

\newpage
\begin{problem}{}{}
% Subject: varc
\textbf{Instructions}

Instructions   

**The passage below is accompanied by a set of questions. Choose the best answer to each question.**
We cannot travel outside our neighbourhood without passports. We must wear the same plain clothes. We must exchange our houses every ten years. We cannot avoid labour. We all go to bed at the same time . . . We have religious freedom, but we cannot deny that the soul dies with the body, since ‘but for the fear of punishment, they would have nothing but contempt for the laws and customs of society'. . . . In More’s time, for much of the population, given the plenty and security on offer, such restraints would not have seemed overly unreasonable. For modern readers, however, Utopia appears to rely upon relentless transparency, the repression of variety, and the curtailment of privacy. Utopia provides security: but at what price? In both its external and internal relations, indeed, it seems perilously dystopian.
Such a conclusion might be fortified by examining selectively the tradition which follows More on these points. This often portrays societies where . . . 'it would be almost impossible for man to be depraved, or wicked'. . . . This is achieved both through institutions and mores, which underpin the common life. . . . The passions are regulated and inequalities of wealth and distinction are minimized. Needs, vanity, and emulation are restrained, often by prizing equality and holding riches in contempt. The desire for public power is curbed. Marriage and sexual intercourse are often controlled: in Tommaso Campanella’s The City of the Sun (1623), the first great literary utopia after More’s, relations are forbidden to men before the age of twenty-one and women before nineteen. Communal child-rearing is normal; for Campanella, this commences at age two. Greater simplicity of life, ‘living according to nature’, is often a result: the desire for simplicity and purity are closely related. People become more alike in appearance, opinion, and outlook than they often have been. Unity, order, and homogeneity thus prevail at the cost of individuality and diversity. This model, as J. C. Davis demonstrates, dominated early modern utopianism. . . . And utopian homogeneity remains a familiar theme well into the twentieth century.
Given these considerations, it is not unreasonable to take as our starting point here the hypothesis that utopia and dystopia evidently share more in common than is often supposed. Indeed, they might be twins, the progeny of the same parents. Insofar as this proves to be the case, my linkage of both here will be uncomfortably close for some readers. Yet we should not mistake this argument for the assertion that all utopias are, or tend to produce, dystopias. Those who defend this proposition will find that their association here is not nearly close enough. For we have only to acknowledge the existence of thousands of successful intentional communities in which a cooperative ethos predominates and where harmony without coercion is the rule to set aside such an assertion. Here the individual’s submersion in the group is consensual (though this concept is not unproblematic). It results not in enslavement but voluntary submission to group norms. Harmony is achieved without . . . harming others.

\vspace{1em}
\textbf{Question 3}

Which sequence of words below best captures the narrative of the passage?

\begin{enumerate}[label=\Alph*.]
    \item Utopia - Security - Homogeneity - Intentional community.
    \item Relentless transparency - Homogeneity - Utopia - Dystopia.
    \item Utopia - Security - Dystopia - Coercion.
    \item Curtailment of privacy - Dystopia - Utopia - Intentional community.
[](https://cracku.in/3-which-sequence-of-words-below-best-captures-the-na-x-cat-2021-slot-1)
\end{enumerate}
\vspace{1em}
\textbf{Solution:}

The passage begins by portraying a **utopian** society. The author then discusses the difference in perspective concerning the underlying elements of such a society. A disagreement originates about the perception of **security** - while the people part of the utopia might find the shackles on their freedom to be reasonable, modern readers perceive this as suppression of heterogeneity and violation of privacy. This is presented with the intention to direct attention towards the tradeoff that exists between security and certain other essential variables. It additionally puts the spotlight on the thin line that exists between a utopia and a dystopia. The author then cites other works in literary history that depict a utopian setting along with certain key attributes that one might stumble upon in such narratives. **Homogeneity** comes across as a prominent idea (a set of beliefs are considered acceptable, and the masses are expected to conform to the same). Towards the end of the discussion, the author reiterates the thin film that separates a utopia from a dystopia. He adds that while many individuals might be tempted to use these two ideas interchangeably, this shouldn't be the case. According to the author, the assertion that "all utopias are, or tend to produce, dystopias" is fallacious. He presents justification concerning the same: there are many utopian settings wherein conformity to doctrines or sacrifice of individuality is intentional - the person voluntarily submits to the group's norms for the greater good. This resonates with the term **intentional community** stated in the options. Option A aptly captures these principal themes.

\end{problem}

\newpage
\begin{problem}{}{}
% Subject: varc
\textbf{Instructions}

Instructions   

**The passage below is accompanied by a set of questions. Choose the best answer to each question.**
We cannot travel outside our neighbourhood without passports. We must wear the same plain clothes. We must exchange our houses every ten years. We cannot avoid labour. We all go to bed at the same time . . . We have religious freedom, but we cannot deny that the soul dies with the body, since ‘but for the fear of punishment, they would have nothing but contempt for the laws and customs of society'. . . . In More’s time, for much of the population, given the plenty and security on offer, such restraints would not have seemed overly unreasonable. For modern readers, however, Utopia appears to rely upon relentless transparency, the repression of variety, and the curtailment of privacy. Utopia provides security: but at what price? In both its external and internal relations, indeed, it seems perilously dystopian.
Such a conclusion might be fortified by examining selectively the tradition which follows More on these points. This often portrays societies where . . . 'it would be almost impossible for man to be depraved, or wicked'. . . . This is achieved both through institutions and mores, which underpin the common life. . . . The passions are regulated and inequalities of wealth and distinction are minimized. Needs, vanity, and emulation are restrained, often by prizing equality and holding riches in contempt. The desire for public power is curbed. Marriage and sexual intercourse are often controlled: in Tommaso Campanella’s The City of the Sun (1623), the first great literary utopia after More’s, relations are forbidden to men before the age of twenty-one and women before nineteen. Communal child-rearing is normal; for Campanella, this commences at age two. Greater simplicity of life, ‘living according to nature’, is often a result: the desire for simplicity and purity are closely related. People become more alike in appearance, opinion, and outlook than they often have been. Unity, order, and homogeneity thus prevail at the cost of individuality and diversity. This model, as J. C. Davis demonstrates, dominated early modern utopianism. . . . And utopian homogeneity remains a familiar theme well into the twentieth century.
Given these considerations, it is not unreasonable to take as our starting point here the hypothesis that utopia and dystopia evidently share more in common than is often supposed. Indeed, they might be twins, the progeny of the same parents. Insofar as this proves to be the case, my linkage of both here will be uncomfortably close for some readers. Yet we should not mistake this argument for the assertion that all utopias are, or tend to produce, dystopias. Those who defend this proposition will find that their association here is not nearly close enough. For we have only to acknowledge the existence of thousands of successful intentional communities in which a cooperative ethos predominates and where harmony without coercion is the rule to set aside such an assertion. Here the individual’s submersion in the group is consensual (though this concept is not unproblematic). It results not in enslavement but voluntary submission to group norms. Harmony is achieved without . . . harming others.

\vspace{1em}
\textbf{Question 4}

All of the following statements can be inferred from the passage EXCEPT that:

\begin{enumerate}[label=\Alph*.]
    \item utopian and dystopian societies are twins, the progeny of the same parents.
    \item it is possible to see utopias as dystopias, with a change in perspective, because one person’s utopia could be seen as another’s dystopia.
    \item many conceptions of utopian societies emphasise the importance of social uniformity and cultural homogeneity.
    \item utopian societies exist in a long tradition of literature dealing with imaginary people practicing imaginary customs, in imaginary worlds.
[](https://cracku.in/4-all-of-the-following-statements-can-be-inferred-fr-x-cat-2021-slot-1)
\end{enumerate}
\vspace{1em}
\textbf{Solution:}

Option A:
_Indeed, they might be twins, the progeny of the same parents. Insofar as this proves to be the case, my linkage of both here will be uncomfortably close for some readers._
**might** be twins. The level of certainty is not absolute. However, Option A goes one step further to assert that they **are** twins and the progeny of the same parents. Hence, A cannot be inferred and is the answer.  
'_Insofar as this proves to be the case_ ' can cause confusion while answering. But note that the case the author is talking about is the level of similarity between the two. Hence, what is being proven is that the two are quite similar to each other, and hence some would presume that they are twins. The excerpt, however, does not support that they actually are.
Option B: The whole passage supports the inference that utopias can be perceived as dystopias by different people. E.g. The author mentions that where some people push for relentless transparency so that they are secure, some people would perceive this as a breach of their privacy. Hence, a utopia for the former would be a dystopia for the latter.
Option C: 
_People become more alike in appearance, opinion, and outlook than they often have been. Unity, order, and homogeneity thus prevail at the cost of individuality and diversity. This model, as J. C. Davis demonstrates, dominated early modern utopianism. . . . And utopian homogeneity remains a familiar theme well into the twentieth century._
Option C is a direct inference from the above excerpt. It has been mentioned that this theme of homogeneity and uniformity dominated early modern utopianism.
Option D: Throughout the passage, the author deals with conceptions of utopian societies as dealt with in literary works. We can infer that utopian societies do exist in literature where the characters practice traditions that the author made up to portray a utopian society.

Instructions   

**The passage below is accompanied by a set of questions. Choose the best answer to each question.**
For the Maya of the Classic period, who lived in Southern Mexico and Central America between 250 and 900 CE, the category of ‘persons’ was not coincident with human beings, as it is for us. That is, human beings were persons - but other, nonhuman entities could be persons, too. . . . In order to explore the slippage of categories between ‘humans’ and ‘persons’, I examined a very specific category of ancient Maya images, found painted in scenes on ceramic vessels. I sought out instances in which faces (some combination of eyes, nose, and mouth) are shown on inanimate objects. . . . Consider my iPhone, which needs to be fed with electricity every night, swaddled in a protective bumper, and enjoys communicating with other fellow-phone-beings. Does it have personhood (if at all) because it is connected to me, drawing this resource from me as an owner or source? For the Maya (who did have plenty of other communicating objects, if not smartphones), the answer was no. Nonhuman persons were not tethered to specific humans, and they did not derive their personhood from a connection with a human. . . . It’s a profoundly democratising way of understanding the world. Humans are not more important persons - we are just one of many kinds of persons who inhabit this world. . . .
The Maya saw personhood as ‘activated’ by experiencing certain bodily needs and through participation in certain social activities. For example, among the faced objects that I examined, persons are marked by personal requirements (such as hunger, tiredness, physical closeness), and by community obligations (communication, interaction, ritual observance). In the images I examined, we see, for instance, faced objects being cradled in humans’ arms; we also see them speaking to humans. These core elements of personhood are both turned inward, what the body or self of a person requires, and outward, what a community expects of the persons who are a part of it, underlining the reciprocal nature of community membership. .
Personhood was a nonbinary proposition for the Maya. Entities were able to be persons while also being something else. The faced objects I looked at indicate that they continue to be functional, doing what objects do (a stone implement continues to chop, an incense burner continues to do its smoky work). Furthermore, the Maya visually depicted many objects in ways that indicated the material category to which they belonged - drawings of the stone implement show that a person-tool is still made of stone. One additional complexity: the incense burner (which would have been made of clay, and decorated with spiky appliques representing the sacred ceiba tree found in this region) is categorised as a person - but also as a tree. With these Maya examples, we are challenged to discard the person/nonperson binary that constitutes our basic ontological outlook. . . . The porousness of boundaries that we have seen in the Maya world points towards the possibility of living with a certain uncategorisability of the world.

\end{problem}

\newpage
\begin{problem}{}{}
% Subject: varc
\textbf{Instructions}

Instructions   

**The passage below is accompanied by a set of questions. Choose the best answer to each question.**
We cannot travel outside our neighbourhood without passports. We must wear the same plain clothes. We must exchange our houses every ten years. We cannot avoid labour. We all go to bed at the same time . . . We have religious freedom, but we cannot deny that the soul dies with the body, since ‘but for the fear of punishment, they would have nothing but contempt for the laws and customs of society'. . . . In More’s time, for much of the population, given the plenty and security on offer, such restraints would not have seemed overly unreasonable. For modern readers, however, Utopia appears to rely upon relentless transparency, the repression of variety, and the curtailment of privacy. Utopia provides security: but at what price? In both its external and internal relations, indeed, it seems perilously dystopian.
Such a conclusion might be fortified by examining selectively the tradition which follows More on these points. This often portrays societies where . . . 'it would be almost impossible for man to be depraved, or wicked'. . . . This is achieved both through institutions and mores, which underpin the common life. . . . The passions are regulated and inequalities of wealth and distinction are minimized. Needs, vanity, and emulation are restrained, often by prizing equality and holding riches in contempt. The desire for public power is curbed. Marriage and sexual intercourse are often controlled: in Tommaso Campanella’s The City of the Sun (1623), the first great literary utopia after More’s, relations are forbidden to men before the age of twenty-one and women before nineteen. Communal child-rearing is normal; for Campanella, this commences at age two. Greater simplicity of life, ‘living according to nature’, is often a result: the desire for simplicity and purity are closely related. People become more alike in appearance, opinion, and outlook than they often have been. Unity, order, and homogeneity thus prevail at the cost of individuality and diversity. This model, as J. C. Davis demonstrates, dominated early modern utopianism. . . . And utopian homogeneity remains a familiar theme well into the twentieth century.
Given these considerations, it is not unreasonable to take as our starting point here the hypothesis that utopia and dystopia evidently share more in common than is often supposed. Indeed, they might be twins, the progeny of the same parents. Insofar as this proves to be the case, my linkage of both here will be uncomfortably close for some readers. Yet we should not mistake this argument for the assertion that all utopias are, or tend to produce, dystopias. Those who defend this proposition will find that their association here is not nearly close enough. For we have only to acknowledge the existence of thousands of successful intentional communities in which a cooperative ethos predominates and where harmony without coercion is the rule to set aside such an assertion. Here the individual’s submersion in the group is consensual (though this concept is not unproblematic). It results not in enslavement but voluntary submission to group norms. Harmony is achieved without . . . harming others.

\vspace{1em}
\textbf{Question 5}

Which one of the following, if true about the Classic Maya, would invalidate the purpose of the iPhone example in the passage?

\begin{enumerate}[label=\Alph*.]
    \item The clay incense burner with spiky appliques was categorised only as a person and not as a tree by the Classic Maya.
    \item Unlike modern societies equipped with mobile phones, the Classic Maya did not have any communicating objects.
    \item Classic Maya songs represent both humans and non-living objects as characters, talking and interacting with each other.
    \item The personhood of the incense burner and the stone chopper was a function of their usefulness to humans.
[](https://cracku.in/5-which-one-of-the-following-if-true-about-the-class-x-cat-2021-slot-1)
\end{enumerate}
\vspace{1em}
\textbf{Solution:}

The author supplements the example of the i-phone with a pertinent question: "_Does it have personhood (if at all) because it is connected to me, drawing this resource from me as an owner or source?_ " He proceeds to then highlight the key takeaway from the example: "_For the Maya (who did have plenty of other communicating objects, if not smartphones), the answer was no. Nonhuman persons were not tethered to specific humans, and they did not derive their personhood from a connection with a human._ " The end idea: the personhood of an object is not a function of its utility or attachment to humans {an object can be categorised as a person based on certain distinct variables aside from its relation to a human}. The only relevant information invalidating this portrayal is in Option D: if the personhood of the incense burner or stone chopper is dependent on their usefulness to humans, the purpose of presenting the example and the associated idea is undermined.
It is unclear how Options A and B invalidate the purpose of the example. Option C aligns with the author's assertions and does not touch upon the idea emphasised by the example. 
Hence, Option D is the correct choice.

\end{problem}

\newpage
\begin{problem}{}{}
% Subject: varc
\textbf{Instructions}

Instructions   

**The passage below is accompanied by a set of questions. Choose the best answer to each question.**
We cannot travel outside our neighbourhood without passports. We must wear the same plain clothes. We must exchange our houses every ten years. We cannot avoid labour. We all go to bed at the same time . . . We have religious freedom, but we cannot deny that the soul dies with the body, since ‘but for the fear of punishment, they would have nothing but contempt for the laws and customs of society'. . . . In More’s time, for much of the population, given the plenty and security on offer, such restraints would not have seemed overly unreasonable. For modern readers, however, Utopia appears to rely upon relentless transparency, the repression of variety, and the curtailment of privacy. Utopia provides security: but at what price? In both its external and internal relations, indeed, it seems perilously dystopian.
Such a conclusion might be fortified by examining selectively the tradition which follows More on these points. This often portrays societies where . . . 'it would be almost impossible for man to be depraved, or wicked'. . . . This is achieved both through institutions and mores, which underpin the common life. . . . The passions are regulated and inequalities of wealth and distinction are minimized. Needs, vanity, and emulation are restrained, often by prizing equality and holding riches in contempt. The desire for public power is curbed. Marriage and sexual intercourse are often controlled: in Tommaso Campanella’s The City of the Sun (1623), the first great literary utopia after More’s, relations are forbidden to men before the age of twenty-one and women before nineteen. Communal child-rearing is normal; for Campanella, this commences at age two. Greater simplicity of life, ‘living according to nature’, is often a result: the desire for simplicity and purity are closely related. People become more alike in appearance, opinion, and outlook than they often have been. Unity, order, and homogeneity thus prevail at the cost of individuality and diversity. This model, as J. C. Davis demonstrates, dominated early modern utopianism. . . . And utopian homogeneity remains a familiar theme well into the twentieth century.
Given these considerations, it is not unreasonable to take as our starting point here the hypothesis that utopia and dystopia evidently share more in common than is often supposed. Indeed, they might be twins, the progeny of the same parents. Insofar as this proves to be the case, my linkage of both here will be uncomfortably close for some readers. Yet we should not mistake this argument for the assertion that all utopias are, or tend to produce, dystopias. Those who defend this proposition will find that their association here is not nearly close enough. For we have only to acknowledge the existence of thousands of successful intentional communities in which a cooperative ethos predominates and where harmony without coercion is the rule to set aside such an assertion. Here the individual’s submersion in the group is consensual (though this concept is not unproblematic). It results not in enslavement but voluntary submission to group norms. Harmony is achieved without . . . harming others.

\vspace{1em}
\textbf{Question 6}

Which one of the following best explains the “additional complexity” that the example of the incense burner illustrates regarding personhood for the Classic Maya?

\begin{enumerate}[label=\Alph*.]
    \item The example adds a new layer to the nonbinary understanding of personhood by bringing in a third category that shares a similar relation with the previous two.
    \item The example provides an exception to the nonbinary understanding of personhood that the passage had hitherto established.
    \item The example adds a new layer to the nonbinary understanding of personhood by bringing in a third category that shares a dissimilar relation with the previous two.
    \item The example complicates the nonbinary understanding of personhood by bringing in the sacred, establishing the porosity of the divine and the profane.
[](https://cracku.in/6-which-one-of-the-following-best-explains-the-addit-x-cat-2021-slot-1)
\end{enumerate}
\vspace{1em}
\textbf{Solution:}

_One additional complexity: the incense burner (which would have been made of clay, and decorated with spiky appliques representing the sacred ceiba tree found in this region) is categorised as a person - but also as a tree._  

The additional complexity that the author talks about here is the addition of another layer in the non-binary understanding of personhood. The incense burner was already classified as a person, but now it has also been classified as a tree. Hence, we have Options A and C. Note that the third category, that is tree, has a relation with the previous two categories. The boundary separating tree and person is porous. And since the incense burner has been categorized as a tree too, the relationship between them is porous too. Hence, we can infer that the third category shares a similar relationship with the previous two categories, and A is the correct answer.
The author is not exemplifying an exception but citing an additional complexity that is present in the definition. Hence, B can be eliminated.
The author does not establish the porosity of the divine and the profane. Hence, Option D is out of the scope of the passage.

\end{problem}

\newpage
\begin{problem}{}{}
% Subject: varc
\textbf{Instructions}

Instructions   

**The passage below is accompanied by a set of questions. Choose the best answer to each question.**
We cannot travel outside our neighbourhood without passports. We must wear the same plain clothes. We must exchange our houses every ten years. We cannot avoid labour. We all go to bed at the same time . . . We have religious freedom, but we cannot deny that the soul dies with the body, since ‘but for the fear of punishment, they would have nothing but contempt for the laws and customs of society'. . . . In More’s time, for much of the population, given the plenty and security on offer, such restraints would not have seemed overly unreasonable. For modern readers, however, Utopia appears to rely upon relentless transparency, the repression of variety, and the curtailment of privacy. Utopia provides security: but at what price? In both its external and internal relations, indeed, it seems perilously dystopian.
Such a conclusion might be fortified by examining selectively the tradition which follows More on these points. This often portrays societies where . . . 'it would be almost impossible for man to be depraved, or wicked'. . . . This is achieved both through institutions and mores, which underpin the common life. . . . The passions are regulated and inequalities of wealth and distinction are minimized. Needs, vanity, and emulation are restrained, often by prizing equality and holding riches in contempt. The desire for public power is curbed. Marriage and sexual intercourse are often controlled: in Tommaso Campanella’s The City of the Sun (1623), the first great literary utopia after More’s, relations are forbidden to men before the age of twenty-one and women before nineteen. Communal child-rearing is normal; for Campanella, this commences at age two. Greater simplicity of life, ‘living according to nature’, is often a result: the desire for simplicity and purity are closely related. People become more alike in appearance, opinion, and outlook than they often have been. Unity, order, and homogeneity thus prevail at the cost of individuality and diversity. This model, as J. C. Davis demonstrates, dominated early modern utopianism. . . . And utopian homogeneity remains a familiar theme well into the twentieth century.
Given these considerations, it is not unreasonable to take as our starting point here the hypothesis that utopia and dystopia evidently share more in common than is often supposed. Indeed, they might be twins, the progeny of the same parents. Insofar as this proves to be the case, my linkage of both here will be uncomfortably close for some readers. Yet we should not mistake this argument for the assertion that all utopias are, or tend to produce, dystopias. Those who defend this proposition will find that their association here is not nearly close enough. For we have only to acknowledge the existence of thousands of successful intentional communities in which a cooperative ethos predominates and where harmony without coercion is the rule to set aside such an assertion. Here the individual’s submersion in the group is consensual (though this concept is not unproblematic). It results not in enslavement but voluntary submission to group norms. Harmony is achieved without . . . harming others.

\vspace{1em}
\textbf{Question 7}

On the basis of the passage, which one of the following worldviews can be inferred to be closest to that of the Classic Maya?

\begin{enumerate}[label=\Alph*.]
    \item A tribe that perceives its hunting weapons as sacred person-artefacts because of their significance to its survival.
    \item A tribe that perceives plants as person-plants because they form an ecosystem and are marked by needs of nutrition.
    \item A futuristic society that perceives robots to be persons as well as robots because of their similarity to humans.
    \item A tribe that perceives its utensils as person-utensils in light of their functionality and bodily needs.
[](https://cracku.in/7-on-the-basis-of-the-passage-which-one-of-the-follo-x-cat-2021-slot-1)
\end{enumerate}
\vspace{1em}
\textbf{Solution:}

The author highlights multiple elements that constitute the Classic Mayan worldview pertaining to personhood: 
  * "_Nonhuman persons were not tethered to specific humans, and they did not derive their personhood from a connection with a human._ "
  * "_The Maya saw personhood as ‘activated’ by experiencing certain bodily needs and through participation in certain social activities. For example, among the faced objects that I examined, persons are marked by personal requirements (such as hunger, tiredness, physical closeness), and by community obligations (communication, interaction, ritual observance)_ "
  * "_Personhood was a nonbinary proposition for the Maya. Entities were able to be persons while also being something else...With these Maya examples, we are challenged to discard the person/nonperson binary that constitutes our basic ontological outlook_ "


Inspecting the options using the above ideas as filters, we notice that Option B is closest to the Classic Mayan worldview. A tribe that "perceives plants as person-plants because they form an ecosystem and are marked by needs of nutrition" acknowledges the two elementary variables for defining personhood - _personal requirements_ and  _community obligations_. 
Options A and C tether the personhood of the objects to their utility to humans - this does not coincide with the Classic Mayan belief. Although Option D mentions bodily needs, the interpretation of the term 'functionality' remains unclear. Hence, we can eliminate Option D.

\end{problem}

\newpage
\begin{problem}{}{}
% Subject: varc
\textbf{Instructions}

Instructions   

**The passage below is accompanied by a set of questions. Choose the best answer to each question.**
We cannot travel outside our neighbourhood without passports. We must wear the same plain clothes. We must exchange our houses every ten years. We cannot avoid labour. We all go to bed at the same time . . . We have religious freedom, but we cannot deny that the soul dies with the body, since ‘but for the fear of punishment, they would have nothing but contempt for the laws and customs of society'. . . . In More’s time, for much of the population, given the plenty and security on offer, such restraints would not have seemed overly unreasonable. For modern readers, however, Utopia appears to rely upon relentless transparency, the repression of variety, and the curtailment of privacy. Utopia provides security: but at what price? In both its external and internal relations, indeed, it seems perilously dystopian.
Such a conclusion might be fortified by examining selectively the tradition which follows More on these points. This often portrays societies where . . . 'it would be almost impossible for man to be depraved, or wicked'. . . . This is achieved both through institutions and mores, which underpin the common life. . . . The passions are regulated and inequalities of wealth and distinction are minimized. Needs, vanity, and emulation are restrained, often by prizing equality and holding riches in contempt. The desire for public power is curbed. Marriage and sexual intercourse are often controlled: in Tommaso Campanella’s The City of the Sun (1623), the first great literary utopia after More’s, relations are forbidden to men before the age of twenty-one and women before nineteen. Communal child-rearing is normal; for Campanella, this commences at age two. Greater simplicity of life, ‘living according to nature’, is often a result: the desire for simplicity and purity are closely related. People become more alike in appearance, opinion, and outlook than they often have been. Unity, order, and homogeneity thus prevail at the cost of individuality and diversity. This model, as J. C. Davis demonstrates, dominated early modern utopianism. . . . And utopian homogeneity remains a familiar theme well into the twentieth century.
Given these considerations, it is not unreasonable to take as our starting point here the hypothesis that utopia and dystopia evidently share more in common than is often supposed. Indeed, they might be twins, the progeny of the same parents. Insofar as this proves to be the case, my linkage of both here will be uncomfortably close for some readers. Yet we should not mistake this argument for the assertion that all utopias are, or tend to produce, dystopias. Those who defend this proposition will find that their association here is not nearly close enough. For we have only to acknowledge the existence of thousands of successful intentional communities in which a cooperative ethos predominates and where harmony without coercion is the rule to set aside such an assertion. Here the individual’s submersion in the group is consensual (though this concept is not unproblematic). It results not in enslavement but voluntary submission to group norms. Harmony is achieved without . . . harming others.

\vspace{1em}
\textbf{Question 8}

Which one of the following, if true, would not undermine the democratising potential of the Classic Maya worldview?

\begin{enumerate}[label=\Alph*.]
    \item They believed that animals like cats and dogs that live in proximity to humans have a more clearly articulated personhood.
    \item While they believed in the personhood of objects and plants, they did not believe in the personhood of rivers and animals.
    \item They understood the stone implement and the incense burner in a purely human form.
    \item They depicted their human healers with physical attributes of local medicinal plants.
[](https://cracku.in/8-which-one-of-the-following-if-true-would-not-under-x-cat-2021-slot-1)
\end{enumerate}
\vspace{1em}
\textbf{Solution:}

The idea concerning the democratising potential can be retraced to the first paragraph wherein the author states: "_Nonhuman persons were not tethered to specific humans, and they did not derive their personhood from a connection with a human. . . . It’s a profoundly democratising way of understanding the world. Humans are not more important persons - we are just one of many kinds of persons who inhabit this world_ "
_Option A_ : Considering proximity as an idea would undermine the portrayal of the Classic Mayan worldview. The author presents the example of the I-phone to convey how the personhood of an object is not a function of its utility or attachment to humans. If true, the statement in A would counter the premise of this example. Hence, we can eliminate this choice. 
_Option B_ : the assessment here is quite similar to Option A; if we create distinctions within the realm of inanimate objects, this will weaken the Mayan worldview. It would diminish the democratising potential of such a viewpoint by introducing specific barriers or criteria for the classification of personhood.
_Option C_ : The claim here runs against the Mayan idea of personhood being nonbinary. Thus, we can eliminate it since it undermines the democratising potential of the Classic Mayan worldview. 
Hence, Option D is the correct choice.

Instructions   

**The passage below is accompanied by a set of questions. Choose the best answer to each question.**
The sleights of hand that conflate consumption with virtue are a central theme in A Thirst for Empire, a sweeping and richly detailed history of tea by the historian Erika Rappaport. How did tea evolve from an obscure “China drink” to a universal beverage imbued with civilising properties? The answer, in brief, revolves around this conflation, not only by profit-motivated marketers but by a wide variety of interest groups. While abundant historical records have allowed the study of how tea itself moved from east to west, Rappaport is focused on the movement of the idea of tea to suit particular purposes.
Beginning in the 1700s, the temperance movement advocated for tea as a pleasure that cheered but did not inebriate, and industrialists soon borrowed this moral argument in advancing their case for free trade in tea (and hence more open markets for their textiles). Factory owners joined in, compelled by the cause of a sober workforce, while Christian missionaries discovered that tea “would soothe any colonial encounter”. During the Second World War, tea service was presented as a social and patriotic activity that uplifted soldiers and calmed refugees.
But it was tea’s consumer-directed marketing by importers and retailers - and later by brands - that most closely portends current trade debates. An early version of the “farm to table” movement was sparked by anti-Chinese sentiment and concerns over trade deficits, as well as by the reality and threat of adulterated tea containing dirt and hedge clippings. Lipton was soon advertising “from the Garden to Tea Cup” supply chains originating in British India and supervised by “educated Englishmen”. While tea marketing always presented direct consumer benefits (health, energy, relaxation), tea drinkers were also assured that they were participating in a larger noble project that advanced the causes of family, nation and civilization. . . .
Rappaport’s treatment of her subject is refreshingly apolitical. Indeed, it is a virtue that readers will be unable to guess her political orientation: both the miracle of markets and capitalism’s dark underbelly are evident in tea’s complex story, as are the complicated effects of British colonialism. . . . Commodity histories are now themselves commodities: recent works investigate cotton, salt, cod, sugar, chocolate, paper and milk. And morality marketing is now a commodity as well, applied to food, “fair trade” apparel and eco-tourism. Yet tea is, Rappaport makes clear, a world apart - an astonishing success story in which tea marketers not only succeeded in conveying a sense of moral elevation to the consumer but also arguably did advance the cause of civilisation and community.
I have been offered tea at a British garden party, a Bedouin campfire, a Turkish carpet shop and a Japanese chashitsu, to name a few settings. In each case the offering was more an idea - friendship, community, respect - than a drink, and in each case the idea then created a reality. It is not a stretch to say that tea marketers have advanced the particularly noble cause of human dialogue and friendship.

\end{problem}

\newpage
\begin{problem}{}{}
% Subject: varc
\textbf{Instructions}

Instructions   

**The passage below is accompanied by a set of questions. Choose the best answer to each question.**
We cannot travel outside our neighbourhood without passports. We must wear the same plain clothes. We must exchange our houses every ten years. We cannot avoid labour. We all go to bed at the same time . . . We have religious freedom, but we cannot deny that the soul dies with the body, since ‘but for the fear of punishment, they would have nothing but contempt for the laws and customs of society'. . . . In More’s time, for much of the population, given the plenty and security on offer, such restraints would not have seemed overly unreasonable. For modern readers, however, Utopia appears to rely upon relentless transparency, the repression of variety, and the curtailment of privacy. Utopia provides security: but at what price? In both its external and internal relations, indeed, it seems perilously dystopian.
Such a conclusion might be fortified by examining selectively the tradition which follows More on these points. This often portrays societies where . . . 'it would be almost impossible for man to be depraved, or wicked'. . . . This is achieved both through institutions and mores, which underpin the common life. . . . The passions are regulated and inequalities of wealth and distinction are minimized. Needs, vanity, and emulation are restrained, often by prizing equality and holding riches in contempt. The desire for public power is curbed. Marriage and sexual intercourse are often controlled: in Tommaso Campanella’s The City of the Sun (1623), the first great literary utopia after More’s, relations are forbidden to men before the age of twenty-one and women before nineteen. Communal child-rearing is normal; for Campanella, this commences at age two. Greater simplicity of life, ‘living according to nature’, is often a result: the desire for simplicity and purity are closely related. People become more alike in appearance, opinion, and outlook than they often have been. Unity, order, and homogeneity thus prevail at the cost of individuality and diversity. This model, as J. C. Davis demonstrates, dominated early modern utopianism. . . . And utopian homogeneity remains a familiar theme well into the twentieth century.
Given these considerations, it is not unreasonable to take as our starting point here the hypothesis that utopia and dystopia evidently share more in common than is often supposed. Indeed, they might be twins, the progeny of the same parents. Insofar as this proves to be the case, my linkage of both here will be uncomfortably close for some readers. Yet we should not mistake this argument for the assertion that all utopias are, or tend to produce, dystopias. Those who defend this proposition will find that their association here is not nearly close enough. For we have only to acknowledge the existence of thousands of successful intentional communities in which a cooperative ethos predominates and where harmony without coercion is the rule to set aside such an assertion. Here the individual’s submersion in the group is consensual (though this concept is not unproblematic). It results not in enslavement but voluntary submission to group norms. Harmony is achieved without . . . harming others.

\vspace{1em}
\textbf{Question 9}

Today, “conflat[ing] consumption with virtue” can be seen in the marketing of:

\begin{enumerate}[label=\Alph*.]
    \item sustainably farmed foods.
    \item natural health supplements.
    \item travel to pristine destinations.
    \item ergonomically designed products.
[](https://cracku.in/9-today-conflating-consumption-with-virtue-can-be-se-x-cat-2021-slot-1)
\end{enumerate}
\vspace{1em}
\textbf{Solution:}

Although mildly subjective, the question tests our understanding of the central idea. Across the passage, we notice the 'conflation' of tea consumption with particular virtues: it was not merely limited to benefits to the consumer but served a greater purpose. The narrative was that by drinking tea, people were/are advancing the cause of civilisation and community {thereby, imparting a sense of moral elevation}. Thus, the welfare highlighted is two-fold: both the consumer and society is benefitted. Towards the end, the author supports the narrative as follows: "_It is not a stretch to say that tea marketers have advanced the particularly noble cause of human dialogue and friendship._ " Any choice showcasing this dual benefit might be the potential answer.
Option A presents sustainably farmed foods; it is easy to identify that any associated marketing mechanics will emphasise the benefit of such food to both people and the environment. Thus, advertisers will make a case for how sustainably farmed food is beneficial not just to the consumer but also to the world at large. This will be equivalent to the 'conflation' that we came across in the passage.
Option B mentions natural health supplements; although we can discern the benefit to the consumer, the benefits to society is hard to perceive. Similarly, Options C and D appear irrelevant; it is difficult to identify what virtues we are conflating here with the subject. 
Hence, of the given choices, Option A appears most appropriate.

\end{problem}

\newpage
\begin{problem}{}{}
% Subject: varc
\textbf{Instructions}

Instructions   

**The passage below is accompanied by a set of questions. Choose the best answer to each question.**
We cannot travel outside our neighbourhood without passports. We must wear the same plain clothes. We must exchange our houses every ten years. We cannot avoid labour. We all go to bed at the same time . . . We have religious freedom, but we cannot deny that the soul dies with the body, since ‘but for the fear of punishment, they would have nothing but contempt for the laws and customs of society'. . . . In More’s time, for much of the population, given the plenty and security on offer, such restraints would not have seemed overly unreasonable. For modern readers, however, Utopia appears to rely upon relentless transparency, the repression of variety, and the curtailment of privacy. Utopia provides security: but at what price? In both its external and internal relations, indeed, it seems perilously dystopian.
Such a conclusion might be fortified by examining selectively the tradition which follows More on these points. This often portrays societies where . . . 'it would be almost impossible for man to be depraved, or wicked'. . . . This is achieved both through institutions and mores, which underpin the common life. . . . The passions are regulated and inequalities of wealth and distinction are minimized. Needs, vanity, and emulation are restrained, often by prizing equality and holding riches in contempt. The desire for public power is curbed. Marriage and sexual intercourse are often controlled: in Tommaso Campanella’s The City of the Sun (1623), the first great literary utopia after More’s, relations are forbidden to men before the age of twenty-one and women before nineteen. Communal child-rearing is normal; for Campanella, this commences at age two. Greater simplicity of life, ‘living according to nature’, is often a result: the desire for simplicity and purity are closely related. People become more alike in appearance, opinion, and outlook than they often have been. Unity, order, and homogeneity thus prevail at the cost of individuality and diversity. This model, as J. C. Davis demonstrates, dominated early modern utopianism. . . . And utopian homogeneity remains a familiar theme well into the twentieth century.
Given these considerations, it is not unreasonable to take as our starting point here the hypothesis that utopia and dystopia evidently share more in common than is often supposed. Indeed, they might be twins, the progeny of the same parents. Insofar as this proves to be the case, my linkage of both here will be uncomfortably close for some readers. Yet we should not mistake this argument for the assertion that all utopias are, or tend to produce, dystopias. Those who defend this proposition will find that their association here is not nearly close enough. For we have only to acknowledge the existence of thousands of successful intentional communities in which a cooperative ethos predominates and where harmony without coercion is the rule to set aside such an assertion. Here the individual’s submersion in the group is consensual (though this concept is not unproblematic). It results not in enslavement but voluntary submission to group norms. Harmony is achieved without . . . harming others.

\vspace{1em}
\textbf{Question 10}

The author of this book review is LEAST likely to support the view that:

\begin{enumerate}[label=\Alph*.]
    \item tea drinking has become a social ritual worldwide.
    \item the ritual of drinking tea promotes congeniality and camaraderie.
    \item tea became the leading drink in Britain in the nineteenth century.
    \item tea drinking was sometimes promoted as a patriotic duty.
[](https://cracku.in/10-the-author-of-this-book-review-is-least-likely-to--x-cat-2021-slot-1)
\end{enumerate}
\vspace{1em}
\textbf{Solution:}

{_During the Second World War, tea service was presented as a social and patriotic activity that uplifted soldiers and calmed refugees._}
{_I have been offered tea at a British garden party, a Bedouin campfire, a Turkish carpet shop and a Japanese chashitsu, to name a few settings. In each case the offering was more an idea - friendship, community, respect - than a drink, and in each case the idea then created a reality. It is not a stretch to say that tea marketers have advanced the particularly noble cause of human dialogue and friendship._}
Options A, B and D have been implied in the above excerpts - we can safely assume that the author will agree to the claims made in these options. Option C, however, has not been stated in the passage. There is no information to deduce that the author will agree with the assertions that tea became the leading drink in Britain in the nineteenth century. Thus, Option C is the correct choice.

\end{problem}

\newpage
\begin{problem}{}{}
% Subject: varc
\textbf{Instructions}

Instructions   

**The passage below is accompanied by a set of questions. Choose the best answer to each question.**
We cannot travel outside our neighbourhood without passports. We must wear the same plain clothes. We must exchange our houses every ten years. We cannot avoid labour. We all go to bed at the same time . . . We have religious freedom, but we cannot deny that the soul dies with the body, since ‘but for the fear of punishment, they would have nothing but contempt for the laws and customs of society'. . . . In More’s time, for much of the population, given the plenty and security on offer, such restraints would not have seemed overly unreasonable. For modern readers, however, Utopia appears to rely upon relentless transparency, the repression of variety, and the curtailment of privacy. Utopia provides security: but at what price? In both its external and internal relations, indeed, it seems perilously dystopian.
Such a conclusion might be fortified by examining selectively the tradition which follows More on these points. This often portrays societies where . . . 'it would be almost impossible for man to be depraved, or wicked'. . . . This is achieved both through institutions and mores, which underpin the common life. . . . The passions are regulated and inequalities of wealth and distinction are minimized. Needs, vanity, and emulation are restrained, often by prizing equality and holding riches in contempt. The desire for public power is curbed. Marriage and sexual intercourse are often controlled: in Tommaso Campanella’s The City of the Sun (1623), the first great literary utopia after More’s, relations are forbidden to men before the age of twenty-one and women before nineteen. Communal child-rearing is normal; for Campanella, this commences at age two. Greater simplicity of life, ‘living according to nature’, is often a result: the desire for simplicity and purity are closely related. People become more alike in appearance, opinion, and outlook than they often have been. Unity, order, and homogeneity thus prevail at the cost of individuality and diversity. This model, as J. C. Davis demonstrates, dominated early modern utopianism. . . . And utopian homogeneity remains a familiar theme well into the twentieth century.
Given these considerations, it is not unreasonable to take as our starting point here the hypothesis that utopia and dystopia evidently share more in common than is often supposed. Indeed, they might be twins, the progeny of the same parents. Insofar as this proves to be the case, my linkage of both here will be uncomfortably close for some readers. Yet we should not mistake this argument for the assertion that all utopias are, or tend to produce, dystopias. Those who defend this proposition will find that their association here is not nearly close enough. For we have only to acknowledge the existence of thousands of successful intentional communities in which a cooperative ethos predominates and where harmony without coercion is the rule to set aside such an assertion. Here the individual’s submersion in the group is consensual (though this concept is not unproblematic). It results not in enslavement but voluntary submission to group norms. Harmony is achieved without . . . harming others.

\vspace{1em}
\textbf{Question 11}

According to this book review, A Thirst for Empire says that, in addition to “profitmotivated marketers”, tea drinking was promoted in Britain by all of the following EXCEPT:

\begin{enumerate}[label=\Alph*.]
    \item manufacturers who were pressing for duty-free imports.
    \item tea drinkers lobbying for product diversity.
    \item the anti-alcohol lobby as a substitute for the consumption of liquor.
    \item factories to instill sobriety in their labour.
[](https://cracku.in/11-according-to-this-book-review-a-thirst-for-empire--x-cat-2021-slot-1)
\end{enumerate}
\vspace{1em}
\textbf{Solution:}

We can refer to the following excerpt to examine the choices: {  _Beginning in the 1700s, the temperance movement advocated for tea as a pleasure that cheered but did not inebriate, and industrialists soon borrowed this moral argument in advancing their case for free trade in tea (and hence more open markets for their textiles). Factory owners joined in, compelled by the cause of a sober workforce, while Christian missionaries discovered that tea “would soothe any colonial encounter”._}
Options A, C and D can be directly inferred from the excerpt. The author does not present any information pertaining to tea drinkers lobbying for product diversity - hence, we can identify this as an incorrect reason. 
Therefore, Option B is the correct choice.

\end{problem}

\newpage
\begin{problem}{}{}
% Subject: varc
\textbf{Instructions}

Instructions   

**The passage below is accompanied by a set of questions. Choose the best answer to each question.**
We cannot travel outside our neighbourhood without passports. We must wear the same plain clothes. We must exchange our houses every ten years. We cannot avoid labour. We all go to bed at the same time . . . We have religious freedom, but we cannot deny that the soul dies with the body, since ‘but for the fear of punishment, they would have nothing but contempt for the laws and customs of society'. . . . In More’s time, for much of the population, given the plenty and security on offer, such restraints would not have seemed overly unreasonable. For modern readers, however, Utopia appears to rely upon relentless transparency, the repression of variety, and the curtailment of privacy. Utopia provides security: but at what price? In both its external and internal relations, indeed, it seems perilously dystopian.
Such a conclusion might be fortified by examining selectively the tradition which follows More on these points. This often portrays societies where . . . 'it would be almost impossible for man to be depraved, or wicked'. . . . This is achieved both through institutions and mores, which underpin the common life. . . . The passions are regulated and inequalities of wealth and distinction are minimized. Needs, vanity, and emulation are restrained, often by prizing equality and holding riches in contempt. The desire for public power is curbed. Marriage and sexual intercourse are often controlled: in Tommaso Campanella’s The City of the Sun (1623), the first great literary utopia after More’s, relations are forbidden to men before the age of twenty-one and women before nineteen. Communal child-rearing is normal; for Campanella, this commences at age two. Greater simplicity of life, ‘living according to nature’, is often a result: the desire for simplicity and purity are closely related. People become more alike in appearance, opinion, and outlook than they often have been. Unity, order, and homogeneity thus prevail at the cost of individuality and diversity. This model, as J. C. Davis demonstrates, dominated early modern utopianism. . . . And utopian homogeneity remains a familiar theme well into the twentieth century.
Given these considerations, it is not unreasonable to take as our starting point here the hypothesis that utopia and dystopia evidently share more in common than is often supposed. Indeed, they might be twins, the progeny of the same parents. Insofar as this proves to be the case, my linkage of both here will be uncomfortably close for some readers. Yet we should not mistake this argument for the assertion that all utopias are, or tend to produce, dystopias. Those who defend this proposition will find that their association here is not nearly close enough. For we have only to acknowledge the existence of thousands of successful intentional communities in which a cooperative ethos predominates and where harmony without coercion is the rule to set aside such an assertion. Here the individual’s submersion in the group is consensual (though this concept is not unproblematic). It results not in enslavement but voluntary submission to group norms. Harmony is achieved without . . . harming others.

\vspace{1em}
\textbf{Question 12}

This book review argues that, according to Rappaport, tea is unlike other “morality” products because it:

\begin{enumerate}[label=\Alph*.]
    \item had an actual beneficial effect on social interaction and society in general.
    \item was actively encouraged by interest groups in the government.
    \item was marketed by a wide range of interest groups.
    \item appealed to a universal group and not just to a niche section of people.
[](https://cracku.in/12-this-book-review-argues-that-according-to-rappapor-x-cat-2021-slot-1)
\end{enumerate}
\vspace{1em}
\textbf{Solution:}

An important clue to analyse the options lies in the following excerpt: {_Yet tea is, Rappaport makes clear, a world apart - an astonishing success story in which tea marketers not only succeeded in conveying a sense of moral elevation to the consumer but also arguably did advance the cause of civilisation and community._} It is emphasised that the moral praise that tea received was not mere talk - tea as a 'morality product' did produce desirable outcomes. Option A is closest to conveying this idea. Options B, C and D fail to highlight the attribute that separates tea from 'other morality products' {as per the passage}.
Hence, Option A is the correct choice.

Instructions   

**The passage below is accompanied by a set of questions. Choose the best answer to each question.**
Cuttlefish are full of personality, as behavioral ecologist Alexandra Schnell found out while researching the cephalopod's potential to display self-control. . . . “Self-control is thought to be the cornerstone of intelligence, as it is an important prerequisite for complex decision-making and planning for the future,” says Schnell . . .
[Schnell's] study used a modified version of the “marshmallow test” . . . During the original marshmallow test, psychologist Walter Mischel presented children between age four and six with one marshmallow. He told them that if they waited 15 minutes and didn’t eat it, he would give them a second marshmallow. A long-term follow-up study showed that the children who waited for the second marshmallow had more success later in life. . . . The cuttlefish version of the experiment looked a lot different. The researchers worked with six cuttlefish under nine months old and presented them with seafood instead of sweets. (Preliminary experiments showed that cuttlefishes’ favorite food is live grass shrimp, while raw prawns are so-so and Asian shore crab is nearly unacceptable.) Since the researchers couldn’t explain to the cuttlefish that they would need to wait for their shrimp, they trained them to recognize certain shapes that indicated when a food item would become available. The symbols were pasted on transparent drawers so that the cuttlefish could see the food that was stored inside. One drawer, labeled with a circle to mean “immediate,” held raw king prawn. Another drawer, labeled with a triangle to mean “delayed,” held live grass shrimp. During a control experiment, square labels meant “never.”
“If their self-control is flexible and I hadn’t just trained them to wait in any context, you would expect the cuttlefish to take the immediate reward [in the control], even if it’s their second preference,” says Schnell . . . and that’s what they did. That showed the researchers that cuttlefish wouldn’t reject the prawns if it was the only food available. In the experimental trials, the cuttlefish didn’t jump on the prawns if the live grass shrimp were labeled with a triangle— many waited for the shrimp drawer to open up. Each time the cuttlefish showed it could wait, the researchers tacked another ten seconds on to the next round of waiting before releasing the shrimp. The longest that a cuttlefish waited was 130 seconds.
Schnell [says] that the cuttlefish usually sat at the bottom of the tank and looked at the two food items while they waited, but sometimes, they would turn away from the king prawn “as if to distract themselves from the temptation of the immediate reward.” In past studies, humans, chimpanzees, parrots and dogs also tried to distract themselves while waiting for a reward.
Not every species can use self-control, but most of the animals that can share another trait in common: long, social lives. Cuttlefish, on the other hand, are solitary creatures that don’t form relationships even with mates or young. . . . “We don’t know if living in a social group is important for complex cognition unless we also show those abilities are lacking in less social species,” says . . . comparative psychologist Jennifer Vonk.

\end{problem}

\newpage
\begin{problem}{}{}
% Subject: varc
\textbf{Instructions}

Instructions   

**The passage below is accompanied by a set of questions. Choose the best answer to each question.**
We cannot travel outside our neighbourhood without passports. We must wear the same plain clothes. We must exchange our houses every ten years. We cannot avoid labour. We all go to bed at the same time . . . We have religious freedom, but we cannot deny that the soul dies with the body, since ‘but for the fear of punishment, they would have nothing but contempt for the laws and customs of society'. . . . In More’s time, for much of the population, given the plenty and security on offer, such restraints would not have seemed overly unreasonable. For modern readers, however, Utopia appears to rely upon relentless transparency, the repression of variety, and the curtailment of privacy. Utopia provides security: but at what price? In both its external and internal relations, indeed, it seems perilously dystopian.
Such a conclusion might be fortified by examining selectively the tradition which follows More on these points. This often portrays societies where . . . 'it would be almost impossible for man to be depraved, or wicked'. . . . This is achieved both through institutions and mores, which underpin the common life. . . . The passions are regulated and inequalities of wealth and distinction are minimized. Needs, vanity, and emulation are restrained, often by prizing equality and holding riches in contempt. The desire for public power is curbed. Marriage and sexual intercourse are often controlled: in Tommaso Campanella’s The City of the Sun (1623), the first great literary utopia after More’s, relations are forbidden to men before the age of twenty-one and women before nineteen. Communal child-rearing is normal; for Campanella, this commences at age two. Greater simplicity of life, ‘living according to nature’, is often a result: the desire for simplicity and purity are closely related. People become more alike in appearance, opinion, and outlook than they often have been. Unity, order, and homogeneity thus prevail at the cost of individuality and diversity. This model, as J. C. Davis demonstrates, dominated early modern utopianism. . . . And utopian homogeneity remains a familiar theme well into the twentieth century.
Given these considerations, it is not unreasonable to take as our starting point here the hypothesis that utopia and dystopia evidently share more in common than is often supposed. Indeed, they might be twins, the progeny of the same parents. Insofar as this proves to be the case, my linkage of both here will be uncomfortably close for some readers. Yet we should not mistake this argument for the assertion that all utopias are, or tend to produce, dystopias. Those who defend this proposition will find that their association here is not nearly close enough. For we have only to acknowledge the existence of thousands of successful intentional communities in which a cooperative ethos predominates and where harmony without coercion is the rule to set aside such an assertion. Here the individual’s submersion in the group is consensual (though this concept is not unproblematic). It results not in enslavement but voluntary submission to group norms. Harmony is achieved without . . . harming others.

\vspace{1em}
\textbf{Question 13}

All of the following constitute a point of difference between the “original” and “modified” versions of the marshmallow test EXCEPT that:

\begin{enumerate}[label=\Alph*.]
    \item the former had human subjects, while the latter had cuttlefish.
    \item the former correlated self-control and future success, while the latter correlated selfcontrol and survival advantages.
    \item the former used verbal communication with its subjects, while the latter had to develop a symbolic means of communication.
    \item the former was performed over a longer time span than the latter.
[](https://cracku.in/13-all-of-the-following-constitute-a-point-of-differe-x-cat-2021-slot-1)
\end{enumerate}
\vspace{1em}
\textbf{Solution:}

Options A, C and D have been explicitly stated in the passage (refer to the second paragraph):
_A_ : ".._.children between age four and six with one marshmallow. He told them that if they waited 15 minutes and didn’t eat it, he would give them a second marshmallow._.." Thus, children were the subject under observation in the original marshmallow experiment, while cuttlefish were studied in the modified version of the same. 
_C_ : "..._Since the researchers couldn’t explain to the cuttlefish that they would need to wait for their shrimp, they trained them to recognize certain shapes that indicated when a food item would become available._.." Option C merely rephrase this excerpt.
_D_ : ".._.A long-term follow-up study showed that the children who waited for the second marshmallow had more success later in life.._." Given that the researchers undertook a long term study to map the successes of children showcasing self-control, we can safely conclude that the cuttlefish-version of the experiment was undertaken over a relatively shorter period. 
Option B cannot be inferred from the discussion - there is no correlation between selfcontrol and survival advantages. Hence, B is the correct choice.

\end{problem}

\newpage
\begin{problem}{}{}
% Subject: varc
\textbf{Instructions}

Instructions   

**The passage below is accompanied by a set of questions. Choose the best answer to each question.**
We cannot travel outside our neighbourhood without passports. We must wear the same plain clothes. We must exchange our houses every ten years. We cannot avoid labour. We all go to bed at the same time . . . We have religious freedom, but we cannot deny that the soul dies with the body, since ‘but for the fear of punishment, they would have nothing but contempt for the laws and customs of society'. . . . In More’s time, for much of the population, given the plenty and security on offer, such restraints would not have seemed overly unreasonable. For modern readers, however, Utopia appears to rely upon relentless transparency, the repression of variety, and the curtailment of privacy. Utopia provides security: but at what price? In both its external and internal relations, indeed, it seems perilously dystopian.
Such a conclusion might be fortified by examining selectively the tradition which follows More on these points. This often portrays societies where . . . 'it would be almost impossible for man to be depraved, or wicked'. . . . This is achieved both through institutions and mores, which underpin the common life. . . . The passions are regulated and inequalities of wealth and distinction are minimized. Needs, vanity, and emulation are restrained, often by prizing equality and holding riches in contempt. The desire for public power is curbed. Marriage and sexual intercourse are often controlled: in Tommaso Campanella’s The City of the Sun (1623), the first great literary utopia after More’s, relations are forbidden to men before the age of twenty-one and women before nineteen. Communal child-rearing is normal; for Campanella, this commences at age two. Greater simplicity of life, ‘living according to nature’, is often a result: the desire for simplicity and purity are closely related. People become more alike in appearance, opinion, and outlook than they often have been. Unity, order, and homogeneity thus prevail at the cost of individuality and diversity. This model, as J. C. Davis demonstrates, dominated early modern utopianism. . . . And utopian homogeneity remains a familiar theme well into the twentieth century.
Given these considerations, it is not unreasonable to take as our starting point here the hypothesis that utopia and dystopia evidently share more in common than is often supposed. Indeed, they might be twins, the progeny of the same parents. Insofar as this proves to be the case, my linkage of both here will be uncomfortably close for some readers. Yet we should not mistake this argument for the assertion that all utopias are, or tend to produce, dystopias. Those who defend this proposition will find that their association here is not nearly close enough. For we have only to acknowledge the existence of thousands of successful intentional communities in which a cooperative ethos predominates and where harmony without coercion is the rule to set aside such an assertion. Here the individual’s submersion in the group is consensual (though this concept is not unproblematic). It results not in enslavement but voluntary submission to group norms. Harmony is achieved without . . . harming others.

\vspace{1em}
\textbf{Question 14}

Which one of the following, if true, would best complement the passage’s findings?

\begin{enumerate}[label=\Alph*.]
    \item Cuttlefish are equally fond of live grass shrimp and raw prawns.
    \item Cuttlefish live in big groups that exhibit sociability.
    \item Cuttlefish wait longer than 100 seconds for the shrimp drawer to open up.
    \item Cuttlefish cannot distinguish between geometrical shapes.
[](https://cracku.in/14-which-one-of-the-following-if-true-would-best-comp-x-cat-2021-slot-1)
\end{enumerate}
\vspace{1em}
\textbf{Solution:}

_Option A_ : If true, this would weaken the findings of the experiments. The methodology used to establish the trait of self-restraint in cuttlefish is based on the premise that cuttlefish prefer one specific kind of food over another. If we demonstrate that cuttlefish are equally fond of live grass shrimp and raw prawn, then the observations made in the study become invalid. Therefore, we can eliminate Option A.
_Option B_ : If true, the finding here would support the comments made in the final paragraph. The author states the following: {_Not every species can use self-control, but most of the animals that can share another trait in common: long, social lives._} If it is proved that cuttlefish fulfils this ancillary criterion of sociability, the primary claim made in the passage is strengthened. 
_Option C_ : If true, the statement here bears no significance to the passage’s findings. [We already know that _"..the longest that a cuttlefish waited was 130 seconds..._ "] 
_Option D_ : If true, this would weaken the findings of the experiments. Similar to Option A, the methodology used to establish the trait of self-restraint in cuttlefish is based on the premise that cuttlefish can distinguish between geometrical shapes. If we demonstrate that this information is false, then the observations made in the study become invalid. Therefore, we can eliminate Option D.
Hence, Option B is the correct choice.

\end{problem}

\newpage
\begin{problem}{}{}
% Subject: varc
\textbf{Instructions}

Instructions   

**The passage below is accompanied by a set of questions. Choose the best answer to each question.**
We cannot travel outside our neighbourhood without passports. We must wear the same plain clothes. We must exchange our houses every ten years. We cannot avoid labour. We all go to bed at the same time . . . We have religious freedom, but we cannot deny that the soul dies with the body, since ‘but for the fear of punishment, they would have nothing but contempt for the laws and customs of society'. . . . In More’s time, for much of the population, given the plenty and security on offer, such restraints would not have seemed overly unreasonable. For modern readers, however, Utopia appears to rely upon relentless transparency, the repression of variety, and the curtailment of privacy. Utopia provides security: but at what price? In both its external and internal relations, indeed, it seems perilously dystopian.
Such a conclusion might be fortified by examining selectively the tradition which follows More on these points. This often portrays societies where . . . 'it would be almost impossible for man to be depraved, or wicked'. . . . This is achieved both through institutions and mores, which underpin the common life. . . . The passions are regulated and inequalities of wealth and distinction are minimized. Needs, vanity, and emulation are restrained, often by prizing equality and holding riches in contempt. The desire for public power is curbed. Marriage and sexual intercourse are often controlled: in Tommaso Campanella’s The City of the Sun (1623), the first great literary utopia after More’s, relations are forbidden to men before the age of twenty-one and women before nineteen. Communal child-rearing is normal; for Campanella, this commences at age two. Greater simplicity of life, ‘living according to nature’, is often a result: the desire for simplicity and purity are closely related. People become more alike in appearance, opinion, and outlook than they often have been. Unity, order, and homogeneity thus prevail at the cost of individuality and diversity. This model, as J. C. Davis demonstrates, dominated early modern utopianism. . . . And utopian homogeneity remains a familiar theme well into the twentieth century.
Given these considerations, it is not unreasonable to take as our starting point here the hypothesis that utopia and dystopia evidently share more in common than is often supposed. Indeed, they might be twins, the progeny of the same parents. Insofar as this proves to be the case, my linkage of both here will be uncomfortably close for some readers. Yet we should not mistake this argument for the assertion that all utopias are, or tend to produce, dystopias. Those who defend this proposition will find that their association here is not nearly close enough. For we have only to acknowledge the existence of thousands of successful intentional communities in which a cooperative ethos predominates and where harmony without coercion is the rule to set aside such an assertion. Here the individual’s submersion in the group is consensual (though this concept is not unproblematic). It results not in enslavement but voluntary submission to group norms. Harmony is achieved without . . . harming others.

\vspace{1em}
\textbf{Question 15}

Which one of the following cannot be inferred from Alexandra Schnell’s experiment?

\begin{enumerate}[label=\Alph*.]
    \item Like human children, cuttlefish are capable of self-control.
    \item Intelligence in a species is impossible without sociability.
    \item Cuttlefish exercise choice when it comes to food.
    \item Cuttlefish exert self-control with the help of diversions.
[](https://cracku.in/15-which-one-of-the-following-cannot-be-inferred-from-x-cat-2021-slot-1)
\end{enumerate}
\vspace{1em}
\textbf{Solution:}

_Option A_ : {_Each time the cuttlefish showed it could wait, the researchers tacked another ten seconds on to the next round of waiting before releasing the shrimp. The longest that a cuttlefish waited was 130 seconds._} The results of Schnell's experiment indicated that cuttlefish exhibit self-restraint. This was shown to be the case with few children in the original marshmallow experiments. Hence, Option A is correct. 
_Option B_ : {_Not every species can use self-control, but most of the animals that can share another trait in common: long, social lives._} The author does not imply causation between intelligence and sociability. He merely highlights these two attributes: self-control is considered indicative of intelligence, and most such organisms showcase social lives as well. Hence, B is a distortion and cannot be inferred. 
_Option C_ : {_Preliminary experiments showed that cuttlefishes’ favorite food is live grass shrimp, while raw prawns are so-so and Asian shore crab is nearly unacceptable._} The above lines depict a preference for certain kinds of food; thus, C is true. 
_Option D_ : {_Schnell [says] that the cuttlefish usually sat at the bottom of the tank and looked at the two food items while they waited, but sometimes, they would turn away from the king prawn “as if to distract themselves from the temptation of the immediate reward._ ”} Exerting self-control through distractions has also been presented as a behavioural trait of the cuttlefish. 
Hence, Option B is the correct choice.

\end{problem}

\newpage
\begin{problem}{}{}
% Subject: varc
\textbf{Instructions}

Instructions   

**The passage below is accompanied by a set of questions. Choose the best answer to each question.**
We cannot travel outside our neighbourhood without passports. We must wear the same plain clothes. We must exchange our houses every ten years. We cannot avoid labour. We all go to bed at the same time . . . We have religious freedom, but we cannot deny that the soul dies with the body, since ‘but for the fear of punishment, they would have nothing but contempt for the laws and customs of society'. . . . In More’s time, for much of the population, given the plenty and security on offer, such restraints would not have seemed overly unreasonable. For modern readers, however, Utopia appears to rely upon relentless transparency, the repression of variety, and the curtailment of privacy. Utopia provides security: but at what price? In both its external and internal relations, indeed, it seems perilously dystopian.
Such a conclusion might be fortified by examining selectively the tradition which follows More on these points. This often portrays societies where . . . 'it would be almost impossible for man to be depraved, or wicked'. . . . This is achieved both through institutions and mores, which underpin the common life. . . . The passions are regulated and inequalities of wealth and distinction are minimized. Needs, vanity, and emulation are restrained, often by prizing equality and holding riches in contempt. The desire for public power is curbed. Marriage and sexual intercourse are often controlled: in Tommaso Campanella’s The City of the Sun (1623), the first great literary utopia after More’s, relations are forbidden to men before the age of twenty-one and women before nineteen. Communal child-rearing is normal; for Campanella, this commences at age two. Greater simplicity of life, ‘living according to nature’, is often a result: the desire for simplicity and purity are closely related. People become more alike in appearance, opinion, and outlook than they often have been. Unity, order, and homogeneity thus prevail at the cost of individuality and diversity. This model, as J. C. Davis demonstrates, dominated early modern utopianism. . . . And utopian homogeneity remains a familiar theme well into the twentieth century.
Given these considerations, it is not unreasonable to take as our starting point here the hypothesis that utopia and dystopia evidently share more in common than is often supposed. Indeed, they might be twins, the progeny of the same parents. Insofar as this proves to be the case, my linkage of both here will be uncomfortably close for some readers. Yet we should not mistake this argument for the assertion that all utopias are, or tend to produce, dystopias. Those who defend this proposition will find that their association here is not nearly close enough. For we have only to acknowledge the existence of thousands of successful intentional communities in which a cooperative ethos predominates and where harmony without coercion is the rule to set aside such an assertion. Here the individual’s submersion in the group is consensual (though this concept is not unproblematic). It results not in enslavement but voluntary submission to group norms. Harmony is achieved without . . . harming others.

\vspace{1em}
\textbf{Question 16}

In which one of the following scenarios would the cuttlefish’s behaviour demonstrate self-control?

\begin{enumerate}[label=\Alph*.]
    \item Asian shore crabs and raw prawns are simultaneously released while a live grass shrimp drawer labelled with a triangle is placed in front of the cuttlefish, to be opened after one minute.
    \item raw prawns are released while a live grass shrimp drawer labelled with a square is placed in front of the cuttlefish.
    \item raw prawns are released while an Asian shore crab drawer labelled with a triangle is placed in front of the cuttlefish, to be opened after one minute.
    \item live grass shrimp are released while two raw prawn drawers labelled with a circle and a triangle respectively are placed in front of the cuttlefish; the triangle-labelled drawer is opened after 50 seconds.
[](https://cracku.in/16-in-which-one-of-the-following-scenarios-would-the--x-cat-2021-slot-1)
\end{enumerate}
\vspace{1em}
\textbf{Solution:}

The question tests our understanding of the experiment stated in the second paragraph. The key highlights of the modified cuttlefish experiment are:
(a) Food choice: 1st preference - live grass shrimp; 2nd preference - raw prawns; 3rd preference - Asian shore crab 
(b) Symbol based training: Circle - immediate availability; Triangle - delayed availability; Square - never available
(c) Observations: in the absence of 1st preference ---> cuttlefish will go for 2nd preference {applies to all such scenarios involving only one food choice}; in the presence of multiple food choices, cuttlefish will wait ---> if 1st pref is available and associated with either Circle or Triangle
Based on the above, we can filter out the given options:
_Option A_ : We know that raw prawns and Asian shore crab are not the primary preference of the cuttlefish. Additionally, we know that live grass shrimp (1st presence) is available for delayed consumption {associated with a Triangle}. If the cuttlefish waits for one minute to consume live grass shrimp and ignores the other two food choices, this definitively showcases that cuttlefish exert self-control. Hence, Option A is a strong candidate for the correct choice since it supplements the experiment's findings. 
_Option B_ : In this case, the cuttlefish will go for the raw prawns since it has been conditioned to understand that the box labelled with Square will never open. This does not contribute to establishing self-control in cuttlefish. Thus, we can reject this choice. 
_Option C_ : In this case, the cuttlefish will go for the raw prawns and avoid the Asian shore crab irrespective of the box it is placed in. This is because raw prawns fare more favourably as a food choice than Asian shore crab. {We know from points (a) and (c) that the cuttlefish will go for the 2nd preference over the 3rd} This does not contribute to establishing self-control in cuttlefish. Thus, we can reject this choice. 
_Option D_ : In this case, the cuttlefish will go for the live grass shrimp and avoid the raw prawns irrespective of the box it is placed in. This is because live grass shrimp fares more favourably as a food choice than raw prawns. {We know from points (a) and (c) that the cuttlefish will go for the 1st preference over the 2nd} This does not contribute to establishing self-control in cuttlefish. Therefore, we can reject this choice.
Hence, the correct answer is Option A.

Instructions   

For the following questions answer them individually

\end{problem}

\newpage
\begin{problem}{}{}
% Subject: varc
\textbf{Instructions}

Instructions   

**The passage below is accompanied by a set of questions. Choose the best answer to each question.**
We cannot travel outside our neighbourhood without passports. We must wear the same plain clothes. We must exchange our houses every ten years. We cannot avoid labour. We all go to bed at the same time . . . We have religious freedom, but we cannot deny that the soul dies with the body, since ‘but for the fear of punishment, they would have nothing but contempt for the laws and customs of society'. . . . In More’s time, for much of the population, given the plenty and security on offer, such restraints would not have seemed overly unreasonable. For modern readers, however, Utopia appears to rely upon relentless transparency, the repression of variety, and the curtailment of privacy. Utopia provides security: but at what price? In both its external and internal relations, indeed, it seems perilously dystopian.
Such a conclusion might be fortified by examining selectively the tradition which follows More on these points. This often portrays societies where . . . 'it would be almost impossible for man to be depraved, or wicked'. . . . This is achieved both through institutions and mores, which underpin the common life. . . . The passions are regulated and inequalities of wealth and distinction are minimized. Needs, vanity, and emulation are restrained, often by prizing equality and holding riches in contempt. The desire for public power is curbed. Marriage and sexual intercourse are often controlled: in Tommaso Campanella’s The City of the Sun (1623), the first great literary utopia after More’s, relations are forbidden to men before the age of twenty-one and women before nineteen. Communal child-rearing is normal; for Campanella, this commences at age two. Greater simplicity of life, ‘living according to nature’, is often a result: the desire for simplicity and purity are closely related. People become more alike in appearance, opinion, and outlook than they often have been. Unity, order, and homogeneity thus prevail at the cost of individuality and diversity. This model, as J. C. Davis demonstrates, dominated early modern utopianism. . . . And utopian homogeneity remains a familiar theme well into the twentieth century.
Given these considerations, it is not unreasonable to take as our starting point here the hypothesis that utopia and dystopia evidently share more in common than is often supposed. Indeed, they might be twins, the progeny of the same parents. Insofar as this proves to be the case, my linkage of both here will be uncomfortably close for some readers. Yet we should not mistake this argument for the assertion that all utopias are, or tend to produce, dystopias. Those who defend this proposition will find that their association here is not nearly close enough. For we have only to acknowledge the existence of thousands of successful intentional communities in which a cooperative ethos predominates and where harmony without coercion is the rule to set aside such an assertion. Here the individual’s submersion in the group is consensual (though this concept is not unproblematic). It results not in enslavement but voluntary submission to group norms. Harmony is achieved without . . . harming others.

\vspace{1em}
\textbf{Question 17}

**The passage given below is followed by four alternate summaries. Choose the option that best captures the essence of the passage.**
McGurk and MacDonald (1976) reported a powerful multisensory illusion occurring with audio-visual speech. They recorded a voice articulating a consonant ‘ba-ba-ba’ and dubbed it with a face articulating another consonant ‘ga-ga-ga’. Even though the acoustic speech signal was well recognized alone, it was heard as another consonant after dubbing with incongruent visual speech i.e., ‘da-da-da’. The illusion, termed as the McGurk effect, has been replicated many times, and it has sparked an abundance of research. The reason for the great impact is that this is a striking demonstration of multisensory integration, where that auditory and visual information is merged into a unified, integrated percept.

\begin{enumerate}[label=\Alph*.]
    \item Visual speech mismatched with auditory speech can result in the perception of an entirely different message: this illusion is known as the McGurk effect.
    \item When the quality of auditory information is poor, the visual information wins over the auditory information.
    \item The McGurk effect which is a demonstration of multisensory integration has been replicated many times.
    \item When the auditory speech signal does not match the visual speech movements, the acoustic speech signal is confusing and integration of the two is imperfect.
[](https://cracku.in/17-the-passage-given-below-is-followed-by-four-altern-x-cat-2021-slot-1)
\end{enumerate}
\vspace{1em}
\textbf{Solution:}

The main points of the paragraph are:
1. A multisensory illusion, dubbing a different visual cue to audio, makes the subject perceive a different sound. (Important point)
2. This illusion is called McGurk effect. (Important point. Related to 1)
3. An impactful subject of research as it demonstrates multisensory integration. (Secondary point. 1 and 2 can stand without this point)
Option A: Covers 1 and 2. Hence, a plausible option.
Option B: It distorts what the author is trying to say. It draws a conclusion out of the results of the study instead of paraphrasing the passage.
Option C: Option C covers only 3. It does not mention 1 and hence is not a good summary.
Option D: Mentions only 1. Not an apt summary.
Hence, the answer is Option A.

\end{problem}

\newpage
\begin{problem}{}{}
% Subject: varc
\textbf{Instructions}

Instructions   

**The passage below is accompanied by a set of questions. Choose the best answer to each question.**
We cannot travel outside our neighbourhood without passports. We must wear the same plain clothes. We must exchange our houses every ten years. We cannot avoid labour. We all go to bed at the same time . . . We have religious freedom, but we cannot deny that the soul dies with the body, since ‘but for the fear of punishment, they would have nothing but contempt for the laws and customs of society'. . . . In More’s time, for much of the population, given the plenty and security on offer, such restraints would not have seemed overly unreasonable. For modern readers, however, Utopia appears to rely upon relentless transparency, the repression of variety, and the curtailment of privacy. Utopia provides security: but at what price? In both its external and internal relations, indeed, it seems perilously dystopian.
Such a conclusion might be fortified by examining selectively the tradition which follows More on these points. This often portrays societies where . . . 'it would be almost impossible for man to be depraved, or wicked'. . . . This is achieved both through institutions and mores, which underpin the common life. . . . The passions are regulated and inequalities of wealth and distinction are minimized. Needs, vanity, and emulation are restrained, often by prizing equality and holding riches in contempt. The desire for public power is curbed. Marriage and sexual intercourse are often controlled: in Tommaso Campanella’s The City of the Sun (1623), the first great literary utopia after More’s, relations are forbidden to men before the age of twenty-one and women before nineteen. Communal child-rearing is normal; for Campanella, this commences at age two. Greater simplicity of life, ‘living according to nature’, is often a result: the desire for simplicity and purity are closely related. People become more alike in appearance, opinion, and outlook than they often have been. Unity, order, and homogeneity thus prevail at the cost of individuality and diversity. This model, as J. C. Davis demonstrates, dominated early modern utopianism. . . . And utopian homogeneity remains a familiar theme well into the twentieth century.
Given these considerations, it is not unreasonable to take as our starting point here the hypothesis that utopia and dystopia evidently share more in common than is often supposed. Indeed, they might be twins, the progeny of the same parents. Insofar as this proves to be the case, my linkage of both here will be uncomfortably close for some readers. Yet we should not mistake this argument for the assertion that all utopias are, or tend to produce, dystopias. Those who defend this proposition will find that their association here is not nearly close enough. For we have only to acknowledge the existence of thousands of successful intentional communities in which a cooperative ethos predominates and where harmony without coercion is the rule to set aside such an assertion. Here the individual’s submersion in the group is consensual (though this concept is not unproblematic). It results not in enslavement but voluntary submission to group norms. Harmony is achieved without . . . harming others.

\vspace{1em}
\textbf{Question 18}

**Five jumbled up sentences, related to a topic, are given below. Four of them can be put together to form a coherent paragraph. Identify the odd one out and key in the number of the sentence as your answer:**
1. There is a dark side to academic research, especially in India, and at its centre is the phenomenon of predatory journals.  
2. But in truth, as long as you pay, you can get anything published.  
3. In look and feel thus, they are exactly like any reputed journal.  
4. They claim to be indexed in the most influential databases, say they possess editorial boards that comprise top scientists and researchers, and claim to have a rigorous peer-review structure.  
5. But a large section of researchers and scientists across the world are at the receiving end of nothing short of an academic publishing scam.
Backspace  
789  
456  
123  
0.-  
Clear All  

Submit
[](https://cracku.in/18-five-jumbled-up-sentences-related-to-a-topic-are-g-x-cat-2021-slot-1)

\begin{enumerate}[label=\Alph*.]
\end{enumerate}
\vspace{1em}
\textbf{Solution:}

The given collection of statements focuses on predatory journals. The author begins by mentioning the subject {Statement 1} in a grim tone. He highlights the kind of claims that these journals state: presence in 'influential databases', quality 'editorial boards' and a 'rigorous peer-review structure' {Statement 4}. On the surface, 'they are exactly like any reputed journal' {Statement 3}. However, the ground reality is starkly different: paying money allows you to publish anything in such predatory journal {Statement 2}. Hence, arrangement 1432 forms a coherent paragraph with predatory journals as the centre of attention. 
Statement 5 deviates from this subject. Although the topic seems to be about academic publishing, the focus is no longer on predatory journals, but instead, it becomes a bit broader. Since there is a mismatch in scope, we discern Statement 5 as the odd one out.

\end{problem}

\newpage
\begin{problem}{}{}
% Subject: varc
\textbf{Instructions}

Instructions   

**The passage below is accompanied by a set of questions. Choose the best answer to each question.**
We cannot travel outside our neighbourhood without passports. We must wear the same plain clothes. We must exchange our houses every ten years. We cannot avoid labour. We all go to bed at the same time . . . We have religious freedom, but we cannot deny that the soul dies with the body, since ‘but for the fear of punishment, they would have nothing but contempt for the laws and customs of society'. . . . In More’s time, for much of the population, given the plenty and security on offer, such restraints would not have seemed overly unreasonable. For modern readers, however, Utopia appears to rely upon relentless transparency, the repression of variety, and the curtailment of privacy. Utopia provides security: but at what price? In both its external and internal relations, indeed, it seems perilously dystopian.
Such a conclusion might be fortified by examining selectively the tradition which follows More on these points. This often portrays societies where . . . 'it would be almost impossible for man to be depraved, or wicked'. . . . This is achieved both through institutions and mores, which underpin the common life. . . . The passions are regulated and inequalities of wealth and distinction are minimized. Needs, vanity, and emulation are restrained, often by prizing equality and holding riches in contempt. The desire for public power is curbed. Marriage and sexual intercourse are often controlled: in Tommaso Campanella’s The City of the Sun (1623), the first great literary utopia after More’s, relations are forbidden to men before the age of twenty-one and women before nineteen. Communal child-rearing is normal; for Campanella, this commences at age two. Greater simplicity of life, ‘living according to nature’, is often a result: the desire for simplicity and purity are closely related. People become more alike in appearance, opinion, and outlook than they often have been. Unity, order, and homogeneity thus prevail at the cost of individuality and diversity. This model, as J. C. Davis demonstrates, dominated early modern utopianism. . . . And utopian homogeneity remains a familiar theme well into the twentieth century.
Given these considerations, it is not unreasonable to take as our starting point here the hypothesis that utopia and dystopia evidently share more in common than is often supposed. Indeed, they might be twins, the progeny of the same parents. Insofar as this proves to be the case, my linkage of both here will be uncomfortably close for some readers. Yet we should not mistake this argument for the assertion that all utopias are, or tend to produce, dystopias. Those who defend this proposition will find that their association here is not nearly close enough. For we have only to acknowledge the existence of thousands of successful intentional communities in which a cooperative ethos predominates and where harmony without coercion is the rule to set aside such an assertion. Here the individual’s submersion in the group is consensual (though this concept is not unproblematic). It results not in enslavement but voluntary submission to group norms. Harmony is achieved without . . . harming others.

\vspace{1em}
\textbf{Question 19}

**The four sentences (labelled 1, 2, 3, 4) below, when properly sequenced would yield a coherent paragraph. Decide on the proper sequencing of the order of the sentences and key in the sequence of the four numbers as your answer:**
1. The work is more than the text, for the text only takes on life, when it is realized and furthermore the realization is by no means independent of the individual disposition of the reader.
2. The convergence of text and reader brings the literary work into existence and this convergence is not to be identified either with the reality of the text or with the individual disposition of the reader.
3. From this polarity it follows that the literary work cannot be completely identical with the text, or with the realization of the text, but in fact must lie halfway between the two.
4. The literary work has two poles, which we might call the artistic and the aesthetic; the artistic refers to the text created by the author, and the aesthetic to the realization accomplished by the reader.
Backspace  
789  
456  
123  
0.-  
Clear All  

Submit
[](https://cracku.in/19-the-four-sentences-labelled-1-2-3-4-below-when-pro-x-cat-2021-slot-1)

\begin{enumerate}[label=\Alph*.]
\end{enumerate}
\vspace{1em}
\textbf{Solution:}

A quick scan of the statements tells us that the statements discuss the nature of literary work - how it originates from the convergence of the text and the disposition of the readers. Note the phrase "From this**polarity** " in statement 3 - this must refer to the distinction presented in statement 4 ["literary work has two **poles** "] Statement 4 shows how the literary work comprises of two ends - (1) the artistic tied to the text, and (2) the aesthetic tied to the reader's realisation of the text. Statement 3 then underlines how the literary work cannot be equated to either of these and must lie somewhere midway netween the two. This idea is carried forward in Statement 1, which elaborates how the literary work is more than the text and becomes alive only when realised by the reader - in a way, this highlights how the literary work takes form due to the interplay between the artistic and the aesthetic. Statement 2 then extends on this idea of "convergence" [shown in 1] and reiterates the idea in 3 that the literary work, despite being an amalgamation of the two components, must not be mistaken with "the text or with the individual disposition of the reader." Hence, we obtain a coherent paragraph from the arrangement 4-3-1-2.

\end{problem}

\newpage
\begin{problem}{}{}
% Subject: varc
\textbf{Instructions}

Instructions   

**The passage below is accompanied by a set of questions. Choose the best answer to each question.**
We cannot travel outside our neighbourhood without passports. We must wear the same plain clothes. We must exchange our houses every ten years. We cannot avoid labour. We all go to bed at the same time . . . We have religious freedom, but we cannot deny that the soul dies with the body, since ‘but for the fear of punishment, they would have nothing but contempt for the laws and customs of society'. . . . In More’s time, for much of the population, given the plenty and security on offer, such restraints would not have seemed overly unreasonable. For modern readers, however, Utopia appears to rely upon relentless transparency, the repression of variety, and the curtailment of privacy. Utopia provides security: but at what price? In both its external and internal relations, indeed, it seems perilously dystopian.
Such a conclusion might be fortified by examining selectively the tradition which follows More on these points. This often portrays societies where . . . 'it would be almost impossible for man to be depraved, or wicked'. . . . This is achieved both through institutions and mores, which underpin the common life. . . . The passions are regulated and inequalities of wealth and distinction are minimized. Needs, vanity, and emulation are restrained, often by prizing equality and holding riches in contempt. The desire for public power is curbed. Marriage and sexual intercourse are often controlled: in Tommaso Campanella’s The City of the Sun (1623), the first great literary utopia after More’s, relations are forbidden to men before the age of twenty-one and women before nineteen. Communal child-rearing is normal; for Campanella, this commences at age two. Greater simplicity of life, ‘living according to nature’, is often a result: the desire for simplicity and purity are closely related. People become more alike in appearance, opinion, and outlook than they often have been. Unity, order, and homogeneity thus prevail at the cost of individuality and diversity. This model, as J. C. Davis demonstrates, dominated early modern utopianism. . . . And utopian homogeneity remains a familiar theme well into the twentieth century.
Given these considerations, it is not unreasonable to take as our starting point here the hypothesis that utopia and dystopia evidently share more in common than is often supposed. Indeed, they might be twins, the progeny of the same parents. Insofar as this proves to be the case, my linkage of both here will be uncomfortably close for some readers. Yet we should not mistake this argument for the assertion that all utopias are, or tend to produce, dystopias. Those who defend this proposition will find that their association here is not nearly close enough. For we have only to acknowledge the existence of thousands of successful intentional communities in which a cooperative ethos predominates and where harmony without coercion is the rule to set aside such an assertion. Here the individual’s submersion in the group is consensual (though this concept is not unproblematic). It results not in enslavement but voluntary submission to group norms. Harmony is achieved without . . . harming others.

\vspace{1em}
\textbf{Question 20}

**The passage given below is followed by four alternate summaries. Choose the option that best captures the essence of the passage.**  
Foreign peacekeepers often exist in a bubble in the poor countries in which they are deployed; they live in posh compounds, drive fancy vehicles, and distance themselves from locals. This may be partially justified as they are outsiders, living in constant fear, performing a job that is emotionally draining. But they are often despised by the locals, and many would like them to leave. A better solution would be bottom-up peacebuilding, which would involve their spending more time working with communities, understanding their grievances and earning their trust, rather than only meeting government officials.

\begin{enumerate}[label=\Alph*.]
    \item Peacekeeping duties would be more effectively performed by local residents given their better understanding, knowledge and rapport with their own communities.
    \item The environment in poor countries has tended to make foreign peacekeeping forces live in enclaves, but it is time to change this scenario.
    \item Extravagant lifestyles and an aloof attitude among the foreigners working as peacekeepers in poor countries have justifiably made them the target of local anger.
    \item Peacekeeping forces in foreign countries have tended to be aloof for valid reasons but would be more effective if they worked more closely with local communities.
[](https://cracku.in/20-the-passage-given-below-is-followed-by-four-altern-x-cat-2021-slot-1)
\end{enumerate}
\vspace{1em}
\textbf{Solution:}

The main points of the paragraph are:
1. The peacekeeping forces often exist in a bubble. Though there are valid reasons behind this, this also results in the locals feeling antipathy towards them.
2. The solution to this problem is to build rapport with the locals too instead of focusing only on the government officials.
Option A: Not implied in the paragraph. The paragraph suggests building relationships with the locals. Appointing only locals as peacekeepers has not been implied.
Option B: This option distorts what is being presented in the paragraph. The paragraph suggests that the bubble is justified sometimes and also suggest measures to counter that. The option implicates the country's environment as being responsible for that bubble, hence the blame is shifted completely. Also, the option fails to mention the antipathy and the measures suggested to counter the bubble.
Option C: This option is distorted. Where the paragraph says that the aloof attitude is justified sometimes, the option blames the peacekeeping forces and their 'extravagant lifestyles' for the antipathy they face. Hence, can be eliminated.
Option D: Option D correctly captures the main points and is the answer.

\end{problem}

\newpage
\begin{problem}{}{}
% Subject: varc
\textbf{Instructions}

Instructions   

**The passage below is accompanied by a set of questions. Choose the best answer to each question.**
We cannot travel outside our neighbourhood without passports. We must wear the same plain clothes. We must exchange our houses every ten years. We cannot avoid labour. We all go to bed at the same time . . . We have religious freedom, but we cannot deny that the soul dies with the body, since ‘but for the fear of punishment, they would have nothing but contempt for the laws and customs of society'. . . . In More’s time, for much of the population, given the plenty and security on offer, such restraints would not have seemed overly unreasonable. For modern readers, however, Utopia appears to rely upon relentless transparency, the repression of variety, and the curtailment of privacy. Utopia provides security: but at what price? In both its external and internal relations, indeed, it seems perilously dystopian.
Such a conclusion might be fortified by examining selectively the tradition which follows More on these points. This often portrays societies where . . . 'it would be almost impossible for man to be depraved, or wicked'. . . . This is achieved both through institutions and mores, which underpin the common life. . . . The passions are regulated and inequalities of wealth and distinction are minimized. Needs, vanity, and emulation are restrained, often by prizing equality and holding riches in contempt. The desire for public power is curbed. Marriage and sexual intercourse are often controlled: in Tommaso Campanella’s The City of the Sun (1623), the first great literary utopia after More’s, relations are forbidden to men before the age of twenty-one and women before nineteen. Communal child-rearing is normal; for Campanella, this commences at age two. Greater simplicity of life, ‘living according to nature’, is often a result: the desire for simplicity and purity are closely related. People become more alike in appearance, opinion, and outlook than they often have been. Unity, order, and homogeneity thus prevail at the cost of individuality and diversity. This model, as J. C. Davis demonstrates, dominated early modern utopianism. . . . And utopian homogeneity remains a familiar theme well into the twentieth century.
Given these considerations, it is not unreasonable to take as our starting point here the hypothesis that utopia and dystopia evidently share more in common than is often supposed. Indeed, they might be twins, the progeny of the same parents. Insofar as this proves to be the case, my linkage of both here will be uncomfortably close for some readers. Yet we should not mistake this argument for the assertion that all utopias are, or tend to produce, dystopias. Those who defend this proposition will find that their association here is not nearly close enough. For we have only to acknowledge the existence of thousands of successful intentional communities in which a cooperative ethos predominates and where harmony without coercion is the rule to set aside such an assertion. Here the individual’s submersion in the group is consensual (though this concept is not unproblematic). It results not in enslavement but voluntary submission to group norms. Harmony is achieved without . . . harming others.

\vspace{1em}
\textbf{Question 21}

**Five jumbled up sentences, related to a topic, are given below. Four of them can be put together to form a coherent paragraph. Identify the odd one out and key in the number of the sentence as your answer:**  
1. The legal status of resources mined in space remains ambiguous; and while the market for asteroid minerals is currently nonexistent, this is likely to change as technical hurdles diminish.  
2. Outer space is a commons, and all of it is open for exploration, however, space law developed in the 1950s and 60s is state-centric and arguably ill-suited to a commercial future.  
3. Laws adopted by the US and Luxembourg are first steps, but they only protect firms from competing claims by their compatriots; a Chinese company will not be bound by US law.  
4. Critics say the US is conferring rights that it has no authority to confer; Russia in particular has condemned this, citing the US’ disrespect for international law.  
5. At issue now is commercial activity, as private firms—rather than nation states — look to space for profit.
Backspace  
789  
456  
123  
0.-  
Clear All  

Submit
[](https://cracku.in/21-five-jumbled-up-sentences-related-to-a-topic-are-g-x-cat-2021-slot-1)

\begin{enumerate}[label=\Alph*.]
\end{enumerate}
\vspace{1em}
\textbf{Solution:}

A brief reading of the sentences suggests that the paragraph is about the inadequacy of laws about commercial activities in space in the wake of rapid technological developments in the same field. All sentences, other than 4, talk about this inadequacy or highlight why the laws are inadequate. Option 4 is out of context here, as it talks about the US disrespecting international law. It does not relate to the inadequacy of space law for commercial activities, and hence, is the answer.

\end{problem}

\newpage
\begin{problem}{}{}
% Subject: varc
\textbf{Instructions}

Instructions   

**The passage below is accompanied by a set of questions. Choose the best answer to each question.**
We cannot travel outside our neighbourhood without passports. We must wear the same plain clothes. We must exchange our houses every ten years. We cannot avoid labour. We all go to bed at the same time . . . We have religious freedom, but we cannot deny that the soul dies with the body, since ‘but for the fear of punishment, they would have nothing but contempt for the laws and customs of society'. . . . In More’s time, for much of the population, given the plenty and security on offer, such restraints would not have seemed overly unreasonable. For modern readers, however, Utopia appears to rely upon relentless transparency, the repression of variety, and the curtailment of privacy. Utopia provides security: but at what price? In both its external and internal relations, indeed, it seems perilously dystopian.
Such a conclusion might be fortified by examining selectively the tradition which follows More on these points. This often portrays societies where . . . 'it would be almost impossible for man to be depraved, or wicked'. . . . This is achieved both through institutions and mores, which underpin the common life. . . . The passions are regulated and inequalities of wealth and distinction are minimized. Needs, vanity, and emulation are restrained, often by prizing equality and holding riches in contempt. The desire for public power is curbed. Marriage and sexual intercourse are often controlled: in Tommaso Campanella’s The City of the Sun (1623), the first great literary utopia after More’s, relations are forbidden to men before the age of twenty-one and women before nineteen. Communal child-rearing is normal; for Campanella, this commences at age two. Greater simplicity of life, ‘living according to nature’, is often a result: the desire for simplicity and purity are closely related. People become more alike in appearance, opinion, and outlook than they often have been. Unity, order, and homogeneity thus prevail at the cost of individuality and diversity. This model, as J. C. Davis demonstrates, dominated early modern utopianism. . . . And utopian homogeneity remains a familiar theme well into the twentieth century.
Given these considerations, it is not unreasonable to take as our starting point here the hypothesis that utopia and dystopia evidently share more in common than is often supposed. Indeed, they might be twins, the progeny of the same parents. Insofar as this proves to be the case, my linkage of both here will be uncomfortably close for some readers. Yet we should not mistake this argument for the assertion that all utopias are, or tend to produce, dystopias. Those who defend this proposition will find that their association here is not nearly close enough. For we have only to acknowledge the existence of thousands of successful intentional communities in which a cooperative ethos predominates and where harmony without coercion is the rule to set aside such an assertion. Here the individual’s submersion in the group is consensual (though this concept is not unproblematic). It results not in enslavement but voluntary submission to group norms. Harmony is achieved without . . . harming others.

\vspace{1em}
\textbf{Question 22}

**The four sentences (labelled 1, 2, 3, 4) below, when properly sequenced would yield a coherent paragraph. Decide on the proper sequencing of the order of the sentences and key in the sequence of the four numbers as your answer:**
1. In the central nervous systems of other animal species, such a comprehensive regeneration of neurons has not yet been proven beyond doubt.  
2. Biologists from the University of Bayreuth have discovered a uniquely rapid form of regeneration in injured neurons and their function in the central nervous system of zebrafish.  
3. They studied the Mauthner cells, which are solely responsible for the escape behaviour of the fish, and previously regarded as incapable of regeneration.  
4. However, their ability to regenerate crucially depends on the location of the injury.
Backspace  
789  
456  
123  
0.-  
Clear All  

Submit
[](https://cracku.in/22-the-four-sentences-labelled-1-2-3-4-below-when-pro-x-cat-2021-slot-1)

\begin{enumerate}[label=\Alph*.]
\end{enumerate}
\vspace{1em}
\textbf{Solution:}

A preliminary read gives us a crude idea about the subject: regeneration of neurons in zebrafish. The discovery is introduced in Statement 2 {'a uniquely rapid form of regeneration in injured neurons and their function in the central nervous system of zebrafish.'} Additional information is presented in statement 3 {what kind of cells were studied? how has the status quo changed?}. Statement 4 presents further clarification concerning the zebrafish {location of the injury is an essential variable} and Statement 1 transitions into an opinion about the regeneration phenomena in other animal species. Hence, arrangement 2341 forms a coherent paragraph.

\end{problem}

\newpage
\begin{problem}{}{}
% Subject: varc
\textbf{Instructions}

Instructions   

**The passage below is accompanied by a set of questions. Choose the best answer to each question.**
We cannot travel outside our neighbourhood without passports. We must wear the same plain clothes. We must exchange our houses every ten years. We cannot avoid labour. We all go to bed at the same time . . . We have religious freedom, but we cannot deny that the soul dies with the body, since ‘but for the fear of punishment, they would have nothing but contempt for the laws and customs of society'. . . . In More’s time, for much of the population, given the plenty and security on offer, such restraints would not have seemed overly unreasonable. For modern readers, however, Utopia appears to rely upon relentless transparency, the repression of variety, and the curtailment of privacy. Utopia provides security: but at what price? In both its external and internal relations, indeed, it seems perilously dystopian.
Such a conclusion might be fortified by examining selectively the tradition which follows More on these points. This often portrays societies where . . . 'it would be almost impossible for man to be depraved, or wicked'. . . . This is achieved both through institutions and mores, which underpin the common life. . . . The passions are regulated and inequalities of wealth and distinction are minimized. Needs, vanity, and emulation are restrained, often by prizing equality and holding riches in contempt. The desire for public power is curbed. Marriage and sexual intercourse are often controlled: in Tommaso Campanella’s The City of the Sun (1623), the first great literary utopia after More’s, relations are forbidden to men before the age of twenty-one and women before nineteen. Communal child-rearing is normal; for Campanella, this commences at age two. Greater simplicity of life, ‘living according to nature’, is often a result: the desire for simplicity and purity are closely related. People become more alike in appearance, opinion, and outlook than they often have been. Unity, order, and homogeneity thus prevail at the cost of individuality and diversity. This model, as J. C. Davis demonstrates, dominated early modern utopianism. . . . And utopian homogeneity remains a familiar theme well into the twentieth century.
Given these considerations, it is not unreasonable to take as our starting point here the hypothesis that utopia and dystopia evidently share more in common than is often supposed. Indeed, they might be twins, the progeny of the same parents. Insofar as this proves to be the case, my linkage of both here will be uncomfortably close for some readers. Yet we should not mistake this argument for the assertion that all utopias are, or tend to produce, dystopias. Those who defend this proposition will find that their association here is not nearly close enough. For we have only to acknowledge the existence of thousands of successful intentional communities in which a cooperative ethos predominates and where harmony without coercion is the rule to set aside such an assertion. Here the individual’s submersion in the group is consensual (though this concept is not unproblematic). It results not in enslavement but voluntary submission to group norms. Harmony is achieved without . . . harming others.

\vspace{1em}
\textbf{Question 23}

**The passage given below is followed by four alternate summaries. Choose the option that best captures the essence of the passage.**  
Developing countries are becoming hotbeds of business innovation in much the same way as Japan did from the 1950s onwards. They are reinventing systems of production and distribution, and experimenting with entirely new business models. Why are countries that were until recently associated with cheap hands now becoming leaders in innovation? Driven by a mixture of ambition and fear they are relentlessly climbing up the value chain. Emerging-market champions have not only proved highly competitive in their own backyards, they are also going global themselves.

\begin{enumerate}[label=\Alph*.]
    \item Competition has driven emerging economies, once suppliers of cheap labour, to become innovators of business models that have enabled them to move up the value chain and go global.
    \item Innovations in production and distribution are helping emerging economies compete with countries to which they once supplied cheap labour.
    \item Developing countries are being forced to invent new business models which challenge the old business models, so they can remain competitive domestically.
    \item Production and distribution models are going through rapid innovations worldwide as developed countries are being challenged by their earlier suppliers from the developing world.
[](https://cracku.in/23-the-passage-given-below-is-followed-by-four-altern-x-cat-2021-slot-1)
\end{enumerate}
\vspace{1em}
\textbf{Solution:}

The main points of the paragraph are:
1. Developing economies are becoming hotbeds of economic innovation.
2. Earlier they used to be associated with cheap labour, but now ambition and fear have made them competitive globally.
Option A: It correctly captures the two main points and hence is the answer. 
Option B: This option is distorted. Business innovations have not been mentioned as the reason why emerging economies have become competitive globally. It has only been mentioned as a factor in close association.
Option C: Again, the paragraph does not mention that the developing economies are being forced to do this in order to stay competitive. This option suggests an element of necessity for the survival of the economies, which is not implied.
Option D: This option is distorted. The passage only mentions innovations in developing economies and not worldwide.

\end{problem}

\newpage
\begin{problem}{}{}
% Subject: varc
\textbf{Instructions}

Instructions   

**The passage below is accompanied by a set of questions. Choose the best answer to each question.**
We cannot travel outside our neighbourhood without passports. We must wear the same plain clothes. We must exchange our houses every ten years. We cannot avoid labour. We all go to bed at the same time . . . We have religious freedom, but we cannot deny that the soul dies with the body, since ‘but for the fear of punishment, they would have nothing but contempt for the laws and customs of society'. . . . In More’s time, for much of the population, given the plenty and security on offer, such restraints would not have seemed overly unreasonable. For modern readers, however, Utopia appears to rely upon relentless transparency, the repression of variety, and the curtailment of privacy. Utopia provides security: but at what price? In both its external and internal relations, indeed, it seems perilously dystopian.
Such a conclusion might be fortified by examining selectively the tradition which follows More on these points. This often portrays societies where . . . 'it would be almost impossible for man to be depraved, or wicked'. . . . This is achieved both through institutions and mores, which underpin the common life. . . . The passions are regulated and inequalities of wealth and distinction are minimized. Needs, vanity, and emulation are restrained, often by prizing equality and holding riches in contempt. The desire for public power is curbed. Marriage and sexual intercourse are often controlled: in Tommaso Campanella’s The City of the Sun (1623), the first great literary utopia after More’s, relations are forbidden to men before the age of twenty-one and women before nineteen. Communal child-rearing is normal; for Campanella, this commences at age two. Greater simplicity of life, ‘living according to nature’, is often a result: the desire for simplicity and purity are closely related. People become more alike in appearance, opinion, and outlook than they often have been. Unity, order, and homogeneity thus prevail at the cost of individuality and diversity. This model, as J. C. Davis demonstrates, dominated early modern utopianism. . . . And utopian homogeneity remains a familiar theme well into the twentieth century.
Given these considerations, it is not unreasonable to take as our starting point here the hypothesis that utopia and dystopia evidently share more in common than is often supposed. Indeed, they might be twins, the progeny of the same parents. Insofar as this proves to be the case, my linkage of both here will be uncomfortably close for some readers. Yet we should not mistake this argument for the assertion that all utopias are, or tend to produce, dystopias. Those who defend this proposition will find that their association here is not nearly close enough. For we have only to acknowledge the existence of thousands of successful intentional communities in which a cooperative ethos predominates and where harmony without coercion is the rule to set aside such an assertion. Here the individual’s submersion in the group is consensual (though this concept is not unproblematic). It results not in enslavement but voluntary submission to group norms. Harmony is achieved without . . . harming others.

\vspace{1em}
\textbf{Question 24}

The four sentences (labelled 1, 2, 3, 4) below, when properly sequenced would yield a coherent paragraph. Decide on the proper sequencing of the order of the sentences and key in the sequence of the four numbers as your answer:  
1. A popular response is the exhortation to plant more trees.  
2. It seems all but certain that global warming will go well above two degrees—quite how high no one knows yet.  
3. Burning them releases it, which is why the scale of forest fires in the Amazon basin last year garnered headlines.  
4. This is because trees sequester carbon by absorbing carbon dioxide.
Backspace  
789  
456  
123  
0.-  
Clear All  

Submit
[](https://cracku.in/24-the-four-sentences-labelled-1-2-3-4-below-when-pro-x-cat-2021-slot-1)

\begin{enumerate}[label=\Alph*.]
\end{enumerate}
\vspace{1em}
\textbf{Solution:}

The given collection of statements appears to correlate trees and global warming. Statements 4 and 3 form a logical block since they emphasise the significance of trees. The author states that trees help in reducing carbon dioxide while burning trees leads to the release of carbon dioxide {'it' in statement 3 refers to carbon dioxide}. Statements 2 and 1 form a pair since 2 introduces a general belief and 1 mentions a 'popular response' to the same. People respond to the news about the increase in temperatures (due to global warming) by suggesting that we should plant more trees. The statement pair 4-3 then justifies this familiar exhortation. Hence, the correct arrangement is 2143. 
[](https://cracku.in/downloads/13893)
  *  
  *

\end{problem}

\newpage
\begin{problem}{}{}
% Subject: varc
\textbf{Instructions}

Instructions   

The passage below is accompanied by a set of questions. Choose the best answer to each question.
It has been said that knowledge, or the problem of knowledge, is the scandal of philosophy. The scandal is philosophy’s apparent inability to show how, when and why we can be sure that we know something or, indeed, that we know anything. Philosopher Michael Williams writes: ‘Is it possible to obtain knowledge at all? This problem is pressing because there are powerful arguments, some very ancient, for the conclusion that it is not . . . Scepticism is the skeleton in Western rationalism’s closet’. While it is not clear that the scandal matters to anyone but philosophers, philosophers point out that it should matter to everyone, at least given a certain conception of knowledge. For, they explain, unless we can ground our claims to knowledge as such, which is to say, distinguish it from mere opinion, superstition, fantasy, wishful thinking, ideology, illusion or delusion, then the actions we take on the basis of presumed knowledge - boarding an airplane, swallowing a pill, finding someone guilty of a crime - will be irrational and unjustifiable.
That is all quite serious-sounding but so also are the rattlings of the skeleton: that is, the sceptic’s contention that we cannot be sure that we know anything - at least not if we think of knowledge as something like having a correct mental representation of reality, and not if we think of reality as something like things-as-they-are-in-themselves, independent of our perceptions, ideas or descriptions. For, the sceptic will note, since reality, under that conception of it, is outside our ken (we cannot catch a glimpse of things-in-themselves around the corner of our own eyes; we cannot form an idea of reality that floats above the processes of our conceiving it), we have no way to compare our mental representations with things-as-they-are-in-themselves and therefore no way to determine whether they are correct or incorrect. Thus the sceptic may repeat (rattling loudly), you cannot be sure you ‘know’ something or anything at all - at least not, he may add (rattling softly before disappearing), if that is the way you conceive ‘knowledge’.
There are a number of ways to handle this situation. The most common is to ignore it. Most people outside the academy - and, indeed, most of us inside it - are unaware of or unperturbed by the philosophical scandal of knowledge and go about our lives without too many epistemic anxieties. We hold our beliefs and presumptive knowledges more or less confidently, usually depending on how we acquired them (I saw it with my own eyes; I heard it on Fox News; a guy at the office told me) and how broadly and strenuously they seem to be shared or endorsed by various relevant people: experts and authorities, friends and family members, colleagues and associates. And we examine our convictions more or less closely, explain them more or less extensively, and defend them more or less vigorously, usually depending on what seems to be at stake for ourselves and/or other people and what resources are available for reassuring ourselves or making our beliefs credible to others (look, it’s right here on the page; add up the figures yourself; I happen to be a heart specialist).

\vspace{1em}
\textbf{Question 1}

The author discusses all of the following arguments in the passage, EXCEPT:

\begin{enumerate}[label=\Alph*.]
    \item sceptics believe that we can never fully know anything, if by “knowing” we mean knowledge of a reality that is independent of the knower.
    \item if we cannot distinguish knowledge from opinion or delusion, we will not be able to justify our actions.
    \item the best way to deal with scepticism about the veracity of knowledge is to ignore it.
    \item philosophers maintain that the scandal of philosophy should be of concern to everyone.
[](https://cracku.in/1-the-author-discusses-all-of-the-following-argument-x-cat-2021-slot-2)
\end{enumerate}
\vspace{1em}
\textbf{Solution:}

The argument in Option A has been discussed in the following excerpt:  

_...the sceptic’s contention that we cannot be sure that we know anything - at least not if we think of knowledge as something like having a correct mental representation of reality, and not if we think of reality as something like things-as-they-are-in-themselves, independent of our perceptions, ideas or descriptions._
The argument in Option B has been discussed in the following excerpt:
__For, they explain, unless we can ground our claims to knowledge as such, which is to say, distinguish it from mere opinion, superstition, fantasy, wishful thinking, ideology, illusion or delusion, then the actions we take on the basis of presumed knowledge - boarding an airplane, swallowing a pill, finding someone guilty of a crime - will be irrational and unjustifiable.__
The argument in Option D has been discussed in the following excerpt:
_While it is not clear that the scandal matters to anyone but philosophers, philosophers point out that it should matter to everyone, at least given a certain conception of knowledge._
The author does says that ignoring the scepticism about the veracity of knowledge is the most common way of dealing with it, not the best way. Hence. Option C has not been discussed.

\end{problem}

\newpage
\begin{problem}{}{}
% Subject: varc
\textbf{Instructions}

Instructions   

The passage below is accompanied by a set of questions. Choose the best answer to each question.
It has been said that knowledge, or the problem of knowledge, is the scandal of philosophy. The scandal is philosophy’s apparent inability to show how, when and why we can be sure that we know something or, indeed, that we know anything. Philosopher Michael Williams writes: ‘Is it possible to obtain knowledge at all? This problem is pressing because there are powerful arguments, some very ancient, for the conclusion that it is not . . . Scepticism is the skeleton in Western rationalism’s closet’. While it is not clear that the scandal matters to anyone but philosophers, philosophers point out that it should matter to everyone, at least given a certain conception of knowledge. For, they explain, unless we can ground our claims to knowledge as such, which is to say, distinguish it from mere opinion, superstition, fantasy, wishful thinking, ideology, illusion or delusion, then the actions we take on the basis of presumed knowledge - boarding an airplane, swallowing a pill, finding someone guilty of a crime - will be irrational and unjustifiable.
That is all quite serious-sounding but so also are the rattlings of the skeleton: that is, the sceptic’s contention that we cannot be sure that we know anything - at least not if we think of knowledge as something like having a correct mental representation of reality, and not if we think of reality as something like things-as-they-are-in-themselves, independent of our perceptions, ideas or descriptions. For, the sceptic will note, since reality, under that conception of it, is outside our ken (we cannot catch a glimpse of things-in-themselves around the corner of our own eyes; we cannot form an idea of reality that floats above the processes of our conceiving it), we have no way to compare our mental representations with things-as-they-are-in-themselves and therefore no way to determine whether they are correct or incorrect. Thus the sceptic may repeat (rattling loudly), you cannot be sure you ‘know’ something or anything at all - at least not, he may add (rattling softly before disappearing), if that is the way you conceive ‘knowledge’.
There are a number of ways to handle this situation. The most common is to ignore it. Most people outside the academy - and, indeed, most of us inside it - are unaware of or unperturbed by the philosophical scandal of knowledge and go about our lives without too many epistemic anxieties. We hold our beliefs and presumptive knowledges more or less confidently, usually depending on how we acquired them (I saw it with my own eyes; I heard it on Fox News; a guy at the office told me) and how broadly and strenuously they seem to be shared or endorsed by various relevant people: experts and authorities, friends and family members, colleagues and associates. And we examine our convictions more or less closely, explain them more or less extensively, and defend them more or less vigorously, usually depending on what seems to be at stake for ourselves and/or other people and what resources are available for reassuring ourselves or making our beliefs credible to others (look, it’s right here on the page; add up the figures yourself; I happen to be a heart specialist).

\vspace{1em}
\textbf{Question 2}

“. . . we cannot catch a glimpse of things-in-themselves around the corner of our own eyes; we cannot form an idea of reality that floats above the processes of our conceiving it . . .” Which one of the following statements best reflects the argument being made in this sentence?

\begin{enumerate}[label=\Alph*.]
    \item Our knowledge of reality floats above our subjective perception of it.
    \item If the reality of things is independent of our eyesight, logically we cannot perceive our perception.
    \item Our knowledge of reality cannot be merged with our process of conceiving it.
    \item If the reality of things is independent of our perception, logically we cannot perceive that reality.
[](https://cracku.in/2-we-cannot-catch-a-glimpse-of-things-in-themselves--x-cat-2021-slot-2)
\end{enumerate}
\vspace{1em}
\textbf{Solution:}

_...and not if we think of reality as something like things-as-they-are-in-themselves,__**independent of our perceptions, ideas or descriptions.**__For, the sceptic will note, since**reality, under that conception of it, is outside our ken** (we cannot catch a glimpse of things-in-themselves around the corner of our own eyes; we cannot form an idea of reality that floats above the processes of our conceiving it), we have no way to compare our mental representations with things-as-they-are-in-themselves and therefore no way to determine whether they are correct or incorrect._
The author is making a logical argument in the sentence given in the question. According to the author, if we say that the reality of objects is independent of our perceptions, then it is out of our ken. Hence, we cannot for this idea of reality, and logically, we would be unable to comprehend it. Option D comes the closest to capturing this point.
Option A is incorrect. The sceptic argues that if reality were independent of individual perceptions, then the reality would float above the processes with which we conceive it. It is being presented as an argument ot negate the viewpoint, while the Option takes it as an established fact.
Perception is more than eyesight. Option B captures only eyesight, and hence, is a distortion.
The process of conceiving reality and our knowledge is not the argument the author presents in the mentioned lines. Hence, Option C is incorrect.

\end{problem}

\newpage
\begin{problem}{}{}
% Subject: varc
\textbf{Instructions}

Instructions   

The passage below is accompanied by a set of questions. Choose the best answer to each question.
It has been said that knowledge, or the problem of knowledge, is the scandal of philosophy. The scandal is philosophy’s apparent inability to show how, when and why we can be sure that we know something or, indeed, that we know anything. Philosopher Michael Williams writes: ‘Is it possible to obtain knowledge at all? This problem is pressing because there are powerful arguments, some very ancient, for the conclusion that it is not . . . Scepticism is the skeleton in Western rationalism’s closet’. While it is not clear that the scandal matters to anyone but philosophers, philosophers point out that it should matter to everyone, at least given a certain conception of knowledge. For, they explain, unless we can ground our claims to knowledge as such, which is to say, distinguish it from mere opinion, superstition, fantasy, wishful thinking, ideology, illusion or delusion, then the actions we take on the basis of presumed knowledge - boarding an airplane, swallowing a pill, finding someone guilty of a crime - will be irrational and unjustifiable.
That is all quite serious-sounding but so also are the rattlings of the skeleton: that is, the sceptic’s contention that we cannot be sure that we know anything - at least not if we think of knowledge as something like having a correct mental representation of reality, and not if we think of reality as something like things-as-they-are-in-themselves, independent of our perceptions, ideas or descriptions. For, the sceptic will note, since reality, under that conception of it, is outside our ken (we cannot catch a glimpse of things-in-themselves around the corner of our own eyes; we cannot form an idea of reality that floats above the processes of our conceiving it), we have no way to compare our mental representations with things-as-they-are-in-themselves and therefore no way to determine whether they are correct or incorrect. Thus the sceptic may repeat (rattling loudly), you cannot be sure you ‘know’ something or anything at all - at least not, he may add (rattling softly before disappearing), if that is the way you conceive ‘knowledge’.
There are a number of ways to handle this situation. The most common is to ignore it. Most people outside the academy - and, indeed, most of us inside it - are unaware of or unperturbed by the philosophical scandal of knowledge and go about our lives without too many epistemic anxieties. We hold our beliefs and presumptive knowledges more or less confidently, usually depending on how we acquired them (I saw it with my own eyes; I heard it on Fox News; a guy at the office told me) and how broadly and strenuously they seem to be shared or endorsed by various relevant people: experts and authorities, friends and family members, colleagues and associates. And we examine our convictions more or less closely, explain them more or less extensively, and defend them more or less vigorously, usually depending on what seems to be at stake for ourselves and/or other people and what resources are available for reassuring ourselves or making our beliefs credible to others (look, it’s right here on the page; add up the figures yourself; I happen to be a heart specialist).

\vspace{1em}
\textbf{Question 3}

According to the last paragraph of the passage, “We hold our beliefs and presumptive knowledges more or less confidently, usually depending on” something. Which one of the following most broadly captures what we depend on?

\begin{enumerate}[label=\Alph*.]
    \item How we come to hold them; how widely they are held in our social circles.
    \item All of the options listed here.
    \item How much of a stake we have in them; what resources there are to support them.
    \item Remaining outside the academy; ignoring epistemic anxieties.
[](https://cracku.in/3-according-to-the-last-paragraph-of-the-passage-we--x-cat-2021-slot-2)
\end{enumerate}
\vspace{1em}
\textbf{Solution:}

_We hold our beliefs and presumptive knowledges more or less confidently, usually**depending on how we acquired them** (I saw it with my own eyes; I heard it on Fox News; a guy at the office told me) **and how broadly and strenuously they seem to be shared or endorsed by various relevant people** : experts and authorities, friends and family members, colleagues and associates._
From the above excerpt, it is clear that held beliefs and presumptive knowledges depend upon how we acquired them, and how strongly they are shared or endorsed by relevant people in our circle. Hence, Option A is the answer.

\end{problem}

\newpage
\begin{problem}{}{}
% Subject: varc
\textbf{Instructions}

Instructions   

The passage below is accompanied by a set of questions. Choose the best answer to each question.
It has been said that knowledge, or the problem of knowledge, is the scandal of philosophy. The scandal is philosophy’s apparent inability to show how, when and why we can be sure that we know something or, indeed, that we know anything. Philosopher Michael Williams writes: ‘Is it possible to obtain knowledge at all? This problem is pressing because there are powerful arguments, some very ancient, for the conclusion that it is not . . . Scepticism is the skeleton in Western rationalism’s closet’. While it is not clear that the scandal matters to anyone but philosophers, philosophers point out that it should matter to everyone, at least given a certain conception of knowledge. For, they explain, unless we can ground our claims to knowledge as such, which is to say, distinguish it from mere opinion, superstition, fantasy, wishful thinking, ideology, illusion or delusion, then the actions we take on the basis of presumed knowledge - boarding an airplane, swallowing a pill, finding someone guilty of a crime - will be irrational and unjustifiable.
That is all quite serious-sounding but so also are the rattlings of the skeleton: that is, the sceptic’s contention that we cannot be sure that we know anything - at least not if we think of knowledge as something like having a correct mental representation of reality, and not if we think of reality as something like things-as-they-are-in-themselves, independent of our perceptions, ideas or descriptions. For, the sceptic will note, since reality, under that conception of it, is outside our ken (we cannot catch a glimpse of things-in-themselves around the corner of our own eyes; we cannot form an idea of reality that floats above the processes of our conceiving it), we have no way to compare our mental representations with things-as-they-are-in-themselves and therefore no way to determine whether they are correct or incorrect. Thus the sceptic may repeat (rattling loudly), you cannot be sure you ‘know’ something or anything at all - at least not, he may add (rattling softly before disappearing), if that is the way you conceive ‘knowledge’.
There are a number of ways to handle this situation. The most common is to ignore it. Most people outside the academy - and, indeed, most of us inside it - are unaware of or unperturbed by the philosophical scandal of knowledge and go about our lives without too many epistemic anxieties. We hold our beliefs and presumptive knowledges more or less confidently, usually depending on how we acquired them (I saw it with my own eyes; I heard it on Fox News; a guy at the office told me) and how broadly and strenuously they seem to be shared or endorsed by various relevant people: experts and authorities, friends and family members, colleagues and associates. And we examine our convictions more or less closely, explain them more or less extensively, and defend them more or less vigorously, usually depending on what seems to be at stake for ourselves and/or other people and what resources are available for reassuring ourselves or making our beliefs credible to others (look, it’s right here on the page; add up the figures yourself; I happen to be a heart specialist).

\vspace{1em}
\textbf{Question 4}

The author of the passage is most likely to support which one of the following statements?

\begin{enumerate}[label=\Alph*.]
    \item The confidence with which we maintain something to be true is usually independent of the source of the alleged truth.
    \item The scandal of philosophy is that we might not know anything at all about reality if we think of reality as independent of our perceptions, ideas or descriptions.
    \item The actions taken on the basis of presumed knowledge are rational and justifiable if we are confident that that knowledge is widely held.
    \item For the sceptic, if we think of reality as independent of our perceptions, ideas or descriptions, we should aim to know that reality independently too.
[](https://cracku.in/4-the-author-of-the-passage-is-most-likely-to-suppor-x-cat-2021-slot-2)
\end{enumerate}
\vspace{1em}
\textbf{Solution:}

_We hold our beliefs and presumptive knowledges more or less confidently, usually depending on how we acquired them..._
As mentioned in the above line, Option A directly contradicts what the author says.
_It has been said that knowledge, or the problem of knowledge, is the scandal of philosophy. The scandal is philosophy’s apparent inability to show how, when and why we can be sure that we know something or, indeed, that we know anything._
The author then goes on to explain that sceptic view is the skeleton in western philosophy's closet when trying to negate this scandal. Thus, it means that the scandal has to do with the sceptic way of thinking.  

In the next paragraph, the sceptic view has been explained, which talks about our inability to grasp reality if we think of it as independent of our perceptions. Hence, the scandal can be construed to be the same. The author is likely to agree with this view. Option B is the answer.
In the last paragraph, the author points out that we defend a viewpoint strongly if we feel that it is held widely in our social circle. However, the author does not allude to the fact that it is appropriate or not. Also, such an argument would bolster the inherent bias we have, and hence would point towards a fault in our decision making. Hence, Option C is incorrect.
Option D is not supported in the passage. The passage presents the sceptic view that if reality were construed as independent of our perceptions, then it would be impossible to grasp reality. It has not been mentioned that in this case, we should aim to study that reality in a similar manner.

Instructions   

The passage below is accompanied by a set of questions. Choose the best answer to each question.
It’s easy to forget that most of the world’s languages are still transmitted orally with no widely established written form. While speech communities are increasingly involved in projects to protect their languages - in print, on air and online - orality is fragile and contributes to linguistic vulnerability. But indigenous languages are about much more than unusual words and intriguing grammar: They function as vehicles for the transmission of cultural traditions, environmental understandings and knowledge about medicinal plants, all at risk when elders die and livelihoods are disrupted.
Both push and pull factors lead to the decline of languages. Through war, famine and natural disasters, whole communities can be destroyed, taking their language with them to the grave, such as the indigenous populations of Tasmania who were wiped out by colonists. More commonly, speakers live on but abandon their language in favor of another vernacular, a widespread process that linguists refer to as “language shift” from which few languages are immune. Such trading up and out of a speech form occurs for complex political, cultural and economic reasons - sometimes voluntary for economic and educational reasons, although often amplified by state coercion or neglect. Welsh, long stigmatized and disparaged by the British state, has rebounded with vigor.
Many speakers of endangered, poorly documented languages have embraced new digital media with excitement. Speakers of previously exclusively oral tongues are turning to the web as a virtual space for languages to live on. Internet technology offers powerful ways for oral traditions and cultural practices to survive, even thrive, among increasingly mobile communities. I have watched as videos of traditional wedding ceremonies and songs are recorded on smartphones in London by Nepali migrants, then uploaded to YouTube and watched an hour later by relatives in remote Himalayan villages . . .
Globalization is regularly, and often uncritically, pilloried as a major threat to linguistic diversity. But in fact, globalization is as much process as it is ideology, certainly when it comes to language. The real forces behind cultural homogenization are unbending beliefs, exchanged through a globalized delivery system, reinforced by the historical monolingualism prevalent in much of the West.
Monolingualism - the condition of being able to speak only one language - is regularly accompanied by a deep-seated conviction in the value of that language over all others. Across the largest economies that make up the G8, being monolingual is still often the norm, with multilingualism appearing unusual and even somewhat exotic. The monolingual mindset stands in sharp contrast to the lived reality of most the world, which throughout its history has been more multilingual than unilingual. Monolingualism, then, not globalization, should be our primary concern.
Multilingualism can help us live in a more connected and more interdependent world. By widening access to technology, globalization can support indigenous and scholarly communities engaged in documenting and protecting our shared linguistic heritage. For the last 5,000 years, the rise and fall of languages was intimately tied to the plow, sword and book. In our digital age, the keyboard, screen and web will play a decisive role in shaping the future linguistic diversity of our species.

\end{problem}

\newpage
\begin{problem}{}{}
% Subject: varc
\textbf{Instructions}

Instructions   

The passage below is accompanied by a set of questions. Choose the best answer to each question.
It has been said that knowledge, or the problem of knowledge, is the scandal of philosophy. The scandal is philosophy’s apparent inability to show how, when and why we can be sure that we know something or, indeed, that we know anything. Philosopher Michael Williams writes: ‘Is it possible to obtain knowledge at all? This problem is pressing because there are powerful arguments, some very ancient, for the conclusion that it is not . . . Scepticism is the skeleton in Western rationalism’s closet’. While it is not clear that the scandal matters to anyone but philosophers, philosophers point out that it should matter to everyone, at least given a certain conception of knowledge. For, they explain, unless we can ground our claims to knowledge as such, which is to say, distinguish it from mere opinion, superstition, fantasy, wishful thinking, ideology, illusion or delusion, then the actions we take on the basis of presumed knowledge - boarding an airplane, swallowing a pill, finding someone guilty of a crime - will be irrational and unjustifiable.
That is all quite serious-sounding but so also are the rattlings of the skeleton: that is, the sceptic’s contention that we cannot be sure that we know anything - at least not if we think of knowledge as something like having a correct mental representation of reality, and not if we think of reality as something like things-as-they-are-in-themselves, independent of our perceptions, ideas or descriptions. For, the sceptic will note, since reality, under that conception of it, is outside our ken (we cannot catch a glimpse of things-in-themselves around the corner of our own eyes; we cannot form an idea of reality that floats above the processes of our conceiving it), we have no way to compare our mental representations with things-as-they-are-in-themselves and therefore no way to determine whether they are correct or incorrect. Thus the sceptic may repeat (rattling loudly), you cannot be sure you ‘know’ something or anything at all - at least not, he may add (rattling softly before disappearing), if that is the way you conceive ‘knowledge’.
There are a number of ways to handle this situation. The most common is to ignore it. Most people outside the academy - and, indeed, most of us inside it - are unaware of or unperturbed by the philosophical scandal of knowledge and go about our lives without too many epistemic anxieties. We hold our beliefs and presumptive knowledges more or less confidently, usually depending on how we acquired them (I saw it with my own eyes; I heard it on Fox News; a guy at the office told me) and how broadly and strenuously they seem to be shared or endorsed by various relevant people: experts and authorities, friends and family members, colleagues and associates. And we examine our convictions more or less closely, explain them more or less extensively, and defend them more or less vigorously, usually depending on what seems to be at stake for ourselves and/or other people and what resources are available for reassuring ourselves or making our beliefs credible to others (look, it’s right here on the page; add up the figures yourself; I happen to be a heart specialist).

\vspace{1em}
\textbf{Question 5}

The author lists all of the following as reasons for the decline or disappearance of a language EXCEPT:

\begin{enumerate}[label=\Alph*.]
    \item a catastrophic event that entirely eliminates a people and their culture.
    \item governments promoting certain languages over others.
    \item the focus on only a few languages as a result of widespread internet use.
    \item people shifting away from their own language to study or work in another language.
[](https://cracku.in/5-the-author-lists-all-of-the-following-as-reasons-f-x-cat-2021-slot-2)
\end{enumerate}
\vspace{1em}
\textbf{Solution:}

_Through war, famine and natural disasters, whole communities can be destroyed, taking their language with them to the grave..._
_Such trading up and out of a speech form occurs for complex political, cultural and economic reasons - sometimes voluntary for economic and educational reasons, although often amplified by_ __state coercion__ _or neglect._
_More commonly, speakers live on but abandon their language in favor of another vernacular, a widespread process that linguists refer to as “language shift” from which few languages  
_
From the above excerpts, Options A, B, and D can be supported. The passage supports that Internet technology has allowed certain endangered languages to thrive. Option C is not supported anywhere in the passage, hence, is the answer.

\end{problem}

\newpage
\begin{problem}{}{}
% Subject: varc
\textbf{Instructions}

Instructions   

The passage below is accompanied by a set of questions. Choose the best answer to each question.
It has been said that knowledge, or the problem of knowledge, is the scandal of philosophy. The scandal is philosophy’s apparent inability to show how, when and why we can be sure that we know something or, indeed, that we know anything. Philosopher Michael Williams writes: ‘Is it possible to obtain knowledge at all? This problem is pressing because there are powerful arguments, some very ancient, for the conclusion that it is not . . . Scepticism is the skeleton in Western rationalism’s closet’. While it is not clear that the scandal matters to anyone but philosophers, philosophers point out that it should matter to everyone, at least given a certain conception of knowledge. For, they explain, unless we can ground our claims to knowledge as such, which is to say, distinguish it from mere opinion, superstition, fantasy, wishful thinking, ideology, illusion or delusion, then the actions we take on the basis of presumed knowledge - boarding an airplane, swallowing a pill, finding someone guilty of a crime - will be irrational and unjustifiable.
That is all quite serious-sounding but so also are the rattlings of the skeleton: that is, the sceptic’s contention that we cannot be sure that we know anything - at least not if we think of knowledge as something like having a correct mental representation of reality, and not if we think of reality as something like things-as-they-are-in-themselves, independent of our perceptions, ideas or descriptions. For, the sceptic will note, since reality, under that conception of it, is outside our ken (we cannot catch a glimpse of things-in-themselves around the corner of our own eyes; we cannot form an idea of reality that floats above the processes of our conceiving it), we have no way to compare our mental representations with things-as-they-are-in-themselves and therefore no way to determine whether they are correct or incorrect. Thus the sceptic may repeat (rattling loudly), you cannot be sure you ‘know’ something or anything at all - at least not, he may add (rattling softly before disappearing), if that is the way you conceive ‘knowledge’.
There are a number of ways to handle this situation. The most common is to ignore it. Most people outside the academy - and, indeed, most of us inside it - are unaware of or unperturbed by the philosophical scandal of knowledge and go about our lives without too many epistemic anxieties. We hold our beliefs and presumptive knowledges more or less confidently, usually depending on how we acquired them (I saw it with my own eyes; I heard it on Fox News; a guy at the office told me) and how broadly and strenuously they seem to be shared or endorsed by various relevant people: experts and authorities, friends and family members, colleagues and associates. And we examine our convictions more or less closely, explain them more or less extensively, and defend them more or less vigorously, usually depending on what seems to be at stake for ourselves and/or other people and what resources are available for reassuring ourselves or making our beliefs credible to others (look, it’s right here on the page; add up the figures yourself; I happen to be a heart specialist).

\vspace{1em}
\textbf{Question 6}

We can infer all of the following about indigenous languages from the passage EXCEPT that:

\begin{enumerate}[label=\Alph*.]
    \item they are repositories of traditional knowledge about the environment and culture.
    \item they are in danger of being wiped out as most can only be transmitted orally.
    \item people are increasingly working on documenting these languages.
    \item their vocabulary and grammatical constructs have been challenging to document.
[](https://cracku.in/6-we-can-infer-all-of-the-following-about-indigenous-x-cat-2021-slot-2)
\end{enumerate}
\vspace{1em}
\textbf{Solution:}

_It’s easy to forget that _most of the world’s languages are still transmitted orally_ with no widely established written form. While speech communities are increasingly involved in projects _to protect their languages - in print,_ on air and online - _orality is fragile and contributes to linguistic vulnerability._ But indigenous languages are about much more than unusual words and intriguing grammar:_They function as vehicles for the transmission of cultural traditions, environmental understandings and knowledge about medicinal plants_ , all at risk when elders die and livelihoods are disrupted._
From the underlined portions of the above excerpt, we can infer options A, B, and C.
The first paragraph does mention that indigenous languages have unusual words and intriguing grammar. But it has not been mentioned as a factor that makes their documentation challenging. Hence, Option D is the answer.

\end{problem}

\newpage
\begin{problem}{}{}
% Subject: varc
\textbf{Instructions}

Instructions   

The passage below is accompanied by a set of questions. Choose the best answer to each question.
It has been said that knowledge, or the problem of knowledge, is the scandal of philosophy. The scandal is philosophy’s apparent inability to show how, when and why we can be sure that we know something or, indeed, that we know anything. Philosopher Michael Williams writes: ‘Is it possible to obtain knowledge at all? This problem is pressing because there are powerful arguments, some very ancient, for the conclusion that it is not . . . Scepticism is the skeleton in Western rationalism’s closet’. While it is not clear that the scandal matters to anyone but philosophers, philosophers point out that it should matter to everyone, at least given a certain conception of knowledge. For, they explain, unless we can ground our claims to knowledge as such, which is to say, distinguish it from mere opinion, superstition, fantasy, wishful thinking, ideology, illusion or delusion, then the actions we take on the basis of presumed knowledge - boarding an airplane, swallowing a pill, finding someone guilty of a crime - will be irrational and unjustifiable.
That is all quite serious-sounding but so also are the rattlings of the skeleton: that is, the sceptic’s contention that we cannot be sure that we know anything - at least not if we think of knowledge as something like having a correct mental representation of reality, and not if we think of reality as something like things-as-they-are-in-themselves, independent of our perceptions, ideas or descriptions. For, the sceptic will note, since reality, under that conception of it, is outside our ken (we cannot catch a glimpse of things-in-themselves around the corner of our own eyes; we cannot form an idea of reality that floats above the processes of our conceiving it), we have no way to compare our mental representations with things-as-they-are-in-themselves and therefore no way to determine whether they are correct or incorrect. Thus the sceptic may repeat (rattling loudly), you cannot be sure you ‘know’ something or anything at all - at least not, he may add (rattling softly before disappearing), if that is the way you conceive ‘knowledge’.
There are a number of ways to handle this situation. The most common is to ignore it. Most people outside the academy - and, indeed, most of us inside it - are unaware of or unperturbed by the philosophical scandal of knowledge and go about our lives without too many epistemic anxieties. We hold our beliefs and presumptive knowledges more or less confidently, usually depending on how we acquired them (I saw it with my own eyes; I heard it on Fox News; a guy at the office told me) and how broadly and strenuously they seem to be shared or endorsed by various relevant people: experts and authorities, friends and family members, colleagues and associates. And we examine our convictions more or less closely, explain them more or less extensively, and defend them more or less vigorously, usually depending on what seems to be at stake for ourselves and/or other people and what resources are available for reassuring ourselves or making our beliefs credible to others (look, it’s right here on the page; add up the figures yourself; I happen to be a heart specialist).

\vspace{1em}
\textbf{Question 7}

From the passage, we can infer that the author is in favour of:

\begin{enumerate}[label=\Alph*.]
    \item greater multilingualism.
    \item “language shifts” across languages.
    \item cultural homogenisation.
    \item an expanded state role in the preservation of languages.
[](https://cracku.in/7-from-the-passage-we-can-infer-that-the-author-is-i-x-cat-2021-slot-2)
\end{enumerate}
\vspace{1em}
\textbf{Solution:}

_Multilingualism can help us live in a more connected and more interdependent world._
Throughout the passage, the author is advocating for greater language diversity and suggesting how to counter the threat to the same. Multilingualism is something that the author has supported as cited above. Hence, Option A is the answer.
_More commonly, speakers live on but abandon their language in favor of another vernacular, a widespread process that linguists refer to as “language shift” from which few languages are immune._
Since the author advocated the preservation of languages, he would likely be against this shift, as it endangers languages. Option B can be eliminated.
Again, the author is pushing for more diversity and preservation of cultures instead of homogenization of the same. Option C can be eliminated.
The author does not push for state intervention in the preservation of languages. Hence, Option D would not be the answer.

\end{problem}

\newpage
\begin{problem}{}{}
% Subject: varc
\textbf{Instructions}

Instructions   

The passage below is accompanied by a set of questions. Choose the best answer to each question.
It has been said that knowledge, or the problem of knowledge, is the scandal of philosophy. The scandal is philosophy’s apparent inability to show how, when and why we can be sure that we know something or, indeed, that we know anything. Philosopher Michael Williams writes: ‘Is it possible to obtain knowledge at all? This problem is pressing because there are powerful arguments, some very ancient, for the conclusion that it is not . . . Scepticism is the skeleton in Western rationalism’s closet’. While it is not clear that the scandal matters to anyone but philosophers, philosophers point out that it should matter to everyone, at least given a certain conception of knowledge. For, they explain, unless we can ground our claims to knowledge as such, which is to say, distinguish it from mere opinion, superstition, fantasy, wishful thinking, ideology, illusion or delusion, then the actions we take on the basis of presumed knowledge - boarding an airplane, swallowing a pill, finding someone guilty of a crime - will be irrational and unjustifiable.
That is all quite serious-sounding but so also are the rattlings of the skeleton: that is, the sceptic’s contention that we cannot be sure that we know anything - at least not if we think of knowledge as something like having a correct mental representation of reality, and not if we think of reality as something like things-as-they-are-in-themselves, independent of our perceptions, ideas or descriptions. For, the sceptic will note, since reality, under that conception of it, is outside our ken (we cannot catch a glimpse of things-in-themselves around the corner of our own eyes; we cannot form an idea of reality that floats above the processes of our conceiving it), we have no way to compare our mental representations with things-as-they-are-in-themselves and therefore no way to determine whether they are correct or incorrect. Thus the sceptic may repeat (rattling loudly), you cannot be sure you ‘know’ something or anything at all - at least not, he may add (rattling softly before disappearing), if that is the way you conceive ‘knowledge’.
There are a number of ways to handle this situation. The most common is to ignore it. Most people outside the academy - and, indeed, most of us inside it - are unaware of or unperturbed by the philosophical scandal of knowledge and go about our lives without too many epistemic anxieties. We hold our beliefs and presumptive knowledges more or less confidently, usually depending on how we acquired them (I saw it with my own eyes; I heard it on Fox News; a guy at the office told me) and how broadly and strenuously they seem to be shared or endorsed by various relevant people: experts and authorities, friends and family members, colleagues and associates. And we examine our convictions more or less closely, explain them more or less extensively, and defend them more or less vigorously, usually depending on what seems to be at stake for ourselves and/or other people and what resources are available for reassuring ourselves or making our beliefs credible to others (look, it’s right here on the page; add up the figures yourself; I happen to be a heart specialist).

\vspace{1em}
\textbf{Question 8}

The author mentions the Welsh language to show that:

\begin{enumerate}[label=\Alph*.]
    \item efforts to integrate Welsh speakers in the English-speaking fold have been fruitless.
    \item languages can revive even after their speakers have gone through a “language shift”.
    \item vulnerable languages can rebound with state effort.
    \item while often pilloried, globalisation can, in fact, support linguistic revival.
[](https://cracku.in/8-the-author-mentions-the-welsh-language-to-show-tha-x-cat-2021-slot-2)
\end{enumerate}
\vspace{1em}
\textbf{Solution:}

_More commonly, speakers live on but abandon their language in favor of another vernacular, a widespread process that linguists refer to as “language shift” from which few languages are immune. Such trading up and out of a speech form occurs for complex political, cultural and economic reasons - sometimes voluntary for economic and educational reasons, although often amplified by state coercion or neglect. Welsh, long stigmatized and disparaged by the British state, has rebounded with vigor._
In the above excerpt, the trading of language for another (language shift) has been mentioned. The author mentions the reasons why this happens, and then the Welsh language is mentioned as an example which has rebounded against the same with vigour. Thus, it is presented as a ray of hope, that a language can be revived even when cultural shift occurs. Option B is the answer.
The example has not been mentioned to spite the efforts that were put to integrate Welsh speakers into English speaking fold. The purpose of the author is to deal with endangered languages, and the subject in Option A would be out of scope.
The role of state effort in revitalising Welsh has not been mentioned, hence, Option C can be eliminated.
The role of globalisation in revitalising Welsh has not been mentioned either. Hence, Option D can be eliminated too.

Instructions   

The passage below is accompanied by a set of questions. Choose the best answer to each question.
I have elaborated . . . a framework for analyzing the contradictory pulls on [Indian] nationalist ideology in its struggle against the dominance of colonialism and the resolution it offered to those contradictions. Briefly, this resolution was built around a separation of the domain of culture into two spheres—the material and the spiritual. It was in the material sphere that the claims of Western civilization were the most powerful. Science, technology, rational forms of economic organization, modern methods of statecraft—these had given the European countries the strength to subjugate the non-European people . . . To overcome this domination, the colonized people had to learn those superior techniques of organizing material life and incorporate them within their own cultures. . . . But this could not mean the imitation of the West in every aspect of life, for then the very distinction between the West and the East would vanish—the self-identity of national culture would itself be threatened. . . .
The discourse of nationalism shows that the material/spiritual distinction was condensed into an analogous, but ideologically far more powerful, dichotomy: that between the outer and the inner. . . . Applying the inner/outer distinction to the matter of concrete day-to-day living separates the social space into ghar and bāhir, the home and the world. The world is the external, the domain of the material; the home represents one’s inner spiritual self, one’s true identity. The world is a treacherous terrain of the pursuit of material interests, where practical considerations reign supreme. It is also typically the domain of the male. The home in its essence must remain unaffected by the profane activities of the material world—and woman is its representation. And so one gets an identification of social roles by gender to correspond with the separation of the social space into ghar and bāhir. . . .
The colonial situation, and the ideological response of nationalism to the critique of Indian tradition, introduced an entirely new substance to [these dichotomies] and effected their transformation. The material/spiritual dichotomy, to which the terms world and home corresponded, had acquired . . . a very special significance in the nationalist mind. The world was where the European power had challenged the non-European people and, by virtue of its superior material culture, had subjugated them. But, the nationalists asserted, it had failed to colonize the inner, essential, identity of the East which lay in its distinctive, and superior, spiritual culture. . . . [I]n the entire phase of the national struggle, the crucial need was to protect, preserve and strengthen the inner core of the national culture, its spiritual essence. . . .
Once we match this new meaning of the home/world dichotomy with the identification of social roles by gender, we get the ideological framework within which nationalism answered the women’s question. It would be a grave error to see in this, as liberals are apt to in their despair at the many marks of social conservatism in nationalist practice, a total rejection of the West. Quite the contrary: the nationalist paradigm in fact supplied an ideological principle of selection.

\end{problem}

\newpage
\begin{problem}{}{}
% Subject: varc
\textbf{Instructions}

Instructions   

The passage below is accompanied by a set of questions. Choose the best answer to each question.
It has been said that knowledge, or the problem of knowledge, is the scandal of philosophy. The scandal is philosophy’s apparent inability to show how, when and why we can be sure that we know something or, indeed, that we know anything. Philosopher Michael Williams writes: ‘Is it possible to obtain knowledge at all? This problem is pressing because there are powerful arguments, some very ancient, for the conclusion that it is not . . . Scepticism is the skeleton in Western rationalism’s closet’. While it is not clear that the scandal matters to anyone but philosophers, philosophers point out that it should matter to everyone, at least given a certain conception of knowledge. For, they explain, unless we can ground our claims to knowledge as such, which is to say, distinguish it from mere opinion, superstition, fantasy, wishful thinking, ideology, illusion or delusion, then the actions we take on the basis of presumed knowledge - boarding an airplane, swallowing a pill, finding someone guilty of a crime - will be irrational and unjustifiable.
That is all quite serious-sounding but so also are the rattlings of the skeleton: that is, the sceptic’s contention that we cannot be sure that we know anything - at least not if we think of knowledge as something like having a correct mental representation of reality, and not if we think of reality as something like things-as-they-are-in-themselves, independent of our perceptions, ideas or descriptions. For, the sceptic will note, since reality, under that conception of it, is outside our ken (we cannot catch a glimpse of things-in-themselves around the corner of our own eyes; we cannot form an idea of reality that floats above the processes of our conceiving it), we have no way to compare our mental representations with things-as-they-are-in-themselves and therefore no way to determine whether they are correct or incorrect. Thus the sceptic may repeat (rattling loudly), you cannot be sure you ‘know’ something or anything at all - at least not, he may add (rattling softly before disappearing), if that is the way you conceive ‘knowledge’.
There are a number of ways to handle this situation. The most common is to ignore it. Most people outside the academy - and, indeed, most of us inside it - are unaware of or unperturbed by the philosophical scandal of knowledge and go about our lives without too many epistemic anxieties. We hold our beliefs and presumptive knowledges more or less confidently, usually depending on how we acquired them (I saw it with my own eyes; I heard it on Fox News; a guy at the office told me) and how broadly and strenuously they seem to be shared or endorsed by various relevant people: experts and authorities, friends and family members, colleagues and associates. And we examine our convictions more or less closely, explain them more or less extensively, and defend them more or less vigorously, usually depending on what seems to be at stake for ourselves and/or other people and what resources are available for reassuring ourselves or making our beliefs credible to others (look, it’s right here on the page; add up the figures yourself; I happen to be a heart specialist).

\vspace{1em}
\textbf{Question 9}

Which one of the following explains the “contradictory pulls” on Indian nationalism?

\begin{enumerate}[label=\Alph*.]
    \item Despite its spiritual superiority, Indian nationalism had to fight against colonial domination.
    \item Despite its fight against colonial domination, Indian nationalism had to borrow from the coloniser in the spiritual sphere.
    \item Despite its scientific and technological inferiority, Indian nationalism had to fight against colonial domination.
    \item Despite its fight against colonial domination, Indian nationalism had to borrow from the coloniser in the material sphere.
[](https://cracku.in/9-which-one-of-the-following-explains-the-contradict-x-cat-2021-slot-2)
\end{enumerate}
\vspace{1em}
\textbf{Solution:}

_I have elaborated . . . a framework for analyzing the contradictory pulls on [Indian] nationalist ideology in its _struggle against the dominance of colonialism_ and the resolution it offered to those contradictions. Briefly, _this resolution was built around a separation of the domain of culture into two spheres—the material and the spiritual._ It was in the material sphere that the claims of Western civilization were the most powerful. Science, technology, rational forms of economic organization, modern methods of statecraft—these had given the European countries the strength to subjugate the non-European people . . . _To overcome this domination, the colonized people had to learn those superior techniques of organizing material life and incorporate them within their own cultures_. . . . But this could not mean the imitation of the West in every aspect of life, for then the very distinction between the West and the East would vanish—the self-identity of national culture would itself be threatened. . . ._
The first paragraph acknowledges that the nationalist ideology was fighting against colonial dominance, and there were certain inherent contradictions in the way this struggle was being carried out. The author says that a method of resolution of these contradictions was to separate material and spiritual domains. This hints at the contradiction present. We can infer from here that this is being done because the nationalists acknowledge that the colonial countries were superior in certain aspects, which allowed them to subjugate non-Europeans, as is also mentioned later in the paragraph. The author also mentions that the colonized people had to learn those superior techniques, instead of all-out rejection of colonialist ideas and the progress they brought with them. Thus, the contradiction was that to overcome colonial dominance, nationalism had to accept that the material ways of the West were superior and incorporate them. Hence, Option D is the answer.

\end{problem}

\newpage
\begin{problem}{}{}
% Subject: varc
\textbf{Instructions}

Instructions   

The passage below is accompanied by a set of questions. Choose the best answer to each question.
It has been said that knowledge, or the problem of knowledge, is the scandal of philosophy. The scandal is philosophy’s apparent inability to show how, when and why we can be sure that we know something or, indeed, that we know anything. Philosopher Michael Williams writes: ‘Is it possible to obtain knowledge at all? This problem is pressing because there are powerful arguments, some very ancient, for the conclusion that it is not . . . Scepticism is the skeleton in Western rationalism’s closet’. While it is not clear that the scandal matters to anyone but philosophers, philosophers point out that it should matter to everyone, at least given a certain conception of knowledge. For, they explain, unless we can ground our claims to knowledge as such, which is to say, distinguish it from mere opinion, superstition, fantasy, wishful thinking, ideology, illusion or delusion, then the actions we take on the basis of presumed knowledge - boarding an airplane, swallowing a pill, finding someone guilty of a crime - will be irrational and unjustifiable.
That is all quite serious-sounding but so also are the rattlings of the skeleton: that is, the sceptic’s contention that we cannot be sure that we know anything - at least not if we think of knowledge as something like having a correct mental representation of reality, and not if we think of reality as something like things-as-they-are-in-themselves, independent of our perceptions, ideas or descriptions. For, the sceptic will note, since reality, under that conception of it, is outside our ken (we cannot catch a glimpse of things-in-themselves around the corner of our own eyes; we cannot form an idea of reality that floats above the processes of our conceiving it), we have no way to compare our mental representations with things-as-they-are-in-themselves and therefore no way to determine whether they are correct or incorrect. Thus the sceptic may repeat (rattling loudly), you cannot be sure you ‘know’ something or anything at all - at least not, he may add (rattling softly before disappearing), if that is the way you conceive ‘knowledge’.
There are a number of ways to handle this situation. The most common is to ignore it. Most people outside the academy - and, indeed, most of us inside it - are unaware of or unperturbed by the philosophical scandal of knowledge and go about our lives without too many epistemic anxieties. We hold our beliefs and presumptive knowledges more or less confidently, usually depending on how we acquired them (I saw it with my own eyes; I heard it on Fox News; a guy at the office told me) and how broadly and strenuously they seem to be shared or endorsed by various relevant people: experts and authorities, friends and family members, colleagues and associates. And we examine our convictions more or less closely, explain them more or less extensively, and defend them more or less vigorously, usually depending on what seems to be at stake for ourselves and/or other people and what resources are available for reassuring ourselves or making our beliefs credible to others (look, it’s right here on the page; add up the figures yourself; I happen to be a heart specialist).

\vspace{1em}
\textbf{Question 10}

On the basis of the information in the passage, all of the following are true about the spiritual/material dichotomy of Indian nationalism EXCEPT that it:

\begin{enumerate}[label=\Alph*.]
    \item represented a continuation of age-old oppositions in Indian culture.
    \item constituted the premise of the ghar/bāhir dichotomy.
    \item was not as ideologically powerful as the inner/outer dichotomy.
    \item helped in safeguarding the identity of Indian nationalism.
[](https://cracku.in/10-on-the-basis-of-the-information-in-the-passage-all-x-cat-2021-slot-2)
\end{enumerate}
\vspace{1em}
\textbf{Solution:}

_The discourse of nationalism shows that the material/spiritual distinction was condensed into an analogous, but ideologically far more powerful, dichotomy: that between the outer and the inner....  
_
The above excerpt shows that the material/spiritual distinction was condensed to form a far more superior dichotomy of the outer and the inner. Thus, the former was the premise for the latter, as well as inferior to the latter. Hence, Options B and C are true.
_To overcome this domination, the colonized people had to learn those superior techniques of organizing material life and incorporate them within their own cultures. . . . But this could not mean the imitation of the West in every aspect of life, for then the very distinction between the West and the East would vanish—the self-identity of national culture would itself be threatened. . ._
From the above excerpt, we can infer that the dichotomy helped save the identity of Indian Nationalism. Option D is also true.
Option A is not true as per the passage, and hence, is the answer.

\end{problem}

\newpage
\begin{problem}{}{}
% Subject: varc
\textbf{Instructions}

Instructions   

The passage below is accompanied by a set of questions. Choose the best answer to each question.
It has been said that knowledge, or the problem of knowledge, is the scandal of philosophy. The scandal is philosophy’s apparent inability to show how, when and why we can be sure that we know something or, indeed, that we know anything. Philosopher Michael Williams writes: ‘Is it possible to obtain knowledge at all? This problem is pressing because there are powerful arguments, some very ancient, for the conclusion that it is not . . . Scepticism is the skeleton in Western rationalism’s closet’. While it is not clear that the scandal matters to anyone but philosophers, philosophers point out that it should matter to everyone, at least given a certain conception of knowledge. For, they explain, unless we can ground our claims to knowledge as such, which is to say, distinguish it from mere opinion, superstition, fantasy, wishful thinking, ideology, illusion or delusion, then the actions we take on the basis of presumed knowledge - boarding an airplane, swallowing a pill, finding someone guilty of a crime - will be irrational and unjustifiable.
That is all quite serious-sounding but so also are the rattlings of the skeleton: that is, the sceptic’s contention that we cannot be sure that we know anything - at least not if we think of knowledge as something like having a correct mental representation of reality, and not if we think of reality as something like things-as-they-are-in-themselves, independent of our perceptions, ideas or descriptions. For, the sceptic will note, since reality, under that conception of it, is outside our ken (we cannot catch a glimpse of things-in-themselves around the corner of our own eyes; we cannot form an idea of reality that floats above the processes of our conceiving it), we have no way to compare our mental representations with things-as-they-are-in-themselves and therefore no way to determine whether they are correct or incorrect. Thus the sceptic may repeat (rattling loudly), you cannot be sure you ‘know’ something or anything at all - at least not, he may add (rattling softly before disappearing), if that is the way you conceive ‘knowledge’.
There are a number of ways to handle this situation. The most common is to ignore it. Most people outside the academy - and, indeed, most of us inside it - are unaware of or unperturbed by the philosophical scandal of knowledge and go about our lives without too many epistemic anxieties. We hold our beliefs and presumptive knowledges more or less confidently, usually depending on how we acquired them (I saw it with my own eyes; I heard it on Fox News; a guy at the office told me) and how broadly and strenuously they seem to be shared or endorsed by various relevant people: experts and authorities, friends and family members, colleagues and associates. And we examine our convictions more or less closely, explain them more or less extensively, and defend them more or less vigorously, usually depending on what seems to be at stake for ourselves and/or other people and what resources are available for reassuring ourselves or making our beliefs credible to others (look, it’s right here on the page; add up the figures yourself; I happen to be a heart specialist).

\vspace{1em}
\textbf{Question 11}

Which one of the following, if true, would weaken the author’s claims in the passage?

\begin{enumerate}[label=\Alph*.]
    \item Indian nationalists rejected the cause of English education for women during the colonial period.
    \item Forces of colonial modernity played an important role in shaping anti-colonial Indian nationalism.
    \item The colonial period saw the hybridisation of Indian culture in all realms as it came in contact with British/European culture.
    \item The Industrial Revolution played a crucial role in shaping the economic prowess of Britain in the eighteenth century.
[](https://cracku.in/11-which-one-of-the-following-if-true-would-weaken-th-x-cat-2021-slot-2)
\end{enumerate}
\vspace{1em}
\textbf{Solution:}

The arguments in the passage are based on the premise that the material and spiritual aspects of culture were different. Hence, even if Indian nationalism accepted the superior material ways of the west, they still would not be giving in to colonial dominance and their identity would be preserved by the spiritual aspect, as it remained unaffected. Hence, to weaken the author's argument, we can give a statement that proves that the spiritual aspect was affected too. Option C does that and is the answer.
Rejecting education for women could have more than one reason. It does not imply that the spiritual part of Indian culture was affected by colonialism. Moreover, the separation of roles according to gender is something that nationalist ideology supported, hence, rejecting education based on gender would not contradict it. Option A can be eliminated.
Option B does not weaken the author's argument since the author already agrees that the forces of colonialist modernity helped shape Indian nationalism, but only in the material aspect.
Option D is unrelated to the argument at hand and can be eliminated too.

\end{problem}

\newpage
\begin{problem}{}{}
% Subject: varc
\textbf{Instructions}

Instructions   

The passage below is accompanied by a set of questions. Choose the best answer to each question.
It has been said that knowledge, or the problem of knowledge, is the scandal of philosophy. The scandal is philosophy’s apparent inability to show how, when and why we can be sure that we know something or, indeed, that we know anything. Philosopher Michael Williams writes: ‘Is it possible to obtain knowledge at all? This problem is pressing because there are powerful arguments, some very ancient, for the conclusion that it is not . . . Scepticism is the skeleton in Western rationalism’s closet’. While it is not clear that the scandal matters to anyone but philosophers, philosophers point out that it should matter to everyone, at least given a certain conception of knowledge. For, they explain, unless we can ground our claims to knowledge as such, which is to say, distinguish it from mere opinion, superstition, fantasy, wishful thinking, ideology, illusion or delusion, then the actions we take on the basis of presumed knowledge - boarding an airplane, swallowing a pill, finding someone guilty of a crime - will be irrational and unjustifiable.
That is all quite serious-sounding but so also are the rattlings of the skeleton: that is, the sceptic’s contention that we cannot be sure that we know anything - at least not if we think of knowledge as something like having a correct mental representation of reality, and not if we think of reality as something like things-as-they-are-in-themselves, independent of our perceptions, ideas or descriptions. For, the sceptic will note, since reality, under that conception of it, is outside our ken (we cannot catch a glimpse of things-in-themselves around the corner of our own eyes; we cannot form an idea of reality that floats above the processes of our conceiving it), we have no way to compare our mental representations with things-as-they-are-in-themselves and therefore no way to determine whether they are correct or incorrect. Thus the sceptic may repeat (rattling loudly), you cannot be sure you ‘know’ something or anything at all - at least not, he may add (rattling softly before disappearing), if that is the way you conceive ‘knowledge’.
There are a number of ways to handle this situation. The most common is to ignore it. Most people outside the academy - and, indeed, most of us inside it - are unaware of or unperturbed by the philosophical scandal of knowledge and go about our lives without too many epistemic anxieties. We hold our beliefs and presumptive knowledges more or less confidently, usually depending on how we acquired them (I saw it with my own eyes; I heard it on Fox News; a guy at the office told me) and how broadly and strenuously they seem to be shared or endorsed by various relevant people: experts and authorities, friends and family members, colleagues and associates. And we examine our convictions more or less closely, explain them more or less extensively, and defend them more or less vigorously, usually depending on what seems to be at stake for ourselves and/or other people and what resources are available for reassuring ourselves or making our beliefs credible to others (look, it’s right here on the page; add up the figures yourself; I happen to be a heart specialist).

\vspace{1em}
\textbf{Question 12}

Which one of the following best describes the liberal perception of Indian nationalism?

\begin{enumerate}[label=\Alph*.]
    \item Indian nationalism’s sophistication resided in its distinction of the material from the spiritual spheres.
    \item Indian nationalist discourses reaffirmed traditional gender roles for Indian women.
    \item Indian nationalism embraced the changes brought about by colonialism in Indian women’s traditional gender roles.
    \item Indian nationalist discourses provided an ideological principle of selection.
[](https://cracku.in/12-which-one-of-the-following-best-describes-the-libe-x-cat-2021-slot-2)
\end{enumerate}
\vspace{1em}
\textbf{Solution:}

_Once we match this new meaning of the home/world dichotomy with the _identification of social roles by gender, we get the ideological framework within which nationalism answered the women’s question_. It would be a grave error to see in this, _as liberals are apt to in their despair at the many marks of social conservatism in nationalist practice_ , a total rejection of the West. Quite the contrary: the nationalist paradigm in fact supplied an ideological principle of selection._
From the above excerpt, we can see that the liberals were concerned over the social conservatism that nationalist practice promoted as an ideological principal of selection, where social roles would be selected according to the gender of the person. Hence, Option B is the answer, as it comes the closest in capturing the liberal perception of the same.
The material/spiritual dichotomy has not been discussed in terms of liberal perspective, hence, Option A is out of the score here.
Option C is incorrect. Indian nationalism did not accept the changes brought about the colonialism, rather, promoted the segregation of gender roles according to their spiritual ideology of home/ world dichotomy.
Option D is contrary to what is mentioned in the passage. The author says that the 'ideological principle of selection' was the actual truth, and the liberal perspective was just contrary to what was actually happening.

Instructions   

The passage below is accompanied by a set of questions. Choose the best answer to each question.
Many people believe that truth conveys power. . . . Hence sticking with the truth is the best strategy for gaining power. Unfortunately, this is just a comforting myth. In fact, truth and power have a far more complicated relationship, because in human society, power means two very different things.  
  
On the one hand, power means having the ability to manipulate objective realities: to hunt animals, to construct bridges, to cure diseases, to build atom bombs. This kind of power is closely tied to truth. If you believe a false physical theory, you won’t be able to build an atom bomb. On the other hand, power also means having the ability to manipulate human beliefs, thereby getting lots of people to cooperate effectively. Building atom bombs requires not just a good understanding of physics, but also the coordinated labor of millions of humans. Planet Earth was conquered by Homo sapiens rather than by chimpanzees or elephants, because we are the only mammals that can cooperate in very large numbers. And large-scale cooperation depends on believing common stories. But these stories need not be true. You can unite millions of people by making them believe in completely fictional stories about God, about race or about economics. The dual nature of power and truth results in the curious fact that we humans know many more truths than any other animal, but we also believe in much more nonsense. . . .  
  
When it comes to uniting people around a common story, fiction actually enjoys three inherent advantages over the truth. First, whereas the truth is universal, fictions tend to be local. Consequently if we want to distinguish our tribe from foreigners, a fictional story will serve as a far better identity marker than a true story. . . . The second huge advantage of fiction over truth has to do with the handicap principle, which says that reliable signals must be costly to the signaler. Otherwise, they can easily be faked by cheaters. . . . If political loyalty is signaled by believing a true story, anyone can fake it. But believing ridiculous and outlandish stories exacts greater cost, and is therefore a better signal of loyalty. . . . Third, and most important, the truth is often painful and disturbing. Hence if you stick to unalloyed reality, few people will follow you. An American presidential candidate who tells the American public the truth, the whole truth and nothing but the truth about American history has a 100 percent guarantee of losing the elections. . . . An uncompromising adherence to the truth is an admirable spiritual practice, but it is not a winning political strategy. . . .
Even if we need to pay some price for deactivating our rational faculties, the advantages of increased social cohesion are often so big that fictional stories routinely triumph over the truth in human history. Scholars have known this for thousands of years, which is why scholars often had to decide whether they served the truth or social harmony. Should they aim to unite people by making sure everyone believes in the same fiction, or should they let people know the truth even at the price of disunity?

\end{problem}

\newpage
\begin{problem}{}{}
% Subject: varc
\textbf{Instructions}

Instructions   

The passage below is accompanied by a set of questions. Choose the best answer to each question.
It has been said that knowledge, or the problem of knowledge, is the scandal of philosophy. The scandal is philosophy’s apparent inability to show how, when and why we can be sure that we know something or, indeed, that we know anything. Philosopher Michael Williams writes: ‘Is it possible to obtain knowledge at all? This problem is pressing because there are powerful arguments, some very ancient, for the conclusion that it is not . . . Scepticism is the skeleton in Western rationalism’s closet’. While it is not clear that the scandal matters to anyone but philosophers, philosophers point out that it should matter to everyone, at least given a certain conception of knowledge. For, they explain, unless we can ground our claims to knowledge as such, which is to say, distinguish it from mere opinion, superstition, fantasy, wishful thinking, ideology, illusion or delusion, then the actions we take on the basis of presumed knowledge - boarding an airplane, swallowing a pill, finding someone guilty of a crime - will be irrational and unjustifiable.
That is all quite serious-sounding but so also are the rattlings of the skeleton: that is, the sceptic’s contention that we cannot be sure that we know anything - at least not if we think of knowledge as something like having a correct mental representation of reality, and not if we think of reality as something like things-as-they-are-in-themselves, independent of our perceptions, ideas or descriptions. For, the sceptic will note, since reality, under that conception of it, is outside our ken (we cannot catch a glimpse of things-in-themselves around the corner of our own eyes; we cannot form an idea of reality that floats above the processes of our conceiving it), we have no way to compare our mental representations with things-as-they-are-in-themselves and therefore no way to determine whether they are correct or incorrect. Thus the sceptic may repeat (rattling loudly), you cannot be sure you ‘know’ something or anything at all - at least not, he may add (rattling softly before disappearing), if that is the way you conceive ‘knowledge’.
There are a number of ways to handle this situation. The most common is to ignore it. Most people outside the academy - and, indeed, most of us inside it - are unaware of or unperturbed by the philosophical scandal of knowledge and go about our lives without too many epistemic anxieties. We hold our beliefs and presumptive knowledges more or less confidently, usually depending on how we acquired them (I saw it with my own eyes; I heard it on Fox News; a guy at the office told me) and how broadly and strenuously they seem to be shared or endorsed by various relevant people: experts and authorities, friends and family members, colleagues and associates. And we examine our convictions more or less closely, explain them more or less extensively, and defend them more or less vigorously, usually depending on what seems to be at stake for ourselves and/or other people and what resources are available for reassuring ourselves or making our beliefs credible to others (look, it’s right here on the page; add up the figures yourself; I happen to be a heart specialist).

\vspace{1em}
\textbf{Question 13}

The central theme of the passage is about the choice between:

\begin{enumerate}[label=\Alph*.]
    \item stories that unite people and those that distinguish groups from each other.
    \item attaining social cohesion and propagating objective truth.
    \item leaders who unknowingly spread fictions and those who intentionally do so.
    \item truth and power.
[](https://cracku.in/13-the-central-theme-of-the-passage-is-about-the-choi-x-cat-2021-slot-2)
\end{enumerate}
\vspace{1em}
\textbf{Solution:}

The author begins that passage by saying that truth does not necessarily carry power. He then goes on to explain that to attain social cohesion, sticking to the truth is not always an optimal strategy. In the last paragraph, the author sums up this trade-off:_  
Even if we need to pay some price for deactivating our rational faculties, the advantages of increased social cohesion are often so big that fictional stories routinely triumph over the truth in human history. Scholars have known this for thousands of years, which is why scholars often had to decide whether they served the truth or social harmony._
Thus, Option B is the answer.
The author is not fixated upon the types of stories, not upon what kind of stories do the leaders propagate. Hence, Options A and C can be eliminated.
Power has been mentioned to indicate that sometimes, absolute truth is not the way forward to achieve maximum utility. The main contention of the author is not the trade-off between truth and power but between truth and social cohesion. Option D can be eliminated too.

\end{problem}

\newpage
\begin{problem}{}{}
% Subject: varc
\textbf{Instructions}

Instructions   

The passage below is accompanied by a set of questions. Choose the best answer to each question.
It has been said that knowledge, or the problem of knowledge, is the scandal of philosophy. The scandal is philosophy’s apparent inability to show how, when and why we can be sure that we know something or, indeed, that we know anything. Philosopher Michael Williams writes: ‘Is it possible to obtain knowledge at all? This problem is pressing because there are powerful arguments, some very ancient, for the conclusion that it is not . . . Scepticism is the skeleton in Western rationalism’s closet’. While it is not clear that the scandal matters to anyone but philosophers, philosophers point out that it should matter to everyone, at least given a certain conception of knowledge. For, they explain, unless we can ground our claims to knowledge as such, which is to say, distinguish it from mere opinion, superstition, fantasy, wishful thinking, ideology, illusion or delusion, then the actions we take on the basis of presumed knowledge - boarding an airplane, swallowing a pill, finding someone guilty of a crime - will be irrational and unjustifiable.
That is all quite serious-sounding but so also are the rattlings of the skeleton: that is, the sceptic’s contention that we cannot be sure that we know anything - at least not if we think of knowledge as something like having a correct mental representation of reality, and not if we think of reality as something like things-as-they-are-in-themselves, independent of our perceptions, ideas or descriptions. For, the sceptic will note, since reality, under that conception of it, is outside our ken (we cannot catch a glimpse of things-in-themselves around the corner of our own eyes; we cannot form an idea of reality that floats above the processes of our conceiving it), we have no way to compare our mental representations with things-as-they-are-in-themselves and therefore no way to determine whether they are correct or incorrect. Thus the sceptic may repeat (rattling loudly), you cannot be sure you ‘know’ something or anything at all - at least not, he may add (rattling softly before disappearing), if that is the way you conceive ‘knowledge’.
There are a number of ways to handle this situation. The most common is to ignore it. Most people outside the academy - and, indeed, most of us inside it - are unaware of or unperturbed by the philosophical scandal of knowledge and go about our lives without too many epistemic anxieties. We hold our beliefs and presumptive knowledges more or less confidently, usually depending on how we acquired them (I saw it with my own eyes; I heard it on Fox News; a guy at the office told me) and how broadly and strenuously they seem to be shared or endorsed by various relevant people: experts and authorities, friends and family members, colleagues and associates. And we examine our convictions more or less closely, explain them more or less extensively, and defend them more or less vigorously, usually depending on what seems to be at stake for ourselves and/or other people and what resources are available for reassuring ourselves or making our beliefs credible to others (look, it’s right here on the page; add up the figures yourself; I happen to be a heart specialist).

\vspace{1em}
\textbf{Question 14}

The author would support none of the following statements about political power EXCEPT that:

\begin{enumerate}[label=\Alph*.]
    \item there are definite advantages to promoting fiction, but there needs to be some limit to a pervasive belief in myths.
    \item while unalloyed truth is not recommended, leaders should stay as close as possible to it.
    \item manipulating people’s beliefs is politically advantageous, but a leader who propagates only myths is likely to lose power.
    \item people cannot handle the unvarnished truth, so leaders retain power by deviating from it.
[](https://cracku.in/14-the-author-would-support-none-of-the-following-sta-x-cat-2021-slot-2)
\end{enumerate}
\vspace{1em}
\textbf{Solution:}

The author does not support that there is a limit to the influence that myths have on people, nor does he support imposing one. Hence, Option A can be eliminated.
_Third, and most important, the truth is often painful and disturbing. Hence if you stick to unalloyed reality, few people will follow you. An American presidential candidate who tells the American public the truth, the whole truth and nothing but the truth about American history has a 100 percent guarantee of losing the elections. . . . An uncompromising adherence to the truth is an admirable spiritual practice, but it is not a winning political strategy. . . ._
Option B is contrary to what is being said in the passage. The author says that untarnished truth is not a good recipe for a political win, hence, the candidate should steer clear of that.
Option C is also contrary to what is being said in the passage. According to the author, not conveying the complete truth will allow a person to stay in power.
Option D is in line with the above excerpt and hence, is the answer.

\end{problem}

\newpage
\begin{problem}{}{}
% Subject: varc
\textbf{Instructions}

Instructions   

The passage below is accompanied by a set of questions. Choose the best answer to each question.
It has been said that knowledge, or the problem of knowledge, is the scandal of philosophy. The scandal is philosophy’s apparent inability to show how, when and why we can be sure that we know something or, indeed, that we know anything. Philosopher Michael Williams writes: ‘Is it possible to obtain knowledge at all? This problem is pressing because there are powerful arguments, some very ancient, for the conclusion that it is not . . . Scepticism is the skeleton in Western rationalism’s closet’. While it is not clear that the scandal matters to anyone but philosophers, philosophers point out that it should matter to everyone, at least given a certain conception of knowledge. For, they explain, unless we can ground our claims to knowledge as such, which is to say, distinguish it from mere opinion, superstition, fantasy, wishful thinking, ideology, illusion or delusion, then the actions we take on the basis of presumed knowledge - boarding an airplane, swallowing a pill, finding someone guilty of a crime - will be irrational and unjustifiable.
That is all quite serious-sounding but so also are the rattlings of the skeleton: that is, the sceptic’s contention that we cannot be sure that we know anything - at least not if we think of knowledge as something like having a correct mental representation of reality, and not if we think of reality as something like things-as-they-are-in-themselves, independent of our perceptions, ideas or descriptions. For, the sceptic will note, since reality, under that conception of it, is outside our ken (we cannot catch a glimpse of things-in-themselves around the corner of our own eyes; we cannot form an idea of reality that floats above the processes of our conceiving it), we have no way to compare our mental representations with things-as-they-are-in-themselves and therefore no way to determine whether they are correct or incorrect. Thus the sceptic may repeat (rattling loudly), you cannot be sure you ‘know’ something or anything at all - at least not, he may add (rattling softly before disappearing), if that is the way you conceive ‘knowledge’.
There are a number of ways to handle this situation. The most common is to ignore it. Most people outside the academy - and, indeed, most of us inside it - are unaware of or unperturbed by the philosophical scandal of knowledge and go about our lives without too many epistemic anxieties. We hold our beliefs and presumptive knowledges more or less confidently, usually depending on how we acquired them (I saw it with my own eyes; I heard it on Fox News; a guy at the office told me) and how broadly and strenuously they seem to be shared or endorsed by various relevant people: experts and authorities, friends and family members, colleagues and associates. And we examine our convictions more or less closely, explain them more or less extensively, and defend them more or less vigorously, usually depending on what seems to be at stake for ourselves and/or other people and what resources are available for reassuring ourselves or making our beliefs credible to others (look, it’s right here on the page; add up the figures yourself; I happen to be a heart specialist).

\vspace{1em}
\textbf{Question 15}

The author implies that, like scholars, successful leaders:

\begin{enumerate}[label=\Alph*.]
    \item today know how to create social cohesion better than in the past.
    \item use myths to attain the first type of power.
    \item know how to balance truth and social unity.
    \item need to leverage both types of power to remain in office.
[](https://cracku.in/15-the-author-implies-that-like-scholars-successful-l-x-cat-2021-slot-2)
\end{enumerate}
\vspace{1em}
\textbf{Solution:}

_Even if we need to pay some price for deactivating our rational faculties, the _advantages of increased social cohesion are often so big that fictional stories routinely triumph over the truth in human history. Scholars have known this for thousands of years_ , which is why scholars often had to decide whether they served the truth or social harmony. Should they aim to unite people by making sure everyone believes in the same fiction, or should they let people know the truth even at the price of disunity?_
In the penultimate paragraph, the author mentions how successful leaders balance truth and social unity to achieve an optimal outcome. The above excerpt shows that scholars have known this for a long time too, and have implemented it. Thus, Option C is the answer. 
That leaders and scholars have improved with time when it comes to achieving social cohesion is not implied in the passage. Option A can be eliminated.
We cannot say that scholars use myths to obtain power as leaders do. Hence, Option B can be eliminated.
We cannot say that scholars use myths to stay in office as leaders do. Option D can be eliminated too.

\end{problem}

\newpage
\begin{problem}{}{}
% Subject: varc
\textbf{Instructions}

Instructions   

The passage below is accompanied by a set of questions. Choose the best answer to each question.
It has been said that knowledge, or the problem of knowledge, is the scandal of philosophy. The scandal is philosophy’s apparent inability to show how, when and why we can be sure that we know something or, indeed, that we know anything. Philosopher Michael Williams writes: ‘Is it possible to obtain knowledge at all? This problem is pressing because there are powerful arguments, some very ancient, for the conclusion that it is not . . . Scepticism is the skeleton in Western rationalism’s closet’. While it is not clear that the scandal matters to anyone but philosophers, philosophers point out that it should matter to everyone, at least given a certain conception of knowledge. For, they explain, unless we can ground our claims to knowledge as such, which is to say, distinguish it from mere opinion, superstition, fantasy, wishful thinking, ideology, illusion or delusion, then the actions we take on the basis of presumed knowledge - boarding an airplane, swallowing a pill, finding someone guilty of a crime - will be irrational and unjustifiable.
That is all quite serious-sounding but so also are the rattlings of the skeleton: that is, the sceptic’s contention that we cannot be sure that we know anything - at least not if we think of knowledge as something like having a correct mental representation of reality, and not if we think of reality as something like things-as-they-are-in-themselves, independent of our perceptions, ideas or descriptions. For, the sceptic will note, since reality, under that conception of it, is outside our ken (we cannot catch a glimpse of things-in-themselves around the corner of our own eyes; we cannot form an idea of reality that floats above the processes of our conceiving it), we have no way to compare our mental representations with things-as-they-are-in-themselves and therefore no way to determine whether they are correct or incorrect. Thus the sceptic may repeat (rattling loudly), you cannot be sure you ‘know’ something or anything at all - at least not, he may add (rattling softly before disappearing), if that is the way you conceive ‘knowledge’.
There are a number of ways to handle this situation. The most common is to ignore it. Most people outside the academy - and, indeed, most of us inside it - are unaware of or unperturbed by the philosophical scandal of knowledge and go about our lives without too many epistemic anxieties. We hold our beliefs and presumptive knowledges more or less confidently, usually depending on how we acquired them (I saw it with my own eyes; I heard it on Fox News; a guy at the office told me) and how broadly and strenuously they seem to be shared or endorsed by various relevant people: experts and authorities, friends and family members, colleagues and associates. And we examine our convictions more or less closely, explain them more or less extensively, and defend them more or less vigorously, usually depending on what seems to be at stake for ourselves and/or other people and what resources are available for reassuring ourselves or making our beliefs credible to others (look, it’s right here on the page; add up the figures yourself; I happen to be a heart specialist).

\vspace{1em}
\textbf{Question 16}

Regarding which one of the following quotes could we argue that the author overemphasises the importance of fiction?

\begin{enumerate}[label=\Alph*.]
    \item “. . . scholars often had to decide whether they served the truth or social harmony. Should they aim to unite people by making sure everyone believes in the same fiction, or should they let people know the truth . . .?”
    \item “On the one hand, power means having the ability to manipulate objective realities: to hunt animals, to construct bridges, to cure diseases, to build atom bombs.”
    \item “Hence sticking with the truth is the best strategy for gaining power. Unfortunately, this is just a comforting myth.”
    \item "In fact, truth and power have a far more complicated relationship, because in human society, power means two very different things."
[](https://cracku.in/16-regarding-which-one-of-the-following-quotes-could--x-cat-2021-slot-2)
\end{enumerate}
\vspace{1em}
\textbf{Solution:}

Option A: The author here emphasizes that the choice between truth and social cohesion is a difficult one for scholars, as it means choosing between truth or uniting everyone using a common narrative. Here, the reach and influence of fiction created by that scholar has been overemphasized, and hence, is the answer.
Option B: There is no overemphasis in this option. Since humans have achieved these feats, and these feats do manipulate the objective reality around us, Option B can be eliminated.
Option C: Here too, the importance of fiction has not been overemphasized, but the importance of truth has been downplayed.
Option D: Option D presents a statement, which is unrelated to the emphasis being placed on the importance of myths.

Instructions   

For the following questions answer them individually

\end{problem}

\newpage
\begin{problem}{}{}
% Subject: varc
\textbf{Instructions}

Instructions   

The passage below is accompanied by a set of questions. Choose the best answer to each question.
It has been said that knowledge, or the problem of knowledge, is the scandal of philosophy. The scandal is philosophy’s apparent inability to show how, when and why we can be sure that we know something or, indeed, that we know anything. Philosopher Michael Williams writes: ‘Is it possible to obtain knowledge at all? This problem is pressing because there are powerful arguments, some very ancient, for the conclusion that it is not . . . Scepticism is the skeleton in Western rationalism’s closet’. While it is not clear that the scandal matters to anyone but philosophers, philosophers point out that it should matter to everyone, at least given a certain conception of knowledge. For, they explain, unless we can ground our claims to knowledge as such, which is to say, distinguish it from mere opinion, superstition, fantasy, wishful thinking, ideology, illusion or delusion, then the actions we take on the basis of presumed knowledge - boarding an airplane, swallowing a pill, finding someone guilty of a crime - will be irrational and unjustifiable.
That is all quite serious-sounding but so also are the rattlings of the skeleton: that is, the sceptic’s contention that we cannot be sure that we know anything - at least not if we think of knowledge as something like having a correct mental representation of reality, and not if we think of reality as something like things-as-they-are-in-themselves, independent of our perceptions, ideas or descriptions. For, the sceptic will note, since reality, under that conception of it, is outside our ken (we cannot catch a glimpse of things-in-themselves around the corner of our own eyes; we cannot form an idea of reality that floats above the processes of our conceiving it), we have no way to compare our mental representations with things-as-they-are-in-themselves and therefore no way to determine whether they are correct or incorrect. Thus the sceptic may repeat (rattling loudly), you cannot be sure you ‘know’ something or anything at all - at least not, he may add (rattling softly before disappearing), if that is the way you conceive ‘knowledge’.
There are a number of ways to handle this situation. The most common is to ignore it. Most people outside the academy - and, indeed, most of us inside it - are unaware of or unperturbed by the philosophical scandal of knowledge and go about our lives without too many epistemic anxieties. We hold our beliefs and presumptive knowledges more or less confidently, usually depending on how we acquired them (I saw it with my own eyes; I heard it on Fox News; a guy at the office told me) and how broadly and strenuously they seem to be shared or endorsed by various relevant people: experts and authorities, friends and family members, colleagues and associates. And we examine our convictions more or less closely, explain them more or less extensively, and defend them more or less vigorously, usually depending on what seems to be at stake for ourselves and/or other people and what resources are available for reassuring ourselves or making our beliefs credible to others (look, it’s right here on the page; add up the figures yourself; I happen to be a heart specialist).

\vspace{1em}
\textbf{Question 17}

**The passage given below is followed by four alternate summaries. Choose the option that best captures the essence of the passage.**
Biologists who publish their research directly to the Web have been labelled as “rogue”, but physicists have been routinely publishing research digitally (“preprints”), prior to submitting in a peer-reviewed journal. Advocates of preprints argue that quick and open dissemination of research speeds up scientific progress and allows for wider access to knowledge. But some journals still don’t accept research previously published as a preprint. Even if the idea of preprints is gaining ground, one of the biggest barriers for biologists is how they would be viewed by members of their conservative research community.

\begin{enumerate}[label=\Alph*.]
    \item One of the advantages of digital preprints of research is they hasten the dissemination process, but these are not accepted by most scientific communities.
    \item Compared to biologists, physicists are less conservative in their acceptance of digital pre-publication of research papers, which allows for faster dissemination of knowledge.
    \item While digital publication of research is gaining popularity in many scientific disciplines, almost all peer-reviewed journals are reluctant to accept papers that have been published before.
    \item Preprints of research are frowned on by some scientific fields as they do not undergo a rigourous reviewing process but are accepted among biologists as a quick way to disseminate information.
[](https://cracku.in/17-the-passage-given-below-is-followed-by-four-altern-x-cat-2021-slot-2)
\end{enumerate}
\vspace{1em}
\textbf{Solution:}

The main points of the paragraph are:  
1. As compared to physicists, biologists are more conservative when it comes to the subject of preprints.  
2. Preprints allow faster dissemination of knowledge.
A: Misses out the comparison between biologists and physicists.
B: Captures both the points appropriately and is the answer.
C: Also misses out the comparison between biologists and physicists.
D: Factually incorrect, physicists and not biologists are open to the idea of preprints.

\end{problem}

\newpage
\begin{problem}{}{}
% Subject: varc
\textbf{Instructions}

Instructions   

The passage below is accompanied by a set of questions. Choose the best answer to each question.
It has been said that knowledge, or the problem of knowledge, is the scandal of philosophy. The scandal is philosophy’s apparent inability to show how, when and why we can be sure that we know something or, indeed, that we know anything. Philosopher Michael Williams writes: ‘Is it possible to obtain knowledge at all? This problem is pressing because there are powerful arguments, some very ancient, for the conclusion that it is not . . . Scepticism is the skeleton in Western rationalism’s closet’. While it is not clear that the scandal matters to anyone but philosophers, philosophers point out that it should matter to everyone, at least given a certain conception of knowledge. For, they explain, unless we can ground our claims to knowledge as such, which is to say, distinguish it from mere opinion, superstition, fantasy, wishful thinking, ideology, illusion or delusion, then the actions we take on the basis of presumed knowledge - boarding an airplane, swallowing a pill, finding someone guilty of a crime - will be irrational and unjustifiable.
That is all quite serious-sounding but so also are the rattlings of the skeleton: that is, the sceptic’s contention that we cannot be sure that we know anything - at least not if we think of knowledge as something like having a correct mental representation of reality, and not if we think of reality as something like things-as-they-are-in-themselves, independent of our perceptions, ideas or descriptions. For, the sceptic will note, since reality, under that conception of it, is outside our ken (we cannot catch a glimpse of things-in-themselves around the corner of our own eyes; we cannot form an idea of reality that floats above the processes of our conceiving it), we have no way to compare our mental representations with things-as-they-are-in-themselves and therefore no way to determine whether they are correct or incorrect. Thus the sceptic may repeat (rattling loudly), you cannot be sure you ‘know’ something or anything at all - at least not, he may add (rattling softly before disappearing), if that is the way you conceive ‘knowledge’.
There are a number of ways to handle this situation. The most common is to ignore it. Most people outside the academy - and, indeed, most of us inside it - are unaware of or unperturbed by the philosophical scandal of knowledge and go about our lives without too many epistemic anxieties. We hold our beliefs and presumptive knowledges more or less confidently, usually depending on how we acquired them (I saw it with my own eyes; I heard it on Fox News; a guy at the office told me) and how broadly and strenuously they seem to be shared or endorsed by various relevant people: experts and authorities, friends and family members, colleagues and associates. And we examine our convictions more or less closely, explain them more or less extensively, and defend them more or less vigorously, usually depending on what seems to be at stake for ourselves and/or other people and what resources are available for reassuring ourselves or making our beliefs credible to others (look, it’s right here on the page; add up the figures yourself; I happen to be a heart specialist).

\vspace{1em}
\textbf{Question 18}

**The four sentences (labelled 1, 2, 3, 4) below, when properly sequenced would yield a coherent paragraph. Decide on the proper sequencing of the order of the sentences and key in the sequence of the four numbers as your answer:**
1. But today there is an epochal challenge to rethink and reconstitute the vision and practice of development as a shared responsibility - a sharing which binds both the agent and the audience, the developed world and the developing, in a bond of shared destiny.  
2. We are at a crossroads now in our vision and practice of development.  
3. This calls for the cultivation of an appropriate ethical mode of being in our lives which enables us to realize this global and planetary situation of shared living and responsibility.  
4. Half a century ago, development began as a hope for a better human possibility, but in the last fifty years, this hope has lost itself in the dreary desert of various kinds of hegemonic applications.
Backspace  
789  
456  
123  
0.-  
Clear All  

Submit
[](https://cracku.in/18-the-four-sentences-labelled-1-2-3-4-below-when-pro-x-cat-2021-slot-2)

\begin{enumerate}[label=\Alph*.]
\end{enumerate}
\vspace{1em}
\textbf{Solution:}

A brief reading of the sentences suggests that the paragraph is about the change needed in the way we go about development. 2 introduces the topic at hand, that this is a watershed moment when it comes to the subject of development. 41 make a mandatory pair, which talks about what the purpose of development was at the beginning and how it needs to be altered to suit the needs of today. 3 then aptly ends the paragraph, suggesting the measures that could be taken to counter the same. Hence, he proper sequence would be 2413.

\end{problem}

\newpage
\begin{problem}{}{}
% Subject: varc
\textbf{Instructions}

Instructions   

The passage below is accompanied by a set of questions. Choose the best answer to each question.
It has been said that knowledge, or the problem of knowledge, is the scandal of philosophy. The scandal is philosophy’s apparent inability to show how, when and why we can be sure that we know something or, indeed, that we know anything. Philosopher Michael Williams writes: ‘Is it possible to obtain knowledge at all? This problem is pressing because there are powerful arguments, some very ancient, for the conclusion that it is not . . . Scepticism is the skeleton in Western rationalism’s closet’. While it is not clear that the scandal matters to anyone but philosophers, philosophers point out that it should matter to everyone, at least given a certain conception of knowledge. For, they explain, unless we can ground our claims to knowledge as such, which is to say, distinguish it from mere opinion, superstition, fantasy, wishful thinking, ideology, illusion or delusion, then the actions we take on the basis of presumed knowledge - boarding an airplane, swallowing a pill, finding someone guilty of a crime - will be irrational and unjustifiable.
That is all quite serious-sounding but so also are the rattlings of the skeleton: that is, the sceptic’s contention that we cannot be sure that we know anything - at least not if we think of knowledge as something like having a correct mental representation of reality, and not if we think of reality as something like things-as-they-are-in-themselves, independent of our perceptions, ideas or descriptions. For, the sceptic will note, since reality, under that conception of it, is outside our ken (we cannot catch a glimpse of things-in-themselves around the corner of our own eyes; we cannot form an idea of reality that floats above the processes of our conceiving it), we have no way to compare our mental representations with things-as-they-are-in-themselves and therefore no way to determine whether they are correct or incorrect. Thus the sceptic may repeat (rattling loudly), you cannot be sure you ‘know’ something or anything at all - at least not, he may add (rattling softly before disappearing), if that is the way you conceive ‘knowledge’.
There are a number of ways to handle this situation. The most common is to ignore it. Most people outside the academy - and, indeed, most of us inside it - are unaware of or unperturbed by the philosophical scandal of knowledge and go about our lives without too many epistemic anxieties. We hold our beliefs and presumptive knowledges more or less confidently, usually depending on how we acquired them (I saw it with my own eyes; I heard it on Fox News; a guy at the office told me) and how broadly and strenuously they seem to be shared or endorsed by various relevant people: experts and authorities, friends and family members, colleagues and associates. And we examine our convictions more or less closely, explain them more or less extensively, and defend them more or less vigorously, usually depending on what seems to be at stake for ourselves and/or other people and what resources are available for reassuring ourselves or making our beliefs credible to others (look, it’s right here on the page; add up the figures yourself; I happen to be a heart specialist).

\vspace{1em}
\textbf{Question 19}

**Five jumbled up sentences, related to a topic, are given below. Four of them can be put together to form a coherent paragraph. Identify the odd one out and key in the number of the sentence as your answer:**  
1. The care with which philosophers examine arguments for and against forms of biotechnology makes this an excellent primer on formulating and assessing moral arguments.  
2. Although most people find at least some forms of genetic engineering disquieting, it is not easy to articulate why: what is wrong with re-engineering our nature?  
3. Breakthroughs in genetics present us with the promise that we will soon be able to prevent a host of debilitating diseases, and the predicament that our newfound genetic knowledge may enable us to enhance our genetic traits.  
4. To grapple with the ethics of enhancement, we need to confront questions that verge on theology, which is why modern philosophers and political theorists tend to shrink from them.  
5. One argument is that the drive for human perfection through genetics is objectionable as it represents a bid for mastery that fails to appreciate the gifts of human powers and achievements.
Backspace  
789  
456  
123  
0.-  
Clear All  

Submit
[](https://cracku.in/19-five-jumbled-up-sentences-related-to-a-topic-are-g-x-cat-2021-slot-2)

\begin{enumerate}[label=\Alph*.]
\end{enumerate}
\vspace{1em}
\textbf{Solution:}

The sentences have been taken from Harvard's Justice, and have been modified considerably. Since a paragraph has not been directly taken here, the better way of elimination here would be to evaluate the major points of each sentence and see which one runs tangent to the discussion at hand. (During the examination, one must try both ways to solve: arranging and eliminating.)
1. Using the debate on biotechnology to evaluate moral arguments.
2. Why is bioengineering disputed?
3. The promise of bioengineering.
4. Ethics of bioengineering based on theology.
5. The theological argument.
We can see here that the last four sentences try to examine why bioengineering is disputed in spite of its huge potential. Then reasons are given about the question on its ethicality, and how it is closely associated with theology on the matter.
1 however runs tangential to the discussion. The main focus is bioengineering while 1 aims to shift the focus and use the debate on the matter as a stepping stone to reach another goal: evaluating/formulating moral arguments. Hence, 1 is the odd one out here.

\end{problem}

\newpage
\begin{problem}{}{}
% Subject: varc
\textbf{Instructions}

Instructions   

The passage below is accompanied by a set of questions. Choose the best answer to each question.
It has been said that knowledge, or the problem of knowledge, is the scandal of philosophy. The scandal is philosophy’s apparent inability to show how, when and why we can be sure that we know something or, indeed, that we know anything. Philosopher Michael Williams writes: ‘Is it possible to obtain knowledge at all? This problem is pressing because there are powerful arguments, some very ancient, for the conclusion that it is not . . . Scepticism is the skeleton in Western rationalism’s closet’. While it is not clear that the scandal matters to anyone but philosophers, philosophers point out that it should matter to everyone, at least given a certain conception of knowledge. For, they explain, unless we can ground our claims to knowledge as such, which is to say, distinguish it from mere opinion, superstition, fantasy, wishful thinking, ideology, illusion or delusion, then the actions we take on the basis of presumed knowledge - boarding an airplane, swallowing a pill, finding someone guilty of a crime - will be irrational and unjustifiable.
That is all quite serious-sounding but so also are the rattlings of the skeleton: that is, the sceptic’s contention that we cannot be sure that we know anything - at least not if we think of knowledge as something like having a correct mental representation of reality, and not if we think of reality as something like things-as-they-are-in-themselves, independent of our perceptions, ideas or descriptions. For, the sceptic will note, since reality, under that conception of it, is outside our ken (we cannot catch a glimpse of things-in-themselves around the corner of our own eyes; we cannot form an idea of reality that floats above the processes of our conceiving it), we have no way to compare our mental representations with things-as-they-are-in-themselves and therefore no way to determine whether they are correct or incorrect. Thus the sceptic may repeat (rattling loudly), you cannot be sure you ‘know’ something or anything at all - at least not, he may add (rattling softly before disappearing), if that is the way you conceive ‘knowledge’.
There are a number of ways to handle this situation. The most common is to ignore it. Most people outside the academy - and, indeed, most of us inside it - are unaware of or unperturbed by the philosophical scandal of knowledge and go about our lives without too many epistemic anxieties. We hold our beliefs and presumptive knowledges more or less confidently, usually depending on how we acquired them (I saw it with my own eyes; I heard it on Fox News; a guy at the office told me) and how broadly and strenuously they seem to be shared or endorsed by various relevant people: experts and authorities, friends and family members, colleagues and associates. And we examine our convictions more or less closely, explain them more or less extensively, and defend them more or less vigorously, usually depending on what seems to be at stake for ourselves and/or other people and what resources are available for reassuring ourselves or making our beliefs credible to others (look, it’s right here on the page; add up the figures yourself; I happen to be a heart specialist).

\vspace{1em}
\textbf{Question 20}

**Five jumbled up sentences, related to a topic, are given below. Four of them can be put together to form a coherent paragraph. Identify the odd one out and key in the number of the sentence as your answer:**  
1. It has taken on a warm, fuzzy glow in the advertising world, where its potential is being widely discussed, and it is being claimed as the undeniable wave of the future.  
2. There is little enthusiasm for this in the scientific arena; for them marketing is not a science, and only a handful of studies have been published in scientific journals.  
3. The new, growing field of neuromarketing attempts to reveal the inner workings of consumer behaviour and is an extension of the study of how choices and decisions are made.  
4. Some see neuromarketing as an attempt to make the "art" of advertising into a science, being used by marketing experts to back up their proposals with some form of real data.  
5. The marketing gurus have already started drawing on psychology in developing tests and theories, and advertising people have borrowed the idea of the focus group from social scientists.
Backspace  
789  
456  
123  
0.-  
Clear All  

Submit
[](https://cracku.in/20-five-jumbled-up-sentences-related-to-a-topic-are-g-x-cat-2021-slot-2)

\begin{enumerate}[label=\Alph*.]
\end{enumerate}
\vspace{1em}
\textbf{Solution:}

A brief reading of the sentences tells us that the paragraph must be about the industry of neuromarketing, which is still in its embryonic phase. 3 can be the opening sentence to the paragraph, as it introduces the topic at hand. All the other sentences need a sentence before them that introduces what is being talked about. 
1,4, and 2 then go on to talk about the opinion of differernt associated parties on the matter. It has taken the advertising industry by storm. Others feel that this 'art' is being masked as a science, and many lack enthusiasm on the matter.
5 however, does not fit in here. The reason is that it talks about 'psychology', which is different from the use of neural science. Even if one is not familiar with the difference, we can see that it goes a step forward to talk about the application of a science, whereas the paragraph is mostly concerned with a growing science and how it is shaping public opinion. Hence, 5 is the odd one out.

\end{problem}

\newpage
\begin{problem}{}{}
% Subject: varc
\textbf{Instructions}

Instructions   

The passage below is accompanied by a set of questions. Choose the best answer to each question.
It has been said that knowledge, or the problem of knowledge, is the scandal of philosophy. The scandal is philosophy’s apparent inability to show how, when and why we can be sure that we know something or, indeed, that we know anything. Philosopher Michael Williams writes: ‘Is it possible to obtain knowledge at all? This problem is pressing because there are powerful arguments, some very ancient, for the conclusion that it is not . . . Scepticism is the skeleton in Western rationalism’s closet’. While it is not clear that the scandal matters to anyone but philosophers, philosophers point out that it should matter to everyone, at least given a certain conception of knowledge. For, they explain, unless we can ground our claims to knowledge as such, which is to say, distinguish it from mere opinion, superstition, fantasy, wishful thinking, ideology, illusion or delusion, then the actions we take on the basis of presumed knowledge - boarding an airplane, swallowing a pill, finding someone guilty of a crime - will be irrational and unjustifiable.
That is all quite serious-sounding but so also are the rattlings of the skeleton: that is, the sceptic’s contention that we cannot be sure that we know anything - at least not if we think of knowledge as something like having a correct mental representation of reality, and not if we think of reality as something like things-as-they-are-in-themselves, independent of our perceptions, ideas or descriptions. For, the sceptic will note, since reality, under that conception of it, is outside our ken (we cannot catch a glimpse of things-in-themselves around the corner of our own eyes; we cannot form an idea of reality that floats above the processes of our conceiving it), we have no way to compare our mental representations with things-as-they-are-in-themselves and therefore no way to determine whether they are correct or incorrect. Thus the sceptic may repeat (rattling loudly), you cannot be sure you ‘know’ something or anything at all - at least not, he may add (rattling softly before disappearing), if that is the way you conceive ‘knowledge’.
There are a number of ways to handle this situation. The most common is to ignore it. Most people outside the academy - and, indeed, most of us inside it - are unaware of or unperturbed by the philosophical scandal of knowledge and go about our lives without too many epistemic anxieties. We hold our beliefs and presumptive knowledges more or less confidently, usually depending on how we acquired them (I saw it with my own eyes; I heard it on Fox News; a guy at the office told me) and how broadly and strenuously they seem to be shared or endorsed by various relevant people: experts and authorities, friends and family members, colleagues and associates. And we examine our convictions more or less closely, explain them more or less extensively, and defend them more or less vigorously, usually depending on what seems to be at stake for ourselves and/or other people and what resources are available for reassuring ourselves or making our beliefs credible to others (look, it’s right here on the page; add up the figures yourself; I happen to be a heart specialist).

\vspace{1em}
\textbf{Question 21}

**The four sentences (labelled 1, 2, 3 and 4) below, when properly sequenced would yield a coherent paragraph. Decide on the proper sequencing of the order of the sentences and key in the sequence of the four numbers as your answer:**  
1. Look forward a few decades to an invention which can end the energy crisis, change the global economy and curb climate change at a stroke: commercial fusion power.  
2. To gain meaningful insights, logic has to be accompanied by asking probing questions of nature through controlled tests, precise observations and clever analysis.  
3. The greatest of all inventions is the über-invention that has provided the insights on which others depend: the modern scientific method.  
4. This invention is inconceivable without the scientific method; it will rest on the application of a diverse range of scientific insights, such as the process transforming hydrogen into helium to release huge amounts of energy.
Backspace  
789  
456  
123  
0.-  
Clear All  

Submit
[](https://cracku.in/21-the-four-sentences-labelled-1-2-3-and-4-below-when-x-cat-2021-slot-2)

\begin{enumerate}[label=\Alph*.]
\end{enumerate}
\vspace{1em}
\textbf{Solution:}

A brief reading of the sentences suggests that the paragraph is about great inventions, focusing on the importance of scientific method and how it forms the foundation for other great inventions. 32 forms an introductory pair that claims that the modern scientific method must be the greatest of all inventions and then talks about its mode of inquiry.
1 then presents an invention that could solve many problems in the future. 4 then claims that the invention would have been impossible if the scientific method did not precede it. Hence, the coherent arrangement is 3214.

\end{problem}

\newpage
\begin{problem}{}{}
% Subject: varc
\textbf{Instructions}

Instructions   

The passage below is accompanied by a set of questions. Choose the best answer to each question.
It has been said that knowledge, or the problem of knowledge, is the scandal of philosophy. The scandal is philosophy’s apparent inability to show how, when and why we can be sure that we know something or, indeed, that we know anything. Philosopher Michael Williams writes: ‘Is it possible to obtain knowledge at all? This problem is pressing because there are powerful arguments, some very ancient, for the conclusion that it is not . . . Scepticism is the skeleton in Western rationalism’s closet’. While it is not clear that the scandal matters to anyone but philosophers, philosophers point out that it should matter to everyone, at least given a certain conception of knowledge. For, they explain, unless we can ground our claims to knowledge as such, which is to say, distinguish it from mere opinion, superstition, fantasy, wishful thinking, ideology, illusion or delusion, then the actions we take on the basis of presumed knowledge - boarding an airplane, swallowing a pill, finding someone guilty of a crime - will be irrational and unjustifiable.
That is all quite serious-sounding but so also are the rattlings of the skeleton: that is, the sceptic’s contention that we cannot be sure that we know anything - at least not if we think of knowledge as something like having a correct mental representation of reality, and not if we think of reality as something like things-as-they-are-in-themselves, independent of our perceptions, ideas or descriptions. For, the sceptic will note, since reality, under that conception of it, is outside our ken (we cannot catch a glimpse of things-in-themselves around the corner of our own eyes; we cannot form an idea of reality that floats above the processes of our conceiving it), we have no way to compare our mental representations with things-as-they-are-in-themselves and therefore no way to determine whether they are correct or incorrect. Thus the sceptic may repeat (rattling loudly), you cannot be sure you ‘know’ something or anything at all - at least not, he may add (rattling softly before disappearing), if that is the way you conceive ‘knowledge’.
There are a number of ways to handle this situation. The most common is to ignore it. Most people outside the academy - and, indeed, most of us inside it - are unaware of or unperturbed by the philosophical scandal of knowledge and go about our lives without too many epistemic anxieties. We hold our beliefs and presumptive knowledges more or less confidently, usually depending on how we acquired them (I saw it with my own eyes; I heard it on Fox News; a guy at the office told me) and how broadly and strenuously they seem to be shared or endorsed by various relevant people: experts and authorities, friends and family members, colleagues and associates. And we examine our convictions more or less closely, explain them more or less extensively, and defend them more or less vigorously, usually depending on what seems to be at stake for ourselves and/or other people and what resources are available for reassuring ourselves or making our beliefs credible to others (look, it’s right here on the page; add up the figures yourself; I happen to be a heart specialist).

\vspace{1em}
\textbf{Question 22}

**The passage given below is followed by four alternate summaries. Choose the option that best captures the essence of the passage.**
Creativity is now viewed as the engine of economic progress. Various organizations are devoted to its study and promotion; there are encyclopedias and handbooks surveying creativity research. But this proliferating success has tended to erode creativity’s stable identity: it has become so invested with value that it has become impossible to police its meaning and the practices that supposedly identify and encourage it. Many people and organizations committed to producing original thoughts now feel that undue obsession with the idea of creativity gets in the way of real creativity.

\begin{enumerate}[label=\Alph*.]
    \item The obsession with original thought, how it can be promoted and researched, has made it impossible for people and organizations to define the concept anymore.
    \item The industry that has built up around researching what comprises and encourages creativity has destroyed the creative process itself.
    \item Creativity has proliferated to the extent that is no longer a stable process, and its mutating identity has stifled the creative process.
    \item The value assigned to creativity today has assumed such proportions that the concept itself has lost its real meaning and this is hampering the engendering of real creativity.
[](https://cracku.in/22-the-passage-given-below-is-followed-by-four-altern-x-cat-2021-slot-2)
\end{enumerate}
\vspace{1em}
\textbf{Solution:}

The main points of the paragraph are:
1. The value of creativity to economic progress has been realised, with serious investment being done to study/promote it.
2. But this success fires back. Policing its meaning can lead to obsession, hampering creativity itself.
A: It is extreme in approach. The paragraph does not imply that it has become impossible to define the concept, but it becomes difficult to practice creativity when it is being forced on oneself.
B: It is also extreme. The obsession hampers, not completely destroys the creative process.
C: This option is a distortion and fails to capture the above points.
D: Comes the closest to capturing the above two points, and hence, is the answer.

\end{problem}

\newpage
\begin{problem}{}{}
% Subject: varc
\textbf{Instructions}

Instructions   

The passage below is accompanied by a set of questions. Choose the best answer to each question.
It has been said that knowledge, or the problem of knowledge, is the scandal of philosophy. The scandal is philosophy’s apparent inability to show how, when and why we can be sure that we know something or, indeed, that we know anything. Philosopher Michael Williams writes: ‘Is it possible to obtain knowledge at all? This problem is pressing because there are powerful arguments, some very ancient, for the conclusion that it is not . . . Scepticism is the skeleton in Western rationalism’s closet’. While it is not clear that the scandal matters to anyone but philosophers, philosophers point out that it should matter to everyone, at least given a certain conception of knowledge. For, they explain, unless we can ground our claims to knowledge as such, which is to say, distinguish it from mere opinion, superstition, fantasy, wishful thinking, ideology, illusion or delusion, then the actions we take on the basis of presumed knowledge - boarding an airplane, swallowing a pill, finding someone guilty of a crime - will be irrational and unjustifiable.
That is all quite serious-sounding but so also are the rattlings of the skeleton: that is, the sceptic’s contention that we cannot be sure that we know anything - at least not if we think of knowledge as something like having a correct mental representation of reality, and not if we think of reality as something like things-as-they-are-in-themselves, independent of our perceptions, ideas or descriptions. For, the sceptic will note, since reality, under that conception of it, is outside our ken (we cannot catch a glimpse of things-in-themselves around the corner of our own eyes; we cannot form an idea of reality that floats above the processes of our conceiving it), we have no way to compare our mental representations with things-as-they-are-in-themselves and therefore no way to determine whether they are correct or incorrect. Thus the sceptic may repeat (rattling loudly), you cannot be sure you ‘know’ something or anything at all - at least not, he may add (rattling softly before disappearing), if that is the way you conceive ‘knowledge’.
There are a number of ways to handle this situation. The most common is to ignore it. Most people outside the academy - and, indeed, most of us inside it - are unaware of or unperturbed by the philosophical scandal of knowledge and go about our lives without too many epistemic anxieties. We hold our beliefs and presumptive knowledges more or less confidently, usually depending on how we acquired them (I saw it with my own eyes; I heard it on Fox News; a guy at the office told me) and how broadly and strenuously they seem to be shared or endorsed by various relevant people: experts and authorities, friends and family members, colleagues and associates. And we examine our convictions more or less closely, explain them more or less extensively, and defend them more or less vigorously, usually depending on what seems to be at stake for ourselves and/or other people and what resources are available for reassuring ourselves or making our beliefs credible to others (look, it’s right here on the page; add up the figures yourself; I happen to be a heart specialist).

\vspace{1em}
\textbf{Question 23}

**The four sentences (labelled 1, 2, 3 and 4) below, when properly sequenced would yield a coherent paragraph. Decide on the proper sequencing of the order of the sentences and key in the sequence of the four numbers as your answer:**  
1. The US has long maintained that the Northwest Passage is an international strait through which its commercial and military vessels have the right to pass without seeking Canada’s permission.  
2. Canada, which officially acquired the group of islands forming the Northwest Passage in 1880, claims sovereignty over all the shipping routes through the Passage.  
3. The dispute could be transitory, however, as scientists speculate that the entire Arctic Ocean will soon be ice-free in summer, so ship owners will not have to ask for permission to sail through any of the Northwest Passage routes.  
4. The US and Canada have never legally settled the question of access through the Passage, but have an agreement whereby the US needs to seek Canada’s consent for any transit.
Backspace  
789  
456  
123  
0.-  
Clear All  

Submit
[](https://cracku.in/23-the-four-sentences-labelled-1-2-3-and-4-below-when-x-cat-2021-slot-2)

\begin{enumerate}[label=\Alph*.]
\end{enumerate}
\vspace{1em}
\textbf{Solution:}

A brief reading of the sentences suggests that the paragraph is about the dispute between the US and Canada over the Northwest passage.
2 is a better opening sentence than 1 as 21 forms a good introduction into the dispute. 2 mentions that Canada claims ownership over the passage, and 1 then mentions the counterclaim the US offers: the passage is an international route.
43 becomes a mandatory pair which follows 21. 4 mentions the historical development in the dispute: it has not been legally settled yet. 3 then hints about the future of this dispute: the dispute itself would vanish once the entire Arctic begins to stay ice-free in summer.
Hence, the coherent arrangement is 2143.

\end{problem}

\newpage
\begin{problem}{}{}
% Subject: varc
\textbf{Instructions}

Instructions   

The passage below is accompanied by a set of questions. Choose the best answer to each question.
It has been said that knowledge, or the problem of knowledge, is the scandal of philosophy. The scandal is philosophy’s apparent inability to show how, when and why we can be sure that we know something or, indeed, that we know anything. Philosopher Michael Williams writes: ‘Is it possible to obtain knowledge at all? This problem is pressing because there are powerful arguments, some very ancient, for the conclusion that it is not . . . Scepticism is the skeleton in Western rationalism’s closet’. While it is not clear that the scandal matters to anyone but philosophers, philosophers point out that it should matter to everyone, at least given a certain conception of knowledge. For, they explain, unless we can ground our claims to knowledge as such, which is to say, distinguish it from mere opinion, superstition, fantasy, wishful thinking, ideology, illusion or delusion, then the actions we take on the basis of presumed knowledge - boarding an airplane, swallowing a pill, finding someone guilty of a crime - will be irrational and unjustifiable.
That is all quite serious-sounding but so also are the rattlings of the skeleton: that is, the sceptic’s contention that we cannot be sure that we know anything - at least not if we think of knowledge as something like having a correct mental representation of reality, and not if we think of reality as something like things-as-they-are-in-themselves, independent of our perceptions, ideas or descriptions. For, the sceptic will note, since reality, under that conception of it, is outside our ken (we cannot catch a glimpse of things-in-themselves around the corner of our own eyes; we cannot form an idea of reality that floats above the processes of our conceiving it), we have no way to compare our mental representations with things-as-they-are-in-themselves and therefore no way to determine whether they are correct or incorrect. Thus the sceptic may repeat (rattling loudly), you cannot be sure you ‘know’ something or anything at all - at least not, he may add (rattling softly before disappearing), if that is the way you conceive ‘knowledge’.
There are a number of ways to handle this situation. The most common is to ignore it. Most people outside the academy - and, indeed, most of us inside it - are unaware of or unperturbed by the philosophical scandal of knowledge and go about our lives without too many epistemic anxieties. We hold our beliefs and presumptive knowledges more or less confidently, usually depending on how we acquired them (I saw it with my own eyes; I heard it on Fox News; a guy at the office told me) and how broadly and strenuously they seem to be shared or endorsed by various relevant people: experts and authorities, friends and family members, colleagues and associates. And we examine our convictions more or less closely, explain them more or less extensively, and defend them more or less vigorously, usually depending on what seems to be at stake for ourselves and/or other people and what resources are available for reassuring ourselves or making our beliefs credible to others (look, it’s right here on the page; add up the figures yourself; I happen to be a heart specialist).

\vspace{1em}
\textbf{Question 24}

**The passage given below is followed by four alternate summaries. Choose the option that best captures the essence of the passage.**
The unlikely alliance of the incumbent industrialist and the distressed unemployed worker is especially powerful amid the debris of corporate bankruptcies and layoffs. In an economic downturn, the capitalist is more likely to focus on costs of the competition emanating from free markets than on the opportunities they create. And the unemployed worker will find many others in a similar condition and with anxieties similar to his, which will make it easier for them to organize together. Using the cover and the political organization provided by the distressed, the capitalist captures the political agenda.

\begin{enumerate}[label=\Alph*.]
    \item In an economic downturn, the capitalists use the anxieties of the unemployed and their political organisation to set the political agenda to suit their economic interests.
    \item The purpose of an unlikely alliance between the industrialist and the unemployed during an economic downturn is to stifle competition in free markets.
    \item An economic downturn creates competition because of which the capitalists capture the political agenda created by the political organisation provided by the unemployed.
    \item An unlikely alliance of the industrialist and the unemployed happens during an economic downturn in which they come together to unite politically and capture the political agenda.
[](https://cracku.in/24-the-passage-given-below-is-followed-by-four-altern-x-cat-2021-slot-2)
\end{enumerate}
\vspace{1em}
\textbf{Solution:}

The main points of the paragraph are:
1. In an economic disaster, the atypical alliance of established industrialist and unemployed workers proves powerful.
2. Anxieties and anticipation lead them to look after their interests.
3. It is the industrialist that benefits the most as he is able to use the latter to achieve his vested political interests.
A: Comes the closest in capturing all three points, and hence, is the answer.
B: Distortion. This purpose has not been mentioned in the passage.
C: Distortion. It has not been mentioned that an economic downturn creates competition. It has been mentioned that during such a disaster, the industrialist is more likely to focus on the downsides emerging from free-market (competition) than the upsides.
D: Distortion. It has been implied that the industrialist manipulates the situation to fulfill his own political agenda, and not that the two parties come together to achieve a single goal.
[](https://cracku.in/downloads/13895)
  *  
  *

\end{problem}

\newpage
\begin{problem}{}{}
% Subject: varc
\textbf{Instructions}

Instructions   

The passage below is accompanied by a set of questions. Choose the best answer to each question.
Starting in 1957, [Noam Chomsky] proclaimed a new doctrine: Language, that most human of all attributes, was innate. The grammatical faculty was built into the infant brain, and your average 3-year-old was not a mere apprentice in the great enterprise of absorbing English from his or her parents, but a “linguistic genius.” Since this message was couched in terms of Chomskyan theoretical linguistics, in discourse so opaque that it was nearly incomprehensible even to some scholars, many people did not hear it. Now, in a brilliant, witty and altogether satisfying book, Mr. Chomsky's colleague Steven Pinker . . . has brought Mr. Chomsky's findings to everyman. In “The Language Instinct” he has gathered persuasive data from such diverse fields as cognitive neuroscience, developmental psychology and speech therapy to make his points, and when he disagrees with Mr. Chomsky he tells you so. . . .
For Mr. Chomsky and Mr. Pinker, somewhere in the human brain there is a complex set of neural circuits that have been programmed with “super-rules” (making up what Mr. Chomsky calls “universal grammar”), and that these rules are unconscious and instinctive. A half-century ago, this would have been pooh-poohed as a “black box” theory, since one could not actually pinpoint this grammatical faculty in a specific part of the brain, or describe its functioning. But now things are different. Neurosurgeons [have now found that this] “black box” is situated in and around Broca’s area, on the left side of the forebrain. . . . 
Unlike Mr. Chomsky, Mr. Pinker firmly places the wiring of the brain for language within the framework of Darwinian natural selection and evolution. He effectively disposes of all claims that intelligent nonhuman primates like chimps have any abilities to learn and use language. It is not that chimps lack the vocal apparatus to speak; it is just that their brains are unable to produce or use grammar. On the other hand, the “language instinct,” when it first appeared among our most distant hominid ancestors, must have given them a selective reproductive advantage over their competitors (including the ancestral chimps). . . .
So according to Mr. Pinker, the roots of language must be in the genes, but there cannot be a “grammar gene” any more than there can be a gene for the heart or any other complex body structure. This proposition will undoubtedly raise the hackles of some behavioral psychologists and anthropologists, for it apparently contradicts the liberal idea that human behavior may be changed for the better by improvements in culture and environment, and it might seem to invite the twin bugaboos of biological determinism and racism. Yet Mr. Pinker stresses one point that should allay such fears. Even though there are 4,000 to 6,000 languages today, they are all sufficiently alike to be considered one language by an extraterrestrial observer. In other words, most of the diversity of the world’s cultures, so beloved to anthropologists, is superficial and minor compared to the similarities. Racial differences are literally only “skin deep.” The fundamental unity of humanity is the theme of Mr. Chomsky's universal grammar, and of this exciting book.

\vspace{1em}
\textbf{Question 1}

Which one of the following statements best summarises the author’s position about Pinker’s book?

\begin{enumerate}[label=\Alph*.]
    \item Culture and environment play a key role in shaping our acquisition of language.
    \item Anatomical developments like the voice box play a key role in determining language acquisition skills.
    \item The evolutionary and deterministic framework of Pinker’s book makes it racist.
    \item The universality of the “language instinct” counters claims that Pinker’s book is racist.
[](https://cracku.in/1-which-one-of-the-following-statements-best-summari-x-cat-2021-slot-3)
\end{enumerate}
\vspace{1em}
\textbf{Solution:}

_The fundamental unity of humanity is the theme of Mr. Chomsky's universal grammar, and of this exciting book._
Throughout the passage, the author seems to support the points made by Mr Pinker. The above line also shows that the opinion of the author towards the book is positive, and the author does not think that the book is racist in any way, but promotes unity and cohesion. Option D captures this point correctly and is the answer.
_So according to Mr. Pinker, the roots of language must be in the genes, but there _cannot be a “grammar gene” any more than there can be a gene for the heart or any other complex body structure_. This proposition will undoubtedly raise the hackles of some behavioral psychologists and anthropologists, _for it apparently contradicts the liberal idea that human behavior may be changed for the better by improvements in culture and environment_...._
_The book does not support that a complex anatomical structure like a 'voice box' plays a key role in determining language acquisition skills. Nor does it support the role of culture and environment in shaping human behaviour Options A and B are eliminated.____  
___
_Option C portrays the book as racist, which is directly in contradiction with the author's opinion. C is eliminated too.  
_

\end{problem}

\newpage
\begin{problem}{}{}
% Subject: varc
\textbf{Instructions}

Instructions   

The passage below is accompanied by a set of questions. Choose the best answer to each question.
Starting in 1957, [Noam Chomsky] proclaimed a new doctrine: Language, that most human of all attributes, was innate. The grammatical faculty was built into the infant brain, and your average 3-year-old was not a mere apprentice in the great enterprise of absorbing English from his or her parents, but a “linguistic genius.” Since this message was couched in terms of Chomskyan theoretical linguistics, in discourse so opaque that it was nearly incomprehensible even to some scholars, many people did not hear it. Now, in a brilliant, witty and altogether satisfying book, Mr. Chomsky's colleague Steven Pinker . . . has brought Mr. Chomsky's findings to everyman. In “The Language Instinct” he has gathered persuasive data from such diverse fields as cognitive neuroscience, developmental psychology and speech therapy to make his points, and when he disagrees with Mr. Chomsky he tells you so. . . .
For Mr. Chomsky and Mr. Pinker, somewhere in the human brain there is a complex set of neural circuits that have been programmed with “super-rules” (making up what Mr. Chomsky calls “universal grammar”), and that these rules are unconscious and instinctive. A half-century ago, this would have been pooh-poohed as a “black box” theory, since one could not actually pinpoint this grammatical faculty in a specific part of the brain, or describe its functioning. But now things are different. Neurosurgeons [have now found that this] “black box” is situated in and around Broca’s area, on the left side of the forebrain. . . . 
Unlike Mr. Chomsky, Mr. Pinker firmly places the wiring of the brain for language within the framework of Darwinian natural selection and evolution. He effectively disposes of all claims that intelligent nonhuman primates like chimps have any abilities to learn and use language. It is not that chimps lack the vocal apparatus to speak; it is just that their brains are unable to produce or use grammar. On the other hand, the “language instinct,” when it first appeared among our most distant hominid ancestors, must have given them a selective reproductive advantage over their competitors (including the ancestral chimps). . . .
So according to Mr. Pinker, the roots of language must be in the genes, but there cannot be a “grammar gene” any more than there can be a gene for the heart or any other complex body structure. This proposition will undoubtedly raise the hackles of some behavioral psychologists and anthropologists, for it apparently contradicts the liberal idea that human behavior may be changed for the better by improvements in culture and environment, and it might seem to invite the twin bugaboos of biological determinism and racism. Yet Mr. Pinker stresses one point that should allay such fears. Even though there are 4,000 to 6,000 languages today, they are all sufficiently alike to be considered one language by an extraterrestrial observer. In other words, most of the diversity of the world’s cultures, so beloved to anthropologists, is superficial and minor compared to the similarities. Racial differences are literally only “skin deep.” The fundamental unity of humanity is the theme of Mr. Chomsky's universal grammar, and of this exciting book.

\vspace{1em}
\textbf{Question 2}

According to the passage, all of the following are true about the language instinct EXCEPT that:

\begin{enumerate}[label=\Alph*.]
    \item all intelligent primates are gifted with it.
    \item developments in neuroscience have increased its acceptance.
    \item it confers an evolutionary reproductive advantage.
    \item not all intelligent primates are gifted with it.
[](https://cracku.in/2-according-to-the-passage-all-of-the-following-are--x-cat-2021-slot-3)
\end{enumerate}
\vspace{1em}
\textbf{Solution:}

_A half-century ago, this would have been pooh-poohed as a “black box” theory, since one could not actually pinpoint this grammatical faculty in a specific part of the brain, or describe its functioning. But now things are different. Neurosurgeons [have now found that this] “black box” is situated in and around Broca’s area, on the left side of the forebrain. . . ._
_On the other hand, the “language instinct,” when it first appeared among our most distant hominid ancestors, must have given them a selective reproductive advantage over their competitors (including the ancestral chimps). . . ._
__He effectively disposes of all claims that intelligent nonhuman primates like chimps have any abilities to learn and use language._  
_
The above excerpts provide support for Options B, C, and D respectively. Option A is in direct contradiction with Option D, and hence, is the answer.

\end{problem}

\newpage
\begin{problem}{}{}
% Subject: varc
\textbf{Instructions}

Instructions   

The passage below is accompanied by a set of questions. Choose the best answer to each question.
Starting in 1957, [Noam Chomsky] proclaimed a new doctrine: Language, that most human of all attributes, was innate. The grammatical faculty was built into the infant brain, and your average 3-year-old was not a mere apprentice in the great enterprise of absorbing English from his or her parents, but a “linguistic genius.” Since this message was couched in terms of Chomskyan theoretical linguistics, in discourse so opaque that it was nearly incomprehensible even to some scholars, many people did not hear it. Now, in a brilliant, witty and altogether satisfying book, Mr. Chomsky's colleague Steven Pinker . . . has brought Mr. Chomsky's findings to everyman. In “The Language Instinct” he has gathered persuasive data from such diverse fields as cognitive neuroscience, developmental psychology and speech therapy to make his points, and when he disagrees with Mr. Chomsky he tells you so. . . .
For Mr. Chomsky and Mr. Pinker, somewhere in the human brain there is a complex set of neural circuits that have been programmed with “super-rules” (making up what Mr. Chomsky calls “universal grammar”), and that these rules are unconscious and instinctive. A half-century ago, this would have been pooh-poohed as a “black box” theory, since one could not actually pinpoint this grammatical faculty in a specific part of the brain, or describe its functioning. But now things are different. Neurosurgeons [have now found that this] “black box” is situated in and around Broca’s area, on the left side of the forebrain. . . . 
Unlike Mr. Chomsky, Mr. Pinker firmly places the wiring of the brain for language within the framework of Darwinian natural selection and evolution. He effectively disposes of all claims that intelligent nonhuman primates like chimps have any abilities to learn and use language. It is not that chimps lack the vocal apparatus to speak; it is just that their brains are unable to produce or use grammar. On the other hand, the “language instinct,” when it first appeared among our most distant hominid ancestors, must have given them a selective reproductive advantage over their competitors (including the ancestral chimps). . . .
So according to Mr. Pinker, the roots of language must be in the genes, but there cannot be a “grammar gene” any more than there can be a gene for the heart or any other complex body structure. This proposition will undoubtedly raise the hackles of some behavioral psychologists and anthropologists, for it apparently contradicts the liberal idea that human behavior may be changed for the better by improvements in culture and environment, and it might seem to invite the twin bugaboos of biological determinism and racism. Yet Mr. Pinker stresses one point that should allay such fears. Even though there are 4,000 to 6,000 languages today, they are all sufficiently alike to be considered one language by an extraterrestrial observer. In other words, most of the diversity of the world’s cultures, so beloved to anthropologists, is superficial and minor compared to the similarities. Racial differences are literally only “skin deep.” The fundamental unity of humanity is the theme of Mr. Chomsky's universal grammar, and of this exciting book.

\vspace{1em}
\textbf{Question 3}

On the basis of the information in the passage, Pinker and Chomsky may disagree with each other on which one of the following points?

\begin{enumerate}[label=\Alph*.]
    \item The language instinct.
    \item The inborn language acquisition skills of humans.
    \item The Darwinian explanatory paradigm for language.
    \item The possibility of a universal grammar.
[](https://cracku.in/3-on-the-basis-of-the-information-in-the-passage-pin-x-cat-2021-slot-3)
\end{enumerate}
\vspace{1em}
\textbf{Solution:}

_Unlike Mr. Chomsky, Mr. Pinker firmly places the wiring of the brain for language within the framework of Darwinian natural selection and evolution._
The passage suggests that Mr. Pinker and Mr. Chomsky agree on almost all topics. However, the above line indicates that they both disagreed on the application of the Darwinian framework to explain language instinct. Where Mr. Pinker was in favour of the same, Mr. Chomsky was against. Hence, Option C is the answer.

\end{problem}

\newpage
\begin{problem}{}{}
% Subject: varc
\textbf{Instructions}

Instructions   

The passage below is accompanied by a set of questions. Choose the best answer to each question.
Starting in 1957, [Noam Chomsky] proclaimed a new doctrine: Language, that most human of all attributes, was innate. The grammatical faculty was built into the infant brain, and your average 3-year-old was not a mere apprentice in the great enterprise of absorbing English from his or her parents, but a “linguistic genius.” Since this message was couched in terms of Chomskyan theoretical linguistics, in discourse so opaque that it was nearly incomprehensible even to some scholars, many people did not hear it. Now, in a brilliant, witty and altogether satisfying book, Mr. Chomsky's colleague Steven Pinker . . . has brought Mr. Chomsky's findings to everyman. In “The Language Instinct” he has gathered persuasive data from such diverse fields as cognitive neuroscience, developmental psychology and speech therapy to make his points, and when he disagrees with Mr. Chomsky he tells you so. . . .
For Mr. Chomsky and Mr. Pinker, somewhere in the human brain there is a complex set of neural circuits that have been programmed with “super-rules” (making up what Mr. Chomsky calls “universal grammar”), and that these rules are unconscious and instinctive. A half-century ago, this would have been pooh-poohed as a “black box” theory, since one could not actually pinpoint this grammatical faculty in a specific part of the brain, or describe its functioning. But now things are different. Neurosurgeons [have now found that this] “black box” is situated in and around Broca’s area, on the left side of the forebrain. . . . 
Unlike Mr. Chomsky, Mr. Pinker firmly places the wiring of the brain for language within the framework of Darwinian natural selection and evolution. He effectively disposes of all claims that intelligent nonhuman primates like chimps have any abilities to learn and use language. It is not that chimps lack the vocal apparatus to speak; it is just that their brains are unable to produce or use grammar. On the other hand, the “language instinct,” when it first appeared among our most distant hominid ancestors, must have given them a selective reproductive advantage over their competitors (including the ancestral chimps). . . .
So according to Mr. Pinker, the roots of language must be in the genes, but there cannot be a “grammar gene” any more than there can be a gene for the heart or any other complex body structure. This proposition will undoubtedly raise the hackles of some behavioral psychologists and anthropologists, for it apparently contradicts the liberal idea that human behavior may be changed for the better by improvements in culture and environment, and it might seem to invite the twin bugaboos of biological determinism and racism. Yet Mr. Pinker stresses one point that should allay such fears. Even though there are 4,000 to 6,000 languages today, they are all sufficiently alike to be considered one language by an extraterrestrial observer. In other words, most of the diversity of the world’s cultures, so beloved to anthropologists, is superficial and minor compared to the similarities. Racial differences are literally only “skin deep.” The fundamental unity of humanity is the theme of Mr. Chomsky's universal grammar, and of this exciting book.

\vspace{1em}
\textbf{Question 4}

From the passage, it can be inferred that all of the following are true about Pinker’s book, “The Language Instinct”, EXCEPT that Pinker:

\begin{enumerate}[label=\Alph*.]
    \item draws extensively from Chomsky’s propositions.
    \item disagrees with Chomsky on certain grounds.
    \item draws from behavioural psychology theories.
    \item writes in a different style from Chomsky.
[](https://cracku.in/4-from-the-passage-it-can-be-inferred-that-all-of-th-x-cat-2021-slot-3)
\end{enumerate}
\vspace{1em}
\textbf{Solution:}

_Since this message was couched in terms of Chomskyan theoretical linguistics, in discourse so opaque that it was nearly incomprehensible even to some scholars, many people did not hear it. Now, in a brilliant, witty and altogether satisfying book, Mr. Chomsky's colleague Steven Pinker . . . has brought Mr. Chomsky's findings to everyman._
From the above excerpt, it is clear that Mr. Pinker's style of writing is much more comprehensible to the common man. Hence, their writing styles are quite different. Also, the above excerpt mentions that the book brings Mr. Chomsky's findings to everyman, hence, it is clear that it draws heavily from the findings. Options A and D are eliminated.
_Unlike Mr. Chomsky, Mr. Pinker firmly places the wiring of the brain for language within the framework of Darwinian natural selection and evolution._
The above excerpt shows that they both disagreed on a certain point. Hence, Option B is eliminated too.
Option C finds no mention in the passage, hence, is the answer.

Instructions   

**The passage below is accompanied by a set of questions. Choose the best answer to each question.**
Today we can hardly conceive of ourselves without an unconscious. Yet between 1700 and 1900, this notion developed as a genuinely original thought. The “unconscious” burst the shell of conventional language, coined as it had been to embody the fleeting ideas and the shifting conceptions of several generations until, finally, it became fixed and defined in specialized terms within the realm of medical psychology and Freudian psychoanalysis.
The vocabulary concerning the soul and the mind increased enormously in the course of the nineteenth century. The enrichments of literary and intellectual language led to an altered understanding of the meanings that underlie time-honored expressions and traditional catchwords. At the same time, once coined, powerful new ideas attracted to themselves a whole host of seemingly unrelated issues, practices, and experiences, creating a peculiar network of preoccupations that as a group had not existed before. The drawn-out attempt to approach and define the unconscious brought together the spiritualist and the psychical researcher of borderline phenomena (such as apparitions, spectral illusions, haunted houses, mediums, trance, automatic writing); the psychiatrist or alienist probing the nature of mental disease, of abnormal ideation, hallucination, delirium, melancholia, mania; the surgeon performing operations with the aid of hypnotism; the magnetizer claiming to correct the disequilibrium in the universal flow of magnetic fluids but who soon came to be regarded as a clever manipulator of the imagination; the physiologist and the physician who puzzled over sleep, dreams, sleepwalking, anesthesia, the influence of the mind on the body in health and disease; the neurologist concerned with the functions of the brain and the physiological basis of mental life; the philosopher interested in the will, the emotions, consciousness, knowledge, imagination and the creative genius; and, last but not least, the psychologist.
Significantly, most if not all of these practices (for example, hypnotism in surgery or psychological magnetism) originated in the waning years of the eighteenth century and during the early decades of the nineteenth century, as did some of the disciplines (such as psychology and psychical research). The majority of topics too were either new or assumed hitherto unknown colors. Thus, before 1790, few if any spoke, in medical terms, of the affinity between creative genius and the hallucinations of the insane . . .
Striving vaguely and independently to give expression to a latent conception, various lines of thought can be brought together by some novel term. The new concept then serves as a kind of resting place or stocktaking in the development of ideas, giving satisfaction and a stimulus for further discussion or speculation. Thus, the massive introduction of the term unconscious by Hartmann in 1869 appeared to focalize many stray thoughts, affording a temporary feeling that a crucial step had been taken forward, a comprehensive knowledge gained, a knowledge that required only further elaboration, explication, and unfolding in order to bring in a bounty of higher understanding. Ultimately, Hartmann’s attempt at defining the unconscious proved fruitless because he extended its reach into every realm of organic and inorganic, spiritual, intellectual, and instinctive existence, severely diluting the precision and compromising the impact of the concept.

\end{problem}

\newpage
\begin{problem}{}{}
% Subject: varc
\textbf{Instructions}

Instructions   

The passage below is accompanied by a set of questions. Choose the best answer to each question.
Starting in 1957, [Noam Chomsky] proclaimed a new doctrine: Language, that most human of all attributes, was innate. The grammatical faculty was built into the infant brain, and your average 3-year-old was not a mere apprentice in the great enterprise of absorbing English from his or her parents, but a “linguistic genius.” Since this message was couched in terms of Chomskyan theoretical linguistics, in discourse so opaque that it was nearly incomprehensible even to some scholars, many people did not hear it. Now, in a brilliant, witty and altogether satisfying book, Mr. Chomsky's colleague Steven Pinker . . . has brought Mr. Chomsky's findings to everyman. In “The Language Instinct” he has gathered persuasive data from such diverse fields as cognitive neuroscience, developmental psychology and speech therapy to make his points, and when he disagrees with Mr. Chomsky he tells you so. . . .
For Mr. Chomsky and Mr. Pinker, somewhere in the human brain there is a complex set of neural circuits that have been programmed with “super-rules” (making up what Mr. Chomsky calls “universal grammar”), and that these rules are unconscious and instinctive. A half-century ago, this would have been pooh-poohed as a “black box” theory, since one could not actually pinpoint this grammatical faculty in a specific part of the brain, or describe its functioning. But now things are different. Neurosurgeons [have now found that this] “black box” is situated in and around Broca’s area, on the left side of the forebrain. . . . 
Unlike Mr. Chomsky, Mr. Pinker firmly places the wiring of the brain for language within the framework of Darwinian natural selection and evolution. He effectively disposes of all claims that intelligent nonhuman primates like chimps have any abilities to learn and use language. It is not that chimps lack the vocal apparatus to speak; it is just that their brains are unable to produce or use grammar. On the other hand, the “language instinct,” when it first appeared among our most distant hominid ancestors, must have given them a selective reproductive advantage over their competitors (including the ancestral chimps). . . .
So according to Mr. Pinker, the roots of language must be in the genes, but there cannot be a “grammar gene” any more than there can be a gene for the heart or any other complex body structure. This proposition will undoubtedly raise the hackles of some behavioral psychologists and anthropologists, for it apparently contradicts the liberal idea that human behavior may be changed for the better by improvements in culture and environment, and it might seem to invite the twin bugaboos of biological determinism and racism. Yet Mr. Pinker stresses one point that should allay such fears. Even though there are 4,000 to 6,000 languages today, they are all sufficiently alike to be considered one language by an extraterrestrial observer. In other words, most of the diversity of the world’s cultures, so beloved to anthropologists, is superficial and minor compared to the similarities. Racial differences are literally only “skin deep.” The fundamental unity of humanity is the theme of Mr. Chomsky's universal grammar, and of this exciting book.

\vspace{1em}
\textbf{Question 5}

Which one of the following statements best describes what the passage is about?

\begin{enumerate}[label=\Alph*.]
    \item The discovery of the unconscious as a part of the human mind.
    \item The growing vocabulary of the soul and the mind, as diverse processes.
    \item The collating of diverse ideas under the single term: unconscious.
    \item The identification of the unconscious as an object of psychical research.
[](https://cracku.in/5-which-one-of-the-following-statements-best-describ-x-cat-2021-slot-3)
\end{enumerate}
\vspace{1em}
\textbf{Solution:}

The passage starts by highlighting that the term 'unconscious', widely held today, came in conception not long ago. With the coining of this term, many unrelated activities/ideas found a common umbrella under which they could be categorized and also allowed them to prosper. The author then writes the following line, which gives us a clear conception of the main theme:
_Thus, the massive introduction of the term unconscious by Hartmann in 1869 appeared to focalize many stray thoughts, affording a temporary feeling that a crucial step had been taken forward, a comprehensive knowledge gained a knowledge that required only further elaboration, explication, and unfolding in order to bring in a bounty of higher understanding._
Thus, the passage is about the assembly of many stray thoughts under the banner of the unconscious. Option C perfectly captures this, and hence, is the answer.
The author does not primarily deal with the unconscious as a part of the mind. Nor does he focus upon the expansion of the vocabulary of the mind and the soul. Thus, Options A and B can be rejected.
'Psychical research' is not the main focus of the passage. The author says that the term allowed certain 'psychic' activities to flourish. He does not focus on the term as an object of psychical research. Hence, Option D can be eliminated too.

\end{problem}

\newpage
\begin{problem}{}{}
% Subject: varc
\textbf{Instructions}

Instructions   

The passage below is accompanied by a set of questions. Choose the best answer to each question.
Starting in 1957, [Noam Chomsky] proclaimed a new doctrine: Language, that most human of all attributes, was innate. The grammatical faculty was built into the infant brain, and your average 3-year-old was not a mere apprentice in the great enterprise of absorbing English from his or her parents, but a “linguistic genius.” Since this message was couched in terms of Chomskyan theoretical linguistics, in discourse so opaque that it was nearly incomprehensible even to some scholars, many people did not hear it. Now, in a brilliant, witty and altogether satisfying book, Mr. Chomsky's colleague Steven Pinker . . . has brought Mr. Chomsky's findings to everyman. In “The Language Instinct” he has gathered persuasive data from such diverse fields as cognitive neuroscience, developmental psychology and speech therapy to make his points, and when he disagrees with Mr. Chomsky he tells you so. . . .
For Mr. Chomsky and Mr. Pinker, somewhere in the human brain there is a complex set of neural circuits that have been programmed with “super-rules” (making up what Mr. Chomsky calls “universal grammar”), and that these rules are unconscious and instinctive. A half-century ago, this would have been pooh-poohed as a “black box” theory, since one could not actually pinpoint this grammatical faculty in a specific part of the brain, or describe its functioning. But now things are different. Neurosurgeons [have now found that this] “black box” is situated in and around Broca’s area, on the left side of the forebrain. . . . 
Unlike Mr. Chomsky, Mr. Pinker firmly places the wiring of the brain for language within the framework of Darwinian natural selection and evolution. He effectively disposes of all claims that intelligent nonhuman primates like chimps have any abilities to learn and use language. It is not that chimps lack the vocal apparatus to speak; it is just that their brains are unable to produce or use grammar. On the other hand, the “language instinct,” when it first appeared among our most distant hominid ancestors, must have given them a selective reproductive advantage over their competitors (including the ancestral chimps). . . .
So according to Mr. Pinker, the roots of language must be in the genes, but there cannot be a “grammar gene” any more than there can be a gene for the heart or any other complex body structure. This proposition will undoubtedly raise the hackles of some behavioral psychologists and anthropologists, for it apparently contradicts the liberal idea that human behavior may be changed for the better by improvements in culture and environment, and it might seem to invite the twin bugaboos of biological determinism and racism. Yet Mr. Pinker stresses one point that should allay such fears. Even though there are 4,000 to 6,000 languages today, they are all sufficiently alike to be considered one language by an extraterrestrial observer. In other words, most of the diversity of the world’s cultures, so beloved to anthropologists, is superficial and minor compared to the similarities. Racial differences are literally only “skin deep.” The fundamental unity of humanity is the theme of Mr. Chomsky's universal grammar, and of this exciting book.

\vspace{1em}
\textbf{Question 6}

“The enrichments of literary and intellectual language led to an altered understanding of the meanings that underlie time-honored expressions and traditional catchwords.” Which one of the following interpretations of this sentence would be closest in meaning to the original?

\begin{enumerate}[label=\Alph*.]
    \item All of the options listed here.
    \item Time-honored expressions and traditional catchwords were enriched by literary and intellectual language.
    \item Literary and intellectual language was altered by time-honored expressions and traditional catchwords.
    \item The meanings of time-honored expressions were changed by innovations in literary and intellectual language.
[](https://cracku.in/6-the-enrichments-of-literary-and-intellectual-langu-x-cat-2021-slot-3)
\end{enumerate}
\vspace{1em}
\textbf{Solution:}

Let us try to break the sentence down and interpret its meaning:
_“The enrichments of literary and intellectual language led to an altered understanding of the meanings that underlie time-honored expressions and traditional catchwords.”_
In simple words |Enrichments of language| led to |change in understanding| of | time-honoured expressions|. 
In the context of the passage, the line means that when the terms related to 'the unconscious' were coined, they enriched the vocabulary of the language and this, in turn, changes the meanings of many old expressions related to this term.
Option D comes the closest in capturing the meaning, and hence, is the answer.
B: The meanings of the catchwords were altered. They were not enriched. Can be eliminated.
C: The catchwords did not cause a change. Their own meaning was changed. Can be eliminated.

\end{problem}

\newpage
\begin{problem}{}{}
% Subject: varc
\textbf{Instructions}

Instructions   

The passage below is accompanied by a set of questions. Choose the best answer to each question.
Starting in 1957, [Noam Chomsky] proclaimed a new doctrine: Language, that most human of all attributes, was innate. The grammatical faculty was built into the infant brain, and your average 3-year-old was not a mere apprentice in the great enterprise of absorbing English from his or her parents, but a “linguistic genius.” Since this message was couched in terms of Chomskyan theoretical linguistics, in discourse so opaque that it was nearly incomprehensible even to some scholars, many people did not hear it. Now, in a brilliant, witty and altogether satisfying book, Mr. Chomsky's colleague Steven Pinker . . . has brought Mr. Chomsky's findings to everyman. In “The Language Instinct” he has gathered persuasive data from such diverse fields as cognitive neuroscience, developmental psychology and speech therapy to make his points, and when he disagrees with Mr. Chomsky he tells you so. . . .
For Mr. Chomsky and Mr. Pinker, somewhere in the human brain there is a complex set of neural circuits that have been programmed with “super-rules” (making up what Mr. Chomsky calls “universal grammar”), and that these rules are unconscious and instinctive. A half-century ago, this would have been pooh-poohed as a “black box” theory, since one could not actually pinpoint this grammatical faculty in a specific part of the brain, or describe its functioning. But now things are different. Neurosurgeons [have now found that this] “black box” is situated in and around Broca’s area, on the left side of the forebrain. . . . 
Unlike Mr. Chomsky, Mr. Pinker firmly places the wiring of the brain for language within the framework of Darwinian natural selection and evolution. He effectively disposes of all claims that intelligent nonhuman primates like chimps have any abilities to learn and use language. It is not that chimps lack the vocal apparatus to speak; it is just that their brains are unable to produce or use grammar. On the other hand, the “language instinct,” when it first appeared among our most distant hominid ancestors, must have given them a selective reproductive advantage over their competitors (including the ancestral chimps). . . .
So according to Mr. Pinker, the roots of language must be in the genes, but there cannot be a “grammar gene” any more than there can be a gene for the heart or any other complex body structure. This proposition will undoubtedly raise the hackles of some behavioral psychologists and anthropologists, for it apparently contradicts the liberal idea that human behavior may be changed for the better by improvements in culture and environment, and it might seem to invite the twin bugaboos of biological determinism and racism. Yet Mr. Pinker stresses one point that should allay such fears. Even though there are 4,000 to 6,000 languages today, they are all sufficiently alike to be considered one language by an extraterrestrial observer. In other words, most of the diversity of the world’s cultures, so beloved to anthropologists, is superficial and minor compared to the similarities. Racial differences are literally only “skin deep.” The fundamental unity of humanity is the theme of Mr. Chomsky's universal grammar, and of this exciting book.

\vspace{1em}
\textbf{Question 7}

Which one of the following sets of words is closest to mapping the main arguments of the passage?

\begin{enumerate}[label=\Alph*.]
    \item Unconscious; Latent conception; Dreams.
    \item Literary language; Unconscious; Insanity.
    \item Language; Unconscious; Psychoanalysis.
    \item Imagination; Magnetism; Psychiatry.
[](https://cracku.in/7-which-one-of-the-following-sets-of-words-is-closes-x-cat-2021-slot-3)
\end{enumerate}
\vspace{1em}
\textbf{Solution:}

Unconscious is the primary focus of the passage. Since D does not have that as a main point, it can be eliminated.
Dreams find a single, small mention as an example in the passage. Hence, Option A can be eliminated too.
Insanity finds a small mention in the passage and is not a main point. Hence, Option B is incorrect.
The author initially deals with how the enrichment of vocabulary on the matter of unconscious has a deep effect and how this later became a subject of psychoanalysis. Hence, Option C is the correct answer.

\end{problem}

\newpage
\begin{problem}{}{}
% Subject: varc
\textbf{Instructions}

Instructions   

The passage below is accompanied by a set of questions. Choose the best answer to each question.
Starting in 1957, [Noam Chomsky] proclaimed a new doctrine: Language, that most human of all attributes, was innate. The grammatical faculty was built into the infant brain, and your average 3-year-old was not a mere apprentice in the great enterprise of absorbing English from his or her parents, but a “linguistic genius.” Since this message was couched in terms of Chomskyan theoretical linguistics, in discourse so opaque that it was nearly incomprehensible even to some scholars, many people did not hear it. Now, in a brilliant, witty and altogether satisfying book, Mr. Chomsky's colleague Steven Pinker . . . has brought Mr. Chomsky's findings to everyman. In “The Language Instinct” he has gathered persuasive data from such diverse fields as cognitive neuroscience, developmental psychology and speech therapy to make his points, and when he disagrees with Mr. Chomsky he tells you so. . . .
For Mr. Chomsky and Mr. Pinker, somewhere in the human brain there is a complex set of neural circuits that have been programmed with “super-rules” (making up what Mr. Chomsky calls “universal grammar”), and that these rules are unconscious and instinctive. A half-century ago, this would have been pooh-poohed as a “black box” theory, since one could not actually pinpoint this grammatical faculty in a specific part of the brain, or describe its functioning. But now things are different. Neurosurgeons [have now found that this] “black box” is situated in and around Broca’s area, on the left side of the forebrain. . . . 
Unlike Mr. Chomsky, Mr. Pinker firmly places the wiring of the brain for language within the framework of Darwinian natural selection and evolution. He effectively disposes of all claims that intelligent nonhuman primates like chimps have any abilities to learn and use language. It is not that chimps lack the vocal apparatus to speak; it is just that their brains are unable to produce or use grammar. On the other hand, the “language instinct,” when it first appeared among our most distant hominid ancestors, must have given them a selective reproductive advantage over their competitors (including the ancestral chimps). . . .
So according to Mr. Pinker, the roots of language must be in the genes, but there cannot be a “grammar gene” any more than there can be a gene for the heart or any other complex body structure. This proposition will undoubtedly raise the hackles of some behavioral psychologists and anthropologists, for it apparently contradicts the liberal idea that human behavior may be changed for the better by improvements in culture and environment, and it might seem to invite the twin bugaboos of biological determinism and racism. Yet Mr. Pinker stresses one point that should allay such fears. Even though there are 4,000 to 6,000 languages today, they are all sufficiently alike to be considered one language by an extraterrestrial observer. In other words, most of the diversity of the world’s cultures, so beloved to anthropologists, is superficial and minor compared to the similarities. Racial differences are literally only “skin deep.” The fundamental unity of humanity is the theme of Mr. Chomsky's universal grammar, and of this exciting book.

\vspace{1em}
\textbf{Question 8}

All of the following statements may be considered valid inferences from the passage, EXCEPT:

\begin{enumerate}[label=\Alph*.]
    \item Without the linguistic developments of the nineteenth century, the growth of understanding of the soul and the mind may not have happened.
    \item Eighteenth century thinkers were the first to perceive a connection between creative genius and insanity.
    \item New conceptions in the nineteenth century could provide new knowledge because of the establishment of fields such as anaesthesiology.
    \item Unrelated practices began to be treated as related to each other, as knowledge of the mind grew in the nineteenth century.
[](https://cracku.in/8-all-of-the-following-statements-may-be-considered--x-cat-2021-slot-3)
\end{enumerate}
\vspace{1em}
\textbf{Solution:}

_The “unconscious” burst the shell of conventional language, coined as it had been to embody the fleeting ideas and the shifting conceptions of several generations until, finally, it became fixed and defined in specialized terms within the realm of medical psychology and Freudian psychoanalysis._
In the passage, the author has clearly outlined the importance of linguistic developments in helping the knowledge of the field grow. Since the option is not extreme in certainty ('may' not have happened), Option A can be inferred.
_Significantly, most if not all of these practices (for example, hypnotism in surgery or psychological magnetism) originated in the waning years of the eighteenth century and during the early decades of the nineteenth century, as did some of the disciplines (such as psychology and psychical research). The majority of topics too were either new or assumed hitherto unknown colors. Thus, before 1790, few if any spoke, in medical terms, of the affinity between creative genius and the hallucinations of the insane . . ._
From the above excerpt, we can infer that the affinity between genius and insanity was not looked into before the 18th century.
_At the same time, once coined, powerful new ideas attracted to themselves a whole host of seemingly unrelated issues, practices, and experiences, creating a peculiar network of preoccupations that as a group had not existed before._
The above excerpt and the examples the author provides after this excerpt can help us infer that as the knowledge of the mind grew, unrelated activities found a common title. Option D can be inferred.
The passage does not imply anywhere that the new conceptions were able to provide new knowledge only because some fields were established. Option C is out of the scope of the passage and cannot be inferred.

Instructions   

The passage below is accompanied by a set of questions. Choose the best answer to each question.
Back in the early 2000s, an awesome thing happened in the New X-Men comics. Our mutant heroes had been battling giant robots called Sentinels for years, but suddenly these mechanical overlords spawned a new threat: Nano-Sentinels! Not content to rule Earth with their metal fists, these tiny robots invaded our bodies at the microscopic level. Infected humans were slowly converted into machines, cell by cell.
Now, a new wave of extremely odd robots is making at least part of the Nano-Sentinels story come true. Using exotic fabrication materials like squishy hydrogels and elastic polymers, researchers are making autonomous devices that are often tiny and that could turn out to be more powerful than an army of Terminators. Some are 1-centimetre blobs that can skate over water. Others are flat sheets that can roll themselves into tubes, or matchstick-sized plastic coils that act as powerful muscles. No, they won’t be invading our bodies and turning us into Sentinels - which I personally find a little disappointing - but some of them could one day swim through our bloodstream to heal us. They could also clean up pollutants in water or fold themselves into different kinds of vehicles for us to drive. . . . 
Unlike a traditional robot, which is made of mechanical parts, these new kinds of robots are made from molecular parts. The principle is the same: both are devices that can move around and do things independently. But a robot made from smart materials might be nothing more than a pink drop of hydrogel. Instead of gears and wires, it’s assembled from two kinds of molecules - some that love water and some that avoid it - which interact to allow the bot to skate on top of a pond.
Sometimes these materials are used to enhance more conventional robots. One team of researchers, for example, has developed a different kind of hydrogel that becomes sticky when exposed to a low-voltage zap of electricity and then stops being sticky when the electricity is switched off. This putty-like gel can be pasted right onto the feet or wheels of a robot. When the robot wants to climb a sheer wall or scoot across the ceiling, it can activate its sticky feet with a few volts. Once it is back on a flat surface again, the robot turns off the adhesive like a light switch. 
Robots that are wholly or partly made of gloop aren’t the future that I was promised in science fiction. But it’s definitely the future I want. I’m especially keen on the nanometre-scale “soft robots” that could one day swim through our bodies. Metin Sitti, a director at the Max Planck Institute for Intelligent Systems in Germany, worked with colleagues to prototype these tiny, synthetic beasts using various stretchy materials, such as simple rubber, and seeding them with magnetic microparticles. They are assembled into a finished shape by applying magnetic fields. The results look like flowers or geometric shapes made from Tinkertoy ball and stick modelling kits. They’re guided through tubes of fluid using magnets, and can even stop and cling to the sides of a tube.

\end{problem}

\newpage
\begin{problem}{}{}
% Subject: varc
\textbf{Instructions}

Instructions   

The passage below is accompanied by a set of questions. Choose the best answer to each question.
Starting in 1957, [Noam Chomsky] proclaimed a new doctrine: Language, that most human of all attributes, was innate. The grammatical faculty was built into the infant brain, and your average 3-year-old was not a mere apprentice in the great enterprise of absorbing English from his or her parents, but a “linguistic genius.” Since this message was couched in terms of Chomskyan theoretical linguistics, in discourse so opaque that it was nearly incomprehensible even to some scholars, many people did not hear it. Now, in a brilliant, witty and altogether satisfying book, Mr. Chomsky's colleague Steven Pinker . . . has brought Mr. Chomsky's findings to everyman. In “The Language Instinct” he has gathered persuasive data from such diverse fields as cognitive neuroscience, developmental psychology and speech therapy to make his points, and when he disagrees with Mr. Chomsky he tells you so. . . .
For Mr. Chomsky and Mr. Pinker, somewhere in the human brain there is a complex set of neural circuits that have been programmed with “super-rules” (making up what Mr. Chomsky calls “universal grammar”), and that these rules are unconscious and instinctive. A half-century ago, this would have been pooh-poohed as a “black box” theory, since one could not actually pinpoint this grammatical faculty in a specific part of the brain, or describe its functioning. But now things are different. Neurosurgeons [have now found that this] “black box” is situated in and around Broca’s area, on the left side of the forebrain. . . . 
Unlike Mr. Chomsky, Mr. Pinker firmly places the wiring of the brain for language within the framework of Darwinian natural selection and evolution. He effectively disposes of all claims that intelligent nonhuman primates like chimps have any abilities to learn and use language. It is not that chimps lack the vocal apparatus to speak; it is just that their brains are unable to produce or use grammar. On the other hand, the “language instinct,” when it first appeared among our most distant hominid ancestors, must have given them a selective reproductive advantage over their competitors (including the ancestral chimps). . . .
So according to Mr. Pinker, the roots of language must be in the genes, but there cannot be a “grammar gene” any more than there can be a gene for the heart or any other complex body structure. This proposition will undoubtedly raise the hackles of some behavioral psychologists and anthropologists, for it apparently contradicts the liberal idea that human behavior may be changed for the better by improvements in culture and environment, and it might seem to invite the twin bugaboos of biological determinism and racism. Yet Mr. Pinker stresses one point that should allay such fears. Even though there are 4,000 to 6,000 languages today, they are all sufficiently alike to be considered one language by an extraterrestrial observer. In other words, most of the diversity of the world’s cultures, so beloved to anthropologists, is superficial and minor compared to the similarities. Racial differences are literally only “skin deep.” The fundamental unity of humanity is the theme of Mr. Chomsky's universal grammar, and of this exciting book.

\vspace{1em}
\textbf{Question 9}

Which one of the following statements best captures the sense of the first paragraph?

\begin{enumerate}[label=\Alph*.]
    \item People who were infected by Nano-Sentinel robots became mutants who were called X-Men.
    \item Tiny sentinels called X-Men infected people, turning them into mutant robot overlords.
    \item None of the options listed here.
    \item The X-Men were mutant heroes who now had to battle tiny robots called Nano-Sentinels
[](https://cracku.in/9-which-one-of-the-following-statements-best-capture-x-cat-2021-slot-3)
\end{enumerate}
\vspace{1em}
\textbf{Solution:}

_Back in the early 2000s, an awesome thing happened in the New X-Men comics. Our mutant heroes had been battling giant robots called Sentinels for years, but suddenly these mechanical overlords spawned a new threat: Nano-Sentinels! Not content to rule Earth with their metal fists, these tiny robots invaded our bodies at the microscopic level. Infected humans were slowly converted into machines, cell by cell_ _._
The first paragraph talks about the X-men comics, in which the mutant heroes, that X-Men, has been battling giant robots called sentinels. But these Sentinels then developed Nano-Sentinels, which could invade bodies at the microscopic level, and the heroes would now have to fight them too. Option D perfectly captures this, and hence, is the answer.
Option A is incorrect. X-men were battling the Sentinels before the invention of Nano-Sentinels. Hence, the origin of X-men is different.
Option B is incorrect. The mechanical overlords made Nano-Sentinels to convert people into machines. It has not been said that the people were converted into the mechanical lords themselves.

\end{problem}

\newpage
\begin{problem}{}{}
% Subject: varc
\textbf{Instructions}

Instructions   

The passage below is accompanied by a set of questions. Choose the best answer to each question.
Starting in 1957, [Noam Chomsky] proclaimed a new doctrine: Language, that most human of all attributes, was innate. The grammatical faculty was built into the infant brain, and your average 3-year-old was not a mere apprentice in the great enterprise of absorbing English from his or her parents, but a “linguistic genius.” Since this message was couched in terms of Chomskyan theoretical linguistics, in discourse so opaque that it was nearly incomprehensible even to some scholars, many people did not hear it. Now, in a brilliant, witty and altogether satisfying book, Mr. Chomsky's colleague Steven Pinker . . . has brought Mr. Chomsky's findings to everyman. In “The Language Instinct” he has gathered persuasive data from such diverse fields as cognitive neuroscience, developmental psychology and speech therapy to make his points, and when he disagrees with Mr. Chomsky he tells you so. . . .
For Mr. Chomsky and Mr. Pinker, somewhere in the human brain there is a complex set of neural circuits that have been programmed with “super-rules” (making up what Mr. Chomsky calls “universal grammar”), and that these rules are unconscious and instinctive. A half-century ago, this would have been pooh-poohed as a “black box” theory, since one could not actually pinpoint this grammatical faculty in a specific part of the brain, or describe its functioning. But now things are different. Neurosurgeons [have now found that this] “black box” is situated in and around Broca’s area, on the left side of the forebrain. . . . 
Unlike Mr. Chomsky, Mr. Pinker firmly places the wiring of the brain for language within the framework of Darwinian natural selection and evolution. He effectively disposes of all claims that intelligent nonhuman primates like chimps have any abilities to learn and use language. It is not that chimps lack the vocal apparatus to speak; it is just that their brains are unable to produce or use grammar. On the other hand, the “language instinct,” when it first appeared among our most distant hominid ancestors, must have given them a selective reproductive advantage over their competitors (including the ancestral chimps). . . .
So according to Mr. Pinker, the roots of language must be in the genes, but there cannot be a “grammar gene” any more than there can be a gene for the heart or any other complex body structure. This proposition will undoubtedly raise the hackles of some behavioral psychologists and anthropologists, for it apparently contradicts the liberal idea that human behavior may be changed for the better by improvements in culture and environment, and it might seem to invite the twin bugaboos of biological determinism and racism. Yet Mr. Pinker stresses one point that should allay such fears. Even though there are 4,000 to 6,000 languages today, they are all sufficiently alike to be considered one language by an extraterrestrial observer. In other words, most of the diversity of the world’s cultures, so beloved to anthropologists, is superficial and minor compared to the similarities. Racial differences are literally only “skin deep.” The fundamental unity of humanity is the theme of Mr. Chomsky's universal grammar, and of this exciting book.

\vspace{1em}
\textbf{Question 10}

Which one of the following scenarios, if false, could be seen as supporting the passage?

\begin{enumerate}[label=\Alph*.]
    \item Robots made from smart materials are likely to become part of our everyday lives in the future.
    \item There are two kinds of molecules used to make some nano-robots: one that reacts positively to water and the other negatively.
    \item Some hydrogels turn sticky when an electric current is passed through them; this potentially has very useful applications.
    \item Nano-Sentinel-like robots are likely to be used to inject people to convert them into robots, cell by cell.
[](https://cracku.in/10-which-one-of-the-following-scenarios-if-false-coul-x-cat-2021-slot-3)
\end{enumerate}
\vspace{1em}
\textbf{Solution:}

We will check which option when proven false will support the passage:
A: Robots becoming a part of everyday life is neither supported nor opposed in the passage. Thus, Option A is not the answer.
B:  _Instead of gears and wires, it’s assembled from two kinds of molecules - some that love water and some that avoid it - which interact to allow the bot to skate on top of a pond._
Option B has been clearly mentioned in the passage. Hence, if it is proven false, it will contradict the passage. Option B is not the answer.
C:  _One team of researchers, for example, has developed a different kind of hydrogel that becomes sticky when exposed to a low-voltage zap of electricity and then stops being sticky when the electricity is switched off._
Option C has been mentioned in the passage. Hence, if it is proven false, it will contradict the passage. Option C is not the answer.
D:  _No, they won’t be invading our bodies and turning us into Sentinels..._
Option D is just the opposite of what has been given in the passage. Hence, if Option D is false, it would support the passage.

\end{problem}

\newpage
\begin{problem}{}{}
% Subject: varc
\textbf{Instructions}

Instructions   

The passage below is accompanied by a set of questions. Choose the best answer to each question.
Starting in 1957, [Noam Chomsky] proclaimed a new doctrine: Language, that most human of all attributes, was innate. The grammatical faculty was built into the infant brain, and your average 3-year-old was not a mere apprentice in the great enterprise of absorbing English from his or her parents, but a “linguistic genius.” Since this message was couched in terms of Chomskyan theoretical linguistics, in discourse so opaque that it was nearly incomprehensible even to some scholars, many people did not hear it. Now, in a brilliant, witty and altogether satisfying book, Mr. Chomsky's colleague Steven Pinker . . . has brought Mr. Chomsky's findings to everyman. In “The Language Instinct” he has gathered persuasive data from such diverse fields as cognitive neuroscience, developmental psychology and speech therapy to make his points, and when he disagrees with Mr. Chomsky he tells you so. . . .
For Mr. Chomsky and Mr. Pinker, somewhere in the human brain there is a complex set of neural circuits that have been programmed with “super-rules” (making up what Mr. Chomsky calls “universal grammar”), and that these rules are unconscious and instinctive. A half-century ago, this would have been pooh-poohed as a “black box” theory, since one could not actually pinpoint this grammatical faculty in a specific part of the brain, or describe its functioning. But now things are different. Neurosurgeons [have now found that this] “black box” is situated in and around Broca’s area, on the left side of the forebrain. . . . 
Unlike Mr. Chomsky, Mr. Pinker firmly places the wiring of the brain for language within the framework of Darwinian natural selection and evolution. He effectively disposes of all claims that intelligent nonhuman primates like chimps have any abilities to learn and use language. It is not that chimps lack the vocal apparatus to speak; it is just that their brains are unable to produce or use grammar. On the other hand, the “language instinct,” when it first appeared among our most distant hominid ancestors, must have given them a selective reproductive advantage over their competitors (including the ancestral chimps). . . .
So according to Mr. Pinker, the roots of language must be in the genes, but there cannot be a “grammar gene” any more than there can be a gene for the heart or any other complex body structure. This proposition will undoubtedly raise the hackles of some behavioral psychologists and anthropologists, for it apparently contradicts the liberal idea that human behavior may be changed for the better by improvements in culture and environment, and it might seem to invite the twin bugaboos of biological determinism and racism. Yet Mr. Pinker stresses one point that should allay such fears. Even though there are 4,000 to 6,000 languages today, they are all sufficiently alike to be considered one language by an extraterrestrial observer. In other words, most of the diversity of the world’s cultures, so beloved to anthropologists, is superficial and minor compared to the similarities. Racial differences are literally only “skin deep.” The fundamental unity of humanity is the theme of Mr. Chomsky's universal grammar, and of this exciting book.

\vspace{1em}
\textbf{Question 11}

Which one of the following statements, if true, would be the most direct extension of the arguments in the passage?

\begin{enumerate}[label=\Alph*.]
    \item Sentinel robots will be used in warfare to cause large-scale destructive mutations amongst civilians.
    \item X-Men may be created by injecting people with mutant nano-gels that will respond to the brain’s magnetic field.
    \item In the future, robots will be used to search and destroy diseases even in the deepest recesses of the human body.
    \item 1-centimetre blobs of gel that have nano-robots in them will be used to send messages.
[](https://cracku.in/11-which-one-of-the-following-statements-if-true-woul-x-cat-2021-slot-3)
\end{enumerate}
\vspace{1em}
\textbf{Solution:}

A: Sentinel robots are just fiction that is mentioned in the passage to introduce the new wave of development that has taken place. Option A is eliminated.
B: The author has introduced X-men as an example only. His arguments are not related to the creation of X-men in any way. Option B can be eliminated.
C: ....._but some of them could one day swim through our bloodstream to heal us._
Throughout the passage, the author is trying to highlight the positives of the new robots. Hence, a direct extension of the argument would be the robots healing us at a microscopic level, as is hinted in the above excerpt. Option C is the answer.
D: Option D, though not entirely incorrect, is not a direct extension of the arguments presented in the passage. Unlike Option C, D has not been hinted at in the passage. Hence, it can be eliminated too.

\end{problem}

\newpage
\begin{problem}{}{}
% Subject: varc
\textbf{Instructions}

Instructions   

The passage below is accompanied by a set of questions. Choose the best answer to each question.
Starting in 1957, [Noam Chomsky] proclaimed a new doctrine: Language, that most human of all attributes, was innate. The grammatical faculty was built into the infant brain, and your average 3-year-old was not a mere apprentice in the great enterprise of absorbing English from his or her parents, but a “linguistic genius.” Since this message was couched in terms of Chomskyan theoretical linguistics, in discourse so opaque that it was nearly incomprehensible even to some scholars, many people did not hear it. Now, in a brilliant, witty and altogether satisfying book, Mr. Chomsky's colleague Steven Pinker . . . has brought Mr. Chomsky's findings to everyman. In “The Language Instinct” he has gathered persuasive data from such diverse fields as cognitive neuroscience, developmental psychology and speech therapy to make his points, and when he disagrees with Mr. Chomsky he tells you so. . . .
For Mr. Chomsky and Mr. Pinker, somewhere in the human brain there is a complex set of neural circuits that have been programmed with “super-rules” (making up what Mr. Chomsky calls “universal grammar”), and that these rules are unconscious and instinctive. A half-century ago, this would have been pooh-poohed as a “black box” theory, since one could not actually pinpoint this grammatical faculty in a specific part of the brain, or describe its functioning. But now things are different. Neurosurgeons [have now found that this] “black box” is situated in and around Broca’s area, on the left side of the forebrain. . . . 
Unlike Mr. Chomsky, Mr. Pinker firmly places the wiring of the brain for language within the framework of Darwinian natural selection and evolution. He effectively disposes of all claims that intelligent nonhuman primates like chimps have any abilities to learn and use language. It is not that chimps lack the vocal apparatus to speak; it is just that their brains are unable to produce or use grammar. On the other hand, the “language instinct,” when it first appeared among our most distant hominid ancestors, must have given them a selective reproductive advantage over their competitors (including the ancestral chimps). . . .
So according to Mr. Pinker, the roots of language must be in the genes, but there cannot be a “grammar gene” any more than there can be a gene for the heart or any other complex body structure. This proposition will undoubtedly raise the hackles of some behavioral psychologists and anthropologists, for it apparently contradicts the liberal idea that human behavior may be changed for the better by improvements in culture and environment, and it might seem to invite the twin bugaboos of biological determinism and racism. Yet Mr. Pinker stresses one point that should allay such fears. Even though there are 4,000 to 6,000 languages today, they are all sufficiently alike to be considered one language by an extraterrestrial observer. In other words, most of the diversity of the world’s cultures, so beloved to anthropologists, is superficial and minor compared to the similarities. Racial differences are literally only “skin deep.” The fundamental unity of humanity is the theme of Mr. Chomsky's universal grammar, and of this exciting book.

\vspace{1em}
\textbf{Question 12}

Which one of the following statements best summarises the central point of the passage?

\begin{enumerate}[label=\Alph*.]
    \item Robots will use nano-robots on their feet and wheels to climb walls or move on ceilings.
    \item Nano-robots made from molecules that react to water have become increasingly useful.
    \item Once the stuff of science fiction, nano-robots now feature in cutting-edge scientific research.
    \item The field of robotics is likely to be featured more and more in comics like the New X-Men.
[](https://cracku.in/12-which-one-of-the-following-statements-best-summari-x-cat-2021-slot-3)
\end{enumerate}
\vspace{1em}
\textbf{Solution:}

The author first introduces an arc of a comic book where nano-robots are used. He then goes on to show how that fiction is increasingly becoming reality. He then goes on to describe the various features present in today's nano-robots. Option C comes the closest in capturing this point, and hence, is the answer.
Option A is just one of the features of the modern nano-robots and is not the focus of the passage.
Hydrophilic and Hydrophobic materials are not the main point of contention here. Option B can be eliminated.
The author uses the example of X-men to introduce the development of technology today. His main contention is not the content of the comic books and how it would be affected by recent developments in technology. Option D can be eliminated too.

Instructions   

The passage below is accompanied by a set of questions. Choose the best answer to each question.
Keeping time accurately comes with a price. The maximum accuracy of a clock is directly related to how much disorder, or entropy, it creates every time it ticks. Natalia Ares at the University of Oxford and her colleagues made this discovery using a tiny clock with an accuracy that can be controlled. The clock consists of a 50-nanometre-thick membrane of silicon nitride, vibrated by an electric current. Each time the membrane moved up and down once and then returned to its original position, the researchers counted a tick, and the regularity of the spacing between the ticks represented the accuracy of the clock. The researchers found that as they increased the clock’s accuracy, the heat produced in the system grew, increasing the entropy of its surroundings by jostling nearby particles . . . “If a clock is more accurate, you are paying for it somehow,” says Ares. In this case, you pay for it by pouring more ordered energy into the clock, which is then converted into entropy. “By measuring time, we are increasing the entropy of the universe,” says Ares. The more entropy there is in the universe, the closer it may be to its eventual demise. “Maybe we should stop measuring time,” says Ares. The scale of the additional entropy is so small, though, that there is no need to worry about its effects, she says. 
The increase in entropy in timekeeping may be related to the “arrow of time”, says Marcus Huber at the Austrian Academy of Sciences in Vienna, who was part of the research team. It has been suggested that the reason that time only flows forward, not in reverse, is that the total amount of entropy in the universe is constantly increasing, creating disorder that cannot be put in order again.
The relationship that the researchers found is a limit on the accuracy of a clock, so it doesn’t mean that a clock that creates the most possible entropy would be maximally accurate - hence a large, inefficient grandfather clock isn’t more precise than an atomic clock. “It’s a bit like fuel use in a car. Just because I’m using more fuel doesn’t mean that I’m going faster or further,” says Huber.
When the researchers compared their results with theoretical models developed for clocks that rely on quantum effects, they were surprised to find that the relationship between accuracy and entropy seemed to be the same for both. . . . We can’t be sure yet that these results are actually universal, though, because there are many types of clocks for which the relationship between accuracy and entropy haven’t been tested. “It’s still unclear how this principle plays out in real devices such as atomic clocks, which push the ultimate quantum limits of accuracy,” says Mark Mitchison at Trinity College Dublin in Ireland. Understanding this relationship could be helpful for designing clocks in the future, particularly those used in quantum computers and other devices where both accuracy and temperature are crucial, says Ares. This finding could also help us understand more generally how the quantum world and the classical world are similar and different in terms of thermodynamics and the passage of time.

\end{problem}

\newpage
\begin{problem}{}{}
% Subject: varc
\textbf{Instructions}

Instructions   

The passage below is accompanied by a set of questions. Choose the best answer to each question.
Starting in 1957, [Noam Chomsky] proclaimed a new doctrine: Language, that most human of all attributes, was innate. The grammatical faculty was built into the infant brain, and your average 3-year-old was not a mere apprentice in the great enterprise of absorbing English from his or her parents, but a “linguistic genius.” Since this message was couched in terms of Chomskyan theoretical linguistics, in discourse so opaque that it was nearly incomprehensible even to some scholars, many people did not hear it. Now, in a brilliant, witty and altogether satisfying book, Mr. Chomsky's colleague Steven Pinker . . . has brought Mr. Chomsky's findings to everyman. In “The Language Instinct” he has gathered persuasive data from such diverse fields as cognitive neuroscience, developmental psychology and speech therapy to make his points, and when he disagrees with Mr. Chomsky he tells you so. . . .
For Mr. Chomsky and Mr. Pinker, somewhere in the human brain there is a complex set of neural circuits that have been programmed with “super-rules” (making up what Mr. Chomsky calls “universal grammar”), and that these rules are unconscious and instinctive. A half-century ago, this would have been pooh-poohed as a “black box” theory, since one could not actually pinpoint this grammatical faculty in a specific part of the brain, or describe its functioning. But now things are different. Neurosurgeons [have now found that this] “black box” is situated in and around Broca’s area, on the left side of the forebrain. . . . 
Unlike Mr. Chomsky, Mr. Pinker firmly places the wiring of the brain for language within the framework of Darwinian natural selection and evolution. He effectively disposes of all claims that intelligent nonhuman primates like chimps have any abilities to learn and use language. It is not that chimps lack the vocal apparatus to speak; it is just that their brains are unable to produce or use grammar. On the other hand, the “language instinct,” when it first appeared among our most distant hominid ancestors, must have given them a selective reproductive advantage over their competitors (including the ancestral chimps). . . .
So according to Mr. Pinker, the roots of language must be in the genes, but there cannot be a “grammar gene” any more than there can be a gene for the heart or any other complex body structure. This proposition will undoubtedly raise the hackles of some behavioral psychologists and anthropologists, for it apparently contradicts the liberal idea that human behavior may be changed for the better by improvements in culture and environment, and it might seem to invite the twin bugaboos of biological determinism and racism. Yet Mr. Pinker stresses one point that should allay such fears. Even though there are 4,000 to 6,000 languages today, they are all sufficiently alike to be considered one language by an extraterrestrial observer. In other words, most of the diversity of the world’s cultures, so beloved to anthropologists, is superficial and minor compared to the similarities. Racial differences are literally only “skin deep.” The fundamental unity of humanity is the theme of Mr. Chomsky's universal grammar, and of this exciting book.

\vspace{1em}
\textbf{Question 13}

None of the following statements can be inferred from the passage EXCEPT that:

\begin{enumerate}[label=\Alph*.]
    \item the arrow of time has not yet been tested for atomic clocks.
    \item quantum computers are likely to produce more heat and, hence, more entropy, because of the emphasis on their clocks' accuracy.
    \item grandfather clocks are likely to produce less heat and, hence, less entropy, because they are not as accurate.
    \item a clock with a 50-nanometre-thick membrane of silicon nitride has been made to vibrate, producing electric currents.
[](https://cracku.in/13-none-of-the-following-statements-can-be-inferred-f-x-cat-2021-slot-3)
\end{enumerate}
\vspace{1em}
\textbf{Solution:}

A: We cannot infer that the 'arrow of time' has not been tested for atomic clocks. Option A can be eliminated.
B: It has been given in the Option that since quantum computers place more emphasis on their clock's accuracy, they would produce more heat.
_The researchers found that as they increased the clock’s accuracy, the heat produced in the system grew, increasing the entropy of its surroundings by jostling nearby particles..._
The passage supports this inference. B is the answer.
C:  _The relationship that the researchers found is a limit on the accuracy of a clock, so it doesn’t mean that a clock that creates the most possible entropy would be maximally accurate - hence a large, inefficient grandfather clock isn’t more precise than an atomic clock_.
The passage gives a specific example of an inefficient grandfather clock. We cannot infer whether all grandfather clocks are efficient or not.
D:  _The clock consists of a 50-nanometre-thick membrane of silicon nitride, vibrated by an electric current._
The clock uses electric current to produce vibrations and not the other way around. Option D can be eliminated.

\end{problem}

\newpage
\begin{problem}{}{}
% Subject: varc
\textbf{Instructions}

Instructions   

The passage below is accompanied by a set of questions. Choose the best answer to each question.
Starting in 1957, [Noam Chomsky] proclaimed a new doctrine: Language, that most human of all attributes, was innate. The grammatical faculty was built into the infant brain, and your average 3-year-old was not a mere apprentice in the great enterprise of absorbing English from his or her parents, but a “linguistic genius.” Since this message was couched in terms of Chomskyan theoretical linguistics, in discourse so opaque that it was nearly incomprehensible even to some scholars, many people did not hear it. Now, in a brilliant, witty and altogether satisfying book, Mr. Chomsky's colleague Steven Pinker . . . has brought Mr. Chomsky's findings to everyman. In “The Language Instinct” he has gathered persuasive data from such diverse fields as cognitive neuroscience, developmental psychology and speech therapy to make his points, and when he disagrees with Mr. Chomsky he tells you so. . . .
For Mr. Chomsky and Mr. Pinker, somewhere in the human brain there is a complex set of neural circuits that have been programmed with “super-rules” (making up what Mr. Chomsky calls “universal grammar”), and that these rules are unconscious and instinctive. A half-century ago, this would have been pooh-poohed as a “black box” theory, since one could not actually pinpoint this grammatical faculty in a specific part of the brain, or describe its functioning. But now things are different. Neurosurgeons [have now found that this] “black box” is situated in and around Broca’s area, on the left side of the forebrain. . . . 
Unlike Mr. Chomsky, Mr. Pinker firmly places the wiring of the brain for language within the framework of Darwinian natural selection and evolution. He effectively disposes of all claims that intelligent nonhuman primates like chimps have any abilities to learn and use language. It is not that chimps lack the vocal apparatus to speak; it is just that their brains are unable to produce or use grammar. On the other hand, the “language instinct,” when it first appeared among our most distant hominid ancestors, must have given them a selective reproductive advantage over their competitors (including the ancestral chimps). . . .
So according to Mr. Pinker, the roots of language must be in the genes, but there cannot be a “grammar gene” any more than there can be a gene for the heart or any other complex body structure. This proposition will undoubtedly raise the hackles of some behavioral psychologists and anthropologists, for it apparently contradicts the liberal idea that human behavior may be changed for the better by improvements in culture and environment, and it might seem to invite the twin bugaboos of biological determinism and racism. Yet Mr. Pinker stresses one point that should allay such fears. Even though there are 4,000 to 6,000 languages today, they are all sufficiently alike to be considered one language by an extraterrestrial observer. In other words, most of the diversity of the world’s cultures, so beloved to anthropologists, is superficial and minor compared to the similarities. Racial differences are literally only “skin deep.” The fundamental unity of humanity is the theme of Mr. Chomsky's universal grammar, and of this exciting book.

\vspace{1em}
\textbf{Question 14}

The author makes all of the following arguments in the passage, EXCEPT that:

\begin{enumerate}[label=\Alph*.]
    \item The relationship between accuracy and entropy may not apply to all clocks.
    \item Researchers found that the heat produced in a system is the price paid for increased accuracy of measurement.
    \item There is no difference in accuracy between an inefficient grandfather clock and an atomic clock.
    \item In designing clocks for quantum computers, both precision and heat have to be taken into account.
[](https://cracku.in/14-the-author-makes-all-of-the-following-arguments-in-x-cat-2021-slot-3)
\end{enumerate}
\vspace{1em}
\textbf{Solution:}

There is an evident confusion between Option B and Option C; however, the official answer key marked Option B as the correct choice. Let us try to rationalise this decision. Options A and D can be understood from the passage:
_Option A_ follows from {_...We can’t be sure yet that these results are actually universal, though, because there are many types of clocks for which the relationship between accuracy and entropy haven’t been tested..._}
_Option D_ follows from {..._Understanding this relationship could be helpful for designing clocks in the future, particularly those used in quantum computers and other devices where both accuracy and temperature are crucial, says Ares..._}
_Option C_ : Pay heed to the following excerpt from the passage - {..._The relationship that the researchers found is a limit on the accuracy of a clock, so it doesn’t mean that a clock that creates the most possible entropy would be maximally accurate - hence a large, inefficient grandfather clock isn’t more precise than an atomic clock. “It’s a bit like fuel use in a car. Just because I’m using more fuel doesn’t mean that I’m going faster or further,” says Huber..._}
A simple correlation is being highlighted: higher accuracy means higher entropy; however, this does not necessarily imply that higher entropy translates to higher accuracy. The example of a grandfather clock is highlighted to emphasise this point: we will come across higher entropy in this case, but it does not mean that the grandfather clock is any more accurate than an atomic clock. In a way, the author tries to point out that the accuracy could very well be similar. This accuracy is not in absolute terms but in the way accuracy is defined by the author earlier in the passage. Thus, in a way, Option C matches the idea conveyed by the author
_Option B_ : Pay heed to the following excerpt from the passage - {.._.The researchers found that as they increased the clock’s accuracy, the heat produced in the system grew, increasing the entropy of its surroundings by jostling nearby particles . . . “If a clock is more accurate, you are paying for it somehow,” says Ares. In this case, you pay for it by pouring more ordered energy into the clock, which is then converted into entropy. “By measuring time, we are increasing the entropy of the universe,” says Ares..._}
The discussion about the price paid appears to be distinct from the earlier segment wherein the author states that when we push for higher accuracy, we will come across more heat. While talking about the cost at which higher accuracy is achieved, the author states that we "pour in" more 'ordered energy' and this subsequently leads to higher entropy. Hence, the focus seems to be on the connection between accuracy and entropy than between heat and its role in creating higher accuracy. We cannot conclusively infer that the "_ordered energy_ " stated in the latter half refers to the "_heat_ " mentioned earlier on. Thus, claiming that heat is the price we pay for generating higher accuracy might be difficult to substantiate. Hence, Option B is distorted.

\end{problem}

\newpage
\begin{problem}{}{}
% Subject: varc
\textbf{Instructions}

Instructions   

The passage below is accompanied by a set of questions. Choose the best answer to each question.
Starting in 1957, [Noam Chomsky] proclaimed a new doctrine: Language, that most human of all attributes, was innate. The grammatical faculty was built into the infant brain, and your average 3-year-old was not a mere apprentice in the great enterprise of absorbing English from his or her parents, but a “linguistic genius.” Since this message was couched in terms of Chomskyan theoretical linguistics, in discourse so opaque that it was nearly incomprehensible even to some scholars, many people did not hear it. Now, in a brilliant, witty and altogether satisfying book, Mr. Chomsky's colleague Steven Pinker . . . has brought Mr. Chomsky's findings to everyman. In “The Language Instinct” he has gathered persuasive data from such diverse fields as cognitive neuroscience, developmental psychology and speech therapy to make his points, and when he disagrees with Mr. Chomsky he tells you so. . . .
For Mr. Chomsky and Mr. Pinker, somewhere in the human brain there is a complex set of neural circuits that have been programmed with “super-rules” (making up what Mr. Chomsky calls “universal grammar”), and that these rules are unconscious and instinctive. A half-century ago, this would have been pooh-poohed as a “black box” theory, since one could not actually pinpoint this grammatical faculty in a specific part of the brain, or describe its functioning. But now things are different. Neurosurgeons [have now found that this] “black box” is situated in and around Broca’s area, on the left side of the forebrain. . . . 
Unlike Mr. Chomsky, Mr. Pinker firmly places the wiring of the brain for language within the framework of Darwinian natural selection and evolution. He effectively disposes of all claims that intelligent nonhuman primates like chimps have any abilities to learn and use language. It is not that chimps lack the vocal apparatus to speak; it is just that their brains are unable to produce or use grammar. On the other hand, the “language instinct,” when it first appeared among our most distant hominid ancestors, must have given them a selective reproductive advantage over their competitors (including the ancestral chimps). . . .
So according to Mr. Pinker, the roots of language must be in the genes, but there cannot be a “grammar gene” any more than there can be a gene for the heart or any other complex body structure. This proposition will undoubtedly raise the hackles of some behavioral psychologists and anthropologists, for it apparently contradicts the liberal idea that human behavior may be changed for the better by improvements in culture and environment, and it might seem to invite the twin bugaboos of biological determinism and racism. Yet Mr. Pinker stresses one point that should allay such fears. Even though there are 4,000 to 6,000 languages today, they are all sufficiently alike to be considered one language by an extraterrestrial observer. In other words, most of the diversity of the world’s cultures, so beloved to anthropologists, is superficial and minor compared to the similarities. Racial differences are literally only “skin deep.” The fundamental unity of humanity is the theme of Mr. Chomsky's universal grammar, and of this exciting book.

\vspace{1em}
\textbf{Question 15}

“It’s a bit like fuel use in a car. Just because I’m using more fuel doesn’t mean that I’m going faster or further . . .” What is the purpose of this example?

\begin{enumerate}[label=\Alph*.]
    \item If you go faster in a car, you will tend to consume more fuel, but the converse is not necessarily true. In the same way, increased entropy does not necessarily mean greater accuracy of a clock.
    \item The further you go in a car, the more fuel you use. In the same way, the faster you go in a car, the less time you use.
    \item If you measure the speed of a car with a grandfather clock, the result will be different than if you measured it with an atomic clock.
    \item The further and faster you go in a car, the greater the amount of fuel you will use, the greater the amount of heat produced and, hence, the greater the entropy.
[](https://cracku.in/15-its-a-bit-like-fuel-use-in-a-car-just-because-im-u-x-cat-2021-slot-3)
\end{enumerate}
\vspace{1em}
\textbf{Solution:}

_The relationship that the researchers found is a limit on the accuracy of a clock, so it doesn’t mean that a clock that creates the most possible entropy would be maximally accurate - hence a large, inefficient grandfather clock isn’t more precise than an atomic clock. “It’s a bit like fuel use in a car. Just because I’m using more fuel doesn’t mean that I’m going faster or further,” says Hube_
In the above excerpt, the author gives an example that though a large, inefficient grandfather clock would produce more entropy, it is not necessarily more precise than an atomic clock. Hence, if a clock produces more entropy, it does not mean that it would be more precise than a clock that produces less entropy. Then the mentioned statement is given as an example. If a car is going faster or further, it will definitely use more fuel. But if a car is using more fuel, then the converse is not true. It could just be possible that the mileage of the car is low. Option A comes the closest to capturing this idea, and hence, is the answer.

\end{problem}

\newpage
\begin{problem}{}{}
% Subject: varc
\textbf{Instructions}

Instructions   

The passage below is accompanied by a set of questions. Choose the best answer to each question.
Starting in 1957, [Noam Chomsky] proclaimed a new doctrine: Language, that most human of all attributes, was innate. The grammatical faculty was built into the infant brain, and your average 3-year-old was not a mere apprentice in the great enterprise of absorbing English from his or her parents, but a “linguistic genius.” Since this message was couched in terms of Chomskyan theoretical linguistics, in discourse so opaque that it was nearly incomprehensible even to some scholars, many people did not hear it. Now, in a brilliant, witty and altogether satisfying book, Mr. Chomsky's colleague Steven Pinker . . . has brought Mr. Chomsky's findings to everyman. In “The Language Instinct” he has gathered persuasive data from such diverse fields as cognitive neuroscience, developmental psychology and speech therapy to make his points, and when he disagrees with Mr. Chomsky he tells you so. . . .
For Mr. Chomsky and Mr. Pinker, somewhere in the human brain there is a complex set of neural circuits that have been programmed with “super-rules” (making up what Mr. Chomsky calls “universal grammar”), and that these rules are unconscious and instinctive. A half-century ago, this would have been pooh-poohed as a “black box” theory, since one could not actually pinpoint this grammatical faculty in a specific part of the brain, or describe its functioning. But now things are different. Neurosurgeons [have now found that this] “black box” is situated in and around Broca’s area, on the left side of the forebrain. . . . 
Unlike Mr. Chomsky, Mr. Pinker firmly places the wiring of the brain for language within the framework of Darwinian natural selection and evolution. He effectively disposes of all claims that intelligent nonhuman primates like chimps have any abilities to learn and use language. It is not that chimps lack the vocal apparatus to speak; it is just that their brains are unable to produce or use grammar. On the other hand, the “language instinct,” when it first appeared among our most distant hominid ancestors, must have given them a selective reproductive advantage over their competitors (including the ancestral chimps). . . .
So according to Mr. Pinker, the roots of language must be in the genes, but there cannot be a “grammar gene” any more than there can be a gene for the heart or any other complex body structure. This proposition will undoubtedly raise the hackles of some behavioral psychologists and anthropologists, for it apparently contradicts the liberal idea that human behavior may be changed for the better by improvements in culture and environment, and it might seem to invite the twin bugaboos of biological determinism and racism. Yet Mr. Pinker stresses one point that should allay such fears. Even though there are 4,000 to 6,000 languages today, they are all sufficiently alike to be considered one language by an extraterrestrial observer. In other words, most of the diversity of the world’s cultures, so beloved to anthropologists, is superficial and minor compared to the similarities. Racial differences are literally only “skin deep.” The fundamental unity of humanity is the theme of Mr. Chomsky's universal grammar, and of this exciting book.

\vspace{1em}
\textbf{Question 16}

Which one of the following sets of words and phrases serves best as keywords of the passage?

\begin{enumerate}[label=\Alph*.]
    \item Electric current; Heat; Quantum effects.
    \item Silicon Nitride; Energy; Grandfather Clock.
    \item Measuring Time; Accuracy; Entropy.
    \item Membrane; Arrow of time; Entropy.
[](https://cracku.in/16-which-one-of-the-following-sets-of-words-and-phras-x-cat-2021-slot-3)
\end{enumerate}
\vspace{1em}
\textbf{Solution:}

_The maximum accuracy of a clock is directly related to how much disorder, or entropy, it creates every time it ticks._
The author highlights in the beginning of the passage that the accuracy associated with measuring time is directly related to how much entropy it creates while ticking. The author then goes on to talk about the relationship between accuracy and entropy, and how quantum mechanics and thermodynamics come in play here. Thus, the main keywords are the measurement of time, accuracy and entropy. Option C is the answer.  

Electric current is just a small part of an example presented in the passage. Option A can be eliminated. The same is the case for Silicon Nitride and Membrane. These are just keywords associated with a particular experiment/example presented in the passage and are not important for the passage as a whole.

Instructions   

For the following questions answer them individually

\end{problem}

\newpage
\begin{problem}{}{}
% Subject: varc
\textbf{Instructions}

Instructions   

The passage below is accompanied by a set of questions. Choose the best answer to each question.
Starting in 1957, [Noam Chomsky] proclaimed a new doctrine: Language, that most human of all attributes, was innate. The grammatical faculty was built into the infant brain, and your average 3-year-old was not a mere apprentice in the great enterprise of absorbing English from his or her parents, but a “linguistic genius.” Since this message was couched in terms of Chomskyan theoretical linguistics, in discourse so opaque that it was nearly incomprehensible even to some scholars, many people did not hear it. Now, in a brilliant, witty and altogether satisfying book, Mr. Chomsky's colleague Steven Pinker . . . has brought Mr. Chomsky's findings to everyman. In “The Language Instinct” he has gathered persuasive data from such diverse fields as cognitive neuroscience, developmental psychology and speech therapy to make his points, and when he disagrees with Mr. Chomsky he tells you so. . . .
For Mr. Chomsky and Mr. Pinker, somewhere in the human brain there is a complex set of neural circuits that have been programmed with “super-rules” (making up what Mr. Chomsky calls “universal grammar”), and that these rules are unconscious and instinctive. A half-century ago, this would have been pooh-poohed as a “black box” theory, since one could not actually pinpoint this grammatical faculty in a specific part of the brain, or describe its functioning. But now things are different. Neurosurgeons [have now found that this] “black box” is situated in and around Broca’s area, on the left side of the forebrain. . . . 
Unlike Mr. Chomsky, Mr. Pinker firmly places the wiring of the brain for language within the framework of Darwinian natural selection and evolution. He effectively disposes of all claims that intelligent nonhuman primates like chimps have any abilities to learn and use language. It is not that chimps lack the vocal apparatus to speak; it is just that their brains are unable to produce or use grammar. On the other hand, the “language instinct,” when it first appeared among our most distant hominid ancestors, must have given them a selective reproductive advantage over their competitors (including the ancestral chimps). . . .
So according to Mr. Pinker, the roots of language must be in the genes, but there cannot be a “grammar gene” any more than there can be a gene for the heart or any other complex body structure. This proposition will undoubtedly raise the hackles of some behavioral psychologists and anthropologists, for it apparently contradicts the liberal idea that human behavior may be changed for the better by improvements in culture and environment, and it might seem to invite the twin bugaboos of biological determinism and racism. Yet Mr. Pinker stresses one point that should allay such fears. Even though there are 4,000 to 6,000 languages today, they are all sufficiently alike to be considered one language by an extraterrestrial observer. In other words, most of the diversity of the world’s cultures, so beloved to anthropologists, is superficial and minor compared to the similarities. Racial differences are literally only “skin deep.” The fundamental unity of humanity is the theme of Mr. Chomsky's universal grammar, and of this exciting book.

\vspace{1em}
\textbf{Question 17}

**The four sentences (labelled 1, 2, 3 and 4) below, when properly sequenced would yield a coherent paragraph. Decide on the proper sequencing of the order of the sentences and key in the sequence of the four numbers as your answer:**  
1. Businesses find automation, such as robotic employees, a big asset in terms of productivity and efficiency.  
2. But in recent years, robotics has had increasing impacts on unemployment, not just of manual labour, as computers are rapidly handling some white-collar and service-sector work.  
3. For years politicians have promised workers that they would bring back their jobs by clamping down on trade, offshoring and immigration.  
4. Economists, based on their research, say that the bigger threat to jobs now is not globalisation but automation.
Backspace  
789  
456  
123  
0.-  
Clear All  

Submit
[](https://cracku.in/17-the-four-sentences-labelled-1-2-3-and-4-below-when-x-cat-2021-slot-3)

\begin{enumerate}[label=\Alph*.]
\end{enumerate}
\vspace{1em}
\textbf{Solution:}

A quick read of the sentences tells us that the paragraph is about the unemployment caused by automation. The passage is best opened by 3, which provides the current state of unemployment. The politicians view globalisation as the factor exacerbating unemployment. 4 contrasts this by saying that expert analysis tells a different story. It is automation that could prove to be a crucial factor. 12 forms a pair, that further elucidate the kind and scope of impact that automation has on jobs. Hence, the correct sequence would be 3412.

\end{problem}

\newpage
\begin{problem}{}{}
% Subject: varc
\textbf{Instructions}

Instructions   

The passage below is accompanied by a set of questions. Choose the best answer to each question.
Starting in 1957, [Noam Chomsky] proclaimed a new doctrine: Language, that most human of all attributes, was innate. The grammatical faculty was built into the infant brain, and your average 3-year-old was not a mere apprentice in the great enterprise of absorbing English from his or her parents, but a “linguistic genius.” Since this message was couched in terms of Chomskyan theoretical linguistics, in discourse so opaque that it was nearly incomprehensible even to some scholars, many people did not hear it. Now, in a brilliant, witty and altogether satisfying book, Mr. Chomsky's colleague Steven Pinker . . . has brought Mr. Chomsky's findings to everyman. In “The Language Instinct” he has gathered persuasive data from such diverse fields as cognitive neuroscience, developmental psychology and speech therapy to make his points, and when he disagrees with Mr. Chomsky he tells you so. . . .
For Mr. Chomsky and Mr. Pinker, somewhere in the human brain there is a complex set of neural circuits that have been programmed with “super-rules” (making up what Mr. Chomsky calls “universal grammar”), and that these rules are unconscious and instinctive. A half-century ago, this would have been pooh-poohed as a “black box” theory, since one could not actually pinpoint this grammatical faculty in a specific part of the brain, or describe its functioning. But now things are different. Neurosurgeons [have now found that this] “black box” is situated in and around Broca’s area, on the left side of the forebrain. . . . 
Unlike Mr. Chomsky, Mr. Pinker firmly places the wiring of the brain for language within the framework of Darwinian natural selection and evolution. He effectively disposes of all claims that intelligent nonhuman primates like chimps have any abilities to learn and use language. It is not that chimps lack the vocal apparatus to speak; it is just that their brains are unable to produce or use grammar. On the other hand, the “language instinct,” when it first appeared among our most distant hominid ancestors, must have given them a selective reproductive advantage over their competitors (including the ancestral chimps). . . .
So according to Mr. Pinker, the roots of language must be in the genes, but there cannot be a “grammar gene” any more than there can be a gene for the heart or any other complex body structure. This proposition will undoubtedly raise the hackles of some behavioral psychologists and anthropologists, for it apparently contradicts the liberal idea that human behavior may be changed for the better by improvements in culture and environment, and it might seem to invite the twin bugaboos of biological determinism and racism. Yet Mr. Pinker stresses one point that should allay such fears. Even though there are 4,000 to 6,000 languages today, they are all sufficiently alike to be considered one language by an extraterrestrial observer. In other words, most of the diversity of the world’s cultures, so beloved to anthropologists, is superficial and minor compared to the similarities. Racial differences are literally only “skin deep.” The fundamental unity of humanity is the theme of Mr. Chomsky's universal grammar, and of this exciting book.

\vspace{1em}
\textbf{Question 18}

**The four sentences (labelled 1, 2, 3, 4) below, when properly sequenced would yield a coherent paragraph. Decide on the proper sequencing of the order of the sentences and key in the sequence of the four numbers as your answer:**  
1. It is regimes of truth that make certain relationships speakable - relationships, like subjectivities, are constituted through discursive formations, which sustain regimes of truth.  
2. Relationships are nothing without the communication that brings them into being; interpersonal communication is connected to knowledge shared by interlocutors, and scholars should attend to relational histories in their analyses.  
3. A Foucauldian approach to relationships goes beyond these conceptions of discourse and history to macrolevel regimes of truth as constituting relationships.  
4. Reconsidering micropractices within relationships that are constituted within and simultaneously contributors to regimes of truth acknowledges the central position of power/knowledge in the constitution of what has come to be considered true and real.
Backspace  
789  
456  
123  
0.-  
Clear All  

Submit
[](https://cracku.in/18-the-four-sentences-labelled-1-2-3-4-below-when-pro-x-cat-2021-slot-3)

\begin{enumerate}[label=\Alph*.]
\end{enumerate}
\vspace{1em}
\textbf{Solution:}

A brief reading of the sentences tells us that the paragraph is about the different conceptions of relationships. 2 explains that communication is an important aspect here, and should be studied properly. 3 mentions a Foucauldian approach, that goes beyond this, and includes macrolevel regimes of truth. 1 then explains why the concept of regimes of truth is relevant here. 4 then aptly concludes the paragraph, implying how the micropractices within the relationships allude the importance of knowledge/power. Thus, the correct sequence would be 2314.

\end{problem}

\newpage
\begin{problem}{}{}
% Subject: varc
\textbf{Instructions}

Instructions   

The passage below is accompanied by a set of questions. Choose the best answer to each question.
Starting in 1957, [Noam Chomsky] proclaimed a new doctrine: Language, that most human of all attributes, was innate. The grammatical faculty was built into the infant brain, and your average 3-year-old was not a mere apprentice in the great enterprise of absorbing English from his or her parents, but a “linguistic genius.” Since this message was couched in terms of Chomskyan theoretical linguistics, in discourse so opaque that it was nearly incomprehensible even to some scholars, many people did not hear it. Now, in a brilliant, witty and altogether satisfying book, Mr. Chomsky's colleague Steven Pinker . . . has brought Mr. Chomsky's findings to everyman. In “The Language Instinct” he has gathered persuasive data from such diverse fields as cognitive neuroscience, developmental psychology and speech therapy to make his points, and when he disagrees with Mr. Chomsky he tells you so. . . .
For Mr. Chomsky and Mr. Pinker, somewhere in the human brain there is a complex set of neural circuits that have been programmed with “super-rules” (making up what Mr. Chomsky calls “universal grammar”), and that these rules are unconscious and instinctive. A half-century ago, this would have been pooh-poohed as a “black box” theory, since one could not actually pinpoint this grammatical faculty in a specific part of the brain, or describe its functioning. But now things are different. Neurosurgeons [have now found that this] “black box” is situated in and around Broca’s area, on the left side of the forebrain. . . . 
Unlike Mr. Chomsky, Mr. Pinker firmly places the wiring of the brain for language within the framework of Darwinian natural selection and evolution. He effectively disposes of all claims that intelligent nonhuman primates like chimps have any abilities to learn and use language. It is not that chimps lack the vocal apparatus to speak; it is just that their brains are unable to produce or use grammar. On the other hand, the “language instinct,” when it first appeared among our most distant hominid ancestors, must have given them a selective reproductive advantage over their competitors (including the ancestral chimps). . . .
So according to Mr. Pinker, the roots of language must be in the genes, but there cannot be a “grammar gene” any more than there can be a gene for the heart or any other complex body structure. This proposition will undoubtedly raise the hackles of some behavioral psychologists and anthropologists, for it apparently contradicts the liberal idea that human behavior may be changed for the better by improvements in culture and environment, and it might seem to invite the twin bugaboos of biological determinism and racism. Yet Mr. Pinker stresses one point that should allay such fears. Even though there are 4,000 to 6,000 languages today, they are all sufficiently alike to be considered one language by an extraterrestrial observer. In other words, most of the diversity of the world’s cultures, so beloved to anthropologists, is superficial and minor compared to the similarities. Racial differences are literally only “skin deep.” The fundamental unity of humanity is the theme of Mr. Chomsky's universal grammar, and of this exciting book.

\vspace{1em}
\textbf{Question 19}

**Five jumbled up sentences, related to a topic, are given below. Four of them can be put together to form a coherent paragraph. Identify the odd one out and key in the number of the sentence as your answer:**  
1. A typical example is Wikipedia, where the overwhelming majority of contributors are male and so the available content is skewed to reflect their interests.  
2. Without diversity of thought and representation, society is left with a distorted picture of future options, which are likely to result in augmenting existing inequalities.  
3. Gross gender inequality in the technology sector is problematic, not only for the industry-wide marginalisation of women, but because technology designs embody the values of their makers.  
4. While redressing unequal representation in the workplace is a step in the right direction, broader social change is needed to address the structural inequalities embedded within the current organisation of work and employment.  
5. If technology merely reflects the perspectives of the male stereotype, then new technologies are unlikely to accommodate the diverse social contexts within which they operate.
Backspace  
789  
456  
123  
0.-  
Clear All  

Submit
[](https://cracku.in/19-five-jumbled-up-sentences-related-to-a-topic-are-g-x-cat-2021-slot-3)

\begin{enumerate}[label=\Alph*.]
\end{enumerate}
\vspace{1em}
\textbf{Solution:}

A brief reading of the sentences suggests that the paragraph must be about the disparity in the representation of different genders. Sentences 1,2,3, and 5 are concerned with the problems that arise when the representation of females is less. 
4, however, runs tangent to the discussion at hand. It talks about 'structural inequalities'. This sentence, if included in the paragraph, would render it incomplete as all the other sentences talk about gender inequality and not structural inequality. Thus, 4 is out of context here.

\end{problem}

\newpage
\begin{problem}{}{}
% Subject: varc
\textbf{Instructions}

Instructions   

The passage below is accompanied by a set of questions. Choose the best answer to each question.
Starting in 1957, [Noam Chomsky] proclaimed a new doctrine: Language, that most human of all attributes, was innate. The grammatical faculty was built into the infant brain, and your average 3-year-old was not a mere apprentice in the great enterprise of absorbing English from his or her parents, but a “linguistic genius.” Since this message was couched in terms of Chomskyan theoretical linguistics, in discourse so opaque that it was nearly incomprehensible even to some scholars, many people did not hear it. Now, in a brilliant, witty and altogether satisfying book, Mr. Chomsky's colleague Steven Pinker . . . has brought Mr. Chomsky's findings to everyman. In “The Language Instinct” he has gathered persuasive data from such diverse fields as cognitive neuroscience, developmental psychology and speech therapy to make his points, and when he disagrees with Mr. Chomsky he tells you so. . . .
For Mr. Chomsky and Mr. Pinker, somewhere in the human brain there is a complex set of neural circuits that have been programmed with “super-rules” (making up what Mr. Chomsky calls “universal grammar”), and that these rules are unconscious and instinctive. A half-century ago, this would have been pooh-poohed as a “black box” theory, since one could not actually pinpoint this grammatical faculty in a specific part of the brain, or describe its functioning. But now things are different. Neurosurgeons [have now found that this] “black box” is situated in and around Broca’s area, on the left side of the forebrain. . . . 
Unlike Mr. Chomsky, Mr. Pinker firmly places the wiring of the brain for language within the framework of Darwinian natural selection and evolution. He effectively disposes of all claims that intelligent nonhuman primates like chimps have any abilities to learn and use language. It is not that chimps lack the vocal apparatus to speak; it is just that their brains are unable to produce or use grammar. On the other hand, the “language instinct,” when it first appeared among our most distant hominid ancestors, must have given them a selective reproductive advantage over their competitors (including the ancestral chimps). . . .
So according to Mr. Pinker, the roots of language must be in the genes, but there cannot be a “grammar gene” any more than there can be a gene for the heart or any other complex body structure. This proposition will undoubtedly raise the hackles of some behavioral psychologists and anthropologists, for it apparently contradicts the liberal idea that human behavior may be changed for the better by improvements in culture and environment, and it might seem to invite the twin bugaboos of biological determinism and racism. Yet Mr. Pinker stresses one point that should allay such fears. Even though there are 4,000 to 6,000 languages today, they are all sufficiently alike to be considered one language by an extraterrestrial observer. In other words, most of the diversity of the world’s cultures, so beloved to anthropologists, is superficial and minor compared to the similarities. Racial differences are literally only “skin deep.” The fundamental unity of humanity is the theme of Mr. Chomsky's universal grammar, and of this exciting book.

\vspace{1em}
\textbf{Question 20}

**The passage given below is followed by four alternate summaries. Choose the option that best captures the essence of the passage.**
People view idleness as a sin and industriousness as a virtue, and in the process have developed an unsatisfactory relationship with their jobs. Work has become a way for them to keep busy, even though many find their work meaningless. In their need for activity people undertake what was once considered work (fishing, gardening) as hobbies. The opposing view is that hard work has made us prosperous and improved our levels of health and education. It has also brought innovation and labour and time-saving devices, which have lessened life’s drudgery.

\begin{enumerate}[label=\Alph*.]
    \item Despite some detractors, hard work is essential in today’s world to enable economic progress, for education and health and to propel innovations that make life easier.
    \item Hard work has overtaken all aspects of our lives and has enabled economic prosperity, but it is important that people reserve their leisure time for some idleness.
    \item Some believe that hard work has been glorified to the extent that it has become meaningless, and led to greater idleness, but it has also had enormous positive impacts on everyday life.
    \item While the idealisation of hard work has propelled people into meaningless jobs and endless activity, it has also led to tremendous social benefits from prosperity and innovation.
[](https://cracku.in/20-the-passage-given-below-is-followed-by-four-altern-x-cat-2021-slot-3)
\end{enumerate}
\vspace{1em}
\textbf{Solution:}

The main points of the passage are:
1. People increasingly view idleness as sin and industriousness as a virtue, pushing them into meaningless jobs.
2. On the other hand, this has also saved us from many of life's drudgeries.
A: Misses out on point 1.
B: A distortion. The author does not advocate idleness. Also, 1 is not covered properly.
C: Incorrect. 'led to greater idleness' is not implied anywhere in the passage.
D: Covers both the points aptly and is the answer.

\end{problem}

\newpage
\begin{problem}{}{}
% Subject: varc
\textbf{Instructions}

Instructions   

The passage below is accompanied by a set of questions. Choose the best answer to each question.
Starting in 1957, [Noam Chomsky] proclaimed a new doctrine: Language, that most human of all attributes, was innate. The grammatical faculty was built into the infant brain, and your average 3-year-old was not a mere apprentice in the great enterprise of absorbing English from his or her parents, but a “linguistic genius.” Since this message was couched in terms of Chomskyan theoretical linguistics, in discourse so opaque that it was nearly incomprehensible even to some scholars, many people did not hear it. Now, in a brilliant, witty and altogether satisfying book, Mr. Chomsky's colleague Steven Pinker . . . has brought Mr. Chomsky's findings to everyman. In “The Language Instinct” he has gathered persuasive data from such diverse fields as cognitive neuroscience, developmental psychology and speech therapy to make his points, and when he disagrees with Mr. Chomsky he tells you so. . . .
For Mr. Chomsky and Mr. Pinker, somewhere in the human brain there is a complex set of neural circuits that have been programmed with “super-rules” (making up what Mr. Chomsky calls “universal grammar”), and that these rules are unconscious and instinctive. A half-century ago, this would have been pooh-poohed as a “black box” theory, since one could not actually pinpoint this grammatical faculty in a specific part of the brain, or describe its functioning. But now things are different. Neurosurgeons [have now found that this] “black box” is situated in and around Broca’s area, on the left side of the forebrain. . . . 
Unlike Mr. Chomsky, Mr. Pinker firmly places the wiring of the brain for language within the framework of Darwinian natural selection and evolution. He effectively disposes of all claims that intelligent nonhuman primates like chimps have any abilities to learn and use language. It is not that chimps lack the vocal apparatus to speak; it is just that their brains are unable to produce or use grammar. On the other hand, the “language instinct,” when it first appeared among our most distant hominid ancestors, must have given them a selective reproductive advantage over their competitors (including the ancestral chimps). . . .
So according to Mr. Pinker, the roots of language must be in the genes, but there cannot be a “grammar gene” any more than there can be a gene for the heart or any other complex body structure. This proposition will undoubtedly raise the hackles of some behavioral psychologists and anthropologists, for it apparently contradicts the liberal idea that human behavior may be changed for the better by improvements in culture and environment, and it might seem to invite the twin bugaboos of biological determinism and racism. Yet Mr. Pinker stresses one point that should allay such fears. Even though there are 4,000 to 6,000 languages today, they are all sufficiently alike to be considered one language by an extraterrestrial observer. In other words, most of the diversity of the world’s cultures, so beloved to anthropologists, is superficial and minor compared to the similarities. Racial differences are literally only “skin deep.” The fundamental unity of humanity is the theme of Mr. Chomsky's universal grammar, and of this exciting book.

\vspace{1em}
\textbf{Question 21}

**Five jumbled up sentences, related to a topic, are given below. Four of them can be put together to form a coherent paragraph. Identify the odd one out and key in the number of the sentence as your answer:**  
1. They often include a foundation course on navigating capitalism with Chinese characteristics and have replaced typical cases from US corporates with a focus on how Western theories apply to China’s buzzing local firms.  
2. The best Chinese business schools look like their Western rivals but are now growing distinct in terms of what they teach and the career boost they offer.  
3. Western schools have enhanced their offerings with double degrees, popular with domestic and overseas students alike—and boosted the prestige of their Chinese partners.  
4. For students, a big draw is the chance to rub shoulders with captains of China’s private sector.  
5. Their business courses now largely cater to the growing demand from China Inc which has become more global, richer and ready to recruit from this sinocentric student body.
Backspace  
789  
456  
123  
0.-  
Clear All  

Submit
[](https://cracku.in/21-five-jumbled-up-sentences-related-to-a-topic-are-g-x-cat-2021-slot-3)

\begin{enumerate}[label=\Alph*.]
\end{enumerate}
\vspace{1em}
\textbf{Solution:}

A brief reading of the sentences tells us that the paragraph is about Chinese business schools and how they stand in comparison to their western counterparts. 2 mentions that thought they have a similar outlook, Chinese business schools have a different curriculum and are also different in what they have to offer. 1, 4, and 5 further talk about the peculiarity of the Chinese schools.
3, however, runs tangent to the discussion. It shifts the focus from Chinese schools and describes western schools. Hence, 3 is out of the context here.

\end{problem}

\newpage
\begin{problem}{}{}
% Subject: varc
\textbf{Instructions}

Instructions   

The passage below is accompanied by a set of questions. Choose the best answer to each question.
Starting in 1957, [Noam Chomsky] proclaimed a new doctrine: Language, that most human of all attributes, was innate. The grammatical faculty was built into the infant brain, and your average 3-year-old was not a mere apprentice in the great enterprise of absorbing English from his or her parents, but a “linguistic genius.” Since this message was couched in terms of Chomskyan theoretical linguistics, in discourse so opaque that it was nearly incomprehensible even to some scholars, many people did not hear it. Now, in a brilliant, witty and altogether satisfying book, Mr. Chomsky's colleague Steven Pinker . . . has brought Mr. Chomsky's findings to everyman. In “The Language Instinct” he has gathered persuasive data from such diverse fields as cognitive neuroscience, developmental psychology and speech therapy to make his points, and when he disagrees with Mr. Chomsky he tells you so. . . .
For Mr. Chomsky and Mr. Pinker, somewhere in the human brain there is a complex set of neural circuits that have been programmed with “super-rules” (making up what Mr. Chomsky calls “universal grammar”), and that these rules are unconscious and instinctive. A half-century ago, this would have been pooh-poohed as a “black box” theory, since one could not actually pinpoint this grammatical faculty in a specific part of the brain, or describe its functioning. But now things are different. Neurosurgeons [have now found that this] “black box” is situated in and around Broca’s area, on the left side of the forebrain. . . . 
Unlike Mr. Chomsky, Mr. Pinker firmly places the wiring of the brain for language within the framework of Darwinian natural selection and evolution. He effectively disposes of all claims that intelligent nonhuman primates like chimps have any abilities to learn and use language. It is not that chimps lack the vocal apparatus to speak; it is just that their brains are unable to produce or use grammar. On the other hand, the “language instinct,” when it first appeared among our most distant hominid ancestors, must have given them a selective reproductive advantage over their competitors (including the ancestral chimps). . . .
So according to Mr. Pinker, the roots of language must be in the genes, but there cannot be a “grammar gene” any more than there can be a gene for the heart or any other complex body structure. This proposition will undoubtedly raise the hackles of some behavioral psychologists and anthropologists, for it apparently contradicts the liberal idea that human behavior may be changed for the better by improvements in culture and environment, and it might seem to invite the twin bugaboos of biological determinism and racism. Yet Mr. Pinker stresses one point that should allay such fears. Even though there are 4,000 to 6,000 languages today, they are all sufficiently alike to be considered one language by an extraterrestrial observer. In other words, most of the diversity of the world’s cultures, so beloved to anthropologists, is superficial and minor compared to the similarities. Racial differences are literally only “skin deep.” The fundamental unity of humanity is the theme of Mr. Chomsky's universal grammar, and of this exciting book.

\vspace{1em}
\textbf{Question 22}

**The passage given below is followed by four alternate summaries. Choose the option that best captures the essence of the passage.**
The human mind is wired to see patterns. Not only does the brain process information as it comes in, it also stores insights from all our past experiences. Every interaction, happy or sad, is catalogued in our memory. Intuition draws from that deep memory well to inform our decisions going forward. In other words, intuitive decisions are based on data, and not contrary to data as many would like to assume. When we subconsciously spot patterns, the body starts firing neurochemicals in both the brain and gut. These “somatic markers” are what give us that instant sense that something is right … or that it’s off. Not only are these automatic processes faster than rational thought, but our intuition draws from decades of diverse qualitative experience (sights, sounds, interactions, etc.) - a wholly human feature that big data alone could never accomplish.

\begin{enumerate}[label=\Alph*.]
    \item Intuition is infinitely richer than big data which is based on rational thought and accomplishes more than what big data can.
    \item Intuitions are automatic processes and are therefore faster than rational thought, and so decisions based on them are better.
    \item Intuition draws from deep memory, and may not be related to data, but to decades of diverse qualitative experience.
    \item Intuitions are neuro-chemical firings based on pattern recognition and draw upon a rich and vast database of experiences.
[](https://cracku.in/22-the-passage-given-below-is-followed-by-four-altern-x-cat-2021-slot-3)
\end{enumerate}
\vspace{1em}
\textbf{Solution:}

The main points of the paragraph are as follows:
1. Intuition draws from a vast array of memories that our brain keeps in store.
2. When our brain recognises a pattern from past memories, neuron firing starts, which gives us the gut feeling of intuition. 
A: Distortion: The passage does not give any detail about big data being based on rational thought.
B: Out of scope. The paragraph does not allude to whether the decisions based on intuition are better or worse.
C: Incorrect: The passage says that intuitive decisions are based on data.
D: Correctly covers the mentioned points and hence, is the answer.

\end{problem}

\newpage
\begin{problem}{}{}
% Subject: varc
\textbf{Instructions}

Instructions   

The passage below is accompanied by a set of questions. Choose the best answer to each question.
Starting in 1957, [Noam Chomsky] proclaimed a new doctrine: Language, that most human of all attributes, was innate. The grammatical faculty was built into the infant brain, and your average 3-year-old was not a mere apprentice in the great enterprise of absorbing English from his or her parents, but a “linguistic genius.” Since this message was couched in terms of Chomskyan theoretical linguistics, in discourse so opaque that it was nearly incomprehensible even to some scholars, many people did not hear it. Now, in a brilliant, witty and altogether satisfying book, Mr. Chomsky's colleague Steven Pinker . . . has brought Mr. Chomsky's findings to everyman. In “The Language Instinct” he has gathered persuasive data from such diverse fields as cognitive neuroscience, developmental psychology and speech therapy to make his points, and when he disagrees with Mr. Chomsky he tells you so. . . .
For Mr. Chomsky and Mr. Pinker, somewhere in the human brain there is a complex set of neural circuits that have been programmed with “super-rules” (making up what Mr. Chomsky calls “universal grammar”), and that these rules are unconscious and instinctive. A half-century ago, this would have been pooh-poohed as a “black box” theory, since one could not actually pinpoint this grammatical faculty in a specific part of the brain, or describe its functioning. But now things are different. Neurosurgeons [have now found that this] “black box” is situated in and around Broca’s area, on the left side of the forebrain. . . . 
Unlike Mr. Chomsky, Mr. Pinker firmly places the wiring of the brain for language within the framework of Darwinian natural selection and evolution. He effectively disposes of all claims that intelligent nonhuman primates like chimps have any abilities to learn and use language. It is not that chimps lack the vocal apparatus to speak; it is just that their brains are unable to produce or use grammar. On the other hand, the “language instinct,” when it first appeared among our most distant hominid ancestors, must have given them a selective reproductive advantage over their competitors (including the ancestral chimps). . . .
So according to Mr. Pinker, the roots of language must be in the genes, but there cannot be a “grammar gene” any more than there can be a gene for the heart or any other complex body structure. This proposition will undoubtedly raise the hackles of some behavioral psychologists and anthropologists, for it apparently contradicts the liberal idea that human behavior may be changed for the better by improvements in culture and environment, and it might seem to invite the twin bugaboos of biological determinism and racism. Yet Mr. Pinker stresses one point that should allay such fears. Even though there are 4,000 to 6,000 languages today, they are all sufficiently alike to be considered one language by an extraterrestrial observer. In other words, most of the diversity of the world’s cultures, so beloved to anthropologists, is superficial and minor compared to the similarities. Racial differences are literally only “skin deep.” The fundamental unity of humanity is the theme of Mr. Chomsky's universal grammar, and of this exciting book.

\vspace{1em}
\textbf{Question 23}

**The passage given below is followed by four alternate summaries. Choose the option that best captures the essence of the passage.**
Brazil’s growth rate has been low, yet most Brazilians say their financial situation has improved, and they expect it to get even better. This is because most incomes are rising fast, with higher minimum wages and very low unemployment. The result is falling inequality and a growing middle class — the result of economic stabilization, improved social security and universal primary education. But despite recent improvements the Brazilian economy is still painfully unequal, with poor Brazilians paying the biggest share of their income in taxes and getting the least back in government services.

\begin{enumerate}[label=\Alph*.]
    \item Economic reforms have benefitted many Brazilians, but they are unaware of the impending problems from rising inequalities in their society.
    \item Good economic indicators have masked the unfair taxation of the poor that is likely to destabilise the Brazilian economy in the next few years.
    \item Most Brazilians feel they have benefitted from recent economic events, but the poor continue to be dealt unfairly by the state.
    \item With rising incomes and falling unemployment, most Brazilians are being misled into thinking that their economy is doing well.
[](https://cracku.in/23-the-passage-given-below-is-followed-by-four-altern-x-cat-2021-slot-3)
\end{enumerate}
\vspace{1em}
\textbf{Solution:}

The main points of the paragraph are:
1. The Brazilian economy has been stagnant, but the popular perception is that the times have improved.
2. The reasons are falling inequality and other important services.
3. Despite this, the economy is brutally unequal.
A: Ignorance on the part of Brazilians is not implied. What the author is saying that though things have improved for the ones who say so, others are still being dealt a rough hand.
B: The paragraph does not imply that the good economic indicators are being used as subterfuge to cover up the prevailing inequality.
C: Comes the closest in capturing the three points. Hence, is the answer.
D: It has a problem similar to Option A. Things have improved for that section of people. They are not ignorant, nor are they being misled into believing something.

\end{problem}

\newpage
\begin{problem}{}{}
% Subject: varc
\textbf{Instructions}

Instructions   

The passage below is accompanied by a set of questions. Choose the best answer to each question.
Starting in 1957, [Noam Chomsky] proclaimed a new doctrine: Language, that most human of all attributes, was innate. The grammatical faculty was built into the infant brain, and your average 3-year-old was not a mere apprentice in the great enterprise of absorbing English from his or her parents, but a “linguistic genius.” Since this message was couched in terms of Chomskyan theoretical linguistics, in discourse so opaque that it was nearly incomprehensible even to some scholars, many people did not hear it. Now, in a brilliant, witty and altogether satisfying book, Mr. Chomsky's colleague Steven Pinker . . . has brought Mr. Chomsky's findings to everyman. In “The Language Instinct” he has gathered persuasive data from such diverse fields as cognitive neuroscience, developmental psychology and speech therapy to make his points, and when he disagrees with Mr. Chomsky he tells you so. . . .
For Mr. Chomsky and Mr. Pinker, somewhere in the human brain there is a complex set of neural circuits that have been programmed with “super-rules” (making up what Mr. Chomsky calls “universal grammar”), and that these rules are unconscious and instinctive. A half-century ago, this would have been pooh-poohed as a “black box” theory, since one could not actually pinpoint this grammatical faculty in a specific part of the brain, or describe its functioning. But now things are different. Neurosurgeons [have now found that this] “black box” is situated in and around Broca’s area, on the left side of the forebrain. . . . 
Unlike Mr. Chomsky, Mr. Pinker firmly places the wiring of the brain for language within the framework of Darwinian natural selection and evolution. He effectively disposes of all claims that intelligent nonhuman primates like chimps have any abilities to learn and use language. It is not that chimps lack the vocal apparatus to speak; it is just that their brains are unable to produce or use grammar. On the other hand, the “language instinct,” when it first appeared among our most distant hominid ancestors, must have given them a selective reproductive advantage over their competitors (including the ancestral chimps). . . .
So according to Mr. Pinker, the roots of language must be in the genes, but there cannot be a “grammar gene” any more than there can be a gene for the heart or any other complex body structure. This proposition will undoubtedly raise the hackles of some behavioral psychologists and anthropologists, for it apparently contradicts the liberal idea that human behavior may be changed for the better by improvements in culture and environment, and it might seem to invite the twin bugaboos of biological determinism and racism. Yet Mr. Pinker stresses one point that should allay such fears. Even though there are 4,000 to 6,000 languages today, they are all sufficiently alike to be considered one language by an extraterrestrial observer. In other words, most of the diversity of the world’s cultures, so beloved to anthropologists, is superficial and minor compared to the similarities. Racial differences are literally only “skin deep.” The fundamental unity of humanity is the theme of Mr. Chomsky's universal grammar, and of this exciting book.

\vspace{1em}
\textbf{Question 24}

**The four sentences (labelled 1, 2, 3 and 4) below, when properly sequenced would yield a coherent paragraph. Decide on the proper sequencing of the order of the sentences and key in the sequence of the four numbers as your answer:**  
1. Restitution of artefacts to original cultures could face legal obstacles, as many Western museums are legally prohibited from disposing off their collections.  
2. This is in response to countries like Nigeria, which are pressurising European museums to return their precious artefacts looted by colonisers in the past.  
3. Museums in Europe today are struggling to come to terms with their colonial legacy, some taking steps to return artefacts but not wanting to lose their prized collections.  
4. Legal hurdles notwithstanding, politicians and institutions in France and Germany would now like to defuse the colonial time bombs, and are now backing the return of part of their holdings.
Backspace  
789  
456  
123  
0.-  
Clear All  

Submit
[](https://cracku.in/24-the-four-sentences-labelled-1-2-3-and-4-below-when-x-cat-2021-slot-3)

\begin{enumerate}[label=\Alph*.]
\end{enumerate}
\vspace{1em}
\textbf{Solution:}

A brief reading of the sentences suggests that the paragraph is about restitution of articles that were taken from colonies. 3 introduces the topic at hand: European museums are trying to come to terms with their colonial legacy by returning artefacts. 2 follows 3, providing the reason why this development has taken place. 1 then further describes the current situation: the legal obstacles this action could face. 4 concludes the paragraph by saying that other than these legal obstacles, institutions in France and Germany are seriously backing the move. Hence, the correct arrangement is 3214.
[](https://cracku.in/downloads/13896)
  *  
  *

\end{problem}

\newpage
\begin{problem}{}{}
% Subject: varc
\textbf{Instructions}

Instructions   

**The passage below is accompanied by a set of questions. Choose the best answer to each question.**
**Comprehension** :
Stoicism was founded in 300 BC by the Greek philosopher Zeno and survived into the Roman era until about AD 300. According to the Stoics, emotions consist of two movements. The first movement is the immediate feeling and other reactions (e.g., physiological response) that occur when a stimulus or event occurs. For instance, consider what could have happened if an army general accused Marcus Aurelius of treason in front of other officers. The first movement for Marcus may have been (internal) surprise and anger in response to this insult, accompanied perhaps by some involuntary physiological and expressive responses such as face flushing and a movement of the eyebrows. The second movement is what one does next about the emotion. Second movement behaviors occur after thinking and are under one’s control. Examples of second movements for Marcus might have included a plot to seek revenge, actions signifying deference and appeasement, or perhaps proceeding as he would have proceeded whether or not this event occurred: continuing to lead the Romans in a way that Marcus Aurelius believed best benefited them. In the Stoic view, choosing a reasoned, unemotional response as the second movement is the only appropriate response.
The Stoics believed that to live the good life and be a good person, we need to free ourselves of nearly all desires such as too much desire for money, power, or sexual gratification. Prior to second movements, we can consider what is important in life. Money, power, and excessive sexual gratification are not important. Character, rationality, and kindness are important. The Epicureans, first associated with the Greek philosopher Epicurus . . . held a similar view, believing that people should enjoy simple pleasures, such as good conversation, friendship, food, and wine, but not be indulgent in these pursuits and not follow passion for those things that hold no real value like power and money. As Oatley (2004) states, “the Epicureans articulated a view—enjoyment of relationship with friends, of things that are real rather than illusory, simple rather than artificially inflated, possible rather than vanishingly unlikely—that is certainly relevant today” . . . In sum, these ancient Greek and Roman philosophers saw emotions, especially strong ones, as potentially dangerous. They viewed emotions as experiences that needed to be [reined] in and controlled.
As Oatley (2004) points out, the Stoic idea bears some similarity to Buddhism. Buddha, living in India in the 6th century BC, argued for cultivating a certain attitude that decreases the probability of (in Stoic terms) destructive second movements. Through meditation and the right attitude, one allows emotions to happen to oneself (it is impossible to prevent this), but one is advised to observe the emotions without necessarily acting on them; one achieves some distance and decides what has value and what does not have value. Additionally, the  
Stoic idea of developing virtue in oneself, of becoming a good person, which the Stoics believed we could do because we have a touch of the divine, laid the foundation for the three monotheistic religions: Judaism, Christianity, and Islam . . . As with Stoicism, tenets of these religions include controlling our emotions lest we engage in sinful behavior.

\vspace{1em}
\textbf{Question 1}

“Through meditation and the right attitude, one allows emotions to happen to oneself (it is impossible to prevent this), but one is advised to observe the emotions without necessarily acting on them; one achieves some distance and decides what has value and what does not have value.” 
In the context of the passage, which one of the following is not a possible implication of the quoted statement?

\begin{enumerate}[label=\Alph*.]
    \item “Meditation and the right attitude”, in this instance, implies an initially passive reception of all experiences.
    \item Meditation allows certain out-of-body experiences that permit us to gain the distance necessary to control our emotions.
    \item The observation of emotions in a distant manner corresponds to the second movement referred to earlier in the passage.
    \item Emotional responses can make it difficult to distinguish valuable experiences from valueless experiences.
[](https://cracku.in/1-through-meditation-and-the-right-attitude-one-allo-x-cat-2022-slot-1)
\end{enumerate}
\vspace{1em}
\textbf{Solution:}

_Option A_ : This is a possible implication because the passage states that through meditation and the right attitude, one allows emotions to happen to themselves, which implies a passive reception of those experiences._  
_
_Option B_ : The passage does not state or imply that meditation allows for certain out-of-body experiences or that it allows an individual to gain distance from their emotions in that way. The passage specifically highlights that through meditation and the right attitude, one allows emotions to happen to themselves and then observes those emotions without necessarily acting on them, which allows one to achieve some distance and decide what has value and what does not have value. It does not mention anything about out-of-body experiences or any other kind of distance that is not achieved through the act of observing emotions.
_Option C_ : This is a possible implication because the passage states that the second movement is what one does next about the emotion and that it occurs after thinking and is under one's control. Observing emotions in a distant manner, as described in the quote, would involve thinking and control and would therefore correspond to the second movement referred to earlier.
_Option D_ : This is a possible implication because the quote mentions that observing emotions allows an individual to decide what has value and what does not have value, implying that emotional responses can make this distinction difficult.
Hence, Option B is the correct answer.

\end{problem}

\newpage
\begin{problem}{}{}
% Subject: varc
\textbf{Instructions}

Instructions   

**The passage below is accompanied by a set of questions. Choose the best answer to each question.**
**Comprehension** :
Stoicism was founded in 300 BC by the Greek philosopher Zeno and survived into the Roman era until about AD 300. According to the Stoics, emotions consist of two movements. The first movement is the immediate feeling and other reactions (e.g., physiological response) that occur when a stimulus or event occurs. For instance, consider what could have happened if an army general accused Marcus Aurelius of treason in front of other officers. The first movement for Marcus may have been (internal) surprise and anger in response to this insult, accompanied perhaps by some involuntary physiological and expressive responses such as face flushing and a movement of the eyebrows. The second movement is what one does next about the emotion. Second movement behaviors occur after thinking and are under one’s control. Examples of second movements for Marcus might have included a plot to seek revenge, actions signifying deference and appeasement, or perhaps proceeding as he would have proceeded whether or not this event occurred: continuing to lead the Romans in a way that Marcus Aurelius believed best benefited them. In the Stoic view, choosing a reasoned, unemotional response as the second movement is the only appropriate response.
The Stoics believed that to live the good life and be a good person, we need to free ourselves of nearly all desires such as too much desire for money, power, or sexual gratification. Prior to second movements, we can consider what is important in life. Money, power, and excessive sexual gratification are not important. Character, rationality, and kindness are important. The Epicureans, first associated with the Greek philosopher Epicurus . . . held a similar view, believing that people should enjoy simple pleasures, such as good conversation, friendship, food, and wine, but not be indulgent in these pursuits and not follow passion for those things that hold no real value like power and money. As Oatley (2004) states, “the Epicureans articulated a view—enjoyment of relationship with friends, of things that are real rather than illusory, simple rather than artificially inflated, possible rather than vanishingly unlikely—that is certainly relevant today” . . . In sum, these ancient Greek and Roman philosophers saw emotions, especially strong ones, as potentially dangerous. They viewed emotions as experiences that needed to be [reined] in and controlled.
As Oatley (2004) points out, the Stoic idea bears some similarity to Buddhism. Buddha, living in India in the 6th century BC, argued for cultivating a certain attitude that decreases the probability of (in Stoic terms) destructive second movements. Through meditation and the right attitude, one allows emotions to happen to oneself (it is impossible to prevent this), but one is advised to observe the emotions without necessarily acting on them; one achieves some distance and decides what has value and what does not have value. Additionally, the  
Stoic idea of developing virtue in oneself, of becoming a good person, which the Stoics believed we could do because we have a touch of the divine, laid the foundation for the three monotheistic religions: Judaism, Christianity, and Islam . . . As with Stoicism, tenets of these religions include controlling our emotions lest we engage in sinful behavior.

\vspace{1em}
\textbf{Question 2}

Which one of the following statements would be an accurate inference from the example of Marcus Aurelius?

\begin{enumerate}[label=\Alph*.]
    \item Marcus Aurelius was humiliated by the accusation of treason in front of the other officers.
    \item Marcus Aurelius was a Stoic whose philosophy survived into the Roman era.
    \item Marcus Aurelius plotted revenge in his quest for justice.
    \item Marcus Aurelius was one of the leaders of the Roman army.
[](https://cracku.in/2-which-one-of-the-following-statements-would-be-an--x-cat-2022-slot-1)
\end{enumerate}
\vspace{1em}
\textbf{Solution:}

Based on the discussion, Option D is the correct answer: the passage describes an example of what might have happened if an army general accused Marcus Aurelius of treason in front of other officers, implying that Marcus Aurelius was a leader in the Roman army. The other options are neither mentioned nor implied in the passage and are therefore not supported by the information provided. 
Option A: is incorrect because the passage does not mention anything about Marcus Aurelius feeling humiliated or embarrassed by the accusation; it only describes the immediate feeling and other reactions that may have occurred in response to the stimulus of the accusation, such as surprise and anger.
Option B: is incorrect since the author does not label Marcus Aurelius as a Stoic or associate him with the philosophy of Stoicism in any way; he only uses him as an example of what might have happened in a specific situation involving an army general accusing him of treason.
Option C: is incorrect since the passage does not state that Marcus Aurelius was plotting revenge or seeking justice; it only mentions that one of the potential second movements that Marcus Aurelius might have chosen in response to the accusation could have been a plot to seek revenge. However, it does not state that this is what actually happened or that it was the only possible second movement that Marcus Aurelius could have chosen.

\end{problem}

\newpage
\begin{problem}{}{}
% Subject: varc
\textbf{Instructions}

Instructions   

**The passage below is accompanied by a set of questions. Choose the best answer to each question.**
**Comprehension** :
Stoicism was founded in 300 BC by the Greek philosopher Zeno and survived into the Roman era until about AD 300. According to the Stoics, emotions consist of two movements. The first movement is the immediate feeling and other reactions (e.g., physiological response) that occur when a stimulus or event occurs. For instance, consider what could have happened if an army general accused Marcus Aurelius of treason in front of other officers. The first movement for Marcus may have been (internal) surprise and anger in response to this insult, accompanied perhaps by some involuntary physiological and expressive responses such as face flushing and a movement of the eyebrows. The second movement is what one does next about the emotion. Second movement behaviors occur after thinking and are under one’s control. Examples of second movements for Marcus might have included a plot to seek revenge, actions signifying deference and appeasement, or perhaps proceeding as he would have proceeded whether or not this event occurred: continuing to lead the Romans in a way that Marcus Aurelius believed best benefited them. In the Stoic view, choosing a reasoned, unemotional response as the second movement is the only appropriate response.
The Stoics believed that to live the good life and be a good person, we need to free ourselves of nearly all desires such as too much desire for money, power, or sexual gratification. Prior to second movements, we can consider what is important in life. Money, power, and excessive sexual gratification are not important. Character, rationality, and kindness are important. The Epicureans, first associated with the Greek philosopher Epicurus . . . held a similar view, believing that people should enjoy simple pleasures, such as good conversation, friendship, food, and wine, but not be indulgent in these pursuits and not follow passion for those things that hold no real value like power and money. As Oatley (2004) states, “the Epicureans articulated a view—enjoyment of relationship with friends, of things that are real rather than illusory, simple rather than artificially inflated, possible rather than vanishingly unlikely—that is certainly relevant today” . . . In sum, these ancient Greek and Roman philosophers saw emotions, especially strong ones, as potentially dangerous. They viewed emotions as experiences that needed to be [reined] in and controlled.
As Oatley (2004) points out, the Stoic idea bears some similarity to Buddhism. Buddha, living in India in the 6th century BC, argued for cultivating a certain attitude that decreases the probability of (in Stoic terms) destructive second movements. Through meditation and the right attitude, one allows emotions to happen to oneself (it is impossible to prevent this), but one is advised to observe the emotions without necessarily acting on them; one achieves some distance and decides what has value and what does not have value. Additionally, the  
Stoic idea of developing virtue in oneself, of becoming a good person, which the Stoics believed we could do because we have a touch of the divine, laid the foundation for the three monotheistic religions: Judaism, Christianity, and Islam . . . As with Stoicism, tenets of these religions include controlling our emotions lest we engage in sinful behavior.

\vspace{1em}
\textbf{Question 3}

Which one of the following statements, if false, could be seen as contradicting the facts/arguments in the passage?

\begin{enumerate}[label=\Alph*.]
    \item Despite practising meditation and cultivating the right attitude, emotions cannot ever be controlled.
    \item The Greek philosopher Zeno survived into the Roman era until about AD 300.
    \item In the Epicurean view, indulging in simple pleasures is not desirable.
    \item In the Stoic view, choosing a reasoned, unemotional response as the first movement is an appropriate response to emotional situations.
[](https://cracku.in/3-which-one-of-the-following-statements-if-false-cou-x-cat-2022-slot-1)
\end{enumerate}
\vspace{1em}
\textbf{Solution:}

In this question, we need to find a statement which is in line with the ideas given in the passage(then, if it is false, it will contradict the passage).
_Option A_ : "_Through meditation and the right attitude, one allows emotions to happen to oneself (it is impossible to prevent this), but one is advised to observe the emotions without necessarily acting on them; one achieves some distance and decides what has value and what does not have value."_
The above excerpt was written to describe the similarities between Stoicism and Buddhism. Since option A is not in line with the above excerpt's idea, it is not the correct option.
_Option B_ : This option can easily be refuted on the basis of the information given in the first line of the passage. It cannot be inferred that Zeno survived into the Roman era until about AD 300.
_Option C_ : "_The Epicureans, first associated with the Greek philosopher Epicurus . . . held a similar view, believing that people should enjoy simple pleasures, such as good conversation, friendship, food, and wine,_**but not be indulgent in these pursuits** and not follow passion for those things that hold no real value like power and money.__"
Since this option reiterates the idea mentioned in the underlined portion of the above excerpt, this is the correct option.
Thus, the correct option is C.

\end{problem}

\newpage
\begin{problem}{}{}
% Subject: varc
\textbf{Instructions}

Instructions   

**The passage below is accompanied by a set of questions. Choose the best answer to each question.**
**Comprehension** :
Stoicism was founded in 300 BC by the Greek philosopher Zeno and survived into the Roman era until about AD 300. According to the Stoics, emotions consist of two movements. The first movement is the immediate feeling and other reactions (e.g., physiological response) that occur when a stimulus or event occurs. For instance, consider what could have happened if an army general accused Marcus Aurelius of treason in front of other officers. The first movement for Marcus may have been (internal) surprise and anger in response to this insult, accompanied perhaps by some involuntary physiological and expressive responses such as face flushing and a movement of the eyebrows. The second movement is what one does next about the emotion. Second movement behaviors occur after thinking and are under one’s control. Examples of second movements for Marcus might have included a plot to seek revenge, actions signifying deference and appeasement, or perhaps proceeding as he would have proceeded whether or not this event occurred: continuing to lead the Romans in a way that Marcus Aurelius believed best benefited them. In the Stoic view, choosing a reasoned, unemotional response as the second movement is the only appropriate response.
The Stoics believed that to live the good life and be a good person, we need to free ourselves of nearly all desires such as too much desire for money, power, or sexual gratification. Prior to second movements, we can consider what is important in life. Money, power, and excessive sexual gratification are not important. Character, rationality, and kindness are important. The Epicureans, first associated with the Greek philosopher Epicurus . . . held a similar view, believing that people should enjoy simple pleasures, such as good conversation, friendship, food, and wine, but not be indulgent in these pursuits and not follow passion for those things that hold no real value like power and money. As Oatley (2004) states, “the Epicureans articulated a view—enjoyment of relationship with friends, of things that are real rather than illusory, simple rather than artificially inflated, possible rather than vanishingly unlikely—that is certainly relevant today” . . . In sum, these ancient Greek and Roman philosophers saw emotions, especially strong ones, as potentially dangerous. They viewed emotions as experiences that needed to be [reined] in and controlled.
As Oatley (2004) points out, the Stoic idea bears some similarity to Buddhism. Buddha, living in India in the 6th century BC, argued for cultivating a certain attitude that decreases the probability of (in Stoic terms) destructive second movements. Through meditation and the right attitude, one allows emotions to happen to oneself (it is impossible to prevent this), but one is advised to observe the emotions without necessarily acting on them; one achieves some distance and decides what has value and what does not have value. Additionally, the  
Stoic idea of developing virtue in oneself, of becoming a good person, which the Stoics believed we could do because we have a touch of the divine, laid the foundation for the three monotheistic religions: Judaism, Christianity, and Islam . . . As with Stoicism, tenets of these religions include controlling our emotions lest we engage in sinful behavior.

\vspace{1em}
\textbf{Question 4}

On the basis of the passage, which one of the following statements can be regarded as true?

\begin{enumerate}[label=\Alph*.]
    \item The Stoics valorised the pursuit of money, power, and sexual gratification.
    \item The Stoic influences can be seen in multiple religions.
    \item The Epicureans believed in controlling all emotions.
    \item There were no Stoics in India at the time of the Roman civilisation.
[](https://cracku.in/4-on-the-basis-of-the-passage-which-one-of-the-follo-x-cat-2022-slot-1)
\end{enumerate}
\vspace{1em}
\textbf{Solution:}

The passage states that "the Stoic idea of developing virtue in oneself, of becoming a good person, which the Stoics believed we could do because we have a touch of the divine, laid the foundation for the three monotheistic religions: Judaism, Christianity, and Islam." This aligns with the claim in Option B. Options A and C are incorrect as the passage states that the Stoics believed in freeing oneself of nearly all desires, including excessive desires for money, power, and sexual gratification [A], and that the Epicureans believed in enjoying simple pleasures but not being indulgent or pursuing things with no real value [not sufficient to validate C]. Similarly, we cannot substantiate the statement in D.
Hence, Option B is the correct answer.

Instructions   

**The passage below is accompanied by a set of questions. Choose the best answer to each question.**
**Comprehension** :
The Chinese have two different concepts of a copy. Fangzhipin . . . are imitations where the difference from the original is obvious. These are small models or copies that can be purchased in a museum shop, for example. The second concept for a copy is fuzhipin . . . They are exact reproductions of the original, which, for the Chinese, are of equal value to the original. It has absolutely no negative connotations. The discrepancy with regard to the understanding of what a copy is has often led to misunderstandings and arguments between China and Western museums. The Chinese often send copies abroad instead of originals, in the firm belief that they are not essentially different from the originals. The rejection that then comes from the Western museums is perceived by the Chinese as an insult. . . .
The Far Eastern notion of identity is also very confusing to the Western observer. The Ise Grand Shrine [in Japan] is 1,300 years old for the millions of Japanese people who go there on pilgrimage every year. But in reality this temple complex is completely rebuilt from scratch every 20 years. . . .
The cathedral of Freiburg Minster in southwest Germany is covered in scaffolding almost all year round. The sandstone from which it is built is a very soft, porous material that does not withstand natural erosion by rain and wind. After a while, it crumbles. As a result, the cathedral is continually being examined for damage, and eroded stones are replaced. And in the cathedral’s dedicated workshop, copies of the damaged sandstone figures are constantly being produced. Of course, attempts are made to preserve the stones from the Middle Ages for as long as possible. But at some point they, too, are removed and replaced with new stones.
Fundamentally, this is the same operation as with the Japanese shrine, except in this case the production of a replica takes place very slowly and over long periods of time. . . . In the field of art as well, the idea of an unassailable original developed historically in the Western world. Back in the 17th century [in the West], excavated artworks from antiquity were treated quite differently from today. They were not restored in a way that was faithful to the original. Instead, there was massive intervention in these works, changing their appearance. . . .
It is probably this intellectual position that explains why Asians have far fewer scruples about cloning than Europeans. The South Korean cloning researcher Hwang Woo-suk, who attracted worldwide attention with his cloning experiments in 2004, is a Buddhist. He found a great deal of support and followers among Buddhists, while Christians called for a ban on human cloning. . . . Hwang legitimised his cloning experiments with his religious affiliation: ‘I am Buddhist, and I have no philosophical problem with cloning. And as you know, the basis of Buddhism is that life is recycled through reincarnation. In some ways, I think, therapeutic cloning restarts the circle of life.’

\end{problem}

\newpage
\begin{problem}{}{}
% Subject: varc
\textbf{Instructions}

Instructions   

**The passage below is accompanied by a set of questions. Choose the best answer to each question.**
**Comprehension** :
Stoicism was founded in 300 BC by the Greek philosopher Zeno and survived into the Roman era until about AD 300. According to the Stoics, emotions consist of two movements. The first movement is the immediate feeling and other reactions (e.g., physiological response) that occur when a stimulus or event occurs. For instance, consider what could have happened if an army general accused Marcus Aurelius of treason in front of other officers. The first movement for Marcus may have been (internal) surprise and anger in response to this insult, accompanied perhaps by some involuntary physiological and expressive responses such as face flushing and a movement of the eyebrows. The second movement is what one does next about the emotion. Second movement behaviors occur after thinking and are under one’s control. Examples of second movements for Marcus might have included a plot to seek revenge, actions signifying deference and appeasement, or perhaps proceeding as he would have proceeded whether or not this event occurred: continuing to lead the Romans in a way that Marcus Aurelius believed best benefited them. In the Stoic view, choosing a reasoned, unemotional response as the second movement is the only appropriate response.
The Stoics believed that to live the good life and be a good person, we need to free ourselves of nearly all desires such as too much desire for money, power, or sexual gratification. Prior to second movements, we can consider what is important in life. Money, power, and excessive sexual gratification are not important. Character, rationality, and kindness are important. The Epicureans, first associated with the Greek philosopher Epicurus . . . held a similar view, believing that people should enjoy simple pleasures, such as good conversation, friendship, food, and wine, but not be indulgent in these pursuits and not follow passion for those things that hold no real value like power and money. As Oatley (2004) states, “the Epicureans articulated a view—enjoyment of relationship with friends, of things that are real rather than illusory, simple rather than artificially inflated, possible rather than vanishingly unlikely—that is certainly relevant today” . . . In sum, these ancient Greek and Roman philosophers saw emotions, especially strong ones, as potentially dangerous. They viewed emotions as experiences that needed to be [reined] in and controlled.
As Oatley (2004) points out, the Stoic idea bears some similarity to Buddhism. Buddha, living in India in the 6th century BC, argued for cultivating a certain attitude that decreases the probability of (in Stoic terms) destructive second movements. Through meditation and the right attitude, one allows emotions to happen to oneself (it is impossible to prevent this), but one is advised to observe the emotions without necessarily acting on them; one achieves some distance and decides what has value and what does not have value. Additionally, the  
Stoic idea of developing virtue in oneself, of becoming a good person, which the Stoics believed we could do because we have a touch of the divine, laid the foundation for the three monotheistic religions: Judaism, Christianity, and Islam . . . As with Stoicism, tenets of these religions include controlling our emotions lest we engage in sinful behavior.

\vspace{1em}
\textbf{Question 5}

Based on the passage, which one of the following copies would a Chinese museum be unlikely to consider as having less value than the original?

\begin{enumerate}[label=\Alph*.]
    \item Pablo Picasso’s painting of Vincent van Gogh’s original painting, bearing Picasso’s signature.
    \item Pablo Picasso’s painting of Vincent van Gogh’s original painting, identical in every respect.
    \item Pablo Picasso’s photograph of Vincent van Gogh’s original painting, printed to exactly the same scale.
    \item Pablo Picasso’s miniaturised, but otherwise faithful and accurate painting of Vincent van Gogh’s original painting.
[](https://cracku.in/5-based-on-the-passage-which-one-of-the-following-co-x-cat-2022-slot-1)
\end{enumerate}
\vspace{1em}
\textbf{Solution:}

The passage discusses cultural differences in the concept of a copy and the value placed on originality, particularly in relation to art and religious buildings. In China, copies (fuzhipin) are considered to be of equal value to the original and do not carry negative connotations, while in the Western world, the idea of an unassailable original has historically held more importance. This difference in perspective has led to misunderstandings and tensions between China and Western museums when Chinese museums send copies abroad.   

Based on the above, a Chinese museum would be unlikely to consider Option B [Pablo Picasso's painting of Vincent van Gogh's original painting, identical in every respect] as having less value than the original. This is because the Chinese concept of a copy (fuzhipin) refers to exact reproductions of the original that are considered to be of equal value to the original and do not carry negative connotations. Contrarily, Option A - a painting of Vincent van Gogh's original painting by Pablo Picasso with Picasso's signature - would not be considered a fuzhipin as it is not an exact reproduction of the original and bears the signature of a different artist. Similarly, Options C and D would also not be considered a fuzhipin since they are not an exact reproduction of the original [but merely different versions/formats].
Hence, Option B is the correct choice.

\end{problem}

\newpage
\begin{problem}{}{}
% Subject: varc
\textbf{Instructions}

Instructions   

**The passage below is accompanied by a set of questions. Choose the best answer to each question.**
**Comprehension** :
Stoicism was founded in 300 BC by the Greek philosopher Zeno and survived into the Roman era until about AD 300. According to the Stoics, emotions consist of two movements. The first movement is the immediate feeling and other reactions (e.g., physiological response) that occur when a stimulus or event occurs. For instance, consider what could have happened if an army general accused Marcus Aurelius of treason in front of other officers. The first movement for Marcus may have been (internal) surprise and anger in response to this insult, accompanied perhaps by some involuntary physiological and expressive responses such as face flushing and a movement of the eyebrows. The second movement is what one does next about the emotion. Second movement behaviors occur after thinking and are under one’s control. Examples of second movements for Marcus might have included a plot to seek revenge, actions signifying deference and appeasement, or perhaps proceeding as he would have proceeded whether or not this event occurred: continuing to lead the Romans in a way that Marcus Aurelius believed best benefited them. In the Stoic view, choosing a reasoned, unemotional response as the second movement is the only appropriate response.
The Stoics believed that to live the good life and be a good person, we need to free ourselves of nearly all desires such as too much desire for money, power, or sexual gratification. Prior to second movements, we can consider what is important in life. Money, power, and excessive sexual gratification are not important. Character, rationality, and kindness are important. The Epicureans, first associated with the Greek philosopher Epicurus . . . held a similar view, believing that people should enjoy simple pleasures, such as good conversation, friendship, food, and wine, but not be indulgent in these pursuits and not follow passion for those things that hold no real value like power and money. As Oatley (2004) states, “the Epicureans articulated a view—enjoyment of relationship with friends, of things that are real rather than illusory, simple rather than artificially inflated, possible rather than vanishingly unlikely—that is certainly relevant today” . . . In sum, these ancient Greek and Roman philosophers saw emotions, especially strong ones, as potentially dangerous. They viewed emotions as experiences that needed to be [reined] in and controlled.
As Oatley (2004) points out, the Stoic idea bears some similarity to Buddhism. Buddha, living in India in the 6th century BC, argued for cultivating a certain attitude that decreases the probability of (in Stoic terms) destructive second movements. Through meditation and the right attitude, one allows emotions to happen to oneself (it is impossible to prevent this), but one is advised to observe the emotions without necessarily acting on them; one achieves some distance and decides what has value and what does not have value. Additionally, the  
Stoic idea of developing virtue in oneself, of becoming a good person, which the Stoics believed we could do because we have a touch of the divine, laid the foundation for the three monotheistic religions: Judaism, Christianity, and Islam . . . As with Stoicism, tenets of these religions include controlling our emotions lest we engage in sinful behavior.

\vspace{1em}
\textbf{Question 6}

Which one of the following scenarios is unlikely to follow from the arguments in the passage?

\begin{enumerate}[label=\Alph*.]
    \item A 17th-century British painter would have no problem adding personal touches when restoring an ancient Roman painting.
    \item A 20th-century Japanese Buddhist monk would value a reconstructed shrine as the original.
    \item A 17th-century French artist who adhered to a Christian worldview would need to be completely true to the original intent of a painting when restoring it.
    \item A 21st-century Christian scientist is likely to oppose cloning because of his philosophical orientation.
[](https://cracku.in/6-which-one-of-the-following-scenarios-is-unlikely-t-x-cat-2022-slot-1)
\end{enumerate}
\vspace{1em}
\textbf{Solution:}

{_Back in the 17th century [in the West], excavated artworks from antiquity were treated quite differently from today. They were not restored in a way that was faithful to the original. Instead, there was massive intervention in these works, changing their appearance._.}
Based on the passage, the scenario in Option C [A 17th-century French artist who adhered to a Christian worldview would need to be completely true to the original intent of a painting when restoring it] is unlikely to follow from the arguments in the passage. The passage mentions that in the 17th century, excavated artworks from antiquity were treated differently from how they are today and were not restored in a way that was faithful to the original. Instead, there was "massive intervention" in these works, changing their appearance. This suggests that the idea of an unassailable original may not have held as much importance in the 17th century as it does today. Therefore, it is unlikely that a 17th-century French artist who adhered to a Christian worldview would necessarily need to be completely true to the original intent of a painting when restoring it. Contrarily, we cannot definitively comment on the other scenarios - A, B, and D. 
Hence, Option C is the correct choice.

\end{problem}

\newpage
\begin{problem}{}{}
% Subject: varc
\textbf{Instructions}

Instructions   

**The passage below is accompanied by a set of questions. Choose the best answer to each question.**
**Comprehension** :
Stoicism was founded in 300 BC by the Greek philosopher Zeno and survived into the Roman era until about AD 300. According to the Stoics, emotions consist of two movements. The first movement is the immediate feeling and other reactions (e.g., physiological response) that occur when a stimulus or event occurs. For instance, consider what could have happened if an army general accused Marcus Aurelius of treason in front of other officers. The first movement for Marcus may have been (internal) surprise and anger in response to this insult, accompanied perhaps by some involuntary physiological and expressive responses such as face flushing and a movement of the eyebrows. The second movement is what one does next about the emotion. Second movement behaviors occur after thinking and are under one’s control. Examples of second movements for Marcus might have included a plot to seek revenge, actions signifying deference and appeasement, or perhaps proceeding as he would have proceeded whether or not this event occurred: continuing to lead the Romans in a way that Marcus Aurelius believed best benefited them. In the Stoic view, choosing a reasoned, unemotional response as the second movement is the only appropriate response.
The Stoics believed that to live the good life and be a good person, we need to free ourselves of nearly all desires such as too much desire for money, power, or sexual gratification. Prior to second movements, we can consider what is important in life. Money, power, and excessive sexual gratification are not important. Character, rationality, and kindness are important. The Epicureans, first associated with the Greek philosopher Epicurus . . . held a similar view, believing that people should enjoy simple pleasures, such as good conversation, friendship, food, and wine, but not be indulgent in these pursuits and not follow passion for those things that hold no real value like power and money. As Oatley (2004) states, “the Epicureans articulated a view—enjoyment of relationship with friends, of things that are real rather than illusory, simple rather than artificially inflated, possible rather than vanishingly unlikely—that is certainly relevant today” . . . In sum, these ancient Greek and Roman philosophers saw emotions, especially strong ones, as potentially dangerous. They viewed emotions as experiences that needed to be [reined] in and controlled.
As Oatley (2004) points out, the Stoic idea bears some similarity to Buddhism. Buddha, living in India in the 6th century BC, argued for cultivating a certain attitude that decreases the probability of (in Stoic terms) destructive second movements. Through meditation and the right attitude, one allows emotions to happen to oneself (it is impossible to prevent this), but one is advised to observe the emotions without necessarily acting on them; one achieves some distance and decides what has value and what does not have value. Additionally, the  
Stoic idea of developing virtue in oneself, of becoming a good person, which the Stoics believed we could do because we have a touch of the divine, laid the foundation for the three monotheistic religions: Judaism, Christianity, and Islam . . . As with Stoicism, tenets of these religions include controlling our emotions lest we engage in sinful behavior.

\vspace{1em}
\textbf{Question 7}

Which one of the following statements does not correctly express the similarity between the Ise Grand Shrine and the cathedral of Freiburg Minster?

\begin{enumerate}[label=\Alph*.]
    \item Both were built as places of worship.
    \item Both can be regarded as very old structures.
    \item Both are continually undergoing restoration.
    \item Both will one day be completely rebuilt.
[](https://cracku.in/7-which-one-of-the-following-statements-does-not-cor-x-cat-2022-slot-1)
\end{enumerate}
\vspace{1em}
\textbf{Solution:}

{_The Ise Grand Shrine [in Japan] is 1,300 years old for the millions of Japanese people who go there on pilgrimage every year. But in reality this temple complex is completely rebuilt from scratch every 20 years._ . . .}
While we know that the cathedral of Freiburg Minster is continually undergoing restoration, the same cannot be said about Ise Grand Shrine - we are told that it is rebuilt periodically, but there is no information not substantiate that it is being continually restored [at least in the same sense as that conveyed in the passage]. The idea of restoration is based on how these monuments are being rebuilt. 
Hence, Option C is the correct choice.

\end{problem}

\newpage
\begin{problem}{}{}
% Subject: varc
\textbf{Instructions}

Instructions   

**The passage below is accompanied by a set of questions. Choose the best answer to each question.**
**Comprehension** :
Stoicism was founded in 300 BC by the Greek philosopher Zeno and survived into the Roman era until about AD 300. According to the Stoics, emotions consist of two movements. The first movement is the immediate feeling and other reactions (e.g., physiological response) that occur when a stimulus or event occurs. For instance, consider what could have happened if an army general accused Marcus Aurelius of treason in front of other officers. The first movement for Marcus may have been (internal) surprise and anger in response to this insult, accompanied perhaps by some involuntary physiological and expressive responses such as face flushing and a movement of the eyebrows. The second movement is what one does next about the emotion. Second movement behaviors occur after thinking and are under one’s control. Examples of second movements for Marcus might have included a plot to seek revenge, actions signifying deference and appeasement, or perhaps proceeding as he would have proceeded whether or not this event occurred: continuing to lead the Romans in a way that Marcus Aurelius believed best benefited them. In the Stoic view, choosing a reasoned, unemotional response as the second movement is the only appropriate response.
The Stoics believed that to live the good life and be a good person, we need to free ourselves of nearly all desires such as too much desire for money, power, or sexual gratification. Prior to second movements, we can consider what is important in life. Money, power, and excessive sexual gratification are not important. Character, rationality, and kindness are important. The Epicureans, first associated with the Greek philosopher Epicurus . . . held a similar view, believing that people should enjoy simple pleasures, such as good conversation, friendship, food, and wine, but not be indulgent in these pursuits and not follow passion for those things that hold no real value like power and money. As Oatley (2004) states, “the Epicureans articulated a view—enjoyment of relationship with friends, of things that are real rather than illusory, simple rather than artificially inflated, possible rather than vanishingly unlikely—that is certainly relevant today” . . . In sum, these ancient Greek and Roman philosophers saw emotions, especially strong ones, as potentially dangerous. They viewed emotions as experiences that needed to be [reined] in and controlled.
As Oatley (2004) points out, the Stoic idea bears some similarity to Buddhism. Buddha, living in India in the 6th century BC, argued for cultivating a certain attitude that decreases the probability of (in Stoic terms) destructive second movements. Through meditation and the right attitude, one allows emotions to happen to oneself (it is impossible to prevent this), but one is advised to observe the emotions without necessarily acting on them; one achieves some distance and decides what has value and what does not have value. Additionally, the  
Stoic idea of developing virtue in oneself, of becoming a good person, which the Stoics believed we could do because we have a touch of the divine, laid the foundation for the three monotheistic religions: Judaism, Christianity, and Islam . . . As with Stoicism, tenets of these religions include controlling our emotions lest we engage in sinful behavior.

\vspace{1em}
\textbf{Question 8}

The value that the modern West assigns to “an unassailable original” has resulted in all of the following EXCEPT:

\begin{enumerate}[label=\Alph*.]
    \item it discourages them from simultaneous displays of multiple copies of a painting.
    \item it allows regular employment for certain craftsmen.
    \item it discourages them from making interventions in ancient art.
    \item it discourages them from carrying out human cloning.
[](https://cracku.in/8-the-value-that-the-modern-west-assigns-to-an-unass-x-cat-2022-slot-1)
\end{enumerate}
\vspace{1em}
\textbf{Solution:}

_Option A_ : The value placed on an unassailable original in the Western world may discourage the simultaneous display of multiple copies of a painting, as the original is considered more valuable and authentic. Hence, Option A is valid.
_Option B_ : This is a valid option because the value placed on the original artwork in the Western world may lead to the regular employment of craftsmen who are responsible for preserving and restoring original works of art. This can include tasks such as examining the artwork for damage and replacing eroded or damaged materials [restoration].
_Option C_ : It is true that the focus on the original in the Western world may discourage interventions in ancient art that would alter the appearance of the original. In the past, ancient artworks were frequently altered during restoration, but this practice has become less common in recent times as the value of preserving the original appearance of the artwork has increased. Thus, Option C is plausible._  
_
_Option D_ : The passage discusses how the idea of an original work of art that cannot be altered developed in the Western world and how this intellectual position has led to different attitudes towards cloning between Europe and Asia. However, the passage does not directly mention that the value placed on an unassailable original has discouraged or influenced attitudes towards human cloning in the Western world.
Hence, Option D is the correct answer.

Instructions   

**The passage below is accompanied by a set of questions. Choose the best answer to each question.**
**Comprehension** :
Stories concerning the Undead have always been with us. From out of the primal darkness of Mankind’s earliest years, come whispers of eerie creatures, not quite alive (or alive in a way which we can understand), yet not quite dead either. These may have been ancient and primitive deities who dwelt deep in the surrounding forests and in remote places, or simply those deceased who refused to remain in their tombs and who wandered about the countryside, physically tormenting and frightening those who were still alive. Mostly they were ill-defined—strange sounds in the night beyond the comforting glow of the fire, or a shape, half-glimpsed in the twilight along the edge of an encampment. They were vague and indistinct, but they were always there with the power to terrify and disturb. They had the power to touch the minds of our early ancestors and to fill them with dread. Such fear formed the basis of the earliest tales although the source and exact nature of such terrors still remained very vague.
And as Mankind became more sophisticated, leaving the gloom of their caves and forming themselves into recognizable communities—towns, cities, whole cultures—so the Undead travelled with them, inhabiting their folklore just as they had in former times. Now they began to take on more definite shapes. They became walking cadavers; the physical embodiment of former deities and things which had existed alongside Man since the Creation. Some still remained vague and ill-defined but, as Mankind strove to explain the horror which it felt towards them, such creatures emerged more readily into the light.
In order to confirm their abnormal status, many of the Undead were often accorded attributes, which defied the natural order of things—the power to transform themselves into other shapes, the ability to sustain themselves by drinking human blood, and the ability to influence human minds across a distance. Such powers—described as supernatural—only [lent] an added dimension to the terror that humans felt regarding them.
And it was only natural, too, that the Undead should become connected with the practice of magic. From very early times, Shamans and witchdoctors had claimed at least some power and control over the spirits of departed ancestors, and this has continued down into more “civilized” times. Formerly, the invisible spirits and forces that thronged around men’s earliest encampments, had spoken “through” the tribal Shamans but now, as entities in their own right, they were subject to magical control and could be physically summoned by a competent sorcerer. However, the relationship between the magician and an Undead creature was often a very tenuous and uncertain one. Some sorcerers might have even become Undead entities once they died, but they might also have been susceptible to the powers of other magicians when they did.
From the Middle Ages and into the Age of Enlightenment, theories of the Undead continued to grow and develop. Their names became more familiar—werewolf, vampire, ghoul—each one certain to strike fear into the hearts of ordinary humans.

\end{problem}

\newpage
\begin{problem}{}{}
% Subject: varc
\textbf{Instructions}

Instructions   

**The passage below is accompanied by a set of questions. Choose the best answer to each question.**
**Comprehension** :
Stoicism was founded in 300 BC by the Greek philosopher Zeno and survived into the Roman era until about AD 300. According to the Stoics, emotions consist of two movements. The first movement is the immediate feeling and other reactions (e.g., physiological response) that occur when a stimulus or event occurs. For instance, consider what could have happened if an army general accused Marcus Aurelius of treason in front of other officers. The first movement for Marcus may have been (internal) surprise and anger in response to this insult, accompanied perhaps by some involuntary physiological and expressive responses such as face flushing and a movement of the eyebrows. The second movement is what one does next about the emotion. Second movement behaviors occur after thinking and are under one’s control. Examples of second movements for Marcus might have included a plot to seek revenge, actions signifying deference and appeasement, or perhaps proceeding as he would have proceeded whether or not this event occurred: continuing to lead the Romans in a way that Marcus Aurelius believed best benefited them. In the Stoic view, choosing a reasoned, unemotional response as the second movement is the only appropriate response.
The Stoics believed that to live the good life and be a good person, we need to free ourselves of nearly all desires such as too much desire for money, power, or sexual gratification. Prior to second movements, we can consider what is important in life. Money, power, and excessive sexual gratification are not important. Character, rationality, and kindness are important. The Epicureans, first associated with the Greek philosopher Epicurus . . . held a similar view, believing that people should enjoy simple pleasures, such as good conversation, friendship, food, and wine, but not be indulgent in these pursuits and not follow passion for those things that hold no real value like power and money. As Oatley (2004) states, “the Epicureans articulated a view—enjoyment of relationship with friends, of things that are real rather than illusory, simple rather than artificially inflated, possible rather than vanishingly unlikely—that is certainly relevant today” . . . In sum, these ancient Greek and Roman philosophers saw emotions, especially strong ones, as potentially dangerous. They viewed emotions as experiences that needed to be [reined] in and controlled.
As Oatley (2004) points out, the Stoic idea bears some similarity to Buddhism. Buddha, living in India in the 6th century BC, argued for cultivating a certain attitude that decreases the probability of (in Stoic terms) destructive second movements. Through meditation and the right attitude, one allows emotions to happen to oneself (it is impossible to prevent this), but one is advised to observe the emotions without necessarily acting on them; one achieves some distance and decides what has value and what does not have value. Additionally, the  
Stoic idea of developing virtue in oneself, of becoming a good person, which the Stoics believed we could do because we have a touch of the divine, laid the foundation for the three monotheistic religions: Judaism, Christianity, and Islam . . . As with Stoicism, tenets of these religions include controlling our emotions lest we engage in sinful behavior.

\vspace{1em}
\textbf{Question 9}

“In order to confirm their abnormal status, many of the Undead were often accorded attributes, which defied the natural order of things . . .”
Which one of the following best expresses the claim made in this statement?

\begin{enumerate}[label=\Alph*.]
    \item Human beings conceptualise the Undead as possessing abnormal features.
    \item The Undead are deified in nature’s order by giving them divine attributes.
    \item The natural attributes of the Undead are rendered abnormal by changing their status.
    \item According the Undead an abnormal status is to reject the natural order of things.
[](https://cracku.in/9-in-order-to-confirm-their-abnormal-status-many-of--x-cat-2022-slot-1)
\end{enumerate}
\vspace{1em}
\textbf{Solution:}

Option A best expresses the claim made in the statement - "_In order to confirm their abnormal status, many of the Undead were often accorded attributes, which defied the natural order of things_. . .” 
The passage states that in order to confirm their abnormal status, the Undead were often given attributes that defied the natural order of things. This suggests that humans conceptualize the Undead as possessing abnormal features to confirm their abnormal status. This differs from Options B and C, which suggest that the Undead are deified or that their natural attributes are rendered abnormal by changing their status. Option D is also inaccurate, as the passage does not mention that giving the Undead an abnormal status is a rejection of the natural order of things.
Hence, Option A is the correct answer.

\end{problem}

\newpage
\begin{problem}{}{}
% Subject: varc
\textbf{Instructions}

Instructions   

**The passage below is accompanied by a set of questions. Choose the best answer to each question.**
**Comprehension** :
Stoicism was founded in 300 BC by the Greek philosopher Zeno and survived into the Roman era until about AD 300. According to the Stoics, emotions consist of two movements. The first movement is the immediate feeling and other reactions (e.g., physiological response) that occur when a stimulus or event occurs. For instance, consider what could have happened if an army general accused Marcus Aurelius of treason in front of other officers. The first movement for Marcus may have been (internal) surprise and anger in response to this insult, accompanied perhaps by some involuntary physiological and expressive responses such as face flushing and a movement of the eyebrows. The second movement is what one does next about the emotion. Second movement behaviors occur after thinking and are under one’s control. Examples of second movements for Marcus might have included a plot to seek revenge, actions signifying deference and appeasement, or perhaps proceeding as he would have proceeded whether or not this event occurred: continuing to lead the Romans in a way that Marcus Aurelius believed best benefited them. In the Stoic view, choosing a reasoned, unemotional response as the second movement is the only appropriate response.
The Stoics believed that to live the good life and be a good person, we need to free ourselves of nearly all desires such as too much desire for money, power, or sexual gratification. Prior to second movements, we can consider what is important in life. Money, power, and excessive sexual gratification are not important. Character, rationality, and kindness are important. The Epicureans, first associated with the Greek philosopher Epicurus . . . held a similar view, believing that people should enjoy simple pleasures, such as good conversation, friendship, food, and wine, but not be indulgent in these pursuits and not follow passion for those things that hold no real value like power and money. As Oatley (2004) states, “the Epicureans articulated a view—enjoyment of relationship with friends, of things that are real rather than illusory, simple rather than artificially inflated, possible rather than vanishingly unlikely—that is certainly relevant today” . . . In sum, these ancient Greek and Roman philosophers saw emotions, especially strong ones, as potentially dangerous. They viewed emotions as experiences that needed to be [reined] in and controlled.
As Oatley (2004) points out, the Stoic idea bears some similarity to Buddhism. Buddha, living in India in the 6th century BC, argued for cultivating a certain attitude that decreases the probability of (in Stoic terms) destructive second movements. Through meditation and the right attitude, one allows emotions to happen to oneself (it is impossible to prevent this), but one is advised to observe the emotions without necessarily acting on them; one achieves some distance and decides what has value and what does not have value. Additionally, the  
Stoic idea of developing virtue in oneself, of becoming a good person, which the Stoics believed we could do because we have a touch of the divine, laid the foundation for the three monotheistic religions: Judaism, Christianity, and Islam . . . As with Stoicism, tenets of these religions include controlling our emotions lest we engage in sinful behavior.

\vspace{1em}
\textbf{Question 10}

Which one of the following observations is a valid conclusion to draw from the statement, “From out of the primal darkness of Mankind’s earliest years, come whispers of eerie creatures, not quite alive (or alive in a way which we can understand), yet not quite dead either.”?

\begin{enumerate}[label=\Alph*.]
    \item Mankind’s early years were marked by a belief in the existence of eerie creatures that were neither quite alive nor dead.
    \item Long ago, eerie creatures used to whisper in the primal darkness that they were not quite dead.
    \item Mankind’s primal years were marked by creatures alive with eerie whispers, but seen only in the darkness.
    \item We can understand the lives of the eerie creatures in Mankind’s early years through their whispers in the darkness.
[](https://cracku.in/10-which-one-of-the-following-observations-is-a-valid-x-cat-2022-slot-1)
\end{enumerate}
\vspace{1em}
\textbf{Solution:}

Option A is a valid conclusion to draw from the lines - "_From out of the primal darkness of Mankind's earliest years, come whispers of eerie creatures, not quite alive (or alive in a way which we can understand), yet not quite dead either._ "
The statement mentions that in Mankind's earliest years, there were whispers of eerie creatures that were not quite alive or dead. This suggests that in these early years, there was a belief in the existence of such creatures. Option A accurately captures this idea. Option B is not a valid conclusion as the given statement does not mention anything about eerie creatures whispering about their own death. Option C is also incorrect since the passage does not mention that the creatures were only seen in the darkness. Similarly, Option D is not a valid conclusion as the given statement offers no information on how we can understand the lives of the eerie creatures. 
Hence, Option A is the correct choice.

\end{problem}

\newpage
\begin{problem}{}{}
% Subject: varc
\textbf{Instructions}

Instructions   

**The passage below is accompanied by a set of questions. Choose the best answer to each question.**
**Comprehension** :
Stoicism was founded in 300 BC by the Greek philosopher Zeno and survived into the Roman era until about AD 300. According to the Stoics, emotions consist of two movements. The first movement is the immediate feeling and other reactions (e.g., physiological response) that occur when a stimulus or event occurs. For instance, consider what could have happened if an army general accused Marcus Aurelius of treason in front of other officers. The first movement for Marcus may have been (internal) surprise and anger in response to this insult, accompanied perhaps by some involuntary physiological and expressive responses such as face flushing and a movement of the eyebrows. The second movement is what one does next about the emotion. Second movement behaviors occur after thinking and are under one’s control. Examples of second movements for Marcus might have included a plot to seek revenge, actions signifying deference and appeasement, or perhaps proceeding as he would have proceeded whether or not this event occurred: continuing to lead the Romans in a way that Marcus Aurelius believed best benefited them. In the Stoic view, choosing a reasoned, unemotional response as the second movement is the only appropriate response.
The Stoics believed that to live the good life and be a good person, we need to free ourselves of nearly all desires such as too much desire for money, power, or sexual gratification. Prior to second movements, we can consider what is important in life. Money, power, and excessive sexual gratification are not important. Character, rationality, and kindness are important. The Epicureans, first associated with the Greek philosopher Epicurus . . . held a similar view, believing that people should enjoy simple pleasures, such as good conversation, friendship, food, and wine, but not be indulgent in these pursuits and not follow passion for those things that hold no real value like power and money. As Oatley (2004) states, “the Epicureans articulated a view—enjoyment of relationship with friends, of things that are real rather than illusory, simple rather than artificially inflated, possible rather than vanishingly unlikely—that is certainly relevant today” . . . In sum, these ancient Greek and Roman philosophers saw emotions, especially strong ones, as potentially dangerous. They viewed emotions as experiences that needed to be [reined] in and controlled.
As Oatley (2004) points out, the Stoic idea bears some similarity to Buddhism. Buddha, living in India in the 6th century BC, argued for cultivating a certain attitude that decreases the probability of (in Stoic terms) destructive second movements. Through meditation and the right attitude, one allows emotions to happen to oneself (it is impossible to prevent this), but one is advised to observe the emotions without necessarily acting on them; one achieves some distance and decides what has value and what does not have value. Additionally, the  
Stoic idea of developing virtue in oneself, of becoming a good person, which the Stoics believed we could do because we have a touch of the divine, laid the foundation for the three monotheistic religions: Judaism, Christianity, and Islam . . . As with Stoicism, tenets of these religions include controlling our emotions lest we engage in sinful behavior.

\vspace{1em}
\textbf{Question 11}

Which one of the following statements best describes what the passage is about?

\begin{enumerate}[label=\Alph*.]
    \item The writer discusses the transition from primitive thinking to the Age of Enlightenment.
    \item The passage discusses the evolution of theories of the Undead from primitive thinking to the Age of Enlightenment.
    \item The passage describes the failure of human beings to fully comprehend their environment.
    \item The writer describes the ways in which the Undead come to be associated with Shamans and the practice of magic.
[](https://cracku.in/11-which-one-of-the-following-statements-best-describ-x-cat-2022-slot-1)
\end{enumerate}
\vspace{1em}
\textbf{Solution:}

The passage underlines that the concept of the Undead, or creatures that are not quite alive or dead, has always been a part of human folklore. In ancient times, the Undead were ill-defined and vague, but as human societies became more sophisticated, the Undead took on more definite shapes and were often associated with supernatural powers, such as the ability to transform, drink blood, and influence human minds. The Undead have also been connected to the practice of magic, and in more recent times, specific names such as werewolf, vampire, and ghoul have become associated with the concept of the Undead. These names are often used to strike fear into the hearts of ordinary humans. The passage suggests that the Undead have evolved and developed over time, and as human societies have advanced, the Undead have become more defined and have gained more specific attributes. Overall, the passage discusses the long-standing presence of the Undead in human folklore and the evolution of their portrayal in various cultures. Option B aptly captures the above idea. 
Option A is incorrect since the passage does not mention the transition from primitive thinking to the Age of Enlightenment. Option C is also inaccurate as the author does not emphasize the failure of human beings "to fully comprehend their environment" [not the focus]. Option D is not a complete description of the passage since it only mentions one aspect of the passage rather than the overall theme of the evolution of the concept of the Undead.
Hence, Option B is the correct choice.

\end{problem}

\newpage
\begin{problem}{}{}
% Subject: varc
\textbf{Instructions}

Instructions   

**The passage below is accompanied by a set of questions. Choose the best answer to each question.**
**Comprehension** :
Stoicism was founded in 300 BC by the Greek philosopher Zeno and survived into the Roman era until about AD 300. According to the Stoics, emotions consist of two movements. The first movement is the immediate feeling and other reactions (e.g., physiological response) that occur when a stimulus or event occurs. For instance, consider what could have happened if an army general accused Marcus Aurelius of treason in front of other officers. The first movement for Marcus may have been (internal) surprise and anger in response to this insult, accompanied perhaps by some involuntary physiological and expressive responses such as face flushing and a movement of the eyebrows. The second movement is what one does next about the emotion. Second movement behaviors occur after thinking and are under one’s control. Examples of second movements for Marcus might have included a plot to seek revenge, actions signifying deference and appeasement, or perhaps proceeding as he would have proceeded whether or not this event occurred: continuing to lead the Romans in a way that Marcus Aurelius believed best benefited them. In the Stoic view, choosing a reasoned, unemotional response as the second movement is the only appropriate response.
The Stoics believed that to live the good life and be a good person, we need to free ourselves of nearly all desires such as too much desire for money, power, or sexual gratification. Prior to second movements, we can consider what is important in life. Money, power, and excessive sexual gratification are not important. Character, rationality, and kindness are important. The Epicureans, first associated with the Greek philosopher Epicurus . . . held a similar view, believing that people should enjoy simple pleasures, such as good conversation, friendship, food, and wine, but not be indulgent in these pursuits and not follow passion for those things that hold no real value like power and money. As Oatley (2004) states, “the Epicureans articulated a view—enjoyment of relationship with friends, of things that are real rather than illusory, simple rather than artificially inflated, possible rather than vanishingly unlikely—that is certainly relevant today” . . . In sum, these ancient Greek and Roman philosophers saw emotions, especially strong ones, as potentially dangerous. They viewed emotions as experiences that needed to be [reined] in and controlled.
As Oatley (2004) points out, the Stoic idea bears some similarity to Buddhism. Buddha, living in India in the 6th century BC, argued for cultivating a certain attitude that decreases the probability of (in Stoic terms) destructive second movements. Through meditation and the right attitude, one allows emotions to happen to oneself (it is impossible to prevent this), but one is advised to observe the emotions without necessarily acting on them; one achieves some distance and decides what has value and what does not have value. Additionally, the  
Stoic idea of developing virtue in oneself, of becoming a good person, which the Stoics believed we could do because we have a touch of the divine, laid the foundation for the three monotheistic religions: Judaism, Christianity, and Islam . . . As with Stoicism, tenets of these religions include controlling our emotions lest we engage in sinful behavior.

\vspace{1em}
\textbf{Question 12}

All of the following statements, if false, could be seen as being in accordance with the passage, EXCEPT:

\begin{enumerate}[label=\Alph*.]
    \item the Undead remained vague and ill-defined, even as Mankind strove to understand the horror they inspired.
    \item the transition from the Middle Ages to the Age of Enlightenment saw new theories of the Undead.
    \item the growing sophistication of Mankind meant that humans stopped believing in the Undead.
    \item the relationship between Shamans and the Undead was believed to be a strong and stable one.
[](https://cracku.in/12-all-of-the-following-statements-if-false-could-be--x-cat-2022-slot-1)
\end{enumerate}
\vspace{1em}
\textbf{Solution:}

The multiple negations indicate that if any of the given statements are false, they could be seen as being consistent with the information provided in the passage [i.e. they do not contradict the information provided in the passage]. Since the question involves "except," we need to find a valid statement based on the information provided [since if this statement is correct, it would not be consistent with the information in the passage]. 
We notice that Option B is consistent with the information provided in the passage: The passage states that as human societies became more sophisticated, the Undead took on more definite shapes and became more defined. It also mentions that from the Middle Ages and into the Age of Enlightenment, theories of the Undead continued to grow and develop. This suggests that the transition from the Middle Ages to the Age of Enlightenment saw new theories of the Undead. There will be inconsistencies if the statement in B is refuted or incorrect.
If the remaining options are false, they will support the discussion in the passage: Option A is not consistent with the passage as it states that the Undead remained vague and ill-defined, even as human societies strove to understand the horror they inspired, while the passage actually states that the Undead became more defined as human societies became more sophisticated. Option C is also incorrect because the passage does not mention that the growing sophistication of Mankind caused humans to stop believing in the Undead. Similarly, Option D is also inconsistent in its current form since the passage does not mention anything about the strength or stability of the relationship between Shamans and the Undead.
Hence, Option B is the correct choice.

Instructions   

**The passage below is accompanied by a set of questions. Choose the best answer to each question.**
**Comprehension** :
Critical theory of technology is a political theory of modernity with a normative dimension. It belongs to a tradition extending from Marx to Foucault and Habermas according to which advances in the formal claims of human rights take center stage while in the background centralization of ever more powerful public institutions and private organizations imposes an authoritarian social order.
Marx attributed this trajectory to the capitalist rationalization of production. Today it marks many institutions besides the factory and every modern political system, including so-called socialist systems. This trajectory arose from the problems of command over a disempowered and deskilled labor force; but everywhere [that] masses are organized - whether it be Foucault’s prisons or Habermas’s public sphere - the same pattern prevails. Technological design and development is shaped by this pattern as the material base of a distinctive social order. Marcuse would later point to a “project” as the basis of what he called rather confusingly “technological rationality.” Releasing technology from this project is a democratic political task.
In accordance with this general line of thought, critical theory of technology regards technologies as an environment rather than as a collection of tools. We live today with and even within technologies that determine our way of life. Along with the constant pressures to build centers of power, many other social values and meanings are inscribed in technological design. A hermeneutics of technology must make explicit the meanings implicit in the devices we use and the rituals they script. Social histories of technologies such as the bicycle, artificial lighting or firearms have made important contributions to this type of analysis. Critical theory of technology attempts to build a methodological approach on the lessons of these histories.
As an environment, technologies shape their inhabitants. In this respect, they are comparable to laws and customs. Each of these institutions can be said to represent those who live under their sway through privileging certain dimensions of their human nature. Laws of property represent the interest in ownership and control. Customs such as parental authority represent the interest of childhood in safety and growth. Similarly, the automobile represents its users in so far as they are interested in mobility. Interests such as these constitute the version of human nature sanctioned by society.
This notion of representation does not imply an eternal human nature. The concept of nature as non-identity in the Frankfurt School suggests an alternative. On these terms, nature is what lies at the limit of history, at the point at which society loses the capacity to imprint its meanings on things and control them effectively. The reference here is, of course, not to the nature of natural science, but to the lived nature in which we find ourselves and which we are. This nature reveals itself as that which cannot be totally encompassed by the machinery of society. For the Frankfurt School, human nature, in all its transcending force, emerges out of a historical context as that context is [depicted] in illicit joys, struggles and pathologies. We can perhaps admit a less romantic . . . conception in which those dimensions of human nature recognized by society are also granted theoretical legitimacy.

\end{problem}

\newpage
\begin{problem}{}{}
% Subject: varc
\textbf{Instructions}

Instructions   

**The passage below is accompanied by a set of questions. Choose the best answer to each question.**
**Comprehension** :
Stoicism was founded in 300 BC by the Greek philosopher Zeno and survived into the Roman era until about AD 300. According to the Stoics, emotions consist of two movements. The first movement is the immediate feeling and other reactions (e.g., physiological response) that occur when a stimulus or event occurs. For instance, consider what could have happened if an army general accused Marcus Aurelius of treason in front of other officers. The first movement for Marcus may have been (internal) surprise and anger in response to this insult, accompanied perhaps by some involuntary physiological and expressive responses such as face flushing and a movement of the eyebrows. The second movement is what one does next about the emotion. Second movement behaviors occur after thinking and are under one’s control. Examples of second movements for Marcus might have included a plot to seek revenge, actions signifying deference and appeasement, or perhaps proceeding as he would have proceeded whether or not this event occurred: continuing to lead the Romans in a way that Marcus Aurelius believed best benefited them. In the Stoic view, choosing a reasoned, unemotional response as the second movement is the only appropriate response.
The Stoics believed that to live the good life and be a good person, we need to free ourselves of nearly all desires such as too much desire for money, power, or sexual gratification. Prior to second movements, we can consider what is important in life. Money, power, and excessive sexual gratification are not important. Character, rationality, and kindness are important. The Epicureans, first associated with the Greek philosopher Epicurus . . . held a similar view, believing that people should enjoy simple pleasures, such as good conversation, friendship, food, and wine, but not be indulgent in these pursuits and not follow passion for those things that hold no real value like power and money. As Oatley (2004) states, “the Epicureans articulated a view—enjoyment of relationship with friends, of things that are real rather than illusory, simple rather than artificially inflated, possible rather than vanishingly unlikely—that is certainly relevant today” . . . In sum, these ancient Greek and Roman philosophers saw emotions, especially strong ones, as potentially dangerous. They viewed emotions as experiences that needed to be [reined] in and controlled.
As Oatley (2004) points out, the Stoic idea bears some similarity to Buddhism. Buddha, living in India in the 6th century BC, argued for cultivating a certain attitude that decreases the probability of (in Stoic terms) destructive second movements. Through meditation and the right attitude, one allows emotions to happen to oneself (it is impossible to prevent this), but one is advised to observe the emotions without necessarily acting on them; one achieves some distance and decides what has value and what does not have value. Additionally, the  
Stoic idea of developing virtue in oneself, of becoming a good person, which the Stoics believed we could do because we have a touch of the divine, laid the foundation for the three monotheistic religions: Judaism, Christianity, and Islam . . . As with Stoicism, tenets of these religions include controlling our emotions lest we engage in sinful behavior.

\vspace{1em}
\textbf{Question 13}

Which one of the following statements best reflects the main argument of the fourth paragraph of the passage?

\begin{enumerate}[label=\Alph*.]
    \item Technology, laws, and customs are comparable, but dissimilar phenomena.
    \item Technological environments privilege certain dimensions of human nature as effectively as laws and customs.
    \item Automobiles represent the interest in mobility present in human nature.
    \item Technology, laws, and customs are not unlike each other if considered as institutions.
[](https://cracku.in/13-which-one-of-the-following-statements-best-reflect-x-cat-2022-slot-1)
\end{enumerate}
\vspace{1em}
\textbf{Solution:}

The fourth paragraph discusses the ways in which technology shapes society and the values and meanings that are inscribed in technological design. It suggests that technology, like laws and customs, represents the interests of those who use it and shapes the version of human nature that is sanctioned by society. The paragraph compares technology to laws and customs, stating that they are similar in the sense that they are institutions that shape the way people live. Therefore, the correct statement that captures the crux of the fourth paragraph is _Option D_ [Technology, laws, and customs are not unlike each other if considered as institutions]._  
_
_Option A_ : While it is true that technology, laws, and customs are comparable phenomena in this sense, the statement does not adequately capture the main point being made in the paragraph.
_Option B_ : is incorrect because it misinterprets the main point of the fourth paragraph. While the paragraph does suggest that technology shapes the version of human nature that is sanctioned by society, it does not directly compare the effectiveness of technological environments and laws and customs in privileging certain dimensions of human nature.
_Option C:_ is only partially true. While the fourth paragraph does mention the idea that technologies represent the interests of those who use them, it does not specifically state that automobiles represent the interest in mobility present in human nature. The paragraph mentions the automobile as an example of a technology that represents its users, but does not explicitly link it to the concept of mobility.
Hence, Option D is the correct choice.

\end{problem}

\newpage
\begin{problem}{}{}
% Subject: varc
\textbf{Instructions}

Instructions   

**The passage below is accompanied by a set of questions. Choose the best answer to each question.**
**Comprehension** :
Stoicism was founded in 300 BC by the Greek philosopher Zeno and survived into the Roman era until about AD 300. According to the Stoics, emotions consist of two movements. The first movement is the immediate feeling and other reactions (e.g., physiological response) that occur when a stimulus or event occurs. For instance, consider what could have happened if an army general accused Marcus Aurelius of treason in front of other officers. The first movement for Marcus may have been (internal) surprise and anger in response to this insult, accompanied perhaps by some involuntary physiological and expressive responses such as face flushing and a movement of the eyebrows. The second movement is what one does next about the emotion. Second movement behaviors occur after thinking and are under one’s control. Examples of second movements for Marcus might have included a plot to seek revenge, actions signifying deference and appeasement, or perhaps proceeding as he would have proceeded whether or not this event occurred: continuing to lead the Romans in a way that Marcus Aurelius believed best benefited them. In the Stoic view, choosing a reasoned, unemotional response as the second movement is the only appropriate response.
The Stoics believed that to live the good life and be a good person, we need to free ourselves of nearly all desires such as too much desire for money, power, or sexual gratification. Prior to second movements, we can consider what is important in life. Money, power, and excessive sexual gratification are not important. Character, rationality, and kindness are important. The Epicureans, first associated with the Greek philosopher Epicurus . . . held a similar view, believing that people should enjoy simple pleasures, such as good conversation, friendship, food, and wine, but not be indulgent in these pursuits and not follow passion for those things that hold no real value like power and money. As Oatley (2004) states, “the Epicureans articulated a view—enjoyment of relationship with friends, of things that are real rather than illusory, simple rather than artificially inflated, possible rather than vanishingly unlikely—that is certainly relevant today” . . . In sum, these ancient Greek and Roman philosophers saw emotions, especially strong ones, as potentially dangerous. They viewed emotions as experiences that needed to be [reined] in and controlled.
As Oatley (2004) points out, the Stoic idea bears some similarity to Buddhism. Buddha, living in India in the 6th century BC, argued for cultivating a certain attitude that decreases the probability of (in Stoic terms) destructive second movements. Through meditation and the right attitude, one allows emotions to happen to oneself (it is impossible to prevent this), but one is advised to observe the emotions without necessarily acting on them; one achieves some distance and decides what has value and what does not have value. Additionally, the  
Stoic idea of developing virtue in oneself, of becoming a good person, which the Stoics believed we could do because we have a touch of the divine, laid the foundation for the three monotheistic religions: Judaism, Christianity, and Islam . . . As with Stoicism, tenets of these religions include controlling our emotions lest we engage in sinful behavior.

\vspace{1em}
\textbf{Question 14}

Which one of the following statements could be inferred as supporting the arguments of the passage?

\begin{enumerate}[label=\Alph*.]
    \item It is not human nature, but human culture that is represented by institutions such as law and custom.
    \item Technologies form the environmental context and shape the contours of human society.
    \item Nature decides the point at which society loses its capacity to control history.
    \item The romantic conception of nature referred to by the passage is the one that requires theoretical legitimacy.
[](https://cracku.in/14-which-one-of-the-following-statements-could-be-inf-x-cat-2022-slot-1)
\end{enumerate}
\vspace{1em}
\textbf{Solution:}

The passage specifically states that "critical theory of technology regards technologies as an environment rather than as a collection of tools," and that "technologies shape their inhabitants" in a way that is similar to laws and customs, which represent certain interests and values of those who live under their sway. There is no mention of the role of nature in shaping society or determining the limits of society's control over history, so options C and D cannot be supporting points. The statement in A wouldn't support the discussion in the passage because the author suggests that institutions such as laws and customs represent certain dimensions of human nature rather than human culture.
Hence, Option B is the correct choice.

\end{problem}

\newpage
\begin{problem}{}{}
% Subject: varc
\textbf{Instructions}

Instructions   

**The passage below is accompanied by a set of questions. Choose the best answer to each question.**
**Comprehension** :
Stoicism was founded in 300 BC by the Greek philosopher Zeno and survived into the Roman era until about AD 300. According to the Stoics, emotions consist of two movements. The first movement is the immediate feeling and other reactions (e.g., physiological response) that occur when a stimulus or event occurs. For instance, consider what could have happened if an army general accused Marcus Aurelius of treason in front of other officers. The first movement for Marcus may have been (internal) surprise and anger in response to this insult, accompanied perhaps by some involuntary physiological and expressive responses such as face flushing and a movement of the eyebrows. The second movement is what one does next about the emotion. Second movement behaviors occur after thinking and are under one’s control. Examples of second movements for Marcus might have included a plot to seek revenge, actions signifying deference and appeasement, or perhaps proceeding as he would have proceeded whether or not this event occurred: continuing to lead the Romans in a way that Marcus Aurelius believed best benefited them. In the Stoic view, choosing a reasoned, unemotional response as the second movement is the only appropriate response.
The Stoics believed that to live the good life and be a good person, we need to free ourselves of nearly all desires such as too much desire for money, power, or sexual gratification. Prior to second movements, we can consider what is important in life. Money, power, and excessive sexual gratification are not important. Character, rationality, and kindness are important. The Epicureans, first associated with the Greek philosopher Epicurus . . . held a similar view, believing that people should enjoy simple pleasures, such as good conversation, friendship, food, and wine, but not be indulgent in these pursuits and not follow passion for those things that hold no real value like power and money. As Oatley (2004) states, “the Epicureans articulated a view—enjoyment of relationship with friends, of things that are real rather than illusory, simple rather than artificially inflated, possible rather than vanishingly unlikely—that is certainly relevant today” . . . In sum, these ancient Greek and Roman philosophers saw emotions, especially strong ones, as potentially dangerous. They viewed emotions as experiences that needed to be [reined] in and controlled.
As Oatley (2004) points out, the Stoic idea bears some similarity to Buddhism. Buddha, living in India in the 6th century BC, argued for cultivating a certain attitude that decreases the probability of (in Stoic terms) destructive second movements. Through meditation and the right attitude, one allows emotions to happen to oneself (it is impossible to prevent this), but one is advised to observe the emotions without necessarily acting on them; one achieves some distance and decides what has value and what does not have value. Additionally, the  
Stoic idea of developing virtue in oneself, of becoming a good person, which the Stoics believed we could do because we have a touch of the divine, laid the foundation for the three monotheistic religions: Judaism, Christianity, and Islam . . . As with Stoicism, tenets of these religions include controlling our emotions lest we engage in sinful behavior.

\vspace{1em}
\textbf{Question 15}

Which one of the following statements contradicts the arguments of the passage?

\begin{enumerate}[label=\Alph*.]
    \item The problems of command over a disempowered and deskilled labour force gave rise to similar patterns of the capitalist rationalisation of production wherever masses were organised.
    \item Marx’s understanding of the capitalist rationalisation of production and Marcuse’s understanding of a “project” of “technological rationality” share theoretical inclinations.
    \item Masses are organised in patterns set by Foucault’s prisons and Habermas’ public sphere.
    \item Paradoxically, the capitalist rationalisation of production is a mark of so-called socialist systems as well.
[](https://cracku.in/15-which-one-of-the-following-statements-contradicts--x-cat-2022-slot-1)
\end{enumerate}
\vspace{1em}
\textbf{Solution:}

_Option A_ : This statement is consistent with the arguments of the passage, which claim that the pattern of the capitalist rationalization of production arises from the problems of command over a disempowered and deskilled labor force and is present in many different contexts, including the factory and socialist systems.  

_Option B_ : This statement is consistent with the passage, which claims that Marx and Marcuse both contribute to the tradition of critical theory of technology, which seeks to understand the ways in which technological systems are shaped by and contribute to the reproduction of social and political hierarchies.
_Option C_ : This statement contradicts an element discussed in the passage. The passage states that the pattern of the capitalist rationalization of production, which is marked by the centralization of power in institutions and organizations and the deskilling of the labor force, arises in many different contexts, including the factory, prisons, and the public sphere. It does not claim that the patterns in these different contexts are set by Foucault's prisons and Habermas' public sphere.
_Option D_ : This statement is consistent with the passage, which claims that the pattern of the capitalist rationalization of production is present in many different contexts, including socialist systems.
Hence, Option C is the correct choice.

\end{problem}

\newpage
\begin{problem}{}{}
% Subject: varc
\textbf{Instructions}

Instructions   

**The passage below is accompanied by a set of questions. Choose the best answer to each question.**
**Comprehension** :
Stoicism was founded in 300 BC by the Greek philosopher Zeno and survived into the Roman era until about AD 300. According to the Stoics, emotions consist of two movements. The first movement is the immediate feeling and other reactions (e.g., physiological response) that occur when a stimulus or event occurs. For instance, consider what could have happened if an army general accused Marcus Aurelius of treason in front of other officers. The first movement for Marcus may have been (internal) surprise and anger in response to this insult, accompanied perhaps by some involuntary physiological and expressive responses such as face flushing and a movement of the eyebrows. The second movement is what one does next about the emotion. Second movement behaviors occur after thinking and are under one’s control. Examples of second movements for Marcus might have included a plot to seek revenge, actions signifying deference and appeasement, or perhaps proceeding as he would have proceeded whether or not this event occurred: continuing to lead the Romans in a way that Marcus Aurelius believed best benefited them. In the Stoic view, choosing a reasoned, unemotional response as the second movement is the only appropriate response.
The Stoics believed that to live the good life and be a good person, we need to free ourselves of nearly all desires such as too much desire for money, power, or sexual gratification. Prior to second movements, we can consider what is important in life. Money, power, and excessive sexual gratification are not important. Character, rationality, and kindness are important. The Epicureans, first associated with the Greek philosopher Epicurus . . . held a similar view, believing that people should enjoy simple pleasures, such as good conversation, friendship, food, and wine, but not be indulgent in these pursuits and not follow passion for those things that hold no real value like power and money. As Oatley (2004) states, “the Epicureans articulated a view—enjoyment of relationship with friends, of things that are real rather than illusory, simple rather than artificially inflated, possible rather than vanishingly unlikely—that is certainly relevant today” . . . In sum, these ancient Greek and Roman philosophers saw emotions, especially strong ones, as potentially dangerous. They viewed emotions as experiences that needed to be [reined] in and controlled.
As Oatley (2004) points out, the Stoic idea bears some similarity to Buddhism. Buddha, living in India in the 6th century BC, argued for cultivating a certain attitude that decreases the probability of (in Stoic terms) destructive second movements. Through meditation and the right attitude, one allows emotions to happen to oneself (it is impossible to prevent this), but one is advised to observe the emotions without necessarily acting on them; one achieves some distance and decides what has value and what does not have value. Additionally, the  
Stoic idea of developing virtue in oneself, of becoming a good person, which the Stoics believed we could do because we have a touch of the divine, laid the foundation for the three monotheistic religions: Judaism, Christianity, and Islam . . . As with Stoicism, tenets of these religions include controlling our emotions lest we engage in sinful behavior.

\vspace{1em}
\textbf{Question 16}

All of the following claims can be inferred from the passage, EXCEPT:

\begin{enumerate}[label=\Alph*.]
    \item the significance of parental authority to children’s safety does not therefore imply that parental authority is a permanent aspect of human nature.
    \item the critical theory of technology argues that, as issues of human rights become more prominent, we lose sight of the ways in which the social order becomes more authoritarian.
    \item analyses of technologies must engage with their social histories to be able to reveal their implicit and explicit meanings for us.
    \item technologies seek to privilege certain dimensions of human nature at a high cost to lived nature.
[](https://cracku.in/16-all-of-the-following-claims-can-be-inferred-from-t-x-cat-2022-slot-1)
\end{enumerate}
\vspace{1em}
\textbf{Solution:}

_Option A_ : The passage states that "laws of property represent the interest in ownership and control. Customs such as parental authority represent the interest of childhood in safety and growth." It then goes on to say that "interests such as these constitute the version of human nature sanctioned by society." This suggests that the concept of human nature is not fixed but rather emerges out of historical context and is shaped by society. Therefore, it can be inferred that the significance of parental authority to children's safety does not imply that parental authority is a permanent aspect of human nature.  

_Option B_ : This claim can be inferred - the author states that "critical theory of technology is a political theory of modernity with a normative dimension" and that it belongs to a tradition "according to which advances in the formal claims of human rights take centre stage while in the background centralization of ever more powerful public institutions and private organizations imposes an authoritarian social order." This suggests that the critical theory of technology argues that as issues of human rights become more prominent, the social order becomes more authoritarian.
_Option C_ : We are told that "a hermeneutics of technology must make explicit the meanings implicit in the devices we use and the rituals they script" and that "social histories of technologies such as the bicycle, artificial lighting or firearms have made important contributions to this type of analysis." This suggests that engaging with the social histories of technologies is necessary to understand their implicit and explicit meanings for us.
_Option D_ : This cannot be inferred from the passage because the passage does not mention any costs or negative consequences of technologies privileging certain dimensions of human nature - it only discusses the idea that technologies represent the interests of their users and shape their behaviour and values; however, no claims are made about the potential negative impacts of this process.
Hence, Option D is the correct choice.

Instructions   

For the following questions answer them individually

\end{problem}

\newpage
\begin{problem}{}{}
% Subject: varc
\textbf{Instructions}

Instructions   

**The passage below is accompanied by a set of questions. Choose the best answer to each question.**
**Comprehension** :
Stoicism was founded in 300 BC by the Greek philosopher Zeno and survived into the Roman era until about AD 300. According to the Stoics, emotions consist of two movements. The first movement is the immediate feeling and other reactions (e.g., physiological response) that occur when a stimulus or event occurs. For instance, consider what could have happened if an army general accused Marcus Aurelius of treason in front of other officers. The first movement for Marcus may have been (internal) surprise and anger in response to this insult, accompanied perhaps by some involuntary physiological and expressive responses such as face flushing and a movement of the eyebrows. The second movement is what one does next about the emotion. Second movement behaviors occur after thinking and are under one’s control. Examples of second movements for Marcus might have included a plot to seek revenge, actions signifying deference and appeasement, or perhaps proceeding as he would have proceeded whether or not this event occurred: continuing to lead the Romans in a way that Marcus Aurelius believed best benefited them. In the Stoic view, choosing a reasoned, unemotional response as the second movement is the only appropriate response.
The Stoics believed that to live the good life and be a good person, we need to free ourselves of nearly all desires such as too much desire for money, power, or sexual gratification. Prior to second movements, we can consider what is important in life. Money, power, and excessive sexual gratification are not important. Character, rationality, and kindness are important. The Epicureans, first associated with the Greek philosopher Epicurus . . . held a similar view, believing that people should enjoy simple pleasures, such as good conversation, friendship, food, and wine, but not be indulgent in these pursuits and not follow passion for those things that hold no real value like power and money. As Oatley (2004) states, “the Epicureans articulated a view—enjoyment of relationship with friends, of things that are real rather than illusory, simple rather than artificially inflated, possible rather than vanishingly unlikely—that is certainly relevant today” . . . In sum, these ancient Greek and Roman philosophers saw emotions, especially strong ones, as potentially dangerous. They viewed emotions as experiences that needed to be [reined] in and controlled.
As Oatley (2004) points out, the Stoic idea bears some similarity to Buddhism. Buddha, living in India in the 6th century BC, argued for cultivating a certain attitude that decreases the probability of (in Stoic terms) destructive second movements. Through meditation and the right attitude, one allows emotions to happen to oneself (it is impossible to prevent this), but one is advised to observe the emotions without necessarily acting on them; one achieves some distance and decides what has value and what does not have value. Additionally, the  
Stoic idea of developing virtue in oneself, of becoming a good person, which the Stoics believed we could do because we have a touch of the divine, laid the foundation for the three monotheistic religions: Judaism, Christianity, and Islam . . . As with Stoicism, tenets of these religions include controlling our emotions lest we engage in sinful behavior.

\vspace{1em}
\textbf{Question 17}

**There is a sentence that is missing in the paragraph below. Look at the paragraph and decide in which blank (option 1, 2, 3, or 4) the following sentence would best fit.**
**Sentence** : Having made citizens more and less knowledgeable than their predecessors, the Internet has proved to be both a blessing and a curse.
**Paragraph** : Never before has a population, nearly all of whom has enjoyed at a least a secondary school education, been exposed to so much information, whether in newspapers and magazines or through YouTube, Google, and Facebook. ___(1)___. Yet it is not clear that people today are more knowledgeable than their barely literate predecessors. Contemporary advances in technology offered more serious and inquisitive students access to realms of knowledge previously unimaginable and unavailable. ___(2)___. But such readily available knowledge leads many more students away from serious study, the reading of actual texts, and toward an inability to write effectively and grammatically. ___(3)___. It has let people choose sources that reinforce their opinions rather than encouraging them to question inherited beliefs. ___(4)___.

\begin{enumerate}[label=\Alph*.]
    \item Option 1
    \item Option 2
    \item Option 3
    \item Option 4
[](https://cracku.in/17-there-is-a-sentence-that-is-missing-in-the-paragra-x-cat-2022-slot-1)
\end{enumerate}
\vspace{1em}
\textbf{Solution:}

The sentence would best fit in Blank 4 because it ties together the ideas presented in the paragraph. The paragraph states that people today have access to a vast amount of information, but it is not clear if they are more knowledgeable as a result. It also mentions that the readily available knowledge can lead students away from serious study and towards an inability to write effectively. The missing sentence introduces the idea that the Internet has both positive and negative effects on knowledge and learning. It fits well in Option 4 because it connects to the idea that the Internet has let people choose sources that reinforce their opinions rather than encouraging them to question inherited beliefs, which is the final point presented in the paragraph.

\end{problem}

\newpage
\begin{problem}{}{}
% Subject: varc
\textbf{Instructions}

Instructions   

**The passage below is accompanied by a set of questions. Choose the best answer to each question.**
**Comprehension** :
Stoicism was founded in 300 BC by the Greek philosopher Zeno and survived into the Roman era until about AD 300. According to the Stoics, emotions consist of two movements. The first movement is the immediate feeling and other reactions (e.g., physiological response) that occur when a stimulus or event occurs. For instance, consider what could have happened if an army general accused Marcus Aurelius of treason in front of other officers. The first movement for Marcus may have been (internal) surprise and anger in response to this insult, accompanied perhaps by some involuntary physiological and expressive responses such as face flushing and a movement of the eyebrows. The second movement is what one does next about the emotion. Second movement behaviors occur after thinking and are under one’s control. Examples of second movements for Marcus might have included a plot to seek revenge, actions signifying deference and appeasement, or perhaps proceeding as he would have proceeded whether or not this event occurred: continuing to lead the Romans in a way that Marcus Aurelius believed best benefited them. In the Stoic view, choosing a reasoned, unemotional response as the second movement is the only appropriate response.
The Stoics believed that to live the good life and be a good person, we need to free ourselves of nearly all desires such as too much desire for money, power, or sexual gratification. Prior to second movements, we can consider what is important in life. Money, power, and excessive sexual gratification are not important. Character, rationality, and kindness are important. The Epicureans, first associated with the Greek philosopher Epicurus . . . held a similar view, believing that people should enjoy simple pleasures, such as good conversation, friendship, food, and wine, but not be indulgent in these pursuits and not follow passion for those things that hold no real value like power and money. As Oatley (2004) states, “the Epicureans articulated a view—enjoyment of relationship with friends, of things that are real rather than illusory, simple rather than artificially inflated, possible rather than vanishingly unlikely—that is certainly relevant today” . . . In sum, these ancient Greek and Roman philosophers saw emotions, especially strong ones, as potentially dangerous. They viewed emotions as experiences that needed to be [reined] in and controlled.
As Oatley (2004) points out, the Stoic idea bears some similarity to Buddhism. Buddha, living in India in the 6th century BC, argued for cultivating a certain attitude that decreases the probability of (in Stoic terms) destructive second movements. Through meditation and the right attitude, one allows emotions to happen to oneself (it is impossible to prevent this), but one is advised to observe the emotions without necessarily acting on them; one achieves some distance and decides what has value and what does not have value. Additionally, the  
Stoic idea of developing virtue in oneself, of becoming a good person, which the Stoics believed we could do because we have a touch of the divine, laid the foundation for the three monotheistic religions: Judaism, Christianity, and Islam . . . As with Stoicism, tenets of these religions include controlling our emotions lest we engage in sinful behavior.

\vspace{1em}
\textbf{Question 18}

**The passage given below is followed by four alternate summaries. Choose the option that best captures the essence of the passage.**
Petitioning is an expeditious democratic tradition, used frequently in prior centuries, by which citizens can bring issues directly to governments. As expressions of collective voice, they support procedural democracy by shaping agendas. They can also recruit citizens to causes, give voice to the voteless, and apply the discipline of rhetorical argument that clarifies a point of view. By contrast, elections are limited in several respects: they involve only a few candidates, and thus fall far short of a representative democracy. Further, voters’ choices are not specific to particular policies or laws, and elections are episodic, whereas the voice of the people needs to be heard and integrated constantly into democratic government.

\begin{enumerate}[label=\Alph*.]
    \item By giving citizens greater control over shaping political and democratic agendas, political petitions are invaluable as they represent an ideal form of a representative democracy.
    \item Citizens become less inclined to petitioning as it enables vocal citizens to shape political agendas, but this needs to change to strengthen democracies today.
    \item Petitioning has been important to democratic functioning, as it supplements the electoral process by enabling ongoing engagement with the government.
    \item Petitioning is definitely more representative of the collective voice, and the functioning of democratic government could improve if we relied more on petitioning rather than holding periodic elections.
[](https://cracku.in/18-the-passage-given-below-is-followed-by-four-altern-x-cat-2022-slot-1)
\end{enumerate}
\vspace{1em}
\textbf{Solution:}

The given passage discusses the role of petitioning in democratic governance; it highlights how petitioning can shape political agendas, recruit citizens to causes, give voice to the voteless, and apply the discipline of rhetorical argument to clarify points of view. The passage also compares petitioning to elections, stating that petitioning supplements the electoral process and enables ongoing engagement with the government. The author does not make a judgment about the relative superiority of petitioning or elections, but rather focuses on the specific ways in which petitioning can contribute to democratic functioning.  _Option C_ accurately summarizes the points discussed above. 
_Option A_ is incorrect because the passage does not state that petitioning is "an ideal form of representative democracy." It only mentions that petitioning can give voice to the voteless and apply the discipline of rhetorical argument, but it does not make a broader claim about the superiority of petitioning over other forms of democracy.
_Option B_ is incorrect because the passage does not suggest that citizens are becoming less inclined to petitioning. In fact, the passage does not address the prevalence of petitioning in contemporary times at all.
_Option D_ is incorrect because the passage does not state that petitioning is more representative of the collective voice than elections. It only mentions that petitioning can give voice to the voteless, but it does not make a comparison to elections in this regard. The passage also does not suggest that the functioning of democratic government would necessarily improve if we relied more on petitioning instead of holding periodic elections; it simply states that petitioning has been important to democratic functioning in the past, and that it supplements the electoral process by enabling ongoing engagement with the government.

\end{problem}

\newpage
\begin{problem}{}{}
% Subject: varc
\textbf{Instructions}

Instructions   

**The passage below is accompanied by a set of questions. Choose the best answer to each question.**
**Comprehension** :
Stoicism was founded in 300 BC by the Greek philosopher Zeno and survived into the Roman era until about AD 300. According to the Stoics, emotions consist of two movements. The first movement is the immediate feeling and other reactions (e.g., physiological response) that occur when a stimulus or event occurs. For instance, consider what could have happened if an army general accused Marcus Aurelius of treason in front of other officers. The first movement for Marcus may have been (internal) surprise and anger in response to this insult, accompanied perhaps by some involuntary physiological and expressive responses such as face flushing and a movement of the eyebrows. The second movement is what one does next about the emotion. Second movement behaviors occur after thinking and are under one’s control. Examples of second movements for Marcus might have included a plot to seek revenge, actions signifying deference and appeasement, or perhaps proceeding as he would have proceeded whether or not this event occurred: continuing to lead the Romans in a way that Marcus Aurelius believed best benefited them. In the Stoic view, choosing a reasoned, unemotional response as the second movement is the only appropriate response.
The Stoics believed that to live the good life and be a good person, we need to free ourselves of nearly all desires such as too much desire for money, power, or sexual gratification. Prior to second movements, we can consider what is important in life. Money, power, and excessive sexual gratification are not important. Character, rationality, and kindness are important. The Epicureans, first associated with the Greek philosopher Epicurus . . . held a similar view, believing that people should enjoy simple pleasures, such as good conversation, friendship, food, and wine, but not be indulgent in these pursuits and not follow passion for those things that hold no real value like power and money. As Oatley (2004) states, “the Epicureans articulated a view—enjoyment of relationship with friends, of things that are real rather than illusory, simple rather than artificially inflated, possible rather than vanishingly unlikely—that is certainly relevant today” . . . In sum, these ancient Greek and Roman philosophers saw emotions, especially strong ones, as potentially dangerous. They viewed emotions as experiences that needed to be [reined] in and controlled.
As Oatley (2004) points out, the Stoic idea bears some similarity to Buddhism. Buddha, living in India in the 6th century BC, argued for cultivating a certain attitude that decreases the probability of (in Stoic terms) destructive second movements. Through meditation and the right attitude, one allows emotions to happen to oneself (it is impossible to prevent this), but one is advised to observe the emotions without necessarily acting on them; one achieves some distance and decides what has value and what does not have value. Additionally, the  
Stoic idea of developing virtue in oneself, of becoming a good person, which the Stoics believed we could do because we have a touch of the divine, laid the foundation for the three monotheistic religions: Judaism, Christianity, and Islam . . . As with Stoicism, tenets of these religions include controlling our emotions lest we engage in sinful behavior.

\vspace{1em}
\textbf{Question 19}

**There is a sentence that is missing in the paragraph below. Look at the paragraph and decide in which blank (option 1, 2, 3, or 4) the following sentence would best fit.**
**Sentence** : Easing the anxiety and pressure of having a “big day” is part of the appeal for many couples who marry in secret.
**Paragraph** : Wedding season is upon us and - after two years of Covid chaos that saw nuptials scaled back- you may think the temptation would be to go all out. ___(1)___. But instead of expanding the guest list, many couples are opting to have entirely secret ceremonies. With Covid case numbers remaining high and the cost of living crisis meaning that many couples are feeling the pinch, it’s no wonder that some are less than eager to send out invites. ___(2)___. Plus, it can’t hurt that in celebrity circles getting married in secret is all the rage. ___(3)___. “I would definitely say that secret weddings are becoming more common,” says Landis Bejar, the founder of a therapy practice, which specialises in helping brides and grooms manage wedding stress. “People are looking for ways to get out of the spotlight and avoid the pomp and circumstance of weddings. ___(4)___. They just want to get to the part where they are married.”

\begin{enumerate}[label=\Alph*.]
    \item Option 1
    \item Option 2
    \item Option 3
    \item Option 4
[](https://cracku.in/19-there-is-a-sentence-that-is-missing-in-the-paragra-x-cat-2022-slot-1)
\end{enumerate}
\vspace{1em}
\textbf{Solution:}

The sentence would best fit Blank 2 because it ties together the ideas presented in the paragraph. The paragraph describes the change in the mindset for wedding celebrations post covid. In the first two lines, the author mentions this mindset. The given sentence will fill the blank 2 as it restates the idea presented in the previous line that many couples are marrying secretly to ease the anxiety and pressure.  

Thus, the correct option is B.

\end{problem}

\newpage
\begin{problem}{}{}
% Subject: varc
\textbf{Instructions}

Instructions   

**The passage below is accompanied by a set of questions. Choose the best answer to each question.**
**Comprehension** :
Stoicism was founded in 300 BC by the Greek philosopher Zeno and survived into the Roman era until about AD 300. According to the Stoics, emotions consist of two movements. The first movement is the immediate feeling and other reactions (e.g., physiological response) that occur when a stimulus or event occurs. For instance, consider what could have happened if an army general accused Marcus Aurelius of treason in front of other officers. The first movement for Marcus may have been (internal) surprise and anger in response to this insult, accompanied perhaps by some involuntary physiological and expressive responses such as face flushing and a movement of the eyebrows. The second movement is what one does next about the emotion. Second movement behaviors occur after thinking and are under one’s control. Examples of second movements for Marcus might have included a plot to seek revenge, actions signifying deference and appeasement, or perhaps proceeding as he would have proceeded whether or not this event occurred: continuing to lead the Romans in a way that Marcus Aurelius believed best benefited them. In the Stoic view, choosing a reasoned, unemotional response as the second movement is the only appropriate response.
The Stoics believed that to live the good life and be a good person, we need to free ourselves of nearly all desires such as too much desire for money, power, or sexual gratification. Prior to second movements, we can consider what is important in life. Money, power, and excessive sexual gratification are not important. Character, rationality, and kindness are important. The Epicureans, first associated with the Greek philosopher Epicurus . . . held a similar view, believing that people should enjoy simple pleasures, such as good conversation, friendship, food, and wine, but not be indulgent in these pursuits and not follow passion for those things that hold no real value like power and money. As Oatley (2004) states, “the Epicureans articulated a view—enjoyment of relationship with friends, of things that are real rather than illusory, simple rather than artificially inflated, possible rather than vanishingly unlikely—that is certainly relevant today” . . . In sum, these ancient Greek and Roman philosophers saw emotions, especially strong ones, as potentially dangerous. They viewed emotions as experiences that needed to be [reined] in and controlled.
As Oatley (2004) points out, the Stoic idea bears some similarity to Buddhism. Buddha, living in India in the 6th century BC, argued for cultivating a certain attitude that decreases the probability of (in Stoic terms) destructive second movements. Through meditation and the right attitude, one allows emotions to happen to oneself (it is impossible to prevent this), but one is advised to observe the emotions without necessarily acting on them; one achieves some distance and decides what has value and what does not have value. Additionally, the  
Stoic idea of developing virtue in oneself, of becoming a good person, which the Stoics believed we could do because we have a touch of the divine, laid the foundation for the three monotheistic religions: Judaism, Christianity, and Islam . . . As with Stoicism, tenets of these religions include controlling our emotions lest we engage in sinful behavior.

\vspace{1em}
\textbf{Question 20}

**The passage given below is followed by four alternate summaries. Choose the option that best captures the essence of the passage.**
It’s not that modern historians of medieval Africa have been ignorant about contacts between Ethiopia and Europe; they just had the power dynamic reversed. The traditional narrative stressed Ethiopia as weak and in trouble in the face of aggression from external forces, so Ethiopia sought military assistance from their fellow Christians to the north. But the real story, buried in plain sight in medieval diplomatic texts, simply had not yet been put together by modern scholars. Recent research pushes scholars of medieval Europe to imagine a much more richly connected medieval world: at the beginning of the so-called Age of Exploration, there is evidence that the kings of Ethiopia were sponsoring their own missions of diplomacy, faith and commerce.

\begin{enumerate}[label=\Alph*.]
    \item Medieval texts have documented how strong connections between the Christian communities of Ethiopia and Europe were invaluable in establishing military and trade links between the two civilisations.
    \item Historians were under the illusion that Ethiopia needed military protection from their neighbours, but in fact the country had close commercial and religious connections with them.
    \item Medieval texts have been ‘cherry-picked’ to promote a view of Ethiopia as weak and in need of Europe’s military help with aggressive neighbours, but recent studies reveal it was a well-connected and outward-looking culture.
    \item Medieval historical sources selectively promoted the narrative that powerful European forces were called on to protect weak African civilisations such as Ethiopia, but this is far from reality.
[](https://cracku.in/20-the-passage-given-below-is-followed-by-four-altern-x-cat-2022-slot-1)
\end{enumerate}
\vspace{1em}
\textbf{Solution:}

The passage touches upon the historical relationship between Ethiopia and Europe during the medieval period. The traditional narrative of this relationship has portrayed Ethiopia as weak and in need of military assistance from Europe, but recent research has revealed that this narrative is not accurate and that Ethiopia was actually a well-connected and outward-looking culture that engaged in missions of diplomacy, faith, and commerce with Europe. The passage also notes that these new findings challenge historians to re-imagine the connections between Ethiopia and Europe during this period and to consider the role of Ethiopia as a more active participant in these relationships. Option C accurately reflects the main points made in the passage [that the traditional narrative of Ethiopia's relationship with Europe is inaccurate]
Option A is incorrect because it suggests that the connections between Ethiopia and Europe were primarily military and commercial in nature, which is not stated in the passage. Similarly, Option B implies that historians had a mistaken view of Ethiopia's relationship with Europe, which is not stated in the passage. The idea in Option D - which suggests that medieval texts were biased in favour of Europe and against Africa - cannot be understood from the discussion.
Hence, Option C is the correct choice.

\end{problem}

\newpage
\begin{problem}{}{}
% Subject: varc
\textbf{Instructions}

Instructions   

**The passage below is accompanied by a set of questions. Choose the best answer to each question.**
**Comprehension** :
Stoicism was founded in 300 BC by the Greek philosopher Zeno and survived into the Roman era until about AD 300. According to the Stoics, emotions consist of two movements. The first movement is the immediate feeling and other reactions (e.g., physiological response) that occur when a stimulus or event occurs. For instance, consider what could have happened if an army general accused Marcus Aurelius of treason in front of other officers. The first movement for Marcus may have been (internal) surprise and anger in response to this insult, accompanied perhaps by some involuntary physiological and expressive responses such as face flushing and a movement of the eyebrows. The second movement is what one does next about the emotion. Second movement behaviors occur after thinking and are under one’s control. Examples of second movements for Marcus might have included a plot to seek revenge, actions signifying deference and appeasement, or perhaps proceeding as he would have proceeded whether or not this event occurred: continuing to lead the Romans in a way that Marcus Aurelius believed best benefited them. In the Stoic view, choosing a reasoned, unemotional response as the second movement is the only appropriate response.
The Stoics believed that to live the good life and be a good person, we need to free ourselves of nearly all desires such as too much desire for money, power, or sexual gratification. Prior to second movements, we can consider what is important in life. Money, power, and excessive sexual gratification are not important. Character, rationality, and kindness are important. The Epicureans, first associated with the Greek philosopher Epicurus . . . held a similar view, believing that people should enjoy simple pleasures, such as good conversation, friendship, food, and wine, but not be indulgent in these pursuits and not follow passion for those things that hold no real value like power and money. As Oatley (2004) states, “the Epicureans articulated a view—enjoyment of relationship with friends, of things that are real rather than illusory, simple rather than artificially inflated, possible rather than vanishingly unlikely—that is certainly relevant today” . . . In sum, these ancient Greek and Roman philosophers saw emotions, especially strong ones, as potentially dangerous. They viewed emotions as experiences that needed to be [reined] in and controlled.
As Oatley (2004) points out, the Stoic idea bears some similarity to Buddhism. Buddha, living in India in the 6th century BC, argued for cultivating a certain attitude that decreases the probability of (in Stoic terms) destructive second movements. Through meditation and the right attitude, one allows emotions to happen to oneself (it is impossible to prevent this), but one is advised to observe the emotions without necessarily acting on them; one achieves some distance and decides what has value and what does not have value. Additionally, the  
Stoic idea of developing virtue in oneself, of becoming a good person, which the Stoics believed we could do because we have a touch of the divine, laid the foundation for the three monotheistic religions: Judaism, Christianity, and Islam . . . As with Stoicism, tenets of these religions include controlling our emotions lest we engage in sinful behavior.

\vspace{1em}
\textbf{Question 21}

**The four sentences (labelled 1, 2, 3 and 4) below, when properly sequenced, would yield a coherent paragraph. Decide on the proper sequencing of the order of the sentences and key in the sequence of the four numbers as your answer:**
1. Some company leaders are basing their decisions on locating offices to foster innovation and growth, as their best-performing inventors suffered the greatest productivity losses when their commutes grew longer.
2. Shorter commutes support innovation by giving employees more time in the office and greater opportunities for in-person collaboration, while removing the physical strain of a long commute.
3. This is not always the case: remote work does not automatically lead to greater creativity and productivity as office water-cooler conversations are also very important for innovation.
4. Some see the link between long commutes and productivity as support for work-from-home scenarios, as many workers have grown accustomed to their commute-free arrangements during the pandemic.
Backspace  
789  
456  
123  
0.-  
Clear All  

Submit
[](https://cracku.in/21-the-four-sentences-labelled-1-2-3-and-4-below-when-x-cat-2022-slot-1)

\begin{enumerate}[label=\Alph*.]
\end{enumerate}
\vspace{1em}
\textbf{Solution:}

The given set of sentences underlines the relationship between commuting and productivity and how it can impact innovation. The idea is that shorter commutes can support innovation by giving employees more time in the office and opportunities for in-person collaboration, while longer commutes can be physically draining and can lead to decreased productivity. This correlation is first highlighted in statement 2, which sets the context for subsequent discussion. We note that statements 4 and 3 form a logical block since 4 presents an idea [work-from-home scenarios to improve productivity] and 3 extends on it [why the idea in 4 might not be effective]. The marker word 'some' in 4 helps us link it with statement 1, which discusses the opinions of "some company leaders." So 1 puts forth the belief of one group of company leaders, and 4 presents us with the opinion held by a few other company leaders. This allows us to form the logical block: [1-4-3]. We can place 2 at the beginning of this structure to obtain a coherent paragraph. Hence, the correct answer is 2143.

\end{problem}

\newpage
\begin{problem}{}{}
% Subject: varc
\textbf{Instructions}

Instructions   

**The passage below is accompanied by a set of questions. Choose the best answer to each question.**
**Comprehension** :
Stoicism was founded in 300 BC by the Greek philosopher Zeno and survived into the Roman era until about AD 300. According to the Stoics, emotions consist of two movements. The first movement is the immediate feeling and other reactions (e.g., physiological response) that occur when a stimulus or event occurs. For instance, consider what could have happened if an army general accused Marcus Aurelius of treason in front of other officers. The first movement for Marcus may have been (internal) surprise and anger in response to this insult, accompanied perhaps by some involuntary physiological and expressive responses such as face flushing and a movement of the eyebrows. The second movement is what one does next about the emotion. Second movement behaviors occur after thinking and are under one’s control. Examples of second movements for Marcus might have included a plot to seek revenge, actions signifying deference and appeasement, or perhaps proceeding as he would have proceeded whether or not this event occurred: continuing to lead the Romans in a way that Marcus Aurelius believed best benefited them. In the Stoic view, choosing a reasoned, unemotional response as the second movement is the only appropriate response.
The Stoics believed that to live the good life and be a good person, we need to free ourselves of nearly all desires such as too much desire for money, power, or sexual gratification. Prior to second movements, we can consider what is important in life. Money, power, and excessive sexual gratification are not important. Character, rationality, and kindness are important. The Epicureans, first associated with the Greek philosopher Epicurus . . . held a similar view, believing that people should enjoy simple pleasures, such as good conversation, friendship, food, and wine, but not be indulgent in these pursuits and not follow passion for those things that hold no real value like power and money. As Oatley (2004) states, “the Epicureans articulated a view—enjoyment of relationship with friends, of things that are real rather than illusory, simple rather than artificially inflated, possible rather than vanishingly unlikely—that is certainly relevant today” . . . In sum, these ancient Greek and Roman philosophers saw emotions, especially strong ones, as potentially dangerous. They viewed emotions as experiences that needed to be [reined] in and controlled.
As Oatley (2004) points out, the Stoic idea bears some similarity to Buddhism. Buddha, living in India in the 6th century BC, argued for cultivating a certain attitude that decreases the probability of (in Stoic terms) destructive second movements. Through meditation and the right attitude, one allows emotions to happen to oneself (it is impossible to prevent this), but one is advised to observe the emotions without necessarily acting on them; one achieves some distance and decides what has value and what does not have value. Additionally, the  
Stoic idea of developing virtue in oneself, of becoming a good person, which the Stoics believed we could do because we have a touch of the divine, laid the foundation for the three monotheistic religions: Judaism, Christianity, and Islam . . . As with Stoicism, tenets of these religions include controlling our emotions lest we engage in sinful behavior.

\vspace{1em}
\textbf{Question 22}

**The four sentences (labelled 1, 2, 3 and 4) below, when properly sequenced, would yield a coherent paragraph. Decide on the proper sequencing of the order of the sentences and key in the sequence of the four numbers as your answer:**
1. The creative element in product design has become of paramount importance as it is one of the few ways a firm or industry can sustain a competitive advantage over its rivals.
2. In fact, the creative element in the value of world industry would be larger still, if we added the contribution of the creative element in other industries, such as the design of tech accessories.
3. The creative industry is receiving a lot of attention today as its growth rate is faster than that of the world economy as a whole.
4. It is for this reason that today’s trade issues are increasingly involving intellectual property, as Western countries have an interest in protecting their revenues along with freeing trade in non-tangibles.
Backspace  
789  
456  
123  
0.-  
Clear All  

Submit
[](https://cracku.in/22-the-four-sentences-labelled-1-2-3-and-4-below-when-x-cat-2022-slot-1)

\begin{enumerate}[label=\Alph*.]
\end{enumerate}
\vspace{1em}
\textbf{Solution:}

A brief reading of the sentences suggests that the paragraph is about the applications of the creative element in the world economy. Sentence 3 gives the introduction by contrasting the growth rate of the world economy and the creative industry. Statement 2 follows 3 by stating that the difference will be more if we consider the contribution of the creative element in other industries. Sentence 1 extends this idea by explaining how the creative element in product design can help sustain a competitive advantage. Sentence 4 concludes the paragraph by stating the implications of the application of the creative element mentioned in sentence 1.
Thus, the correct order will be 3-2-1-4.

\end{problem}

\newpage
\begin{problem}{}{}
% Subject: varc
\textbf{Instructions}

Instructions   

**The passage below is accompanied by a set of questions. Choose the best answer to each question.**
**Comprehension** :
Stoicism was founded in 300 BC by the Greek philosopher Zeno and survived into the Roman era until about AD 300. According to the Stoics, emotions consist of two movements. The first movement is the immediate feeling and other reactions (e.g., physiological response) that occur when a stimulus or event occurs. For instance, consider what could have happened if an army general accused Marcus Aurelius of treason in front of other officers. The first movement for Marcus may have been (internal) surprise and anger in response to this insult, accompanied perhaps by some involuntary physiological and expressive responses such as face flushing and a movement of the eyebrows. The second movement is what one does next about the emotion. Second movement behaviors occur after thinking and are under one’s control. Examples of second movements for Marcus might have included a plot to seek revenge, actions signifying deference and appeasement, or perhaps proceeding as he would have proceeded whether or not this event occurred: continuing to lead the Romans in a way that Marcus Aurelius believed best benefited them. In the Stoic view, choosing a reasoned, unemotional response as the second movement is the only appropriate response.
The Stoics believed that to live the good life and be a good person, we need to free ourselves of nearly all desires such as too much desire for money, power, or sexual gratification. Prior to second movements, we can consider what is important in life. Money, power, and excessive sexual gratification are not important. Character, rationality, and kindness are important. The Epicureans, first associated with the Greek philosopher Epicurus . . . held a similar view, believing that people should enjoy simple pleasures, such as good conversation, friendship, food, and wine, but not be indulgent in these pursuits and not follow passion for those things that hold no real value like power and money. As Oatley (2004) states, “the Epicureans articulated a view—enjoyment of relationship with friends, of things that are real rather than illusory, simple rather than artificially inflated, possible rather than vanishingly unlikely—that is certainly relevant today” . . . In sum, these ancient Greek and Roman philosophers saw emotions, especially strong ones, as potentially dangerous. They viewed emotions as experiences that needed to be [reined] in and controlled.
As Oatley (2004) points out, the Stoic idea bears some similarity to Buddhism. Buddha, living in India in the 6th century BC, argued for cultivating a certain attitude that decreases the probability of (in Stoic terms) destructive second movements. Through meditation and the right attitude, one allows emotions to happen to oneself (it is impossible to prevent this), but one is advised to observe the emotions without necessarily acting on them; one achieves some distance and decides what has value and what does not have value. Additionally, the  
Stoic idea of developing virtue in oneself, of becoming a good person, which the Stoics believed we could do because we have a touch of the divine, laid the foundation for the three monotheistic religions: Judaism, Christianity, and Islam . . . As with Stoicism, tenets of these religions include controlling our emotions lest we engage in sinful behavior.

\vspace{1em}
\textbf{Question 23}

**The passage given below is followed by four alternate summaries. Choose the option that best captures the essence of the passage.**
All that we think we know about how life hangs together is really some kind of illusion that we have perpetrated on ourselves because of our limited vision. What appear to be inanimate objects such as stones turn out not only to be alive in the same way that we are, but also in many infinitesimal ways to be affected by stimuli just as humans are. The distinction between animate and inanimate simply cannot be made when you enter the world of quantum mechanics and try to determine how those apparent subatomic particles, of which you and everything else in our universe is composed, are all tied together. The point is that physics and metaphysics show there is a pattern to the universe that goes beyond our capacity to grasp it with our brains.

\begin{enumerate}[label=\Alph*.]
    \item The effect of stimuli is similar in inanimate objects when compared to animate objects or living beings.
    \item Quantum physics indicates that an astigmatic view of reality results in erroneous assumptions about the universe.
    \item The inanimate world is both sentient and cognizant like its animate counterpart.
    \item Arbitrary distinctions between inanimate and animate objects disappear at the scale at which quantum mechanics works.
[](https://cracku.in/23-the-passage-given-below-is-followed-by-four-altern-x-cat-2022-slot-1)
\end{enumerate}
\vspace{1em}
\textbf{Solution:}

The passage is about the limitations of our understanding of the universe and how it is connected. It suggests that our understanding of the distinction between animate and inanimate objects may be flawed because of our limited perspective, and that quantum mechanics reveals a different understanding of this distinction; it also suggests that there may be a pattern to the universe that we cannot fully grasp with our brains. Option D correctly captures these points.
Options A and B are incorrect because the passage does not state that the effect of stimuli is similar in inanimate objects compared to living beings or mention astigmatism or an erroneous view of reality. Option C is not understood since the passage does not suggest that inanimate objects are sentient or cognizant; it only highlights that the distinction between animate and inanimate objects disappears at the quantum scale.
Hence, Option D is the correct choice.

\end{problem}

\newpage
\begin{problem}{}{}
% Subject: varc
\textbf{Instructions}

Instructions   

**The passage below is accompanied by a set of questions. Choose the best answer to each question.**
**Comprehension** :
Stoicism was founded in 300 BC by the Greek philosopher Zeno and survived into the Roman era until about AD 300. According to the Stoics, emotions consist of two movements. The first movement is the immediate feeling and other reactions (e.g., physiological response) that occur when a stimulus or event occurs. For instance, consider what could have happened if an army general accused Marcus Aurelius of treason in front of other officers. The first movement for Marcus may have been (internal) surprise and anger in response to this insult, accompanied perhaps by some involuntary physiological and expressive responses such as face flushing and a movement of the eyebrows. The second movement is what one does next about the emotion. Second movement behaviors occur after thinking and are under one’s control. Examples of second movements for Marcus might have included a plot to seek revenge, actions signifying deference and appeasement, or perhaps proceeding as he would have proceeded whether or not this event occurred: continuing to lead the Romans in a way that Marcus Aurelius believed best benefited them. In the Stoic view, choosing a reasoned, unemotional response as the second movement is the only appropriate response.
The Stoics believed that to live the good life and be a good person, we need to free ourselves of nearly all desires such as too much desire for money, power, or sexual gratification. Prior to second movements, we can consider what is important in life. Money, power, and excessive sexual gratification are not important. Character, rationality, and kindness are important. The Epicureans, first associated with the Greek philosopher Epicurus . . . held a similar view, believing that people should enjoy simple pleasures, such as good conversation, friendship, food, and wine, but not be indulgent in these pursuits and not follow passion for those things that hold no real value like power and money. As Oatley (2004) states, “the Epicureans articulated a view—enjoyment of relationship with friends, of things that are real rather than illusory, simple rather than artificially inflated, possible rather than vanishingly unlikely—that is certainly relevant today” . . . In sum, these ancient Greek and Roman philosophers saw emotions, especially strong ones, as potentially dangerous. They viewed emotions as experiences that needed to be [reined] in and controlled.
As Oatley (2004) points out, the Stoic idea bears some similarity to Buddhism. Buddha, living in India in the 6th century BC, argued for cultivating a certain attitude that decreases the probability of (in Stoic terms) destructive second movements. Through meditation and the right attitude, one allows emotions to happen to oneself (it is impossible to prevent this), but one is advised to observe the emotions without necessarily acting on them; one achieves some distance and decides what has value and what does not have value. Additionally, the  
Stoic idea of developing virtue in oneself, of becoming a good person, which the Stoics believed we could do because we have a touch of the divine, laid the foundation for the three monotheistic religions: Judaism, Christianity, and Islam . . . As with Stoicism, tenets of these religions include controlling our emotions lest we engage in sinful behavior.

\vspace{1em}
\textbf{Question 24}

**The four sentences (labelled 1, 2, 3 and 4) below, when properly sequenced, would yield a coherent paragraph. Decide on the proper sequencing of the order of the sentences and key in the sequence of the four numbers as your answer:**
1. Fish skin collagen has excellent thermo-stability and tensile strength making it ideal for use as bandage that adheres to the skin and adjusts to body movements.
2. Collagen, one of the main structural proteins in connective tissues in the human body, is well known for promoting skin regeneration.
3. Fish skin swims in here as diseases and bacteria that affect fish are different from most human pathogens.
4. The risk of introducing disease agents into other species through the use of pig and cow collagen proteins for wound healing has inhibited its broader applications in the medical field.
Backspace  
789  
456  
123  
0.-  
Clear All  

Submit
[](https://cracku.in/24-the-four-sentences-labelled-1-2-3-and-4-below-when-x-cat-2022-slot-1)

\begin{enumerate}[label=\Alph*.]
\end{enumerate}
\vspace{1em}
\textbf{Solution:}

A brief reading of the sentences suggests that the paragraph is about the applications of a structural protein called Collagen. Sentence 2 introduces the paragraph by defining Collagen. 2-4 is an obvious link, as sentence 2 ends by stating that Collagen promotes skin regeneration, and sentence 4 explains why it is not widely used in the medical field despite this. Sentence 3 extends the idea given in sentence 4 by contrasting the bacteria affecting fish skins with human pathogens. Sentence 1 follows sentence 3 by providing the applications of fish skin collagen.
Thus, the correct order will be 2-4-3-1.
[](https://cracku.in/downloads/17290)
  *  
  *

\end{problem}

\newpage
\begin{problem}{}{}
% Subject: varc
\textbf{Instructions}

Instructions   

**The passage below is accompanied by a set of questions. Choose the best answer to each question.**
[Octopuses are] misfits in their own extended families . . . They belong to the Mollusca class Cephalopoda. But they don’t look like their cousins at all. Other molluscs include sea snails, sea slugs, bivalves - most are shelled invertebrates with a dorsal foot. Cephalopods are all arms, and can be as tiny as 1 centimetre and as large at 30 feet. Some of them have brains the size of a walnut, which is large for an invertebrate. . . .
It makes sense for these molluscs to have added protection in the form of a higher cognition; they don’t have a shell covering them, and pretty much everything feeds on cephalopods, including humans. But how did cephalopods manage to secure their own invisibility cloak? Cephalopods fire from multiple cylinders to achieve this in varying degrees from species to species. There are four main catalysts - chromatophores, iridophores, papillae and leucophores. . . .
[Chromatophores] are organs on their bodies that contain pigment sacs, which have red, yellow and brown pigment granules. These sacs have a network of radial muscles, meaning muscles arranged in a circle radiating outwards. These are connected to the brain by a nerve. When the cephalopod wants to change colour, the brain carries an electrical impulse through the nerve to the muscles that expand outwards, pulling open the sacs to display the colours on the skin. Why these three colours? Because these are the colours the light reflects at the depths they live in (the rest is absorbed before it reaches those depths). . . .
Well, what about other colours? Cue the iridophores. Think of a second level of skin that has thin stacks of cells. These can reflect light back at different wavelengths. . . . It’s using the same properties that we’ve seen in hologram stickers, or rainbows on puddles of oil. You move your head and you see a different colour. The sticker isn’t doing anything but reflecting light - it’s your movement that’s changing the appearance of the colour. This property of holograms, oil and other such surfaces is called “iridescence”. . . .
Papillae are sections of the skin that can be deformed to make a texture bumpy. Even humans possess them (goosebumps) but cannot use them in the manner that cephalopods can. For instance, the use of these cells is how an octopus can wrap itself over a rock and appear jagged or how a squid or cuttlefish can imitate the look of a coral reef by growing miniature towers on its skin. It actually matches the texture of the substrate it chooses.
Finally, the leucophores: According to a paper, published in Nature, cuttlefish and octopuses possess an additional type of reflector cell called a leucophore. They are cells that scatter full spectrum light so that they appear white in a similar way that a polar bear’s fur appears white. Leucophores will also reflect any filtered light shown on them . . . If the water appears blue at a certain depth, the octopuses and cuttlefish can appear blue; if the water appears green, they appear green, and so on and so forth.

\vspace{1em}
\textbf{Question 1}

Based on the passage, it can be inferred that camouflaging techniques in an octopus are most dissimilar to those in:

\begin{enumerate}[label=\Alph*.]
    \item sea snails
    \item cuttlefish
    \item polar bears
    \item squids
[](https://cracku.in/1-based-on-the-passage-it-can-be-inferred-that-camou-x-cat-2022-slot-2)
\end{enumerate}
\vspace{1em}
\textbf{Solution:}

The author discusses camouflaging techniques in an octopus vis-a-vis other organisms such as cuttlefish, squids and polar bears: { For instance, the use of these cells is how an octopus can wrap itself over a rock and appear jagged or how a **squid** or **cuttlefish** can imitate the look of a coral reef by growing miniature towers on its skin. It actually matches the texture of the substrate it chooses...Finally, the leucophores: According to a paper, published in Nature, cuttlefish and octopuses possess an additional type of reflector cell called a leucophore. They are cells that scatter full spectrum light so that they appear white in a similar way that a **polar bear’s fur** appears white}
However, note that no such discussion on sea snails is presented; hence, Option A is the correct choice.

\end{problem}

\newpage
\begin{problem}{}{}
% Subject: varc
\textbf{Instructions}

Instructions   

**The passage below is accompanied by a set of questions. Choose the best answer to each question.**
[Octopuses are] misfits in their own extended families . . . They belong to the Mollusca class Cephalopoda. But they don’t look like their cousins at all. Other molluscs include sea snails, sea slugs, bivalves - most are shelled invertebrates with a dorsal foot. Cephalopods are all arms, and can be as tiny as 1 centimetre and as large at 30 feet. Some of them have brains the size of a walnut, which is large for an invertebrate. . . .
It makes sense for these molluscs to have added protection in the form of a higher cognition; they don’t have a shell covering them, and pretty much everything feeds on cephalopods, including humans. But how did cephalopods manage to secure their own invisibility cloak? Cephalopods fire from multiple cylinders to achieve this in varying degrees from species to species. There are four main catalysts - chromatophores, iridophores, papillae and leucophores. . . .
[Chromatophores] are organs on their bodies that contain pigment sacs, which have red, yellow and brown pigment granules. These sacs have a network of radial muscles, meaning muscles arranged in a circle radiating outwards. These are connected to the brain by a nerve. When the cephalopod wants to change colour, the brain carries an electrical impulse through the nerve to the muscles that expand outwards, pulling open the sacs to display the colours on the skin. Why these three colours? Because these are the colours the light reflects at the depths they live in (the rest is absorbed before it reaches those depths). . . .
Well, what about other colours? Cue the iridophores. Think of a second level of skin that has thin stacks of cells. These can reflect light back at different wavelengths. . . . It’s using the same properties that we’ve seen in hologram stickers, or rainbows on puddles of oil. You move your head and you see a different colour. The sticker isn’t doing anything but reflecting light - it’s your movement that’s changing the appearance of the colour. This property of holograms, oil and other such surfaces is called “iridescence”. . . .
Papillae are sections of the skin that can be deformed to make a texture bumpy. Even humans possess them (goosebumps) but cannot use them in the manner that cephalopods can. For instance, the use of these cells is how an octopus can wrap itself over a rock and appear jagged or how a squid or cuttlefish can imitate the look of a coral reef by growing miniature towers on its skin. It actually matches the texture of the substrate it chooses.
Finally, the leucophores: According to a paper, published in Nature, cuttlefish and octopuses possess an additional type of reflector cell called a leucophore. They are cells that scatter full spectrum light so that they appear white in a similar way that a polar bear’s fur appears white. Leucophores will also reflect any filtered light shown on them . . . If the water appears blue at a certain depth, the octopuses and cuttlefish can appear blue; if the water appears green, they appear green, and so on and so forth.

\vspace{1em}
\textbf{Question 2}

All of the following are reasons for octopuses being “misfits” EXCEPT that they:

\begin{enumerate}[label=\Alph*.]
    \item are consumed by humans and other animals.
    \item do not possess an outer protective shell.
    \item exhibit higher intelligence than other molluscs.
    \item have several arms
[](https://cracku.in/2-all-of-the-following-are-reasons-for-octopuses-bei-x-cat-2022-slot-2)
\end{enumerate}
\vspace{1em}
\textbf{Solution:}

{_But _they don’t look like their cousins_ at all. Other molluscs include sea snails, sea slugs, bivalves - _most are shelled invertebrates with a dorsal foot_. Cephalopods are _all arms_ , and can be as tiny as 1 centimetre and as large at 30 feet. Some of them have brains the size of a walnut, which is large for an invertebrate. . . .It _makes sense for these molluscs to have added protection in the form of a higher cognition_ ; _they don’t have a shell covering them_ , and pretty much everything feeds on cephalopods, including humans._} 
We can understand that B, C and D are true - they are points of dissimilarities between octopuses and other molluscs [we are told that they have multiple appendages instead of a single dorsal foot and lack shells; they also have higher intelligence]. However, there is no information on whether humans consume molluscs like sea snails or not [furthermore, the author does not use this fact to claim that octopuses are distinct in this regard]. 
Hence, Option A is the correct choice.

\end{problem}

\newpage
\begin{problem}{}{}
% Subject: varc
\textbf{Instructions}

Instructions   

**The passage below is accompanied by a set of questions. Choose the best answer to each question.**
[Octopuses are] misfits in their own extended families . . . They belong to the Mollusca class Cephalopoda. But they don’t look like their cousins at all. Other molluscs include sea snails, sea slugs, bivalves - most are shelled invertebrates with a dorsal foot. Cephalopods are all arms, and can be as tiny as 1 centimetre and as large at 30 feet. Some of them have brains the size of a walnut, which is large for an invertebrate. . . .
It makes sense for these molluscs to have added protection in the form of a higher cognition; they don’t have a shell covering them, and pretty much everything feeds on cephalopods, including humans. But how did cephalopods manage to secure their own invisibility cloak? Cephalopods fire from multiple cylinders to achieve this in varying degrees from species to species. There are four main catalysts - chromatophores, iridophores, papillae and leucophores. . . .
[Chromatophores] are organs on their bodies that contain pigment sacs, which have red, yellow and brown pigment granules. These sacs have a network of radial muscles, meaning muscles arranged in a circle radiating outwards. These are connected to the brain by a nerve. When the cephalopod wants to change colour, the brain carries an electrical impulse through the nerve to the muscles that expand outwards, pulling open the sacs to display the colours on the skin. Why these three colours? Because these are the colours the light reflects at the depths they live in (the rest is absorbed before it reaches those depths). . . .
Well, what about other colours? Cue the iridophores. Think of a second level of skin that has thin stacks of cells. These can reflect light back at different wavelengths. . . . It’s using the same properties that we’ve seen in hologram stickers, or rainbows on puddles of oil. You move your head and you see a different colour. The sticker isn’t doing anything but reflecting light - it’s your movement that’s changing the appearance of the colour. This property of holograms, oil and other such surfaces is called “iridescence”. . . .
Papillae are sections of the skin that can be deformed to make a texture bumpy. Even humans possess them (goosebumps) but cannot use them in the manner that cephalopods can. For instance, the use of these cells is how an octopus can wrap itself over a rock and appear jagged or how a squid or cuttlefish can imitate the look of a coral reef by growing miniature towers on its skin. It actually matches the texture of the substrate it chooses.
Finally, the leucophores: According to a paper, published in Nature, cuttlefish and octopuses possess an additional type of reflector cell called a leucophore. They are cells that scatter full spectrum light so that they appear white in a similar way that a polar bear’s fur appears white. Leucophores will also reflect any filtered light shown on them . . . If the water appears blue at a certain depth, the octopuses and cuttlefish can appear blue; if the water appears green, they appear green, and so on and so forth.

\vspace{1em}
\textbf{Question 3}

Which one of the following statements is not true about the camouflaging ability of Cephalopods?

\begin{enumerate}[label=\Alph*.]
    \item Cephalopods can blend into the colour of their surroundings.
    \item Cephalopods can change their texture.
    \item Cephalopods can change their colour.
    \item Cephalopods can take on the colour of their predator.
[](https://cracku.in/3-which-one-of-the-following-statements-is-not-true--x-cat-2022-slot-2)
\end{enumerate}
\vspace{1em}
\textbf{Solution:}

Options A and C: The discussion on Chromatophores, iridophores and leucophores sufficiently supports the statements here.
Option B: We can infer this from the following excerpt - {_Papillae are sections of the skin that can be deformed to make a texture bumpy... For instance, the use of these cells is how an octopus can wrap itself over a rock and appear jagged or how a squid or cuttlefish can imitate the look of a coral reef by growing miniature towers on its skin. It actually matches the texture of the substrate it chooses_.}
Option D is not discussed anywhere in the passage.

\end{problem}

\newpage
\begin{problem}{}{}
% Subject: varc
\textbf{Instructions}

Instructions   

**The passage below is accompanied by a set of questions. Choose the best answer to each question.**
[Octopuses are] misfits in their own extended families . . . They belong to the Mollusca class Cephalopoda. But they don’t look like their cousins at all. Other molluscs include sea snails, sea slugs, bivalves - most are shelled invertebrates with a dorsal foot. Cephalopods are all arms, and can be as tiny as 1 centimetre and as large at 30 feet. Some of them have brains the size of a walnut, which is large for an invertebrate. . . .
It makes sense for these molluscs to have added protection in the form of a higher cognition; they don’t have a shell covering them, and pretty much everything feeds on cephalopods, including humans. But how did cephalopods manage to secure their own invisibility cloak? Cephalopods fire from multiple cylinders to achieve this in varying degrees from species to species. There are four main catalysts - chromatophores, iridophores, papillae and leucophores. . . .
[Chromatophores] are organs on their bodies that contain pigment sacs, which have red, yellow and brown pigment granules. These sacs have a network of radial muscles, meaning muscles arranged in a circle radiating outwards. These are connected to the brain by a nerve. When the cephalopod wants to change colour, the brain carries an electrical impulse through the nerve to the muscles that expand outwards, pulling open the sacs to display the colours on the skin. Why these three colours? Because these are the colours the light reflects at the depths they live in (the rest is absorbed before it reaches those depths). . . .
Well, what about other colours? Cue the iridophores. Think of a second level of skin that has thin stacks of cells. These can reflect light back at different wavelengths. . . . It’s using the same properties that we’ve seen in hologram stickers, or rainbows on puddles of oil. You move your head and you see a different colour. The sticker isn’t doing anything but reflecting light - it’s your movement that’s changing the appearance of the colour. This property of holograms, oil and other such surfaces is called “iridescence”. . . .
Papillae are sections of the skin that can be deformed to make a texture bumpy. Even humans possess them (goosebumps) but cannot use them in the manner that cephalopods can. For instance, the use of these cells is how an octopus can wrap itself over a rock and appear jagged or how a squid or cuttlefish can imitate the look of a coral reef by growing miniature towers on its skin. It actually matches the texture of the substrate it chooses.
Finally, the leucophores: According to a paper, published in Nature, cuttlefish and octopuses possess an additional type of reflector cell called a leucophore. They are cells that scatter full spectrum light so that they appear white in a similar way that a polar bear’s fur appears white. Leucophores will also reflect any filtered light shown on them . . . If the water appears blue at a certain depth, the octopuses and cuttlefish can appear blue; if the water appears green, they appear green, and so on and so forth.

\vspace{1em}
\textbf{Question 4}

Based on the passage, we can infer that all of the following statements, if true, would weaken the camouflaging adeptness of Cephalopods EXCEPT:

\begin{enumerate}[label=\Alph*.]
    \item the number of chromatophores in Cephalopods is half the number of iridophores and leucophores.
    \item the temperature of water at the depths at which Cephalopods reside renders the transmission of neural signals difficult.
    \item light reflects the colours red, green, and yellow at the depths at which Cephalopods reside.
    \item the hydrostatic pressure at the depths at which Cephalopods reside renders radial muscle movements difficult.
[](https://cracku.in/4-based-on-the-passage-we-can-infer-that-all-of-the--x-cat-2022-slot-2)
\end{enumerate}
\vspace{1em}
\textbf{Solution:}

Let us evaluate the choices individually:
_Option A_ : [**the number of chromatophores in Cephalopods is half the number of iridophores and leucophores**] If true, this does not undermine the camouflaging adeptness of Cephalopods primarily because each of chromatophores, iridophores and leucophores have specific [somewhat independent] roles to play and it is unclear how their quantity would directly impact the camouflaging capacity. Even if the number of chromatophores is fewer than the other types, the octopus can still maintain its camouflaging adeptness. 
_Option B_ : [**the temperature of water at the depths at which Cephalopods reside renders the transmission of neural signals difficult.**] If true, this would limit the camouflaging adeptness primarily because the underlying mechanism is being restricted/impacted. {When the cephalopod wants to change colour, the brain carries an electrical impulse through the nerve to the muscles that expand outwards, pulling open the sacs to display the colours on the skin.}
_Option C_ : [**light reflects the colours red, green, and yellow at the depths at which Cephalopods reside.**] If true, this would limit the camouflaging adeptness primarily because the underlying mechanism is being restricted/impacted. If the colour scheme is distinct, it would undermine the observations presented below: {_[Chromatophores] are organs on their bodies that contain pigment sacs, which have red, yellow and brown pigment granules...Why these three colours? Because these are the colours the light reflects at the depths they live in (the rest is absorbed before it reaches those depths)_}
_Option D_ : [**the hydrostatic pressure at the depths at which Cephalopods reside renders radial muscle movements difficult.**] If true, this would limit the camouflaging adeptness primarily because the underlying mechanism is being restricted/impacted. {These sacs have a network of radial muscles, meaning muscles arranged in a circle radiating outwards.}
Hence, Option A is the correct choice.

Instructions   

**The passage below is accompanied by a set of questions. Choose the best answer to each question.**
When we teach engineering problems now, we ask students to come to a single “best” solution defined by technical ideals like low cost, speed to build, and ability to scale. This way of teaching primes students to believe that their decision-making is purely objective, as it is grounded in math and science. This is known as technical-social dualism, the idea that the technical and social dimensions of engineering problems are readily separable and remain distinct throughout the problem-definition and solution process.
Nontechnical parameters such as access to a technology, cultural relevancy or potential harms are deemed political and invalid in this way of learning. But those technical ideals are at their core social and political choices determined by a dominant culture focused on economic growth for the most privileged segments of society. By choosing to downplay public welfare as a critical parameter for engineering design, we risk creating a culture of disengagement from societal concerns amongst engineers that is antithetical to the ethical code of engineering.
In my field of medical devices, ignoring social dimensions has real consequences. . . . Most FDA-approved drugs are incorrectly dosed for people assigned female at birth, leading to unexpected adverse reactions. This is because they have been inadequately represented in clinical trials.
Beyond physical failings, subjective beliefs treated as facts by those in decision-making roles can encode social inequities. For example, spirometers, routinely used devices that measure lung capacity, still have correction factors that automatically assume smaller lung capacity in Black and Asian individuals. These racially based adjustments are derived from research done by eugenicists who thought these racial differences were biologically determined and who considered nonwhite people as inferior. These machines ignore the influence of social and environmental factors on lung capacity.
Many technologies for systemically marginalized people have not been built because they were not deemed important such as better early diagnostics and treatment for diseases like endometriosis, a disease that afflicts 10 percent of people with uteruses. And we hardly question whether devices are built sustainably, which has led to a crisis of medical waste and health care accounting for 10 percent of U.S. greenhouse gas emissions.
Social justice must be made core to the way engineers are trained. Some universities are working on this. . . . Engineers taught this way will be prepared to think critically about what problems we choose to solve, how we do so responsibly and how we build teams that challenge our ways of thinking.
Individual engineering professors are also working to embed societal needs in their pedagogy. Darshan Karwat at the University of Arizona developed activist engineering to challenge engineers to acknowledge their full moral and social responsibility through practical self-reflection. Khalid Kadir at the University of California, Berkeley, created the popular course Engineering, Environment, and Society that teaches engineers how to engage in place-based knowledge, an understanding of the people, context and history, to design better technical approaches in collaboration with communities. When we design and build with equity and justice in mind, we craft better solutions that respond to the complexities of entrenched systemic problems.

\end{problem}

\newpage
\begin{problem}{}{}
% Subject: varc
\textbf{Instructions}

Instructions   

**The passage below is accompanied by a set of questions. Choose the best answer to each question.**
[Octopuses are] misfits in their own extended families . . . They belong to the Mollusca class Cephalopoda. But they don’t look like their cousins at all. Other molluscs include sea snails, sea slugs, bivalves - most are shelled invertebrates with a dorsal foot. Cephalopods are all arms, and can be as tiny as 1 centimetre and as large at 30 feet. Some of them have brains the size of a walnut, which is large for an invertebrate. . . .
It makes sense for these molluscs to have added protection in the form of a higher cognition; they don’t have a shell covering them, and pretty much everything feeds on cephalopods, including humans. But how did cephalopods manage to secure their own invisibility cloak? Cephalopods fire from multiple cylinders to achieve this in varying degrees from species to species. There are four main catalysts - chromatophores, iridophores, papillae and leucophores. . . .
[Chromatophores] are organs on their bodies that contain pigment sacs, which have red, yellow and brown pigment granules. These sacs have a network of radial muscles, meaning muscles arranged in a circle radiating outwards. These are connected to the brain by a nerve. When the cephalopod wants to change colour, the brain carries an electrical impulse through the nerve to the muscles that expand outwards, pulling open the sacs to display the colours on the skin. Why these three colours? Because these are the colours the light reflects at the depths they live in (the rest is absorbed before it reaches those depths). . . .
Well, what about other colours? Cue the iridophores. Think of a second level of skin that has thin stacks of cells. These can reflect light back at different wavelengths. . . . It’s using the same properties that we’ve seen in hologram stickers, or rainbows on puddles of oil. You move your head and you see a different colour. The sticker isn’t doing anything but reflecting light - it’s your movement that’s changing the appearance of the colour. This property of holograms, oil and other such surfaces is called “iridescence”. . . .
Papillae are sections of the skin that can be deformed to make a texture bumpy. Even humans possess them (goosebumps) but cannot use them in the manner that cephalopods can. For instance, the use of these cells is how an octopus can wrap itself over a rock and appear jagged or how a squid or cuttlefish can imitate the look of a coral reef by growing miniature towers on its skin. It actually matches the texture of the substrate it chooses.
Finally, the leucophores: According to a paper, published in Nature, cuttlefish and octopuses possess an additional type of reflector cell called a leucophore. They are cells that scatter full spectrum light so that they appear white in a similar way that a polar bear’s fur appears white. Leucophores will also reflect any filtered light shown on them . . . If the water appears blue at a certain depth, the octopuses and cuttlefish can appear blue; if the water appears green, they appear green, and so on and so forth.

\vspace{1em}
\textbf{Question 5}

We can infer that the author would approve of a more evolved engineering pedagogy that includes all of the following EXCEPT:

\begin{enumerate}[label=\Alph*.]
    \item making considerations of environmental sustainability intrinsic to the development of technological solutions.
    \item a more responsible approach to technical design and problem-solving than a focus on speed in developing and bringing to scale.
    \item design that is based on the needs of communities using local knowledge and responding to local priorities.
    \item moving towards technical-social dualism where social community needs are incorporated in problem-definition and solutions.
[](https://cracku.in/5-we-can-infer-that-the-author-would-approve-of-a-mo-x-cat-2022-slot-2)
\end{enumerate}
\vspace{1em}
\textbf{Solution:}

Based on the passage, it is likely that the author would approve of all of the options except for Option D, which is moving towards technical-social dualism. Technical-social dualism is described in the passage as the idea that the technical and social dimensions of engineering problems are readily separable and remain distinct throughout the problem-definition and solution process. The passage criticizes this approach, arguing that it ignores the social dimensions of engineering problems and leads to a focus on technical ideals such as cost and efficiency at the expense of broader societal concerns. Therefore, it is unlikely that the author would approve of moving towards technical-social dualism.
Option A is likely to be included because the passage mentions the need for engineers to be aware of the potential impacts of their work on different groups of people, including the environment. It is suggested that ignoring these factors can result in technologies that are not sustainable, which can contribute to a crisis of medical waste and health care, accounting for 10% of U.S. greenhouse gas emissions.
Option B is likely to be included because the passage discusses the consequences of ignoring social dimensions in engineering, such as physical failures and the perpetuation of social inequities. It suggests that a more responsible approach to technical design and problem-solving would consider the full range of stakeholders and the potential impacts of a technology on different groups of people.
Option C is likely to be included because the passage emphasizes the importance of considering social justice in engineering education and practice. It mentions courses focusing on place-based knowledge and community engagement as examples of efforts to incorporate social justice into engineering education. Such an approach would involve designing technologies that are responsive to the needs of communities, using local knowledge and taking into account local priorities.
Hence, Option D is the correct choice.

\end{problem}

\newpage
\begin{problem}{}{}
% Subject: varc
\textbf{Instructions}

Instructions   

**The passage below is accompanied by a set of questions. Choose the best answer to each question.**
[Octopuses are] misfits in their own extended families . . . They belong to the Mollusca class Cephalopoda. But they don’t look like their cousins at all. Other molluscs include sea snails, sea slugs, bivalves - most are shelled invertebrates with a dorsal foot. Cephalopods are all arms, and can be as tiny as 1 centimetre and as large at 30 feet. Some of them have brains the size of a walnut, which is large for an invertebrate. . . .
It makes sense for these molluscs to have added protection in the form of a higher cognition; they don’t have a shell covering them, and pretty much everything feeds on cephalopods, including humans. But how did cephalopods manage to secure their own invisibility cloak? Cephalopods fire from multiple cylinders to achieve this in varying degrees from species to species. There are four main catalysts - chromatophores, iridophores, papillae and leucophores. . . .
[Chromatophores] are organs on their bodies that contain pigment sacs, which have red, yellow and brown pigment granules. These sacs have a network of radial muscles, meaning muscles arranged in a circle radiating outwards. These are connected to the brain by a nerve. When the cephalopod wants to change colour, the brain carries an electrical impulse through the nerve to the muscles that expand outwards, pulling open the sacs to display the colours on the skin. Why these three colours? Because these are the colours the light reflects at the depths they live in (the rest is absorbed before it reaches those depths). . . .
Well, what about other colours? Cue the iridophores. Think of a second level of skin that has thin stacks of cells. These can reflect light back at different wavelengths. . . . It’s using the same properties that we’ve seen in hologram stickers, or rainbows on puddles of oil. You move your head and you see a different colour. The sticker isn’t doing anything but reflecting light - it’s your movement that’s changing the appearance of the colour. This property of holograms, oil and other such surfaces is called “iridescence”. . . .
Papillae are sections of the skin that can be deformed to make a texture bumpy. Even humans possess them (goosebumps) but cannot use them in the manner that cephalopods can. For instance, the use of these cells is how an octopus can wrap itself over a rock and appear jagged or how a squid or cuttlefish can imitate the look of a coral reef by growing miniature towers on its skin. It actually matches the texture of the substrate it chooses.
Finally, the leucophores: According to a paper, published in Nature, cuttlefish and octopuses possess an additional type of reflector cell called a leucophore. They are cells that scatter full spectrum light so that they appear white in a similar way that a polar bear’s fur appears white. Leucophores will also reflect any filtered light shown on them . . . If the water appears blue at a certain depth, the octopuses and cuttlefish can appear blue; if the water appears green, they appear green, and so on and so forth.

\vspace{1em}
\textbf{Question 6}

All of the following are examples of the negative outcomes of focusing on technical ideals in the medical sphere EXCEPT the:

\begin{enumerate}[label=\Alph*.]
    \item neglect of research and development of medical technologies for the diagnosis and treatment of diseases that typically afflict marginalised communities.
    \item continuing calibration of medical devices based on past racial biases that have remained unadjusted for changes.
    \item exclusion of non-privileged groups in clinical trials which leads to incorrect drug dosages.
    \item incorrect assignment of people as female at birth which has resulted in faulty drug interventions.
[](https://cracku.in/6-all-of-the-following-are-examples-of-the-negative--x-cat-2022-slot-2)
\end{enumerate}
\vspace{1em}
\textbf{Solution:}

Option D is not mentioned in the passage as a negative outcome of focusing on technical ideals in the medical sphere. The passage specifically mentions that "most FDA-approved drugs are incorrectly dosed for people assigned female at birth, leading to unexpected adverse reactions. This is because they have been inadequately represented in clinical trials," which suggests that the incorrect dosing of drugs for people assigned female at birth is a consequence of inadequate representation in clinical trials, rather than a result of focusing on technical ideals.
The other options are all mentioned in the passage as negative outcomes of focusing on technical ideals in the medical sphere. Option A refers to the passage's mention of the lack of technologies for "systemically marginalized people" such as those with endometriosis. Option B relates to the discussion on spirometers that have correction factors that assume smaller lung capacity in Black and Asian individuals based on research by eugenicists. Option C is tied to the discussion on "most FDA-approved drugs" being incorrectly dosed for people assigned female at birth due to inadequate representation in clinical trials.
Hence, Option D is the correct choice.

\end{problem}

\newpage
\begin{problem}{}{}
% Subject: varc
\textbf{Instructions}

Instructions   

**The passage below is accompanied by a set of questions. Choose the best answer to each question.**
[Octopuses are] misfits in their own extended families . . . They belong to the Mollusca class Cephalopoda. But they don’t look like their cousins at all. Other molluscs include sea snails, sea slugs, bivalves - most are shelled invertebrates with a dorsal foot. Cephalopods are all arms, and can be as tiny as 1 centimetre and as large at 30 feet. Some of them have brains the size of a walnut, which is large for an invertebrate. . . .
It makes sense for these molluscs to have added protection in the form of a higher cognition; they don’t have a shell covering them, and pretty much everything feeds on cephalopods, including humans. But how did cephalopods manage to secure their own invisibility cloak? Cephalopods fire from multiple cylinders to achieve this in varying degrees from species to species. There are four main catalysts - chromatophores, iridophores, papillae and leucophores. . . .
[Chromatophores] are organs on their bodies that contain pigment sacs, which have red, yellow and brown pigment granules. These sacs have a network of radial muscles, meaning muscles arranged in a circle radiating outwards. These are connected to the brain by a nerve. When the cephalopod wants to change colour, the brain carries an electrical impulse through the nerve to the muscles that expand outwards, pulling open the sacs to display the colours on the skin. Why these three colours? Because these are the colours the light reflects at the depths they live in (the rest is absorbed before it reaches those depths). . . .
Well, what about other colours? Cue the iridophores. Think of a second level of skin that has thin stacks of cells. These can reflect light back at different wavelengths. . . . It’s using the same properties that we’ve seen in hologram stickers, or rainbows on puddles of oil. You move your head and you see a different colour. The sticker isn’t doing anything but reflecting light - it’s your movement that’s changing the appearance of the colour. This property of holograms, oil and other such surfaces is called “iridescence”. . . .
Papillae are sections of the skin that can be deformed to make a texture bumpy. Even humans possess them (goosebumps) but cannot use them in the manner that cephalopods can. For instance, the use of these cells is how an octopus can wrap itself over a rock and appear jagged or how a squid or cuttlefish can imitate the look of a coral reef by growing miniature towers on its skin. It actually matches the texture of the substrate it chooses.
Finally, the leucophores: According to a paper, published in Nature, cuttlefish and octopuses possess an additional type of reflector cell called a leucophore. They are cells that scatter full spectrum light so that they appear white in a similar way that a polar bear’s fur appears white. Leucophores will also reflect any filtered light shown on them . . . If the water appears blue at a certain depth, the octopuses and cuttlefish can appear blue; if the water appears green, they appear green, and so on and so forth.

\vspace{1em}
\textbf{Question 7}

In this passage, the author is making the claim that:

\begin{enumerate}[label=\Alph*.]
    \item the objective of best solutions in engineering has shifted the focus of pedagogy from humanism and social obligations to technological perfection.
    \item engineering students today are taught to focus on objective technical outcomes, independent of the social dimensions of their work.
    \item engineering students today are trained to be non-subjective in their reasoning as this best enables them to develop much-needed universal solutions.
    \item technical-social dualism has emerged as a technique for engineering students to incorporate social considerations into their technical problem-solving processes.
[](https://cracku.in/7-in-this-passage-the-author-is-making-the-claim-tha-x-cat-2022-slot-2)
\end{enumerate}
\vspace{1em}
\textbf{Solution:}

The passage discusses the concept of technical-social dualism, which separates the technical and social dimensions of engineering problems and can result in a focus on technical ideals such as cost and efficiency at the expense of broader societal concerns. The passage states that this way of teaching "primes students to believe that their decision-making is purely objective, as it is grounded in math and science" and that "nontechnical parameters such as access to a technology, cultural relevancy or potential harms are deemed political and invalid." These statements support the claim that engineering students are taught to focus on objective technical outcomes, independent of the social dimensions of their work. Option B best captures the above understanding. The other answer choices do not accurately reflect the content or the main argument of the passage.

\end{problem}

\newpage
\begin{problem}{}{}
% Subject: varc
\textbf{Instructions}

Instructions   

**The passage below is accompanied by a set of questions. Choose the best answer to each question.**
[Octopuses are] misfits in their own extended families . . . They belong to the Mollusca class Cephalopoda. But they don’t look like their cousins at all. Other molluscs include sea snails, sea slugs, bivalves - most are shelled invertebrates with a dorsal foot. Cephalopods are all arms, and can be as tiny as 1 centimetre and as large at 30 feet. Some of them have brains the size of a walnut, which is large for an invertebrate. . . .
It makes sense for these molluscs to have added protection in the form of a higher cognition; they don’t have a shell covering them, and pretty much everything feeds on cephalopods, including humans. But how did cephalopods manage to secure their own invisibility cloak? Cephalopods fire from multiple cylinders to achieve this in varying degrees from species to species. There are four main catalysts - chromatophores, iridophores, papillae and leucophores. . . .
[Chromatophores] are organs on their bodies that contain pigment sacs, which have red, yellow and brown pigment granules. These sacs have a network of radial muscles, meaning muscles arranged in a circle radiating outwards. These are connected to the brain by a nerve. When the cephalopod wants to change colour, the brain carries an electrical impulse through the nerve to the muscles that expand outwards, pulling open the sacs to display the colours on the skin. Why these three colours? Because these are the colours the light reflects at the depths they live in (the rest is absorbed before it reaches those depths). . . .
Well, what about other colours? Cue the iridophores. Think of a second level of skin that has thin stacks of cells. These can reflect light back at different wavelengths. . . . It’s using the same properties that we’ve seen in hologram stickers, or rainbows on puddles of oil. You move your head and you see a different colour. The sticker isn’t doing anything but reflecting light - it’s your movement that’s changing the appearance of the colour. This property of holograms, oil and other such surfaces is called “iridescence”. . . .
Papillae are sections of the skin that can be deformed to make a texture bumpy. Even humans possess them (goosebumps) but cannot use them in the manner that cephalopods can. For instance, the use of these cells is how an octopus can wrap itself over a rock and appear jagged or how a squid or cuttlefish can imitate the look of a coral reef by growing miniature towers on its skin. It actually matches the texture of the substrate it chooses.
Finally, the leucophores: According to a paper, published in Nature, cuttlefish and octopuses possess an additional type of reflector cell called a leucophore. They are cells that scatter full spectrum light so that they appear white in a similar way that a polar bear’s fur appears white. Leucophores will also reflect any filtered light shown on them . . . If the water appears blue at a certain depth, the octopuses and cuttlefish can appear blue; if the water appears green, they appear green, and so on and so forth.

\vspace{1em}
\textbf{Question 8}

The author gives all of the following reasons for why marginalised people are systematically discriminated against in technology-related interventions EXCEPT:

\begin{enumerate}[label=\Alph*.]
    \item “But those technical ideals are at their core social and political choices determined by a dominant culture focused on economic growth for the most privileged segments of society.”
    \item “And we hardly question whether devices are built sustainably, which has led to a crisis of medical waste and health care accounting for 10 percent of U.S. greenhouse gas emissions.”
    \item “These racially based adjustments are derived from research done by eugenicists who thought these racial differences were biologically determined and who considered nonwhite people as inferior.”
    \item “Beyond physical failings, subjective beliefs treated as facts by those in decision-making roles can encode social inequities.”
[](https://cracku.in/8-the-author-gives-all-of-the-following-reasons-for--x-cat-2022-slot-2)
\end{enumerate}
\vspace{1em}
\textbf{Solution:}

_Option A_ : [**Correct**] The passage states that technical ideals, such as cost and efficiency, are often determined by a dominant culture that prioritizes economic growth for the most privileged segments of society. This can result in technologies and interventions that are not designed with the needs and concerns of marginalized groups in mind, leading to systemic discrimination against these groups.
_Option B_ : [**Incorrect**] The passage does not mention anything about sustainability or medical waste contributing to greenhouse gas emissions as a reason for the systematic discrimination of marginalized people in technology-related interventions.
_Option C_ : [**Correct**] The passage mentions that certain technologies, such as spirometers, have correction factors that assume smaller lung capacity in Black and Asian individuals based on research by eugenicists who believed in racial hierarchies and considered nonwhite people as inferior. This is an example of how subjective beliefs can be treated as facts and encoded into technologies, leading to social inequities.
_Option D_ : [**Correct**] The passage discusses how subjective beliefs treated as facts by those in decision-making roles can result in physical failures, such as incorrect dosing of drugs for people assigned female at birth, and also encode social inequities, such as the correction factors on spirometers that assume smaller lung capacity in Black and Asian individuals.
Hence, Option B is the correct choice.

Instructions   

**The passage below is accompanied by a set of questions. Choose the best answer to each question.**
We begin with the emergence of the philosophy of the social sciences as an arena of thought and as a set of social institutions. The two characterisations overlap but are not congruent. Academic disciplines are social institutions. . . . My view is that institutions are all those social entities that organise action: they link acting individuals into social structures. There are various kinds of institutions. Hegelians and Marxists emphasise universal institutions such as the family, rituals, governance, economy and the military. These are mostly institutions that just grew. Perhaps in some imaginary beginning of time they spontaneously appeared. In their present incarnations, however, they are very much the product of conscious attempts to mould and plan them. We have family law, established and disestablished churches, constitutions and laws, including those governing the economy and the military. Institutions deriving from statute, like joint-stock companies are formal by contrast with informal ones such as friendships. There are some institutions that come in both informal and formal variants, as well as in mixed ones. Consider the fact that the stock exchange and the black market are both market institutions, one formal one not. Consider further that there are many features of the work of the stock exchange that rely on informal, noncodifiable agreements, not least the language used for communication. To be precise, mixtures are the norm . . . From constitutions at the top to by-laws near the bottom we are always adding to, or tinkering with, earlier institutions, the grown and the designed are intertwined.
It is usual in social thought to treat culture and tradition as different from, although alongside, institutions. The view taken here is different. Culture and tradition are sub-sets of institutions analytically isolated for explanatory or expository purposes. Some social scientists have taken all institutions, even purely local ones, to be entities that satisfy basic human needs - under local conditions . . . Others differed and declared any structure of reciprocal roles and norms an institution. Most of these differences are differences of emphasis rather than disagreements. Let us straddle all these versions and present institutions very generally . . . as structures that serve to coordinate the actions of individuals. . . . Institutions themselves then have no aims or purpose other than those given to them by actors or used by actors to explain them . . .
Language is the formative institution for social life and for science . . . Both formal and informal language is involved, naturally grown or designed. (Language is all of these to varying degrees.) Languages are paradigms of institutions or, from another perspective, nested sets of institutions. Syntax, semantics, lexicon and alphabet/character-set are all institutions within the larger institutional framework of a written language. Natural languages are typical examples of what Ferguson called ‘the result of human action, but not the execution of any human design’[;] reformed natural languages and artificial languages introduce design into their modifications or refinements of natural language. Above all, languages are paradigms of institutional tools that function to coordinate.

\end{problem}

\newpage
\begin{problem}{}{}
% Subject: varc
\textbf{Instructions}

Instructions   

**The passage below is accompanied by a set of questions. Choose the best answer to each question.**
[Octopuses are] misfits in their own extended families . . . They belong to the Mollusca class Cephalopoda. But they don’t look like their cousins at all. Other molluscs include sea snails, sea slugs, bivalves - most are shelled invertebrates with a dorsal foot. Cephalopods are all arms, and can be as tiny as 1 centimetre and as large at 30 feet. Some of them have brains the size of a walnut, which is large for an invertebrate. . . .
It makes sense for these molluscs to have added protection in the form of a higher cognition; they don’t have a shell covering them, and pretty much everything feeds on cephalopods, including humans. But how did cephalopods manage to secure their own invisibility cloak? Cephalopods fire from multiple cylinders to achieve this in varying degrees from species to species. There are four main catalysts - chromatophores, iridophores, papillae and leucophores. . . .
[Chromatophores] are organs on their bodies that contain pigment sacs, which have red, yellow and brown pigment granules. These sacs have a network of radial muscles, meaning muscles arranged in a circle radiating outwards. These are connected to the brain by a nerve. When the cephalopod wants to change colour, the brain carries an electrical impulse through the nerve to the muscles that expand outwards, pulling open the sacs to display the colours on the skin. Why these three colours? Because these are the colours the light reflects at the depths they live in (the rest is absorbed before it reaches those depths). . . .
Well, what about other colours? Cue the iridophores. Think of a second level of skin that has thin stacks of cells. These can reflect light back at different wavelengths. . . . It’s using the same properties that we’ve seen in hologram stickers, or rainbows on puddles of oil. You move your head and you see a different colour. The sticker isn’t doing anything but reflecting light - it’s your movement that’s changing the appearance of the colour. This property of holograms, oil and other such surfaces is called “iridescence”. . . .
Papillae are sections of the skin that can be deformed to make a texture bumpy. Even humans possess them (goosebumps) but cannot use them in the manner that cephalopods can. For instance, the use of these cells is how an octopus can wrap itself over a rock and appear jagged or how a squid or cuttlefish can imitate the look of a coral reef by growing miniature towers on its skin. It actually matches the texture of the substrate it chooses.
Finally, the leucophores: According to a paper, published in Nature, cuttlefish and octopuses possess an additional type of reflector cell called a leucophore. They are cells that scatter full spectrum light so that they appear white in a similar way that a polar bear’s fur appears white. Leucophores will also reflect any filtered light shown on them . . . If the water appears blue at a certain depth, the octopuses and cuttlefish can appear blue; if the water appears green, they appear green, and so on and so forth.

\vspace{1em}
\textbf{Question 9}

“Consider the fact that the stock exchange and the black market are both market institutions, one formal one not.” Which one of the following statements best explains this quote, in the context of the passage?

\begin{enumerate}[label=\Alph*.]
    \item Market instruments can be formally traded in the stock exchange and informally traded in the black market.
    \item The stock exchange and the black market are both organised to function by rules.
    \item The stock exchange and the black market are examples of how, even within the same domain, different kinds of institutions can co-exist.
    \item The stock exchange and the black market are both dependent on the market to survive.
[](https://cracku.in/9-consider-the-fact-that-the-stock-exchange-and-the--x-cat-2022-slot-2)
\end{enumerate}
\vspace{1em}
\textbf{Solution:}

The author states that the stock exchange and the black market are both market institutions, but notes that the stock exchange is formal, while the black market is not. This suggests that the two institutions differ in terms of how they are structured and operate, and the author suggests that this difference is an example of how different kinds of institutions can co-exist within the same domain. Option C aptly presents this idea. Contrarily, the ideas mentioned in Options A, B and D are not supported by the information in the passage. 
Hence, Option C is the correct choice.

\end{problem}

\newpage
\begin{problem}{}{}
% Subject: varc
\textbf{Instructions}

Instructions   

**The passage below is accompanied by a set of questions. Choose the best answer to each question.**
[Octopuses are] misfits in their own extended families . . . They belong to the Mollusca class Cephalopoda. But they don’t look like their cousins at all. Other molluscs include sea snails, sea slugs, bivalves - most are shelled invertebrates with a dorsal foot. Cephalopods are all arms, and can be as tiny as 1 centimetre and as large at 30 feet. Some of them have brains the size of a walnut, which is large for an invertebrate. . . .
It makes sense for these molluscs to have added protection in the form of a higher cognition; they don’t have a shell covering them, and pretty much everything feeds on cephalopods, including humans. But how did cephalopods manage to secure their own invisibility cloak? Cephalopods fire from multiple cylinders to achieve this in varying degrees from species to species. There are four main catalysts - chromatophores, iridophores, papillae and leucophores. . . .
[Chromatophores] are organs on their bodies that contain pigment sacs, which have red, yellow and brown pigment granules. These sacs have a network of radial muscles, meaning muscles arranged in a circle radiating outwards. These are connected to the brain by a nerve. When the cephalopod wants to change colour, the brain carries an electrical impulse through the nerve to the muscles that expand outwards, pulling open the sacs to display the colours on the skin. Why these three colours? Because these are the colours the light reflects at the depths they live in (the rest is absorbed before it reaches those depths). . . .
Well, what about other colours? Cue the iridophores. Think of a second level of skin that has thin stacks of cells. These can reflect light back at different wavelengths. . . . It’s using the same properties that we’ve seen in hologram stickers, or rainbows on puddles of oil. You move your head and you see a different colour. The sticker isn’t doing anything but reflecting light - it’s your movement that’s changing the appearance of the colour. This property of holograms, oil and other such surfaces is called “iridescence”. . . .
Papillae are sections of the skin that can be deformed to make a texture bumpy. Even humans possess them (goosebumps) but cannot use them in the manner that cephalopods can. For instance, the use of these cells is how an octopus can wrap itself over a rock and appear jagged or how a squid or cuttlefish can imitate the look of a coral reef by growing miniature towers on its skin. It actually matches the texture of the substrate it chooses.
Finally, the leucophores: According to a paper, published in Nature, cuttlefish and octopuses possess an additional type of reflector cell called a leucophore. They are cells that scatter full spectrum light so that they appear white in a similar way that a polar bear’s fur appears white. Leucophores will also reflect any filtered light shown on them . . . If the water appears blue at a certain depth, the octopuses and cuttlefish can appear blue; if the water appears green, they appear green, and so on and so forth.

\vspace{1em}
\textbf{Question 10}

All of the following inferences from the passage are false, EXCEPT:

\begin{enumerate}[label=\Alph*.]
    \item as concepts, “culture” and “tradition” have no analytical, explanatory or expository power, especially when they are treated in isolation.
    \item the institution of friendship cannot be found in the institution of joint-stock companies because the first is an informal institution, while the second is a formal one.
    \item institutions like the family, rituals, governance, economy, and the military are natural and cannot be consciously modified.
    \item “natural language” refers to that stage of language development where no conscious human intent is evident in the formation of language.
[](https://cracku.in/10-all-of-the-following-inferences-from-the-passage-a-x-cat-2022-slot-2)
\end{enumerate}
\vspace{1em}
\textbf{Solution:}

Based on the passage, the correct inference is: Option D - **"natural language" refers to that stage of language development where no conscious human intent is evident in the formation of language.**
The passage states that natural languages are "the result of human action, but not the execution of any human design," suggesting that they are not consciously designed or modified. Artificial languages and reformed natural languages, on the other hand, are described as introducing "design into their modifications or refinements," indicating that they are consciously modified.
The passage does not support the other inferences. In fact, it suggests that culture and tradition can be understood as subsets of institutions and have analytical, explanatory, and expository power when they are studied in this context. It also suggests that there are both informal and formal institutions, and that many institutions are a mixture of the two, rather than being mutually exclusive categories. The passage also does not imply that institutions like the family, rituals, governance, economy, and the military are natural and cannot be consciously modified; rather, it suggests that these institutions are the product of conscious attempts to mold and plan them, and that they can be modified through processes such as family law, established and disestablished churches, constitutions and laws, and so on [Options A, B and C].

\end{problem}

\newpage
\begin{problem}{}{}
% Subject: varc
\textbf{Instructions}

Instructions   

**The passage below is accompanied by a set of questions. Choose the best answer to each question.**
[Octopuses are] misfits in their own extended families . . . They belong to the Mollusca class Cephalopoda. But they don’t look like their cousins at all. Other molluscs include sea snails, sea slugs, bivalves - most are shelled invertebrates with a dorsal foot. Cephalopods are all arms, and can be as tiny as 1 centimetre and as large at 30 feet. Some of them have brains the size of a walnut, which is large for an invertebrate. . . .
It makes sense for these molluscs to have added protection in the form of a higher cognition; they don’t have a shell covering them, and pretty much everything feeds on cephalopods, including humans. But how did cephalopods manage to secure their own invisibility cloak? Cephalopods fire from multiple cylinders to achieve this in varying degrees from species to species. There are four main catalysts - chromatophores, iridophores, papillae and leucophores. . . .
[Chromatophores] are organs on their bodies that contain pigment sacs, which have red, yellow and brown pigment granules. These sacs have a network of radial muscles, meaning muscles arranged in a circle radiating outwards. These are connected to the brain by a nerve. When the cephalopod wants to change colour, the brain carries an electrical impulse through the nerve to the muscles that expand outwards, pulling open the sacs to display the colours on the skin. Why these three colours? Because these are the colours the light reflects at the depths they live in (the rest is absorbed before it reaches those depths). . . .
Well, what about other colours? Cue the iridophores. Think of a second level of skin that has thin stacks of cells. These can reflect light back at different wavelengths. . . . It’s using the same properties that we’ve seen in hologram stickers, or rainbows on puddles of oil. You move your head and you see a different colour. The sticker isn’t doing anything but reflecting light - it’s your movement that’s changing the appearance of the colour. This property of holograms, oil and other such surfaces is called “iridescence”. . . .
Papillae are sections of the skin that can be deformed to make a texture bumpy. Even humans possess them (goosebumps) but cannot use them in the manner that cephalopods can. For instance, the use of these cells is how an octopus can wrap itself over a rock and appear jagged or how a squid or cuttlefish can imitate the look of a coral reef by growing miniature towers on its skin. It actually matches the texture of the substrate it chooses.
Finally, the leucophores: According to a paper, published in Nature, cuttlefish and octopuses possess an additional type of reflector cell called a leucophore. They are cells that scatter full spectrum light so that they appear white in a similar way that a polar bear’s fur appears white. Leucophores will also reflect any filtered light shown on them . . . If the water appears blue at a certain depth, the octopuses and cuttlefish can appear blue; if the water appears green, they appear green, and so on and so forth.

\vspace{1em}
\textbf{Question 11}

In the first paragraph of the passage, what are the two “characterisations” that are seen as overlapping but not congruent?

\begin{enumerate}[label=\Alph*.]
    \item “an arena of thought” and “academic disciplines".
    \item “individuals” and “social structures”.
    \item “the philosophy of the social sciences” and “a set of social institutions”.
    \item “academic disciplines” and “institutions”.
[](https://cracku.in/11-in-the-first-paragraph-of-the-passage-what-are-the-x-cat-2022-slot-2)
\end{enumerate}
\vspace{1em}
\textbf{Solution:}

{_We begin with the emergence of the philosophy of the social sciences as an arena of thought and as a set of social institutions. The two characterisations overlap but are not congruent. Academic disciplines are social institutions. . ._}
In the first paragraph of the passage, the two "characterisations" that are seen as overlapping but not congruent are the philosophy of the social sciences as an arena of thought [academic discipline] and as a set of social institutions. The author suggests that these two characterizations overlap in the sense that they both involve the study of social phenomena, but they are not congruent because the philosophy of the social sciences is a field of study that encompasses a broad range of ideas and theories, while social institutions are specific structures that organize and coordinate social action. Although they are related, they are not identical and can be understood as distinct but overlapping aspects of the social sciences. Option D correctly represents the above.

\end{problem}

\newpage
\begin{problem}{}{}
% Subject: varc
\textbf{Instructions}

Instructions   

**The passage below is accompanied by a set of questions. Choose the best answer to each question.**
[Octopuses are] misfits in their own extended families . . . They belong to the Mollusca class Cephalopoda. But they don’t look like their cousins at all. Other molluscs include sea snails, sea slugs, bivalves - most are shelled invertebrates with a dorsal foot. Cephalopods are all arms, and can be as tiny as 1 centimetre and as large at 30 feet. Some of them have brains the size of a walnut, which is large for an invertebrate. . . .
It makes sense for these molluscs to have added protection in the form of a higher cognition; they don’t have a shell covering them, and pretty much everything feeds on cephalopods, including humans. But how did cephalopods manage to secure their own invisibility cloak? Cephalopods fire from multiple cylinders to achieve this in varying degrees from species to species. There are four main catalysts - chromatophores, iridophores, papillae and leucophores. . . .
[Chromatophores] are organs on their bodies that contain pigment sacs, which have red, yellow and brown pigment granules. These sacs have a network of radial muscles, meaning muscles arranged in a circle radiating outwards. These are connected to the brain by a nerve. When the cephalopod wants to change colour, the brain carries an electrical impulse through the nerve to the muscles that expand outwards, pulling open the sacs to display the colours on the skin. Why these three colours? Because these are the colours the light reflects at the depths they live in (the rest is absorbed before it reaches those depths). . . .
Well, what about other colours? Cue the iridophores. Think of a second level of skin that has thin stacks of cells. These can reflect light back at different wavelengths. . . . It’s using the same properties that we’ve seen in hologram stickers, or rainbows on puddles of oil. You move your head and you see a different colour. The sticker isn’t doing anything but reflecting light - it’s your movement that’s changing the appearance of the colour. This property of holograms, oil and other such surfaces is called “iridescence”. . . .
Papillae are sections of the skin that can be deformed to make a texture bumpy. Even humans possess them (goosebumps) but cannot use them in the manner that cephalopods can. For instance, the use of these cells is how an octopus can wrap itself over a rock and appear jagged or how a squid or cuttlefish can imitate the look of a coral reef by growing miniature towers on its skin. It actually matches the texture of the substrate it chooses.
Finally, the leucophores: According to a paper, published in Nature, cuttlefish and octopuses possess an additional type of reflector cell called a leucophore. They are cells that scatter full spectrum light so that they appear white in a similar way that a polar bear’s fur appears white. Leucophores will also reflect any filtered light shown on them . . . If the water appears blue at a certain depth, the octopuses and cuttlefish can appear blue; if the water appears green, they appear green, and so on and so forth.

\vspace{1em}
\textbf{Question 12}

Which of the following statements best represents the essence of the passage?

\begin{enumerate}[label=\Alph*.]
    \item It is usual in social thought to treat culture and tradition as different from institutions.
    \item Language is the fundamental formal institution for social life and for science.
    \item The stock exchange and the black market are both market institutions.
    \item Institutions are structures that serve to coordinate the actions of individuals.
[](https://cracku.in/12-which-of-the-following-statements-best-represents--x-cat-2022-slot-2)
\end{enumerate}
\vspace{1em}
\textbf{Solution:}

The passage discusses the concept of institutions within the philosophy of the social sciences. Institutions are seen as structures that coordinate the actions of individuals and can be either formal or informal; in this regard, the author discusses various types of institutions, including universal institutions, formal institutions, and informal institutions. The author also argues that culture and tradition can be seen as subsets of institutions, and that language is a particularly important institution that shapes social life and science. The author notes that there can be differences in the way that institutions are understood and emphasized by different social scientists but that these differences are often matters of emphasis rather than fundamental disagreements. Primarily, the passage underlines that these institutions are structures that coordinate the actions of individuals. Option D best captures the essence of the abovementioned elements. 
_Option A_ is incorrect because it only partially reflects what the author says about culture and tradition. While the author does discuss culture and tradition in relation to institutions, they are not treated as being completely separate from institutions. Instead, the author argues that culture and tradition can be seen as subsets of institutions, meaning they are part of institutions rather than completely separate from them.  _Option B_ exaggerates the emphasis the author places on language as an institution. While the author does discuss language as an important institution, he does not suggest that it is the only or even the most fundamental institution.  _Option C_ , while true, does not capture the main focus of the passage, which is on the concept of institutions in general rather than on specific examples of institutions.

Instructions   

**The passage below is accompanied by a set of questions. Choose the best answer to each question.**
Humans today make music. Think beyond all the qualifications that might trail after this bald statement: that only certain humans make music, that extensive training is involved, that many societies distinguish musical specialists from nonmusicians, that in today’s societies most listen to music rather than making it, and so forth. These qualifications, whatever their local merit, are moot in the face of the overarching truth that making music, considered from a cognitive and psychological vantage, is the province of all those who perceive and experience what is made. We are, almost all of us, musicians — everyone who can entrain (not necessarily dance) to a beat, who can recognize a repeated tune (not necessarily sing it), who can distinguish one instrument or one singing voice from another. I will often use an antique word, recently revived, to name this broader musical experience. Humans are musicking creatures. . . .
The set of capacities that enables musicking is a principal marker of modern humanity. There is nothing polemical in this assertion except a certain insistence, which will figure often in what follows, that musicking be included in our thinking about fundamental human commonalities. Capacities involved in musicking are many and take shape in complicated ways, arising from innate dispositions . . . Most of these capacities overlap with nonmusical ones, though a few may be distinct and dedicated to musical perception and production. In the area of overlap, linguistic capacities seem to be particularly important, and humans are (in principle) language-makers in addition to music-makers — speaking creatures as well as musicking ones. 
Humans are symbol-makers too, a feature tightly bound up with language, not so tightly with music. The species Cassirer dubbed Homo symbolicus cannot help but tangle musicking in webs of symbolic thought and expression, habitually making it a component of behavioral complexes that form such expression. But in fundamental features musicking is neither language-like nor symbol-like, and from these differences come many clues to its ancient emergence.
If musicking is a primary, shared trait of modern humans, then to describe its emergence must be to detail the coalescing of that modernity. This took place, archaeologists are clear, over a very long durée: at least 50,000 years or so, more likely something closer to 200,000, depending in part on what that coalescence is taken to comprise. If we look back 20,000 years, a small portion of this long period, we reach the lives of humans whose musical capacities were probably little different from our own. As we look farther back we reach horizons where this similarity can no longer hold — perhaps 40,000 years ago, perhaps 70,000, perhaps 100,000. But we never cross a line before which all the cognitive capacities recruited in modern musicking abruptly disappear. Unless we embrace the incredible notion that music sprang forth in full-blown glory, its emergence will have to be tracked in gradualist terms across a long period.
This is one general feature of a history of music’s emergence . . . The history was at once sociocultural and biological . . . The capacities recruited in musicking are many, so describing its emergence involves following several or many separate strands.

\end{problem}

\newpage
\begin{problem}{}{}
% Subject: varc
\textbf{Instructions}

Instructions   

**The passage below is accompanied by a set of questions. Choose the best answer to each question.**
[Octopuses are] misfits in their own extended families . . . They belong to the Mollusca class Cephalopoda. But they don’t look like their cousins at all. Other molluscs include sea snails, sea slugs, bivalves - most are shelled invertebrates with a dorsal foot. Cephalopods are all arms, and can be as tiny as 1 centimetre and as large at 30 feet. Some of them have brains the size of a walnut, which is large for an invertebrate. . . .
It makes sense for these molluscs to have added protection in the form of a higher cognition; they don’t have a shell covering them, and pretty much everything feeds on cephalopods, including humans. But how did cephalopods manage to secure their own invisibility cloak? Cephalopods fire from multiple cylinders to achieve this in varying degrees from species to species. There are four main catalysts - chromatophores, iridophores, papillae and leucophores. . . .
[Chromatophores] are organs on their bodies that contain pigment sacs, which have red, yellow and brown pigment granules. These sacs have a network of radial muscles, meaning muscles arranged in a circle radiating outwards. These are connected to the brain by a nerve. When the cephalopod wants to change colour, the brain carries an electrical impulse through the nerve to the muscles that expand outwards, pulling open the sacs to display the colours on the skin. Why these three colours? Because these are the colours the light reflects at the depths they live in (the rest is absorbed before it reaches those depths). . . .
Well, what about other colours? Cue the iridophores. Think of a second level of skin that has thin stacks of cells. These can reflect light back at different wavelengths. . . . It’s using the same properties that we’ve seen in hologram stickers, or rainbows on puddles of oil. You move your head and you see a different colour. The sticker isn’t doing anything but reflecting light - it’s your movement that’s changing the appearance of the colour. This property of holograms, oil and other such surfaces is called “iridescence”. . . .
Papillae are sections of the skin that can be deformed to make a texture bumpy. Even humans possess them (goosebumps) but cannot use them in the manner that cephalopods can. For instance, the use of these cells is how an octopus can wrap itself over a rock and appear jagged or how a squid or cuttlefish can imitate the look of a coral reef by growing miniature towers on its skin. It actually matches the texture of the substrate it chooses.
Finally, the leucophores: According to a paper, published in Nature, cuttlefish and octopuses possess an additional type of reflector cell called a leucophore. They are cells that scatter full spectrum light so that they appear white in a similar way that a polar bear’s fur appears white. Leucophores will also reflect any filtered light shown on them . . . If the water appears blue at a certain depth, the octopuses and cuttlefish can appear blue; if the water appears green, they appear green, and so on and so forth.

\vspace{1em}
\textbf{Question 13}

Which one of the following sets of terms best serves as keywords to the passage?

\begin{enumerate}[label=\Alph*.]
    \item Musicking; Cognitive psychology; Antique; Symbol-makers; Modernity.
    \item Humans; Capacities; Language; Symbols; Modernity.
    \item Humans; Musicking; Linguistic capacities; Symbol-making; Modern humanity.
    \item Humans; Psychological vantage; Musicking; Cassirer; Emergence of music.
[](https://cracku.in/13-which-one-of-the-following-sets-of-terms-best-serv-x-cat-2022-slot-2)
\end{enumerate}
\vspace{1em}
\textbf{Solution:}

The passage is about the idea that making music is a fundamental and universal aspect of the human experience, and that it has a long history that is both sociocultural and biological in nature. The author suggests that the capacity for musicking, or making music, is innate in humans and is closely related to other capacities such as language and symbol-making. The author also acknowledges that there are variations in how musical capacities are expressed and developed among different cultures, but maintains that all humans possess these capacities to some extent. The author suggests that the emergence of music can be traced back to at least 50,000 years ago, and that it likely developed gradually over a longer period of time. The author also notes that the emergence of music involved the recruitment of many cognitive capacities, and that understanding this process involves following multiple strands. In this regard, Option C includes all the keywords central to the discussion. 
The other choices do not include all of the relevant terms. Option A includes some of the terms mentioned in the passage, but omits important ones, such as humans and modern humanity. Option B includes humans and capacities but omits important terms such as musicking and symbol-making. Option D includes humans and the emergence of music but omits important terms such as musicking and linguistic capacities.
Hence, Option C is the correct choice.

\end{problem}

\newpage
\begin{problem}{}{}
% Subject: varc
\textbf{Instructions}

Instructions   

**The passage below is accompanied by a set of questions. Choose the best answer to each question.**
[Octopuses are] misfits in their own extended families . . . They belong to the Mollusca class Cephalopoda. But they don’t look like their cousins at all. Other molluscs include sea snails, sea slugs, bivalves - most are shelled invertebrates with a dorsal foot. Cephalopods are all arms, and can be as tiny as 1 centimetre and as large at 30 feet. Some of them have brains the size of a walnut, which is large for an invertebrate. . . .
It makes sense for these molluscs to have added protection in the form of a higher cognition; they don’t have a shell covering them, and pretty much everything feeds on cephalopods, including humans. But how did cephalopods manage to secure their own invisibility cloak? Cephalopods fire from multiple cylinders to achieve this in varying degrees from species to species. There are four main catalysts - chromatophores, iridophores, papillae and leucophores. . . .
[Chromatophores] are organs on their bodies that contain pigment sacs, which have red, yellow and brown pigment granules. These sacs have a network of radial muscles, meaning muscles arranged in a circle radiating outwards. These are connected to the brain by a nerve. When the cephalopod wants to change colour, the brain carries an electrical impulse through the nerve to the muscles that expand outwards, pulling open the sacs to display the colours on the skin. Why these three colours? Because these are the colours the light reflects at the depths they live in (the rest is absorbed before it reaches those depths). . . .
Well, what about other colours? Cue the iridophores. Think of a second level of skin that has thin stacks of cells. These can reflect light back at different wavelengths. . . . It’s using the same properties that we’ve seen in hologram stickers, or rainbows on puddles of oil. You move your head and you see a different colour. The sticker isn’t doing anything but reflecting light - it’s your movement that’s changing the appearance of the colour. This property of holograms, oil and other such surfaces is called “iridescence”. . . .
Papillae are sections of the skin that can be deformed to make a texture bumpy. Even humans possess them (goosebumps) but cannot use them in the manner that cephalopods can. For instance, the use of these cells is how an octopus can wrap itself over a rock and appear jagged or how a squid or cuttlefish can imitate the look of a coral reef by growing miniature towers on its skin. It actually matches the texture of the substrate it chooses.
Finally, the leucophores: According to a paper, published in Nature, cuttlefish and octopuses possess an additional type of reflector cell called a leucophore. They are cells that scatter full spectrum light so that they appear white in a similar way that a polar bear’s fur appears white. Leucophores will also reflect any filtered light shown on them . . . If the water appears blue at a certain depth, the octopuses and cuttlefish can appear blue; if the water appears green, they appear green, and so on and so forth.

\vspace{1em}
\textbf{Question 14}

“Think beyond all the qualifications that might trail after this bald statement . . .” In the context of the passage, what is the author trying to communicate in this quoted extract?

\begin{enumerate}[label=\Alph*.]
    \item A bald statement is one that is trailed by a series of qualifying clarifications and caveats.
    \item A bald statement is one that requires no qualifications to infer its meaning.
    \item Although there may be many caveats and other considerations, the statement is essentially true.
    \item Thinking beyond qualifications allows us to give free reign to musical expressions.
[](https://cracku.in/14-think-beyond-all-the-qualifications-that-might-tra-x-cat-2022-slot-2)
\end{enumerate}
\vspace{1em}
\textbf{Solution:}

In this context, the author is trying to communicate that although there may be various qualifications and considerations that might trail after the statement that "humans today make music," the statement is fundamentally true. The author suggests that almost all humans are musicians to some extent, and that the capacity for making music is a fundamental aspect of the human experience. The author is urging readers to consider this statement without getting bogged down in the various qualifications and considerations that might be attached to it, and to recognize its underlying truth. Option C aptly captures this idea.
_Option A_ : [Incorrect] The phrase "trail after" does not necessarily imply that a bald statement is followed by a series of qualifying clarifications and caveats. Rather, it simply means that something follows after something else. In this case, the author is suggesting that there may be various qualifications and considerations that follow after the statement that "humans today make music," but is not implying that the statement itself is trailed by a series of clarifications and caveats.
_Option B_ : [Incorrect] The phrase "bald statement" does not necessarily imply that a statement requires no qualifications to infer its meaning. Rather, it simply means that the statement is presented in a straightforward and unembellished way. In this case, the author is suggesting that the statement that "humans today make music" is presented in a bald and straightforward way, but is not implying that the statement itself requires no qualifications or considerations.  

_Option D_ : [Incorrect] The phrase "give free reign to" does not accurately describe the author's intention in this context. The author is not suggesting that readers should allow musical expressions to be unrestricted or uncontrolled, but rather that they should consider the statement that "humans today make music" without getting bogged down in various qualifications and considerations.
Hence, Option C is the correct choice.

\end{problem}

\newpage
\begin{problem}{}{}
% Subject: varc
\textbf{Instructions}

Instructions   

**The passage below is accompanied by a set of questions. Choose the best answer to each question.**
[Octopuses are] misfits in their own extended families . . . They belong to the Mollusca class Cephalopoda. But they don’t look like their cousins at all. Other molluscs include sea snails, sea slugs, bivalves - most are shelled invertebrates with a dorsal foot. Cephalopods are all arms, and can be as tiny as 1 centimetre and as large at 30 feet. Some of them have brains the size of a walnut, which is large for an invertebrate. . . .
It makes sense for these molluscs to have added protection in the form of a higher cognition; they don’t have a shell covering them, and pretty much everything feeds on cephalopods, including humans. But how did cephalopods manage to secure their own invisibility cloak? Cephalopods fire from multiple cylinders to achieve this in varying degrees from species to species. There are four main catalysts - chromatophores, iridophores, papillae and leucophores. . . .
[Chromatophores] are organs on their bodies that contain pigment sacs, which have red, yellow and brown pigment granules. These sacs have a network of radial muscles, meaning muscles arranged in a circle radiating outwards. These are connected to the brain by a nerve. When the cephalopod wants to change colour, the brain carries an electrical impulse through the nerve to the muscles that expand outwards, pulling open the sacs to display the colours on the skin. Why these three colours? Because these are the colours the light reflects at the depths they live in (the rest is absorbed before it reaches those depths). . . .
Well, what about other colours? Cue the iridophores. Think of a second level of skin that has thin stacks of cells. These can reflect light back at different wavelengths. . . . It’s using the same properties that we’ve seen in hologram stickers, or rainbows on puddles of oil. You move your head and you see a different colour. The sticker isn’t doing anything but reflecting light - it’s your movement that’s changing the appearance of the colour. This property of holograms, oil and other such surfaces is called “iridescence”. . . .
Papillae are sections of the skin that can be deformed to make a texture bumpy. Even humans possess them (goosebumps) but cannot use them in the manner that cephalopods can. For instance, the use of these cells is how an octopus can wrap itself over a rock and appear jagged or how a squid or cuttlefish can imitate the look of a coral reef by growing miniature towers on its skin. It actually matches the texture of the substrate it chooses.
Finally, the leucophores: According to a paper, published in Nature, cuttlefish and octopuses possess an additional type of reflector cell called a leucophore. They are cells that scatter full spectrum light so that they appear white in a similar way that a polar bear’s fur appears white. Leucophores will also reflect any filtered light shown on them . . . If the water appears blue at a certain depth, the octopuses and cuttlefish can appear blue; if the water appears green, they appear green, and so on and so forth.

\vspace{1em}
\textbf{Question 15}

Based on the passage, which one of the following statements is a valid argument about the emergence of music/musicking?

\begin{enumerate}[label=\Alph*.]
    \item Anyone who can perceive and experience music must be considered capable of musicking.
    \item Although musicking is not language-like, it shares the quality of being a form of expression.
    \item 20,000 years ago, human musical capacities were not very different from what they are today.
    \item All musical work is located in the overlap between linguistic capacity and music production.
[](https://cracku.in/15-based-on-the-passage-which-one-of-the-following-st-x-cat-2022-slot-2)
\end{enumerate}
\vspace{1em}
\textbf{Solution:}

_Option A_ : This is incorrect because it overgeneralizes the ideas presented in the passage. The passage does not state that all humans are musicians or capable of musicking, only that making music is a universal aspect of the human experience and that the capacity for musicking is innate in humans.
_Option B_ : The passage does state that musicking is a form of expression, but there is no discussion supporting the claim "musicking is not language-like." Thus, we can eliminate this choice.
_Option C_ : The passage states - {"..._if we look back 20,000 years, a small portion of this long period, we reach the lives of humans whose musical capacities were probably little different from our own.._."} This suggests that the musical capacities of humans 20,000 years ago were not significantly different from those of humans today.
_Option D_ : This is not supported by the passage. The passage does not state that all musical work is located in the overlap between linguistic capacity and music production. In fact, the passage states that "_most of these capacities overlap with nonmusical ones, though a few may be distinct and dedicated to musical perception and production_ ," suggesting that there are some capacities specifically dedicated to musical perception and production that do not overlap with nonmusical ones.
Hence, the correct answer is Option C.

\end{problem}

\newpage
\begin{problem}{}{}
% Subject: varc
\textbf{Instructions}

Instructions   

**The passage below is accompanied by a set of questions. Choose the best answer to each question.**
[Octopuses are] misfits in their own extended families . . . They belong to the Mollusca class Cephalopoda. But they don’t look like their cousins at all. Other molluscs include sea snails, sea slugs, bivalves - most are shelled invertebrates with a dorsal foot. Cephalopods are all arms, and can be as tiny as 1 centimetre and as large at 30 feet. Some of them have brains the size of a walnut, which is large for an invertebrate. . . .
It makes sense for these molluscs to have added protection in the form of a higher cognition; they don’t have a shell covering them, and pretty much everything feeds on cephalopods, including humans. But how did cephalopods manage to secure their own invisibility cloak? Cephalopods fire from multiple cylinders to achieve this in varying degrees from species to species. There are four main catalysts - chromatophores, iridophores, papillae and leucophores. . . .
[Chromatophores] are organs on their bodies that contain pigment sacs, which have red, yellow and brown pigment granules. These sacs have a network of radial muscles, meaning muscles arranged in a circle radiating outwards. These are connected to the brain by a nerve. When the cephalopod wants to change colour, the brain carries an electrical impulse through the nerve to the muscles that expand outwards, pulling open the sacs to display the colours on the skin. Why these three colours? Because these are the colours the light reflects at the depths they live in (the rest is absorbed before it reaches those depths). . . .
Well, what about other colours? Cue the iridophores. Think of a second level of skin that has thin stacks of cells. These can reflect light back at different wavelengths. . . . It’s using the same properties that we’ve seen in hologram stickers, or rainbows on puddles of oil. You move your head and you see a different colour. The sticker isn’t doing anything but reflecting light - it’s your movement that’s changing the appearance of the colour. This property of holograms, oil and other such surfaces is called “iridescence”. . . .
Papillae are sections of the skin that can be deformed to make a texture bumpy. Even humans possess them (goosebumps) but cannot use them in the manner that cephalopods can. For instance, the use of these cells is how an octopus can wrap itself over a rock and appear jagged or how a squid or cuttlefish can imitate the look of a coral reef by growing miniature towers on its skin. It actually matches the texture of the substrate it chooses.
Finally, the leucophores: According to a paper, published in Nature, cuttlefish and octopuses possess an additional type of reflector cell called a leucophore. They are cells that scatter full spectrum light so that they appear white in a similar way that a polar bear’s fur appears white. Leucophores will also reflect any filtered light shown on them . . . If the water appears blue at a certain depth, the octopuses and cuttlefish can appear blue; if the water appears green, they appear green, and so on and so forth.

\vspace{1em}
\textbf{Question 16}

Which one of the following statements, if true, would weaken the author’s claim that humans are musicking creatures?

\begin{enumerate}[label=\Alph*.]
    \item Nonmusical capacities are of far greater consequence to human survival than the capacity for music.
    \item From a cognitive and psychological vantage, musicking arises from unconscious dispositions, not conscious ones.
    \item As musicking is neither language-like nor symbol-like, it is a much older form of expression.
    \item Musical capacities are primarily socio-cultural, which explains the wide diversity of musical forms.
[](https://cracku.in/16-which-one-of-the-following-statements-if-true-woul-x-cat-2022-slot-2)
\end{enumerate}
\vspace{1em}
\textbf{Solution:}

Let us examine the given statements:
_Option A_ : [**Nonmusical capacities are of far greater consequence to human survival than the capacity for music.**] It is unclear how this ties into the discussion since the survival aspect is not touched upon or implied in the passage. Thus, the claim, if true, does little to undermine the author's claim that humans are musicking creatures. 
_Option B_ : [**From a cognitive and psychological vantage, musicking arises from unconscious dispositions, not conscious ones.**] This does not weaken the author's claim that musicking is a universal aspect of the human experience - while the author does suggest that musicking arises from innate dispositions, the fact that these dispositions may be unconscious rather than conscious does not weaken the overall argument.
_Option C_ : [**As musicking is neither language-like nor symbol-like, it is a much older form of expression.**] This, again, does not contradict the author's claim in the passage. Though the author suggests that musicking is distinct from language and symbol-making, this does not necessarily mean that it is a much older form of expression. In fact, the author notes that the emergence of musicking can be traced back to at least 50,000 years ago, which is relatively recent in evolutionary terms. Even if this were true, it does not undermine the author's claim.
_Option D_ : [**Musical capacities are primarily socio-cultural, which explains the wide diversity of musical forms**.] If true, this directly contradicts the author's claim that musicking is a universal aspect of the human experience. The author argues that the capacity for musicking is innate in all humans and that it has a long history that is both sociocultural and biological in nature. However, if musical capacities are primarily socio-cultural, as suggested in Option D, this would mean that musicking is largely shaped by cultural and social factors rather than being a fundamental aspect of the human experience. This would greatly weaken the author's overall argument that all humans are musicking creatures. Hence, Option D is the correct choice.

Instructions   

For the following questions answer them individually

\end{problem}

\newpage
\begin{problem}{}{}
% Subject: varc
\textbf{Instructions}

Instructions   

**The passage below is accompanied by a set of questions. Choose the best answer to each question.**
[Octopuses are] misfits in their own extended families . . . They belong to the Mollusca class Cephalopoda. But they don’t look like their cousins at all. Other molluscs include sea snails, sea slugs, bivalves - most are shelled invertebrates with a dorsal foot. Cephalopods are all arms, and can be as tiny as 1 centimetre and as large at 30 feet. Some of them have brains the size of a walnut, which is large for an invertebrate. . . .
It makes sense for these molluscs to have added protection in the form of a higher cognition; they don’t have a shell covering them, and pretty much everything feeds on cephalopods, including humans. But how did cephalopods manage to secure their own invisibility cloak? Cephalopods fire from multiple cylinders to achieve this in varying degrees from species to species. There are four main catalysts - chromatophores, iridophores, papillae and leucophores. . . .
[Chromatophores] are organs on their bodies that contain pigment sacs, which have red, yellow and brown pigment granules. These sacs have a network of radial muscles, meaning muscles arranged in a circle radiating outwards. These are connected to the brain by a nerve. When the cephalopod wants to change colour, the brain carries an electrical impulse through the nerve to the muscles that expand outwards, pulling open the sacs to display the colours on the skin. Why these three colours? Because these are the colours the light reflects at the depths they live in (the rest is absorbed before it reaches those depths). . . .
Well, what about other colours? Cue the iridophores. Think of a second level of skin that has thin stacks of cells. These can reflect light back at different wavelengths. . . . It’s using the same properties that we’ve seen in hologram stickers, or rainbows on puddles of oil. You move your head and you see a different colour. The sticker isn’t doing anything but reflecting light - it’s your movement that’s changing the appearance of the colour. This property of holograms, oil and other such surfaces is called “iridescence”. . . .
Papillae are sections of the skin that can be deformed to make a texture bumpy. Even humans possess them (goosebumps) but cannot use them in the manner that cephalopods can. For instance, the use of these cells is how an octopus can wrap itself over a rock and appear jagged or how a squid or cuttlefish can imitate the look of a coral reef by growing miniature towers on its skin. It actually matches the texture of the substrate it chooses.
Finally, the leucophores: According to a paper, published in Nature, cuttlefish and octopuses possess an additional type of reflector cell called a leucophore. They are cells that scatter full spectrum light so that they appear white in a similar way that a polar bear’s fur appears white. Leucophores will also reflect any filtered light shown on them . . . If the water appears blue at a certain depth, the octopuses and cuttlefish can appear blue; if the water appears green, they appear green, and so on and so forth.

\vspace{1em}
\textbf{Question 17}

**The four sentences (labelled 1, 2, 3 and 4) below, when properly sequenced, would yield a coherent paragraph. Decide on the proper sequencing of the order of the sentences and key in the sequence of the four numbers as your answer:**
1. The trajectory of cheerfulness through the self is linked to the history of the word ‘cheer’ which comes from an Old French meaning ‘face’.
2. Translations of the Bible into vernacular languages, expanded the noun ‘cheer’ into the more abstract ‘cheerful-ness’, something that circulates as an emotional and social quality defining the self and a moral community.
3. When you take on a cheerful expression, no matter what the state of your soul, your cheerfulness moves into the self: the interior of the self is changed by the power of cheer.
4. People in the medieval ‘Canterbury Tales’ have a ‘piteous’ or a ‘sober’ cheer; ‘cheer’ is an expression and a body part, lying at the intersection of emotions and physiognomy.
Backspace  
789  
456  
123  
0.-  
Clear All  

Submit
[](https://cracku.in/17-the-four-sentences-labelled-1-2-3-and-4-below-when-x-cat-2022-slot-2)

\begin{enumerate}[label=\Alph*.]
\end{enumerate}
\vspace{1em}
\textbf{Solution:}

The given set of sentences talks about the concept of cheerfulness and how it can affect a person's internal state. It discusses the history of the word "cheer," which originally referred to a person's facial expression, and how it came to be associated with an abstract concept of positive emotion. The passage also notes that translations of the Bible into vernacular languages helped to expand the meaning of "cheer" to include a sense of cheerfulness or positivity as a quality of the self and a moral community. We note that sentences 1, 2 and 4 are part of the timeline the author traces, while statement 3 is a broad claim that would serve as an apt introduction. The idea of cheerfulness and self that is mentioned in 3 is linked to the " trajectory of cheerfulness through the self" highlighted in statement 1, allowing us to form the pair [3-1]. Note that 4 discusses how cheer came to be perceived as a trait at the intersection of emotion and a physical feature [as highlighted in 1 - where it has been categorised as a facial feature]. Sentence 2 then extends on this idea by how cheerfulness, over time, came to be considered as "an emotional and social quality." Hence, the correct arrangement is 3142.

\end{problem}

\newpage
\begin{problem}{}{}
% Subject: varc
\textbf{Instructions}

Instructions   

**The passage below is accompanied by a set of questions. Choose the best answer to each question.**
[Octopuses are] misfits in their own extended families . . . They belong to the Mollusca class Cephalopoda. But they don’t look like their cousins at all. Other molluscs include sea snails, sea slugs, bivalves - most are shelled invertebrates with a dorsal foot. Cephalopods are all arms, and can be as tiny as 1 centimetre and as large at 30 feet. Some of them have brains the size of a walnut, which is large for an invertebrate. . . .
It makes sense for these molluscs to have added protection in the form of a higher cognition; they don’t have a shell covering them, and pretty much everything feeds on cephalopods, including humans. But how did cephalopods manage to secure their own invisibility cloak? Cephalopods fire from multiple cylinders to achieve this in varying degrees from species to species. There are four main catalysts - chromatophores, iridophores, papillae and leucophores. . . .
[Chromatophores] are organs on their bodies that contain pigment sacs, which have red, yellow and brown pigment granules. These sacs have a network of radial muscles, meaning muscles arranged in a circle radiating outwards. These are connected to the brain by a nerve. When the cephalopod wants to change colour, the brain carries an electrical impulse through the nerve to the muscles that expand outwards, pulling open the sacs to display the colours on the skin. Why these three colours? Because these are the colours the light reflects at the depths they live in (the rest is absorbed before it reaches those depths). . . .
Well, what about other colours? Cue the iridophores. Think of a second level of skin that has thin stacks of cells. These can reflect light back at different wavelengths. . . . It’s using the same properties that we’ve seen in hologram stickers, or rainbows on puddles of oil. You move your head and you see a different colour. The sticker isn’t doing anything but reflecting light - it’s your movement that’s changing the appearance of the colour. This property of holograms, oil and other such surfaces is called “iridescence”. . . .
Papillae are sections of the skin that can be deformed to make a texture bumpy. Even humans possess them (goosebumps) but cannot use them in the manner that cephalopods can. For instance, the use of these cells is how an octopus can wrap itself over a rock and appear jagged or how a squid or cuttlefish can imitate the look of a coral reef by growing miniature towers on its skin. It actually matches the texture of the substrate it chooses.
Finally, the leucophores: According to a paper, published in Nature, cuttlefish and octopuses possess an additional type of reflector cell called a leucophore. They are cells that scatter full spectrum light so that they appear white in a similar way that a polar bear’s fur appears white. Leucophores will also reflect any filtered light shown on them . . . If the water appears blue at a certain depth, the octopuses and cuttlefish can appear blue; if the water appears green, they appear green, and so on and so forth.

\vspace{1em}
\textbf{Question 18}

**There is a sentence that is missing in the paragraph below. Look at the paragraph and decide in which blank (option 1, 2, 3, or 4) the following sentence would best fit.**
**Sentence** : Most were first-time users of a tablet and a digital app.
**Paragraph** : Aage Badhein’s USP lies in the ethnographic research that constituted the foundation of its development process. Customizations based on learning directly from potential users were critical to making this self-paced app suitable for both a literate and non-literate audience. ___(1)___ The user interface caters to a Hindi-speaking audience who have minimal to no experience with digital services and devices. ___(2)___ The content and functionality of the app are suitable for a wide audience. This includes youth preparing for an independent role in life or a student ready to create a strong foundation of financial management early in her life. ___(3)___ Household members desirous of improving their family’s financial strength to reach their aspirations can also benefit. We piloted Aage Badhein in early 2021 with over 400 women from rural areas. ___(4)___ The digital solution generated a large amount of interest in the communities.

\begin{enumerate}[label=\Alph*.]
    \item Option 1
    \item Option 2
    \item Option 3
    \item Option 4
[](https://cracku.in/18-there-is-a-sentence-that-is-missing-in-the-paragra-x-cat-2022-slot-2)
\end{enumerate}
\vspace{1em}
\textbf{Solution:}

The sentence "**Most were first-time users of a tablet and a digital app** " would fit best in blank 4 because it provides information about the specific group of people who were involved in the pilot study for the app. The preceding sentence mentions that the app was piloted in early 2021 with over 400 women from rural areas, and the added sentence provides further detail about the specific experiences that these users had with technology. This information helps to clarify that the pilot study for the app involved a group of people who may be unfamiliar with digital devices and services, and it also helps to emphasize the potential impact that the app may have had on these users.
Placing the sentence in blank 4 is logical because it provides relevant and specific information about the pilot study for the app, which is the topic being discussed in that part of the paragraph. By contrast, placing the sentence in any of the other blanks would be less logical because it would be introducing information that is not directly related to the main topic being discussed in those parts of the paragraph.
Hence, Option D is the correct choice.

\end{problem}

\newpage
\begin{problem}{}{}
% Subject: varc
\textbf{Instructions}

Instructions   

**The passage below is accompanied by a set of questions. Choose the best answer to each question.**
[Octopuses are] misfits in their own extended families . . . They belong to the Mollusca class Cephalopoda. But they don’t look like their cousins at all. Other molluscs include sea snails, sea slugs, bivalves - most are shelled invertebrates with a dorsal foot. Cephalopods are all arms, and can be as tiny as 1 centimetre and as large at 30 feet. Some of them have brains the size of a walnut, which is large for an invertebrate. . . .
It makes sense for these molluscs to have added protection in the form of a higher cognition; they don’t have a shell covering them, and pretty much everything feeds on cephalopods, including humans. But how did cephalopods manage to secure their own invisibility cloak? Cephalopods fire from multiple cylinders to achieve this in varying degrees from species to species. There are four main catalysts - chromatophores, iridophores, papillae and leucophores. . . .
[Chromatophores] are organs on their bodies that contain pigment sacs, which have red, yellow and brown pigment granules. These sacs have a network of radial muscles, meaning muscles arranged in a circle radiating outwards. These are connected to the brain by a nerve. When the cephalopod wants to change colour, the brain carries an electrical impulse through the nerve to the muscles that expand outwards, pulling open the sacs to display the colours on the skin. Why these three colours? Because these are the colours the light reflects at the depths they live in (the rest is absorbed before it reaches those depths). . . .
Well, what about other colours? Cue the iridophores. Think of a second level of skin that has thin stacks of cells. These can reflect light back at different wavelengths. . . . It’s using the same properties that we’ve seen in hologram stickers, or rainbows on puddles of oil. You move your head and you see a different colour. The sticker isn’t doing anything but reflecting light - it’s your movement that’s changing the appearance of the colour. This property of holograms, oil and other such surfaces is called “iridescence”. . . .
Papillae are sections of the skin that can be deformed to make a texture bumpy. Even humans possess them (goosebumps) but cannot use them in the manner that cephalopods can. For instance, the use of these cells is how an octopus can wrap itself over a rock and appear jagged or how a squid or cuttlefish can imitate the look of a coral reef by growing miniature towers on its skin. It actually matches the texture of the substrate it chooses.
Finally, the leucophores: According to a paper, published in Nature, cuttlefish and octopuses possess an additional type of reflector cell called a leucophore. They are cells that scatter full spectrum light so that they appear white in a similar way that a polar bear’s fur appears white. Leucophores will also reflect any filtered light shown on them . . . If the water appears blue at a certain depth, the octopuses and cuttlefish can appear blue; if the water appears green, they appear green, and so on and so forth.

\vspace{1em}
\textbf{Question 19}

**The four sentences (labelled 1, 2, 3 and 4) below, when properly sequenced, would yield a coherent paragraph. Decide on the proper sequencing of the order of the sentences and key in the sequence of the four numbers as your answer:**
1. Women may prioritize cooking because they feel they alone are responsible for mediating a toxic and unhealthy food system.
2. Food is commonly framed through the lens of individual choice: you can choose to eat healthily.
3. This is particularly so in a neoliberal context where the state has transferred the responsibility for food onto individual consumers.
4. The individualized framing of choice appeals to a popular desire to experience agency, but draws away from the structural obstacles that stratify individual food choices.
Backspace  
789  
456  
123  
0.-  
Clear All  

Submit
[](https://cracku.in/19-the-four-sentences-labelled-1-2-3-and-4-below-when-x-cat-2022-slot-2)

\begin{enumerate}[label=\Alph*.]
\end{enumerate}
\vspace{1em}
\textbf{Solution:}

The statements here discuss how the way we think about food often centres around individual choice, rather than acknowledging the larger societal and structural factors that influence food choices. It also mentions the role of neoliberalism in shifting the responsibility for food onto individual consumers and how this may disproportionately impact women, who may feel a sense of responsibility to navigate and "mediate" an unhealthy food system. Sentences 2 and 4 create a logical block since 2 highlights the framing of food through the "lens of individual choice", and 4 extends on the impact of such a framing ["appeals to a popular desire to experience agency, but draws away from the structural obstacles that stratify individual food choices"]. "This" in statement 3 refers to the particular outcome of framing underlined in 4, allowing us to create a logical block: [2-4-3]. Statement 1 can be placed at the end of this arrangement to emphasise how women are perhaps disproportionately impacted by the elements discussed in 2-4-3. Hence, the correct answer is 2431.

\end{problem}

\newpage
\begin{problem}{}{}
% Subject: varc
\textbf{Instructions}

Instructions   

**The passage below is accompanied by a set of questions. Choose the best answer to each question.**
[Octopuses are] misfits in their own extended families . . . They belong to the Mollusca class Cephalopoda. But they don’t look like their cousins at all. Other molluscs include sea snails, sea slugs, bivalves - most are shelled invertebrates with a dorsal foot. Cephalopods are all arms, and can be as tiny as 1 centimetre and as large at 30 feet. Some of them have brains the size of a walnut, which is large for an invertebrate. . . .
It makes sense for these molluscs to have added protection in the form of a higher cognition; they don’t have a shell covering them, and pretty much everything feeds on cephalopods, including humans. But how did cephalopods manage to secure their own invisibility cloak? Cephalopods fire from multiple cylinders to achieve this in varying degrees from species to species. There are four main catalysts - chromatophores, iridophores, papillae and leucophores. . . .
[Chromatophores] are organs on their bodies that contain pigment sacs, which have red, yellow and brown pigment granules. These sacs have a network of radial muscles, meaning muscles arranged in a circle radiating outwards. These are connected to the brain by a nerve. When the cephalopod wants to change colour, the brain carries an electrical impulse through the nerve to the muscles that expand outwards, pulling open the sacs to display the colours on the skin. Why these three colours? Because these are the colours the light reflects at the depths they live in (the rest is absorbed before it reaches those depths). . . .
Well, what about other colours? Cue the iridophores. Think of a second level of skin that has thin stacks of cells. These can reflect light back at different wavelengths. . . . It’s using the same properties that we’ve seen in hologram stickers, or rainbows on puddles of oil. You move your head and you see a different colour. The sticker isn’t doing anything but reflecting light - it’s your movement that’s changing the appearance of the colour. This property of holograms, oil and other such surfaces is called “iridescence”. . . .
Papillae are sections of the skin that can be deformed to make a texture bumpy. Even humans possess them (goosebumps) but cannot use them in the manner that cephalopods can. For instance, the use of these cells is how an octopus can wrap itself over a rock and appear jagged or how a squid or cuttlefish can imitate the look of a coral reef by growing miniature towers on its skin. It actually matches the texture of the substrate it chooses.
Finally, the leucophores: According to a paper, published in Nature, cuttlefish and octopuses possess an additional type of reflector cell called a leucophore. They are cells that scatter full spectrum light so that they appear white in a similar way that a polar bear’s fur appears white. Leucophores will also reflect any filtered light shown on them . . . If the water appears blue at a certain depth, the octopuses and cuttlefish can appear blue; if the water appears green, they appear green, and so on and so forth.

\vspace{1em}
\textbf{Question 20}

**The four sentences (labelled 1, 2, 3 and 4) below, when properly sequenced, would yield a coherent paragraph. Decide on the proper sequencing of the order of the sentences and key in the sequence of the four numbers as your answer:**
1. From chemical pollutants in the environment to the damming of rivers to invasive species transported through global trade and travel, every environmental issue is different and there is no single tech solution that can solve this crisis.
2. Discourse on the threat of environmental collapse revolves around cutting down emissions, but biodiversity loss and ecosystem collapse are caused by myriad and diverse reasons.
3. This would require legislation that recognises the rights of future generations and other species that allows the judiciary to uphold a much higher standard of environmental protection than currently possible.
4. Clearly, our environmental crisis requires large political solutions, not minor technological ones, so, instead of focusing on infinite growth, we could consider a path of stable-state economies, while preserving markets and healthy competition.
Backspace  
789  
456  
123  
0.-  
Clear All  

Submit
[](https://cracku.in/20-the-four-sentences-labelled-1-2-3-and-4-below-when-x-cat-2022-slot-2)

\begin{enumerate}[label=\Alph*.]
\end{enumerate}
\vspace{1em}
\textbf{Solution:}

The statements discuss the complexity of environmental issues and the need for political solutions rather than technological ones to address the threat of environmental collapse. The discussion focuses on the various causes of biodiversity loss and ecosystem collapse, such as chemical pollutants, damming of rivers, and invasive species, and highlights that there is no single technological solution to address these issues. The suggestion is that we should consider moving towards stable-state economies, which preserve markets and competition while also recognizing the rights of future generations and other species and allowing for stronger environmental protections. Statements 1 and 2 both address the complexity and diverse nature of environmental issues. Statement 2 notes that while discourse often focuses on emissions, there are many other factors contributing to environmental collapse, including biodiversity loss and ecosystem collapse. Statement 1 then adds to this by outlining specific examples of diverse causes of environmental problems, including chemical pollutants, damming of rivers, and invasive species. Statements 4 and 3 offer suggestions to combat the issues highlighted in 2 and 1. Hence, the correct arrangement is 2143.

\end{problem}

\newpage
\begin{problem}{}{}
% Subject: varc
\textbf{Instructions}

Instructions   

**The passage below is accompanied by a set of questions. Choose the best answer to each question.**
[Octopuses are] misfits in their own extended families . . . They belong to the Mollusca class Cephalopoda. But they don’t look like their cousins at all. Other molluscs include sea snails, sea slugs, bivalves - most are shelled invertebrates with a dorsal foot. Cephalopods are all arms, and can be as tiny as 1 centimetre and as large at 30 feet. Some of them have brains the size of a walnut, which is large for an invertebrate. . . .
It makes sense for these molluscs to have added protection in the form of a higher cognition; they don’t have a shell covering them, and pretty much everything feeds on cephalopods, including humans. But how did cephalopods manage to secure their own invisibility cloak? Cephalopods fire from multiple cylinders to achieve this in varying degrees from species to species. There are four main catalysts - chromatophores, iridophores, papillae and leucophores. . . .
[Chromatophores] are organs on their bodies that contain pigment sacs, which have red, yellow and brown pigment granules. These sacs have a network of radial muscles, meaning muscles arranged in a circle radiating outwards. These are connected to the brain by a nerve. When the cephalopod wants to change colour, the brain carries an electrical impulse through the nerve to the muscles that expand outwards, pulling open the sacs to display the colours on the skin. Why these three colours? Because these are the colours the light reflects at the depths they live in (the rest is absorbed before it reaches those depths). . . .
Well, what about other colours? Cue the iridophores. Think of a second level of skin that has thin stacks of cells. These can reflect light back at different wavelengths. . . . It’s using the same properties that we’ve seen in hologram stickers, or rainbows on puddles of oil. You move your head and you see a different colour. The sticker isn’t doing anything but reflecting light - it’s your movement that’s changing the appearance of the colour. This property of holograms, oil and other such surfaces is called “iridescence”. . . .
Papillae are sections of the skin that can be deformed to make a texture bumpy. Even humans possess them (goosebumps) but cannot use them in the manner that cephalopods can. For instance, the use of these cells is how an octopus can wrap itself over a rock and appear jagged or how a squid or cuttlefish can imitate the look of a coral reef by growing miniature towers on its skin. It actually matches the texture of the substrate it chooses.
Finally, the leucophores: According to a paper, published in Nature, cuttlefish and octopuses possess an additional type of reflector cell called a leucophore. They are cells that scatter full spectrum light so that they appear white in a similar way that a polar bear’s fur appears white. Leucophores will also reflect any filtered light shown on them . . . If the water appears blue at a certain depth, the octopuses and cuttlefish can appear blue; if the water appears green, they appear green, and so on and so forth.

\vspace{1em}
\textbf{Question 21}

**The passage given below is followed by four alternate summaries. Choose the option that best captures the essence of the passage.**
There's a common idea that museum artworks are somehow timeless objects available to admire for generations to come. But many are objects of decay. Even the most venerable Old Master paintings don't escape: pigments discolour, varnishes crack, canvases warp. This challenging fact of art-world life is down to something that sounds more like a thread from a morality tale: inherent vice. Damien Hirst's iconic shark floating in a tank - entitled The Physical Impossibility of Death in the Mind of Someone Living - is a work that put a spotlight on inherent vice. When he made it in 1991, Hirst got himself in a pickle by not using the right kind of pickle to preserve the giant fish. The result was that the shark began to decompose quite quickly - its preserving liquid clouding, the skin wrinkling, and an unpleasant smell wafting from the tank.

\begin{enumerate}[label=\Alph*.]
    \item Museums are left with the moral responsibility of restoring and preserving the artworks since artists cannot preserve their works beyond their life.
    \item Museums have to guard timeless art treasures from intrinsic defects such as the deterioration of paint, polish and canvas.
    \item The role of museums has evolved to ensure that the artworks are preserved forever in addition to guarding and displaying them.
    \item Artworks may not last forever; they may deteriorate with time, and the challenge is to slow down their degeneration.
[](https://cracku.in/21-the-passage-given-below-is-followed-by-four-altern-x-cat-2022-slot-2)
\end{enumerate}
\vspace{1em}
\textbf{Solution:}

The passage is about the 'inherent vice' or the natural tendency of certain artworks to deteriorate over time due to various factors such as discolouration of pigments, cracking of varnishes, and warping of canvases. The passage also mentions an example of Damien Hirst's artwork, The Physical Impossibility of Death in the Mind of Someone Living, which began to decompose quickly due to the use of the wrong preserving liquid. In this regard, Option D offers an apt summary of the passage because it accurately captures the main idea of the passage, which is that artworks may not last forever and may deteriorate with time. Option A is incorrect because the passage does not mention any moral responsibility of museums to restore and preserve artworks. Similarly, Option B can be eliminated since the passage does not specifically mention museums guarding art treasures from intrinsic defects. Option C is also inaccurate because the discussion does not present the evolution of the role of museums in preserving artworks forever.
Hence, Option D is the correct choice.

\end{problem}

\newpage
\begin{problem}{}{}
% Subject: varc
\textbf{Instructions}

Instructions   

**The passage below is accompanied by a set of questions. Choose the best answer to each question.**
[Octopuses are] misfits in their own extended families . . . They belong to the Mollusca class Cephalopoda. But they don’t look like their cousins at all. Other molluscs include sea snails, sea slugs, bivalves - most are shelled invertebrates with a dorsal foot. Cephalopods are all arms, and can be as tiny as 1 centimetre and as large at 30 feet. Some of them have brains the size of a walnut, which is large for an invertebrate. . . .
It makes sense for these molluscs to have added protection in the form of a higher cognition; they don’t have a shell covering them, and pretty much everything feeds on cephalopods, including humans. But how did cephalopods manage to secure their own invisibility cloak? Cephalopods fire from multiple cylinders to achieve this in varying degrees from species to species. There are four main catalysts - chromatophores, iridophores, papillae and leucophores. . . .
[Chromatophores] are organs on their bodies that contain pigment sacs, which have red, yellow and brown pigment granules. These sacs have a network of radial muscles, meaning muscles arranged in a circle radiating outwards. These are connected to the brain by a nerve. When the cephalopod wants to change colour, the brain carries an electrical impulse through the nerve to the muscles that expand outwards, pulling open the sacs to display the colours on the skin. Why these three colours? Because these are the colours the light reflects at the depths they live in (the rest is absorbed before it reaches those depths). . . .
Well, what about other colours? Cue the iridophores. Think of a second level of skin that has thin stacks of cells. These can reflect light back at different wavelengths. . . . It’s using the same properties that we’ve seen in hologram stickers, or rainbows on puddles of oil. You move your head and you see a different colour. The sticker isn’t doing anything but reflecting light - it’s your movement that’s changing the appearance of the colour. This property of holograms, oil and other such surfaces is called “iridescence”. . . .
Papillae are sections of the skin that can be deformed to make a texture bumpy. Even humans possess them (goosebumps) but cannot use them in the manner that cephalopods can. For instance, the use of these cells is how an octopus can wrap itself over a rock and appear jagged or how a squid or cuttlefish can imitate the look of a coral reef by growing miniature towers on its skin. It actually matches the texture of the substrate it chooses.
Finally, the leucophores: According to a paper, published in Nature, cuttlefish and octopuses possess an additional type of reflector cell called a leucophore. They are cells that scatter full spectrum light so that they appear white in a similar way that a polar bear’s fur appears white. Leucophores will also reflect any filtered light shown on them . . . If the water appears blue at a certain depth, the octopuses and cuttlefish can appear blue; if the water appears green, they appear green, and so on and so forth.

\vspace{1em}
\textbf{Question 22}

**The passage given below is followed by four alternate summaries. Choose the option that best captures the essence of the passage.**
Today, many of the debates about behavioural control in the age of big data echo Cold War-era anxieties about brainwashing, insidious manipulation and repression in the ‘technological society’. In his book Psychopolitics, Han warns of the sophisticated use of targeted online content, enabling ‘influence to take place on a pre-reflexive level’. On our current trajectory, “freedom will prove to have been merely an interlude.” The fear is that the digital age has not liberated us but exposed us, by offering up our private lives to machine-learning algorithms that can process masses of personal and behavioural data. In a world of influencers and digital entrepreneurs, it’s not easy to imagine the resurgence of a culture engendered through disconnect and disaffiliation, but concerns over the threat of online targeting, polarisation and big data have inspired recent polemics about the need to rediscover solitude and disconnect.

\begin{enumerate}[label=\Alph*.]
    \item Rather than freeing us, digital technology is enslaving us by collecting personal information and influencing our online behaviour.
    \item With big data making personal information freely available, the debate on the nature of freedom and the need for privacy has resurfaced.
    \item The role of technology in influencing public behaviour is reminiscent of the manner in which behaviour was manipulated during the Cold War.
    \item The notion of freedom and privacy is at stake in a world where artificial intelligence is capable of influencing behaviour through data gathered online.
[](https://cracku.in/22-the-passage-given-below-is-followed-by-four-altern-x-cat-2022-slot-2)
\end{enumerate}
\vspace{1em}
\textbf{Solution:}

The passage discusses the ways in which big data and targeted online content can potentially influence and manipulate behaviour, leading to concerns over freedom and privacy in the digital age. This is reflected in the statements that "behavioural control" in the age of big data echoes Cold War-era anxieties about "brainwashing" and "repression," and that the use of targeted online content can enable "influence to take place on a pre-reflexive level." The passage also mentions the fear that the digital age has not liberated us, but rather exposed us by making personal and behavioural data available to machine-learning algorithms. Option B accurately reflects this central theme of the passage by stating that the debate on the nature of freedom and privacy has resurfaced due to the availability of personal information through big data. Option A is incorrect because it goes beyond the scope of the passage by stating that digital technology is "enslaving" us, which is not explicitly stated in the text. Similarly, Option C is inaccurate since the author only mentions the Cold War as a reference point for similar debates on behavioural control, but does not focus on the Cold War itself. Option D is wrong because the passage does not mention artificial intelligence specifically, but rather machine-learning algorithms.  

Hence, Option B is the correct choice.

\end{problem}

\newpage
\begin{problem}{}{}
% Subject: varc
\textbf{Instructions}

Instructions   

**The passage below is accompanied by a set of questions. Choose the best answer to each question.**
[Octopuses are] misfits in their own extended families . . . They belong to the Mollusca class Cephalopoda. But they don’t look like their cousins at all. Other molluscs include sea snails, sea slugs, bivalves - most are shelled invertebrates with a dorsal foot. Cephalopods are all arms, and can be as tiny as 1 centimetre and as large at 30 feet. Some of them have brains the size of a walnut, which is large for an invertebrate. . . .
It makes sense for these molluscs to have added protection in the form of a higher cognition; they don’t have a shell covering them, and pretty much everything feeds on cephalopods, including humans. But how did cephalopods manage to secure their own invisibility cloak? Cephalopods fire from multiple cylinders to achieve this in varying degrees from species to species. There are four main catalysts - chromatophores, iridophores, papillae and leucophores. . . .
[Chromatophores] are organs on their bodies that contain pigment sacs, which have red, yellow and brown pigment granules. These sacs have a network of radial muscles, meaning muscles arranged in a circle radiating outwards. These are connected to the brain by a nerve. When the cephalopod wants to change colour, the brain carries an electrical impulse through the nerve to the muscles that expand outwards, pulling open the sacs to display the colours on the skin. Why these three colours? Because these are the colours the light reflects at the depths they live in (the rest is absorbed before it reaches those depths). . . .
Well, what about other colours? Cue the iridophores. Think of a second level of skin that has thin stacks of cells. These can reflect light back at different wavelengths. . . . It’s using the same properties that we’ve seen in hologram stickers, or rainbows on puddles of oil. You move your head and you see a different colour. The sticker isn’t doing anything but reflecting light - it’s your movement that’s changing the appearance of the colour. This property of holograms, oil and other such surfaces is called “iridescence”. . . .
Papillae are sections of the skin that can be deformed to make a texture bumpy. Even humans possess them (goosebumps) but cannot use them in the manner that cephalopods can. For instance, the use of these cells is how an octopus can wrap itself over a rock and appear jagged or how a squid or cuttlefish can imitate the look of a coral reef by growing miniature towers on its skin. It actually matches the texture of the substrate it chooses.
Finally, the leucophores: According to a paper, published in Nature, cuttlefish and octopuses possess an additional type of reflector cell called a leucophore. They are cells that scatter full spectrum light so that they appear white in a similar way that a polar bear’s fur appears white. Leucophores will also reflect any filtered light shown on them . . . If the water appears blue at a certain depth, the octopuses and cuttlefish can appear blue; if the water appears green, they appear green, and so on and so forth.

\vspace{1em}
\textbf{Question 23}

**The passage given below is followed by four alternate summaries. Choose the option that best captures the essence of the passage.**
Several of the world’s earliest cities were organised along egalitarian lines. In some regions, urban populations governed themselves for centuries without any indication of the temples and palaces that would later emerge; in others, temples and palaces never emerged at all, and there is simply no evidence of a class of administrators or any other sort of ruling stratum. It would seem that the mere fact of urban life does not, necessarily, imply any particular form of political organization, and never did. Far from resigning us to inequality, the picture that is now emerging of humanity’s past may open our eyes to egalitarian possibilities we otherwise would have never considered.

\begin{enumerate}[label=\Alph*.]
    \item The lack of hierarchical administration in ancient cities can be deduced by the absence of religious and regal structures such as temples and palaces.
    \item Contrary to our assumption that urban settlements have always involved hierarchical political and administrative structures, ancient cities were not organised in this way.
    \item The emergence of a class of administrators and ruling stratum transformed the egalitarian urban life of ancient cities to the hierarchical civic organisations of today.
    \item We now have the evidence in support of the existence of an egalitarian urban life in some ancient cities, where political and civic organisation was far less hierarchical.
[](https://cracku.in/23-the-passage-given-below-is-followed-by-four-altern-x-cat-2022-slot-2)
\end{enumerate}
\vspace{1em}
\textbf{Solution:}

The passage is about the political and civic organization of ancient cities. It states that some ancient cities were organized along egalitarian lines, without any indication of temples or palaces (which suggests a lack of a ruling class or administrators), and that in other cities, temples and palaces never emerged at all. Option D correctly summarizes this information by stating that there was evidence of an egalitarian urban life in some ancient cities, where the political and civic organization was less hierarchical. Option A is incorrect because it only mentions the absence of temples and palaces, but does not mention the fact that some ancient cities were organized along egalitarian lines. On a similar note, Option B presents an exaggeration by suggesting that 'all' ancient cities were organized along egalitarian lines, which is not stated in the passage. Option C is also inaccurate since it asserts that ancient cities were transformed from egalitarian to hierarchical, but the passage only states that some ancient cities were egalitarian and does not mention any transformation. 
Hence, Option D is the correct choice.

\end{problem}

\newpage
\begin{problem}{}{}
% Subject: varc
\textbf{Instructions}

Instructions   

**The passage below is accompanied by a set of questions. Choose the best answer to each question.**
[Octopuses are] misfits in their own extended families . . . They belong to the Mollusca class Cephalopoda. But they don’t look like their cousins at all. Other molluscs include sea snails, sea slugs, bivalves - most are shelled invertebrates with a dorsal foot. Cephalopods are all arms, and can be as tiny as 1 centimetre and as large at 30 feet. Some of them have brains the size of a walnut, which is large for an invertebrate. . . .
It makes sense for these molluscs to have added protection in the form of a higher cognition; they don’t have a shell covering them, and pretty much everything feeds on cephalopods, including humans. But how did cephalopods manage to secure their own invisibility cloak? Cephalopods fire from multiple cylinders to achieve this in varying degrees from species to species. There are four main catalysts - chromatophores, iridophores, papillae and leucophores. . . .
[Chromatophores] are organs on their bodies that contain pigment sacs, which have red, yellow and brown pigment granules. These sacs have a network of radial muscles, meaning muscles arranged in a circle radiating outwards. These are connected to the brain by a nerve. When the cephalopod wants to change colour, the brain carries an electrical impulse through the nerve to the muscles that expand outwards, pulling open the sacs to display the colours on the skin. Why these three colours? Because these are the colours the light reflects at the depths they live in (the rest is absorbed before it reaches those depths). . . .
Well, what about other colours? Cue the iridophores. Think of a second level of skin that has thin stacks of cells. These can reflect light back at different wavelengths. . . . It’s using the same properties that we’ve seen in hologram stickers, or rainbows on puddles of oil. You move your head and you see a different colour. The sticker isn’t doing anything but reflecting light - it’s your movement that’s changing the appearance of the colour. This property of holograms, oil and other such surfaces is called “iridescence”. . . .
Papillae are sections of the skin that can be deformed to make a texture bumpy. Even humans possess them (goosebumps) but cannot use them in the manner that cephalopods can. For instance, the use of these cells is how an octopus can wrap itself over a rock and appear jagged or how a squid or cuttlefish can imitate the look of a coral reef by growing miniature towers on its skin. It actually matches the texture of the substrate it chooses.
Finally, the leucophores: According to a paper, published in Nature, cuttlefish and octopuses possess an additional type of reflector cell called a leucophore. They are cells that scatter full spectrum light so that they appear white in a similar way that a polar bear’s fur appears white. Leucophores will also reflect any filtered light shown on them . . . If the water appears blue at a certain depth, the octopuses and cuttlefish can appear blue; if the water appears green, they appear green, and so on and so forth.

\vspace{1em}
\textbf{Question 24}

**There is a sentence that is missing in the paragraph below. Look at the paragraph and decide in which blank (option 1, 2, 3, or 4) the following sentence would best fit.**
**Sentence** : This was years in the making but fast-tracked during the pandemic, when "people started being more mindful about their food", he explained.
**Paragraph** : For millennia, ghee has been a venerated staple of the subcontinental diet, but it fell out of favour a few decades ago when saturated fats were largely considered to be unhealthy. ___(1)___ But more recently, as the thinking around saturated fats is shifting globally, Indians are finding their own way back to this ingredient that is so integral to their cuisine. ___(2)___ For Karmakar, a renewed interest in ghee is emblematic of a return-to-basics movement in India. ___(3)___ This movement is also part of an overall trend towards "slow food". In keeping with the movement's philosophy, ghee can be produced locally (even at home) and has inextricable cultural ties. ___(4)___ At a basic level, ghee is a type of clarified butter believed to have originated in India as a way to preserve butter from going rancid in the hot climate.

\begin{enumerate}[label=\Alph*.]
    \item Option 1
    \item Option 2
    \item Option 3
    \item Option 4
[](https://cracku.in/24-there-is-a-sentence-that-is-missing-in-the-paragra-x-cat-2022-slot-2)
\end{enumerate}
\vspace{1em}
\textbf{Solution:}

We can easily eliminate Options A and B since the given sentence "**This was years in the making but fast-tracked during the pandemic, when "people started being more mindful about their food", he explained** " uses the pronoun "he," which must refer to "Karmakar." The sentence should take the spot of blank 3 since it fits in well with the idea that ghee is part of a "return-to-basics movement in India", as it mentions that the renewed interest in ghee has been developing for a while but was accelerated during the pandemic when people became more conscious about their food choices. Additionally, the idea that this movement is part of an overall trend towards "slow food" is also mentioned in the sentence, as it talks about people being mindful of their food. By contrast, placing the sentence in any of the other blanks would be less logical because it would be introducing information that is not directly related to the main topic being discussed in those parts of the paragraph.
Hence, Option C is the correct choice. 
[](https://cracku.in/downloads/17298)
  *  
  *

\end{problem}

\newpage
\begin{problem}{}{}
% Subject: varc
\textbf{Instructions}

Instructions   

A set of questions accompanies the passage below. Choose the best answer to each question.
Interpretations of the Indian past . . . were inevitably influenced by colonial concerns and interests, and also by prevalent European ideas about history, civilization and the Orient. Orientalist scholars studied the languages and the texts with selected Indian scholars, but made little attempt to understand the worldview of those who were teaching them. The readings, therefore, are something of a disjuncture from the traditional ways of looking at the Indian past. . . .
Orientalism [which we can understand broadly as Western perceptions of the Orient] fuelled the fantasy and the freedom sought by European Romanticism, particularly in its opposition to the more disciplined Neo-Classicism. The cultures of Asia were seen as bringing a new Romantic paradigm. Another Renaissance was anticipated through an acquaintance with the Orient, and this, it was thought, would be different from the earlier Greek Renaissance. It was believed that this Oriental Renaissance would liberate European thought and literature from the increasing focus on discipline and rationality that had followed from the earlier Enlightenment. . . . [The Romantic English poets, Wordsworth and Coleridge,] were apprehensive of the changes introduced by industrialization and turned to nature and to fantasies of the Orient.
However, this enthusiasm gradually changed, to conform with the emphasis later in the nineteenth century on the innate superiority of European civilization. Oriental civilizations were now seen as having once been great but currently in decline. The various phases of Orientalism tended to mould European understanding of the Indian past into a particular pattern. . . . There was an attempt to formulate Indian culture as uniform, such formulations being derived from texts that were given priority. The so-called ‘discovery’ of India was largely through selected literature in Sanskrit. This interpretation tended to emphasize non-historical aspects of Indian culture, for example, the idea of an unchanging continuity of society and religion over 3,000 years; and it was believed that the Indian pattern of life was so concerned with metaphysics and the subtleties of religious belief that little attention was given to the more tangible aspects.
German Romanticism endorsed this image of India, and it became the mystic land for many Europeans, where even the most ordinary actions were imbued with a complex symbolism. This was the genesis of the idea of the spiritual east, and also, incidentally, the refuge of European intellectuals seeking to distance themselves from the changing patterns of their own societies. A dichotomy in values was maintained, Indian values being described as ‘spiritual’ and European values as ‘materialistic’, with little attempt to juxtapose these values with the reality of Indian society. This theme has been even more firmly endorsed by a section of Indian opinion during the last hundred years.
It was a consolation to the Indian intelligentsia for its perceived inability to counter the technical superiority of the west, a superiority viewed as having enabled Europe to colonize Asia and other parts of the world. At the height of anti-colonial nationalism it acted as a salve for having been made a colony of Britain.

\vspace{1em}
\textbf{Question 1}

It can be inferred from the passage that to gain a more accurate view of a nation’s history and culture, scholars should do all of the following EXCEPT:

\begin{enumerate}[label=\Alph*.]
    \item develop an oppositional framework to grasp cultural differences.
    \item examine their own beliefs and biases.
    \item read widely in the country’s literature.
    \item examine the complex reality of that nation’s society.
[](https://cracku.in/1-it-can-be-inferred-from-the-passage-that-to-gain-a-x-cat-2022-slot-3)
\end{enumerate}
\vspace{1em}
\textbf{Solution:}

Option A:  __"_ There was an attempt to formulate Indian culture as uniform, such formulations being derived from texts that were given priority... A _dichotomy in values was maintained, Indian values being described as 'spiritual' and European values as 'materialistic', with little attempt to juxtapose these values with the reality of Indian society_..._"
It can be understood from the above lines of the passage that the author did not approve(even criticize) the position where one should develop an oppositional framework to grasp cultural differences. Thus, this will not bear any fruit in getting a more accurate view of one nation's history and culture. This is the correct option.   

Option B: Throughout the passage, the author criticized the framework the Englishmen adopted to understand India's culture. Thus, the author will support this view that can help give a more accurate picture of a nation's history and culture; hence, this is not the correct option.
Option C: Reading widely into a country's literature without any selection bias will give a more encompassing view of the nation's history; hence, this is also not the correct option.
Option D: This can be rejected on the same ground as option B.
Thus, the correct option is A.

\end{problem}

\newpage
\begin{problem}{}{}
% Subject: varc
\textbf{Instructions}

Instructions   

A set of questions accompanies the passage below. Choose the best answer to each question.
Interpretations of the Indian past . . . were inevitably influenced by colonial concerns and interests, and also by prevalent European ideas about history, civilization and the Orient. Orientalist scholars studied the languages and the texts with selected Indian scholars, but made little attempt to understand the worldview of those who were teaching them. The readings, therefore, are something of a disjuncture from the traditional ways of looking at the Indian past. . . .
Orientalism [which we can understand broadly as Western perceptions of the Orient] fuelled the fantasy and the freedom sought by European Romanticism, particularly in its opposition to the more disciplined Neo-Classicism. The cultures of Asia were seen as bringing a new Romantic paradigm. Another Renaissance was anticipated through an acquaintance with the Orient, and this, it was thought, would be different from the earlier Greek Renaissance. It was believed that this Oriental Renaissance would liberate European thought and literature from the increasing focus on discipline and rationality that had followed from the earlier Enlightenment. . . . [The Romantic English poets, Wordsworth and Coleridge,] were apprehensive of the changes introduced by industrialization and turned to nature and to fantasies of the Orient.
However, this enthusiasm gradually changed, to conform with the emphasis later in the nineteenth century on the innate superiority of European civilization. Oriental civilizations were now seen as having once been great but currently in decline. The various phases of Orientalism tended to mould European understanding of the Indian past into a particular pattern. . . . There was an attempt to formulate Indian culture as uniform, such formulations being derived from texts that were given priority. The so-called ‘discovery’ of India was largely through selected literature in Sanskrit. This interpretation tended to emphasize non-historical aspects of Indian culture, for example, the idea of an unchanging continuity of society and religion over 3,000 years; and it was believed that the Indian pattern of life was so concerned with metaphysics and the subtleties of religious belief that little attention was given to the more tangible aspects.
German Romanticism endorsed this image of India, and it became the mystic land for many Europeans, where even the most ordinary actions were imbued with a complex symbolism. This was the genesis of the idea of the spiritual east, and also, incidentally, the refuge of European intellectuals seeking to distance themselves from the changing patterns of their own societies. A dichotomy in values was maintained, Indian values being described as ‘spiritual’ and European values as ‘materialistic’, with little attempt to juxtapose these values with the reality of Indian society. This theme has been even more firmly endorsed by a section of Indian opinion during the last hundred years.
It was a consolation to the Indian intelligentsia for its perceived inability to counter the technical superiority of the west, a superiority viewed as having enabled Europe to colonize Asia and other parts of the world. At the height of anti-colonial nationalism it acted as a salve for having been made a colony of Britain.

\vspace{1em}
\textbf{Question 2}

It can be inferred from the passage that the author is not likely to support the view that:

\begin{enumerate}[label=\Alph*.]
    \item India’s culture has evolved over the centuries.
    \item the Orientalist view of Asia fired the imagination of some Western poets.
    \item India became a colony although it matched the technical knowledge of the West.
    \item Indian culture acknowledges the material aspects of life.
[](https://cracku.in/2-it-can-be-inferred-from-the-passage-that-the-autho-x-cat-2022-slot-3)
\end{enumerate}
\vspace{1em}
\textbf{Solution:}

Option A: Since the author criticized the uniform view adopted by the Englishmen to understand Indian culture, the author will support the argument presented in this option. Thus, this is not the correct option.
Option B: This can be understood from the starting and concluding lines of the second paragraph, and hence is not the correct option.
Option C: "_It was a consolation to the Indian intelligentsia _for its perceived inability to counter the technical superiority of the west_ , a superiority viewed as having enabled Europe to colonize Asia and other parts of the world_."
Although it can be inferred from the above excerpt that the Indians underestimated their culture and knowledge, it cannot be inferred from this excerpt that they matched the technical understanding of the west.
Thus, this view will not be supported by the author and hence, C is the correct option.
Option D: "_...it was believed that the Indian pattern of life was so concerned with metaphysics and the subtleties of religious belief that little attention was given to**the more tangible aspects**._"
From the above excerpt, it can be inferred that the author disapproved of the Orientalist ignorance of the Indian view towards the materialistic(tangible) aspects. Thus, the author will agree with the view that the Indian culture acknowledges the material aspects of life.
Thus, the correct option is C.

\end{problem}

\newpage
\begin{problem}{}{}
% Subject: varc
\textbf{Instructions}

Instructions   

A set of questions accompanies the passage below. Choose the best answer to each question.
Interpretations of the Indian past . . . were inevitably influenced by colonial concerns and interests, and also by prevalent European ideas about history, civilization and the Orient. Orientalist scholars studied the languages and the texts with selected Indian scholars, but made little attempt to understand the worldview of those who were teaching them. The readings, therefore, are something of a disjuncture from the traditional ways of looking at the Indian past. . . .
Orientalism [which we can understand broadly as Western perceptions of the Orient] fuelled the fantasy and the freedom sought by European Romanticism, particularly in its opposition to the more disciplined Neo-Classicism. The cultures of Asia were seen as bringing a new Romantic paradigm. Another Renaissance was anticipated through an acquaintance with the Orient, and this, it was thought, would be different from the earlier Greek Renaissance. It was believed that this Oriental Renaissance would liberate European thought and literature from the increasing focus on discipline and rationality that had followed from the earlier Enlightenment. . . . [The Romantic English poets, Wordsworth and Coleridge,] were apprehensive of the changes introduced by industrialization and turned to nature and to fantasies of the Orient.
However, this enthusiasm gradually changed, to conform with the emphasis later in the nineteenth century on the innate superiority of European civilization. Oriental civilizations were now seen as having once been great but currently in decline. The various phases of Orientalism tended to mould European understanding of the Indian past into a particular pattern. . . . There was an attempt to formulate Indian culture as uniform, such formulations being derived from texts that were given priority. The so-called ‘discovery’ of India was largely through selected literature in Sanskrit. This interpretation tended to emphasize non-historical aspects of Indian culture, for example, the idea of an unchanging continuity of society and religion over 3,000 years; and it was believed that the Indian pattern of life was so concerned with metaphysics and the subtleties of religious belief that little attention was given to the more tangible aspects.
German Romanticism endorsed this image of India, and it became the mystic land for many Europeans, where even the most ordinary actions were imbued with a complex symbolism. This was the genesis of the idea of the spiritual east, and also, incidentally, the refuge of European intellectuals seeking to distance themselves from the changing patterns of their own societies. A dichotomy in values was maintained, Indian values being described as ‘spiritual’ and European values as ‘materialistic’, with little attempt to juxtapose these values with the reality of Indian society. This theme has been even more firmly endorsed by a section of Indian opinion during the last hundred years.
It was a consolation to the Indian intelligentsia for its perceived inability to counter the technical superiority of the west, a superiority viewed as having enabled Europe to colonize Asia and other parts of the world. At the height of anti-colonial nationalism it acted as a salve for having been made a colony of Britain.

\vspace{1em}
\textbf{Question 3}

In the context of the passage, all of the following statements are true EXCEPT:

\begin{enumerate}[label=\Alph*.]
    \item Indian texts influenced Orientalist scholars.
    \item Orientalist scholarship influenced Indians.
    \item India’s spiritualism served as a salve for European colonisers.
    \item Orientalists’ understanding of Indian history was linked to colonial concerns.
[](https://cracku.in/3-in-the-context-of-the-passage-all-of-the-following-x-cat-2022-slot-3)
\end{enumerate}
\vspace{1em}
\textbf{Solution:}

"_It was a consolation to the Indian intelligentsia for its perceived inability to counter the technical superiority of the west, a superiority viewed as having enabled Europe to colonize Asia and other parts of the world. At the height of anti-colonial nationalism it acted as a salve for having been made a colony of Britain._ "
The author mentioned the reference to being a salve in the last paragraph of the passage(above excerpt). The above excerpt was not regarding colonisers; rather, it refers to the Indian intelligentsia (intellectuals or highly educated people as a group). Thus, it can be inferred that option C is not the correct inference and hence is the correct option.
Throughout the passage, it can be inferred that the Orientalist scholars’ understanding of Indian history and culture was selective, uniform, generalized, and biased. They viewed the Indian culture largely through the lenses of limited and selected literature in Sanskrit. Thus, option A can be inferred and is not the correct option.
"_A dichotomy in values was maintained, Indian values being described as ‘spiritual’ and European values as ‘materialistic’, with little attempt to juxtapose these values with the reality of Indian society. This theme has been even more firmly endorsed by a section of Indian opinion during the last hundred years."_
From the above lines, option B can also be inferred.
Thus, the correct option is C.

\end{problem}

\newpage
\begin{problem}{}{}
% Subject: varc
\textbf{Instructions}

Instructions   

A set of questions accompanies the passage below. Choose the best answer to each question.
Interpretations of the Indian past . . . were inevitably influenced by colonial concerns and interests, and also by prevalent European ideas about history, civilization and the Orient. Orientalist scholars studied the languages and the texts with selected Indian scholars, but made little attempt to understand the worldview of those who were teaching them. The readings, therefore, are something of a disjuncture from the traditional ways of looking at the Indian past. . . .
Orientalism [which we can understand broadly as Western perceptions of the Orient] fuelled the fantasy and the freedom sought by European Romanticism, particularly in its opposition to the more disciplined Neo-Classicism. The cultures of Asia were seen as bringing a new Romantic paradigm. Another Renaissance was anticipated through an acquaintance with the Orient, and this, it was thought, would be different from the earlier Greek Renaissance. It was believed that this Oriental Renaissance would liberate European thought and literature from the increasing focus on discipline and rationality that had followed from the earlier Enlightenment. . . . [The Romantic English poets, Wordsworth and Coleridge,] were apprehensive of the changes introduced by industrialization and turned to nature and to fantasies of the Orient.
However, this enthusiasm gradually changed, to conform with the emphasis later in the nineteenth century on the innate superiority of European civilization. Oriental civilizations were now seen as having once been great but currently in decline. The various phases of Orientalism tended to mould European understanding of the Indian past into a particular pattern. . . . There was an attempt to formulate Indian culture as uniform, such formulations being derived from texts that were given priority. The so-called ‘discovery’ of India was largely through selected literature in Sanskrit. This interpretation tended to emphasize non-historical aspects of Indian culture, for example, the idea of an unchanging continuity of society and religion over 3,000 years; and it was believed that the Indian pattern of life was so concerned with metaphysics and the subtleties of religious belief that little attention was given to the more tangible aspects.
German Romanticism endorsed this image of India, and it became the mystic land for many Europeans, where even the most ordinary actions were imbued with a complex symbolism. This was the genesis of the idea of the spiritual east, and also, incidentally, the refuge of European intellectuals seeking to distance themselves from the changing patterns of their own societies. A dichotomy in values was maintained, Indian values being described as ‘spiritual’ and European values as ‘materialistic’, with little attempt to juxtapose these values with the reality of Indian society. This theme has been even more firmly endorsed by a section of Indian opinion during the last hundred years.
It was a consolation to the Indian intelligentsia for its perceived inability to counter the technical superiority of the west, a superiority viewed as having enabled Europe to colonize Asia and other parts of the world. At the height of anti-colonial nationalism it acted as a salve for having been made a colony of Britain.

\vspace{1em}
\textbf{Question 4}

Which one of the following styles of research is most similar to the Orientalist scholars’ method of understanding Indian history and culture?

\begin{enumerate}[label=\Alph*.]
    \item Studying artefacts excavated at a palace to understand the lifestyle of those who lived there.
    \item Reading 18th century accounts by travellers to India to see how they viewed Indian life and culture of the time.
    \item Reading about the life of early American settlers and later waves of migration to understand the evolution of American culture.
    \item Analysing Hollywood action movies that depict violence and sex to understand contemporary America.
[](https://cracku.in/4-which-one-of-the-following-styles-of-research-is-m-x-cat-2022-slot-3)
\end{enumerate}
\vspace{1em}
\textbf{Solution:}

"_There was an attempt to formulate Indian culture as uniform, such formulations being derived from texts that were given priority. The so-called ‘discovery’ of India was largely through selected literature in Sanskrit. This interpretation tended to emphasize non-historical aspects of Indian culture, for example, the idea of an unchanging continuity of society and religion over 3,000 years_ "
From the above excerpt of the passage, it can be inferred that the Orientalist scholars’ method of understanding Indian history and culture was selective, uniform, generalized, and biased. They viewed the Indian culture largely through the lenses of limited and selected literature in Sanskrit.
Thus, we need to select an option which resembles the same approach.
Out of the four options, only option D uses a very limited understanding(of selected American movies) to form a generalized view of a nation.
Thus, option D is the correct option.

Instructions   

**The passage below is accompanied by a set of questions. Choose the best answer to each question.**
Sociologists working in the Chicago School tradition have focused on how rapid or dramatic social change causes increases in crime. Just as Durkheim, Marx, Toennies, and other European sociologists thought that the rapid changes produced by industrialization and urbanization produced crime and disorder, so too did the Chicago School theorists. The location of the University of Chicago provided an excellent opportunity for Park, Burgess, and McKenzie to study the social ecology of the city. Shaw and McKay found . . . that areas of the city characterized by high levels of social disorganization had higher rates of crime and delinquency.
In the 1920s and 1930s Chicago, like many American cities, experienced considerable immigration. Rapid population growth is a disorganizing influence, but growth resulting from in-migration of very different people is particularly disruptive. Chicago’s in-migrants were both native-born whites and blacks from rural areas and small towns, and foreign immigrants. The heavy industry of cities like Chicago, Detroit, and Pittsburgh drew those seeking opportunities and new lives. Farmers and villagers from America’s hinterland, like their European cousins of whom Durkheim wrote, moved in large numbers into cities. At the start of the twentieth century, Americans were predominately a rural population, but by the century’s mid-point, most lived in urban areas. The social lives of these migrants, as well as those already living in the cities they moved to, were disrupted by the differences between urban and rural life. According to social disorganization theory, until the social ecology of the ‘‘new place’’ can adapt, this rapid change is a criminogenic influence. But most rural migrants, and even many of the foreign immigrants to the city, looked like and eventually spoke the same language as the natives of the cities into which they moved. These similarities allowed for more rapid social integration for these migrants than was the case for African Americans and most foreign immigrants.
In these same decades, America experienced what has been called ‘‘the great migration’’: the massive movement of African Americans out of the rural South and into northern (and some southern) cities. The scale of this migration is one of the most dramatic in human history. These migrants, unlike their white counterparts, were not integrated into the cities they now called home. In fact, most American cities at the end of the twentieth century were characterized by high levels of racial residential segregation . . . Failure to integrate these immigrants, coupled with other forces of social disorganization such as crowding, poverty, and illness, caused crime rates to climb in the cities, particularly in the segregated wards and neighbourhoods where the migrants were forced to live.
Foreign immigrants during this period did not look as dramatically different from the rest of the population as blacks did, but the migrants from eastern and southern Europe who came to American cities did not speak English, and were frequently Catholic, while the native born were mostly Protestant. The combination of rapid population growth with the diversity of those moving into the cities created what the Chicago School sociologists called social disorganization.

\end{problem}

\newpage
\begin{problem}{}{}
% Subject: varc
\textbf{Instructions}

Instructions   

A set of questions accompanies the passage below. Choose the best answer to each question.
Interpretations of the Indian past . . . were inevitably influenced by colonial concerns and interests, and also by prevalent European ideas about history, civilization and the Orient. Orientalist scholars studied the languages and the texts with selected Indian scholars, but made little attempt to understand the worldview of those who were teaching them. The readings, therefore, are something of a disjuncture from the traditional ways of looking at the Indian past. . . .
Orientalism [which we can understand broadly as Western perceptions of the Orient] fuelled the fantasy and the freedom sought by European Romanticism, particularly in its opposition to the more disciplined Neo-Classicism. The cultures of Asia were seen as bringing a new Romantic paradigm. Another Renaissance was anticipated through an acquaintance with the Orient, and this, it was thought, would be different from the earlier Greek Renaissance. It was believed that this Oriental Renaissance would liberate European thought and literature from the increasing focus on discipline and rationality that had followed from the earlier Enlightenment. . . . [The Romantic English poets, Wordsworth and Coleridge,] were apprehensive of the changes introduced by industrialization and turned to nature and to fantasies of the Orient.
However, this enthusiasm gradually changed, to conform with the emphasis later in the nineteenth century on the innate superiority of European civilization. Oriental civilizations were now seen as having once been great but currently in decline. The various phases of Orientalism tended to mould European understanding of the Indian past into a particular pattern. . . . There was an attempt to formulate Indian culture as uniform, such formulations being derived from texts that were given priority. The so-called ‘discovery’ of India was largely through selected literature in Sanskrit. This interpretation tended to emphasize non-historical aspects of Indian culture, for example, the idea of an unchanging continuity of society and religion over 3,000 years; and it was believed that the Indian pattern of life was so concerned with metaphysics and the subtleties of religious belief that little attention was given to the more tangible aspects.
German Romanticism endorsed this image of India, and it became the mystic land for many Europeans, where even the most ordinary actions were imbued with a complex symbolism. This was the genesis of the idea of the spiritual east, and also, incidentally, the refuge of European intellectuals seeking to distance themselves from the changing patterns of their own societies. A dichotomy in values was maintained, Indian values being described as ‘spiritual’ and European values as ‘materialistic’, with little attempt to juxtapose these values with the reality of Indian society. This theme has been even more firmly endorsed by a section of Indian opinion during the last hundred years.
It was a consolation to the Indian intelligentsia for its perceived inability to counter the technical superiority of the west, a superiority viewed as having enabled Europe to colonize Asia and other parts of the world. At the height of anti-colonial nationalism it acted as a salve for having been made a colony of Britain.

\vspace{1em}
\textbf{Question 5}

Which one of the following sets of words/phrases best encapsulates the issues discussed in the passage?

\begin{enumerate}[label=\Alph*.]
    \item Chicago School; Native-born Whites; European immigrants; Poverty
    \item Chicago School; Social organisation; Migration; Crime
    \item Durkheim; Marx; Toennies; Shaw
    \item Rapid population growth; Heavy industry; Segregation; Crime
[](https://cracku.in/5-which-one-of-the-following-sets-of-wordsphrases-be-x-cat-2022-slot-3)
\end{enumerate}
\vspace{1em}
\textbf{Solution:}

The passage starts by stating about the sociologists working in the Chicago school tradition on the causality between social disorganization and crime. Then the author describes the immigration experienced in American cities in the 1920s and 1930s. The author then gives the reason why this led to an increase in crime rates(some examples being _failure to integrate these immigrants, coupled with other forces of social disorganization, such as crowding, poverty, and illness, caused crime rates to climb in the cities, particularly in the segregated wards and neighborhoods where the migrants were forced to live._)
Both options A and C should be eliminated because, in these options, the words social disorganization/ organization or crimes are missing. Compared with D, B is better because the term heavy industry is not a keyword of the passage. Also, more than population growth, migration is the primary reason behind social disorganization.
Thus, the correct option is B.

\end{problem}

\newpage
\begin{problem}{}{}
% Subject: varc
\textbf{Instructions}

Instructions   

A set of questions accompanies the passage below. Choose the best answer to each question.
Interpretations of the Indian past . . . were inevitably influenced by colonial concerns and interests, and also by prevalent European ideas about history, civilization and the Orient. Orientalist scholars studied the languages and the texts with selected Indian scholars, but made little attempt to understand the worldview of those who were teaching them. The readings, therefore, are something of a disjuncture from the traditional ways of looking at the Indian past. . . .
Orientalism [which we can understand broadly as Western perceptions of the Orient] fuelled the fantasy and the freedom sought by European Romanticism, particularly in its opposition to the more disciplined Neo-Classicism. The cultures of Asia were seen as bringing a new Romantic paradigm. Another Renaissance was anticipated through an acquaintance with the Orient, and this, it was thought, would be different from the earlier Greek Renaissance. It was believed that this Oriental Renaissance would liberate European thought and literature from the increasing focus on discipline and rationality that had followed from the earlier Enlightenment. . . . [The Romantic English poets, Wordsworth and Coleridge,] were apprehensive of the changes introduced by industrialization and turned to nature and to fantasies of the Orient.
However, this enthusiasm gradually changed, to conform with the emphasis later in the nineteenth century on the innate superiority of European civilization. Oriental civilizations were now seen as having once been great but currently in decline. The various phases of Orientalism tended to mould European understanding of the Indian past into a particular pattern. . . . There was an attempt to formulate Indian culture as uniform, such formulations being derived from texts that were given priority. The so-called ‘discovery’ of India was largely through selected literature in Sanskrit. This interpretation tended to emphasize non-historical aspects of Indian culture, for example, the idea of an unchanging continuity of society and religion over 3,000 years; and it was believed that the Indian pattern of life was so concerned with metaphysics and the subtleties of religious belief that little attention was given to the more tangible aspects.
German Romanticism endorsed this image of India, and it became the mystic land for many Europeans, where even the most ordinary actions were imbued with a complex symbolism. This was the genesis of the idea of the spiritual east, and also, incidentally, the refuge of European intellectuals seeking to distance themselves from the changing patterns of their own societies. A dichotomy in values was maintained, Indian values being described as ‘spiritual’ and European values as ‘materialistic’, with little attempt to juxtapose these values with the reality of Indian society. This theme has been even more firmly endorsed by a section of Indian opinion during the last hundred years.
It was a consolation to the Indian intelligentsia for its perceived inability to counter the technical superiority of the west, a superiority viewed as having enabled Europe to colonize Asia and other parts of the world. At the height of anti-colonial nationalism it acted as a salve for having been made a colony of Britain.

\vspace{1em}
\textbf{Question 6}

A fundamental conclusion by the author is that:

\begin{enumerate}[label=\Alph*.]
    \item according to European sociologists, crime in America is mainly in Chicago.
    \item the best circumstances for crime to flourish are when there are severe racial disparities.
    \item to prevent crime, it is important to maintain social order through maintaining social segregation.
    \item rapid population growth and demographic diversity give rise to social disorganization that can feed the growth of crime.
[](https://cracku.in/6-a-fundamental-conclusion-by-the-author-is-that-x-cat-2022-slot-3)
\end{enumerate}
\vspace{1em}
\textbf{Solution:}

The above passage is a case study of how rapid or dramatic social change causes relate to increasing crime growth in Chicago. The passage focus on the effects of social disorganization on crime in Chicago.   

Option A: There is no comparison of crime in Chicago with that of crime in other states. Thus, this is not the correct answer.
Option B: The passage focus on the effects of social disorganization on crime in Chicago. It is not specific only to the racial aspect. Thus, this is not the correct option.
Option C: This is a distortion of the passage's main idea. Thus, this is not the correct option.
Option D: This option aptly describes the main conclusion of the passage and hence, is the correct option.
The correct option is D.

\end{problem}

\newpage
\begin{problem}{}{}
% Subject: varc
\textbf{Instructions}

Instructions   

A set of questions accompanies the passage below. Choose the best answer to each question.
Interpretations of the Indian past . . . were inevitably influenced by colonial concerns and interests, and also by prevalent European ideas about history, civilization and the Orient. Orientalist scholars studied the languages and the texts with selected Indian scholars, but made little attempt to understand the worldview of those who were teaching them. The readings, therefore, are something of a disjuncture from the traditional ways of looking at the Indian past. . . .
Orientalism [which we can understand broadly as Western perceptions of the Orient] fuelled the fantasy and the freedom sought by European Romanticism, particularly in its opposition to the more disciplined Neo-Classicism. The cultures of Asia were seen as bringing a new Romantic paradigm. Another Renaissance was anticipated through an acquaintance with the Orient, and this, it was thought, would be different from the earlier Greek Renaissance. It was believed that this Oriental Renaissance would liberate European thought and literature from the increasing focus on discipline and rationality that had followed from the earlier Enlightenment. . . . [The Romantic English poets, Wordsworth and Coleridge,] were apprehensive of the changes introduced by industrialization and turned to nature and to fantasies of the Orient.
However, this enthusiasm gradually changed, to conform with the emphasis later in the nineteenth century on the innate superiority of European civilization. Oriental civilizations were now seen as having once been great but currently in decline. The various phases of Orientalism tended to mould European understanding of the Indian past into a particular pattern. . . . There was an attempt to formulate Indian culture as uniform, such formulations being derived from texts that were given priority. The so-called ‘discovery’ of India was largely through selected literature in Sanskrit. This interpretation tended to emphasize non-historical aspects of Indian culture, for example, the idea of an unchanging continuity of society and religion over 3,000 years; and it was believed that the Indian pattern of life was so concerned with metaphysics and the subtleties of religious belief that little attention was given to the more tangible aspects.
German Romanticism endorsed this image of India, and it became the mystic land for many Europeans, where even the most ordinary actions were imbued with a complex symbolism. This was the genesis of the idea of the spiritual east, and also, incidentally, the refuge of European intellectuals seeking to distance themselves from the changing patterns of their own societies. A dichotomy in values was maintained, Indian values being described as ‘spiritual’ and European values as ‘materialistic’, with little attempt to juxtapose these values with the reality of Indian society. This theme has been even more firmly endorsed by a section of Indian opinion during the last hundred years.
It was a consolation to the Indian intelligentsia for its perceived inability to counter the technical superiority of the west, a superiority viewed as having enabled Europe to colonize Asia and other parts of the world. At the height of anti-colonial nationalism it acted as a salve for having been made a colony of Britain.

\vspace{1em}
\textbf{Question 7}

Which one of the following is not a valid inference from the passage?

\begin{enumerate}[label=\Alph*.]
    \item The failure to integrate in-migrants, along with social problems like poverty, was a significant reason for the rise in crime in American cities.
    \item According to social disorganisation theory, the social integration of African American migrants into Chicago was slower because they were less organised.
    \item The differences between urban and rural lifestyles were crucial factors in the disruption experienced by migrants to American cities.
    \item According to social disorganisation theory, fast-paced social change provides fertile ground for the rapid growth of crime.
[](https://cracku.in/7-which-one-of-the-following-is-not-a-valid-inferenc-x-cat-2022-slot-3)
\end{enumerate}
\vspace{1em}
\textbf{Solution:}

"_Failure to integrate these immigrants, coupled with other forces of social disorganization such as crowding, poverty, and illness, caused crime rates to climb in the cities, particularly in the segregated wards and neighborhoods where the migrants were forced to live._ "
From the above excerpt of the penultimate paragraph of the passage, it can be inferred that poverty and the rise in disorganization contributed to the increase in crime in American cities. Thus, option A is a valid inference and hence can be eliminated.
"_The social lives of these migrants, as well as those already living in cities they moved to, were disrupted by the differences between urban and rural life . According to social disorganization theory, until the social ecology of the ‘‘new place’’ can adapt, this rapid change is a criminogenic influence._ "
From the above excerpt of the second passage, it can be inferred that the difference between urban and rural life contributed to the disruption experienced by the migrants. Also, this rapid change contributes to the rise in crime in these cities. Thus, options C and D are also valid and hence cannot be the answer.
"_These migrants, unlike their white counterparts, were not integrated into the cities they now called home. In fact, most American cities at the end of the twentieth century was characterized by high levels of racial residential segregation._ "
Although the African American migrants faced a high level of racial segregation, it is nowhere mentioned that it was because they were less organised. Thus, option B is not true in the scope of the passage and hence is the correct option.
Thus, the correct option is B.

\end{problem}

\newpage
\begin{problem}{}{}
% Subject: varc
\textbf{Instructions}

Instructions   

A set of questions accompanies the passage below. Choose the best answer to each question.
Interpretations of the Indian past . . . were inevitably influenced by colonial concerns and interests, and also by prevalent European ideas about history, civilization and the Orient. Orientalist scholars studied the languages and the texts with selected Indian scholars, but made little attempt to understand the worldview of those who were teaching them. The readings, therefore, are something of a disjuncture from the traditional ways of looking at the Indian past. . . .
Orientalism [which we can understand broadly as Western perceptions of the Orient] fuelled the fantasy and the freedom sought by European Romanticism, particularly in its opposition to the more disciplined Neo-Classicism. The cultures of Asia were seen as bringing a new Romantic paradigm. Another Renaissance was anticipated through an acquaintance with the Orient, and this, it was thought, would be different from the earlier Greek Renaissance. It was believed that this Oriental Renaissance would liberate European thought and literature from the increasing focus on discipline and rationality that had followed from the earlier Enlightenment. . . . [The Romantic English poets, Wordsworth and Coleridge,] were apprehensive of the changes introduced by industrialization and turned to nature and to fantasies of the Orient.
However, this enthusiasm gradually changed, to conform with the emphasis later in the nineteenth century on the innate superiority of European civilization. Oriental civilizations were now seen as having once been great but currently in decline. The various phases of Orientalism tended to mould European understanding of the Indian past into a particular pattern. . . . There was an attempt to formulate Indian culture as uniform, such formulations being derived from texts that were given priority. The so-called ‘discovery’ of India was largely through selected literature in Sanskrit. This interpretation tended to emphasize non-historical aspects of Indian culture, for example, the idea of an unchanging continuity of society and religion over 3,000 years; and it was believed that the Indian pattern of life was so concerned with metaphysics and the subtleties of religious belief that little attention was given to the more tangible aspects.
German Romanticism endorsed this image of India, and it became the mystic land for many Europeans, where even the most ordinary actions were imbued with a complex symbolism. This was the genesis of the idea of the spiritual east, and also, incidentally, the refuge of European intellectuals seeking to distance themselves from the changing patterns of their own societies. A dichotomy in values was maintained, Indian values being described as ‘spiritual’ and European values as ‘materialistic’, with little attempt to juxtapose these values with the reality of Indian society. This theme has been even more firmly endorsed by a section of Indian opinion during the last hundred years.
It was a consolation to the Indian intelligentsia for its perceived inability to counter the technical superiority of the west, a superiority viewed as having enabled Europe to colonize Asia and other parts of the world. At the height of anti-colonial nationalism it acted as a salve for having been made a colony of Britain.

\vspace{1em}
\textbf{Question 8}

The author notes that, “At the start of the twentieth century, Americans were predominately a rural population, but by the century’s mid-point most lived in urban areas.” Which one of the following statements, if true, does not contradict this statement?

\begin{enumerate}[label=\Alph*.]
    \item Economists have found that throughout the twentieth century, the size of the labour force in America has always been largest in rural areas.
    \item A population census conducted in 1952 showed that more Americans lived in rural areas than in urban ones.
    \item The estimation of per capita income in America in the mid-twentieth century primarily required data from rural areas.
    \item Demographic transition in America in the twentieth century is strongly marked by an out-migration from rural areas.
[](https://cracku.in/8-the-author-notes-that-at-the-start-of-the-twentiet-x-cat-2022-slot-3)
\end{enumerate}
\vspace{1em}
\textbf{Solution:}

Option A: If the workforce size is the largest in the rural area throughout the twenty-first century, then it directly contradicts the theory of the migration of most of the population from rural areas to urban areas. The sample space of the workforce population does not tally with this theory of migration. Thus, this is not the correct option.
Option B: If this population census of 1952 is accurate, it directly nullifies the above migration theory. Thus, this is not the correct option.
Option C: If the data for the estimation of per capita income in the mid-twentieth century primarily required data from the rural areas, then it is not possible that the majority of the population lives in the urban areas. Thus, this is not the correct option.
Option D: If this option is true, it will strengthen the theory of the intermigration of people from rural to urban areas. Thus, this is the correct option.
Thus, the correct option is D.

Instructions   

**The passage below is accompanied by a set of questions. Choose the best answer to each question.**
Nature has all along yielded her flesh to humans. First, we took nature’s materials as food, fibers, and shelter. Then we learned to extract raw materials from her biosphere to create our own new synthetic materials. Now Bios is yielding us her mind—we are taking her logic.
Clockwork logic—the logic of the machines—will only build simple contraptions. Truly complex systems such as a cell, a meadow, an economy, or a brain (natural or artificial) require a rigorous nontechnological logic. We now see that no logic except bio-logic can assemble a thinking device, or even a workable system of any magnitude.
It is an astounding discovery that one can extract the logic of Bios out of biology and have something useful. Although many philosophers in the past have suspected one could abstract the laws of life and apply them elsewhere, it wasn’t until the complexity of computers and human-made systems became as complicated as living things, that it was possible to prove this. It’s eerie how much of life can be transferred. So far, some of the traits of the living that have successfully been transported to mechanical systems are: self-replication, self-governance, limited self-repair, mild evolution, and partial learning.
We have reason to believe yet more can be synthesized and made into something new. Yet at the same time that the logic of Bios is being imported into machines, the logic of Technos is being imported into life. The root of bioengineering is the desire to control the organic long enough to improve it. Domesticated plants and animals are examples of technos-logic applied to life. The wild aromatic root of the Queen Anne’s lace weed has been fine-tuned over generations by selective herb gatherers until it has evolved into a sweet carrot of the garden; the udders of wild bovines have been selectively enlarged in an “unnatural” way to satisfy humans rather than calves. Milk cows and carrots, therefore, are human inventions as much as steam engines and gunpowder are. But milk cows and carrots are more indicative of the kind of inventions humans will make in the future: products that are grown rather than manufactured.
Genetic engineering is precisely what cattle breeders do when they select better strains ofHolsteins, only bioengineers employ more precise and powerful control. While carrot and milk cow breeders had to rely on diffuse organic evolution, modern genetic engineers can use directed artificial evolution—purposeful design—which greatly accelerates improvements.
The overlap of the mechanical and the lifelike increases year by year. Part of this bionic convergence is a matter of words. The meanings of “mechanical” and “life” are both stretching until all complicated things can be perceived as machines, and all self-sustaining machines can be perceived as alive. Yet beyond semantics, two concrete trends are happening: (1)Human-made things are behaving more lifelike, and (2) Life is becoming more engineered. The apparent veil between the organic and the manufactured has crumpled to reveal that the two really are, and have always been, of one being.

\end{problem}

\newpage
\begin{problem}{}{}
% Subject: varc
\textbf{Instructions}

Instructions   

A set of questions accompanies the passage below. Choose the best answer to each question.
Interpretations of the Indian past . . . were inevitably influenced by colonial concerns and interests, and also by prevalent European ideas about history, civilization and the Orient. Orientalist scholars studied the languages and the texts with selected Indian scholars, but made little attempt to understand the worldview of those who were teaching them. The readings, therefore, are something of a disjuncture from the traditional ways of looking at the Indian past. . . .
Orientalism [which we can understand broadly as Western perceptions of the Orient] fuelled the fantasy and the freedom sought by European Romanticism, particularly in its opposition to the more disciplined Neo-Classicism. The cultures of Asia were seen as bringing a new Romantic paradigm. Another Renaissance was anticipated through an acquaintance with the Orient, and this, it was thought, would be different from the earlier Greek Renaissance. It was believed that this Oriental Renaissance would liberate European thought and literature from the increasing focus on discipline and rationality that had followed from the earlier Enlightenment. . . . [The Romantic English poets, Wordsworth and Coleridge,] were apprehensive of the changes introduced by industrialization and turned to nature and to fantasies of the Orient.
However, this enthusiasm gradually changed, to conform with the emphasis later in the nineteenth century on the innate superiority of European civilization. Oriental civilizations were now seen as having once been great but currently in decline. The various phases of Orientalism tended to mould European understanding of the Indian past into a particular pattern. . . . There was an attempt to formulate Indian culture as uniform, such formulations being derived from texts that were given priority. The so-called ‘discovery’ of India was largely through selected literature in Sanskrit. This interpretation tended to emphasize non-historical aspects of Indian culture, for example, the idea of an unchanging continuity of society and religion over 3,000 years; and it was believed that the Indian pattern of life was so concerned with metaphysics and the subtleties of religious belief that little attention was given to the more tangible aspects.
German Romanticism endorsed this image of India, and it became the mystic land for many Europeans, where even the most ordinary actions were imbued with a complex symbolism. This was the genesis of the idea of the spiritual east, and also, incidentally, the refuge of European intellectuals seeking to distance themselves from the changing patterns of their own societies. A dichotomy in values was maintained, Indian values being described as ‘spiritual’ and European values as ‘materialistic’, with little attempt to juxtapose these values with the reality of Indian society. This theme has been even more firmly endorsed by a section of Indian opinion during the last hundred years.
It was a consolation to the Indian intelligentsia for its perceived inability to counter the technical superiority of the west, a superiority viewed as having enabled Europe to colonize Asia and other parts of the world. At the height of anti-colonial nationalism it acted as a salve for having been made a colony of Britain.

\vspace{1em}
\textbf{Question 9}

Which one of the following sets of words/phrases best serves as keywords to thepassage?

\begin{enumerate}[label=\Alph*.]
    \item Complex systems; Carrots; Milk cows; Convergence; Technos-logic
    \item Nature; Computers; Carrots; Milk cows; Genetic engineering
    \item Nature; Bios; Technos; Self-repair; Holsteins
    \item Complex systems; Bio-logic; Bioengineering; Technos-logic; Convergence
[](https://cracku.in/9-which-one-of-the-following-sets-of-wordsphrases-be-x-cat-2022-slot-3)
\end{enumerate}
\vspace{1em}
\textbf{Solution:}

The starting two paragraphs discuss the complexity of the biosphere and how it is impossible to build a thinking device without bio-logic. In the next paragraphs, the author describes how with the increasing complexity of human-made systems(not until it was comparable to living things), it has become possible to transfer these traits into mechanical systems. Examples of these are bioengineering and genetic engineering. Then in the concluding paragraph, the author discusses about the convergence of these two logics(Biologic and Techno logic).
Options B and C do not talk about the conclusion of the passage(convergence of the logics), and hence can be eliminated. Out of options A and D, we should select the option with bio-logic and techno-logic instead of carrots and cows, because the broader idea is about bio and techno, not carrots and cows.
Thus, the correct option is D.

\end{problem}

\newpage
\begin{problem}{}{}
% Subject: varc
\textbf{Instructions}

Instructions   

A set of questions accompanies the passage below. Choose the best answer to each question.
Interpretations of the Indian past . . . were inevitably influenced by colonial concerns and interests, and also by prevalent European ideas about history, civilization and the Orient. Orientalist scholars studied the languages and the texts with selected Indian scholars, but made little attempt to understand the worldview of those who were teaching them. The readings, therefore, are something of a disjuncture from the traditional ways of looking at the Indian past. . . .
Orientalism [which we can understand broadly as Western perceptions of the Orient] fuelled the fantasy and the freedom sought by European Romanticism, particularly in its opposition to the more disciplined Neo-Classicism. The cultures of Asia were seen as bringing a new Romantic paradigm. Another Renaissance was anticipated through an acquaintance with the Orient, and this, it was thought, would be different from the earlier Greek Renaissance. It was believed that this Oriental Renaissance would liberate European thought and literature from the increasing focus on discipline and rationality that had followed from the earlier Enlightenment. . . . [The Romantic English poets, Wordsworth and Coleridge,] were apprehensive of the changes introduced by industrialization and turned to nature and to fantasies of the Orient.
However, this enthusiasm gradually changed, to conform with the emphasis later in the nineteenth century on the innate superiority of European civilization. Oriental civilizations were now seen as having once been great but currently in decline. The various phases of Orientalism tended to mould European understanding of the Indian past into a particular pattern. . . . There was an attempt to formulate Indian culture as uniform, such formulations being derived from texts that were given priority. The so-called ‘discovery’ of India was largely through selected literature in Sanskrit. This interpretation tended to emphasize non-historical aspects of Indian culture, for example, the idea of an unchanging continuity of society and religion over 3,000 years; and it was believed that the Indian pattern of life was so concerned with metaphysics and the subtleties of religious belief that little attention was given to the more tangible aspects.
German Romanticism endorsed this image of India, and it became the mystic land for many Europeans, where even the most ordinary actions were imbued with a complex symbolism. This was the genesis of the idea of the spiritual east, and also, incidentally, the refuge of European intellectuals seeking to distance themselves from the changing patterns of their own societies. A dichotomy in values was maintained, Indian values being described as ‘spiritual’ and European values as ‘materialistic’, with little attempt to juxtapose these values with the reality of Indian society. This theme has been even more firmly endorsed by a section of Indian opinion during the last hundred years.
It was a consolation to the Indian intelligentsia for its perceived inability to counter the technical superiority of the west, a superiority viewed as having enabled Europe to colonize Asia and other parts of the world. At the height of anti-colonial nationalism it acted as a salve for having been made a colony of Britain.

\vspace{1em}
\textbf{Question 10}

The author claims that, “Part of this bionic convergence is a matter of words”. Which one of the following statements best expresses the point being made by the author?

\begin{enumerate}[label=\Alph*.]
    \item “Bios” and “Technos” are both convergent forms of logic, but they generate meanings about the world that are mutually exclusive.
    \item “Mechanical” and “life” are words from different logical systems and are, therefore,fundamentally incompatible in meaning.
    \item A bionic convergence indicates the meeting ground of genetic engineering and artificial intelligence.
    \item “Mechanical” and “life” were earlier seen as opposite in meaning, but the difference between the two is increasingly blurred.
[](https://cracku.in/10-the-author-claims-that-part-of-this-bionic-converg-x-cat-2022-slot-3)
\end{enumerate}
\vspace{1em}
\textbf{Solution:}

_“The overlap of the mechanical and the lifelike increases year by year.**_Part of this bionic convergence is a matter of words_**. The meanings of “mechanical” and “life” are both stretching until all complicated things can be perceived as machines, and all self-sustaining machines can be perceived as alive._”
From the above line, the author tries to show the increasing similarities between ‘mechanical’ and ‘lifelike’ with the passage of time. He states that this increase in similarities will continue till the meanings and the perception of the words become synonymous.
Option A: This option states the opposite of what the author tried to convey and hence is not the correct option.
Option B: This option is distorted and can be rejected on the same grounds as option A.
Option C: This is a distorted inference, and the author did not use the above statement to show the meeting grounds of ‘genetic engineering’ and ‘mechanical engineering’. Thus, this is not the correct option.
Option D: This option aptly expresses the point made by the author in the last paragraph, and hence is the correct option.
Thus, the correct option is D.

\end{problem}

\newpage
\begin{problem}{}{}
% Subject: varc
\textbf{Instructions}

Instructions   

A set of questions accompanies the passage below. Choose the best answer to each question.
Interpretations of the Indian past . . . were inevitably influenced by colonial concerns and interests, and also by prevalent European ideas about history, civilization and the Orient. Orientalist scholars studied the languages and the texts with selected Indian scholars, but made little attempt to understand the worldview of those who were teaching them. The readings, therefore, are something of a disjuncture from the traditional ways of looking at the Indian past. . . .
Orientalism [which we can understand broadly as Western perceptions of the Orient] fuelled the fantasy and the freedom sought by European Romanticism, particularly in its opposition to the more disciplined Neo-Classicism. The cultures of Asia were seen as bringing a new Romantic paradigm. Another Renaissance was anticipated through an acquaintance with the Orient, and this, it was thought, would be different from the earlier Greek Renaissance. It was believed that this Oriental Renaissance would liberate European thought and literature from the increasing focus on discipline and rationality that had followed from the earlier Enlightenment. . . . [The Romantic English poets, Wordsworth and Coleridge,] were apprehensive of the changes introduced by industrialization and turned to nature and to fantasies of the Orient.
However, this enthusiasm gradually changed, to conform with the emphasis later in the nineteenth century on the innate superiority of European civilization. Oriental civilizations were now seen as having once been great but currently in decline. The various phases of Orientalism tended to mould European understanding of the Indian past into a particular pattern. . . . There was an attempt to formulate Indian culture as uniform, such formulations being derived from texts that were given priority. The so-called ‘discovery’ of India was largely through selected literature in Sanskrit. This interpretation tended to emphasize non-historical aspects of Indian culture, for example, the idea of an unchanging continuity of society and religion over 3,000 years; and it was believed that the Indian pattern of life was so concerned with metaphysics and the subtleties of religious belief that little attention was given to the more tangible aspects.
German Romanticism endorsed this image of India, and it became the mystic land for many Europeans, where even the most ordinary actions were imbued with a complex symbolism. This was the genesis of the idea of the spiritual east, and also, incidentally, the refuge of European intellectuals seeking to distance themselves from the changing patterns of their own societies. A dichotomy in values was maintained, Indian values being described as ‘spiritual’ and European values as ‘materialistic’, with little attempt to juxtapose these values with the reality of Indian society. This theme has been even more firmly endorsed by a section of Indian opinion during the last hundred years.
It was a consolation to the Indian intelligentsia for its perceived inability to counter the technical superiority of the west, a superiority viewed as having enabled Europe to colonize Asia and other parts of the world. At the height of anti-colonial nationalism it acted as a salve for having been made a colony of Britain.

\vspace{1em}
\textbf{Question 11}

The author claims that, “The apparent veil between the organic and the manufactured has crumpled to reveal that the two really are, and have always been, of one being.”Which one of the following statements best expresses the point being made by the author here?

\begin{enumerate}[label=\Alph*.]
    \item Organic reality has crumpled under the veil of manufacturing, rendering the apparent and the real as the same being.
    \item The crumpling of the organic veil between apparent and manufactured reality reveals them to have the same being.
    \item Scientific advances are making it increasingly difficult to distinguish between organic reality and manufactured reality.
    \item Apparent reality and organic reality are distinguished by the fact that the former is manufactured.
[](https://cracku.in/11-the-author-claims-that-the-apparent-veil-between-t-x-cat-2022-slot-3)
\end{enumerate}
\vspace{1em}
\textbf{Solution:}

"_Yet beyond semantics, two concrete trends are happening: (1)Human-made things are behaving more lifelike, and (2) Life is becoming more engineered._The apparent veil between the organic and the manufactured has crumpled to reveal that the two really are, and have always been, of one being_.._"
The main argument made by the author in the last paragraph is regarding the increasing similarities between manufactured and organic(lifelike) reality. According to the author, the growing similarities(because of the scientific advances) have distorted the understanding of the realities and have made us think that perhaps these two are and have always been the same.
Option A: This is a distorted inference. It is not that the Organic reality has crumpled under the veil of manufacturing; instead, their meanings are converging mutually. Thus, this is not the correct option.
Option B: This is again a distorted inference. It is not the organic veil that has crumpled; instead, it is the apparent veil. Similarly, in the second half of the option, the organic reality is replaced with the apparent reality.
Option C: This option aptly expresses the main point of the author and is the correct option.   

Option D: The author nowhere stated or implied this, and hence this option can be easily eliminated.
Thus, the correct option is D.

\end{problem}

\newpage
\begin{problem}{}{}
% Subject: varc
\textbf{Instructions}

Instructions   

A set of questions accompanies the passage below. Choose the best answer to each question.
Interpretations of the Indian past . . . were inevitably influenced by colonial concerns and interests, and also by prevalent European ideas about history, civilization and the Orient. Orientalist scholars studied the languages and the texts with selected Indian scholars, but made little attempt to understand the worldview of those who were teaching them. The readings, therefore, are something of a disjuncture from the traditional ways of looking at the Indian past. . . .
Orientalism [which we can understand broadly as Western perceptions of the Orient] fuelled the fantasy and the freedom sought by European Romanticism, particularly in its opposition to the more disciplined Neo-Classicism. The cultures of Asia were seen as bringing a new Romantic paradigm. Another Renaissance was anticipated through an acquaintance with the Orient, and this, it was thought, would be different from the earlier Greek Renaissance. It was believed that this Oriental Renaissance would liberate European thought and literature from the increasing focus on discipline and rationality that had followed from the earlier Enlightenment. . . . [The Romantic English poets, Wordsworth and Coleridge,] were apprehensive of the changes introduced by industrialization and turned to nature and to fantasies of the Orient.
However, this enthusiasm gradually changed, to conform with the emphasis later in the nineteenth century on the innate superiority of European civilization. Oriental civilizations were now seen as having once been great but currently in decline. The various phases of Orientalism tended to mould European understanding of the Indian past into a particular pattern. . . . There was an attempt to formulate Indian culture as uniform, such formulations being derived from texts that were given priority. The so-called ‘discovery’ of India was largely through selected literature in Sanskrit. This interpretation tended to emphasize non-historical aspects of Indian culture, for example, the idea of an unchanging continuity of society and religion over 3,000 years; and it was believed that the Indian pattern of life was so concerned with metaphysics and the subtleties of religious belief that little attention was given to the more tangible aspects.
German Romanticism endorsed this image of India, and it became the mystic land for many Europeans, where even the most ordinary actions were imbued with a complex symbolism. This was the genesis of the idea of the spiritual east, and also, incidentally, the refuge of European intellectuals seeking to distance themselves from the changing patterns of their own societies. A dichotomy in values was maintained, Indian values being described as ‘spiritual’ and European values as ‘materialistic’, with little attempt to juxtapose these values with the reality of Indian society. This theme has been even more firmly endorsed by a section of Indian opinion during the last hundred years.
It was a consolation to the Indian intelligentsia for its perceived inability to counter the technical superiority of the west, a superiority viewed as having enabled Europe to colonize Asia and other parts of the world. At the height of anti-colonial nationalism it acted as a salve for having been made a colony of Britain.

\vspace{1em}
\textbf{Question 12}

None of the following statements is implied by the arguments of the passage, EXCEPT:

\begin{enumerate}[label=\Alph*.]
    \item historically, philosophers have known that the laws of life can be abstracted and applied elsewhere.
    \item genetic engineers and bioengineers are the same insofar as they both seek to force evolution in an artificial way.
    \item the biological realm is as complex as the mechanical one; which is why the logic of Bios is being imported into machines.
    \item purposeful design represents the pinnacle of scientific expertise in the service of human betterment and civilisational progress.
[](https://cracku.in/12-none-of-the-following-statements-is-implied-by-the-x-cat-2022-slot-3)
\end{enumerate}
\vspace{1em}
\textbf{Solution:}

_"Although many philosophers in the past**have suspected** one could abstract the laws of life and apply them elsewhere, it wasn't until the complexity of computers and human-made systems became as complicated as living things that it was possible to prove this._"
Option A can be easily rejected from the above excerpt from the passage. Also, it can be inferred that now(not before), since the complexity of computers and human-made systems are comparable, the logic of Bios can be applied to machines. Although option C seems to convey the same meaning, it generalises the complexity and is a distorted inference.   

The author has nowhere mentioned or implied in the passage that purposeful design represents the pinnacle of scientific expertise in the service of human betterment and civilisational progress. Thus, option D can also be rejected.
_"Genetic engineering is precisely what cattle breeders do when they select better strains of Holsteins, only bioengineers employ more precise and powerful control._ _While carrot and milk cow breeders had to rely on diffuse organic evolution, modern genetic engineers can use directed artificial evolution—purposeful design—which greatly accelerates improvements._ "
From the above excerpt from the penultimate paragraph, it can be inferred that although genetic engineering has less control over the products than bioengineering, they both try to evolve the product artificially. Thus, option B can be inferred from the passage.
Thus, the correct option is B.

Instructions   

**The passage below is accompanied by a set of questions. Choose the best answer to each question.**
As software improves, the people using it become less likely to sharpen their own know-how. Applications that offer lots of prompts and tips are often to blame; simpler, less solicitous programs push people harder to think, act and learn.
Ten years ago, information scientists at Utrecht University in the Netherlands had a group of people carry out complicated analytical and planning tasks using either rudimentary software that provided no assistance or sophisticated software that offered a great deal of aid. The researchers found that the people using the simple software developed better strategies, made fewer mistakes and developed a deeper aptitude for the work. The people using the more advanced software, meanwhile, would often “aimlessly click around” when confronted with a tricky problem. The supposedly helpful software actually short-circuited their thinking and learning.
[According to] philosopher Hubert Dreyfus . . . . our skills get sharper only through practice, when we use them regularly to overcome different sorts of difficult challenges. The goal of modern software, by contrast, is to ease our way through such challenges. Arduous, painstaking work is exactly what programmers are most eager to automate—after all, that is where the immediate efficiency gains tend to lie. In other words, a fundamental tension ripples between the interests of the people doing the automation and the interests of the people doing the work.
Nevertheless, automation’s scope continues to widen. With the rise of electronic health records, physicians increasingly rely on software templates to guide them through patient exams. The programs incorporate valuable checklists and alerts, but they also make medicine more routinized and formulaic—and distance doctors from their patients. . . . Harvard Medical School professor Beth Lown, in a 2012 journal article . . . warned that when doctors become“screen-driven,” following a computer’s prompts rather than “the patient’s narrative thread,” their thinking can become constricted. In the worst cases, they may miss important diagnostic signals. . . .
In a recent paper published in the journal Diagnosis, three medical researchers . . . examined the misdiagnosis of Thomas Eric Duncan, the first person to die of Ebola in the U.S., at Texas Health Presbyterian Hospital Dallas. They argue that the digital templates used by the hospital’s clinicians to record patient information probably helped to induce a kind of tunnel vision. “These highly constrained tools,” the researchers write, “are optimized for data capture but at the expense of sacrificing their utility for appropriate triage and diagnosis, leading users to miss the forest for the trees.” Medical software, they write, is no “replacement for basic history-taking, examination skills, and critical thinking.” . . .
There is an alternative. In “human-centred automation,” the talents of people take precedence. . . . In this model, software plays an essential but secondary role. It takes over routine functions that a human operator has already mastered, issues alerts when unexpected situations arise, provides fresh information that expands the operator’s perspective and counters the biases that often distort human thinking. The technology becomes the expert's partner, not the expert’s replacement.

\end{problem}

\newpage
\begin{problem}{}{}
% Subject: varc
\textbf{Instructions}

Instructions   

A set of questions accompanies the passage below. Choose the best answer to each question.
Interpretations of the Indian past . . . were inevitably influenced by colonial concerns and interests, and also by prevalent European ideas about history, civilization and the Orient. Orientalist scholars studied the languages and the texts with selected Indian scholars, but made little attempt to understand the worldview of those who were teaching them. The readings, therefore, are something of a disjuncture from the traditional ways of looking at the Indian past. . . .
Orientalism [which we can understand broadly as Western perceptions of the Orient] fuelled the fantasy and the freedom sought by European Romanticism, particularly in its opposition to the more disciplined Neo-Classicism. The cultures of Asia were seen as bringing a new Romantic paradigm. Another Renaissance was anticipated through an acquaintance with the Orient, and this, it was thought, would be different from the earlier Greek Renaissance. It was believed that this Oriental Renaissance would liberate European thought and literature from the increasing focus on discipline and rationality that had followed from the earlier Enlightenment. . . . [The Romantic English poets, Wordsworth and Coleridge,] were apprehensive of the changes introduced by industrialization and turned to nature and to fantasies of the Orient.
However, this enthusiasm gradually changed, to conform with the emphasis later in the nineteenth century on the innate superiority of European civilization. Oriental civilizations were now seen as having once been great but currently in decline. The various phases of Orientalism tended to mould European understanding of the Indian past into a particular pattern. . . . There was an attempt to formulate Indian culture as uniform, such formulations being derived from texts that were given priority. The so-called ‘discovery’ of India was largely through selected literature in Sanskrit. This interpretation tended to emphasize non-historical aspects of Indian culture, for example, the idea of an unchanging continuity of society and religion over 3,000 years; and it was believed that the Indian pattern of life was so concerned with metaphysics and the subtleties of religious belief that little attention was given to the more tangible aspects.
German Romanticism endorsed this image of India, and it became the mystic land for many Europeans, where even the most ordinary actions were imbued with a complex symbolism. This was the genesis of the idea of the spiritual east, and also, incidentally, the refuge of European intellectuals seeking to distance themselves from the changing patterns of their own societies. A dichotomy in values was maintained, Indian values being described as ‘spiritual’ and European values as ‘materialistic’, with little attempt to juxtapose these values with the reality of Indian society. This theme has been even more firmly endorsed by a section of Indian opinion during the last hundred years.
It was a consolation to the Indian intelligentsia for its perceived inability to counter the technical superiority of the west, a superiority viewed as having enabled Europe to colonize Asia and other parts of the world. At the height of anti-colonial nationalism it acted as a salve for having been made a colony of Britain.

\vspace{1em}
\textbf{Question 13}

In the Ebola misdiagnosis case, we can infer that doctors probably missed the forest for the trees because:

\begin{enumerate}[label=\Alph*.]
    \item they were led by the data processed by digital templates
    \item the data collected were not sufficient for appropriate triage.
    \item the digital templates forced them to acquire tunnel vision.
    \item they used the wrong type of digital templates for the case.
[](https://cracku.in/13-in-the-ebola-misdiagnosis-case-we-can-infer-that-d-x-cat-2022-slot-3)
\end{enumerate}
\vspace{1em}
\textbf{Solution:}

"_In a recent paper published in the journal Diagnosis, three medical researchers . . . examined the misdiagnosis of Thomas Eric Duncan, the first person to die of Ebola in the U.S., at Texas Health Presbyterian Hospital Dallas. They argue that the digital templates used by the hospital’s clinicians to record patient information probably helped to induce a kind of tunnel vision._ "
From the above expert, we can infer that the misdiagnosis of the ebola patient could have been caused by the digital templates used. The information stored in the templates may have helped induce tunnel vision, so the diagnosis could not capture the virus.
"_"These highly constrained tools", the researchers write, "are optimized for data capture but at the expense of sacrificing their utility for appropriate triage and diagnosis, leading users to miss the forest for the trees". Medical software, they write, is no "replacement for basic history-taking, examination skills, and critical thinking.". . ._ "
The subsequent excerpt gives information about such tools, _which are primarily used to capture/store data at the expense of appropriate diagnosis_ , which leads the users to miss the more important information[to miss the forest for the trees].
Thus, the major culprit, in this case, is the less important data captured by the doctors, which led them to think in the wrong direction.
Since only option A puts the onus on the data processed by the digital templates, this is the correct option.
Options B, C, and D can be eliminated on the basis of the explanation provided above.

\end{problem}

\newpage
\begin{problem}{}{}
% Subject: varc
\textbf{Instructions}

Instructions   

A set of questions accompanies the passage below. Choose the best answer to each question.
Interpretations of the Indian past . . . were inevitably influenced by colonial concerns and interests, and also by prevalent European ideas about history, civilization and the Orient. Orientalist scholars studied the languages and the texts with selected Indian scholars, but made little attempt to understand the worldview of those who were teaching them. The readings, therefore, are something of a disjuncture from the traditional ways of looking at the Indian past. . . .
Orientalism [which we can understand broadly as Western perceptions of the Orient] fuelled the fantasy and the freedom sought by European Romanticism, particularly in its opposition to the more disciplined Neo-Classicism. The cultures of Asia were seen as bringing a new Romantic paradigm. Another Renaissance was anticipated through an acquaintance with the Orient, and this, it was thought, would be different from the earlier Greek Renaissance. It was believed that this Oriental Renaissance would liberate European thought and literature from the increasing focus on discipline and rationality that had followed from the earlier Enlightenment. . . . [The Romantic English poets, Wordsworth and Coleridge,] were apprehensive of the changes introduced by industrialization and turned to nature and to fantasies of the Orient.
However, this enthusiasm gradually changed, to conform with the emphasis later in the nineteenth century on the innate superiority of European civilization. Oriental civilizations were now seen as having once been great but currently in decline. The various phases of Orientalism tended to mould European understanding of the Indian past into a particular pattern. . . . There was an attempt to formulate Indian culture as uniform, such formulations being derived from texts that were given priority. The so-called ‘discovery’ of India was largely through selected literature in Sanskrit. This interpretation tended to emphasize non-historical aspects of Indian culture, for example, the idea of an unchanging continuity of society and religion over 3,000 years; and it was believed that the Indian pattern of life was so concerned with metaphysics and the subtleties of religious belief that little attention was given to the more tangible aspects.
German Romanticism endorsed this image of India, and it became the mystic land for many Europeans, where even the most ordinary actions were imbued with a complex symbolism. This was the genesis of the idea of the spiritual east, and also, incidentally, the refuge of European intellectuals seeking to distance themselves from the changing patterns of their own societies. A dichotomy in values was maintained, Indian values being described as ‘spiritual’ and European values as ‘materialistic’, with little attempt to juxtapose these values with the reality of Indian society. This theme has been even more firmly endorsed by a section of Indian opinion during the last hundred years.
It was a consolation to the Indian intelligentsia for its perceived inability to counter the technical superiority of the west, a superiority viewed as having enabled Europe to colonize Asia and other parts of the world. At the height of anti-colonial nationalism it acted as a salve for having been made a colony of Britain.

\vspace{1em}
\textbf{Question 14}

In the context of the passage, all of the following can be considered examples of human-centered automation EXCEPT:

\begin{enumerate}[label=\Alph*.]
    \item medical software that provides optional feedback on the doctor’s analysis of the medical situation.
    \item a smart-home system that changes the temperature as instructed by the resident.
    \item software that auto-completes text when the user writes an email.
    \item software that offers interpretations when requested by the human operator.
[](https://cracku.in/14-in-the-context-of-the-passage-all-of-the-following-x-cat-2022-slot-3)
\end{enumerate}
\vspace{1em}
\textbf{Solution:}

"_There is an alternative. In “human-centred automation,” the talents of people takeprecedence. . . . In this model, software plays an essential but secondary role. It takes over routine functions that a human operator has already mastered, issues alerts when unexpected situations arise, provides fresh information that expands the operator’s perspective and counters the biases that often distort human thinking. The technology becomes the expert’s partner, not the expert’s replacement._ "
The above excerpt from the passage defines and applies the human-centred approach. This model should have humans as the primary mind, and the software's rule should be restricted only to assistance.
Option A: Since the role of the software is only specified to the feedback on the doctor's analysis, this is a perfect example of the hum-centred approach. Thus, this is not the correct option.
Option B: In this option, too, the role of technology is dependent on the instructions provided by the resident(human), and hence, it is not the correct option.
Option C: Since the software, in this case, operates on its own(auto-completion), it does not take account of human talent and thinking and hence, is not an example of human-centred automation. Thus, this is the correct option.
Option D: In this case, the software only works or provides assistance when the user requests, and hence, it is not the correct option.
Thus, the correct option is C.

\end{problem}

\newpage
\begin{problem}{}{}
% Subject: varc
\textbf{Instructions}

Instructions   

A set of questions accompanies the passage below. Choose the best answer to each question.
Interpretations of the Indian past . . . were inevitably influenced by colonial concerns and interests, and also by prevalent European ideas about history, civilization and the Orient. Orientalist scholars studied the languages and the texts with selected Indian scholars, but made little attempt to understand the worldview of those who were teaching them. The readings, therefore, are something of a disjuncture from the traditional ways of looking at the Indian past. . . .
Orientalism [which we can understand broadly as Western perceptions of the Orient] fuelled the fantasy and the freedom sought by European Romanticism, particularly in its opposition to the more disciplined Neo-Classicism. The cultures of Asia were seen as bringing a new Romantic paradigm. Another Renaissance was anticipated through an acquaintance with the Orient, and this, it was thought, would be different from the earlier Greek Renaissance. It was believed that this Oriental Renaissance would liberate European thought and literature from the increasing focus on discipline and rationality that had followed from the earlier Enlightenment. . . . [The Romantic English poets, Wordsworth and Coleridge,] were apprehensive of the changes introduced by industrialization and turned to nature and to fantasies of the Orient.
However, this enthusiasm gradually changed, to conform with the emphasis later in the nineteenth century on the innate superiority of European civilization. Oriental civilizations were now seen as having once been great but currently in decline. The various phases of Orientalism tended to mould European understanding of the Indian past into a particular pattern. . . . There was an attempt to formulate Indian culture as uniform, such formulations being derived from texts that were given priority. The so-called ‘discovery’ of India was largely through selected literature in Sanskrit. This interpretation tended to emphasize non-historical aspects of Indian culture, for example, the idea of an unchanging continuity of society and religion over 3,000 years; and it was believed that the Indian pattern of life was so concerned with metaphysics and the subtleties of religious belief that little attention was given to the more tangible aspects.
German Romanticism endorsed this image of India, and it became the mystic land for many Europeans, where even the most ordinary actions were imbued with a complex symbolism. This was the genesis of the idea of the spiritual east, and also, incidentally, the refuge of European intellectuals seeking to distance themselves from the changing patterns of their own societies. A dichotomy in values was maintained, Indian values being described as ‘spiritual’ and European values as ‘materialistic’, with little attempt to juxtapose these values with the reality of Indian society. This theme has been even more firmly endorsed by a section of Indian opinion during the last hundred years.
It was a consolation to the Indian intelligentsia for its perceived inability to counter the technical superiority of the west, a superiority viewed as having enabled Europe to colonize Asia and other parts of the world. At the height of anti-colonial nationalism it acted as a salve for having been made a colony of Britain.

\vspace{1em}
\textbf{Question 15}

From the passage, we can infer that the author is apprehensive about the use of sophisticated automation for all of the following reasons EXCEPT that:

\begin{enumerate}[label=\Alph*.]
    \item it stunts the development of its users.
    \item it could mislead people.
    \item computers could replace humans.
    \item it stops users from exercising their minds.
[](https://cracku.in/15-from-the-passage-we-can-infer-that-the-author-is-a-x-cat-2022-slot-3)
\end{enumerate}
\vspace{1em}
\textbf{Solution:}

Option A: From the findings of the information scientists at Utrecht University research(2nd paragraph), it can be concluded that the excessive usage of sophisticated software stunted the thinking and learning of its users. Thus, this is not the correct option.
Option B: From the penultimate paragraph of the passage, it can be inferred that the overemphasis on the data by these 'highly constrained tools' can lead the users astray from their desired target. Thus, this option is also not the correct option.
Option C: Nowhere in the passage is it mentioned or implied that the software can replace human beings. On the contrary, the example of the research paper in the journal Diagnosis points to the limitation of these 'sophisticated' softwares. Thus, this is the correct option.
Option D: This can be inferred from the second and third paragraphs of the passage.
Thus, the correct option is C.

\end{problem}

\newpage
\begin{problem}{}{}
% Subject: varc
\textbf{Instructions}

Instructions   

A set of questions accompanies the passage below. Choose the best answer to each question.
Interpretations of the Indian past . . . were inevitably influenced by colonial concerns and interests, and also by prevalent European ideas about history, civilization and the Orient. Orientalist scholars studied the languages and the texts with selected Indian scholars, but made little attempt to understand the worldview of those who were teaching them. The readings, therefore, are something of a disjuncture from the traditional ways of looking at the Indian past. . . .
Orientalism [which we can understand broadly as Western perceptions of the Orient] fuelled the fantasy and the freedom sought by European Romanticism, particularly in its opposition to the more disciplined Neo-Classicism. The cultures of Asia were seen as bringing a new Romantic paradigm. Another Renaissance was anticipated through an acquaintance with the Orient, and this, it was thought, would be different from the earlier Greek Renaissance. It was believed that this Oriental Renaissance would liberate European thought and literature from the increasing focus on discipline and rationality that had followed from the earlier Enlightenment. . . . [The Romantic English poets, Wordsworth and Coleridge,] were apprehensive of the changes introduced by industrialization and turned to nature and to fantasies of the Orient.
However, this enthusiasm gradually changed, to conform with the emphasis later in the nineteenth century on the innate superiority of European civilization. Oriental civilizations were now seen as having once been great but currently in decline. The various phases of Orientalism tended to mould European understanding of the Indian past into a particular pattern. . . . There was an attempt to formulate Indian culture as uniform, such formulations being derived from texts that were given priority. The so-called ‘discovery’ of India was largely through selected literature in Sanskrit. This interpretation tended to emphasize non-historical aspects of Indian culture, for example, the idea of an unchanging continuity of society and religion over 3,000 years; and it was believed that the Indian pattern of life was so concerned with metaphysics and the subtleties of religious belief that little attention was given to the more tangible aspects.
German Romanticism endorsed this image of India, and it became the mystic land for many Europeans, where even the most ordinary actions were imbued with a complex symbolism. This was the genesis of the idea of the spiritual east, and also, incidentally, the refuge of European intellectuals seeking to distance themselves from the changing patterns of their own societies. A dichotomy in values was maintained, Indian values being described as ‘spiritual’ and European values as ‘materialistic’, with little attempt to juxtapose these values with the reality of Indian society. This theme has been even more firmly endorsed by a section of Indian opinion during the last hundred years.
It was a consolation to the Indian intelligentsia for its perceived inability to counter the technical superiority of the west, a superiority viewed as having enabled Europe to colonize Asia and other parts of the world. At the height of anti-colonial nationalism it acted as a salve for having been made a colony of Britain.

\vspace{1em}
\textbf{Question 16}

It can be inferred that in the Utrecht University experiment, one group of people was“aimlessly clicking around” because:

\begin{enumerate}[label=\Alph*.]
    \item they did not have the skill-set to address complicated tasks.
    \item they were hoping that the software would help carry out the tasks.
    \item the other group was carrying out the tasks more efficiently.
    \item they wanted to avoid making mistakes.
[](https://cracku.in/16-it-can-be-inferred-that-in-the-utrecht-university--x-cat-2022-slot-3)
\end{enumerate}
\vspace{1em}
\textbf{Solution:}

"_The researchers found that the people using the simple software developed better strategies, made fewer mistakes and developed a deeper aptitude for the work._The people using the more advanced software, meanwhile, would often “**aimlessly click around** ” when confronted with a tricky problem._ The supposedly helpful software actually short-circuited their thinking and learning._"
The above excerpt gives the findings of the Utrecht University experiment. The two study groups(one assisted by simple software and the other by a more sophisticated one) show contrasting behaviours. When confronted with a tricky problem, the one with the advanced software would often aimlessly click around the screen to solve the problem. This shows the dependent behaviour of the user on the software. In other words, the users, rather than trying to develop a strategy for the problem, were expecting it to get done by the software.   

Option A: Nowhere in the excerpt was the competency of the users questioned. Instead, it was the effect of the dependency on the software being tested. Thus, this option cannot be inferred.
Option B: The users expected the software to help with the tricky problems. This was described by the phrase "aimlessly click around" in the above excerpt. Thus, this is the correct option.
Option C: The phrase "aimlessly click around" was not used to contrast the strategies adopted by the two study groups; hence, this is not the correct option.
Option D: This again cannot be inferred from the above excerpt.
Thus, the correct option is B.

Instructions   

For the following questions answer them individually

\end{problem}

\newpage
\begin{problem}{}{}
% Subject: varc
\textbf{Instructions}

Instructions   

A set of questions accompanies the passage below. Choose the best answer to each question.
Interpretations of the Indian past . . . were inevitably influenced by colonial concerns and interests, and also by prevalent European ideas about history, civilization and the Orient. Orientalist scholars studied the languages and the texts with selected Indian scholars, but made little attempt to understand the worldview of those who were teaching them. The readings, therefore, are something of a disjuncture from the traditional ways of looking at the Indian past. . . .
Orientalism [which we can understand broadly as Western perceptions of the Orient] fuelled the fantasy and the freedom sought by European Romanticism, particularly in its opposition to the more disciplined Neo-Classicism. The cultures of Asia were seen as bringing a new Romantic paradigm. Another Renaissance was anticipated through an acquaintance with the Orient, and this, it was thought, would be different from the earlier Greek Renaissance. It was believed that this Oriental Renaissance would liberate European thought and literature from the increasing focus on discipline and rationality that had followed from the earlier Enlightenment. . . . [The Romantic English poets, Wordsworth and Coleridge,] were apprehensive of the changes introduced by industrialization and turned to nature and to fantasies of the Orient.
However, this enthusiasm gradually changed, to conform with the emphasis later in the nineteenth century on the innate superiority of European civilization. Oriental civilizations were now seen as having once been great but currently in decline. The various phases of Orientalism tended to mould European understanding of the Indian past into a particular pattern. . . . There was an attempt to formulate Indian culture as uniform, such formulations being derived from texts that were given priority. The so-called ‘discovery’ of India was largely through selected literature in Sanskrit. This interpretation tended to emphasize non-historical aspects of Indian culture, for example, the idea of an unchanging continuity of society and religion over 3,000 years; and it was believed that the Indian pattern of life was so concerned with metaphysics and the subtleties of religious belief that little attention was given to the more tangible aspects.
German Romanticism endorsed this image of India, and it became the mystic land for many Europeans, where even the most ordinary actions were imbued with a complex symbolism. This was the genesis of the idea of the spiritual east, and also, incidentally, the refuge of European intellectuals seeking to distance themselves from the changing patterns of their own societies. A dichotomy in values was maintained, Indian values being described as ‘spiritual’ and European values as ‘materialistic’, with little attempt to juxtapose these values with the reality of Indian society. This theme has been even more firmly endorsed by a section of Indian opinion during the last hundred years.
It was a consolation to the Indian intelligentsia for its perceived inability to counter the technical superiority of the west, a superiority viewed as having enabled Europe to colonize Asia and other parts of the world. At the height of anti-colonial nationalism it acted as a salve for having been made a colony of Britain.

\vspace{1em}
\textbf{Question 17}

The passage given below is followed by four alternate summaries. Choose the option that best captures the essence of the passage.
Tamsin Blanchard, curator of Fashion Open Studio, an initiative by a campaign group showcasing the work of ethical designers says, “We're all drawn to an exquisite piece of embroidery, a colourful textile or even a style of dressing that might have originated from another heritage. [But] this magpie mentality, where all of culture and history is up for grabs as 'inspiration', has accelerated since the proliferation of social media...Where once a fashion student might research the history and traditions of a particular item of clothing with care and respect, we now have a world where images are lifted from image libraries without a care for their cultural significance. It's easier than ever to steal a motif or a craft technique and transfer it on to a piece of clothing that is either mass produced or appears on a runway without credit or compensation to their original communities."

\begin{enumerate}[label=\Alph*.]
    \item Copying an embroidery design or pattern of textile from native communities who own them is tantamount to stealing, and they need to be compensated.
    \item Media has encouraged mass production; images are copied effortlessly without care or concern for the interests of ethnic communities.
    \item Taking fashion ideas from any cultural group without their consent is a form of appropriation without giving due credit, compensation, and respect.
    \item Cultural collaboration is the need of the hour. Beautiful design ideas of indigenous people need to be showcased and shared worldwide.
[](https://cracku.in/17-the-passage-given-below-is-followed-by-four-altern-x-cat-2022-slot-3)
\end{enumerate}
\vspace{1em}
\textbf{Solution:}

The main points of the paragraph are:
i) The copying of fashion ideas unique to particular cultures or heritages is rising in this age of social media.
ii) The original communities are not credited and compensated when their unique ideas are used.
Option A: This is a distorted option. It is generalizing that copying a fashion idea is tantamount to stealing(not specifying whether it is done with or without the consent of the original communities.).Thus, this is not the correct option.
Option B: Again, this is a very general and extreme option. Also, it is a distorted inference that the media has encouraged mass production. Thus, this is also not the correct option.
Option C: Since this includes both the main points, this is the correct option.
Option D: This is a distorted option and does not include the main ideas of the paragraph.
Thus, the correct option is C.

\end{problem}

\newpage
\begin{problem}{}{}
% Subject: varc
\textbf{Instructions}

Instructions   

A set of questions accompanies the passage below. Choose the best answer to each question.
Interpretations of the Indian past . . . were inevitably influenced by colonial concerns and interests, and also by prevalent European ideas about history, civilization and the Orient. Orientalist scholars studied the languages and the texts with selected Indian scholars, but made little attempt to understand the worldview of those who were teaching them. The readings, therefore, are something of a disjuncture from the traditional ways of looking at the Indian past. . . .
Orientalism [which we can understand broadly as Western perceptions of the Orient] fuelled the fantasy and the freedom sought by European Romanticism, particularly in its opposition to the more disciplined Neo-Classicism. The cultures of Asia were seen as bringing a new Romantic paradigm. Another Renaissance was anticipated through an acquaintance with the Orient, and this, it was thought, would be different from the earlier Greek Renaissance. It was believed that this Oriental Renaissance would liberate European thought and literature from the increasing focus on discipline and rationality that had followed from the earlier Enlightenment. . . . [The Romantic English poets, Wordsworth and Coleridge,] were apprehensive of the changes introduced by industrialization and turned to nature and to fantasies of the Orient.
However, this enthusiasm gradually changed, to conform with the emphasis later in the nineteenth century on the innate superiority of European civilization. Oriental civilizations were now seen as having once been great but currently in decline. The various phases of Orientalism tended to mould European understanding of the Indian past into a particular pattern. . . . There was an attempt to formulate Indian culture as uniform, such formulations being derived from texts that were given priority. The so-called ‘discovery’ of India was largely through selected literature in Sanskrit. This interpretation tended to emphasize non-historical aspects of Indian culture, for example, the idea of an unchanging continuity of society and religion over 3,000 years; and it was believed that the Indian pattern of life was so concerned with metaphysics and the subtleties of religious belief that little attention was given to the more tangible aspects.
German Romanticism endorsed this image of India, and it became the mystic land for many Europeans, where even the most ordinary actions were imbued with a complex symbolism. This was the genesis of the idea of the spiritual east, and also, incidentally, the refuge of European intellectuals seeking to distance themselves from the changing patterns of their own societies. A dichotomy in values was maintained, Indian values being described as ‘spiritual’ and European values as ‘materialistic’, with little attempt to juxtapose these values with the reality of Indian society. This theme has been even more firmly endorsed by a section of Indian opinion during the last hundred years.
It was a consolation to the Indian intelligentsia for its perceived inability to counter the technical superiority of the west, a superiority viewed as having enabled Europe to colonize Asia and other parts of the world. At the height of anti-colonial nationalism it acted as a salve for having been made a colony of Britain.

\vspace{1em}
\textbf{Question 18}

The four sentences (labelled 1, 2, 3 and 4) below, when properly sequenced, would yield a coherent paragraph. Decide on the proper sequencing of the order of the sentences and key in the sequence of the four numbers as your answer:
1. If I wanted to sit indoors and read, or play Sonic the Hedgehog on a red-hot SegaMega Drive, I would often be made to feel guilty about not going outside to “enjoy it while it lasts”.  
2. My mum, quite reasonably, wanted me and my sister out of the house, in the sun.  
3. Tales of my mum’s idyllic-sounding childhood in the Sussex countryside, where trees were climbed by 8 am and streams navigated by lunchtime, were passed down to us like folklore.  
4. To an introverted kid, that felt like a threat - and the feeling has stayed with me.
Backspace  
789  
456  
123  
0.-  
Clear All  

Submit
[](https://cracku.in/18-the-four-sentences-labelled-1-2-3-and-4-below-when-x-cat-2022-slot-3)

\begin{enumerate}[label=\Alph*.]
\end{enumerate}
\vspace{1em}
\textbf{Solution:}

A brief reading of sentences 1 and 4 tells us that they form a pair. In sentence 1, the author described how he/she was made to feel guilty about not going outside while staying Indoors. Statement 4 describes how it felt like a threat for an introverted kid. Thus, 1-4 is a pair.  
Now, out of sentences 2 and 3, 2 initiates the topic of discussion[how the author was impelled by her mother to go in the sun]. Statement 3 gives the information about the tales and folklore given to the author to impel. Thus, 2-3 is a pair.  
1 is a direct follow-up to 3; hence, the correct order is 2-3-1-4.

\end{problem}

\newpage
\begin{problem}{}{}
% Subject: varc
\textbf{Instructions}

Instructions   

A set of questions accompanies the passage below. Choose the best answer to each question.
Interpretations of the Indian past . . . were inevitably influenced by colonial concerns and interests, and also by prevalent European ideas about history, civilization and the Orient. Orientalist scholars studied the languages and the texts with selected Indian scholars, but made little attempt to understand the worldview of those who were teaching them. The readings, therefore, are something of a disjuncture from the traditional ways of looking at the Indian past. . . .
Orientalism [which we can understand broadly as Western perceptions of the Orient] fuelled the fantasy and the freedom sought by European Romanticism, particularly in its opposition to the more disciplined Neo-Classicism. The cultures of Asia were seen as bringing a new Romantic paradigm. Another Renaissance was anticipated through an acquaintance with the Orient, and this, it was thought, would be different from the earlier Greek Renaissance. It was believed that this Oriental Renaissance would liberate European thought and literature from the increasing focus on discipline and rationality that had followed from the earlier Enlightenment. . . . [The Romantic English poets, Wordsworth and Coleridge,] were apprehensive of the changes introduced by industrialization and turned to nature and to fantasies of the Orient.
However, this enthusiasm gradually changed, to conform with the emphasis later in the nineteenth century on the innate superiority of European civilization. Oriental civilizations were now seen as having once been great but currently in decline. The various phases of Orientalism tended to mould European understanding of the Indian past into a particular pattern. . . . There was an attempt to formulate Indian culture as uniform, such formulations being derived from texts that were given priority. The so-called ‘discovery’ of India was largely through selected literature in Sanskrit. This interpretation tended to emphasize non-historical aspects of Indian culture, for example, the idea of an unchanging continuity of society and religion over 3,000 years; and it was believed that the Indian pattern of life was so concerned with metaphysics and the subtleties of religious belief that little attention was given to the more tangible aspects.
German Romanticism endorsed this image of India, and it became the mystic land for many Europeans, where even the most ordinary actions were imbued with a complex symbolism. This was the genesis of the idea of the spiritual east, and also, incidentally, the refuge of European intellectuals seeking to distance themselves from the changing patterns of their own societies. A dichotomy in values was maintained, Indian values being described as ‘spiritual’ and European values as ‘materialistic’, with little attempt to juxtapose these values with the reality of Indian society. This theme has been even more firmly endorsed by a section of Indian opinion during the last hundred years.
It was a consolation to the Indian intelligentsia for its perceived inability to counter the technical superiority of the west, a superiority viewed as having enabled Europe to colonize Asia and other parts of the world. At the height of anti-colonial nationalism it acted as a salve for having been made a colony of Britain.

\vspace{1em}
\textbf{Question 19}

There is a sentence that is missing in the paragraph below. Look at the paragraph and decide in which blank (option 1, 2, 3, or 4) the following sentence would best fit.
Sentence: When people socially learn from each other, they often learn without understanding why what they’re copying—the beliefs and behaviours and technologies and know-how—works.
Paragraph: ___(1)___. The dual-inheritance theory ….says....that inheritance is itself an evolutionary system. It has variation. What makes us a new kind of animal, and so different and successful as a species, is we rely heavily on social learning, to the point where socially acquired information is effectively a second line of inheritance, the first being our genes…. ___(2)___. People tend to home in on who seems to be the smartest or most successful person around, as well as what everybody seems to be doing—the majority of people have something worth learning. ___(3)___. When you repeat this process over time, you can get, around the world, cultural packages—beliefs or behaviours or technology or other solutions—that are adapted to the local conditions. People have different psychologies, effectively. ___(4)___.

\begin{enumerate}[label=\Alph*.]
    \item Option 1
    \item Option 2
    \item Option 3
    \item Option 4
[](https://cracku.in/19-there-is-a-sentence-that-is-missing-in-the-paragra-x-cat-2022-slot-3)
\end{enumerate}
\vspace{1em}
\textbf{Solution:}

The sentence would best fit Blank 2 because it ties together the ideas presented in the paragraph. The paragraph describes the dual inheritance theory. In the first two lines, the author states the theory. The given sentence extends the idea and gives the implication of the theory. After blank 2, the author exemplifies the implication given in the sentence(problem sentence) in the following line.  

Thus, the correct option is B.

\end{problem}

\newpage
\begin{problem}{}{}
% Subject: varc
\textbf{Instructions}

Instructions   

A set of questions accompanies the passage below. Choose the best answer to each question.
Interpretations of the Indian past . . . were inevitably influenced by colonial concerns and interests, and also by prevalent European ideas about history, civilization and the Orient. Orientalist scholars studied the languages and the texts with selected Indian scholars, but made little attempt to understand the worldview of those who were teaching them. The readings, therefore, are something of a disjuncture from the traditional ways of looking at the Indian past. . . .
Orientalism [which we can understand broadly as Western perceptions of the Orient] fuelled the fantasy and the freedom sought by European Romanticism, particularly in its opposition to the more disciplined Neo-Classicism. The cultures of Asia were seen as bringing a new Romantic paradigm. Another Renaissance was anticipated through an acquaintance with the Orient, and this, it was thought, would be different from the earlier Greek Renaissance. It was believed that this Oriental Renaissance would liberate European thought and literature from the increasing focus on discipline and rationality that had followed from the earlier Enlightenment. . . . [The Romantic English poets, Wordsworth and Coleridge,] were apprehensive of the changes introduced by industrialization and turned to nature and to fantasies of the Orient.
However, this enthusiasm gradually changed, to conform with the emphasis later in the nineteenth century on the innate superiority of European civilization. Oriental civilizations were now seen as having once been great but currently in decline. The various phases of Orientalism tended to mould European understanding of the Indian past into a particular pattern. . . . There was an attempt to formulate Indian culture as uniform, such formulations being derived from texts that were given priority. The so-called ‘discovery’ of India was largely through selected literature in Sanskrit. This interpretation tended to emphasize non-historical aspects of Indian culture, for example, the idea of an unchanging continuity of society and religion over 3,000 years; and it was believed that the Indian pattern of life was so concerned with metaphysics and the subtleties of religious belief that little attention was given to the more tangible aspects.
German Romanticism endorsed this image of India, and it became the mystic land for many Europeans, where even the most ordinary actions were imbued with a complex symbolism. This was the genesis of the idea of the spiritual east, and also, incidentally, the refuge of European intellectuals seeking to distance themselves from the changing patterns of their own societies. A dichotomy in values was maintained, Indian values being described as ‘spiritual’ and European values as ‘materialistic’, with little attempt to juxtapose these values with the reality of Indian society. This theme has been even more firmly endorsed by a section of Indian opinion during the last hundred years.
It was a consolation to the Indian intelligentsia for its perceived inability to counter the technical superiority of the west, a superiority viewed as having enabled Europe to colonize Asia and other parts of the world. At the height of anti-colonial nationalism it acted as a salve for having been made a colony of Britain.

\vspace{1em}
\textbf{Question 20}

The passage given below is followed by four alternate summaries. Choose the option that best captures the essence of the passage.
To defend the sequence of alphabetisation may seem bizarre, so obvious is its application that it is hard to imagine a reference, catalogue or listing without it. But alphabetical order was not an immediate consequence of the alphabet itself. In the Middle Ages, deference for ecclesiastical tradition left scholars reluctant to categorise things according to the alphabet — to do so would be a rejection of the divine order. The rediscovery of the ancient Greek and Roman classics necessitated more efficient ways of ordering, searching and referencing texts. Government bureaucracy in the 16th and 17th centuries quickened the advance of alphabetical order, bringing with it pigeonholes, notebooks and card indexes.

\begin{enumerate}[label=\Alph*.]
    \item Unlike the alphabet, once the efficacy of the alphabetic sequence became apparent to scholars and administrators, its use became widespread.
    \item The alphabetic order took several centuries to gain common currency because of religious beliefs and a lack of appreciation of its efficacy in the ordering of things.
    \item The ban on the use by scholars of any form of categorisation - but the divinely ordained one - delayed the adoption of the alphabetic sequence by several centuries.
    \item While adoption of the written alphabet was easily accomplished, it took scholars several centuries to accept the alphabetic sequence as a useful tool in their work.
[](https://cracku.in/20-the-passage-given-below-is-followed-by-four-altern-x-cat-2022-slot-3)
\end{enumerate}
\vspace{1em}
\textbf{Solution:}

The main ideas of the paragraph are:
i) The alphabetical order did not directly follow the discovery of alphabets.
ii) Scholars were reluctant to categorize the alphabet in the middle ages because of the fear of rejection of the divine order.
iii) Only after the rediscovery of Greek and Roman classics and the Government bureaucracy in later centuries did the categorization happen.
Option A misses capturing the point of why scholars were reluctant to categorize things according to the alphabet.  
Option C is factually incorrect in mentioning the ban on the use. Option D can be eliminated on the same grounds as option A.   

Option B captures the essence of the main points most aptly and hence, is the best answer.
Thus, the correct option is B.

\end{problem}

\newpage
\begin{problem}{}{}
% Subject: varc
\textbf{Instructions}

Instructions   

A set of questions accompanies the passage below. Choose the best answer to each question.
Interpretations of the Indian past . . . were inevitably influenced by colonial concerns and interests, and also by prevalent European ideas about history, civilization and the Orient. Orientalist scholars studied the languages and the texts with selected Indian scholars, but made little attempt to understand the worldview of those who were teaching them. The readings, therefore, are something of a disjuncture from the traditional ways of looking at the Indian past. . . .
Orientalism [which we can understand broadly as Western perceptions of the Orient] fuelled the fantasy and the freedom sought by European Romanticism, particularly in its opposition to the more disciplined Neo-Classicism. The cultures of Asia were seen as bringing a new Romantic paradigm. Another Renaissance was anticipated through an acquaintance with the Orient, and this, it was thought, would be different from the earlier Greek Renaissance. It was believed that this Oriental Renaissance would liberate European thought and literature from the increasing focus on discipline and rationality that had followed from the earlier Enlightenment. . . . [The Romantic English poets, Wordsworth and Coleridge,] were apprehensive of the changes introduced by industrialization and turned to nature and to fantasies of the Orient.
However, this enthusiasm gradually changed, to conform with the emphasis later in the nineteenth century on the innate superiority of European civilization. Oriental civilizations were now seen as having once been great but currently in decline. The various phases of Orientalism tended to mould European understanding of the Indian past into a particular pattern. . . . There was an attempt to formulate Indian culture as uniform, such formulations being derived from texts that were given priority. The so-called ‘discovery’ of India was largely through selected literature in Sanskrit. This interpretation tended to emphasize non-historical aspects of Indian culture, for example, the idea of an unchanging continuity of society and religion over 3,000 years; and it was believed that the Indian pattern of life was so concerned with metaphysics and the subtleties of religious belief that little attention was given to the more tangible aspects.
German Romanticism endorsed this image of India, and it became the mystic land for many Europeans, where even the most ordinary actions were imbued with a complex symbolism. This was the genesis of the idea of the spiritual east, and also, incidentally, the refuge of European intellectuals seeking to distance themselves from the changing patterns of their own societies. A dichotomy in values was maintained, Indian values being described as ‘spiritual’ and European values as ‘materialistic’, with little attempt to juxtapose these values with the reality of Indian society. This theme has been even more firmly endorsed by a section of Indian opinion during the last hundred years.
It was a consolation to the Indian intelligentsia for its perceived inability to counter the technical superiority of the west, a superiority viewed as having enabled Europe to colonize Asia and other parts of the world. At the height of anti-colonial nationalism it acted as a salve for having been made a colony of Britain.

\vspace{1em}
\textbf{Question 21}

The four sentences (labelled 1, 2, 3 and 4) below, when properly sequenced, would yield a coherent paragraph. Decide on the proper sequencing of the order of the sentences and key in the sequence of the four numbers as your answer:
1. The more we are able to accept that our achievements are largely out of our control, the easier it becomes to understand that our failures, and those of others, are too.  
2. But the raft of recent books about the limits of merit is an important correction to the arrogance of contemporary entitlement and an opportunity to reassert the importance of luck, or grace, in our thinking.  
3. Meritocracy as an organising principle is an inevitable function of a free society, as we are designed to see our achievements as worthy of reward.  
4. And that in turn should increase our humility and the respect with which we treat our fellow citizens, helping ultimately to build a more compassionate society.
Backspace  
789  
456  
123  
0.-  
Clear All  

Submit
[](https://cracku.in/21-the-four-sentences-labelled-1-2-3-and-4-below-when-x-cat-2022-slot-3)

\begin{enumerate}[label=\Alph*.]
\end{enumerate}
\vspace{1em}
\textbf{Solution:}

A brief reading of the sentences suggests that the paragraph is about the limits to considering meritocracy as an organising principle in a free society. Statement 3 introduces the topic at hand by stating the inevitability of meritocracy. Statement 2 initiates the main idea by citing the general idea given in recent books about the limits of merit and its applicability. Thus, 3-2 will form a pair.   
Statement 1 extends this idea by giving the benefits of understanding the limits, and statement 4 follows this by applying this to build a more compassionate society.
Thus, the correct sequence will be 3-2-1-4.

\end{problem}

\newpage
\begin{problem}{}{}
% Subject: varc
\textbf{Instructions}

Instructions   

A set of questions accompanies the passage below. Choose the best answer to each question.
Interpretations of the Indian past . . . were inevitably influenced by colonial concerns and interests, and also by prevalent European ideas about history, civilization and the Orient. Orientalist scholars studied the languages and the texts with selected Indian scholars, but made little attempt to understand the worldview of those who were teaching them. The readings, therefore, are something of a disjuncture from the traditional ways of looking at the Indian past. . . .
Orientalism [which we can understand broadly as Western perceptions of the Orient] fuelled the fantasy and the freedom sought by European Romanticism, particularly in its opposition to the more disciplined Neo-Classicism. The cultures of Asia were seen as bringing a new Romantic paradigm. Another Renaissance was anticipated through an acquaintance with the Orient, and this, it was thought, would be different from the earlier Greek Renaissance. It was believed that this Oriental Renaissance would liberate European thought and literature from the increasing focus on discipline and rationality that had followed from the earlier Enlightenment. . . . [The Romantic English poets, Wordsworth and Coleridge,] were apprehensive of the changes introduced by industrialization and turned to nature and to fantasies of the Orient.
However, this enthusiasm gradually changed, to conform with the emphasis later in the nineteenth century on the innate superiority of European civilization. Oriental civilizations were now seen as having once been great but currently in decline. The various phases of Orientalism tended to mould European understanding of the Indian past into a particular pattern. . . . There was an attempt to formulate Indian culture as uniform, such formulations being derived from texts that were given priority. The so-called ‘discovery’ of India was largely through selected literature in Sanskrit. This interpretation tended to emphasize non-historical aspects of Indian culture, for example, the idea of an unchanging continuity of society and religion over 3,000 years; and it was believed that the Indian pattern of life was so concerned with metaphysics and the subtleties of religious belief that little attention was given to the more tangible aspects.
German Romanticism endorsed this image of India, and it became the mystic land for many Europeans, where even the most ordinary actions were imbued with a complex symbolism. This was the genesis of the idea of the spiritual east, and also, incidentally, the refuge of European intellectuals seeking to distance themselves from the changing patterns of their own societies. A dichotomy in values was maintained, Indian values being described as ‘spiritual’ and European values as ‘materialistic’, with little attempt to juxtapose these values with the reality of Indian society. This theme has been even more firmly endorsed by a section of Indian opinion during the last hundred years.
It was a consolation to the Indian intelligentsia for its perceived inability to counter the technical superiority of the west, a superiority viewed as having enabled Europe to colonize Asia and other parts of the world. At the height of anti-colonial nationalism it acted as a salve for having been made a colony of Britain.

\vspace{1em}
\textbf{Question 22}

The four sentences (labelled 1, 2, 3 and 4) below, when properly sequenced, would yield a coherent paragraph. Decide on the proper sequencing of the order of the sentences and key in the sequence of the four numbers as your answer:
1. Various industrial sectors including retail, transit systems, enterprises, educational institutions, event organizing, finance, travel etc. have now started leveraging these beacons solutions to track and communicate with their customers.  
2. A beacon fixed on to a shop wall enables the retailer to assess the proximity of the customer, and come up with a much targeted or personalized communication like offers, discounts and combos on products in each shelf.  
3. Smartphones or other mobile devices can capture the beacon signals, and distance can be estimated by measuring received signal strength.  
4. Beacons are tiny and inexpensive, micro-location-based technology devices that can send radio frequency signals and notify nearby Bluetooth devices of their presence and transmit information.
Backspace  
789  
456  
123  
0.-  
Clear All  

Submit
[](https://cracku.in/22-the-four-sentences-labelled-1-2-3-and-4-below-when-x-cat-2022-slot-3)

\begin{enumerate}[label=\Alph*.]
\end{enumerate}
\vspace{1em}
\textbf{Solution:}

A brief reading of the sentences suggests that the paragraph is about beacons and their applications. Statement 4 introduces the topic by giving the definition and its utility. Statement 3 sheds light on the working of beacons. Thus, statement 3 will follow statement 4.
Sentences 1 and 2 discuss the utility of these beacons in different industrial sectors, with statement 2 explaining how the beacon would help retailers.
Thus, the correct order will be 4-3-1-2.

\end{problem}

\newpage
\begin{problem}{}{}
% Subject: varc
\textbf{Instructions}

Instructions   

A set of questions accompanies the passage below. Choose the best answer to each question.
Interpretations of the Indian past . . . were inevitably influenced by colonial concerns and interests, and also by prevalent European ideas about history, civilization and the Orient. Orientalist scholars studied the languages and the texts with selected Indian scholars, but made little attempt to understand the worldview of those who were teaching them. The readings, therefore, are something of a disjuncture from the traditional ways of looking at the Indian past. . . .
Orientalism [which we can understand broadly as Western perceptions of the Orient] fuelled the fantasy and the freedom sought by European Romanticism, particularly in its opposition to the more disciplined Neo-Classicism. The cultures of Asia were seen as bringing a new Romantic paradigm. Another Renaissance was anticipated through an acquaintance with the Orient, and this, it was thought, would be different from the earlier Greek Renaissance. It was believed that this Oriental Renaissance would liberate European thought and literature from the increasing focus on discipline and rationality that had followed from the earlier Enlightenment. . . . [The Romantic English poets, Wordsworth and Coleridge,] were apprehensive of the changes introduced by industrialization and turned to nature and to fantasies of the Orient.
However, this enthusiasm gradually changed, to conform with the emphasis later in the nineteenth century on the innate superiority of European civilization. Oriental civilizations were now seen as having once been great but currently in decline. The various phases of Orientalism tended to mould European understanding of the Indian past into a particular pattern. . . . There was an attempt to formulate Indian culture as uniform, such formulations being derived from texts that were given priority. The so-called ‘discovery’ of India was largely through selected literature in Sanskrit. This interpretation tended to emphasize non-historical aspects of Indian culture, for example, the idea of an unchanging continuity of society and religion over 3,000 years; and it was believed that the Indian pattern of life was so concerned with metaphysics and the subtleties of religious belief that little attention was given to the more tangible aspects.
German Romanticism endorsed this image of India, and it became the mystic land for many Europeans, where even the most ordinary actions were imbued with a complex symbolism. This was the genesis of the idea of the spiritual east, and also, incidentally, the refuge of European intellectuals seeking to distance themselves from the changing patterns of their own societies. A dichotomy in values was maintained, Indian values being described as ‘spiritual’ and European values as ‘materialistic’, with little attempt to juxtapose these values with the reality of Indian society. This theme has been even more firmly endorsed by a section of Indian opinion during the last hundred years.
It was a consolation to the Indian intelligentsia for its perceived inability to counter the technical superiority of the west, a superiority viewed as having enabled Europe to colonize Asia and other parts of the world. At the height of anti-colonial nationalism it acted as a salve for having been made a colony of Britain.

\vspace{1em}
\textbf{Question 23}

There is a sentence that is missing in the paragraph below. Look at the paragraph and decide in which blank (option 1, 2, 3, or 4) the following sentence would best fit.
Sentence: This has meant a lot of uncertainty around what a wide-scale return to office might look like in practice.
Paragraph: Bringing workers back to their desks has been a rocky road for employers and employees alike. The evolution of the pandemic has meant that best-laid plans have often not materialised. ___(1)___ The flow of workers back into offices has been more of a trickle than a steady stream. ___(2)___ Yet while plenty of companies are still working through their new policies, some employees across the globe are now back at their desks, whether on a full-time or hybrid basis. ___(3)___ That means we’re beginning to get some clarity on what return-to-office means - what’s working, as well as what has yet to be settled. ___(4)___

\begin{enumerate}[label=\Alph*.]
    \item Option 1
    \item Option 2
    \item Option 3
    \item Option 4
[](https://cracku.in/23-there-is-a-sentence-that-is-missing-in-the-paragra-x-cat-2022-slot-3)
\end{enumerate}
\vspace{1em}
\textbf{Solution:}

The sentence would best fit in Blank 2 because it ties together the ideas presented in the paragraph. The paragraph describes the problems in getting back the employees in the office. In the first two lines, the author mentions this problem. The given sentence gives the effect of the problems and hence will occupy the blank 2. Also, after blank 2, the author changes the flow of the idea and starts describing the present scenario regarding the steps taken by the companies to face this issue.
Thus, the correct option is B.

\end{problem}

\newpage
\begin{problem}{}{}
% Subject: varc
\textbf{Instructions}

Instructions   

A set of questions accompanies the passage below. Choose the best answer to each question.
Interpretations of the Indian past . . . were inevitably influenced by colonial concerns and interests, and also by prevalent European ideas about history, civilization and the Orient. Orientalist scholars studied the languages and the texts with selected Indian scholars, but made little attempt to understand the worldview of those who were teaching them. The readings, therefore, are something of a disjuncture from the traditional ways of looking at the Indian past. . . .
Orientalism [which we can understand broadly as Western perceptions of the Orient] fuelled the fantasy and the freedom sought by European Romanticism, particularly in its opposition to the more disciplined Neo-Classicism. The cultures of Asia were seen as bringing a new Romantic paradigm. Another Renaissance was anticipated through an acquaintance with the Orient, and this, it was thought, would be different from the earlier Greek Renaissance. It was believed that this Oriental Renaissance would liberate European thought and literature from the increasing focus on discipline and rationality that had followed from the earlier Enlightenment. . . . [The Romantic English poets, Wordsworth and Coleridge,] were apprehensive of the changes introduced by industrialization and turned to nature and to fantasies of the Orient.
However, this enthusiasm gradually changed, to conform with the emphasis later in the nineteenth century on the innate superiority of European civilization. Oriental civilizations were now seen as having once been great but currently in decline. The various phases of Orientalism tended to mould European understanding of the Indian past into a particular pattern. . . . There was an attempt to formulate Indian culture as uniform, such formulations being derived from texts that were given priority. The so-called ‘discovery’ of India was largely through selected literature in Sanskrit. This interpretation tended to emphasize non-historical aspects of Indian culture, for example, the idea of an unchanging continuity of society and religion over 3,000 years; and it was believed that the Indian pattern of life was so concerned with metaphysics and the subtleties of religious belief that little attention was given to the more tangible aspects.
German Romanticism endorsed this image of India, and it became the mystic land for many Europeans, where even the most ordinary actions were imbued with a complex symbolism. This was the genesis of the idea of the spiritual east, and also, incidentally, the refuge of European intellectuals seeking to distance themselves from the changing patterns of their own societies. A dichotomy in values was maintained, Indian values being described as ‘spiritual’ and European values as ‘materialistic’, with little attempt to juxtapose these values with the reality of Indian society. This theme has been even more firmly endorsed by a section of Indian opinion during the last hundred years.
It was a consolation to the Indian intelligentsia for its perceived inability to counter the technical superiority of the west, a superiority viewed as having enabled Europe to colonize Asia and other parts of the world. At the height of anti-colonial nationalism it acted as a salve for having been made a colony of Britain.

\vspace{1em}
\textbf{Question 24}

The passage given below is followed by four alternate summaries. Choose the option that best captures the essence of the passage.
“It does seem to me that the job of comedy is to offend, or have the potential to offend, and it cannot be drained of that potential,” Rowan Atkinson said of cancel culture.“Every joke has a victim. That’s the definition of a joke. Someone or something or an idea is made to look ridiculous.” The Netflix star continued, “I think you’ve got to be very, very careful about saying what you’re allowed to make jokes about. You’ve always got to kick up? Really?” He added, “There are lots of extremely smug and self-satisfied people in what would be deemed lower down in society, who also deserve to be pulled up. In a proper free society, you should be allowed to make jokes about absolutely anything.”

\begin{enumerate}[label=\Alph*.]
    \item All jokes target someone and one should be able to joke about anyone in the society, which is inconsistent with cancel culture.
    \item Every joke needs a victim and one needs to include people from lower down the society and not just the upper class.
    \item Victims of jokes must not only be politicians and royalty, but also arrogant people from lower classes should be mentioned by comedians.
    \item Cancel culture does not understand the role and duty of comedians, which is to deride and mock everyone.
[](https://cracku.in/24-the-passage-given-below-is-followed-by-four-altern-x-cat-2022-slot-3)
\end{enumerate}
\vspace{1em}
\textbf{Solution:}

The main ideas of the passage are:
i) The job of a joke is to offend its target(victim) irrespective of its status.
ii) The cancel culture deems it inappropriate to joke about people deemed lower in society.
Option A: This option includes both the main points and hence is the correct answer.
Option B: This is a distorted option. The ideas in the paragraph are not intended to persuade to include people from the lower class in the joke. Thus, this is not the correct answer.
Option C: Again, this is a distorted option and can be eliminated based on the explanation given in option B.
Option D: This is also a distorted option, as nowhere in the passage the duties of a comedian are mentioned. Thus, this is not the correct option.
Thus, the correct option is A.
[](https://cracku.in/downloads/17293)
  *  
  *

\end{problem}

\newpage
\begin{problem}{}{}
% Subject: varc
\textbf{Instructions}

Instructions   

**The passage below is accompanied by four questions. Based on the passage, choose the best answer for each question.**
For early postcolonial literature, the world of the novel was often the nation. Postcolonial novels were usually [concerned with] national questions. Sometimes the whole story of the novel was taken as an allegory of the nation, whether India or Tanzania. This was important for supporting anti-colonial nationalism, but could also be limiting - land-focused and inward-looking.
My new book “Writing Ocean Worlds” explores another kind of world of the novel: not the village or nation, but the Indian Ocean world. The book describes a set of novels in which the Indian Ocean is at the centre of the story. It focuses on the novelists Amitav Ghosh, Abdulrazak Gurnah, Lindsey Collen and Joseph Conrad [who have] centred the Indian Ocean world in the majority of their novels. . . . Their work reveals a world that is outward-looking - full of movement, border-crossing and south-south interconnection. They are all very different - from colonially inclined (Conrad) to radically anti-capitalist (Collen), but together draw on and shape a wider sense of Indian Ocean space through themes, images, metaphors and language. This has the effect of remapping the world in the reader’s mind, as centred in the interconnected global south. . . .
The Indian Ocean world is a term used to describe the very long-lasting connections among the coasts of East Africa, the Arab coasts, and South and East Asia. These connections were made possible by the geography of the Indian Ocean. For much of history, travel by sea was much easier than by land, which meant that port cities very far apart were often more easily connected to each other than to much closer inland cities. Historical and archaeological evidence suggests that what we now call globalisation first appeared in the Indian Ocean. This is the interconnected oceanic world referenced and produced by the novels in my book. . . .
For their part Ghosh, Gurnah, Collen and even Conrad reference a different set of histories and geographies than the ones most commonly found in fiction in English. Those [commonly found ones] are mostly centred in Europe or the US, assume a background of Christianity and whiteness, and mention places like Paris and New York. The novels in [my] book highlight instead a largely Islamic space, feature characters of colour and centralise the ports of Malindi, Mombasa, Aden, Java and Bombay. . . . It is a densely imagined, richly sensory image of a southern cosmopolitan culture which provides for an enlarged sense of place in the world.
This remapping is particularly powerful for the representation of Africa. In the fiction, sailors and travellers are not all European. . . . African, as well as Indian and Arab characters, are traders, nakhodas (dhow ship captains), runaways, villains, missionaries and activists. This does not mean that Indian Ocean Africa is romanticised. Migration is often a matter of force; travel is portrayed as abandonment rather than adventure, freedoms are kept from women and slavery is rife. What it does mean is that the African part of the Indian Ocean world plays an active role in its long, rich history and therefore in that of the wider world.

\vspace{1em}
\textbf{Question 1}

Which one of the following statements is not true about migration in the Indian Ocean world?

\begin{enumerate}[label=\Alph*.]
    \item The Indian Ocean world’s migration networks connected the global north with the global south.
    \item Geographical location rather than geographical proximity determined the choice of destination for migrants.
    \item The Indian Ocean world’s migration networks were shaped by religious and commercial histories of the region.
    \item Migration in the Indian Ocean world was an ambivalent experience.
[](https://cracku.in/1-which-one-of-the-following-statements-is-not-true--x-cat-2023-slot-1)
\end{enumerate}
\vspace{1em}
\textbf{Solution:}

The passage focuses on the interconnectedness within the global south in the context of the Indian Ocean world's migration networks. It emphasizes historical connections between the coasts of East Africa, the Arab coasts, and South and East Asia. The passage does not specifically highlight migration networks connecting the Indian Ocean world with the global north. Instead, it underscores the significance of geographical location, religious histories, and commercial interactions within the region, pointing to a more localized and regional perspective on migration. Therefore, **Option A** is not true according to the passage. Additionally, the passage mentioned the Indian Ocean as “ _a term used to describe the very long-lasting connections among_** _the coasts of East Africa, the Arab coasts, and South and East Asia_** _.”_ and not north and south.
**Option B** is correct as the passage mentions that for much of history, travel by sea in the Indian Ocean was easier than by land, emphasizing the importance of geographical location.
Option C is correct as the passage indicates that the novels in the book draw on and shape a wider sense of Indian Ocean space through themes, images, metaphors, and language, including religious and commercial aspects.
Option D is correct as the passage notes that migration is often portrayed as abandonment rather than adventure, indicating a complex and ambivalent nature of the migration experience in the Indian Ocean world.

\end{problem}

\newpage
\begin{problem}{}{}
% Subject: varc
\textbf{Instructions}

Instructions   

**The passage below is accompanied by four questions. Based on the passage, choose the best answer for each question.**
For early postcolonial literature, the world of the novel was often the nation. Postcolonial novels were usually [concerned with] national questions. Sometimes the whole story of the novel was taken as an allegory of the nation, whether India or Tanzania. This was important for supporting anti-colonial nationalism, but could also be limiting - land-focused and inward-looking.
My new book “Writing Ocean Worlds” explores another kind of world of the novel: not the village or nation, but the Indian Ocean world. The book describes a set of novels in which the Indian Ocean is at the centre of the story. It focuses on the novelists Amitav Ghosh, Abdulrazak Gurnah, Lindsey Collen and Joseph Conrad [who have] centred the Indian Ocean world in the majority of their novels. . . . Their work reveals a world that is outward-looking - full of movement, border-crossing and south-south interconnection. They are all very different - from colonially inclined (Conrad) to radically anti-capitalist (Collen), but together draw on and shape a wider sense of Indian Ocean space through themes, images, metaphors and language. This has the effect of remapping the world in the reader’s mind, as centred in the interconnected global south. . . .
The Indian Ocean world is a term used to describe the very long-lasting connections among the coasts of East Africa, the Arab coasts, and South and East Asia. These connections were made possible by the geography of the Indian Ocean. For much of history, travel by sea was much easier than by land, which meant that port cities very far apart were often more easily connected to each other than to much closer inland cities. Historical and archaeological evidence suggests that what we now call globalisation first appeared in the Indian Ocean. This is the interconnected oceanic world referenced and produced by the novels in my book. . . .
For their part Ghosh, Gurnah, Collen and even Conrad reference a different set of histories and geographies than the ones most commonly found in fiction in English. Those [commonly found ones] are mostly centred in Europe or the US, assume a background of Christianity and whiteness, and mention places like Paris and New York. The novels in [my] book highlight instead a largely Islamic space, feature characters of colour and centralise the ports of Malindi, Mombasa, Aden, Java and Bombay. . . . It is a densely imagined, richly sensory image of a southern cosmopolitan culture which provides for an enlarged sense of place in the world.
This remapping is particularly powerful for the representation of Africa. In the fiction, sailors and travellers are not all European. . . . African, as well as Indian and Arab characters, are traders, nakhodas (dhow ship captains), runaways, villains, missionaries and activists. This does not mean that Indian Ocean Africa is romanticised. Migration is often a matter of force; travel is portrayed as abandonment rather than adventure, freedoms are kept from women and slavery is rife. What it does mean is that the African part of the Indian Ocean world plays an active role in its long, rich history and therefore in that of the wider world.

\vspace{1em}
\textbf{Question 2}

On the basis of the nature of the relationship between the items in each pair below, choose the odd pair out:

\begin{enumerate}[label=\Alph*.]
    \item Indian Ocean novels : Outward-looking
    \item Postcolonial novels : Border-crossing
    \item Indian Ocean world : Slavery
    \item Postcolonial novels : Anti-colonial nationalism
[](https://cracku.in/2-on-the-basis-of-the-nature-of-the-relationship-bet-x-cat-2023-slot-1)
\end{enumerate}
\vspace{1em}
\textbf{Solution:}

Options A, C and D have the following format World/Novels : Characteristic of that particular world/novel. Option B is the odd one out as the characteristic of Border-crossing does not belong to the Postcolonial novels world. 
From the passage, we can infer that the Indian Ocean novels were "outward-looking - full of movement, border-crossing and south-south interconnection". At the same time, they showcased elements of the global south like Slavery, Forced Migration etc. Hence, A and C showcase valid elements of the Indian Ocean Novels World.
On the other hand, postcolonial novels were usually [concerned with] national questions, land-focused and inward-looking. They featured anti-colonial nationalism. Hence, option D is also a valid Theme:Characteristic combination.
However, we note that Border-crossing is an element of the Indian Ocean novel world and not the Postcolonial novel world. Hence, option B is not a valid combination and thus is the odd one out.

\end{problem}

\newpage
\begin{problem}{}{}
% Subject: varc
\textbf{Instructions}

Instructions   

**The passage below is accompanied by four questions. Based on the passage, choose the best answer for each question.**
For early postcolonial literature, the world of the novel was often the nation. Postcolonial novels were usually [concerned with] national questions. Sometimes the whole story of the novel was taken as an allegory of the nation, whether India or Tanzania. This was important for supporting anti-colonial nationalism, but could also be limiting - land-focused and inward-looking.
My new book “Writing Ocean Worlds” explores another kind of world of the novel: not the village or nation, but the Indian Ocean world. The book describes a set of novels in which the Indian Ocean is at the centre of the story. It focuses on the novelists Amitav Ghosh, Abdulrazak Gurnah, Lindsey Collen and Joseph Conrad [who have] centred the Indian Ocean world in the majority of their novels. . . . Their work reveals a world that is outward-looking - full of movement, border-crossing and south-south interconnection. They are all very different - from colonially inclined (Conrad) to radically anti-capitalist (Collen), but together draw on and shape a wider sense of Indian Ocean space through themes, images, metaphors and language. This has the effect of remapping the world in the reader’s mind, as centred in the interconnected global south. . . .
The Indian Ocean world is a term used to describe the very long-lasting connections among the coasts of East Africa, the Arab coasts, and South and East Asia. These connections were made possible by the geography of the Indian Ocean. For much of history, travel by sea was much easier than by land, which meant that port cities very far apart were often more easily connected to each other than to much closer inland cities. Historical and archaeological evidence suggests that what we now call globalisation first appeared in the Indian Ocean. This is the interconnected oceanic world referenced and produced by the novels in my book. . . .
For their part Ghosh, Gurnah, Collen and even Conrad reference a different set of histories and geographies than the ones most commonly found in fiction in English. Those [commonly found ones] are mostly centred in Europe or the US, assume a background of Christianity and whiteness, and mention places like Paris and New York. The novels in [my] book highlight instead a largely Islamic space, feature characters of colour and centralise the ports of Malindi, Mombasa, Aden, Java and Bombay. . . . It is a densely imagined, richly sensory image of a southern cosmopolitan culture which provides for an enlarged sense of place in the world.
This remapping is particularly powerful for the representation of Africa. In the fiction, sailors and travellers are not all European. . . . African, as well as Indian and Arab characters, are traders, nakhodas (dhow ship captains), runaways, villains, missionaries and activists. This does not mean that Indian Ocean Africa is romanticised. Migration is often a matter of force; travel is portrayed as abandonment rather than adventure, freedoms are kept from women and slavery is rife. What it does mean is that the African part of the Indian Ocean world plays an active role in its long, rich history and therefore in that of the wider world.

\vspace{1em}
\textbf{Question 3}

All of the following statements, if true, would weaken the passage’s claim about the relationship between mainstream English-language fiction and Indian Ocean novels EXCEPT:

\begin{enumerate}[label=\Alph*.]
    \item the depiction of Africa in most Indian Ocean novels is driven by a postcolonial nostalgia for an idyllic past
    \item the depiction of Africa in most Indian Ocean novels is driven by an Orientalist imagination of its cultural crudeness.
    \item very few mainstream English-language novels have historically been set in American and European metropolitan centres.
    \item most mainstream English-language novels have historically privileged the Christian, white, male experience of travel and adventure.
[](https://cracku.in/3-all-of-the-following-statements-if-true-would-weak-x-cat-2023-slot-1)
\end{enumerate}
\vspace{1em}
\textbf{Solution:}

_“For their part Ghosh, Gurnah, Collen and even Conrad reference a different set of histories and geographies than the ones most commonly found in fiction in English._**_Those [commonly found ones] are mostly centred in Europe or the US, assume a background of Christianity and whiteness, and mention places like Paris and New York._**_“_
The passage argues that the novels discussed in "Writing Ocean Worlds" diverge from the common representations found in English fiction, which often center on Europe or the US, assume a background of Christianity and whiteness, and mention places like Paris and New York. If **Option D** were true, it would support the passage's claim rather than weaken it. Therefore, Option D is the correct answer.
Through the passage, the author claims that the Indian Ocean novels provide a more realistic picture of the Indian Ocean space, particularly in the representation of Africa. The author claims that the depiction is more authentic and free from Eurocentricity that is seen in other novels.
Option A weakens the passage by contradicting these claims and suggesting that the depiction of Africa is influenced by postcolonial nostalgia.
Option B weakens the passage by suggesting a potential bias or negative stereotyping in the portrayal of Africa in Indian Ocean novels.
Option C weakens the author's claim by disputing that there is eurocentric perspective in other novels.

\end{problem}

\newpage
\begin{problem}{}{}
% Subject: varc
\textbf{Instructions}

Instructions   

**The passage below is accompanied by four questions. Based on the passage, choose the best answer for each question.**
For early postcolonial literature, the world of the novel was often the nation. Postcolonial novels were usually [concerned with] national questions. Sometimes the whole story of the novel was taken as an allegory of the nation, whether India or Tanzania. This was important for supporting anti-colonial nationalism, but could also be limiting - land-focused and inward-looking.
My new book “Writing Ocean Worlds” explores another kind of world of the novel: not the village or nation, but the Indian Ocean world. The book describes a set of novels in which the Indian Ocean is at the centre of the story. It focuses on the novelists Amitav Ghosh, Abdulrazak Gurnah, Lindsey Collen and Joseph Conrad [who have] centred the Indian Ocean world in the majority of their novels. . . . Their work reveals a world that is outward-looking - full of movement, border-crossing and south-south interconnection. They are all very different - from colonially inclined (Conrad) to radically anti-capitalist (Collen), but together draw on and shape a wider sense of Indian Ocean space through themes, images, metaphors and language. This has the effect of remapping the world in the reader’s mind, as centred in the interconnected global south. . . .
The Indian Ocean world is a term used to describe the very long-lasting connections among the coasts of East Africa, the Arab coasts, and South and East Asia. These connections were made possible by the geography of the Indian Ocean. For much of history, travel by sea was much easier than by land, which meant that port cities very far apart were often more easily connected to each other than to much closer inland cities. Historical and archaeological evidence suggests that what we now call globalisation first appeared in the Indian Ocean. This is the interconnected oceanic world referenced and produced by the novels in my book. . . .
For their part Ghosh, Gurnah, Collen and even Conrad reference a different set of histories and geographies than the ones most commonly found in fiction in English. Those [commonly found ones] are mostly centred in Europe or the US, assume a background of Christianity and whiteness, and mention places like Paris and New York. The novels in [my] book highlight instead a largely Islamic space, feature characters of colour and centralise the ports of Malindi, Mombasa, Aden, Java and Bombay. . . . It is a densely imagined, richly sensory image of a southern cosmopolitan culture which provides for an enlarged sense of place in the world.
This remapping is particularly powerful for the representation of Africa. In the fiction, sailors and travellers are not all European. . . . African, as well as Indian and Arab characters, are traders, nakhodas (dhow ship captains), runaways, villains, missionaries and activists. This does not mean that Indian Ocean Africa is romanticised. Migration is often a matter of force; travel is portrayed as abandonment rather than adventure, freedoms are kept from women and slavery is rife. What it does mean is that the African part of the Indian Ocean world plays an active role in its long, rich history and therefore in that of the wider world.

\vspace{1em}
\textbf{Question 4}

All of the following claims contribute to the “remapping” discussed by the passage,  
EXCEPT:

\begin{enumerate}[label=\Alph*.]
    \item Indian Ocean novels have gone beyond the specifics of national concerns to explore rich regional pasts.
    \item the world of early international trade and commerce was not the sole domain of white Europeans.
    \item cosmopolitanism originated in the West and travelled to the East through globalisation.
    \item the global south, as opposed to the global north, was the first centre of globalisation.
[](https://cracku.in/4-all-of-the-following-claims-contribute-to-the-rema-x-cat-2023-slot-1)
\end{enumerate}
\vspace{1em}
\textbf{Solution:}

Option C is the correct answer because it contradicts the idea of "remapping" discussed in the passage. The passage emphasizes that the novels under consideration challenge the common representations found in English fiction, particularly those centered in the West. Option C, suggesting that cosmopolitanism originated in the West and traveled to the East through globalization, aligns with the conventional Western-centric narrative rather than the passage's argument of reshaping perspectives and centralizing the interconnected global south, particularly the Indian Ocean world, as a key space in the reimagined literary landscape.
Option A aligns with the passage's discussion of the novels focusing on the Indian Ocean world, contributing to the "remapping" beyond national concerns.
Option B aligns with the passage's emphasis on the interconnected Indian Ocean world, challenging the Eurocentric perspective on trade and commerce.
Option D supports the passage's claim that historical evidence suggests that globalization first appeared in the Indian Ocean, contributing to the "remapping" of the world's historical and geographical perspectives.

Instructions   

**The passage below is accompanied by four questions. Based on the passage, choose the best answer for each question.**
Many human phenomena and characteristics - such as behaviors, beliefs, economies, genes, incomes, life expectancies, and other things - are influenced both by geographic factors and by non-geographic factors. Geographic factors mean physical and biological factors tied to geographic location, including climate, the distributions of wild plant and animal species, soils, and topography. Non-geographic factors include those factors subsumed under the term culture, other factors subsumed under the term history, and decisions by individual people. . . .
[T]he differences between the current economies of North and South Korea . . . cannot be attributed to the modest environmental differences between [them] . . . They are instead due entirely to the different [government] policies . . . At the opposite extreme, the Inuit and other traditional peoples living north of the Arctic Circle developed warm fur clothes but no agriculture, while equatorial lowland peoples around the world never developed warm fur clothes but often did develop agriculture. The explanation is straightforwardly geographic, rather than a cultural or historical quirk unrelated to geography. . . . Aboriginal Australia remained the sole continent occupied only by hunter/gatherers and with no indigenous farming or herding . . . [Here the] explanation is biogeographic: the Australian continent has no domesticable native animal species and few domesticable native plant species. Instead, the crops and domestic animals that now make Australia a food and wool exporter are all non-native (mainly Eurasian) species such as sheep, wheat, and grapes, brought to Australia by overseas colonists.
Today, no scholar would be silly enough to deny that culture, history, and individual choices play a big role in many human phenomena. Scholars don’t react to cultural, historical, and individual-agent explanations by denouncing “cultural determinism,” “historical determinism,” or “individual determinism,” and then thinking no further. But many scholars do react to any explanation invoking some geographic role, by denouncing “geographic determinism” . . .
Several reasons may underlie this widespread but nonsensical view. One reason is that some geographic explanations advanced a century ago were racist, thereby causing all geographic explanations to become tainted by racist associations in the minds of many scholars other than geographers. But many genetic, historical, psychological, and anthropological explanations advanced a century ago were also racist, yet the validity of newer non-racist genetic etc. explanations is widely accepted today.
Another reason for reflex rejection of geographic explanations is that historians have a tradition, in their discipline, of stressing the role of contingency (a favorite word among historians) based on individual decisions and chance. Often that view is warranted . . . But often, too, that view is unwarranted. The development of warm fur clothes among the Inuit living north of the Arctic Circle was not because one influential Inuit leader persuaded other Inuit in 1783 to adopt warm fur clothes, for no good environmental reason.
A third reason is that geographic explanations usually depend on detailed technical facts of geography and other fields of scholarship . . . Most historians and economists don’t acquire that detailed knowledge as part of the professional training.

\end{problem}

\newpage
\begin{problem}{}{}
% Subject: varc
\textbf{Instructions}

Instructions   

**The passage below is accompanied by four questions. Based on the passage, choose the best answer for each question.**
For early postcolonial literature, the world of the novel was often the nation. Postcolonial novels were usually [concerned with] national questions. Sometimes the whole story of the novel was taken as an allegory of the nation, whether India or Tanzania. This was important for supporting anti-colonial nationalism, but could also be limiting - land-focused and inward-looking.
My new book “Writing Ocean Worlds” explores another kind of world of the novel: not the village or nation, but the Indian Ocean world. The book describes a set of novels in which the Indian Ocean is at the centre of the story. It focuses on the novelists Amitav Ghosh, Abdulrazak Gurnah, Lindsey Collen and Joseph Conrad [who have] centred the Indian Ocean world in the majority of their novels. . . . Their work reveals a world that is outward-looking - full of movement, border-crossing and south-south interconnection. They are all very different - from colonially inclined (Conrad) to radically anti-capitalist (Collen), but together draw on and shape a wider sense of Indian Ocean space through themes, images, metaphors and language. This has the effect of remapping the world in the reader’s mind, as centred in the interconnected global south. . . .
The Indian Ocean world is a term used to describe the very long-lasting connections among the coasts of East Africa, the Arab coasts, and South and East Asia. These connections were made possible by the geography of the Indian Ocean. For much of history, travel by sea was much easier than by land, which meant that port cities very far apart were often more easily connected to each other than to much closer inland cities. Historical and archaeological evidence suggests that what we now call globalisation first appeared in the Indian Ocean. This is the interconnected oceanic world referenced and produced by the novels in my book. . . .
For their part Ghosh, Gurnah, Collen and even Conrad reference a different set of histories and geographies than the ones most commonly found in fiction in English. Those [commonly found ones] are mostly centred in Europe or the US, assume a background of Christianity and whiteness, and mention places like Paris and New York. The novels in [my] book highlight instead a largely Islamic space, feature characters of colour and centralise the ports of Malindi, Mombasa, Aden, Java and Bombay. . . . It is a densely imagined, richly sensory image of a southern cosmopolitan culture which provides for an enlarged sense of place in the world.
This remapping is particularly powerful for the representation of Africa. In the fiction, sailors and travellers are not all European. . . . African, as well as Indian and Arab characters, are traders, nakhodas (dhow ship captains), runaways, villains, missionaries and activists. This does not mean that Indian Ocean Africa is romanticised. Migration is often a matter of force; travel is portrayed as abandonment rather than adventure, freedoms are kept from women and slavery is rife. What it does mean is that the African part of the Indian Ocean world plays an active role in its long, rich history and therefore in that of the wider world.

\vspace{1em}
\textbf{Question 5}

All of the following are advanced by the author as reasons why non-geographers disregard geographic influences on human phenomena EXCEPT their:

\begin{enumerate}[label=\Alph*.]
    \item lingering impressions of past geographic analyses that were politically offensive.
    \item belief in the central role of humans, unrelated to physical surroundings, in influencing phenomena.
    \item disciplinary training which typically does not include technical knowledge of geography.
    \item dismissal of explanations that involve geographical causes for human behaviour.
[](https://cracku.in/5-all-of-the-following-are-advanced-by-the-author-as-x-cat-2023-slot-1)
\end{enumerate}
\vspace{1em}
\textbf{Solution:}

Option D is not explicitly presented by the author as a reason why non-geographers disregard geographic influences. The author suggests that scholars often react negatively to explanations involving a geographic role by denouncing "geographic determinism." However, the specific idea of dismissal is not explicitly outlined in the passage.
The other options on the other hand, can be inferred from the passage:
Option A can be inferred from the following lines: _“ One reason is that some geographic explanations advanced a century ago were racist, thereby causing all geographic explanations to become tainted.”_
Option B can be inferred from the following lines:_“Another reason for reflex rejection of geographic explanations is that historians have a tradition, in their discipline, of stressing the role of contingency (a favorite word among historians) based on individual decisions and chance.”_
Option C can be inferred from the last paragraph of the passage:” _Geographic explanations usually depend on detailed technical facts of geography and other fields of scholarship . . . Most historians and economists don’t acquire that detailed knowledge as part of the professional training.”_

\end{problem}

\newpage
\begin{problem}{}{}
% Subject: varc
\textbf{Instructions}

Instructions   

**The passage below is accompanied by four questions. Based on the passage, choose the best answer for each question.**
For early postcolonial literature, the world of the novel was often the nation. Postcolonial novels were usually [concerned with] national questions. Sometimes the whole story of the novel was taken as an allegory of the nation, whether India or Tanzania. This was important for supporting anti-colonial nationalism, but could also be limiting - land-focused and inward-looking.
My new book “Writing Ocean Worlds” explores another kind of world of the novel: not the village or nation, but the Indian Ocean world. The book describes a set of novels in which the Indian Ocean is at the centre of the story. It focuses on the novelists Amitav Ghosh, Abdulrazak Gurnah, Lindsey Collen and Joseph Conrad [who have] centred the Indian Ocean world in the majority of their novels. . . . Their work reveals a world that is outward-looking - full of movement, border-crossing and south-south interconnection. They are all very different - from colonially inclined (Conrad) to radically anti-capitalist (Collen), but together draw on and shape a wider sense of Indian Ocean space through themes, images, metaphors and language. This has the effect of remapping the world in the reader’s mind, as centred in the interconnected global south. . . .
The Indian Ocean world is a term used to describe the very long-lasting connections among the coasts of East Africa, the Arab coasts, and South and East Asia. These connections were made possible by the geography of the Indian Ocean. For much of history, travel by sea was much easier than by land, which meant that port cities very far apart were often more easily connected to each other than to much closer inland cities. Historical and archaeological evidence suggests that what we now call globalisation first appeared in the Indian Ocean. This is the interconnected oceanic world referenced and produced by the novels in my book. . . .
For their part Ghosh, Gurnah, Collen and even Conrad reference a different set of histories and geographies than the ones most commonly found in fiction in English. Those [commonly found ones] are mostly centred in Europe or the US, assume a background of Christianity and whiteness, and mention places like Paris and New York. The novels in [my] book highlight instead a largely Islamic space, feature characters of colour and centralise the ports of Malindi, Mombasa, Aden, Java and Bombay. . . . It is a densely imagined, richly sensory image of a southern cosmopolitan culture which provides for an enlarged sense of place in the world.
This remapping is particularly powerful for the representation of Africa. In the fiction, sailors and travellers are not all European. . . . African, as well as Indian and Arab characters, are traders, nakhodas (dhow ship captains), runaways, villains, missionaries and activists. This does not mean that Indian Ocean Africa is romanticised. Migration is often a matter of force; travel is portrayed as abandonment rather than adventure, freedoms are kept from women and slavery is rife. What it does mean is that the African part of the Indian Ocean world plays an active role in its long, rich history and therefore in that of the wider world.

\vspace{1em}
\textbf{Question 6}

The author criticises scholars who are not geographers for all of the following reasons EXCEPT:

\begin{enumerate}[label=\Alph*.]
    \item their rejection of the role of biogeographic factors in social and cultural phenomena.
    \item their outdated interpretations of past cultural and historical phenomena.
    \item the importance they place on the role of individual decisions when studying human phenomena
    \item their labelling of geographic explanations as deterministic.
[](https://cracku.in/6-the-author-criticises-scholars-who-are-not-geograp-x-cat-2023-slot-1)
\end{enumerate}
\vspace{1em}
\textbf{Solution:}

The passage does not explicitly mention the criticism of scholars for having outdated interpretations of past cultural and historical phenomena. The primary focus of the author's criticism, as discussed in the passage, centers on scholars' tendencies to dismiss geographic factors, label geographic explanations as deterministic, and associate geographic analyses with past racism.Therefore Option B is the correct answer.

\end{problem}

\newpage
\begin{problem}{}{}
% Subject: varc
\textbf{Instructions}

Instructions   

**The passage below is accompanied by four questions. Based on the passage, choose the best answer for each question.**
For early postcolonial literature, the world of the novel was often the nation. Postcolonial novels were usually [concerned with] national questions. Sometimes the whole story of the novel was taken as an allegory of the nation, whether India or Tanzania. This was important for supporting anti-colonial nationalism, but could also be limiting - land-focused and inward-looking.
My new book “Writing Ocean Worlds” explores another kind of world of the novel: not the village or nation, but the Indian Ocean world. The book describes a set of novels in which the Indian Ocean is at the centre of the story. It focuses on the novelists Amitav Ghosh, Abdulrazak Gurnah, Lindsey Collen and Joseph Conrad [who have] centred the Indian Ocean world in the majority of their novels. . . . Their work reveals a world that is outward-looking - full of movement, border-crossing and south-south interconnection. They are all very different - from colonially inclined (Conrad) to radically anti-capitalist (Collen), but together draw on and shape a wider sense of Indian Ocean space through themes, images, metaphors and language. This has the effect of remapping the world in the reader’s mind, as centred in the interconnected global south. . . .
The Indian Ocean world is a term used to describe the very long-lasting connections among the coasts of East Africa, the Arab coasts, and South and East Asia. These connections were made possible by the geography of the Indian Ocean. For much of history, travel by sea was much easier than by land, which meant that port cities very far apart were often more easily connected to each other than to much closer inland cities. Historical and archaeological evidence suggests that what we now call globalisation first appeared in the Indian Ocean. This is the interconnected oceanic world referenced and produced by the novels in my book. . . .
For their part Ghosh, Gurnah, Collen and even Conrad reference a different set of histories and geographies than the ones most commonly found in fiction in English. Those [commonly found ones] are mostly centred in Europe or the US, assume a background of Christianity and whiteness, and mention places like Paris and New York. The novels in [my] book highlight instead a largely Islamic space, feature characters of colour and centralise the ports of Malindi, Mombasa, Aden, Java and Bombay. . . . It is a densely imagined, richly sensory image of a southern cosmopolitan culture which provides for an enlarged sense of place in the world.
This remapping is particularly powerful for the representation of Africa. In the fiction, sailors and travellers are not all European. . . . African, as well as Indian and Arab characters, are traders, nakhodas (dhow ship captains), runaways, villains, missionaries and activists. This does not mean that Indian Ocean Africa is romanticised. Migration is often a matter of force; travel is portrayed as abandonment rather than adventure, freedoms are kept from women and slavery is rife. What it does mean is that the African part of the Indian Ocean world plays an active role in its long, rich history and therefore in that of the wider world.

\vspace{1em}
\textbf{Question 7}

All of the following can be inferred from the passage EXCEPT:

\begin{enumerate}[label=\Alph*.]
    \item individual dictat and contingency were not the causal factors for the use of fur clothing in some very cold climates.
    \item agricultural practices changed drastically in the Australian continent after it was colonised.
    \item while most human phenomena result from culture and individual choice, some have bio-geographic origins.
    \item several academic studies of human phenomena in the past involved racist interpretations.
[](https://cracku.in/7-all-of-the-following-can-be-inferred-from-the-pass-x-cat-2023-slot-1)
\end{enumerate}
\vspace{1em}
\textbf{Solution:}

Option C cannot be directly inferred from the passage. The passage does discuss the influence of both geographic factors (such as biogeography) and non-geographic factors (culture, history, individual decisions) on human phenomena. However, the passage does not explicitly quantify or compare the prevalence of these influences by stating that _"most human phenomena result from culture and individual choice."_
Option A: The author discusses the development of warm fur clothes among the Inuit living north of the Arctic Circle and asserts that it was not due to a specific individual decision or historical contingency in 1783 but rather a response to environmental factors.
Option B: The author discusses the current state of agricultural practices in Australia, stating that the crops and domestic animals that make Australia a food and wool exporter today are all non-native species (mainly Eurasian) brought to Australia by overseas colonists. The use of the term "non-native" implies a change in agricultural practices from what was originally present in the Australian continent. 
Option D: The author mentions that some geographic explanations advanced a century ago were racist, causing all geographic explanations to become tainted by racist associations in the minds of many scholars.

\end{problem}

\newpage
\begin{problem}{}{}
% Subject: varc
\textbf{Instructions}

Instructions   

**The passage below is accompanied by four questions. Based on the passage, choose the best answer for each question.**
For early postcolonial literature, the world of the novel was often the nation. Postcolonial novels were usually [concerned with] national questions. Sometimes the whole story of the novel was taken as an allegory of the nation, whether India or Tanzania. This was important for supporting anti-colonial nationalism, but could also be limiting - land-focused and inward-looking.
My new book “Writing Ocean Worlds” explores another kind of world of the novel: not the village or nation, but the Indian Ocean world. The book describes a set of novels in which the Indian Ocean is at the centre of the story. It focuses on the novelists Amitav Ghosh, Abdulrazak Gurnah, Lindsey Collen and Joseph Conrad [who have] centred the Indian Ocean world in the majority of their novels. . . . Their work reveals a world that is outward-looking - full of movement, border-crossing and south-south interconnection. They are all very different - from colonially inclined (Conrad) to radically anti-capitalist (Collen), but together draw on and shape a wider sense of Indian Ocean space through themes, images, metaphors and language. This has the effect of remapping the world in the reader’s mind, as centred in the interconnected global south. . . .
The Indian Ocean world is a term used to describe the very long-lasting connections among the coasts of East Africa, the Arab coasts, and South and East Asia. These connections were made possible by the geography of the Indian Ocean. For much of history, travel by sea was much easier than by land, which meant that port cities very far apart were often more easily connected to each other than to much closer inland cities. Historical and archaeological evidence suggests that what we now call globalisation first appeared in the Indian Ocean. This is the interconnected oceanic world referenced and produced by the novels in my book. . . .
For their part Ghosh, Gurnah, Collen and even Conrad reference a different set of histories and geographies than the ones most commonly found in fiction in English. Those [commonly found ones] are mostly centred in Europe or the US, assume a background of Christianity and whiteness, and mention places like Paris and New York. The novels in [my] book highlight instead a largely Islamic space, feature characters of colour and centralise the ports of Malindi, Mombasa, Aden, Java and Bombay. . . . It is a densely imagined, richly sensory image of a southern cosmopolitan culture which provides for an enlarged sense of place in the world.
This remapping is particularly powerful for the representation of Africa. In the fiction, sailors and travellers are not all European. . . . African, as well as Indian and Arab characters, are traders, nakhodas (dhow ship captains), runaways, villains, missionaries and activists. This does not mean that Indian Ocean Africa is romanticised. Migration is often a matter of force; travel is portrayed as abandonment rather than adventure, freedoms are kept from women and slavery is rife. What it does mean is that the African part of the Indian Ocean world plays an active role in its long, rich history and therefore in that of the wider world.

\vspace{1em}
\textbf{Question 8}

The examples of the Inuit and Aboriginal Australians are offered in the passage to show:

\begin{enumerate}[label=\Alph*.]
    \item human resourcefulness across cultures in adapting to their surroundings.
    \item how physical circumstances can dictate human behaviour and cultures.
    \item that despite geographical isolation, traditional societies were self-sufficient and adaptive.
    \item how environmental factors lead to comparatively divergent paths in livelihoods and development.
[](https://cracku.in/8-the-examples-of-the-inuit-and-aboriginal-australia-x-cat-2023-slot-1)
\end{enumerate}
\vspace{1em}
\textbf{Solution:}

Option B is the correct answer because the passage uses examples like the Inuit and Aboriginal Australians to illustrate the influence of physical circumstances, particularly environmental factors, on human behavior and cultural practices. The discussion about the development of warm fur clothes among the Inuit due to the Arctic environment and the absence of indigenous farming in Aboriginal Australia because of the lack of domesticable native species underscores how physical circumstances dictate certain aspects of human behavior and shape cultural adaptations. Therefore, Option B accurately captures the main idea conveyed by the examples provided in the passage.
Option A is not explicitly emphasized in the passage; the focus is more on how environmental factors influence behavior and cultures.
Option C: The passage doesn't explicitly highlight self-sufficiency but rather the impact of specific environmental factors on the development of societies.
Option D is not entirely incorrect, but Option B more precisely captures the emphasis on physical circumstances dictating human behavior and cultures in the context of the examples provided in the passage.

Instructions   

The passage below is accompanied by four questions. Based on the passage, choose the best answer for each question.
RESIDENTS of Lozère, a hilly department in southern France, recite complaints familiar to many rural corners of Europe. In remote hamlets and villages, with names such as Le Bacon and Le Bacon Vieux, mayors grumble about a lack of local schools, jobs, or phone and internet connections. Farmers of grazing animals add another concern: the return of wolves. Eradicated from France last century, the predators are gradually creeping back to more forests and hillsides. “The wolf must be taken in hand,” said an aspiring parliamentarian, Francis Palombi, when pressed by voters in an election campaign early this summer. Tourists enjoy visiting a wolf park in Lozère, but farmers fret over their livestock and their livelihoods. .  
. .
As early as the ninth century, the royal office of the Luparii—wolf-catchers—was created in France to tackle the predators. Those official hunters (and others) completed their job in the 1930s, when the last wolf disappeared from the mainland. Active hunting and improved technology such as rifles in the 19th century, plus the use of poison such as strychnine later on, caused the population collapse. But in the early 1990s the animals reappeared. They crossed the Alps from Italy, upsetting sheep farmers on the French side of the border. Wolves have since spread to areas such as Lozère, delighting environmentalists, who see the predators’ presence as a sign of wider ecological health. Farmers, who say the wolves cause the deaths of thousands of sheep and other grazing animals, are less cheerful. They grumble that green activists and politically correct urban types have allowed the return of an old enemy.
Various factors explain the changes of the past few decades. Rural depopulation is part of the story. In Lozère, for example, farming and a once-flourishing mining industry supported a population of over 140,000 residents in the mid-19th century. Today the department has fewer than 80,000 people, many in its towns. As humans withdraw, forests are expanding. In France, between 1990 and 2015, forest cover increased by an average of 102,000 hectares each year, as more fields were given over to trees. Now, nearly one-third of mainland France is covered by woodland of some sort. The decline of hunting as a sport also means more forests fall quiet. In the mid-to-late 20th century over 2m hunters regularly spent winter weekends tramping in woodland, seeking boars, birds and other prey. Today the Fédération Nationale des Chasseurs, the national body, claims 1.1m people hold hunting licences, though the number of active hunters is probably lower. The mostly protected status of the wolf in Europe—hunting them is now forbidden, other than when occasional culls are sanctioned by the state—plus the efforts of NGOs to track and count the animals, also contribute to the recovery of wolf populations.
As the lupine population of Europe spreads westwards, with occasional reports of wolves seen closer to urban areas, expect to hear of more clashes between farmers and those who celebrate the predators’ return. Farmers’ losses are real, but are not the only economic story. Tourist venues, such as parks where wolves are kept and the animals’ spread is discussed, also generate income and jobs in rural areas.

\end{problem}

\newpage
\begin{problem}{}{}
% Subject: varc
\textbf{Instructions}

Instructions   

**The passage below is accompanied by four questions. Based on the passage, choose the best answer for each question.**
For early postcolonial literature, the world of the novel was often the nation. Postcolonial novels were usually [concerned with] national questions. Sometimes the whole story of the novel was taken as an allegory of the nation, whether India or Tanzania. This was important for supporting anti-colonial nationalism, but could also be limiting - land-focused and inward-looking.
My new book “Writing Ocean Worlds” explores another kind of world of the novel: not the village or nation, but the Indian Ocean world. The book describes a set of novels in which the Indian Ocean is at the centre of the story. It focuses on the novelists Amitav Ghosh, Abdulrazak Gurnah, Lindsey Collen and Joseph Conrad [who have] centred the Indian Ocean world in the majority of their novels. . . . Their work reveals a world that is outward-looking - full of movement, border-crossing and south-south interconnection. They are all very different - from colonially inclined (Conrad) to radically anti-capitalist (Collen), but together draw on and shape a wider sense of Indian Ocean space through themes, images, metaphors and language. This has the effect of remapping the world in the reader’s mind, as centred in the interconnected global south. . . .
The Indian Ocean world is a term used to describe the very long-lasting connections among the coasts of East Africa, the Arab coasts, and South and East Asia. These connections were made possible by the geography of the Indian Ocean. For much of history, travel by sea was much easier than by land, which meant that port cities very far apart were often more easily connected to each other than to much closer inland cities. Historical and archaeological evidence suggests that what we now call globalisation first appeared in the Indian Ocean. This is the interconnected oceanic world referenced and produced by the novels in my book. . . .
For their part Ghosh, Gurnah, Collen and even Conrad reference a different set of histories and geographies than the ones most commonly found in fiction in English. Those [commonly found ones] are mostly centred in Europe or the US, assume a background of Christianity and whiteness, and mention places like Paris and New York. The novels in [my] book highlight instead a largely Islamic space, feature characters of colour and centralise the ports of Malindi, Mombasa, Aden, Java and Bombay. . . . It is a densely imagined, richly sensory image of a southern cosmopolitan culture which provides for an enlarged sense of place in the world.
This remapping is particularly powerful for the representation of Africa. In the fiction, sailors and travellers are not all European. . . . African, as well as Indian and Arab characters, are traders, nakhodas (dhow ship captains), runaways, villains, missionaries and activists. This does not mean that Indian Ocean Africa is romanticised. Migration is often a matter of force; travel is portrayed as abandonment rather than adventure, freedoms are kept from women and slavery is rife. What it does mean is that the African part of the Indian Ocean world plays an active role in its long, rich history and therefore in that of the wider world.

\vspace{1em}
\textbf{Question 9}

Which one of the following has NOT contributed to the growing wolf population in Lozère?

\begin{enumerate}[label=\Alph*.]
    \item A decline in the rural population of Lozère.
    \item An increase in woodlands and forest cover in Lozère.
    \item The shutting down of the royal office of the Luparii.
    \item The granting of a protected status to wolves in Europe.
[](https://cracku.in/9-which-one-of-the-following-has-not-contributed-to--x-cat-2023-slot-1)
\end{enumerate}
\vspace{1em}
\textbf{Solution:}

The passage mentions that as early as the ninth century, the royal office of the Luparii, or wolf-catchers, was created in France to tackle the predators. However, this office became redundant as it had finished it’s job (kill the last wolf). So the resurgence of the wolfs can’t be attributed to it shutting down. The other options on the other hand, can be clearly inferred.
Option A: _“Various factors explain the changes of the past few decades. Rural depopulation is part of the story. In Lozère, for example, farming and a once-flourishing mining industry supported a population of over 140,000 residents in the mid-19th century. Today the department has fewer than 80,000 people, many in its towns. “_
Option B: “ _As humans withdraw, forests are expanding. In France, between 1990 and 2015, forest cover increased by an average of 102,000 hectares each year, as more fields were given over to trees. Now, nearly one-third of mainland France is covered by woodland of some sort. “_
Option D: “ _The mostly protected status of the wolf in Europe—hunting them is now forbidden, other than when occasional culls are sanctioned by the state—plus the efforts of NGOs to track and count the animals, also contribute to the recovery of wolf populations.”_

\end{problem}

\newpage
\begin{problem}{}{}
% Subject: varc
\textbf{Instructions}

Instructions   

**The passage below is accompanied by four questions. Based on the passage, choose the best answer for each question.**
For early postcolonial literature, the world of the novel was often the nation. Postcolonial novels were usually [concerned with] national questions. Sometimes the whole story of the novel was taken as an allegory of the nation, whether India or Tanzania. This was important for supporting anti-colonial nationalism, but could also be limiting - land-focused and inward-looking.
My new book “Writing Ocean Worlds” explores another kind of world of the novel: not the village or nation, but the Indian Ocean world. The book describes a set of novels in which the Indian Ocean is at the centre of the story. It focuses on the novelists Amitav Ghosh, Abdulrazak Gurnah, Lindsey Collen and Joseph Conrad [who have] centred the Indian Ocean world in the majority of their novels. . . . Their work reveals a world that is outward-looking - full of movement, border-crossing and south-south interconnection. They are all very different - from colonially inclined (Conrad) to radically anti-capitalist (Collen), but together draw on and shape a wider sense of Indian Ocean space through themes, images, metaphors and language. This has the effect of remapping the world in the reader’s mind, as centred in the interconnected global south. . . .
The Indian Ocean world is a term used to describe the very long-lasting connections among the coasts of East Africa, the Arab coasts, and South and East Asia. These connections were made possible by the geography of the Indian Ocean. For much of history, travel by sea was much easier than by land, which meant that port cities very far apart were often more easily connected to each other than to much closer inland cities. Historical and archaeological evidence suggests that what we now call globalisation first appeared in the Indian Ocean. This is the interconnected oceanic world referenced and produced by the novels in my book. . . .
For their part Ghosh, Gurnah, Collen and even Conrad reference a different set of histories and geographies than the ones most commonly found in fiction in English. Those [commonly found ones] are mostly centred in Europe or the US, assume a background of Christianity and whiteness, and mention places like Paris and New York. The novels in [my] book highlight instead a largely Islamic space, feature characters of colour and centralise the ports of Malindi, Mombasa, Aden, Java and Bombay. . . . It is a densely imagined, richly sensory image of a southern cosmopolitan culture which provides for an enlarged sense of place in the world.
This remapping is particularly powerful for the representation of Africa. In the fiction, sailors and travellers are not all European. . . . African, as well as Indian and Arab characters, are traders, nakhodas (dhow ship captains), runaways, villains, missionaries and activists. This does not mean that Indian Ocean Africa is romanticised. Migration is often a matter of force; travel is portrayed as abandonment rather than adventure, freedoms are kept from women and slavery is rife. What it does mean is that the African part of the Indian Ocean world plays an active role in its long, rich history and therefore in that of the wider world.

\vspace{1em}
\textbf{Question 10}

The inhabitants of Lozère have to grapple with all of the following problems, EXCEPT:

\begin{enumerate}[label=\Alph*.]
    \item lack of educational facilities.
    \item poor rural communication infrastructure.
    \item livestock losses.
    \item decline in the number of hunting licences.
[](https://cracku.in/10-the-inhabitants-of-lozere-have-to-grapple-with-all-x-cat-2023-slot-1)
\end{enumerate}
\vspace{1em}
\textbf{Solution:}

Considering the first paragraph: “ _RESIDENTS of Lozère, a hilly department in southern France, recite complaints familiar to many rural corners of Europe. In remote hamlets and villages, with names such as Le Bacon and Le Bacon Vieux,_**_mayors grumble about a lack of local schools, jobs, or phone and internet connections._**_Farmers of grazing animals add another concern: the return of wolves. Eradicated from France last century, the predators are gradually creeping back to more forests and hillsides. “The wolf must be taken in hand,” said an aspiring parliamentarian, Francis Palombi, when pressed by voters in an election campaign early this summer. Tourists enjoy visiting a wolf park in Lozère,_**_but farmers fret over their livestock and their livelihoods_** _. .”_
Options A, B and C can be clearly inferred from the highlighted part.
The passage mentions that the number of people holding hunting licenses is still high but the number of people who still actively hunt is low. So Option D which states that there is decline in the number of hunting licences is incorrect.

\end{problem}

\newpage
\begin{problem}{}{}
% Subject: varc
\textbf{Instructions}

Instructions   

**The passage below is accompanied by four questions. Based on the passage, choose the best answer for each question.**
For early postcolonial literature, the world of the novel was often the nation. Postcolonial novels were usually [concerned with] national questions. Sometimes the whole story of the novel was taken as an allegory of the nation, whether India or Tanzania. This was important for supporting anti-colonial nationalism, but could also be limiting - land-focused and inward-looking.
My new book “Writing Ocean Worlds” explores another kind of world of the novel: not the village or nation, but the Indian Ocean world. The book describes a set of novels in which the Indian Ocean is at the centre of the story. It focuses on the novelists Amitav Ghosh, Abdulrazak Gurnah, Lindsey Collen and Joseph Conrad [who have] centred the Indian Ocean world in the majority of their novels. . . . Their work reveals a world that is outward-looking - full of movement, border-crossing and south-south interconnection. They are all very different - from colonially inclined (Conrad) to radically anti-capitalist (Collen), but together draw on and shape a wider sense of Indian Ocean space through themes, images, metaphors and language. This has the effect of remapping the world in the reader’s mind, as centred in the interconnected global south. . . .
The Indian Ocean world is a term used to describe the very long-lasting connections among the coasts of East Africa, the Arab coasts, and South and East Asia. These connections were made possible by the geography of the Indian Ocean. For much of history, travel by sea was much easier than by land, which meant that port cities very far apart were often more easily connected to each other than to much closer inland cities. Historical and archaeological evidence suggests that what we now call globalisation first appeared in the Indian Ocean. This is the interconnected oceanic world referenced and produced by the novels in my book. . . .
For their part Ghosh, Gurnah, Collen and even Conrad reference a different set of histories and geographies than the ones most commonly found in fiction in English. Those [commonly found ones] are mostly centred in Europe or the US, assume a background of Christianity and whiteness, and mention places like Paris and New York. The novels in [my] book highlight instead a largely Islamic space, feature characters of colour and centralise the ports of Malindi, Mombasa, Aden, Java and Bombay. . . . It is a densely imagined, richly sensory image of a southern cosmopolitan culture which provides for an enlarged sense of place in the world.
This remapping is particularly powerful for the representation of Africa. In the fiction, sailors and travellers are not all European. . . . African, as well as Indian and Arab characters, are traders, nakhodas (dhow ship captains), runaways, villains, missionaries and activists. This does not mean that Indian Ocean Africa is romanticised. Migration is often a matter of force; travel is portrayed as abandonment rather than adventure, freedoms are kept from women and slavery is rife. What it does mean is that the African part of the Indian Ocean world plays an active role in its long, rich history and therefore in that of the wider world.

\vspace{1em}
\textbf{Question 11}

Which one of the following statements, if true, would weaken the author’s claims?

\begin{enumerate}[label=\Alph*.]
    \item Having migrated out in the last century, wolves are now returning to Lozère.
    \item Unemployment concerns the residents of Lozère.
    \item Wolf attacks on tourists in Lozère are on the rise.
    \item The old mining sites of Lozère are now being used as grazing pastures for sheep.
[](https://cracku.in/11-which-one-of-the-following-statements-if-true-woul-x-cat-2023-slot-1)
\end{enumerate}
\vspace{1em}
\textbf{Solution:}

The author's claims seem to be focused on the conflicts between farmers and the return of wolves, the economic implications, and the coexistence challenges. If wolf attacks on tourists were on the rise, it might shift the narrative and suggest a broader safety concern beyond the impact on farmers, potentially weakening the author's emphasis on the positive economic aspects of wolf-related tourism. Therefore Option C, if true, would weaken the author’s argument.
Option A supports the author's claims about the return of wolves to Lozère.
Option B is not directly related to the author's claims about conflicts between farmers and wolves or the economic implications of wolf-related tourism.
Option D , if true, would not necessarily weaken the author's claims but might be seen as providing additional information about land use in Lozère.

\end{problem}

\newpage
\begin{problem}{}{}
% Subject: varc
\textbf{Instructions}

Instructions   

**The passage below is accompanied by four questions. Based on the passage, choose the best answer for each question.**
For early postcolonial literature, the world of the novel was often the nation. Postcolonial novels were usually [concerned with] national questions. Sometimes the whole story of the novel was taken as an allegory of the nation, whether India or Tanzania. This was important for supporting anti-colonial nationalism, but could also be limiting - land-focused and inward-looking.
My new book “Writing Ocean Worlds” explores another kind of world of the novel: not the village or nation, but the Indian Ocean world. The book describes a set of novels in which the Indian Ocean is at the centre of the story. It focuses on the novelists Amitav Ghosh, Abdulrazak Gurnah, Lindsey Collen and Joseph Conrad [who have] centred the Indian Ocean world in the majority of their novels. . . . Their work reveals a world that is outward-looking - full of movement, border-crossing and south-south interconnection. They are all very different - from colonially inclined (Conrad) to radically anti-capitalist (Collen), but together draw on and shape a wider sense of Indian Ocean space through themes, images, metaphors and language. This has the effect of remapping the world in the reader’s mind, as centred in the interconnected global south. . . .
The Indian Ocean world is a term used to describe the very long-lasting connections among the coasts of East Africa, the Arab coasts, and South and East Asia. These connections were made possible by the geography of the Indian Ocean. For much of history, travel by sea was much easier than by land, which meant that port cities very far apart were often more easily connected to each other than to much closer inland cities. Historical and archaeological evidence suggests that what we now call globalisation first appeared in the Indian Ocean. This is the interconnected oceanic world referenced and produced by the novels in my book. . . .
For their part Ghosh, Gurnah, Collen and even Conrad reference a different set of histories and geographies than the ones most commonly found in fiction in English. Those [commonly found ones] are mostly centred in Europe or the US, assume a background of Christianity and whiteness, and mention places like Paris and New York. The novels in [my] book highlight instead a largely Islamic space, feature characters of colour and centralise the ports of Malindi, Mombasa, Aden, Java and Bombay. . . . It is a densely imagined, richly sensory image of a southern cosmopolitan culture which provides for an enlarged sense of place in the world.
This remapping is particularly powerful for the representation of Africa. In the fiction, sailors and travellers are not all European. . . . African, as well as Indian and Arab characters, are traders, nakhodas (dhow ship captains), runaways, villains, missionaries and activists. This does not mean that Indian Ocean Africa is romanticised. Migration is often a matter of force; travel is portrayed as abandonment rather than adventure, freedoms are kept from women and slavery is rife. What it does mean is that the African part of the Indian Ocean world plays an active role in its long, rich history and therefore in that of the wider world.

\vspace{1em}
\textbf{Question 12}

The author presents a possible economic solution to an existing issue facing Lozère that takes into account the divergent and competing interests of:

\begin{enumerate}[label=\Alph*.]
    \item politicians and farmers.
    \item environmentalists and politicians.
    \item farmers and environmentalists.
    \item tourists and environmentalists.
[](https://cracku.in/12-the-author-presents-a-possible-economic-solution-t-x-cat-2023-slot-1)
\end{enumerate}
\vspace{1em}
\textbf{Solution:}

_“As the lupine population of Europe spreads westwards, with occasional reports of wolves seen closer to urban areas, expect to hear of more clashes between farmers and those who celebrate the predators’ return. Farmers’ losses are real, but are not the only economic story. Tourist venues, such as parks where wolves are kept and the animals’ spread is discussed, also generate income and jobs in rural areas.”_
The passage mentions that farmers in Lozère are concerned about the return of wolves causing losses in livestock. On the other hand, environmentalists view the presence of wolves as a sign of wider ecological health. The suggested economic solution involves tourist venues related to wolves, such as parks, which not only address the economic concerns of farmers by generating income but also align with the interests of environmentalists who appreciate the return of the predators. Therefore, Option C accurately captures the collaboration between farmers and environmentalists in the proposed solution.

Instructions   

The passage below is accompanied by four questions. Based on the passage, choose the best answer for each question.
[Fifty] years after its publication in English [in 1972], and just a year since [Marshall] Sahlins himself died—we may ask: why did [his essay] “Original Affluent Society” have such an impact, and how has it fared since? . . . Sahlins’s principal argument was simple but counterintuitive: before being driven into marginal environments by colonial powers, hunter-gatherers, or foragers, were not engaged in a desperate struggle for meager survival. Quite the contrary, they satisfied their needs with far less work than people in agricultural and industrial societies, leaving them more time to use as they wished. Hunters, he quipped, keep bankers’ hours. Refusing to maximize, many were “more concerned with games of chance than with chances of game.” . . . The so-called Neolithic Revolution, rather than improving life, imposed a harsher work regime and set in motion the long history of growing inequality . . .
Moreover, foragers had other options. The contemporary Hadza of Tanzania, who had long been surrounded by farmers, knew they had alternatives and rejected them. To Sahlins, this showed that foragers are not simply examples of human diversity or victimhood but something more profound: they demonstrated that societies make real choices. Culture, a way of living oriented around a distinctive set of values, manifests a fundamental principle of collective self-determination. . . .
But the point [of the essay] is not so much the empirical validity of the data—the real interest for most readers, after all, is not in foragers either today or in the Paleolithic—but rather its conceptual challenge to contemporary economic life and bourgeois individualism. The empirical served a philosophical and political project, a thought experiment and stimulus to the imagination of possibilities.
With its title’s nod toward The Affluent Society (1958), economist John Kenneth Galbraith’s famously skeptical portrait of America’s postwar prosperity and inequality, and dripping with New Left contempt for consumerism, “The Original Affluent Society” brought this critical perspective to bear on the contemporary world. It did so through the classic anthropological move of showing that radical alternatives to the readers’ lives really exist. If the capitalist world seeks wealth through ever greater material production to meet infinitely expansive desires, foraging societies follow “the Zen road to affluence”: not by getting more, but by wanting less. If it seems that foragers have been left behind by “progress,” this is due only to the ethnocentric self-congratulation of the West. Rather than accumulate material goods, these societies are guided by other values: leisure, mobility, and above all, freedom. . . .
Viewed in today’s context, of course, not every aspect of the essay has aged well. While acknowledging the violence of colonialism, racism, and dispossession, it does not thematize them as heavily as we might today. Rebuking evolutionary anthropologists for treating present-day foragers as “left behind” by progress, it too can succumb to the temptation to use them as proxies for the Paleolithic. Yet these characteristics should not distract us from appreciating Sahlins’s effort to show that if we want to conjure new possibilities, we need to learn about actually inhabitable worlds.

\end{problem}

\newpage
\begin{problem}{}{}
% Subject: varc
\textbf{Instructions}

Instructions   

**The passage below is accompanied by four questions. Based on the passage, choose the best answer for each question.**
For early postcolonial literature, the world of the novel was often the nation. Postcolonial novels were usually [concerned with] national questions. Sometimes the whole story of the novel was taken as an allegory of the nation, whether India or Tanzania. This was important for supporting anti-colonial nationalism, but could also be limiting - land-focused and inward-looking.
My new book “Writing Ocean Worlds” explores another kind of world of the novel: not the village or nation, but the Indian Ocean world. The book describes a set of novels in which the Indian Ocean is at the centre of the story. It focuses on the novelists Amitav Ghosh, Abdulrazak Gurnah, Lindsey Collen and Joseph Conrad [who have] centred the Indian Ocean world in the majority of their novels. . . . Their work reveals a world that is outward-looking - full of movement, border-crossing and south-south interconnection. They are all very different - from colonially inclined (Conrad) to radically anti-capitalist (Collen), but together draw on and shape a wider sense of Indian Ocean space through themes, images, metaphors and language. This has the effect of remapping the world in the reader’s mind, as centred in the interconnected global south. . . .
The Indian Ocean world is a term used to describe the very long-lasting connections among the coasts of East Africa, the Arab coasts, and South and East Asia. These connections were made possible by the geography of the Indian Ocean. For much of history, travel by sea was much easier than by land, which meant that port cities very far apart were often more easily connected to each other than to much closer inland cities. Historical and archaeological evidence suggests that what we now call globalisation first appeared in the Indian Ocean. This is the interconnected oceanic world referenced and produced by the novels in my book. . . .
For their part Ghosh, Gurnah, Collen and even Conrad reference a different set of histories and geographies than the ones most commonly found in fiction in English. Those [commonly found ones] are mostly centred in Europe or the US, assume a background of Christianity and whiteness, and mention places like Paris and New York. The novels in [my] book highlight instead a largely Islamic space, feature characters of colour and centralise the ports of Malindi, Mombasa, Aden, Java and Bombay. . . . It is a densely imagined, richly sensory image of a southern cosmopolitan culture which provides for an enlarged sense of place in the world.
This remapping is particularly powerful for the representation of Africa. In the fiction, sailors and travellers are not all European. . . . African, as well as Indian and Arab characters, are traders, nakhodas (dhow ship captains), runaways, villains, missionaries and activists. This does not mean that Indian Ocean Africa is romanticised. Migration is often a matter of force; travel is portrayed as abandonment rather than adventure, freedoms are kept from women and slavery is rife. What it does mean is that the African part of the Indian Ocean world plays an active role in its long, rich history and therefore in that of the wider world.

\vspace{1em}
\textbf{Question 13}

We can infer that Sahlins's main goal in writing his essay was to:

\begin{enumerate}[label=\Alph*.]
    \item put forth the view that, despite egalitarian origins, economic progress brings greater inequality and social hierarchies.
    \item highlight the fact that while we started off as a fairly contented egalitarian people, we have progressively degenerated into materialism.
    \item hold a mirror to an acquisitive society, with examples of other communities that have chosen successfully to be non-materialistic
    \item counter Galbraith’s pessimistic view of the inevitability of a capitalist trajectory for economic growth.
[](https://cracku.in/13-we-can-infer-that-sahlinss-main-goal-in-writing-hi-x-cat-2023-slot-1)
\end{enumerate}
\vspace{1em}
\textbf{Solution:}

The passage emphasizes that Marshall Sahlins's main goal in writing his essay was to hold a mirror to an acquisitive society (contemporary economic life and bourgeois individualism). The essay accomplishes this by providing examples of foraging societies that made real choices to prioritize values such as leisure, mobility, and freedom over material accumulation. Sahlins contrasts the Zen road to affluence, where affluence is achieved by wanting less, with the capitalist pursuit of wealth through material production and consumerism. Therefore, Sahlins's goal, as portrayed in the passage, aligns with the idea of presenting examples of communities that have successfully chosen non-materialistic paths as a critique of acquisitive societies. So, Option C is the correct answer.
Option A: While Sahlins's essay acknowledges growing inequality and social hierarchies resulting from the Neolithic Revolution, it is more focused on contrasting foraging societies with contemporary economic life.
Option B: The primary emphasis is on showcasing foraging societies' choices and values rather than asserting a progressive degeneration of society.
Option D: Even though Sahlins's essay critiques aspects of contemporary economic views, its primary focus is not explicitly countering Galbraith's pessimistic view but rather presenting alternative possibilities through examples of non-materialistic societies.

\end{problem}

\newpage
\begin{problem}{}{}
% Subject: varc
\textbf{Instructions}

Instructions   

**The passage below is accompanied by four questions. Based on the passage, choose the best answer for each question.**
For early postcolonial literature, the world of the novel was often the nation. Postcolonial novels were usually [concerned with] national questions. Sometimes the whole story of the novel was taken as an allegory of the nation, whether India or Tanzania. This was important for supporting anti-colonial nationalism, but could also be limiting - land-focused and inward-looking.
My new book “Writing Ocean Worlds” explores another kind of world of the novel: not the village or nation, but the Indian Ocean world. The book describes a set of novels in which the Indian Ocean is at the centre of the story. It focuses on the novelists Amitav Ghosh, Abdulrazak Gurnah, Lindsey Collen and Joseph Conrad [who have] centred the Indian Ocean world in the majority of their novels. . . . Their work reveals a world that is outward-looking - full of movement, border-crossing and south-south interconnection. They are all very different - from colonially inclined (Conrad) to radically anti-capitalist (Collen), but together draw on and shape a wider sense of Indian Ocean space through themes, images, metaphors and language. This has the effect of remapping the world in the reader’s mind, as centred in the interconnected global south. . . .
The Indian Ocean world is a term used to describe the very long-lasting connections among the coasts of East Africa, the Arab coasts, and South and East Asia. These connections were made possible by the geography of the Indian Ocean. For much of history, travel by sea was much easier than by land, which meant that port cities very far apart were often more easily connected to each other than to much closer inland cities. Historical and archaeological evidence suggests that what we now call globalisation first appeared in the Indian Ocean. This is the interconnected oceanic world referenced and produced by the novels in my book. . . .
For their part Ghosh, Gurnah, Collen and even Conrad reference a different set of histories and geographies than the ones most commonly found in fiction in English. Those [commonly found ones] are mostly centred in Europe or the US, assume a background of Christianity and whiteness, and mention places like Paris and New York. The novels in [my] book highlight instead a largely Islamic space, feature characters of colour and centralise the ports of Malindi, Mombasa, Aden, Java and Bombay. . . . It is a densely imagined, richly sensory image of a southern cosmopolitan culture which provides for an enlarged sense of place in the world.
This remapping is particularly powerful for the representation of Africa. In the fiction, sailors and travellers are not all European. . . . African, as well as Indian and Arab characters, are traders, nakhodas (dhow ship captains), runaways, villains, missionaries and activists. This does not mean that Indian Ocean Africa is romanticised. Migration is often a matter of force; travel is portrayed as abandonment rather than adventure, freedoms are kept from women and slavery is rife. What it does mean is that the African part of the Indian Ocean world plays an active role in its long, rich history and therefore in that of the wider world.

\vspace{1em}
\textbf{Question 14}

The author mentions Tanzania’s Hadza community to illustrate:

\begin{enumerate}[label=\Alph*.]
    \item that hunter-gatherer communities’ subsistence-level techniques equipped them to survive well into contemporary times.
    \item how pre-agrarian societies did not hamper the emergence of more advanced agrarian practices in contiguous communities.
    \item that forager communities’ lifestyles derived not from ignorance about alternatives, but from their own choice.
    \item how two vastly different ways of living and working were able to coexist in proximity for centuries.
[](https://cracku.in/14-the-author-mentions-tanzanias-hadza-community-to-i-x-cat-2023-slot-1)
\end{enumerate}
\vspace{1em}
\textbf{Solution:}

Option C is the correct answer because the passage uses the example of Tanzania's Hadza community to illustrate that forager communities, like the Hadza, do not conform to a simple narrative of human diversity or victimhood. Instead, they actively make choices about their way of life. The passage mentions that the Hadza, despite being surrounded by farmers, knew they had alternatives and consciously rejected them. This example serves to emphasize that forager communities are not constrained by ignorance about alternatives; rather, their lifestyles derive from their own choices. Therefore, Option C accurately captures the essence of the Hadza illustration in the passage.
Option A: The passage doesn't specifically highlight the survival techniques of hunter-gatherer communities into contemporary times, but rather emphasizes their choices and values.
Option B: The passage doesn't discuss the Hadza community in the context of agrarian practices in contiguous communities, making this option irrelevant to the illustration.
Option D: The passage does not suggest that the Hadza community coexisted with vastly different ways of living and working for centuries

\end{problem}

\newpage
\begin{problem}{}{}
% Subject: varc
\textbf{Instructions}

Instructions   

**The passage below is accompanied by four questions. Based on the passage, choose the best answer for each question.**
For early postcolonial literature, the world of the novel was often the nation. Postcolonial novels were usually [concerned with] national questions. Sometimes the whole story of the novel was taken as an allegory of the nation, whether India or Tanzania. This was important for supporting anti-colonial nationalism, but could also be limiting - land-focused and inward-looking.
My new book “Writing Ocean Worlds” explores another kind of world of the novel: not the village or nation, but the Indian Ocean world. The book describes a set of novels in which the Indian Ocean is at the centre of the story. It focuses on the novelists Amitav Ghosh, Abdulrazak Gurnah, Lindsey Collen and Joseph Conrad [who have] centred the Indian Ocean world in the majority of their novels. . . . Their work reveals a world that is outward-looking - full of movement, border-crossing and south-south interconnection. They are all very different - from colonially inclined (Conrad) to radically anti-capitalist (Collen), but together draw on and shape a wider sense of Indian Ocean space through themes, images, metaphors and language. This has the effect of remapping the world in the reader’s mind, as centred in the interconnected global south. . . .
The Indian Ocean world is a term used to describe the very long-lasting connections among the coasts of East Africa, the Arab coasts, and South and East Asia. These connections were made possible by the geography of the Indian Ocean. For much of history, travel by sea was much easier than by land, which meant that port cities very far apart were often more easily connected to each other than to much closer inland cities. Historical and archaeological evidence suggests that what we now call globalisation first appeared in the Indian Ocean. This is the interconnected oceanic world referenced and produced by the novels in my book. . . .
For their part Ghosh, Gurnah, Collen and even Conrad reference a different set of histories and geographies than the ones most commonly found in fiction in English. Those [commonly found ones] are mostly centred in Europe or the US, assume a background of Christianity and whiteness, and mention places like Paris and New York. The novels in [my] book highlight instead a largely Islamic space, feature characters of colour and centralise the ports of Malindi, Mombasa, Aden, Java and Bombay. . . . It is a densely imagined, richly sensory image of a southern cosmopolitan culture which provides for an enlarged sense of place in the world.
This remapping is particularly powerful for the representation of Africa. In the fiction, sailors and travellers are not all European. . . . African, as well as Indian and Arab characters, are traders, nakhodas (dhow ship captains), runaways, villains, missionaries and activists. This does not mean that Indian Ocean Africa is romanticised. Migration is often a matter of force; travel is portrayed as abandonment rather than adventure, freedoms are kept from women and slavery is rife. What it does mean is that the African part of the Indian Ocean world plays an active role in its long, rich history and therefore in that of the wider world.

\vspace{1em}
\textbf{Question 15}

The author of the passage mentions Galbraith’s “The Affluent Society” to:

\begin{enumerate}[label=\Alph*.]
    \item show how Galbraith’s theories refute Sahlins’s thesis on the contentment of pre-hunter-gatherer communities.
    \item contrast the materialist nature of contemporary growth paths with the pacifist content ways of living among the foragers.
    \item document the influence of Galbraith’s cynical views on modern consumerism on Sahlins’s analysis of pre-historic societies.
    \item show how Sahlins’s views complemented Galbraith’s criticism of the consumerism and inequality of contemporary society.
[](https://cracku.in/15-the-author-of-the-passage-mentions-galbraiths-the--x-cat-2023-slot-1)
\end{enumerate}
\vspace{1em}
\textbf{Solution:}

The passage explicitly mentions that Sahlins's essay, "The Original Affluent Society," brought a critical perspective to contemporary consumerism and inequality, echoing the themes found in John Kenneth Galbraith's work, "The Affluent Society." The passage notes that Sahlins's essay contrasts the values of foraging societies with the capitalist pursuit of wealth, and it suggests that the essay complements Galbraith's skeptical portrait of postwar prosperity and inequality. Therefore, Option D accurately reflects the information presented in the passage regarding the relationship between Sahlins's views and Galbraith's criticism of contemporary society.
The passage does not suggest that Galbraith’s theories refute Sahlins’s thesis but rather highlights their complementarity (Option A) nor does it focus on contrasting foragers' ways of living with Galbraith's views on contemporary growth paths (Option B).
The passage does not document the influence of Galbraith’s views on Sahlins’s analysis; instead, it emphasizes how Sahlins's essay complements Galbraith’s critical perspective on contemporary society. Therefore Option C is incorrect too.

\end{problem}

\newpage
\begin{problem}{}{}
% Subject: varc
\textbf{Instructions}

Instructions   

**The passage below is accompanied by four questions. Based on the passage, choose the best answer for each question.**
For early postcolonial literature, the world of the novel was often the nation. Postcolonial novels were usually [concerned with] national questions. Sometimes the whole story of the novel was taken as an allegory of the nation, whether India or Tanzania. This was important for supporting anti-colonial nationalism, but could also be limiting - land-focused and inward-looking.
My new book “Writing Ocean Worlds” explores another kind of world of the novel: not the village or nation, but the Indian Ocean world. The book describes a set of novels in which the Indian Ocean is at the centre of the story. It focuses on the novelists Amitav Ghosh, Abdulrazak Gurnah, Lindsey Collen and Joseph Conrad [who have] centred the Indian Ocean world in the majority of their novels. . . . Their work reveals a world that is outward-looking - full of movement, border-crossing and south-south interconnection. They are all very different - from colonially inclined (Conrad) to radically anti-capitalist (Collen), but together draw on and shape a wider sense of Indian Ocean space through themes, images, metaphors and language. This has the effect of remapping the world in the reader’s mind, as centred in the interconnected global south. . . .
The Indian Ocean world is a term used to describe the very long-lasting connections among the coasts of East Africa, the Arab coasts, and South and East Asia. These connections were made possible by the geography of the Indian Ocean. For much of history, travel by sea was much easier than by land, which meant that port cities very far apart were often more easily connected to each other than to much closer inland cities. Historical and archaeological evidence suggests that what we now call globalisation first appeared in the Indian Ocean. This is the interconnected oceanic world referenced and produced by the novels in my book. . . .
For their part Ghosh, Gurnah, Collen and even Conrad reference a different set of histories and geographies than the ones most commonly found in fiction in English. Those [commonly found ones] are mostly centred in Europe or the US, assume a background of Christianity and whiteness, and mention places like Paris and New York. The novels in [my] book highlight instead a largely Islamic space, feature characters of colour and centralise the ports of Malindi, Mombasa, Aden, Java and Bombay. . . . It is a densely imagined, richly sensory image of a southern cosmopolitan culture which provides for an enlarged sense of place in the world.
This remapping is particularly powerful for the representation of Africa. In the fiction, sailors and travellers are not all European. . . . African, as well as Indian and Arab characters, are traders, nakhodas (dhow ship captains), runaways, villains, missionaries and activists. This does not mean that Indian Ocean Africa is romanticised. Migration is often a matter of force; travel is portrayed as abandonment rather than adventure, freedoms are kept from women and slavery is rife. What it does mean is that the African part of the Indian Ocean world plays an active role in its long, rich history and therefore in that of the wider world.

\vspace{1em}
\textbf{Question 16}

The author of the passage criticises Sahlins’s essay for its:

\begin{enumerate}[label=\Alph*.]
    \item cursory treatment of the effects of racism and colonialism on societies.
    \item outdated values regarding present-day foragers versus ancient foraging communities.
    \item critique of anthropologists who disparage the choices of foragers in today’s society
    \item failure to supplement its thesis with robust empirical data.
[](https://cracku.in/16-the-author-of-the-passage-criticises-sahlinss-essa-x-cat-2023-slot-1)
\end{enumerate}
\vspace{1em}
\textbf{Solution:}

"_Viewed in today’s context, of course, not every aspect of the essay has aged well. While acknowledging the violence of colonialism, racism, and dispossession, it does not thematize them as heavily as we might today."_
Option A is the correct answer because the passage explicitly mentions that, when viewed in today's context, not every aspect of Sahlins's essay has aged well, and it acknowledges that the essay does not thematize issues like racism, colonialism, and dispossession as heavily as might be expected today. The term "cursory treatment" suggests that the essay provides only a brief or superficial examination of the effects of racism and colonialism on societies, and the passage criticizes this aspect of the essay for not giving these important issues more comprehensive attention.

Instructions   

For the following questions answer them individually

\end{problem}

\newpage
\begin{problem}{}{}
% Subject: varc
\textbf{Instructions}

Instructions   

**The passage below is accompanied by four questions. Based on the passage, choose the best answer for each question.**
For early postcolonial literature, the world of the novel was often the nation. Postcolonial novels were usually [concerned with] national questions. Sometimes the whole story of the novel was taken as an allegory of the nation, whether India or Tanzania. This was important for supporting anti-colonial nationalism, but could also be limiting - land-focused and inward-looking.
My new book “Writing Ocean Worlds” explores another kind of world of the novel: not the village or nation, but the Indian Ocean world. The book describes a set of novels in which the Indian Ocean is at the centre of the story. It focuses on the novelists Amitav Ghosh, Abdulrazak Gurnah, Lindsey Collen and Joseph Conrad [who have] centred the Indian Ocean world in the majority of their novels. . . . Their work reveals a world that is outward-looking - full of movement, border-crossing and south-south interconnection. They are all very different - from colonially inclined (Conrad) to radically anti-capitalist (Collen), but together draw on and shape a wider sense of Indian Ocean space through themes, images, metaphors and language. This has the effect of remapping the world in the reader’s mind, as centred in the interconnected global south. . . .
The Indian Ocean world is a term used to describe the very long-lasting connections among the coasts of East Africa, the Arab coasts, and South and East Asia. These connections were made possible by the geography of the Indian Ocean. For much of history, travel by sea was much easier than by land, which meant that port cities very far apart were often more easily connected to each other than to much closer inland cities. Historical and archaeological evidence suggests that what we now call globalisation first appeared in the Indian Ocean. This is the interconnected oceanic world referenced and produced by the novels in my book. . . .
For their part Ghosh, Gurnah, Collen and even Conrad reference a different set of histories and geographies than the ones most commonly found in fiction in English. Those [commonly found ones] are mostly centred in Europe or the US, assume a background of Christianity and whiteness, and mention places like Paris and New York. The novels in [my] book highlight instead a largely Islamic space, feature characters of colour and centralise the ports of Malindi, Mombasa, Aden, Java and Bombay. . . . It is a densely imagined, richly sensory image of a southern cosmopolitan culture which provides for an enlarged sense of place in the world.
This remapping is particularly powerful for the representation of Africa. In the fiction, sailors and travellers are not all European. . . . African, as well as Indian and Arab characters, are traders, nakhodas (dhow ship captains), runaways, villains, missionaries and activists. This does not mean that Indian Ocean Africa is romanticised. Migration is often a matter of force; travel is portrayed as abandonment rather than adventure, freedoms are kept from women and slavery is rife. What it does mean is that the African part of the Indian Ocean world plays an active role in its long, rich history and therefore in that of the wider world.

\vspace{1em}
\textbf{Question 17}

**There is a sentence that is missing in the paragraph below. Look at the paragraph and decide where (option 1, 2, 3, or 4) the following sentence would best fit.**
Sentence: The discovery helps to explain archeological similarities between the Paleolithic peoples of China, Japan, and the Americas.
Paragraph: The researchers also uncovered an unexpected genetic link between Native Americans and Japanese people. ___(1)___. During the deglaciation period, another group branched out from northern coastal China and travelled to Japan. ___(2)___. "We were surprised to find that this ancestral source also contributed to the Japanese gene pool, especially the indigenous Ainus," says Li. ___(3)___. They shared similarities in how they crafted stemmed projectile points for arrowheads and spears. ___(4)___. "This suggests that the Pleistocene connection among the Americas, China, and Japan was not confined to culture but also to genetics," says senior author Qing-Peng Kong, an evolutionary geneticist at the Chinese Academy of Sciences.

\begin{enumerate}[label=\Alph*.]
    \item Option 2
    \item Option 3
    \item Option 1
    \item Option 4
[](https://cracku.in/17-there-is-a-sentence-that-is-missing-in-the-paragra-x-cat-2023-slot-1)
\end{enumerate}
\vspace{1em}
\textbf{Solution:}

The sentence best fits in Blank 3 because it logically follows the mention of the unexpected genetic link between Native Americans and Japanese people mentioned in the sentence preceding Blank 3. After establishing this genetic connection, the sentence provides additional context by explaining the archaeological implications of the discovery. It suggests that the shared genetic link has archaeological manifestations, leading to similarities in the Paleolithic peoples of China, Japan, and the Americas. Placing the sentence here helps to connect the genetic findings to broader archaeological and cultural aspects.
Blank 1: The sentence before Blank 1 introduces the unexpected genetic link between Native Americans and Japanese people. Placing this sentence in Blank 1 would disrupt the logical progression of information, as Blank 1 should provide information that directly connects with or follows from the mention of the unexpected genetic link, and the sentence is more relevant to explaining the broader context after the genetic link has been introduced.
Blank 2: Putting the sentence here makes no sense as there has been no mention of any discovery before.
Blank 4: Placing the sentence here would disrupt the flow of the passage as before Blank 4 the passage mentions how there were similarities between them and After Blank 4, the passage states “ _This suggests..”._ “This” clearly refers to the similarities mentioned before Blank 4.

\end{problem}

\newpage
\begin{problem}{}{}
% Subject: varc
\textbf{Instructions}

Instructions   

**The passage below is accompanied by four questions. Based on the passage, choose the best answer for each question.**
For early postcolonial literature, the world of the novel was often the nation. Postcolonial novels were usually [concerned with] national questions. Sometimes the whole story of the novel was taken as an allegory of the nation, whether India or Tanzania. This was important for supporting anti-colonial nationalism, but could also be limiting - land-focused and inward-looking.
My new book “Writing Ocean Worlds” explores another kind of world of the novel: not the village or nation, but the Indian Ocean world. The book describes a set of novels in which the Indian Ocean is at the centre of the story. It focuses on the novelists Amitav Ghosh, Abdulrazak Gurnah, Lindsey Collen and Joseph Conrad [who have] centred the Indian Ocean world in the majority of their novels. . . . Their work reveals a world that is outward-looking - full of movement, border-crossing and south-south interconnection. They are all very different - from colonially inclined (Conrad) to radically anti-capitalist (Collen), but together draw on and shape a wider sense of Indian Ocean space through themes, images, metaphors and language. This has the effect of remapping the world in the reader’s mind, as centred in the interconnected global south. . . .
The Indian Ocean world is a term used to describe the very long-lasting connections among the coasts of East Africa, the Arab coasts, and South and East Asia. These connections were made possible by the geography of the Indian Ocean. For much of history, travel by sea was much easier than by land, which meant that port cities very far apart were often more easily connected to each other than to much closer inland cities. Historical and archaeological evidence suggests that what we now call globalisation first appeared in the Indian Ocean. This is the interconnected oceanic world referenced and produced by the novels in my book. . . .
For their part Ghosh, Gurnah, Collen and even Conrad reference a different set of histories and geographies than the ones most commonly found in fiction in English. Those [commonly found ones] are mostly centred in Europe or the US, assume a background of Christianity and whiteness, and mention places like Paris and New York. The novels in [my] book highlight instead a largely Islamic space, feature characters of colour and centralise the ports of Malindi, Mombasa, Aden, Java and Bombay. . . . It is a densely imagined, richly sensory image of a southern cosmopolitan culture which provides for an enlarged sense of place in the world.
This remapping is particularly powerful for the representation of Africa. In the fiction, sailors and travellers are not all European. . . . African, as well as Indian and Arab characters, are traders, nakhodas (dhow ship captains), runaways, villains, missionaries and activists. This does not mean that Indian Ocean Africa is romanticised. Migration is often a matter of force; travel is portrayed as abandonment rather than adventure, freedoms are kept from women and slavery is rife. What it does mean is that the African part of the Indian Ocean world plays an active role in its long, rich history and therefore in that of the wider world.

\vspace{1em}
\textbf{Question 18}

There is a sentence that is missing in the paragraph below. Look at the paragraph and decide where (option 1, 2, 3, or 4) the following sentence would best fit.
**Sentence** : This philosophical cut at one’s core beliefs, values, and way of life is difficult enough.
**Paragraph** : The experience of reading philosophy is often disquieting. When reading philosophy, the values around which one has heretofore organised one’s life may come to look provincial, flatly wrong, or even evil. ___(1)___. When beliefs previously held as truths are rendered implausible, new beliefs, values, and ways of living may be required. ___(2)___. What’s worse, philosophers admonish each other to remain unsutured until such time as a defensible new answer is revealed or constructed. Sometimes, philosophical writing is even strictly critical in that it does not even attempt to provide an alternative after tearing down a cultural or conceptual citadel. ___(3)___. The reader of philosophy must be prepared for the possibility of this experience. While reading philosophy can help one clarify one’s values, and even make one self-conscious for the first time of the fact that there are good reasons for believing what one believes, it can also generate unremediated doubt that is difficult to live with. ___(4)___.

\begin{enumerate}[label=\Alph*.]
    \item Option 1
    \item Option 4
    \item Option 3
    \item Option 2
[](https://cracku.in/18-there-is-a-sentence-that-is-missing-in-the-paragra-x-cat-2023-slot-1)
\end{enumerate}
\vspace{1em}
\textbf{Solution:}

The sentence best fits in Blank 2 because it provides a continuation and elaboration on the disquieting experience mentioned before Blank 2. Before Blank 2, the passage introduces the disquieting nature of reading philosophy, and the sentence builds on that by explaining that the philosophical examination challenges the values around which one has organized their life, making them appear provincial, flatly wrong, or even evil. The sentence in question serves to articulate the difficulty and discomfort associated with this profound philosophical scrutiny, logically following the initial statement about the disquieting experience of reading philosophy.
We can also see that the sentence ends with _“...and way of life is difficult enough.”_ This is immediately followed by the passage highlighting what’s worse _“What’s worse, philosophers admonish each other to remain unsutured until such time as a defensible new answer is revealed or constructed.”_. So we can see that Sentence 2 serves as a preceding sentence to this.

\end{problem}

\newpage
\begin{problem}{}{}
% Subject: varc
\textbf{Instructions}

Instructions   

**The passage below is accompanied by four questions. Based on the passage, choose the best answer for each question.**
For early postcolonial literature, the world of the novel was often the nation. Postcolonial novels were usually [concerned with] national questions. Sometimes the whole story of the novel was taken as an allegory of the nation, whether India or Tanzania. This was important for supporting anti-colonial nationalism, but could also be limiting - land-focused and inward-looking.
My new book “Writing Ocean Worlds” explores another kind of world of the novel: not the village or nation, but the Indian Ocean world. The book describes a set of novels in which the Indian Ocean is at the centre of the story. It focuses on the novelists Amitav Ghosh, Abdulrazak Gurnah, Lindsey Collen and Joseph Conrad [who have] centred the Indian Ocean world in the majority of their novels. . . . Their work reveals a world that is outward-looking - full of movement, border-crossing and south-south interconnection. They are all very different - from colonially inclined (Conrad) to radically anti-capitalist (Collen), but together draw on and shape a wider sense of Indian Ocean space through themes, images, metaphors and language. This has the effect of remapping the world in the reader’s mind, as centred in the interconnected global south. . . .
The Indian Ocean world is a term used to describe the very long-lasting connections among the coasts of East Africa, the Arab coasts, and South and East Asia. These connections were made possible by the geography of the Indian Ocean. For much of history, travel by sea was much easier than by land, which meant that port cities very far apart were often more easily connected to each other than to much closer inland cities. Historical and archaeological evidence suggests that what we now call globalisation first appeared in the Indian Ocean. This is the interconnected oceanic world referenced and produced by the novels in my book. . . .
For their part Ghosh, Gurnah, Collen and even Conrad reference a different set of histories and geographies than the ones most commonly found in fiction in English. Those [commonly found ones] are mostly centred in Europe or the US, assume a background of Christianity and whiteness, and mention places like Paris and New York. The novels in [my] book highlight instead a largely Islamic space, feature characters of colour and centralise the ports of Malindi, Mombasa, Aden, Java and Bombay. . . . It is a densely imagined, richly sensory image of a southern cosmopolitan culture which provides for an enlarged sense of place in the world.
This remapping is particularly powerful for the representation of Africa. In the fiction, sailors and travellers are not all European. . . . African, as well as Indian and Arab characters, are traders, nakhodas (dhow ship captains), runaways, villains, missionaries and activists. This does not mean that Indian Ocean Africa is romanticised. Migration is often a matter of force; travel is portrayed as abandonment rather than adventure, freedoms are kept from women and slavery is rife. What it does mean is that the African part of the Indian Ocean world plays an active role in its long, rich history and therefore in that of the wider world.

\vspace{1em}
\textbf{Question 19}

Five jumbled up sentences (labelled 1, 2, 3, 4 and 5), related to a topic, are given below. Four of them can be put together to form a coherent paragraph. Identify the odd sentence and key in the number of that sentence as your answer.  
1. Having an appreciation for the workings of another person’s mind is considered a prerequisite for natural language acquisition, strategic social interaction, reflexive thought, and moral judgment.  
2. It is a ‘theory of mind’ though some scholars prefer to call it ‘mentalizing’ or ‘mindreading’, which is important for the development of one's cognitive abilities.  
3. Though we must speculate about its evolutionary origin, we do have indications that the capacity evolved sometime in the last few million years.  
4. This capacity develops from early beginnings in the first year of life to the adult’s fast and often effortless understanding of others’ thoughts, feelings, and intentions.  
5. One of the most fascinating human capacities is the ability to perceive and interpret other people’s behaviour in terms of their mental states.
Backspace  
789  
456  
123  
0.-  
Clear All  

Submit
[](https://cracku.in/19-five-jumbled-up-sentences-labelled-1-2-3-4-and-5-r-x-cat-2023-slot-1)

\begin{enumerate}[label=\Alph*.]
\end{enumerate}
\vspace{1em}
\textbf{Solution:}

Sentence 2 is the odd one out because it introduces a term, "theory of mind," and discusses scholars' preferences for alternative terms like "mentalizing" or "mindreading." Unlike the other sentences, which focus on explaining and elaborating on the concept of understanding others' mental states, sentence 2 provides more of a meta-discussion about the terminology used to describe this capacity rather than directly contributing to the explanation of the topic. The other sentences contribute to the substantive discussion of the capacity to perceive and interpret other people's behavior in terms of their mental states, making sentence 2 the odd one out in the context of forming a coherent paragraph.

\end{problem}

\newpage
\begin{problem}{}{}
% Subject: varc
\textbf{Instructions}

Instructions   

**The passage below is accompanied by four questions. Based on the passage, choose the best answer for each question.**
For early postcolonial literature, the world of the novel was often the nation. Postcolonial novels were usually [concerned with] national questions. Sometimes the whole story of the novel was taken as an allegory of the nation, whether India or Tanzania. This was important for supporting anti-colonial nationalism, but could also be limiting - land-focused and inward-looking.
My new book “Writing Ocean Worlds” explores another kind of world of the novel: not the village or nation, but the Indian Ocean world. The book describes a set of novels in which the Indian Ocean is at the centre of the story. It focuses on the novelists Amitav Ghosh, Abdulrazak Gurnah, Lindsey Collen and Joseph Conrad [who have] centred the Indian Ocean world in the majority of their novels. . . . Their work reveals a world that is outward-looking - full of movement, border-crossing and south-south interconnection. They are all very different - from colonially inclined (Conrad) to radically anti-capitalist (Collen), but together draw on and shape a wider sense of Indian Ocean space through themes, images, metaphors and language. This has the effect of remapping the world in the reader’s mind, as centred in the interconnected global south. . . .
The Indian Ocean world is a term used to describe the very long-lasting connections among the coasts of East Africa, the Arab coasts, and South and East Asia. These connections were made possible by the geography of the Indian Ocean. For much of history, travel by sea was much easier than by land, which meant that port cities very far apart were often more easily connected to each other than to much closer inland cities. Historical and archaeological evidence suggests that what we now call globalisation first appeared in the Indian Ocean. This is the interconnected oceanic world referenced and produced by the novels in my book. . . .
For their part Ghosh, Gurnah, Collen and even Conrad reference a different set of histories and geographies than the ones most commonly found in fiction in English. Those [commonly found ones] are mostly centred in Europe or the US, assume a background of Christianity and whiteness, and mention places like Paris and New York. The novels in [my] book highlight instead a largely Islamic space, feature characters of colour and centralise the ports of Malindi, Mombasa, Aden, Java and Bombay. . . . It is a densely imagined, richly sensory image of a southern cosmopolitan culture which provides for an enlarged sense of place in the world.
This remapping is particularly powerful for the representation of Africa. In the fiction, sailors and travellers are not all European. . . . African, as well as Indian and Arab characters, are traders, nakhodas (dhow ship captains), runaways, villains, missionaries and activists. This does not mean that Indian Ocean Africa is romanticised. Migration is often a matter of force; travel is portrayed as abandonment rather than adventure, freedoms are kept from women and slavery is rife. What it does mean is that the African part of the Indian Ocean world plays an active role in its long, rich history and therefore in that of the wider world.

\vspace{1em}
\textbf{Question 20}

**Five jumbled up sentences (labelled 1, 2, 3, 4 and 5), related to a topic, are given below. Four of them can be put together to form a coherent paragraph. Identify the odd sentence and key in the number of that sentence as your answer.**
1. In English, there is no systematic rule for the naming of numbers; after ten, we have "eleven" and "twelve" and then the teens: "thirteen", "fourteen", "fifteen" and so on.
2. Even more confusingly, some English words invert the numbers they refer to: the word "fourteen" puts the four first, even though it appears last.
3. It can take children a while to learn all these words, and understand that "fourteen" is different from "forty".
4. For multiples of 10, English speakers switch to a different pattern: "twenty", "thirty", "forty" and so on.
5. If you didn't know the word for "eleven", you would be unable to just guess it - you might come up with something like "one-teen".
Backspace  
789  
456  
123  
0.-  
Clear All  

Submit
[](https://cracku.in/20-five-jumbled-up-sentences-labelled-1-2-3-4-and-5-r-x-cat-2023-slot-1)

\begin{enumerate}[label=\Alph*.]
\end{enumerate}
\vspace{1em}
\textbf{Solution:}

Sentence 3 is the odd one out because it introduces a different topic compared to the other sentences. While the other sentences focus on the naming patterns of numbers in English, especially the irregularities and variations in the system, sentence 3 shifts the focus to the learning process of children and their understanding of the differences between numbers like "fourteen" and "forty." The other sentences contribute to the discussion about the intricacies of English number naming, making sentence 3 less aligned with the central theme of the paragraph.

\end{problem}

\newpage
\begin{problem}{}{}
% Subject: varc
\textbf{Instructions}

Instructions   

**The passage below is accompanied by four questions. Based on the passage, choose the best answer for each question.**
For early postcolonial literature, the world of the novel was often the nation. Postcolonial novels were usually [concerned with] national questions. Sometimes the whole story of the novel was taken as an allegory of the nation, whether India or Tanzania. This was important for supporting anti-colonial nationalism, but could also be limiting - land-focused and inward-looking.
My new book “Writing Ocean Worlds” explores another kind of world of the novel: not the village or nation, but the Indian Ocean world. The book describes a set of novels in which the Indian Ocean is at the centre of the story. It focuses on the novelists Amitav Ghosh, Abdulrazak Gurnah, Lindsey Collen and Joseph Conrad [who have] centred the Indian Ocean world in the majority of their novels. . . . Their work reveals a world that is outward-looking - full of movement, border-crossing and south-south interconnection. They are all very different - from colonially inclined (Conrad) to radically anti-capitalist (Collen), but together draw on and shape a wider sense of Indian Ocean space through themes, images, metaphors and language. This has the effect of remapping the world in the reader’s mind, as centred in the interconnected global south. . . .
The Indian Ocean world is a term used to describe the very long-lasting connections among the coasts of East Africa, the Arab coasts, and South and East Asia. These connections were made possible by the geography of the Indian Ocean. For much of history, travel by sea was much easier than by land, which meant that port cities very far apart were often more easily connected to each other than to much closer inland cities. Historical and archaeological evidence suggests that what we now call globalisation first appeared in the Indian Ocean. This is the interconnected oceanic world referenced and produced by the novels in my book. . . .
For their part Ghosh, Gurnah, Collen and even Conrad reference a different set of histories and geographies than the ones most commonly found in fiction in English. Those [commonly found ones] are mostly centred in Europe or the US, assume a background of Christianity and whiteness, and mention places like Paris and New York. The novels in [my] book highlight instead a largely Islamic space, feature characters of colour and centralise the ports of Malindi, Mombasa, Aden, Java and Bombay. . . . It is a densely imagined, richly sensory image of a southern cosmopolitan culture which provides for an enlarged sense of place in the world.
This remapping is particularly powerful for the representation of Africa. In the fiction, sailors and travellers are not all European. . . . African, as well as Indian and Arab characters, are traders, nakhodas (dhow ship captains), runaways, villains, missionaries and activists. This does not mean that Indian Ocean Africa is romanticised. Migration is often a matter of force; travel is portrayed as abandonment rather than adventure, freedoms are kept from women and slavery is rife. What it does mean is that the African part of the Indian Ocean world plays an active role in its long, rich history and therefore in that of the wider world.

\vspace{1em}
\textbf{Question 21}

**The four sentences (labelled 1, 2, 3 and 4) given below, when properly sequenced, would yield a coherent paragraph. Decide on the proper sequencing of the order of the sentences and key in the sequence of the four numbers as your answer.**
1. What precisely are the “unusual elements” that make a particular case so attractive to a certain kind of audience?
2 . It might be a particularly savage or unfathomable level of depravity, very often it has something to do with the precise amount of mystery involved.
3. Unsolved, and perhaps unsolvable cases offer something that “ordinary” murder doesn’t.
4. Why are some crimes destined for perpetual re-examination and others locked into permanent obscurity?
Backspace  
789  
456  
123  
0.-  
Clear All  

Submit
[](https://cracku.in/21-the-four-sentences-labelled-1-2-3-and-4-given-belo-x-cat-2023-slot-1)

\begin{enumerate}[label=\Alph*.]
\end{enumerate}
\vspace{1em}
\textbf{Solution:}

The correct order is 4-1-2-3.
Sentence 4 introduces the central question about the enduring fascination with certain crimes and the perpetual re-examination of some cases.
Sentence 1 builds on this question by specifically asking about the "unusual elements" that make a particular case attractive to a certain audience, providing a more focused inquiry.
Sentence 2 follows by suggesting possible reasons, such as a particularly savage or unfathomable level of depravity, and the role of mystery in drawing attention to these cases.
Sentence 3 then generalizes the idea, stating that unsolved and perhaps unsolvable cases offer something unique that "ordinary" murder cases don't, emphasizing the enduring allure of unresolved mysteries in the realm of crime.

\end{problem}

\newpage
\begin{problem}{}{}
% Subject: varc
\textbf{Instructions}

Instructions   

**The passage below is accompanied by four questions. Based on the passage, choose the best answer for each question.**
For early postcolonial literature, the world of the novel was often the nation. Postcolonial novels were usually [concerned with] national questions. Sometimes the whole story of the novel was taken as an allegory of the nation, whether India or Tanzania. This was important for supporting anti-colonial nationalism, but could also be limiting - land-focused and inward-looking.
My new book “Writing Ocean Worlds” explores another kind of world of the novel: not the village or nation, but the Indian Ocean world. The book describes a set of novels in which the Indian Ocean is at the centre of the story. It focuses on the novelists Amitav Ghosh, Abdulrazak Gurnah, Lindsey Collen and Joseph Conrad [who have] centred the Indian Ocean world in the majority of their novels. . . . Their work reveals a world that is outward-looking - full of movement, border-crossing and south-south interconnection. They are all very different - from colonially inclined (Conrad) to radically anti-capitalist (Collen), but together draw on and shape a wider sense of Indian Ocean space through themes, images, metaphors and language. This has the effect of remapping the world in the reader’s mind, as centred in the interconnected global south. . . .
The Indian Ocean world is a term used to describe the very long-lasting connections among the coasts of East Africa, the Arab coasts, and South and East Asia. These connections were made possible by the geography of the Indian Ocean. For much of history, travel by sea was much easier than by land, which meant that port cities very far apart were often more easily connected to each other than to much closer inland cities. Historical and archaeological evidence suggests that what we now call globalisation first appeared in the Indian Ocean. This is the interconnected oceanic world referenced and produced by the novels in my book. . . .
For their part Ghosh, Gurnah, Collen and even Conrad reference a different set of histories and geographies than the ones most commonly found in fiction in English. Those [commonly found ones] are mostly centred in Europe or the US, assume a background of Christianity and whiteness, and mention places like Paris and New York. The novels in [my] book highlight instead a largely Islamic space, feature characters of colour and centralise the ports of Malindi, Mombasa, Aden, Java and Bombay. . . . It is a densely imagined, richly sensory image of a southern cosmopolitan culture which provides for an enlarged sense of place in the world.
This remapping is particularly powerful for the representation of Africa. In the fiction, sailors and travellers are not all European. . . . African, as well as Indian and Arab characters, are traders, nakhodas (dhow ship captains), runaways, villains, missionaries and activists. This does not mean that Indian Ocean Africa is romanticised. Migration is often a matter of force; travel is portrayed as abandonment rather than adventure, freedoms are kept from women and slavery is rife. What it does mean is that the African part of the Indian Ocean world plays an active role in its long, rich history and therefore in that of the wider world.

\vspace{1em}
\textbf{Question 22}

The four sentences (labelled 1, 2, 3 and 4) given below, when properly sequenced, would yield a coherent paragraph. Decide on the proper sequencing of the order of the sentences and key in the sequence of the four numbers as your answer.
1. Algorithms hosted on the internet are accessed by many, so biases in AI models have resulted in much larger impact, adversely affecting far larger groups of people.  
2. Though “algorithmic bias” is the popular term, the foundation of such bias is not in algorithms, but in the data; algorithms are not biased, data is, as algorithms merely reflect persistent patterns that are present in the training data.  
3. Despite their widespread impact, it is relatively easier to fix AI biases than human-generated biases, as it is simpler to identify the former than to try to make people unlearn behaviors learnt over generations.  
4. The impact of biased decisions made by humans is localised and geographically confined, but with the advent of AI, the impact of such decisions is spread over a much wider scale.
Backspace  
789  
456  
123  
0.-  
Clear All  

Submit
[](https://cracku.in/22-the-four-sentences-labelled-1-2-3-and-4-given-belo-x-cat-2023-slot-1)

\begin{enumerate}[label=\Alph*.]
\end{enumerate}
\vspace{1em}
\textbf{Solution:}

4-1-2-3 is the correct order.
Sentence 4 introduces the idea that biased decisions made by humans have a localized impact, but with AI's emergence, the impact scale becomes much broader. This sets the context for the discussion to be followed.
Now, if we consider Sentences 1 and 2 we can see that Sentence 1 is expanding on the impact of biased decisions in the context of AI, highlighting that algorithms hosted on the internet, accessed by many, result in more significant adverse effects on larger groups of people.Sentence 2 is providing a clarification on the term "algorithmic bias," emphasizing that the foundation of bias lies in the data rather than the algorithms themselves. It explains that algorithms reflect persistent patterns present in the training data. Therefore we can say that Sentence 2 must be following 1.
Sentence 3 is pointing out that despite the widespread impact of AI biases, it is comparatively easier to fix them than human-generated biases, as it is simpler to identify and address biases in algorithms than to make people unlearn behaviors learned over generations. This logically follows the point made in Sentence 2. Therefore Sentence 3 must be following Sentence 2.
Therefore the correct order is 4-1-2-3.

\end{problem}

\newpage
\begin{problem}{}{}
% Subject: varc
\textbf{Instructions}

Instructions   

**The passage below is accompanied by four questions. Based on the passage, choose the best answer for each question.**
For early postcolonial literature, the world of the novel was often the nation. Postcolonial novels were usually [concerned with] national questions. Sometimes the whole story of the novel was taken as an allegory of the nation, whether India or Tanzania. This was important for supporting anti-colonial nationalism, but could also be limiting - land-focused and inward-looking.
My new book “Writing Ocean Worlds” explores another kind of world of the novel: not the village or nation, but the Indian Ocean world. The book describes a set of novels in which the Indian Ocean is at the centre of the story. It focuses on the novelists Amitav Ghosh, Abdulrazak Gurnah, Lindsey Collen and Joseph Conrad [who have] centred the Indian Ocean world in the majority of their novels. . . . Their work reveals a world that is outward-looking - full of movement, border-crossing and south-south interconnection. They are all very different - from colonially inclined (Conrad) to radically anti-capitalist (Collen), but together draw on and shape a wider sense of Indian Ocean space through themes, images, metaphors and language. This has the effect of remapping the world in the reader’s mind, as centred in the interconnected global south. . . .
The Indian Ocean world is a term used to describe the very long-lasting connections among the coasts of East Africa, the Arab coasts, and South and East Asia. These connections were made possible by the geography of the Indian Ocean. For much of history, travel by sea was much easier than by land, which meant that port cities very far apart were often more easily connected to each other than to much closer inland cities. Historical and archaeological evidence suggests that what we now call globalisation first appeared in the Indian Ocean. This is the interconnected oceanic world referenced and produced by the novels in my book. . . .
For their part Ghosh, Gurnah, Collen and even Conrad reference a different set of histories and geographies than the ones most commonly found in fiction in English. Those [commonly found ones] are mostly centred in Europe or the US, assume a background of Christianity and whiteness, and mention places like Paris and New York. The novels in [my] book highlight instead a largely Islamic space, feature characters of colour and centralise the ports of Malindi, Mombasa, Aden, Java and Bombay. . . . It is a densely imagined, richly sensory image of a southern cosmopolitan culture which provides for an enlarged sense of place in the world.
This remapping is particularly powerful for the representation of Africa. In the fiction, sailors and travellers are not all European. . . . African, as well as Indian and Arab characters, are traders, nakhodas (dhow ship captains), runaways, villains, missionaries and activists. This does not mean that Indian Ocean Africa is romanticised. Migration is often a matter of force; travel is portrayed as abandonment rather than adventure, freedoms are kept from women and slavery is rife. What it does mean is that the African part of the Indian Ocean world plays an active role in its long, rich history and therefore in that of the wider world.

\vspace{1em}
\textbf{Question 23}

The passage given below is followed by four alternate summaries. Choose the option that best captures the essence of the passage.
Manipulating information was a feature of history long before modern journalism established rules of integrity. A record dates back to ancient Rome, when Antony met Cleopatra and his political enemy Octavian launched a smear campaign against him with “short, sharp slogans written upon coins.” The perpetrator became the first Roman Emperor and “fake news had allowed Octavian to hack the republican system once and for all”. But the 21st century has seen the weaponization of information on an unprecedented scale. Powerful new technology makes the fabrication of content simple, and social networks amplify falsehoods peddled by States, populist politicians, and dishonest corporate entities. The platforms have become fertile ground for computational propaganda, ‘trolling’ and ‘troll armies’.

\begin{enumerate}[label=\Alph*.]
    \item Disinformation, which is mediated by technology today, is not new and has existed since ancient times.
    \item People need to become critical of what they read, since historically, weaponization of information has led to corruption.
    \item Use of misinformation for attaining power, a practice that is as old as the Octavian era, is currently fueled by technology.
    \item Octavian used fake news to manipulate people and attain power and influence, just as people do today
[](https://cracku.in/23-the-passage-given-below-is-followed-by-four-altern-x-cat-2023-slot-1)
\end{enumerate}
\vspace{1em}
\textbf{Solution:}

The passage discusses the historical use of misinformation for political purposes, dating back to ancient Rome with Octavian's smear campaign against Antony. It then highlights how the 21st century has seen an unprecedented scale of information weaponization, facilitated by powerful technology and amplified through social networks. Option C effectively conveys the continuity of using misinformation for power throughout history, now fueled by modern technology. Therefore, Option C is the correct answer.
Option A: While the passage acknowledges the historical aspect of disinformation, it emphasizes the unprecedented scale in the 21st century, which is not captured in this option.
Option B focuses on the need for critical reading without explicitly highlighting the historical context and the weaponization of information for power.
Option D does not emphasize the broader historical and contemporary context of misinformation for political purposes.

\end{problem}

\newpage
\begin{problem}{}{}
% Subject: varc
\textbf{Instructions}

Instructions   

**The passage below is accompanied by four questions. Based on the passage, choose the best answer for each question.**
For early postcolonial literature, the world of the novel was often the nation. Postcolonial novels were usually [concerned with] national questions. Sometimes the whole story of the novel was taken as an allegory of the nation, whether India or Tanzania. This was important for supporting anti-colonial nationalism, but could also be limiting - land-focused and inward-looking.
My new book “Writing Ocean Worlds” explores another kind of world of the novel: not the village or nation, but the Indian Ocean world. The book describes a set of novels in which the Indian Ocean is at the centre of the story. It focuses on the novelists Amitav Ghosh, Abdulrazak Gurnah, Lindsey Collen and Joseph Conrad [who have] centred the Indian Ocean world in the majority of their novels. . . . Their work reveals a world that is outward-looking - full of movement, border-crossing and south-south interconnection. They are all very different - from colonially inclined (Conrad) to radically anti-capitalist (Collen), but together draw on and shape a wider sense of Indian Ocean space through themes, images, metaphors and language. This has the effect of remapping the world in the reader’s mind, as centred in the interconnected global south. . . .
The Indian Ocean world is a term used to describe the very long-lasting connections among the coasts of East Africa, the Arab coasts, and South and East Asia. These connections were made possible by the geography of the Indian Ocean. For much of history, travel by sea was much easier than by land, which meant that port cities very far apart were often more easily connected to each other than to much closer inland cities. Historical and archaeological evidence suggests that what we now call globalisation first appeared in the Indian Ocean. This is the interconnected oceanic world referenced and produced by the novels in my book. . . .
For their part Ghosh, Gurnah, Collen and even Conrad reference a different set of histories and geographies than the ones most commonly found in fiction in English. Those [commonly found ones] are mostly centred in Europe or the US, assume a background of Christianity and whiteness, and mention places like Paris and New York. The novels in [my] book highlight instead a largely Islamic space, feature characters of colour and centralise the ports of Malindi, Mombasa, Aden, Java and Bombay. . . . It is a densely imagined, richly sensory image of a southern cosmopolitan culture which provides for an enlarged sense of place in the world.
This remapping is particularly powerful for the representation of Africa. In the fiction, sailors and travellers are not all European. . . . African, as well as Indian and Arab characters, are traders, nakhodas (dhow ship captains), runaways, villains, missionaries and activists. This does not mean that Indian Ocean Africa is romanticised. Migration is often a matter of force; travel is portrayed as abandonment rather than adventure, freedoms are kept from women and slavery is rife. What it does mean is that the African part of the Indian Ocean world plays an active role in its long, rich history and therefore in that of the wider world.

\vspace{1em}
\textbf{Question 24}

The passage given below is followed by four alternate summaries. Choose the option that best captures the essence of the passage.
Colonialism is not a modern phenomenon. World history is full of examples of one society gradually expanding by incorporating adjacent territory and settling its people on newly conquered territory. In the sixteenth century, colonialism changed decisively because of technological developments in navigation that began to connect more remote parts of the world. The modern European colonial project emerged when it became possible to move large numbers of people across the ocean and to maintain political control in spite of geographical dispersion. The term colonialism is used to describe the process of European settlement, violent dispossession and political  
domination over the rest of the world, including the Americas, Australia, and parts of Africa and Asia.

\begin{enumerate}[label=\Alph*.]
    \item As a result of developments in navigation technology, European colonialism led to the displacement of indigenous populations and global political changes in the 16th century.
    \item Colonialism, conceptualized in the 16th century, allowed colonizers to expand their territories, establish settlements, and exercise political power.
    \item Technological advancements in navigation in the 16th century, transformed colonialism, enabling Europeans to establish settlements and exert political dominance over distant regions.
    \item Colonialism surged in the 16th century due to advancements in navigation, enabling British settlements abroad and global dominance.
[](https://cracku.in/24-the-passage-given-below-is-followed-by-four-altern-x-cat-2023-slot-1)
\end{enumerate}
\vspace{1em}
\textbf{Solution:}

Option C is the correct answer because it accurately captures the main idea of the passage. It highlights how technological advancements in navigation during the sixteenth century transformed colonialism by enabling Europeans to establish settlements and exert political dominance over distant regions, including the Americas, Australia, and parts of Africa and Asia.
Option A focuses on the displacement of indigenous populations, which is not the central point of the passage.
While Option B mentions the expansion of territories and political power, it does not emphasize the technological advancements in navigation.
Option D introduces the concept of British settlements, which is narrower than the broader context of European colonialism discussed in the passage.
[](https://cracku.in/downloads/18891)
  *  
  *

\end{problem}

\newpage
\begin{problem}{}{}
% Subject: varc
\textbf{Instructions}

Instructions   

The passage below is accompanied by four questions. Based on the passage, choose the best answer for each question.
The Positivists, anxious to stake out their claim for history as a science, contributed the weight of their influence to the cult of facts. First ascertain the facts, said the positivists, then draw your conclusions from them. . . . This is what may [be] called the common-sense view of history. History consists of a corpus of ascertained facts. The facts are available to the historian in documents, inscriptions, and so on . . . [Sir George Clark] contrasted the "hard core of facts" in history with the surrounding pulp of disputable interpretation forgetting perhaps that the pulpy part of the fruit is more rewarding than the hard core. . . . It recalls the favourite dictum of the great liberal journalist C. P. Scott: "Facts are sacred, opinion is free.". . .
What is a historical fact? . . . According to the common-sense view, there are certain basic facts which are the same for all historians and which form, so to speak, the backbone of history—the fact, for example, that the Battle of Hastings was fought in 1066. But this view calls for two observations. In the first place, it is not with facts like these that the historian is primarily concerned. It is no doubt important to know that the great battle was fought in 1066 and not in 1065 or 1067, and that it was fought at Hastings and not at Eastbourne or Brighton. The historian must not get these things wrong. But [to] praise a historian for his accuracy is like praising an architect for using well-seasoned timber or properly mixed concrete in his building. It is a necessary condition of his work, but not his essential function. It is precisely for matters of this kind that the historian is entitled to rely on what have been called the "auxiliary sciences" of history—archaeology, epigraphy, numismatics, chronology, and so forth. . . .
The second observation is that the necessity to establish these basic facts rests not on any quality in the facts themselves, but on an apriori decision of the historian. In spite of C. P. Scott's motto, every journalist knows today that the most effective way to influence opinion is by the selection and arrangement of the appropriate facts. It used to be said that facts speak for themselves. This is, of course, untrue. The facts speak only when the historian calls on them: it is he who decides to which facts to give the floor, and in what order or context. . . . The only reason why we are interested to know that the battle was fought at Hastings in 1066 is that historians regard it as a major historical event. . . . Professor Talcott Parsons once called [science] "a selective system of cognitive orientations to reality." It might perhaps have been put more simply. But history is, among other things, that. The historian is necessarily selective. The belief in a hard core of historical facts existing objectively and independently of the interpretation of the historian is a preposterous fallacy, but one which it is very hard to eradicate.

\vspace{1em}
\textbf{Question 1}

According to this passage, which one of the following statements best describes the significance of archaeology for historians?

\begin{enumerate}[label=\Alph*.]
    \item Archaeology helps historians to locate the oldest civilisations in history.
    \item Archaeology helps historians to ascertain factual accuracy.
    \item Archaeology helps historians to carry out their primary duty.
    \item Archaeology helps historians to interpret historical facts.
[](https://cracku.in/1-according-to-this-passage-which-one-of-the-followi-x-cat-2023-slot-2)
\end{enumerate}
\vspace{1em}
\textbf{Solution:}

The passage suggests that historians can rely on disciplines such as archaeology, among others, to establish basic facts. The relevant portion of the passage is:
_"But [to] praise a historian for his accuracy is like praising an architect for using well-seasoned timber or properly mixed concrete in his building. It is a necessary condition of his work, but not his essential function. It is precisely for matters of this kind that the historian is entitled to rely on what have been called the 'auxiliary sciences' of history—archaeology, epigraphy, numismatics, chronology, and so forth."_
In this context, the "auxiliary sciences" are mentioned as tools that historians can use to ensure the accuracy of basic facts. Archaeology is included in this list, suggesting that it helps historians in ascertaining factual accuracy by providing evidence from material remains, artifacts, and other archaeological findings. Therefore, Option B correctly reflects the role of archaeology in supporting historians in their pursuit of factual accuracy.

\end{problem}

\newpage
\begin{problem}{}{}
% Subject: varc
\textbf{Instructions}

Instructions   

The passage below is accompanied by four questions. Based on the passage, choose the best answer for each question.
The Positivists, anxious to stake out their claim for history as a science, contributed the weight of their influence to the cult of facts. First ascertain the facts, said the positivists, then draw your conclusions from them. . . . This is what may [be] called the common-sense view of history. History consists of a corpus of ascertained facts. The facts are available to the historian in documents, inscriptions, and so on . . . [Sir George Clark] contrasted the "hard core of facts" in history with the surrounding pulp of disputable interpretation forgetting perhaps that the pulpy part of the fruit is more rewarding than the hard core. . . . It recalls the favourite dictum of the great liberal journalist C. P. Scott: "Facts are sacred, opinion is free.". . .
What is a historical fact? . . . According to the common-sense view, there are certain basic facts which are the same for all historians and which form, so to speak, the backbone of history—the fact, for example, that the Battle of Hastings was fought in 1066. But this view calls for two observations. In the first place, it is not with facts like these that the historian is primarily concerned. It is no doubt important to know that the great battle was fought in 1066 and not in 1065 or 1067, and that it was fought at Hastings and not at Eastbourne or Brighton. The historian must not get these things wrong. But [to] praise a historian for his accuracy is like praising an architect for using well-seasoned timber or properly mixed concrete in his building. It is a necessary condition of his work, but not his essential function. It is precisely for matters of this kind that the historian is entitled to rely on what have been called the "auxiliary sciences" of history—archaeology, epigraphy, numismatics, chronology, and so forth. . . .
The second observation is that the necessity to establish these basic facts rests not on any quality in the facts themselves, but on an apriori decision of the historian. In spite of C. P. Scott's motto, every journalist knows today that the most effective way to influence opinion is by the selection and arrangement of the appropriate facts. It used to be said that facts speak for themselves. This is, of course, untrue. The facts speak only when the historian calls on them: it is he who decides to which facts to give the floor, and in what order or context. . . . The only reason why we are interested to know that the battle was fought at Hastings in 1066 is that historians regard it as a major historical event. . . . Professor Talcott Parsons once called [science] "a selective system of cognitive orientations to reality." It might perhaps have been put more simply. But history is, among other things, that. The historian is necessarily selective. The belief in a hard core of historical facts existing objectively and independently of the interpretation of the historian is a preposterous fallacy, but one which it is very hard to eradicate.

\vspace{1em}
\textbf{Question 2}

All of the following, if true, can weaken the passage’s claim that facts do not speak for themselves, EXCEPT:

\begin{enumerate}[label=\Alph*.]
    \item the truth value of a fact is independent of the historian who expresses it.
    \item a fact, by its very nature, is objective and universal, irrespective of the context in which it is placed.
    \item facts, like truth, can be relative: what is fact for person X may not be so for person Y.
    \item the order in which a series of facts is presented does not have any bearing on the production of meaning.
[](https://cracku.in/2-all-of-the-following-if-true-can-weaken-the-passag-x-cat-2023-slot-2)
\end{enumerate}
\vspace{1em}
\textbf{Solution:}

Option C is the correct answer because it aligns with the passage's perspective that the interpretation of facts is subjective and can be influenced by different perspectives. The passage argues that historians play a crucial role in selecting and interpreting facts, and Option C supports this idea by suggesting that facts, like truth, can be relative.
If facts are relative, it means that what one person considers a fact may not be viewed the same way by another person. This relativity of facts supports the notion that the historian's interpretation and perspective play a significant role in determining what is considered a fact. Therefore, Option C, if true, reinforces the passage's claim that facts are not entirely objective and independent of the historian's perspective, and it does not weaken the passage's argument.

\end{problem}

\newpage
\begin{problem}{}{}
% Subject: varc
\textbf{Instructions}

Instructions   

The passage below is accompanied by four questions. Based on the passage, choose the best answer for each question.
The Positivists, anxious to stake out their claim for history as a science, contributed the weight of their influence to the cult of facts. First ascertain the facts, said the positivists, then draw your conclusions from them. . . . This is what may [be] called the common-sense view of history. History consists of a corpus of ascertained facts. The facts are available to the historian in documents, inscriptions, and so on . . . [Sir George Clark] contrasted the "hard core of facts" in history with the surrounding pulp of disputable interpretation forgetting perhaps that the pulpy part of the fruit is more rewarding than the hard core. . . . It recalls the favourite dictum of the great liberal journalist C. P. Scott: "Facts are sacred, opinion is free.". . .
What is a historical fact? . . . According to the common-sense view, there are certain basic facts which are the same for all historians and which form, so to speak, the backbone of history—the fact, for example, that the Battle of Hastings was fought in 1066. But this view calls for two observations. In the first place, it is not with facts like these that the historian is primarily concerned. It is no doubt important to know that the great battle was fought in 1066 and not in 1065 or 1067, and that it was fought at Hastings and not at Eastbourne or Brighton. The historian must not get these things wrong. But [to] praise a historian for his accuracy is like praising an architect for using well-seasoned timber or properly mixed concrete in his building. It is a necessary condition of his work, but not his essential function. It is precisely for matters of this kind that the historian is entitled to rely on what have been called the "auxiliary sciences" of history—archaeology, epigraphy, numismatics, chronology, and so forth. . . .
The second observation is that the necessity to establish these basic facts rests not on any quality in the facts themselves, but on an apriori decision of the historian. In spite of C. P. Scott's motto, every journalist knows today that the most effective way to influence opinion is by the selection and arrangement of the appropriate facts. It used to be said that facts speak for themselves. This is, of course, untrue. The facts speak only when the historian calls on them: it is he who decides to which facts to give the floor, and in what order or context. . . . The only reason why we are interested to know that the battle was fought at Hastings in 1066 is that historians regard it as a major historical event. . . . Professor Talcott Parsons once called [science] "a selective system of cognitive orientations to reality." It might perhaps have been put more simply. But history is, among other things, that. The historian is necessarily selective. The belief in a hard core of historical facts existing objectively and independently of the interpretation of the historian is a preposterous fallacy, but one which it is very hard to eradicate.

\vspace{1em}
\textbf{Question 3}

If the author of the passage were to write a book on the Battle of Hastings along the lines of his/her own reasoning, the focus of the historical account would be on:

\begin{enumerate}[label=\Alph*.]
    \item exploring the socio-political and economic factors that led to the Battle.
    \item producing a detailed timeline of the various events that led to the Battle.
    \item deriving historical facts from the relevant documents and inscriptions.
    \item providing a nuanced interpretation by relying on the auxiliary sciences.
[](https://cracku.in/3-if-the-author-of-the-passage-were-to-write-a-book--x-cat-2023-slot-2)
\end{enumerate}
\vspace{1em}
\textbf{Solution:}

The passage suggests that while establishing basic facts is necessary, the essential function of historians goes beyond this. It emphasizes the selective and interpretive nature of historical writing, where historians are expected to go deeper into understanding the context and motivations behind historical events.
Option A, which involves exploring the socio-political and economic factors that led to the Battle of Hastings, is in line with the idea of providing a nuanced interpretation. This option indicates a focus on understanding the underlying causes and influences that shaped the historical event, reflecting a more comprehensive and contextual approach to historical writing.
Option B: While timelines are important, the author suggests that the historian's essential function goes beyond establishing basic chronological facts.
Option C: The author acknowledges the importance of basic facts but argues that historians must go beyond mere fact-finding.
Option D: While the author acknowledges the use of auxiliary sciences, it also suggests that the historian's focus should go beyond relying solely on these sciences for basic facts, emphasizing the historian's selective and interpretive role in presenting historical events.

\end{problem}

\newpage
\begin{problem}{}{}
% Subject: varc
\textbf{Instructions}

Instructions   

The passage below is accompanied by four questions. Based on the passage, choose the best answer for each question.
The Positivists, anxious to stake out their claim for history as a science, contributed the weight of their influence to the cult of facts. First ascertain the facts, said the positivists, then draw your conclusions from them. . . . This is what may [be] called the common-sense view of history. History consists of a corpus of ascertained facts. The facts are available to the historian in documents, inscriptions, and so on . . . [Sir George Clark] contrasted the "hard core of facts" in history with the surrounding pulp of disputable interpretation forgetting perhaps that the pulpy part of the fruit is more rewarding than the hard core. . . . It recalls the favourite dictum of the great liberal journalist C. P. Scott: "Facts are sacred, opinion is free.". . .
What is a historical fact? . . . According to the common-sense view, there are certain basic facts which are the same for all historians and which form, so to speak, the backbone of history—the fact, for example, that the Battle of Hastings was fought in 1066. But this view calls for two observations. In the first place, it is not with facts like these that the historian is primarily concerned. It is no doubt important to know that the great battle was fought in 1066 and not in 1065 or 1067, and that it was fought at Hastings and not at Eastbourne or Brighton. The historian must not get these things wrong. But [to] praise a historian for his accuracy is like praising an architect for using well-seasoned timber or properly mixed concrete in his building. It is a necessary condition of his work, but not his essential function. It is precisely for matters of this kind that the historian is entitled to rely on what have been called the "auxiliary sciences" of history—archaeology, epigraphy, numismatics, chronology, and so forth. . . .
The second observation is that the necessity to establish these basic facts rests not on any quality in the facts themselves, but on an apriori decision of the historian. In spite of C. P. Scott's motto, every journalist knows today that the most effective way to influence opinion is by the selection and arrangement of the appropriate facts. It used to be said that facts speak for themselves. This is, of course, untrue. The facts speak only when the historian calls on them: it is he who decides to which facts to give the floor, and in what order or context. . . . The only reason why we are interested to know that the battle was fought at Hastings in 1066 is that historians regard it as a major historical event. . . . Professor Talcott Parsons once called [science] "a selective system of cognitive orientations to reality." It might perhaps have been put more simply. But history is, among other things, that. The historian is necessarily selective. The belief in a hard core of historical facts existing objectively and independently of the interpretation of the historian is a preposterous fallacy, but one which it is very hard to eradicate.

\vspace{1em}
\textbf{Question 4}

All of the following describe the “common-sense view” of history, EXCEPT:

\begin{enumerate}[label=\Alph*.]
    \item real history can be found in ancient engravings and archival documents.
    \item only the positivist methods can lead to credible historical knowledge.
    \item history is like science: a selective system of cognitive orientations to reality.
    \item history can be objective like the sciences if it is derived from historical facts.
[](https://cracku.in/4-all-of-the-following-describe-the-common-sense-vie-x-cat-2023-slot-2)
\end{enumerate}
\vspace{1em}
\textbf{Solution:}

“ _According to the common-sense view, there are certain basic facts which are the same for all historians and which form, so to speak, the backbone of history:_
The common-sense view of history, as described in the passage, emphasizes that history consists of a body of ascertained facts found in documents and engravings. In contrast, Option C suggests that history is like science, framing it as a selective system of cognitive orientations to reality. This characterization does not align with the common-sense view, which is more concerned with factual accuracy and tangible historical evidence rather than presenting history as a selective system akin to science. Therefore Option C is the correct answer.
Option A accurately reflects the common-sense view as described in the passage.
The common-sense view includes a broader acceptance of historical methods beyond positivism, as mentioned in the passage (option B) and as per the passage, involves a fallacy in believing that historical facts are objective and independent of interpretation (Option D).

Instructions   

The passage below is accompanied by four questions. Based on the passage, choose the best answer for each question.
Umberto Eco, an Italian writer, was right when he said the language of Europe is translation. Netflix and other deep-pocketed global firms speak it well. Just as the EU employs a small army of translators and interpreters to turn intricate laws or impassioned speeches of Romanian MEPs into the EU’s 24 official languages, so do the likes of Netflix. It now offers dubbing in 34 languages and subtitling in a few more. . . . 
The economics of European productions are more appealing, too. American audiences are more willing than before to give dubbed or subtitled viewing a chance. This means shows such as “Lupin”, a French crime caper on Netflix, can become global hits. . . . In 2015, about 75% of Netflix’s original content was American; now the figure is half, according to Ampere, a media-analysis company. Netflix has about 100 productions under way in Europe, which is more than big public broadcasters in France or Germany. . . .
Not everything works across borders. Comedy sometimes struggles. Whodunits and bloodthirsty maelstroms between arch Romans and uppity tribesmen have a more universal appeal. Some do it better than others. Barbarians aside, German television is not always built for export, says one executive, being polite. A bigger problem is that national broadcasters still dominate. Streaming services, such as Netflix or Disney+, account for about a third of all viewing hours, even in markets where they are well-established. Europe is an ageing continent. The generation of teens staring at phones is outnumbered by their elders who prefer to gawp at the box.
In Brussels and national capitals, the prospect of Netflix as a cultural hegemon is seen as a threat. “Cultural sovereignty” is the watchword of European executives worried that the Americans will eat their lunch. To be fair, Netflix content sometimes seems stuck in an uncanny valley somewhere in the mid-Atlantic, with local quirks stripped out. Netflix originals tend to have fewer specific cultural references than shows produced by domestic rivals, according to Enders, a market analyst. The company used to have an imperial model of commissioning, with executives in Los Angeles cooking up ideas French people might like. Now Netflix has offices across Europe. But ultimately the big decisions rest with American executives. This makes European politicians nervous.
They should not be. An irony of European integration is that it is often American companies that facilitate it. Google Translate makes European newspapers comprehensible, even if a little clunky, for the continent’s non-polyglots. American social-media companies make it easier for Europeans to talk politics across borders. (That they do not always like to hear what they say about each other is another matter.) Now Netflix and friends pump the same content into homes across a continent, making culture a cross-border endeavour, too. If Europeans are to share a currency, bail each other out in times of financial need and share vaccines in a pandemic, then they need to have something in common—even if it is just bingeing on the same series. Watching fictitious northern and southern Europeans tear each other apart 2,000 years ago beats doing so in reality.

\end{problem}

\newpage
\begin{problem}{}{}
% Subject: varc
\textbf{Instructions}

Instructions   

The passage below is accompanied by four questions. Based on the passage, choose the best answer for each question.
The Positivists, anxious to stake out their claim for history as a science, contributed the weight of their influence to the cult of facts. First ascertain the facts, said the positivists, then draw your conclusions from them. . . . This is what may [be] called the common-sense view of history. History consists of a corpus of ascertained facts. The facts are available to the historian in documents, inscriptions, and so on . . . [Sir George Clark] contrasted the "hard core of facts" in history with the surrounding pulp of disputable interpretation forgetting perhaps that the pulpy part of the fruit is more rewarding than the hard core. . . . It recalls the favourite dictum of the great liberal journalist C. P. Scott: "Facts are sacred, opinion is free.". . .
What is a historical fact? . . . According to the common-sense view, there are certain basic facts which are the same for all historians and which form, so to speak, the backbone of history—the fact, for example, that the Battle of Hastings was fought in 1066. But this view calls for two observations. In the first place, it is not with facts like these that the historian is primarily concerned. It is no doubt important to know that the great battle was fought in 1066 and not in 1065 or 1067, and that it was fought at Hastings and not at Eastbourne or Brighton. The historian must not get these things wrong. But [to] praise a historian for his accuracy is like praising an architect for using well-seasoned timber or properly mixed concrete in his building. It is a necessary condition of his work, but not his essential function. It is precisely for matters of this kind that the historian is entitled to rely on what have been called the "auxiliary sciences" of history—archaeology, epigraphy, numismatics, chronology, and so forth. . . .
The second observation is that the necessity to establish these basic facts rests not on any quality in the facts themselves, but on an apriori decision of the historian. In spite of C. P. Scott's motto, every journalist knows today that the most effective way to influence opinion is by the selection and arrangement of the appropriate facts. It used to be said that facts speak for themselves. This is, of course, untrue. The facts speak only when the historian calls on them: it is he who decides to which facts to give the floor, and in what order or context. . . . The only reason why we are interested to know that the battle was fought at Hastings in 1066 is that historians regard it as a major historical event. . . . Professor Talcott Parsons once called [science] "a selective system of cognitive orientations to reality." It might perhaps have been put more simply. But history is, among other things, that. The historian is necessarily selective. The belief in a hard core of historical facts existing objectively and independently of the interpretation of the historian is a preposterous fallacy, but one which it is very hard to eradicate.

\vspace{1em}
\textbf{Question 5}

Based on information provided in the passage, all of the following are true, EXCEPT:

\begin{enumerate}[label=\Alph*.]
    \item European television productions have the potential to become global hits.
    \item Netflix has been able to transform itself into a truly European entity.
    \item only half of Netflix’s original programming in the EU is now produced in America.
    \item national broadcasters dominate in the EU in terms of total television viewing hours.
[](https://cracku.in/5-based-on-information-provided-in-the-passage-all-o-x-cat-2023-slot-2)
\end{enumerate}
\vspace{1em}
\textbf{Solution:}

Option B is not supported by the information in the passage. While the passage mentions that Netflix has offices across Europe, it also notes that the big decisions still rest with American executives. The passage suggests that Netflix's content might still exhibit a somewhat mid-Atlantic quality, and the company's executive decisions remain under the control of Americans, making it less accurate to claim that Netflix has completely transformed into a truly European entity.
Option A is supported by the passage, which mentions shows like "Lupin," a French crime caper on Netflix, becoming global hits.
Option C: The passage specifically mentions that, according to Ampere, a media-analysis company, in 2015, about 75% of Netflix’s original content was American, but now the figure is half. This indicates a shift in the geographical distribution of Netflix's original programming, with a decrease in the proportion of American content.
Option D is true as the passage mentions that streaming services like Netflix account for about a third of all viewing hours, challenging the dominance of national broadcasters.

\end{problem}

\newpage
\begin{problem}{}{}
% Subject: varc
\textbf{Instructions}

Instructions   

The passage below is accompanied by four questions. Based on the passage, choose the best answer for each question.
The Positivists, anxious to stake out their claim for history as a science, contributed the weight of their influence to the cult of facts. First ascertain the facts, said the positivists, then draw your conclusions from them. . . . This is what may [be] called the common-sense view of history. History consists of a corpus of ascertained facts. The facts are available to the historian in documents, inscriptions, and so on . . . [Sir George Clark] contrasted the "hard core of facts" in history with the surrounding pulp of disputable interpretation forgetting perhaps that the pulpy part of the fruit is more rewarding than the hard core. . . . It recalls the favourite dictum of the great liberal journalist C. P. Scott: "Facts are sacred, opinion is free.". . .
What is a historical fact? . . . According to the common-sense view, there are certain basic facts which are the same for all historians and which form, so to speak, the backbone of history—the fact, for example, that the Battle of Hastings was fought in 1066. But this view calls for two observations. In the first place, it is not with facts like these that the historian is primarily concerned. It is no doubt important to know that the great battle was fought in 1066 and not in 1065 or 1067, and that it was fought at Hastings and not at Eastbourne or Brighton. The historian must not get these things wrong. But [to] praise a historian for his accuracy is like praising an architect for using well-seasoned timber or properly mixed concrete in his building. It is a necessary condition of his work, but not his essential function. It is precisely for matters of this kind that the historian is entitled to rely on what have been called the "auxiliary sciences" of history—archaeology, epigraphy, numismatics, chronology, and so forth. . . .
The second observation is that the necessity to establish these basic facts rests not on any quality in the facts themselves, but on an apriori decision of the historian. In spite of C. P. Scott's motto, every journalist knows today that the most effective way to influence opinion is by the selection and arrangement of the appropriate facts. It used to be said that facts speak for themselves. This is, of course, untrue. The facts speak only when the historian calls on them: it is he who decides to which facts to give the floor, and in what order or context. . . . The only reason why we are interested to know that the battle was fought at Hastings in 1066 is that historians regard it as a major historical event. . . . Professor Talcott Parsons once called [science] "a selective system of cognitive orientations to reality." It might perhaps have been put more simply. But history is, among other things, that. The historian is necessarily selective. The belief in a hard core of historical facts existing objectively and independently of the interpretation of the historian is a preposterous fallacy, but one which it is very hard to eradicate.

\vspace{1em}
\textbf{Question 6}

The author sees the rise of Netflix in Europe as:

\begin{enumerate}[label=\Alph*.]
    \item a unifying force.
    \item filling an entertainment gap.
    \item a looming cultural threat.
    \item an economic threat.
[](https://cracku.in/6-the-author-sees-the-rise-of-netflix-in-europe-as-x-cat-2023-slot-2)
\end{enumerate}
\vspace{1em}
\textbf{Solution:}

Option A is the correct answer because the passage emphasizes that the rise of Netflix in Europe is seen as a unifying force. The author notes that Netflix and similar streaming services, by pumping the same content into homes across the continent, contribute to making culture a cross-border endeavor. This is described as a shared experience among Europeans, as they binge-watch the same series. The idea is that having a common cultural experience, facilitated by platforms like Netflix, can be a unifying factor among the diverse populations of Europe. Therefore, the rise of Netflix is portrayed in a positive light as a force that brings people together through shared cultural consumption.

\end{problem}

\newpage
\begin{problem}{}{}
% Subject: varc
\textbf{Instructions}

Instructions   

The passage below is accompanied by four questions. Based on the passage, choose the best answer for each question.
The Positivists, anxious to stake out their claim for history as a science, contributed the weight of their influence to the cult of facts. First ascertain the facts, said the positivists, then draw your conclusions from them. . . . This is what may [be] called the common-sense view of history. History consists of a corpus of ascertained facts. The facts are available to the historian in documents, inscriptions, and so on . . . [Sir George Clark] contrasted the "hard core of facts" in history with the surrounding pulp of disputable interpretation forgetting perhaps that the pulpy part of the fruit is more rewarding than the hard core. . . . It recalls the favourite dictum of the great liberal journalist C. P. Scott: "Facts are sacred, opinion is free.". . .
What is a historical fact? . . . According to the common-sense view, there are certain basic facts which are the same for all historians and which form, so to speak, the backbone of history—the fact, for example, that the Battle of Hastings was fought in 1066. But this view calls for two observations. In the first place, it is not with facts like these that the historian is primarily concerned. It is no doubt important to know that the great battle was fought in 1066 and not in 1065 or 1067, and that it was fought at Hastings and not at Eastbourne or Brighton. The historian must not get these things wrong. But [to] praise a historian for his accuracy is like praising an architect for using well-seasoned timber or properly mixed concrete in his building. It is a necessary condition of his work, but not his essential function. It is precisely for matters of this kind that the historian is entitled to rely on what have been called the "auxiliary sciences" of history—archaeology, epigraphy, numismatics, chronology, and so forth. . . .
The second observation is that the necessity to establish these basic facts rests not on any quality in the facts themselves, but on an apriori decision of the historian. In spite of C. P. Scott's motto, every journalist knows today that the most effective way to influence opinion is by the selection and arrangement of the appropriate facts. It used to be said that facts speak for themselves. This is, of course, untrue. The facts speak only when the historian calls on them: it is he who decides to which facts to give the floor, and in what order or context. . . . The only reason why we are interested to know that the battle was fought at Hastings in 1066 is that historians regard it as a major historical event. . . . Professor Talcott Parsons once called [science] "a selective system of cognitive orientations to reality." It might perhaps have been put more simply. But history is, among other things, that. The historian is necessarily selective. The belief in a hard core of historical facts existing objectively and independently of the interpretation of the historian is a preposterous fallacy, but one which it is very hard to eradicate.

\vspace{1em}
\textbf{Question 7}

Based only on information provided in the passage, which one of the following hypothetical Netflix shows would be most successful with audiences across the EU?

\begin{enumerate}[label=\Alph*.]
    \item An Italian comedy show hosted by an international star.
    \item An original German TV science fiction production.
    \item A murder mystery drama set in North Africa and France.
    \item A trans-Atlantic romantic drama set in Europe and America.
[](https://cracku.in/7-based-only-on-information-provided-in-the-passage--x-cat-2023-slot-2)
\end{enumerate}
\vspace{1em}
\textbf{Solution:}

The passage mentions that certain genres, particularly murder mysteries and dramatic conflicts like "bloodthirsty maelstroms between arch Romans and uppity tribesmen," have a more universal appeal. This suggests that themes involving suspense, mystery, and conflicts can transcend cultural boundaries and be attractive to a broader audience.
Therefore, a murder mystery drama set in North Africa and France aligns with the passage's implication that such themes have a more universal appeal, making it likely to be successful with audiences across the diverse countries of the EU. So Option C is the correct answer.
Option A is incorrect as the passage states that comedy does not travel well.
Option B: While science fiction may have a global appeal, the passage emphasizes genres like murder mysteries having a more universal appeal.
Option D: The passage doesn't discuss the appeal of romantic dramas, and the trans-Atlantic setting may not necessarily align with the passage's suggestion that certain themes work better across borders within Europe.

\end{problem}

\newpage
\begin{problem}{}{}
% Subject: varc
\textbf{Instructions}

Instructions   

The passage below is accompanied by four questions. Based on the passage, choose the best answer for each question.
The Positivists, anxious to stake out their claim for history as a science, contributed the weight of their influence to the cult of facts. First ascertain the facts, said the positivists, then draw your conclusions from them. . . . This is what may [be] called the common-sense view of history. History consists of a corpus of ascertained facts. The facts are available to the historian in documents, inscriptions, and so on . . . [Sir George Clark] contrasted the "hard core of facts" in history with the surrounding pulp of disputable interpretation forgetting perhaps that the pulpy part of the fruit is more rewarding than the hard core. . . . It recalls the favourite dictum of the great liberal journalist C. P. Scott: "Facts are sacred, opinion is free.". . .
What is a historical fact? . . . According to the common-sense view, there are certain basic facts which are the same for all historians and which form, so to speak, the backbone of history—the fact, for example, that the Battle of Hastings was fought in 1066. But this view calls for two observations. In the first place, it is not with facts like these that the historian is primarily concerned. It is no doubt important to know that the great battle was fought in 1066 and not in 1065 or 1067, and that it was fought at Hastings and not at Eastbourne or Brighton. The historian must not get these things wrong. But [to] praise a historian for his accuracy is like praising an architect for using well-seasoned timber or properly mixed concrete in his building. It is a necessary condition of his work, but not his essential function. It is precisely for matters of this kind that the historian is entitled to rely on what have been called the "auxiliary sciences" of history—archaeology, epigraphy, numismatics, chronology, and so forth. . . .
The second observation is that the necessity to establish these basic facts rests not on any quality in the facts themselves, but on an apriori decision of the historian. In spite of C. P. Scott's motto, every journalist knows today that the most effective way to influence opinion is by the selection and arrangement of the appropriate facts. It used to be said that facts speak for themselves. This is, of course, untrue. The facts speak only when the historian calls on them: it is he who decides to which facts to give the floor, and in what order or context. . . . The only reason why we are interested to know that the battle was fought at Hastings in 1066 is that historians regard it as a major historical event. . . . Professor Talcott Parsons once called [science] "a selective system of cognitive orientations to reality." It might perhaps have been put more simply. But history is, among other things, that. The historian is necessarily selective. The belief in a hard core of historical facts existing objectively and independently of the interpretation of the historian is a preposterous fallacy, but one which it is very hard to eradicate.

\vspace{1em}
\textbf{Question 8}

Which one of the following research findings would weaken the author’s conclusion in the final paragraph?

\begin{enumerate}[label=\Alph*.]
    \item Research shows there is a wide variance in the popularity and viewing of Netflix shows across different EU countries.
    \item Research shows that Netflix has been gradually losing market share to other streaming television service providers.
    \item Research shows that Netflix hits produced in France are very popular with North American audiences.
    \item Research shows that older women across the EU enjoy watching romantic comedies on Netflix, whereas younger women prefer historical fiction dramas.
[](https://cracku.in/8-which-one-of-the-following-research-findings-would-x-cat-2023-slot-2)
\end{enumerate}
\vspace{1em}
\textbf{Solution:}

The author concludes that the rise of Netflix in Europe is seen as a unifying force, emphasizing shared experiences through common series. If research were to show a wide variance in the popularity and viewing of Netflix shows across different EU countries (Option A), it would suggest that the impact of Netflix on cultural unity is not as consistent or unifying as the author implies. The wide variance could indicate that cultural preferences or barriers within different EU countries limit the effectiveness of Netflix as a unifying force across the entire region.Therefore Option A is the correct answer.
Option B: While this finding could have implications for Netflix's business, it doesn't necessarily address the cultural unifying aspect mentioned in the author's conclusion.
Option C: This finding doesn't directly relate to the author's conclusion about Netflix serving as a unifying force within Eur
Option D: While this finding highlights audience preferences, it doesn't directly address the author's conclusion regarding the cross-border unifying role of Netflix in Europe.

Instructions   

The passage below is accompanied by four questions. Based on the passage, choose the best answer for each question.
The Second Hand September campaign, led by Oxfam . . . seeks to encourage shopping at local organisations and charities as alternatives to fast fashion brands such as Primark and Boohoo in the name of saving our planet. As innocent as mindless scrolling through online shops may seem, such consumers are unintentionally—or perhaps even knowingly —contributing to an industry that uses more energy than aviation. . . .
Brits buy more garments than any other country in Europe, so it comes as no shock that many of those clothes end up in UK landfills each year: 300,000 tonnes of them, to be exact. This waste of clothing is destructive to our planet, releasing greenhouse gasses as clothes are burnt as well as bleeding toxins and dyes into the surrounding soil and water. As ecologist Chelsea Rochman bluntly put it, “The mismanagement of our waste has even come back to haunt us on our dinner plate.”
It’s not surprising, then, that people are scrambling for a solution, the most common of which is second-hand shopping. Retailers selling consigned clothing are currently expanding at a rapid rate . . . If everyone bought just one used item in a year, it would save 449 million lbs of waste, equivalent to the weight of 1 million Polar bears. “Thrifting” has increasingly become a trendy practice. London is home to many second-hand, or more commonly coined ‘vintage’, shops across the city from Bayswater to Brixton.
So you’re cool and you care about the planet; you’ve killed two birds with one stone. But do people simply purchase a second-hand item, flash it on Instagram with #vintage and call it a day without considering whether what they are doing is actually effective?
According to a study commissioned by Patagonia, for instance, older clothes shed more microfibres. These can end up in our rivers and seas after just one wash due to the worn material, thus contributing to microfibre pollution. To break it down, the amount of microfibres released by laundering 100,000 fleece jackets is equivalent to as many as 11,900 plastic grocery bags, and up to 40 per cent of that ends up in our oceans. . . . So where does this leave second-hand consumers? [They would be well advised to buy] high-quality items that shed less and last longer [as this] combats both microfibre pollution and excess garments ending up in landfills. . . .
Luxury brands would rather not circulate their latest season stock around the globe to be sold at a cheaper price, which is why companies like ThredUP, a US fashion resale marketplace, have not yet caught on in the UK. There will always be a market for consignment but there is also a whole generation of people who have been taught that only buying new products is the norm; second-hand luxury goods are not in their psyche. Ben Whitaker, director at Liquidation Firm B-Stock, told Prospect that unless recycling becomes cost-effective and filters into mass production, with the right technology to partner it, “high-end retailers would rather put brand before sustainability.”

\end{problem}

\newpage
\begin{problem}{}{}
% Subject: varc
\textbf{Instructions}

Instructions   

The passage below is accompanied by four questions. Based on the passage, choose the best answer for each question.
The Positivists, anxious to stake out their claim for history as a science, contributed the weight of their influence to the cult of facts. First ascertain the facts, said the positivists, then draw your conclusions from them. . . . This is what may [be] called the common-sense view of history. History consists of a corpus of ascertained facts. The facts are available to the historian in documents, inscriptions, and so on . . . [Sir George Clark] contrasted the "hard core of facts" in history with the surrounding pulp of disputable interpretation forgetting perhaps that the pulpy part of the fruit is more rewarding than the hard core. . . . It recalls the favourite dictum of the great liberal journalist C. P. Scott: "Facts are sacred, opinion is free.". . .
What is a historical fact? . . . According to the common-sense view, there are certain basic facts which are the same for all historians and which form, so to speak, the backbone of history—the fact, for example, that the Battle of Hastings was fought in 1066. But this view calls for two observations. In the first place, it is not with facts like these that the historian is primarily concerned. It is no doubt important to know that the great battle was fought in 1066 and not in 1065 or 1067, and that it was fought at Hastings and not at Eastbourne or Brighton. The historian must not get these things wrong. But [to] praise a historian for his accuracy is like praising an architect for using well-seasoned timber or properly mixed concrete in his building. It is a necessary condition of his work, but not his essential function. It is precisely for matters of this kind that the historian is entitled to rely on what have been called the "auxiliary sciences" of history—archaeology, epigraphy, numismatics, chronology, and so forth. . . .
The second observation is that the necessity to establish these basic facts rests not on any quality in the facts themselves, but on an apriori decision of the historian. In spite of C. P. Scott's motto, every journalist knows today that the most effective way to influence opinion is by the selection and arrangement of the appropriate facts. It used to be said that facts speak for themselves. This is, of course, untrue. The facts speak only when the historian calls on them: it is he who decides to which facts to give the floor, and in what order or context. . . . The only reason why we are interested to know that the battle was fought at Hastings in 1066 is that historians regard it as a major historical event. . . . Professor Talcott Parsons once called [science] "a selective system of cognitive orientations to reality." It might perhaps have been put more simply. But history is, among other things, that. The historian is necessarily selective. The belief in a hard core of historical facts existing objectively and independently of the interpretation of the historian is a preposterous fallacy, but one which it is very hard to eradicate.

\vspace{1em}
\textbf{Question 9}

The central idea of the passage would be undermined if:

\begin{enumerate}[label=\Alph*.]
    \item customers bought all their clothes online.
    \item clothes were not thrown and burnt in landfills
    \item second-hand stores sold only high-quality clothes.
    \item Primark and Boohoo recycled their clothes for vintage stores
[](https://cracku.in/9-the-central-idea-of-the-passage-would-be-undermine-x-cat-2023-slot-2)
\end{enumerate}
\vspace{1em}
\textbf{Solution:}

The central idea of the passage is the promotion of sustainable shopping practices, particularly second-hand shopping, as a means to combat the detrimental environmental effects of the fashion industry. But, the passage also discusses the need for consumers to be mindful of the environmental impact of their clothing choices, **opting for high-quality items** that last longer and shed fewer microfibers.  
The passage argues that opting for second clothing might not always be beneficial for the environment by highlighting the microfibre pollution that they can potentially cause. Now, if the second-hand clothes being sold were only of higher quality, it would take care of this problem (_[They would be well advised to buy] high-quality items that shed less and last longer_** _[as this] combats both microfibre pollution and excess garments ending up in landfills”_** _)  
_ So, the correct answer is Option C.
Option A is more about the purchasing channel than the nature of the clothes so it does not necessarily undermine the central idea of the passage.
Option B supports the central idea by reducing environmental harm.
Option D could align with the sustainability goal and support the central idea, so it doesn't necessarily undermine it.

\end{problem}

\newpage
\begin{problem}{}{}
% Subject: varc
\textbf{Instructions}

Instructions   

The passage below is accompanied by four questions. Based on the passage, choose the best answer for each question.
The Positivists, anxious to stake out their claim for history as a science, contributed the weight of their influence to the cult of facts. First ascertain the facts, said the positivists, then draw your conclusions from them. . . . This is what may [be] called the common-sense view of history. History consists of a corpus of ascertained facts. The facts are available to the historian in documents, inscriptions, and so on . . . [Sir George Clark] contrasted the "hard core of facts" in history with the surrounding pulp of disputable interpretation forgetting perhaps that the pulpy part of the fruit is more rewarding than the hard core. . . . It recalls the favourite dictum of the great liberal journalist C. P. Scott: "Facts are sacred, opinion is free.". . .
What is a historical fact? . . . According to the common-sense view, there are certain basic facts which are the same for all historians and which form, so to speak, the backbone of history—the fact, for example, that the Battle of Hastings was fought in 1066. But this view calls for two observations. In the first place, it is not with facts like these that the historian is primarily concerned. It is no doubt important to know that the great battle was fought in 1066 and not in 1065 or 1067, and that it was fought at Hastings and not at Eastbourne or Brighton. The historian must not get these things wrong. But [to] praise a historian for his accuracy is like praising an architect for using well-seasoned timber or properly mixed concrete in his building. It is a necessary condition of his work, but not his essential function. It is precisely for matters of this kind that the historian is entitled to rely on what have been called the "auxiliary sciences" of history—archaeology, epigraphy, numismatics, chronology, and so forth. . . .
The second observation is that the necessity to establish these basic facts rests not on any quality in the facts themselves, but on an apriori decision of the historian. In spite of C. P. Scott's motto, every journalist knows today that the most effective way to influence opinion is by the selection and arrangement of the appropriate facts. It used to be said that facts speak for themselves. This is, of course, untrue. The facts speak only when the historian calls on them: it is he who decides to which facts to give the floor, and in what order or context. . . . The only reason why we are interested to know that the battle was fought at Hastings in 1066 is that historians regard it as a major historical event. . . . Professor Talcott Parsons once called [science] "a selective system of cognitive orientations to reality." It might perhaps have been put more simply. But history is, among other things, that. The historian is necessarily selective. The belief in a hard core of historical facts existing objectively and independently of the interpretation of the historian is a preposterous fallacy, but one which it is very hard to eradicate.

\vspace{1em}
\textbf{Question 10}

According to the author, companies like ThredUP have not caught on in the UK for all of the following reasons EXCEPT that:

\begin{enumerate}[label=\Alph*.]
    \item recycling is currently not financially attractive for luxury brands.
    \item luxury brands want to maintain their brand image.
    \item luxury brands do not like their product to be devalued.
    \item the British don’t buy second-hand clothing.
[](https://cracku.in/10-according-to-the-author-companies-like-thredup-hav-x-cat-2023-slot-2)
\end{enumerate}
\vspace{1em}
\textbf{Solution:}

Option D is the correct answer because the passage does not mention or suggest that the British don't buy second-hand clothing. Instead, the passage discusses challenges related to luxury brands and their reluctance to circulate their latest season stock globally at a cheaper price. The reasons mentioned include the financial aspect(Option A), concerns about brand image(Option B), and the desire to avoid devaluing their products(Option D). Therefore, the passage does not attribute the slow adoption of companies like ThredUP in the UK to the British not buying second-hand clothing.

\end{problem}

\newpage
\begin{problem}{}{}
% Subject: varc
\textbf{Instructions}

Instructions   

The passage below is accompanied by four questions. Based on the passage, choose the best answer for each question.
The Positivists, anxious to stake out their claim for history as a science, contributed the weight of their influence to the cult of facts. First ascertain the facts, said the positivists, then draw your conclusions from them. . . . This is what may [be] called the common-sense view of history. History consists of a corpus of ascertained facts. The facts are available to the historian in documents, inscriptions, and so on . . . [Sir George Clark] contrasted the "hard core of facts" in history with the surrounding pulp of disputable interpretation forgetting perhaps that the pulpy part of the fruit is more rewarding than the hard core. . . . It recalls the favourite dictum of the great liberal journalist C. P. Scott: "Facts are sacred, opinion is free.". . .
What is a historical fact? . . . According to the common-sense view, there are certain basic facts which are the same for all historians and which form, so to speak, the backbone of history—the fact, for example, that the Battle of Hastings was fought in 1066. But this view calls for two observations. In the first place, it is not with facts like these that the historian is primarily concerned. It is no doubt important to know that the great battle was fought in 1066 and not in 1065 or 1067, and that it was fought at Hastings and not at Eastbourne or Brighton. The historian must not get these things wrong. But [to] praise a historian for his accuracy is like praising an architect for using well-seasoned timber or properly mixed concrete in his building. It is a necessary condition of his work, but not his essential function. It is precisely for matters of this kind that the historian is entitled to rely on what have been called the "auxiliary sciences" of history—archaeology, epigraphy, numismatics, chronology, and so forth. . . .
The second observation is that the necessity to establish these basic facts rests not on any quality in the facts themselves, but on an apriori decision of the historian. In spite of C. P. Scott's motto, every journalist knows today that the most effective way to influence opinion is by the selection and arrangement of the appropriate facts. It used to be said that facts speak for themselves. This is, of course, untrue. The facts speak only when the historian calls on them: it is he who decides to which facts to give the floor, and in what order or context. . . . The only reason why we are interested to know that the battle was fought at Hastings in 1066 is that historians regard it as a major historical event. . . . Professor Talcott Parsons once called [science] "a selective system of cognitive orientations to reality." It might perhaps have been put more simply. But history is, among other things, that. The historian is necessarily selective. The belief in a hard core of historical facts existing objectively and independently of the interpretation of the historian is a preposterous fallacy, but one which it is very hard to eradicate.

\vspace{1em}
\textbf{Question 11}

Based on the passage, we can infer that the opposite of fast fashion, ‘slow fashion’, would most likely refer to clothes that:

\begin{enumerate}[label=\Alph*.]
    \item are of high quality and long lasting.
    \item do not bleed toxins and dyes.
    \item are sold by genuine vintage stores.
    \item do not shed microfibres.
[](https://cracku.in/11-based-on-the-passage-we-can-infer-that-the-opposit-x-cat-2023-slot-2)
\end{enumerate}
\vspace{1em}
\textbf{Solution:}

Option A is the correct answer because the passage emphasizes the environmental issues associated with fast fashion, including the wasteful disposal of garments in landfills. The opposite of this disposable and rapid turnover nature of fast fashion would be a more sustainable and durable approach, which aligns with the idea of "slow fashion."
The passage suggests that buying high-quality items that last longer is a way to combat the negative environmental impact of the fashion industry. Therefore, 'slow fashion' can be inferred to refer to clothes that are of high quality and long-lasting, promoting a more sustainable and environmentally friendly approach to fashion consumption.

\end{problem}

\newpage
\begin{problem}{}{}
% Subject: varc
\textbf{Instructions}

Instructions   

The passage below is accompanied by four questions. Based on the passage, choose the best answer for each question.
The Positivists, anxious to stake out their claim for history as a science, contributed the weight of their influence to the cult of facts. First ascertain the facts, said the positivists, then draw your conclusions from them. . . . This is what may [be] called the common-sense view of history. History consists of a corpus of ascertained facts. The facts are available to the historian in documents, inscriptions, and so on . . . [Sir George Clark] contrasted the "hard core of facts" in history with the surrounding pulp of disputable interpretation forgetting perhaps that the pulpy part of the fruit is more rewarding than the hard core. . . . It recalls the favourite dictum of the great liberal journalist C. P. Scott: "Facts are sacred, opinion is free.". . .
What is a historical fact? . . . According to the common-sense view, there are certain basic facts which are the same for all historians and which form, so to speak, the backbone of history—the fact, for example, that the Battle of Hastings was fought in 1066. But this view calls for two observations. In the first place, it is not with facts like these that the historian is primarily concerned. It is no doubt important to know that the great battle was fought in 1066 and not in 1065 or 1067, and that it was fought at Hastings and not at Eastbourne or Brighton. The historian must not get these things wrong. But [to] praise a historian for his accuracy is like praising an architect for using well-seasoned timber or properly mixed concrete in his building. It is a necessary condition of his work, but not his essential function. It is precisely for matters of this kind that the historian is entitled to rely on what have been called the "auxiliary sciences" of history—archaeology, epigraphy, numismatics, chronology, and so forth. . . .
The second observation is that the necessity to establish these basic facts rests not on any quality in the facts themselves, but on an apriori decision of the historian. In spite of C. P. Scott's motto, every journalist knows today that the most effective way to influence opinion is by the selection and arrangement of the appropriate facts. It used to be said that facts speak for themselves. This is, of course, untrue. The facts speak only when the historian calls on them: it is he who decides to which facts to give the floor, and in what order or context. . . . The only reason why we are interested to know that the battle was fought at Hastings in 1066 is that historians regard it as a major historical event. . . . Professor Talcott Parsons once called [science] "a selective system of cognitive orientations to reality." It might perhaps have been put more simply. But history is, among other things, that. The historian is necessarily selective. The belief in a hard core of historical facts existing objectively and independently of the interpretation of the historian is a preposterous fallacy, but one which it is very hard to eradicate.

\vspace{1em}
\textbf{Question 12}

The act of “thrifting”, as described in the passage, can be considered ironic because it:

\begin{enumerate}[label=\Alph*.]
    \item has created environmental problems.
    \item is not cost-effective for retailers
    \item offers luxury clothing at cut-rate prices.
    \item is an anti-consumerist attitude.
[](https://cracku.in/12-the-act-of-thrifting-as-described-in-the-passage-c-x-cat-2023-slot-2)
\end{enumerate}
\vspace{1em}
\textbf{Solution:}

The irony of "thrifting," as discussed in the passage, is rooted in its unintended environmental consequences. While thrifting is commonly associated with sustainable and eco-friendly practices, the passage highlights a potential environmental issue linked to the act. Specifically, the passage mentions a study commissioned by Patagonia that reveals older clothes, often found in second-hand stores, tend to shed more microfibers. These microfibers can end up in rivers and oceans, contributing to microfiber pollution. Therefore, the seemingly environmentally conscious act of thrifting, aimed at reducing waste, may inadvertently result in environmental problems through the shedding of microfibers during the washing of older garments.Therefore Option A is the correct answer

Instructions   

The passage below is accompanied by four questions. Based on the passage, choose the best answer for each question.
Over the past four centuries liberalism has been so successful that it has driven all its opponents off the battlefield. Now it is disintegrating, destroyed by a mix of hubris and internal contradictions, according to Patrick Deneen, a professor of politics at the University of Notre Dame. . . . Equality of opportunity has produced a new meritocratic aristocracy that has all the aloofness of the old aristocracy with none of its sense of noblesse oblige. Democracy has degenerated into a theatre of the absurd. And technological advances are reducing ever more areas of work into meaningless drudgery. “The gap between liberalism’s claims about itself and the lived reality of the citizenry” is now so wide that “the lie can no longer be accepted,” Mr Deneen writes. What better proof of this than the vision of 1,000 private planes whisking their occupants to Davos to discuss the question of “creating a shared future in a fragmented world”? . . .
Deneen does an impressive job of capturing the current mood of disillusionment, echoing leftwing complaints about rampant commercialism, right-wing complaints about narcissistic and bullying students, and general worries about atomisation and selfishness. But when he concludes that all this adds up to a failure of liberalism, is his argument convincing? . . . He argues that the essence of liberalism lies in freeing individuals from constraints. In fact, liberalism contains a wide range of intellectual traditions which provide different answers to the question of how to trade off the relative claims of rights and responsibilities, individual expression and social ties. . . . liberals experimented with a range of ideas from devolving power from the centre to creating national education systems.
Mr Deneen’s fixation on the essence of liberalism leads to the second big problem of his book: his failure to recognise liberalism’s ability to reform itself and address its internal problems. The late 19th century saw America suffering from many of the problems that are reappearing today, including the creation of a business aristocracy, the rise of vast companies, the corruption of politics and the sense that society was dividing into winners and losers. But a wide variety of reformers, working within the liberal tradition, tackled these problems head on. Theodore Roosevelt took on the trusts. Progressives cleaned up government corruption. University reformers modernised academic syllabuses and built ladders of opportunity. Rather than dying, liberalism reformed itself.
Mr Deneen is right to point out that the record of liberalism in recent years has been dismal. He is also right to assert that the world has much to learn from the premodern notions of liberty as self-mastery and self-denial. The biggest enemy of liberalism is not so much atomisation but old-fashioned greed, as members of the Davos elite pile their plates ever higher with perks and share options. But he is wrong to argue that the only way for people to liberate themselves from the contradictions of liberalism is “liberation from liberalism itself”. The best way to read “Why Liberalism Failed” is not as a funeral oration but as a call to action: up your game, or else.

\end{problem}

\newpage
\begin{problem}{}{}
% Subject: varc
\textbf{Instructions}

Instructions   

The passage below is accompanied by four questions. Based on the passage, choose the best answer for each question.
The Positivists, anxious to stake out their claim for history as a science, contributed the weight of their influence to the cult of facts. First ascertain the facts, said the positivists, then draw your conclusions from them. . . . This is what may [be] called the common-sense view of history. History consists of a corpus of ascertained facts. The facts are available to the historian in documents, inscriptions, and so on . . . [Sir George Clark] contrasted the "hard core of facts" in history with the surrounding pulp of disputable interpretation forgetting perhaps that the pulpy part of the fruit is more rewarding than the hard core. . . . It recalls the favourite dictum of the great liberal journalist C. P. Scott: "Facts are sacred, opinion is free.". . .
What is a historical fact? . . . According to the common-sense view, there are certain basic facts which are the same for all historians and which form, so to speak, the backbone of history—the fact, for example, that the Battle of Hastings was fought in 1066. But this view calls for two observations. In the first place, it is not with facts like these that the historian is primarily concerned. It is no doubt important to know that the great battle was fought in 1066 and not in 1065 or 1067, and that it was fought at Hastings and not at Eastbourne or Brighton. The historian must not get these things wrong. But [to] praise a historian for his accuracy is like praising an architect for using well-seasoned timber or properly mixed concrete in his building. It is a necessary condition of his work, but not his essential function. It is precisely for matters of this kind that the historian is entitled to rely on what have been called the "auxiliary sciences" of history—archaeology, epigraphy, numismatics, chronology, and so forth. . . .
The second observation is that the necessity to establish these basic facts rests not on any quality in the facts themselves, but on an apriori decision of the historian. In spite of C. P. Scott's motto, every journalist knows today that the most effective way to influence opinion is by the selection and arrangement of the appropriate facts. It used to be said that facts speak for themselves. This is, of course, untrue. The facts speak only when the historian calls on them: it is he who decides to which facts to give the floor, and in what order or context. . . . The only reason why we are interested to know that the battle was fought at Hastings in 1066 is that historians regard it as a major historical event. . . . Professor Talcott Parsons once called [science] "a selective system of cognitive orientations to reality." It might perhaps have been put more simply. But history is, among other things, that. The historian is necessarily selective. The belief in a hard core of historical facts existing objectively and independently of the interpretation of the historian is a preposterous fallacy, but one which it is very hard to eradicate.

\vspace{1em}
\textbf{Question 13}

The author of the passage refers to “the Davos elite” to illustrate his views on:

\begin{enumerate}[label=\Alph*.]
    \item the unlikelihood of a return to the liberalism of the past as long as the rich continue to benefit from the decline in liberal values.
    \item the way the debate around liberalism has been captured by the rich who have managed to insulate themselves from economic hardships.
    \item the hypocrisy of the liberal rich, who profess to subscribe to liberal values while cornering most of the wealth.
    \item the fact that the rise in liberalism had led to a greater interest in shared futures from unlikely social classes.
[](https://cracku.in/13-the-author-of-the-passage-refers-to-the-davos-elit-x-cat-2023-slot-2)
\end{enumerate}
\vspace{1em}
\textbf{Solution:}

Option C is the correct answer because the mention of "the Davos elite" in the passage serves to illustrate the perceived hypocrisy of wealthy individuals who profess to adhere to liberal values while simultaneously amassing the majority of the wealth. The author points to the contradiction between the elite's participation in events like discussions on creating a shared future and their exclusive access to privileges, symbolized by the use of private planes. This highlights the critique that the liberal rich, represented by the Davos elite, may not align their actions with the egalitarian ideals they claim to support, emphasizing a discrepancy between rhetoric and behavior.
Option A is incorrect because the passage does not explicitly connect the decline in liberal values to the rich benefiting, but rather to internal contradictions and hubris.
Option B is incorrect as the passage does not focus on how the debate around liberalism is captured by the rich; instead, it critiques the actions of the Davos elite.  
Option D is incorrect because the passage does not suggest that the rise in liberalism has led to greater interest in shared futures from unlikely social classes; rather, it critiques the behavior of the Davos elite as hypocritical.

\end{problem}

\newpage
\begin{problem}{}{}
% Subject: varc
\textbf{Instructions}

Instructions   

The passage below is accompanied by four questions. Based on the passage, choose the best answer for each question.
The Positivists, anxious to stake out their claim for history as a science, contributed the weight of their influence to the cult of facts. First ascertain the facts, said the positivists, then draw your conclusions from them. . . . This is what may [be] called the common-sense view of history. History consists of a corpus of ascertained facts. The facts are available to the historian in documents, inscriptions, and so on . . . [Sir George Clark] contrasted the "hard core of facts" in history with the surrounding pulp of disputable interpretation forgetting perhaps that the pulpy part of the fruit is more rewarding than the hard core. . . . It recalls the favourite dictum of the great liberal journalist C. P. Scott: "Facts are sacred, opinion is free.". . .
What is a historical fact? . . . According to the common-sense view, there are certain basic facts which are the same for all historians and which form, so to speak, the backbone of history—the fact, for example, that the Battle of Hastings was fought in 1066. But this view calls for two observations. In the first place, it is not with facts like these that the historian is primarily concerned. It is no doubt important to know that the great battle was fought in 1066 and not in 1065 or 1067, and that it was fought at Hastings and not at Eastbourne or Brighton. The historian must not get these things wrong. But [to] praise a historian for his accuracy is like praising an architect for using well-seasoned timber or properly mixed concrete in his building. It is a necessary condition of his work, but not his essential function. It is precisely for matters of this kind that the historian is entitled to rely on what have been called the "auxiliary sciences" of history—archaeology, epigraphy, numismatics, chronology, and so forth. . . .
The second observation is that the necessity to establish these basic facts rests not on any quality in the facts themselves, but on an apriori decision of the historian. In spite of C. P. Scott's motto, every journalist knows today that the most effective way to influence opinion is by the selection and arrangement of the appropriate facts. It used to be said that facts speak for themselves. This is, of course, untrue. The facts speak only when the historian calls on them: it is he who decides to which facts to give the floor, and in what order or context. . . . The only reason why we are interested to know that the battle was fought at Hastings in 1066 is that historians regard it as a major historical event. . . . Professor Talcott Parsons once called [science] "a selective system of cognitive orientations to reality." It might perhaps have been put more simply. But history is, among other things, that. The historian is necessarily selective. The belief in a hard core of historical facts existing objectively and independently of the interpretation of the historian is a preposterous fallacy, but one which it is very hard to eradicate.

\vspace{1em}
\textbf{Question 14}

All of the following statements are evidence of the decline of liberalism today, EXCEPT:

\begin{enumerate}[label=\Alph*.]
    \item “And technological advances are reducing ever more areas of work into meaningless drudgery.”
    \item “. . . the creation of a business aristocracy, the rise of vast companies . . .”
    \item “Democracy has degenerated into a theatre of the absurd.”
    \item “‘The gap between liberalism’s claims about itself and the lived reality of the citizenry’ is now so wide that ‘the lie can no longer be accepted,’. . .”
[](https://cracku.in/14-all-of-the-following-statements-are-evidence-of-th-x-cat-2023-slot-2)
\end{enumerate}
\vspace{1em}
\textbf{Solution:}

All the options, other than A, are direct signs of declining or ineffective liberalism.   
Option B:  _Creation of business aristocracy_ , the author in the first paragraph says that liberalism promoted a meritocratic aristocracy and then went ahead to argue why the meritocratic aristocracy is not a good replacement of the old aristocracy. Creation of a business aristocracy and the rise of vast companies are against the ideals of liberalism.   
Option C:  _Democracy has degenerated into a theatre of the absurd,_ this clearly shows the non-functionality of liberalism and is a pretty valid argument for the decline of liberalism.   
Option D:  _The gap between liberalism’s claims about itself and the lived reality of the citizenry’ is now so wide that ‘the lie can no longer be accepted,_ this lines says that the gap between want liberalism asked us to do and what is actually different are two very different thing. This too can be an evidence of liberalism's decline.   
  
Option A:  _And technological advances are reducing ever more areas of work into meaningless drudgery,_ while this line does talk about the technological advancement in a negative sense, it does not necessarily provide evidence of the decline of liberalism per se. Instead, it highlights a potential consequence or critique within the context of technological advances. The negative impact of technology on certain types of work might be seen as a challenge that needs to be addressed within the liberal framework rather than direct evidence of the decline of liberalism.  
The same challenge could be seen at a time when liberalism was prospering and thus is not an evidence of its decline.

\end{problem}

\newpage
\begin{problem}{}{}
% Subject: varc
\textbf{Instructions}

Instructions   

The passage below is accompanied by four questions. Based on the passage, choose the best answer for each question.
The Positivists, anxious to stake out their claim for history as a science, contributed the weight of their influence to the cult of facts. First ascertain the facts, said the positivists, then draw your conclusions from them. . . . This is what may [be] called the common-sense view of history. History consists of a corpus of ascertained facts. The facts are available to the historian in documents, inscriptions, and so on . . . [Sir George Clark] contrasted the "hard core of facts" in history with the surrounding pulp of disputable interpretation forgetting perhaps that the pulpy part of the fruit is more rewarding than the hard core. . . . It recalls the favourite dictum of the great liberal journalist C. P. Scott: "Facts are sacred, opinion is free.". . .
What is a historical fact? . . . According to the common-sense view, there are certain basic facts which are the same for all historians and which form, so to speak, the backbone of history—the fact, for example, that the Battle of Hastings was fought in 1066. But this view calls for two observations. In the first place, it is not with facts like these that the historian is primarily concerned. It is no doubt important to know that the great battle was fought in 1066 and not in 1065 or 1067, and that it was fought at Hastings and not at Eastbourne or Brighton. The historian must not get these things wrong. But [to] praise a historian for his accuracy is like praising an architect for using well-seasoned timber or properly mixed concrete in his building. It is a necessary condition of his work, but not his essential function. It is precisely for matters of this kind that the historian is entitled to rely on what have been called the "auxiliary sciences" of history—archaeology, epigraphy, numismatics, chronology, and so forth. . . .
The second observation is that the necessity to establish these basic facts rests not on any quality in the facts themselves, but on an apriori decision of the historian. In spite of C. P. Scott's motto, every journalist knows today that the most effective way to influence opinion is by the selection and arrangement of the appropriate facts. It used to be said that facts speak for themselves. This is, of course, untrue. The facts speak only when the historian calls on them: it is he who decides to which facts to give the floor, and in what order or context. . . . The only reason why we are interested to know that the battle was fought at Hastings in 1066 is that historians regard it as a major historical event. . . . Professor Talcott Parsons once called [science] "a selective system of cognitive orientations to reality." It might perhaps have been put more simply. But history is, among other things, that. The historian is necessarily selective. The belief in a hard core of historical facts existing objectively and independently of the interpretation of the historian is a preposterous fallacy, but one which it is very hard to eradicate.

\vspace{1em}
\textbf{Question 15}

The author of the passage is likely to disagree with all of the following statements, EXCEPT:

\begin{enumerate}[label=\Alph*.]
    \item claims about liberalism’s disintegration are exaggerated and misunderstand its core features.
    \item if we accept that liberalism is a dying ideal, we must work to find a viable substitute.
    \item liberalism was the dominant ideal in the past century, but it had to reform itself to remain so.
    \item the essence of liberalism lies in greater individual self-expression and freedoms.
[](https://cracku.in/15-the-author-of-the-passage-is-likely-to-disagree-wi-x-cat-2023-slot-2)
\end{enumerate}
\vspace{1em}
\textbf{Solution:}

The author is likely to agree with the statement in Option C, as it aligns with the author's argument in the passage that liberalism has historically reformed itself in the face of challenges. The author emphasizes the ability of liberalism to address internal problems and reform rather than attributing its success to being the dominant ideal in the past century.
Option A: While the author emphasizes the historical ability of liberalism to reform itself, there is an acknowledgment of its current challenges and failures. Therefore, the author might not fully endorse the idea that claims about liberalism's disintegration are merely exaggerated misunderstandings.
Option B: The author may disagree with this statement as it suggests accepting liberalism's decline and seeking a substitute, which contradicts the call to reform liberalism.
Option D:The author may disagree with this statement as it simplifies the essence of liberalism, which the author argues encompasses various intellectual traditions and responses to the trade-off between rights and responsibilities.

\end{problem}

\newpage
\begin{problem}{}{}
% Subject: varc
\textbf{Instructions}

Instructions   

The passage below is accompanied by four questions. Based on the passage, choose the best answer for each question.
The Positivists, anxious to stake out their claim for history as a science, contributed the weight of their influence to the cult of facts. First ascertain the facts, said the positivists, then draw your conclusions from them. . . . This is what may [be] called the common-sense view of history. History consists of a corpus of ascertained facts. The facts are available to the historian in documents, inscriptions, and so on . . . [Sir George Clark] contrasted the "hard core of facts" in history with the surrounding pulp of disputable interpretation forgetting perhaps that the pulpy part of the fruit is more rewarding than the hard core. . . . It recalls the favourite dictum of the great liberal journalist C. P. Scott: "Facts are sacred, opinion is free.". . .
What is a historical fact? . . . According to the common-sense view, there are certain basic facts which are the same for all historians and which form, so to speak, the backbone of history—the fact, for example, that the Battle of Hastings was fought in 1066. But this view calls for two observations. In the first place, it is not with facts like these that the historian is primarily concerned. It is no doubt important to know that the great battle was fought in 1066 and not in 1065 or 1067, and that it was fought at Hastings and not at Eastbourne or Brighton. The historian must not get these things wrong. But [to] praise a historian for his accuracy is like praising an architect for using well-seasoned timber or properly mixed concrete in his building. It is a necessary condition of his work, but not his essential function. It is precisely for matters of this kind that the historian is entitled to rely on what have been called the "auxiliary sciences" of history—archaeology, epigraphy, numismatics, chronology, and so forth. . . .
The second observation is that the necessity to establish these basic facts rests not on any quality in the facts themselves, but on an apriori decision of the historian. In spite of C. P. Scott's motto, every journalist knows today that the most effective way to influence opinion is by the selection and arrangement of the appropriate facts. It used to be said that facts speak for themselves. This is, of course, untrue. The facts speak only when the historian calls on them: it is he who decides to which facts to give the floor, and in what order or context. . . . The only reason why we are interested to know that the battle was fought at Hastings in 1066 is that historians regard it as a major historical event. . . . Professor Talcott Parsons once called [science] "a selective system of cognitive orientations to reality." It might perhaps have been put more simply. But history is, among other things, that. The historian is necessarily selective. The belief in a hard core of historical facts existing objectively and independently of the interpretation of the historian is a preposterous fallacy, but one which it is very hard to eradicate.

\vspace{1em}
\textbf{Question 16}

The author of the passage faults Deneen’s conclusions for all of the following reasons, EXCEPT:

\begin{enumerate}[label=\Alph*.]
    \item its repeated harking back to premodern notions of liberty.
    \item its failure to note historical instances in which the process of declining liberalism has managed to reverse itself.
    \item its extreme pessimism about the future of liberalism today and predictions of an ultimate decline.
    \item its very narrow definition of liberalism limited to individual freedoms.
[](https://cracku.in/16-the-author-of-the-passage-faults-deneens-conclusio-x-cat-2023-slot-2)
\end{enumerate}
\vspace{1em}
\textbf{Solution:}

Option A is the correct answer because the passage does not explicitly identify Deneen's repeated emphasis on premodern notions of liberty as a reason for faulting his conclusions. While the passage criticizes Deneen for his extreme pessimism about the future of liberalism, his narrow definition of liberalism limited to individual freedoms, and his fixation on the essence of liberalism, it does not specifically address his tendency to hark back to premodern notions. Therefore, Option A stands out as an exception, as it does not align with the explicitly stated reasons for faulting Deneen's conclusions in the passage.

Instructions   

For the following questions answer them individually

\end{problem}

\newpage
\begin{problem}{}{}
% Subject: varc
\textbf{Instructions}

Instructions   

The passage below is accompanied by four questions. Based on the passage, choose the best answer for each question.
The Positivists, anxious to stake out their claim for history as a science, contributed the weight of their influence to the cult of facts. First ascertain the facts, said the positivists, then draw your conclusions from them. . . . This is what may [be] called the common-sense view of history. History consists of a corpus of ascertained facts. The facts are available to the historian in documents, inscriptions, and so on . . . [Sir George Clark] contrasted the "hard core of facts" in history with the surrounding pulp of disputable interpretation forgetting perhaps that the pulpy part of the fruit is more rewarding than the hard core. . . . It recalls the favourite dictum of the great liberal journalist C. P. Scott: "Facts are sacred, opinion is free.". . .
What is a historical fact? . . . According to the common-sense view, there are certain basic facts which are the same for all historians and which form, so to speak, the backbone of history—the fact, for example, that the Battle of Hastings was fought in 1066. But this view calls for two observations. In the first place, it is not with facts like these that the historian is primarily concerned. It is no doubt important to know that the great battle was fought in 1066 and not in 1065 or 1067, and that it was fought at Hastings and not at Eastbourne or Brighton. The historian must not get these things wrong. But [to] praise a historian for his accuracy is like praising an architect for using well-seasoned timber or properly mixed concrete in his building. It is a necessary condition of his work, but not his essential function. It is precisely for matters of this kind that the historian is entitled to rely on what have been called the "auxiliary sciences" of history—archaeology, epigraphy, numismatics, chronology, and so forth. . . .
The second observation is that the necessity to establish these basic facts rests not on any quality in the facts themselves, but on an apriori decision of the historian. In spite of C. P. Scott's motto, every journalist knows today that the most effective way to influence opinion is by the selection and arrangement of the appropriate facts. It used to be said that facts speak for themselves. This is, of course, untrue. The facts speak only when the historian calls on them: it is he who decides to which facts to give the floor, and in what order or context. . . . The only reason why we are interested to know that the battle was fought at Hastings in 1066 is that historians regard it as a major historical event. . . . Professor Talcott Parsons once called [science] "a selective system of cognitive orientations to reality." It might perhaps have been put more simply. But history is, among other things, that. The historian is necessarily selective. The belief in a hard core of historical facts existing objectively and independently of the interpretation of the historian is a preposterous fallacy, but one which it is very hard to eradicate.

\vspace{1em}
\textbf{Question 17}

There is a sentence that is missing in the paragraph below. Look at the paragraph and decide where (option 1, 2, 3, or 4) the following sentence would best fit.
Sentence: Dualism was long held as the defining feature of developing countries in contrast to developed countries, where frontier technologies and high productivity were assumed to prevail.
Paragraph: ___(1)___. At the core of development economics lies the idea of ‘productive dualism’: that poor countries’ economies are split between a narrow ‘modern’ sector that uses advanced technologies and a larger ‘traditional’ sector characterized by very low productivity.___(2)___. While this distinction between developing and advanced economies may have made some sense in the 1950s and 1960s, it no longer appears to be very relevant. A combination of forces have produced a widening gap between the winners and those left behind.___(3)___. Convergence between poor and rich parts of the economy was arrested and regional disparities widened.___(4)___. As a result, policymakers in advanced economies are now grappling with the same questions that have long preoccupied developing economies: mainly how to close the gap with the more advanced parts of the economy.

\begin{enumerate}[label=\Alph*.]
    \item Option 1
    \item Option 2
    \item Option 3
    \item Option 4
[](https://cracku.in/17-there-is-a-sentence-that-is-missing-in-the-paragra-x-cat-2023-slot-2)
\end{enumerate}
\vspace{1em}
\textbf{Solution:}

The sentence best fits in Blank 2 because it directly elaborates on the concept of 'productive dualism' introduced in before Blank 2. It provides additional information about the distinction between the "modern" and "traditional" sectors in poor countries' economies and highlights the historical perspective that contrasts developing and developed countries. This sentence helps set the stage for the subsequent discussion about the relevance of this distinction in the present context.

\end{problem}

\newpage
\begin{problem}{}{}
% Subject: varc
\textbf{Instructions}

Instructions   

The passage below is accompanied by four questions. Based on the passage, choose the best answer for each question.
The Positivists, anxious to stake out their claim for history as a science, contributed the weight of their influence to the cult of facts. First ascertain the facts, said the positivists, then draw your conclusions from them. . . . This is what may [be] called the common-sense view of history. History consists of a corpus of ascertained facts. The facts are available to the historian in documents, inscriptions, and so on . . . [Sir George Clark] contrasted the "hard core of facts" in history with the surrounding pulp of disputable interpretation forgetting perhaps that the pulpy part of the fruit is more rewarding than the hard core. . . . It recalls the favourite dictum of the great liberal journalist C. P. Scott: "Facts are sacred, opinion is free.". . .
What is a historical fact? . . . According to the common-sense view, there are certain basic facts which are the same for all historians and which form, so to speak, the backbone of history—the fact, for example, that the Battle of Hastings was fought in 1066. But this view calls for two observations. In the first place, it is not with facts like these that the historian is primarily concerned. It is no doubt important to know that the great battle was fought in 1066 and not in 1065 or 1067, and that it was fought at Hastings and not at Eastbourne or Brighton. The historian must not get these things wrong. But [to] praise a historian for his accuracy is like praising an architect for using well-seasoned timber or properly mixed concrete in his building. It is a necessary condition of his work, but not his essential function. It is precisely for matters of this kind that the historian is entitled to rely on what have been called the "auxiliary sciences" of history—archaeology, epigraphy, numismatics, chronology, and so forth. . . .
The second observation is that the necessity to establish these basic facts rests not on any quality in the facts themselves, but on an apriori decision of the historian. In spite of C. P. Scott's motto, every journalist knows today that the most effective way to influence opinion is by the selection and arrangement of the appropriate facts. It used to be said that facts speak for themselves. This is, of course, untrue. The facts speak only when the historian calls on them: it is he who decides to which facts to give the floor, and in what order or context. . . . The only reason why we are interested to know that the battle was fought at Hastings in 1066 is that historians regard it as a major historical event. . . . Professor Talcott Parsons once called [science] "a selective system of cognitive orientations to reality." It might perhaps have been put more simply. But history is, among other things, that. The historian is necessarily selective. The belief in a hard core of historical facts existing objectively and independently of the interpretation of the historian is a preposterous fallacy, but one which it is very hard to eradicate.

\vspace{1em}
\textbf{Question 18}

There is a sentence that is missing in the paragraph below. Look at the paragraph and decide where (option 1, 2, 3, or 4) the following sentence would best fit.
Sentence: And probably much earlier, moving the documentation for kissing back 1,000 years compared to what was acknowledged in the scientific community.
Paragraph: Research has hypothesised that the earliest evidence of human lip kissing originated in a very specific geographical location in South Asia 3,500 years ago.___(1)___. From there it may have spread to other regions, simultaneously accelerating the spread of the herpes simplex virus 1. According to Dr Troels Pank Arbøll and Dr Sophie Lund Rasmussen, who in a new article in the journal Science draw on a range of written sources from the earliest Mesopotamian societies, kissing was already a well-established practice 4,500 years ago in the Middle East.___(2)___. In ancient Mesopotamia, people wrote in cuneiform script on clay tablets.___(3)___. Many thousands of these clay tablets have survived to this day, and they contain clear examples that kissing was considered a part of romantic intimacy in ancient times.___(4)___. “Kissing could also have been part of friendships and family members' relations," says Dr Troels Pank Arbøll, an expert on the history of medicine in Mesopotamia.

\begin{enumerate}[label=\Alph*.]
    \item Option 3
    \item Option 4
    \item Option 1
    \item Option 2
[](https://cracku.in/18-there-is-a-sentence-that-is-missing-in-the-paragra-x-cat-2023-slot-2)
\end{enumerate}
\vspace{1em}
\textbf{Solution:}

The sentence best fits in Blank 2. Placing it there provides additional information about the timeline and challenges the previously acknowledged timeline in the scientific community, creating a logical flow in the paragraph.
The sentence does not fit in Blank 1 as it does not add anything to the Sentence preceding Blank 1. Placing the sentence in Blank 3 would disrupt the flow of the passage as there is a direct link between the sentence preceding and the sentence following Blank 3. (_script on clay tablets………Many thousands of these clay tablets”.)_
Similarly placing the Sentence in Blank 4 would disrupt the flow of the passage.

\end{problem}

\newpage
\begin{problem}{}{}
% Subject: varc
\textbf{Instructions}

Instructions   

The passage below is accompanied by four questions. Based on the passage, choose the best answer for each question.
The Positivists, anxious to stake out their claim for history as a science, contributed the weight of their influence to the cult of facts. First ascertain the facts, said the positivists, then draw your conclusions from them. . . . This is what may [be] called the common-sense view of history. History consists of a corpus of ascertained facts. The facts are available to the historian in documents, inscriptions, and so on . . . [Sir George Clark] contrasted the "hard core of facts" in history with the surrounding pulp of disputable interpretation forgetting perhaps that the pulpy part of the fruit is more rewarding than the hard core. . . . It recalls the favourite dictum of the great liberal journalist C. P. Scott: "Facts are sacred, opinion is free.". . .
What is a historical fact? . . . According to the common-sense view, there are certain basic facts which are the same for all historians and which form, so to speak, the backbone of history—the fact, for example, that the Battle of Hastings was fought in 1066. But this view calls for two observations. In the first place, it is not with facts like these that the historian is primarily concerned. It is no doubt important to know that the great battle was fought in 1066 and not in 1065 or 1067, and that it was fought at Hastings and not at Eastbourne or Brighton. The historian must not get these things wrong. But [to] praise a historian for his accuracy is like praising an architect for using well-seasoned timber or properly mixed concrete in his building. It is a necessary condition of his work, but not his essential function. It is precisely for matters of this kind that the historian is entitled to rely on what have been called the "auxiliary sciences" of history—archaeology, epigraphy, numismatics, chronology, and so forth. . . .
The second observation is that the necessity to establish these basic facts rests not on any quality in the facts themselves, but on an apriori decision of the historian. In spite of C. P. Scott's motto, every journalist knows today that the most effective way to influence opinion is by the selection and arrangement of the appropriate facts. It used to be said that facts speak for themselves. This is, of course, untrue. The facts speak only when the historian calls on them: it is he who decides to which facts to give the floor, and in what order or context. . . . The only reason why we are interested to know that the battle was fought at Hastings in 1066 is that historians regard it as a major historical event. . . . Professor Talcott Parsons once called [science] "a selective system of cognitive orientations to reality." It might perhaps have been put more simply. But history is, among other things, that. The historian is necessarily selective. The belief in a hard core of historical facts existing objectively and independently of the interpretation of the historian is a preposterous fallacy, but one which it is very hard to eradicate.

\vspace{1em}
\textbf{Question 19}

Five jumbled up sentences (labelled 1, 2, 3, 4 and 5), related to a topic, are given below. Four of them can be put together to form a coherent paragraph. Identify the odd sentence and key in the number of that sentence as your answer.
1. The banning of Northern Lights could be considered a precursor to censoring books for “moral”, world view or religious reasons.  
2. Attempts to ban books are attempts to silence authors who have summoned immense courage in telling their stories.  
3. Now the banning and challenging of books in the US has escalated to an unprecedented level.  
4. The widely acclaimed fantasy novel Northern Lights was banned in some parts of the US, and was the second most challenged book in the US.  
5. The American Library Association documented an unparalleled number of reported book challenges in 2022, about 2,500 unique titles.
Backspace  
789  
456  
123  
0.-  
Clear All  

Submit
[](https://cracku.in/19-five-jumbled-up-sentences-labelled-1-2-3-4-and-5-r-x-cat-2023-slot-2)

\begin{enumerate}[label=\Alph*.]
\end{enumerate}
\vspace{1em}
\textbf{Solution:}

Sentence 2 is the odd one out because it introduces a broader statement about attempts to ban books and silence authors, while the other sentences are specifically focused on the banning and challenging of books in the US. Sentences 1, 3, 4, and 5 collectively form a coherent paragraph discussing the banning of Northern Lights and the escalation of book challenges in the US, while Sentence 2 introduces a different perspective that doesn't directly contribute to the flow of the paragraph regarding the specific incidents and trends mentioned in the other sentences.

\end{problem}

\newpage
\begin{problem}{}{}
% Subject: varc
\textbf{Instructions}

Instructions   

The passage below is accompanied by four questions. Based on the passage, choose the best answer for each question.
The Positivists, anxious to stake out their claim for history as a science, contributed the weight of their influence to the cult of facts. First ascertain the facts, said the positivists, then draw your conclusions from them. . . . This is what may [be] called the common-sense view of history. History consists of a corpus of ascertained facts. The facts are available to the historian in documents, inscriptions, and so on . . . [Sir George Clark] contrasted the "hard core of facts" in history with the surrounding pulp of disputable interpretation forgetting perhaps that the pulpy part of the fruit is more rewarding than the hard core. . . . It recalls the favourite dictum of the great liberal journalist C. P. Scott: "Facts are sacred, opinion is free.". . .
What is a historical fact? . . . According to the common-sense view, there are certain basic facts which are the same for all historians and which form, so to speak, the backbone of history—the fact, for example, that the Battle of Hastings was fought in 1066. But this view calls for two observations. In the first place, it is not with facts like these that the historian is primarily concerned. It is no doubt important to know that the great battle was fought in 1066 and not in 1065 or 1067, and that it was fought at Hastings and not at Eastbourne or Brighton. The historian must not get these things wrong. But [to] praise a historian for his accuracy is like praising an architect for using well-seasoned timber or properly mixed concrete in his building. It is a necessary condition of his work, but not his essential function. It is precisely for matters of this kind that the historian is entitled to rely on what have been called the "auxiliary sciences" of history—archaeology, epigraphy, numismatics, chronology, and so forth. . . .
The second observation is that the necessity to establish these basic facts rests not on any quality in the facts themselves, but on an apriori decision of the historian. In spite of C. P. Scott's motto, every journalist knows today that the most effective way to influence opinion is by the selection and arrangement of the appropriate facts. It used to be said that facts speak for themselves. This is, of course, untrue. The facts speak only when the historian calls on them: it is he who decides to which facts to give the floor, and in what order or context. . . . The only reason why we are interested to know that the battle was fought at Hastings in 1066 is that historians regard it as a major historical event. . . . Professor Talcott Parsons once called [science] "a selective system of cognitive orientations to reality." It might perhaps have been put more simply. But history is, among other things, that. The historian is necessarily selective. The belief in a hard core of historical facts existing objectively and independently of the interpretation of the historian is a preposterous fallacy, but one which it is very hard to eradicate.

\vspace{1em}
\textbf{Question 20}

**Five jumbled up sentences (labelled 1, 2, 3, 4 and 5), related to a topic, are given below. Four of them can be put together to form a coherent paragraph. Identify the odd sentence and key in the number of that sentence as your answer.**
1. Self-care particularly links to loneliness, behavioural problems, and negative academic outcomes.
2. “Latchkey children” refers to children who routinely return home from school to empty homes and take care of themselves for extended periods of time.
3. Although self-care generally points to negative outcomes, it is important to consider that the bulk of research has yet to track long-term consequences.
4. In research and practice, the phrase “children in self-care” has come to replace "latchkey" in an effort to more accurately reflect the nature of their circumstances.
5. Although parents might believe that self-care would be beneficial for development, recent research has found quite the opposite.
Backspace  
789  
456  
123  
0.-  
Clear All  

Submit
[](https://cracku.in/20-five-jumbled-up-sentences-labelled-1-2-3-4-and-5-r-x-cat-2023-slot-2)

\begin{enumerate}[label=\Alph*.]
\end{enumerate}
\vspace{1em}
\textbf{Solution:}

Sentence 3 is the odd one out because it introduces a different perspective. While the other sentences discuss the negative outcomes associated with self-care in children, Sentence 3 suggests a need to consider that the bulk of research has yet to track long-term consequences. This sentence introduces a more neutral or balanced viewpoint that doesn't align with the general theme of the other sentences, which focus on the negative aspects of self-care for children.
The remaining sentences (1, 2, 4, and 5) can be put together to form a coherent paragraph discussing the concept of self-care in relation to children, particularly those referred to as "Latchkey children."

\end{problem}

\newpage
\begin{problem}{}{}
% Subject: varc
\textbf{Instructions}

Instructions   

The passage below is accompanied by four questions. Based on the passage, choose the best answer for each question.
The Positivists, anxious to stake out their claim for history as a science, contributed the weight of their influence to the cult of facts. First ascertain the facts, said the positivists, then draw your conclusions from them. . . . This is what may [be] called the common-sense view of history. History consists of a corpus of ascertained facts. The facts are available to the historian in documents, inscriptions, and so on . . . [Sir George Clark] contrasted the "hard core of facts" in history with the surrounding pulp of disputable interpretation forgetting perhaps that the pulpy part of the fruit is more rewarding than the hard core. . . . It recalls the favourite dictum of the great liberal journalist C. P. Scott: "Facts are sacred, opinion is free.". . .
What is a historical fact? . . . According to the common-sense view, there are certain basic facts which are the same for all historians and which form, so to speak, the backbone of history—the fact, for example, that the Battle of Hastings was fought in 1066. But this view calls for two observations. In the first place, it is not with facts like these that the historian is primarily concerned. It is no doubt important to know that the great battle was fought in 1066 and not in 1065 or 1067, and that it was fought at Hastings and not at Eastbourne or Brighton. The historian must not get these things wrong. But [to] praise a historian for his accuracy is like praising an architect for using well-seasoned timber or properly mixed concrete in his building. It is a necessary condition of his work, but not his essential function. It is precisely for matters of this kind that the historian is entitled to rely on what have been called the "auxiliary sciences" of history—archaeology, epigraphy, numismatics, chronology, and so forth. . . .
The second observation is that the necessity to establish these basic facts rests not on any quality in the facts themselves, but on an apriori decision of the historian. In spite of C. P. Scott's motto, every journalist knows today that the most effective way to influence opinion is by the selection and arrangement of the appropriate facts. It used to be said that facts speak for themselves. This is, of course, untrue. The facts speak only when the historian calls on them: it is he who decides to which facts to give the floor, and in what order or context. . . . The only reason why we are interested to know that the battle was fought at Hastings in 1066 is that historians regard it as a major historical event. . . . Professor Talcott Parsons once called [science] "a selective system of cognitive orientations to reality." It might perhaps have been put more simply. But history is, among other things, that. The historian is necessarily selective. The belief in a hard core of historical facts existing objectively and independently of the interpretation of the historian is a preposterous fallacy, but one which it is very hard to eradicate.

\vspace{1em}
\textbf{Question 21}

The four sentences (labelled 1, 2, 3 and 4) given below, when properly sequenced, would yield a coherent paragraph. Decide on the proper sequencing of the order of the sentences and key in the sequence of the four numbers as your answer.
1. Like the ants that make up a colony, no single neuron holds complex information like self-awareness, hope or pride.  
2. Although the human brain is not yet understood enough to identify the mechanism by which emergence functions, most neurobiologists agree that complex interconnections among the parts give rise to qualities that belong only to the whole.  
3. Nonetheless, the sum of all neurons in the nervous system generate complex human emotions like fear and joy, none of which can be attributed to a single neuron.  
4. Human consciousness is often called an emergent property of the human brain.
Backspace  
789  
456  
123  
0.-  
Clear All  

Submit
[](https://cracku.in/21-the-four-sentences-labelled-1-2-3-and-4-given-belo-x-cat-2023-slot-2)

\begin{enumerate}[label=\Alph*.]
\end{enumerate}
\vspace{1em}
\textbf{Solution:}

The correct order is 4-1-3-2.
Sentence 4 introduces the main idea that human consciousness is often referred to as an emergent property of the human brain. This gives context for the discussion to be followed.
Now if we consider Sentences 1 and 3 we can see that Sentence 1 ends with _“_** _no single neuron_** _holds complex information like self-awareness, hope or pride.”_ Sentence 3 builds on the idea in Sentence 1 by stating that, nonetheless, the collective activity of all neurons in the nervous system generates complex human emotions like fear and joy “** _Nonetheless_** _, the sum of all neurons …..”_. So we can infer that Sentence 3 must be following Sentence 1. Sentence 3 also supports the notion that emergent properties arise from the interaction of individual components.
Finally, Sentence 2 concludes the paragraph by explaining that although the exact mechanism of emergence in the human brain is not fully understood, neurobiologists agree that complex interconnections among the parts give rise to qualities specific to the whole. This sentence wraps up the discussion and reinforces the concept introduced Sentence 4.
Therefore the correct order is 4-1-3-2.

\end{problem}

\newpage
\begin{problem}{}{}
% Subject: varc
\textbf{Instructions}

Instructions   

The passage below is accompanied by four questions. Based on the passage, choose the best answer for each question.
The Positivists, anxious to stake out their claim for history as a science, contributed the weight of their influence to the cult of facts. First ascertain the facts, said the positivists, then draw your conclusions from them. . . . This is what may [be] called the common-sense view of history. History consists of a corpus of ascertained facts. The facts are available to the historian in documents, inscriptions, and so on . . . [Sir George Clark] contrasted the "hard core of facts" in history with the surrounding pulp of disputable interpretation forgetting perhaps that the pulpy part of the fruit is more rewarding than the hard core. . . . It recalls the favourite dictum of the great liberal journalist C. P. Scott: "Facts are sacred, opinion is free.". . .
What is a historical fact? . . . According to the common-sense view, there are certain basic facts which are the same for all historians and which form, so to speak, the backbone of history—the fact, for example, that the Battle of Hastings was fought in 1066. But this view calls for two observations. In the first place, it is not with facts like these that the historian is primarily concerned. It is no doubt important to know that the great battle was fought in 1066 and not in 1065 or 1067, and that it was fought at Hastings and not at Eastbourne or Brighton. The historian must not get these things wrong. But [to] praise a historian for his accuracy is like praising an architect for using well-seasoned timber or properly mixed concrete in his building. It is a necessary condition of his work, but not his essential function. It is precisely for matters of this kind that the historian is entitled to rely on what have been called the "auxiliary sciences" of history—archaeology, epigraphy, numismatics, chronology, and so forth. . . .
The second observation is that the necessity to establish these basic facts rests not on any quality in the facts themselves, but on an apriori decision of the historian. In spite of C. P. Scott's motto, every journalist knows today that the most effective way to influence opinion is by the selection and arrangement of the appropriate facts. It used to be said that facts speak for themselves. This is, of course, untrue. The facts speak only when the historian calls on them: it is he who decides to which facts to give the floor, and in what order or context. . . . The only reason why we are interested to know that the battle was fought at Hastings in 1066 is that historians regard it as a major historical event. . . . Professor Talcott Parsons once called [science] "a selective system of cognitive orientations to reality." It might perhaps have been put more simply. But history is, among other things, that. The historian is necessarily selective. The belief in a hard core of historical facts existing objectively and independently of the interpretation of the historian is a preposterous fallacy, but one which it is very hard to eradicate.

\vspace{1em}
\textbf{Question 22}

The four sentences (labelled 1, 2, 3 and 4) given below, when properly sequenced, would yield a coherent paragraph. Decide on the proper sequencing of the order of the sentences and key in the sequence of the four numbers as your answer.
1. Contemporary African writing like ‘The Bottled Leopard’ voices this theme using two children and two backgrounds to juxtapose two varying cultures.  
2. Chukwuemeka Ike explores the conflict, and casts the Western tradition as condescending, enveloping and unaccommodating towards local African practice.  
3. However, their views contradict the reality, for a rich and sustaining local African cultural ethos exists for all who care, to see and experience.  
4. Western Christian concepts tend to deny or feign ignorance about the existence of a genuine and enduring indigenous African tradition.
Backspace  
789  
456  
123  
0.-  
Clear All  

Submit
[](https://cracku.in/22-the-four-sentences-labelled-1-2-3-and-4-given-belo-x-cat-2023-slot-2)

\begin{enumerate}[label=\Alph*.]
\end{enumerate}
\vspace{1em}
\textbf{Solution:}

The correct order is 4-3-1-2.
Sentence 4 introduces the conflict between Western Christian concepts and indigenous African tradition.This gives context for the discussion to be followed.
Now we can see that Sentence 3 starts with _“However….”_ which implies that this sentence is challenging the views presented in Sentence 4 by asserting the existence of a rich indigenous African cultural ethos. Therefore Sentence 3 must be following Sentence 4.
Sentence 2 introduces Chuwkuemeka Ike and provides additional information about Chukwuemeka Ike's exploration of the conflict, adding depth to the discussion. Sentence 1 introduces an example of contemporary African writing that explores the conflict. It makes sense that first we give an example and then elaborate on it. Sentence 1 introduces the book; it is a logical inferences that Ike must be the author of this book.
Therefore Sentence 2 must be following Sentence 1.  
Therefore the correct order is 4-3-1-2.

\end{problem}

\newpage
\begin{problem}{}{}
% Subject: varc
\textbf{Instructions}

Instructions   

The passage below is accompanied by four questions. Based on the passage, choose the best answer for each question.
The Positivists, anxious to stake out their claim for history as a science, contributed the weight of their influence to the cult of facts. First ascertain the facts, said the positivists, then draw your conclusions from them. . . . This is what may [be] called the common-sense view of history. History consists of a corpus of ascertained facts. The facts are available to the historian in documents, inscriptions, and so on . . . [Sir George Clark] contrasted the "hard core of facts" in history with the surrounding pulp of disputable interpretation forgetting perhaps that the pulpy part of the fruit is more rewarding than the hard core. . . . It recalls the favourite dictum of the great liberal journalist C. P. Scott: "Facts are sacred, opinion is free.". . .
What is a historical fact? . . . According to the common-sense view, there are certain basic facts which are the same for all historians and which form, so to speak, the backbone of history—the fact, for example, that the Battle of Hastings was fought in 1066. But this view calls for two observations. In the first place, it is not with facts like these that the historian is primarily concerned. It is no doubt important to know that the great battle was fought in 1066 and not in 1065 or 1067, and that it was fought at Hastings and not at Eastbourne or Brighton. The historian must not get these things wrong. But [to] praise a historian for his accuracy is like praising an architect for using well-seasoned timber or properly mixed concrete in his building. It is a necessary condition of his work, but not his essential function. It is precisely for matters of this kind that the historian is entitled to rely on what have been called the "auxiliary sciences" of history—archaeology, epigraphy, numismatics, chronology, and so forth. . . .
The second observation is that the necessity to establish these basic facts rests not on any quality in the facts themselves, but on an apriori decision of the historian. In spite of C. P. Scott's motto, every journalist knows today that the most effective way to influence opinion is by the selection and arrangement of the appropriate facts. It used to be said that facts speak for themselves. This is, of course, untrue. The facts speak only when the historian calls on them: it is he who decides to which facts to give the floor, and in what order or context. . . . The only reason why we are interested to know that the battle was fought at Hastings in 1066 is that historians regard it as a major historical event. . . . Professor Talcott Parsons once called [science] "a selective system of cognitive orientations to reality." It might perhaps have been put more simply. But history is, among other things, that. The historian is necessarily selective. The belief in a hard core of historical facts existing objectively and independently of the interpretation of the historian is a preposterous fallacy, but one which it is very hard to eradicate.

\vspace{1em}
\textbf{Question 23}

The passage given below is followed by four alternate summaries. Choose the option that best captures the essence of the passage.
People spontaneously create counterfactual alternatives to reality when they think “if only” or “what if” and imagine how the past could have been different. The mind computes counterfactuals for many reasons. Counterfactuals explain the past and prepare for the future, they implicate various relations including causal ones, and they affect intentions and decisions. They modulate emotions such as regret and relief, and they support moral judgments such as blame. The ability to create counterfactuals develops throughout childhood and contributes to reasoning about other people's beliefs, including their false beliefs.

\begin{enumerate}[label=\Alph*.]
    \item Counterfactuals help people to prepare for the future by understanding intentions and making decisions.
    \item People create counterfactual alternatives to reality for various reasons, including reasoning about other people's beliefs.
    \item Counterfactual alternatives to reality are created for a variety of reasons and is part of one's developmental process.
    \item Counterfactual thinking helps to reverse past and future actions and reason out false beliefs.
[](https://cracku.in/23-the-passage-given-below-is-followed-by-four-altern-x-cat-2023-slot-2)
\end{enumerate}
\vspace{1em}
\textbf{Solution:}

The passage discusses the phenomenon of counterfactual thinking, highlighting that people spontaneously create counterfactual alternatives to reality for various reasons. These reasons include explaining the past, preparing for the future, implicating various relations (including causal ones), affecting emotions, and supporting moral judgments. Additionally, the passage mentions that the ability to create counterfactuals develops throughout childhood and contributes to reasoning about other people's beliefs. Option C effectively encompasses these key points, making it the most accurate summary of the passage.
Option A focuses primarily on the preparation for the future aspect, neglecting the broader reasons for creating counterfactual alternatives.
Option B does not emphasize the developmental aspect and various reasons for creating 
Option D inaccurately suggests that counterfactual thinking helps reverse past and future actions, which is not the main point of the passage, and it oversimplifies the role of counterfactuals.

\end{problem}

\newpage
\begin{problem}{}{}
% Subject: varc
\textbf{Instructions}

Instructions   

The passage below is accompanied by four questions. Based on the passage, choose the best answer for each question.
The Positivists, anxious to stake out their claim for history as a science, contributed the weight of their influence to the cult of facts. First ascertain the facts, said the positivists, then draw your conclusions from them. . . . This is what may [be] called the common-sense view of history. History consists of a corpus of ascertained facts. The facts are available to the historian in documents, inscriptions, and so on . . . [Sir George Clark] contrasted the "hard core of facts" in history with the surrounding pulp of disputable interpretation forgetting perhaps that the pulpy part of the fruit is more rewarding than the hard core. . . . It recalls the favourite dictum of the great liberal journalist C. P. Scott: "Facts are sacred, opinion is free.". . .
What is a historical fact? . . . According to the common-sense view, there are certain basic facts which are the same for all historians and which form, so to speak, the backbone of history—the fact, for example, that the Battle of Hastings was fought in 1066. But this view calls for two observations. In the first place, it is not with facts like these that the historian is primarily concerned. It is no doubt important to know that the great battle was fought in 1066 and not in 1065 or 1067, and that it was fought at Hastings and not at Eastbourne or Brighton. The historian must not get these things wrong. But [to] praise a historian for his accuracy is like praising an architect for using well-seasoned timber or properly mixed concrete in his building. It is a necessary condition of his work, but not his essential function. It is precisely for matters of this kind that the historian is entitled to rely on what have been called the "auxiliary sciences" of history—archaeology, epigraphy, numismatics, chronology, and so forth. . . .
The second observation is that the necessity to establish these basic facts rests not on any quality in the facts themselves, but on an apriori decision of the historian. In spite of C. P. Scott's motto, every journalist knows today that the most effective way to influence opinion is by the selection and arrangement of the appropriate facts. It used to be said that facts speak for themselves. This is, of course, untrue. The facts speak only when the historian calls on them: it is he who decides to which facts to give the floor, and in what order or context. . . . The only reason why we are interested to know that the battle was fought at Hastings in 1066 is that historians regard it as a major historical event. . . . Professor Talcott Parsons once called [science] "a selective system of cognitive orientations to reality." It might perhaps have been put more simply. But history is, among other things, that. The historian is necessarily selective. The belief in a hard core of historical facts existing objectively and independently of the interpretation of the historian is a preposterous fallacy, but one which it is very hard to eradicate.

\vspace{1em}
\textbf{Question 24}

The passage given below is followed by four alternate summaries. Choose the option that best captures the essence of the passage.
Heatwaves are becoming longer, frequent and intense due to climate change. The impacts of extreme heat are unevenly experienced; with older people and young children, those with pre-existing medical conditions and on low incomes significantly more vulnerable. Adaptation to heatwaves is a significant public policy concern. Research conducted among at-risk people in the UK reveals that even vulnerable people do not perceive themselves as at risk of extreme heat; therefore, early warnings of extreme heat events do not perform as intended. This suggests that understanding how extreme heat is narrated is very important. The news media play a central role in this process and can help warn people about the potential danger, as well as about impacts on infrastructure and society.

\begin{enumerate}[label=\Alph*.]
    \item Protection from heat waves is important but current reports and public policies seem ineffective.
    \item People are vulnerable to heatwaves caused due to climate change, measures taken are ineffective.
    \item Heatwaves pose an enormous risk; the media plays a pivotal role in alerting people to this danger.
    \item News stories help in warning about heatwaves, but they have to become more effective.
[](https://cracku.in/24-the-passage-given-below-is-followed-by-four-altern-x-cat-2023-slot-2)
\end{enumerate}
\vspace{1em}
\textbf{Solution:}

The passage discusses the increasing frequency and intensity of heatwaves due to climate change, with vulnerable groups experiencing uneven impacts. It emphasizes that adaptation to heatwaves is a significant public policy concern. The research findings suggest that even vulnerable individuals may not perceive themselves as at risk of extreme heat, highlighting the importance of understanding how extreme heat is narrated. The passage specifically mentions the central role of the news media in warning people about the potential danger of heatwaves and their impacts on infrastructure and society. Option C effectively conveys the primary focus on heatwaves posing a substantial risk and the critical role of the media in alerting the public to this danger.
Option A implies a general importance of protection without specifically highlighting the role of the media in alerting people to the risks of heatwaves.
Option B acknowledges the vulnerability to heatwaves but it does not emphasize the role of the media in alerting people and suggests a broader critique of measures taken.
Option D mentions the need for news stories to become more effective but does not emphasize the central role of the media in alerting people to the risk of heatwaves, as the passage does.
[](https://cracku.in/downloads/18893)
  *  
  *

\end{problem}

\newpage
\begin{problem}{}{}
% Subject: varc
\textbf{Instructions}

Instructions   

The passage below is accompanied by four questions. Based on the passage, choose the best answer for each question.
In 2006, the Met [art museum in the US] agreed to return the Euphronios krater, a masterpiece Greek urn that had been a museum draw since 1972. In 2007, the Getty [art museum in the US] agreed to return 40 objects to Italy, including a marble Aphrodite, in the midst of looting scandals. And in December, Sotheby’s and a private owner agreed to return an ancient Khmer statue of a warrior, pulled from auction two years before, to Cambodia.
Cultural property, or patrimony, laws limit the transfer of cultural property outside the source country’s territory, including outright export prohibitions and national ownership laws. Most art historians, archaeologists, museum officials and policymakers portray cultural property laws in general as invaluable tools for counteracting the ugly legacy of Western cultural imperialism.
During the late 19th and early 20th century — an era former Met director Thomas Hoving called “the age of piracy” — American and European art museums acquired antiquities by hook or by crook, from grave robbers or souvenir collectors, bounty from digs and ancient sites in impoverished but art-rich source countries. Patrimony laws were intended to protect future archaeological discoveries against Western imperialist designs. . . .
I surveyed 90 countries with one or more archaeological sites on UNESCO’s World Heritage Site list, and my study shows that in most cases the number of discovered sites diminishes sharply after a country passes a cultural property law. There are 222 archaeological sites listed for those 90 countries. When you look into the history of the sites, you see that all but 21 were discovered before the passage of cultural property laws. . . .
Strict cultural patrimony laws are popular in most countries. But the downside may be that they reduce incentives for foreign governments, nongovernmental organizations and educational institutions to invest in overseas exploration because their efforts will not necessarily be rewarded by opportunities to hold, display and study what is uncovered. To the extent that source countries can fund their own archaeological projects, artifacts and sites may still be discovered. . . . The survey has far-reaching implications. It suggests that source countries, particularly in the developing world, should narrow their cultural property laws so that they can reap the benefits of new archaeological discoveries, which typically increase tourism and enhance cultural pride. This does not mean these nations should abolish restrictions on foreign excavation and foreign claims to artifacts.
China provides an interesting alternative approach for source nations eager for foreign archaeological investment. From 1935 to 2003, China had a restrictive cultural property law that prohibited foreign ownership of Chinese cultural artifacts. In those years, China’s most significant archaeological discovery occurred by chance, in 1974, when peasant farmers accidentally uncovered ranks of buried terra cotta warriors, which are part of Emperor Qin’s spectacular tomb system.
In 2003, the Chinese government switched course, dropping its cultural property law and embracing collaborative international archaeological research. Since then, China has nominated 11 archaeological sites for inclusion in the World Heritage Site list, including eight in 2013, the most ever for China.

\vspace{1em}
\textbf{Question 1}

Which one of the following statements best expresses the paradox of patrimony laws?

\begin{enumerate}[label=\Alph*.]
    \item They were aimed at protecting cultural property, but instead reduced business for auctioneers like Sotheby’s.
    \item They were intended to protect cultural property, but instead resulted in the neglect of historical sites.
    \item They were intended to protect cultural property, but instead resulted in the withholding of national treasure from museums.
    \item They were aimed at protecting cultural property, but instead reduced new archaeological discoveries.
[](https://cracku.in/1-which-one-of-the-following-statements-best-express-x-cat-2023-slot-3)
\end{enumerate}
\vspace{1em}
\textbf{Solution:}

The primary purpose of patrimony laws is stated in the passage as being "aimed at protecting cultural property," implying that the intention behind these laws is to preserve and safeguard a country's cultural heritage. However, the paradox lies in the unintended consequence of these laws, as highlighted in the passage. The author argues that, despite the good intentions of protecting cultural property, the strict implementation of patrimony laws has led to a reduction in new archaeological discoveries. This reduction is attributed to diminished incentives for foreign entities, such as governments, NGOs, and educational institutions, to invest in overseas archaeological exploration. In other words, the very laws designed to protect cultural property end up hindering the process of making new archaeological discoveries. This underscores the tension between preserving cultural heritage and the potential negative impact on the exploration and understanding of that heritage. Option D aptly captures this point.

\end{problem}

\newpage
\begin{problem}{}{}
% Subject: varc
\textbf{Instructions}

Instructions   

The passage below is accompanied by four questions. Based on the passage, choose the best answer for each question.
In 2006, the Met [art museum in the US] agreed to return the Euphronios krater, a masterpiece Greek urn that had been a museum draw since 1972. In 2007, the Getty [art museum in the US] agreed to return 40 objects to Italy, including a marble Aphrodite, in the midst of looting scandals. And in December, Sotheby’s and a private owner agreed to return an ancient Khmer statue of a warrior, pulled from auction two years before, to Cambodia.
Cultural property, or patrimony, laws limit the transfer of cultural property outside the source country’s territory, including outright export prohibitions and national ownership laws. Most art historians, archaeologists, museum officials and policymakers portray cultural property laws in general as invaluable tools for counteracting the ugly legacy of Western cultural imperialism.
During the late 19th and early 20th century — an era former Met director Thomas Hoving called “the age of piracy” — American and European art museums acquired antiquities by hook or by crook, from grave robbers or souvenir collectors, bounty from digs and ancient sites in impoverished but art-rich source countries. Patrimony laws were intended to protect future archaeological discoveries against Western imperialist designs. . . .
I surveyed 90 countries with one or more archaeological sites on UNESCO’s World Heritage Site list, and my study shows that in most cases the number of discovered sites diminishes sharply after a country passes a cultural property law. There are 222 archaeological sites listed for those 90 countries. When you look into the history of the sites, you see that all but 21 were discovered before the passage of cultural property laws. . . .
Strict cultural patrimony laws are popular in most countries. But the downside may be that they reduce incentives for foreign governments, nongovernmental organizations and educational institutions to invest in overseas exploration because their efforts will not necessarily be rewarded by opportunities to hold, display and study what is uncovered. To the extent that source countries can fund their own archaeological projects, artifacts and sites may still be discovered. . . . The survey has far-reaching implications. It suggests that source countries, particularly in the developing world, should narrow their cultural property laws so that they can reap the benefits of new archaeological discoveries, which typically increase tourism and enhance cultural pride. This does not mean these nations should abolish restrictions on foreign excavation and foreign claims to artifacts.
China provides an interesting alternative approach for source nations eager for foreign archaeological investment. From 1935 to 2003, China had a restrictive cultural property law that prohibited foreign ownership of Chinese cultural artifacts. In those years, China’s most significant archaeological discovery occurred by chance, in 1974, when peasant farmers accidentally uncovered ranks of buried terra cotta warriors, which are part of Emperor Qin’s spectacular tomb system.
In 2003, the Chinese government switched course, dropping its cultural property law and embracing collaborative international archaeological research. Since then, China has nominated 11 archaeological sites for inclusion in the World Heritage Site list, including eight in 2013, the most ever for China.

\vspace{1em}
\textbf{Question 2}

It can be inferred from the passage that archaeological sites are considered important by some source countries because they:

\begin{enumerate}[label=\Alph*.]
    \item are a symbol of Western imperialism.
    \item are subject to strict patrimony laws.
    \item generate funds for future discoveries.
    \item give a boost to the tourism sector.
[](https://cracku.in/2-it-can-be-inferred-from-the-passage-that-archaeolo-x-cat-2023-slot-3)
\end{enumerate}
\vspace{1em}
\textbf{Solution:}

The author suggests that archaeological sites are important to some source countries because they can reap benefits from new archaeological discoveries, and one of the mentioned benefits is that such discoveries typically increase tourism. The passage emphasizes the economic and cultural advantages associated with tourism, which includes enhancing cultural pride and potentially attracting visitors to explore archaeological sites. Option D correctly presents this point. None of the other choices can be considered as valid inferences.

\end{problem}

\newpage
\begin{problem}{}{}
% Subject: varc
\textbf{Instructions}

Instructions   

The passage below is accompanied by four questions. Based on the passage, choose the best answer for each question.
In 2006, the Met [art museum in the US] agreed to return the Euphronios krater, a masterpiece Greek urn that had been a museum draw since 1972. In 2007, the Getty [art museum in the US] agreed to return 40 objects to Italy, including a marble Aphrodite, in the midst of looting scandals. And in December, Sotheby’s and a private owner agreed to return an ancient Khmer statue of a warrior, pulled from auction two years before, to Cambodia.
Cultural property, or patrimony, laws limit the transfer of cultural property outside the source country’s territory, including outright export prohibitions and national ownership laws. Most art historians, archaeologists, museum officials and policymakers portray cultural property laws in general as invaluable tools for counteracting the ugly legacy of Western cultural imperialism.
During the late 19th and early 20th century — an era former Met director Thomas Hoving called “the age of piracy” — American and European art museums acquired antiquities by hook or by crook, from grave robbers or souvenir collectors, bounty from digs and ancient sites in impoverished but art-rich source countries. Patrimony laws were intended to protect future archaeological discoveries against Western imperialist designs. . . .
I surveyed 90 countries with one or more archaeological sites on UNESCO’s World Heritage Site list, and my study shows that in most cases the number of discovered sites diminishes sharply after a country passes a cultural property law. There are 222 archaeological sites listed for those 90 countries. When you look into the history of the sites, you see that all but 21 were discovered before the passage of cultural property laws. . . .
Strict cultural patrimony laws are popular in most countries. But the downside may be that they reduce incentives for foreign governments, nongovernmental organizations and educational institutions to invest in overseas exploration because their efforts will not necessarily be rewarded by opportunities to hold, display and study what is uncovered. To the extent that source countries can fund their own archaeological projects, artifacts and sites may still be discovered. . . . The survey has far-reaching implications. It suggests that source countries, particularly in the developing world, should narrow their cultural property laws so that they can reap the benefits of new archaeological discoveries, which typically increase tourism and enhance cultural pride. This does not mean these nations should abolish restrictions on foreign excavation and foreign claims to artifacts.
China provides an interesting alternative approach for source nations eager for foreign archaeological investment. From 1935 to 2003, China had a restrictive cultural property law that prohibited foreign ownership of Chinese cultural artifacts. In those years, China’s most significant archaeological discovery occurred by chance, in 1974, when peasant farmers accidentally uncovered ranks of buried terra cotta warriors, which are part of Emperor Qin’s spectacular tomb system.
In 2003, the Chinese government switched course, dropping its cultural property law and embracing collaborative international archaeological research. Since then, China has nominated 11 archaeological sites for inclusion in the World Heritage Site list, including eight in 2013, the most ever for China.

\vspace{1em}
\textbf{Question 3}

Which one of the following statements, if true, would undermine the central idea of the passage?

\begin{enumerate}[label=\Alph*.]
    \item Affluent archaeologically-rich source countries can afford to carry out their own excavations.
    \item Museums established in economically deprived archaeologically-rich source countries can display the antiques discovered there.
    \item UNESCO finances archaeological research in poor, but archaeologically-rich source countries.
    \item Western countries will have to apologise to countries for looting their cultural property in the past century.
[](https://cracku.in/3-which-one-of-the-following-statements-if-true-woul-x-cat-2023-slot-3)
\end{enumerate}
\vspace{1em}
\textbf{Solution:}

The central idea of the passage is that strict cultural property laws, although popular, may reduce incentives for foreign entities to invest in overseas archaeological exploration. The passage suggests that this reduction in incentives could be detrimental to archaeological discoveries and, consequently, to the tourism and cultural pride of source countries.
Among the given options, the only statement that would undermine the central idea is presented in Option C - it introduces the idea that there is external financial support for archaeological research in these countries. If true, then the lack of discoveries could potentially be attributed to a completely different factor/variable that the author might have failed to account for.

\end{problem}

\newpage
\begin{problem}{}{}
% Subject: varc
\textbf{Instructions}

Instructions   

The passage below is accompanied by four questions. Based on the passage, choose the best answer for each question.
In 2006, the Met [art museum in the US] agreed to return the Euphronios krater, a masterpiece Greek urn that had been a museum draw since 1972. In 2007, the Getty [art museum in the US] agreed to return 40 objects to Italy, including a marble Aphrodite, in the midst of looting scandals. And in December, Sotheby’s and a private owner agreed to return an ancient Khmer statue of a warrior, pulled from auction two years before, to Cambodia.
Cultural property, or patrimony, laws limit the transfer of cultural property outside the source country’s territory, including outright export prohibitions and national ownership laws. Most art historians, archaeologists, museum officials and policymakers portray cultural property laws in general as invaluable tools for counteracting the ugly legacy of Western cultural imperialism.
During the late 19th and early 20th century — an era former Met director Thomas Hoving called “the age of piracy” — American and European art museums acquired antiquities by hook or by crook, from grave robbers or souvenir collectors, bounty from digs and ancient sites in impoverished but art-rich source countries. Patrimony laws were intended to protect future archaeological discoveries against Western imperialist designs. . . .
I surveyed 90 countries with one or more archaeological sites on UNESCO’s World Heritage Site list, and my study shows that in most cases the number of discovered sites diminishes sharply after a country passes a cultural property law. There are 222 archaeological sites listed for those 90 countries. When you look into the history of the sites, you see that all but 21 were discovered before the passage of cultural property laws. . . .
Strict cultural patrimony laws are popular in most countries. But the downside may be that they reduce incentives for foreign governments, nongovernmental organizations and educational institutions to invest in overseas exploration because their efforts will not necessarily be rewarded by opportunities to hold, display and study what is uncovered. To the extent that source countries can fund their own archaeological projects, artifacts and sites may still be discovered. . . . The survey has far-reaching implications. It suggests that source countries, particularly in the developing world, should narrow their cultural property laws so that they can reap the benefits of new archaeological discoveries, which typically increase tourism and enhance cultural pride. This does not mean these nations should abolish restrictions on foreign excavation and foreign claims to artifacts.
China provides an interesting alternative approach for source nations eager for foreign archaeological investment. From 1935 to 2003, China had a restrictive cultural property law that prohibited foreign ownership of Chinese cultural artifacts. In those years, China’s most significant archaeological discovery occurred by chance, in 1974, when peasant farmers accidentally uncovered ranks of buried terra cotta warriors, which are part of Emperor Qin’s spectacular tomb system.
In 2003, the Chinese government switched course, dropping its cultural property law and embracing collaborative international archaeological research. Since then, China has nominated 11 archaeological sites for inclusion in the World Heritage Site list, including eight in 2013, the most ever for China.

\vspace{1em}
\textbf{Question 4}

From the passage we can infer that the author is likely to advise poor, but archaeologically-rich source countries to do all of the following, EXCEPT:

\begin{enumerate}[label=\Alph*.]
    \item fund institutes in other countries to undertake archaeological exploration in the source country reaping the benefits of cutting-edge techniques.
    \item allow foreign countries to analyse and exhibit the archaeological finds made in the source country.
    \item to find ways to motivate other countries to finance archaeological explorations in their country.
    \item adopt China’s strategy of dropping its cultural property laws and carrying out archaeological research through international collaboration.
[](https://cracku.in/4-from-the-passage-we-can-infer-that-the-author-is-l-x-cat-2023-slot-3)
\end{enumerate}
\vspace{1em}
\textbf{Solution:}

Let us examine the given choices - 
Option A: The passage emphasises the benefits of international collaboration and suggests that source countries could reap the benefits of new archaeological discoveries through such collaboration. Funding institutes in other countries to undertake exploration may not align with this perspective, as it involves outsourcing the exploration rather than fostering collaboration within the source country. Thus, Option A is an unlikely recommendation. 
Option B: This aligns with the idea of international collaboration and suggests that the author might advise source countries to permit foreign analysis and exhibition of archaeological finds. The passage does suggest that strict cultural property laws may hinder opportunities to hold, display, and study uncovered artefacts.
Option C: We are told that strict cultural property laws might reduce incentives for foreign governments, NGOs, and educational institutions to invest in overseas exploration. Therefore, the author would support a proactive approach to encourage these entities to invest in expeditions in source countries. 
Option D: The passage highlights China's shift in strategy from strict cultural property laws to embracing collaborative international archaeological research. The author suggests that this approach has led to an increase in archaeological discoveries and recognition. Hence, this would be a relevant recommendation.
Hence, Option A is the correct choice.

Instructions   

The passage below is accompanied by four questions. Based on the passage, choose the best answer for each question.
Steven Pinker’s new book, “Rationality: What It Is, Why It Seems Scarce, Why It Matters,” offers a pragmatic dose of measured optimism, presenting rationality as a fragile but achievable ideal in personal and civic life. . . . Pinker’s ambition to illuminate such a crucial topic offers the welcome prospect of a return to sanity. . . . It’s no small achievement to make formal logic, game theory, statistics and Bayesian reasoning delightful topics full of charm and relevance.
It’s also plausible to believe that a wider application of the rational tools he analyzes would improve the world in important ways. His primer on statistics and scientific uncertainty is particularly timely and should be required reading before consuming any news about the [COVID] pandemic. More broadly, he argues that less media coverage of shocking but vanishingly rare events, from shark attacks to adverse vaccine reactions, would help prevent dangerous overreactions, fatalism and the diversion of finite resources away from solvable but less-dramatic issues, like malnutrition in the developing world.
It’s a reasonable critique, and Pinker is not the first to make it. But analyzing the political economy of journalism — its funding structures, ownership concentration and increasing reliance on social media shares — would have given a fuller picture of why so much coverage is so misguided and what we might do about it.
Pinker’s main focus is the sort of conscious, sequential reasoning that can track the steps in a geometric proof or an argument in formal logic. Skill in this domain maps directly onto the navigation of many real-world problems, and Pinker shows how greater mastery of the tools of rationality can improve decision-making in medical, legal, financial and many other contexts in which we must act on uncertain and shifting information. . . .
Despite the undeniable power of the sort of rationality he describes, many of the deepest insights in the history of science, math, music and art strike their originators in moments of epiphany. From the 19th-century chemist Friedrich August Kekulé’s discovery of the structure of benzene to any of Mozart’s symphonies, much extraordinary human achievement is not a product of conscious, sequential reasoning. Even Plato’s Socrates — who anticipated many of Pinker’s points by nearly 2,500 years, showing the virtue of knowing what you do not know and examining all premises in arguments, not simply trusting speakers’ authority or charisma — attributed many of his most profound insights to dreams and visions. Conscious reasoning is helpful in sorting the wheat from the chaff, but it would be interesting to consider the hidden aquifers that make much of the grain grow in the first place.
The role of moral and ethical education in promoting rational behavior is also underexplored. Pinker recognizes that rationality “is not just a cognitive virtue but a moral one.” But this profoundly important point, one subtly explored by ancient Greek philosophers like Plato and Aristotle, doesn’t really get developed. This is a shame, since possessing the right sort of moral character is arguably a precondition for using rationality in beneficial ways.

\end{problem}

\newpage
\begin{problem}{}{}
% Subject: varc
\textbf{Instructions}

Instructions   

The passage below is accompanied by four questions. Based on the passage, choose the best answer for each question.
In 2006, the Met [art museum in the US] agreed to return the Euphronios krater, a masterpiece Greek urn that had been a museum draw since 1972. In 2007, the Getty [art museum in the US] agreed to return 40 objects to Italy, including a marble Aphrodite, in the midst of looting scandals. And in December, Sotheby’s and a private owner agreed to return an ancient Khmer statue of a warrior, pulled from auction two years before, to Cambodia.
Cultural property, or patrimony, laws limit the transfer of cultural property outside the source country’s territory, including outright export prohibitions and national ownership laws. Most art historians, archaeologists, museum officials and policymakers portray cultural property laws in general as invaluable tools for counteracting the ugly legacy of Western cultural imperialism.
During the late 19th and early 20th century — an era former Met director Thomas Hoving called “the age of piracy” — American and European art museums acquired antiquities by hook or by crook, from grave robbers or souvenir collectors, bounty from digs and ancient sites in impoverished but art-rich source countries. Patrimony laws were intended to protect future archaeological discoveries against Western imperialist designs. . . .
I surveyed 90 countries with one or more archaeological sites on UNESCO’s World Heritage Site list, and my study shows that in most cases the number of discovered sites diminishes sharply after a country passes a cultural property law. There are 222 archaeological sites listed for those 90 countries. When you look into the history of the sites, you see that all but 21 were discovered before the passage of cultural property laws. . . .
Strict cultural patrimony laws are popular in most countries. But the downside may be that they reduce incentives for foreign governments, nongovernmental organizations and educational institutions to invest in overseas exploration because their efforts will not necessarily be rewarded by opportunities to hold, display and study what is uncovered. To the extent that source countries can fund their own archaeological projects, artifacts and sites may still be discovered. . . . The survey has far-reaching implications. It suggests that source countries, particularly in the developing world, should narrow their cultural property laws so that they can reap the benefits of new archaeological discoveries, which typically increase tourism and enhance cultural pride. This does not mean these nations should abolish restrictions on foreign excavation and foreign claims to artifacts.
China provides an interesting alternative approach for source nations eager for foreign archaeological investment. From 1935 to 2003, China had a restrictive cultural property law that prohibited foreign ownership of Chinese cultural artifacts. In those years, China’s most significant archaeological discovery occurred by chance, in 1974, when peasant farmers accidentally uncovered ranks of buried terra cotta warriors, which are part of Emperor Qin’s spectacular tomb system.
In 2003, the Chinese government switched course, dropping its cultural property law and embracing collaborative international archaeological research. Since then, China has nominated 11 archaeological sites for inclusion in the World Heritage Site list, including eight in 2013, the most ever for China.

\vspace{1em}
\textbf{Question 5}

According to the author, for Pinker as well as the ancient Greek philosophers, rational thinking involves all of the following EXCEPT:

\begin{enumerate}[label=\Alph*.]
    \item an awareness of underlying assumptions in an argument and gaps in one’s own knowledge
    \item the belief that the ability to reason logically encompasses an ethical and moral dimension.
    \item the primacy of conscious sequential reasoning as the basis for seminal human achievements.
    \item arriving at independent conclusions irrespective of who is presenting the argument.
[](https://cracku.in/5-according-to-the-author-for-pinker-as-well-as-the--x-cat-2023-slot-3)
\end{enumerate}
\vspace{1em}
\textbf{Solution:}

Based on the discussion, the option that is NOT associated with Pinker's view of rational thinking (as well as that of the ancient Greek philosophers) is Option C - the passage suggests that while sequential reasoning is valuable, many profound human achievements come from moments of epiphany or insight rather than solely from conscious, sequential reasoning.
In relation to this thought, we are told that the emphasis on rational thought involves an understanding of the gaps in one’s own knowledge [Option A] and also ‘arriving at independent conclusions’ [Option D]: {“ _Even Plato’s Socrates — who anticipated many of Pinker’s points by nearly 2,500 years, showing the virtue of knowing what you do not know and examining all premises in arguments, not simply trusting speakers’ authority or charisma..._ ”}
Towards the end of the passage, we are informed of an ethical and moral dimension [Option B] to rationality, which the author asserts that Pinker considers but does not elaborate on.
Hence, Option C is the correct choice.

\end{problem}

\newpage
\begin{problem}{}{}
% Subject: varc
\textbf{Instructions}

Instructions   

The passage below is accompanied by four questions. Based on the passage, choose the best answer for each question.
In 2006, the Met [art museum in the US] agreed to return the Euphronios krater, a masterpiece Greek urn that had been a museum draw since 1972. In 2007, the Getty [art museum in the US] agreed to return 40 objects to Italy, including a marble Aphrodite, in the midst of looting scandals. And in December, Sotheby’s and a private owner agreed to return an ancient Khmer statue of a warrior, pulled from auction two years before, to Cambodia.
Cultural property, or patrimony, laws limit the transfer of cultural property outside the source country’s territory, including outright export prohibitions and national ownership laws. Most art historians, archaeologists, museum officials and policymakers portray cultural property laws in general as invaluable tools for counteracting the ugly legacy of Western cultural imperialism.
During the late 19th and early 20th century — an era former Met director Thomas Hoving called “the age of piracy” — American and European art museums acquired antiquities by hook or by crook, from grave robbers or souvenir collectors, bounty from digs and ancient sites in impoverished but art-rich source countries. Patrimony laws were intended to protect future archaeological discoveries against Western imperialist designs. . . .
I surveyed 90 countries with one or more archaeological sites on UNESCO’s World Heritage Site list, and my study shows that in most cases the number of discovered sites diminishes sharply after a country passes a cultural property law. There are 222 archaeological sites listed for those 90 countries. When you look into the history of the sites, you see that all but 21 were discovered before the passage of cultural property laws. . . .
Strict cultural patrimony laws are popular in most countries. But the downside may be that they reduce incentives for foreign governments, nongovernmental organizations and educational institutions to invest in overseas exploration because their efforts will not necessarily be rewarded by opportunities to hold, display and study what is uncovered. To the extent that source countries can fund their own archaeological projects, artifacts and sites may still be discovered. . . . The survey has far-reaching implications. It suggests that source countries, particularly in the developing world, should narrow their cultural property laws so that they can reap the benefits of new archaeological discoveries, which typically increase tourism and enhance cultural pride. This does not mean these nations should abolish restrictions on foreign excavation and foreign claims to artifacts.
China provides an interesting alternative approach for source nations eager for foreign archaeological investment. From 1935 to 2003, China had a restrictive cultural property law that prohibited foreign ownership of Chinese cultural artifacts. In those years, China’s most significant archaeological discovery occurred by chance, in 1974, when peasant farmers accidentally uncovered ranks of buried terra cotta warriors, which are part of Emperor Qin’s spectacular tomb system.
In 2003, the Chinese government switched course, dropping its cultural property law and embracing collaborative international archaeological research. Since then, China has nominated 11 archaeological sites for inclusion in the World Heritage Site list, including eight in 2013, the most ever for China.

\vspace{1em}
\textbf{Question 6}

The author endorses Pinker’s views on the importance of logical reasoning as it:

\begin{enumerate}[label=\Alph*.]
    \item provides a moral compass for resolving important ethical dilemmas.
    \item focuses public attention on real issues like development rather than sensational events.
    \item equips people with the ability to tackle challenging practical problems.
    \item helps people to gain expertise in statistics and other scientific disciplines.
[](https://cracku.in/6-the-author-endorses-pinkers-views-on-the-importanc-x-cat-2023-slot-3)
\end{enumerate}
\vspace{1em}
\textbf{Solution:}

The passage emphasises Pinker's focus on sequential reasoning and the tools of rationality, suggesting that greater mastery of these tools can improve decision-making in various practical contexts where individuals must act on ‘uncertain and shifting information.’ The author’s endorsement or support for Pinker’s work is centred on the idea that logical reasoning “equips people with the ability to tackle challenging practical problems” [Option C]. 
Option A is incorrect - while the author acknowledges that rationality is seen by Pinker as a moral virtue, he adds that this role of moral and ethical education is underexplored in Pinker's work. Option B presents a very specific use case of Pinker’s views and fails to capture the broader message. Option D is similarly limited in scope - the emphasis is more on the broader applicability of rationality in decision-making.
Hence, Option C is the correct choice.

\end{problem}

\newpage
\begin{problem}{}{}
% Subject: varc
\textbf{Instructions}

Instructions   

The passage below is accompanied by four questions. Based on the passage, choose the best answer for each question.
In 2006, the Met [art museum in the US] agreed to return the Euphronios krater, a masterpiece Greek urn that had been a museum draw since 1972. In 2007, the Getty [art museum in the US] agreed to return 40 objects to Italy, including a marble Aphrodite, in the midst of looting scandals. And in December, Sotheby’s and a private owner agreed to return an ancient Khmer statue of a warrior, pulled from auction two years before, to Cambodia.
Cultural property, or patrimony, laws limit the transfer of cultural property outside the source country’s territory, including outright export prohibitions and national ownership laws. Most art historians, archaeologists, museum officials and policymakers portray cultural property laws in general as invaluable tools for counteracting the ugly legacy of Western cultural imperialism.
During the late 19th and early 20th century — an era former Met director Thomas Hoving called “the age of piracy” — American and European art museums acquired antiquities by hook or by crook, from grave robbers or souvenir collectors, bounty from digs and ancient sites in impoverished but art-rich source countries. Patrimony laws were intended to protect future archaeological discoveries against Western imperialist designs. . . .
I surveyed 90 countries with one or more archaeological sites on UNESCO’s World Heritage Site list, and my study shows that in most cases the number of discovered sites diminishes sharply after a country passes a cultural property law. There are 222 archaeological sites listed for those 90 countries. When you look into the history of the sites, you see that all but 21 were discovered before the passage of cultural property laws. . . .
Strict cultural patrimony laws are popular in most countries. But the downside may be that they reduce incentives for foreign governments, nongovernmental organizations and educational institutions to invest in overseas exploration because their efforts will not necessarily be rewarded by opportunities to hold, display and study what is uncovered. To the extent that source countries can fund their own archaeological projects, artifacts and sites may still be discovered. . . . The survey has far-reaching implications. It suggests that source countries, particularly in the developing world, should narrow their cultural property laws so that they can reap the benefits of new archaeological discoveries, which typically increase tourism and enhance cultural pride. This does not mean these nations should abolish restrictions on foreign excavation and foreign claims to artifacts.
China provides an interesting alternative approach for source nations eager for foreign archaeological investment. From 1935 to 2003, China had a restrictive cultural property law that prohibited foreign ownership of Chinese cultural artifacts. In those years, China’s most significant archaeological discovery occurred by chance, in 1974, when peasant farmers accidentally uncovered ranks of buried terra cotta warriors, which are part of Emperor Qin’s spectacular tomb system.
In 2003, the Chinese government switched course, dropping its cultural property law and embracing collaborative international archaeological research. Since then, China has nominated 11 archaeological sites for inclusion in the World Heritage Site list, including eight in 2013, the most ever for China.

\vspace{1em}
\textbf{Question 7}

The author mentions Kekulé’s discovery of the structure of benzene and Mozart’s symphonies to illustrate the point that:

\begin{enumerate}[label=\Alph*.]
    \item great innovations across various fields can stem from flashes of intuition and are not always propelled by logical thinking.
    \item Pinker’s conclusions on sequential reasoning are belied by European achievements which, in the past, were more rooted in unconscious bursts of genius.
    \item it is not just the creative arts, but also scientific fields that have benefitted from flashes of creativity.
    \item unlike the sciences, human achievements in other fields are a mix of logical reasoning and spontaneous epiphanies.
[](https://cracku.in/7-the-author-mentions-kekules-discovery-of-the-struc-x-cat-2023-slot-3)
\end{enumerate}
\vspace{1em}
\textbf{Solution:}

In the case of Kekulé, the discovery of the benzene structure reportedly came to him in a dream, showcasing how creative insights can emerge unexpectedly and unconsciously. Similarly, Mozart's symphonies, considered masterpieces of classical music, are often seen as products of his musical genius and creative intuition. Therefore, the examples support the notion that groundbreaking achievements in both scientific and artistic domains may involve moments of inspiration, intuition, or epiphany, challenging the idea that all significant accomplishments are the result of conscious and sequential reasoning. This aligns with the broader point that while conscious reasoning is valuable, there are also subconscious and intuitive processes at play in the generation of innovative ideas and creations. Option A correctly captures this idea.

\end{problem}

\newpage
\begin{problem}{}{}
% Subject: varc
\textbf{Instructions}

Instructions   

The passage below is accompanied by four questions. Based on the passage, choose the best answer for each question.
In 2006, the Met [art museum in the US] agreed to return the Euphronios krater, a masterpiece Greek urn that had been a museum draw since 1972. In 2007, the Getty [art museum in the US] agreed to return 40 objects to Italy, including a marble Aphrodite, in the midst of looting scandals. And in December, Sotheby’s and a private owner agreed to return an ancient Khmer statue of a warrior, pulled from auction two years before, to Cambodia.
Cultural property, or patrimony, laws limit the transfer of cultural property outside the source country’s territory, including outright export prohibitions and national ownership laws. Most art historians, archaeologists, museum officials and policymakers portray cultural property laws in general as invaluable tools for counteracting the ugly legacy of Western cultural imperialism.
During the late 19th and early 20th century — an era former Met director Thomas Hoving called “the age of piracy” — American and European art museums acquired antiquities by hook or by crook, from grave robbers or souvenir collectors, bounty from digs and ancient sites in impoverished but art-rich source countries. Patrimony laws were intended to protect future archaeological discoveries against Western imperialist designs. . . .
I surveyed 90 countries with one or more archaeological sites on UNESCO’s World Heritage Site list, and my study shows that in most cases the number of discovered sites diminishes sharply after a country passes a cultural property law. There are 222 archaeological sites listed for those 90 countries. When you look into the history of the sites, you see that all but 21 were discovered before the passage of cultural property laws. . . .
Strict cultural patrimony laws are popular in most countries. But the downside may be that they reduce incentives for foreign governments, nongovernmental organizations and educational institutions to invest in overseas exploration because their efforts will not necessarily be rewarded by opportunities to hold, display and study what is uncovered. To the extent that source countries can fund their own archaeological projects, artifacts and sites may still be discovered. . . . The survey has far-reaching implications. It suggests that source countries, particularly in the developing world, should narrow their cultural property laws so that they can reap the benefits of new archaeological discoveries, which typically increase tourism and enhance cultural pride. This does not mean these nations should abolish restrictions on foreign excavation and foreign claims to artifacts.
China provides an interesting alternative approach for source nations eager for foreign archaeological investment. From 1935 to 2003, China had a restrictive cultural property law that prohibited foreign ownership of Chinese cultural artifacts. In those years, China’s most significant archaeological discovery occurred by chance, in 1974, when peasant farmers accidentally uncovered ranks of buried terra cotta warriors, which are part of Emperor Qin’s spectacular tomb system.
In 2003, the Chinese government switched course, dropping its cultural property law and embracing collaborative international archaeological research. Since then, China has nominated 11 archaeological sites for inclusion in the World Heritage Site list, including eight in 2013, the most ever for China.

\vspace{1em}
\textbf{Question 8}

The author refers to the ancient Greek philosophers to:

\begin{enumerate}[label=\Alph*.]
    \item show how dreams and visions have for centuries influenced subconscious behaviour and pathbreaking inventions.
    \item indicate the various similarities between their thinking and Pinker’s conclusions.
    \item reveal gaps in Pinker’s discussion of the importance of ethical considerations in rational behaviour.
    \item highlight the influence of their thinking on the development of Pinker’s arguments.
[](https://cracku.in/8-the-author-refers-to-the-ancient-greek-philosopher-x-cat-2023-slot-3)
\end{enumerate}
\vspace{1em}
\textbf{Solution:}

In the passage, the author notes that Pinker recognises rationality as both “a cognitive and moral virtue.” However, the author points out that this "profoundly important" connection between rationality and morality is not thoroughly developed in Pinker's book. By bringing up the ancient Greek philosophers who, according to the text, subtly explored the connection between moral character and rationality, the author is implying that Pinker's work could benefit from a more in-depth consideration of the ethical dimension of rational behaviour. Option C accurately reflects this point - none of the other choices correctly capture the intention behind mentioning the Greek philosophers.

Instructions   

**The passage below is accompanied by four questions. Based on the passage, choose the best answer for each question.**  
Understanding romantic aesthetics is not a simple undertaking for reasons that are internal to the nature of the subject. Distinguished scholars, such as Arthur Lovejoy, Northrop Frye and Isaiah Berlin, have remarked on the notorious challenges facing any attempt to define romanticism. Lovejoy, for example, claimed that romanticism is “the scandal of literary history and criticism” . . . The main difficulty in studying the romantics, according to him, is the lack of any “single real entity, or type of entity” that the concept “romanticism” designates. Lovejoy concluded, “the word ‘romantic’ has come to mean so many things that, by itself, it means nothing” . . .
The more specific task of characterizing romantic aesthetics adds to these difficulties an air of paradox. Conventionally, “aesthetics” refers to a theory concerning beauty and art or the branch of philosophy that studies these topics. However, many of the romantics rejected the identification of aesthetics with a circumscribed domain of human life that is separated from the practical and theoretical domains of life. The most characteristic romantic commitment is to the idea that the character of art and beauty and of our engagement with them should shape all aspects of human life. Being fundamental to human existence, beauty and art should be a central ingredient not only in a philosophical or artistic life, but also in the lives of ordinary men and women. Another challenge for any attempt to characterize romantic aesthetics lies in the fact that most of the romantics were poets and artists whose views of art and beauty are, for the most part, to be found not in developed theoretical accounts, but in fragments, aphorisms and poems, which are often more elusive and suggestive than conclusive.
Nevertheless, in spite of these challenges the task of characterizing romantic aesthetics is neither impossible nor undesirable, as numerous thinkers responding to Lovejoy’s radical skepticism have noted. While warning against a reductive definition of romanticism, Berlin, for example, still heralded the need for a general characterization: “[Although] one does have a certain sympathy with Lovejoy’s despair…[he is] in this instance mistaken. There was a romantic movement…and it is important to discover what it is” . . .
Recent attempts to characterize romanticism and to stress its contemporary relevance follow this path. Instead of overlooking the undeniable differences between the variety of romanticisms of different nations that Lovejoy had stressed, such studies attempt to characterize romanticism, not in terms of a single definition, a specific time, or a specific place, but in terms of “particular philosophical questions and concerns” . . .
While the German, British and French romantics are all considered, the central protagonists in the following are the German romantics. Two reasons explain this focus: first, because it  
has paved the way for the other romanticisms, German romanticism has a pride of place among the different national romanticisms . . . Second, the aesthetic outlook that was developed in Germany roughly between 1796 and 1801-02 — the period that corresponds to the heyday of what is known as “Early Romanticism” . . .— offers the most philosophical expression of romanticism since it is grounded primarily in the epistemological, metaphysical, ethical, and political concerns that the German romantics discerned in the aftermath of Kant’s philosophy.

\end{problem}

\newpage
\begin{problem}{}{}
% Subject: varc
\textbf{Instructions}

Instructions   

The passage below is accompanied by four questions. Based on the passage, choose the best answer for each question.
In 2006, the Met [art museum in the US] agreed to return the Euphronios krater, a masterpiece Greek urn that had been a museum draw since 1972. In 2007, the Getty [art museum in the US] agreed to return 40 objects to Italy, including a marble Aphrodite, in the midst of looting scandals. And in December, Sotheby’s and a private owner agreed to return an ancient Khmer statue of a warrior, pulled from auction two years before, to Cambodia.
Cultural property, or patrimony, laws limit the transfer of cultural property outside the source country’s territory, including outright export prohibitions and national ownership laws. Most art historians, archaeologists, museum officials and policymakers portray cultural property laws in general as invaluable tools for counteracting the ugly legacy of Western cultural imperialism.
During the late 19th and early 20th century — an era former Met director Thomas Hoving called “the age of piracy” — American and European art museums acquired antiquities by hook or by crook, from grave robbers or souvenir collectors, bounty from digs and ancient sites in impoverished but art-rich source countries. Patrimony laws were intended to protect future archaeological discoveries against Western imperialist designs. . . .
I surveyed 90 countries with one or more archaeological sites on UNESCO’s World Heritage Site list, and my study shows that in most cases the number of discovered sites diminishes sharply after a country passes a cultural property law. There are 222 archaeological sites listed for those 90 countries. When you look into the history of the sites, you see that all but 21 were discovered before the passage of cultural property laws. . . .
Strict cultural patrimony laws are popular in most countries. But the downside may be that they reduce incentives for foreign governments, nongovernmental organizations and educational institutions to invest in overseas exploration because their efforts will not necessarily be rewarded by opportunities to hold, display and study what is uncovered. To the extent that source countries can fund their own archaeological projects, artifacts and sites may still be discovered. . . . The survey has far-reaching implications. It suggests that source countries, particularly in the developing world, should narrow their cultural property laws so that they can reap the benefits of new archaeological discoveries, which typically increase tourism and enhance cultural pride. This does not mean these nations should abolish restrictions on foreign excavation and foreign claims to artifacts.
China provides an interesting alternative approach for source nations eager for foreign archaeological investment. From 1935 to 2003, China had a restrictive cultural property law that prohibited foreign ownership of Chinese cultural artifacts. In those years, China’s most significant archaeological discovery occurred by chance, in 1974, when peasant farmers accidentally uncovered ranks of buried terra cotta warriors, which are part of Emperor Qin’s spectacular tomb system.
In 2003, the Chinese government switched course, dropping its cultural property law and embracing collaborative international archaeological research. Since then, China has nominated 11 archaeological sites for inclusion in the World Heritage Site list, including eight in 2013, the most ever for China.

\vspace{1em}
\textbf{Question 9}

The main difficulty in studying romanticism is the:

\begin{enumerate}[label=\Alph*.]
    \item elusive and suggestive nature of romantic aesthetics.
    \item lack of clear conceptual contours of the domain.
    \item controversial and scandalous history of romantic literature.
    \item absence of written accounts by romantic poets and artists.
[](https://cracku.in/9-the-main-difficulty-in-studying-romanticism-is-the-x-cat-2023-slot-3)
\end{enumerate}
\vspace{1em}
\textbf{Solution:}

“ _The main difficulty in studying the romantics, according to him, is the lack of any “single real entity, or type of entity” that the concept “romanticism” designates.”_
Option B is the correct answer because it accurately captures the main difficulty highlighted in the passage when studying romanticism. The passage, particularly referencing Arthur Lovejoy, emphasizes the challenge posed by the lack of clear conceptual contours or a single, cohesive entity associated with the term "romanticism." Lovejoy's assertion that romanticism is the "scandal of literary history and criticism" underscores the difficulty in defining the boundaries and characteristics of this literary and artistic movement.
The elusive and suggestive nature of romantic aesthetics (Option A) is mentioned in the passage as a challenge, but it is not identified as the main difficulty. Similarly, the controversial and scandalous history of romantic literature (Option C) is not specified as the primary obstacle. The absence of written accounts by romantic poets and artists (Option D) is acknowledged as a challenge, but the primary difficulty highlighted in the passage is the lack of clear conceptual contours associated with romanticism.

\end{problem}

\newpage
\begin{problem}{}{}
% Subject: varc
\textbf{Instructions}

Instructions   

The passage below is accompanied by four questions. Based on the passage, choose the best answer for each question.
In 2006, the Met [art museum in the US] agreed to return the Euphronios krater, a masterpiece Greek urn that had been a museum draw since 1972. In 2007, the Getty [art museum in the US] agreed to return 40 objects to Italy, including a marble Aphrodite, in the midst of looting scandals. And in December, Sotheby’s and a private owner agreed to return an ancient Khmer statue of a warrior, pulled from auction two years before, to Cambodia.
Cultural property, or patrimony, laws limit the transfer of cultural property outside the source country’s territory, including outright export prohibitions and national ownership laws. Most art historians, archaeologists, museum officials and policymakers portray cultural property laws in general as invaluable tools for counteracting the ugly legacy of Western cultural imperialism.
During the late 19th and early 20th century — an era former Met director Thomas Hoving called “the age of piracy” — American and European art museums acquired antiquities by hook or by crook, from grave robbers or souvenir collectors, bounty from digs and ancient sites in impoverished but art-rich source countries. Patrimony laws were intended to protect future archaeological discoveries against Western imperialist designs. . . .
I surveyed 90 countries with one or more archaeological sites on UNESCO’s World Heritage Site list, and my study shows that in most cases the number of discovered sites diminishes sharply after a country passes a cultural property law. There are 222 archaeological sites listed for those 90 countries. When you look into the history of the sites, you see that all but 21 were discovered before the passage of cultural property laws. . . .
Strict cultural patrimony laws are popular in most countries. But the downside may be that they reduce incentives for foreign governments, nongovernmental organizations and educational institutions to invest in overseas exploration because their efforts will not necessarily be rewarded by opportunities to hold, display and study what is uncovered. To the extent that source countries can fund their own archaeological projects, artifacts and sites may still be discovered. . . . The survey has far-reaching implications. It suggests that source countries, particularly in the developing world, should narrow their cultural property laws so that they can reap the benefits of new archaeological discoveries, which typically increase tourism and enhance cultural pride. This does not mean these nations should abolish restrictions on foreign excavation and foreign claims to artifacts.
China provides an interesting alternative approach for source nations eager for foreign archaeological investment. From 1935 to 2003, China had a restrictive cultural property law that prohibited foreign ownership of Chinese cultural artifacts. In those years, China’s most significant archaeological discovery occurred by chance, in 1974, when peasant farmers accidentally uncovered ranks of buried terra cotta warriors, which are part of Emperor Qin’s spectacular tomb system.
In 2003, the Chinese government switched course, dropping its cultural property law and embracing collaborative international archaeological research. Since then, China has nominated 11 archaeological sites for inclusion in the World Heritage Site list, including eight in 2013, the most ever for China.

\vspace{1em}
\textbf{Question 10}

According to the romantics, aesthetics:

\begin{enumerate}[label=\Alph*.]
    \item should be confined to a specific domain separate from the practical and theoretical aspects of life.
    \item is primarily the concern of philosophers and artists, rather than of ordinary people.
    \item is widely considered to be irrelevant to human existence.
    \item permeates all aspects of human life, philosophical and mundane.
[](https://cracku.in/10-according-to-the-romantics-aesthetics-x-cat-2023-slot-3)
\end{enumerate}
\vspace{1em}
\textbf{Solution:}

“ _The most characteristic romantic commitment is to the idea that the character of art and beauty and of our engagement with them should shape all aspects of human life. Being fundamental to human existence, beauty and art should be a central ingredient not only in a philosophical or artistic life, but also in the lives of ordinary men and women.”_
According to the passage, the romantics rejected the idea of confining aesthetics to a specific domain separate from practical and theoretical aspects of life. Instead, they believed that aesthetics, encompassing the character of art and beauty, should permeate all aspects of human existence, not only in philosophical or artistic lives but also in the lives of ordinary men and women.Therefore Option D is the correct answer.

\end{problem}

\newpage
\begin{problem}{}{}
% Subject: varc
\textbf{Instructions}

Instructions   

The passage below is accompanied by four questions. Based on the passage, choose the best answer for each question.
In 2006, the Met [art museum in the US] agreed to return the Euphronios krater, a masterpiece Greek urn that had been a museum draw since 1972. In 2007, the Getty [art museum in the US] agreed to return 40 objects to Italy, including a marble Aphrodite, in the midst of looting scandals. And in December, Sotheby’s and a private owner agreed to return an ancient Khmer statue of a warrior, pulled from auction two years before, to Cambodia.
Cultural property, or patrimony, laws limit the transfer of cultural property outside the source country’s territory, including outright export prohibitions and national ownership laws. Most art historians, archaeologists, museum officials and policymakers portray cultural property laws in general as invaluable tools for counteracting the ugly legacy of Western cultural imperialism.
During the late 19th and early 20th century — an era former Met director Thomas Hoving called “the age of piracy” — American and European art museums acquired antiquities by hook or by crook, from grave robbers or souvenir collectors, bounty from digs and ancient sites in impoverished but art-rich source countries. Patrimony laws were intended to protect future archaeological discoveries against Western imperialist designs. . . .
I surveyed 90 countries with one or more archaeological sites on UNESCO’s World Heritage Site list, and my study shows that in most cases the number of discovered sites diminishes sharply after a country passes a cultural property law. There are 222 archaeological sites listed for those 90 countries. When you look into the history of the sites, you see that all but 21 were discovered before the passage of cultural property laws. . . .
Strict cultural patrimony laws are popular in most countries. But the downside may be that they reduce incentives for foreign governments, nongovernmental organizations and educational institutions to invest in overseas exploration because their efforts will not necessarily be rewarded by opportunities to hold, display and study what is uncovered. To the extent that source countries can fund their own archaeological projects, artifacts and sites may still be discovered. . . . The survey has far-reaching implications. It suggests that source countries, particularly in the developing world, should narrow their cultural property laws so that they can reap the benefits of new archaeological discoveries, which typically increase tourism and enhance cultural pride. This does not mean these nations should abolish restrictions on foreign excavation and foreign claims to artifacts.
China provides an interesting alternative approach for source nations eager for foreign archaeological investment. From 1935 to 2003, China had a restrictive cultural property law that prohibited foreign ownership of Chinese cultural artifacts. In those years, China’s most significant archaeological discovery occurred by chance, in 1974, when peasant farmers accidentally uncovered ranks of buried terra cotta warriors, which are part of Emperor Qin’s spectacular tomb system.
In 2003, the Chinese government switched course, dropping its cultural property law and embracing collaborative international archaeological research. Since then, China has nominated 11 archaeological sites for inclusion in the World Heritage Site list, including eight in 2013, the most ever for China.

\vspace{1em}
\textbf{Question 11}

Which one of the following statements is NOT supported by the passage?

\begin{enumerate}[label=\Alph*.]
    \item Characterising romantic aesthetics is both possible and desirable, despite the challenges involved.
    \item Recent studies on romanticism seek to refute the differences between national romanticisms.
    \item Romantic aesthetics are primarily expressed through fragments, aphorisms, and poems.
    \item Many romantics rejected the idea of aesthetics as a domain separate from other aspects of life.
[](https://cracku.in/11-which-one-of-the-following-statements-is-not-suppo-x-cat-2023-slot-3)
\end{enumerate}
\vspace{1em}
\textbf{Solution:}

The passage suggests that recent studies on romanticism do not seek to overlook the differences between national romanticisms but rather attempt to characterize romanticism in terms of particular philosophical questions and concerns. The focus is on understanding romanticism without ignoring the diversity among different national expressions of it. Therefore Option B is the correct answer as it is not supported by the passage.
Option A: Although the passage acknowledges the challenges in characterizing romantic aesthetics, it also argues that it is not impossible or undesirable. Scholars recognize the difficulty but still emphasize the importance of discovering the nature of romanticism.  
Option C: The passage mentions that the views of romantics on art and beauty are often found in fragments, aphorisms, and poems rather than in fully developed theoretical accounts. This emphasizes the elusive and suggestive nature of their expressions.
Option D: The passage notes that many romantics rejected the identification of aesthetics with a circumscribed domain of human life that is separated from the practical and theoretical domains of life. Instead, they believed that the character of art and beauty should shape all aspects of human life.

\end{problem}

\newpage
\begin{problem}{}{}
% Subject: varc
\textbf{Instructions}

Instructions   

The passage below is accompanied by four questions. Based on the passage, choose the best answer for each question.
In 2006, the Met [art museum in the US] agreed to return the Euphronios krater, a masterpiece Greek urn that had been a museum draw since 1972. In 2007, the Getty [art museum in the US] agreed to return 40 objects to Italy, including a marble Aphrodite, in the midst of looting scandals. And in December, Sotheby’s and a private owner agreed to return an ancient Khmer statue of a warrior, pulled from auction two years before, to Cambodia.
Cultural property, or patrimony, laws limit the transfer of cultural property outside the source country’s territory, including outright export prohibitions and national ownership laws. Most art historians, archaeologists, museum officials and policymakers portray cultural property laws in general as invaluable tools for counteracting the ugly legacy of Western cultural imperialism.
During the late 19th and early 20th century — an era former Met director Thomas Hoving called “the age of piracy” — American and European art museums acquired antiquities by hook or by crook, from grave robbers or souvenir collectors, bounty from digs and ancient sites in impoverished but art-rich source countries. Patrimony laws were intended to protect future archaeological discoveries against Western imperialist designs. . . .
I surveyed 90 countries with one or more archaeological sites on UNESCO’s World Heritage Site list, and my study shows that in most cases the number of discovered sites diminishes sharply after a country passes a cultural property law. There are 222 archaeological sites listed for those 90 countries. When you look into the history of the sites, you see that all but 21 were discovered before the passage of cultural property laws. . . .
Strict cultural patrimony laws are popular in most countries. But the downside may be that they reduce incentives for foreign governments, nongovernmental organizations and educational institutions to invest in overseas exploration because their efforts will not necessarily be rewarded by opportunities to hold, display and study what is uncovered. To the extent that source countries can fund their own archaeological projects, artifacts and sites may still be discovered. . . . The survey has far-reaching implications. It suggests that source countries, particularly in the developing world, should narrow their cultural property laws so that they can reap the benefits of new archaeological discoveries, which typically increase tourism and enhance cultural pride. This does not mean these nations should abolish restrictions on foreign excavation and foreign claims to artifacts.
China provides an interesting alternative approach for source nations eager for foreign archaeological investment. From 1935 to 2003, China had a restrictive cultural property law that prohibited foreign ownership of Chinese cultural artifacts. In those years, China’s most significant archaeological discovery occurred by chance, in 1974, when peasant farmers accidentally uncovered ranks of buried terra cotta warriors, which are part of Emperor Qin’s spectacular tomb system.
In 2003, the Chinese government switched course, dropping its cultural property law and embracing collaborative international archaeological research. Since then, China has nominated 11 archaeological sites for inclusion in the World Heritage Site list, including eight in 2013, the most ever for China.

\vspace{1em}
\textbf{Question 12}

According to the passage, recent studies on romanticism avoid “a single definition, a specific time, or a specific place” because they:

\begin{enumerate}[label=\Alph*.]
    \item understand that the variety of romanticisms renders a general analysis impossible.
    \item prefer to highlight the paradox of romantic aesthetics as a concept.
    \item prefer to focus on the fundamental concerns of the romantics.
    \item seek to discredit Lovejoy’s scepticism regarding romanticism.
[](https://cracku.in/12-according-to-the-passage-recent-studies-on-romanti-x-cat-2023-slot-3)
\end{enumerate}
\vspace{1em}
\textbf{Solution:}

Option C is the correct answer because it accurately reflects the passage's explanation of why recent studies on romanticism avoid seeking "a single definition, a specific time, or a specific place." According to the passage, these studies opt to characterize romanticism in terms of "particular philosophical questions and concerns" rather than attempting to provide a singular, all-encompassing definition. The reason for this approach is to delve into the fundamental concerns of the romantics, recognizing that romanticism is a complex and multifaceted movement with diverse expressions in different nations and contexts.
The passage suggests that romanticism is not easily confined to a single, universally applicable definition due to the variety of romanticisms in different nations. Instead of discrediting or refuting Lovejoy's skepticism (Option D), recent studies acknowledge the challenges and complexities of defining romanticism but seek to understand it by focusing on the core philosophical questions and concerns that were central to the romantics' worldview.

Instructions   

The passage below is accompanied by four questions. Based on the passage, choose the best answer for each question.
Comprehension:  
The biggest challenge [The Nutmeg’s Curse by Ghosh] throws down is to the prevailing understanding of when the climate crisis started. Most of us have accepted . . . that it started with the widespread use of coal at the beginning of the Industrial Age in the 18th century and worsened with the mass adoption of oil and natural gas in the 20th. Ghosh takes this history at least three centuries back, to the start of European colonialism in the 15th century. He [starts] the book with a 1621 massacre by Dutch invaders determined to impose a monopoly on nutmeg cultivation and trade in the Banda islands in today’s Indonesia. Not only do the Dutch systematically depopulate the islands through genocide, they also try their best to bring nutmeg cultivation into plantation mode. These are the two points to which Ghosh returns through examples from around the world. One, how European colonialists decimated not only indigenous populations but also indigenous understanding of the relationship between humans and Earth. Two, how this was an invasion not only of humans but of the Earth itself, and how this continues to the present day by looking at nature as a ‘resource’ to exploit. . . .
We know we are facing more frequent and more severe heatwaves, storms, floods, droughts and wildfires due to climate change. We know our expansion through deforestation, dam building, canal cutting - in short, terraforming, the word Ghosh uses - has brought us repeated disasters . . . Are these the responses of an angry Gaia who has finally had enough? By using the word ‘curse’ in the title, the author makes it clear that he thinks so. I use the pronoun ‘who’ knowingly, because Ghosh has quoted many non-European sources to enquire into the relationship between humans and the world around them so that he can question the prevalent way of looking at Earth as an inert object to be exploited to the maximum.
As Ghosh’s text, notes and bibliography show once more, none of this is new. There have always been challenges to the way European colonialists looked at other civilisations and at Earth. It is just that the invaders and their myriad backers in the fields of economics, politics, anthropology, philosophy, literature, technology, physics, chemistry, biology have dominated global intellectual discourse. . . .
There are other points of view that we can hear today if we listen hard enough. Those observing global climate negotiations know about the Latin American way of looking at Earth as Pachamama (Earth Mother). They also know how such a framing is just provided lip service and is ignored in the substantive portions of the negotiations. In The Nutmeg’s Curse, Ghosh explains why. He shows the extent of the vested interest in the oil economy - not only for oil-exporting countries, but also for a superpower like the US that controls oil drilling, oil prices and oil movement around the world. Many of us know power utilities are sabotaging decentralised solar power generation today because it hits their revenues and control. And how the other points of view are so often drowned out.

\end{problem}

\newpage
\begin{problem}{}{}
% Subject: varc
\textbf{Instructions}

Instructions   

The passage below is accompanied by four questions. Based on the passage, choose the best answer for each question.
In 2006, the Met [art museum in the US] agreed to return the Euphronios krater, a masterpiece Greek urn that had been a museum draw since 1972. In 2007, the Getty [art museum in the US] agreed to return 40 objects to Italy, including a marble Aphrodite, in the midst of looting scandals. And in December, Sotheby’s and a private owner agreed to return an ancient Khmer statue of a warrior, pulled from auction two years before, to Cambodia.
Cultural property, or patrimony, laws limit the transfer of cultural property outside the source country’s territory, including outright export prohibitions and national ownership laws. Most art historians, archaeologists, museum officials and policymakers portray cultural property laws in general as invaluable tools for counteracting the ugly legacy of Western cultural imperialism.
During the late 19th and early 20th century — an era former Met director Thomas Hoving called “the age of piracy” — American and European art museums acquired antiquities by hook or by crook, from grave robbers or souvenir collectors, bounty from digs and ancient sites in impoverished but art-rich source countries. Patrimony laws were intended to protect future archaeological discoveries against Western imperialist designs. . . .
I surveyed 90 countries with one or more archaeological sites on UNESCO’s World Heritage Site list, and my study shows that in most cases the number of discovered sites diminishes sharply after a country passes a cultural property law. There are 222 archaeological sites listed for those 90 countries. When you look into the history of the sites, you see that all but 21 were discovered before the passage of cultural property laws. . . .
Strict cultural patrimony laws are popular in most countries. But the downside may be that they reduce incentives for foreign governments, nongovernmental organizations and educational institutions to invest in overseas exploration because their efforts will not necessarily be rewarded by opportunities to hold, display and study what is uncovered. To the extent that source countries can fund their own archaeological projects, artifacts and sites may still be discovered. . . . The survey has far-reaching implications. It suggests that source countries, particularly in the developing world, should narrow their cultural property laws so that they can reap the benefits of new archaeological discoveries, which typically increase tourism and enhance cultural pride. This does not mean these nations should abolish restrictions on foreign excavation and foreign claims to artifacts.
China provides an interesting alternative approach for source nations eager for foreign archaeological investment. From 1935 to 2003, China had a restrictive cultural property law that prohibited foreign ownership of Chinese cultural artifacts. In those years, China’s most significant archaeological discovery occurred by chance, in 1974, when peasant farmers accidentally uncovered ranks of buried terra cotta warriors, which are part of Emperor Qin’s spectacular tomb system.
In 2003, the Chinese government switched course, dropping its cultural property law and embracing collaborative international archaeological research. Since then, China has nominated 11 archaeological sites for inclusion in the World Heritage Site list, including eight in 2013, the most ever for China.

\vspace{1em}
\textbf{Question 13}

On the basis of information in the passage, which one of the following is NOT a reason for the failure of policies seeking to address climate change?

\begin{enumerate}[label=\Alph*.]
    \item The greed of organisations benefiting from non-renewable energy resources.
    \item The global dominance of oil economies and international politics built around it.
    \item The marginalised status of non-European ways of looking at nature and the environment.
    \item The decentralised characteristic of renewable energy resources like solar power.
[](https://cracku.in/13-on-the-basis-of-information-in-the-passage-which-o-x-cat-2023-slot-3)
\end{enumerate}
\vspace{1em}
\textbf{Solution:}

“ _Those observing global climate negotiations know about the Latin American way of looking at Earth as Pachamama (Earth Mother)._**_They also know how such a framing is just provided lip service and is ignored in the substantive portions of the negotiations._**_In The Nutmeg’s Curse, Ghosh explains why. He shows_** _the extent of the vested interest in the oil economy - not only for oil-exporting countries, but also for a superpower like the US that controls oil drilling, oil prices and oil movement around the world._**_Many of us know_** _power utilities are sabotaging decentralised solar power generation today because it hits their revenues and control.”_**
From the highlighted part we can clearly infer Options A, B and C. The passage does not suggest that the decentralised characteristic of renewable energy resources like solar power is a reason for the failure of climate change policies. Instead, it mentions that power utilities may be sabotaging decentralized solar power generation because it affects their revenues and control, but it does not frame the decentralised nature of renewable energy as a cause for failure. Therefore Option D is not a reason for the failure of policies seeking to address climate change.

\end{problem}

\newpage
\begin{problem}{}{}
% Subject: varc
\textbf{Instructions}

Instructions   

The passage below is accompanied by four questions. Based on the passage, choose the best answer for each question.
In 2006, the Met [art museum in the US] agreed to return the Euphronios krater, a masterpiece Greek urn that had been a museum draw since 1972. In 2007, the Getty [art museum in the US] agreed to return 40 objects to Italy, including a marble Aphrodite, in the midst of looting scandals. And in December, Sotheby’s and a private owner agreed to return an ancient Khmer statue of a warrior, pulled from auction two years before, to Cambodia.
Cultural property, or patrimony, laws limit the transfer of cultural property outside the source country’s territory, including outright export prohibitions and national ownership laws. Most art historians, archaeologists, museum officials and policymakers portray cultural property laws in general as invaluable tools for counteracting the ugly legacy of Western cultural imperialism.
During the late 19th and early 20th century — an era former Met director Thomas Hoving called “the age of piracy” — American and European art museums acquired antiquities by hook or by crook, from grave robbers or souvenir collectors, bounty from digs and ancient sites in impoverished but art-rich source countries. Patrimony laws were intended to protect future archaeological discoveries against Western imperialist designs. . . .
I surveyed 90 countries with one or more archaeological sites on UNESCO’s World Heritage Site list, and my study shows that in most cases the number of discovered sites diminishes sharply after a country passes a cultural property law. There are 222 archaeological sites listed for those 90 countries. When you look into the history of the sites, you see that all but 21 were discovered before the passage of cultural property laws. . . .
Strict cultural patrimony laws are popular in most countries. But the downside may be that they reduce incentives for foreign governments, nongovernmental organizations and educational institutions to invest in overseas exploration because their efforts will not necessarily be rewarded by opportunities to hold, display and study what is uncovered. To the extent that source countries can fund their own archaeological projects, artifacts and sites may still be discovered. . . . The survey has far-reaching implications. It suggests that source countries, particularly in the developing world, should narrow their cultural property laws so that they can reap the benefits of new archaeological discoveries, which typically increase tourism and enhance cultural pride. This does not mean these nations should abolish restrictions on foreign excavation and foreign claims to artifacts.
China provides an interesting alternative approach for source nations eager for foreign archaeological investment. From 1935 to 2003, China had a restrictive cultural property law that prohibited foreign ownership of Chinese cultural artifacts. In those years, China’s most significant archaeological discovery occurred by chance, in 1974, when peasant farmers accidentally uncovered ranks of buried terra cotta warriors, which are part of Emperor Qin’s spectacular tomb system.
In 2003, the Chinese government switched course, dropping its cultural property law and embracing collaborative international archaeological research. Since then, China has nominated 11 archaeological sites for inclusion in the World Heritage Site list, including eight in 2013, the most ever for China.

\vspace{1em}
\textbf{Question 14}

Which one of the following, if true, would make the reviewer’s choice of the pronoun “who” for Gaia inappropriate?

\begin{enumerate}[label=\Alph*.]
    \item Modern western science discovers new evidence for the Earth being an inanimate object.
    \item There is a direct cause-effect relationship between human activities and global climate change.
    \item Ghosh’s book has a different title: “The Nutmeg’s Revenge”.
    \item Non-European societies have perceived the Earth as a non-living source of all resources.
[](https://cracku.in/14-which-one-of-the-following-if-true-would-make-the--x-cat-2023-slot-3)
\end{enumerate}
\vspace{1em}
\textbf{Solution:}

If non-European societies perceive the Earth as a non-living source of all resources, it contradicts the personification implied by the use of "who" for Gaia.
In the context of the passage, the author uses the word "curse" in the title, and the pronoun "who" is used for Gaia, suggesting a perspective that sees Earth as a living, conscious entity. If non-European societies do not share this perspective and view the Earth as a non-living source of resources, it challenges the appropriateness of using the pronoun "who" for Gaia in the context of their beliefs. This would make the reviewer's choice of the pronoun "who" inappropriate, given the differing perspectives on the nature of the Earth. Therefore Option D is the correct answer.

\end{problem}

\newpage
\begin{problem}{}{}
% Subject: varc
\textbf{Instructions}

Instructions   

The passage below is accompanied by four questions. Based on the passage, choose the best answer for each question.
In 2006, the Met [art museum in the US] agreed to return the Euphronios krater, a masterpiece Greek urn that had been a museum draw since 1972. In 2007, the Getty [art museum in the US] agreed to return 40 objects to Italy, including a marble Aphrodite, in the midst of looting scandals. And in December, Sotheby’s and a private owner agreed to return an ancient Khmer statue of a warrior, pulled from auction two years before, to Cambodia.
Cultural property, or patrimony, laws limit the transfer of cultural property outside the source country’s territory, including outright export prohibitions and national ownership laws. Most art historians, archaeologists, museum officials and policymakers portray cultural property laws in general as invaluable tools for counteracting the ugly legacy of Western cultural imperialism.
During the late 19th and early 20th century — an era former Met director Thomas Hoving called “the age of piracy” — American and European art museums acquired antiquities by hook or by crook, from grave robbers or souvenir collectors, bounty from digs and ancient sites in impoverished but art-rich source countries. Patrimony laws were intended to protect future archaeological discoveries against Western imperialist designs. . . .
I surveyed 90 countries with one or more archaeological sites on UNESCO’s World Heritage Site list, and my study shows that in most cases the number of discovered sites diminishes sharply after a country passes a cultural property law. There are 222 archaeological sites listed for those 90 countries. When you look into the history of the sites, you see that all but 21 were discovered before the passage of cultural property laws. . . .
Strict cultural patrimony laws are popular in most countries. But the downside may be that they reduce incentives for foreign governments, nongovernmental organizations and educational institutions to invest in overseas exploration because their efforts will not necessarily be rewarded by opportunities to hold, display and study what is uncovered. To the extent that source countries can fund their own archaeological projects, artifacts and sites may still be discovered. . . . The survey has far-reaching implications. It suggests that source countries, particularly in the developing world, should narrow their cultural property laws so that they can reap the benefits of new archaeological discoveries, which typically increase tourism and enhance cultural pride. This does not mean these nations should abolish restrictions on foreign excavation and foreign claims to artifacts.
China provides an interesting alternative approach for source nations eager for foreign archaeological investment. From 1935 to 2003, China had a restrictive cultural property law that prohibited foreign ownership of Chinese cultural artifacts. In those years, China’s most significant archaeological discovery occurred by chance, in 1974, when peasant farmers accidentally uncovered ranks of buried terra cotta warriors, which are part of Emperor Qin’s spectacular tomb system.
In 2003, the Chinese government switched course, dropping its cultural property law and embracing collaborative international archaeological research. Since then, China has nominated 11 archaeological sites for inclusion in the World Heritage Site list, including eight in 2013, the most ever for China.

\vspace{1em}
\textbf{Question 15}

All of the following can be inferred from the reviewer’s discussion of “The Nutmeg’s Curse”, EXCEPT:

\begin{enumerate}[label=\Alph*.]
    \item the history of climate change is deeply intertwined with the history of colonialism.
    \item the contemporary dominant perception of nature and the environment was put in place by processes of colonialism.
    \item environmental preservation policy makers can learn a lot from non-European and/or pre-colonial societies.
    \item academic discourses have always served the function of raising awareness about environmental preservation.
[](https://cracku.in/15-all-of-the-following-can-be-inferred-from-the-revi-x-cat-2023-slot-3)
\end{enumerate}
\vspace{1em}
\textbf{Solution:}

Option D cannot be directly drawn from the passage. The passage discusses the historical perspective on climate change presented in "The Nutmeg's Curse" and emphasizes the impact of colonialism on the contemporary dominant perception of nature and the environment. It suggests that there are alternative views from non-European and/or pre-colonial societies that can provide insights for environmental preservation policy makers. However, the passage does not explicitly state that academic discourses have always served the function of raising awareness about environmental preservation.
The passage connects the history of climate change with colonialism (Option A), highlights that colonial processes shaped the contemporary perception of nature and the environment (Option B) and suggests that non-European and/or pre-colonial societies hold valuable insights for environmental preservation policy makers (Option C).

\end{problem}

\newpage
\begin{problem}{}{}
% Subject: varc
\textbf{Instructions}

Instructions   

The passage below is accompanied by four questions. Based on the passage, choose the best answer for each question.
In 2006, the Met [art museum in the US] agreed to return the Euphronios krater, a masterpiece Greek urn that had been a museum draw since 1972. In 2007, the Getty [art museum in the US] agreed to return 40 objects to Italy, including a marble Aphrodite, in the midst of looting scandals. And in December, Sotheby’s and a private owner agreed to return an ancient Khmer statue of a warrior, pulled from auction two years before, to Cambodia.
Cultural property, or patrimony, laws limit the transfer of cultural property outside the source country’s territory, including outright export prohibitions and national ownership laws. Most art historians, archaeologists, museum officials and policymakers portray cultural property laws in general as invaluable tools for counteracting the ugly legacy of Western cultural imperialism.
During the late 19th and early 20th century — an era former Met director Thomas Hoving called “the age of piracy” — American and European art museums acquired antiquities by hook or by crook, from grave robbers or souvenir collectors, bounty from digs and ancient sites in impoverished but art-rich source countries. Patrimony laws were intended to protect future archaeological discoveries against Western imperialist designs. . . .
I surveyed 90 countries with one or more archaeological sites on UNESCO’s World Heritage Site list, and my study shows that in most cases the number of discovered sites diminishes sharply after a country passes a cultural property law. There are 222 archaeological sites listed for those 90 countries. When you look into the history of the sites, you see that all but 21 were discovered before the passage of cultural property laws. . . .
Strict cultural patrimony laws are popular in most countries. But the downside may be that they reduce incentives for foreign governments, nongovernmental organizations and educational institutions to invest in overseas exploration because their efforts will not necessarily be rewarded by opportunities to hold, display and study what is uncovered. To the extent that source countries can fund their own archaeological projects, artifacts and sites may still be discovered. . . . The survey has far-reaching implications. It suggests that source countries, particularly in the developing world, should narrow their cultural property laws so that they can reap the benefits of new archaeological discoveries, which typically increase tourism and enhance cultural pride. This does not mean these nations should abolish restrictions on foreign excavation and foreign claims to artifacts.
China provides an interesting alternative approach for source nations eager for foreign archaeological investment. From 1935 to 2003, China had a restrictive cultural property law that prohibited foreign ownership of Chinese cultural artifacts. In those years, China’s most significant archaeological discovery occurred by chance, in 1974, when peasant farmers accidentally uncovered ranks of buried terra cotta warriors, which are part of Emperor Qin’s spectacular tomb system.
In 2003, the Chinese government switched course, dropping its cultural property law and embracing collaborative international archaeological research. Since then, China has nominated 11 archaeological sites for inclusion in the World Heritage Site list, including eight in 2013, the most ever for China.

\vspace{1em}
\textbf{Question 16}

Which one of the following best explains the primary purpose of the discussion of the colonisation of the Banda islands in “The Nutmeg’s Curse”?

\begin{enumerate}[label=\Alph*.]
    \item To illustrate the role played by the cultivation of certain crops in the plantation mode in contributing to climate change.
    \item To illustrate the first instance in history when the processes responsible for climate change were initiated.
    \item To illustrate how systemic violence against the colonised constituted the cornerstone of colonialism.
    \item To illustrate how colonialism represented and perpetuated the mindset that has led to climate change.
[](https://cracku.in/16-which-one-of-the-following-best-explains-the-prima-x-cat-2023-slot-3)
\end{enumerate}
\vspace{1em}
\textbf{Solution:}

_“These are the two points to which Ghosh returns through examples from around the world. One, how European colonialists decimated not only indigenous populations but also indigenous understanding of the relationship between humans and Earth. Two, how this was an invasion not only of humans but of the Earth itself, and how this continues to the present day by looking at nature as a ‘resource’ to exploit”_
The passage discusses how the colonization of the Banda islands, as presented in "The Nutmeg’s Curse," is used to illustrate the broader idea that colonialism played a significant role in shaping the mindset and practices that have led to climate change. The exploitation of both indigenous populations and the Earth's resources during colonialism is portrayed as a key factor in perpetuating the mindset that views nature as a resource to be exploited, contributing to the environmental challenges faced today. Therefore, the primary purpose of discussing the colonization of the Banda islands is to highlight how colonialism represented and perpetuated the mindset that has led to climate change. Therefore Option D is the correct answer.

Instructions   

For the following questions answer them individually

\end{problem}

\newpage
\begin{problem}{}{}
% Subject: varc
\textbf{Instructions}

Instructions   

The passage below is accompanied by four questions. Based on the passage, choose the best answer for each question.
In 2006, the Met [art museum in the US] agreed to return the Euphronios krater, a masterpiece Greek urn that had been a museum draw since 1972. In 2007, the Getty [art museum in the US] agreed to return 40 objects to Italy, including a marble Aphrodite, in the midst of looting scandals. And in December, Sotheby’s and a private owner agreed to return an ancient Khmer statue of a warrior, pulled from auction two years before, to Cambodia.
Cultural property, or patrimony, laws limit the transfer of cultural property outside the source country’s territory, including outright export prohibitions and national ownership laws. Most art historians, archaeologists, museum officials and policymakers portray cultural property laws in general as invaluable tools for counteracting the ugly legacy of Western cultural imperialism.
During the late 19th and early 20th century — an era former Met director Thomas Hoving called “the age of piracy” — American and European art museums acquired antiquities by hook or by crook, from grave robbers or souvenir collectors, bounty from digs and ancient sites in impoverished but art-rich source countries. Patrimony laws were intended to protect future archaeological discoveries against Western imperialist designs. . . .
I surveyed 90 countries with one or more archaeological sites on UNESCO’s World Heritage Site list, and my study shows that in most cases the number of discovered sites diminishes sharply after a country passes a cultural property law. There are 222 archaeological sites listed for those 90 countries. When you look into the history of the sites, you see that all but 21 were discovered before the passage of cultural property laws. . . .
Strict cultural patrimony laws are popular in most countries. But the downside may be that they reduce incentives for foreign governments, nongovernmental organizations and educational institutions to invest in overseas exploration because their efforts will not necessarily be rewarded by opportunities to hold, display and study what is uncovered. To the extent that source countries can fund their own archaeological projects, artifacts and sites may still be discovered. . . . The survey has far-reaching implications. It suggests that source countries, particularly in the developing world, should narrow their cultural property laws so that they can reap the benefits of new archaeological discoveries, which typically increase tourism and enhance cultural pride. This does not mean these nations should abolish restrictions on foreign excavation and foreign claims to artifacts.
China provides an interesting alternative approach for source nations eager for foreign archaeological investment. From 1935 to 2003, China had a restrictive cultural property law that prohibited foreign ownership of Chinese cultural artifacts. In those years, China’s most significant archaeological discovery occurred by chance, in 1974, when peasant farmers accidentally uncovered ranks of buried terra cotta warriors, which are part of Emperor Qin’s spectacular tomb system.
In 2003, the Chinese government switched course, dropping its cultural property law and embracing collaborative international archaeological research. Since then, China has nominated 11 archaeological sites for inclusion in the World Heritage Site list, including eight in 2013, the most ever for China.

\vspace{1em}
\textbf{Question 17}

There is a sentence that is missing in the paragraph below. Look at the paragraph and decide where (option 1, 2, 3, or 4) the following sentence would best fit.  
Sentence: Beyond undermining the monopoly of the State on the use of force, armed conflict also creates an environment that can enable organized crime to prosper.
Paragraph: ___(1)___. Linkages between illicit arms, organized crime, and armed conflict can reinforce one another while also escalating and prolonging violence and eroding governance.___(2)___. Financial gains from crime can lengthen or intensify armed conflicts by creating revenue streams for non-State armed groups (NSAGs).___(3)___. In this context, when hostilities cease and parties to a conflict move towards a peaceful resolution, the widespread availability of surplus arms and ammunition can contribute to a situation of ‘criminalized peace’ that obstructs sustainable peacebuilding efforts.___(4)___.

\begin{enumerate}[label=\Alph*.]
    \item Option 4
    \item Option 3
    \item Option 1
    \item Option 2
[](https://cracku.in/17-there-is-a-sentence-that-is-missing-in-the-paragra-x-cat-2023-slot-3)
\end{enumerate}
\vspace{1em}
\textbf{Solution:}

The Sentence best fits in Blank 3 because it directly relates to the theme introduced in before Blank 3. Before Blank 3 the passage discusses the linkages between illicit arms, organized crime, and armed conflict, and the provided sentence explains how armed conflict creates an environment favorable for organized crime to thrive. Therefore, Blank 3 is the appropriate placement for the sentence, enhancing the coherence and thematic continuity of the paragraph.
Blank 1 and 4 can be easily eliminated as the Sentence in question does not serve as a proper introductory or a concluding statement.
The provided sentence introduces a different aspect, focusing on the broader impact of armed conflict on organized crime rather than the financial dynamics discussed after Blank 2.

\end{problem}

\newpage
\begin{problem}{}{}
% Subject: varc
\textbf{Instructions}

Instructions   

The passage below is accompanied by four questions. Based on the passage, choose the best answer for each question.
In 2006, the Met [art museum in the US] agreed to return the Euphronios krater, a masterpiece Greek urn that had been a museum draw since 1972. In 2007, the Getty [art museum in the US] agreed to return 40 objects to Italy, including a marble Aphrodite, in the midst of looting scandals. And in December, Sotheby’s and a private owner agreed to return an ancient Khmer statue of a warrior, pulled from auction two years before, to Cambodia.
Cultural property, or patrimony, laws limit the transfer of cultural property outside the source country’s territory, including outright export prohibitions and national ownership laws. Most art historians, archaeologists, museum officials and policymakers portray cultural property laws in general as invaluable tools for counteracting the ugly legacy of Western cultural imperialism.
During the late 19th and early 20th century — an era former Met director Thomas Hoving called “the age of piracy” — American and European art museums acquired antiquities by hook or by crook, from grave robbers or souvenir collectors, bounty from digs and ancient sites in impoverished but art-rich source countries. Patrimony laws were intended to protect future archaeological discoveries against Western imperialist designs. . . .
I surveyed 90 countries with one or more archaeological sites on UNESCO’s World Heritage Site list, and my study shows that in most cases the number of discovered sites diminishes sharply after a country passes a cultural property law. There are 222 archaeological sites listed for those 90 countries. When you look into the history of the sites, you see that all but 21 were discovered before the passage of cultural property laws. . . .
Strict cultural patrimony laws are popular in most countries. But the downside may be that they reduce incentives for foreign governments, nongovernmental organizations and educational institutions to invest in overseas exploration because their efforts will not necessarily be rewarded by opportunities to hold, display and study what is uncovered. To the extent that source countries can fund their own archaeological projects, artifacts and sites may still be discovered. . . . The survey has far-reaching implications. It suggests that source countries, particularly in the developing world, should narrow their cultural property laws so that they can reap the benefits of new archaeological discoveries, which typically increase tourism and enhance cultural pride. This does not mean these nations should abolish restrictions on foreign excavation and foreign claims to artifacts.
China provides an interesting alternative approach for source nations eager for foreign archaeological investment. From 1935 to 2003, China had a restrictive cultural property law that prohibited foreign ownership of Chinese cultural artifacts. In those years, China’s most significant archaeological discovery occurred by chance, in 1974, when peasant farmers accidentally uncovered ranks of buried terra cotta warriors, which are part of Emperor Qin’s spectacular tomb system.
In 2003, the Chinese government switched course, dropping its cultural property law and embracing collaborative international archaeological research. Since then, China has nominated 11 archaeological sites for inclusion in the World Heritage Site list, including eight in 2013, the most ever for China.

\vspace{1em}
\textbf{Question 18}

**There is a sentence that is missing in the paragraph below. Look at the paragraph and decide where (option 1, 2, 3, or 4) the following sentence would best fit.****  
**
**Sentence:** For theoretical purposes, arguments may be considered as freestanding entities, abstracted from their contexts of use in actual human activities.
**Paragraph** : ___(1)___. An argument can be defined as a complex symbolic structure where some parts, known as the premises, offer support to another part, the conclusion. Alternatively, an argument can be viewed as a complex speech act consisting of one or more acts of premising (which assert propositions in favor of the conclusion), an act of concluding, and a stated or implicit marker (“hence”, “therefore”) that indicates that the conclusion follows from the premises.___(2)___. The relation of support between premises and conclusion can be cashed out in different ways: the premises may guarantee the truth of the conclusion, or make its truth more probable; the premises may imply the conclusion; the premises may make the conclusion more acceptable (or assertible).___(3)___. But depending on one’s explanatory goals, there is also much to be gained from considering arguments as they in fact occur in human communicative practices.___(4)___.

\begin{enumerate}[label=\Alph*.]
    \item Option 4
    \item Option 2
    \item Option 1
    \item Option 3
[](https://cracku.in/18-there-is-a-sentence-that-is-missing-in-the-paragra-x-cat-2023-slot-3)
\end{enumerate}
\vspace{1em}
\textbf{Solution:}

The sentence is most fitting in Blank 3 because it follows the idea presented in the previous sentences. Before Blank 3 the passage discusses the various ways to understand the relation of support between premises and conclusions, presenting different perspectives on how arguments can be analysed or interpreted. The sentence in question adds another layer to this discussion by suggesting that, for theoretical purposes, arguments can be considered independently of their real-world contexts. This placement enhances the overall flow of the paragraph by introducing a theoretical perspective on the nature of arguments.

\end{problem}

\newpage
\begin{problem}{}{}
% Subject: varc
\textbf{Instructions}

Instructions   

The passage below is accompanied by four questions. Based on the passage, choose the best answer for each question.
In 2006, the Met [art museum in the US] agreed to return the Euphronios krater, a masterpiece Greek urn that had been a museum draw since 1972. In 2007, the Getty [art museum in the US] agreed to return 40 objects to Italy, including a marble Aphrodite, in the midst of looting scandals. And in December, Sotheby’s and a private owner agreed to return an ancient Khmer statue of a warrior, pulled from auction two years before, to Cambodia.
Cultural property, or patrimony, laws limit the transfer of cultural property outside the source country’s territory, including outright export prohibitions and national ownership laws. Most art historians, archaeologists, museum officials and policymakers portray cultural property laws in general as invaluable tools for counteracting the ugly legacy of Western cultural imperialism.
During the late 19th and early 20th century — an era former Met director Thomas Hoving called “the age of piracy” — American and European art museums acquired antiquities by hook or by crook, from grave robbers or souvenir collectors, bounty from digs and ancient sites in impoverished but art-rich source countries. Patrimony laws were intended to protect future archaeological discoveries against Western imperialist designs. . . .
I surveyed 90 countries with one or more archaeological sites on UNESCO’s World Heritage Site list, and my study shows that in most cases the number of discovered sites diminishes sharply after a country passes a cultural property law. There are 222 archaeological sites listed for those 90 countries. When you look into the history of the sites, you see that all but 21 were discovered before the passage of cultural property laws. . . .
Strict cultural patrimony laws are popular in most countries. But the downside may be that they reduce incentives for foreign governments, nongovernmental organizations and educational institutions to invest in overseas exploration because their efforts will not necessarily be rewarded by opportunities to hold, display and study what is uncovered. To the extent that source countries can fund their own archaeological projects, artifacts and sites may still be discovered. . . . The survey has far-reaching implications. It suggests that source countries, particularly in the developing world, should narrow their cultural property laws so that they can reap the benefits of new archaeological discoveries, which typically increase tourism and enhance cultural pride. This does not mean these nations should abolish restrictions on foreign excavation and foreign claims to artifacts.
China provides an interesting alternative approach for source nations eager for foreign archaeological investment. From 1935 to 2003, China had a restrictive cultural property law that prohibited foreign ownership of Chinese cultural artifacts. In those years, China’s most significant archaeological discovery occurred by chance, in 1974, when peasant farmers accidentally uncovered ranks of buried terra cotta warriors, which are part of Emperor Qin’s spectacular tomb system.
In 2003, the Chinese government switched course, dropping its cultural property law and embracing collaborative international archaeological research. Since then, China has nominated 11 archaeological sites for inclusion in the World Heritage Site list, including eight in 2013, the most ever for China.

\vspace{1em}
\textbf{Question 19}

Five jumbled up sentences (labelled 1, 2, 3, 4 and 5), related to a topic, are given below. Four of them can be put together to form a coherent paragraph. Identify the odd sentence and key in the number of that sentence as your answer.
1. Boa Senior, who lived through the 2004 tsunami, the Japanese occupation and diseases brought by British settlers, was the last native of the island chain who was fluent in Bo.  
2. The indigenous population has been steadily collapsing since the island chain was colonised by British settlers in 1858 and used for most of the following 100 years as a colonial penal colony.  
3. Taking its name from a now-extinct tribe, Bo is one of the 10 Great Andamanese languages, which are thought to date back to pre-Neolithic human settlement of south-east Asia.  
4. The last speaker of an ancient tribal language has died in the Andaman Islands, breaking a 65,000-year link to one of the world's oldest cultures.  
5. Though the language has been closely studied by researchers of linguistic history, Boa Senior spent the last few years of her life unable to converse with anyone in her mother tongue.
Backspace  
789  
456  
123  
0.-  
Clear All  

Submit
[](https://cracku.in/19-five-jumbled-up-sentences-labelled-1-2-3-4-and-5-r-x-cat-2023-slot-3)

\begin{enumerate}[label=\Alph*.]
\end{enumerate}
\vspace{1em}
\textbf{Solution:}

Sentence 2 doesn't fit well with the others, which focus on the language Bo and its last native speaker, Boa Senior. The other sentences provide information about the language, its history, and the last fluent speaker, creating a coherent narrative. Sentence 2 introduces a different topic about the decline of the indigenous population without directly contributing to the discussion about the tribal language and its last speaker.

\end{problem}

\newpage
\begin{problem}{}{}
% Subject: varc
\textbf{Instructions}

Instructions   

The passage below is accompanied by four questions. Based on the passage, choose the best answer for each question.
In 2006, the Met [art museum in the US] agreed to return the Euphronios krater, a masterpiece Greek urn that had been a museum draw since 1972. In 2007, the Getty [art museum in the US] agreed to return 40 objects to Italy, including a marble Aphrodite, in the midst of looting scandals. And in December, Sotheby’s and a private owner agreed to return an ancient Khmer statue of a warrior, pulled from auction two years before, to Cambodia.
Cultural property, or patrimony, laws limit the transfer of cultural property outside the source country’s territory, including outright export prohibitions and national ownership laws. Most art historians, archaeologists, museum officials and policymakers portray cultural property laws in general as invaluable tools for counteracting the ugly legacy of Western cultural imperialism.
During the late 19th and early 20th century — an era former Met director Thomas Hoving called “the age of piracy” — American and European art museums acquired antiquities by hook or by crook, from grave robbers or souvenir collectors, bounty from digs and ancient sites in impoverished but art-rich source countries. Patrimony laws were intended to protect future archaeological discoveries against Western imperialist designs. . . .
I surveyed 90 countries with one or more archaeological sites on UNESCO’s World Heritage Site list, and my study shows that in most cases the number of discovered sites diminishes sharply after a country passes a cultural property law. There are 222 archaeological sites listed for those 90 countries. When you look into the history of the sites, you see that all but 21 were discovered before the passage of cultural property laws. . . .
Strict cultural patrimony laws are popular in most countries. But the downside may be that they reduce incentives for foreign governments, nongovernmental organizations and educational institutions to invest in overseas exploration because their efforts will not necessarily be rewarded by opportunities to hold, display and study what is uncovered. To the extent that source countries can fund their own archaeological projects, artifacts and sites may still be discovered. . . . The survey has far-reaching implications. It suggests that source countries, particularly in the developing world, should narrow their cultural property laws so that they can reap the benefits of new archaeological discoveries, which typically increase tourism and enhance cultural pride. This does not mean these nations should abolish restrictions on foreign excavation and foreign claims to artifacts.
China provides an interesting alternative approach for source nations eager for foreign archaeological investment. From 1935 to 2003, China had a restrictive cultural property law that prohibited foreign ownership of Chinese cultural artifacts. In those years, China’s most significant archaeological discovery occurred by chance, in 1974, when peasant farmers accidentally uncovered ranks of buried terra cotta warriors, which are part of Emperor Qin’s spectacular tomb system.
In 2003, the Chinese government switched course, dropping its cultural property law and embracing collaborative international archaeological research. Since then, China has nominated 11 archaeological sites for inclusion in the World Heritage Site list, including eight in 2013, the most ever for China.

\vspace{1em}
\textbf{Question 20}

Five jumbled up sentences (labelled 1, 2, 3, 4 and 5), related to a topic, are given below. Four of them can be put together to form a coherent paragraph. Identify the odd sentence and key in the number of that sentence as your answer.
1. Although hard skills have traditionally ruled the roost, some companies are moving away from choosing prospective hires based on technical abilities alone.  
2. Companies are shaking off the old definition of an ideal candidate and ditching the idea of looking for the singularly perfect candidate altogether.  
3. Now, some job descriptions are frequently asking for candidates to demonstrate soft skills, such as leadership or teamwork.  
4. That’s not to say that practical know-how is no longer required - some jobs still call for highly specific expertise  
5. The move towards prioritising soft skills “is a natural response to three years of the pandemic” says a senior recruiter at Cenlar FSB.
Backspace  
789  
456  
123  
0.-  
Clear All  

Submit
[](https://cracku.in/20-five-jumbled-up-sentences-labelled-1-2-3-4-and-5-r-x-cat-2023-slot-3)

\begin{enumerate}[label=\Alph*.]
\end{enumerate}
\vspace{1em}
\textbf{Solution:}

Sentence 2 is the odd one out because it does not directly contribute to the theme of a shift in hiring priorities towards soft skills. Sentences 1, 3, 4, and 5 all discuss changes in hiring practices, the importance of soft skills, and the response to the pandemic, creating a coherent narrative. In contrast, Sentence 2 introduces a more general idea about redefining the ideal candidate without specifically addressing the shift towards prioritizing soft skills, making it the odd one out in the context of the paragraph's focus.

\end{problem}

\newpage
\begin{problem}{}{}
% Subject: varc
\textbf{Instructions}

Instructions   

The passage below is accompanied by four questions. Based on the passage, choose the best answer for each question.
In 2006, the Met [art museum in the US] agreed to return the Euphronios krater, a masterpiece Greek urn that had been a museum draw since 1972. In 2007, the Getty [art museum in the US] agreed to return 40 objects to Italy, including a marble Aphrodite, in the midst of looting scandals. And in December, Sotheby’s and a private owner agreed to return an ancient Khmer statue of a warrior, pulled from auction two years before, to Cambodia.
Cultural property, or patrimony, laws limit the transfer of cultural property outside the source country’s territory, including outright export prohibitions and national ownership laws. Most art historians, archaeologists, museum officials and policymakers portray cultural property laws in general as invaluable tools for counteracting the ugly legacy of Western cultural imperialism.
During the late 19th and early 20th century — an era former Met director Thomas Hoving called “the age of piracy” — American and European art museums acquired antiquities by hook or by crook, from grave robbers or souvenir collectors, bounty from digs and ancient sites in impoverished but art-rich source countries. Patrimony laws were intended to protect future archaeological discoveries against Western imperialist designs. . . .
I surveyed 90 countries with one or more archaeological sites on UNESCO’s World Heritage Site list, and my study shows that in most cases the number of discovered sites diminishes sharply after a country passes a cultural property law. There are 222 archaeological sites listed for those 90 countries. When you look into the history of the sites, you see that all but 21 were discovered before the passage of cultural property laws. . . .
Strict cultural patrimony laws are popular in most countries. But the downside may be that they reduce incentives for foreign governments, nongovernmental organizations and educational institutions to invest in overseas exploration because their efforts will not necessarily be rewarded by opportunities to hold, display and study what is uncovered. To the extent that source countries can fund their own archaeological projects, artifacts and sites may still be discovered. . . . The survey has far-reaching implications. It suggests that source countries, particularly in the developing world, should narrow their cultural property laws so that they can reap the benefits of new archaeological discoveries, which typically increase tourism and enhance cultural pride. This does not mean these nations should abolish restrictions on foreign excavation and foreign claims to artifacts.
China provides an interesting alternative approach for source nations eager for foreign archaeological investment. From 1935 to 2003, China had a restrictive cultural property law that prohibited foreign ownership of Chinese cultural artifacts. In those years, China’s most significant archaeological discovery occurred by chance, in 1974, when peasant farmers accidentally uncovered ranks of buried terra cotta warriors, which are part of Emperor Qin’s spectacular tomb system.
In 2003, the Chinese government switched course, dropping its cultural property law and embracing collaborative international archaeological research. Since then, China has nominated 11 archaeological sites for inclusion in the World Heritage Site list, including eight in 2013, the most ever for China.

\vspace{1em}
\textbf{Question 21}

The four sentences (labelled 1, 2, 3 and 4) given below, when properly sequenced, would yield a coherent paragraph. Decide on the proper sequencing of the order of the sentences and key in the sequence of the four numbers as your answer.
1. Centuries later formal learning is still mostly based on reading, even with the widespread use of other possible education-affecting technologies such as film, radio, and television.  
2. One of the immediate and recognisable impacts of the printing press was on how people learned; in the scribal culture it primarily involved listening, so memorization was paramount.  
3. The transformation of learners from listeners to readers was a complex social and cultural phenomenon, and it was not until the industrial era that the concept of universal literacy took root.  
4. The printing press shifted the learning process, as listening and memorisation gradually gave way to reading and learning no longer required the presence of a mentor; it could be done privately.
Backspace  
789  
456  
123  
0.-  
Clear All  

Submit
[](https://cracku.in/21-the-four-sentences-labelled-1-2-3-and-4-given-belo-x-cat-2023-slot-3)

\begin{enumerate}[label=\Alph*.]
\end{enumerate}
\vspace{1em}
\textbf{Solution:}

The correct order is 2-4-3-1
Sentence 2 introduces the immediate impact of the printing press on how people learned in the scribal culture, emphasizing the shift from listening and memorization to a new learning process.
Sentence 4 builds on this by explaining the transformation brought about by the printing press, highlighting the shift from listening and memorization to reading and the newfound ability to learn privately.
Sentence 3 provides additional context, emphasizing that the transformation from listeners to readers was a complex social and cultural phenomenon that became more widespread during the industrial era.
Sentence 1 concludes the paragraph by noting that despite technological advancements such as film, radio, and television, centuries later, formal learning is still predominantly based on reading. This brings the historical context to the present, concluding the narrative cohesively.

\end{problem}

\newpage
\begin{problem}{}{}
% Subject: varc
\textbf{Instructions}

Instructions   

The passage below is accompanied by four questions. Based on the passage, choose the best answer for each question.
In 2006, the Met [art museum in the US] agreed to return the Euphronios krater, a masterpiece Greek urn that had been a museum draw since 1972. In 2007, the Getty [art museum in the US] agreed to return 40 objects to Italy, including a marble Aphrodite, in the midst of looting scandals. And in December, Sotheby’s and a private owner agreed to return an ancient Khmer statue of a warrior, pulled from auction two years before, to Cambodia.
Cultural property, or patrimony, laws limit the transfer of cultural property outside the source country’s territory, including outright export prohibitions and national ownership laws. Most art historians, archaeologists, museum officials and policymakers portray cultural property laws in general as invaluable tools for counteracting the ugly legacy of Western cultural imperialism.
During the late 19th and early 20th century — an era former Met director Thomas Hoving called “the age of piracy” — American and European art museums acquired antiquities by hook or by crook, from grave robbers or souvenir collectors, bounty from digs and ancient sites in impoverished but art-rich source countries. Patrimony laws were intended to protect future archaeological discoveries against Western imperialist designs. . . .
I surveyed 90 countries with one or more archaeological sites on UNESCO’s World Heritage Site list, and my study shows that in most cases the number of discovered sites diminishes sharply after a country passes a cultural property law. There are 222 archaeological sites listed for those 90 countries. When you look into the history of the sites, you see that all but 21 were discovered before the passage of cultural property laws. . . .
Strict cultural patrimony laws are popular in most countries. But the downside may be that they reduce incentives for foreign governments, nongovernmental organizations and educational institutions to invest in overseas exploration because their efforts will not necessarily be rewarded by opportunities to hold, display and study what is uncovered. To the extent that source countries can fund their own archaeological projects, artifacts and sites may still be discovered. . . . The survey has far-reaching implications. It suggests that source countries, particularly in the developing world, should narrow their cultural property laws so that they can reap the benefits of new archaeological discoveries, which typically increase tourism and enhance cultural pride. This does not mean these nations should abolish restrictions on foreign excavation and foreign claims to artifacts.
China provides an interesting alternative approach for source nations eager for foreign archaeological investment. From 1935 to 2003, China had a restrictive cultural property law that prohibited foreign ownership of Chinese cultural artifacts. In those years, China’s most significant archaeological discovery occurred by chance, in 1974, when peasant farmers accidentally uncovered ranks of buried terra cotta warriors, which are part of Emperor Qin’s spectacular tomb system.
In 2003, the Chinese government switched course, dropping its cultural property law and embracing collaborative international archaeological research. Since then, China has nominated 11 archaeological sites for inclusion in the World Heritage Site list, including eight in 2013, the most ever for China.

\vspace{1em}
\textbf{Question 22}

The four sentences (labelled 1, 2, 3 and 4) given below, when properly sequenced, would yield a coherent paragraph. Decide on the proper sequencing of the order of the sentences and key in the sequence of the four numbers as your answer.
1. Veena Sahajwalla, a materials scientist at the University of New South Wales, believes there is a new way of solving this problem.  
2. Her vision is for automated drones and robots to pick out components, put them into a small furnace and smelt them at specific temperatures to extract the metals one by one before they are sent off to manufacturers for reuse.  
3. E-waste contains huge quantities of valuable metals, ceramics and plastics that could be salvaged and recycled, although currently not enough of it is.  
4. She plans to build microfactories that can tease apart the tangle of materials in mobile phones, computers and other e-waste.
Backspace  
789  
456  
123  
0.-  
Clear All  

Submit
[](https://cracku.in/22-the-four-sentences-labelled-1-2-3-and-4-given-belo-x-cat-2023-slot-3)

\begin{enumerate}[label=\Alph*.]
\end{enumerate}
\vspace{1em}
\textbf{Solution:}

Sentence 3 introduces the issue, highlighting that e-waste contains valuable materials that are not being efficiently recycled. This sets context for the discussion to be followed. 
Now if we consider Sentences 1, 2 and 4, we can see that Sentence 2 and 4 are referencing a “She” and “Her”. Considering these 3 sentences, they can only reference Veena Sahajwalla. Therefore Sentence 1 must follow Sentence 3.
Now Sentence 1 states that she has a new way of solving the problem. Sentence 4 provides details about Sahajwalla's plan, stating that she intends to build microfactories to extract valuable materials from e-waste. Therefore Sentence 4 must be following Sentence 1. Sentence 2 logically follows Sentence 4 by elaborating on Sahajwalla's vision, explaining how automated drones and robots will be utilized in the process of extracting metals from e-waste, providing a clear and detailed picture of the proposed solution.
Therefore the correct order is 3-1-4-2.

\end{problem}

\newpage
\begin{problem}{}{}
% Subject: varc
\textbf{Instructions}

Instructions   

The passage below is accompanied by four questions. Based on the passage, choose the best answer for each question.
In 2006, the Met [art museum in the US] agreed to return the Euphronios krater, a masterpiece Greek urn that had been a museum draw since 1972. In 2007, the Getty [art museum in the US] agreed to return 40 objects to Italy, including a marble Aphrodite, in the midst of looting scandals. And in December, Sotheby’s and a private owner agreed to return an ancient Khmer statue of a warrior, pulled from auction two years before, to Cambodia.
Cultural property, or patrimony, laws limit the transfer of cultural property outside the source country’s territory, including outright export prohibitions and national ownership laws. Most art historians, archaeologists, museum officials and policymakers portray cultural property laws in general as invaluable tools for counteracting the ugly legacy of Western cultural imperialism.
During the late 19th and early 20th century — an era former Met director Thomas Hoving called “the age of piracy” — American and European art museums acquired antiquities by hook or by crook, from grave robbers or souvenir collectors, bounty from digs and ancient sites in impoverished but art-rich source countries. Patrimony laws were intended to protect future archaeological discoveries against Western imperialist designs. . . .
I surveyed 90 countries with one or more archaeological sites on UNESCO’s World Heritage Site list, and my study shows that in most cases the number of discovered sites diminishes sharply after a country passes a cultural property law. There are 222 archaeological sites listed for those 90 countries. When you look into the history of the sites, you see that all but 21 were discovered before the passage of cultural property laws. . . .
Strict cultural patrimony laws are popular in most countries. But the downside may be that they reduce incentives for foreign governments, nongovernmental organizations and educational institutions to invest in overseas exploration because their efforts will not necessarily be rewarded by opportunities to hold, display and study what is uncovered. To the extent that source countries can fund their own archaeological projects, artifacts and sites may still be discovered. . . . The survey has far-reaching implications. It suggests that source countries, particularly in the developing world, should narrow their cultural property laws so that they can reap the benefits of new archaeological discoveries, which typically increase tourism and enhance cultural pride. This does not mean these nations should abolish restrictions on foreign excavation and foreign claims to artifacts.
China provides an interesting alternative approach for source nations eager for foreign archaeological investment. From 1935 to 2003, China had a restrictive cultural property law that prohibited foreign ownership of Chinese cultural artifacts. In those years, China’s most significant archaeological discovery occurred by chance, in 1974, when peasant farmers accidentally uncovered ranks of buried terra cotta warriors, which are part of Emperor Qin’s spectacular tomb system.
In 2003, the Chinese government switched course, dropping its cultural property law and embracing collaborative international archaeological research. Since then, China has nominated 11 archaeological sites for inclusion in the World Heritage Site list, including eight in 2013, the most ever for China.

\vspace{1em}
\textbf{Question 23}

The passage given below is followed by four alternate summaries. Choose the option that best captures the essence of the passage.
The weight of society’s expectations is hardly a new phenomenon but it has become particularly draining over recent decades, perhaps because expectations themselves are so multifarious and contradictory. The perfectionism of the 1950s was rooted in the norms of mass culture and captured in famous advertising images of the ideal white American family that now seem self-satirising. In that era, perfectionism meant seamlessly conforming to values, behaviour and appearance: chiselled confidence for men, demure graciousness for women. The perfectionist was under pressure to look like everyone else, only more so. The perfectionists of today, by contrast, feel an obligation to stand out through their idiosyncratic style and wit if they are to gain a foothold in the attention economy.

\begin{enumerate}[label=\Alph*.]
    \item The image of perfectionism is reflected in and perpetuated by the media; and people do their best to adhere to these ideals.
    \item Though long-standing, the pressure to appear perfect and thereby attract attention, has evolved over time from one of conformism to one of non-conformism.
    \item The pressure to appear perfect has been the cause of tension and conflict because the idea itself has been in a state of flux and hard to define.
    \item The desire to attract attention is so deep-rooted in individual consciousness that people are willing to go to any lengths to achieve it.
[](https://cracku.in/23-the-passage-given-below-is-followed-by-four-altern-x-cat-2023-slot-3)
\end{enumerate}
\vspace{1em}
\textbf{Solution:}

The passage contrasts the 1950s, where perfectionism meant conforming to societal norms, with contemporary times, where individuals feel pressured to stand out and gain attention through unique style and wit. This evolution from conformism to non-conformism is a key point in the passage, making Option B the most accurate summary. It effectively captures the changing nature of perfectionism in response to multifarious and contradictory societal expectations over the decades.
Option A focuses on the media's role and people adhering to ideals, which is not the primary emphasis of the passage that highlights the evolution of societal expectations over time.
Option C suggests tension and conflict related to the changing idea of perfection, but the passage emphasises the historical shift in expectations rather than conflict.
Option D overgeneralizes by stating that people are willing to go to any lengths to attract attention, which is not explicitly supported by the passage that highlights the changing nature of perfectionism.

\end{problem}

\newpage
\begin{problem}{}{}
% Subject: varc
\textbf{Instructions}

Instructions   

The passage below is accompanied by four questions. Based on the passage, choose the best answer for each question.
In 2006, the Met [art museum in the US] agreed to return the Euphronios krater, a masterpiece Greek urn that had been a museum draw since 1972. In 2007, the Getty [art museum in the US] agreed to return 40 objects to Italy, including a marble Aphrodite, in the midst of looting scandals. And in December, Sotheby’s and a private owner agreed to return an ancient Khmer statue of a warrior, pulled from auction two years before, to Cambodia.
Cultural property, or patrimony, laws limit the transfer of cultural property outside the source country’s territory, including outright export prohibitions and national ownership laws. Most art historians, archaeologists, museum officials and policymakers portray cultural property laws in general as invaluable tools for counteracting the ugly legacy of Western cultural imperialism.
During the late 19th and early 20th century — an era former Met director Thomas Hoving called “the age of piracy” — American and European art museums acquired antiquities by hook or by crook, from grave robbers or souvenir collectors, bounty from digs and ancient sites in impoverished but art-rich source countries. Patrimony laws were intended to protect future archaeological discoveries against Western imperialist designs. . . .
I surveyed 90 countries with one or more archaeological sites on UNESCO’s World Heritage Site list, and my study shows that in most cases the number of discovered sites diminishes sharply after a country passes a cultural property law. There are 222 archaeological sites listed for those 90 countries. When you look into the history of the sites, you see that all but 21 were discovered before the passage of cultural property laws. . . .
Strict cultural patrimony laws are popular in most countries. But the downside may be that they reduce incentives for foreign governments, nongovernmental organizations and educational institutions to invest in overseas exploration because their efforts will not necessarily be rewarded by opportunities to hold, display and study what is uncovered. To the extent that source countries can fund their own archaeological projects, artifacts and sites may still be discovered. . . . The survey has far-reaching implications. It suggests that source countries, particularly in the developing world, should narrow their cultural property laws so that they can reap the benefits of new archaeological discoveries, which typically increase tourism and enhance cultural pride. This does not mean these nations should abolish restrictions on foreign excavation and foreign claims to artifacts.
China provides an interesting alternative approach for source nations eager for foreign archaeological investment. From 1935 to 2003, China had a restrictive cultural property law that prohibited foreign ownership of Chinese cultural artifacts. In those years, China’s most significant archaeological discovery occurred by chance, in 1974, when peasant farmers accidentally uncovered ranks of buried terra cotta warriors, which are part of Emperor Qin’s spectacular tomb system.
In 2003, the Chinese government switched course, dropping its cultural property law and embracing collaborative international archaeological research. Since then, China has nominated 11 archaeological sites for inclusion in the World Heritage Site list, including eight in 2013, the most ever for China.

\vspace{1em}
\textbf{Question 24}

The passage given below is followed by four alternate summaries. Choose the option that best captures the essence of the passage.
Gradually, life for the island’s birds is improving. Antarctic prions and white-headed petrels, which also nest in burrows, had managed to cling on in some sites while pests were on the island. Their numbers are now increasing. “It’s fantastic and so exciting,” Shaw says. As birds return to breed, they also poo. This adds nutrients to the soil, which in turn helps the plants to grow back stronger. Tall plants then help burrowing birds hide from predatory skuas. “It’s this wonderful feedback loop,” Shaw says. Today, the “pretty paddock” that Houghton first experienced has been transformed. “The tussock is over your head, and you’re dodging all these penguin tunnels,” she says. The orchids and tiny herb that had been protected by fencing have started turning up all over the place.

\begin{enumerate}[label=\Alph*.]
    \item There is an increasing number of predatory birds and plants on the island despite the presence of pests which is a positive development.
    \item In the absence of pests, life on the island is now protected, and there has been a revival of a variety of birds and plants.
    \item Flowering plants, herbs and birds are now being protected on this wonderful Antarctic island.
    \item There is a huge positive transformation of the ecosystem of the island when brought under environmental protection.
[](https://cracku.in/24-the-passage-given-below-is-followed-by-four-altern-x-cat-2023-slot-3)
\end{enumerate}
\vspace{1em}
\textbf{Solution:}

The passage describes how life on the island is gradually improving, particularly for birds like Antarctic prions and white-headed petrels. It highlights that these birds are increasing in numbers as pests are controlled on the island. The absence of pests allows the birds to return, breed, and contribute to the ecosystem positively. The passage also mentions how bird droppings add nutrients to the soil, supporting plant growth. Overall, Option B effectively conveys the central theme of the passage - the revival and improvement of life on the island due to the absence of pests and the positive impact on birds and plants.
Option A incorrectly suggests an increase in the number of predatory birds, which contradicts the positive developments mentioned in the passage.  
Option C is incorrect as the passage doesn't explicitly state that this protection is the primary cause of the positive changes.
Option D is more general and doesn't specifically address the absence of pests as a crucial factor in the positive transformation mentioned in the passage.
[](https://cracku.in/downloads/18892)
  *  
  *

\end{problem}

\newpage
\begin{problem}{}{}
% Subject: varc
\textbf{Instructions}

Instructions   

**The passage below is accompanied by four questions. Based on the passage, choose the best answer for each question.**
. . . [T]he idea of craftsmanship is not simply nostalgic. . . . Crafts require distinct skills, an allround approach to work that involves the whole product, rather than individual parts, and an attitude that necessitates devotion to the job and a focus on the communal interest. The concept of craft emphasises the human touch and individual judgment.
Essentially, the crafts concept seems to run against the preponderant ethos of management studies which, as the academics note, have long prioritised efficiency and consistency. . . . Craft skills were portrayed as being primitive and traditionalist.
The contrast between artisanship and efficiency first came to the fore in the 19th century when British manufacturers suddenly faced competition from across the Atlantic as firms developed the “American system” using standardised parts. . . . the worldwide success of the Singer sewing machine showed the potential of a mass-produced device. This process created its own reaction, first in the form of the Arts and Crafts movement of the late 19th century, and then again in the “small is beautiful” movement of the 1970s. A third crafts movement is emerging as people become aware of the environmental impact of conventional industry.
There are two potential markets for those who practise crafts. The first stems from the existence of consumers who are willing to pay a premium price for goods that are deemed to be of extra quality. . . . The second market lies in those consumers who wish to use their purchases to support local workers, or to reduce their environmental impact by taking goods to craftspeople to be mended, or recycled.
For workers, the appeal of craftsmanship is that it allows them the autonomy to make creative choices, and thus makes a job far more satisfying. In that sense, it could offer hope for the overall labour market. Let the machines automate dull and repetitive tasks and let workers focus purely on their skills, judgment and imagination. As a current example, the academics cite the “agile” manifesto in the software sector, an industry at the heart of technological change. The pioneers behind the original agile manifesto promised to prioritise “individuals and interactions over processes and tools”. By bringing together experts from different teams, agile working is designed to improve creativity.
But the broader question is whether crafts can create a lot more jobs than they do today. Demand for crafted products may rise but will it be easy to retrain workers in sectors that might get automated (such as truck drivers) to take advantage? In a world where products and services often have to pass through regulatory hoops, large companies will usually have the advantage.
History also suggests that the link between crafts and creativity is not automatic. Medieval craft guilds were monopolies which resisted new entrants. They were also highly hierarchical with young men required to spend long periods as apprentices and journeymen before they could set up on their own; by that time the innovative spirit may have been knocked out of them. Craft workers can thrive in the modern era, but only if they don’t get too organised.

\vspace{1em}
\textbf{Question 1}

We can infer from the passage that medieval crafts guilds resembled mass production in that both

\begin{enumerate}[label=\Alph*.]
    \item did not necessarily promote creativity.
    \item discouraged innovation by restricting entry through strict rules
    \item did not always employ egalitarian production processes.
    \item focused excessively on product quality
[](https://cracku.in/1-we-can-infer-from-the-passage-that-medieval-crafts-x-cat-2024-slot-1)
\end{enumerate}
\vspace{1em}
\textbf{Solution:}

Option A is the correct answer.   
  
The passage mentions that medieval craft guilds were monopolies that resisted new entrants and required long apprenticeship periods, which could stifle innovative spirit. Similarly, the ethos of mass production prioritizes efficiency and consistency, which often comes at the expense of creativity. Therefore, Option A could be inferred.
Option B: While it is true that medieval guilds restricted entry through strict rules, the passage does not suggest that mass production involves such restrictions. 
Option C: The passage does not discuss whether mass production or medieval guilds employed egalitarian processes. 
Option D: The passage does not mention that medieval guilds or mass production focused excessively on product quality.

\end{problem}

\newpage
\begin{problem}{}{}
% Subject: varc
\textbf{Instructions}

Instructions   

**The passage below is accompanied by four questions. Based on the passage, choose the best answer for each question.**
. . . [T]he idea of craftsmanship is not simply nostalgic. . . . Crafts require distinct skills, an allround approach to work that involves the whole product, rather than individual parts, and an attitude that necessitates devotion to the job and a focus on the communal interest. The concept of craft emphasises the human touch and individual judgment.
Essentially, the crafts concept seems to run against the preponderant ethos of management studies which, as the academics note, have long prioritised efficiency and consistency. . . . Craft skills were portrayed as being primitive and traditionalist.
The contrast between artisanship and efficiency first came to the fore in the 19th century when British manufacturers suddenly faced competition from across the Atlantic as firms developed the “American system” using standardised parts. . . . the worldwide success of the Singer sewing machine showed the potential of a mass-produced device. This process created its own reaction, first in the form of the Arts and Crafts movement of the late 19th century, and then again in the “small is beautiful” movement of the 1970s. A third crafts movement is emerging as people become aware of the environmental impact of conventional industry.
There are two potential markets for those who practise crafts. The first stems from the existence of consumers who are willing to pay a premium price for goods that are deemed to be of extra quality. . . . The second market lies in those consumers who wish to use their purchases to support local workers, or to reduce their environmental impact by taking goods to craftspeople to be mended, or recycled.
For workers, the appeal of craftsmanship is that it allows them the autonomy to make creative choices, and thus makes a job far more satisfying. In that sense, it could offer hope for the overall labour market. Let the machines automate dull and repetitive tasks and let workers focus purely on their skills, judgment and imagination. As a current example, the academics cite the “agile” manifesto in the software sector, an industry at the heart of technological change. The pioneers behind the original agile manifesto promised to prioritise “individuals and interactions over processes and tools”. By bringing together experts from different teams, agile working is designed to improve creativity.
But the broader question is whether crafts can create a lot more jobs than they do today. Demand for crafted products may rise but will it be easy to retrain workers in sectors that might get automated (such as truck drivers) to take advantage? In a world where products and services often have to pass through regulatory hoops, large companies will usually have the advantage.
History also suggests that the link between crafts and creativity is not automatic. Medieval craft guilds were monopolies which resisted new entrants. They were also highly hierarchical with young men required to spend long periods as apprentices and journeymen before they could set up on their own; by that time the innovative spirit may have been knocked out of them. Craft workers can thrive in the modern era, but only if they don’t get too organised.

\vspace{1em}
\textbf{Question 2}

Which one of the following statements is NOT inconsistent with the views stated in the passage?

\begin{enumerate}[label=\Alph*.]
    \item We need to support the crafts; only then can we retain the creativity intrinsic to their production.
    \item Creativity in the crafts could be stifled if the market for artisan goods becomes too organised.
    \item The Arts and Crafts movement was initially inspired by the “American system” of production.
    \item The agile movement in software is a throwback to the tenets of the medieval crafts guilds
[](https://cracku.in/2-which-one-of-the-following-statements-is-not-incon-x-cat-2024-slot-1)
\end{enumerate}
\vspace{1em}
\textbf{Solution:}

Option B is the correct answer.   
  
The passage mentions that medieval craft guilds were monopolistic and hierarchical, resisting new entrants and imposing long apprenticeships, which could stifle innovation. It also warns that modern craft workers “can thrive... only if they don’t get too organised,” supporting option B.
Option A: This is an extreme interpretation. The passage does not argue that supporting crafts is the only way to retain the creativity intrinsic to their production.
Option C: The Arts and Crafts movement is described as a reaction against the “American system” and the rise of mass production, not an inspiration drawn from it.
Option D: Agile movement is praised for prioritizing creativity and collaboration, whereas medieval craft guilds are described as hierarchical and restrictive, which stifled innovation. These are contrasting ideas. Therefore, agile movement cannot be a throwback to the tenets (principles) of medieval craft guilds)

\end{problem}

\newpage
\begin{problem}{}{}
% Subject: varc
\textbf{Instructions}

Instructions   

**The passage below is accompanied by four questions. Based on the passage, choose the best answer for each question.**
. . . [T]he idea of craftsmanship is not simply nostalgic. . . . Crafts require distinct skills, an allround approach to work that involves the whole product, rather than individual parts, and an attitude that necessitates devotion to the job and a focus on the communal interest. The concept of craft emphasises the human touch and individual judgment.
Essentially, the crafts concept seems to run against the preponderant ethos of management studies which, as the academics note, have long prioritised efficiency and consistency. . . . Craft skills were portrayed as being primitive and traditionalist.
The contrast between artisanship and efficiency first came to the fore in the 19th century when British manufacturers suddenly faced competition from across the Atlantic as firms developed the “American system” using standardised parts. . . . the worldwide success of the Singer sewing machine showed the potential of a mass-produced device. This process created its own reaction, first in the form of the Arts and Crafts movement of the late 19th century, and then again in the “small is beautiful” movement of the 1970s. A third crafts movement is emerging as people become aware of the environmental impact of conventional industry.
There are two potential markets for those who practise crafts. The first stems from the existence of consumers who are willing to pay a premium price for goods that are deemed to be of extra quality. . . . The second market lies in those consumers who wish to use their purchases to support local workers, or to reduce their environmental impact by taking goods to craftspeople to be mended, or recycled.
For workers, the appeal of craftsmanship is that it allows them the autonomy to make creative choices, and thus makes a job far more satisfying. In that sense, it could offer hope for the overall labour market. Let the machines automate dull and repetitive tasks and let workers focus purely on their skills, judgment and imagination. As a current example, the academics cite the “agile” manifesto in the software sector, an industry at the heart of technological change. The pioneers behind the original agile manifesto promised to prioritise “individuals and interactions over processes and tools”. By bringing together experts from different teams, agile working is designed to improve creativity.
But the broader question is whether crafts can create a lot more jobs than they do today. Demand for crafted products may rise but will it be easy to retrain workers in sectors that might get automated (such as truck drivers) to take advantage? In a world where products and services often have to pass through regulatory hoops, large companies will usually have the advantage.
History also suggests that the link between crafts and creativity is not automatic. Medieval craft guilds were monopolies which resisted new entrants. They were also highly hierarchical with young men required to spend long periods as apprentices and journeymen before they could set up on their own; by that time the innovative spirit may have been knocked out of them. Craft workers can thrive in the modern era, but only if they don’t get too organised.

\vspace{1em}
\textbf{Question 3}

The author questions the ability of crafts to create substantial employment opportunities presently because

\begin{enumerate}[label=\Alph*.]
    \item the low scale of crafts production will not be able to absorb the mass of redundant labour.
    \item regulatory requirements could make it difficult for small crafts outfits to compete.
    \item workers made redundant by automation are unlikely to opt for crafts-related work.
    \item crafts guilds tend to resist new entrants and are unlikely to accept large numbers of trainees.
[](https://cracku.in/3-the-author-questions-the-ability-of-crafts-to-crea-x-cat-2024-slot-1)
\end{enumerate}
\vspace{1em}
\textbf{Solution:}

Option B is the correct answer.   
  
The author says _"In a world where products and services often have to pass through regulatory hoops, large companies will usually have the advantage."_ From this, we can infer that the author doubts whether crafts can create substantial employment opportunities because smaller craft businesses may struggle to compete with larger companies due to regulatory barriers._  
_  
_Option A:_ The passage does not focus on the low scale of crafts production as the primary obstacle. Instead, the author emphasizes regulatory challenges and not the scale._  
  
Option C:_ The passage does not argue that workers wouldn’t want to pursue crafts-related work, just that retraining them for these roles might be difficult._  
  
Option D:_ The passage mentions that craft guilds resisted new entrants, but it does not suggest that they are unlikely to accept large number of trainees. The passage suggests that craft workers can thrive in the modern era, but the challenge lies in how modern crafts are organized and their potential to scale up in a competitive market.

\end{problem}

\newpage
\begin{problem}{}{}
% Subject: varc
\textbf{Instructions}

Instructions   

**The passage below is accompanied by four questions. Based on the passage, choose the best answer for each question.**
. . . [T]he idea of craftsmanship is not simply nostalgic. . . . Crafts require distinct skills, an allround approach to work that involves the whole product, rather than individual parts, and an attitude that necessitates devotion to the job and a focus on the communal interest. The concept of craft emphasises the human touch and individual judgment.
Essentially, the crafts concept seems to run against the preponderant ethos of management studies which, as the academics note, have long prioritised efficiency and consistency. . . . Craft skills were portrayed as being primitive and traditionalist.
The contrast between artisanship and efficiency first came to the fore in the 19th century when British manufacturers suddenly faced competition from across the Atlantic as firms developed the “American system” using standardised parts. . . . the worldwide success of the Singer sewing machine showed the potential of a mass-produced device. This process created its own reaction, first in the form of the Arts and Crafts movement of the late 19th century, and then again in the “small is beautiful” movement of the 1970s. A third crafts movement is emerging as people become aware of the environmental impact of conventional industry.
There are two potential markets for those who practise crafts. The first stems from the existence of consumers who are willing to pay a premium price for goods that are deemed to be of extra quality. . . . The second market lies in those consumers who wish to use their purchases to support local workers, or to reduce their environmental impact by taking goods to craftspeople to be mended, or recycled.
For workers, the appeal of craftsmanship is that it allows them the autonomy to make creative choices, and thus makes a job far more satisfying. In that sense, it could offer hope for the overall labour market. Let the machines automate dull and repetitive tasks and let workers focus purely on their skills, judgment and imagination. As a current example, the academics cite the “agile” manifesto in the software sector, an industry at the heart of technological change. The pioneers behind the original agile manifesto promised to prioritise “individuals and interactions over processes and tools”. By bringing together experts from different teams, agile working is designed to improve creativity.
But the broader question is whether crafts can create a lot more jobs than they do today. Demand for crafted products may rise but will it be easy to retrain workers in sectors that might get automated (such as truck drivers) to take advantage? In a world where products and services often have to pass through regulatory hoops, large companies will usually have the advantage.
History also suggests that the link between crafts and creativity is not automatic. Medieval craft guilds were monopolies which resisted new entrants. They were also highly hierarchical with young men required to spend long periods as apprentices and journeymen before they could set up on their own; by that time the innovative spirit may have been knocked out of them. Craft workers can thrive in the modern era, but only if they don’t get too organised.

\vspace{1em}
\textbf{Question 4}

The most recent revival in interest in the crafts is a result of the emergence of all of the following EXCEPT:

\begin{enumerate}[label=\Alph*.]
    \item support for individual creations as opposed to mass-produced objects
    \item concerns about the environmental impact of mass production.
    \item a niche market for discerning buyers of quality products.
    \item a greater interest in buying locally produced goods.
[](https://cracku.in/4-the-most-recent-revival-in-interest-in-the-crafts--x-cat-2024-slot-1)
\end{enumerate}
\vspace{1em}
\textbf{Solution:}

Option A is the correct answer.   
  
While the passage mentions the premium for high-quality crafted goods and the contrast between mass production and craftsmanship, it does not mention that there is support for "individual creations" as opposed to mass-produced goods. The focus is on quality and sustainability rather than explicitly advocating for individuality in products.
Option B: This is mentioned. The passage talks about consumers wanting to reduce environmental impact by supporting local workers or recycling goods, which reflects growing concerns over mass production.
Option C: This is mentioned. The passage describes a market of consumers who pay a premium for high-quality, hand-crafted goods.
Option D: This is mentioned. The passage highlights a market where consumers buy from local workers to support the community and reduce environmental impact.

Instructions   

For the following questions answer them individually

\end{problem}

\newpage
\begin{problem}{}{}
% Subject: varc
\textbf{Instructions}

Instructions   

**The passage below is accompanied by four questions. Based on the passage, choose the best answer for each question.**
. . . [T]he idea of craftsmanship is not simply nostalgic. . . . Crafts require distinct skills, an allround approach to work that involves the whole product, rather than individual parts, and an attitude that necessitates devotion to the job and a focus on the communal interest. The concept of craft emphasises the human touch and individual judgment.
Essentially, the crafts concept seems to run against the preponderant ethos of management studies which, as the academics note, have long prioritised efficiency and consistency. . . . Craft skills were portrayed as being primitive and traditionalist.
The contrast between artisanship and efficiency first came to the fore in the 19th century when British manufacturers suddenly faced competition from across the Atlantic as firms developed the “American system” using standardised parts. . . . the worldwide success of the Singer sewing machine showed the potential of a mass-produced device. This process created its own reaction, first in the form of the Arts and Crafts movement of the late 19th century, and then again in the “small is beautiful” movement of the 1970s. A third crafts movement is emerging as people become aware of the environmental impact of conventional industry.
There are two potential markets for those who practise crafts. The first stems from the existence of consumers who are willing to pay a premium price for goods that are deemed to be of extra quality. . . . The second market lies in those consumers who wish to use their purchases to support local workers, or to reduce their environmental impact by taking goods to craftspeople to be mended, or recycled.
For workers, the appeal of craftsmanship is that it allows them the autonomy to make creative choices, and thus makes a job far more satisfying. In that sense, it could offer hope for the overall labour market. Let the machines automate dull and repetitive tasks and let workers focus purely on their skills, judgment and imagination. As a current example, the academics cite the “agile” manifesto in the software sector, an industry at the heart of technological change. The pioneers behind the original agile manifesto promised to prioritise “individuals and interactions over processes and tools”. By bringing together experts from different teams, agile working is designed to improve creativity.
But the broader question is whether crafts can create a lot more jobs than they do today. Demand for crafted products may rise but will it be easy to retrain workers in sectors that might get automated (such as truck drivers) to take advantage? In a world where products and services often have to pass through regulatory hoops, large companies will usually have the advantage.
History also suggests that the link between crafts and creativity is not automatic. Medieval craft guilds were monopolies which resisted new entrants. They were also highly hierarchical with young men required to spend long periods as apprentices and journeymen before they could set up on their own; by that time the innovative spirit may have been knocked out of them. Craft workers can thrive in the modern era, but only if they don’t get too organised.

\vspace{1em}
\textbf{Question 5}

**The passage given below is followed by four alternate summaries. Choose the option that best captures the essence of the passage.**
Scientific research shows that many animals are very intelligent and have sensory and motor abilities that dwarf ours. Dogs are able to detect diseases such as cancer and diabetes and warn humans of impending heart attacks and strokes. Elephants, whales, hippopotamuses, giraffes, and alligators use low-frequency sounds to communicate over long distances, often miles. Many animals also display wide-ranging emotions, including joy, happiness, empathy, compassion, grief, and even resentment and embarrassment. It’s not surprising that animals share many emotions with us because we also share brain structures, located in the limbic system, that are the seat of our emotions.

\begin{enumerate}[label=\Alph*.]
    \item The advanced sensory and motor abilities of animals is the reason why they can display wide-ranging emotions
    \item The similarity in brain structure explains why animals show emotions typically associated with humans
    \item Animals can show emotions which are typically associated with humans.
    \item Animals are more intelligent than us in sensing danger and detecting diseases.
[](https://cracku.in/5-the-passage-given-below-is-followed-by-four-altern-x-cat-2024-slot-1)
\end{enumerate}
\vspace{1em}
\textbf{Solution:}

Option B is the correct answer.   
  
The paragraph states that animals share many emotions with humans, such as joy, happiness, empathy, and grief, because of shared brain structures, particularly in the limbic system, which is responsible for emotions in both humans and animals. This is the key point that ties together animals' intelligence and emotional capacity, as discussed in the passage.
Option A: The passage does not attribute emotions to sensory and motor abilities. The emphasis is on brain structures, not sensory abilities.
Option C: While the passage states that animals share emotions with humans, this option fails to capture the reason behind this, i.e. shared brain structures. It misses the point that makes the emotional similarity possible, which is central to the passage's message.
Option D: The passage discusses animals' sensory abilities but does not suggest that their intelligence is superior to humans'.

Instructions   

**The passage below is accompanied by four questions. Based on the passage, choose the best answer for each question.**
Oftentimes, when economists cross borders, they are less interested in learning from others than in invading their garden plots. Gary Becker, for instance, pioneered the idea of human capital. To do so, he famously tackled topics like crime and domesticity, applying methods honed in the study of markets to domains of nonmarket life. He projected economics outward into new realms: for example, by revealing the extent to which humans calculate marginal utilities when choosing their spouses or stealing from neighbors. At the same time, he did not let other ways of thinking enter his own economic realm: for example, he did not borrow from anthropology or history or let observations of nonmarket economics inform his homo economicus. Becker was a picture of the imperial economist in the heyday of the discipline’s bravura.
Times have changed for the once almighty discipline. Economics has been taken to task, within and beyond its ramparts. Some economists have reached out, imported, borrowed, and collaborated—been less imperial, more open. Consider Thomas Piketty and his outreach to historians. The booming field of behavioral economics—the fusion of economics and social psychology—is another case. Having spawned active subfields, like judgment, decisionmaking and a turn to experimentation, the field aims to go beyond the caricature of Rational Man to explain how humans make decisions….
It is important to underscore how this flips the way we think about economics. For generations, economists have presumed that people have interests—“preferences,” in the neoclassical argot—that get revealed in the course of peoples’ choices. Interests come before actions and determine them. If you are hungry, you buy lunch; if you are cold, you get a sweater. If you only have so much money and can’t afford to deal with both your growling stomach and your shivering, which need you choose to meet using your scarce savings reveals your preference.
Psychologists take one look at this simple formulation and shake their heads. Increasingly, even some mainstream economists have to admit that homo economicus doesn’t always behave like the textbook maximizer; irrational behavior can’t simply be waved away as extraeconomic expressions of passions over interests, and thus the domain of other disciplines…. This is one place where the humanist can help the economist. If narrative economics is going to help us understand how rivals duke it out, who wins and who loses, we are going to need much more than lessons from epidemiological studies of viruses or intracranial stimuli.
Above all, we need politics and institutions. Shiller [the Nobel prize winning economist] connects perceptions of narratives to changes in behavior and thence to social outcomes. He completes a circle that was key to behavioral economics and brings in storytelling to make sense of how perceptions get framed. This cycle (perception to behavior to society) was once mediated or dominated by institutions: the political parties, lobby groups, and media organizations that played a vital role in legitimating, representing, and excluding interests. Yet institutions have been stripped from Shiller’s account, to reveal a bare dynamic of emotions and economics, without the intermediating place of politics.

\end{problem}

\newpage
\begin{problem}{}{}
% Subject: varc
\textbf{Instructions}

Instructions   

**The passage below is accompanied by four questions. Based on the passage, choose the best answer for each question.**
. . . [T]he idea of craftsmanship is not simply nostalgic. . . . Crafts require distinct skills, an allround approach to work that involves the whole product, rather than individual parts, and an attitude that necessitates devotion to the job and a focus on the communal interest. The concept of craft emphasises the human touch and individual judgment.
Essentially, the crafts concept seems to run against the preponderant ethos of management studies which, as the academics note, have long prioritised efficiency and consistency. . . . Craft skills were portrayed as being primitive and traditionalist.
The contrast between artisanship and efficiency first came to the fore in the 19th century when British manufacturers suddenly faced competition from across the Atlantic as firms developed the “American system” using standardised parts. . . . the worldwide success of the Singer sewing machine showed the potential of a mass-produced device. This process created its own reaction, first in the form of the Arts and Crafts movement of the late 19th century, and then again in the “small is beautiful” movement of the 1970s. A third crafts movement is emerging as people become aware of the environmental impact of conventional industry.
There are two potential markets for those who practise crafts. The first stems from the existence of consumers who are willing to pay a premium price for goods that are deemed to be of extra quality. . . . The second market lies in those consumers who wish to use their purchases to support local workers, or to reduce their environmental impact by taking goods to craftspeople to be mended, or recycled.
For workers, the appeal of craftsmanship is that it allows them the autonomy to make creative choices, and thus makes a job far more satisfying. In that sense, it could offer hope for the overall labour market. Let the machines automate dull and repetitive tasks and let workers focus purely on their skills, judgment and imagination. As a current example, the academics cite the “agile” manifesto in the software sector, an industry at the heart of technological change. The pioneers behind the original agile manifesto promised to prioritise “individuals and interactions over processes and tools”. By bringing together experts from different teams, agile working is designed to improve creativity.
But the broader question is whether crafts can create a lot more jobs than they do today. Demand for crafted products may rise but will it be easy to retrain workers in sectors that might get automated (such as truck drivers) to take advantage? In a world where products and services often have to pass through regulatory hoops, large companies will usually have the advantage.
History also suggests that the link between crafts and creativity is not automatic. Medieval craft guilds were monopolies which resisted new entrants. They were also highly hierarchical with young men required to spend long periods as apprentices and journeymen before they could set up on their own; by that time the innovative spirit may have been knocked out of them. Craft workers can thrive in the modern era, but only if they don’t get too organised.

\vspace{1em}
\textbf{Question 6}

We can infer from the passage that the term '‘homo economicus” refers to someone who

\begin{enumerate}[label=\Alph*.]
    \item is not influenced by the preferences and choices of others.
    \item believes in borrowing and collaborating with other disciplines in their work.
    \item makes rational decisions based on their own preferences.
    \item maximises their opportunities based on nonmarket choices.
[](https://cracku.in/6-we-can-infer-from-the-passage-that-the-term-homo-e-x-cat-2024-slot-1)
\end{enumerate}
\vspace{1em}
\textbf{Solution:}

Option C is the correct answer.   
  
The passage suggests that the idea of 'homo economicus' assumes that people have “preferences” that determine their choices. Deciding whether to spend limited money on food or warmth illustrates this rational decision-making based on individual preferences. This aligns with option C
Option A: The passage does not suggest that 'homo economicus' is not influenced by others’ preferences. It only states that this model assumes people have their own preferences that determine their actions.
Option B: This is wrong as the passage states that economists like Gary Becker, associated with the homo economicus model, did not borrow or collaborate with other disciplines.
Option D: The passage mentions applying economic reasoning to nonmarket domains but does not tie this specifically to homo economicus.

\end{problem}

\newpage
\begin{problem}{}{}
% Subject: varc
\textbf{Instructions}

Instructions   

**The passage below is accompanied by four questions. Based on the passage, choose the best answer for each question.**
. . . [T]he idea of craftsmanship is not simply nostalgic. . . . Crafts require distinct skills, an allround approach to work that involves the whole product, rather than individual parts, and an attitude that necessitates devotion to the job and a focus on the communal interest. The concept of craft emphasises the human touch and individual judgment.
Essentially, the crafts concept seems to run against the preponderant ethos of management studies which, as the academics note, have long prioritised efficiency and consistency. . . . Craft skills were portrayed as being primitive and traditionalist.
The contrast between artisanship and efficiency first came to the fore in the 19th century when British manufacturers suddenly faced competition from across the Atlantic as firms developed the “American system” using standardised parts. . . . the worldwide success of the Singer sewing machine showed the potential of a mass-produced device. This process created its own reaction, first in the form of the Arts and Crafts movement of the late 19th century, and then again in the “small is beautiful” movement of the 1970s. A third crafts movement is emerging as people become aware of the environmental impact of conventional industry.
There are two potential markets for those who practise crafts. The first stems from the existence of consumers who are willing to pay a premium price for goods that are deemed to be of extra quality. . . . The second market lies in those consumers who wish to use their purchases to support local workers, or to reduce their environmental impact by taking goods to craftspeople to be mended, or recycled.
For workers, the appeal of craftsmanship is that it allows them the autonomy to make creative choices, and thus makes a job far more satisfying. In that sense, it could offer hope for the overall labour market. Let the machines automate dull and repetitive tasks and let workers focus purely on their skills, judgment and imagination. As a current example, the academics cite the “agile” manifesto in the software sector, an industry at the heart of technological change. The pioneers behind the original agile manifesto promised to prioritise “individuals and interactions over processes and tools”. By bringing together experts from different teams, agile working is designed to improve creativity.
But the broader question is whether crafts can create a lot more jobs than they do today. Demand for crafted products may rise but will it be easy to retrain workers in sectors that might get automated (such as truck drivers) to take advantage? In a world where products and services often have to pass through regulatory hoops, large companies will usually have the advantage.
History also suggests that the link between crafts and creativity is not automatic. Medieval craft guilds were monopolies which resisted new entrants. They were also highly hierarchical with young men required to spend long periods as apprentices and journeymen before they could set up on their own; by that time the innovative spirit may have been knocked out of them. Craft workers can thrive in the modern era, but only if they don’t get too organised.

\vspace{1em}
\textbf{Question 7}

“Times have changed for the once almighty discipline.” We can infer from this statement and the associated paragraph that the author is being

\begin{enumerate}[label=\Alph*.]
    \item sarcastic about how economists, who earlier shunned other disciplines, are now beginning to incorporate them in their analyses
    \item disparaging of economists’ inability to precisely predict market behaviour, and are now borrowing from other disciplines to remedy this
    \item judgemental about the ability of economic tools to accurately manage crises leading to the downfall of this lofty science.
    \item critical of economists’ openly borrowing and collaborating across disciplines to explain how humans make decisions
[](https://cracku.in/7-times-have-changed-for-the-once-almighty-disciplin-x-cat-2024-slot-1)
\end{enumerate}
\vspace{1em}
\textbf{Solution:}

Option A is the correct answer.   
  
The tone of the phrase “almighty discipline” suggests mild sarcasm regarding economics' earlier dominance and self-containment. The statement reflects a shift in economics from being insular and self-assured to being more open to interdisciplinary borrowing. The author describes earlier economists, like Gary Becker, as imperialistic and unwilling to incorporate ideas from other disciplines. In contrast, contemporary economists like Thomas Piketty are portrayed as reaching out and collaborating.
Option B: The passage mentions that economists are now collaborating with other disciplines, but it does not suggest this is due to their inability to predict market behaviour. The focus is on how the discipline has shifted over time, not any predictive failures.
Option C: The author does not criticize economic tools for managing crises or say that economics as a discipline has fallen. The passage highlights a transformation toward openness, not a downfall.
Option D: The author is not critical of economists collaborating with other disciplines. Instead, the passage presents this shift as a positive change, contrasting it with the earlier insularity of the discipline.

\end{problem}

\newpage
\begin{problem}{}{}
% Subject: varc
\textbf{Instructions}

Instructions   

**The passage below is accompanied by four questions. Based on the passage, choose the best answer for each question.**
. . . [T]he idea of craftsmanship is not simply nostalgic. . . . Crafts require distinct skills, an allround approach to work that involves the whole product, rather than individual parts, and an attitude that necessitates devotion to the job and a focus on the communal interest. The concept of craft emphasises the human touch and individual judgment.
Essentially, the crafts concept seems to run against the preponderant ethos of management studies which, as the academics note, have long prioritised efficiency and consistency. . . . Craft skills were portrayed as being primitive and traditionalist.
The contrast between artisanship and efficiency first came to the fore in the 19th century when British manufacturers suddenly faced competition from across the Atlantic as firms developed the “American system” using standardised parts. . . . the worldwide success of the Singer sewing machine showed the potential of a mass-produced device. This process created its own reaction, first in the form of the Arts and Crafts movement of the late 19th century, and then again in the “small is beautiful” movement of the 1970s. A third crafts movement is emerging as people become aware of the environmental impact of conventional industry.
There are two potential markets for those who practise crafts. The first stems from the existence of consumers who are willing to pay a premium price for goods that are deemed to be of extra quality. . . . The second market lies in those consumers who wish to use their purchases to support local workers, or to reduce their environmental impact by taking goods to craftspeople to be mended, or recycled.
For workers, the appeal of craftsmanship is that it allows them the autonomy to make creative choices, and thus makes a job far more satisfying. In that sense, it could offer hope for the overall labour market. Let the machines automate dull and repetitive tasks and let workers focus purely on their skills, judgment and imagination. As a current example, the academics cite the “agile” manifesto in the software sector, an industry at the heart of technological change. The pioneers behind the original agile manifesto promised to prioritise “individuals and interactions over processes and tools”. By bringing together experts from different teams, agile working is designed to improve creativity.
But the broader question is whether crafts can create a lot more jobs than they do today. Demand for crafted products may rise but will it be easy to retrain workers in sectors that might get automated (such as truck drivers) to take advantage? In a world where products and services often have to pass through regulatory hoops, large companies will usually have the advantage.
History also suggests that the link between crafts and creativity is not automatic. Medieval craft guilds were monopolies which resisted new entrants. They were also highly hierarchical with young men required to spend long periods as apprentices and journeymen before they could set up on their own; by that time the innovative spirit may have been knocked out of them. Craft workers can thrive in the modern era, but only if they don’t get too organised.

\vspace{1em}
\textbf{Question 8}

The author critiques Schiller’s approach to behavioural economics for

\begin{enumerate}[label=\Alph*.]
    \item ignoring the marginal role that media and politics play in influencing people’s behaviour.
    \item denigrating the role of institutions while creating a link between behavioural economics and perceptions
    \item linking emotions and rational behaviour without considering the mediation of social institutions.
    \item relying excessively on storytelling as the main influence on the formation of perceptions.
[](https://cracku.in/8-the-author-critiques-schillers-approach-to-behavio-x-cat-2024-slot-1)
\end{enumerate}
\vspace{1em}
\textbf{Solution:}

Option C is the correct answer.   
  
The author critiques Schiller for focusing on the direct relationship between emotions, perceptions, and behaviour while excluding the mediating role of institutions such as political parties, media organizations, and lobby groups. The passage states that these institutions historically played a vital role in framing perceptions and legitimizing interests, which Schiller’s narrative leaves out.
Option A: The passage suggests that media and politics are key intermediaries in shaping perceptions and behaviour and not a marginal role.
Option B: The passage does not suggest that Shiller actively denigrates or diminishes the role of institutions. The issue is his omission of institutions, not his critique of their importance.
Option D: The passage mentions that storytelling is part of Schiller’s approach but does not critique him for relying excessively on it. The critique is about the exclusion of institutions, not the use of storytelling.

\end{problem}

\newpage
\begin{problem}{}{}
% Subject: varc
\textbf{Instructions}

Instructions   

**The passage below is accompanied by four questions. Based on the passage, choose the best answer for each question.**
. . . [T]he idea of craftsmanship is not simply nostalgic. . . . Crafts require distinct skills, an allround approach to work that involves the whole product, rather than individual parts, and an attitude that necessitates devotion to the job and a focus on the communal interest. The concept of craft emphasises the human touch and individual judgment.
Essentially, the crafts concept seems to run against the preponderant ethos of management studies which, as the academics note, have long prioritised efficiency and consistency. . . . Craft skills were portrayed as being primitive and traditionalist.
The contrast between artisanship and efficiency first came to the fore in the 19th century when British manufacturers suddenly faced competition from across the Atlantic as firms developed the “American system” using standardised parts. . . . the worldwide success of the Singer sewing machine showed the potential of a mass-produced device. This process created its own reaction, first in the form of the Arts and Crafts movement of the late 19th century, and then again in the “small is beautiful” movement of the 1970s. A third crafts movement is emerging as people become aware of the environmental impact of conventional industry.
There are two potential markets for those who practise crafts. The first stems from the existence of consumers who are willing to pay a premium price for goods that are deemed to be of extra quality. . . . The second market lies in those consumers who wish to use their purchases to support local workers, or to reduce their environmental impact by taking goods to craftspeople to be mended, or recycled.
For workers, the appeal of craftsmanship is that it allows them the autonomy to make creative choices, and thus makes a job far more satisfying. In that sense, it could offer hope for the overall labour market. Let the machines automate dull and repetitive tasks and let workers focus purely on their skills, judgment and imagination. As a current example, the academics cite the “agile” manifesto in the software sector, an industry at the heart of technological change. The pioneers behind the original agile manifesto promised to prioritise “individuals and interactions over processes and tools”. By bringing together experts from different teams, agile working is designed to improve creativity.
But the broader question is whether crafts can create a lot more jobs than they do today. Demand for crafted products may rise but will it be easy to retrain workers in sectors that might get automated (such as truck drivers) to take advantage? In a world where products and services often have to pass through regulatory hoops, large companies will usually have the advantage.
History also suggests that the link between crafts and creativity is not automatic. Medieval craft guilds were monopolies which resisted new entrants. They were also highly hierarchical with young men required to spend long periods as apprentices and journeymen before they could set up on their own; by that time the innovative spirit may have been knocked out of them. Craft workers can thrive in the modern era, but only if they don’t get too organised.

\vspace{1em}
\textbf{Question 9}

In the first paragraph the author is making the point that economists like Becker

\begin{enumerate}[label=\Alph*.]
    \item benefitted from the application of their principles and concepts to non-economic phenomena.
    \item had begun to borrow concepts from other disciplines but were averse to the latter applying economic principles.
    \item used economics to analyse non-market behaviour, without incorporating perspectives from other areas of inquiry.
    \item tended to guard their discipline from poaching by academics from other subject areas.
[](https://cracku.in/9-in-the-first-paragraph-the-author-is-making-the-po-x-cat-2024-slot-1)
\end{enumerate}
\vspace{1em}
\textbf{Solution:}

Option C is the correct answer. 
The first paragraph describes how Becker projected economics into non-market domains like crime and domesticity, analyzing them with economic tools. It also states that he did not borrow ideas or perspectives from fields such as anthropology or history. Option C aligns with this.
Option A: The passage mentions that Becker applied economic principles to non-market phenomena like crime and domesticity, but it does not state or imply that he benefitted from this.
Option B: The passage explicitly contradicts this idea, stating: _“At the same time, he did not let other ways of thinking enter his own economic realm: for example, he did not borrow from anthropology or history.”_ Becker did not borrow from other disciplines, so this option is incorrect.
Option D: The passage does not state that Becker actively guarded economics against outside influence. Instead, it focuses on Becker’s one-sided application of economics to other areas, without accepting external perspectives.

Instructions   

For the following questions answer them individually

\end{problem}

\newpage
\begin{problem}{}{}
% Subject: varc
\textbf{Instructions}

Instructions   

**The passage below is accompanied by four questions. Based on the passage, choose the best answer for each question.**
. . . [T]he idea of craftsmanship is not simply nostalgic. . . . Crafts require distinct skills, an allround approach to work that involves the whole product, rather than individual parts, and an attitude that necessitates devotion to the job and a focus on the communal interest. The concept of craft emphasises the human touch and individual judgment.
Essentially, the crafts concept seems to run against the preponderant ethos of management studies which, as the academics note, have long prioritised efficiency and consistency. . . . Craft skills were portrayed as being primitive and traditionalist.
The contrast between artisanship and efficiency first came to the fore in the 19th century when British manufacturers suddenly faced competition from across the Atlantic as firms developed the “American system” using standardised parts. . . . the worldwide success of the Singer sewing machine showed the potential of a mass-produced device. This process created its own reaction, first in the form of the Arts and Crafts movement of the late 19th century, and then again in the “small is beautiful” movement of the 1970s. A third crafts movement is emerging as people become aware of the environmental impact of conventional industry.
There are two potential markets for those who practise crafts. The first stems from the existence of consumers who are willing to pay a premium price for goods that are deemed to be of extra quality. . . . The second market lies in those consumers who wish to use their purchases to support local workers, or to reduce their environmental impact by taking goods to craftspeople to be mended, or recycled.
For workers, the appeal of craftsmanship is that it allows them the autonomy to make creative choices, and thus makes a job far more satisfying. In that sense, it could offer hope for the overall labour market. Let the machines automate dull and repetitive tasks and let workers focus purely on their skills, judgment and imagination. As a current example, the academics cite the “agile” manifesto in the software sector, an industry at the heart of technological change. The pioneers behind the original agile manifesto promised to prioritise “individuals and interactions over processes and tools”. By bringing together experts from different teams, agile working is designed to improve creativity.
But the broader question is whether crafts can create a lot more jobs than they do today. Demand for crafted products may rise but will it be easy to retrain workers in sectors that might get automated (such as truck drivers) to take advantage? In a world where products and services often have to pass through regulatory hoops, large companies will usually have the advantage.
History also suggests that the link between crafts and creativity is not automatic. Medieval craft guilds were monopolies which resisted new entrants. They were also highly hierarchical with young men required to spend long periods as apprentices and journeymen before they could set up on their own; by that time the innovative spirit may have been knocked out of them. Craft workers can thrive in the modern era, but only if they don’t get too organised.

\vspace{1em}
\textbf{Question 10}

**There is a sentence that is missing in the paragraph below. Look at the paragraph and decide where (option 1, 2, 3, or 4) the following sentence would best fit.**
**Sentence** : Comprehending a wide range of emotions, Renaissance music nevertheless portrayed all emotions in a balanced and moderate fashion.   
  
**Paragraph** : A volume of translated Italian madrigals were published in London during the year of 1588. This sudden public interest facilitated a surge of English Madrigal writing as well as a spurt of other secular music writing and publication. ___(1)___. This music boom lasted for thirty years and was as much a golden age of music as British literature was with Shakespeare and Queen Elizabeth I. ___(2)___. The rebirth in both literature and music originated in Italy and migrated to England; the English madrigal became more humorous and lighter in England as compared to Italy. Renaissance music was mostly polyphonic in texture. ___(3)___. Extreme use of and contrasts in dynamics, rhythm, and tone colour do not occur. ___(4)___. The rhythms in Renaissance music tend to have a smooth, soft flow instead of a sharp, well-defined pulse of accents.

\begin{enumerate}[label=\Alph*.]
    \item Option 1
    \item Option 3
    \item Option 4
    \item Option 2
[](https://cracku.in/10-there-is-a-sentence-that-is-missing-in-the-paragra-x-cat-2024-slot-1)
\end{enumerate}
\vspace{1em}
\textbf{Solution:}

The sentence would best fit in B) Option 3.  
  
Option 3 is the most logical position because it continues the description of Renaissance music’s characteristics, which was mentioned in the sentence before blank 3. The missing sentence fits perfectly here, as it discusses how Renaissance music handles emotions in a balanced and moderate manner, aligning with the flow and tone of the preceding sentences.  
  
The sentences around Option 1 and Option 2 discuss the historical context of English madrigals and their cultural significance, so adding the sentence there would disrupt the historical narrative.  

Option 4 may appear plausible, but the sentence discussing the balanced and moderate nature of music after asserting that extremes do not occur would be redundant. The given sentence should precede the one addressing extremes to maintain logical flow.

\end{problem}

\newpage
\begin{problem}{}{}
% Subject: varc
\textbf{Instructions}

Instructions   

**The passage below is accompanied by four questions. Based on the passage, choose the best answer for each question.**
. . . [T]he idea of craftsmanship is not simply nostalgic. . . . Crafts require distinct skills, an allround approach to work that involves the whole product, rather than individual parts, and an attitude that necessitates devotion to the job and a focus on the communal interest. The concept of craft emphasises the human touch and individual judgment.
Essentially, the crafts concept seems to run against the preponderant ethos of management studies which, as the academics note, have long prioritised efficiency and consistency. . . . Craft skills were portrayed as being primitive and traditionalist.
The contrast between artisanship and efficiency first came to the fore in the 19th century when British manufacturers suddenly faced competition from across the Atlantic as firms developed the “American system” using standardised parts. . . . the worldwide success of the Singer sewing machine showed the potential of a mass-produced device. This process created its own reaction, first in the form of the Arts and Crafts movement of the late 19th century, and then again in the “small is beautiful” movement of the 1970s. A third crafts movement is emerging as people become aware of the environmental impact of conventional industry.
There are two potential markets for those who practise crafts. The first stems from the existence of consumers who are willing to pay a premium price for goods that are deemed to be of extra quality. . . . The second market lies in those consumers who wish to use their purchases to support local workers, or to reduce their environmental impact by taking goods to craftspeople to be mended, or recycled.
For workers, the appeal of craftsmanship is that it allows them the autonomy to make creative choices, and thus makes a job far more satisfying. In that sense, it could offer hope for the overall labour market. Let the machines automate dull and repetitive tasks and let workers focus purely on their skills, judgment and imagination. As a current example, the academics cite the “agile” manifesto in the software sector, an industry at the heart of technological change. The pioneers behind the original agile manifesto promised to prioritise “individuals and interactions over processes and tools”. By bringing together experts from different teams, agile working is designed to improve creativity.
But the broader question is whether crafts can create a lot more jobs than they do today. Demand for crafted products may rise but will it be easy to retrain workers in sectors that might get automated (such as truck drivers) to take advantage? In a world where products and services often have to pass through regulatory hoops, large companies will usually have the advantage.
History also suggests that the link between crafts and creativity is not automatic. Medieval craft guilds were monopolies which resisted new entrants. They were also highly hierarchical with young men required to spend long periods as apprentices and journeymen before they could set up on their own; by that time the innovative spirit may have been knocked out of them. Craft workers can thrive in the modern era, but only if they don’t get too organised.

\vspace{1em}
\textbf{Question 11}

**There is a sentence that is missing in the paragraph below. Look at the paragraph and decide where (option 1, 2, 3, or 4) the following sentence would best fit.**
**Sentence** : Understanding central Asia’s role helps developments make more sense not only across Asia but in Europe, the Americas and Africa.
**Paragraph** : The nations of the Silk Roads are sometimes called ‘developing countries’, but they are actually some of the world’s most highly developed countries, the very crossroads of civilization, in advanced states of disrepair. ___(1)___. These countries lie at the centre of global affairs: they have since the beginning of history. Running across the spine of Asia, they form a web of connections fanning out in every direction, routes along which pilgrims and warriors, nomads and merchants have travelled, goods and produce have been bought and sold, and ideas exchanged, adapted and refined. ___(2)___ .They have carried not only prosperity, but also death and violence, disease and disaster. ___(3)___. The Silk Roads are the world’s central nervous system, connecting otherwise far-flung peoples and places…. ___(4)___. It allows us to see patterns and links, causes and effects that remain invisible if one looks only at Europe, or North America.

\begin{enumerate}[label=\Alph*.]
    \item Option 3
    \item Option 2
    \item Option 1
    \item Option 4
[](https://cracku.in/11-there-is-a-sentence-that-is-missing-in-the-paragra-x-cat-2024-slot-1)
\end{enumerate}
\vspace{1em}
\textbf{Solution:}

We can observe that the only spot where the given sentence would fit is Blank 4. This sentence in question emphasises the broader, global significance of Central Asia's role; it reflects the idea that understanding the Silk Roads provides clarity on historical and global interconnectedness. 
Since the given sentence shifts focus from describing the historical and cultural importance of the Silk Roads to their global implications, which is broader than the current discussion, placing it in any of Blanks 1, 2, or 3 interrupts the flow between the statements outlining certain features linked to the Silk Roads. For example, Blank 1 is a poor fit since the succeeding sentence mentions “These countries,” which refers to the “developing countries” mentioned in the sentence preceding the blank. Similarly, “They” in the sentence after Blank 2 refers to routes/Silk Roads described earlier. We can eliminate Blank 3 using a similar logic; the sentence after this blank establishes the metaphor of the Silk Roads as a central nervous system, emphasizing their global connectivity. The given sentence doesn’t align with the focus on the Silk Roads’ tangible impacts or the metaphorical depiction of their connectivity; therefore, placing it in Blank 3 would shift the focus prematurely to understanding their role globally.

\end{problem}

\newpage
\begin{problem}{}{}
% Subject: varc
\textbf{Instructions}

Instructions   

**The passage below is accompanied by four questions. Based on the passage, choose the best answer for each question.**
. . . [T]he idea of craftsmanship is not simply nostalgic. . . . Crafts require distinct skills, an allround approach to work that involves the whole product, rather than individual parts, and an attitude that necessitates devotion to the job and a focus on the communal interest. The concept of craft emphasises the human touch and individual judgment.
Essentially, the crafts concept seems to run against the preponderant ethos of management studies which, as the academics note, have long prioritised efficiency and consistency. . . . Craft skills were portrayed as being primitive and traditionalist.
The contrast between artisanship and efficiency first came to the fore in the 19th century when British manufacturers suddenly faced competition from across the Atlantic as firms developed the “American system” using standardised parts. . . . the worldwide success of the Singer sewing machine showed the potential of a mass-produced device. This process created its own reaction, first in the form of the Arts and Crafts movement of the late 19th century, and then again in the “small is beautiful” movement of the 1970s. A third crafts movement is emerging as people become aware of the environmental impact of conventional industry.
There are two potential markets for those who practise crafts. The first stems from the existence of consumers who are willing to pay a premium price for goods that are deemed to be of extra quality. . . . The second market lies in those consumers who wish to use their purchases to support local workers, or to reduce their environmental impact by taking goods to craftspeople to be mended, or recycled.
For workers, the appeal of craftsmanship is that it allows them the autonomy to make creative choices, and thus makes a job far more satisfying. In that sense, it could offer hope for the overall labour market. Let the machines automate dull and repetitive tasks and let workers focus purely on their skills, judgment and imagination. As a current example, the academics cite the “agile” manifesto in the software sector, an industry at the heart of technological change. The pioneers behind the original agile manifesto promised to prioritise “individuals and interactions over processes and tools”. By bringing together experts from different teams, agile working is designed to improve creativity.
But the broader question is whether crafts can create a lot more jobs than they do today. Demand for crafted products may rise but will it be easy to retrain workers in sectors that might get automated (such as truck drivers) to take advantage? In a world where products and services often have to pass through regulatory hoops, large companies will usually have the advantage.
History also suggests that the link between crafts and creativity is not automatic. Medieval craft guilds were monopolies which resisted new entrants. They were also highly hierarchical with young men required to spend long periods as apprentices and journeymen before they could set up on their own; by that time the innovative spirit may have been knocked out of them. Craft workers can thrive in the modern era, but only if they don’t get too organised.

\vspace{1em}
\textbf{Question 12}

**Five jumbled up sentences (labelled 1, 2, 3, 4 and 5), related to a topic, are given below. Four of them can be put together to form a coherent paragraph. Identify the odd sentence and key in the number of that sentence as your answer.**
1. Urbanites also have more and better options for getting around: Uber is ubiquitous; easy-to-rent dockless bicycles are spreading; battery-powered scooters will be next.  
2. When more people use buses or trains the service usually improves because public-transport agencies run more buses and trains.  
3. Worsening services on public transport, terrorist attacks in some urban metros and a rise in fares have been blamed for this trend.  
4. It seems more likely that public transport is being squeezed structurally as people’s need to travel is diminishing as a result of smartphones, videoconferencing, online shopping and so on.  
5. There has been a puzzling decline in the use of urban public transport in many countries in the west, despite the growth in urban populations and rising employment.
Backspace  
789  
456  
123  
0.-  
Clear All  

Submit
[](https://cracku.in/12-five-jumbled-up-sentences-labelled-1-2-3-4-and-5-r-x-cat-2024-slot-1)

\begin{enumerate}[label=\Alph*.]
\end{enumerate}
\vspace{1em}
\textbf{Solution:}

**Sentence 2 is the odd one out.  
  
The sequence 5-3-4-1 forms a coherent paragraph.**
Sentence 5 introduces the issue of the puzzling decline in public transport despite growing urban populations and employment.  
Sentence 3, then explains some of the reasons for the decline, like worsening services, terrorist attacks, and rising fares.  
Sentence 4 adds a structural explanation, stating that public transport is being squeezed as people travel less due to technology.  
Sentence 1 provides more details about the alternatives to public transport, which are contributing to the decline in usage.  
  
Sentence 2 talks about how more people using buses or trains improves the service, but this idea is not directly relevant to the overall discussion in the passage, which focuses on the decline of public transport usage. It shifts the focus toward an increase in usage improving services, which contradicts the central theme of decline.

Instructions   

**The passage below is accompanied by four questions. Based on the passage, choose the best answer for each question.**
Landing in Australia, the British colonists weren’t much impressed with the small-bodied, slender-snooted marsupials called bandicoots. “Their muzzle, which is much too long, gives them an air exceedingly stupid,” one naturalist noted in 1805. They nicknamed one type the “zebra rat” because of its black-striped rump.
Silly-looking or not, though, the zebra rat—the smallest bandicoot, more commonly known today as the western barred bandicoot—exhibited a genius for survival in the harsh outback, where its ancestors had persisted for some 26 million years. Its births were triggered by rainfall in the bone-dry desert. It carried its breath-mint-size babies in a backward-facing pouch so mothers could forage for food and dig shallow, camouflaged shelters.
Still, these adaptations did not prepare the western barred bandicoot for the colonial-era transformation of its ecosystem, particularly the onslaught of imported British animals, from cattle and rabbits that damaged delicate desert vegetation to ravenous house cats that soon developed a taste for bandicoots. Several of the dozen-odd bandicoot species went extinct, and by the 1940s the western barred bandicoot, whose original range stretched across much of the continent, persisted only on two predator-free islands in Shark Bay, off Australia’s western coast.
“Our isolated fauna had simply not been exposed to these predators,” says Reece Pedler, an ecologist with the Wild Deserts conservation program.
Now Wild Deserts is using descendants of those few thousand island survivors, called Shark Bay bandicoots, in a new effort to seed a mainland bandicoot revival. They’ve imported 20 bandicoots to a preserve on the edge of the Strzelecki Desert, in the remote interior of New South Wales. This sanctuary is a challenging place, desolate much of the year, with one of the world’s most mercurial rainfall patterns—relentless droughts followed by sudden drenching floods.
The imported bandicoots occupy two fenced “exclosures,” cleared of invasive rabbits (courtesy of Pedler’s sheepdog) and of feral cats (which slunk off once the rabbits disappeared). A third fenced area contains the program’s Wild Training Zone, where two other rare marsupials (bilbies, a larger type of bandicoot, and mulgaras, a somewhat fearsome fuzzball known for sucking the brains out of prey) currently share terrain with controlled numbers of cats, learning to evade them. It’s unclear whether the Shark Bay bandicoots, which are perhaps even more predator-naive than their now-extinct mainland bandicoot kin, will be able to make that kind of breakthrough.
For now, though, a recent surge of rainfall has led to a bandicoot joey boom, raising the Wild Deserts population to about 100, with other sanctuaries adding to that number. There are also signs of rebirth in the landscape itself. With their constant digging, the bandicoots trap moisture and allow for seed germination so the cattle-damaged desert can restore itself.
They have a new nickname—a flattering one, this time. “We call them ecosystem engineers,” Pedler says.

\end{problem}

\newpage
\begin{problem}{}{}
% Subject: varc
\textbf{Instructions}

Instructions   

**The passage below is accompanied by four questions. Based on the passage, choose the best answer for each question.**
. . . [T]he idea of craftsmanship is not simply nostalgic. . . . Crafts require distinct skills, an allround approach to work that involves the whole product, rather than individual parts, and an attitude that necessitates devotion to the job and a focus on the communal interest. The concept of craft emphasises the human touch and individual judgment.
Essentially, the crafts concept seems to run against the preponderant ethos of management studies which, as the academics note, have long prioritised efficiency and consistency. . . . Craft skills were portrayed as being primitive and traditionalist.
The contrast between artisanship and efficiency first came to the fore in the 19th century when British manufacturers suddenly faced competition from across the Atlantic as firms developed the “American system” using standardised parts. . . . the worldwide success of the Singer sewing machine showed the potential of a mass-produced device. This process created its own reaction, first in the form of the Arts and Crafts movement of the late 19th century, and then again in the “small is beautiful” movement of the 1970s. A third crafts movement is emerging as people become aware of the environmental impact of conventional industry.
There are two potential markets for those who practise crafts. The first stems from the existence of consumers who are willing to pay a premium price for goods that are deemed to be of extra quality. . . . The second market lies in those consumers who wish to use their purchases to support local workers, or to reduce their environmental impact by taking goods to craftspeople to be mended, or recycled.
For workers, the appeal of craftsmanship is that it allows them the autonomy to make creative choices, and thus makes a job far more satisfying. In that sense, it could offer hope for the overall labour market. Let the machines automate dull and repetitive tasks and let workers focus purely on their skills, judgment and imagination. As a current example, the academics cite the “agile” manifesto in the software sector, an industry at the heart of technological change. The pioneers behind the original agile manifesto promised to prioritise “individuals and interactions over processes and tools”. By bringing together experts from different teams, agile working is designed to improve creativity.
But the broader question is whether crafts can create a lot more jobs than they do today. Demand for crafted products may rise but will it be easy to retrain workers in sectors that might get automated (such as truck drivers) to take advantage? In a world where products and services often have to pass through regulatory hoops, large companies will usually have the advantage.
History also suggests that the link between crafts and creativity is not automatic. Medieval craft guilds were monopolies which resisted new entrants. They were also highly hierarchical with young men required to spend long periods as apprentices and journeymen before they could set up on their own; by that time the innovative spirit may have been knocked out of them. Craft workers can thrive in the modern era, but only if they don’t get too organised.

\vspace{1em}
\textbf{Question 13}

According to the text, the western barred bandicoots now have a flattering name because they have

\begin{enumerate}[label=\Alph*.]
    \item aided in altering an arid environment.
    \item led a revival in preserving the species.
    \item grown fivefold in terms of population
    \item led to a surge and increase of rainfall
[](https://cracku.in/13-according-to-the-text-the-western-barred-bandicoot-x-cat-2024-slot-1)
\end{enumerate}
\vspace{1em}
\textbf{Solution:}

Option A is the correct answer.   
  
The passage mentions how bandicoots contribute to the ecosystem by digging, which traps moisture and aids in seed germination. This activity helps restore the desert ecosystem damaged by cattle. Their role as "ecosystem engineers" stems from these positive environmental impacts.
Option B: Although efforts to preserve the bandicoot species are ongoing, the passage does not link their new nickname to these efforts. The name "ecosystem engineers" specifically reflects their environmental contributions rather than conservation measures.
Option C: While the passage mentions a population increase due to rainfall, the new nickname is unrelated to this growth. Instead, it is tied to their environmental engineering role.
Option D: This is a wrong interpretation because there is no mention in the passage of the bandicoots affecting rainfall.

\end{problem}

\newpage
\begin{problem}{}{}
% Subject: varc
\textbf{Instructions}

Instructions   

**The passage below is accompanied by four questions. Based on the passage, choose the best answer for each question.**
. . . [T]he idea of craftsmanship is not simply nostalgic. . . . Crafts require distinct skills, an allround approach to work that involves the whole product, rather than individual parts, and an attitude that necessitates devotion to the job and a focus on the communal interest. The concept of craft emphasises the human touch and individual judgment.
Essentially, the crafts concept seems to run against the preponderant ethos of management studies which, as the academics note, have long prioritised efficiency and consistency. . . . Craft skills were portrayed as being primitive and traditionalist.
The contrast between artisanship and efficiency first came to the fore in the 19th century when British manufacturers suddenly faced competition from across the Atlantic as firms developed the “American system” using standardised parts. . . . the worldwide success of the Singer sewing machine showed the potential of a mass-produced device. This process created its own reaction, first in the form of the Arts and Crafts movement of the late 19th century, and then again in the “small is beautiful” movement of the 1970s. A third crafts movement is emerging as people become aware of the environmental impact of conventional industry.
There are two potential markets for those who practise crafts. The first stems from the existence of consumers who are willing to pay a premium price for goods that are deemed to be of extra quality. . . . The second market lies in those consumers who wish to use their purchases to support local workers, or to reduce their environmental impact by taking goods to craftspeople to be mended, or recycled.
For workers, the appeal of craftsmanship is that it allows them the autonomy to make creative choices, and thus makes a job far more satisfying. In that sense, it could offer hope for the overall labour market. Let the machines automate dull and repetitive tasks and let workers focus purely on their skills, judgment and imagination. As a current example, the academics cite the “agile” manifesto in the software sector, an industry at the heart of technological change. The pioneers behind the original agile manifesto promised to prioritise “individuals and interactions over processes and tools”. By bringing together experts from different teams, agile working is designed to improve creativity.
But the broader question is whether crafts can create a lot more jobs than they do today. Demand for crafted products may rise but will it be easy to retrain workers in sectors that might get automated (such as truck drivers) to take advantage? In a world where products and services often have to pass through regulatory hoops, large companies will usually have the advantage.
History also suggests that the link between crafts and creativity is not automatic. Medieval craft guilds were monopolies which resisted new entrants. They were also highly hierarchical with young men required to spend long periods as apprentices and journeymen before they could set up on their own; by that time the innovative spirit may have been knocked out of them. Craft workers can thrive in the modern era, but only if they don’t get too organised.

\vspace{1em}
\textbf{Question 14}

Which one of the following options does NOT represent the characteristics of the western barred bandicoot?

\begin{enumerate}[label=\Alph*.]
    \item Shallow diggers having an elongated muzzle
    \item Smallest black striped marsupial that uses camouflage and dig
    \item Look of a rat but with a baby pouch and a slender snout
    \item Long thin nose, black striped back, pouch for joeys
[](https://cracku.in/14-which-one-of-the-following-options-does-not-repres-x-cat-2024-slot-1)
\end{enumerate}
\vspace{1em}
\textbf{Solution:}

Option B is the correct answer.   
  
The passage does not mention the bandicoots' use of camouflage as a survival technique. While their shelters may be camouflaged (hidden), the bandicoots themselves are not described as using camouflage directly.
Option A: This is correct. The passage mentions their long, slender snouts and their digging behaviour, which allows them to create shallow shelters in the desert.
Option C: This is correct. The nickname "zebra rat" comes from their appearance, and their slender snout and backward-facing pouch for carrying joeys are described in the passage.
Option D: This is also correct. The passage mentions these features as characteristics of the western barred bandicoot.

\end{problem}

\newpage
\begin{problem}{}{}
% Subject: varc
\textbf{Instructions}

Instructions   

**The passage below is accompanied by four questions. Based on the passage, choose the best answer for each question.**
. . . [T]he idea of craftsmanship is not simply nostalgic. . . . Crafts require distinct skills, an allround approach to work that involves the whole product, rather than individual parts, and an attitude that necessitates devotion to the job and a focus on the communal interest. The concept of craft emphasises the human touch and individual judgment.
Essentially, the crafts concept seems to run against the preponderant ethos of management studies which, as the academics note, have long prioritised efficiency and consistency. . . . Craft skills were portrayed as being primitive and traditionalist.
The contrast between artisanship and efficiency first came to the fore in the 19th century when British manufacturers suddenly faced competition from across the Atlantic as firms developed the “American system” using standardised parts. . . . the worldwide success of the Singer sewing machine showed the potential of a mass-produced device. This process created its own reaction, first in the form of the Arts and Crafts movement of the late 19th century, and then again in the “small is beautiful” movement of the 1970s. A third crafts movement is emerging as people become aware of the environmental impact of conventional industry.
There are two potential markets for those who practise crafts. The first stems from the existence of consumers who are willing to pay a premium price for goods that are deemed to be of extra quality. . . . The second market lies in those consumers who wish to use their purchases to support local workers, or to reduce their environmental impact by taking goods to craftspeople to be mended, or recycled.
For workers, the appeal of craftsmanship is that it allows them the autonomy to make creative choices, and thus makes a job far more satisfying. In that sense, it could offer hope for the overall labour market. Let the machines automate dull and repetitive tasks and let workers focus purely on their skills, judgment and imagination. As a current example, the academics cite the “agile” manifesto in the software sector, an industry at the heart of technological change. The pioneers behind the original agile manifesto promised to prioritise “individuals and interactions over processes and tools”. By bringing together experts from different teams, agile working is designed to improve creativity.
But the broader question is whether crafts can create a lot more jobs than they do today. Demand for crafted products may rise but will it be easy to retrain workers in sectors that might get automated (such as truck drivers) to take advantage? In a world where products and services often have to pass through regulatory hoops, large companies will usually have the advantage.
History also suggests that the link between crafts and creativity is not automatic. Medieval craft guilds were monopolies which resisted new entrants. They were also highly hierarchical with young men required to spend long periods as apprentices and journeymen before they could set up on their own; by that time the innovative spirit may have been knocked out of them. Craft workers can thrive in the modern era, but only if they don’t get too organised.

\vspace{1em}
\textbf{Question 15}

The text uses the word ‘exclosures’ because Wild Deserts has adopted a measure of

\begin{enumerate}[label=\Alph*.]
    \item restoring cattle damaged deserts to green landscapes.
    \item ridding the main desert of feral cats and large bilbies
    \item excluding animals to make the islands predator-free.
    \item barring the entry of invasive species.
[](https://cracku.in/15-the-text-uses-the-word-exclosures-because-wild-des-x-cat-2024-slot-1)
\end{enumerate}
\vspace{1em}
\textbf{Solution:}

Option D is the correct answer. 
The 'exclosures' are mentioned as fenced areas cleared of invasive rabbits and feral cats. The term "exclosures" points to the intentional exclusion of these invasive species to create a safe environment for the bandicoots and other native animals.  
  
Option A: While the bandicoots help in restoring the cattle damaged landscape, the term 'exclosure' does not relate to it.  
  
Option B: The exclosures are cleared of feral cats, but the passage does not mention removing large bilbies, which are actually part of the controlled environment. The exclosures themselves are specifically to protect the bandicoots from cats and rabbits, not bilbies.  
  
Option C: The exclosures are not about making an area entirely predator-free for all species. Instead, the purpose is to create controlled environments where invasive species like rabbits and feral cats are removed. Predators are still present in the Wild Training Zone, where bandicoots and other marsupials learn to coexist and evade predators.

\end{problem}

\newpage
\begin{problem}{}{}
% Subject: varc
\textbf{Instructions}

Instructions   

**The passage below is accompanied by four questions. Based on the passage, choose the best answer for each question.**
. . . [T]he idea of craftsmanship is not simply nostalgic. . . . Crafts require distinct skills, an allround approach to work that involves the whole product, rather than individual parts, and an attitude that necessitates devotion to the job and a focus on the communal interest. The concept of craft emphasises the human touch and individual judgment.
Essentially, the crafts concept seems to run against the preponderant ethos of management studies which, as the academics note, have long prioritised efficiency and consistency. . . . Craft skills were portrayed as being primitive and traditionalist.
The contrast between artisanship and efficiency first came to the fore in the 19th century when British manufacturers suddenly faced competition from across the Atlantic as firms developed the “American system” using standardised parts. . . . the worldwide success of the Singer sewing machine showed the potential of a mass-produced device. This process created its own reaction, first in the form of the Arts and Crafts movement of the late 19th century, and then again in the “small is beautiful” movement of the 1970s. A third crafts movement is emerging as people become aware of the environmental impact of conventional industry.
There are two potential markets for those who practise crafts. The first stems from the existence of consumers who are willing to pay a premium price for goods that are deemed to be of extra quality. . . . The second market lies in those consumers who wish to use their purchases to support local workers, or to reduce their environmental impact by taking goods to craftspeople to be mended, or recycled.
For workers, the appeal of craftsmanship is that it allows them the autonomy to make creative choices, and thus makes a job far more satisfying. In that sense, it could offer hope for the overall labour market. Let the machines automate dull and repetitive tasks and let workers focus purely on their skills, judgment and imagination. As a current example, the academics cite the “agile” manifesto in the software sector, an industry at the heart of technological change. The pioneers behind the original agile manifesto promised to prioritise “individuals and interactions over processes and tools”. By bringing together experts from different teams, agile working is designed to improve creativity.
But the broader question is whether crafts can create a lot more jobs than they do today. Demand for crafted products may rise but will it be easy to retrain workers in sectors that might get automated (such as truck drivers) to take advantage? In a world where products and services often have to pass through regulatory hoops, large companies will usually have the advantage.
History also suggests that the link between crafts and creativity is not automatic. Medieval craft guilds were monopolies which resisted new entrants. They were also highly hierarchical with young men required to spend long periods as apprentices and journeymen before they could set up on their own; by that time the innovative spirit may have been knocked out of them. Craft workers can thrive in the modern era, but only if they don’t get too organised.

\vspace{1em}
\textbf{Question 16}

Which one of the following statements provides a gist of this passage?

\begin{enumerate}[label=\Alph*.]
    \item The onslaught of animals, such as cattle, rabbits and housecats, brought in by the British led to the extinction of the western barred bandicoot.
    \item The negligent attitude of the British colonists towards these bandicoots evidenced by the names given to them led to their annihilation.
    \item Marsupials are going extinct due to the colonial era transformation of the ecosystem which also destroyed natural vegetation
    \item A type of bandicoots was nearly wiped out by invasive species but rescuers now pin hopes on a remnant island population.
[](https://cracku.in/16-which-one-of-the-following-statements-provides-a-g-x-cat-2024-slot-1)
\end{enumerate}
\vspace{1em}
\textbf{Solution:}

Option D is the correct answer.  

Option D captures the main idea of the passage. It reflects the near-extinction of the western barred bandicoot due to invasive species and highlights the conservation efforts using survivors from Shark Bay islands.
Option A: This is not entirely true. The western barred bandicoot did not go extinct; instead, it survived in small numbers on two predator-free islands. This option incorrectly asserts total extinction and ignores the ongoing efforts to revive the species.
Option B: This is a distortion. While the colonists' negligence and the nicknames they gave reflect their disregard, the passage clearly attributes the near-extinction of bandicoots to ecological disruptions caused by invasive species, not merely the colonists' attitudes.
Option C: This generalizes the issue and does not focus specifically on the western barred bandicoot, the subject of the passage. Furthermore, it does not highlight the rescue efforts which are central to the passage.

Instructions   

For the following questions answer them individually

\end{problem}

\newpage
\begin{problem}{}{}
% Subject: varc
\textbf{Instructions}

Instructions   

**The passage below is accompanied by four questions. Based on the passage, choose the best answer for each question.**
. . . [T]he idea of craftsmanship is not simply nostalgic. . . . Crafts require distinct skills, an allround approach to work that involves the whole product, rather than individual parts, and an attitude that necessitates devotion to the job and a focus on the communal interest. The concept of craft emphasises the human touch and individual judgment.
Essentially, the crafts concept seems to run against the preponderant ethos of management studies which, as the academics note, have long prioritised efficiency and consistency. . . . Craft skills were portrayed as being primitive and traditionalist.
The contrast between artisanship and efficiency first came to the fore in the 19th century when British manufacturers suddenly faced competition from across the Atlantic as firms developed the “American system” using standardised parts. . . . the worldwide success of the Singer sewing machine showed the potential of a mass-produced device. This process created its own reaction, first in the form of the Arts and Crafts movement of the late 19th century, and then again in the “small is beautiful” movement of the 1970s. A third crafts movement is emerging as people become aware of the environmental impact of conventional industry.
There are two potential markets for those who practise crafts. The first stems from the existence of consumers who are willing to pay a premium price for goods that are deemed to be of extra quality. . . . The second market lies in those consumers who wish to use their purchases to support local workers, or to reduce their environmental impact by taking goods to craftspeople to be mended, or recycled.
For workers, the appeal of craftsmanship is that it allows them the autonomy to make creative choices, and thus makes a job far more satisfying. In that sense, it could offer hope for the overall labour market. Let the machines automate dull and repetitive tasks and let workers focus purely on their skills, judgment and imagination. As a current example, the academics cite the “agile” manifesto in the software sector, an industry at the heart of technological change. The pioneers behind the original agile manifesto promised to prioritise “individuals and interactions over processes and tools”. By bringing together experts from different teams, agile working is designed to improve creativity.
But the broader question is whether crafts can create a lot more jobs than they do today. Demand for crafted products may rise but will it be easy to retrain workers in sectors that might get automated (such as truck drivers) to take advantage? In a world where products and services often have to pass through regulatory hoops, large companies will usually have the advantage.
History also suggests that the link between crafts and creativity is not automatic. Medieval craft guilds were monopolies which resisted new entrants. They were also highly hierarchical with young men required to spend long periods as apprentices and journeymen before they could set up on their own; by that time the innovative spirit may have been knocked out of them. Craft workers can thrive in the modern era, but only if they don’t get too organised.

\vspace{1em}
\textbf{Question 17}

**The passage given below is followed by four alternate summaries. Choose the option that best captures the essence of the passage.**
Cartographers design and create maps to communicate information about phenomena located somewhere on our planet. In the past, cartographers did not worry too much about who was going to read their maps. Although some simple “usability” research was done—like comparing whether circle or bar symbols worked best—cartographers knew how to make maps. This has changed now, however, due to all kinds of societal and technological developments. Today, map readers are more demanding—mostly because of the tools they use to read maps. Cartographers, who are also influenced by these trends, are now more interested in seeing if their products are efficient, effective, and appreciated.

\begin{enumerate}[label=\Alph*.]
    \item Today, cartographers also need to look into the usability of maps because of the new technological developments.
    \item Modern mapmakers evaluate a map’s effectiveness efficiency and satisfaction of the user through a series of experiments
    \item Maps are being used for a variety of reasons and therefore map readers have become more demanding
    \item New technological developments have prompted cartographers to experiment with their maps by applying these new innovations.
[](https://cracku.in/17-the-passage-given-below-is-followed-by-four-altern-x-cat-2024-slot-1)
\end{enumerate}
\vspace{1em}
\textbf{Solution:}

Option A is the correct answer.   
  
The passage emphasizes that cartographers now should pay attention to the usability of maps due to the evolving expectations of map readers. The key point is that technological developments have made users more demanding, leading cartographers to focus on how efficient, effective, and appreciated their maps are.
Option B: While it is true that cartographers are focused on usability, the passage does not mention specific experiments or evaluation methods.
Option C: This option suggests that maps are being used for a variety of reasons, which is not mentioned in the passage. The focus of the passage is on the demanding nature of modern map readers and not on the reasons for which maps are used.
Option D: While new technological developments are mentioned, the passage does not state that cartographers are experimenting with these innovations in their maps.

\end{problem}

\newpage
\begin{problem}{}{}
% Subject: varc
\textbf{Instructions}

Instructions   

**The passage below is accompanied by four questions. Based on the passage, choose the best answer for each question.**
. . . [T]he idea of craftsmanship is not simply nostalgic. . . . Crafts require distinct skills, an allround approach to work that involves the whole product, rather than individual parts, and an attitude that necessitates devotion to the job and a focus on the communal interest. The concept of craft emphasises the human touch and individual judgment.
Essentially, the crafts concept seems to run against the preponderant ethos of management studies which, as the academics note, have long prioritised efficiency and consistency. . . . Craft skills were portrayed as being primitive and traditionalist.
The contrast between artisanship and efficiency first came to the fore in the 19th century when British manufacturers suddenly faced competition from across the Atlantic as firms developed the “American system” using standardised parts. . . . the worldwide success of the Singer sewing machine showed the potential of a mass-produced device. This process created its own reaction, first in the form of the Arts and Crafts movement of the late 19th century, and then again in the “small is beautiful” movement of the 1970s. A third crafts movement is emerging as people become aware of the environmental impact of conventional industry.
There are two potential markets for those who practise crafts. The first stems from the existence of consumers who are willing to pay a premium price for goods that are deemed to be of extra quality. . . . The second market lies in those consumers who wish to use their purchases to support local workers, or to reduce their environmental impact by taking goods to craftspeople to be mended, or recycled.
For workers, the appeal of craftsmanship is that it allows them the autonomy to make creative choices, and thus makes a job far more satisfying. In that sense, it could offer hope for the overall labour market. Let the machines automate dull and repetitive tasks and let workers focus purely on their skills, judgment and imagination. As a current example, the academics cite the “agile” manifesto in the software sector, an industry at the heart of technological change. The pioneers behind the original agile manifesto promised to prioritise “individuals and interactions over processes and tools”. By bringing together experts from different teams, agile working is designed to improve creativity.
But the broader question is whether crafts can create a lot more jobs than they do today. Demand for crafted products may rise but will it be easy to retrain workers in sectors that might get automated (such as truck drivers) to take advantage? In a world where products and services often have to pass through regulatory hoops, large companies will usually have the advantage.
History also suggests that the link between crafts and creativity is not automatic. Medieval craft guilds were monopolies which resisted new entrants. They were also highly hierarchical with young men required to spend long periods as apprentices and journeymen before they could set up on their own; by that time the innovative spirit may have been knocked out of them. Craft workers can thrive in the modern era, but only if they don’t get too organised.

\vspace{1em}
\textbf{Question 18}

**There is a sentence that is missing in the paragraph below. Look at the paragraph and decide where (option 1, 2, 3, or 4) the following sentence would best fit.**
**Sentence:** The brain isn’t organized the way you might set up your home office or bathroom medicine cabinet.
**Paragraph** : ___(1)___. You can’t just put things anywhere you want to. The evolved architecture of the brain is haphazard and disjointed, and incorporates multiple systems, each of which has a mind of its own. ___(2)___. Evolution doesn’t design things and it doesn’t build systems—it settles on systems that, historically, conveyed a survival benefit. There is no overarching, grand planner engineering the systems so that they work harmoniously together. ___(3)___. The brain is more like a big, old house with piecemeal renovations done on every floor, and less like new construction. ___(4)___.

\begin{enumerate}[label=\Alph*.]
    \item Option 4
    \item Option 1
    \item Option 2
    \item Option 3
[](https://cracku.in/18-there-is-a-sentence-that-is-missing-in-the-paragra-x-cat-2024-slot-1)
\end{enumerate}
\vspace{1em}
\textbf{Solution:}

Let’s evaluate each option:**_  
  
Option 1_** : The given sentence works well as an introduction. It sets the stage by comparing the brain’s organization to familiar spaces like a home office or medicine cabinet. This analogy smoothly transitions into the next sentence about the brain's "haphazard and disjointed" architecture, making it a natural fit here.**_  
  
Option 2_** : Inserting the sentence here would disrupt the flow. This part focuses on the brain's evolved systems and how evolution works. The analogy to home offices feels out of place and does not connect coherently with the subsequent sentences._**  
  
Option 3**_ : This gap already includes a metaphor—the brain as "a big, old house with piecemeal renovations." Adding the sentence here would be redundant, as it repeats the idea of disorganization without adding new insight._**  
  
Option 4**_ : This section discusses how evolution doesn’t design harmonious systems. The sentence about the brain’s organization doesn’t align well with this abstract explanation, making it feel disjointed and out of context.  
**  
Option 1 is the best fit** as it introduces the idea clearly, aligns with the paragraph's theme, and transitions seamlessly into the discussion of the brain's complexity.

\end{problem}

\newpage
\begin{problem}{}{}
% Subject: varc
\textbf{Instructions}

Instructions   

**The passage below is accompanied by four questions. Based on the passage, choose the best answer for each question.**
. . . [T]he idea of craftsmanship is not simply nostalgic. . . . Crafts require distinct skills, an allround approach to work that involves the whole product, rather than individual parts, and an attitude that necessitates devotion to the job and a focus on the communal interest. The concept of craft emphasises the human touch and individual judgment.
Essentially, the crafts concept seems to run against the preponderant ethos of management studies which, as the academics note, have long prioritised efficiency and consistency. . . . Craft skills were portrayed as being primitive and traditionalist.
The contrast between artisanship and efficiency first came to the fore in the 19th century when British manufacturers suddenly faced competition from across the Atlantic as firms developed the “American system” using standardised parts. . . . the worldwide success of the Singer sewing machine showed the potential of a mass-produced device. This process created its own reaction, first in the form of the Arts and Crafts movement of the late 19th century, and then again in the “small is beautiful” movement of the 1970s. A third crafts movement is emerging as people become aware of the environmental impact of conventional industry.
There are two potential markets for those who practise crafts. The first stems from the existence of consumers who are willing to pay a premium price for goods that are deemed to be of extra quality. . . . The second market lies in those consumers who wish to use their purchases to support local workers, or to reduce their environmental impact by taking goods to craftspeople to be mended, or recycled.
For workers, the appeal of craftsmanship is that it allows them the autonomy to make creative choices, and thus makes a job far more satisfying. In that sense, it could offer hope for the overall labour market. Let the machines automate dull and repetitive tasks and let workers focus purely on their skills, judgment and imagination. As a current example, the academics cite the “agile” manifesto in the software sector, an industry at the heart of technological change. The pioneers behind the original agile manifesto promised to prioritise “individuals and interactions over processes and tools”. By bringing together experts from different teams, agile working is designed to improve creativity.
But the broader question is whether crafts can create a lot more jobs than they do today. Demand for crafted products may rise but will it be easy to retrain workers in sectors that might get automated (such as truck drivers) to take advantage? In a world where products and services often have to pass through regulatory hoops, large companies will usually have the advantage.
History also suggests that the link between crafts and creativity is not automatic. Medieval craft guilds were monopolies which resisted new entrants. They were also highly hierarchical with young men required to spend long periods as apprentices and journeymen before they could set up on their own; by that time the innovative spirit may have been knocked out of them. Craft workers can thrive in the modern era, but only if they don’t get too organised.

\vspace{1em}
\textbf{Question 19}

**The passage given below is followed by four alternate summaries. Choose the option that best captures the essence of the passage.**
Certain codes may, of course, be so widely distributed in a specific language community or culture, and be learned at so early an age, that they appear not to be constructed - the effect of an articulation between sign and referent - but to be ‘naturally’ given. Simple visual signs appear to have achieved a ‘near-universality’ in this sense: though evidence remains that even apparently ‘natural’ visual codes are culture specific. However, this does not mean that no codes have intervened; rather, that the codes have been profoundly naturalized. The operation of naturalized codes reveals not the transparency and ‘naturalness’ of language but the depth, the habituation and the near-universality of the codes in use. They produce apparently ‘natural’ recognitions. This has the (ideological) effect of concealing the practices of coding which are present.

\begin{enumerate}[label=\Alph*.]
    \item Learning linguistic and visual signs at an early age makes all such codes appear natural. This naturalization of codes is the effect of ideology.
    \item Not all codes are natural but certain codes are naturalized and made to appear universal. Ideology aims to hide the mechanism of coding behind signs.
    \item Language and visual signs are codes. However, some of the codes are so widespread that they not only seem naturally given but also hide the mechanism of coding behind the signs.
    \item All codes, linguistic and visual, have a natural origin but some are so widespread that they become universal. This is what hides the mechanism of coding behind signs.
[](https://cracku.in/19-the-passage-given-below-is-followed-by-four-altern-x-cat-2024-slot-1)
\end{enumerate}
\vspace{1em}
\textbf{Solution:}

Option C is the correct answer. 
This option captures the main idea of the passage that some codes, like language and visual signs, are so commonly used that they appear natural and conceal the process of how they were created.
Option A: The passage does not suggest that early learning is why codes appear natural. The cause-and-effect relationship is incorrectly stated here. 
Option B: This option misinterprets two key aspects of the passage. First, the idea that certain codes are "made to appear universal" is somewhat misleading because the passage doesn't claim that codes are deliberately made universal; instead, it describes how codes, through habituation and widespread use, come to feel "natural". Second, the phrase "Ideology aims to hide the mechanism of coding" is not supported by the passage. The passage suggests that the naturalization of codes leads to the illusion of transparency and naturalness, which conceals the mechanisms of coding, but it doesn't explicitly discuss ideology as a force that intentionally hides these mechanisms
Option D: This option is incorrect because the passage doesn’t claim that all codes have a natural origin. It states that codes become naturalized through use, not that they were naturally originating from the start.

\end{problem}

\newpage
\begin{problem}{}{}
% Subject: varc
\textbf{Instructions}

Instructions   

**The passage below is accompanied by four questions. Based on the passage, choose the best answer for each question.**
. . . [T]he idea of craftsmanship is not simply nostalgic. . . . Crafts require distinct skills, an allround approach to work that involves the whole product, rather than individual parts, and an attitude that necessitates devotion to the job and a focus on the communal interest. The concept of craft emphasises the human touch and individual judgment.
Essentially, the crafts concept seems to run against the preponderant ethos of management studies which, as the academics note, have long prioritised efficiency and consistency. . . . Craft skills were portrayed as being primitive and traditionalist.
The contrast between artisanship and efficiency first came to the fore in the 19th century when British manufacturers suddenly faced competition from across the Atlantic as firms developed the “American system” using standardised parts. . . . the worldwide success of the Singer sewing machine showed the potential of a mass-produced device. This process created its own reaction, first in the form of the Arts and Crafts movement of the late 19th century, and then again in the “small is beautiful” movement of the 1970s. A third crafts movement is emerging as people become aware of the environmental impact of conventional industry.
There are two potential markets for those who practise crafts. The first stems from the existence of consumers who are willing to pay a premium price for goods that are deemed to be of extra quality. . . . The second market lies in those consumers who wish to use their purchases to support local workers, or to reduce their environmental impact by taking goods to craftspeople to be mended, or recycled.
For workers, the appeal of craftsmanship is that it allows them the autonomy to make creative choices, and thus makes a job far more satisfying. In that sense, it could offer hope for the overall labour market. Let the machines automate dull and repetitive tasks and let workers focus purely on their skills, judgment and imagination. As a current example, the academics cite the “agile” manifesto in the software sector, an industry at the heart of technological change. The pioneers behind the original agile manifesto promised to prioritise “individuals and interactions over processes and tools”. By bringing together experts from different teams, agile working is designed to improve creativity.
But the broader question is whether crafts can create a lot more jobs than they do today. Demand for crafted products may rise but will it be easy to retrain workers in sectors that might get automated (such as truck drivers) to take advantage? In a world where products and services often have to pass through regulatory hoops, large companies will usually have the advantage.
History also suggests that the link between crafts and creativity is not automatic. Medieval craft guilds were monopolies which resisted new entrants. They were also highly hierarchical with young men required to spend long periods as apprentices and journeymen before they could set up on their own; by that time the innovative spirit may have been knocked out of them. Craft workers can thrive in the modern era, but only if they don’t get too organised.

\vspace{1em}
\textbf{Question 20}

**Five jumbled up sentences (labelled 1, 2, 3, 4 and 5), related to a topic, are given below. Four of them can be put together to form a coherent paragraph. Identify the odd sentence and key in the number of that sentence as your answer.**
1. Animals have an interest in fulfilling their basic needs, but also in avoiding suffering, and thus we ought to extend moral consideration.  
2. Singer viewed himself as a utilitarian, and presents a direct moral theory concerning animal rights, in contrast to indirect positions, such as welfarist views.  
3. He argued for extending moral consideration to animals because, similar to humans, animals have certain significant interests.  
4. The event that publicly announced animal rights as a legitimate issue within contemporary philosophy was Peter Singer’s Animal Liberation text in 1975.  
5. As such, we ought to view their interests alongside and equal to human interests, which results in humans having direct moral duties towards animals.
Backspace  
789  
456  
123  
0.-  
Clear All  

Submit
[](https://cracku.in/20-five-jumbled-up-sentences-labelled-1-2-3-4-and-5-r-x-cat-2024-slot-1)

\begin{enumerate}[label=\Alph*.]
\end{enumerate}
\vspace{1em}
\textbf{Solution:}

**Sentence 1 is the odd sentence.**  
  
The correct sequence is 4-2-3-5.**  
  
Sentence 4** introduces the historical significance of Peter Singer’s _Animal Liberation_ , which set the stage for the discussion of animal rights in contemporary philosophy.  
**Sentence 2** expands on Singer’s philosophy, explaining his utilitarian view and how it contrasts with other indirect moral views on animal rights.  
**Sentence 3** then explains the core of Singer’s argument, that animals, like humans, have significant interests and deserve moral consideration.**  
Sentence 5** concludes Singer's argument by stating that humans should treat animal interests equally to human interests, leading to moral duties towards animals.  
**  
Sentence 1 is the odd one out** because it does not focus on Peter Singer's views on animal rights and the moral obligations humans have towards animals, which is the main theme of the paragraph.

Instructions   

**The passage below is accompanied by four questions. Based on the passage, choose the best answer for each question.**
In the summer of 2022, subscribers to the US streaming service HBO MAX were alarmed to discover that dozens of the platform’s offerings - from the Covid-themed heist thriller Locked Down to the recent remake of The Witches - had been quietly removed from the service . . . The news seemed like vindication to those who had long warned that streaming was more about controlling access to the cultural commons than expanding it, as did reports (since denied by the show’s creators) that Netflix had begun editing old episodes of Stranger Things to retroactively improve their visual effects.
What’s less clear is whether the commonly prescribed cure for these cultural ills - a return to the material pleasures of physical media - is the right one. While the makers of Blu-ray discs claim they have a shelf life of 100 years, such statistics remain largely theoretical until they come to pass, and are dependent on storage conditions, not to mention the continued availability of playback equipment. The humble DVD has already proved far less resilient, with many early releases already beginning to deteriorate in quality Digital movie purchases provide even less security. Any film “bought” on iTunes could disappear if you move to another territory with a different rights agreement and try to redownload it. It’s a bold new frontier in the commodification of art: the birth of the product recall. After a man took to Twitter to bemoan losing access to Cars 2 after moving from Canada to Australia, Apple clarified that users who downloaded films to their devices would retain permanent access to those downloads, even if they relocated to a hemisphere where the [content was] subject to a different set of rights agreements. Thanks to the company’s ironclad digital rights management technology, however, such files cannot be moved or backed up, locking you into watching with your Apple account.
Anyone who does manage to acquire Digital Rights Management free (DRM-free) copies of their favourite films must nonetheless grapple with ever-changing file format standards, not to mention data decay - the gradual process by which electronic information slowly but surely corrupts. Only the regular migration of files from hard drive to hard drive can delay the inevitable, in a sisyphean battle against the ravages of digital time.
In a sense, none of this is new. Charlie Chaplin burned the negative of his 1926 film A Woman of the Sea as a tax write-off. Many more films have been lost through accident, negligence or plain indifference. During a heatwave in July 1937, a Fox film vault in New Jersey burned down, destroying a majority of the silent films produced by the studio.
Back then, at least, cinema was defined by its ephemerality: the sense that a film was as good as gone once it left your local cinema. Today, with film studios keen to stress the breadth of their back catalogues (or to put in Hollywood terms, the value of their IPs), audiences may start to wonder why those same studios seem happy to set the vault alight themselves if it’ll help next quarter’s numbers.

\end{problem}

\newpage
\begin{problem}{}{}
% Subject: varc
\textbf{Instructions}

Instructions   

**The passage below is accompanied by four questions. Based on the passage, choose the best answer for each question.**
. . . [T]he idea of craftsmanship is not simply nostalgic. . . . Crafts require distinct skills, an allround approach to work that involves the whole product, rather than individual parts, and an attitude that necessitates devotion to the job and a focus on the communal interest. The concept of craft emphasises the human touch and individual judgment.
Essentially, the crafts concept seems to run against the preponderant ethos of management studies which, as the academics note, have long prioritised efficiency and consistency. . . . Craft skills were portrayed as being primitive and traditionalist.
The contrast between artisanship and efficiency first came to the fore in the 19th century when British manufacturers suddenly faced competition from across the Atlantic as firms developed the “American system” using standardised parts. . . . the worldwide success of the Singer sewing machine showed the potential of a mass-produced device. This process created its own reaction, first in the form of the Arts and Crafts movement of the late 19th century, and then again in the “small is beautiful” movement of the 1970s. A third crafts movement is emerging as people become aware of the environmental impact of conventional industry.
There are two potential markets for those who practise crafts. The first stems from the existence of consumers who are willing to pay a premium price for goods that are deemed to be of extra quality. . . . The second market lies in those consumers who wish to use their purchases to support local workers, or to reduce their environmental impact by taking goods to craftspeople to be mended, or recycled.
For workers, the appeal of craftsmanship is that it allows them the autonomy to make creative choices, and thus makes a job far more satisfying. In that sense, it could offer hope for the overall labour market. Let the machines automate dull and repetitive tasks and let workers focus purely on their skills, judgment and imagination. As a current example, the academics cite the “agile” manifesto in the software sector, an industry at the heart of technological change. The pioneers behind the original agile manifesto promised to prioritise “individuals and interactions over processes and tools”. By bringing together experts from different teams, agile working is designed to improve creativity.
But the broader question is whether crafts can create a lot more jobs than they do today. Demand for crafted products may rise but will it be easy to retrain workers in sectors that might get automated (such as truck drivers) to take advantage? In a world where products and services often have to pass through regulatory hoops, large companies will usually have the advantage.
History also suggests that the link between crafts and creativity is not automatic. Medieval craft guilds were monopolies which resisted new entrants. They were also highly hierarchical with young men required to spend long periods as apprentices and journeymen before they could set up on their own; by that time the innovative spirit may have been knocked out of them. Craft workers can thrive in the modern era, but only if they don’t get too organised.

\vspace{1em}
\textbf{Question 21}

Which one of the following statements about art best captures the arguments made in the passage?

\begin{enumerate}[label=\Alph*.]
    \item In the age of online subscription services, it is time to change our understanding of classic works of art being primarily immutable and easily available to the public
    \item As art is increasingly created, stored and distributed digitally, access to it is counterintuitively likely to be made more difficult by the rapid churn in technology and the whims of host platforms.
    \item Accepting retroactive changes to works of art is dangerous because it will encourage creators to not put enough effort into the original attempt, given that they can always edit or update their work later
    \item Works of art belong to the cultural commons and hence must remain available in perpetuity, irrespective of who pays for access to them.
[](https://cracku.in/21-which-one-of-the-following-statements-about-art-be-x-cat-2024-slot-1)
\end{enumerate}
\vspace{1em}
\textbf{Solution:}

Option B is the correct answer.   
  
The passage argues that, despite the advances in digital distribution and storage of films (via streaming services or digital purchases), access to art is becoming more fragile. It mentions how content can disappear from platforms, how digital files deteriorate over time, and how rights agreements can limit access based on geographical location. The idea that technology and platform control lead to difficulties in maintaining access captured in Option B 
Option A: The passage does not advocate changing the understanding of art as immutable or easily available. Hence, this is wrong.
Option C: This is an overstatement of the passage’s argument. While the passage briefly touches on the idea of retroactive changes to works like Stranger Things, it does not present these changes as inherently "dangerous." 
Option D: The passage does not argue for the availability of art in the cultural commons in perpetuity. Instead, it highlights how access is controlled by platforms and technological challenges rather than making a broader ideological statement about cultural commons.

\end{problem}

\newpage
\begin{problem}{}{}
% Subject: varc
\textbf{Instructions}

Instructions   

**The passage below is accompanied by four questions. Based on the passage, choose the best answer for each question.**
. . . [T]he idea of craftsmanship is not simply nostalgic. . . . Crafts require distinct skills, an allround approach to work that involves the whole product, rather than individual parts, and an attitude that necessitates devotion to the job and a focus on the communal interest. The concept of craft emphasises the human touch and individual judgment.
Essentially, the crafts concept seems to run against the preponderant ethos of management studies which, as the academics note, have long prioritised efficiency and consistency. . . . Craft skills were portrayed as being primitive and traditionalist.
The contrast between artisanship and efficiency first came to the fore in the 19th century when British manufacturers suddenly faced competition from across the Atlantic as firms developed the “American system” using standardised parts. . . . the worldwide success of the Singer sewing machine showed the potential of a mass-produced device. This process created its own reaction, first in the form of the Arts and Crafts movement of the late 19th century, and then again in the “small is beautiful” movement of the 1970s. A third crafts movement is emerging as people become aware of the environmental impact of conventional industry.
There are two potential markets for those who practise crafts. The first stems from the existence of consumers who are willing to pay a premium price for goods that are deemed to be of extra quality. . . . The second market lies in those consumers who wish to use their purchases to support local workers, or to reduce their environmental impact by taking goods to craftspeople to be mended, or recycled.
For workers, the appeal of craftsmanship is that it allows them the autonomy to make creative choices, and thus makes a job far more satisfying. In that sense, it could offer hope for the overall labour market. Let the machines automate dull and repetitive tasks and let workers focus purely on their skills, judgment and imagination. As a current example, the academics cite the “agile” manifesto in the software sector, an industry at the heart of technological change. The pioneers behind the original agile manifesto promised to prioritise “individuals and interactions over processes and tools”. By bringing together experts from different teams, agile working is designed to improve creativity.
But the broader question is whether crafts can create a lot more jobs than they do today. Demand for crafted products may rise but will it be easy to retrain workers in sectors that might get automated (such as truck drivers) to take advantage? In a world where products and services often have to pass through regulatory hoops, large companies will usually have the advantage.
History also suggests that the link between crafts and creativity is not automatic. Medieval craft guilds were monopolies which resisted new entrants. They were also highly hierarchical with young men required to spend long periods as apprentices and journeymen before they could set up on their own; by that time the innovative spirit may have been knocked out of them. Craft workers can thrive in the modern era, but only if they don’t get too organised.

\vspace{1em}
\textbf{Question 22}

Which one of the following statements, if true, would best invalidate the main argument of the passage?

\begin{enumerate}[label=\Alph*.]
    \item Recent research has irrefutably proven that Blu-Ray discs have a shelf life of at least 100 years.
    \item Studios and streaming services have committed to giving customers perpetual and platform independent access to the original digital content they have paid for.
    \item When moving to a different geographical location, customers can easily use Virtual Private Networks (VPNs) to bypass geo-blocking and regain access to their content on any streaming service.
    \item Improved cloud storage services have made it possible for movie collections to now be preserved in perpetuity, without the need to keep migrating the files.
[](https://cracku.in/22-which-one-of-the-following-statements-if-true-woul-x-cat-2024-slot-1)
\end{enumerate}
\vspace{1em}
\textbf{Solution:}

Option B is the correct answer.   
  
Option B would invalidate the main argument because it directly addresses the issue raised in the passage, i.e., the lack of permanent access to digital content. The passage highlights concerns about the temporary and restricted nature of digital ownership. If studios and streaming services committed to providing perpetual and platform-independent access, it would resolve the problem of content being removed or restricted, making the author's argument about the instability of digital media irrelevant.
Option A: This would not invalidate the argument because the passage mentions that Blu-ray discs have a theoretical shelf life but acknowledges that their durability depends on storage conditions and the availability of playback equipment.
Option C: This option doesn't directly invalidate the argument either. While VPNs might help users bypass geo-restrictions, it doesn't address the broader issue of digital ownership and the fragility of digital rights, especially the fear of losing access to content due to different rights agreements. The passage focuses on the unreliability and restrictions of digital ownership, not just geographical access.
Option D: While this option touches on the potential for preserving digital content, it doesn't directly address the problem raised in the passage: the lack of permanent, independent access to content.

\end{problem}

\newpage
\begin{problem}{}{}
% Subject: varc
\textbf{Instructions}

Instructions   

**The passage below is accompanied by four questions. Based on the passage, choose the best answer for each question.**
. . . [T]he idea of craftsmanship is not simply nostalgic. . . . Crafts require distinct skills, an allround approach to work that involves the whole product, rather than individual parts, and an attitude that necessitates devotion to the job and a focus on the communal interest. The concept of craft emphasises the human touch and individual judgment.
Essentially, the crafts concept seems to run against the preponderant ethos of management studies which, as the academics note, have long prioritised efficiency and consistency. . . . Craft skills were portrayed as being primitive and traditionalist.
The contrast between artisanship and efficiency first came to the fore in the 19th century when British manufacturers suddenly faced competition from across the Atlantic as firms developed the “American system” using standardised parts. . . . the worldwide success of the Singer sewing machine showed the potential of a mass-produced device. This process created its own reaction, first in the form of the Arts and Crafts movement of the late 19th century, and then again in the “small is beautiful” movement of the 1970s. A third crafts movement is emerging as people become aware of the environmental impact of conventional industry.
There are two potential markets for those who practise crafts. The first stems from the existence of consumers who are willing to pay a premium price for goods that are deemed to be of extra quality. . . . The second market lies in those consumers who wish to use their purchases to support local workers, or to reduce their environmental impact by taking goods to craftspeople to be mended, or recycled.
For workers, the appeal of craftsmanship is that it allows them the autonomy to make creative choices, and thus makes a job far more satisfying. In that sense, it could offer hope for the overall labour market. Let the machines automate dull and repetitive tasks and let workers focus purely on their skills, judgment and imagination. As a current example, the academics cite the “agile” manifesto in the software sector, an industry at the heart of technological change. The pioneers behind the original agile manifesto promised to prioritise “individuals and interactions over processes and tools”. By bringing together experts from different teams, agile working is designed to improve creativity.
But the broader question is whether crafts can create a lot more jobs than they do today. Demand for crafted products may rise but will it be easy to retrain workers in sectors that might get automated (such as truck drivers) to take advantage? In a world where products and services often have to pass through regulatory hoops, large companies will usually have the advantage.
History also suggests that the link between crafts and creativity is not automatic. Medieval craft guilds were monopolies which resisted new entrants. They were also highly hierarchical with young men required to spend long periods as apprentices and journeymen before they could set up on their own; by that time the innovative spirit may have been knocked out of them. Craft workers can thrive in the modern era, but only if they don’t get too organised.

\vspace{1em}
\textbf{Question 23}

Which of the following statements is suggested by the sentence “Back then, at least, cinema was defined by its ephemerality: the sense that a film was as good as gone once it left your local cinema”?

\begin{enumerate}[label=\Alph*.]
    \item Around a century ago, people were more accepting of not having access to films once they left the local cinema.
    \item Today, films are expected to be available for a long time, since they are no longer tied solely to their stay at the local cinema
    \item Cinema is now no longer as ephemeral as it used to be earlier, because the technology used for creating and preserving films has improved manifold
    \item Presently, there is no reason why film studios should remove access to films once they have left the local cinema
[](https://cracku.in/23-which-of-the-following-statements-is-suggested-by--x-cat-2024-slot-1)
\end{enumerate}
\vspace{1em}
\textbf{Solution:}

Option B is the correct answer. 
The passage contrasts the past and present by mentioning that, in the past, films were considered “as good as gone” once they left the cinema, meaning they were ephemeral and not readily accessible afterwards. Whereas, today's audience expects ongoing access to films well beyond their initial cinema run, thanks to technological advancements like streaming services and digital media. This shift in expectations is what the passage implies when referencing the previous era's ephemerality versus today's more lasting availability.
Option A: This is not the main point. While it may be true that people accepted films as temporary, the passage emphasizes today's expectations rather than discussing past acceptance.
Option C: The passage does not mention technology improvements. It focuses more on audience expectations or belief that films should now be available beyond just the cinema.
Option D: While the passage suggests that audiences may expect films to remain accessible, it does not claim there is no reason for studios to remove access. The passage acknowledges that financial motives may lead to films being removed from platforms.

\end{problem}

\newpage
\begin{problem}{}{}
% Subject: varc
\textbf{Instructions}

Instructions   

**The passage below is accompanied by four questions. Based on the passage, choose the best answer for each question.**
. . . [T]he idea of craftsmanship is not simply nostalgic. . . . Crafts require distinct skills, an allround approach to work that involves the whole product, rather than individual parts, and an attitude that necessitates devotion to the job and a focus on the communal interest. The concept of craft emphasises the human touch and individual judgment.
Essentially, the crafts concept seems to run against the preponderant ethos of management studies which, as the academics note, have long prioritised efficiency and consistency. . . . Craft skills were portrayed as being primitive and traditionalist.
The contrast between artisanship and efficiency first came to the fore in the 19th century when British manufacturers suddenly faced competition from across the Atlantic as firms developed the “American system” using standardised parts. . . . the worldwide success of the Singer sewing machine showed the potential of a mass-produced device. This process created its own reaction, first in the form of the Arts and Crafts movement of the late 19th century, and then again in the “small is beautiful” movement of the 1970s. A third crafts movement is emerging as people become aware of the environmental impact of conventional industry.
There are two potential markets for those who practise crafts. The first stems from the existence of consumers who are willing to pay a premium price for goods that are deemed to be of extra quality. . . . The second market lies in those consumers who wish to use their purchases to support local workers, or to reduce their environmental impact by taking goods to craftspeople to be mended, or recycled.
For workers, the appeal of craftsmanship is that it allows them the autonomy to make creative choices, and thus makes a job far more satisfying. In that sense, it could offer hope for the overall labour market. Let the machines automate dull and repetitive tasks and let workers focus purely on their skills, judgment and imagination. As a current example, the academics cite the “agile” manifesto in the software sector, an industry at the heart of technological change. The pioneers behind the original agile manifesto promised to prioritise “individuals and interactions over processes and tools”. By bringing together experts from different teams, agile working is designed to improve creativity.
But the broader question is whether crafts can create a lot more jobs than they do today. Demand for crafted products may rise but will it be easy to retrain workers in sectors that might get automated (such as truck drivers) to take advantage? In a world where products and services often have to pass through regulatory hoops, large companies will usually have the advantage.
History also suggests that the link between crafts and creativity is not automatic. Medieval craft guilds were monopolies which resisted new entrants. They were also highly hierarchical with young men required to spend long periods as apprentices and journeymen before they could set up on their own; by that time the innovative spirit may have been knocked out of them. Craft workers can thrive in the modern era, but only if they don’t get too organised.

\vspace{1em}
\textbf{Question 24}

“Netflix had begun editing old episodes of Stranger Things to retroactively improve their visual effects.” What is the purpose of this example used in the passage?

\begin{enumerate}[label=\Alph*.]
    \item To show that streaming services are controlling access to the cultural commons rather than expanding it.
    \item To show how unsubstantiated reports are leading to an increase in the level of distrust towards streaming services
    \item To show a practice that justifies the fears of people who feel streaming services cannot be trusted to be custodians of cultural artefacts like film.
    \item To show that art in the digital age, specifically film, is no longer sacrosanct, and may be changed to suit changing tastes or technology
[](https://cracku.in/24-netflix-had-begun-editing-old-episodes-of-stranger-x-cat-2024-slot-1)
\end{enumerate}
\vspace{1em}
\textbf{Solution:}

Option C is the correct answer. 
The passage highlights that the practice of streaming services, like Netflix, editing old episodes of Stranger Things retroactively raised concerns. Altering a popular show’s content without consent or transparency supports the concern that such platforms can tamper with or erase parts of culture at their discretion rather than preserve them as they were originally created.
Option A: This option doesn't fully address the specific concern raised by the Stranger Things editing example. The example focuses on altering the content, not controlling access. 
Option B: The example concerns the possibility and practice of editing content on streaming platforms, not whether unsubstantiated reports cause distrust.
Option D: The example of Stranger Things is not necessarily about changing films to suit new tastes or technology. Rather, it’s about the platform’s ability to alter existing content without transparency or input from the original creators. 
[](https://cracku.in/downloads/19318)
  *  
  *

\end{problem}

\newpage
\begin{problem}{}{}
% Subject: varc
\textbf{Instructions}

Instructions   

For the following questions answer them individually

\vspace{1em}
\textbf{Question 1}

**There is a sentence that is missing in the paragraph below. Look at the paragraph and decide where (option 1, 2, 3, or 4) the following sentence would best fit.**
**Sentence** : Science has officially crowned us superior to our early-rising brethren.
**Paragraph** : My fellow night owls, grab a strong cup of coffee and gather around: I have great news. ___(1)___. For a long time, our kind has been unfairly maligned. Stereotyped as lazy and undisciplined. Told we ought to be morning larks. Advised to go to bed early so we can wake before 5am and run a marathon before breakfast like all high-flyers seem to do. Now, however, we are having the last laugh. ___(2)___. It may be a tad more complicated than that. A study published last week, which you may have already seen while scrolling at 1am, suggests that staying up late could be good for brain power. ___(3)___. Is this study a thinly veiled PR exercise conducted by a caffeine-pill company? Nope, it’s legit. ___(4)___. Research led by academics at Imperial College London studied data on more than 26,000 people and found that “self-declared ‘night owls’ generally tend to have higher cognitive scores”.

\begin{enumerate}[label=\Alph*.]
    \item Option 4
    \item Option 3
    \item Option 1
    \item Option 2
[](https://cracku.in/1-there-is-a-sentence-that-is-missing-in-the-paragra-x-cat-2024-slot-2)
\end{enumerate}
\vspace{1em}
\textbf{Solution:}

**D) Option 2 is the correct answer.  
**  
The sentence before option 2 states "Now, however, we are having the last laugh". The reason they are having the last laugh is given in the sentence,  _"Science has officially crowned us superior to our early-rising brethren."_ Therefore, it fits here perfectly forming a coherent paragraph.  
  
_A) Option 4_ : Placing the sentence here would seem a bit disconnected. The sentence before and after Blank 4 flows continuously and placing the sentence here would break this flow  
  
 _B) Option 3_ : This placement would disrupt the flow of the sentence before and after Blank 3. The study mentioned before, blank 3, is the moment where the night owl's position is validated, so introducing the idea of "superiority" will be out of place and feel redundant.  
  
_C) Option 1_ : It may seem plausible initially, but it would not be right to talk about the claim without first mentioning the negative stereotypes, making this placement less effective.

\end{problem}

\newpage
\begin{problem}{}{}
% Subject: varc
\textbf{Instructions}

Instructions   

For the following questions answer them individually

\vspace{1em}
\textbf{Question 2}

**Five jumbled up sentences (labelled 1, 2, 3, 4 and 5), related to a topic, are given below. Four of them can be put together to form a coherent paragraph. Identify the odd sentence and key in the number of that sentence as your answer.**
1. No known real researcher of human behaviour would say that gender is all nature or all nurture.  
2. The evidence for a biological basis for gender certainly doesn’t mean we should be complacent in the face of sexism.  
3. Many people are uncomfortable with the idea that gender is not purely a social construct.  
4. Despite this empirical truth, researchers who study the biological basis of gender often face political pushback.  
5. There’s a political preference for gender to be only a reflection of social factors and so entirely malleable.
Backspace  
789  
456  
123  
0.-  
Clear All  

Submit
[](https://cracku.in/2-five-jumbled-up-sentences-labelled-1-2-3-4-and-5-r-x-cat-2024-slot-2)

\begin{enumerate}[label=\Alph*.]
\end{enumerate}
\vspace{1em}
\textbf{Solution:}

Sentence 2 is the odd one out.
Sequence 1435 forms a coherent paragraph:  
**  
Sentence 1** introduces the debate about gender being influenced by both biological and social factors. It sets the tone for the paragraph, addressing the complexity of gender.**  
Sentence 4** follows logically from Sentence 1. After introducing the complexity of gender, it suggests that, even with the acknowledgment of both nature and nurture, researchers studying gender's biological basis still face political challenges.  
**Sentence 3** builds on Sentence 4, highlighting why researchers studying the biological aspects of gender face resistance. The public discomfort with viewing gender as more than just a social construct is explained.**  
Sentence 5** connects to Sentence 3 by elaborating on the political preference for viewing gender purely as a social construct, reinforcing the challenges faced by researchers advocating for a more nuanced view of gender.  
**  
Sentence 2 doesn't fit in the context of the paragraph**. The rest of the sentences are focused on the complexities of gender as a mix of nature and nurture and the political and societal challenges researchers face. Sentence 2 shifts the focus to the issue of sexism, which, while related to gender, does not directly tie into the previous sentences discussing the nature vs. nurture debate and the political resistance to biological explanations.

\end{problem}

\newpage
\begin{problem}{}{}
% Subject: varc
\textbf{Instructions}

Instructions   

For the following questions answer them individually

\vspace{1em}
\textbf{Question 3}

**There is a sentence that is missing in the paragraph below. Look at the paragraph and decide where (option 1, 2, 3, or 4) the following sentence would best fit.**
**Sentence** : [T]he Europeans did not invent globalization.
**Paragraph** : The first phase of globalization occurred long before the introduction of either steam or electric power…Chinese consumers at all social levels consumed vast quantities of spices, fragrant woods and unusual plants. The peoples of Southeast Asia who lived in forests gave up their traditional livelihoods and completely reoriented their economies to supply Chinese consumers….___(1)___. These exchanges of the year 1000 opened some of the routes through which goods and peoples continued to travel after Columbus traversed the mid- Atlantic. ___(2)___. Yet the world of 1000 differed from that of 1492 in important ways….the travellers who encountered one another in the year 1000 were much closer technologically. ___(3)___. They changed and augmented what was already there since 1000. ___(4)___. If globalization hadn’t yet begun, Europeans wouldn’t have been able to penetrate the markets in so many places as quickly as they did after 1492.

\begin{enumerate}[label=\Alph*.]
    \item Option 1
    \item Option 4
    \item Option 2
    \item Option 3
[](https://cracku.in/3-there-is-a-sentence-that-is-missing-in-the-paragra-x-cat-2024-slot-2)
\end{enumerate}
\vspace{1em}
\textbf{Solution:}

We can attempt to place the given sentence by breaking the passage down into three segments -
The discussion till Blank (2) can be considered the first segment: it introduces the idea that globalization was already underway well before the industrial era. It highlights the economic interconnectedness between Chinese consumers and Southeast Asian suppliers, emphasizing that trade networks and economic reorientations were already shaping societies. The missing sentence does not fit here in Blanks (1) or (2) because the focus is on early economic exchange rather than the role of Europeans in globalization.
The discussion between Blanks (2) and (3) can be labelled as the second segment: it explores how trade routes developed in the year 1000, predating Columbus’s journey. It sets up a comparison between different historical periods - 1000 and 1492 - but the idea this comparison is leading to is still unclear. 
This is where the final segment [the discussion after Blank (3)] fits in. It delivers the passage’s key argument: European involvement in global trade was not a beginning but a continuation of pre-existing systems. It reinforces the idea that the economic structures and trade routes established before 1492 were instrumental in enabling European expansion. The given sentence would best fit Blank (3) because we make a declaration ["The Europeans did not invent globalization"] and then follow up with information that reinforces this idea [“They changed and augmented what was already there since 1000”]. 
Hence, Option D is the correct choice.

\end{problem}

\newpage
\begin{problem}{}{}
% Subject: varc
\textbf{Instructions}

Instructions   

For the following questions answer them individually

\vspace{1em}
\textbf{Question 4}

**The passage given below is followed by four alternate summaries. Choose the option that best captures the essence of the passage.**
Different from individuals, states conduct warfare operations using the DIME model— “diplomacy, information, military, and economics.” Most states do everything they can to inflict pain and confusion on their enemies before deploying the military. In fact, attacks on vectors of information are a well-worn tactic of war and usually are the first target when the charge begins. It’s common for telecom data and communications networks to be routinely monitored by governments, which is why the open data policies of the web are so concerning to many advocates of privacy and human rights. With the worldwide adoption of social media, more governments are getting involved in low-grade information warfare through the use of cyber troops. According to a study by the Oxford Internet Institute in 2020, cyber troops are “government or political party actors tasked with manipulating public opinion online.” The Oxford research group was able to identify 81 countries with active cyber troop operations, utilizing many different strategies to spread false information, including spending millions on online advertising.

\begin{enumerate}[label=\Alph*.]
    \item Following the DIME model, many governments have taken advantage of open data policies of the web to deploy cyber troops who manipulate domestic public opinion, using advertising and other strategies to spread false information.
    \item Governments primarily use the DIME model to deploy cyber troops who practise lowgrade information warfare, seeking to manipulate public opinion with the objective of inflicting pain and confusion on their enemies.
    \item Using the DIME model, together with military operations, many governments simultaneously conduct information warfare with the help of cyber troops and routinely monitor telecom data and communications networks.
    \item As part of conducting information warfare as per the DIME model, many governments routinely monitor telecom data and communications networks, and use cyber troops on social media to manipulate public opinion.
[](https://cracku.in/4-the-passage-given-below-is-followed-by-four-altern-x-cat-2024-slot-2)
\end{enumerate}
\vspace{1em}
\textbf{Solution:}

The passage starts by explaining how states use the DIME model in warfare, often targeting information systems first to destabilize enemies. Governments monitor communication networks, raising privacy concerns. With social media's rise, many countries employ "cyber troops" to manipulate public opinion online, spreading disinformation through tactics like paid ads. Option D best captures all these points. 
_Option A_ : This option incorrectly focuses on open data policies being used to deploy cyber troops, which the passage does not emphasize.
_Option B_ : This option fails to address the key part of the passage, which is monitoring telecom data and networks.
_Option C_ : The passage does not indicate that governments conduct warfare simultaneously with military forces; rather, they act before involving the military. We can infer this from the statement,_"Most states do everything they can to inflict pain and confusion on their enemies before deploying the military."_

\end{problem}

\newpage
\begin{problem}{}{}
% Subject: varc
\textbf{Instructions}

Instructions   

For the following questions answer them individually

\vspace{1em}
\textbf{Question 5}

**The passage given below is followed by four alternate summaries. Choose the option that best captures the essence of the passage.**
John Cleese told Fox News Digital that comedians do not have the freedom to be funny in 2022. “There’s always been limitations on what they’re allowed to say,” Cleese said. “I think it’s particularly worrying at the moment because you can only create in an atmosphere of freedom, where you’re not checking everything you say critically before you move on. What you have to be able to do is to build without knowing where you’re going because you’ve never been there before. That’s what creativity is — you have to be allowed to build. And a lot of comedians now are sitting there and when they think of something, they say something like, ‘Can I get away with it? I don’t think so. So and so got into trouble, and he said that, oh, she said that.’ You see what I mean? And that’s the death of creativity.”

\begin{enumerate}[label=\Alph*.]
    \item Comedians must not check what they think and say. They must go where no one has gone before.
    \item Creativity and critical thinking cannot work together. Comedians must first be creative, and later be critical.
    \item Comedians are being prevented from saying what they want and that is the death of this art form.
    \item Freedom and creativity are essential for comedy. Fear about offending people hinders originality.
[](https://cracku.in/5-the-passage-given-below-is-followed-by-four-altern-x-cat-2024-slot-2)
\end{enumerate}
\vspace{1em}
\textbf{Solution:}

Option D best captures the essence of the passage.   
  
John Cleese argues that comedians need freedom to be creative and that fear of offending people or worrying about the consequences of what they say hinders their ability to be original. He highlights how modern comedians often second-guess themselves, which stifles their creativity. Option D reflects this idea, which stresses the importance of freedom and creativity in comedy, while warning against the fear that stifles innovation.  
  
_Option A_ : While Cleese advocates for freedom in comedy, he doesn't say that comedians _"must go where no one has gone before,_ " which is more of an extreme interpretation than the essence of the passage.  
  
_Option B_ : This option focuses too much on the relationship between creativity and critical thinking. The passage is more about how fear of offending hinders creativity, not about creativity being incompatible with critical thinking.  
  
_Option C_ : This focuses on the _"death of the art form,"_ but Cleese's main point is about how fear of repercussions impacts creativity, not about the art form dying as such.

\end{problem}

\newpage
\begin{problem}{}{}
% Subject: varc
\textbf{Instructions}

Instructions   

For the following questions answer them individually

\vspace{1em}
\textbf{Question 6}

**There is a sentence that is missing in the paragraph below. Look at the paragraph and decide where (option 1, 2, 3, or 4) the following sentence would best fit.****  
  
Sentence** : Yet each day the flock produced eggs with calcareous shells though they apparently had not ingested any calcium from land which was entirely lacking in limestone.
**Paragraph** : Early in this century a young Breton schoolboy who preparing himself for a scientific career began to notice a strange fact about hens in his father's poultry yard. ___(1) ___. As they scratched the soil they constantly seemed to be pecking at specks of mica, a siliceous material dotting the ground. ___(2)___. No one could explain to Louis Kervran why the chickens selected the mica, or why each time a bird was killed for the family cooking pot no trace of the mica could be found in its gizzard. ___(3) ___. It took Kervran many years to establish that the chickens were transmuting one element into another. ___(4)___.

\begin{enumerate}[label=\Alph*.]
    \item Option 1
    \item Option 3
    \item Option 2
    \item Option 4
[](https://cracku.in/6-there-is-a-sentence-that-is-missing-in-the-paragra-x-cat-2024-slot-2)
\end{enumerate}
\vspace{1em}
\textbf{Solution:}

Evaluation of the options:
Option 1: The sentence starting with "Yet each day..." would not fit here. This part is introducing Kervran’s initial observations about the hens. Inserting this sentence would feel disjointed because the paragraph hasn't yet explained the significance of the eggs.
Option 2: The sentence wouldn't fit here either. The focus here is on what the hens are doing (pecking mica), and the sentence about eggs feels disconnected from this point in the narrative. 
Option 3: The sentence about the eggs fits perfectly here because it explains a key element of the puzzle: the hens are producing eggs with calcareous shells despite not having an obvious source of calcium. This provides a crucial piece of information that explains why Kervran was puzzled in the first place. It ties directly to the anomaly he is investigating. 
Option 4: It won't fit here as the paragraph has already shifted to Kervran's conclusion about the chickens transmuting one element to into another. Inserting the sentence about the eggs here would introduce new information too late, disrupting the focus on Kervran's final findings. 
Therefore, The sentence fits best in Option 3. It helps to explain Kervran's puzzling observations and contributes to the mystery that he's trying to solve, all while maintaining a smooth narrative flow. The other options either disrupt the flow or introduce new information too soon or too late.

Instructions   

**The passage below is accompanied by four questions. Based on the passage, choose the best answer for each question.**
The history of any major technological or industrial advance is inevitably shadowed by a less predictable history of unintended consequences and secondary effects — what economists sometimes call “externalities.” Sometimes those consequences are innocuous ones, or even beneficial. Gutenberg invents the printing press, and literacy rates rise, which causes a significant part of the reading public to require spectacles for the first time, which creates a surge of investment in lens-making across Europe, which leads to the invention of the telescope and the microscope.
Oftentimes the secondary effects seem to belong to an entirely different sphere of society. When Willis Carrier hit upon the idea of air-conditioning, the technology was primarily intended for industrial use: ensuring cool, dry air for factories that required low-humidity environments. But…it touched off one of the largest migrations in the history of the United States, enabling the rise of metropolitan areas like Phoenix and Las Vegas that barely existed when Carrier first started tinkering with the idea in the early 1900s.
Sometimes the unintended consequence comes about when consumers use an invention in a surprising way. Edison famously thought his phonograph, which he sometimes called “the talking machine,” would primarily be used to take dictation….But then later innovators… discovered a much larger audience willing to pay for musical recordings made on descendants of Edison’s original invention. In other cases, the original innovation comes into the world disguised as a plaything…the way the animatronic dolls of the mid-1700s inspired Jacquard to invent the first “programmable” loom and Charles Babbage to invent the first machine that fit the modern definition of a computer, setting the stage for the revolution in programmable technology that would transform the 21st century in countless ways.
We live under the gathering storm of modern history’s most momentous unintended consequence….carbon-based climate change. Imagine the vast sweep of inventors whose ideas started the Industrial Revolution, all the entrepreneurs and scientists and hobbyists who had a hand in bringing it about. Line up a thousand of them and ask them all what they had been hoping to do with their work. Not one would say that their intent had been to deposit enough carbon in the atmosphere to create a greenhouse effect that trapped heat at the surface of the planet. And yet here we are.
Ethyl (leaded fuel) and Freon belonged to the same general class of secondary effect: innovations whose unintended consequences stem from some kind of waste by-product that they emit. But the potential health threats of Ethyl (unleaded fuel) were visible in the 1920s, unlike, say, the long-term effects of atmospheric carbon build up in the early days of the Industrial Revolution….
Indeed, it is reasonable to see CFCs (chlorofluorocarbons) as a forerunner of the kind of threat we will most likely face in the coming decades, as it becomes increasingly possible for individuals or small groups to create new scientific advances — through chemistry or biotechnology or materials science — setting off unintended consequences that reverberate on a global scale.

\end{problem}

\newpage
\begin{problem}{}{}
% Subject: varc
\textbf{Instructions}

Instructions   

For the following questions answer them individually

\vspace{1em}
\textbf{Question 7}

The author lists all of the following examples as “externalities” of major technical advances EXCEPT:

\begin{enumerate}[label=\Alph*.]
    \item build-up of chlorofluorocarbons in the atmosphere
    \item cooling and de-humidifying of factories through air-conditioning
    \item application of the Jacquard loom to modern IT programming
    \item extension of the phonograph to large-scale recording of music
[](https://cracku.in/7-the-author-lists-all-of-the-following-examples-as--x-cat-2024-slot-2)
\end{enumerate}
\vspace{1em}
\textbf{Solution:}

Option B is the correct answer. 
This is not an externality because it was the original intended use of air-conditioning. The passage mentions that Carrier invented air-conditioning to ensure cool, dry air for factories with low-humidity requirements.
The other options are externalities:
Option A: This is an unintended consequence. CFCs were initially used in refrigeration and air-conditioning, but their long-term environmental impact (ozone depletion) was not anticipated.
Option C: The Jacquard loom was originally a mechanical device for weaving patterns in fabric, but it led to the development of programmable machines, which had far-reaching effects on modern computing, which was also an unintended consequence of the loom's invention.
Option D: The passage states that the phonograph was initially designed for dictation but was adapted for music recording, which was an unintended consequence.

\end{problem}

\newpage
\begin{problem}{}{}
% Subject: varc
\textbf{Instructions}

Instructions   

For the following questions answer them individually

\vspace{1em}
\textbf{Question 8}

Which of the following best conveys the main point of the first paragraph?

\begin{enumerate}[label=\Alph*.]
    \item The secondary effects of most major technological advances in the past, especially if they were unintended, have turned out to be beneficial.
    \item The full impact of technological advances cannot be estimated in the short run as the ripple effects often extend far beyond the original intent.
    \item It is important to judge an invention not by its immediate outcomes, but by the holistic impact of its secondary effects.
    \item The entire impact of a technological advance should be evaluated by the boost its secondary effects gives to generating further technological advances.
[](https://cracku.in/8-which-of-the-following-best-conveys-the-main-point-x-cat-2024-slot-2)
\end{enumerate}
\vspace{1em}
\textbf{Solution:}

Option B is the correct answer.
  
  
Option A: While some secondary effects may be beneficial, the focus of the paragraph is more on the unpredictability and far-reaching nature of these effects, rather than their inherent benefit.
Option C: The paragraph doesn't advocate judging inventions by their secondary effects. Instead, it mentions that these effects are unpredictable and sometimes surprising. The main point is their unpredictability, not how to judge an invention.
Option D: The paragraph does not suggest that the impact of a technological advance should be evaluated by the boost its secondary effects give to generating further technological advances.

\end{problem}

\newpage
\begin{problem}{}{}
% Subject: varc
\textbf{Instructions}

Instructions   

For the following questions answer them individually

\vspace{1em}
\textbf{Question 9}

Carrier, Babbage, and Edison are mentioned in the passage to illustrate the author’s point that

\begin{enumerate}[label=\Alph*.]
    \item the secondary effect of past inventions mostly resulted in the creation of new inventions.
    \item these inventors could not have visualised the eventual impact of their inventions on society.
    \item despite the original intention, the unintended consequences of their inventions were largely beneficial.
    \item inventions typically end up being used for entirely different purposes than the intended ones.
[](https://cracku.in/9-carrier-babbage-and-edison-are-mentioned-in-the-pa-x-cat-2024-slot-2)
\end{enumerate}
\vspace{1em}
\textbf{Solution:}

Option B is the correct answer. 
The author mentions Carrier, Babbage, and Edison to emphasize that the inventors' original intentions were not related to the unexpected societal impacts their inventions had:
  * Carrier created air-conditioning for industrial use, but it triggered a mass migration to cities like Phoenix and Las Vegas. 
  * Babbage's invention of a programmable loom and Edison’s phonograph were originally intended for specific purposes (textile weaving and dictation, respectively). Still, they led to far-reaching technological developments in computing and music industries.   
Therefore, we can infer that the inventors did not anticipate the full consequences of their inventions.


Option A: The secondary effects are shown as surprising or leading to unforeseen societal changes, rather than leading to more inventions.
Option C: The passage doesn't claim that the unintended consequences were largely beneficial. While some consequences may have been beneficial (like the telescope and microscope from the printing press), others (like climate change ) have been harmful.
Option D: This is close, but it's not the main point. The intended purpose of the inventions may have been different from their actual use, but the author’s primary argument is about how the inventors could not have predicted the full societal impacts of their inventions, rather than focusing on how inventions end up being used for different purposes.

\end{problem}

\newpage
\begin{problem}{}{}
% Subject: varc
\textbf{Instructions}

Instructions   

For the following questions answer them individually

\vspace{1em}
\textbf{Question 10}

We can assume that the author would support all of the following views EXCEPT:

\begin{enumerate}[label=\Alph*.]
    \item While technological advances in the past have had innocuous or beneficial outcomes, more recent advances have the potential to be more threatening globally.
    \item The by-products of leaded fuel, rather than the fuel itself, were responsible for the build-up of carbon-related gases in the atmosphere.
    \item It has become far easier for people today to bring out innovations with dire worldwide consequences than it was earlier.
    \item The emissions caused by the large-scale use of leaded fuel ought to have been addressed earlier than they were.
[](https://cracku.in/10-we-can-assume-that-the-author-would-support-all-of-x-cat-2024-slot-2)
\end{enumerate}
\vspace{1em}
\textbf{Solution:}

Option A is the correct answer. 
The author does not imply that recent advances are more threatening than past ones. Instead, he suggests that the nature of technological progress (with more individuals and smaller groups able to innovate) has changed, leading to new risks. The focus is not on comparing "past vs. recent" as more threatening but on the unexpected global impacts of all technological advances. Thus, option A misrepresents the author's view.
Option B: The author would agree with this because the passage explains that Ethyl (leaded fuel) and Freon are examples of innovations whose unintended consequences stem from by-products they emitted (such as lead from fuel and chemicals from Freon), which had secondary effects on health and the environment.
Option C: The author would support this as he mentions how individuals or small groups can now create innovations that have global impacts, particularly in fields like biotechnology and chemistry, which was less true in the past.
Option D: The author suggests that the health threats of leaded fuel were visible earlier (in the 1920s), implying that they should have been addressed sooner. Hence, the author would support this

Instructions   

For the following questions answer them individually

\end{problem}

\newpage
\begin{problem}{}{}
% Subject: varc
\textbf{Instructions}

Instructions   

For the following questions answer them individually

\vspace{1em}
\textbf{Question 11}

**The passage given below is followed by four alternate summaries. Choose the option that best captures the essence of the passage.**
Recent important scientific findings have emerged from crossing the boundaries of scientific fields. They stem from physicists collaborating with biologists, sociologists and others, to answer questions about our world. But physicists and their potential collaborators often find their cultures out of sync. For one, physicists often discard a lot of information while extracting broad patterns; for other scientists, information is not readily disposed. Further, many non-physicists are uncomfortable with mathematical models. Still, the desire to work on something new and different is real, and there are clear benefits from the collision of views.

\begin{enumerate}[label=\Alph*.]
    \item Despite differences in their research styles, physicists’ research collaborations with scholars from other disciplines have yielded important research findings.
    \item Large data sets and mathematical models in physics research combined with the research methods of non-physicist collaborators have yielded important scientific findings.
    \item The desire to diversify their research and answer important questions has led to several collaborations between physicists and other social scientists.
    \item Physicists have successfully buried their differences on research methods applied in other fields in their desire to find answers to baffling scientific questions.
[](https://cracku.in/11-the-passage-given-below-is-followed-by-four-altern-x-cat-2024-slot-2)
\end{enumerate}
\vspace{1em}
\textbf{Solution:}

Option A is the correct answer.   
  
Option A captures the core idea that their collaborations have led to valuable scientific discoveries despite differences in research methods between physicists and other scientists. The passage emphasizes how these contrasting approaches still lead to productive outcomes, demonstrating the benefits of cross-disciplinary work.
Option B: While large data sets and mathematical models are mentioned, this option incorrectly focuses on "large data sets and mathematical models" as the main contributor, which is not the main point of the passage.
Option C: This is partially true, but the passage does not emphasize the "desire to diversify" research or focus on social scientists. It is about the collaboration of different scientific fields, not specifically social science.
Option D: This is inaccurate because the passage does not state that physicists have "buried" their differences; rather, it says that their differences exist, but the collaboration is still valuable. This is an extreme interpretation.

\end{problem}

\newpage
\begin{problem}{}{}
% Subject: varc
\textbf{Instructions}

Instructions   

For the following questions answer them individually

\vspace{1em}
\textbf{Question 12}

Five jumbled up sentences (labelled 1, 2, 3, 4 and 5), related to a topic, are given below. Four of them can be put together to form a coherent paragraph. Identify the odd sentence and key in the number of that sentence as your answer.
1. The UK is a world leader in developing cultivated meat and the approval of a cultivated pet food is an important milestone.  
2. If we’re to realise the full potential benefits of cultivated meat the government must invest in research and infrastructure.  
3. The first UK applications for cultivated meat produced for humans remain under assessment with the Food Standards Agency.  
4. The previous UK government had been looking at fast-tracking the approval of cultivated meat for human consumption.  
5. It underscores the potential for new innovation to help reduce the negative impacts of intensive animal agriculture.
Backspace  
789  
456  
123  
0.-  
Clear All  

Submit
[](https://cracku.in/12-five-jumbled-up-sentences-labelled-1-2-3-4-and-5-r-x-cat-2024-slot-2)

\begin{enumerate}[label=\Alph*.]
\end{enumerate}
\vspace{1em}
\textbf{Solution:}

**Sentence 4 is the odd one out.  
**  
Sentences 1, 2, 3, and 5 all focus on the current status and future prospects of cultivated meat in the UK, such as the leadership in developing cultivated meat, regulatory assessments, and the importance of government investment in research.  
  
Sentence 4, however, refers to the previous UK government's efforts to fast-track approval of cultivated meat which is disconnected from the other sentences, as the focus shifts to a past action rather than the current developments or future needs in the industry.

Instructions   

**The passage below is accompanied by four questions. Based on the passage, choose the best answer for each question.**
The job of a peer reviewer is thankless. Collectively, academics spend around 70 million hours every year evaluating each other’s manuscripts on the behalf of scholarly journals — and they usually receive no monetary compensation and little if any recognition for their effort. Some do it as a way to keep abreast with developments in their field; some simply see it as a duty to the discipline. Either way, academic publishing would likely crumble without them.
In recent years, some scientists have begun posting their reviews online, mainly to claim credit for their work. Sites like Publons allow researchers to either share entire referee reports or simply list the journals for whom they’ve carried out a review…. The rise of Publons suggests that academics are increasingly placing value on the work of peer review and asking others, such as grant funders, to do the same. While that’s vital in the publish-or-perish culture of academia, there’s also immense value in the data underlying peer review. Sharing peer review data could help journals stamp out fraud, inefficiency, and systemic bias in academic publishing.….
Peer review data could also help root out bias. Last year, a study based on peer review data for nearly 24,000 submissions to the biomedical journal eLife found that women and non- Westerners were vastly underrepresented among peer reviewers. Only around one in every five reviewers was female, and less than two percent of reviewers were based in developing countries…. Openly publishing peer review data could perhaps also help journals address another problem in academic publishing: fraudulent peer reviews. For instance, a minority of authors have been known to use phony email addresses to pose as an outside expert and review their own manuscripts.…
Opponents of open peer review commonly argue that confidentiality is vital to the integrity of the review process; referees may be less critical of manuscripts if their reports are published, especially if they are revealing their identities by signing them. Some also hold concerns that open reviewing may deter referees from agreeing to judge manuscripts in the first place, or that they’ll take longer to do so out of fear of scrutiny….
Even when the content of reviews and the identity of reviewers can’t be shared publicly, perhaps journals could share the data with outside researchers for study. Or they could release other figures that wouldn’t compromise the anonymity of reviews but that might answer important questions about how long the reviewing process takes, how many researchers editors have to reach out to on average to find one who will carry out the work, and the geographic distribution of peer reviewers.
Of course, opening up data underlying the reviewing process will not fix peer review entirely, and there may be instances in which there are valid reasons to keep the content of peer reviews hidden and the identity of the referees confidential. But the norm should shift from opacity in all cases to opacity only when necessary.

\end{problem}

\newpage
\begin{problem}{}{}
% Subject: varc
\textbf{Instructions}

Instructions   

For the following questions answer them individually

\vspace{1em}
\textbf{Question 13}

According to the passage, which of the following is the only reason NOT given in favour of making peer review data public?

\begin{enumerate}[label=\Alph*.]
    \item It will deal with peer review fraud such as authors publishing bogus reviews of their work.
    \item It would highlight the gender and race biases currently existing in the selection of reviewers.
    \item It could address various inefficiencies and fraudulent practices that continue in academic publishing process.
    \item It can tackle the problem of selecting appropriately qualified reviewers for academic writing.
[](https://cracku.in/13-according-to-the-passage-which-of-the-following-is-x-cat-2024-slot-2)
\end{enumerate}
\vspace{1em}
\textbf{Solution:}

Option D is the correct answer.
The passage provides several arguments in favour of making peer review data public, but it does not mention that making this data public would help in selecting appropriately qualified reviewers for academic writing.
Option A: The passage mentions that openly publishing peer review data could help journals address fraudulent peer reviews, such as authors using fake email addresses to review their own manuscripts. Therefore, this is a reason given in favour of making the data public.
Option B: The passage discusses a study showing that women and non-Westerners are underrepresented among peer reviewers. Publishing peer review data could help highlight these gender and race biases, which is one of the reasons for making the data public.
Option C: The passage suggests that sharing peer-review data could help journals tackle issues like fraud, inefficiency, and systemic bias in the publishing process, which is another reason given in favour of making the data public.

\end{problem}

\newpage
\begin{problem}{}{}
% Subject: varc
\textbf{Instructions}

Instructions   

For the following questions answer them individually

\vspace{1em}
\textbf{Question 14}

All of the following are listed as reasons why academics choose to review other scholars’ work EXCEPT:

\begin{enumerate}[label=\Alph*.]
    \item It helps them keep current with cutting-edge ideas in their academic disciplines.
    \item Some use this as an opportunity to publicise their own review work.
    \item It is seen as a form of service to the academic community.
    \item It is seen as an opportunity to expand their influence in the academic community.
[](https://cracku.in/14-all-of-the-following-are-listed-as-reasons-why-aca-x-cat-2024-slot-2)
\end{enumerate}
\vspace{1em}
\textbf{Solution:}

Option D is the correct answer. 
This is not mentioned as a reason in the passage. The focus is on staying informed, contributing to the field, and publicizing their own work, not on expanding influence. Therefore, option D is the correct answer as it is the only reason not mentioned in the passage.
_Option A_ : The passage mentions that some academics review work to _"keep abreast with developments in their field,"_ which aligns with the idea that reviewing helps them stay updated with cutting-edge ideas. This is a valid reason.
_Option B_ : The passage also mentions that some scientists post their reviews online _"mainly to claim credit for their work,_ " which indicates that some view reviewing as an opportunity to publicize their contributions. This is a valid reason
_Option C_ : The passage states that some view reviewing as _"a duty to the discipline,"_ which is a form of service to the academic community. Therefore, option C is also a valid reason.

\end{problem}

\newpage
\begin{problem}{}{}
% Subject: varc
\textbf{Instructions}

Instructions   

For the following questions answer them individually

\vspace{1em}
\textbf{Question 15}

Based on the passage we can infer that the author would most probably support

\begin{enumerate}[label=\Alph*.]
    \item more careful screening to ensure the recruitment of content-familiar peer reviewers.
    \item preserving the anonymity of reviewers to protect them from criticism.
    \item publicising peer review data rather than the publication of actual reviews.
    \item greater transparency across the peer review process in academic publishing.
[](https://cracku.in/15-based-on-the-passage-we-can-infer-that-the-author--x-cat-2024-slot-2)
\end{enumerate}
\vspace{1em}
\textbf{Solution:}

Option D is the correct answer. 
_"from opacity in all cases to opacity only when necessary,"_ implying support for greater transparency in peer review.
_Option A_ : The passage mentions concerns about reviewer selection, such as gender and geographic imbalances, but it does not suggest that careful screening of content-familiar reviewers is a priority.
_Option B_ : The passage mentioned about maintaining confidentiality, but it does not strongly argue in favour of preserving anonymity. 
_Option C_ : The author does not fully endorse the publication of actual reviews or reviewer identities in every situation. While the author advocates for sharing peer review data, they also acknowledge that, in some cases, the content of reviews and reviewers' identities may need to remain confidential.

\end{problem}

\newpage
\begin{problem}{}{}
% Subject: varc
\textbf{Instructions}

Instructions   

For the following questions answer them individually

\vspace{1em}
\textbf{Question 16}

According to the passage, some are opposed to making peer reviews public for all the following reasons EXCEPT that it

\begin{enumerate}[label=\Alph*.]
    \item makes reviewers reluctant to review manuscripts, especially if these are critical of the submitted work.
    \item leaves the reviewers unexposed to unwarranted and unjustified criticism or comments from others.
    \item deters reviewers from producing honest, if critical, reviews that are vital to the sound publishing process.
    \item delays the manuscript evaluation process as reviewers would take longer to write their reviews.
[](https://cracku.in/16-according-to-the-passage-some-are-opposed-to-makin-x-cat-2024-slot-2)
\end{enumerate}
\vspace{1em}
\textbf{Solution:}

Let's evaluate the options:  
  
_Option A_ : The passage mentions this concern:_"referees may be less critical of manuscripts if their reports are published, especially if they are revealing their identities by signing them."  
  
__Option B_ : The passage does not mention that one of the reasons to oppose open peer review is to protect reviewers from unwarranted and unjustified criticism. Instead, the passage discusses the concern that reviewers may avoid giving critical feedback if their identities and reports are made public, not because they want to avoid unjust criticism. The key concern is the fear of justified criticism rather than avoiding unjust or unwarranted criticism.  
_  
__Option C_ : The passage mentions that reviewers may be less critical in their reports if they fear their reviews will be published. This might prevent reviewers from being honest and offering critical assessments, which are crucial for the integrity of the publishing process.  
  
_Option D_ : Another concern raised in the passage is that reviewers might take longer to submit their reviews if they know they will be publicly scrutinized. This delay in the review process is seen as a disadvantage of open peer review.  
  
Therefore,**option B is the correct answer** because the passage does not mention concerns about leaving reviewers unexposed to unwarranted criticism.

Instructions   

**The passage below is accompanied by four questions. Based on the passage, choose the best answer for each question.  
**  
[S]pices were a global commodity centuries before European voyages. There was a complex chain of relations, yet consumers had little knowledge of producers and vice versa. Desire for spices helped fuel European colonial empires to create political, military and commercial networks under a single power.
Historians know a fair amount about the supply of spices in Europe during the medieval period - the origins, methods of transportation, the prices - but less about demand. Why go to such extraordinary efforts to procure expensive products from exotic lands? Still, demand was great enough to inspire the voyages of Christopher Columbus and Vasco Da Gama, launching the first fateful wave of European colonialism. . . .
So, why were spices so highly prized in Europe in the centuries from about 1000 to 1500? One widely disseminated explanation for medieval demand for spices was that they covered the taste of spoiled meat. . . . Medieval purchasers consumed meat much fresher than what the average city-dweller in the developed world of today has at hand. However, refrigeration was not available, and some hot spices have been shown to serve as an anti-bacterial agent. Salting, smoking or drying meat were other means of preservation. Most spices used in cooking began as medical ingredients, and throughout the Middle Ages spices were used as both medicines and condiments. Above all, medieval recipes involve the combination of medical and culinary lore in order to balance food's humeral properties and prevent disease. Most spices were hot and dry and so appropriate in sauces to counteract the moist and wet properties supposedly possessed by most meat and fish. . . .
Where spices came from was known in a vague sense centuries before the voyages of Columbus. Just how vague may be judged by looking at medieval world maps . . . To the medieval European imagination, the East was exotic and alluring. Medieval maps often placed India close to the so-called Earthly Paradise, the Garden of Eden described in the Bible.
Geographical knowledge has a lot to do with the perceptions of spices’ relative scarcity and the reasons for their high prices. An example of the varying notions of scarcity is the conflicting information about how pepper is harvested. As far back as the 7th century Europeans thought that pepper in India grew on trees "guarded" by serpents that would bite and poison anyone who attempted to gather the fruit. The only way to harvest pepper was to burn the trees, which would drive the snakes underground. Of course, this bit of lore would explain the shriveled black peppercorns, but not white, pink or other colors.
Spices never had the enduring allure or power of gold and silver or the commercial potential of new products such as tobacco, indigo or sugar. But the taste for spices did continue for a while beyond the Middle Ages. As late as the 17th century, the English and the Dutch were struggling for control of the Spice Islands: Dutch New Amsterdam, or New York, was exchanged by the British for one of the Moluccan Islands where nutmeg was grown.

\end{problem}

\newpage
\begin{problem}{}{}
% Subject: varc
\textbf{Instructions}

Instructions   

For the following questions answer them individually

\vspace{1em}
\textbf{Question 17}

It can be inferred that all of the following contributed to a decline in the allure of spices, EXCEPT:

\begin{enumerate}[label=\Alph*.]
    \item the development of refrigeration techniques.
    \item increase in the availability of spices.
    \item changes in the system of medical treatment.
    \item changes in European cuisine.
[](https://cracku.in/17-it-can-be-inferred-that-all-of-the-following-contr-x-cat-2024-slot-2)
\end{enumerate}
\vspace{1em}
\textbf{Solution:}

Option B is the correct answer. 
The demand for spices was not necessarily tied to their availability. The passage states that medieval Europeans had limited geographical knowledge of where spices came from and were highly prized despite their relative scarcity. Spices were used for culinary and medicinal purposes, and the demand was driven by cultural and medical factors rather than availability.
Option A: The passage states that spices were used partly due to the unavailability of refrigeration, and some spices served as antibacterial agents. With the development of refrigeration techniques, the necessity for spices to preserve or enhance food properties would diminish, leading to a decline in their appeal.
Option C: The passage mentions how spices were used for medicinal properties in medieval times as part of humeral balancing. As medical science evolved, such practices would likely fall out of favour, reducing their significance.
Option D: The passage states that spices played a significant role in medieval cooking, balancing humeral properties in food. Changes in European tastes or culinary practices would have contributed to a decline in the demand for spices.

\end{problem}

\newpage
\begin{problem}{}{}
% Subject: varc
\textbf{Instructions}

Instructions   

For the following questions answer them individually

\vspace{1em}
\textbf{Question 18}

In the context of the passage, the people who heard the story of pepper trees being guarded by snakes would be least likely to arrive at the conclusion that

\begin{enumerate}[label=\Alph*.]
    \item this is why pepper is so hot.
    \item pepper is costly for good reason.
    \item it is not advisable to go to India to harvest the pepper themselves.
    \item it is no surprise that the pepper supply is so limited.
[](https://cracku.in/18-in-the-context-of-the-passage-the-people-who-heard-x-cat-2024-slot-2)
\end{enumerate}
\vspace{1em}
\textbf{Solution:}

Option A is the correct answer. 
This is a wrong conclusion because it confuses the physical heat from the harvesting process with the pepper's actual spiciness. The process of using fire for harvesting could not be related to the spiciness (being hot) of the pepper.
Option B: The story of snakes and burning trees implies that pepper was difficult and dangerous to harvest, which could explain why it was costly. This is a reasonable conclusion based on the story.
Option C: Given the danger described in the myth (snakes and burning trees), it's logical that people would conclude it's not advisable to go to India to harvest pepper themselves. This could be a conclusion from the story.
Option D: The story suggests that harvesting pepper is difficult and dangerous, which might lead people to think that the supply is limited. Based on the information provided in the myth, this is a reasonable conclusion.

\end{problem}

\newpage
\begin{problem}{}{}
% Subject: varc
\textbf{Instructions}

Instructions   

For the following questions answer them individually

\vspace{1em}
\textbf{Question 19}

In the context of the passage, which one of the following conclusions CANNOT be reached?

\begin{enumerate}[label=\Alph*.]
    \item The spice trade was a driver of colonial expansion.
    \item India was colonised for its spices and gold.
    \item Tobacco was more marketable than spices.
    \item Colonialism was motivated by the demand for spices.
[](https://cracku.in/19-in-the-context-of-the-passage-which-one-of-the-fol-x-cat-2024-slot-2)
\end{enumerate}
\vspace{1em}
\textbf{Solution:}

Option B is the correct answer.**  
  
** While spices were a major part of European trade with the East, the passage does claim that gold was a motivation for colonizing India. The main focus of the passage is on spices, not gold. Hence, the conclusion that India was colonized for both spices and gold cannot be definitively drawn from the passage.  
  
_Option A_ : The passage hints that the desire for spices played a significant role in driving European colonial expansion. Therefore, this conclusion can be reached from the passage.  
  
_Option C_ : The passage briefly mentions that spices never had the same enduring allure or commercial potential as gold, silver, tobacco, indigo, or sugar. From this, we can infer that tobacco was more marketable than spices at certain points in history.   
 _  
Option D_ : In the passage, the desire for spices is described as one of the major factors that led to European colonialism. Therefore, this conclusion could be reached.

\end{problem}

\newpage
\begin{problem}{}{}
% Subject: varc
\textbf{Instructions}

Instructions   

For the following questions answer them individually

\vspace{1em}
\textbf{Question 20}

If a trader brought white peppercorns from India to medieval Europe, all of the following are unlikely to happen, EXCEPT:

\begin{enumerate}[label=\Alph*.]
    \item medieval maps would be used as navigational aids.
    \item Europeans would doubt the story of pepper harvesting.
    \item the price of spices would decrease.
    \item pepper would no longer be considered exotic.
[](https://cracku.in/20-if-a-trader-brought-white-peppercorns-from-india-t-x-cat-2024-slot-2)
\end{enumerate}
\vspace{1em}
\textbf{Solution:}

Option B is the correct answer.  
  
If white peppercorns were brought to Europe, Europeans would likely doubt the myth of harvesting pepper by burning trees, as white peppercorns would not be burnt like the black peppercorns described in the story. This inconsistency would make them question the accuracy of the myth._  
  
Option A_ : Medieval maps were mentioned as symbolic and inaccurate, making them impractical for navigation. Traders did not rely on these maps for precise geographical guidance.  
  
_Option C_ : This outcome is unlikely. The passage explains that spices were scarce and expensive; thus, bringing white pepper would not significantly affect its price, since it remained a rare and highly valued commodity. 
_Option D_ : Even if a trader brought white peppercorns, pepper would still be viewed as exotic by Europeans due to its rarity and the long journey it took to reach Europe. So, this outcome is also unlikely.

Instructions   

**The passage below is accompanied by four questions. Based on the passage, choose the best answer for each question.**
Carnivores that recognise humans as a means to get food, are a different story. As they become more reliant on human food they might find at campsites or in rubbish bins, they become less avoidant of humans. Losing that instinctive fear response puts them into more situations where they could get into an altercation with a human, which often results in that bear being put down by humans. “A fed bear is a dead bear,” says Servheen, referring to a common saying among biologists and conservationists. Predatory or predation-related attacks are quite rare, only accounting for 17% of attacks in North America since 1955. They occur when a carnivore views a human as prey and hunts it like it would any other animal it uses for food. (. . .)
Then there are animal attacks provoked by people taking pictures with them or feeding them in natural settings such as national parks which often end with animals being euthanised out of precaution. “Eventually, that animal becomes habituated to people, and [then] bad things happen to the animal. And the folks who initially wanted to make that connection don’t necessarily realise that,” says Christine Wilkinson, a postdoctoral researcher at UC Berkeley, California, who’s been studying coyote-human conflicts.
After conducting countless postmortems on all types of carnivore-human attacks spanning 75 years, Penteriani’s team believes 50% could have been avoided if humans reacted differently. A 2017 study co-authored by Penteriani found that engaging in risky behaviour around large carnivores increases the likelihood of an attack.
Two of the most common risky behaviours are parents leaving their children to play outside unattended and walking an unleashed dog, according to the study. Wilkinson says 66% of coyote attacks involve a dog. “[People] end up in a situation where their dog is being chased, or their dog chases a coyote, or maybe they’re walking their dog near a den that’s marked, and the coyote wants to escort them away,” says Wilkinson.
Experts believe climate change also plays a part in the escalation of human-carnivore conflicts, but the correlation still needs to be ironed out. “As finite resources become scarcer, carnivores and people are coming into more frequent contact, which means that more conflict could occur,” says Jen Miller, international programme specialist for the US Fish & Wildlife Service. For example, she says, there was an uptick in lion attacks in western India during a drought when lions and people were relying on the same water sources.
(. . .) The likelihood of human-carnivore conflicts appears to be higher in areas of low-income countries dominated by vast rural landscapes and farmland, according to Penteriani’s research. “There are a lot of working landscapes in the Global South that are really heterogeneous, that are interspersed with carnivore habitats, forests and savannahs, which creates a lot more opportunity for these encounters, just statistically,” says Wilkinson.

\end{problem}

\newpage
\begin{problem}{}{}
% Subject: varc
\textbf{Instructions}

Instructions   

For the following questions answer them individually

\vspace{1em}
\textbf{Question 21}

According to the passage, what is a significant factor that contributes to the habituation of carnivores to human presence?

\begin{enumerate}[label=\Alph*.]
    \item The natural aggression exhibited by carnivores, exacerbated by human interference, particularly when they are safeguarding their offspring or food sources.
    \item The increased scarcity of resources due to climate change, forcing carnivores to venture outside their natural habitats in search of sustenance.
    \item The predatory perception of humans as potential prey within the carnivores’ food chain.
    \item The reduction in carnivores’ instinctive fear response, resulting from their reliance upon human-provided food.
[](https://cracku.in/21-according-to-the-passage-what-is-a-significant-fac-x-cat-2024-slot-2)
\end{enumerate}
\vspace{1em}
\textbf{Solution:}

Option D is the correct answer.  
  

The passage states that:_“As they become more reliant on human food they might find at campsites or in rubbish bins, they become less avoidant of humans. Losing that instinctive fear response puts them into more situations where they could get into an altercation with a human, which often results in that bear being put down by humans."_
This means that the factor contributing most to carnivores' habituation to human presence is the reduction in instinctive fear response and reliance on human food sources, which is evident in option D.
_Option A_ : The passage does not mention their natural aggression. It focuses more on losing fear due to food reliance on human food.  
  
_Option B_ : The passage mentions climate change as a possible reason for increased conflict, but it doesn't suggest it directly contributes to loss of fear or habituation.  
  
_Option C_ : The passage does not mention about the predatory perception of humans as potential prey, hence eliminated.

\end{problem}

\newpage
\begin{problem}{}{}
% Subject: varc
\textbf{Instructions}

Instructions   

For the following questions answer them individually

\vspace{1em}
\textbf{Question 22}

Given the insights provided by Penteriani’s research and Wilkinson’s statement, which of the following conclusions can be drawn about the relationship between landscape heterogeneity and human-carnivore conflicts?

\begin{enumerate}[label=\Alph*.]
    \item Low-income countries with vast, contiguous wilderness areas are less prone to human-carnivore conflicts because these areas lack the human presence necessary for such encounters.
    \item Landscape heterogeneity, characterized by a mix of farmland and natural habitats, inherently reduces the chances of human-carnivore conflicts by providing more refuge for wildlife away from human activity.
    \item Homogeneous landscapes with uniform agricultural practices are more likely to experience high rates of human-carnivore conflicts due to the predictability of resources.
    \item The diversity and interspersion of working landscapes with carnivore habitats in rural areas increase the statistical probability of encounters between humans and carnivores.
[](https://cracku.in/22-given-the-insights-provided-by-penterianis-researc-x-cat-2024-slot-2)
\end{enumerate}
\vspace{1em}
\textbf{Solution:}

Option D is the correct answer. The passage mentions that landscape heterogeneity (a mix of farmland, forests, and carnivore habitats ) in rural areas of low-income countries creates more opportunities for human-carnivore encounters.   
  
As Penteriani’s research shows, such landscapes increase the statistical probability of these conflicts because the areas are interspersed with human and carnivore habitats. This aligns with option D._  
  
Option A_ : This is inconsistent with the passage. The passage states,_"The likelihood of human-carnivore conflicts appears to be higher in areas of low-income countries dominated by vast rural landscapes and farmland"._ Therefore, it is not less prone, rather more prone as per the passage.  
  
_Option B_ : The passage does not claim that landscape heterogeneity inherently decreases the chances of human-carnivore conflict. Instead, it states that diversity increases the likelihood of encounters rather than reducing it.  
  
_Option C_ : The passage does not state that homogeneous landscapes cause high rates of conflict due to predictability. Instead, it suggests that landscape heterogeneity increases encounters.

\end{problem}

\newpage
\begin{problem}{}{}
% Subject: varc
\textbf{Instructions}

Instructions   

For the following questions answer them individually

\vspace{1em}
\textbf{Question 23}

Which of the following statements, if false, would be inconsistent with the concerns raised in the passage regarding the drivers of carnivore-human conflicts?

\begin{enumerate}[label=\Alph*.]
    \item Climate change has had negligible effects on the frequency of carnivore-human interactions in affected regions.
    \item Predatory attacks by carnivores are a common occurrence and have steadily increased over the past few decades.
    \item Carnivores lose their instinctive fear of humans, when consistently exposed to human food sources
    \item Human efforts to avoid risky behaviours around large carnivores have proven effective in reducing conflict incidents.
[](https://cracku.in/23-which-of-the-following-statements-if-false-would-b-x-cat-2024-slot-2)
\end{enumerate}
\vspace{1em}
\textbf{Solution:}

This is an official CAT 2024 Question, and the marked answer is according to the official answer key. We disagree with this answer.
Let us look at falsifying these statements: 
Option A: Climate change has had negligible effects on the frequency of carnivore-human interactions in affected regions. 
False version 1: Climate change has had no effect on the frequency of carnivore-human interactions in affected regions. - Inconsistent with the passage
False version 2: Climate change has had a lot of effect on the frequency of carnivore-human interactions in affected regions. - Consistent
As one of the versions contradicts the passage, we can say that the statement if false is inconsistent.
Option B:
Predatory attacks by carnivores are a common occurrence and have steadily increased over the past few decades.  
False Version: Predatory attacks by carnivores are a rare occurrence and have steadily increased over the past few decades.
This is consistent with the passage.
Option C: Carnivores lose their instinctive fear of humans, when consistently exposed to human food sources.  
False Version: Carnivores do not lose their instinctive fear of humans, when consistently exposed to human food sources. -- irrelevant to the passage as the author speaks on reliance on human food and not exposure to human food "sources".
Option D: Human efforts to avoid risky behaviours around large carnivores have proven effective in reducing conflict incidents.  
False version: Human efforts to avoid risky behaviours around large carnivores have not proven effective in reducing conflict incidents. -- beyond the scope of the passage.

\end{problem}

\newpage
\begin{problem}{}{}
% Subject: varc
\textbf{Instructions}

Instructions   

For the following questions answer them individually

\vspace{1em}
\textbf{Question 24}

According to the passage, which of the following scenarios would MOST likely exacerbate the frequency of carnivore-human conflicts?

\begin{enumerate}[label=\Alph*.]
    \item Implementing 'food waste' management strategies to prevent wild animals being attracted to human food sources.
    \item Addressing the impact of climate change on the availability of resources for wildlife.
    \item Attempting to photograph wild animals from within secured viewing areas in national parks and protected zones.
    \item Unleashing dogs by pet owners in areas with known high concentrations of large carnivores.
[](https://cracku.in/24-according-to-the-passage-which-of-the-following-sc-x-cat-2024-slot-2)
\end{enumerate}
\vspace{1em}
\textbf{Solution:}

Option D is the correct answer.
The passage mentions that 66% of coyote attacks involve a dog, which can either provoke a carnivore or escalate a dangerous situation when the dog chases a carnivore or vice versa. In areas with large carnivores, unleashing dogs can increase the likelihood of encounters and conflicts. Therefore, option D is the most likely scenario to exacerbate carnivore-human conflicts.  
  
Why the other options are less likely to exacerbate the conflicts:  
  
_Option A_ : Preventing wild animals from being attracted to human food sources would actually reduce carnivore-human conflicts by keeping animals from becoming habituated to humans. This would prevent potential issues, not exacerbate them.  
  
_Option B_ : The passage suggests that climate change could increase the frequency of human-carnivore encounters due to scarcity of resources. Therefore, addressing climate change would likely help prevent the issue rather than exacerbate it.  
  
_Option C_ : According to the passage, photographing wild animals in secured viewing areas is not a major driver of carnivore-human conflicts. Conflicts typically arise from behaviours that encourage animals to approach humans or interact in risky ways, not from observation in protected zones. Therefore, this is unlikely to exacerbate conflicts.
[](https://cracku.in/downloads/19317)
  *  
  *

\end{problem}

\newpage
\begin{problem}{}{}
% Subject: varc
\textbf{Instructions}

Instructions   

For the following questions answer them individually

\vspace{1em}
\textbf{Question 1}

**There is a sentence that is missing in the paragraph below. Look at the paragraph and decide where (option 1, 2, 3, or 4) the following sentence would best fit.**
**Sentence** : This reality is putting stress on employees who have to pay for transport, desk lunches, more childcare, clothing and that after-work socialisation - costs they haven’t incurred for nearly two years.
**Paragraph** : ___(1)___. Prices are rising at their fastest rate in 40 years; consequently, return-to-office-related costs have shot up - think petrol and food, for instance.___(2)___. Yet wages haven’t kept up with inflation - even despite the salary growth many workers have enjoyed during a favourable pandemic labour market. ___(3)___.This is especially jarring for workers who were able to save during remote work, when these expenditures weren’t a factor. ___(4)___. In April 2022, Umus, a London university lecturer, told BBC Worklife that they were spending nearly a quarter of what they made every day on return-to-work costs.

\begin{enumerate}[label=\Alph*.]
    \item Option 4
    \item Option 3
    \item Option 2
    \item Option 1
[](https://cracku.in/1-there-is-a-sentence-that-is-missing-in-the-paragra-x-cat-2024-slot-3)
\end{enumerate}
\vspace{1em}
\textbf{Solution:}

The given sentence focuses on the impact of these rising costs on employees - an idea that hasn’t yet been introduced here. Placing this in _Blank 1_ would prematurely shift the focus from a general economic observation to the employee-specific consequences, disrupting the logical flow. Similarly, _Blank 2_ would be a poor choice since it would interrupt the argument about wages not keeping pace with inflation. The broader context of inflation and wages needs to be fully established before narrowing the discussion to employee struggles. 
_Blank 3_ could be a good fit. We observe that the preceding sentence sets up the problem: wages are lagging behind rising costs. The subsequent sentence emphasises the stark contrast between remote work savings and current financial stress; the anecdote from the London University lecturer is further linked to this. Therefore, _Blank 4_ is also unsuitable. By placing the given sentence in Blank 3, we bridge the two ideas presented above by specifying how these rising costs and stagnant wages directly impact employees. It introduces the tangible pressures employees face, which makes the subsequent sentence about savings during remote work even more impactful.
Hence, Option B is the correct choice.

Instructions   

**The passage below is accompanied by four questions. Based on the passage, choose the best answer for each question.**
Fears of artificial intelligence (AI) have haunted humanity since the very beginning of the computer age. Hitherto, these fears focused on machines using physical means to kill, enslave or replace people. But over the past couple of years, new AI tools have emerged that threaten the survival of human civilisation from an unexpected direction. AI has gained some remarkable abilities to manipulate and generate language, whether with words, sounds or images. AI has thereby hacked the operating system of our civilisation.
Language is the stuff almost all human culture is made of. Human rights, for example, aren’t inscribed in our DNA. Rather, they are cultural artefacts we created by telling stories and writing laws. Gods aren’t physical realities. Rather, they are cultural artefacts we created by inventing myths and writing scriptures….What would happen once a non-human intelligence becomes better than the average human at telling stories, composing melodies, drawing images, and writing laws and scriptures? When people think about Chatgpt and other new AI tools, they are often drawn to examples like schoolchildren using AI to write their essays. What will happen to the school system when kids do that? But this kind of question misses the big picture. Forget about school essays. Think of the next American presidential race in 2024, and try to imagine the impact of AI tools that can be made to mass-produce political content, fake news stories and scriptures for new cults…
Through its mastery of language, AI could even form intimate relationships with people, and use the power of intimacy to change our opinions and worldviews. Although there is no indication that AI has any consciousness or feelings of its own, to foster fake intimacy with humans, it is enough if the AI can make them feel emotionally attached to it….
What will happen to the course of history when AI takes over culture, and begins producing stories, melodies, laws and religions? Previous tools like the printing press and radio helped spread the cultural ideas of humans, but they never created new cultural ideas of their own. AI is fundamentally different. AI can create completely new ideas, completely new culture….Of course, the new power of AI could be used for good purposes as well. I won’t dwell on this because the people who develop AI talk about it enough….
We can still regulate the new AI tools, but we must act quickly. Whereas nukes cannot invent more powerful nukes, AI can make exponentially more powerful AI.… Unregulated AI deployments would create social chaos, which would benefit autocrats and ruin democracies. Democracy is a conversation, and conversations rely on language. When AI hacks language, it could destroy our ability to have meaningful conversations, thereby destroying democracy …. And the first regulation I would suggest is to make it mandatory for AI to disclose that it is an AI. If I am having a conversation with someone, and I cannot tell whether it is a human or an AI—that’s the end of democracy. This text has been generated by a human. Or has it?

\end{problem}

\newpage
\begin{problem}{}{}
% Subject: varc
\textbf{Instructions}

Instructions   

For the following questions answer them individually

\vspace{1em}
\textbf{Question 2}

The author identifies all of the following as dire outcomes of the capture of language by AI EXCEPT that it could

\begin{enumerate}[label=\Alph*.]
    \item spawn a completely new culture through its ability to create new ideas and opinions.
    \item out-strip human creativity and endeavours in the spheres such as art and music and,in the formulation of laws.
    \item eventually subvert democratic processes through the mass creation and spread offake political content and news.
    \item apply its mastery of language to create strong emotional ties which could exacerbate the polarization of political views.
[](https://cracku.in/2-the-author-identifies-all-of-the-following-as-dire-x-cat-2024-slot-3)
\end{enumerate}
\vspace{1em}
\textbf{Solution:}

Let us evaluate the given choices - 
_Option A_ : The creation of new culture and ideas is a central theme of the passage, highlighted as a potential outcome of AI's linguistic capabilities: [“..._AI can create completely new ideas, completely new culture_ …”]
_Option B_ : The author hints that AI could surpass human creativity in areas like storytelling, composing music, and drafting laws or scriptures: [“..._What would happen once a non-human intelligence becomes better than the average human at telling stories, composing melodies, drawing images, and writing laws and scriptures?_...”]
_Option C_ : The threat to democracy through the mass production of fake news and political content is a major concern raised by the author: [“..._Think of the next American presidential race in 2024, and try to imagine the impact of AI tools that can be made to mass-produce political content, fake-news stories_ …”]
_Option D_ : The passage does discuss AI's ability to create emotional connections with individuals: [“..._form intimate relationships with people, and use the power of intimacy to change our opinions and worldviews_ …”] However, it does not explicitly connect this capability to ‘exacerbating the polarization’ of political views. The focus is on fostering fake intimacy to influence opinions, not specifically on worsening political polarization.
Hence, Option D is the correct choice.

\end{problem}

\newpage
\begin{problem}{}{}
% Subject: varc
\textbf{Instructions}

Instructions   

For the following questions answer them individually

\vspace{1em}
\textbf{Question 3}

The author terms language “the operating system of our civilization” for all the following reasons EXCEPT that it

\begin{enumerate}[label=\Alph*.]
    \item can influence political views and opinions as it engenders close emotional ties among people.
    \item is the basis of AI tools like ChatGPT which can be used to generate academic content and opinion.
    \item is fundamental to the articulation and spread of human values and culture in our society.
    \item has laid the foundation for the creation of cultural artefacts through writing and telling of stories.
[](https://cracku.in/3-the-author-terms-language-the-operating-system-of--x-cat-2024-slot-3)
\end{enumerate}
\vspace{1em}
\textbf{Solution:}

The author emphasises that language is foundational to human culture and civilisation because it:
  * articulates and spreads human values and culture (as noted in _Option C_).
  * lays the groundwork for creating cultural artefacts through storytelling and laws (as noted in _Option D_).
  * influences political views and fosters emotional ties (as noted in _Option A_).


_Option B_ , however, does not align with the rationale behind the “operating system” metaphor: the author does not suggest that language is the "operating system" of civilisation because it underpins AI tools. Instead, the passage treats AI tools like ChatGPT as leveraging language's existing role in civilization.

\end{problem}

\newpage
\begin{problem}{}{}
% Subject: varc
\textbf{Instructions}

Instructions   

For the following questions answer them individually

\vspace{1em}
\textbf{Question 4}

We can infer that the author is most likely to agree with which of the following statements?

\begin{enumerate}[label=\Alph*.]
    \item People’s fears of the dangers of students using ChatGPT and other new AI tools are unfounded.
    \item The commonly expressed fear that future AI developments will fatally harm humans is unfounded.
    \item Apart from its drawbacks, AI tools have been beneficial in boosting technological and industrial advance worldwide.
    \item One of the biggest casualties from the spread of unregulated AI is likely to be the democratic process.
[](https://cracku.in/4-we-can-infer-that-the-author-is-most-likely-to-agr-x-cat-2024-slot-3)
\end{enumerate}
\vspace{1em}
\textbf{Solution:}

Let us evaluate the choices based on the information in the passage - 
_Option A_ : The author does not dismiss fears about students using AI but instead deems such concerns possibly trivial compared to AI’s larger societal threats.
_Option B_ : The author doesn’t argue that fears of AI harming humans physically are unfounded, but shifts focus to the linguistic and cultural dangers AI presents. It’s unclear whether he will explicitly support the view presented here. 
_Option C_ : Though the author briefly acknowledges that AI can be used for good, this idea is not a central focus of the passage; he also does not emphasise technological or industrial benefits. Therefore, it’s unclear whether the author will support his view. 
_Option D_ : The concern stated here has been clearly underlined in the passage. We are informed of the threat AI poses to democracy through its manipulation of language and ability to generate misinformation, fake intimacy, and propaganda. The author explicitly states that democracy relies on meaningful human conversations, which are undermined when AI becomes indistinguishable from humans. Therefore, this option aligns most closely with the author's argument.
Hence, Option D is the correct choice.

\end{problem}

\newpage
\begin{problem}{}{}
% Subject: varc
\textbf{Instructions}

Instructions   

For the following questions answer them individually

\vspace{1em}
\textbf{Question 5}

The tone of the passage could best be described as

\begin{enumerate}[label=\Alph*.]
    \item cautionary, because the author lays out some adverse effects of the proliferation of unregulated AI tools.
    \item prescient, as the author analyses the future impact of the use of new AI tools on crucial areas of our society and culture.
    \item alarmist, because the passage discusses scenarios of the influence of new AI tools on language and human emotions.
    \item quizzical, as the passage poses several questions, concluding with the question of whether or not the passage content has been generated by AI.
[](https://cracku.in/5-the-tone-of-the-passage-could-best-be-described-as-x-cat-2024-slot-3)
\end{enumerate}
\vspace{1em}
\textbf{Solution:}

The passage clearly highlights the potential risks associated with AI's linguistic capabilities, urging action to regulate its use. While the passage does present hypothetical scenarios and employs rhetorical questions, its tone is primarily grounded in a warning, not sensationalism or mere speculation. This points us towards _Option A_. The author systematically explains the dangers of unregulated AI tools, particularly their capacity to manipulate language and influence human culture. The warnings are thoughtful and aim to provoke awareness and a sense of urgency without excessive dramatisation. The tone is serious and measured, which aligns with a cautionary style.
The remaining tones do not describe the discussion appropriately. For instance, consider _Option B_ : though the passage discusses AI's future implications, the focus is less on prediction and more on warning about what could happen if action is not taken. "Prescient" implies a focus on foresight and vision, but the passage emphasises immediate concerns and actionable advice. Similarly, the passage is not “alarmist” (_Option C_), as it avoids overly exaggerated or emotional claims. It uses logical arguments and examples rather than fearmongering. _Option D_ is also a poor fit: although the passage concludes with a rhetorical question, this is a stylistic device rather than a defining characteristic of the tone. The primary goal is to issue a warning, not to leave the reader in a state of curiosity or wonder, making “quizzical” an incorrect characterisation.
Hence, Option A is the best choice.

Instructions   

For the following questions answer them individually

\end{problem}

\newpage
\begin{problem}{}{}
% Subject: varc
\textbf{Instructions}

Instructions   

For the following questions answer them individually

\vspace{1em}
\textbf{Question 6}

**There is a sentence that is missing in the paragraph below. Look at the paragraph and decide where (option 1, 2, 3, or 4) the following sentence would best fit.**
**Sentence** : Many have had to leave their homes behind, with more than 1.3 million people being displaced due to the drought.
**Passage** : Somalia has been dealing with an enormous humanitarian catastrophe, driven by the longest and most severe drought the country has experienced in at least 40 years. ___(1)___. Five consecutive rainy seasons have failed, causing more than 8 million people - almost half of the country’s population - to experience acute food insecurity. ___(2)___. More than 43,000 people are believed to have lost their lives, with half of the lives lost likely being children under five. The damage the drought has caused is far-reaching. ___(3)___. Farmers have lost all their agricultural income, while pastoralists have lost more than 3 million livestock, impoverishing entire communities, and leaving them on the brink of famine. ___(4)___. Some, like the pastoralists, may never be able to go back as their livelihoods have been irreversibly wiped out.

\begin{enumerate}[label=\Alph*.]
    \item Option 4
    \item Option 2
    \item Option 3
    \item Option 1
[](https://cracku.in/6-there-is-a-sentence-that-is-missing-in-the-paragra-x-cat-2024-slot-3)
\end{enumerate}
\vspace{1em}
\textbf{Solution:}

A useful strategy for determining where the given sentence best fits is to look for places where the flow of ideas feels awkward or disconnected. When we examine the sentences around Blank 1, we see that they are well-linked, moving smoothly from the broader issue of the drought to its immediate consequence, food insecurity. The same logic applies to Blanks 2 and 3 - there are no noticeable disruptions in the flow of ideas here. The structure remains clear and coherent even without inserting the given sentence into these blanks.
However, when we look at Blank 4, we notice a discrepancy. The sentence before Blank 4 mentions farmers and pastoralists, and the sentence following it continues discussing pastoralists, but there’s a slight gap in the connection. If the sentences were closely linked, the author could have used pronouns or rephrased the second sentence to refer to pastoralists more naturally. Additionally, there’s a shift in focus: the paragraph moves from talking about "leaving them on the brink of famine" to "never being able to go back as their livelihoods have been irreversibly wiped out." This jump in ideas suggests the need for a transitional sentence to bridge the two concepts - enter the given sentence. Moreover, the "Many... Some..." structure works well here. First, the author mentions the larger group of farmers and pastoralists and how they were displaced. Then, the sentence narrows the focus to a smaller portion of this group, highlighting how, even for those who were not displaced, their livelihoods were so affected that they may never recover. This progression makes the insertion of the given sentence at Blank 4 both logical and necessary for maintaining coherence.
Hence, Option A is the correct choice.

\end{problem}

\newpage
\begin{problem}{}{}
% Subject: varc
\textbf{Instructions}

Instructions   

For the following questions answer them individually

\vspace{1em}
\textbf{Question 7}

**Five jumbled-up sentences (labelled 1, 2, 3, 4 and 5) related to a topic are given below. Four of them can be put together to form a coherent paragraph. Identify the odd sentence and key in the number of that sentence as your answer.**
1. Part of the appeal of forecasting is not just that it seems to work, but that you don’t seem to need specialized expertise to succeed at it.
2. The tight connection between forecasting and building a model of the world helps explain why so much of the early interest in the idea came from the intelligence community.
3. This was true even though the latter had access to classified intelligence.
4. One frequently cited study found that accurate forecasters’ predictions of geopolitical events, when aggregated using standard scientific methods, were more accurate than the forecasts of members of the US intelligence community who answered the same questions in a confidential prediction market.
5. The aggregated opinions of non-experts doing forecasting have proven to be a better guide to the future than the aggregated opinions of experts.
Backspace  
789  
456  
123  
0.-  
Clear All  

Submit
[](https://cracku.in/7-five-jumbled-up-sentences-labelled-1-2-3-4-and-5-r-x-cat-2024-slot-3)

\begin{enumerate}[label=\Alph*.]
\end{enumerate}
\vspace{1em}
\textbf{Solution:}

Here, Sentences 1, 4, and 5 discuss the general theme of how non-experts or forecasters without specialized expertise can make accurate predictions, often outperforming experts. Sentence 3 adds context by highlighting that this success occurs even when experts have access to classified information.
Contrarily, Sentence 2 shifts the focus to the “intelligence community's interest in forecasting models” rather than continuing the discussion on the accuracy and success of forecasters versus experts. This makes it unrelated to the main flow of the paragraph.
Let us examine the points presented in each sentence to further understand how the statements link. Sentence 1 introduces the main idea that forecasting is appealing because it works even without requiring specialized expertise. Building on this idea, Sentence 5 highlights that non-experts can often outperform experts in forecasting; this, in a way, relates to the appeal of forecasting introduced in the first sentence. Sentence 4 provides evidence to support the claim made in Sentence 5. It refers to a specific study where non-experts outperformed experts in predicting geopolitical events, solidifying the argument that non-expert forecasting can be more reliable. Sentence 3 strengthens the previous claim by highlighting that the experts, despite having access to classified intelligence, were still outperformed by non-experts. In this manner, the arrangement 1-5-4-3 renders a coherent paragraph. 
Hence, Sentence 2 is the odd one out.

Instructions   

**The passage below is accompanied by four questions. Based on the passage, choose the best answer for each question.**
Moutai has been the global booze sensation of the decade. A bottle of its Flying Fairy, which sold in the 1980s for the equivalent of a dollar, now retails for $400. Moutai’s listed shares have soared by almost 600% in the past five years, outpacing the likes of Amazon ...
It does this while disregarding every Western marketing mantra. It is not global, has meagre digital sales and does not appeal to millennials. It scores pitifully on environmental, social and government measures. In the Boy Scout world of Western business, it would leave a bad taste in more ways than one.
Moutai owes its intoxicating success to three factors—not all of them easy to emulate. First, it profits from Chinese nationalism. Moutai is known as the “national liquor”. It was used to raise spirits and disinfect wounds in Mao’s Long March. It was Premier Zhou Enlai’s favourite tipple, shared with Richard Nixon in 1972. Its centuries-old craftsmanship—it is distilled eight times and stored for years in earthenware jars—is a source of national pride. It also claims to be hangover-proof, which would make it an invention to rival gunpowder ...
Second, it chose to serve China’s super-rich rather than its middle class. Markets are littered with the corpses of firms that could not compete in the cut-throat battle for Chinese middle-class wallets. And the country’s premium market is massive—at 73m-strong, bigger than the population of France, notes Euan McLeish of Bernstein, an investment firm, and still less crowded with prestige brands than advanced economies. Moutai is to these well-heeled drinkers what vintage champagne is to the rest of the world ...
Third, Moutai looks beyond affluent millennials and digital natives. The elderly and the middle-aged, it found, can be just as lucrative. Its biggest market now is (male) drinkers in their mid-30s. Many have no siblings, thanks to four decades of China’s one-child policy—which also means their elderly parents can splash out on weddings and banquets. Moutai is often a guest of honour.
Moutai has succeeded thanks to nationalism, elitism and ageism, in other words—not in spite of this unholy trinity. But it faces risks. The government is its largest shareholder—and a meddlesome one. It appears to want prices to remain stable. Exorbitantly priced booze is at odds with its professed socialist ideals. Yet minority investors—including many foreign funds—lament that Moutai’s wholesale price is a third of what it sells for in shops. Raising it could boost the company’s profits further. Instead, in what some see as a travesty of corporate governance, its majority owner has plans to set up its own sales channel ...
In the long run, its biggest risk may be millennials. As they grow older, health concerns, work-life balance and the desire for more wholesome pursuits than binge-drinking may curb the“Ganbei!” toasting culture [heavy drinking] on which so much of the demand for Moutai rests. For the time being, though, the party goes on.

\end{problem}

\newpage
\begin{problem}{}{}
% Subject: varc
\textbf{Instructions}

Instructions   

For the following questions answer them individually

\vspace{1em}
\textbf{Question 8}

The phrase “would make it an invention to rival gunpowder” has been used in the passage in a sense that is

\begin{enumerate}[label=\Alph*.]
    \item literal
    \item substantive
    \item metaphorical
    \item synonymical
[](https://cracku.in/8-the-phrase-would-make-it-an-invention-to-rival-gun-x-cat-2024-slot-3)
\end{enumerate}
\vspace{1em}
\textbf{Solution:}

The statement in the question draws a comparison between Moutai’s claimed hangover-proof quality and a groundbreaking historical invention (gunpowder). The use of "would make it" suggests a hypothetical scenario, and the comparison is meant to emphasise significance, not an actual invention. This indicates that the comparison is metaphorical/figurative (_Option C_), conveying the liquor’s potential impact on culture and society; hence, we can eliminate _Option A_. 
_Option B_ also implies that the comparison is based on actual substance or tangible qualities. Though the claim about being hangover-proof is substantive, the comparison to gunpowder is not grounded in tangible, measurable terms but rather in its significance. Similarly, _Option D_ suggests that the phrase equates Moutai directly with gunpowder, which is incorrect: the phrase does not treat the two as synonyms but draws a symbolic comparison.
Hence, Option C is the correct choice.

\end{problem}

\newpage
\begin{problem}{}{}
% Subject: varc
\textbf{Instructions}

Instructions   

For the following questions answer them individually

\vspace{1em}
\textbf{Question 9}

Which one of the following is both a reason for Moutai’s success as well as a possible threat to that success?

\begin{enumerate}[label=\Alph*.]
    \item Chinese love of liquor filled celebration.
    \item Government involvement in its business.
    \item Its appeal to the rich.
    \item Its appeal to the older age group.
[](https://cracku.in/9-which-one-of-the-following-is-both-a-reason-for-mo-x-cat-2024-slot-3)
\end{enumerate}
\vspace{1em}
\textbf{Solution:}

To identify the factor that is both a reason for Moutai’s success and a potential threat, we could consider aspects that currently drive demand while also holding the potential to hinder it in the future. Looking at the given choices, _Option A_ may be a contributing factor to Moutai’s success, as cultural practices around drinking play a role in its popularity. However, the passage does not explicitly frame this as a threat, especially since it is deeply ingrained in Chinese society, making this option less fitting. _Option C_ is also a significant reason for Moutai's success, as it targets the super-wealthy, but the passage does not indicate that this would become a threat in the future. _Option B_ is presented more as a hurdle to Moutai’s success, while the explicit benefits remain to be discussed. 
On the other hand, we can deduce that _Option D_ is both a key reason for Moutai’s success - by tapping into the spending power of older consumers - and a potential threat, as the younger generations, with different health-conscious lifestyles, could move away from the heavy drinking culture that has driven Moutai's demand. The passage specifically highlights this generational shift as a risk to Moutai’s long-term success.
Hence, Option D is the best choice.

\end{problem}

\newpage
\begin{problem}{}{}
% Subject: varc
\textbf{Instructions}

Instructions   

For the following questions answer them individually

\vspace{1em}
\textbf{Question 10}

In the context of the passage, it is most likely that the author refers to Moutai’s marketing strategy as “the unholy trinity” because

\begin{enumerate}[label=\Alph*.]
    \item there is nothing holy about marketing techniques for liquor.
    \item it profits from Chinese nationalist feelings.
    \item it contradicts the Western strategy of marketing.
    \item it exposes the firm to long term risks.
[](https://cracku.in/10-in-the-context-of-the-passage-it-is-most-likely-th-x-cat-2024-slot-3)
\end{enumerate}
\vspace{1em}
\textbf{Solution:}

The author uses the phrase “the unholy trinity” to describe Moutai’s marketing strategy, which relies on three factors: nationalism, elitism, and ageism. The word “unholy” implies that these factors may be unconventional or controversial, which could be why they are described as such. The author presents this idea in the context of Western marketing, suggesting that Moutai might have succeeded due to these factors, even though they seem to defy conventional Western norms. Thus, the phrase reflects how Moutai’s marketing strategy is in stark contrast to Western business practices (_Option C_). 
Contrarily, Options A, B, and D either do not address the contrast with Western strategies or misinterpret the focus of the passage. For instance, _Option A_ is not the best choice because the phrase “unholy trinity” isn’t a comment on the morality of marketing liquor itself but rather on the controversial nature of Moutai’s specific strategy. _Options B and D_ focus on tangential aspects that are irrelevant to the question.

\end{problem}

\newpage
\begin{problem}{}{}
% Subject: varc
\textbf{Instructions}

Instructions   

For the following questions answer them individually

\vspace{1em}
\textbf{Question 11}

In the context of the passage, we can infer that to succeed in the liquor industry in China, a marketing firm must consider all of the following factors affecting the Chinese liquor market EXCEPT that

\begin{enumerate}[label=\Alph*.]
    \item there is money to be made from marketing to the middle class.
    \item the government may control the pricing of products.
    \item there are few competitors to meet the demands of high end liquor consumers.
    \item the competition for winning over the middle class is very stiff.
[](https://cracku.in/11-in-the-context-of-the-passage-we-can-infer-that-to-x-cat-2024-slot-3)
\end{enumerate}
\vspace{1em}
\textbf{Solution:}

Let's evaluate the given choices and check if they are consistent with the information in the passage - 
_Option A_ : According to the passage, this is not true. The author emphasises that Moutai deliberately avoids targeting the middle class, as the competition in that market is intense, and instead focuses on the super-rich. This suggests that marketing to the middle class is not as lucrative or straightforward as it might appear.
_Option B_ : The passage mentions that the Chinese government is Moutai's largest shareholder and potentially plays a role in controlling the prices of products, which is a key factor for any firm in this market to consider. Furthermore, the author discusses how the government might further interfere with the operations in this space: [“... _in what some see as a travesty of corporate governance, its majority owner has plans to set up its own sales channel ..._ ”]
_Option C_ : The author highlights that the premium market in China is still growing and not overcrowded with luxury brands, indicating a potential opportunity for firms targeting high-end consumers.
_Option D_ : The passage states that many firms have failed in their attempt to cater to the middle class due to fierce competition. This is a valid consideration for any firm looking to enter the market.
Hence, Option A is the correct choice.

Instructions   

For the following questions answer them individually

\end{problem}

\newpage
\begin{problem}{}{}
% Subject: varc
\textbf{Instructions}

Instructions   

For the following questions answer them individually

\vspace{1em}
\textbf{Question 12}

**There is a sentence that is missing in the paragraph below. Look at the paragraph and decide where (option 1, 2, 3, or 4) the following sentence would best fit.**
**Sentence** : Taken outside the village of Trang Bang on June 8, 1972, the picture captured the trauma and indiscriminate violence of a conflict that claimed, by some estimates, a million or more civilian lives.
**Paragraph** : The horrifying photograph of children fleeing a deadly napalm attack has become a defining image not only of the Vietnam War but the 20th century. ___(1)___. Dark smoke billowing behind them, the young subjects’ faces are painted with a mixture of terror, pain and confusion. ___(2)___. Soldiers from the South Vietnamese Army’s 25th Division follow helplessly behind. ___(3)___. The picture was officially titled “The Terror of War,” but the photo is better known by the nickname given to the naked 9-year-old at its centre: “Napalm Girl”. ___(4)___.

\begin{enumerate}[label=\Alph*.]
    \item Option 4
    \item Option 1
    \item Option 3
    \item Option 2
[](https://cracku.in/12-there-is-a-sentence-that-is-missing-in-the-paragra-x-cat-2024-slot-3)
\end{enumerate}
\vspace{1em}
\textbf{Solution:}

To determine where the given sentence fits best, we must consider the logical flow of ideas between sentences. The sentence in question provides important historical and contextual information about the photograph. The author begins by describing the image and ends the paragraph by mentioning the official title of the photo. While sentences offering historical context are typically placed at the beginning, inserting the sentence in either Blank 1 or 2 would create a disruption in the flow of the image's description. This suggests that the sentence must fit in either Blank 3 or 4. Offering contextual information after mentioning the title at the end seems odd; thus, we can eliminate Blank 4. 
When we refocus on Blank 3, we see that by this point in the paragraph, the author has just described the soldiers from the South Vietnamese Army’s 25th Division following behind the children. This sets up a moment of helplessness that emphasises the chaos and brutality of the scene. The sentence about the trauma and violence of the war logically follows the description of the soldiers because it explains the broader implications of the photo’s content, providing a context for the violence and its far-reaching consequences. After discussing the soldiers, the sentence connects the personal and immediate suffering in the photograph with the larger scope of the Vietnam War’s impact, which then leads into the official title of the photograph. The official title, "The Terror of War," is directly linked to the broader theme of violence and destruction that the sentence describes. Thus, the given sentence naturally sets the stage for introducing the official title and the iconic nickname, "Napalm Girl," ensuring a smooth and coherent transition.
Hence, Option C is the correct choice.

Instructions   

**The passage below is accompanied by four questions. Based on the passage, choose the best answer for each question.**
Languages become endangered and die out for many reasons. Sadly, the physical annihilation of communities of native speakers of a language is all too often the cause of language extinction. In North America, European colonists brought death and destruction to many Native American communities. This was followed by US federal policies restricting the use of indigenous languages, including the removal of native children from their communities to federal boarding schools where native languages and cultural practices were prohibited. As many as 75 percent of the languages spoken in the territories that became the United States have gone extinct, with slightly better language survival rates in Central and South America ...
Even without physical annihilation and prohibitions against language use, the language of the "dominant" cultures may drive other languages into extinction; young people see education, jobs, culture and technology associated with the dominant language and focus their attention on that language. The largest language "killers" are English, Spanish, Portuguese, French, Russian, Hindi, and Chinese, all of which have privileged status as dominant languages threatening minority languages.
When we lose a language, we lose the worldview, culture and knowledge of the people who spoke it, constituting a loss to all humanity. People around the world live in direct contact with their native environment, their habitat. When the language they speak goes extinct, the rest of humanity loses their knowledge of that environment, their wisdom about the relationship between local plants and illness, their philosophical and religious beliefs, as well as their native cultural expression (in music, visual art and poetry) that has enriched both the speakers of that language and others who would have encountered that culture ...
As educators deeply immersed in the liberal arts, we believe that educating students broadly in all facets of language and culture ... yields immense rewards. Some individuals educated in the liberal arts tradition will pursue advanced study in linguistics and become actively engaged in language preservation, setting out for the Amazon, for example, with video recording equipment to interview the last surviving elders in a community to record and document a language spoken by no children.
Certainly, though, the vast majority of students will not pursue this kind of activity. For these students, a liberal arts education is absolutely critical from the twin perspectives of language extinction and global citizenship. When students study languages other than their own, they are sensitized to the existence of different cultural perspectives and practices. With such an education, students are more likely to be able to articulate insights into their own cultural biases, be more empathetic to individuals of other cultures, communicate successfully across linguistic and cultural differences, consider and resolve questions in a way that reflects multiple cultural perspectives, and, ultimately extend support to people, programs, practices, and policies that support the preservation of endangered languages.
There is ample evidence that such preservation can work in languages spiraling toward extinction. For example, Navajo, Cree, and Inuit communities have established schools in which these languages are the language of instruction, and the number of speakers of each has increased.

\end{problem}

\newpage
\begin{problem}{}{}
% Subject: varc
\textbf{Instructions}

Instructions   

For the following questions answer them individually

\vspace{1em}
\textbf{Question 13}

In the context of the passage, which one of the following hypothetical scenarios, if true, is NOT an example of the kind of loss that occurs when a language becomes extinct?

\begin{enumerate}[label=\Alph*.]
    \item The Nicobarese language describes 20 different moods of the ocean. By the time the last speaker is educated in a Central Board school, they will have forgotten their language.
    \item The Lamkangs of Manipur have only 3 remaining native speakers of the language. When they die, we will lose one more group from the government list of indigenous tribes.
    \item The Andamanese language has a word to describe someone who has lost a step-sister. When the language dies, we will lose the concept of the word and the emotions it evokes.
    \item The Inuits of Alaska have 35 different words to describe the texture of snow. When the language becomes extinct, we will lose that understanding of nature.
[](https://cracku.in/13-in-the-context-of-the-passage-which-one-of-the-fol-x-cat-2024-slot-3)
\end{enumerate}
\vspace{1em}
\textbf{Solution:}

The question asks us to identify the hypothetical scenario that does not reflect the kind of loss described in the passage. The passage discusses cultural, ecological, and intellectual losses caused by language extinction, including the loss of:
  * Unique cultural expressions (e.g., music, art, and emotions tied to language).
  * Knowledge about the environment (e.g., relationships between plants and illness).
  * Worldviews and philosophical insights.


We need to identify a scenario that doesn’t align with these themes. Let us examine the options based on this understanding -
_Option A_ : This scenario reflects the loss of unique cultural knowledge - in this case, the ability to describe the “20 different moods of the ocean,” which likely represents detailed ecological and environmental understanding.
_Option B_ : This scenario focuses on some form of administrative or statistical change in a government list; the focus is not on the cultural, ecological, or intellectual loss emphasised in the passage. Therefore, the option does not reflect the deeper, humanity-wide loss described in the passage.
_Option C_ : This scenario aligns with the loss of unique cultural concepts and the emotional depth tied to a word or phrase. In the author’s perspective, losing this concept would diminish humanity's understanding of the complexity of human relationships.
_Option D_ : This scenario reflects the loss of ecological knowledge about snow textures, which likely has practical implications for living in the Arctic environment. This aligns with the passage's discussion of losing environmental wisdom when languages die.
Hence, Option B is the correct choice.

\end{problem}

\newpage
\begin{problem}{}{}
% Subject: varc
\textbf{Instructions}

Instructions   

For the following questions answer them individually

\vspace{1em}
\textbf{Question 14}

Which one of the following hypothetical scenarios, if true, would most strongly undermine the central ideas of the passage?

\begin{enumerate}[label=\Alph*.]
    \item Most liberal arts students will pursue jobs in publishing and human resource management rather than doctorates in linguistics.
    \item A liberal arts education requires that, in addition to being fluent in English, students gain fluency in two of the top five most spoken languages globally.
    \item Schools that teach endangered languages can preserve the language only for a generation.
    \item Recording a dying language that has only a few remaining speakers freezes it in time: it stops evolving further.
[](https://cracku.in/14-which-one-of-the-following-hypothetical-scenarios--x-cat-2024-slot-3)
\end{enumerate}
\vspace{1em}
\textbf{Solution:}

The central idea of the passage is that endangered languages, as carriers of unique cultural perspectives and human knowledge, must be preserved to benefit humanity. The author argues that liberal arts education plays a vital role in this preservation, both by fostering global citizenship and by encouraging some individuals to directly engage in preservation efforts. Liberal arts education, as framed in the passage, sensitises students to cultural diversity and equips them with the tools to support endangered languages and cultures.
We observe that _Option B_ , however, directly undermines this central idea by redefining the focus of liberal arts education. Requiring fluency in two of the most widely spoken global languages (e.g., English, Spanish, Mandarin) would prioritise dominant languages rather than endangered ones. This hypothetical scenario shifts resources and attention away from the preservation of linguistic diversity, which is central to the passage's argument. Such a requirement would reinforce the dominance of already powerful languages, the very phenomenon identified as a major “language killer” in the passage. By institutionalizing the focus on dominant languages, it would erode the argument that liberal arts education fosters support for endangered languages and cultural preservation, ultimately weakening the role of liberal arts in addressing language extinction.
In contrast, the remaining choices either align with the discussion or do not serve as strong counterarguments to the points presented in the passage. For instance, _Option A_ acknowledges that most liberal arts students will not directly engage in language preservation but does not challenge the broader idea that liberal arts education fosters empathy and support for endangered languages. _Option C_ limits the long-term success of language preservation but does not negate its immediate benefits or the potential for renewal in subsequent generations. Similarly, _Option D_ highlights a limitation of recording dying languages but does not undermine the broader argument that documentation is a valuable and necessary tool in preservation.

\end{problem}

\newpage
\begin{problem}{}{}
% Subject: varc
\textbf{Instructions}

Instructions   

For the following questions answer them individually

\vspace{1em}
\textbf{Question 15}

It can be inferred from the passage that it is likely South America had a slightly better language survival rate than North America for all of the following reasons EXCEPT:

\begin{enumerate}[label=\Alph*.]
    \item European colonists allowed children of native speakers to stay at home with their families.
    \item the colonial government was unable to mainstream the locals.
    \item not many native speakers were killed by European colonists.
    \item locals were provided job opportunities in the colonial administration.
[](https://cracku.in/15-it-can-be-inferred-from-the-passage-that-it-is-lik-x-cat-2024-slot-3)
\end{enumerate}
\vspace{1em}
\textbf{Solution:}

The passage highlights the widespread extinction of indigenous languages in North America due to colonisation, physical annihilation, and assimilation policies while noting slightly better survival rates in Central and South America. It implies that language survival may be influenced by factors such as social policies, cultural integration, and the extent of physical and cultural displacement. 
Evaluating the choices, we note that  _Option A_ is plausible since allowing children to stay with families would help preserve native languages, unlike the North American policy of removing children to boarding schools, as discussed in the passage. _Option B_ is also reasonable, as less effective assimilation efforts by colonial governments could lead to better language retention. _Option C_ also aligns with the passage’s context, as it discusses physical annihilation as a significant driver of language extinction, but the survival rate being slightly better in South America could suggest marginally less physical annihilation. 
However, _Option D_ is problematic; while providing locals with jobs in the colonial administration might seem like a factor that supports language retention, this scenario is not consistent with the passage. The passage emphasises that dominant languages often replace indigenous ones through socio-economic pressures, and employment in colonial administration would likely reinforce the use of the dominant language rather than preserve native languages. 
Hence, Option D is the best choice.

\end{problem}

\newpage
\begin{problem}{}{}
% Subject: varc
\textbf{Instructions}

Instructions   

For the following questions answer them individually

\vspace{1em}
\textbf{Question 16}

The author believes that a liberal arts education combined with participation in language preservation empower students in all of the following ways EXCEPT that they will

\begin{enumerate}[label=\Alph*.]
    \item overcome cultural barriers to communication.
    \item learn different languages.
    \item establish schools to preserve languages spiralling towards extinction.
    \item develop a better understanding of their own culture.
[](https://cracku.in/16-the-author-believes-that-a-liberal-arts-education--x-cat-2024-slot-3)
\end{enumerate}
\vspace{1em}
\textbf{Solution:}

The author discusses how a liberal arts education, combined with participation in language preservation efforts, empowers students in several significant ways. He highlights that such an education broadens students' cultural understanding, helps them communicate across linguistic barriers (_Option A_), and enables them to gain insights into both their own and others' cultures (_Option D_).  _Option B_ is an implicit aspect of this discussion. The passage also suggests that some students may even become involved in active language preservation, though it emphasises that the majority might not pursue this path. Contrarily, _Option C_ is not explicitly mentioned or understood: while the passage provides examples of communities that have established such schools, it does not suggest that students themselves will take on this role.

Instructions   

For the following questions answer them individually

\end{problem}

\newpage
\begin{problem}{}{}
% Subject: varc
\textbf{Instructions}

Instructions   

For the following questions answer them individually

\vspace{1em}
\textbf{Question 17}

**The passage given below is followed by four alternate summaries. Choose the option that best captures the essence of the passage.**

\begin{enumerate}[label=\Alph*.]
    \item The tradwife’s commitment to outdated gender roles and retro fashion critiques the superficiality of today's societal ideals.
    \item By promoting an idealized past, the tradwife exposes the artifice of contemporary values and mocks societal norms.
    \item The tradwife, with her vintage dress and traditional roles, highlights the superficiality of modern life and challenges current societal norms.
    \item The tradwife’s vintage dress and adherence to traditional roles reveal the artificial nature of modern life and its superficial values.
[](https://cracku.in/17-the-passage-given-below-is-followed-by-four-altern-x-cat-2024-slot-3)
\end{enumerate}
\vspace{1em}
\textbf{Solution:}

**Option C is the correct answer.** This option best captures the essence of the passage. The tradwife’s embrace of traditional roles and vintage fashion naturally contrasts with and highlights the superficiality of modern life. The passage emphasizes that she is not overtly critiquing society but rather embodying a regressive ideal that challenges contemporary norms. This fits with how the passage describes her actions—her existence itself is a challenge to modern societal expectations.
_Option A:_ The passage does not suggest that the tradwife is actively critiquing modern society. Instead, the tradwife embodies a lifestyle and set of ideals that are in contrast to contemporary values. The critique comes from others' reactions to her behavior, not from her intentions. 
_Option B_ : The passage does not show her as actively trying to expose or mock anything. She simply lives in a way that contrasts with modern expectations. The "mockery" is a reaction from others, not her goal. Therefore, this option overstates her intentions._  
  
Option D_ : This focuses on the tradwife's dress and adherence to traditional roles as the means of revealing modern life’s artificial nature. While the tradwife does embody these traditional values, the passage is more focused on how her actions highlight societal superficiality rather than "revealing" it. Additionally, the passage suggests that her behaviour is not an overt revelation but rather something that others react to—making this option somewhat distorted.

\end{problem}

\newpage
\begin{problem}{}{}
% Subject: varc
\textbf{Instructions}

Instructions   

For the following questions answer them individually

\vspace{1em}
\textbf{Question 18}

**The passage given below is followed by four alternate summaries. Choose the option that best captures the essence of the passage.**
Lyric poetry is a genre of private meditation rather than public commitment. The impulse in Marxism toward changing a society deemed unacceptable in its basic design would seem to place demands on lyric poetry that such poetry, with its tendency toward the personal, the small scale, and the idiosyncratic, could never answer. There is within Marxism, however, also a strand of thought that would locate in lyric poetry alternative modes of perception and description that call forth a vision of worlds at odds with a repressive reality or that draw attention to the workings of ideology within the hegemonic culture. The poetic imagination may indeed deflect larger social concerns, but it may also be implicitly critical and utopian.

\begin{enumerate}[label=\Alph*.]
    \item The focus of lyric poetry is largely personal while that of Marxism is bringing change in society. Unless the difference is resolved, poetry will remain largely utopian.
    \item Marxism has internal contradictions due to which one strand of Marxism sees no merit in lyric poetry while another appreciates the alternative modes of perception in poetry.
    \item The focus of lyric poetry as personal may not seem compatible with Marxism. However, it is possible to envisage lyric poetry as a symbol of resistance against an oppressive culture.
    \item Marxism makes unreasonable demands on lyric poetry. However, lyric poetry has its own merits that are largely ignored by Marxism due to its personal nature.
[](https://cracku.in/18-the-passage-given-below-is-followed-by-four-altern-x-cat-2024-slot-3)
\end{enumerate}
\vspace{1em}
\textbf{Solution:}

The passage contrasts the personal, introspective nature of lyric poetry with Marxism's outward focus on societal change. At first glance, the characteristics of lyric poetry - being personal, small-scale, and idiosyncratic - seem incompatible with Marxism's demand for a transformative critique of an unjust society. However, the passage introduces a nuanced perspective within Marxist thought, which recognises lyric poetry as having an implicit critical and utopian function. Through its imaginative and alternative modes of perception, lyric poetry can challenge dominant ideologies and suggest a vision of resistance to oppression. Thus, while lyric poetry does not directly engage with large-scale social concerns, it can still align with Marxist ideals by offering subtle forms of critique and hope for change. _Option C_ most effectively captures this idea. 
None of the other choices present a valid interpretation of the passage. For instance, _Option A_ oversimplifies the issue by portraying lyric poetry as merely ‘utopian,” ignoring its critical and resistant potential as described in the passage. It misses the nuanced compatibility suggested between lyric poetry and Marxism. _Option B_ inaccurately suggests that Marxism has “internal contradictions,” which is not the focus of the passage. Similarly, _Option D_ frames Marxism as dismissive of lyric poetry, ignoring the strand of Marxist thought that finds value in poetry's alternative perspectives; this misrepresentation makes it a weaker summary.

\end{problem}

\newpage
\begin{problem}{}{}
% Subject: varc
\textbf{Instructions}

Instructions   

For the following questions answer them individually

\vspace{1em}
\textbf{Question 19}

**Five jumbled-up sentences (labelled 1, 2, 3, 4 and 5) related to a topic are given below. Four of them can be put together to form a coherent paragraph. Identify the odd sentence and key in the number of that sentence as your answer.**
2. This fetal warm-up act—the soldering of neural connections before the eyes actually function—is crucial to the performance of the visual system.
3. The reasons for this paring back of synapses is a mystery, but synaptic pruning is thought to sharpen and reinforce the “correct” synapses, while removing the weak and unnecessary ones.
4. Neural connections between the eyes and the brain are formed long before birth, establishing the wiring and the circuitry that allow a child to begin visualizing the world the minute she emerges from the womb.
5. During this rehearsal period, synapses—points of chemical connection—between nerve cells are generated in great excess, only to be pruned back during later development.
Backspace  
789  
456  
123  
0.-  
Clear All  

Submit
[](https://cracku.in/19-five-jumbled-up-sentences-labelled-1-2-3-4-and-5-r-x-cat-2024-slot-3)

\begin{enumerate}[label=\Alph*.]
\end{enumerate}
\vspace{1em}
\textbf{Solution:}

A good starting point here would be to link sentences 5 and 3 based on the idea of “pruning.” Sentence 5 talks about the synaptic development that occurs during the early period described in sentences 4 and 2. It describes how synapses are initially overproduced, a concept that sets up the next part of the process: synaptic pruning. Sentence 3 continues this discussion by offering an explanation for why synapses are pruned: to reinforce the necessary connections and eliminate the weaker ones. This helps to fine-tune the brain’s wiring, ensuring that the important pathways are retained while the less useful ones are discarded. Together, 5-3 ties in with the broader purpose of neural wiring presented in Sentence 2: it emphasises that the creation of neural connections in the fetus is a preparatory stage and that these connections are essential for the future performance of the visual system. Sentence 4 sets the context of the paragraph: it mentions the formation of neural connections before birth, specifically focusing on the connections between the eyes and the brain. 
While Sentence 1 provides a technical description of synapse creation and function, it doesn’t fit the broader context set by sentences 4, 2, 5, and 3. These sentences focus on brain development (especially related to vision) and the process of synapse formation and pruning in a developmental context. Sentence 1 focuses on the structural aspect of synapse formation (the “specialized structures” at the terminal end of an axon), which is more about how a synapse functions at a microscopic level rather than the larger developmental or functional role synapses play in early brain development, as discussed in the other sentences.
Therefore, Sentence 1 is the odd one out.

\end{problem}

\newpage
\begin{problem}{}{}
% Subject: varc
\textbf{Instructions}

Instructions   

For the following questions answer them individually

\vspace{1em}
\textbf{Question 20}

**The passage given below is followed by four alternate summaries. Choose the option that best captures the essence of the passage.**
Humans have managed to tweak the underlying biology of various plants and animals to produce high-tech crops and microbes. But regulating these entities is complicated, as the framework of policies and procedures are outdated and not flexible enough to adapt to emerging technology. The question is whether regulation will ever be able to keep up with human innovation, to regulate living things, which are apt to be unpredictable and unique; to capture all the potential risks when new biological entities are introduced, or when they pass on variations of their genes?

\begin{enumerate}[label=\Alph*.]
    \item The mercurial nature of biological entities calls for scientists to shape the regulations governing emerging technology, with regular calibration to handle variations in the field.
    \item A new framework of rules and procedures for regulating the most recent research emerging from biotechnology is urgently needed, to keep up with this rapidly changing discipline.
    \item Current regulation of biotechnology is outdated, but it is debatable if we can create a framework, imaginative and flexible, to cover all contingencies in this fast-changing area.
    \item The problem with formulating regulation for innovation in the scientific arena it that it is impossible to imagine the outcomes or risks related to the outcomes of all the research.
[](https://cracku.in/20-the-passage-given-below-is-followed-by-four-altern-x-cat-2024-slot-3)
\end{enumerate}
\vspace{1em}
\textbf{Solution:}

The passage discusses the challenges of regulating innovations in biotechnology. It highlights two key issues: (i) current regulations are outdated and insufficient for new biological advancements, and (ii) the unpredictable nature of living entities and rapid technological changes make it difficult to design a regulatory framework that anticipates all risks and contingencies. The core question is whether regulation can ever keep pace with the unpredictable and unique outcomes of biotechnology. _Option C_ correctly touches upon these ideas.
The other options are either too narrow or miss the essence. _Option A_ focuses on calibration and scientists' roles, which is not the central concern of the passage. Similarly, _Option B_ highlights the urgency for new rules but omits the skepticism about achieving comprehensive regulation. _Option D_ overemphasises the impossibility of imagining all risks without addressing the need for flexibility in regulations.
Hence, Option C is the best choice.

Instructions   

**The passage below is accompanied by four questions. Based on the passage, choose the best answer for each question.**
There is a group in the space community who view the solar system not as an opportunity to expand human potential but as a nature preserve, forever the provenance of an elite group of scientists and their sanitary robotic probes. These planetary protection advocates [call] for avoiding “harmful contamination” of celestial bodies. Under this regime, NASA incurs great expense sterilizing robotic probes in order to prevent the contamination of entirely theoretical biospheres ...
Transporting bacteria would matter if Mars were the vital world once imagined by astronomers who mistook optical illusions for canals. Nobody wants to expose Martians to measles, but sadly, robotic exploration reveals a bleak, rusted landscape, lacking oxygen and flooded with radiation ready to sterilize any Earthly microbes. Simple life might exist underground, or down at the bottom of a deep canyon, but it has been very hard to find with robots. . . . The upsides from human exploration and development of Mars clearly outweigh the welfare of purely speculative Martian fungi ...
The other likely targets of human exploration, development, and settlement, our moon and the asteroids, exist in a desiccated, radiation-soaked realm of hard vacuum and extreme temperature variations that would kill nearly anything. It’s also important to note that many international competitors will ignore the demands of these protection extremists in any case. For example, China recently sent a terrarium to the moon and germinated a plant seed—with, unsurprisingly, no protest from its own scientific community. In contrast, when it was recently revealed that a researcher had surreptitiously smuggled super-resilient microscopic tardigrades aboard the ill-fated Israeli Beresheet lunar probe, a firestorm was unleashed within the space community ...
NASA’s previous human exploration efforts made no serious attempt at sterility, with little notice. As the Mars expert Robert Zubrin noted in the National Review, U.S. lunar landings did not leave the campsites cleaner than they found it. Apollo’s bacteria-infested litter included bags of feces. Forcing NASA’s proposed Mars exploration to do better, scrubbing everything and hauling out all the trash, would destroy NASA’s human exploration budget and encroach on the agency’s other directorates, too. Getting future astronauts off Mars is enough of a challenge, without trying to tote weeks of waste along as well.
A reasonable compromise is to continue on the course laid out by the U.S. government and the National Research Council, which proposed a system of zones on Mars, some for science only, some for habitation, and some for resource exploitation. This approach minimizes contamination, maximizes scientific exploration ... Mars presents a stark choice of diverging human futures. We can turn inward, pursuing ever more limited futures while we await whichever natural or manmade disaster will eradicate our species and life on Earth. Alternatively, we can choose to propel our biosphere further into the solar system, simultaneously protecting our home planet and providing a backup plan for the only life we know exists in the universe. Are the lives on Earth worth less than some hypothetical microbe lurking under Martian rocks?

\end{problem}

\newpage
\begin{problem}{}{}
% Subject: varc
\textbf{Instructions}

Instructions   

For the following questions answer them individually

\vspace{1em}
\textbf{Question 21}

The author is unlikely to disagree with any of the following EXCEPT:

\begin{enumerate}[label=\Alph*.]
    \item the proposal for a zonal segregation of the Martian landscape into regions for different purposes.
    \item that while NASA’s earlier missions were not ideal in their approach to space contamination, they likely did no grave damage.
    \item space contamination should be minimised until the possibility of life on the astronomical body being explored is ruled out.
    \item the exorbitant costs of continuing to keep the space environment pristine may be unsustainable.
[](https://cracku.in/21-the-author-is-unlikely-to-disagree-with-any-of-the-x-cat-2024-slot-3)
\end{enumerate}
\vspace{1em}
\textbf{Solution:}

The phrase “unlikely to disagree + EXCEPT” can seem tricky to interpret. In simple terms, the question requires us to find a statement the author will disagree with. Let us inspect the choices - 
_Option A_ : The author supports the proposal for zonal segregation as a reasonable compromise, balancing scientific exploration with human settlement.
_Option B_ : The author agrees that NASA’s earlier missions did not prioritise contamination but implies they caused no significant harm.
_Option C_ : This viewpoint reflects a cautious approach to space exploration. The author dismisses concerns about hypothetical extraterrestrial life as speculative and prioritises human exploration and development over minimising contamination. Therefore, he’s likely to disagree with this position. 
_Option D_ : In the passage, the author argues that the costs of maintaining strict planetary protection measures are excessive and could undermine future exploration efforts. This is consistent with his stance. 
Hence, Option C is the best choice.

\end{problem}

\newpage
\begin{problem}{}{}
% Subject: varc
\textbf{Instructions}

Instructions   

For the following questions answer them individually

\vspace{1em}
\textbf{Question 22}

The author mentions all of the following reasons to dismiss concerns about contaminating Mars EXCEPT:

\begin{enumerate}[label=\Alph*.]
    \item the lack of evidence of living organisms on Mars makes possible contamination from earthly microbes a moot point.
    \item efforts to contain contamination on Mars are likely to be derailed as competitor countries may not follow similar restrictions.
    \item the use of similar probes on astronomical bodies like the moon have had little effect on the environment.
    \item earlier explorations have already contaminated pristine space environments.
[](https://cracku.in/22-the-author-mentions-all-of-the-following-reasons-t-x-cat-2024-slot-3)
\end{enumerate}
\vspace{1em}
\textbf{Solution:}

The passage discusses the debate surrounding planetary protection policies, particularly the concerns about contaminating Mars with Earth-based microbes. The author argues against these concerns, citing several reasons why the risk of contamination should not hinder human exploration and development of Mars. These reasons include:
  * the lack of evidence for life on Mars (describes Mars as a “bleak, rusted landscape” with no confirmed life) [_Option A_]
  * the disregard for such protocols by international competitors (China’s lenient approach to planetary protection) [_Option B_]
  * the historical precedent of contamination from earlier human missions (Apollo missions left waste on the Moon) [_Option D_]


On the other hand, _Option C_ is not presented as a valid reason. The author does not specifically argue that probes have had “little effect” on the Moon's environment but instead focuses on human waste and contamination from earlier human missions, not robotic probes.

\end{problem}

\newpage
\begin{problem}{}{}
% Subject: varc
\textbf{Instructions}

Instructions   

For the following questions answer them individually

\vspace{1em}
\textbf{Question 23}

The author’s overall tone in the first paragraph can be described as

\begin{enumerate}[label=\Alph*.]
    \item sceptical about the excessive efforts to sanitise planets where life has not yet been proven to exist.
    \item equivocal about the reasons extended by the group of scientists seeking to limit space exploration.
    \item indifferent to the elitism of a few scientists aiming to corner space exploration.
    \item approving of the amount of money NASA spends to restrict the spread of contamination in space.
[](https://cracku.in/23-the-authors-overall-tone-in-the-first-paragraph-ca-x-cat-2024-slot-3)
\end{enumerate}
\vspace{1em}
\textbf{Solution:}

The first paragraph critiques the stringent planetary protection policies advocated by a group of scientists who aim to prevent biological contamination of celestial bodies. The author portrays these efforts as excessive, particularly given the lack of evidence for extraterrestrial life, and highlights the significant financial burden these measures place on space agencies like NASA. _Option A_ accurately reflects this scepticism, as the author questions the need to sterilize planets where life has not been proven to exist. 
_Option B_ is incorrect because the author is not equivocal (i.e., ambiguous or undecided); instead, he expresses a clear stance against these strict protocols. _Option C_ is also inaccurate, as the author is not indifferent to elitism but rather critiques the scientists’ restrictive approach. Similarly, _Option D_ can be eliminated because the author does not approve of NASA's spending on sterilization but views it as an unnecessary expense.

\end{problem}

\newpage
\begin{problem}{}{}
% Subject: varc
\textbf{Instructions}

Instructions   

For the following questions answer them individually

\vspace{1em}
\textbf{Question 24}

The contrasting reactions to the Chinese and Israeli “contaminations” of lunar space

\begin{enumerate}[label=\Alph*.]
    \item are valid as the contamination of the lunar environment from animal sources is far greater than from plants.
    \item are evidence of China’s reasonable approach towards space contamination.
    \item indicate that national scientists may have different sensitivities to issues of biosphere protection.
    \item reveal global biases prevalent in attitudes towards different countries.
[](https://cracku.in/24-the-contrasting-reactions-to-the-chinese-and-israe-x-cat-2024-slot-3)
\end{enumerate}
\vspace{1em}
\textbf{Solution:}

The passage highlights contrasting reactions to two instances of potential contamination of the lunar environment: China’s germination of a plant seed on the Moon, which elicited little controversy, and Israel’s accidental release of tardigrades aboard the Beresheet probe, which sparked significant backlash within the space community. This contrast underscores differences in how national or regional scientific communities respond to issues of planetary protection. _Option C_ most closely reflects this idea. 
The passage does not suggest that contamination from animals is inherently more harmful than from plants, as suggested in _Option A_. Similarly, _Option B_ inaccurately implies that the passage endorses China’s approach as inherently “reasonable,” which it does not. _Option D_ diverges a bit from the discussion by emphasising global biases against specific countries, but the passage provides no evidence of such biases, focusing instead on scientific reactions. 
[](https://cracku.in/downloads/19319)
  *  
  *

\end{problem}
